% !TeX program = xelatex
% 数学 LaTeX 排版 Skill v4.0.0；版式保持 v3 金标准. XeLaTeX 编译至少两次. 不分发字体.
% WZ-LOCKED-CLASS-BEGIN
\begin{filecontents*}[overwrite]{elegantbook-original-adapter.cls}
\NeedsTeXFormat{LaTeX2e}
\ProvidesClass{elegantbook-original-adapter}[2026/09/23 v3.0 fixed mathematical book environments]
\LoadClass[10pt,letterpaper,oneside]{book}
\RequirePackage[UTF8,scheme=plain,fontset=fandol]{ctex}
\RequirePackage{fontspec}
\setmainfont{cmunrm.otf}[BoldFont=cmunbx.otf,ItalicFont=cmunti.otf,BoldItalicFont=cmunbi.otf]
\setCJKfamilyfont{sourceheader}[ItalicFont=FandolKai-Regular.otf]{FandolKai-Regular.otf}
\RequirePackage{amsmath,amssymb,mathrsfs}
\RequirePackage{amsthm,etoolbox}
\RequirePackage{xcolor,geometry,setspace,fancyhdr,enumitem,needspace,xurl}
\definecolor{structurecolor}{HTML}{880E4F}
\definecolor{main}{HTML}{9C27B0}
\geometry{paperwidth=612bp,paperheight=792bp,left=20mm,right=20mm,top=95.04bp,bottom=20mm,headheight=12pt,headsep=25pt,footskip=30pt}
\setstretch{1.6}
\setlength{\parindent}{2em}
\setlength{\parskip}{0pt}
\setlist{nosep,leftmargin=2em}
\AtBeginDocument{\setlength{\abovedisplayskip}{4pt plus 1pt minus 1pt}\setlength{\belowdisplayskip}{4pt plus 1pt minus 1pt}\setlength{\abovedisplayshortskip}{2pt}\setlength{\belowdisplayshortskip}{3pt}}
\fancyhf{}
\fancyhead[L]{{\itshape\CJKfamily{sourceheader}\rightmark}}
\fancyhead[R]{{\itshape\CJKfamily{sourceheader}\leftmark}}
\fancyfoot[C]{\thepage}
\renewcommand{\headrulewidth}{0.4pt}
\renewcommand{\headrule}{{\color{structurecolor}\hrule height\headrulewidth width\headwidth\vskip-\headrulewidth}}
\pagestyle{fancy}
\fancypagestyle{cover}{\fancyhf{}\fancyfoot[C]{\thepage}\renewcommand{\headrulewidth}{0pt}\renewcommand{\headrule}{}}
\RequirePackage{hyperref}
\hypersetup{unicode=true,colorlinks=true,linkcolor=structurecolor,urlcolor=structurecolor,citecolor=structurecolor,linktoc=section,bookmarksnumbered=true,bookmarksopen=true,pdfborder={0 0 0},pdfstartview=FitH}
\setcounter{tocdepth}{1}
\renewcommand*\l@section{\@dottedtocline{1}{1.5em}{2.4em}}
\renewcommand{\thesection}{\thechapter.\arabic{section}}
\newcommand{\originalchapter}[1]{\par\addvspace{14pt}\Needspace{5\baselineskip}\refstepcounter{chapter}\setcounter{section}{0}\addcontentsline{toc}{chapter}{\protect\numberline{\thechapter}#1}\markboth{\thechapter\quad #1}{}\begingroup\centering\color{structurecolor}\bfseries\fontsize{14.4}{20}\selectfont\thechapter\quad #1\par\endgroup\nobreak\vspace{8pt}}
\newcommand{\originalsection}[1]{\par\addvspace{9pt}\Needspace{4\baselineskip}\refstepcounter{section}\addcontentsline{toc}{section}{\protect\hspace*{1.5em}\protect\numberline{\protect\textcolor{structurecolor}{\thesection}}#1}\markright{\thesection\quad #1}\noindent{\fontsize{12}{17}\selectfont\color{structurecolor}\thesection\quad}{\bfseries\fontsize{12}{17}\selectfont #1}\par\nobreak\vspace{4pt}}

% Semantic environments are registered once here, not re-designed in manuscripts.
\newtheoremstyle{wzstatement}{6pt}{5pt}
  {\normalfont\kaishu}{0pt}{\normalfont\bfseries\color{structurecolor}}
  {}{.7em}{\thmname{#1}\thmnumber{ #2}\thmnote{\normalfont\ (#3)}}
\newtheoremstyle{wzdefinition}{6pt}{5pt}
  {\normalfont\songti}{0pt}{\normalfont\bfseries\color{structurecolor}}
  {}{.7em}{\thmname{#1}\thmnumber{ #2}\thmnote{\normalfont\ (#3)}}
\newtheoremstyle{wzproblem}{7pt}{5pt}
  {\normalfont\songti}{2em}{\normalfont\bfseries\color{main}}
  {}{.7em}{\thmname{#1}\thmnumber{ #2}}
\newtheoremstyle{wzremark}{4pt}{4pt}
  {\normalfont\songti}{0pt}{\normalfont\bfseries}
  {}{.7em}{\thmname{#1}}
\theoremstyle{wzstatement}
\newtheorem{theorem}{定理}[chapter]
\newtheorem{lemma}[theorem]{引理}
\newtheorem{proposition}[theorem]{命题}
\newtheorem{corollary}[theorem]{推论}
\theoremstyle{wzdefinition}
\newtheorem{definition}[theorem]{定义}
\theoremstyle{wzproblem}
\newtheorem{example}{例题}[chapter]
\newtheorem{exercise}{习题}[chapter]
\theoremstyle{wzremark}
\newtheorem*{remark}{注}
\renewcommand{\proofname}{证明}
% Retain the native amsthm QED stack; only adjust the fixed typography.
\expandafter\patchcmd\csname\string\proof\endcsname{\normalfont \topsep6\p@\@plus6\p@\relax}
  {\normalfont\songti \topsep5\p@\@plus1\p@\relax}
  {}{\ClassError{elegantbook-original-adapter}{Cannot patch proof spacing}{Check the amsthm version.}}
\expandafter\patchcmd\csname\string\proof\endcsname{\itshape}{\normalfont\bfseries}
  {}{\ClassError{elegantbook-original-adapter}{Cannot patch proof heading}{Check the amsthm version.}}
\expandafter\patchcmd\csname\string\proof\endcsname{\ignorespaces}
  {\setlength{\parindent}{2em}\ignorespaces}
  {}{\ClassError{elegantbook-original-adapter}{Cannot patch proof paragraphs}{Check the amsthm version.}}
\AtBeginEnvironment{proof}{\par\Needspace{3\baselineskip}}
\renewcommand{\qedsymbol}{\ensuremath{\square}}
\newenvironment{solution}{\begin{proof}[解答]}{\end{proof}}
\newcommand{\wzcheckchapter}{\ifnum\value{chapter}<1
  \ClassError{elegantbook-original-adapter}{Numbered mathematics before chapter 1}{Move it after the first originalchapter.}\fi}

\newcommand{\wzregisterguard}[1]{\AtBeginEnvironment{#1}{\wzcheckchapter\par\Needspace{3\baselineskip}}}
\forcsvlist{\wzregisterguard}{theorem,lemma,proposition,corollary,definition,example,exercise}
% Front matter and bibliography are not numbered mathematical chapters.
\newcommand{\originalpreface}{\par\addvspace{10pt}\Needspace{5\baselineskip}
  \noindent{\bfseries\fontsize{12}{17}\selectfont 前言}\par\nobreak\vspace{4pt}}
\newcommand{\originalcontents}{\par\addvspace{14pt}
  \begin{center}{\color{structurecolor}\bfseries\fontsize{14.4}{20}\selectfont 目录}\end{center}
  \vspace{-10pt}\@starttoc{toc}}
\renewcommand{\bibname}{参考文献}
\renewenvironment{thebibliography}[1]{
  \par\addvspace{12pt}\Needspace{3\baselineskip}\phantomsection
  \addcontentsline{toc}{chapter}{\bibname}
  \markboth{\bibname}{}
  \noindent{\color{structurecolor}\bfseries\fontsize{14.4}{20}\selectfont\bibname}\par\nobreak\vspace{6pt}
  \list{\@biblabel{\@arabic\c@enumiv}}
    {\settowidth\labelwidth{\@biblabel{#1}}\leftmargin\labelwidth
     \advance\leftmargin\labelsep\usecounter{enumiv}
     \let\p@enumiv\@empty\renewcommand\theenumiv{\@arabic\c@enumiv}
     \setlength{\itemsep}{3pt}\setlength{\parsep}{0pt}}
  \sloppy\clubpenalty4000\widowpenalty4000\sfcode`\.=1000\relax
}{\def\@noitemerr{\@latex@warning{Empty bibliography}}\endlist}

\numberwithin{equation}{chapter}
\renewcommand{\theequation}{\thechapter.\arabic{equation}}
\allowdisplaybreaks[2]
\emergencystretch=2em
\clubpenalty=6000
\widowpenalty=6000
\raggedbottom
\endinput
\end{filecontents*}
% WZ-LOCKED-CLASS-END
\documentclass{elegantbook-original-adapter}
% ===== 元数据内容槽 =====
\newcommand{\BookTitle}{常微分方程与动力系统}
\newcommand{\BookAuthor}{Haynes Miller, Arthur Mattuck, Daniel Rothman 原课程; 中文重编}
\newcommand{\BookDate}{MIT 18.03 Spring 2010; 12.006J Fall 2022}
\hypersetup{pdftitle={常微分方程与动力系统: MIT OCW 中文重编本},pdfauthor={MIT OCW 原课程作者; 中文重编}}
\usepackage{tikz,pgfplots,booktabs,longtable}
\pgfplotsset{compat=1.18}
\usetikzlibrary{arrows.meta,calc,positioning}
\newcommand{\R}{\mathbb R}
\newcommand{\C}{\mathbb C}
\newcommand{\N}{\mathbb N}
\newcommand{\dd}{\,\mathrm d}
\newcommand{\ee}{\mathrm e}
\newcommand{\ii}{\mathrm i}
\newcommand{\eps}{\varepsilon}
\newcommand{\Lap}{\mathcal L}
\newcommand{\tr}{\operatorname{tr}}
\newcommand{\diag}{\operatorname{diag}}
\begin{document}
\frontmatter
\begin{center}
{\color{structurecolor}\bfseries\fontsize{18}{25}\selectfont\BookTitle\par}
{\bfseries MIT OpenCourseWare 中文重编本\par}
\BookAuthor\par
\BookDate\par
\end{center}
\originalpreface
常微分方程把变化规律写成时间导数与当前状态之间的关系. 能求出一个公式只是第一层目标; 初值是否决定唯一运动、解能延续多久、扰动会不会放大、长期运动由什么结构控制, 同样属于方程本身的问题. 本册面向掌握一元微积分、基本多元微积分及矩阵运算的读者, 按初值问题、解析方法、线性响应、系统与非线性动力学的依赖顺序编排.

主要依据为 MIT 18.03 Differential Equations, Spring 2010 的官方课堂讲义及 Miller 的补充讲义、Mattuck 的 Notes and Exercises. 动力系统部分选用 Daniel Rothman 的 MIT 12.006J Nonlinear Dynamics: Chaos, Fall 2022. 原件逐文件页数、作者标注、年份与许可例外见随书来源清单. 本书合并相同数学内容的课堂记录与补充讲义, 不把同一资源的多个链接算成多份讲义. 第1--13章保留既有数学讲义, 其例题与章节末练习为本册整理或另设; 后续课程题库、作业、习题课及考试按原题编号另编, 公开题面完整中文表述, 官方译审与编者补解分别注明.

数值算法只给与解的连续依赖有关的误差衔接, 其完整算法分析参见本仓库数值分析册. 局部存在唯一性在这里给出可直接使用的定理与积分迭代证明, 不重复拓扑册的一般不动点理论. Fourier 级数仅处理线性周期响应所需的情形. 流体偏微分方程、分形测度的完整理论以及混沌系统的全局分类不在本册范围. 此编排是应用数学课程补充, 不作为2027年任何招生考纲独立新增学科的证明.

MIT OCW 默认许可为 \href{https://creativecommons.org/licenses/by-nc-sa/4.0/}{CC BY-NC-SA 4.0}; 中文改编以相同许可供非商业使用并以相同方式共享. 原作者与机构依课程及文件署名登记, 改动包括翻译、重排、补充证明与自行绘图, 不表示 MIT 或作者认可本项目. 单独版权声明优先于默认许可; 含第三方保留内容的整份文件仅留官方链接, 不随仓库及源码 ZIP 复制. 英文原件只在 sources 出处档案, 不嵌入中文阅读 PDF, 亦不收入中文源码 ZIP. 本册图形用 TikZ 独立构造. 官方条款见 \href{https://ocw.mit.edu/pages/privacy-and-terms-of-use/}{MIT OCW 使用条款}. 核查日期为2026-10-08.
\clearpage
\phantomsection
\addcontentsline{toc}{chapter}{原课程信息与课历}
\markboth{原课程信息与课历}{}
\noindent{\color{structurecolor}\bfseries 原课程信息与课历}\par

\noindent\textbf{课程身份与先修.} MIT 18.03 Differential Equations, Spring 2010, 为数学系本科课程, 原课教师为 Haynes Miller 与 Arthur Mattuck. 每周三次一小时讲课、两次一小时习题课. 先修18.01一元微积分, 同修18.02多元微积分; 允许与18.03同时学习18.02. 本册的矩阵记号与必要计算在正文中说明.

\noindent\textbf{原课程教材.} C. Edwards 与 D. Penney, \textit{Elementary Differential Equations with Boundary Value Problems}, 第6版, Prentice Hall, 2003, ISBN 9780136006138; 大纲亦允许第5版, ISBN 9780131457744. 原课还使用 Arthur Mattuck 的 \textit{18.03 Notes and Exercises} 和 Haynes Miller 的 \textit{18.03 Supplementary Notes}. 后文的 EP 表示该商业教材, Notes 表示 Mattuck 讲义中的题号, SN 表示 Miller 补充讲义. 商业教材的题号定位不等于公开题面, 本册没有按猜测复原它们.

\noindent\textbf{原课程大纲.} 原课学习常微分方程及物理系统建模: 一阶方程的解析、图解与数值方法; 高阶线性方程, 尤其常系数二阶方程; 待定系数与常数变易; 复数、复指数以及正弦和指数输入; 振动、阻尼、共振与频率响应; Fourier 级数及周期解; 阶跃、脉冲、卷积与 Laplace 方法; 矩阵、特征值、特征向量与一阶系统; 非线性自治系统的平衡分析及相平面. 本册第1--13章沿数学依赖重组该主线, 并补入已注明的理论内容.

\noindent\textbf{教学组织.} 原课堂通过构造、求解和解释方程训练主动参与; 学生应准备积极参加课堂. 第一次习题课发给投票卡片, 此后每次讲课携带, 用来对课堂问题投票; 意见不一致时讨论, 课末可有一分钟短反馈. 习题课以小组讨论、解题和答疑为主, 鼓励尽早、经常提问, 助教另有办公时间. 原课另设由有经验的本科生值班的辅导室, 可在此完成作业, 考前增派人员. 互动数学小程序 Mathlets 在课堂及每次作业中使用; 下文保留原题所要求的观察任务, 以解析论证和本册独立图形说明数学结论, 不伪称运行过2010年的 Java 小程序.

\noindent\textbf{作业与考核.} 九次作业均分两部分: 第一部分指定教材或 Notes 题目, 第二部分提供更具综合性的公开题面. 应随讲次练习, 而非截止日前一次完成. 原作业说明允许合作讨论, 要求各自独立写出答案并在首页列合作者姓名; 因截止后立即公布答案, 原说明不接受延期. 常规作业按100原始分计, 四讲各24分, 整洁与表达4分; 作业6、7对应三讲, 每讲32分, 表达4分. 学期总评采用下表的885分权重, 与单次作业的100分原始尺度区分.
\begin{center}
\begin{tabular}{lr}\toprule 项目 & 总评点数 \\\midrule
九次作业 & 225 \\
三次一小时期中考试 & 300 \\
一次三小时综合期末考试 & 360 \\\midrule
合计 & 885 \\\bottomrule
\end{tabular}
\end{center}
当期期中试卷题头日期分别为2010年2月24日、3月17日、4月23日. 当期期末题头为2010年5月18日9:00--12:00. 期末没有公开官方答案. 官网练习期末题头写2010, 所链接答案题头写 Spring 2008; 两者若有内容差异, 本册在相应解答中分别说明, 不仅凭链接位置配对.

\noindent\textbf{原课十项核心技能.}
原大纲要求学生个人熟练掌握这些技能, 因其会在其他 MIT 课程中使用; 技能清单向以18.03为先修的任课教师广泛提供. 大纲当时记载140门课程将18.03列为先修或同修, 此数字是2010课程大纲的历史陈述, 不是当前课程数量统计.
\begin{enumerate}
\item 建模得到一阶方程, 用方向场与等斜线理解解, 用 Euler 方法近似.
\item 用积分因子或常数变易求解一阶线性方程.
\item 熟练运算复数和复指数.
\item 求解输入为指数乘多项式的常系数二阶初值问题; 正弦输入时求增益及相位差.
\item 计算 Fourier 系数, 以级数求线性方程的周期解.
\item 用脉冲描述突变, 求单位脉冲响应, 以卷积表示一般输入响应.
\item 使用 Laplace 表求权函数或单位脉冲响应及线性初值解, 联系极点、阻尼与频率响应.
\item 求特征值、特征向量与矩阵指数, 解一阶系统并联系高阶方程.
\item 从迹与行列式重建二维线性自治系统相图.
\item 通过平衡点附近分析判断二维非线性自治系统的定性行为.
\end{enumerate}

\noindent\textbf{公开题面的编排与解答身份.} 正文之后收录七个 Notes 题库单元的217个主编号、九次作业的公开自带题面、22次有公开题面的习题课、三次当期期中与一期末及四套练习考. 第一部分引用 Notes 者按题号及指定小问回指题库, 同一题不再算一个新题; EP 中未公开题面者只保留中文作业指定和准确定位. 每题保留原编号和原条件, 解答标明官方答案译审、译审并补证或编者补解. 编者补解为 AI 辅助形成的独立解答, 与官方答案分开标注, 由另一独立第二读者核算; 同模型实例的交叉审阅不冒充外部专家审稿. 原题或原答的可证实错漏另注明, 不静默改题. 末级小问、出处、重复回指及补解状态见随书逐题清单.

\noindent\textbf{原课程课历.} 以下完整保留官网的讲课 L、习题课 R 顺序、主题、首次引入的技能及作业发布/提交节点. 官网课历没有逐行给出公历日期; 表中不补造日期. 期中I、III的官网行未标 L 编号; L19明确是期中II. R6、R12、R20与R26是复习节点, 无单独公开习题文件; 其余22份题面实际存在, 包括 R23.
\begin{longtable}{p{.065\textwidth}p{.22\textwidth}p{.47\textwidth}p{.105\textwidth}}
\toprule 节点 & 主题 & 首次引入的技能与概念 & 作业与考试 \\\midrule\endhead
\multicolumn{4}{l}{\textbf{一、一阶微分方程}}\\
R1 & 自然增长与分离变量 & 指数增长及捕获建模; 增长率; 分离变量; 通解与特解; 合并积分常数; $\ln|y|$ 及消去; 补回损失解; 用常数满足初值 & -- \\
L1 & 方向场与存在唯一性 & 方向场、积分曲线、等斜线、漏斗; 隐式解; 导数无穷导致不能延续 & 发布作业1 \\
R2 & 方向场与分界曲线 & 分离轨线及解的极值 & -- \\
L2 & 数值方法 & Euler 方法 & -- \\
L3 & 线性方程与建模 & 一阶线性方程; 系统与信号视角; 银行账户; RC 电路; 常输入的分离变量解 & -- \\
R3 & Euler 方法及线性模型 & 混盐问题 & -- \\
L4 & 积分因子 & 齐次方程与零输入; 积分因子; 暂态; 扩散模型及耦合常数 & -- \\
R4 & 一阶线性方程 & 正弦输入 & -- \\
L5 & 复数与单位根 & 复数; 单位根 & 提交1; 发布2 \\
L6 & 复指数与正弦函数 & 复指数; 正弦的振幅、角频率及相位滞后 & -- \\
L7 & 指数及正弦输入的线性响应 & 一阶线性系统的指数/正弦响应; 正弦输入的复方程; 作业中的半衰期 & -- \\
R5 & 复数与复指数 & 原栏为空 & -- \\
L8 & 自治方程、相线与稳定性 & 自治方程; 相线; 稳定性; $e^{k(t-t_0)}$ 与 $ce^{kt}$ & 提交2; 发布3 \\
L9 & 线性与非线性 & 解不能延续的情形 & -- \\
R6 & 期中I复习 & 原栏为空 & 无独立题面 \\
-- & 期中I & 原栏为空 & 一小时考试 \\
\multicolumn{4}{l}{\textbf{二、二阶线性方程}}\\
R7 & 二阶方程的解 & 简谐振子; 初值; 齐次叠加 & -- \\
L11 & 模态与特征多项式 & 弹簧、质量及阻尼器; 一般二阶线性方程; 特征多项式; 实根解 & -- \\
L12 & 振动与阻尼 & 复根; 欠/过/临界阻尼; 复数替换与实解提取; 暂态; 根图 & -- \\
R8 & 二阶常系数齐次方程 & 一般正弦响应; 归一化解 & -- \\
L13 & 指数响应与弹簧驱动 & 受迫系统; 叠加; 指数响应公式; 复数替换; 正弦输入的正弦响应 & -- \\
R9 & 指数与正弦输入 & 原栏为空 & -- \\
L14 & 复增益与阻尼器驱动 & 增益、相位滞后及复增益 & 提交3; 发布4 \\
L15 & 算子、待定系数与共振 & 微分算子; 共振; 待定系数 & -- \\
R10 & 增益、共振与待定系数 & 原栏为空 & -- \\
L16 & 频率响应 & 频率响应 & -- \\
R11 & 频率响应 & 一阶频率响应 & -- \\
L17 & LTI 系统与 RLC 电路 & RLC 电路; 时间平移不变性 & 提交4; 发布5 \\
L18 & 工程应用 & 阻尼比 & -- \\
R12 & 期中II复习 & 原栏为空 & 无独立题面 \\
L19 & 期中II & 原栏为空 & 一小时考试 \\
\multicolumn{4}{l}{\textbf{三、Fourier 级数及变换响应}}\\
R13 & Fourier 级数入门 & 周期函数 & -- \\
L20 & Fourier 级数 & Fourier 级数; 正交性; Fourier 积分 & -- \\
L21 & Fourier 级数的运算 & 方波; 分段连续; 三角恒等式、线性组合及平移 & -- \\
R14 & Fourier 级数 & 不同周期 & -- \\
L22 & 周期解与共振 & 级数求导与积分; 谐波响应; Fourier 级数的振幅--相位表达 & -- \\
R15 & Fourier 谐波响应 & 原栏为空 & -- \\
L23 & 阶跃及脉冲 & 阶跃; $\delta$; 正则函数与奇异函数; 广义函数及导数 & 提交5; 发布6 \\
L24 & 阶跃响应与脉冲响应 & 单位阶跃/脉冲响应; 静止初值; 一阶、二阶响应 & -- \\
R16 & 阶跃、脉冲及其响应 & 原栏为空 & -- \\
L25 & 卷积 & 脉冲后的初值; 时间不变与 $D$、时间平移的交换; 卷积乘积; 初值响应 $w*q$ & -- \\
R17 & 卷积 & $\delta$ 为卷积单位元 & -- \\
L26 & Laplace 基本性质 & Laplace 变换; 收敛域; $\Lap[t^n]$; $s$ 平移; 正弦、余弦变换; 时域与变换域 & 提交6; 发布7 \\
L27 & Laplace 解方程 & $\Lap[\delta]$; 时间导数规则; 逆变换; 部分分式与遮盖法; 一阶非静止初值 & -- \\
R18 & Laplace 变换 & 用变换求单位阶跃响应 & -- \\
L28 & 二阶方程与配方 & $s$ 求导规则; 二阶方程 & -- \\
R19 & Laplace 方法II & 原栏为空 & -- \\
L29 & 极点图 & 权函数、传递函数及其变换关系; 时间平移; 极点及长期行为 & 提交7; 发布8 \\
L30 & 传递函数与频率响应 & 稳定性; 传递函数与增益 & -- \\
R20 & 期中III复习 & 原栏为空 & 无独立题面 \\
-- & 期中III & 原栏为空 & 一小时考试 \\
\multicolumn{4}{l}{\textbf{四、一阶方程组}}\\
L32 & 线性系统与矩阵 & 一阶线性系统; 消元; 矩阵; 反消元及伴随矩阵 & -- \\
R21 & 一阶线性系统 & 原栏为空 & -- \\
L33 & 特征值与特征向量 & 行列式; 特征值; 特征向量; 初值 & -- \\
R22 & 特征值与特征向量 & 解与轨道的区别 & -- \\
L34 & 复特征值或重根 & 特征值与系数; 复特征值; 重根; 缺陷情形与完整特征基 & 提交8; 发布9 \\
L35 & 线性系统相平面 & 迹--行列式平面; 稳定性 & -- \\
R23 & 线性相图 & 线性相图随参数变形 & -- \\
L36 & 正常模态与矩阵指数 & 矩阵指数; 解耦系统; 指数律 & -- \\
R24 & 矩阵指数 & 非齐次系统, 常输入 & -- \\
L37 & 非线性系统 & 非线性自治系统; 向量场; 相图; 平衡点; 平衡附近线性化; Jacobian 矩阵 & 提交9 \\
L38 & 平衡附近线性化及非线性单摆 & 非线性单摆; 飞机长周期振荡; Tacoma Narrows 桥 & -- \\
R25 & 自治系统 & 捕食--被捕食系统 & -- \\
L39 & 线性方法的局限 & 结构稳定; 极限环; 奇异吸引子 & -- \\
R26 & 课程复习 & 原栏为空 & 无独立题面 \\
-- & 综合期末 & 原栏为空 & 三小时考试 \\
\bottomrule
\end{longtable}


来源: \href{https://ocw.mit.edu/courses/18-03-differential-equations-spring-2010/pages/syllabus/}{原课程大纲}, \href{https://ocw.mit.edu/courses/18-03-differential-equations-spring-2010/pages/calendar/}{原课程课历}, \href{https://ocw.mit.edu/courses/18-03-differential-equations-spring-2010/pages/assignments/}{作业目录}, \href{https://ocw.mit.edu/courses/18-03-differential-equations-spring-2010/pages/recitations/}{习题课目录}, \href{https://ocw.mit.edu/courses/18-03-differential-equations-spring-2010/pages/exams/}{考试目录}. 课程信息仅说明2010年的实际安排, 不是2027年招生考纲的认定.

\clearpage
\originalcontents
\clearpage
\mainmatter
\originalchapter{初值问题与解的基本理论}
\originalsection{方程、状态与积分形式}
质量为 $m$ 的物体受力 $F(t,x,x')$ 时满足 $mx''=F(t,x,x')$. 若只记录位置 $x$, 此刻的位置不足以决定下一刻运动; 再记录速度 $v=x'$, 状态 $(x,v)$ 才满足一阶系统 $x'=v$, $v'=F(t,x,v)/m$. 因而高阶方程与一阶系统使用同一套初值理论, 初值包含各阶导数而非一个任意常数.
\begin{definition}
设 $D\subset\R\times\R^n$ 为开集, $f:D\to\R^n$ 连续. 初值问题是
\begin{equation}\label{eq:ivp}
x'=f(t,x),\qquad x(t_0)=x_0,\qquad (t_0,x_0)\in D.
\end{equation}
其解为包含 $t_0$ 的区间 $I$ 上的 $C^1$ 函数 $x:I\to\R^n$, 满足 $(t,x(t))\in D$ 和上述等式. 解的定义域是答案的一部分. 若任何同初值的解都与它在交集上相等, 称其唯一.
\end{definition}
由微积分基本定理, \eqref{eq:ivp} 等价于 $x(t)=x_0+\int_{t_0}^t f(s,x(s))\dd s$. 积分形式消除了未知导数, 适合构造解和估计两条解之间的距离. 方向场则在平面 $(t,x)$ 的每一点画斜率 $f(t,x)$; 等斜线由 $f(t,x)=c$ 给出, 它通常不是解曲线.
\begin{example}
求 $x'=x^2$, $x(0)=a$ 的最大解区间.
\end{example}
\begin{solution}
$a=0$ 时有恒零解. $a\ne0$ 时分离变量得 $-1/x=t-1/a$, 因而 $x(t)=a/(1-at)$. 对 $a>0$, 包含零的最大区间为 $(-\infty,1/a)$; 对 $a<0$, 为 $(1/a,\infty)$. 分母为零处不可能用一个有限的连续值延续, 尽管右端 $x^2$ 处处光滑. 局部存在与全时存在是两个不同结论.
\end{solution}
\originalsection{存在唯一性与延续}
\begin{theorem}[局部存在唯一性]\label{thm:picard}
若 $f$ 连续, 且在 $(t_0,x_0)$ 的某个闭柱体 $K=[t_0-a,t_0+a]\times\overline B(x_0,b)\subset D$ 上对状态满足一致 Lipschitz 条件 $\|f(t,x)-f(t,y)\|\le L\|x-y\|$, 则初值问题在某个 $[t_0-h,t_0+h]$ 上有唯一解. $f$ 对状态为 $C^1$ 是足够条件.
\end{theorem}
\begin{proof}
取 $M=\max_K\|f\|$ 和 $h>0$ 使 $h\le a$, $Mh\le b$, $q=Lh<1$. 若 $M=0$ 或 $L=0$, 对应限制可以省去. 在闭集 $X=\{u\in C([t_0-h,t_0+h],\R^n):\|u-x_0\|_\infty\le b\}$ 上定义
$Tu(t)=x_0+\int_{t_0}^t f(s,u(s))\dd s$. 有 $T(X)\subset X$ 及 $\|Tu-Tv\|_\infty\le q\|u-v\|_\infty$. 从 $u_0=x_0$ 开始迭代 $u_{k+1}=Tu_k$, 则
$\|u_{k+1}-u_k\|_\infty\le q^k\|u_1-u_0\|_\infty$. 几何级数使 $u_k$ 一致 Cauchy, 完备性给出极限 $u\in X$. $T$ 连续, 所以 $Tu=u$; 积分形式说明 $u$ 是解. 两个不动点之差至多为自身距离的 $q$ 倍, 故相等.

任一同初值解在足够小邻域中落在柱体内, 因而局部重合. 在所选小区间内若要离开柱体, 由积分式有 $\|x(t)-x_0\|\le M|t-t_0|\le b$, 并可在需要时把 $h$ 再缩小使 $Mh<b$, 排除首次离开. 因而此小区间上的解唯一. 仅在初值附近的条件不保证未来沿整条轨道仍唯一; 全域的唯一粘合在下述定理的全域局部 Lipschitz 假设下成立. 当 $D_xf$ 连续时, 缩小柱体后其范数有界, 对线段积分就给出所需 Lipschitz 常数.
\end{proof}
连续性本身不能保证唯一性. 方程 $x'=2\sqrt{|x|}$, $x(0)=0$ 既有零解, 也有对任意 $c\ge0$ 定义的 $x(t)=0$ ($t\le c$), $x(t)=(t-c)^2$ ($t\ge c$). 在连接处两侧导数均为零. 唯一性失效正发生在状态方向不满足局部 Lipschitz 的点.
\begin{theorem}[最大解与延续]\label{thm:continuation}
若 $f$ 在 $D$ 上连续且对状态局部 Lipschitz, 则每个初值有唯一最大解, 定义在某个开区间 $(\alpha,\beta)$. 若 $\beta<\infty$ 且终段图像包含在一个紧集 $K\subset D$ 中, 解可延续越过 $\beta$, 与最大性矛盾.
\end{theorem}
\begin{proof}
把全部同初值的局部解按唯一性粘合, 其区间并集是包含 $t_0$ 的区间, 端点若属于定义域就还能局部延续, 故最大区间是开区间. 在紧集 $K$ 上 $\|f\|\le M$, 因而 $\|x(t)-x(s)\|\le M|t-s|$; 当 $t\uparrow\beta$ 时有有限极限 $x_\beta$. 图像的极限点 $(\beta,x_\beta)$ 属于 $K\subset D$, 从此点构造局部解并粘合. 两段在连接处导数都趋向 $f(\beta,x_\beta)$, 所以得到 $C^1$ 延续.
\end{proof}
特别地, 对定义在全空间的自治方程, 有界轨道不会在有限时间爆破. 若解接近 $D$ 的边界, 即使状态有界也可能不能继续, 所以不能删去紧集包含在 $D$ 内的条件.
\originalsection{连续依赖与误差衔接}
\begin{lemma}[Gronwall 不等式]\label{lem:gronwall}
若连续函数 $u\ge0$ 在 $[0,T]$ 上满足 $u(t)\le A+L\int_0^t u(s)\dd s$, 其中 $A,L\ge0$, 则 $u(t)\le A\ee^{Lt}$.
\end{lemma}
\begin{proof}
令 $v(t)=A+L\int_0^t u(s)\dd s$. 当 $A>0$ 时, $v'\le Lv$, 从而 $(\ee^{-Lt}v)'\le0$, 得 $u\le v\le A\ee^{Lt}$. $A=0$ 时先用任意 $A=\eps>0$ 作上界, 再令 $\eps\downarrow0$.
\end{proof}
若两个解在同一 Lipschitz 区域中, 初值差为 $d$, 右端之差的上界为 $\eta$, 则对 $t\ge t_0$ 有
\begin{equation}\label{eq:dependence}
\|x(t)-\widetilde x(t)\|\le d\ee^{L(t-t_0)}+\frac{\eta}{L}(\ee^{L(t-t_0)}-1).
\end{equation}
$L=0$ 时后一项解释为 $\eta(t-t_0)$. 证明方法是对差的积分式取范数, 再对辅助积分函数乘积分因子. 具体令 $u(t)=\|x(t)-\widetilde x(t)\|$, $v(t)=d+\eta(t-t_0)+L\int_{t_0}^t u(s)\dd s$, 则 $u\le v$, $v'\le\eta+Lv$, 积分后得到所写上界. 这说明固定有限时间内的扰动控制, 没有声称时间趋于无穷时误差仍小.

对于 $C^1$ 自治向量场 $f$, 解映射写为 $\phi_t(x_0)$. 先在一条轨道周围取紧邻域并利用有限次局部存在与连续依赖, 可使充分邻近的初值在同一紧时间段内存在. 在该共同时间段内, 对初值的导数 $Y(t)=D\phi_t(x_0)$ 满足
\begin{equation}\label{eq:variational}
Y'=Df(\phi_t(x_0))Y,\qquad Y(0)=I.
\end{equation}
为说明此处微分不是形式运算, 令初值增量为 $h$, 用 \eqref{eq:dependence} 得轨道增量为 $O(\|h\|)$. 连续导数的积分 Taylor 公式将向量场增量写成 $Df(\phi_t)\Delta x+r$, 其中紧时间段上 $\sup\|r\|/\|h\|\to0$. 从轨道增量的积分式减去线性系统 $Yh$ 的积分式, 再用引理 \ref{lem:gronwall}, 得余项一致为 $o(\|h\|)$. 这就证明了导数及变分方程.

Euler 等数值方法用离散点近似积分轨道. 若一条可微近似轨道满足 $\widetilde x'=f(t,\widetilde x)+r(t)$, 则 \eqref{eq:dependence} 直接给出残差到全局误差的转换. 阶数、刚性和步长稳定域的完整讨论属于数值分析册; 在此不由图像吻合替代误差证明.
\originalsection{习题与解答}
\begin{exercise}求 $x'=-x^3$, $x(0)=a$ 的最大解, 并比较正向和反向延续.\end{exercise}
\begin{solution}若 $a=0$ 则解恒零. 否则 $x=a/\sqrt{1+2a^2t}$, 最大区间为 $(-1/(2a^2),\infty)$. 正向趋于零而反向有限时间爆破.\end{solution}
\begin{exercise}设 $|f(t,x)|\le A+B|x|$ 在 $\R^2$ 上成立, 且满足局部存在唯一性条件. 证明每个解全时存在.\end{exercise}
\begin{solution}对任意有限 $T$, 从积分形式和 Gronwall 得 $|x(t)|\le(|x_0|+AT)\ee^{BT}$ ($|t-t_0|\le T$), 反向时间作相同估计. 有限端点前图像位于全空间的紧柱体内, 定理 \ref{thm:continuation} 排除有限最大端点.\end{solution}

\originalchapter{一阶方程的解析方法}
\originalsection{分离变量与变量代换}
对 $y'=a(t)b(y)$, 在 $b(y)\ne0$ 的区间上, 选择原函数 $B'(y)=1/b(y)$, 则 $B(y(t))=\int a(t)\dd t+C$. 隐函数能否在初值附近解为 $y(t)$, 由 $B'(y_0)\ne0$ 保证. 先除以 $b(y)$ 会删掉它的零点所对应的常值解; 必须先登记这些解.
\begin{example}求 $y'=y(1-y)$ 的所有初值解 $y(0)=y_0$.\end{example}
\begin{solution}
$y_0=0,1$ 时分别为常值解. 其他情形由 $1/[y(1-y)]=1/y+1/(1-y)$ 得 $\log|y/(1-y)|=t+C$, 从而
\[
y(t)=\frac{y_0\ee^t}{1-y_0+y_0\ee^t}.
\]
取包含零且分母非零的最大区间. $0<y_0<1$ 时全时存在, 正向趋于1; $y_0>1$ 时反向有限时间爆破; $y_0<0$ 时正向有限时间趋于负无穷. 有理式形式也覆盖两个常值解, 但这是代回原方程后确认的, 不是除法步骤自动保留的.
\end{solution}
对于齐次型 $y'=F(y/t)$, $t\ne0$ 时令 $v=y/t$, 得 $t v'=F(v)-v$. 这里“齐次”指右端的比值结构, 不等于线性齐次方程. Bernoulli 方程 $y'+p(t)y=q(t)y^\gamma$, $\gamma\ne1$, 在幂函数有定义且 $y\ne0$ 的区间上令 $z=y^{1-\gamma}$, 变为 $z'+(1-\gamma)pz=(1-\gamma)q$. 每次变换都应记下逆变换的符号与定义域.
\begin{example}解 $y'+y=t y^2$, $y(0)=1$.\end{example}
\begin{solution}
令 $z=1/y$, 得 $z'-z=-t$. 一个特解为 $t+1$, 故 $z=t+1+C\ee^t$. 初值得 $C=0$, 所以 $y=1/(t+1)$, 最大区间为 $(-1,\infty)$. 原方程另有零解, 此初值不选它. 当 $z$ 穿过零时不能使用 $y=1/z$ 延续.
\end{solution}
\originalsection{一阶线性方程}
\begin{theorem}[积分因子]\label{thm:firstlinear}
设 $p,q$ 在区间 $I$ 上连续, $t_0\in I$. 问题 $y'+p(t)y=q(t)$, $y(t_0)=y_0$ 的唯一解为
\begin{equation}\label{eq:firstlinear}
y(t)=\ee^{-P(t)}\left[y_0+\int_{t_0}^t\ee^{P(s)}q(s)\dd s\right],\qquad P(t)=\int_{t_0}^t p(s)\dd s.
\end{equation}
\end{theorem}
\begin{proof}
选取 $\mu$ 的目的在于把左端变成 $(\mu y)'$. 比较系数须满足 $\mu'=p\mu$, 可以取 $\mu=\ee^P$, 它处处非零. 乘以 $\mu$ 后积分, 使用 $P(t_0)=0$, 得公式. 微分公式核验方程及初值. 两个解之差满足 $(\mu z)'=0$, 零初值迫使 $z=0$.
\end{proof}
公式中的齐次项描述初值的遗留, 积分项累积外部输入. 若 $p=k>0$, 记输入为 $q(t)$, 则 $y(t)=\ee^{-k(t-t_0)}y_0+\int_{t_0}^t\ee^{-k(t-s)}q(s)\dd s$. 较早输入乘的权重较小, 这是后面脉冲响应的第一种实例.
\begin{example}
容器中有 $V$ 升充分混合的溶液. 每分钟流入并流出 $r$ 升, 输入浓度为 $c(t)$. 写出并求解盐量 $M(t)$ 的方程.
\end{example}
\begin{solution}
盐流入率为 $rc(t)$, 流出浓度为 $M/V$, 故 $M'+(r/V)M=rc(t)$. 当 $c(t)=c_0$ 为常数时, $M(t)=Vc_0+(M(0)-Vc_0)\ee^{-rt/V}$. 体积不变和瞬时充分混合是这个模型的条件; 若流入流出量不同, 应先算 $V(t)$, 不能照抄常系数模型.
\end{solution}
\originalsection{恰当方程与隐式解}
方程 $M(t,y)\dd t+N(t,y)\dd y=0$ 沿曲线表示 $M+Ny'=0$. 若存在势函数 $\Psi$ 使 $\Psi_t=M$, $\Psi_y=N$, 则解落在等值线 $\Psi=C$ 上. 反过来只有在 $N\ne0$ 的点附近才能从等值线得到图像形式的解.
\begin{proposition}
若 $M,N\in C^1$ 定义在开矩形上, 则存在上述势函数当且仅当 $M_y=N_t$.
\end{proposition}
\begin{proof}
必要性来自二阶混合偏导相等. 充分性可直接构造
$\Psi(t,y)=\int_{t_*}^{t}M(s,y_*)\dd s+\int_{y_*}^{y}N(t,v)\dd v$.
对 $y$ 求导得 $N$. 对 $t$ 求导, 用 $N_t=M_y$, 得 $M(t,y_*)+\int_{y_*}^{y}M_y(t,v)\dd v=M(t,y)$. 矩形条件保证积分路径始终在定义域中.
\end{proof}
一般区域中局部相容不自动给全局势函数. 本册只使用矩形或直接构造的势函数, 不把拓扑条件隐藏在符号计算里.
\begin{example}解 $(2ty+1)\dd t+(t^2+2y)\dd y=0$, 过点 $(0,1)$.\end{example}
\begin{solution}
$M_y=2t=N_t$. 对 $t$ 积分得 $\Psi=t^2y+t+g(y)$, 与 $\Psi_y=t^2+g'(y)=t^2+2y$ 比较得 $g=y^2$. 初值给 $C=1$. 初值附近选根
$y=(-t^2+\sqrt{t^4-4t+4})/2$; 这条支满足 $y(0)=1$. 被开方数 $t^4-4t+4=1+(t-1)^2(t^2+2t+3)>0$, 且所选支上 $N=t^2+2y=\sqrt{t^4-4t+4}>0$. 因而该支实际在整个 $\R$ 上光滑, 是最大解; 另一个代数根不满足此初值.
\end{solution}
若恰当性失败, 可以尝试积分因子. 只依赖 $t$ 的因子 $\mu$ 必须满足 $(\mu M)_y=(\mu N)_t$, 即 $\mu'/\mu=(M_y-N_t)/N$. 右端若只依赖 $t$ 且 $N\ne0$, 积分即可构造 $\mu$. 只依赖 $y$ 的情形相应为 $\mu_y/\mu=(N_t-M_y)/M$.

\originalsection{Riccati 方程与已知特解}
Riccati 方程 $y'=a(t)y^2+b(t)y+c(t)$ 的二次项阻止线性叠加, 但若已知特解 $y_1$, 可把这一特解作为变量代换的基点. 在 $y\ne y_1$ 的区间令 $y=y_1+1/v$. 消去特解方程得 $v'+(2ay_1+b)v=-a$, 是线性方程. $y=y_1$ 这一分支必须另记; $v$ 的零点对应原变量爆破.
\begin{example}求 $y'=y^2-2/t^2$, $t>0$, $y(1)=0$ 的解.\end{example}
\begin{solution}
试幂函数 $y_1=k/t$, 得 $-k=k^2-2$, 可取 $k=1$. 令 $y=1/t+1/v$, 得 $v'+2v/t=-1$, 积分因子为 $t^2$. 因此 $v=-t/3+C/t^2$. 初值要求 $v(1)=-1$, 得 $C=-2/3$, 所以
\[
y(t)=\frac1t-\frac{3t^2}{t^3+2}=\frac{2(1-t^3)}{t(t^3+2)}.
\]
在 $(0,\infty)$ 分母不为零, 这是指定定义域内的最大解. 第一次代换的构造方向来自二次差 $y^2-y_1^2=2y_1(y-y_1)+(y-y_1)^2$: 倒数恰能把剩下的平方项化为常数输入.
\end{solution}

\originalsection{习题与解答}
\begin{exercise}求 $ty'=y+t\ee^{-y/t}$, $t>0$, $y(1)=0$ 的解.\end{exercise}
\begin{solution}令 $v=y/t$, 得 $tv'=\ee^{-v}$, 所以 $(\ee^v)'=1/t$. 初值给 $\ee^v=1+\log t$, 故 $y=t\log(1+\log t)$, 最大区间为 $(\ee^{-1},\infty)$.\end{solution}
\begin{exercise}解 $y'+2ty=2t$, $y(0)=a$, 并说明不同初值解的差.\end{exercise}
\begin{solution}积分因子为 $\ee^{t^2}$, 得 $y=1+(a-1)\ee^{-t^2}$. 两个解之差为 $(a-b)\ee^{-t^2}$, 正反两方向均趋于零, 因为此方程并非自治方程, 不能从此推出自治流的平衡点同时正反向吸引.\end{solution}

\originalchapter{自治方程与一维相线}
\originalsection{平衡点与单调运动}
对于 $x'=f(x)$, 时间不显含于右端. 一条解的时间平移仍为解; 唯一性给出流的复合规律 $\phi_{t+s}(x)=\phi_t(\phi_s(x))$, 只要各段都有定义. 平衡点 $a$ 满足 $f(a)=0$, 对应恒定轨道. 一维非平衡轨道在 $f$ 符号不变的区间内严格单调.
\begin{proposition}
设 $f$ 局部 Lipschitz. 若相邻零点为 $a<b$, 且 $(a,b)$ 内 $f>0$, 则从其中出发的解正向趋向 $b$, 反向趋向 $a$, 并且全时存在. 若 $f<0$, 方向相反.
\end{proposition}
\begin{proof}
唯一性使轨道不能在有限时刻穿越平衡点: 若到达 $b$, 从到达时刻反向唯一解是恒 $b$, 与初值矛盾. 轨道被限制在 $[a,b]$, 延续定理给出全时存在. 单调有界性给正向极限 $\ell$. 若 $\ell<b$, 连续性使 $f$ 在 $\ell$ 附近有正下界 $c$, 最终 $x'\ge c$ 与收敛矛盾. 反向同理.
\end{proof}
\begin{definition}
自治系统的平衡点 $a$ 若满足: 对任意 $\eps>0$, 存在 $\delta>0$, 使 $\|x(0)-a\|<\delta$ 的解对所有 $t\ge0$ 存在并满足 $\|x(t)-a\|<\eps$, 则称 Lyapunov 稳定. 若另有一个邻域内的解均趋于 $a$, 称渐近稳定. 吸引只要求趋于 $a$, 不应代替稳定的量词.
\end{definition}
\begin{theorem}[一维判别]
设 $f\in C^1$, $f(a)=0$. 若 $f'(a)<0$, 则 $a$ 渐近稳定; 若 $f'(a)>0$, 则不稳定. 当 $f'(a)=0$ 时必须考察两侧符号.
\end{theorem}
\begin{proof}
$f'(a)<0$ 时, 足够小邻域内有 $(x-a)f(x)<0$. 两侧轨道都向 $a$ 运动, 因而不会离开任意足够小的初始区间. 单调极限论证说明其极限为 $a$, 这同时证明稳定与吸引. 若导数为正, 选固定小边界 $a+\eps$; 任意从 $(a,a+\eps)$ 出发的解都远离 $a$, 且必到达边界, 否则极限会是内部零点. 任意小扰动可逃出固定邻域, 所以不稳定.
\end{proof}
$x'=-x^3$ 中原点渐近稳定, $x'=x^3$ 中原点不稳定; $x'=x^2$ 中左侧吸引而右侧排斥, 仍不稳定. 三者线性化都为 $x'=0$. 把零特征值当作“稳定”会混淆这些行为.
\begin{center}
\begin{tikzpicture}[>=Stealth,x=1.4cm]
\draw[->] (-2.4,0)--(2.4,0) node[right]{$x$};
\fill (-1,0) circle(2pt) node[below=5pt]{$0$}; \fill (1,0) circle(2pt) node[below=5pt]{$1$};
\draw[->,thick](-1.5,0)--(-2,0);\draw[->,thick](-.5,0)--(.3,0);\draw[->,thick](2,0)--(1.5,0);
\node[above=8pt] at (-1,0){排斥};\node[above=8pt] at (1,0){吸引};
\end{tikzpicture}

$ x'=x(1-x)$ 的相线. 箭头表示时间增加方向.
\end{center}
在负半轴上右端为负, 因而轨道向左; 相线的每支箭头都由对应区间的符号确定.
\originalsection{参数、阈值与分岔}
参数改变时, 平衡点的数目或稳定性发生变化称为分岔. 三个基本模型分别是
\[
 x'=\mu-x^2,\qquad x'=\mu x-x^2,\qquad x'=\mu x-x^3.
\]
第一式在 $\mu<0$ 时无平衡点, $\mu>0$ 时有 $\pm\sqrt\mu$; 正根吸引、负根排斥, 是鞍结分岔. $\mu=0$ 时半稳定. 第二式有 $0,\mu$ 两条平衡分支, 导数分别为 $\mu,-\mu$, 穿越零时交换稳定性, 是跨临界分岔. 第三式在 $\mu\le0$ 时只有零点, $\mu>0$ 时又生出 $\pm\sqrt\mu$, 两个新点吸引而原点排斥, 是超临界叉形分岔. 对称性 $f(-x)=-f(x)$ 解释两条新分支成对出现; 一般扰动破坏对称性后不能期望同一分岔图.
\begin{example}
分析带恒定捕获的种群模型 $x'=rx(1-x/K)-h$, 其中 $r,K>0$, $h\ge0$, 生物状态取 $x\ge0$.
\end{example}
\begin{solution}
令 $u=x/K$, $\tau=rt$, $H=h/(rK)$, 得 $\dd u/\dd\tau=u(1-u)-H$. 平衡点为 $u_\pm=(1\pm\sqrt{1-4H})/2$, 存在条件为 $H\le1/4$. $H<1/4$ 时小根不稳定, 大根渐近稳定. 当 $H=1/4$ 时, 方程成为 $u'=-(u-1/2)^2$. 从上方出发的解渐近趋于 $1/2$, 从下方出发则向零下降, 所以临界点半稳定而非 Lyapunov 稳定. 对 $H<1/4$, 超过小根的初值趋向大根, 低于小根的初值向零减少. $H>1/4$ 时所有非负状态都下降, 因为右端至多为 $1/4-H<0$. 此模型在零处仍给负导数, 不能自动把负种群当作真实状态; 生物解释应在首次到达零时停止, 或另行设定捕获关闭机制. 数学方程的延续与物理模型的适用区间必须区分.
\end{solution}
\originalsection{比较与标量不等式}
\begin{proposition}[比较原理]
设 $f$ 在 $\R^2$ 上连续, 并对状态在每个紧柱体上有一致的 Lipschitz 常数, $u$ 为 $C^1$ 函数, $u'\le f(t,u)$, 而 $v'=f(t,v)$, $u(t_0)\le v(t_0)$. 在共同存在的正向区间内有 $u\le v$.
\end{proposition}
\begin{proof}
在任意共同紧时间段上取覆盖两图像的有限 Lipschitz 常数 $L$. 令 $w=(u-v)_+$, 则它绝对连续, 几乎处处满足 $w'\le Lw$, 初值为零. 在 $u-v>0$ 处由 Lipschitz 得该不等式, 在 $u-v<0$ 处 $w'=0$, 在零点几乎处处正部导数亦满足它. 乘 $\ee^{-L(t-t_0)}$ 或应用 Gronwall 得 $w=0$. 紧时间段任意, 故结论成立.
\end{proof}
这个结果可用来夹住不能显式积分的方程. 例如 $x'=-x+\sin t$, $|\sin t|\le1$, 与 $z'=-z\pm1$ 比较得到最终有界区间; 无须知道特殊解的相位.

\originalsection{流行病模型的阈值}
将总人口归一化为1, 用 $S,I,R$ 表示易感、感染和移出人群的比例. 假设接触感染率为 $\beta SI$, 感染者以速率 $\gamma I$ 移出, $\beta,\gamma>0$, 得 SIR 系统
\[
S'=-\beta SI,\qquad I'=\beta SI-\gamma I,\qquad R'=\gamma I.
\]
总和 $S+I+R$ 守恒, 各边界的向量场不向负区域, 因而非负单纯形正向不变. 对 $S_0>0,I_0>0$, 指数积分式使 $S,I$ 在有限时间严格正, 且轨道在紧集内全正向存在. $S$ 单调下降, $R$ 单调上升; $I$ 可以先升后降, 判据是 $S$ 是否大于 $\gamma/\beta$.
\begin{proposition}
若 $S_0>0,I_0>0$, 则 $I(t)\to0$, $S(t)\to S_\infty>0$, 且
\begin{equation}\label{eq:sirfinal}
\log\frac{S_\infty}{S_0}=-\frac\beta\gamma(S_0+I_0-S_\infty).
\end{equation}
若 $S_0>\gamma/\beta$, 感染比例先上升, 在 $S=\gamma/\beta$ 时取得唯一峰值, 再趋零.
\end{proposition}
\begin{proof}
$R'=\gamma I$ 给 $\int_0^\infty I\dd t=(R_\infty-R_0)/\gamma<\infty$. 因右端在单纯形有界, $I'$ 有界, 故 $I$ 一致连续. 非负一致连续可积函数趋零: 否则有无限次取值至少为 $\eps$, 一致连续性给每次周围固定长度区间上至少为 $\eps/2$, 可选不相交子列, 与积分有限矛盾. 从 $(\log S)'=-\beta I$ 得 $S_\infty=S_0\exp[-\beta\int_0^\infty I\dd t]>0$. 将 $R_\infty=1-S_\infty$ 和 $R_0=1-S_0-I_0$ 代入得到 \eqref{eq:sirfinal}.

若 $S_0>\gamma/\beta$ 而永不降到阈值, 则 $I$ 一直不减且至少为 $I_0>0$, 与积分有限矛盾. $S$ 严格下降, 所以只穿过阈值一次. 穿过前 $I'>0$, 穿过后 $I'<0$, 得唯一峰值.
\end{proof}
消去时间也能得到不变量 $I+S-(\gamma/\beta)\log S$. 若初值超过阈值, 记 $s_c=\gamma/\beta$, 峰值为
$I_{\max}=I_0+S_0-s_c-s_c\log(S_0/s_c)$.
阈值量 $\beta S_0/\gamma$ 比较初始感染增长与移出. 模型不包含出生、死亡、再次易感或时变接触率, 因而这些变化不能仅靠替换一个数字解释. 这里把 Miller 附录A中的流行病主题补成一个闭合的动力学例子.

\originalsection{习题与解答}
\begin{exercise}判定 $x'=x(x-1)(2-x)$ 的全部平衡点稳定性及区间 $(0,1)$ 的轨道极限.\end{exercise}
\begin{solution}符号从左至右依次为正、负、正、负, 所以 $0,2$ 渐近稳定, $1$ 不稳定. $(0,1)$ 内解正向趋于0, 反向趋于1.\end{solution}
\begin{exercise}分析 $x'=\mu x+x^3$ 的平衡点, 说明与超临界叉形模型的区别.\end{exercise}
\begin{solution}$\mu<0$ 时零点吸引, 两个非零点 $\pm\sqrt{-\mu}$ 的导数为 $-2\mu>0$, 均排斥. $\mu>0$ 时只剩不稳定的零点; $\mu=0$ 时 $x'=x^3$ 仍不稳定. 非零不稳定分支位于参数负侧, 是亚临界叉形的基本模型.\end{solution}

\originalchapter{高阶线性方程与解空间}
\originalsection{线性结构与基本解}
正规 $n$ 阶线性方程为
\begin{equation}\label{eq:nlinear}
y^{(n)}+a_{n-1}(t)y^{(n-1)}+\cdots+a_0(t)y=g(t).
\end{equation}
假设各系数和 $g$ 在区间 $I$ 上连续. 令状态 $X=(y,y',\ldots,y^{(n-1)})^{\mathsf T}$, 可把方程改成连续系数的一阶线性系统. 此右端在状态上为线性, 在每个紧时间段上满足线性增长界; 第1章理论给出整个 $I$ 上的唯一初值解. 若最高阶系数原来不是1, 须先除以它, 其零点处正规化失效, 本节定理不能直接用于跨越该点; 个别原方程的解仍可能光滑跨过, 如 $ty'=y$ 的解 $y=Ct$.
\begin{theorem}\label{thm:space}
\eqref{eq:nlinear} 的齐次解组成 $n$ 维实向量空间. 在 $t_0$ 处规定 $n$ 个初值的映射是此空间到 $\R^n$ 的线性同构. 任一非齐次解等于一个特解加一个齐次解.
\end{theorem}
\begin{proof}
线性算子使齐次解空间对线性组合封闭. 初值映射线性; 唯一性给单射, 存在性给满射. 规定初值为各标准基向量所得的 $n$ 个解构成基. 若 $y_p$ 是特解, 两解之差满足齐次方程; 反之 $y_p+y_h$ 满足同一非齐次方程.
\end{proof}
“通解有 $n$ 个常数”在此来自初值同构, 不是数积分次数的口号. 若给出的 $n$ 个函数初值向量不独立, 即使带了 $n$ 个常数, 也不能覆盖所有解.
\originalsection{Wronskian 与降阶}
\begin{definition}
函数 $y_1,\ldots,y_n$ 的 Wronskian 是 $W(t)=\det(y_j^{(i-1)}(t))_{1\le i,j\le n}$.
\end{definition}
\begin{theorem}[Abel 公式]
若这些函数解同一齐次方程 \eqref{eq:nlinear}, 则
\[
W(t)=W(t_0)\exp\left(-\int_{t_0}^t a_{n-1}(s)\dd s\right).
\]
因此它们为基本解组当且仅当某一个点的 Wronskian 非零.
\end{theorem}
\begin{proof}
对行列式按行微分. 前 $n-1$ 行的导数等于下一行, 对应行列式含两条相同的行, 贡献为零. 最后一行导数按方程展开, 除 $-a_{n-1}$ 乘原末行外, 其他项都与已有行重复, 因而 $W'=-a_{n-1}W$. 应用一阶线性公式. 非零 Wronskian 等价于初值向量独立, 再用定理 \ref{thm:space} 即得基本解组结论.
\end{proof}
对任意可微函数, “Wronskian 恒为零”不能不加条件地当作线性相关判据. 在这里有效的原因是它们解同一个正规线性方程.

对二阶方程 $y''+p y'+q y=0$, 若已知在某区间上非零的解 $y_1$, 令 $y=v y_1$. 消去 $v$ 项得 $v''y_1+v'(2y_1'+py_1)=0$, 即
\[
v'=C\frac{\ee^{-\int p\dd t}}{y_1^2},\qquad
 y_2=y_1\int\frac{\ee^{-\int p\dd t}}{y_1^2}\dd t.
\]
因为 $W(y_1,y_2)=\ee^{-\int p\dd t}\ne0$, 新解与原解独立. 有零点的 $y_1$ 可分区间使用公式, 再用正规方程初值理论延续, 不能直接除以零.
\originalsection{常系数方程}
令 $D=\dd/\dd t$, $P(z)=z^n+a_{n-1}z^{n-1}+\cdots+a_0$. 因 $D\ee^{\lambda t}=\lambda\ee^{\lambda t}$, 有 $P(D)\ee^{\lambda t}=P(\lambda)\ee^{\lambda t}$. 这使特征根给出指数解.
\begin{theorem}
若 $\lambda$ 是 $P$ 的 $m$ 重根, 则 $\ee^{\lambda t},t\ee^{\lambda t},\ldots,t^{m-1}\ee^{\lambda t}$ 是复齐次解. 对所有不同根集合这些函数得到复基本解组. 实系数时, 共轭根 $\alpha\pm\ii\beta$ ($\beta>0$) 对应实解 $t^k\ee^{\alpha t}\cos\beta t$, $t^k\ee^{\alpha t}\sin\beta t$.
\end{theorem}
\begin{proof}
指数移位恒等式为 $P(D)(\ee^{\lambda t}v)=\ee^{\lambda t}P(D+\lambda)v$, 先对 $D$ 用乘积法则, 再对多项式逐项展开. 写 $P(z)=(z-\lambda)^m Q(z)$, 则 $P(D+\lambda)=D^mQ(D+\lambda)$, 所以次数小于 $m$ 的多项式被消去.

证明全部解独立: 若 $\sum_j\ee^{\lambda_j t}p_j(t)=0$, 其中 $\deg p_j<m_j$, 固定 $j$ 并施加 $\prod_{k\ne j}(D-\lambda_k)^{m_k}$. 其余项消失, 剩下 $\ee^{\lambda_jt}\prod_{k\ne j}(D+\lambda_j-\lambda_k)^{m_k}p_j=0$. 每个 $D+c$, $c\ne0$, 在有限维多项式空间上可逆, 因其矩阵三角且对角为 $c$. 因而 $p_j=0$. 总数是 $n$, 由解空间维数得到基本解组. 共轭解取实虚部完成实形式.
\end{proof}
\begin{example}求 $y'''-3y''+3y'-y=0$, $y(0)=1$, $y'(0)=0$, $y''(0)=0$.\end{example}
\begin{solution}
$P(z)=(z-1)^3$, 所以 $y=\ee^t(A+Bt+Ct^2)$. 初值依次给 $A=1$, $A+B=0$, $A+2B+2C=0$, 得 $y=\ee^t(1-t+t^2/2)$. 重根产生的多项式因子不能漏掉.
\end{solution}
Euler 型方程 $t^2y''+at y'+by=0$, $t>0$, 令 $s=\log t$, $u(s)=y(\ee^s)$, 则 $ty'=u'$, $t^2y''=u''-u'$, 得 $u''+(a-1)u'+bu=0$. 因而可用常系数理论处理, 但 $t=0$ 是变换及正规性同时失效的位置.
\originalsection{常数变易与非齐次方程}
对 $y''+py'+qy=g$, 取基本解组 $y_1,y_2$, 设 $y=u_1y_1+u_2y_2$. 为避免冗余选择规定 $u_1'y_1+u_2'y_2=0$. 代回方程得到 $u_1'y_1'+u_2'y_2'=g$, 所以
\[
u_1'=-\frac{y_2g}{W},\qquad u_2'=\frac{y_1g}{W},\qquad
 y_p=-y_1\int_{t_0}^t\frac{y_2g}{W}\dd s+y_2\int_{t_0}^t\frac{y_1g}{W}\dd s.
\]
积分中的 $y_j,g,W$ 均在 $s$ 处取值. 这给零位置、零速度初值的特解; 可直接微分核查. 方法依赖 $W\ne0$, 正是基本解组所保证的条件.
\begin{example}求 $y''+y=\sec t$, $y(0)=y'(0)=0$, 在 $(-\pi/2,\pi/2)$ 上的解.\end{example}
\begin{solution}
取 $y_1=\cos t$, $y_2=\sin t$, $W=1$. 有 $u_1'=-\tan t$, $u_2'=1$, 并取零积分常数. 故 $u_1=\log\cos t$, $u_2=t$, 解为 $y=\cos t\log\cos t+t\sin t$. 正规右端只在指定区间连续, 不用这个表达式声称跨越奇点的解.
\end{solution}
\originalsection{习题与解答}
\begin{exercise}求 $y''+2y'+5y=0$ 的通解并说明长期行为.\end{exercise}
\begin{solution}根为 $-1\pm2\ii$, 故 $y=\ee^{-t}(A\cos2t+B\sin2t)$. 每个解及其导数正向趋于零, 振荡角频率是2, 衰减指数是1.\end{solution}
\begin{exercise}求 $t^2y''-ty'+y=0$, $t>0$, 并计算基本解组的 Wronskian.\end{exercise}
\begin{solution}变换后为 $u''-2u'+u=0$, 得 $y_1=t$, $y_2=t\log t$, $W=t$. 正规方程中 $p=-1/t$, Abel 公式也给 $W=Ct$.\end{solution}

\originalchapter{强迫振动与线性响应}
\originalsection{阻尼振子与能量}
弹簧恢复力 $-kx$, 黏性阻力 $-cx'$, 外力 $F(t)$ 给出 $mx''+cx'+kx=F(t)$, 其中 $m,k>0$, $c\ge0$. 除以质量后写成
\begin{equation}\label{eq:oscillator}
x''+2\zeta\omega_0x'+\omega_0^2x=f(t),\quad
 \omega_0=\sqrt{k/m},\quad \zeta=\frac{c}{2\sqrt{mk}}.
\end{equation}
$\omega_0$ 是无阻尼固有角频率, $\zeta$ 是无量纲阻尼比. 特征根为 $-\zeta\omega_0\pm\omega_0\sqrt{\zeta^2-1}$. $0\le\zeta<1$ 时欠阻尼, $\zeta=1$ 时临界阻尼, $\zeta>1$ 时过阻尼. 临界情形含 $t\ee^{-\omega_0t}$, 不是两个相同指数乘常数.

无外力时取能量 $E=(m(x')^2+kx^2)/2$, 得 $E'=-c(x')^2\le0$. 这验证阻尼的耗散方向. $c=0$ 时能量守恒, 振动一般不衰减; $c>0$ 时特征根实部全负, 所有齐次运动趋于零. 能量单调本身不立即说明其极限为零, 后面将给出不变集论证.
\originalsection{指数输入与共振}
\begin{proposition}[指数响应]
对常系数 $P(D)y=\ee^{\lambda t}$, 若 $P(\lambda)\ne0$, 有特解 $y_p=\ee^{\lambda t}/P(\lambda)$. 若 $\lambda$ 是 $m$ 重根, 则有特解 $y_p=t^m\ee^{\lambda t}/P^{(m)}(\lambda)$.
\end{proposition}
\begin{proof}
第一式直接代入. 重根情形用指数移位恒等式, 令 $P(z)=(z-\lambda)^mQ(z)$. 因为 $D^mQ(D+\lambda)t^m=m!Q(\lambda)=P^{(m)}(\lambda)$, 第二式的左端正好为 $\ee^{\lambda t}$.
\end{proof}
待定系数法把指数、三角函数和多项式的有限维微分封闭空间作为候选. 与齐次解重合时乘足够次的 $t$, 本质就是上述移位公式. 它不适用于任意右端, 此时应使用常数变易或卷积.
\begin{example}求 $y''+y=\cos t$, $y(0)=y'(0)=0$.\end{example}
\begin{solution}
复输入 $\ee^{\ii t}$ 中 $P(z)=z^2+1$, $\ii$ 为简单根, 得复特解 $t\ee^{\ii t}/(2\ii)$. 取实部为 $t\sin t/2$, 已满足两个零初值. 振幅线性增长来自外力频率与固有频率重合, 并非“所有周期输入都使振动变大”.
\end{solution}
\begin{example}设 $\omega\ge0$, $\omega\ne\omega_0$. 求 $x''+\omega_0^2x=F\cos\omega t$, 零初值的解, 并解释接近共振时的拍频.\end{example}
\begin{solution}
特解为 $F\cos\omega t/(\omega_0^2-\omega^2)$, 配合初值得
\[
x=\frac{F(\cos\omega t-\cos\omega_0t)}{\omega_0^2-\omega^2}
=-\frac{2F\sin((\omega+\omega_0)t/2)\sin((\omega-\omega_0)t/2)}{\omega_0^2-\omega^2}.
\]
接近但不等于共振时, 慢因子调制快振动, 产生拍. 令 $\omega\to\omega_0$ 且固定 $t$ 可得共振解 $Ft\sin\omega_0t/(2\omega_0)$. 此极限对所有时间不一致, 不能把非共振有界解误当作共振也有界.
\end{solution}
\originalsection{复增益、相位与频率响应}
设 $P$ 为实系数且 $P(\ii\omega)\ne0$. 输入 $\operatorname{Re}(F\ee^{\ii\omega t})$ 的正弦特解是 $\operatorname{Re}(H(\ii\omega)F\ee^{\ii\omega t})$, 其中 $H(s)=1/P(s)$ 为传递函数. $|H(\ii\omega)|$ 是增益, $\arg H(\ii\omega)$ 是相位变化. 取实部与常系数微分算子交换, 所以使用复数没有改变实方程.

对 \eqref{eq:oscillator},
\[
|H(\ii\omega)|=\frac1{\sqrt{(\omega_0^2-\omega^2)^2+4\zeta^2\omega_0^2\omega^2}}.
\]
相位为 $-\arg(\omega_0^2-\omega^2+2\ii\zeta\omega_0\omega)$, 象限不能由一个不带象限的反正切决定. 阻尼正时瞬态趋于零, 这一特解才是长期稳态. 阻尼为零时, 初值携带的自由振动不消失.

令 $u=\omega^2$, 对增益分母平方求导, 得 $2(u-\omega_0^2)+4\zeta^2\omega_0^2$. 当 $0<\zeta<1/\sqrt2$ 时峰值频率为 $\omega_r=\omega_0\sqrt{1-2\zeta^2}$; 阻尼更大时峰值在零频率处. 因而“最大振幅频率”通常不等于欠阻尼自由振动频率 $\omega_0\sqrt{1-\zeta^2}$.

电路中的串联 RLC 模型以电荷 $q$ 为状态, 电流为 $q'$, 满足 $Lq''+Rq'+q/C=V(t)$. 与机械振子对应的是 $L,m$; $R,c$; $1/C,k$. 输入若通过移动支座或黏性连接产生, 右端可能含输入的导数, 传递函数的分子随之改变, 不能统一写为 $1/P$.
\originalsection{输入输出与叠加}
一个零初始状态的线性时不变系统对输入线性, 并与时间平移相容. “零初始状态”使叠加清楚: 非零初值产生的自由响应应单独加上. 对双输入分别求特解后相加, 是线性算子的性质; 对非线性方程一般不成立.
\begin{example}分析一阶滤波器 $y'+ay=af(t)$, $a>0$ 的频率响应.\end{example}
\begin{solution}
$H(s)=a/(s+a)$, 增益为 $a/\sqrt{a^2+\omega^2}$, 相位为 $-\arctan(\omega/a)$. 慢变化输入几乎原样通过, 高频被压低, 因而是低通滤波器. 输入 $f=1$, 零初值时 $y=1-\ee^{-at}$; 阶跃最终增益 $H(0)=1$ 与频率零极限一致.
\end{solution}

\originalsection{多项式强迫与瞬态识别}
对多项式 $q(t)$, 候选特解可选为相同或更高次数的多项式. 例如 $P(0)\ne0$ 时, $P(D)$ 在次数不超过 $m$ 的多项式空间上矩阵为三角形, 对角全为 $P(0)$, 因而可逆; 次数不超过 $m$ 的输入恰有一个此空间内的特解. $0$ 为特征根时应先乘相应次数的 $t$.
\begin{example}求 $y''+2y'+y=t^2$, $y(0)=y'(0)=0$.\end{example}
\begin{solution}
令 $y_p=At^2+Bt+C$, 比较系数得 $A=1$, $4A+B=0$, $2A+2B+C=0$, 所以 $y_p=t^2-4t+6$. 加上齐次项 $(c_1+c_2t)\ee^{-t}$, 初值给 $c_1=-6$, $c_2=-2$. 故 $y=t^2-4t+6-(6+2t)\ee^{-t}$. 输入非周期, 长期输出也非有界常量, 但自由响应仍衰减.
\end{solution}
在 $0\le\zeta<1$, $A\ne0$ 的欠阻尼自由响应 $x=A\ee^{-\zeta\omega_0t}\cos(\omega_dt-\phi)$ 中, $\omega_d=\omega_0\sqrt{1-\zeta^2}$. 同号相邻峰值相隔 $T_d=2\pi/\omega_d$, 前一个同号峰的绝对值与后一个的比为 $\ee^{\zeta\omega_0T_d}$, 对数减量
$\delta=2\pi\zeta/\sqrt{1-\zeta^2}$.
极值条件含一个固定相位偏移, 不影响相邻同号峰的时间差与比值. 若测得自由衰减的 $T_d$ 与 $\delta$, 可反推 $\zeta=\delta/\sqrt{4\pi^2+\delta^2}$ 及 $\omega_0=2\pi/[T_d\sqrt{1-\zeta^2}]$. 这些关系以理想线性黏性阻尼为前提, 不应用于一般非线性耗散的任意波形.

\originalsection{习题与解答}
\begin{exercise}求 $y''-2y'+y=\ee^t$ 的通解.\end{exercise}
\begin{solution}特征根1为二重根, $P''(1)=2$, 特解为 $t^2\ee^t/2$. 故 $y=\ee^t(A+Bt+t^2/2)$.\end{solution}
\begin{exercise}对 $x''+2x'+2x=\cos t$, 求长期稳态及振幅.\end{exercise}
\begin{solution}$P(\ii)=1+2\ii$, $H(\ii)=(1-2\ii)/5$. 稳态为 $(\cos t+2\sin t)/5$, 振幅为 $1/\sqrt5$. 齐次根为 $-1\pm\ii$, 因而任意初值的瞬态衰减.\end{solution}

\originalchapter{Laplace 变换、卷积与脉冲}
\originalsection{变换与导数公式}
\begin{definition}
若 $f:[0,\infty)\to\C$ 局部分段连续且满足 $|f(t)|\le M\ee^{at}$ (足够大 $t$), 则其单边 Laplace 变换为 $F(s)=\Lap f(s)=\int_0^\infty\ee^{-st}f(t)\dd t$, 在 $\operatorname{Re}s>a$ 绝对收敛. 收敛半平面与公式共同描述变换.
\end{definition}
有限段积分因局部可积而存在, 尾部可用 $M\ee^{-(\operatorname{Re}s-a)t}$ 控制. 相同估计加一个 $t^k$ 因子仍可积, 所以在严格内部的闭子半平面上可对 $s$ 求导, $F^{(k)}(s)=\Lap((-t)^kf)(s)$. 因而变换在其收敛半平面内解析.
\begin{proposition}
若 $f$ 局部绝对连续, $f'$ 局部分段连续, 且 $f,f'$ 为指数阶, 则在共同收敛半平面中 $\Lap f'=sF-f(0+)$. 对满足相应条件的 $n$ 阶导数,
\[
\Lap f^{(n)}=s^nF-s^{n-1}f(0+)-\cdots-f^{(n-1)}(0+).
\]
\end{proposition}
\begin{proof}
在 $[0,R]$ 上分部积分: $\int_0^R\ee^{-st}f' =\ee^{-sR}f(R)-f(0+)+s\int_0^R\ee^{-st}f$. 当实部超过指数阶, 边界尾项趋于零. 对更高阶反复使用该公式即可. 跳跃函数不能直接应用这个普通导数公式, 必须另计跳跃产生的脉冲.
\end{proof}
基本表可从积分得到 $\Lap1=1/s$, $\Lap\ee^{at}=1/(s-a)$, $\Lap\cos\omega t=s/(s^2+\omega^2)$, $\Lap\sin\omega t=\omega/(s^2+\omega^2)$. $\Lap t^n=n!/s^{n+1}$ 来自反复分部积分. 线性性和指数移位 $\Lap(\ee^{at}f)=F(s-a)$ 使表的适用范围扩大.

这里使用的逆变换在局部分段连续指数阶函数中唯一, 相等是除跳跃点取值外的几乎处处相等. 一个证明途径如下. 若变换差 $g$ 为零, 选实数 $b$ 严格大于其指数阶并使 $b>0$. $F(b+n)=0$ 对所有 $n\ge0$ 成立. 代换 $u=\ee^{-t}$ 得 $\int_0^1u^n h(u)\dd u=0$, $h(u)=u^{b-1}g(-\log u)\in L^1(0,1)$. 多项式在连续函数中一致稠密, 故 $h$ 对全部连续测试函数积分为零. 用连续函数逼近区间示性函数, 得每个区间积分为零, 再由不定积分的绝对连续性得 $h=0$ 几乎处处. 此处使用的一致多项式逼近和积分基本性质是分析先修; 具体有理式逆变换仍逐项代回核验.
\originalsection{初值问题与部分分式}
\begin{example}解 $y''+3y'+2y=1$, $y(0)=0$, $y'(0)=1$.\end{example}
\begin{solution}
记 $Y=\Lap y$, 导数公式给 $(s^2+3s+2)Y=1+1/s$. 因此 $Y=(s+1)/[s(s+1)(s+2)]=1/[s(s+2)]$, 得 $y=(1-\ee^{-2t})/2$. 初值项不可遗漏; 分子中的消去由此特定初速度与输入共同决定.
\end{solution}
对二次不可约分母应先配平方, 例如 $s^2+2as+b=(s+a)^2+\omega^2$ ($\omega^2=b-a^2>0$). 重复极点的逆变换含 $t$ 的幂; 这与特征方程的重根完全一致.
\originalsection{延迟、阶跃与卷积}
设 $H(t)$ 为 $t>0$ 时1、$t<0$ 时0的阶跃, 单点值不影响普通积分. 对 $a>0$,
\[
\Lap\{H(t-a)f(t-a)\}=\ee^{-as}F(s).
\]
这由 $u=t-a$ 代换得出. $H(t-a)f(t)$ 的变换一般不是 $\ee^{-as}F(s)$, 必须先把有效输入改写为延迟变量的函数.
\begin{definition}
两个因果函数的卷积为 $(f*g)(t)=\int_0^t f(t-u)g(u)\dd u$, $t\ge0$. 因果指 $t<0$ 时延拓为零.
\end{definition}
\begin{theorem}[卷积公式]
若 $f,g$ 局部分段连续且为指数阶, 则在足够右侧半平面 $\Lap(f*g)=FG$.
\end{theorem}
\begin{proof}
足够大的 $\sigma=\operatorname{Re}s$ 使 $\int_0^\infty\ee^{-\sigma t}|f(t)|\dd t$ 与 $g$ 的对应积分有限. 因而三角区域上的二重积分绝对可积, 可交换积分. 令 $v=t-u$, 得
\[
\int_0^\infty\int_0^t\ee^{-st}f(t-u)g(u)\dd u\dd t
=\int_0^\infty\int_0^\infty\ee^{-s(v+u)}f(v)g(u)\dd v\dd u=FG.
\]
所用交换依靠绝对可积, 不是把发散的两个表达式形式相乘.
\end{proof}
\begin{example}求 $y''+y=H(t-\pi)$, 零初值的解.\end{example}
\begin{solution}
$Y=\ee^{-\pi s}/[s(s^2+1)]$. 因 $1/[s(s^2+1)]=1/s-s/(s^2+1)$, 得 $y=H(t-\pi)[1-\cos(t-\pi)]$. 在开关点位置及速度连续, 加速度发生有限跳跃. 也可用脉冲响应 $k(t)=\sin t$ 卷积输入, 直接积分得到同一结果.
\end{solution}
\originalsection{脉冲与 Green 函数}
Dirac 脉冲 $\delta_a$ 不是取无穷值的普通函数. 它对光滑紧支撑测试函数 $\varphi$ 的作用为 $\langle\delta_a,\varphi\rangle=\varphi(a)$; 分布导数定义为 $\langle T',\varphi\rangle=-\langle T,\varphi'\rangle$. 从定义分部积分得 $H'=\delta_0$. 在因果约定下取 $\Lap\delta_0=1$, $\Lap\delta_a=\ee^{-as}$ ($a>0$), 零时刻质量全部计入, 不采用半质量约定.

若 $y$ 在 $a$ 左右光滑, 有跳跃 $[y]_a$, 则分布导数为常规分段导数加 $[y]_a\delta_a$. 再求导可得二阶导数中的 $[y']_a\delta_a+[y]_a\delta_a'$. 对 $my''+cy'+ky=J\delta_a$, 比较 $\delta_a'$ 系数先得 $[y]_a=0$, 再比较 $\delta_a$ 得 $m[y']_a=J$. 若先保留位置跳跃, 脉冲系数为 $m[y']_a+c[y]_a$, 在 $[y]_a=0$ 后才化为所写结论. 瞬时冲量改变速度, 不使位置跳跃.

对 $y''+py'+qy=f$ 的常系数版本, 因果脉冲响应为 $k(t)=H(t)k_+(t)$, 其中 $k_+$ 解齐次方程并满足 $k_+(0)=0$, $k_+'(0)=1$. 上述跳跃计算给 $(D^2+pD+q)k=\delta_0$. 零初值响应 $y=k*f$, 其 Laplace 变换为 $F/P(s)$. 这也可直接在卷积积分中微分两次证明, 所以卷积的意义不依赖先记住变换表.
\begin{example}求 $y''+2y'+2y=3\delta_2$, 零初值的解.\end{example}
\begin{solution}
$k_+(t)=\ee^{-t}\sin t$, 所以 $y=3H(t-2)\ee^{-(t-2)}\sin(t-2)$. 位置连续, 速度跳跃为3, 与单位质量的冲量条件一致. 普通微分方程只在 $t\ne2$ 处成立, 在全区间应按分布意义理解.
\end{solution}
\originalsection{极点与初终值}
对严格真有理的 $F(s)$, 部分分式给时间项 $t^{m-1}\ee^{\lambda t}$, 极点 $\lambda$ 的实部控制衰减, 虚部控制振荡. 零点消去后的实际极点才决定该输出; 不能根据未约分分母断言输入输出中必有某个模式.

若 $f$ 在零右连续且为指数阶, 由 $u=st$ 和控制收敛得 $\lim_{s\to+\infty}sF(s)=f(0+)$: 远尾用指数阶控制, 近零用连续性. 对有理变换, 取由部分分式得到的连续时间函数作为逆变换代表. 若约分后的 $sF(s)$ 的全部极点都在严格左半平面, 则 $f$ 为常数加衰减模式, 所以 $\lim_{t\to\infty}f(t)=\lim_{s\to0}sF(s)$. 这是终值定理在本册所需的充分版本. $F=1/(s^2+1)$ 时右侧为零而 $f=\sin t$ 无终值; 缺少极点条件不能使用公式.
\originalsection{习题与解答}
\begin{exercise}解 $y'+y=H(t-1)(t-1)$, $y(0)=0$.\end{exercise}
\begin{solution}$Y=\ee^{-s}/[s^2(s+1)]$, 其中 $1/[s^2(s+1)]=1/s^2-1/s+1/(s+1)$. 因而 $y=H(t-1)[t-2+\ee^{-(t-1)}]$.\end{solution}
\begin{exercise}求 $(\ee^{-t}*\ee^{-2t})(t)$, 并验证卷积变换.\end{exercise}
\begin{solution}积分为 $\ee^{-t}\int_0^t\ee^{-u}\dd u=\ee^{-t}-\ee^{-2t}$. 变换为 $1/(s+1)-1/(s+2)=1/[(s+1)(s+2)]$.\end{solution}

\originalchapter{周期输入与 Fourier 方法}
\originalsection{正交展开与系数}
周期输入通常由多个频率组成. 对 $2\pi$ 周期函数, 使用 $\sin nt,\cos nt$ 的目的是让常系数算子逐频率作用. 在实内积 $\langle f,g\rangle=\int_{-\pi}^{\pi}fg\dd t$ 下, 不同频率正交, $\|1\|^2=2\pi$, $\|\cos nt\|^2=\|\sin nt\|^2=\pi$ ($n\ge1$). 用积化和差并积分即可核验这些值.
\begin{definition}
局部可积的 $2\pi$ 周期函数的 Fourier 系数是
\[
a_0=\frac1\pi\int_{-\pi}^{\pi}f(t)\dd t,\quad
 a_n=\frac1\pi\int_{-\pi}^{\pi}f(t)\cos nt\dd t,\quad
 b_n=\frac1\pi\int_{-\pi}^{\pi}f(t)\sin nt\dd t.
\]
其形式级数为 $a_0/2+\sum_{n\ge1}(a_n\cos nt+b_n\sin nt)$. 周期 $T$ 时把 $nt$ 替换为 $2\pi nt/T$, 系数归一化相应为 $2/T$.
\end{definition}
形式展开不自动保证收敛. 本章所用的收敛条件为每周期内可分为有限个区间, 且在每个闭子区间上有单侧 $C^1$ 延拓, 即函数和一阶导数都具有有限单侧极限. 这比只要求开区间内 $C^1$ 更强, 保证断点附近的单侧导数有界.
\begin{theorem}[分段光滑函数的 Fourier 收敛]
在上述条件下, Fourier 部分和在任一点 $t$ 收敛到 $[f(t+)+f(t-)]/2$; 连续点即为函数值.
\end{theorem}
\begin{proof}
由三角恒等式, 部分和可写成 $S_N(t)=(2\pi)^{-1}\int_{-\pi}^{\pi}f(t-u)D_N(u)\dd u$, 其中 $D_N(u)=\sin((N+1/2)u)/\sin(u/2)$, 且积分为 $2\pi$. 将正负 $u$ 配对并减去 $m=(f(t+)+f(t-))/2$, 得
\[
S_N(t)-m=\frac1{2\pi}\int_0^\pi
\frac{f(t-u)+f(t+u)-2m}{\sin(u/2)}\sin((N+1/2)u)\dd u.
\]
分段 $C^1$ 使分子在 $u\downarrow0$ 时为 $O(u)$, 所以前面分式是可积的分段连续函数 $g(u)$. 任意这样的可积函数可以用有限区间阶梯函数作 $L^1$ 逼近: 在各连续段内部用一致连续性分割, 断点附近则用总长度控制积分. 阶梯函数与 $\sin(ku)$ 的积分为有限个 $O(1/k)$ 项; $L^1$ 误差对振荡积分的贡献至多为该误差. 先令 $k=N+1/2\to\infty$, 再令逼近误差趋零, 得结论.
\end{proof}
跳跃点的平均值解释了为什么不能直接在方波的跳跃处要求级数等于选定单点值. 部分和在跳跃附近有 Gibbs 超调, 点态收敛不代表整段一致收敛. 本册不把示意波形的超调数值当作严格估计.
\begin{example}展开 $f(t)=1$ ($0<t<\pi$), $f(t)=-1$ ($-\pi<t<0$), 周期为 $2\pi$.\end{example}
\begin{solution}
$f$ 为奇函数, 故 $a_n=0$. 计算得 $b_n=(2/\pi)\int_0^\pi\sin nt\dd t=2(1-(-1)^n)/(\pi n)$. 因此
$f(t)=\frac4\pi\sum_{j\ge0}\sin((2j+1)t)/(2j+1)$ 在连续点成立, 跳跃点和为零.
\end{solution}
\originalsection{逐频率响应的合法性}
对 $x''+cx'+kx=f(t)$, $c,k>0$, 采用标准复 Fourier 系数 $\widehat f_n=(2\pi)^{-1}\int_{-\pi}^{\pi}f(t)\ee^{-\ii nt}\dd t$ 时, 第 $n$ 频率的复系数乘 $1/(k-n^2+\ii cn)$. 形式叠加容易, 核查它是否确为解更重要. 对有跳跃的输入, 可以先从周期初值问题构造解, 再从解的系数验证频率关系, 这样无需未经证明地逐项二次求导.
\begin{proposition}\label{prop:periodic}
常系数系统 $X'=AX+g(t)$ 中, 若 $g$ 是周期 $T$ 的分段连续函数且 $A$ 的全部特征值实部负, 则有唯一连续、分段 $C^1$ 的 $T$ 周期解, 所有其他解趋向它.
\end{proposition}
\begin{proof}
常数变易给 $X(T)=\ee^{AT}X(0)+b$, $b=\int_0^T\ee^{A(T-s)}g(s)\dd s$. 矩阵 $\ee^{AT}$ 的特征值模都小于1, 故 $I-\ee^{AT}$ 可逆, 唯一选择 $X(0)=(I-\ee^{AT})^{-1}b$. 由周期输入和唯一性, 该解重复每周期. 任一其他解之差为 $\ee^{At}v$, Jordan 形式中每项为多项式乘实部为负的指数, 故趋于零. 矩阵指数的构造和谱性质在第8章证明; 本命题也可在二阶情形直接用第4章特征根证明.
\end{proof}
对上述二阶方程, $x,x'$ 连续且周期, $x''$ 分段连续. 分部积分时周期边界项消失, 第 $n$ 个 Fourier 系数满足 $(k-n^2+\ii cn)\widehat x_n=\widehat f_n$. 若输入分段 $C^1$, 则 $\widehat f_n=O(1/n)$, 这是各光滑段分部积分并保留有限跳跃项的结果. 因而 $\widehat x_n=O(1/n^3)$, $x$ 和 $x'$ 的 Fourier 级数绝对一致收敛. 不需要断言二阶导数级数也绝对收敛; 方程在输入的连续段按普通意义成立, 在跳跃点按分段或分布意义解释.
\begin{example}求方波输入下 $x''+2x'+2x=f(t)$ 的周期响应级数.\end{example}
\begin{solution}
对奇数 $n$, 输入项为 $4\sin nt/(\pi n)$, 它写成 $\operatorname{Re}(A_n\ee^{\ii nt})$ 时的复振幅是 $A_n=-4\ii/(\pi n)$; 标准 Fourier 系数则为 $\widehat f_n=A_n/2=-2\ii/(\pi n)$. 记 $a=2-n^2$, $b=2n$, 则输出为
\[
x(t)=\sum_{\substack{n\ge1\\ n\text{为奇数}}}\frac4{\pi n}
\frac{a\sin nt-b\cos nt}{a^2+b^2}.
\]
根 $-1\pm\ii$ 使周期解唯一且吸引. 每项幅度为 $O(n^{-3})$, 因而高频输入被显著平滑. 在 $t=0$ 的输出一般不为零, 周期响应不同于从静止开始的响应.
\end{solution}
\originalsection{周期解与无阻尼共振}
对 $x''+\omega_0^2x=f(t)$, 若某频率 $n\Omega$ 与 $\omega_0$ 相等, 那一 Fourier 系数非零时产生随 $t$ 增大的共振项, 不存在相同周期的解. 这可不靠无穷级数证明: 若 $x$ 和 $x'$ 都周期, 将方程乘 $\cos\omega_0t$ 或 $\sin\omega_0t$ 并在完整周期积分两次分部积分, 左边为零, 因而右边对应系数必须为零.

即使没有受迫共振, 无阻尼齐次运动不会消失, 所以周期特解通常不是所有初值的长期极限. 若 $\omega_0T$ 不是 $2\pi$ 的整数倍, $I-\ee^{AT}$ 仍可逆, 有唯一 $T$ 周期特解; 若恰为整数倍而共振系数为零, 周期齐次解使周期解不唯一.
\originalsection{习题与解答}
\begin{exercise}求 $f(t)=t$, $-\pi<t<\pi$, 周期延拓的 Fourier 系数.\end{exercise}
\begin{solution}奇性给 $a_n=0$, 分部积分给 $b_n=2(-1)^{n+1}/n$. 级数在区间内部收敛到 $t$, 在 $\pm\pi$ 的跳跃处收敛到零.\end{solution}
\begin{exercise}对 $y'+ay=f(t)$, $a>0$, 周期 $T$, 给出周期解的初值, 不使用 Fourier 展开.\end{exercise}
\begin{solution}令 $b=\int_0^T\ee^{-a(T-s)}f(s)\dd s$, 周期条件为 $y(0)=\ee^{-aT}y(0)+b$, 故 $y(0)=b/(1-\ee^{-aT})$. 任意解与它之差为 $C\ee^{-at}$.\end{solution}

\originalchapter{线性方程组与矩阵指数}
\originalsection{基本矩阵与常数变易}
\begin{definition}
对 $X'=A(t)X$, 若矩阵 $\Phi(t)$ 的各列组成一组线性独立解, 称其为基本矩阵. 一般非齐次系统是 $X'=A(t)X+g(t)$.
\end{definition}
在某点独立的列, 不可能在另一点突然相关: 若有常数向量 $c\ne0$ 使 $\Phi(t_1)c=0$, 则解 $\Phi(t)c$ 的零初值唯一性迫使它处处为零, 与初始独立矛盾. 所以基本矩阵整个区间可逆.
\begin{theorem}
连续系数系统的解为
\begin{equation}\label{eq:matrixvariation}
X(t)=\Phi(t)\Phi(t_0)^{-1}X_0+\Phi(t)\int_{t_0}^t\Phi(s)^{-1}g(s)\dd s.
\end{equation}
\end{theorem}
\begin{proof}
设 $X=\Phi c(t)$, 代入得 $\Phi'c+\Phi c'=A\Phi c+g$, 消去齐次部分后 $c'=\Phi^{-1}g$. 积分并使用初值即得公式. 矩阵乘法不交换, 因此顺序不可任意调换.
\end{proof}
\begin{proposition}[Liouville 公式]\label{prop:liouville}
基本矩阵满足 $\det\Phi(t)=\det\Phi(t_0)\exp(\int_{t_0}^t\tr A(s)\dd s)$.
\end{proposition}
\begin{proof}
按列展开行列式导数, 各列导数为 $A$ 乘该列. 在线性代数中该和等于 $\tr A\det\Phi$: 可先对标准基计算 $A$ 的作用, 非对角项导致两列重复, 只有对角项保留, 再换到当前列基. 因而行列式满足一阶线性方程, 解得公式.
\end{proof}
\originalsection{矩阵指数与谱}
\begin{definition}
对常矩阵 $A$, 定义 $\ee^{At}=\sum_{k=0}^\infty A^kt^k/k!$.
\end{definition}
取任一相容矩阵范数, 各项范数由 $\|A\|^k|t|^k/k!$ 控制, 所以紧时间区间上一致绝对收敛. 导数级数也同样收敛, 得 $(\ee^{At})'=A\ee^{At}$, $\ee^{A0}=I$. 初值唯一性给 $\ee^{A(t+s)}=\ee^{At}\ee^{As}$, 以及逆矩阵 $\ee^{-At}$. 因而 $\Phi=\ee^{At}$ 是基本矩阵.

若 $Av=\lambda v$, 则 $\ee^{At}v=\ee^{\lambda t}v$. 若 $A=SJS^{-1}$, 则 $\ee^{At}=S\ee^{Jt}S^{-1}$. 对 Jordan 块 $J=\lambda I+N$, $N^m=0$, 有 $\ee^{Jt}=\ee^{\lambda t}(I+tN+\cdots+t^{m-1}N^{m-1}/(m-1)!)$. 本册把有限维矩阵的特征分解与 Jordan 标准形作为线性代数先修, 此处完整给出它们对方程的作用. 对不交换的 $A,B$, 一般 $\ee^{A+B}\ne\ee^A\ee^B$, 不能把变系数系统直接写成 $\exp(\int A)$.
\begin{example}求 $X'=\begin{pmatrix}2&1\\0&2\end{pmatrix}X$, $X(0)=(a,b)^{\mathsf T}$.\end{example}
\begin{solution}
$A=2I+N$, $N^2=0$, 所以 $\ee^{At}=\ee^{2t}\begin{pmatrix}1&t\\0&1\end{pmatrix}$, $X=\ee^{2t}(a+bt,b)^{\mathsf T}$. 只有一条特征向量方向, 但仍需两个独立模式. 缺陷矩阵的第二模式是广义特征向量产生的 $t$ 项.
\end{solution}
\begin{example}对 $A=\begin{pmatrix}\alpha&-\beta\\\beta&\alpha\end{pmatrix}$, 求流并判定旋转方向.\end{example}
\begin{solution}
写 $A=\alpha I+\beta J$, $J^2=-I$, 得 $\ee^{At}=\ee^{\alpha t}\begin{pmatrix}\cos\beta t&-\sin\beta t\\\sin\beta t&\cos\beta t\end{pmatrix}$. $\beta>0$ 时从正 $x$ 轴向正 $y$ 轴逆时针旋转; $\alpha$ 控制半径增长或衰减.
\end{solution}
\originalsection{模态与耦合振动}
两个质量相同为 $m$ 的物体各接外侧刚度 $k$ 的弹簧, 中间用刚度 $\kappa$ 的弹簧连接. 位移满足
\[
m\begin{pmatrix}x_1''\\x_2''\end{pmatrix}
=-\begin{pmatrix}k+\kappa&-\kappa\\-\kappa&k+\kappa\end{pmatrix}
\begin{pmatrix}x_1\\x_2\end{pmatrix}.
\]
同相方向 $(1,1)$ 的特征刚度为 $k$, 反相方向 $(1,-1)$ 的特征刚度为 $k+2\kappa$. 令 $u=(x_1+x_2)/2$, $v=(x_1-x_2)/2$, 得 $mu''+ku=0$, $mv''+(k+2\kappa)v=0$. 这两个正常模态将耦合运动分解为独立振子. 改变坐标不是忽略耦合, 而是选取力矩阵的特征方向.
\originalsection{稳定性与输入}
\begin{theorem}\label{thm:linearstable}
常系数系统 $X'=AX$ 的原点渐近稳定当且仅当全部特征值实部严格负. 它 Lyapunov 稳定当且仅当全部实部非正, 且实部为零的特征值的 Jordan 块均为一阶.
\end{theorem}
\begin{proof}
Jordan 公式显示: 严格负的指数压过任意固定次数多项式, 故存在 $M,\gamma>0$ 使 $\|\ee^{At}\|\le M\ee^{-\gamma t}$, 同时得到稳定与收敛. 零实部的一阶块有界, 与负实部块合在一起使半群有界, 所以小初值始终保持小. 若有正实部, 在相应实不变子空间内取特征运动, 任意小非零初值最终逃出固定球. 若零实部块有非零幂零部分, 取最高广义特征方向得多项式增长, 也不稳定. 若存在零实部的一阶块, 则其非零运动不趋于零, 因而不能渐近稳定. 相似变换只改变固定范数常数, 不改变这些量词结论.
\end{proof}
稳定矩阵的非齐次解满足 $\|X(t)\|\le M\ee^{-\gamma t}\|X_0\|+M\int_0^t\ee^{-\gamma(t-s)}\|g(s)\|\dd s$. 有界输入给有界状态; 输入趋零时, 将积分拆成固定旧时段和最终小输入时段可证明状态趋零. 相同结论不适用于只有 Lyapunov 稳定、没有衰减的系统, 如无阻尼振子会对共振输入积累能量.
\originalsection{习题与解答}
\begin{exercise}求 $X'=\begin{pmatrix}1&2\\2&1\end{pmatrix}X$ 的通解.\end{exercise}
\begin{solution}特征值3、$-1$, 对应 $(1,1)$、$(1,-1)$. 因而 $X=c_1\ee^{3t}(1,1)^{\mathsf T}+c_2\ee^{-t}(1,-1)^{\mathsf T}$; 只有初值在第二条直线上才正向趋零.\end{solution}
\begin{exercise}给出全部特征值为零而原点不稳定的二维系统, 并解释机制.\end{exercise}
\begin{solution}取 $A=\begin{pmatrix}0&1\\0&0\end{pmatrix}$, 有 $x=x_0+y_0t$, $y=y_0$. 任意小 $y_0\ne0$ 都使 $x$ 最终无界; 非平凡 Jordan 块产生线性增长.\end{solution}

\originalchapter{平面系统的相图}
\originalsection{轨道与零流线}
自治平面系统 $x'=f(x,y)$, $y'=g(x,y)$ 的轨道是状态平面内的定向曲线. 在 $f\ne0$ 的局部可写 $\dd y/\dd x=g/f$, 但这个商在竖直运动处失效, 且丢掉了速度大小. 零流线 $f=0$ 和 $g=0$ 分别是水平或竖直分量为零的点集, 不是一般的不变轨道; 两者交点才是平衡点.

唯一性说明不同轨道不能在非平衡点相交, 但同一轨道可以周期重复. 穿过某点再次到达同一点, 由自治性与唯一性得到周期运动. 平衡点旁箭头很短也不意味着轨道能在有限时间到达它.
\originalsection{迹、行列式与线性分类}
对 $X'=AX$, $A$ 为实 $2\times2$ 矩阵, 设 $\tau=\tr A$, $\Delta=\det A$. 特征方程为 $\lambda^2-\tau\lambda+\Delta=0$. 因而:
\begin{itemize}
\item $\Delta<0$: 两个实根异号, 原点为鞍点. 负根特征线是稳定方向, 正根特征线是不稳定方向.
\item $\Delta>0$, $\tau^2>4\Delta$: 实根同号, $\tau<0$ 为稳定结点, $\tau>0$ 为不稳定结点.
\item $\Delta>0$, $\tau^2<4\Delta$: 共轭复根实部 $\tau/2$, 为稳定或不稳定焦点; $\tau=0$ 时为中心.
\item $\tau^2=4\Delta>0$: 重实根, 标量矩阵给星形结点, 非标量情形可能给 Jordan 模式, 都按根的符号判断吸引或排斥.
\end{itemize}
$\Delta=0$ 时至少有零特征值, 可能有一条平衡点直线, 不在上述孤立类型中. 稳定区域 $\tau<0$, $\Delta>0$ 同时包含结点与焦点.
\begin{center}
\begin{tikzpicture}[>=Stealth,scale=.8]
\begin{scope}
\draw[->](-2,0)--(2,0)node[right]{$x$};\draw[->](0,-1.8)--(0,1.8)node[above]{$y$};
\draw[->,thick](.3,0)--(1.5,0);\draw[->,thick](-.3,0)--(-1.5,0);
\draw[->,thick](0,1.5)--(0,.3);\draw[->,thick](0,-1.5)--(0,-.3);
\draw[domain=.38:1.6,smooth,variable=\u]plot({\u},{.55/\u});
\draw[->,thick](.85,.65)--(1.12,.49);
\node at (0,-2.4){$x'=x,\ y'=-y$};
\end{scope}
\begin{scope}[xshift=6cm]
\draw[->](-2,0)--(2,0)node[right]{$x$};\draw[->](0,-1.8)--(0,1.8)node[above]{$y$};
\draw[domain=0:13,samples=180,smooth,variable=\u]plot({1.6*exp(-.15*\u)*cos(\u r)},{1.6*exp(-.15*\u)*sin(\u r)});
\draw[->,thick](.93,.78)--(.64,.99);
\node at (0,-2.4){$x'=-.15x-y,\ y'=x-.15y$};
\end{scope}
\end{tikzpicture}

左: 鞍点. 右: 向内的逆时针螺旋, 为稳定焦点.
\end{center}
图只展示所给系统的定性方向, 相似变换后的相图可能被拉伸倾斜, 不应把图中直角当作所有特征向量都正交.
\originalsection{非线性系统与线性化}
若 $F\in C^1$, $F(a)=0$, 令 $u=X-a$, Taylor 公式给 $u'=Au+R(u)$, $A=DF(a)$, $\|R(u)\|/\|u\|\to0$. 在小邻域中线性部分是首阶近似, 但随时间累积的误差不能仅凭一次 Taylor 展开判定. 第10章将证明 $A$ 全部特征值实部负时非线性平衡点渐近稳定.

一般双曲平衡点指没有零实部特征值的平衡点. 局部拓扑分类还可调用 Hartman--Grobman 定理: $C^1$ 向量场在双曲平衡点附近的轨道结构与其线性化局部拓扑等价, 保持时间方向. 本册不证明这一一般拓扑共轭定理, 不靠它替代后文稳定性的直接证明; 鞍点附近稳定与不稳定方向在非线性系统中通常弯成曲线, 不能把它们等同于整条特征直线. 非双曲情形不在该定理条件内.
\begin{example}分析单摆 $\theta'=v$, $v'=-\sin\theta-\gamma v$, $\gamma\ge0$ 的平衡点.\end{example}
\begin{solution}
平衡点为 $(k\pi,0)$. Jacobian 为 $\begin{pmatrix}0&1\\-\cos k\pi&-\gamma\end{pmatrix}$. $k$ 奇数时行列式为 $-1$, 线性化为鞍点. $k$ 偶数且 $\gamma>0$ 时迹负、行列式1, 原点型平衡渐近稳定; $0<\gamma<2$ 为焦点, $\gamma>2$ 为结点, $\gamma=2$ 为重根. $\gamma=0$ 时线性化为中心, 不能只由此断言非线性稳定; 能量 $v^2/2+1-\cos\theta$ 的局部正定性在下一章给出证明.
\end{solution}
\originalsection{守恒量与捕食模型}
Lotka--Volterra 模型为 $x'=x(a-by)$, $y'=y(-c+dx)$, $a,b,c,d>0$. 坐标轴是不可穿越的不变边界; 正初值保持正, 因为 $x=x_0\exp\int(a-by)\dd t$, $y=y_0\exp\int(-c+dx)\dd t$. 正平衡点为 $(c/d,a/b)$.
\begin{proposition}
在正象限内 $V(x,y)=dx-c\log x+by-a\log y$ 沿轨道恒定. 非平衡的正象限轨道为周期曲线.
\end{proposition}
\begin{proof}
微分得
$V'=(d-c/x)x(a-by)+(b-a/y)y(-c+dx)=0$.
$V$ 的 Hessian 为 $\diag(c/x^2,a/y^2)$, 严格正定, 在正平衡点取得唯一最小值; 接近轴或趋于无穷时 $V\to\infty$. 每个高于最小值的等值线是围住平衡点的光滑闭曲线: 梯度只在最小点为零, 严格凸性使每条从该点出发的射线与等值线恰交一次. 等值线上向量场不为零, 相切且速度有正下界. 用弧长参数 $s$ 写 $\dd s/\dd t=v(s)>0$, 绕行时间 $\int\dd s/v(s)$ 有限, 所以轨道周期. 守恒量也将轨道限制在正象限紧集, 保证全时存在.
\end{proof}
这里的闭轨道形成整族, 不属于孤立极限环. 线性化中心与非线性闭轨道相容, 真正支持周期性的是守恒量及其等值线几何.

\originalsection{单摆的周期与分离轨道}
无阻尼单摆的守恒量为 $E=v^2/2+1-\cos\theta$. 在展开的 $(\theta,v)$ 平面中, 每个能量 $0<h<2$ 围绕下垂点形成闭曲线; 转角最大值 $\theta_m$ 满足 $1-\cos\theta_m=h$. 运动从 $\theta=0$ 到 $\theta_m$ 的时间可由 $v=\sqrt{2(h-1+\cos\theta)}$ 积分. 设 $k=\sin(\theta_m/2)=\sqrt{h/2}$, 用 $\sin(\theta/2)=k\sin\varphi$ 消去端点平方根, 得周期
\[
T(h)=4\int_0^{\pi/2}\frac{\dd\varphi}{\sqrt{1-k^2\sin^2\varphi}}.
\]
当 $h\downarrow0$, 分母一致趋于1, 所以 $T\to2\pi$, 与小振幅线性化一致. 当 $h\uparrow2$, 被积函数单调趋于 $1/\cos\varphi$, 其积分发散, 单调收敛给 $T\to\infty$. 非线性摆的周期随振幅增大, 不能把线性近似周期用于接近翻转的运动.

能量 $h=2$ 时, 在 $-\pi<\theta<\pi$ 的上支有 $v=2\cos(\theta/2)$. 对应显式轨道为 $\theta(t)=4\arctan\ee^{t-t_*}-\pi$, $v(t)=2\operatorname{sech}(t-t_*)$. 直接求导验证 $\theta'=v$, $v'=-\sin\theta$. 当 $t\to\pm\infty$, 它分别趋于相邻的两个上竖平衡点; 达到平衡需要无限时间. 这一分离轨道把往复摆动与完整旋转的能量区域分开. $h>2$ 时速度不为零, 转角在展开平面不断增加或减少; 若把角度模 $2\pi$ 识别, 状态空间是圆柱, 与展开平面内“闭曲线”的说法应区分.

\originalsection{习题与解答}
\begin{exercise}分类 $A=\begin{pmatrix}0&1\\-2&-3\end{pmatrix}$ 的原点并求特征方向.\end{exercise}
\begin{solution}迹 $-3$, 行列式2, 根为 $-1,-2$, 是稳定结点. 方向分别为 $(1,-1)$、$(1,-2)$; 一般轨道正向以较慢的 $-1$ 方向趋向原点, 在纯 $-2$ 特征线上则保持快方向.\end{solution}
\begin{exercise}系统 $x'=-y-x(x^2+y^2)$, $y'=x-y(x^2+y^2)$ 的线性化是中心. 判断原点实际稳定性.\end{exercise}
\begin{solution}对 $r>0$, $r'=-r^3$, $\theta'=1$, 因而 $r(t)=r_0/\sqrt{1+2r_0^2t}$ 正向趋零且不增. 原点渐近稳定, 非线性三次项改变了中心的长期行为.\end{solution}

\originalchapter{Lyapunov 方法与耗散}
\originalsection{正定函数与稳定性}
相图能看二维运动, 但在更高维或无法显式求解时, 一个沿轨道单调的标量函数更有用. 对 $X'=F(X)$, 平衡点移到0, 沿 $C^1$ 函数 $V$ 的导数为 $\dot V(X)=\nabla V(X)\cdot F(X)$.
\begin{definition}
在原点邻域中, 若 $V(0)=0$, 且所有 $X\ne0$ 都有 $V(X)>0$, 称 $V$ 正定. 若 $V(X)\to\infty$ 当 $\|X\|\to\infty$, 称其径向无界.
\end{definition}
\begin{theorem}[Lyapunov 稳定性定理]\label{thm:lyapunov}
设 $F$ 局部 Lipschitz, $F(0)=0$. 若存在邻域内 $C^1$ 正定函数 $V$ 且 $\dot V\le0$, 则原点稳定. 若 $\dot V(X)<0$ 对非零点成立, 则原点渐近稳定. 若正定性与导数负定条件在全空间成立且 $V$ 径向无界, 则结论为全局渐近稳定.
\end{theorem}
\begin{proof}
取闭球 $\overline B_\eps$ 完全位于该邻域. 球面上 $V$ 的最小值 $m>0$. 由 $V$ 在零处连续, 取 $\delta<\eps$ 使 $\|X_0\|<\delta$ 时 $V(X_0)<m$. 沿轨道 $V$ 不增, 所以无法首次到达球面. 更精确地, 它被限制在 $\overline B_\eps$ 内的闭子水平集 $V\le V(X_0)$, 与球面分离, 延续定理保证全正向存在, 得稳定.

若导数负定, 取上述紧子水平集中的轨道. 对任意 $\eta>0$, 在紧环带 $\eta\le\|X\|\le\eps$ 上 $-\dot V$ 有正下界. 若轨道始终在这个环带内, $V$ 会最终变成负数, 矛盾, 因而它进入任意小球. 结合已证明的稳定性: 给定目标半径 $\rho$, 先选使进入后永不离开 $B_\rho$ 的稳定半径 $\eta$, 再证明轨道最终进入 $B_\eta$, 就得趋于零. 全空间径向无界时每个子水平集紧, 同一论证覆盖任意初值.
\end{proof}
\begin{example}证明 $x'=-x-y^3$, $y'=x^3-y$ 的原点全局渐近稳定.\end{example}
\begin{solution}
选 $V=(x^4+y^4)/4$, 因交叉项次数匹配,
$\dot V=x^3(-x-y^3)+y^3(x^3-y)=-x^4-y^4=-4V$. $V$ 正定、径向无界且导数负定, 定理适用. 事实上 $V(t)=V(0)\ee^{-4t}$, 可直接得到状态的指数趋零上界.
\end{solution}
\originalsection{线性化稳定性的证明}
\begin{theorem}\label{thm:nonlinearstable}
设 $F\in C^1$ 且 $F(0)=0$. 若 $A=DF(0)$ 的全部特征值实部严格负, 则原点局部渐近稳定, 并且在足够小邻域有指数衰减估计.
\end{theorem}
\begin{proof}
由定理 \ref{thm:linearstable}, $\|\ee^{At}\|\le M\ee^{-\gamma t}$. 定义对称矩阵
$P=\int_0^\infty\ee^{A^{\mathsf T}t}\ee^{At}\dd t$. 积分绝对收敛; 对非零 $X$, $X^{\mathsf T}PX=\int_0^\infty\|\ee^{At}X\|^2\dd t>0$, 故 $P$ 正定. 微分被积矩阵并在零与无穷处取边界得 $A^{\mathsf T}P+PA=-I$.

写 $F(X)=AX+R(X)$, $R=o(\|X\|)$, 取 $V=X^{\mathsf T}PX$. 有
$\dot V=-\|X\|^2+2X^{\mathsf T}PR(X)$. 足够小球内 $\|R(X)\|\le\|X\|/(4\|P\|)$, 故 $\dot V\le-\|X\|^2/2\le-V/(2\lambda_{\max}(P))$. 前述子水平集保证轨道保持在此球中, 所以 $V$ 指数下降. 使用 $\lambda_{\min}\|X\|^2\le V\le\lambda_{\max}\|X\|^2$ 转成状态指数估计. 这个证明明确控制了非线性余项, 而不把局部 Taylor 近似直接延伸到无限时间.
\end{proof}
\originalsection{LaSalle 不变原理}
能量导数常只有半负定, 如阻尼单摆的 $\dot E=-\gamma v^2$. 零导数集合含整条 $v=0$ 线, 但该线上除平衡点之外的点会立即离开, 所以不能把它都当作长期极限.
\begin{definition}
轨道的 $\omega$ 极限集是所有 $t_j\to\infty$ 时 $X(t_j)$ 的极限点组成的集合. 集合正向不变指从其中出发的轨道正向始终留在集合内; 完全不变指还可沿其整条轨道反向留在集合内.
\end{definition}
\begin{theorem}[LaSalle 的紧集版本]\label{thm:lasalle}
设 $K$ 为紧的正向不变集, $F$ 在其邻域局部 Lipschitz, $V\in C^1$ 且 $\dot V\le0$ 于 $K$. 每条从 $K$ 出发的轨道到集合 $E=\{X\in K:\dot V(X)=0\}$ 中最大完全不变子集 $M$ 的距离趋于零.
\end{theorem}
\begin{proof}
轨道在紧集内全正向存在. 单调有界的 $V(X(t))$ 有极限 $\ell$. 紧性使 $\omega$ 集非空紧. 对其任一点 $p=\lim X(t_j)$, 任意固定正时间 $s$ 的连续依赖给 $\phi_s(p)=\lim X(t_j+s)\in\omega$. 对负时间, 可在 $K$ 邻域取一致局部存在长度, 从 $X(t_j)$ 反向到 $X(t_j-s)$ 并用连续依赖取极限; 重复这个过程得到 $p$ 的全反向轨道也在 $\omega$. 因而 $\omega$ 完全不变. $V$ 在 $\omega$ 上恒为 $\ell$, 沿其中轨道求导得 $\dot V=0$, 故 $\omega\subset M$.

若到 $M$ 的距离不趋零, 有 $t_j\to\infty$ 使距离至少为某个 $\eps>0$. 紧性给出收敛子列, 其极限属于 $\omega\subset M$, 与距离下界矛盾.
\end{proof}
\begin{example}证明有阻尼单摆在下垂平衡点附近渐近稳定, 同时解释无阻尼情形.\end{example}
\begin{solution}
取 $E=v^2/2+1-\cos\theta$, 在 $|\theta|<\pi$ 内以 $(0,0)$ 为唯一零点. 对 $0<h<2$, 选包含原点的子水平集 $E\le h$ 分支, 它紧且与 $\theta=\pm\pi$ 分离. $E'=-\gamma v^2$ 使其正向不变. 当 $\gamma>0$, $E'=0$ 要求 $v=0$; 要在此线一直停留又要求 $v'=-\sin\theta=0$, 在该集内只有原点. LaSalle 给趋于原点, 正定能量给稳定. 若 $\gamma=0$, 能量守恒, 原点稳定但非零小能量轨道不会趋于零.
\end{solution}
\originalsection{体积收缩与吸引范围}
$C^1$ 流的变分矩阵满足 \eqref{eq:variational}, 用 Liouville 公式得
\[
\det D\phi_t(X_0)=\exp\int_0^t\operatorname{div}F(\phi_s(X_0))\dd s.
\]
变换变量公式因此把小体积的变化与散度联系. 散度恒负表示体积收缩, 不能单独推出所有轨道有界或都趋向一个点. 例如 $x'=x$, $y'=-2y$ 的散度为 $-1$, 但一般轨道的 $x$ 分量增长; 一个方向的扩张可与更大的另一方向收缩共存.

平衡点的吸引域是所有正向趋向它的初值集合. 局部 Lyapunov 函数给其中一部分可验证子水平集, 不保证已算出整个吸引域. 多稳态系统中的边界可由鞍点的稳定轨道分隔, 仅画若干数值轨道不足以证明全域边界.
\originalsection{习题与解答}
\begin{exercise}证明 $x'=-x+x^3$ 中原点的吸引域恰为 $(-1,1)$.\end{exercise}
\begin{solution}平衡点为 $-1,0,1$. 相线显示 $(-1,0)$ 上右端正, $(0,1)$ 上负, 内部轨道趋零. $\pm1$ 恒定, 外侧轨道向外, 不趋零. 因而边界是两个不稳定平衡点.\end{solution}
\begin{exercise}对 $x'=y$, $y'=-x-y-y^3$, 证明原点全局渐近稳定.\end{exercise}
\begin{solution}取 $V=(x^2+y^2)/2$, $\dot V=-y^2-y^4\le0$. 每个子水平集紧且不变. 导数为零时 $y=0$, 在该线保持又要求 $y'=-x=0$, 所以最大不变子集为原点. LaSalle 给全局吸引, 正定 $V$ 给稳定.\end{solution}

\originalchapter{极限环、返回映射与分岔}
\originalsection{孤立周期轨道}
\begin{definition}
自治系统的周期轨道若在某邻域中没有其他周期轨道, 称极限环. 若附近轨道到它的距离正向趋零且对轨道集合稳定, 称稳定极限环. 稳定性比较的是到整条轨道的距离, 不是每条邻近解与某一固定相位点的距离.
\end{definition}
极坐标在 $r>0$ 上满足 $r'=(xx'+yy')/r$, $\theta'=(xy'-yx')/r^2$. 它把径向吸引和绕行角速度分开, 适用于有旋转结构的系统.
\begin{example}分析 $x'=\mu x-y-x(x^2+y^2)$, $y'=x+\mu y-y(x^2+y^2)$.\end{example}
\begin{solution}
对 $r>0$ 有 $r'=\mu r-r^3$, $\theta'=1$. $\mu<0$ 时所有半径正向趋零; $\mu=0$ 时 $r'=-r^3$, 原点仍渐近稳定, 但衰减是代数速度. $\mu>0$ 时原点排斥, 圆 $r=\sqrt\mu$ 为周期 $2\pi$ 的极限环. $0<r<\sqrt\mu$ 向外, $r>\sqrt\mu$ 向内, 所有非零轨道到该圆趋零. 径向线性化导数 $\mu-3r^2=-2\mu<0$ 验证吸引. 圆内外的半径单调性排除其他周期轨道, 所以该圆确实孤立. 这是超临界 Hopf 分岔的正规模型.
\end{solution}
\begin{center}
\begin{tikzpicture}[>=Stealth,scale=1.05]
\draw[->](-2,0)--(2,0)node[right]{$x$};\draw[->](0,-2)--(0,2)node[above]{$y$};
\draw[thick,color=structurecolor](0,0)circle(1.2);
\draw[->,thick](1.2,0)arc(0:45:1.2);
\draw[->](.5,.45)--(.8,.72);\draw[->](1.65,1.0)--(1.25,.76);
\node[anchor=west]at(1.25,-1.1){$r=\sqrt\mu$};\fill(0,0)circle(1.8pt);
\end{tikzpicture}

$\mu>0$ 的径向吸引与逆时针绕行示意. 内外箭头只表示径向趋势.
\end{center}
一般 Hopf 定理还需光滑性、特征值横穿虚轴及非退化非线性系数等条件. 仅在某参数处出现纯虚根不能保证生出极限环; 本册证明的是此正规模型, 不把它当作一般 Hopf 定理的完整证明.
\originalsection{Poincar\'e 返回映射}
选周期轨道上的非平衡点 $p$, 取一条横截线 $\Sigma$ 与向量场横交. 附近初值在接近一周期的时间再次穿过 $\Sigma$, 定义返回映射 $P$. 横交性使穿越时刻由隐函数定理局部确定; 对 $C^1$ 向量场, 初值可微依赖使 $P$ 为 $C^1$. 周期轨道对应返回映射的固定点.
\begin{proposition}
一维返回映射的固定点 $z_*$ 若 $|P'(z_*)|<1$, 则局部吸引且稳定; 若 $|P'(z_*)|>1$, 则不稳定. 导数模等于1时, 线性判据无结论.
\end{proposition}
\begin{proof}
第一种情形在小区间内 $|P'|\le q<1$, 均值定理给 $|P(z)-z_*|\le q|z-z_*|$, 迭代得几何衰减且区间不变. 第二种情形取 $|P'|\ge q>1$ 的小区间, 在轨道留于其中时每次偏离至少放大 $q$ 倍, 任意非零小扰动最终离开. 从截面到相邻一次绕行的时间与位置连续依赖, 因而离散稳定性转为周期轨道集合的稳定性.
\end{proof}
\begin{proposition}[平面周期轨道的乘子]\label{prop:multiplier}
设 $\gamma(t)$ 为 $C^1$ 平面向量场的周期 $T$ 轨道. 返回映射在固定点的导数为
$P'(z_*)=\exp\int_0^T\operatorname{div}F(\gamma(t))\dd t$.
\end{proposition}
\begin{proof}
变分矩阵 $Y(T)$ 的一个特征值为1, 因为切向量 $F(\gamma(0))$ 经一周期回到自身: 对流的时间平移关系求导即可. 另一个特征值等于 $\det Y(T)$, Liouville 公式给出所写指数. 返回映射导数是在横向商空间上的诱导作用, 穿越时刻变化只添一个切向分量, 不改变横向商. 因此另一特征值即为 $P'(z_*)$.
\end{proof}
在上例 $\operatorname{div}F=2\mu-4r^2=-2\mu$ 于圆上, 故乘子为 $\ee^{-4\pi\mu}<1$, 与径向线性化一致.
\originalsection{排除周期轨道}
\begin{theorem}[Bendixson 判据]
若 $F=(f,g)\in C^1$ 在单连通平面区域 $U$ 上, $\operatorname{div}F$ 处处严格同号, 则 $U$ 内没有非恒定周期轨道.
\end{theorem}
\begin{proof}
若有周期轨道, 唯一性和最小周期使它是一条简单闭曲线, 且速度不为零. 单连通性使所围内部也位于 $U$. 向量场沿边界切向, 通量为零. Green 公式给 $0=\oint F\cdot n\dd s=\iint\operatorname{div}F\dd A$, 与散度严格同号矛盾. 此处使用平面 Jordan 曲线与 Green 公式作为微积分和拓扑先修.
\end{proof}
Dulac 扩展是选 $C^1$ 标量 $B$ 使 $\operatorname{div}(BF)$ 严格同号. 沿原轨道 $BF$ 仍切向, 同一通量证明适用. 区域若有孔, 曲线围成的内部可能不在区域内, 此论证就不能使用.
\begin{example}证明 $x'=x(1-x-ay)$, $y'=y(1-y-bx)$ ($a,b>0$) 在正象限内没有周期轨道.\end{example}
\begin{solution}
取 $B=1/(xy)$, 它在正象限 $C^1$. 有 $BF=((1-x-ay)/y,(1-y-bx)/x)$, 散度为 $-1/y-1/x<0$. 正象限单连通, Dulac 判据排除其内部周期轨道. 边界行为另需相线分析, 不包含在此结论中.
\end{solution}
\originalsection{平面极限与高维限制}
Poincar\'e--Bendixson 定理的常用版本为: $C^1$ 平面流的一条正向轨道若留在紧集内, 且其 $\omega$ 极限集不含平衡点, 则该极限集是一条周期轨道. 这是本册引用而不展开完整横截面拓扑证明的定理; 条件中的紧性与排除平衡点均不可省. 它可以在有不变环域且环域内无平衡点时保证至少一条周期轨道, 不能自动保证唯一性或稳定性.

平面自治流的横截面是一维, 唯一性限制轨道的交错. 三维流的截面可为二维, 时间周期强迫的二维系统也可通过增加相位变成三维自治系统. 因此不能把平面定理用于三维 Lorenz 模型或二维非自治振子来排除混沌. 周期受迫振子可按每个强迫周期采样, 得到二维离散返回映射.

\originalsection{Van der Pol 振子}
Van der Pol 方程 $x''-\mu(1-x^2)x'+x=0$, $\mu>0$, 将小振幅下的负阻尼与大振幅下的正阻尼结合. 令 $v=x'$, 则 $x'=v$, $v'=\mu(1-x^2)v-x$. 能量 $E=(x^2+v^2)/2$ 满足 $E'=\mu(1-x^2)v^2$. 在 $|x|<1$ 的运动片段中它增长, 在 $|x|>1$ 的片段中它下降; 一次片段的符号并不证明整个周期的总能量变化.

因 $E'\le\mu v^2\le2\mu E$, 有 $E(t)\le E(0)\ee^{2\mu t}$, 所以每个解全正向存在. 原点 Jacobian 为 $\begin{pmatrix}0&1\\-1&\mu\end{pmatrix}$, 线性特征根实部正或两正根, 表明小扰动具有增长模式. 为研究闭轨道, 令 $F(x)=\mu(x^3/3-x)$, $z=v+F(x)$, 得 Liénard 平面形式 $x'=z-F(x)$, $z'=-x$.

这里阻尼函数 $u(x)=\mu(x^2-1)$ 为偶连续函数, 恢复力 $g(x)=x$ 为奇函数且在 $x>0$ 时为正, 积分 $\int_0^xg(s)\dd s=x^2/2\to\infty$; $F=\int_0^xu(s)\dd s$. 在此模型中, $F$ 为奇函数, 在 $0<x<\sqrt3$ 为负, 在 $x>\sqrt3$ 为正且严格递增并趋于无穷. 标准 Liénard 定理据这些条件给出唯一吸引极限环. 原 MIT 笔记 LC 在 PDF 第5页引用的 Levinson--Smith 版本以偶阻尼、奇恢复力、积分符号及增长条件表达该结论. 本册不证明这一全局唯一性定理, 也不从 $E'$ 的局部符号替代它; 这里只核验本模型满足条件, 说明极限环的自持振动机制. 对 $\mu\to0$ 的幅值近似或 $\mu$ 很大的松弛振荡, 还需平均法或奇异摄动的误差理论, 不作为已证明的定量结论.

\originalsection{习题与解答}
\begin{exercise}对 $r'=r(1-r^2)$, $\theta'=2$, 求极限环周期与横向乘子.\end{exercise}
\begin{solution}环为 $r=1$, 周期 $T=\pi$. 径向扰动满足 $\rho'=-2\rho$, 所以乘子 $\ee^{-2\pi}$. 原点不稳定, 所有非零正向轨道半径趋于1.\end{solution}
\begin{exercise}说明负散度为什么不能排除三维流中的周期轨道.\end{exercise}
\begin{solution}取平面稳定环正规模型 $x'=x-y-xr^2$, $y'=x+y-yr^2$, 并附加 $z'=-3z$, $r^2=x^2+y^2$. 三维散度为 $-1-4r^2<0$, 但 $r=1,z=0$ 是周期轨道. 平面 Green 公式用围域通量的证明不适用于空间闭曲线.\end{solution}

\originalchapter{幂级数解与边值问题}
\originalsection{普通点的级数构造}
解析方法不限于初等函数. 若方程系数解析, 在普通点附近可把未知解展开为幂级数, 用方程递推系数, 但还需证明所得级数确实收敛.
\begin{theorem}
若矩阵 $A(t)$ 与向量 $b(t)$ 在 $|t|<R$ 有收敛幂级数, 则系统 $X'=A(t)X+b(t)$ 的初值解在零附近解析. 递推公式为
\[
(n+1)X_{n+1}=\sum_{k=0}^n A_kX_{n-k}+b_n,\qquad X_0=X(0).
\]
\end{theorem}
\begin{proof}
形式代入后比较 $t^n$ 系数得到递推, 因而系数唯一. 任取 $0<\rho<R$, 收敛性给出某个 $M>0$ 使 $\|A_k\|,\|b_k\|\le M\rho^{-k}$. 用正系数标量级数 $v$ 支配, 令
$v'=M(v+1)/(1-t/\rho)$, $v(0)=v_0\ge\|X_0\|$. 其解为 $v+1=(v_0+1)(1-t/\rho)^{-M\rho}$, 在 $|t|<\rho$ 解析且 Taylor 系数非负. 比较递推并归纳得 $\|X_n\|\le v_n$. 因此 $X$ 的级数在此圆盘绝对收敛, 内部紧子圆盘可逐项微分并乘系数级数, 故所得函数解原方程. 初值唯一性把它与真实解识别. $\rho<R$ 任意, 可覆盖系数共同收敛圆盘.
\end{proof}
\begin{example}求 $y''+ty=0$, $y(0)=a$, $y'(0)=b$ 的级数解.\end{example}
\begin{solution}
写 $y=\sum c_nt^n$, 则 $c_0=a,c_1=b,c_2=0$, 对 $n\ge1$,
$c_{n+2}=-c_{n-1}/[(n+2)(n+1)]$. 因而
\[
y=a\left(1-\frac{t^3}{6}+\frac{t^6}{180}-\cdots\right)
+b\left(t-\frac{t^4}{12}+\frac{t^7}{504}-\cdots\right).
\]
把它改为一阶系统后系数为多项式, 可在任意有限 $R$ 应用定理, 所以两级数处处收敛. 前几个项不是完整解的替代, 递推与收敛证明才确定全部解.
\end{solution}
\originalsection{正则奇点与 Frobenius 代入}
若 $y''+p(t)y'+q(t)y=0$ 的 $p,q$ 在零有奇点, 而 $tp(t)$ 和 $t^2q(t)$ 在零解析, 称零为正则奇点. 尝试 $y=t^r\sum_{n\ge0}c_nt^n$, $c_0\ne0$, 其中 $t>0$ 时 $t^r=\ee^{r\log t}$ 取实幂或复幂. 最低次项给指标方程 $r(r-1)+p_0r+q_0=0$, $p_0=(tp)(0)$, $q_0=(t^2q)(0)$.

此方程给候选指数, 不自动给两个独立纯幂级数解. 根差为整数或重根时可能出现对数项; 本册不使用未证明的一般 Frobenius 完备性定理, 而在下例把递推、收敛和第二解直接做完.
\begin{example}分析 $t^2y''+ty'+t^2y=0$ 在 $t>0$ 靠近零处的解.\end{example}
\begin{solution}
指标方程为 $r^2=0$, 重根 $r=0$. 取 $c_0=1$, 比较系数得 $c_1=0$, $n^2c_n+c_{n-2}=0$ ($n\ge2$), 所以
\[
y_1(t)=\sum_{k=0}^\infty\frac{(-1)^kt^{2k}}{4^k(k!)^2}.
\]
相邻绝对项之比为 $t^2/[4(k+1)^2]\to0$, 级数处处收敛, 可微分代回验证. 在零附近 $y_1=1+O(t^2)$, 不为零. 用第4章降阶公式, 因 $p=1/t$, 可取
$y_2=y_1\int\dd t/(t y_1^2)$. 被积函数为 $1/t+O(t)$, 因此 $y_2=y_1\log t+O(t^2)y_1$ (适当取积分常数). Wronskian 为 $1/t\ne0$, 两解构成局部基本解组. 第二解在零有对数奇性, 正规初值理论不能以 $t=0$ 为普通初始时刻应用. 这里 $y_1$ 是零阶 Bessel 函数的一个标准规范化解, 名称不承担证明作用.
\end{solution}
\originalsection{初值与边值的区别}
以下取 $a<b$, 实连续 $f$ 与 $q$. 边值问题把条件放在不同位置, 如 $-y''+q(t)y=f(t)$, $y(a)=y(b)=0$. 初值唯一性不能保证此问题有唯一解, 因为起始导数未知, 要调节它使末端条件成立.
\begin{proposition}
若实连续 $q\ge0$ 于 $[a,b]$, 则上述 Dirichlet 边值问题至多有一个 $C^2$ 解.
\end{proposition}
\begin{proof}
两个解之差 $z$ 满足 $-z''+qz=0$, 零端点. 乘 $z$ 并积分分部, 得 $\int_a^b[(z')^2+qz^2]\dd t=0$. 两项非负, $z'\equiv0$, 再由端点得 $z=0$.
\end{proof}
对 $-y''=f$ 可直接构造 Green 核
\[
G(t,s)=\begin{cases}(t-a)(b-s)/(b-a),&t\le s,\\(s-a)(b-t)/(b-a),&s\le t.\end{cases}
\]
它两端为零, 在 $t=s$ 连续, 一阶导数跳跃为 $-1$, 所以 $-\partial_t^2G=\delta_s$. 因此 $y(t)=\int_a^bG(t,s)f(s)\dd s$ 解边值问题. 也可分成两段普通积分求导验证; Green 核与脉冲响应的共同点是单位集中输入, 区别在于所满足的边界条件.
\begin{example}求 $y''+\lambda y=0$, $y(0)=y(\pi)=0$ 的非零解存在条件.\end{example}
\begin{solution}
$\lambda<0$ 时双曲函数解代端点只能为零; $\lambda=0$ 时线性解也只能为零. $\lambda>0$ 时 $y=A\sin\sqrt\lambda t$, 末端要求 $\sin\sqrt\lambda\pi=0$. 所以特征值为 $\lambda=n^2$, $n=1,2,\ldots$, 每个对应 $A\sin nt$. $n$ 个半波来自边界约束, 而不是初值方程突然失去唯一性.
\end{solution}
\originalsection{共振时的相容条件}
取正整数 $n$ 和连续 $f$. 对 $y''+n^2y=f(t)$, $y(0)=y(\pi)=0$, 必要条件为 $\int_0^\pi f(t)\sin nt\dd t=0$, 由乘 $\sin nt$ 并积分分部两次得到. 条件也充分: 取初始导数零的特解
$y_p(t)=\int_0^t\sin(n(t-s))f(s)\dd s/n$.
它满足 $y_p(0)=0$, 而 $y_p(\pi)=(-1)^{n+1}\int_0^\pi f(s)\sin ns\dd s/n=0$. 所有解为 $y_p+C\sin nt$, 故存在但不唯一. 若条件不满足则无解. 这是边值共振与强迫振动共振的对应: 右端在核方向上的成分阻碍解存在.
\originalsection{习题与解答}
\begin{exercise}求 $y''-ty'=0$, $y(0)=0$, $y'(0)=1$ 的积分与级数形式.\end{exercise}
\begin{solution}令 $v=y'$, 得 $v=\ee^{t^2/2}$, 所以 $y=\int_0^t\ee^{s^2/2}\dd s=\sum_{k\ge0}t^{2k+1}/[(2k+1)2^kk!]$. 处处收敛, 且符合初值.\end{solution}
\begin{exercise}判定 $y''+y=1$, $y(0)=y(\pi)=0$ 是否有解.\end{exercise}
\begin{solution}相容积分 $\int_0^\pi\sin t\dd t=2\ne0$, 因而无解. 直接通解 $y=1+A\cos t+B\sin t$ 也显示 $A=-1$ 后 $y(\pi)=2$.\end{solution}

\originalchapter{Lorenz 模型与离散动力学}
\originalsection{三维耗散模型}
Lorenz 方程为
\begin{equation}\label{eq:lorenz}
x'=\sigma(y-x),\quad y'=rx-y-xz,\quad z'=xy-bz,
\qquad \sigma,r,b>0.
\end{equation}
它可由对流的低维截断得到, 本册将它作为已经给定的三维常微分方程研究, 不复制流体方程的完整推导. 参数 $r$ 控制驱动, $\sigma,b$ 控制相应时间尺度. 右端为多项式, 局部存在唯一性立即成立, 全正向存在则仍需界.
\begin{proposition}
Lorenz 系统的所有解全正向存在, 并最终进入一个有界椭球. 相空间体积按因子 $\ee^{-(\sigma+1+b)t}$ 收缩.
\end{proposition}
\begin{proof}
取 $V=rx^2+\sigma y^2+\sigma(z-2r)^2$, 展开导数时三次项与 $xy$ 项相消, 得
\[
V'=-2\sigma\{rx^2+y^2+b(z-r)^2-br^2\}.
\]
记 $Q=rx^2+y^2+b(z-r)^2$, $C=\max(1,\sigma,2\sigma/b)$. 因 $(z-2r)^2\le2(z-r)^2+2r^2$, 有 $V\le CQ+2\sigma r^2$. 故
$V'\le-(2\sigma/C)V+D$, $D=4\sigma^2r^2/C+2\sigma br^2$.
积分因子得 $V(t)\le V(0)\ee^{-at}+(D/a)(1-\ee^{-at})$, $a=2\sigma/C$. 有限时段内状态有界, 延续定理给全正向存在. 任何 $V\le R$ 且 $R>D/a$ 的椭球在边界上导数负, 正向不变, 所有轨道最终进入其中. 散度为 $-(\sigma+1+b)$, Liouville 公式给体积因子.
\end{proof}
这个界与体积收缩不证明每个轨道混沌. 平衡点、周期轨道或其他长期状态仍需按参数研究.
\originalsection{平衡点与失稳阈值}
平衡点总有 $O=(0,0,0)$. 当 $r>1$ 时另有
$C_\pm=(\pm\sqrt{b(r-1)},\pm\sqrt{b(r-1)},r-1)$.
原点线性化含根 $-b$ 及二次方程 $\lambda^2+(\sigma+1)\lambda+\sigma(1-r)=0$. 对 $0<r<1$ 两根实部负, 原点局部渐近稳定; $r>1$ 时一根正一根负, 原点具有增长方向. $r=1$ 出现零根, 线性稳定判据不决定临界点的非线性结论.

在 $C_\pm$ 处直接计算 Jacobian 的特征多项式得
\[
\lambda^3+(\sigma+b+1)\lambda^2+b(\sigma+r)\lambda+2\sigma b(r-1).
\]
\begin{lemma}[三次多项式判据]
若 $A,B,C>0$, 则 $z^3+Az^2+Bz+C$ 的全部根实部严格负, 当且仅当 $AB>C$.
\end{lemma}
\begin{proof}
当 $C=AB$, 多项式为 $(z+A)(z^2+B)$, 含虚根. 固定 $A,B>0$, 从 $C=0$ 开始增加 $C$. $C=0$ 时零根的导数 $\dd z/\dd C=-1/B<0$, 另外两根由 $z^2+Az+B=0$ 给出且实部负, 所以足够小正 $C$ 时稳定. 根连续依赖系数, 要离开左半平面必须穿越虚轴. 代 $z=\ii\omega$ 得 $C=A\omega^2$, $\omega(B-\omega^2)=0$. 对 $C>0$ 不可能 $\omega=0$, 唯一穿越值是 $C=AB$. 在该值 $z=\ii\sqrt B$ 处,
$\operatorname{Re}(\dd z/\dd C)=\operatorname{Re}[-1/(-2B+2A\ii\sqrt B)]>0$, 所以穿越后两根进入右半平面. 无其他穿越值, 故对全部 $C>AB$ 不稳定. 根的连续性可由有限维多项式根连续性这一代数先修保证; 穿越处根简单, 隐函数求导合法.
\end{proof}
对 $r>1$, 判据化为
$(\sigma+b+1)(\sigma+r)>2\sigma(r-1)$.
若 $\sigma\le b+1$, 此条件对所有 $r>1$ 成立. 若 $\sigma>b+1$, 稳定范围是
\[
1<r<r_H=\frac{\sigma(\sigma+b+3)}{\sigma-b-1}.
\]
经典参数 $\sigma=10,b=8/3$ 给 $r_H=470/19\approx24.74$. 超过此值两个非零平衡点失去线性稳定性, 这不是已经证明系统出现混沌. 一般 Hopf 分岔的方向还需非线性系数分析, 本册不从一个特征根计算越级推断.
\originalsection{离散迭代与固定点}
采样或返回映射产生 $u_{n+1}=F(u_n)$. 它的线性稳定性由导数的模决定, 对应连续时间中的指数增长因子. 差分方程是一种数学对象, 不是将 ODE 数值步长随意变大后得到的同一精确动力系统.
\begin{proposition}
$C^1$ 一维迭代中, 固定点 $u_*$ 若 $|F'(u_*)|<1$, 则局部渐近稳定; 若 $|F'(u_*)|>1$, 则不稳定.
\end{proposition}
\begin{proof}在小邻域中将导数模夹在1以下或以上, 应用均值定理和迭代估计, 与第11章返回映射证明完全相同.\end{proof}
\begin{example}分析 logistic 迭代 $u_{n+1}=a u_n(1-u_n)$, $0<a\le4$, 状态在 $[0,1]$.\end{example}
\begin{solution}
$F([0,1])\subset[0,a/4]\subset[0,1]$. 固定点为0及 $u_*=1-1/a$ (后者在 $a\ge1$ 时属于状态区间). 在零处导数为 $a$, $0<a<1$ 吸引, $a>1$ 排斥. 非零点导数为 $2-a$, 在 $1<a<3$ 吸引. $a=1$ 时 $0<u<1$ 的迭代严格下降, 极限只能是零, 所以零点在状态区间内仍渐近稳定. $a=3$ 时令 $q=u-2/3$, 一次映射为 $q\mapsto-q-3q^2$, 二次映射为 $q\mapsto q-18q^3-27q^4$. 足够小非零 $q$ 经二次映射保持符号且模变小, 极限由固定点方程只能为零; 中间奇数次也连续趋零, 所以该临界固定点仍局部渐近稳定, 尽管线性判据不再够用. 解 $F(F(u))=u$ 并去掉两个固定点因子, 得二周期候选
\[
u_\pm=\frac{a+1\pm\sqrt{(a-3)(a+1)}}{2a},\qquad a>3.
\]
它们互相映射, 二周期乘子 $F'(u_+)F'(u_-)=-a^2+2a+4$. 因而 $3<a<1+\sqrt6$ 时该二周期吸引. 这证明第一次倍周期后的明确区间, 没有凭图像证明整个无限倍周期级联.
\end{solution}
线性常系数差分方程 $u_{n+1}-a u_n=g_n$ 的解为 $u_n=a^nu_0+\sum_{k=0}^{n-1}a^{n-1-k}g_k$. 直接归纳即可核验, 它对应连续卷积响应. 单边 $Z$ 变换定义为 $U(z)=\sum_{n\ge0}u_nz^{-n}$, 在收敛区域有 $\mathcal Z(u_{n+1})=z(U-u_0)$. 离散稳定极点在单位圆内, 连续稳定极点在左半平面. 对精确采样 $X(nh)$, $z=\ee^{sh}$ 将后者映到前者; 不同 $s$ 的虚部相差 $2\pi/h$ 会给相同采样因子, 这是采样不能识别所有连续频率的原因.
\originalsection{伸长、折叠与 Lyapunov 指数}
沿离散轨道的线性扰动为 $\delta u_n\approx\prod_{j=0}^{n-1}F'(u_j)\delta u_0$. 若极限存在, 定义最大一维 Lyapunov 指数
$\lambda=\lim_{n\to\infty}n^{-1}\sum_{j=0}^{n-1}\log|F'(u_j)|$; 遇导数零应按扩展实值解释. 连续流则对变分矩阵定义 $\limsup_{t\to\infty}t^{-1}\log\|D\phi_t v\|$ 的方向增长率. 正指数表示线性扰动的长期指数增长, 不能单独当作每条轨道“混沌”的完整定义.

在 $a=4$ 时代 $u=\sin^2\pi\theta$ 得 $F(u)=\sin^2(2\pi\theta)$, 与倍角映射 $\theta\mapsto2\theta\pmod1$ 半共轭. 它是半共轭而非一一共轭, 因为 $\theta$ 与 $1-\theta$ 给同一 $u$. 对二进制符号展开, 倍角将小数点后的首位移去, 除去二进制双表示的点后, 初始第 $k$ 位的差异在 $k-1$ 次迭代后进入首位; 再迭代一次该位被移去. 这种位移说明小差异可能随迭代放大, 不能推出任意两点在固定步数内必定分离; 半共轭还可能将对称点映到同一状态. 这解释伸长和折叠的机制, 同时这个系统仍有固定点及周期轨道, 不能声称所有初值均呈同样随机行为.
\originalsection{习题与解答}
\begin{exercise}求 $u_{n+1}=u_n/2+1$, $u_0=c$ 的解与单边 $Z$ 变换.\end{exercise}
\begin{solution}$u_n=2+(c-2)2^{-n}$, 所以 $U(z)=2z/(z-1)+(c-2)z/(z-1/2)$, 在 $|z|>1$ 收敛. 极点1对应常量输入带来的终值2, 当 $c\ne2$ 时自由扰动极点 $1/2$ 对应稳定模式; $c=2$ 时此项约去, 输出恒为2.\end{solution}
\begin{exercise}在 Lorenz 原点取 $r<1$, 为什么严格负散度并不是稳定性的完整证明?\end{exercise}
\begin{solution}负散度仅给三个增长率之和为负, 仍可能有一个正根. 稳定须核验二次特征多项式常数项 $\sigma(1-r)>0$ 及正一次系数, 再加根 $-b<0$. $r>1$ 时散度相同但有一个正根, 直接给出反例.\end{solution}

% MIT 18.03 Spring 2010. 原题编号保留; 解答与原题分离.
\originalchapter{题库一: 一阶微分方程}
\originalsection{1A: 入门与变量分离}
\begin{exercise}\label{cw:bank-1A-1}
原题1A-1. 验证下列方程具有所列解, 其中 $c_1,c_2,a$ 为常数:
\begin{enumerate}[label=\alph*)]
\item $y''-2y'+y=0$, $y=c_1e^x+c_2xe^x$.
\item $xy'+y=x\sin x$, $y=(\sin x+a)/x-\cos x$.
\end{enumerate}
\end{exercise}
\begin{solution}
官方答案译审. (a) $y'=e^x(c_1+c_2x+c_2)$, $y''=e^x(c_1+c_2x+2c_2)$, 代入即为零. (b) 在 $x\ne0$ 上, $xy=\sin x+a-x\cos x$, 故 $(xy)'=x\sin x$.
\end{solution}
\begin{exercise}\label{cw:bank-1A-2}
原题1A-2. 下列函数族实际依赖几个任意常数? $a,b,c,d,k,c_i$ 均为常数, 表面出现的常数不一定独立.
\begin{enumerate}[label=\alph*)]
\item $c_1e^{kx}$. \item $c_1e^{x+a}$. \item $c_1+c_2\cos2x+c_3\cos^2x$. \item $\ln(ax+b)+\ln(cx+d)$.
\end{enumerate}
\end{exercise}
\begin{solution}
官方答案译审. 依次为 $2,1,2,3$. (a) 非零成员的振幅与指数率可独立改变. (b) 合并为 $(c_1e^a)e^x$. (c) 用 $\cos2x=2\cos^2x-1$, 得 $(c_1-c_2)+(2c_2+c_3)\cos^2x$. (d) 在两个对数均有定义的区间, 合并为 $\ln[acx^2+(ad+bc)x+bd]$, 只涉及三个组合参数; 这三个组合在一般参数点能独立变化. 个别退化成员不会增加整个函数族的参数数目.
\end{solution}
\begin{exercise}\label{cw:bank-1A-3}
原题1A-3. 写出含定积分的显式解:
\begin{enumerate}[label=\alph*)]
\item $y'=1/(y^2\ln x)$, $y(2)=0$.
\item $y'=ye^x/x$, $y(1)=1$.
\end{enumerate}
\end{exercise}
\begin{solution}
依据官方答案补核. (a) 原题初值处右端未定义, 因而没有包含 $x=2$ 的经典解. 官方的分离计算给出形式曲线
\[
 y(x)=\sqrt[3]{3\int_2^x\frac{dt}{\ln t}}.
\]
它在 $2$ 附近连续, 在 $x\ne2$ 的两侧分别满足方程, 但 $y'(x)$ 在 $2$ 处无界, 不能称为此初值问题的经典解. (b) $\ln y=\int_1^x e^t/t\,dt$, 故 $y=\exp(\int_1^x e^t/t\,dt)$, $x>0$.
\end{solution}
\begin{exercise}\label{cw:bank-1A-4}
原题1A-4. 解初值问题:
\begin{enumerate}[label=\alph*)]
\item $y'=(xy+x)/y$, $y(2)=0$.
\item $du/dt=\sin t\cos^2u$, $u(0)=0$.
\end{enumerate}
\end{exercise}
\begin{solution}
依据官方答案补核. (a) 初值处方程未定义, 无经典解. 在 $y\ne0,-1$ 上分离得 $y-\ln|y+1|=x^2/2+C$; 经过 $(2,0)$ 的形式关系是 $y-\ln|y+1|=x^2/2-2$. 左端在 $y=0$ 附近为 $y^2/2+O(y^3)$, 故这里只产生 $x\ge2$ 一侧的分支, 其起点切线竖直. (b) $\sec^2u\,du=\sin t\,dt$, 从初值积分得 $\tan u=1-\cos t$, 因而 $u=\arctan(1-\cos t)$.
\end{solution}
\begin{exercise}\label{cw:bank-1A-5}
原题1A-5. 用变量分离求通解:
\begin{enumerate}[label=\alph*)]
\item $(y^2-2y)\,dx+x^2\,dy=0$.
\item $x\,dv/dx=\sqrt{1-v^2}$.
\item $y'=((y-1)/(x+1))^2$.
\item $dx/dt=\sqrt{1+x}/(t^2+4)$.
\end{enumerate}
\end{exercise}
\begin{solution}
官方答案译审并补齐分离时遗漏的常数解. (a) $\frac12\ln|(y-2)/y|=1/x+C$, 故 $y=2/(1-Ce^{2/x})$, 并有 $y=0$; 以上取 $x\ne0$ 的解区间. (b) $\arcsin v=\ln|x|+C$, 即 $v=\sin(\ln|x|+C)$, 必须限制在 $\cos(\ln|x|+C)\ge0$ 的区间, 因为原式平方根非负. 常数解 $v=\pm1$ 也须保留, 并可在端点与非常数分支粘接. (c) $-1/(y-1)=-1/(x+1)+C$, 即 $y=1+(x+1)/(1-C(x+1))$; 另有 $y=1$, 且 $x\ne-1$. (d) $2\sqrt{1+x}=\frac12\arctan(t/2)+C$, 所以 $x=[\frac14\arctan(t/2)+C]^2-1$, 要求括号非负; 另有 $x=-1$, 可在非负分支的零点粘接.
\end{solution}
\originalsection{1B: 一阶方程的标准方法}
\begin{exercise}\label{cw:bank-1B-1}
原题1B-1. 判定是否恰当, 对恰当方程求通解:
\begin{enumerate}[label=\alph*)]
\item $3x^2y\,dx+(x^3+y^3)\,dy=0$.
\item $(x^2-y^2)\,dx+(y^2-x^2)\,dy=0$.
\item $ve^{uv}\,du+ue^{uv}\,dv=0$.
\item $2xy\,dx-x^2\,dy=0$.
\end{enumerate}
\end{exercise}
\begin{solution}
官方答案译审. 比较 $M_y$ 与 $N_x$: (a) 均为 $3x^2$, 势函数 $x^3y+y^4/4$, 解为 $x^3y+y^4/4=C$. (b) 分别为 $-2y,-2x$, 不恰当. (c) 均为 $(1+uv)e^{uv}$, 势函数 $e^{uv}$, 解为 $e^{uv}=C$, 等价于 $uv=C_1$. (d) 分别为 $2x,-2x$, 不恰当.
\end{solution}
\begin{exercise}\label{cw:bank-1B-2}
原题1B-2. 求积分因子并解方程:
\begin{enumerate}[label=\alph*)]
\item $2x\,dx+(x^2/y)\,dy=0$.
\item $y\,dx-(x+y)\,dy=0$, $y(1)=1$.
\item $(t^2+4)\,dt+t\,dx=x\,dt$.
\item $u(du-dv)+v(du+dv)=0$, $v(0)=1$.
\end{enumerate}
\end{exercise}
\begin{solution}
官方答案译审. (a) 乘 $y$ 后为 $d(x^2y)=0$, 故 $x^2y=C$, 取原式有定义处. (b) 乘 $y^{-2}$ 后为 $d(x/y)-d(\ln|y|)=0$, 初值给 $x/y-\ln y=1$, 即 $x=y(1+\ln y)$, $y>0$. (c) 乘 $t^{-2}$, 得 $d(x/t)=(1+4/t^2)dt$, 所以 $x=t^2-4+Ct$, $t\ne0$. (d) 乘 $(u^2+v^2)^{-1}$, 在初值附近以 $u=r\sin\theta,v=r\cos\theta$ 表示, 方程变为 $d\ln r+d\theta=0$. 初值给 $\ln r+\theta=0$, 即 $r=e^{-\theta}$. 角度须连续选取; 这也给出经过 $(0,1)$ 的螺线.
\end{solution}
\begin{exercise}\label{cw:bank-1B-3}
原题1B-3. 解齐次方程:
\begin{enumerate}[label=\alph*)]
\item $y'=(2y-x)/(y+4x)$.
\item $dw/du=2uw/(u^2-w^2)$.
\item $xy\,dy-y^2\,dx=x\sqrt{x^2-y^2}\,dx$.
\end{enumerate}
\end{exercise}
\begin{solution}
官方答案译审并核对定义域. (a) 令 $z=y/x$, 则 $xz'=-(z+1)^2/(z+4)$. 积分得 $\ln|z+1|-3/(z+1)=-\ln|x|+C$, 即 $\ln|x+y|-3x/(x+y)=C$. 分离遗漏的 $y=-x$ 也是解. (b) 令 $z=w/u$, 则 $uz'=z(1+z^2)/(1-z^2)$; 积分得 $\ln|z/(1+z^2)|=\ln|u|+C$, 故 $w/(u^2+w^2)=C$. 零解含在其中, 原方程排除 $u^2=w^2$. (c) 在 $x>0$ 令 $z=y/x$, 则 $xz z'=\sqrt{1-z^2}$, 得 $\sqrt{1-z^2}=C-\ln x$. 因而 $y=\pm x\sqrt{1-(C-\ln x)^2}$, 经典非边界分支须 $0<C-\ln x<1$, 即 $|y|<|x|$ 且 $y\ne0$. 等号 $C-\ln x=1$ 对应 $y=0$ 的竖直切线端点, 不能包含为可微的 $y(x)$ 解; $C-\ln x=0$ 可与边界解相接. 在 $x<0$ 上相应为 $\sqrt{1-z^2}=C+\ln|x|$, 同样在严格内侧 $0<C+\ln|x|<1$ 取经典非边界分支. 边界直线 $y=\pm x$ 也满足原微分关系.
\end{solution}
\begin{exercise}\label{cw:bank-1B-4}
原题1B-4. 证明代换 $u=y/x^n$ 对适当的 $n$ 可把 $y'=(4+xy^2)/(x^2y)$ 化为变量可分离方程, 并求解. 提示: 先用 $u,x$ 表示 $y,y'$.
\end{exercise}
\begin{solution}
官方答案译审. 写成 $y'=y/x+4/(x^2y)$. 取 $n=1$, $y=xu$, 则 $xu'=4/(x^3u)$, 即 $u u'=4/x^4$. 所以 $u^2=C-8/(3x^3)$, $y^2=Cx^2-8/(3x)$. 取右端为正的区间和任一固定符号平方根, 且 $x,y\ne0$.
\end{solution}
\begin{exercise}\label{cw:bank-1B-5}
原题1B-5. 求通解或满足初值的解:
\begin{enumerate}[label=\alph*)]
\item $xy'+2y=x$.
\item $dx/dt-x\tan t=t/\cos t$, $x(0)=0$.
\item $(x^2-1)y'=1-2xy$.
\item $3v\,dt=t(dt-dv)$, $v(1)=1/4$.
\end{enumerate}
\end{exercise}
\begin{solution}
官方答案译审. (a) 乘积分因子 $x^2$: $(x^2y)'=x^2$, 得 $y=x/3+C/x^2$. (b) $(x\cos t)'=t$, 得 $x=t^2/(2\cos t)$, 初值区间 $|t|<\pi/2$. (c) $[(x^2-1)y]'=1$, 得 $y=(x+C)/(x^2-1)$, 可消去的奇点需另检查原式. (d) $v'+3v/t=1$, 故 $(t^3v)'=t^3$, $v=t/4+C/t^3$. 初值使 $C=0$.
\end{solution}
\begin{exercise}\label{cw:bank-1B-6}
原题1B-6. 设 $a>0$, $r(t)\to0$ ($t\to\infty$). 证明 $x'+ax=r(t)$ 的每个解均趋于零. 提示: 可用洛必达法则.
\end{exercise}
\begin{solution}
官方思路译审并补齐极限论证. 积分因子给
\[
 x(t)=e^{-a(t-t_0)}x(t_0)+\int_{t_0}^t e^{-a(t-s)}r(s)\,ds.
\]
给定 $\varepsilon>0$, 取 $T$ 使 $s\ge T$ 时 $|r(s)|\le a\varepsilon$. 积分的 $[t_0,T]$ 部分随 $t\to\infty$ 趋零, $[T,t]$ 部分绝对值至多 $\varepsilon(1-e^{-a(t-T)})$. 所以 $\limsup|x(t)|\le\varepsilon$, 再令 $\varepsilon\to0$. 这避免在积分分子可能不趋于无穷时误用洛必达法则.
\end{solution}
\begin{exercise}\label{cw:bank-1B-7}
原题1B-7. 解 $y'=y/(y^3+x)$. 提示: 考察 $dx/dy$.
\end{exercise}
\begin{solution}
官方答案译审. 对非零解交换变量: $dx/dy-x/y=y^2$. 乘 $1/y$ 得 $d(x/y)/dy=y$, 故 $x=y^3/2+Cy$. 此关系只在 $y^3+x\ne0$ 且能局部表示为 $y(x)$ 时成立. 另有 $y=0$ ($x\ne0$).
\end{solution}
\begin{exercise}\label{cw:bank-1B-8}
原题1B-8. Bernoulli 方程 $y'+p(x)y=q(x)y^n$, $n\ne1$. 证明代换 $u=y^{1-n}$ 将其化为线性方程. 提示: 先除以 $y^n$.
\end{exercise}
\begin{solution}
官方答案译审. 在实幂有定义且 $y\ne0$ 的区间, $u'=(1-n)y^{-n}y'$, 所以 $u'+(1-n)pu=(1-n)q$. 代换前须另检查零解; 它不一定属于实幂方程的定义域.
\end{solution}
\begin{exercise}\label{cw:bank-1B-9}
原题1B-9. 用1B-8的方法解:
\begin{enumerate}[label=\alph*)]
\item $y'+y=2xy^2$. \item $x^2y'-y^3=xy$.
\end{enumerate}
\end{exercise}
\begin{solution}
官方答案译审. (a) $u=1/y$ 满足 $u'-u=-2x$, 故 $u=2x+2+Ce^x$, $y=(2x+2+Ce^x)^{-1}$; 另有 $y=0$. (b) $u=y^{-2}$ 满足 $u'+2u/x=-2/x^2$, 得 $(x^2u)'=-2$, 故 $y=\pm x/\sqrt{C-2x}$, 在 $C-2x>0$, $x\ne0$ 上取固定分支; 另有 $y=0$.
\end{solution}
\begin{exercise}\label{cw:bank-1B-10}
原题1B-10. Riccati 方程 $y'=A(x)+B(x)y+C(x)y^2$ 一般不能用初等函数求解.
\begin{enumerate}[label=\alph*)]
\item 已知一个解 $y_1$, 证明通解可写成 $y=y_1+u$, 其中 $u$ 满足某个 Bernoulli 方程.
\item 据此解 $y'=1-x^2+y^2$.
\end{enumerate}
\end{exercise}
\begin{solution}
官方答案译审. (a) 代入并减去 $y_1$ 的方程, 得 $u'=(B+2Cy_1)u+Cu^2$. (b) $y_1=x$ 是特解. 若 $u\ne0$, 令 $w=1/u$, 得 $w'+2xw=-1$, 故
\[
 y=x+\frac{e^{x^2}}{K-\int_0^x e^{t^2}\,dt}.
\]
分母非零时有效; 原特解 $y=x$ 也要保留.
\end{solution}
\begin{exercise}\label{cw:bank-1B-11}
原题1B-11. 解下列二阶自治方程. 自治表示独立变量没有显式出现在方程中.
\begin{enumerate}[label=\alph*)]
\item $y''=a^2y$. \item $yy''=(y')^2$. \item $y''=y'(1+3y^2)$, $y(0)=1$, $y'(0)=2$.
\end{enumerate}
\end{exercise}
\begin{solution}
官方答案译审并补齐退化解. 以 $v=y'$ 为 $y$ 的函数, 则 $y''=v\,dv/dy$. (a) 积分得 $v^2=a^2y^2+C$, 最终 $y=Ae^{ax}+Be^{-ax}$ ($a\ne0$); $a=0$ 时 $y=Ax+B$. (b) 在 $y\ne0$ 时 $(y'/y)'=0$, 所以 $y=Ae^{kx}$. 零解也满足; 非零解不能有限时达到零, 因而没有额外跨零分支. (c) $dv/dy=1+3y^2$, 得 $v=y+y^3+C$; 初值给 $C=0$. 分离后 $y/\sqrt{1+y^2}=e^x/\sqrt2$, 因而 $y=e^x/\sqrt{2-e^{2x}}$, 最大初值区间为 $x<\frac12\ln2$.
\end{solution}
\begin{exercise}\label{cw:bank-1B-12}
原题1B-12. 只说明各式的类型和求解方法, 不必求解; 某些式可用多种方法.
\begin{enumerate}[label=\arabic*.]
\item $(x^3+y)dx+xdy=0$. \item $dy/dt+2ty-e^{-t}=0$.
\item $y'=(x^2-y^2)/(5xy)$. \item $(1+2p)dq+(2-q)dp=0$.
\item $\cos x\,dy=(y\sin x+e^x)dx$. \item $x\tan y\,y'=-1$.
\item $y'=y/x+1/y$. \item $dv/du=e^{2u+3v}$.
\item $xy'=y+xe^{y/x}$. \item $xy'-y=x^2\sin x$.
\item $y'=(x+e^y)^{-1}$. \item $y'+2y/x-y^2/x=0$.
\item $dx/dy=-x(2x^2y+\cos y)/(3x^2y^2+\sin y)$.
\item $y'+3y=e^{-3t}$. \item $x\,dy/dx-y=\sqrt{x^2+y^2}$.
\item $(y'-1)/x^2=1$. \item $xy'-2y+y^2=x^4$.
\item $y''=y(y+1)/y'$. \item $t\,ds/dt=s(1-\ln t+\ln s)$.
\item $dy/dx=(3-2y)/(2x+y+1)$. \item $x^2y'+xy+y^2=0$.
\item $y'\tan(x+y)=1-\tan(x+y)$. \item $y\,ds-3s\,dy=y^4dy$.
\item $du=-(1+u\cos^2t)/(t\cos^2t)\,dt$.
\item $y'+y^2+(2x+1)y+1+x+x^2=0$. \item $y''+x^2y'+3x^3=\sin x$.
\end{enumerate}
\end{exercise}
\begin{solution}
官方答案译审. 按题号依次为:
\begin{enumerate}[label=\arabic*.]
\item 恰当, 或以 $y$ 为未知量的线性式.
\item 线性, 积分因子 $e^{t^2}$.
\item 齐次, 令 $y=xz$.
\item 可分离; 也分别对 $p$ 或 $q$ 线性.
\item 恰当, 或线性.
\item 可分离.
\item Bernoulli, $n=-1$, 令 $u=y^2$.
\item 可分离.
\item 齐次, 令 $z=y/x$.
\item 线性, 积分因子 $1/x$.
\item 交换变量后 $dx/dy=x+e^y$, 线性.
\item 可分离或 Bernoulli.
\item 写成 $(3x^2y^2+\sin y)dx+x(2x^2y+\cos y)dy=0$, 恰当.
\item 对 $t$ 的线性式.
\item 在 $x>0$ 或 $x<0$ 分别令 $y=xz$, 化为可分离式.
\item 直接积分.
\item Riccati, 已知特解 $y=x^2$, 再令 $y=x^2+u$.
\item 自治降阶, 令 $v=y'$, 得 $v^2dv/dy=y(y+1)$.
\item 在 $s,t>0$ 上令 $s=tz$, 利用 $\ln s-\ln t=\ln z$ 分离.
\item 写成 $(2y-3)dx+(2x+y+1)dy=0$, 恰当.
\item Bernoulli, $n=2$.
\item 令 $u=x+y$, 则 $u'=\cot u$, 可分离.
\item 以 $y$ 为独立变量、$s$ 为未知函数的线性式.
\item 以 $t$ 为独立变量、$u$ 为未知函数的线性式.
\item Riccati, 特解 $y=-x$; 或令 $u=x+y$, 得 $u'+u^2+u=0$.
\item 令 $v=y'$, 得一阶线性式 $v'+x^2v=\sin x-3x^3$.
\end{enumerate}
\end{solution}
\originalsection{1C: 图解与数值方法}
\begin{exercise}\label{cw:bank-1C-1}
原题1C-1. 对各式选约五条等倾线绘方向场, 在两坐标均为 $[-2,2]$ 的正方形内画若干积分曲线, 并完成附加要求.
\begin{enumerate}[label=\alph*)]
\item $y'=-y/x$: 求精确解并比较.
\item $y'=2x+y$: 找出图形同时是等倾线的解, 并用计算验证.
\item $y'=x-y$: 同(b).
\item $y'=x^2+y^2-1$.
\item $y'=1/(x+y)$: 改用 $[-3,3]^2$, 画经过 $(0,0),(-1,1),(0,-2)$ 的积分曲线. 它们会越过 $y=-x-1$ 吗? 用解曲线不相交原理解释.
\end{enumerate}
\end{exercise}
\begin{solution}
官方答案译审并补核奇点. 等倾线分别为 (a) $y=-cx$, (b) $y=c-2x$, (c) $y=x-c$, (d) $x^2+y^2=1+c$, (e) $y=-x+1/c$. (a) $y=C/x$. (b) 解族 $y=Ce^x-2x-2$, $C=0$ 给直线 $y=-2x-2$, 其斜率和右端都为 $-2$. (c) $y=Ce^{-x}+x-1$, $C=0$ 给等倾线解 $y=x-1$.

(e) 令 $z=x+y$, 则 $z'=1+1/z$, 因而 $x=z-\ln|z+1|+C$, 即 $y-\ln|x+y+1|=C_1$. 经过三点的形式常数依次为 $0,1,-2$. 前两点位于 $x+y=0$ 的奇线, 没有以该点为初值的经典解; 图中画的是该点附近的几何分支, 起点有竖直切线. $y=-x-1$ 本身是经典解, 其附近右端光滑, 所以其他经典分支不能与它相交.
\begin{center}\begin{tikzpicture}\begin{axis}[width=6.8cm,height=6.8cm,axis lines=middle,xmin=-2,xmax=2,ymin=-2,ymax=2,xlabel=$x$,ylabel=$y$,title={(a)},clip=true]
\addplot[black,dashed,domain=-2:2] {--2*x};
\addplot[black,dashed,domain=-2:2] {--1*x};
\addplot[black,dashed,domain=-2:2] {-0*x};
\addplot[black,dashed,domain=-2:2] {-1*x};
\addplot[black,dashed,domain=-2:2] {-2*x};
\draw[gray] (axis cs:-2.08485,-1.91515)--(axis cs:-1.91515,-2.08485);
\draw[gray] (axis cs:-2.09600,-1.42800)--(axis cs:-1.90400,-1.57200);
\draw[gray] (axis cs:-2.10733,-0.94633)--(axis cs:-1.89267,-1.05367);
\draw[gray] (axis cs:-2.11642,-0.47090)--(axis cs:-1.88358,-0.52910);
\draw[gray] (axis cs:-2.12000,0.00000)--(axis cs:-1.88000,0.00000);
\draw[gray] (axis cs:-2.11642,0.47090)--(axis cs:-1.88358,0.52910);
\draw[gray] (axis cs:-2.10733,0.94633)--(axis cs:-1.89267,1.05367);
\draw[gray] (axis cs:-2.09600,1.42800)--(axis cs:-1.90400,1.57200);
\draw[gray] (axis cs:-2.08485,1.91515)--(axis cs:-1.91515,2.08485);
\draw[gray] (axis cs:-1.57200,-1.90400)--(axis cs:-1.42800,-2.09600);
\draw[gray] (axis cs:-1.58485,-1.41515)--(axis cs:-1.41515,-1.58485);
\draw[gray] (axis cs:-1.59985,-0.93344)--(axis cs:-1.40015,-1.06656);
\draw[gray] (axis cs:-1.61384,-0.46205)--(axis cs:-1.38616,-0.53795);
\draw[gray] (axis cs:-1.62000,0.00000)--(axis cs:-1.38000,0.00000);
\draw[gray] (axis cs:-1.61384,0.46205)--(axis cs:-1.38616,0.53795);
\draw[gray] (axis cs:-1.59985,0.93344)--(axis cs:-1.40015,1.06656);
\draw[gray] (axis cs:-1.58485,1.41515)--(axis cs:-1.41515,1.58485);
\draw[gray] (axis cs:-1.57200,1.90400)--(axis cs:-1.42800,2.09600);
\draw[gray] (axis cs:-1.05367,-1.89267)--(axis cs:-0.94633,-2.10733);
\draw[gray] (axis cs:-1.06656,-1.40015)--(axis cs:-0.93344,-1.59985);
\draw[gray] (axis cs:-1.08485,-0.91515)--(axis cs:-0.91515,-1.08485);
\draw[gray] (axis cs:-1.10733,-0.44633)--(axis cs:-0.89267,-0.55367);
\draw[gray] (axis cs:-1.12000,0.00000)--(axis cs:-0.88000,0.00000);
\draw[gray] (axis cs:-1.10733,0.44633)--(axis cs:-0.89267,0.55367);
\draw[gray] (axis cs:-1.08485,0.91515)--(axis cs:-0.91515,1.08485);
\draw[gray] (axis cs:-1.06656,1.40015)--(axis cs:-0.93344,1.59985);
\draw[gray] (axis cs:-1.05367,1.89267)--(axis cs:-0.94633,2.10733);
\draw[gray] (axis cs:-0.52910,-1.88358)--(axis cs:-0.47090,-2.11642);
\draw[gray] (axis cs:-0.53795,-1.38616)--(axis cs:-0.46205,-1.61384);
\draw[gray] (axis cs:-0.55367,-0.89267)--(axis cs:-0.44633,-1.10733);
\draw[gray] (axis cs:-0.58485,-0.41515)--(axis cs:-0.41515,-0.58485);
\draw[gray] (axis cs:-0.62000,0.00000)--(axis cs:-0.38000,0.00000);
\draw[gray] (axis cs:-0.58485,0.41515)--(axis cs:-0.41515,0.58485);
\draw[gray] (axis cs:-0.55367,0.89267)--(axis cs:-0.44633,1.10733);
\draw[gray] (axis cs:-0.53795,1.38616)--(axis cs:-0.46205,1.61384);
\draw[gray] (axis cs:-0.52910,1.88358)--(axis cs:-0.47090,2.11642);
\draw[gray] (axis cs:0.47090,-2.11642)--(axis cs:0.52910,-1.88358);
\draw[gray] (axis cs:0.46205,-1.61384)--(axis cs:0.53795,-1.38616);
\draw[gray] (axis cs:0.44633,-1.10733)--(axis cs:0.55367,-0.89267);
\draw[gray] (axis cs:0.41515,-0.58485)--(axis cs:0.58485,-0.41515);
\draw[gray] (axis cs:0.38000,0.00000)--(axis cs:0.62000,0.00000);
\draw[gray] (axis cs:0.41515,0.58485)--(axis cs:0.58485,0.41515);
\draw[gray] (axis cs:0.44633,1.10733)--(axis cs:0.55367,0.89267);
\draw[gray] (axis cs:0.46205,1.61384)--(axis cs:0.53795,1.38616);
\draw[gray] (axis cs:0.47090,2.11642)--(axis cs:0.52910,1.88358);
\draw[gray] (axis cs:0.94633,-2.10733)--(axis cs:1.05367,-1.89267);
\draw[gray] (axis cs:0.93344,-1.59985)--(axis cs:1.06656,-1.40015);
\draw[gray] (axis cs:0.91515,-1.08485)--(axis cs:1.08485,-0.91515);
\draw[gray] (axis cs:0.89267,-0.55367)--(axis cs:1.10733,-0.44633);
\draw[gray] (axis cs:0.88000,0.00000)--(axis cs:1.12000,0.00000);
\draw[gray] (axis cs:0.89267,0.55367)--(axis cs:1.10733,0.44633);
\draw[gray] (axis cs:0.91515,1.08485)--(axis cs:1.08485,0.91515);
\draw[gray] (axis cs:0.93344,1.59985)--(axis cs:1.06656,1.40015);
\draw[gray] (axis cs:0.94633,2.10733)--(axis cs:1.05367,1.89267);
\draw[gray] (axis cs:1.42800,-2.09600)--(axis cs:1.57200,-1.90400);
\draw[gray] (axis cs:1.41515,-1.58485)--(axis cs:1.58485,-1.41515);
\draw[gray] (axis cs:1.40015,-1.06656)--(axis cs:1.59985,-0.93344);
\draw[gray] (axis cs:1.38616,-0.53795)--(axis cs:1.61384,-0.46205);
\draw[gray] (axis cs:1.38000,0.00000)--(axis cs:1.62000,0.00000);
\draw[gray] (axis cs:1.38616,0.53795)--(axis cs:1.61384,0.46205);
\draw[gray] (axis cs:1.40015,1.06656)--(axis cs:1.59985,0.93344);
\draw[gray] (axis cs:1.41515,1.58485)--(axis cs:1.58485,1.41515);
\draw[gray] (axis cs:1.42800,2.09600)--(axis cs:1.57200,1.90400);
\draw[gray] (axis cs:1.91515,-2.08485)--(axis cs:2.08485,-1.91515);
\draw[gray] (axis cs:1.90400,-1.57200)--(axis cs:2.09600,-1.42800);
\draw[gray] (axis cs:1.89267,-1.05367)--(axis cs:2.10733,-0.94633);
\draw[gray] (axis cs:1.88358,-0.52910)--(axis cs:2.11642,-0.47090);
\draw[gray] (axis cs:1.88000,0.00000)--(axis cs:2.12000,0.00000);
\draw[gray] (axis cs:1.88358,0.52910)--(axis cs:2.11642,0.47090);
\draw[gray] (axis cs:1.89267,1.05367)--(axis cs:2.10733,0.94633);
\draw[gray] (axis cs:1.90400,1.57200)--(axis cs:2.09600,1.42800);
\draw[gray] (axis cs:1.91515,2.08485)--(axis cs:2.08485,1.91515);
\addplot[main,domain=-2.00:-0.08,samples=70] {-1.00/x};
\addplot[main,domain=0.08:2.00,samples=70] {-1.00/x};
\addplot[main,domain=-2.00:-0.08,samples=70] {-0.40/x};
\addplot[main,domain=0.08:2.00,samples=70] {-0.40/x};
\addplot[main,domain=-2.00:-0.08,samples=70] {0.00/x};
\addplot[main,domain=0.08:2.00,samples=70] {0.00/x};
\addplot[main,domain=-2.00:-0.08,samples=70] {0.40/x};
\addplot[main,domain=0.08:2.00,samples=70] {0.40/x};
\addplot[main,domain=-2.00:-0.08,samples=70] {1.00/x};
\addplot[main,domain=0.08:2.00,samples=70] {1.00/x};
\end{axis}\end{tikzpicture}\end{center}
\begin{center}\begin{tikzpicture}\begin{axis}[width=6.8cm,height=6.8cm,axis lines=middle,xmin=-2,xmax=2,ymin=-2,ymax=2,xlabel=$x$,ylabel=$y$,title={(b)},clip=true]
\addplot[black,dashed,domain=-2:2] {-2-2*x};
\addplot[black,dashed,domain=-2:2] {-1-2*x};
\addplot[black,dashed,domain=-2:2] {0-2*x};
\addplot[black,dashed,domain=-2:2] {1-2*x};
\addplot[black,dashed,domain=-2:2] {2-2*x};
\draw[gray] (axis cs:-2.01973,-1.88163)--(axis cs:-1.98027,-2.11837);
\draw[gray] (axis cs:-2.02147,-1.38194)--(axis cs:-1.97853,-1.61806);
\draw[gray] (axis cs:-2.02353,-0.88233)--(axis cs:-1.97647,-1.11767);
\draw[gray] (axis cs:-2.02603,-0.38286)--(axis cs:-1.97397,-0.61714);
\draw[gray] (axis cs:-2.02910,0.11642)--(axis cs:-1.97090,-0.11642);
\draw[gray] (axis cs:-2.03297,0.61538)--(axis cs:-1.96703,0.38462);
\draw[gray] (axis cs:-2.03795,1.11384)--(axis cs:-1.96205,0.88616);
\draw[gray] (axis cs:-2.04457,1.61142)--(axis cs:-1.95543,1.38858);
\draw[gray] (axis cs:-2.05367,2.10733)--(axis cs:-1.94633,1.89267);
\draw[gray] (axis cs:-1.52353,-1.88233)--(axis cs:-1.47647,-2.11767);
\draw[gray] (axis cs:-1.52603,-1.38286)--(axis cs:-1.47397,-1.61714);
\draw[gray] (axis cs:-1.52910,-0.88358)--(axis cs:-1.47090,-1.11642);
\draw[gray] (axis cs:-1.53297,-0.38462)--(axis cs:-1.46703,-0.61538);
\draw[gray] (axis cs:-1.53795,0.11384)--(axis cs:-1.46205,-0.11384);
\draw[gray] (axis cs:-1.54457,0.61142)--(axis cs:-1.45543,0.38858);
\draw[gray] (axis cs:-1.55367,1.10733)--(axis cs:-1.44633,0.89267);
\draw[gray] (axis cs:-1.56656,1.59985)--(axis cs:-1.43344,1.40015);
\draw[gray] (axis cs:-1.58485,2.08485)--(axis cs:-1.41515,1.91515);
\draw[gray] (axis cs:-1.02910,-1.88358)--(axis cs:-0.97090,-2.11642);
\draw[gray] (axis cs:-1.03297,-1.38462)--(axis cs:-0.96703,-1.61538);
\draw[gray] (axis cs:-1.03795,-0.88616)--(axis cs:-0.96205,-1.11384);
\draw[gray] (axis cs:-1.04457,-0.38858)--(axis cs:-0.95543,-0.61142);
\draw[gray] (axis cs:-1.05367,0.10733)--(axis cs:-0.94633,-0.10733);
\draw[gray] (axis cs:-1.06656,0.59985)--(axis cs:-0.93344,0.40015);
\draw[gray] (axis cs:-1.08485,1.08485)--(axis cs:-0.91515,0.91515);
\draw[gray] (axis cs:-1.10733,1.55367)--(axis cs:-0.89267,1.44633);
\draw[gray] (axis cs:-1.12000,2.00000)--(axis cs:-0.88000,2.00000);
\draw[gray] (axis cs:-0.53795,-1.88616)--(axis cs:-0.46205,-2.11384);
\draw[gray] (axis cs:-0.54457,-1.38858)--(axis cs:-0.45543,-1.61142);
\draw[gray] (axis cs:-0.55367,-0.89267)--(axis cs:-0.44633,-1.10733);
\draw[gray] (axis cs:-0.56656,-0.40015)--(axis cs:-0.43344,-0.59985);
\draw[gray] (axis cs:-0.58485,0.08485)--(axis cs:-0.41515,-0.08485);
\draw[gray] (axis cs:-0.60733,0.55367)--(axis cs:-0.39267,0.44633);
\draw[gray] (axis cs:-0.62000,1.00000)--(axis cs:-0.38000,1.00000);
\draw[gray] (axis cs:-0.60733,1.44633)--(axis cs:-0.39267,1.55367);
\draw[gray] (axis cs:-0.58485,1.91515)--(axis cs:-0.41515,2.08485);
\draw[gray] (axis cs:-0.05367,-1.89267)--(axis cs:0.05367,-2.10733);
\draw[gray] (axis cs:-0.06656,-1.40015)--(axis cs:0.06656,-1.59985);
\draw[gray] (axis cs:-0.08485,-0.91515)--(axis cs:0.08485,-1.08485);
\draw[gray] (axis cs:-0.10733,-0.44633)--(axis cs:0.10733,-0.55367);
\draw[gray] (axis cs:-0.12000,0.00000)--(axis cs:0.12000,0.00000);
\draw[gray] (axis cs:-0.10733,0.44633)--(axis cs:0.10733,0.55367);
\draw[gray] (axis cs:-0.08485,0.91515)--(axis cs:0.08485,1.08485);
\draw[gray] (axis cs:-0.06656,1.40015)--(axis cs:0.06656,1.59985);
\draw[gray] (axis cs:-0.05367,1.89267)--(axis cs:0.05367,2.10733);
\draw[gray] (axis cs:0.41515,-1.91515)--(axis cs:0.58485,-2.08485);
\draw[gray] (axis cs:0.39267,-1.44633)--(axis cs:0.60733,-1.55367);
\draw[gray] (axis cs:0.38000,-1.00000)--(axis cs:0.62000,-1.00000);
\draw[gray] (axis cs:0.39267,-0.55367)--(axis cs:0.60733,-0.44633);
\draw[gray] (axis cs:0.41515,-0.08485)--(axis cs:0.58485,0.08485);
\draw[gray] (axis cs:0.43344,0.40015)--(axis cs:0.56656,0.59985);
\draw[gray] (axis cs:0.44633,0.89267)--(axis cs:0.55367,1.10733);
\draw[gray] (axis cs:0.45543,1.38858)--(axis cs:0.54457,1.61142);
\draw[gray] (axis cs:0.46205,1.88616)--(axis cs:0.53795,2.11384);
\draw[gray] (axis cs:0.88000,-2.00000)--(axis cs:1.12000,-2.00000);
\draw[gray] (axis cs:0.89267,-1.55367)--(axis cs:1.10733,-1.44633);
\draw[gray] (axis cs:0.91515,-1.08485)--(axis cs:1.08485,-0.91515);
\draw[gray] (axis cs:0.93344,-0.59985)--(axis cs:1.06656,-0.40015);
\draw[gray] (axis cs:0.94633,-0.10733)--(axis cs:1.05367,0.10733);
\draw[gray] (axis cs:0.95543,0.38858)--(axis cs:1.04457,0.61142);
\draw[gray] (axis cs:0.96205,0.88616)--(axis cs:1.03795,1.11384);
\draw[gray] (axis cs:0.96703,1.38462)--(axis cs:1.03297,1.61538);
\draw[gray] (axis cs:0.97090,1.88358)--(axis cs:1.02910,2.11642);
\draw[gray] (axis cs:1.41515,-2.08485)--(axis cs:1.58485,-1.91515);
\draw[gray] (axis cs:1.43344,-1.59985)--(axis cs:1.56656,-1.40015);
\draw[gray] (axis cs:1.44633,-1.10733)--(axis cs:1.55367,-0.89267);
\draw[gray] (axis cs:1.45543,-0.61142)--(axis cs:1.54457,-0.38858);
\draw[gray] (axis cs:1.46205,-0.11384)--(axis cs:1.53795,0.11384);
\draw[gray] (axis cs:1.46703,0.38462)--(axis cs:1.53297,0.61538);
\draw[gray] (axis cs:1.47090,0.88358)--(axis cs:1.52910,1.11642);
\draw[gray] (axis cs:1.47397,1.38286)--(axis cs:1.52603,1.61714);
\draw[gray] (axis cs:1.47647,1.88233)--(axis cs:1.52353,2.11767);
\draw[gray] (axis cs:1.94633,-2.10733)--(axis cs:2.05367,-1.89267);
\draw[gray] (axis cs:1.95543,-1.61142)--(axis cs:2.04457,-1.38858);
\draw[gray] (axis cs:1.96205,-1.11384)--(axis cs:2.03795,-0.88616);
\draw[gray] (axis cs:1.96703,-0.61538)--(axis cs:2.03297,-0.38462);
\draw[gray] (axis cs:1.97090,-0.11642)--(axis cs:2.02910,0.11642);
\draw[gray] (axis cs:1.97397,0.38286)--(axis cs:2.02603,0.61714);
\draw[gray] (axis cs:1.97647,0.88233)--(axis cs:2.02353,1.11767);
\draw[gray] (axis cs:1.97853,1.38194)--(axis cs:2.02147,1.61806);
\draw[gray] (axis cs:1.98027,1.88163)--(axis cs:2.01973,2.11837);
\addplot[main,domain=-2:2,samples=70] {-1.00*exp(x)-2*x-2};
\addplot[main,domain=-2:2,samples=70] {0.00*exp(x)-2*x-2};
\addplot[main,domain=-2:2,samples=70] {0.30*exp(x)-2*x-2};
\addplot[main,domain=-2:2,samples=70] {1.00*exp(x)-2*x-2};
\addplot[main,domain=-2:2,samples=70] {2.00*exp(x)-2*x-2};
\end{axis}\end{tikzpicture}\end{center}
\begin{center}\begin{tikzpicture}\begin{axis}[width=6.8cm,height=6.8cm,axis lines=middle,xmin=-2,xmax=2,ymin=-2,ymax=2,xlabel=$x$,ylabel=$y$,title={(c)},clip=true]
\addplot[black,dashed,domain=-2:2] {x--2};
\addplot[black,dashed,domain=-2:2] {x--1};
\addplot[black,dashed,domain=-2:2] {x-0};
\addplot[black,dashed,domain=-2:2] {x-1};
\addplot[black,dashed,domain=-2:2] {x-2};
\draw[gray] (axis cs:-2.12000,-2.00000)--(axis cs:-1.88000,-2.00000);
\draw[gray] (axis cs:-2.10733,-1.44633)--(axis cs:-1.89267,-1.55367);
\draw[gray] (axis cs:-2.08485,-0.91515)--(axis cs:-1.91515,-1.08485);
\draw[gray] (axis cs:-2.06656,-0.40015)--(axis cs:-1.93344,-0.59985);
\draw[gray] (axis cs:-2.05367,0.10733)--(axis cs:-1.94633,-0.10733);
\draw[gray] (axis cs:-2.04457,0.61142)--(axis cs:-1.95543,0.38858);
\draw[gray] (axis cs:-2.03795,1.11384)--(axis cs:-1.96205,0.88616);
\draw[gray] (axis cs:-2.03297,1.61538)--(axis cs:-1.96703,1.38462);
\draw[gray] (axis cs:-2.02910,2.11642)--(axis cs:-1.97090,1.88358);
\draw[gray] (axis cs:-1.60733,-2.05367)--(axis cs:-1.39267,-1.94633);
\draw[gray] (axis cs:-1.62000,-1.50000)--(axis cs:-1.38000,-1.50000);
\draw[gray] (axis cs:-1.60733,-0.94633)--(axis cs:-1.39267,-1.05367);
\draw[gray] (axis cs:-1.58485,-0.41515)--(axis cs:-1.41515,-0.58485);
\draw[gray] (axis cs:-1.56656,0.09985)--(axis cs:-1.43344,-0.09985);
\draw[gray] (axis cs:-1.55367,0.60733)--(axis cs:-1.44633,0.39267);
\draw[gray] (axis cs:-1.54457,1.11142)--(axis cs:-1.45543,0.88858);
\draw[gray] (axis cs:-1.53795,1.61384)--(axis cs:-1.46205,1.38616);
\draw[gray] (axis cs:-1.53297,2.11538)--(axis cs:-1.46703,1.88462);
\draw[gray] (axis cs:-1.08485,-2.08485)--(axis cs:-0.91515,-1.91515);
\draw[gray] (axis cs:-1.10733,-1.55367)--(axis cs:-0.89267,-1.44633);
\draw[gray] (axis cs:-1.12000,-1.00000)--(axis cs:-0.88000,-1.00000);
\draw[gray] (axis cs:-1.10733,-0.44633)--(axis cs:-0.89267,-0.55367);
\draw[gray] (axis cs:-1.08485,0.08485)--(axis cs:-0.91515,-0.08485);
\draw[gray] (axis cs:-1.06656,0.59985)--(axis cs:-0.93344,0.40015);
\draw[gray] (axis cs:-1.05367,1.10733)--(axis cs:-0.94633,0.89267);
\draw[gray] (axis cs:-1.04457,1.61142)--(axis cs:-0.95543,1.38858);
\draw[gray] (axis cs:-1.03795,2.11384)--(axis cs:-0.96205,1.88616);
\draw[gray] (axis cs:-0.56656,-2.09985)--(axis cs:-0.43344,-1.90015);
\draw[gray] (axis cs:-0.58485,-1.58485)--(axis cs:-0.41515,-1.41515);
\draw[gray] (axis cs:-0.60733,-1.05367)--(axis cs:-0.39267,-0.94633);
\draw[gray] (axis cs:-0.62000,-0.50000)--(axis cs:-0.38000,-0.50000);
\draw[gray] (axis cs:-0.60733,0.05367)--(axis cs:-0.39267,-0.05367);
\draw[gray] (axis cs:-0.58485,0.58485)--(axis cs:-0.41515,0.41515);
\draw[gray] (axis cs:-0.56656,1.09985)--(axis cs:-0.43344,0.90015);
\draw[gray] (axis cs:-0.55367,1.60733)--(axis cs:-0.44633,1.39267);
\draw[gray] (axis cs:-0.54457,2.11142)--(axis cs:-0.45543,1.88858);
\draw[gray] (axis cs:-0.05367,-2.10733)--(axis cs:0.05367,-1.89267);
\draw[gray] (axis cs:-0.06656,-1.59985)--(axis cs:0.06656,-1.40015);
\draw[gray] (axis cs:-0.08485,-1.08485)--(axis cs:0.08485,-0.91515);
\draw[gray] (axis cs:-0.10733,-0.55367)--(axis cs:0.10733,-0.44633);
\draw[gray] (axis cs:-0.12000,0.00000)--(axis cs:0.12000,0.00000);
\draw[gray] (axis cs:-0.10733,0.55367)--(axis cs:0.10733,0.44633);
\draw[gray] (axis cs:-0.08485,1.08485)--(axis cs:0.08485,0.91515);
\draw[gray] (axis cs:-0.06656,1.59985)--(axis cs:0.06656,1.40015);
\draw[gray] (axis cs:-0.05367,2.10733)--(axis cs:0.05367,1.89267);
\draw[gray] (axis cs:0.45543,-2.11142)--(axis cs:0.54457,-1.88858);
\draw[gray] (axis cs:0.44633,-1.60733)--(axis cs:0.55367,-1.39267);
\draw[gray] (axis cs:0.43344,-1.09985)--(axis cs:0.56656,-0.90015);
\draw[gray] (axis cs:0.41515,-0.58485)--(axis cs:0.58485,-0.41515);
\draw[gray] (axis cs:0.39267,-0.05367)--(axis cs:0.60733,0.05367);
\draw[gray] (axis cs:0.38000,0.50000)--(axis cs:0.62000,0.50000);
\draw[gray] (axis cs:0.39267,1.05367)--(axis cs:0.60733,0.94633);
\draw[gray] (axis cs:0.41515,1.58485)--(axis cs:0.58485,1.41515);
\draw[gray] (axis cs:0.43344,2.09985)--(axis cs:0.56656,1.90015);
\draw[gray] (axis cs:0.96205,-2.11384)--(axis cs:1.03795,-1.88616);
\draw[gray] (axis cs:0.95543,-1.61142)--(axis cs:1.04457,-1.38858);
\draw[gray] (axis cs:0.94633,-1.10733)--(axis cs:1.05367,-0.89267);
\draw[gray] (axis cs:0.93344,-0.59985)--(axis cs:1.06656,-0.40015);
\draw[gray] (axis cs:0.91515,-0.08485)--(axis cs:1.08485,0.08485);
\draw[gray] (axis cs:0.89267,0.44633)--(axis cs:1.10733,0.55367);
\draw[gray] (axis cs:0.88000,1.00000)--(axis cs:1.12000,1.00000);
\draw[gray] (axis cs:0.89267,1.55367)--(axis cs:1.10733,1.44633);
\draw[gray] (axis cs:0.91515,2.08485)--(axis cs:1.08485,1.91515);
\draw[gray] (axis cs:1.46703,-2.11538)--(axis cs:1.53297,-1.88462);
\draw[gray] (axis cs:1.46205,-1.61384)--(axis cs:1.53795,-1.38616);
\draw[gray] (axis cs:1.45543,-1.11142)--(axis cs:1.54457,-0.88858);
\draw[gray] (axis cs:1.44633,-0.60733)--(axis cs:1.55367,-0.39267);
\draw[gray] (axis cs:1.43344,-0.09985)--(axis cs:1.56656,0.09985);
\draw[gray] (axis cs:1.41515,0.41515)--(axis cs:1.58485,0.58485);
\draw[gray] (axis cs:1.39267,0.94633)--(axis cs:1.60733,1.05367);
\draw[gray] (axis cs:1.38000,1.50000)--(axis cs:1.62000,1.50000);
\draw[gray] (axis cs:1.39267,2.05367)--(axis cs:1.60733,1.94633);
\draw[gray] (axis cs:1.97090,-2.11642)--(axis cs:2.02910,-1.88358);
\draw[gray] (axis cs:1.96703,-1.61538)--(axis cs:2.03297,-1.38462);
\draw[gray] (axis cs:1.96205,-1.11384)--(axis cs:2.03795,-0.88616);
\draw[gray] (axis cs:1.95543,-0.61142)--(axis cs:2.04457,-0.38858);
\draw[gray] (axis cs:1.94633,-0.10733)--(axis cs:2.05367,0.10733);
\draw[gray] (axis cs:1.93344,0.40015)--(axis cs:2.06656,0.59985);
\draw[gray] (axis cs:1.91515,0.91515)--(axis cs:2.08485,1.08485);
\draw[gray] (axis cs:1.89267,1.44633)--(axis cs:2.10733,1.55367);
\draw[gray] (axis cs:1.88000,2.00000)--(axis cs:2.12000,2.00000);
\addplot[main,domain=-2:2,samples=70] {-1.00*exp(-x)+x-1};
\addplot[main,domain=-2:2,samples=70] {0.00*exp(-x)+x-1};
\addplot[main,domain=-2:2,samples=70] {0.30*exp(-x)+x-1};
\addplot[main,domain=-2:2,samples=70] {1.00*exp(-x)+x-1};
\addplot[main,domain=-2:2,samples=70] {2.00*exp(-x)+x-1};
\end{axis}\end{tikzpicture}\end{center}
\begin{center}\begin{tikzpicture}\begin{axis}[width=6.8cm,height=6.8cm,axis lines=middle,xmin=-2,xmax=2,ymin=-2,ymax=2,xlabel=$x$,ylabel=$y$,title={(d)},clip=true]
\addplot[black,dashed,domain=0:360,samples=80] ({0.5*cos(x)},{0.5*sin(x)});
\addplot[black,dashed,domain=0:360,samples=80] ({1*cos(x)},{1*sin(x)});
\addplot[black,dashed,domain=0:360,samples=80] ({1.4*cos(x)},{1.4*sin(x)});
\addplot[black,dashed,domain=0:360,samples=80] ({1.7*cos(x)},{1.7*sin(x)});
\addplot[black,dashed,domain=0:360,samples=80] ({2*cos(x)},{2*sin(x)});
\draw[gray] (axis cs:-2.01697,-2.11879)--(axis cs:-1.98303,-1.88121);
\draw[gray] (axis cs:-2.02245,-1.61788)--(axis cs:-1.97755,-1.38212);
\draw[gray] (axis cs:-2.02910,-1.11642)--(axis cs:-1.97090,-0.88358);
\draw[gray] (axis cs:-2.03529,-0.61469)--(axis cs:-1.96471,-0.38531);
\draw[gray] (axis cs:-2.03795,-0.11384)--(axis cs:-1.96205,0.11384);
\draw[gray] (axis cs:-2.03529,0.38531)--(axis cs:-1.96471,0.61469);
\draw[gray] (axis cs:-2.02910,0.88358)--(axis cs:-1.97090,1.11642);
\draw[gray] (axis cs:-2.02245,1.38212)--(axis cs:-1.97755,1.61788);
\draw[gray] (axis cs:-2.01697,1.88121)--(axis cs:-1.98303,2.11879);
\draw[gray] (axis cs:-1.52245,-2.11788)--(axis cs:-1.47755,-1.88212);
\draw[gray] (axis cs:-1.53297,-1.61538)--(axis cs:-1.46703,-1.38462);
\draw[gray] (axis cs:-1.54874,-1.10966)--(axis cs:-1.45126,-0.89034);
\draw[gray] (axis cs:-1.56656,-0.59985)--(axis cs:-1.43344,-0.40015);
\draw[gray] (axis cs:-1.57496,-0.09370)--(axis cs:-1.42504,0.09370);
\draw[gray] (axis cs:-1.56656,0.40015)--(axis cs:-1.43344,0.59985);
\draw[gray] (axis cs:-1.54874,0.89034)--(axis cs:-1.45126,1.10966);
\draw[gray] (axis cs:-1.53297,1.38462)--(axis cs:-1.46703,1.61538);
\draw[gray] (axis cs:-1.52245,1.88212)--(axis cs:-1.47755,2.11788);
\draw[gray] (axis cs:-1.02910,-2.11642)--(axis cs:-0.97090,-1.88358);
\draw[gray] (axis cs:-1.04874,-1.60966)--(axis cs:-0.95126,-1.39034);
\draw[gray] (axis cs:-1.08485,-1.08485)--(axis cs:-0.91515,-0.91515);
\draw[gray] (axis cs:-1.11642,-0.52910)--(axis cs:-0.88358,-0.47090);
\draw[gray] (axis cs:-1.12000,0.00000)--(axis cs:-0.88000,0.00000);
\draw[gray] (axis cs:-1.11642,0.47090)--(axis cs:-0.88358,0.52910);
\draw[gray] (axis cs:-1.08485,0.91515)--(axis cs:-0.91515,1.08485);
\draw[gray] (axis cs:-1.04874,1.39034)--(axis cs:-0.95126,1.60966);
\draw[gray] (axis cs:-1.02910,1.88358)--(axis cs:-0.97090,2.11642);
\draw[gray] (axis cs:-0.53529,-2.11469)--(axis cs:-0.46471,-1.88531);
\draw[gray] (axis cs:-0.56656,-1.59985)--(axis cs:-0.43344,-1.40015);
\draw[gray] (axis cs:-0.61642,-1.02910)--(axis cs:-0.38358,-0.97090);
\draw[gray] (axis cs:-0.60733,-0.44633)--(axis cs:-0.39267,-0.55367);
\draw[gray] (axis cs:-0.59600,0.07200)--(axis cs:-0.40400,-0.07200);
\draw[gray] (axis cs:-0.60733,0.55367)--(axis cs:-0.39267,0.44633);
\draw[gray] (axis cs:-0.61642,0.97090)--(axis cs:-0.38358,1.02910);
\draw[gray] (axis cs:-0.56656,1.40015)--(axis cs:-0.43344,1.59985);
\draw[gray] (axis cs:-0.53529,1.88531)--(axis cs:-0.46471,2.11469);
\draw[gray] (axis cs:-0.03795,-2.11384)--(axis cs:0.03795,-1.88616);
\draw[gray] (axis cs:-0.07496,-1.59370)--(axis cs:0.07496,-1.40630);
\draw[gray] (axis cs:-0.12000,-1.00000)--(axis cs:0.12000,-1.00000);
\draw[gray] (axis cs:-0.09600,-0.42800)--(axis cs:0.09600,-0.57200);
\draw[gray] (axis cs:-0.08485,0.08485)--(axis cs:0.08485,-0.08485);
\draw[gray] (axis cs:-0.09600,0.57200)--(axis cs:0.09600,0.42800);
\draw[gray] (axis cs:-0.12000,1.00000)--(axis cs:0.12000,1.00000);
\draw[gray] (axis cs:-0.07496,1.40630)--(axis cs:0.07496,1.59370);
\draw[gray] (axis cs:-0.03795,1.88616)--(axis cs:0.03795,2.11384);
\draw[gray] (axis cs:0.46471,-2.11469)--(axis cs:0.53529,-1.88531);
\draw[gray] (axis cs:0.43344,-1.59985)--(axis cs:0.56656,-1.40015);
\draw[gray] (axis cs:0.38358,-1.02910)--(axis cs:0.61642,-0.97090);
\draw[gray] (axis cs:0.39267,-0.44633)--(axis cs:0.60733,-0.55367);
\draw[gray] (axis cs:0.40400,0.07200)--(axis cs:0.59600,-0.07200);
\draw[gray] (axis cs:0.39267,0.55367)--(axis cs:0.60733,0.44633);
\draw[gray] (axis cs:0.38358,0.97090)--(axis cs:0.61642,1.02910);
\draw[gray] (axis cs:0.43344,1.40015)--(axis cs:0.56656,1.59985);
\draw[gray] (axis cs:0.46471,1.88531)--(axis cs:0.53529,2.11469);
\draw[gray] (axis cs:0.97090,-2.11642)--(axis cs:1.02910,-1.88358);
\draw[gray] (axis cs:0.95126,-1.60966)--(axis cs:1.04874,-1.39034);
\draw[gray] (axis cs:0.91515,-1.08485)--(axis cs:1.08485,-0.91515);
\draw[gray] (axis cs:0.88358,-0.52910)--(axis cs:1.11642,-0.47090);
\draw[gray] (axis cs:0.88000,0.00000)--(axis cs:1.12000,0.00000);
\draw[gray] (axis cs:0.88358,0.47090)--(axis cs:1.11642,0.52910);
\draw[gray] (axis cs:0.91515,0.91515)--(axis cs:1.08485,1.08485);
\draw[gray] (axis cs:0.95126,1.39034)--(axis cs:1.04874,1.60966);
\draw[gray] (axis cs:0.97090,1.88358)--(axis cs:1.02910,2.11642);
\draw[gray] (axis cs:1.47755,-2.11788)--(axis cs:1.52245,-1.88212);
\draw[gray] (axis cs:1.46703,-1.61538)--(axis cs:1.53297,-1.38462);
\draw[gray] (axis cs:1.45126,-1.10966)--(axis cs:1.54874,-0.89034);
\draw[gray] (axis cs:1.43344,-0.59985)--(axis cs:1.56656,-0.40015);
\draw[gray] (axis cs:1.42504,-0.09370)--(axis cs:1.57496,0.09370);
\draw[gray] (axis cs:1.43344,0.40015)--(axis cs:1.56656,0.59985);
\draw[gray] (axis cs:1.45126,0.89034)--(axis cs:1.54874,1.10966);
\draw[gray] (axis cs:1.46703,1.38462)--(axis cs:1.53297,1.61538);
\draw[gray] (axis cs:1.47755,1.88212)--(axis cs:1.52245,2.11788);
\draw[gray] (axis cs:1.98303,-2.11879)--(axis cs:2.01697,-1.88121);
\draw[gray] (axis cs:1.97755,-1.61788)--(axis cs:2.02245,-1.38212);
\draw[gray] (axis cs:1.97090,-1.11642)--(axis cs:2.02910,-0.88358);
\draw[gray] (axis cs:1.96471,-0.61469)--(axis cs:2.03529,-0.38531);
\draw[gray] (axis cs:1.96205,-0.11384)--(axis cs:2.03795,0.11384);
\draw[gray] (axis cs:1.96471,0.38531)--(axis cs:2.03529,0.61469);
\draw[gray] (axis cs:1.97090,0.88358)--(axis cs:2.02910,1.11642);
\draw[gray] (axis cs:1.97755,1.38212)--(axis cs:2.02245,1.61788);
\draw[gray] (axis cs:1.98303,1.88121)--(axis cs:2.01697,2.11879);
\addplot[main] coordinates {(0.0000,-1.5000) (-0.0083,-1.5105) (-0.0383,-1.5508) (-0.0683,-1.5951) (-0.0983,-1.6440) (-0.1283,-1.6981) (-0.1583,-1.7583) (-0.1883,-1.8254) (-0.2183,-1.9007) (-0.2483,-1.9856) (-0.2783,-2.0816)};
\addplot[main] coordinates {(0.0000,-1.5000) (0.0084,-1.4897) (0.0384,-1.4547) (0.0684,-1.4225) (0.0984,-1.3929) (0.1284,-1.3654) (0.1584,-1.3399) (0.1884,-1.3161) (0.2184,-1.2938) (0.2484,-1.2728) (0.2784,-1.2529) (0.3084,-1.2339) (0.3384,-1.2158) (0.3684,-1.1983) (0.3984,-1.1814) (0.4284,-1.1650) (0.4584,-1.1489) (0.4884,-1.1332) (0.5184,-1.1176) (0.5484,-1.1021) (0.5784,-1.0866) (0.6084,-1.0711) (0.6384,-1.0556) (0.6684,-1.0398) (0.6984,-1.0239) (0.7284,-1.0076) (0.7584,-0.9911) (0.7884,-0.9742) (0.8184,-0.9569) (0.8484,-0.9391) (0.8784,-0.9207) (0.9084,-0.9019) (0.9384,-0.8824) (0.9684,-0.8623) (0.9984,-0.8415) (1.0284,-0.8200) (1.0584,-0.7977) (1.0884,-0.7746) (1.1184,-0.7506) (1.1484,-0.7257) (1.1784,-0.6999) (1.2084,-0.6730) (1.2384,-0.6451) (1.2684,-0.6160) (1.2984,-0.5858) (1.3284,-0.5543) (1.3584,-0.5214) (1.3884,-0.4872) (1.4184,-0.4515) (1.4484,-0.4143) (1.4784,-0.3753) (1.5084,-0.3346) (1.5384,-0.2921) (1.5684,-0.2475) (1.5984,-0.2008) (1.6284,-0.1517) (1.6584,-0.1002) (1.6884,-0.0460) (1.7184,0.0110) (1.7484,0.0712) (1.7784,0.1348) (1.8084,0.2022) (1.8384,0.2736) (1.8684,0.3496) (1.8984,0.4306) (1.9284,0.5171) (1.9584,0.6100) (1.9884,0.7099) (2.0000,0.7507)};
\addplot[main] coordinates {(0.0000,-0.5000) (-0.0093,-0.4930) (-0.0393,-0.4700) (-0.0693,-0.4464) (-0.0993,-0.4223) (-0.1293,-0.3977) (-0.1593,-0.3728) (-0.1893,-0.3476) (-0.2193,-0.3222) (-0.2493,-0.2967) (-0.2793,-0.2713) (-0.3093,-0.2459) (-0.3393,-0.2207) (-0.3693,-0.1957) (-0.3993,-0.1712) (-0.4293,-0.1471) (-0.4593,-0.1236) (-0.4893,-0.1007) (-0.5193,-0.0786) (-0.5493,-0.0573) (-0.5793,-0.0369) (-0.6093,-0.0175) (-0.6393,0.0008) (-0.6693,0.0180) (-0.6993,0.0339) (-0.7293,0.0485) (-0.7593,0.0618) (-0.7893,0.0737) (-0.8193,0.0841) (-0.8493,0.0930) (-0.8793,0.1003) (-0.9093,0.1060) (-0.9393,0.1100) (-0.9693,0.1123) (-0.9993,0.1129) (-1.0293,0.1116) (-1.0593,0.1085) (-1.0893,0.1036) (-1.1193,0.0967) (-1.1493,0.0878) (-1.1793,0.0770) (-1.2093,0.0640) (-1.2393,0.0490) (-1.2693,0.0317) (-1.2993,0.0122) (-1.3293,-0.0096) (-1.3593,-0.0338) (-1.3893,-0.0606) (-1.4193,-0.0899) (-1.4493,-0.1220) (-1.4793,-0.1569) (-1.5093,-0.1948) (-1.5393,-0.2359) (-1.5693,-0.2803) (-1.5993,-0.3284) (-1.6293,-0.3804) (-1.6593,-0.4365) (-1.6893,-0.4971) (-1.7193,-0.5627) (-1.7493,-0.6336) (-1.7793,-0.7106) (-1.8093,-0.7941) (-1.8393,-0.8851) (-1.8693,-0.9844) (-1.8993,-1.0933) (-1.9293,-1.2131) (-1.9593,-1.3456) (-1.9893,-1.4929) (-2.0000,-1.5497)};
\addplot[main] coordinates {(0.0000,-0.5000) (0.0092,-0.5069) (0.0392,-0.5288) (0.0692,-0.5500) (0.0992,-0.5704) (0.1292,-0.5899) (0.1592,-0.6085) (0.1892,-0.6261) (0.2192,-0.6428) (0.2492,-0.6584) (0.2792,-0.6730) (0.3092,-0.6866) (0.3392,-0.6990) (0.3692,-0.7103) (0.3992,-0.7206) (0.4292,-0.7296) (0.4592,-0.7376) (0.4892,-0.7443) (0.5192,-0.7500) (0.5492,-0.7544) (0.5792,-0.7577) (0.6092,-0.7598) (0.6392,-0.7608) (0.6692,-0.7606) (0.6992,-0.7592) (0.7292,-0.7567) (0.7592,-0.7530) (0.7892,-0.7481) (0.8192,-0.7420) (0.8492,-0.7348) (0.8792,-0.7263) (0.9092,-0.7167) (0.9392,-0.7059) (0.9692,-0.6939) (0.9992,-0.6807) (1.0292,-0.6662) (1.0592,-0.6505) (1.0892,-0.6335) (1.1192,-0.6152) (1.1492,-0.5956) (1.1792,-0.5747) (1.2092,-0.5523) (1.2392,-0.5286) (1.2692,-0.5034) (1.2992,-0.4767) (1.3292,-0.4485) (1.3592,-0.4186) (1.3892,-0.3871) (1.4192,-0.3538) (1.4492,-0.3187) (1.4792,-0.2817) (1.5092,-0.2426) (1.5392,-0.2014) (1.5692,-0.1580) (1.5992,-0.1121) (1.6292,-0.0637) (1.6592,-0.0126) (1.6892,0.0415) (1.7192,0.0988) (1.7492,0.1596) (1.7792,0.2241) (1.8092,0.2927) (1.8392,0.3657) (1.8692,0.4438) (1.8992,0.5274) (1.9292,0.6172) (1.9592,0.7139) (1.9892,0.8184) (2.0000,0.8580)};
\addplot[main] coordinates {(0.0000,0.0000) (-0.0001,0.0001) (-0.0011,0.0011) (-0.0111,0.0111) (-0.0411,0.0411) (-0.0711,0.0709) (-0.1011,0.1004) (-0.1311,0.1296) (-0.1611,0.1583) (-0.1911,0.1865) (-0.2211,0.2140) (-0.2511,0.2408) (-0.2811,0.2667) (-0.3111,0.2918) (-0.3411,0.3158) (-0.3711,0.3388) (-0.4011,0.3606) (-0.4311,0.3813) (-0.4611,0.4007) (-0.4911,0.4189) (-0.5211,0.4357) (-0.5511,0.4512) (-0.5811,0.4653) (-0.6111,0.4779) (-0.6411,0.4892) (-0.6711,0.4989) (-0.7011,0.5072) (-0.7311,0.5140) (-0.7611,0.5193) (-0.7911,0.5231) (-0.8211,0.5253) (-0.8511,0.5260) (-0.8811,0.5252) (-0.9111,0.5229) (-0.9411,0.5190) (-0.9711,0.5136) (-1.0011,0.5066) (-1.0311,0.4981) (-1.0611,0.4879) (-1.0911,0.4762) (-1.1211,0.4629) (-1.1511,0.4479) (-1.1811,0.4313) (-1.2111,0.4131) (-1.2411,0.3931) (-1.2711,0.3714) (-1.3011,0.3479) (-1.3311,0.3225) (-1.3611,0.2953) (-1.3911,0.2661) (-1.4211,0.2349) (-1.4511,0.2016) (-1.4811,0.1661) (-1.5111,0.1283) (-1.5411,0.0881) (-1.5711,0.0453) (-1.6011,-0.0002) (-1.6311,-0.0486) (-1.6611,-0.1001) (-1.6911,-0.1548) (-1.7211,-0.2132) (-1.7511,-0.2754) (-1.7811,-0.3418) (-1.8111,-0.4129) (-1.8411,-0.4890) (-1.8711,-0.5708) (-1.9011,-0.6589) (-1.9311,-0.7540) (-1.9611,-0.8571) (-1.9911,-0.9692) (-2.0000,-1.0044)};
\addplot[main] coordinates {(0.0000,0.0000) (0.0001,-0.0001) (0.0011,-0.0011) (0.0111,-0.0111) (0.0411,-0.0411) (0.0711,-0.0709) (0.1011,-0.1004) (0.1311,-0.1296) (0.1611,-0.1583) (0.1911,-0.1865) (0.2211,-0.2140) (0.2511,-0.2408) (0.2811,-0.2667) (0.3111,-0.2918) (0.3411,-0.3158) (0.3711,-0.3388) (0.4011,-0.3606) (0.4311,-0.3813) (0.4611,-0.4007) (0.4911,-0.4189) (0.5211,-0.4357) (0.5511,-0.4512) (0.5811,-0.4653) (0.6111,-0.4779) (0.6411,-0.4892) (0.6711,-0.4989) (0.7011,-0.5072) (0.7311,-0.5140) (0.7611,-0.5193) (0.7911,-0.5231) (0.8211,-0.5253) (0.8511,-0.5260) (0.8811,-0.5252) (0.9111,-0.5229) (0.9411,-0.5190) (0.9711,-0.5136) (1.0011,-0.5066) (1.0311,-0.4981) (1.0611,-0.4879) (1.0911,-0.4762) (1.1211,-0.4629) (1.1511,-0.4479) (1.1811,-0.4313) (1.2111,-0.4131) (1.2411,-0.3931) (1.2711,-0.3714) (1.3011,-0.3479) (1.3311,-0.3225) (1.3611,-0.2953) (1.3911,-0.2661) (1.4211,-0.2349) (1.4511,-0.2016) (1.4811,-0.1661) (1.5111,-0.1283) (1.5411,-0.0881) (1.5711,-0.0453) (1.6011,0.0002) (1.6311,0.0486) (1.6611,0.1001) (1.6911,0.1548) (1.7211,0.2132) (1.7511,0.2754) (1.7811,0.3418) (1.8111,0.4129) (1.8411,0.4890) (1.8711,0.5708) (1.9011,0.6589) (1.9311,0.7540) (1.9611,0.8571) (1.9911,0.9692) (2.0000,1.0044)};
\addplot[main] coordinates {(0.0000,0.5000) (-0.0092,0.5069) (-0.0392,0.5288) (-0.0692,0.5500) (-0.0992,0.5704) (-0.1292,0.5899) (-0.1592,0.6085) (-0.1892,0.6261) (-0.2192,0.6428) (-0.2492,0.6584) (-0.2792,0.6730) (-0.3092,0.6866) (-0.3392,0.6990) (-0.3692,0.7103) (-0.3992,0.7206) (-0.4292,0.7296) (-0.4592,0.7376) (-0.4892,0.7443) (-0.5192,0.7500) (-0.5492,0.7544) (-0.5792,0.7577) (-0.6092,0.7598) (-0.6392,0.7608) (-0.6692,0.7606) (-0.6992,0.7592) (-0.7292,0.7567) (-0.7592,0.7530) (-0.7892,0.7481) (-0.8192,0.7420) (-0.8492,0.7348) (-0.8792,0.7263) (-0.9092,0.7167) (-0.9392,0.7059) (-0.9692,0.6939) (-0.9992,0.6807) (-1.0292,0.6662) (-1.0592,0.6505) (-1.0892,0.6335) (-1.1192,0.6152) (-1.1492,0.5956) (-1.1792,0.5747) (-1.2092,0.5523) (-1.2392,0.5286) (-1.2692,0.5034) (-1.2992,0.4767) (-1.3292,0.4485) (-1.3592,0.4186) (-1.3892,0.3871) (-1.4192,0.3538) (-1.4492,0.3187) (-1.4792,0.2817) (-1.5092,0.2426) (-1.5392,0.2014) (-1.5692,0.1580) (-1.5992,0.1121) (-1.6292,0.0637) (-1.6592,0.0126) (-1.6892,-0.0415) (-1.7192,-0.0988) (-1.7492,-0.1596) (-1.7792,-0.2241) (-1.8092,-0.2927) (-1.8392,-0.3657) (-1.8692,-0.4438) (-1.8992,-0.5274) (-1.9292,-0.6172) (-1.9592,-0.7139) (-1.9892,-0.8184) (-2.0000,-0.8580)};
\addplot[main] coordinates {(0.0000,0.5000) (0.0093,0.4930) (0.0393,0.4700) (0.0693,0.4464) (0.0993,0.4223) (0.1293,0.3977) (0.1593,0.3728) (0.1893,0.3476) (0.2193,0.3222) (0.2493,0.2967) (0.2793,0.2713) (0.3093,0.2459) (0.3393,0.2207) (0.3693,0.1957) (0.3993,0.1712) (0.4293,0.1471) (0.4593,0.1236) (0.4893,0.1007) (0.5193,0.0786) (0.5493,0.0573) (0.5793,0.0369) (0.6093,0.0175) (0.6393,-0.0008) (0.6693,-0.0180) (0.6993,-0.0339) (0.7293,-0.0485) (0.7593,-0.0618) (0.7893,-0.0737) (0.8193,-0.0841) (0.8493,-0.0930) (0.8793,-0.1003) (0.9093,-0.1060) (0.9393,-0.1100) (0.9693,-0.1123) (0.9993,-0.1129) (1.0293,-0.1116) (1.0593,-0.1085) (1.0893,-0.1036) (1.1193,-0.0967) (1.1493,-0.0878) (1.1793,-0.0770) (1.2093,-0.0640) (1.2393,-0.0490) (1.2693,-0.0317) (1.2993,-0.0122) (1.3293,0.0096) (1.3593,0.0338) (1.3893,0.0606) (1.4193,0.0899) (1.4493,0.1220) (1.4793,0.1569) (1.5093,0.1948) (1.5393,0.2359) (1.5693,0.2803) (1.5993,0.3284) (1.6293,0.3804) (1.6593,0.4365) (1.6893,0.4971) (1.7193,0.5627) (1.7493,0.6336) (1.7793,0.7106) (1.8093,0.7941) (1.8393,0.8851) (1.8693,0.9844) (1.8993,1.0933) (1.9293,1.2131) (1.9593,1.3456) (1.9893,1.4929) (2.0000,1.5497)};
\addplot[main] coordinates {(0.0000,1.5000) (-0.0084,1.4897) (-0.0384,1.4547) (-0.0684,1.4225) (-0.0984,1.3929) (-0.1284,1.3654) (-0.1584,1.3399) (-0.1884,1.3161) (-0.2184,1.2938) (-0.2484,1.2728) (-0.2784,1.2529) (-0.3084,1.2339) (-0.3384,1.2158) (-0.3684,1.1983) (-0.3984,1.1814) (-0.4284,1.1650) (-0.4584,1.1489) (-0.4884,1.1332) (-0.5184,1.1176) (-0.5484,1.1021) (-0.5784,1.0866) (-0.6084,1.0711) (-0.6384,1.0556) (-0.6684,1.0398) (-0.6984,1.0239) (-0.7284,1.0076) (-0.7584,0.9911) (-0.7884,0.9742) (-0.8184,0.9569) (-0.8484,0.9391) (-0.8784,0.9207) (-0.9084,0.9019) (-0.9384,0.8824) (-0.9684,0.8623) (-0.9984,0.8415) (-1.0284,0.8200) (-1.0584,0.7977) (-1.0884,0.7746) (-1.1184,0.7506) (-1.1484,0.7257) (-1.1784,0.6999) (-1.2084,0.6730) (-1.2384,0.6451) (-1.2684,0.6160) (-1.2984,0.5858) (-1.3284,0.5543) (-1.3584,0.5214) (-1.3884,0.4872) (-1.4184,0.4515) (-1.4484,0.4143) (-1.4784,0.3753) (-1.5084,0.3346) (-1.5384,0.2921) (-1.5684,0.2475) (-1.5984,0.2008) (-1.6284,0.1517) (-1.6584,0.1002) (-1.6884,0.0460) (-1.7184,-0.0110) (-1.7484,-0.0712) (-1.7784,-0.1348) (-1.8084,-0.2022) (-1.8384,-0.2736) (-1.8684,-0.3496) (-1.8984,-0.4306) (-1.9284,-0.5171) (-1.9584,-0.6100) (-1.9884,-0.7099) (-2.0000,-0.7507)};
\addplot[main] coordinates {(0.0000,1.5000) (0.0083,1.5105) (0.0383,1.5508) (0.0683,1.5951) (0.0983,1.6440) (0.1283,1.6981) (0.1583,1.7583) (0.1883,1.8254) (0.2183,1.9007) (0.2483,1.9856) (0.2783,2.0816)};
\end{axis}\end{tikzpicture}\end{center}
\begin{center}\begin{tikzpicture}\begin{axis}[width=6.8cm,height=6.8cm,axis lines=middle,xmin=-3,xmax=3,ymin=-3,ymax=3,xlabel=$x$,ylabel=$y$,title={(e)},clip=true]
\addplot[black,dashed,domain=-3:3] {-x+1/-2};
\addplot[black,dashed,domain=-3:3] {-x+1/-1};
\addplot[black,dashed,domain=-3:3] {-x+1/0.5};
\addplot[black,dashed,domain=-3:3] {-x+1/1};
\addplot[black,dashed,domain=-3:3] {-x+1/2};
\draw[gray] (axis cs:-3.11837,-2.98027)--(axis cs:-2.88163,-3.01973);
\draw[gray] (axis cs:-3.11788,-2.22755)--(axis cs:-2.88212,-2.27245);
\draw[gray] (axis cs:-3.11714,-1.47397)--(axis cs:-2.88286,-1.52603);
\draw[gray] (axis cs:-3.11595,-0.71908)--(axis cs:-2.88405,-0.78092);
\draw[gray] (axis cs:-3.11384,0.03795)--(axis cs:-2.88616,-0.03795);
\draw[gray] (axis cs:-3.10966,0.79874)--(axis cs:-2.89034,0.70126);
\draw[gray] (axis cs:-3.09985,1.56656)--(axis cs:-2.90015,1.43344);
\draw[gray] (axis cs:-3.07200,2.34600)--(axis cs:-2.92800,2.15400);
\draw[gray] (axis cs:-2.36788,-2.97755)--(axis cs:-2.13212,-3.02245);
\draw[gray] (axis cs:-2.36714,-2.22397)--(axis cs:-2.13286,-2.27603);
\draw[gray] (axis cs:-2.36595,-1.46908)--(axis cs:-2.13405,-1.53092);
\draw[gray] (axis cs:-2.36384,-0.71205)--(axis cs:-2.13616,-0.78795);
\draw[gray] (axis cs:-2.35966,0.04874)--(axis cs:-2.14034,-0.04874);
\draw[gray] (axis cs:-2.34985,0.81656)--(axis cs:-2.15015,0.68344);
\draw[gray] (axis cs:-2.32200,1.59600)--(axis cs:-2.17800,1.40400);
\draw[gray] (axis cs:-2.32200,2.90400)--(axis cs:-2.17800,3.09600);
\draw[gray] (axis cs:-1.61714,-2.97397)--(axis cs:-1.38286,-3.02603);
\draw[gray] (axis cs:-1.61595,-2.21908)--(axis cs:-1.38405,-2.28092);
\draw[gray] (axis cs:-1.61384,-1.46205)--(axis cs:-1.38616,-1.53795);
\draw[gray] (axis cs:-1.60966,-0.70126)--(axis cs:-1.39034,-0.79874);
\draw[gray] (axis cs:-1.59985,0.06656)--(axis cs:-1.40015,-0.06656);
\draw[gray] (axis cs:-1.57200,0.84600)--(axis cs:-1.42800,0.65400);
\draw[gray] (axis cs:-1.57200,2.15400)--(axis cs:-1.42800,2.34600);
\draw[gray] (axis cs:-1.59985,2.93344)--(axis cs:-1.40015,3.06656);
\draw[gray] (axis cs:-0.86595,-2.96908)--(axis cs:-0.63405,-3.03092);
\draw[gray] (axis cs:-0.86384,-2.21205)--(axis cs:-0.63616,-2.28795);
\draw[gray] (axis cs:-0.85966,-1.45126)--(axis cs:-0.64034,-1.54874);
\draw[gray] (axis cs:-0.84985,-0.68344)--(axis cs:-0.65015,-0.81656);
\draw[gray] (axis cs:-0.82200,0.09600)--(axis cs:-0.67800,-0.09600);
\draw[gray] (axis cs:-0.82200,1.40400)--(axis cs:-0.67800,1.59600);
\draw[gray] (axis cs:-0.84985,2.18344)--(axis cs:-0.65015,2.31656);
\draw[gray] (axis cs:-0.85966,2.95126)--(axis cs:-0.64034,3.04874);
\draw[gray] (axis cs:-0.11384,-2.96205)--(axis cs:0.11384,-3.03795);
\draw[gray] (axis cs:-0.10966,-2.20126)--(axis cs:0.10966,-2.29874);
\draw[gray] (axis cs:-0.09985,-1.43344)--(axis cs:0.09985,-1.56656);
\draw[gray] (axis cs:-0.07200,-0.65400)--(axis cs:0.07200,-0.84600);
\draw[gray] (axis cs:-0.07200,0.65400)--(axis cs:0.07200,0.84600);
\draw[gray] (axis cs:-0.09985,1.43344)--(axis cs:0.09985,1.56656);
\draw[gray] (axis cs:-0.10966,2.20126)--(axis cs:0.10966,2.29874);
\draw[gray] (axis cs:-0.11384,2.96205)--(axis cs:0.11384,3.03795);
\draw[gray] (axis cs:0.64034,-2.95126)--(axis cs:0.85966,-3.04874);
\draw[gray] (axis cs:0.65015,-2.18344)--(axis cs:0.84985,-2.31656);
\draw[gray] (axis cs:0.67800,-1.40400)--(axis cs:0.82200,-1.59600);
\draw[gray] (axis cs:0.67800,-0.09600)--(axis cs:0.82200,0.09600);
\draw[gray] (axis cs:0.65015,0.68344)--(axis cs:0.84985,0.81656);
\draw[gray] (axis cs:0.64034,1.45126)--(axis cs:0.85966,1.54874);
\draw[gray] (axis cs:0.63616,2.21205)--(axis cs:0.86384,2.28795);
\draw[gray] (axis cs:0.63405,2.96908)--(axis cs:0.86595,3.03092);
\draw[gray] (axis cs:1.40015,-2.93344)--(axis cs:1.59985,-3.06656);
\draw[gray] (axis cs:1.42800,-2.15400)--(axis cs:1.57200,-2.34600);
\draw[gray] (axis cs:1.42800,-0.84600)--(axis cs:1.57200,-0.65400);
\draw[gray] (axis cs:1.40015,-0.06656)--(axis cs:1.59985,0.06656);
\draw[gray] (axis cs:1.39034,0.70126)--(axis cs:1.60966,0.79874);
\draw[gray] (axis cs:1.38616,1.46205)--(axis cs:1.61384,1.53795);
\draw[gray] (axis cs:1.38405,2.21908)--(axis cs:1.61595,2.28092);
\draw[gray] (axis cs:1.38286,2.97397)--(axis cs:1.61714,3.02603);
\draw[gray] (axis cs:2.17800,-2.90400)--(axis cs:2.32200,-3.09600);
\draw[gray] (axis cs:2.17800,-1.59600)--(axis cs:2.32200,-1.40400);
\draw[gray] (axis cs:2.15015,-0.81656)--(axis cs:2.34985,-0.68344);
\draw[gray] (axis cs:2.14034,-0.04874)--(axis cs:2.35966,0.04874);
\draw[gray] (axis cs:2.13616,0.71205)--(axis cs:2.36384,0.78795);
\draw[gray] (axis cs:2.13405,1.46908)--(axis cs:2.36595,1.53092);
\draw[gray] (axis cs:2.13286,2.22397)--(axis cs:2.36714,2.27603);
\draw[gray] (axis cs:2.13212,2.97755)--(axis cs:2.36788,3.02245);
\draw[gray] (axis cs:2.92800,-2.34600)--(axis cs:3.07200,-2.15400);
\draw[gray] (axis cs:2.90015,-1.56656)--(axis cs:3.09985,-1.43344);
\draw[gray] (axis cs:2.89034,-0.79874)--(axis cs:3.10966,-0.70126);
\draw[gray] (axis cs:2.88616,-0.03795)--(axis cs:3.11384,0.03795);
\draw[gray] (axis cs:2.88405,0.71908)--(axis cs:3.11595,0.78092);
\draw[gray] (axis cs:2.88286,1.47397)--(axis cs:3.11714,1.52603);
\draw[gray] (axis cs:2.88212,2.22755)--(axis cs:3.11788,2.27245);
\draw[gray] (axis cs:2.88163,2.98027)--(axis cs:3.11837,3.01973);
\addplot[main] coordinates {(-3.0993,0.4812) (-3.0667,0.4687) (-3.0340,0.4561) (-3.0010,0.4433) (-2.9679,0.4303) (-2.9347,0.4171) (-2.9013,0.4038) (-2.8676,0.3903) (-2.8338,0.3766) (-2.7998,0.3627) (-2.7657,0.3486) (-2.7313,0.3343) (-2.6967,0.3198) (-2.6619,0.3051) (-2.6268,0.2902) (-2.5916,0.2750) (-2.5561,0.2596) (-2.5204,0.2440) (-2.4844,0.2281) (-2.4482,0.2120) (-2.4117,0.1956) (-2.3749,0.1789) (-2.3379,0.1620) (-2.3005,0.1448) (-2.2629,0.1272) (-2.2249,0.1094) (-2.1866,0.0912) (-2.1480,0.0726) (-2.1090,0.0538) (-2.0697,0.0345) (-2.0300,0.0149) (-1.9899,-0.0051) (-1.9494,-0.0255) (-1.9084,-0.0463) (-1.8671,-0.0676) (-1.8252,-0.0893) (-1.7829,-0.1116) (-1.7401,-0.1343) (-1.6967,-0.1575) (-1.6528,-0.1814) (-1.6083,-0.2057) (-1.5632,-0.2307) (-1.5174,-0.2564) (-1.4710,-0.2827) (-1.4239,-0.3097) (-1.3760,-0.3375) (-1.3273,-0.3661) (-1.2778,-0.3955) (-1.2274,-0.4258) (-1.1761,-0.4571) (-1.1237,-0.4893) (-1.0702,-0.5227) (-1.0157,-0.5572) (-0.9598,-0.5929) (-0.9027,-0.6299) (-0.8441,-0.6684) (-0.7840,-0.7084) (-0.7222,-0.7501) (-0.6586,-0.7936) (-0.5931,-0.8390) (-0.5253,-0.8867) (-0.4552,-0.9367) (-0.3825,-0.9894) (-0.3068,-1.0449) (-0.2278,-1.1038) (-0.1452,-1.1663) (-0.0584,-1.2330) (0.0332,-1.3045) (0.1302,-1.3815) (0.2337,-1.4649) (0.3448,-1.5558) (0.4650,-1.6560) (0.5964,-1.7672) (0.7417,-1.8924) (0.9050,-2.0356) (1.0922,-2.2027) (1.3131,-2.4035) (1.5846,-2.6549) (1.9415,-2.9917)};
\addplot[main] coordinates {(2.0465,-2.9965) (1.8742,-2.8143) (1.7301,-2.6601) (1.6065,-2.5266) (1.4988,-2.4088) (1.4034,-2.3035) (1.3181,-2.2082) (1.2411,-2.1212) (1.1710,-2.0412) (1.1069,-1.9671) (1.0479,-1.8981) (0.9934,-1.8335) (0.9427,-1.7729) (0.8956,-1.7158) (0.8515,-1.6617) (0.8102,-1.6104) (0.7714,-1.5616) (0.7349,-1.5151) (0.7004,-1.4706) (0.6679,-1.4281) (0.6370,-1.3873) (0.6078,-1.3480) (0.5800,-1.3103) (0.5537,-1.2739) (0.5286,-1.2389) (0.5047,-1.2049) (0.4819,-1.1722) (0.4601,-1.1404) (0.4393,-1.1096) (0.4195,-1.0798) (0.4005,-1.0508) (0.3823,-1.0226) (0.3649,-0.9952) (0.3482,-0.9686) (0.3322,-0.9426) (0.3169,-0.9173) (0.3022,-0.8926) (0.2881,-0.8685) (0.2745,-0.8450) (0.2615,-0.8220) (0.2491,-0.7995) (0.2371,-0.7775) (0.2255,-0.7560) (0.2145,-0.7350) (0.2039,-0.7143) (0.1936,-0.6941) (0.1838,-0.6743) (0.1744,-0.6549) (0.1653,-0.6359) (0.1566,-0.6172) (0.1483,-0.5988) (0.1403,-0.5808) (0.1325,-0.5631) (0.1251,-0.5457) (0.1180,-0.5286) (0.1112,-0.5118) (0.1047,-0.4953) (0.0984,-0.4790) (0.0924,-0.4630) (0.0866,-0.4473) (0.0811,-0.4318) (0.0759,-0.4165) (0.0708,-0.4015) (0.0660,-0.3867) (0.0614,-0.3721) (0.0570,-0.3577) (0.0528,-0.3435) (0.0488,-0.3295) (0.0450,-0.3157) (0.0414,-0.3021) (0.0379,-0.2887) (0.0347,-0.2754) (0.0316,-0.2624) (0.0287,-0.2495) (0.0259,-0.2367) (0.0233,-0.2241) (0.0209,-0.2117) (0.0186,-0.1994) (0.0165,-0.1873) (0.0145,-0.1754) (0.0127,-0.1635) (0.0110,-0.1518) (0.0094,-0.1403) (0.0080,-0.1288) (0.0066,-0.1175) (0.0055,-0.1064) (0.0044,-0.0953) (0.0035,-0.0844) (0.0026,-0.0736) (0.0019,-0.0629) (0.0013,-0.0523) (0.0009,-0.0418) (0.0005,-0.0315) (0.0002,-0.0212) (0.0001,-0.0111) (0.0000,-0.0010)};
\addplot[main] coordinates {(0.0000,0.0010) (0.0008,0.0406) (0.0032,0.0786) (0.0069,0.1153) (0.0119,0.1506) (0.0182,0.1848) (0.0255,0.2178) (0.0340,0.2498) (0.0434,0.2808) (0.0537,0.3108) (0.0649,0.3400) (0.0770,0.3683) (0.0898,0.3959) (0.1034,0.4227) (0.1177,0.4489) (0.1326,0.4743) (0.1482,0.4991) (0.1643,0.5234) (0.1811,0.5470) (0.1984,0.5701) (0.2162,0.5927) (0.2345,0.6148) (0.2533,0.6364) (0.2725,0.6576) (0.2922,0.6783) (0.3123,0.6986) (0.3328,0.7184) (0.3537,0.7379) (0.3750,0.7571) (0.3966,0.7758) (0.4186,0.7943) (0.4409,0.8124) (0.4635,0.8301) (0.4864,0.8476) (0.5097,0.8647) (0.5332,0.8816) (0.5570,0.8982) (0.5811,0.9145) (0.6054,0.9306) (0.6300,0.9464) (0.6548,0.9619) (0.6799,0.9773) (0.7052,0.9923) (0.7307,1.0072) (0.7565,1.0219) (0.7824,1.0363) (0.8086,1.0505) (0.8350,1.0645) (0.8615,1.0784) (0.8883,1.0920) (0.9152,1.1055) (0.9423,1.1188) (0.9696,1.1319) (0.9971,1.1448) (1.0247,1.1576) (1.0525,1.1702) (1.0804,1.1827) (1.1085,1.1950) (1.1367,1.2071) (1.1651,1.2191) (1.1936,1.2310) (1.2223,1.2427) (1.2511,1.2543) (1.2800,1.2658) (1.3091,1.2771) (1.3383,1.2883) (1.3676,1.2994) (1.3971,1.3103) (1.4266,1.3212) (1.4563,1.3319) (1.4861,1.3425) (1.5160,1.3530) (1.5460,1.3634) (1.5761,1.3737) (1.6063,1.3838) (1.6366,1.3939) (1.6671,1.4039) (1.6976,1.4137) (1.7282,1.4235) (1.7589,1.4332) (1.7897,1.4428) (1.8206,1.4523) (1.8516,1.4617) (1.8827,1.4710) (1.9138,1.4803) (1.9451,1.4894) (1.9764,1.4985) (2.0078,1.5075) (2.0393,1.5164) (2.0709,1.5252) (2.1025,1.5339) (2.1342,1.5426) (2.1660,1.5512) (2.1979,1.5598) (2.2298,1.5682) (2.2618,1.5766) (2.2939,1.5849) (2.3261,1.5931) (2.3583,1.6013) (2.3906,1.6094)};
\addplot[main] coordinates {(-3.0697,1.0345) (-3.0300,1.0149) (-2.9899,0.9949) (-2.9494,0.9745) (-2.9084,0.9537) (-2.8671,0.9324) (-2.8252,0.9107) (-2.7829,0.8884) (-2.7401,0.8657) (-2.6967,0.8425) (-2.6528,0.8186) (-2.6083,0.7943) (-2.5632,0.7693) (-2.5174,0.7436) (-2.4710,0.7173) (-2.4239,0.6903) (-2.3760,0.6625) (-2.3273,0.6339) (-2.2778,0.6045) (-2.2274,0.5742) (-2.1761,0.5429) (-2.1237,0.5107) (-2.0702,0.4773) (-2.0157,0.4428) (-1.9598,0.4071) (-1.9027,0.3701) (-1.8441,0.3316) (-1.7840,0.2916) (-1.7222,0.2499) (-1.6586,0.2064) (-1.5931,0.1610) (-1.5253,0.1133) (-1.4552,0.0633) (-1.3825,0.0106) (-1.3068,-0.0449) (-1.2278,-0.1038) (-1.1452,-0.1663) (-1.0584,-0.2330) (-0.9668,-0.3045) (-0.8698,-0.3815) (-0.7663,-0.4649) (-0.6552,-0.5558) (-0.5350,-0.6560) (-0.4036,-0.7672) (-0.2583,-0.8924) (-0.0950,-1.0356) (0.0922,-1.2027) (0.3131,-1.4035) (0.5846,-1.6549) (0.9415,-1.9917) (1.4731,-2.5032)};
\addplot[main] coordinates {(1.9325,-2.9125) (1.5372,-2.5072) (1.2596,-2.2196) (1.0465,-1.9965) (0.8742,-1.8143) (0.7301,-1.6601) (0.6065,-1.5266) (0.4988,-1.4088) (0.4034,-1.3035) (0.3181,-1.2082) (0.2411,-1.1212) (0.1710,-1.0412) (0.1069,-0.9671) (0.0479,-0.8981) (-0.0066,-0.8335) (-0.0573,-0.7729) (-0.1044,-0.7158) (-0.1485,-0.6617) (-0.1898,-0.6104) (-0.2286,-0.5616) (-0.2651,-0.5151) (-0.2996,-0.4706) (-0.3321,-0.4281) (-0.3630,-0.3873) (-0.3922,-0.3480) (-0.4200,-0.3103) (-0.4463,-0.2739) (-0.4714,-0.2389) (-0.4953,-0.2049) (-0.5181,-0.1722) (-0.5399,-0.1404) (-0.5607,-0.1096) (-0.5805,-0.0798) (-0.5995,-0.0508) (-0.6177,-0.0226) (-0.6351,0.0048) (-0.6518,0.0314) (-0.6678,0.0574) (-0.6831,0.0827) (-0.6978,0.1074) (-0.7119,0.1315) (-0.7255,0.1550) (-0.7385,0.1780) (-0.7509,0.2005) (-0.7629,0.2225) (-0.7745,0.2440) (-0.7855,0.2650) (-0.7961,0.2857) (-0.8064,0.3059) (-0.8162,0.3257) (-0.8256,0.3451) (-0.8347,0.3641) (-0.8434,0.3828) (-0.8517,0.4012) (-0.8597,0.4192) (-0.8675,0.4369) (-0.8749,0.4543) (-0.8820,0.4714) (-0.8888,0.4882) (-0.8953,0.5047) (-0.9016,0.5210) (-0.9076,0.5370) (-0.9134,0.5527) (-0.9189,0.5682) (-0.9241,0.5835) (-0.9292,0.5985) (-0.9340,0.6133) (-0.9386,0.6279) (-0.9430,0.6423) (-0.9472,0.6565) (-0.9512,0.6705) (-0.9550,0.6843) (-0.9586,0.6979) (-0.9621,0.7113) (-0.9653,0.7246) (-0.9684,0.7376) (-0.9713,0.7505) (-0.9741,0.7633) (-0.9767,0.7759) (-0.9791,0.7883) (-0.9814,0.8006) (-0.9835,0.8127) (-0.9855,0.8246) (-0.9873,0.8365) (-0.9890,0.8482) (-0.9906,0.8597) (-0.9920,0.8712) (-0.9934,0.8825) (-0.9945,0.8936) (-0.9956,0.9047) (-0.9965,0.9156) (-0.9974,0.9264) (-0.9981,0.9371) (-0.9987,0.9477) (-0.9991,0.9582) (-0.9995,0.9685) (-0.9998,0.9788) (-0.9999,0.9889) (-1.0000,0.9990)};
\addplot[main] coordinates {(-1.0000,1.0010) (-0.9992,1.0406) (-0.9968,1.0786) (-0.9931,1.1153) (-0.9881,1.1506) (-0.9818,1.1848) (-0.9745,1.2178) (-0.9660,1.2498) (-0.9566,1.2808) (-0.9463,1.3108) (-0.9351,1.3400) (-0.9230,1.3683) (-0.9102,1.3959) (-0.8966,1.4227) (-0.8823,1.4489) (-0.8674,1.4743) (-0.8518,1.4991) (-0.8357,1.5234) (-0.8189,1.5470) (-0.8016,1.5701) (-0.7838,1.5927) (-0.7655,1.6148) (-0.7467,1.6364) (-0.7275,1.6576) (-0.7078,1.6783) (-0.6877,1.6986) (-0.6672,1.7184) (-0.6463,1.7379) (-0.6250,1.7571) (-0.6034,1.7758) (-0.5814,1.7943) (-0.5591,1.8124) (-0.5365,1.8301) (-0.5136,1.8476) (-0.4903,1.8647) (-0.4668,1.8816) (-0.4430,1.8982) (-0.4189,1.9145) (-0.3946,1.9306) (-0.3700,1.9464) (-0.3452,1.9619) (-0.3201,1.9773) (-0.2948,1.9923) (-0.2693,2.0072) (-0.2435,2.0219) (-0.2176,2.0363) (-0.1914,2.0505) (-0.1650,2.0645) (-0.1385,2.0784) (-0.1117,2.0920) (-0.0848,2.1055) (-0.0577,2.1188) (-0.0304,2.1319) (-0.0029,2.1448) (0.0247,2.1576) (0.0525,2.1702) (0.0804,2.1827) (0.1085,2.1950) (0.1367,2.2071) (0.1651,2.2191) (0.1936,2.2310) (0.2223,2.2427) (0.2511,2.2543) (0.2800,2.2658) (0.3091,2.2771) (0.3383,2.2883) (0.3676,2.2994) (0.3971,2.3103) (0.4266,2.3212) (0.4563,2.3319) (0.4861,2.3425) (0.5160,2.3530) (0.5460,2.3634) (0.5761,2.3737) (0.6063,2.3838) (0.6366,2.3939) (0.6671,2.4039) (0.6976,2.4137) (0.7282,2.4235) (0.7589,2.4332) (0.7897,2.4428) (0.8206,2.4523) (0.8516,2.4617) (0.8827,2.4710) (0.9138,2.4803) (0.9451,2.4894) (0.9764,2.4985) (1.0078,2.5075) (1.0393,2.5164) (1.0709,2.5252) (1.1025,2.5339) (1.1342,2.5426) (1.1660,2.5512) (1.1979,2.5598) (1.2298,2.5682) (1.2618,2.5766) (1.2939,2.5849) (1.3261,2.5931) (1.3583,2.6013) (1.3906,2.6094)};
\addplot[main] coordinates {(-1.6931,-1.3069) (-1.6629,-1.3170) (-1.6326,-1.3272) (-1.6022,-1.3375) (-1.5717,-1.3479) (-1.5411,-1.3584) (-1.5103,-1.3691) (-1.4795,-1.3798) (-1.4485,-1.3907) (-1.4174,-1.4017) (-1.3862,-1.4128) (-1.3549,-1.4240) (-1.3234,-1.4354) (-1.2918,-1.4469) (-1.2601,-1.4585) (-1.2282,-1.4703) (-1.1962,-1.4822) (-1.1641,-1.4942) (-1.1318,-1.5064) (-1.0993,-1.5188) (-1.0667,-1.5313) (-1.0340,-1.5439) (-1.0010,-1.5567) (-0.9679,-1.5697) (-0.9347,-1.5829) (-0.9013,-1.5962) (-0.8676,-1.6097) (-0.8338,-1.6234) (-0.7998,-1.6373) (-0.7657,-1.6514) (-0.7313,-1.6657) (-0.6967,-1.6802) (-0.6619,-1.6949) (-0.6268,-1.7098) (-0.5916,-1.7250) (-0.5561,-1.7404) (-0.5204,-1.7560) (-0.4844,-1.7719) (-0.4482,-1.7880) (-0.4117,-1.8044) (-0.3749,-1.8211) (-0.3379,-1.8380) (-0.3005,-1.8552) (-0.2629,-1.8728) (-0.2249,-1.8906) (-0.1866,-1.9088) (-0.1480,-1.9274) (-0.1090,-1.9462) (-0.0697,-1.9655) (-0.0300,-1.9851) (0.0101,-2.0051) (0.0506,-2.0255) (0.0916,-2.0463) (0.1329,-2.0676) (0.1748,-2.0893) (0.2171,-2.1116) (0.2599,-2.1343) (0.3033,-2.1575) (0.3472,-2.1814) (0.3917,-2.2057) (0.4368,-2.2307) (0.4826,-2.2564) (0.5290,-2.2827) (0.5761,-2.3097) (0.6240,-2.3375) (0.6727,-2.3661) (0.7222,-2.3955) (0.7726,-2.4258) (0.8239,-2.4571) (0.8763,-2.4893) (0.9298,-2.5227) (0.9843,-2.5572) (1.0402,-2.5929) (1.0973,-2.6299) (1.1559,-2.6684) (1.2160,-2.7084) (1.2778,-2.7501) (1.3414,-2.7936) (1.4069,-2.8390) (1.4747,-2.8867) (1.5448,-2.9367) (1.6175,-2.9894) (1.6932,-3.0449)};
\addplot[main] coordinates {(2.4195,-3.0798) (2.4005,-3.0508) (2.3823,-3.0226) (2.3649,-2.9952) (2.3482,-2.9686) (2.3322,-2.9426) (2.3169,-2.9173) (2.3022,-2.8926) (2.2881,-2.8685) (2.2745,-2.8450) (2.2615,-2.8220) (2.2491,-2.7995) (2.2371,-2.7775) (2.2255,-2.7560) (2.2145,-2.7350) (2.2039,-2.7143) (2.1936,-2.6941) (2.1838,-2.6743) (2.1744,-2.6549) (2.1653,-2.6359) (2.1566,-2.6172) (2.1483,-2.5988) (2.1403,-2.5808) (2.1325,-2.5631) (2.1251,-2.5457) (2.1180,-2.5286) (2.1112,-2.5118) (2.1047,-2.4953) (2.0984,-2.4790) (2.0924,-2.4630) (2.0866,-2.4473) (2.0811,-2.4318) (2.0759,-2.4165) (2.0708,-2.4015) (2.0660,-2.3867) (2.0614,-2.3721) (2.0570,-2.3577) (2.0528,-2.3435) (2.0488,-2.3295) (2.0450,-2.3157) (2.0414,-2.3021) (2.0379,-2.2887) (2.0347,-2.2754) (2.0316,-2.2624) (2.0287,-2.2495) (2.0259,-2.2367) (2.0233,-2.2241) (2.0209,-2.2117) (2.0186,-2.1994) (2.0165,-2.1873) (2.0145,-2.1754) (2.0127,-2.1635) (2.0110,-2.1518) (2.0094,-2.1403) (2.0080,-2.1288) (2.0066,-2.1175) (2.0055,-2.1064) (2.0044,-2.0953) (2.0035,-2.0844) (2.0026,-2.0736) (2.0019,-2.0629) (2.0013,-2.0523) (2.0009,-2.0418) (2.0005,-2.0315) (2.0002,-2.0212) (2.0001,-2.0111) (2.0000,-2.0010)};
\addplot[main] coordinates {(2.0000,-1.9990) (2.0008,-1.9594) (2.0032,-1.9214) (2.0069,-1.8847) (2.0119,-1.8494) (2.0182,-1.8152) (2.0255,-1.7822) (2.0340,-1.7502) (2.0434,-1.7192) (2.0537,-1.6892) (2.0649,-1.6600) (2.0770,-1.6317) (2.0898,-1.6041) (2.1034,-1.5773) (2.1177,-1.5511) (2.1326,-1.5257) (2.1482,-1.5009) (2.1643,-1.4766) (2.1811,-1.4530) (2.1984,-1.4299) (2.2162,-1.4073) (2.2345,-1.3852) (2.2533,-1.3636) (2.2725,-1.3424) (2.2922,-1.3217) (2.3123,-1.3014) (2.3328,-1.2816) (2.3537,-1.2621) (2.3750,-1.2429) (2.3966,-1.2242) (2.4186,-1.2057) (2.4409,-1.1876) (2.4635,-1.1699) (2.4864,-1.1524) (2.5097,-1.1353) (2.5332,-1.1184) (2.5570,-1.1018) (2.5811,-1.0855) (2.6054,-1.0694) (2.6300,-1.0536) (2.6548,-1.0381) (2.6799,-1.0227) (2.7052,-1.0077) (2.7307,-0.9928) (2.7565,-0.9781) (2.7824,-0.9637) (2.8086,-0.9495) (2.8350,-0.9355) (2.8615,-0.9216) (2.8883,-0.9080) (2.9152,-0.8945) (2.9423,-0.8812) (2.9696,-0.8681) (2.9971,-0.8552) (3.0247,-0.8424) (3.0525,-0.8298) (3.0804,-0.8173)};
\addplot[dashed,domain=-3:3] {-x-1};
\addplot[only marks,mark=*,black] coordinates {(0,0)(-1,1)(0,-2)};
\end{axis}\end{tikzpicture}\end{center}
图中短线段按各方程斜率独立绘制; (d)的曲线用数值积分描绘, 不作为解析证明.\end{solution}
\begin{exercise}\label{cw:bank-1C-2}
原题1C-2. 在第一象限重点绘 $y'=-y/(x^2+y^2)$ 的方向场. 结合图形和方程说明满足 $y(0)=1$ 的解在 $x>0$ 时 (a) 单调递减; (b) 始终为正.
\end{exercise}
\begin{solution}
官方答案译审. 正半平面斜率为负, 且等倾线 $x^2+y^2+y/c=0$ 为圆弧. 在 $x>0$ 的任何点, 右端在 $y=0$ 附近光滑, 零函数是解. 唯一性使起初为正的解不能在有限正 $x$ 达到或穿越零; 因而 $y>0$, 继而 $y'<0$. 解在任意有限正区间内有界且远离唯一奇点 $(0,0)$, 可以继续延拓.
\begin{center}\begin{tikzpicture}\begin{axis}[width=6cm,height=6cm,axis lines=middle,xmin=0,xmax=2,ymin=0,ymax=2,xlabel=$x$,ylabel=$y$]
\draw[gray] (axis cs:0.20975,0.33050)--(axis cs:0.29025,0.16950);
\draw[gray] (axis cs:0.20230,0.57632)--(axis cs:0.29770,0.42368);
\draw[gray] (axis cs:0.19238,0.81914)--(axis cs:0.30762,0.68086);
\draw[gray] (axis cs:0.18446,1.06168)--(axis cs:0.31554,0.93832);
\draw[gray] (axis cs:0.17866,1.30487)--(axis cs:0.32134,1.19513);
\draw[gray] (axis cs:0.17449,1.54898)--(axis cs:0.32551,1.45102);
\draw[gray] (axis cs:0.17147,1.79397)--(axis cs:0.32853,1.70603);
\draw[gray] (axis cs:0.16925,2.03975)--(axis cs:0.33075,1.96025);
\draw[gray] (axis cs:0.42972,0.30622)--(axis cs:0.57028,0.19378);
\draw[gray] (axis cs:0.43636,0.56364)--(axis cs:0.56364,0.43636);
\draw[gray] (axis cs:0.43387,0.81105)--(axis cs:0.56613,0.68895);
\draw[gray] (axis cs:0.42972,1.05622)--(axis cs:0.57028,0.94378);
\draw[gray] (axis cs:0.42591,1.30110)--(axis cs:0.57409,1.19890);
\draw[gray] (axis cs:0.42283,1.54630)--(axis cs:0.57717,1.45370);
\draw[gray] (axis cs:0.42042,1.79204)--(axis cs:0.57958,1.70796);
\draw[gray] (axis cs:0.41857,2.03832)--(axis cs:0.58143,1.96168);
\draw[gray] (axis cs:0.66644,0.28343)--(axis cs:0.83356,0.21657);
\draw[gray] (axis cs:0.67335,0.54717)--(axis cs:0.82665,0.45283);
\draw[gray] (axis cs:0.67512,0.79992)--(axis cs:0.82488,0.70008);
\draw[gray] (axis cs:0.67420,1.04851)--(axis cs:0.82580,0.95149);
\draw[gray] (axis cs:0.67243,1.29563)--(axis cs:0.82757,1.20437);
\draw[gray] (axis cs:0.67059,1.54235)--(axis cs:0.82941,1.45765);
\draw[gray] (axis cs:0.66895,1.78913)--(axis cs:0.83105,1.71087);
\draw[gray] (axis cs:0.66757,2.03613)--(axis cs:0.83243,1.96387);
\draw[gray] (axis cs:0.91239,0.27061)--(axis cs:1.08761,0.22939);
\draw[gray] (axis cs:0.91644,0.53343)--(axis cs:1.08356,0.46657);
\draw[gray] (axis cs:0.91886,0.78895)--(axis cs:1.08114,0.71105);
\draw[gray] (axis cs:0.91950,1.04025)--(axis cs:1.08050,0.95975);
\draw[gray] (axis cs:0.91911,1.28946)--(axis cs:1.08089,1.21054);
\draw[gray] (axis cs:0.91828,1.53772)--(axis cs:1.08172,1.46228);
\draw[gray] (axis cs:0.91734,1.78561)--(axis cs:1.08266,1.71439);
\draw[gray] (axis cs:0.91644,2.03343)--(axis cs:1.08356,1.96657);
\draw[gray] (axis cs:1.16105,0.26369)--(axis cs:1.33895,0.23631);
\draw[gray] (axis cs:1.16324,0.52393)--(axis cs:1.33676,0.47607);
\draw[gray] (axis cs:1.16513,0.77995)--(axis cs:1.33487,0.72005);
\draw[gray] (axis cs:1.16616,1.03272)--(axis cs:1.33384,0.96728);
\draw[gray] (axis cs:1.16644,1.28343)--(axis cs:1.33356,1.21657);
\draw[gray] (axis cs:1.16625,1.53295)--(axis cs:1.33375,1.46705);
\draw[gray] (axis cs:1.16582,1.78185)--(axis cs:1.33418,1.71815);
\draw[gray] (axis cs:1.16531,2.03045)--(axis cs:1.33469,1.96955);
\draw[gray] (axis cs:1.41052,0.25967)--(axis cs:1.58948,0.24033);
\draw[gray] (axis cs:1.41175,0.51765)--(axis cs:1.58825,0.48235);
\draw[gray] (axis cs:1.41304,0.77319)--(axis cs:1.58696,0.72681);
\draw[gray] (axis cs:1.41398,1.02647)--(axis cs:1.58602,0.97353);
\draw[gray] (axis cs:1.41448,1.27804)--(axis cs:1.58552,1.22196);
\draw[gray] (axis cs:1.41462,1.52846)--(axis cs:1.58538,1.47154);
\draw[gray] (axis cs:1.41452,1.77816)--(axis cs:1.58548,1.72184);
\draw[gray] (axis cs:1.41428,2.02743)--(axis cs:1.58572,1.97257);
\draw[gray] (axis cs:1.66029,0.25718)--(axis cs:1.83971,0.24282);
\draw[gray] (axis cs:1.66101,0.51343)--(axis cs:1.83899,0.48657);
\draw[gray] (axis cs:1.66187,0.76823)--(axis cs:1.83813,0.73177);
\draw[gray] (axis cs:1.66261,1.02151)--(axis cs:1.83739,0.97849);
\draw[gray] (axis cs:1.66312,1.27348)--(axis cs:1.83688,1.22652);
\draw[gray] (axis cs:1.66339,1.52446)--(axis cs:1.83661,1.47554);
\draw[gray] (axis cs:1.66346,1.77472)--(axis cs:1.83654,1.72528);
\draw[gray] (axis cs:1.66341,2.02452)--(axis cs:1.83659,1.97548);
\draw[gray] (axis cs:1.91017,0.25553)--(axis cs:2.08983,0.24447);
\draw[gray] (axis cs:1.91062,0.51052)--(axis cs:2.08938,0.48948);
\draw[gray] (axis cs:1.91119,0.76460)--(axis cs:2.08881,0.73540);
\draw[gray] (axis cs:1.91175,1.01765)--(axis cs:2.08825,0.98235);
\draw[gray] (axis cs:1.91219,1.26973)--(axis cs:2.08781,1.23027);
\draw[gray] (axis cs:1.91249,1.52100)--(axis cs:2.08751,1.47900);
\draw[gray] (axis cs:1.91264,1.77165)--(axis cs:2.08736,1.72835);
\draw[gray] (axis cs:1.91269,2.02183)--(axis cs:2.08731,1.97817);
\end{axis}\end{tikzpicture}\end{center}
\end{solution}
\begin{exercise}\label{cw:bank-1C-3}
原题1C-3. 设 $y'=x-y$, $y(0)=1$.
\begin{enumerate}[label=\alph*)]
\item 用步长 $h=0.1$ 的 Euler 法列表估计 $y(0.1),y(0.2),y(0.3)$. 对 $y(0.3)$ 的估计偏高还是偏低? 为什么?
\item 用相同步长的改进 Euler 法(Heun 法)估计 $y(0.1)$, 说明它是否朝正确方向修正了(a)的值.
\end{enumerate}
\end{exercise}
\begin{solution}
官方答案译审. (a) $y_{j+1}=y_j+0.1(x_j-y_j)$, 得
\begin{center}\begin{tabular}{c|rrrr}$x_j$&0&0.1&0.2&0.3\\\hline$y_j$&1&0.9&0.82&0.758\end{tabular}\end{center}
精确解 $y=2e^{-x}+x-1$ 有 $y''=2e^{-x}>0$, 切线步骤位于曲线下方, 因而估计偏低. (b) 预测值为 $0.9$, 两端斜率为 $-1,-0.8$, 所以校正值 $1+0.1(-1-0.8)/2=0.91$. 精确值 $0.909674\ldots$, 修正方向向上是正确的.
\end{solution}
\begin{exercise}\label{cw:bank-1C-4}
原题1C-4. 对 $y'=x+y^2$, $y(0)=1$, 用步长 $0.1$ 的 Euler 法列表估计 $y(0.1),y(0.2),y(0.3)$. 最后的估计偏高还是偏低?
\end{exercise}
\begin{solution}
官方答案译审. 用 $y_{j+1}=y_j+0.1(x_j+y_j^2)$ 得
\begin{center}\begin{tabular}{c|rrrr}$x_j$&0&0.1&0.2&0.3\\\hline$y_j$&1&1.1&1.231&1.4025361\end{tabular}\end{center}
在此区间 $y,y'>0$, $y''=1+2yy'>0$. 每步的切线低于真曲线, 而右端对正 $y$ 单调递增, 低估会沿步骤传播, 所以最后偏低.
\end{solution}
\begin{exercise}\label{cw:bank-1C-5}
原题1C-5. 对 $y'=x-y$, $y(0)=1$, 证明 Euler 法以 $h=1/n$ 计算的 $y(1)$ 随 $n\to\infty$ 收敛到精确值. 先用归纳法证明第 $j$ 步值为 $y_j=2(1-h)^j-1+jh$. 精确解为 $y=2e^{-x}+x-1$.
\end{exercise}
\begin{solution}
官方答案译审. $j=0$ 时公式给 $1$. 若第 $j$ 步成立, 则
\[
 y_{j+1}=(1-h)y_j+jh^2=2(1-h)^{j+1}-1+(j+1)h.
\]
于是 $j=n$, $h=1/n$ 时 $y_n=2(1-1/n)^n\to2/e$, 等于精确解的 $y(1)$.
\end{solution}
\begin{exercise}\label{cw:bank-1C-6}
原题1C-6. 对 $y'=f(x)$, $y(0)=y_0$, 用四阶 Runge--Kutta 法进行 $n$ 步计算 $y(2nh)$, 每步长度为 $2h$. 精确值为 $y_0+\int_0^{2nh}f(x)dx$. 证明数值结果等于以间距 $h$ 用复合 Simpson 法计算该积分所得的值.
\end{exercise}
\begin{solution}
官方答案译审. 在 $x_j=2jh$ 的一步, 四个斜率是 $f(x_j),f(x_j+h),f(x_j+h),f(x_j+2h)$, 因而增量为 $h[f(x_j)+4f(x_j+h)+f(x_j+2h)]/3$. 对 $j=0,\ldots,n-1$ 相加正是复合 Simpson 公式. 原题的 $h$ 是半个 Runge--Kutta 步长.
\end{solution}
\begin{exercise}\label{cw:bank-1C-7}
原题1C-7. 为保证 $a(x)y'+b(x)y=c(x)$, $y(x_0)=y_0$ 在 $[x_0-h,x_0+h]$ 上有唯一解, 存在唯一性定理要求 $a,b,c$ 满足什么条件?
\end{exercise}
\begin{solution}
官方答案译审. 一个常用充分条件是 $a,b,c$ 在 $x_0$ 附近连续且 $a(x_0)\ne0$. 缩小区间使 $a\ne0$, 标准化后的 $y'=-(b/a)y+c/a$ 及其对 $y$ 的偏导 $-b/a$ 连续, 所以存在唯一解. 更直接的条件是 $b/a,c/a$ 在该区间连续, 不必要求三个分子分母分别连续.
\end{solution}
\originalsection{1D: 几何与物理应用}
\begin{exercise}\label{cw:bank-1D-1}
原题1D-1. 将几何条件写成微分方程, 求所有满足条件的曲线 $y=y(x)$:
\begin{enumerate}[label=\alph*)]
\item 每条切线在第一象限内的线段, 被切点平分.
\item 每条法线从曲线到 $x$ 轴之间的线段长度恒为1.
\item 每条法线从曲线到 $x$ 轴之间的线段, 被 $y$ 轴平分.
\item 固定 $a$, 曲线在 $a$ 与 $x$ 之间的有向面积与 $y(x)-y(a)$ 成正比.
\end{enumerate}
\end{exercise}
\begin{solution}
官方答案译审. (a) 切点 $(x,y)$ 是两轴截点 $(2x,0),(0,2y)$ 的中点, 故 $y'=-y/x$, 得 $xy=C$, 在第一象限取 $C>0$. (b) 法线斜率为 $-1/y'$, 它与 $x$ 轴的交点为 $(x+yy',0)$, 故长度平方 $y^2(1+(y')^2)=1$. 分离得 $(x-C)^2+y^2=1$, 取可作为函数的圆弧; 另有水平线 $y=\pm1$. (c) 平分条件给法线另一端横坐标 $-x$, 故法线斜率 $y/(2x)=-1/y'$, 得 $yy'=-2x$, 所以 $x^2+y^2/2=C$. (d) 写成 $\int_a^x y(t)dt=k[y(x)-y(a)]$, 对 $x$ 求导得 $y=ky'$. $k\ne0$ 时 $y=Ce^{x/k}$; $k=0$ 只给零函数.
\end{solution}
\begin{exercise}\label{cw:bank-1D-2}
原题1D-2. 对每个曲线族: (i) 求以它为积分曲线的微分方程; (ii) 求正交轨线; (iii) 同图画出原族和正交族.
\begin{enumerate}[label=\alph*)]
\item $y$ 轴截距为斜率两倍的直线.
\item 指数曲线 $y=ce^x$.
\item 双曲线 $x^2-y^2=c$.
\item 圆心在 $y$ 轴且与 $x$ 轴相切的圆.
\end{enumerate}
\end{exercise}
\begin{solution}
官方答案译审. 正交方向将有限非零斜率 $m$ 换成 $-1/m$, 竖直/水平处用隐式曲线理解.
(a) $y=m(x+2)$, 原方程 $y'=y/(x+2)$; 正交式 $yy'=-(x+2)$, 得 $(x+2)^2+y^2=C$. (b) 原方程 $y'=y$, 正交式 $yy'=-1$, 得 $y^2+2x=C$. (c) 原式 $yy'=x$, 正交式 $y'=-y/x$, 得 $xy=C$. (d) 原族 $x^2+y^2=2ay$, 消去 $a$ 得 $y'=2xy/(x^2-y^2)$. 正交式 $y'=(y^2-x^2)/(2xy)$. 令 $y=xz$ 积分, 得 $x^2+y^2=Cx$: 圆心在 $x$ 轴并与 $y$ 轴相切的圆.
\begin{center}\begin{tikzpicture}
\begin{axis}[width=6.5cm,height=6.5cm,axis lines=middle,xmin=-4,xmax=1,ymin=-2,ymax=2,title={(a)},clip=true]
\addplot[main,domain=-4:1] {-1*(x+2)};
\addplot[main,domain=-4:1] {-0.5*(x+2)};
\addplot[main,domain=-4:1] {0.5*(x+2)};
\addplot[main,domain=-4:1] {1*(x+2)};
\addplot[black,dashed,domain=0:360,samples=90] ({-2+0.5*cos(x)},{0.5*sin(x)});
\addplot[black,dashed,domain=0:360,samples=90] ({-2+1*cos(x)},{1*sin(x)});
\addplot[black,dashed,domain=0:360,samples=90] ({-2+1.5*cos(x)},{1.5*sin(x)});
\end{axis}\end{tikzpicture}\quad
\begin{tikzpicture}\begin{axis}[width=6.5cm,height=6.5cm,axis lines=middle,xmin=-2,xmax=2,ymin=-2,ymax=2,title={(b)},clip=true]
\addplot[main,domain=-2:2] {-1*exp(x)};
\addplot[main,domain=-2:2] {-0.5*exp(x)};
\addplot[main,domain=-2:2] {0.5*exp(x)};
\addplot[main,domain=-2:2] {1*exp(x)};
\addplot[black,dashed,domain=-2:2,samples=80] ({(-1-x*x)/2},{x});
\addplot[black,dashed,domain=-2:2,samples=80] ({(0-x*x)/2},{x});
\addplot[black,dashed,domain=-2:2,samples=80] ({(1-x*x)/2},{x});
\addplot[black,dashed,domain=-2:2,samples=80] ({(2-x*x)/2},{x});
\end{axis}\end{tikzpicture}
\begin{tikzpicture}\begin{axis}[width=6.5cm,height=6.5cm,axis lines=middle,xmin=-2,xmax=2,ymin=-2,ymax=2,title={(c)},clip=true]
\addplot[main,domain=-2:2] ({0.5*cosh(x)},{0.5*sinh(x)});\addplot[main,domain=-2:2] ({-0.5*cosh(x)},{0.5*sinh(x)});\addplot[main,domain=-2:2] ({0.5*sinh(x)},{0.5*cosh(x)});\addplot[main,domain=-2:2] ({0.5*sinh(x)},{-0.5*cosh(x)});
\addplot[main,domain=-2:2] ({1*cosh(x)},{1*sinh(x)});\addplot[main,domain=-2:2] ({-1*cosh(x)},{1*sinh(x)});\addplot[main,domain=-2:2] ({1*sinh(x)},{1*cosh(x)});\addplot[main,domain=-2:2] ({1*sinh(x)},{-1*cosh(x)});
\addplot[black,dashed,domain=-2:-0.1] {-1/x};\addplot[black,dashed,domain=0.1:2] {-1/x};
\addplot[black,dashed,domain=-2:-0.1] {-0.5/x};\addplot[black,dashed,domain=0.1:2] {-0.5/x};
\addplot[black,dashed,domain=-2:-0.1] {0.5/x};\addplot[black,dashed,domain=0.1:2] {0.5/x};
\addplot[black,dashed,domain=-2:-0.1] {1/x};\addplot[black,dashed,domain=0.1:2] {1/x};
\end{axis}\end{tikzpicture}\quad
\begin{tikzpicture}\begin{axis}[width=6.5cm,height=6.5cm,axis lines=middle,xmin=-2,xmax=2,ymin=-2,ymax=2,title={(d)},clip=true]
\addplot[main,domain=0:360,samples=90] ({abs(-1)*cos(x)},{-1+abs(-1)*sin(x)});\addplot[black,dashed,domain=0:360,samples=90] ({-1+abs(-1)*cos(x)},{abs(-1)*sin(x)});
\addplot[main,domain=0:360,samples=90] ({abs(-0.5)*cos(x)},{-0.5+abs(-0.5)*sin(x)});\addplot[black,dashed,domain=0:360,samples=90] ({-0.5+abs(-0.5)*cos(x)},{abs(-0.5)*sin(x)});
\addplot[main,domain=0:360,samples=90] ({abs(0.5)*cos(x)},{0.5+abs(0.5)*sin(x)});\addplot[black,dashed,domain=0:360,samples=90] ({0.5+abs(0.5)*cos(x)},{abs(0.5)*sin(x)});
\addplot[main,domain=0:360,samples=90] ({abs(1)*cos(x)},{1+abs(1)*sin(x)});\addplot[black,dashed,domain=0:360,samples=90] ({1+abs(1)*cos(x)},{abs(1)*sin(x)});
\end{axis}\end{tikzpicture}\end{center}
实线表示原族, 虚线表示正交族.
\end{solution}
\begin{exercise}\label{cw:bank-1D-3}
原题1D-3. 容器有 $V$ 升盐水, 初浓度 $c_0$ 克/升. 浓度 $c_1$ 的盐水以 $k$ 升/分流入并立即均匀混合, 混合液以同速流出. 设盐量 $x(t)$, 浓度 $c=x/V$.
\begin{enumerate}[label=\alph*)]
\item 建立 $x$ 的方程及初值.
\item 流入纯水时求解, 记 $a=k/V$.
\item $c_1$ 恒定时求 $c(t)$ 及其极限, 解释极限.
\item 改设 $c_1(t)=c_0e^{-\alpha t}$, $\alpha>0$, $a\ne\alpha$. 求 $c(t)$, 并令 $\alpha=0$ 与(c)比较.
\end{enumerate}
\end{exercise}
\begin{solution}
官方答案译审. (a) 流入盐速率 $kc_1$, 流出速率 $kx/V$, 故 $x'+ax=kc_1$, $x(0)=Vc_0$. (b) $c=c_0e^{-at}$. (c) 积分因子给 $c=c_1+(c_0-c_1)e^{-at}\to c_1$: 长期换水使原盐水被输入盐水替换. (d) $c'+ac=ac_0e^{-\alpha t}$, 故
\[
 c(t)=\frac{ac_0}{a-\alpha}e^{-\alpha t}-\frac{\alpha c_0}{a-\alpha}e^{-at}.
\]
$\alpha=0$ 给 $c=c_0$, 与(c)取 $c_1=c_0$ 一致. 若补取 $a=\alpha$, 则 $c=c_0(1+at)e^{-at}$.
\end{solution}
\begin{exercise}\label{cw:bank-1D-4}
原题1D-4. 放射性物质 $A$ 衰变为 $B$, $B$ 再衰变为 $C$.
\begin{enumerate}[label=\alph*)]
\item 衰变常数分别为 $\lambda_1,\lambda_2$ (定义为 $\ln2$/半衰期), 初始量分别为 $A_0,B_0$. 建立并解 $B(t)$ 的方程, 假设 $\lambda_1\ne\lambda_2$.
\item 设 $\lambda_1=1,\lambda_2=2$, 求 $B$ 达到最大值的时间.
\end{enumerate}
\end{exercise}
\begin{solution}
官方答案译审并补核边界. (a) $A=A_0e^{-\lambda_1t}$, $B'+\lambda_2B=\lambda_1A_0e^{-\lambda_1t}$, $B(0)=B_0$, 所以
\[
 B=\frac{\lambda_1A_0}{\lambda_2-\lambda_1}e^{-\lambda_1t}+\left(B_0-\frac{\lambda_1A_0}{\lambda_2-\lambda_1}\right)e^{-\lambda_2t}.
\]
(b) $B=A_0e^{-t}+(B_0-A_0)e^{-2t}$. 对非负初量, 若 $A_0>2B_0$, 导数零点为 $t_* =\ln[2(A_0-B_0)/A_0]>0$, 导数先正后负. 若 $A_0\le2B_0$, $B$ 从 $t=0$ 起递减或立即开始递减, 最大值在 $t=0$. $A_0=B_0=0$ 时处处取最大值零.
\end{solution}
\begin{exercise}\label{cw:bank-1D-5}
原题1D-5. Newton 冷却律说温度变化率与物体和环境的温差成正比. 一锅沸水在 $t=0$ 离火, 5分钟后为 $80^\circ\mathrm C$, 室温 $20^\circ\mathrm C$. 建模并求何时降到 $60^\circ\mathrm C$.
\end{exercise}
\begin{solution}
官方答案译审. $T'=-k(T-20)$, $T(0)=100$, 故 $T=20+80e^{-kt}$. 由 $T(5)=80$ 得 $e^{-5k}=3/4$, $k=\ln(4/3)/5$. 降到60时 $e^{-kt}=1/2$, 所以 $t=5\ln2/\ln(4/3)\approx12.05$ 分钟.
\end{solution}
\begin{exercise}\label{cw:bank-1D-6}
原题1D-6. 质量 $m$ 的物体在重力下穿过空气下落, 重力加速度为 $g$. 求速度 $v(t)$ 与终端速度, 分别假设 (a) 阻力为 $kv$ (小速度); (b) 阻力为 $kv^2$ (高速度). (b)可记 $a=\sqrt{gm/k}$.
\end{exercise}
\begin{solution}
官方答案以静止释放为初值; 原题未指定初速, 此处同时给初速参数. 取向下为正. (a) $mv'=mg-kv$, 故 $v=mg/k+(v_0-mg/k)e^{-kt/m}$, 终端速度 $mg/k$. (b) 在向下运动阶段 $mv'=mg-kv^2$. 对 $0\le v_0<a$, 分离得
\[
 v=a\tanh\left(\frac{g}{a}t+\operatorname{arctanh}(v_0/a)\right).
\]
$v_0=0$ 时为 $a\tanh(gt/a)$, 终端速度 $a$. $v_0=a$ 为常数解, $v_0>a$ 时用 $a\coth(gt/a+\operatorname{arccoth}(v_0/a))$, 极限仍为 $a$.
\end{solution}
\begin{exercise}\label{cw:bank-1D-7}
原题1D-7. 受载缆索悬于两个支点, 最低点为 $Q$. 弧段 $QP$ 受其总载荷 $W$、$Q$ 处恒定水平张力 $T_Q$、$P$ 处张力 $T$. $s$ 为 $QP$ 弧长, $x,y$ 是以 $Q$ 附近水平、竖直方向为坐标的曲线参数.
\begin{enumerate}[label=\alph*)]
\item 证明 $dx/T_Q=dy/W=ds/T$.
\item 自重均匀时证明 (i) $y''=k\sqrt{1+(y')^2}$; (ii) $y=\sqrt{s^2+c^2}+c_1$.
\item 缆索自重可忽略、载荷按水平距离均匀分布时, 求悬索桥缆索 $y(x)$.
\item 缆索自重可忽略, 悬挂等间距且密集的竖杆, 杆底端在同一水平线. 以此线为 $x$ 轴, 证明 Marseille 幕帘曲线满足 $y''=k^2y$.
\end{enumerate}
\end{exercise}
\begin{solution}
官方答案译审. (a) 切线角 $\phi$ 满足 $T\cos\phi=T_Q$, $T\sin\phi=W$, 而 $\cos\phi=dx/ds$, $\sin\phi=dy/ds$, 即所求比例. (b) 每单位弧长重量为 $\rho$, 则 $W=\rho s$, $y'=W/T_Q=ks$, $k=\rho/T_Q$. 求导给(i). 又 $ds/dy=T/W=\sqrt{T_Q^2+\rho^2s^2}/(\rho s)$, 积分得 $y=\sqrt{s^2+(T_Q/\rho)^2}+c_1$. 相应 $y(x)=k^{-1}\cosh(k(x-x_Q))+C$. (c) 水平载荷密度 $\lambda$ 时 $W=\lambda(x-x_Q)$, 故 $y'=\lambda(x-x_Q)/T_Q$, $y=\lambda(x-x_Q)^2/(2T_Q)+C$. (d) 竖杆重量与长度 $y$ 成正比, 密集极限的水平载荷密度为 $\lambda y$. 因而 $W'=\lambda y$, $y''=W'/T_Q=(\lambda/T_Q)y=k^2y$. 最低点对称的解为 $C\cosh(k(x-x_Q))$.
\end{solution}
\originalsection{1E: 一阶自治方程}
\begin{exercise}\label{cw:bank-1E-1}
原题1E-1. 对 $dx/dt=f(x)$: (i) 在水平 $x$ 轴标出平衡点和运动箭头, 判定稳定、不稳定或半稳定, 并用虚线画 $f(x)$ 说明符号来源; (ii) 在 $tx$ 平面画典型解, 包括常数解和非常数解的渐近行为.
\begin{enumerate}[label=\alph*)]
\item $x'=x^2+2x$. \item $x'=-(x-1)^2$. \item $x'=2x-x^2$. \item $x'=(2-x)^3$.
\end{enumerate}
\end{exercise}
\begin{solution}
官方答案译审. (a) 符号在 $(-\infty,-2),(-2,0),(0,\infty)$ 为 $+,-,+$, 所以 $-2$ 稳定, $0$ 不稳定. (b) 除 $1$ 外均为负, $1$ 从右吸引、从左排斥, 半稳定. (c) 符号在 $(-\infty,0),(0,2),(2,\infty)$ 为 $-,+,-$, 所以 $0$ 不稳定, $2$ 稳定. (d) $x<2$ 时正、$x>2$ 时负, $2$ 稳定. 典型解可由分离求得: (a) $x=2Ce^{2t}/(1-Ce^{2t})$; (b) $x=1+1/(t+C)$; (c) $x=2/(1+Ce^{-2t})$; (d) $x=2\pm(2t+C)^{-1/2}$. 各式只取其最大有定义区间, 并补回常数解. 发散分支可能在有限时间爆破, 不能把所有解都画成全时间存在.
\begin{center}
\begin{tikzpicture}\begin{axis}[width=6cm,height=4.2cm,axis lines=middle,xmin=-3,xmax=1,ymin=-3,ymax=3,xlabel=$x$,ylabel=$f(x)$,title={(a)},clip=true]\addplot[dashed,domain=-3:1,samples=70] {x*x+2*x};
\draw[-{Stealth},main,thick] (axis cs:-2.70,0)--(axis cs:-2.45,0);
\draw[-{Stealth},main,thick] (axis cs:-1.00,0)--(axis cs:-1.25,0);
\draw[-{Stealth},main,thick] (axis cs:0.70,0)--(axis cs:0.95,0);
\fill[main] (axis cs:-2,0) circle (2pt);
\fill[main] (axis cs:0,0) circle (2pt);
\end{axis}\end{tikzpicture}\quad\begin{tikzpicture}\begin{axis}[width=6cm,height=4.2cm,axis lines=middle,xmin=0,xmax=2,ymin=-3,ymax=4,xlabel=$t$,ylabel=$x$,clip=true]
\addplot[main,domain=0:2,samples=80] {-2};
\addplot[main,domain=0:2,samples=80] {0};
\addplot[main,domain=0:2,samples=80] {2*(-0.3)*exp(2*x)/(1+0.3*exp(2*x))};
\addplot[main,domain=0:1.49,samples=80] {2*0.05*exp(2*x)/(1-0.05*exp(2*x))};
\addplot[main,domain=0:2,samples=80] {4*exp(2*x)/(1-2*exp(2*x))};
\end{axis}\end{tikzpicture}
\begin{tikzpicture}\begin{axis}[width=6cm,height=4.2cm,axis lines=middle,xmin=-1,xmax=3,ymin=-3,ymax=3,xlabel=$x$,ylabel=$f(x)$,title={(b)},clip=true]\addplot[dashed,domain=-1:3,samples=70] {-(x-1)^2};
\draw[-{Stealth},main,thick] (axis cs:-0.70,0)--(axis cs:-0.95,0);
\draw[-{Stealth},main,thick] (axis cs:0.30,0)--(axis cs:0.05,0);
\draw[-{Stealth},main,thick] (axis cs:2.70,0)--(axis cs:2.45,0);
\fill[main] (axis cs:1,0) circle (2pt);
\end{axis}\end{tikzpicture}\quad\begin{tikzpicture}\begin{axis}[width=6cm,height=4.2cm,axis lines=middle,xmin=0,xmax=2,ymin=-3,ymax=4,xlabel=$t$,ylabel=$x$,clip=true]
\addplot[main,domain=0:2,samples=80] {1};
\addplot[main,domain=0:2,samples=80] {1+1/(x+1)};
\addplot[main,domain=0:2,samples=80] {1+1/(x-3)};
\end{axis}\end{tikzpicture}
\begin{tikzpicture}\begin{axis}[width=6cm,height=4.2cm,axis lines=middle,xmin=-1,xmax=3,ymin=-3,ymax=3,xlabel=$x$,ylabel=$f(x)$,title={(c)},clip=true]\addplot[dashed,domain=-1:3,samples=70] {2*x-x*x};
\draw[-{Stealth},main,thick] (axis cs:-0.70,0)--(axis cs:-0.95,0);
\draw[-{Stealth},main,thick] (axis cs:1.00,0)--(axis cs:1.25,0);
\draw[-{Stealth},main,thick] (axis cs:2.70,0)--(axis cs:2.45,0);
\fill[main] (axis cs:0,0) circle (2pt);
\fill[main] (axis cs:2,0) circle (2pt);
\end{axis}\end{tikzpicture}\quad\begin{tikzpicture}\begin{axis}[width=6cm,height=4.2cm,axis lines=middle,xmin=0,xmax=2,ymin=-3,ymax=4,xlabel=$t$,ylabel=$x$,clip=true]
\addplot[main,domain=0:2,samples=80] {0};
\addplot[main,domain=0:2,samples=80] {2};
\addplot[main,domain=0:2,samples=80] {2/(1+3*exp(-2*x))};
\addplot[main,domain=0:2,samples=80] {2/(1-0.5*exp(-2*x))};
\addplot[main,domain=0:0.545,samples=100] {2/(1-3*exp(-2*x))};
\end{axis}\end{tikzpicture}
\begin{tikzpicture}\begin{axis}[width=6cm,height=4.2cm,axis lines=middle,xmin=0,xmax=4,ymin=-3,ymax=3,xlabel=$x$,ylabel=$f(x)$,title={(d)},clip=true]\addplot[dashed,domain=0:4,samples=70] {(2-x)^3};
\draw[-{Stealth},main,thick] (axis cs:0.30,0)--(axis cs:0.55,0);
\draw[-{Stealth},main,thick] (axis cs:1.50,0)--(axis cs:1.75,0);
\draw[-{Stealth},main,thick] (axis cs:3.70,0)--(axis cs:3.45,0);
\fill[main] (axis cs:2,0) circle (2pt);
\end{axis}\end{tikzpicture}\quad\begin{tikzpicture}\begin{axis}[width=6cm,height=4.2cm,axis lines=middle,xmin=0,xmax=2,ymin=-3,ymax=4,xlabel=$t$,ylabel=$x$,clip=true]
\addplot[main,domain=0:2,samples=80] {2};
\addplot[main,domain=0:2,samples=80] {2+1/sqrt(2*x+0.4)};
\addplot[main,domain=0:2,samples=80] {2-1/sqrt(2*x+0.4)};
\end{axis}\end{tikzpicture}
\end{center}\end{solution}
\originalchapter{题库二: 高阶线性微分方程}
\originalsection{2A: 二阶线性方程的一般性质}
\begin{exercise}\label{cw:bank-2A-1}
原题2A-1. 设 $L=D^2+p(x)D+q(x)$, 方程为 $Ly=r(x)$. (a) 对任意二次可微函数 $u_1,u_2$ 和常数 $c$, 验证 $L(u_1+u_2)=Lu_1+Lu_2$, $L(cu_1)=cLu_1$. (b) 若 $y_p$ 是一个特解, 证明全部解恰为 $y=y_c+y_p$, 其中 $Ly_c=0$.
\end{exercise}
\begin{solution}
官方答案译审. (a) 求导、乘以固定函数及相加都保持加法和数乘, 逐项展开即得两式. (b) 若 $Ly_c=0$, 则 $L(y_c+y_p)=r$; 反之任意解 $y$ 满足 $L(y-y_p)=r-r=0$, 所以 $y-y_p$ 是齐次解.
\end{solution}
\begin{exercise}\label{cw:bank-2A-2}
原题2A-2. (a) 消去常数, 求通解为 $y=c_1e^x+c_2e^{2x}$ 的二阶齐次线性方程. (b) 验证任意初值 $y(x_0)=y_0$, $y\prime(x_0)=v_0$ 都能实现.
\end{exercise}
\begin{solution}
官方答案译审. (a) $y\prime-y=c_2e^{2x}$, 再求导得 $y\prime\prime-3y\prime+2y=0$. (b) 初值系数矩阵行列式 $e^{3x_0}\ne0$, 故有唯一 $c_1,c_2$; 具体 $c_2=e^{-2x_0}(v_0-y_0)$, $c_1=e^{-x_0}(2y_0-v_0)$.
\end{solution}
\begin{exercise}\label{cw:bank-2A-3}
原题2A-3. (a) 消去常数, 求通解为 $y=c_1x+c_2x^2$ 的二阶齐次线性方程. (b) 证明它没有满足 $y(0)=1,y\prime(0)=1$ 的解. (c) 为什么不违背二阶线性方程存在定理?
\end{exercise}
\begin{solution}
官方答案译审. (a) $y\prime\prime=2c_2$, 消去得 $x^2y\prime\prime-2xy\prime+2y=0$. (b) 在原方程 $x=0$ 处直接得到 $2y(0)=0$, 与初值矛盾. (c) 标准式系数 $-2/x,2/x^2$ 在0不连续, 不满足定理条件.
\end{solution}
\begin{exercise}\label{cw:bank-2A-4}
原题2A-4. 考虑 $y\prime\prime+p(x)y\prime+q(x)y=0$. (a) 若 $p,q$ 在全实轴连续, 证明图形在某处与 $x$ 轴相切的解恒为零. (b) 写出以 $x^2$ 为解的上述形式方程, 说明它为什么不违背(a).
\end{exercise}
\begin{solution}
官方答案译审. (a) 相切点 $x_0$ 给 $y(x_0)=y\prime(x_0)=0$. 零解有相同初值, 全实轴上的唯一性给 $y\equiv0$. (b) $xy\prime\prime-y\prime=0$, 即 $y\prime\prime-y\prime/x=0$ ($x\ne0$); 标准式系数在0不连续.
\end{solution}
\begin{exercise}\label{cw:bank-2A-5}
原题2A-5. 计算 Wronskian, 证明下列函数对线性无关: (a) $e^{m_1x},e^{m_2x}$, $m_1\ne m_2$; (b) $e^{mx},xe^{mx}$; (b)允许 $m=0$ 吗?
\end{exercise}
\begin{solution}
官方答案译审. (a) $W=(m_2-m_1)e^{(m_1+m_2)x}\ne0$. (b) $W=e^{2mx}\ne0$, 包括 $m=0$, 此时函数对为 $1,x$.
\end{solution}
\begin{exercise}\label{cw:bank-2A-6}
原题2A-6. 设 $y_1=x^2,y_2=x|x|$, 画 $y_2$ 的图. (a) 证明 $W(y_1,y_2)\equiv0$. (b) 证明两函数在任何含0的区间 $(a,b)$ 上不线性相关, 并说明为什么不违背解函数的 Wronskian 判别定理.
\end{exercise}
\begin{solution}
官方答案译审. $y_2$ 在右半轴为 $x^2$, 左半轴为 $-x^2$, 图由这两条抛物线弧拼成. 两侧各有 $W=0$, 在0也为0. 若 $c_1y_1+c_2y_2=0$, 右侧给 $c_1+c_2=0$, 左侧给 $c_1-c_2=0$, 所以两常数都零. 判别定理要求二者是同一正则线性方程的解; $y_2$ 在0没有二阶导数, 不满足要求.\begin{center}\begin{tikzpicture}\begin{axis}[width=6cm,height=4cm,axis lines=middle,xmin=-2,xmax=2,ymin=-4,ymax=4,xlabel=$x$,ylabel=$y_2$]\addplot[main,domain=-2:0] {-x^2};\addplot[main,domain=0:2] {x^2};\end{axis}\end{tikzpicture}\end{center}
\end{solution}
\begin{exercise}\label{cw:bank-2A-7}
原题2A-7. 设 $y_1,y_2$ 是 $y\prime\prime+py\prime+qy=0$ 的解. (a) 证明 $W\prime=-pW$. (b) 推出 $p=0$ 时 $W$ 恒定. (c) 对 $y\prime\prime+k^2y=0$, $k\ne0$, 直接验证(b), 其解为 $c_1\sin kx+c_2\cos kx$.
\end{exercise}
\begin{solution}
官方答案译审. (a) 对 $W=y_1y_2\prime-y_1\prime y_2$ 求导, 得 $W\prime=y_1y_2\prime\prime-y_1\prime\prime y_2=-p(y_1y_2\prime-y_1\prime y_2)$. (b) $W\prime=0$. (c) 取两个解的系数分别为 $(a,b),(c,d)$, 直接得到 $W=k(bc-ad)$, 与 $x$ 无关.
\end{solution}
\originalsection{2B: 降阶法}
\begin{exercise}\label{cw:bank-2B-1}
原题2B-1. 已知 $y_1=e^x$ 解 $y\prime\prime-2y\prime+y=0$, 分别 (a) 令 $y_2=ue^x$ 代入; (b) 先由2A-7求 $W$ 再求 $y_2$; (c) 用公式 $y_2=y_1\int e^{-\int p dx}/y_1^2\,dx$, 求第二解. (d) 解释结果的可能差异, 写最一般的第二解.
\end{exercise}
\begin{solution}
官方答案译审. (a) 代入得 $e^xu\prime\prime=0$, 故 $y_2=(Ax+B)e^x$, $A\ne0$. (b) $W\prime=2W$ 给 $W=Ke^{2x}$, 而 $W=e^x(y_2\prime-y_2)$, 故 $(e^{-x}y_2)\prime=K$, 给同一式. (c) 被积函数为 $e^{2x}/e^{2x}=1$, 得 $xe^x$ 加常数倍 $e^x$. (d) 非零倍数及加上第一解不改变独立性, 最一般式为 $(Ax+B)e^x$, $A\ne0$.
\end{solution}
\begin{exercise}\label{cw:bank-2B-2}
原题2B-2. 证明2B-1(c)的一般降阶公式给出的 $y_2$ 与 $y_1$ 线性无关, 可计算其 Wronskian.
\end{exercise}
\begin{solution}
官方答案译审. 设 $I\prime=e^{-\int p dx}/y_1^2$, $y_2=y_1I$, 则 $W=y_1^2I\prime=e^{-\int p dx}\ne0$. 公式先在 $y_1\ne0$ 的区间使用; 第二解可按正则方程延拓.
\end{solution}
\begin{exercise}\label{cw:bank-2B-3}
原题2B-3. 已知 $y_1=x$ 解 $x^2y\prime\prime+2xy\prime-2y=0$, 用 $y_2=xu$ 降阶求第二解.
\end{exercise}
\begin{solution}
官方答案译审. 代入得 $x^3u\prime\prime+4x^2u\prime=0$, 令 $v=u\prime$, 则 $v\prime=-4v/x$, $v=Cx^{-4}$. 再积分 $u=A+Bx^{-3}$, 第二解可取 $x^{-2}$, $x\ne0$.
\end{solution}
\begin{exercise}\label{cw:bank-2B-4}
原题2B-4. 已知 $y_1=x$ 为解, 求 $(1-x^2)y\prime\prime-2xy\prime+2y=0$ 在 $(-1,1)$ 上的通解.
\end{exercise}
\begin{solution}
官方答案译审. $e^{-\int p dx}=(1-x^2)^{-1}$, 且 $1/[x^2(1-x^2)]=1/x^2+1/[2(1-x)]+1/[2(1+x)]$. 公式给第二解 $-1+\frac{x}{2}\ln\frac{1+x}{1-x}$. 因而 $y=Ax+B[1-\frac{x}{2}\ln\frac{1+x}{1-x}]$. 该式在0平滑, 可直接验证并由唯一性延拓穿过0.
\end{solution}
\originalsection{2C: 常系数二阶方程}
\begin{exercise}\label{cw:bank-2C-1}
原题2C-1. 求通解或给定初值解: (a) $y\prime\prime-3y\prime+2y=0$; (b) $y\prime\prime+2y\prime-3y=0$, $y(0)=1,y\prime(0)=-1$; (c) $y\prime\prime+2y\prime+2y=0$; (d) $y\prime\prime-2y\prime+5y=0$, $y(0)=1,y\prime(0)=-1$; (e) $y\prime\prime-4y\prime+4y=0$, $y(0)=1,y\prime(0)=1$.
\end{exercise}
\begin{solution}
官方答案译审. 特征根依次为 $(1,2),(1,-3),(-1\pm i),(1\pm2i),(2,2)$. 因而 (a) $Ae^x+Be^{2x}$; (b) 初值给 $A+B=1,A-3B=-1$, 得 $(e^x+e^{-3x})/2$; (c) $e^{-x}(A\cos x+B\sin x)$; (d) $e^x(\cos2x-\sin2x)$; (e) $(1-x)e^{2x}$.
\end{solution}
\begin{exercise}\label{cw:bank-2C-2}
原题2C-2. 用 Wronskian 证明 $e^{ax}\cos bx,e^{ax}\sin bx$ 线性无关, 并说明 $a,b$ 的限制.
\end{exercise}
\begin{solution}
官方答案译审. $W=be^{2ax}$. 所以任意实 $a$ 都可, 必须 $b\ne0$; 若 $b=0$, 第二函数恒零.
\end{solution}
\begin{exercise}\label{cw:bank-2C-3}
原题2C-3. 考虑 $y\prime\prime+cy\prime+4y=0$, $c$ 为常数. 说明哪些 $c$ 满足: (a) 存在振荡解; (b) 所有非零解为衰减振荡.
\end{exercise}
\begin{solution}
官方答案译审. 特征根 $(-c\pm\sqrt{c^2-16})/2$. (a) 需非实根, 即 $-4<c<4$. (b) 再要求实部 $-c/2<0$, 即 $0<c<4$.
\end{solution}
\begin{exercise}\label{cw:bank-2C-4}
原题2C-4. Euler 等维方程 $x^2y\prime\prime+pxy\prime+qy=0$, $p,q$ 为常数. (a) 证明 $x=e^t$ 化为常系数方程. (b) 据此解 $x^2y\prime\prime+xy\prime+y=0$.
\end{exercise}
\begin{solution}
官方答案译审. 令 $Y(t)=y(e^t)$, 则 $xy\prime=\dot Y$, $x^2y\prime\prime=\ddot Y-\dot Y$, 所以 $\ddot Y+(p-1)\dot Y+qY=0$. (b) $\ddot Y+Y=0$, 得 $y=A\cos\ln x+B\sin\ln x$, $x>0$; 在负半轴以 $\ln|x|$ 给同样公式.
\end{solution}
\begin{exercise}\label{cw:bank-2C-5}
原题2C-5. 阻尼弹簧质量系统 $mx\prime\prime+cx\prime+kx=0$ 何时临界阻尼, 即恰处于振荡与非振荡的边界?
\end{exercise}
\begin{solution}
官方答案译审. 对 $m,k>0,c\ge0$, 特征式 $mr^2+cr+k=0$ 的判别式必须为零, 即 $c^2=4mk$, $c=2\sqrt{mk}$.
\end{solution}
\begin{exercise}\label{cw:bank-2C-6}
原题2C-6. 摆长 $L$, 质量 $m$, 切向阻力为 $mc\,d\alpha/dt$. 说明小摆角 $\alpha$ 近似服从常系数二阶方程. 无阻尼 $c=0$ 时用 $L,m,g$ 表示周期.
\end{exercise}
\begin{solution}
官方答案译审. 切向 Newton 方程为 $mL\alpha\prime\prime=-mg\sin\alpha-mc\alpha\prime$, 故 $\alpha\prime\prime+(c/L)\alpha\prime+(g/L)\alpha=0$ (用 $\sin\alpha\approx\alpha$). 无阻尼时角频率 $\sqrt{g/L}$, 周期 $2\pi\sqrt{L/g}$, 不依赖质量.
\end{solution}
\begin{exercise}\label{cw:bank-2C-7}
原题2C-7. 只写待定系数法所需的特解试探形式: (a) $y\prime\prime+2y\prime+2y=x+e^x$; (b) $y\prime\prime-4y\prime=\cos2x$; (c) $y\prime\prime+4y=3\cos2x$; (d) $y\prime\prime-2y\prime+y=3e^x$; (e) $y\prime\prime-3y\prime+2y=e^{-x}+3e^{2x}$; (f) $y\prime\prime-6y\prime+9y=2xe^{3x}$.
\end{exercise}
\begin{solution}
官方答案译审. 依次取 (a) $Ax+B+Ce^x$; (b) $A\cos2x+B\sin2x$; (c) $x(A\cos2x+B\sin2x)$; (d) $Ax^2e^x$; (e) $Ae^{-x}+Bxe^{2x}$; (f) $x^2(Ax+B)e^{3x}$. 每个共振项乘 $x$ 的次数等于对应特征根的重数.
\end{solution}
\begin{exercise}\label{cw:bank-2C-8}
原题2C-8. 求通解或初值解: (a) $y\prime\prime-6y\prime+5y=e^x$; (b) $y\prime\prime+4y=2\cos x$, $y(0)=0,y\prime(0)=1$; (c) $y\prime\prime+y\prime+y=2xe^x$; (d) $y\prime\prime-y=x^2$, $y(0)=0,y\prime(0)=-1$.
\end{exercise}
\begin{solution}
官方答案译审. (a) 特征根1,5, 试 $Axe^x$, 代入给 $-4Ae^x=e^x$, 故 $y=C_1e^x+C_2e^{5x}-xe^x/4$. (b) 特解 $2\cos x/3$, 初值得 $y=-2\cos2x/3+\sin2x/2+2\cos x/3$. (c) 试 $e^x(Ax+B)$, 移位后为 $e^x[(D^2+3D+3)(Ax+B)]$, 给 $A=2/3,B=-2/3$, 因而 $y=e^{-x/2}[C_1\cos(\sqrt3x/2)+C_2\sin(\sqrt3x/2)]+2(x-1)e^x/3$. (d) 特解 $-x^2-2$, 初值解 $y=e^x/2+3e^{-x}/2-x^2-2$.
\end{solution}
\begin{exercise}\label{cw:bank-2C-9}
原题2C-9. 对 $Ly=y\prime\prime+py\prime+qy=r$, (a) 若 $Ly_i=r_i$ ($i=1,2$), 证明 $y_1+y_2$ 是输入 $r_1+r_2$ 的特解. (b) 据此求 $y\prime\prime+2y\prime+2y=2x+\cos x$ 的特解.
\end{exercise}
\begin{solution}
官方答案译审. (a) 线性性给 $L(y_1+y_2)=r_1+r_2$. (b) 输入 $2x$ 试 $Ax+B$, 得 $A=1,B=-1$; 输入 $\cos x$ 试 $a\cos x+b\sin x$, 得 $a+2b=1,b-2a=0$, 所以 $y_p=x-1+\cos x/5+2\sin x/5$.
\end{solution}
\begin{exercise}\label{cw:bank-2C-10}
原题2C-10. 串联 RLC 电路满足 $Lq\prime\prime+Rq\prime+q/C=E$, $Li\prime\prime+Ri\prime+i/C=E\prime$, $i=q\prime$. $L,R,C$ 为电感、电阻、电容, $E$ 为电动势. (a) $R=E=0,L\ne0$ 时证明电荷周期变化并求周期. (b) $E=0$ 时电流振荡所需关系是什么? (c) $R=0,E=E_0\sin\omega t$ 时, $\omega$ 接近哪个 $\omega_0$ 会出现大振幅? 说明理由.
\end{exercise}
\begin{solution}
官方答案译审. 取物理参数 $L,C>0$. (a) $q\prime\prime+q/(LC)=0$, 周期 $2\pi\sqrt{LC}$. (b) 特征判别式 $R^2-4L/C<0$, 即 $R^2<4L/C$. (c) 自然频率 $\omega_0=1/\sqrt{LC}$; 强迫项与齐次频率相同时发生共振, 特解含 $t\sin\omega_0t$ 或 $t\cos\omega_0t$, 近共振时响应系数含 $(\omega_0^2-\omega^2)^{-1}$.
\end{solution}
\originalsection{2D: 参数变易法}
\begin{exercise}\label{cw:bank-2D-1}
原题2D-1. 用参数变易法求特解: (a) $y\prime\prime+y=\tan x$; (b) $y\prime\prime+2y\prime-3y=e^{-x}$; (c) $y\prime\prime+4y=\sec^22x$.
\end{exercise}
\begin{solution}
官方答案译审. 公式 $v_1\prime=-y_2r/W,v_2\prime=y_1r/W$. (a) 取 $(y_1,y_2)=(\cos x,\sin x)$, $W=1$, 得 $v_1=\sin x-\ln|\sec x+\tan x|$, $v_2=-\cos x$, 所以 $y_p=-\cos x\ln|\sec x+\tan x|$. (b) 取 $(e^x,e^{-3x})$, $W=-4e^{-2x}$, 得 $v_1\prime=e^{-2x}/4,v_2\prime=-e^{2x}/4$, 所以 $y_p=-e^{-x}/4$. (c) 取 $(\cos2x,\sin2x)$, $W=2$, 得 $v_1=-\sec2x/4$, $v_2=\ln|\sec2x+\tan2x|/4$, 所以 $y_p=-1/4+\sin2x\ln|\sec2x+\tan2x|/4$. 各解取三角函数右端无奇点的区间.
\end{solution}
\begin{exercise}\label{cw:bank-2D-2}
原题2D-2. Bessel 方程 $x^2y\prime\prime+xy\prime+(x^2-p^2)y=0$ 在 $p=1/2,x>0$ 时有独立解 $\sin x/\sqrt x,\cos x/\sqrt x$. 求 $x^2y\prime\prime+xy\prime+(x^2-1/4)y=x^{3/2}\cos x$ 的通解.
\end{exercise}
\begin{solution}
官方答案译审. 标准式右端 $r=\cos x/\sqrt x$, 取给定解的 $W=-1/x$. 参数变易给 $v_1\prime=\cos^2x$, $v_2\prime=-\sin x\cos x$, 可取 $v_1=x/2+\sin2x/4$, $v_2=\cos2x/4$. 组合后特解为 $\sqrt x\sin x/2$ 加一个齐次项, 故通解 $y=[A\sin x+B\cos x]/\sqrt x+\sqrt x\sin x/2$.
\end{solution}
\begin{exercise}\label{cw:bank-2D-3}
原题2D-3. 对 $y\prime\prime+py\prime+qy=r$, (a) 证明参数变易特解可写 $y_p(x)=\int_a^x [y_1(t)y_2(x)-y_2(t)y_1(x)]r(t)/W(t)\,dt$. (b) 使用不定积分产生任意积分常数, 为什么不影响结果?
\end{exercise}
\begin{solution}
官方答案译审. (a) 将 $v_1=-\int_a^x y_2r/W$, $v_2=\int_a^x y_1r/W$ 代入 $y_p=v_1y_1+v_2y_2$, 合并积分即得. (b) 常数变化使特解增加 $Ay_1+By_2$, 它属于齐次解, 所以通解集合不变.
\end{solution}
\begin{exercise}\label{cw:bank-2D-4}
原题2D-4. 什么时候必须使用参数变易法, 而不能用待定系数法求特解?
\end{exercise}
\begin{solution}
官方答案译审. 通常的待定系数法适用于常系数线性方程, 右端为有限个多项式乘指数、正弦、余弦之和. 若系数变动或右端不属于这类有限求导封闭的函数族, 该法没有相应有限试探形式; 已知齐次基本解时可用参数变易法. 这不排除特殊方程另有更简便方法.
\end{solution}
\originalsection{2E: 复数}
\begin{exercise}\label{cw:bank-2E-1}
原题2E-1. 化为极坐标形式: (a) $-1+i$; (b) $\sqrt3-i$.
\end{exercise}
\begin{solution}
官方答案译审. 模长和辐角分别给 (a) $\sqrt2e^{3\pi i/4}$; (b) $2e^{-\pi i/6}$. 辐角可加 $2\pi$ 的整数倍.
\end{solution}
\begin{exercise}\label{cw:bank-2E-2}
原题2E-2. 用两种方法把 $(1-i)/(1+i)$ 化为 $a+bi$: 全程直角坐标运算, 以及先化极坐标运算. 证明结果一致.
\end{exercise}
\begin{solution}
官方答案译审. 乘共轭分母得 $(1-i)^2/2=-i$. 极坐标法得 $(\sqrt2e^{-\pi i/4})/(\sqrt2e^{\pi i/4})=e^{-\pi i/2}=-i$.
\end{solution}
\begin{exercise}\label{cw:bank-2E-3}
原题2E-3. 证明复平面两点 $z_1,z_2$ 的距离为 $|z_2-z_1|$.
\end{exercise}
\begin{solution}
编者补解. 若 $z_j=x_j+iy_j$, 欧氏距离为 $\sqrt{(x_2-x_1)^2+(y_2-y_1)^2}$, 恰是差数的模.
\end{solution}
\begin{exercise}\label{cw:bank-2E-4}
原题2E-4. 对任意复数 $z,w$, 证明 (a) $\overline{z+w}=\bar z+\bar w$; (b) $\overline{zw}=\bar z\bar w$.
\end{exercise}
\begin{solution}
官方答案译审. 写 $z=a+bi,w=c+di$. (a) 两边都为 $(a+c)-(b+d)i$. (b) 两边都为 $(ac-bd)-(ad+bc)i$.
\end{solution}
\begin{exercise}\label{cw:bank-2E-5}
原题2E-5. 实系数多项式 $f$ 若以 $a+ib$ 为零点, 用2E-4证明 $a-ib$ 也为零点, 因此非实根成共轭对.
\end{exercise}
\begin{solution}
编者补解. 对 $f(z)=\sum c_jz^j$, $c_j\in\R$, 共轭法则给 $f(\bar z)=\overline{f(z)}$. 当 $f(z)=0$, 右端为零.
\end{solution}
\begin{exercise}\label{cw:bank-2E-6}
原题2E-6. 用 Euler 定义 $e^{i\theta}=\cos\theta+i\sin\theta$ 和三角加法公式, 证明 $e^{i\theta}e^{i\phi}=e^{i(\theta+\phi)}$.
\end{exercise}
\begin{solution}
编者补解. 展开乘积, 实部为 $\cos\theta\cos\phi-\sin\theta\sin\phi=\cos(\theta+\phi)$, 虚部为 $\sin\theta\cos\phi+\cos\theta\sin\phi=\sin(\theta+\phi)$.
\end{solution}
\begin{exercise}\label{cw:bank-2E-7}
原题2E-7. 分别用极坐标与 De Moivre 公式、二项式定理计算: (a) $(1-i)^4$; (b) $(1+i\sqrt3)^3$.
\end{exercise}
\begin{solution}
官方答案译审. (a) $(\sqrt2e^{-\pi i/4})^4=-4$; 二项式展开为 $1-4i-6+4i+1=-4$. (b) $(2e^{\pi i/3})^3=-8$; 展开为 $1+3i\sqrt3-9-3i\sqrt3=-8$.
\end{solution}
\begin{exercise}\label{cw:bank-2E-8}
原题2E-8. 用 De Moivre 公式和二项式定理, 以 $\cos\theta,\sin\theta$ 表示 $\cos3\theta,\sin3\theta$.
\end{exercise}
\begin{solution}
编者补解. $(\cos\theta+i\sin\theta)^3=\cos3\theta+i\sin3\theta$. 比较实虚部给 $\cos3\theta=\cos^3\theta-3\cos\theta\sin^2\theta$, $\sin3\theta=3\cos^2\theta\sin\theta-\sin^3\theta$.
\end{solution}
\begin{exercise}\label{cw:bank-2E-9}
原题2E-9. 把1的六个六次方根写成 $a+bi$.
\end{exercise}
\begin{solution}
官方答案译审. 根为 $e^{ik\pi/3}$ ($k=0,\ldots,5$), 即 $1,-1,(1\pm i\sqrt3)/2,(-1\pm i\sqrt3)/2$.
\end{solution}
\begin{exercise}\label{cw:bank-2E-10}
原题2E-10. 解 $x^4+16=0$.
\end{exercise}
\begin{solution}
官方答案译审. $-16=16e^{i(\pi+2k\pi)}$, 根为 $2e^{i(\pi/4+k\pi/2)}$, $k=0,1,2,3$, 即 $\sqrt2(\pm1\pm i)$, 两符号独立.
\end{solution}
\begin{exercise}\label{cw:bank-2E-11}
原题2E-11. 解 $x^4+2x^2+4=0$, 将四根写成极坐标形式和 $a+bi$ 形式.
\end{exercise}
\begin{solution}
编者补解. 令 $u=x^2$, 得 $u=-1\pm i\sqrt3=2e^{\pm2\pi i/3}$. 所以 $x=\sqrt2e^{i\pi/3},\sqrt2e^{2\pi i/3},\sqrt2e^{4\pi i/3},\sqrt2e^{5\pi i/3}$, 即 $\pm1/\sqrt2\pm i\sqrt{3/2}$, 两符号独立.
\end{solution}
\begin{exercise}\label{cw:bank-2E-12}
原题2E-12. Notes C第一页以三次方程 $x^3-6x+2=0$ 为例, 记 $A=\sqrt[3]{-1+\sqrt{-7}}$, $B=\sqrt[3]{-1-\sqrt{-7}}$. 将 $A,B$ 明确写成 $a+bi$, 并证明 $A+B$ 为实数且为该方程的根.
\end{exercise}
\begin{solution}
编者补解; 题面引用由官方Notes C第1页展开. 复立方根须选相容分支. 取 $\theta=\arg(-1+i\sqrt7)\in(\pi/2,\pi)$, 其模为 $2\sqrt2$, 令 $A=\sqrt2[\cos(\theta/3)+i\sin(\theta/3)]$, $B=\sqrt2[\cos(\theta/3)-i\sin(\theta/3)]$. 这就是各自的 $a+bi$ 形式; $AB=2$, $A^3+B^3=-2$. 令 $x=A+B\in\R$, 则 $x^3=A^3+B^3+3ABx=-2+6x$. 另外两组共轭相容分支也给其余实根; 任意独立选分支并不保证 $AB=2$.
\end{solution}
\begin{exercise}\label{cw:bank-2E-13}
原题2E-13. 证明 Notes C公式(16)的复指数法则: $e^{z+w}=e^ze^w$, 其中 $e^{a+ib}=e^a(\cos b+i\sin b)$.
\end{exercise}
\begin{solution}
编者补解. 写 $z=a+ib,w=c+id$. 实指数法则及2E-6给 $e^ze^w=e^{a+c}e^{ib}e^{id}=e^{a+c}e^{i(b+d)}=e^{z+w}$.
\end{solution}
\begin{exercise}\label{cw:bank-2E-14}
原题2E-14. 利用 $\sin x=(e^{ix}-e^{-ix})/(2i)$ 和二项式定理, 把 $\sin^4x$ 表为 $\cos4x,\cos2x$ 的组合. 为什么不应出现 $\sin4x,\sin2x$?
\end{exercise}
\begin{solution}
官方答案译审. 展开为 $(e^{4ix}-4e^{2ix}+6-4e^{-2ix}+e^{-4ix})/16$, 所以 $\sin^4x=3/8-\cos2x/2+\cos4x/8$. 原函数为偶函数, 两个正弦项为奇函数, 且线性无关, 故它们的系数都为零.
\end{solution}
\begin{exercise}\label{cw:bank-2E-15}
原题2E-15. 用复指数求 $\int e^{2x}\sin x\,dx$.
\end{exercise}
\begin{solution}
官方答案译审. 取 $\int e^{(2+i)x}dx=e^{(2+i)x}/(2+i)$ 的虚部, 得 $e^{2x}(2\sin x-\cos x)/5+C$.
\end{solution}
\begin{exercise}\label{cw:bank-2E-16}
原题2E-16. 证明 (a) $\cos x=(e^{ix}+e^{-ix})/2$; (b) $\sin x=(e^{ix}-e^{-ix})/(2i)$.
\end{exercise}
\begin{solution}
官方答案译审. Euler 公式给 $e^{ix}=\cos x+i\sin x$, $e^{-ix}=\cos x-i\sin x$; 相加、相减分别得到所需式.
\end{solution}
\begin{exercise}\label{cw:bank-2E-17}
原题2E-17. 由复指数定义和 $D(u+iv)=Du+iDv$, 推导 $D(e^{(a+ib)x})=(a+ib)e^{(a+ib)x}$.
\end{exercise}
\begin{solution}
编者补解. 对 $e^{ax}(\cos bx+i\sin bx)$ 逐项求导, 得 $e^{ax}[(a\cos bx-b\sin bx)+i(a\sin bx+b\cos bx)]$, 等于 $(a+ib)e^{(a+ib)x}$.
\end{solution}
\begin{exercise}\label{cw:bank-2E-18}
原题2E-18. 在单位圆上定位并用初等几何写出1的三个立方根的 $a+bi$ 形式.
\end{exercise}
\begin{solution}
编者补解. 三个点等距分布, 圆心角为 $0,2\pi/3,4\pi/3$, 组成等边三角形. 后两点横坐标 $-1/2$, 由单位圆方程得纵坐标 $\pm\sqrt3/2$. 所以为 $1,(-1\pm i\sqrt3)/2$.
\end{solution}
\originalsection{2F: 线性算子与高阶方程}
\begin{exercise}\label{cw:bank-2F-1}
原题2F-1. 求通解: (a) $(D-2)^3(D^2+2D+2)y=0$; (b) $(D^8+2D^4+1)y=0$; (c) $y^{(4)}+y=0$; (d) $y^{(4)}-8y\prime\prime+16y=0$; (e) $y^{(6)}-y=0$; (f) $y^{(4)}+16y=0$. (e),(f)可分别用2E-9、2E-10.
\end{exercise}
\begin{solution}
官方答案译审并补正(b). 每个重数 $m$ 的根 $\lambda$ 对应 $x^je^{\lambda x}$, $0\le j<m$, 共轭根合并为实三角形式. (a) $e^{2x}(A+Bx+Cx^2)+e^{-x}(E\cos x+F\sin x)$. (b) 原文是加号; 官方答案却按 $D^8-2D^4+1$ 分解, 须补正. 本题 $(D^4+1)^2$, 根 $\pm\alpha\pm i\alpha$ 各二重, $\alpha=1/\sqrt2$, 故\[y=e^{\alpha x}[(A+Bx)\cos\alpha x+(C+Ex)\sin\alpha x]+e^{-\alpha x}[(F+Gx)\cos\alpha x+(H+Jx)\sin\alpha x].\](c) 同四个根为单根, 去掉(b)中的各 $x$ 项. (d) 特征式 $(r^2-4)^2$, 得 $(A+Bx)e^{2x}+(C+Ex)e^{-2x}$. (e) 六次单位根给 $Ae^x+Be^{-x}+e^{x/2}(C\cos\beta x+E\sin\beta x)+e^{-x/2}(F\cos\beta x+G\sin\beta x)$, $\beta=\sqrt3/2$. (f) $e^{\sqrt2x}(A\cos\sqrt2x+B\sin\sqrt2x)+e^{-\sqrt2x}(C\cos\sqrt2x+E\sin\sqrt2x)$.
\end{solution}
\begin{exercise}\label{cw:bank-2F-2}
原题2F-2. 求 $y^{(4)}-16y=0$ 满足 $y(0)=y\prime(0)=0,y(\pi)=1$, 且存在常数 $K$ 使所有 $x>0$ 有 $|y(x)|<K$ 的解.
\end{exercise}
\begin{solution}
官方答案译审. 根为 $2,-2,2i,-2i$. 有界性排除 $e^{2x}$, 初值得 $y=C(e^{-2x}+\sin2x-\cos2x)$. 由 $y(\pi)=1$ 得 $C=(e^{-2\pi}-1)^{-1}$.
\end{solution}
\begin{exercise}\label{cw:bank-2F-3}
原题2F-3. 求通解: (a) $(D^3-D^2+2D-2)y=0$; (b) $(D^3+D^2-2)y=0$; (c) $y^{(3)}-2y\prime-4=0$; (d) $y^{(4)}+2y\prime\prime+4y=0$. 提示: 整系数首一多项式的整数根整除常数项.
\end{exercise}
\begin{solution}
官方答案译审并补正(c). (a) 特征式 $(r-1)(r^2+2)$, 得 $Ae^x+B\cos\sqrt2x+C\sin\sqrt2x$. (b) $(r-1)(r^2+2r+2)$, 得 $Ae^x+e^{-x}(B\cos x+C\sin x)$. (c) 原文常数项为 $-4$, 官方扫描答案按 $-4y$ 求解, 与原题不同. 原题令 $v=y\prime$, 得 $v\prime\prime-2v=4$, 特解 $v=-2$, 所以 $y=A+Be^{\sqrt2x}+Ce^{-\sqrt2x}-2x$. 若作者意图是 $y^{(3)}-2y\prime-4y=0$, 则扫描答案 $Ae^{2x}+e^{-x}(B\cos x+C\sin x)$ 适用. (d) 令 $u=r^2$, $u=-1\pm i\sqrt3$, 根由2E-11为 $\pm\alpha\pm i\beta$, $\alpha=1/\sqrt2,\beta=\sqrt{3/2}$. 得 $e^{\alpha x}(A\cos\beta x+B\sin\beta x)+e^{-\alpha x}(C\cos\beta x+E\sin\beta x)$.
\end{solution}
\begin{exercise}\label{cw:bank-2F-4}
原题2F-4. 两个串接弹簧的质量及弹性常数均为1, 位移 $x_1,x_2$ 满足 $x_1\prime\prime+2x_1-x_2=0$, $x_2\prime\prime+x_2-x_1=0$. (a) 消去 $x_1$ 得 $x_2$ 的四阶方程. (b) 求通解.
\end{exercise}
\begin{solution}
官方答案译审. (a) $x_1=x_2\prime\prime+x_2$, 代入第一式得 $x_2^{(4)}+3x_2\prime\prime+x_2=0$. (b) 令 $\omega_\pm^2=(3\pm\sqrt5)/2$, 则 $x_2=A\cos\omega_+t+B\sin\omega_+t+C\cos\omega_-t+E\sin\omega_-t$, $x_1=x_2\prime\prime+x_2$, 所以各频率项分别乘 $1-\omega_\pm^2$.
\end{solution}
\begin{exercise}\label{cw:bank-2F-5}
原题2F-5. 设 $y=e^{2x}\cos x$, 用算子公式求 $y\prime\prime$.
\end{exercise}
\begin{solution}
官方答案译审. 指数移位给 $D^2(e^{2x}\cos x)=e^{2x}(D+2)^2\cos x=e^{2x}(3\cos x-4\sin x)$.
\end{solution}
\begin{exercise}\label{cw:bank-2F-6}
原题2F-6. 用算子方法, 尽可能用复指数求特解: (a) $(D^2+1)y=4e^x$; (b) $y^{(3)}+y\prime\prime-y\prime+2y=2\cos x$; (c) $y\prime\prime-2y\prime+4y=e^x\cos x$; (d) $y\prime\prime-6y\prime+9y=e^{3x}$.
\end{exercise}
\begin{solution}
官方答案译审. (a) $P(1)=2$, 得 $y_p=2e^x$. (b) $P(i)=1-2i$, 取 $2e^{ix}/(1-2i)$ 的实部, 得 $2\cos x/5-4\sin x/5$. (c) $P(1+i)=2$, 得 $y_p=e^x\cos x/2$. (d) $(D-3)^2(e^{3x}u)=e^{3x}u\prime\prime$, 取 $u=x^2/2$, 得 $y_p=x^2e^{3x}/2$.
\end{solution}
\begin{exercise}\label{cw:bank-2F-7}
原题2F-7. 对 $y\prime+ay=f(x)$ 假设 $y_p=e^{-ax}u$, 用指数移位公式求特解.
\end{exercise}
\begin{solution}
官方答案译审. $(D+a)e^{-ax}u=e^{-ax}u\prime=f$, 故 $u=\int e^{ax}f(x)dx$, $y_p=e^{-ax}\int e^{ax}f(x)dx$.
\end{solution}
\originalsection{2G: 常系数线性方程的稳定性}
\begin{exercise}\label{cw:bank-2G-1}
原题2G-1. 对 $y\prime\prime+2y\prime+cy=0$: (i) 判定哪些 $c$ 对应两实根、重根、共轭非实根; (ii) 两实根时判定两负、两正或异号; (iii) 在 $c$ 轴标记各区间; (iv) 标记所有解趋于零的稳定区间.
\end{exercise}
\begin{solution}
官方答案译审. 根为 $-1\pm\sqrt{1-c}$. $c<1$ 两实根, $c=1$ 重根 $-1$, $c>1$ 非实共轭且实部 $-1$. 对 $c<0$ 两根异号; $c=0$ 为 $0,-2$; $0<c<1$ 两根都负; 不会出现两正根. 稳定区间为 $c>0$.\begin{center}\begin{tikzpicture}\draw[-{Stealth}] (-4,0)--(4,0) node[right] {$c$};\draw (0,-.12)--(0,.12) node[above] {$0$};\draw (2,-.12)--(2,.12) node[above] {$1$};\node at (-2,-.5) {异号实根};\node at (1,-.5) {两负实根};\node at (3,-.5) {非实根};\draw[main,thick] (0,-1)--(3.8,-1);\node[main] at (2,-1.4) {稳定 $c>0$};\end{tikzpicture}\end{center}
\end{solution}
\begin{exercise}\label{cw:bank-2G-2}
原题2G-2. 证明二阶稳定性系数判据的必要方向: $a_0y\prime\prime+a_1y\prime+a_2y=0$, 可先乘 $-1$ 使 $a_0>0$; 稳定意味着全部解趋于零, 证明 $a_1,a_2>0$.
\end{exercise}
\begin{solution}
官方答案译审. 实根情形稳定要求 $r_1,r_2<0$, 而 $a_1/a_0=-(r_1+r_2)>0$, $a_2/a_0=r_1r_2>0$. 共轭根 $\alpha\pm i\beta$ 情形要求 $\alpha<0$, 所以 $a_1/a_0=-2\alpha>0$, $a_2/a_0=\alpha^2+\beta^2>0$. 重根同样成立.
\end{solution}
\begin{exercise}\label{cw:bank-2G-3}
原题2G-3. 证明二阶系数稳定判据的充分方向: $a_0,a_1,a_2>0$ 时全部解趋于零, 可用求根公式并特别检查两实根情况.
\end{exercise}
\begin{solution}
官方答案译审. 若判别式负, 实部为 $-a_1/(2a_0)<0$. 若非负, 因 $a_1^2-4a_0a_2<a_1^2$, 两根 $(-a_1\pm\sqrt{a_1^2-4a_0a_2})/(2a_0)$ 都负. 重根解多项式乘衰减指数仍趋于零.
\end{solution}
\begin{exercise}\label{cw:bank-2G-4}
原题2G-4. (a) 证明稳定的实系数高阶齐次常系数方程, 所有系数非零且同号. 可将实多项式按实根和共轭对分解. (b) 若全部特征根实, 证明上述条件反过来也充分. (c) 用 $y^{(3)}+y\prime\prime+y\prime+6y=0$ 说明一般情形逆命题不成立.
\end{exercise}
\begin{solution}
编者补解. (a) 取首项系数正. 稳定根的实部负, 实线性因子是 $r+a$, $a>0$, 共轭因子是 $r^2-2\alpha r+\alpha^2+\beta^2$, $\alpha<0$. 各因子系数全正, 乘积也全正. (b) 正系数多项式在 $r\ge0$ 严格正, 无非负实根; 全部根实遂全部负. (c) $r=-2$ 是根, 分解为 $(r+2)(r^2-r+3)$, 后两根为 $(1\pm i\sqrt{11})/2$, 实部正, 所以不稳定.
\end{solution}
\begin{exercise}\label{cw:bank-2G-5}
原题2G-5. (a) 证明 $n=2$ 的 Routh--Hurwitz 条件与二阶正系数判据相同. (b) 用该条件求 $y^{(3)}+y\prime\prime+y\prime+cy=0$ 稳定的全部 $c$.
\end{exercise}
\begin{solution}
编者补解. 首项归一后, 二阶 $r^2+a_1r+a_2$ 的 Hurwitz 主子式为 $\Delta_1=a_1$, $\Delta_2=a_1a_2$, 二者正等价于 $a_1,a_2>0$. 三阶 $r^3+a_1r^2+a_2r+a_3$ 的主子式为 $a_1,a_1a_2-a_3,a_3(a_1a_2-a_3)$. 本题为 $1,1-c,c(1-c)$, 全正等价于 $0<c<1$.
\end{solution}
\begin{exercise}\label{cw:bank-2G-6}
原题2G-6. 以输入 $r(t)=At$ 作用于 $ay\prime\prime+by\prime+cy=r(t)$. (a) 假设 $a,b,c>0$, 用待定系数法求稳态特解并写成 $K(t-d)$. (b) 把 $d$ 解释为时间滞后, 分别用弹簧系统 $m,b,k$ 和电路 $R,L,C$ 表示它.
\end{exercise}
\begin{solution}
编者补解. 试 $y_p=Kt+B$, 得 $cK=A$, $bK+cB=0$. 所以 $y_p=(A/c)(t-b/c)$, $K=A/c,d=b/c$. 二阶正系数使齐次项衰减, 因而该特解是渐近稳态. 弹簧为 $d=b/k$; 电荷电路为 $d=RC$ (因 $a=L,b=R,c=1/C$). $A=0$ 时零特解的延迟参数不唯一, $b/c$ 仍是系统的形式滞后.
\end{solution}
\originalsection{2H: 脉冲响应与卷积}
\begin{exercise}\label{cw:bank-2H-1}
原题2H-1. 求 $y\prime\prime-k^2y=f(t)$ 的单位脉冲响应 $w(t)$.
\end{exercise}
\begin{solution}
官方答案译审. 因果脉冲响应满足 $w(0)=0,w\prime(0+)=1$, 齐次解给 $w(t)=\sinh(kt)/k$, $t\ge0$. $k=0$ 时取极限 $w=t$; $t<0$ 置零.
\end{solution}
\begin{exercise}\label{cw:bank-2H-2}
原题2H-2. (a) 求 $y\prime\prime-(a+b)y\prime+aby=f(t)$ 的单位脉冲响应. (b) 证明 $b\to a$ 时它趋于重根响应 $te^{at}$; 可令 $b=a+h$.
\end{exercise}
\begin{solution}
编者补解. (a) $a\ne b$ 时 $w=(e^{bt}-e^{at})/(b-a)$, 其初值为0, 导数初值为1. (b) $w=e^{at}(e^{ht}-1)/h\to te^{at}$, 即指数函数对参数的导数.
\end{solution}
\begin{exercise}\label{cw:bank-2H-3}
原题2H-3. (a) 用脉冲响应卷积求 $y\prime\prime+4y\prime+4y=e^{-2x}$, $y(0)=y\prime(0)=0$, $x\ge0$, 再用待定系数法校验. (b) 改将右端设为 $e^{-2x}$ ($0\le x\le1$), 此后为0, 求解.
\end{exercise}
\begin{solution}
官方答案译审(a), 编者补解(b). (a) $w(x)=xe^{-2x}$, 所以 $y=e^{-2x}\int_0^x(x-t)dt=x^2e^{-2x}/2$. 用 $y=e^{-2x}u$ 得 $u\prime\prime=1$ 也给此式. (b) 积分上限变为 $m=\min(x,1)$, 得 $y=e^{-2x}(xm-m^2/2)$, 即 $x\le1$ 时 $x^2e^{-2x}/2$, $x>1$ 时 $(x-1/2)e^{-2x}$.
\end{solution}
\begin{exercise}\label{cw:bank-2H-4}
原题2H-4. 设 $\phi(x)=\int_0^x(2x+3t)^2dt$, 分别 (a) 用 Leibniz 公式; (b) 先积分再求导, 计算 $\phi\prime$.
\end{exercise}
\begin{solution}
官方答案译审. (a) $\phi\prime=25x^2+\int_0^x4(2x+3t)dt=25x^2+14x^2=39x^2$. (b) 积分得 $\phi=[(2x+3t)^3/9]_0^x=13x^3$, 求导仍为 $39x^2$.
\end{solution}
\begin{exercise}\label{cw:bank-2H-5}
原题2H-5. 用 Leibniz 公式直接验证下列给定解. (a) 原文方程为 $y\prime\prime+y=f(x)$, $y(0)=y\prime(0)=0$, 给定 $y_p=\frac1k\int_0^x\sin k(x-t)f(t)dt$. (b) $y\prime\prime-2ky\prime+k^2y=f(x)$, 相同零初值, 给定 $y_p=\int_0^x(x-t)e^{k(x-t)}f(t)dt$.
\end{exercise}
\begin{solution}
编者补解并说明原文勘误. 编者核对发现(a)原文的方程与核不一致. 对 $k\ne0$, 该核导数满足 $w(0)=0,w\prime(0)=1,w\prime\prime=-k^2w$, 所以 Leibniz 给 $y_p\prime=\int_0^x\cos k(x-t)f(t)dt$, $y_p\prime\prime=f(x)-k^2y_p$. 因而应验证的是 $y\prime\prime+k^2y=f$; 原文只在 $k^2=1$ 时成立. $k=0$ 取核 $x-t$. (b) 令 $w(z)=ze^{kz}$, 则 $w(0)=0,w\prime(0)=1$ 且 $w\prime\prime-2kw\prime+k^2w=0$. 两次 Leibniz 求导产生端点项 $f(x)$, 即得方程; 两个零初值也由积分上限为0得到.
\end{solution}
\begin{exercise}\label{cw:bank-2H-6}
原题2H-6. 求显式卷积: (a) $e^{ax}*e^{ax}$; (b) $1*x$; (c) $x*x^2$.
\end{exercise}
\begin{solution}
编者补解. 按 $(f*g)(x)=\int_0^xf(t)g(x-t)dt$: (a) $\int_0^xe^{at}e^{a(x-t)}dt=xe^{ax}$; (b) $\int_0^x(x-t)dt=x^2/2$; (c) $\int_0^xt(x-t)^2dt=x^4(1/2-2/3+1/4)=x^4/12$.
\end{solution}
\begin{exercise}\label{cw:bank-2H-7}
原题2H-7. 求 Notes I例7并说明理由. 例7: 银行以恒定连续利率 $r$ 计息, $t=0$ 存入 $A_0$ 到 $t$ 时增长为 $A_0e^{rt}$. 从第0日起, 一位街头艺人以速率 $d(t)$ 存入每日收入, 不取款. 给出 $t=x$ 时账户金额的近似卷积表达式.
\end{exercise}
\begin{solution}
编者补解; 题面由官方Notes I第7页例7展开. 时间 $t$ 的小额存款 $d(t)dt$ 到 $x$ 增值为 $d(t)e^{r(x-t)}dt$. 相加并取连续极限, 得 $M(x)=\int_0^xd(t)e^{r(x-t)}dt=(d*e^{r\cdot})(x)$. 若另有初始余额 $M_0$, 再加 $M_0e^{rx}$. 离散每日存款时相应为求和; 积分是按连续存款率作的近似.
\end{solution}
\begin{exercise}\label{cw:bank-2H-8}
原题2H-8. 证明 $y\prime+ay=r(x)$, $y(0)=0$ 的解是 $y_p=e^{-ax}*r(x)$, 分别 (a) 用 Leibniz 公式; (b) 用一阶积分因子法与定积分.
\end{exercise}
\begin{solution}
编者补解. (a) $y=\int_0^xe^{-a(x-t)}r(t)dt$, 故 $y\prime=r(x)-ay$, 且 $y(0)=0$. (b) $(e^{ax}y)\prime=e^{ax}r(x)$, 从0积分后乘 $e^{-ax}$, 得同一式.
\end{solution}
\begin{exercise}\label{cw:bank-2H-9}
原题2H-9. 设 $y_c=c_1u(x)+c_2v(x)$ 为 $y\prime\prime+p(x)y\prime+q(x)y=0$ 的通解. 以 $W(t)=u(t)v\prime(t)-u\prime(t)v(t)$, $g(x,t)=[u(t)v(x)-v(t)u(x)]/W(t)$ 定义 $y(x)=\int_0^xg(x,t)f(t)dt$. 用 Leibniz 公式证明它解 $y\prime\prime+py\prime+qy=f$, $y(0)=y\prime(0)=0$.
\end{exercise}
\begin{solution}
编者补解. 在正则区间内 $W\ne0$. 对固定 $t$, $g$ 作为 $x$ 的函数满足齐次方程, 且 $g(t,t)=0$, $g_x(t,t)=1$. 因而 $y\prime=\int_0^xg_x(x,t)f(t)dt$, $y\prime\prime=f(x)+\int_0^xg_{xx}(x,t)f(t)dt$. 代入后积分项因齐次关系消失, 得右端 $f(x)$; 积分上限为0给两个初值.
\end{solution}
\originalchapter{题库三: Laplace 变换}
\originalsection{3A: 基本性质与公式}
\begin{exercise}\label{cw:bank-3A-1}
原题3A-1. 由定义证明 $\Lap(t)=1/s^2$, $s>0$.
\end{exercise}
\begin{solution}
官方答案译审. $\int_0^\infty te^{-st}dt=[-te^{-st}/s]_0^\infty+\frac1s\int_0^\infty e^{-st}dt=1/s^2$. $s>0$ 保证端点项消失且积分收敛.
\end{solution}
\begin{exercise}\label{cw:bank-3A-2}
原题3A-2. 假设 $\Lap(e^{\alpha t})=1/(s-\alpha)$ 对复数 $\alpha$ 也成立, 并用 $\Lap(u+iv)=\Lap u+i\Lap v$, 推导 $\Lap(e^{at}\cos bt)$ 与 $\Lap(e^{at}\sin bt)$.
\end{exercise}
\begin{solution}
官方答案译审. $\Lap(e^{(a+ib)t})=1/(s-a-ib)=[(s-a)+ib]/[(s-a)^2+b^2]$. 比较实虚部得到 $(s-a)/[(s-a)^2+b^2]$ 和 $b/[(s-a)^2+b^2]$, $s>a$.
\end{solution}
\begin{exercise}\label{cw:bank-3A-3}
原题3A-3. 求逆变换, (c),(e)使用部分分式: (a) $1/(s/2+3)$; (b) $3/(s^2+4)$; (c) $1/(s^2-4)$; (d) $(1+2s)/s^3$; (e) $1/(s^4-9s^2)$.
\end{exercise}
\begin{solution}
官方答案译审. (a) $1/(s/2+3)=2/(s+6)$, 故 $2e^{-6t}$. (b) $3\sin2t/2$. (c) $[1/(s-2)-1/(s+2)]/4$ 给 $\sinh2t/2$. (d) $1/s^3+2/s^2$ 给 $t^2/2+2t$. (e) 分解为 $-1/(9s^2)+1/[54(s-3)]-1/[54(s+3)]$, 给 $-t/9+\sinh3t/27$.
\end{solution}
\begin{exercise}\label{cw:bank-3A-4}
原题3A-4. 由变换定义、已知 $\Lap(\cos at)$ 和分部积分, 推导 $\Lap(\sin at)$.
\end{exercise}
\begin{solution}
官方答案译审. $I=\int_0^\infty e^{-st}\sin at\,dt$ 分部积分为 $I=(a/s)\int_0^\infty e^{-st}\cos at\,dt=(a/s)s/(s^2+a^2)=a/(s^2+a^2)$, $s>0$.
\end{solution}
\begin{exercise}\label{cw:bank-3A-5}
原题3A-5. (a) 用三角恒等式求 $\Lap(\cos^2at),\Lap(\sin^2at)$. (b) 相加核对结果, 说明应等于什么.
\end{exercise}
\begin{solution}
官方答案译审. 由平方降幂公式, 分别为 $1/(2s)+s/[2(s^2+4a^2)]$ 与 $1/(2s)-s/[2(s^2+4a^2)]$. 相加为 $1/s=\Lap(1)$, 与 $\cos^2+\sin^2=1$ 一致.
\end{solution}
\begin{exercise}\label{cw:bank-3A-6}
原题3A-6. (a) 用 $\int_0^\infty e^{-x^2}dx=\sqrt\pi/2$ 证明 $\Lap(1/\sqrt t)=\sqrt{\pi/s}$, $s>0$. 提示: 积分中作变量替换, 视 $s$ 为常数. (b) 推出 $\Lap(\sqrt t)=\sqrt\pi/(2s^{3/2})$.
\end{exercise}
\begin{solution}
官方答案译审. (a) 令 $x=\sqrt{st}$, $dt=2x\,dx/s$, 积分变为 $2s^{-1/2}\int_0^\infty e^{-x^2}dx=\sqrt\pi/\sqrt s$. (b) 分部积分给 $\int_0^\infty e^{-st}\sqrt t\,dt=(2s)^{-1}\int_0^\infty e^{-st}t^{-1/2}dt$, 得所需式. 边界项 $e^{-st}\sqrt t$ 两端均趋零.
\end{solution}
\begin{exercise}\label{cw:bank-3A-7}
原题3A-7. 证明 $\Lap(e^{t^2})$ 不在任何形如 $s>a$ 的区间存在, 即积分对每个实数 $s$ 都发散.
\end{exercise}
\begin{solution}
官方答案译审. 当 $t\ge\max(0,s)$ 时 $t^2-st\ge0$, 因而 $e^{t^2-st}\ge1$. 尾积分下界为无限长区间上的1的积分, 必发散.
\end{solution}
\begin{exercise}\label{cw:bank-3A-8}
原题3A-8. 哪些实数 $k$ 使 $\Lap(1/t^k)$ 存在? 从定义分析收敛性.
\end{exercise}
\begin{solution}
官方答案译审. 对 $s>0$, 无穷端指数衰减克服任意幂, 只须检查0端 $\int_0^1t^{-k}dt$, 收敛当且仅当 $k<1$. 所以恰为 $k<1$.
\end{solution}
\begin{exercise}\label{cw:bank-3A-9}
原题3A-9. 查变换公式求: (a) $\Lap(e^{-t}\sin3t)$; (b) $\Lap(e^{2t}(t^2-3t+2))$.
\end{exercise}
\begin{solution}
官方答案译审. 用指数移位, (a) $3/[(s+1)^2+9]$, $s>-1$; (b) $2/(s-2)^3-3/(s-2)^2+2/(s-2)$, $s>2$.
\end{solution}
\begin{exercise}\label{cw:bank-3A-10}
原题3A-10. 求逆变换: (a) $3/(s-2)^4$; (b) $1/[s(s-2)]$; (c) $(s+1)/(s^2-4s+5)$.
\end{exercise}
\begin{solution}
官方答案译审. (a) 用 $\Lap(t^3)=6/s^4$ 得 $t^3e^{2t}/2$. (b) $\frac12[1/(s-2)-1/s]$ 给 $(e^{2t}-1)/2$. (c) 分母为 $(s-2)^2+1$, 分子 $(s-2)+3$, 得 $e^{2t}(\cos t+3\sin t)$.
\end{solution}
\originalsection{3B: 导数公式与初值问题}
\begin{exercise}\label{cw:bank-3B-1}
原题3B-1. 用 Laplace 变换解: (a) $y\prime-y=e^{3t}$, $y(0)=1$; (b) $y\prime\prime-3y\prime+2y=0$, $y(0)=y\prime(0)=1$; (c) $y\prime\prime+4y=\sin t$, $y(0)=1,y\prime(0)=0$; (d) $y\prime\prime-2y\prime+2y=2e^t$, $y(0)=0,y\prime(0)=1$; (e) $y\prime\prime-2y\prime+y=e^t$, $y(0)=1,y\prime(0)=0$.
\end{exercise}
\begin{solution}
官方答案译审. 记 $Y=\Lap y$. 导数公式给\begin{align*}(a)\quad &(s-1)Y=1+1/(s-3),&Y&=\tfrac1{2(s-1)}+\tfrac1{2(s-3)};\\(b)\quad &(s^2-3s+2)Y=s-2,&Y&=1/(s-1);\\(c)\quad &(s^2+4)Y=s+1/(s^2+1),&Y&=\frac{s}{s^2+4}+\frac1{3(s^2+1)}-\frac1{3(s^2+4)};\\(d)\quad &[(s-1)^2+1]Y=1+2/(s-1),&Y&=\frac2{s-1}-\frac{2(s-1)}{(s-1)^2+1}+\frac1{(s-1)^2+1};\\(e)\quad &(s-1)^2Y=s-2+1/(s-1),&Y&=\frac1{s-1}-\frac1{(s-1)^2}+\frac1{(s-1)^3}.\end{align*}依次逆变换得 (a) $(e^t+e^{3t})/2$; (b) $e^t$; (c) $\cos2t+\sin t/3-\sin2t/6$; (d) $e^t(2-2\cos t+\sin t)$; (e) $e^t(1-t+t^2/2)$.
\end{solution}
\begin{exercise}\label{cw:bank-3B-2}
原题3B-2. 不查书或笔记, 由定义推导 $\Lap(f\prime)$ 与 $\Lap f$ 的关系, 说明对 $f,f\prime$ 的假设.
\end{exercise}
\begin{solution}
官方答案译审. 可假设 $f$ 连续且分段 $C^1$, $f\prime$ 分段连续, 二者在无穷端为指数阶. 对充分大的 $s$, 分部积分给 $\int_0^\infty e^{-st}f\prime(t)dt=[e^{-st}f(t)]_0^\infty+s\int_0^\infty e^{-st}f(t)dt=sF(s)-f(0+)$. 分段点处 $f$ 连续, 内部端点项相消; 无穷端指数阶保证边界项趋零.
\end{solution}
\begin{exercise}\label{cw:bank-3B-3}
原题3B-3. 用公式和变换表求: (a) $\Lap(t\cos bt)$; (b) $\Lap(t^ne^{kt})$, 给两种方法; (c) $\Lap(e^{at}t\sin t)$.
\end{exercise}
\begin{solution}
官方答案译审并补正. (a) 原扫描答案分子写为 $b^2-s^2$, 负号有误; 由 $-D_s[s/(s^2+b^2)]=(s^2-b^2)/(s^2+b^2)^2$. (b) $n$ 为非负整数, 指数移位给 $n!/(s-k)^{n+1}$; 另一方法是 $(-D_s)^n[1/(s-k)]$, 给同值. (c) 对 $1/(s^2+1)$ 求负导数后移位, 得 $2(s-a)/[(s-a)^2+1]^2$.
\end{solution}
\begin{exercise}\label{cw:bank-3B-4}
原题3B-4. 求逆变换: (a) $s/(s^2+1)^2$; (b) $1/(s^2+1)^2$.
\end{exercise}
\begin{solution}
官方答案译审. (a) 因 $\Lap(t\sin t)=2s/(s^2+1)^2$, 得 $t\sin t/2$. (b) $\Lap(t\cos t)=(s^2-1)/(s^2+1)^2$, 与 $\Lap(\sin t)=1/(s^2+1)$ 相减, 得逆变换 $(\sin t-t\cos t)/2$.
\end{solution}
\begin{exercise}\label{cw:bank-3B-5}
原题3B-5. 不查书或笔记, 推导 (a) $\Lap(e^{at}f(t))=F(s-a)$; (b) $\Lap(tf(t))=-F\prime(s)$.
\end{exercise}
\begin{solution}
官方答案译审. (a) 定义积分的指数合并为 $e^{-(s-a)t}$ 即得. (b) 对定义积分在 $s$ 下求导: $F\prime(s)=\int_0^\infty(-t)e^{-st}f(t)dt$. 若 $f$ 为指数阶且分段连续, 在收敛半平面内取一个较小的闭邻域, 可用 $Ct e^{-\eta t}$ 控制导数尾部, 因而允许积分下求导.
\end{solution}
\begin{exercise}\label{cw:bank-3B-6}
原题3B-6. 若 $y\prime\prime+ty=0$, $y(0)=1,y\prime(0)=0$, $Y(s)=\Lap y$, 求 $Y$ 满足的微分方程. 此 $y$ 是一种 Airy 函数.
\end{exercise}
\begin{solution}
官方答案译审. $\Lap(y\prime\prime)=s^2Y-s$, $\Lap(ty)=-Y\prime$, 所以 $Y\prime-s^2Y=-s$. 原题只要求这个关于 $s$ 的方程; 作为 Laplace 变换还须具有相应无穷端行为 $Y(s)\sim1/s$.
\end{solution}
\originalsection{3C: 不连续函数}
\begin{exercise}\label{cw:bank-3C-1}
原题3C-1. 尽量利用已知函数及变换表求各函数的 Laplace 变换, 尽可能用单位阶跃 $u(t)$, 并画图: (a) $f=1$ ($0\le t\le1$), $f=-1$ ($1<t\le2$), 其余为0; (b) $f=t$ ($0\le t\le1$), $f=2-t$ ($1\le t\le2$), 其余为0; (c) $f=|\sin t|$, $t\ge0$.
\end{exercise}
\begin{solution}
官方答案译审. (a) $f=1-2u(t-1)+u(t-2)$, $F=(1-2e^{-s}+e^{-2s})/s$. (b) $f=t-2(t-1)u(t-1)+(t-2)u(t-2)$, $F=(1-2e^{-s}+e^{-2s})/s^2$. (c) 周期 $\pi$, 分段积分后求几何级数给 $F=(1+e^{-\pi s})/[(s^2+1)(1-e^{-\pi s})]$, $s>0$.\begin{center}\begin{tikzpicture}\begin{axis}[width=4.4cm,height=3.8cm,axis lines=middle,xmin=0,xmax=3,ymin=-1.2,ymax=1.2,title={(a)},xlabel=$t$]\addplot[main] coordinates {(0,1)(1,1)};\addplot[main] coordinates {(1,-1)(2,-1)};\addplot[main] coordinates {(2,0)(3,0)};\end{axis}\end{tikzpicture}\begin{tikzpicture}\begin{axis}[width=4.4cm,height=3.8cm,axis lines=middle,xmin=0,xmax=3,ymin=0,ymax=1.2,title={(b)},xlabel=$t$]\addplot[main] coordinates {(0,0)(1,1)(2,0)(3,0)};\end{axis}\end{tikzpicture}\begin{tikzpicture}\begin{axis}[width=4.4cm,height=3.8cm,axis lines=middle,xmin=0,xmax=7,ymin=0,ymax=1.2,title={(c)},xlabel=$t$]\addplot[main,domain=0:7,samples=160] {abs(sin(deg(x)))};\end{axis}\end{tikzpicture}\end{center}变换不受有限个跳点处单独取值影响.
\end{solution}
\begin{exercise}\label{cw:bank-3C-2}
原题3C-2. 求逆变换: (a) $e^{-s}/(s^2+3s+2)$; (b) $(e^{-s}-e^{-3s})/s$, 并画(b)的图.
\end{exercise}
\begin{solution}
官方答案译审. (a) 无移位部分分解为 $1/(s+1)-1/(s+2)$, 所以为 $u(t-1)[e^{-(t-1)}-e^{-2(t-1)}]$. (b) $u(t-1)-u(t-3)$, 即1至3之间高度1的矩形脉冲.\begin{center}\begin{tikzpicture}\begin{axis}[width=6cm,height=3.5cm,axis lines=middle,xmin=0,xmax=4,ymin=0,ymax=1.3,xlabel=$t$]\addplot[main] coordinates {(0,0)(1,0)};\addplot[main] coordinates {(1,1)(3,1)};\addplot[main] coordinates {(3,0)(4,0)};\end{axis}\end{tikzpicture}\end{center}
\end{solution}
\begin{exercise}\label{cw:bank-3C-3}
原题3C-3. 方波 $f(t)=1$ 在 $2n\le t\le2n+1$ ($n=0,1,\ldots$), 其余为0. 分别 (a) 直接从定义; (b) 写为无限级数, 逐项变换再求和, 求 $\Lap f$.
\end{exercise}
\begin{solution}
官方答案译审. (a) $F=\sum_{n\ge0}\int_{2n}^{2n+1}e^{-st}dt=\frac{1-e^{-s}}s\sum_{n\ge0}e^{-2ns}=1/[s(1+e^{-s})]$, $s>0$. (b) $f=\sum_{n\ge0}[u(t-2n)-u(t-2n-1)]$. 每个方括号非负且有界, 可以对 $s>0$ 用单调收敛逐项积分; 变换后得到(a)中的同一几何级数.
\end{solution}
\begin{exercise}\label{cw:bank-3C-4}
原题3C-4. 用 Laplace 变换解 $y\prime\prime+2y\prime+2y=h(t)$, $y(0)=0,y\prime(0)=1$, 其中 $h=1$ ($\pi\le t\le2\pi$), 其余为0. 用分段形式写解.
\end{exercise}
\begin{solution}
官方答案译审. $h=u(t-\pi)-u(t-2\pi)$, 故 $Y=[1+(e^{-\pi s}-e^{-2\pi s})/s]/[(s+1)^2+1]$. 记阶跃响应 $g(v)=[1-e^{-v}(\cos v+\sin v)]/2$. 逆变换并分段为\[y(t)=\begin{cases}e^{-t}\sin t,&0\le t<\pi,\\e^{-t}\sin t+g(t-\pi),&\pi\le t<2\pi,\\e^{-t}\sin t+g(t-\pi)-g(t-2\pi),&t\ge2\pi.\end{cases}\]此处每段端点的 $y,y\prime$ 连续.
\end{solution}
\begin{exercise}\label{cw:bank-3C-5}
原题3C-5. 解 $y\prime\prime-3y\prime+2y=u(t)t$, $y(0)=1,y\prime(0)=0$, 右端为斜坡函数.
\end{exercise}
\begin{solution}
官方答案译审. $Y=(s-3)/[(s-1)(s-2)]+1/[s^2(s-1)(s-2)]$. 部分分式为 $Y=1/(s-1)-3/[4(s-2)]+3/(4s)+1/(2s^2)$, 所以 $y=e^t-3e^{2t}/4+3/4+t/2$, $t\ge0$.
\end{solution}
\originalsection{3D: 卷积与 Dirac 脉冲}
\begin{exercise}\label{cw:bank-3D-1}
原题3D-1. 解 $y\prime\prime+2y\prime+y=\delta(t)+u(t-1)$, $y(0)=0,y\prime(0)=1$, 并按 $0\le t\le1$ 与 $t>1$ 分段写解.
\end{exercise}
\begin{solution}
官方答案译审. 原文初导数写 $y\prime(0)=1$; 在原点另有单位脉冲的单侧变换约定下, 将该初值解释为脉冲前 $y\prime(0-)=1$. 原点单位脉冲使 $y\prime(0+)=2$. 因而 $(s+1)^2Y=2+e^{-s}/s$. 逆变换给 $2te^{-t}+u(t-1)[1-te^{-(t-1)}]$, 即\[y=\begin{cases}2te^{-t},&0\le t\le1,\\2te^{-t}+1-te^{-(t-1)},&t>1.\end{cases}\]原扫描在将后段合并为 $1+(2-e)te^{-t}$ 时指数常数辨认易错, 此式保留移位形式便于直接核对.
\end{solution}
\begin{exercise}\label{cw:bank-3D-2}
原题3D-2. 解 $y\prime\prime+y=r(t)$, $y(0)=0,y\prime(0)=1$, $r=1$ ($0\le t\le\pi$), 其余为0; 用分段形式写解.
\end{exercise}
\begin{solution}
官方答案译审. $r=1-u(t-\pi)$, 所以 $Y=1/(s^2+1)+(1-e^{-\pi s})/[s(s^2+1)]$. 逆变换为 $\sin t+1-\cos t-u(t-\pi)[1-\cos(t-\pi)]$, 即 $0\le t\le\pi$ 时 $1+\sin t-\cos t$, $t>\pi$ 时 $\sin t-2\cos t$.
\end{solution}
\begin{exercise}\label{cw:bank-3D-3}
原题3D-3. 若 $f(t+c)=f(t)$, $c>0$, 称其周期为 $c$. (a) 证明 $F(s)=\frac1{1-e^{-cs}}\int_0^ce^{-st}f(t)dt$. (b) 用该式重做3C-3.
\end{exercise}
\begin{solution}
官方答案译审. 对分段连续周期函数, $s>0$ 时绝对收敛. 按 $[nc,(n+1)c]$ 分割, 令 $t=nc+v$, 周期性给 $F=\sum_{n\ge0}e^{-ncs}\int_0^ce^{-sv}f(v)dv$, 求几何级数即得. (b) $c=2$, 单周期积分为 $(1-e^{-s})/s$, 故 $F=1/[s(1+e^{-s})]$.
\end{solution}
\begin{exercise}\label{cw:bank-3D-4}
原题3D-4. 用卷积求逆变换, 最终答案不得保留卷积符号: (a) $s/[(s+1)(s^2+4)]$; (b) $1/(s^2+1)^2$.
\end{exercise}
\begin{solution}
官方答案译审. (a) 为 $e^{-t}*\cos2t$, 即 $e^{-t}\int_0^te^u\cos2u\,du=(\cos2t+2\sin2t-e^{-t})/5$. (b) 为 $\sin t*\sin t$, 用积化和差 $\sin(t-u)\sin u=[\cos(t-2u)-\cos t]/2$, 积分得 $(\sin t-t\cos t)/2$.
\end{solution}
\begin{exercise}\label{cw:bank-3D-5}
原题3D-5. 假设 $f(t)=0$ ($t\le0$). 先由卷积形式非正式说明 $\delta*f=f$, 再用 $\delta$ 的定义证明. 原课本把它记为 $\delta_0$.
\end{exercise}
\begin{solution}
官方答案译审. 形式上 $\int\delta(u)f(t-u)du$ 只取 $u=0$ 的值, 所以等于 $f(t)$. 严格说 $\delta$ 是分布: 对测试函数 $\varphi$, $\langle\delta,\varphi\rangle=\varphi(0)$. 将 $f$ 作为因果函数延拓, 分布卷积满足 $\langle\delta*f,\varphi\rangle=\langle f,\varphi\rangle$, 因为 $\delta$ 的作用将平移参数固定为0. 若以单位面积窄脉冲逼近, 在 $f$ 的连续点也得到上述积分极限; 跳点处不应凭普通函数单点值规定分布等式.
\end{solution}
\begin{exercise}\label{cw:bank-3D-6}
原题3D-6. 直接在卷积积分中换元, 证明 $f*g=g*f$.
\end{exercise}
\begin{solution}
官方答案译审. $(f*g)(t)=\int_0^tf(u)g(t-u)du$. 令 $v=t-u$, 反转上下限后为 $\int_0^tg(v)f(t-v)dv=(g*f)(t)$.
\end{solution}
\begin{exercise}\label{cw:bank-3D-7}
原题3D-7. 用 Laplace 变换和卷积证明 $y\prime\prime+k^2y=r(t)$, $y(0)=y\prime(0)=0$ 的解为 $y(t)=\frac1k\int_0^tr(u)\sin k(t-u)du$.
\end{exercise}
\begin{solution}
官方答案译审. $Y=R/(s^2+k^2)$, 而 $\Lap^{-1}[1/(s^2+k^2)]=\sin kt/k$, 卷积定理给所需式. $k=0$ 时连续极限为 $\int_0^t(t-u)r(u)du$.
\end{solution}
\begin{exercise}\label{cw:bank-3D-8}
原题3D-8. 用 Laplace 变换和卷积证明 $y\prime\prime+ay\prime+by=r(t)$, $y(0)=y\prime(0)=0$ 的解为 $y(t)=\int_0^tw(t-u)r(u)du$, 其中 $w$ 满足齐次方程及 $w(0)=0,w\prime(0)=1$. 核 $w(t-u)$ 为算子 $D^2+aD+b$ 的 Green 函数.
\end{exercise}
\begin{solution}
官方答案译审. 零初值给 $Y=R/(s^2+as+b)$. 对 $w$ 的初值问题变换, 得 $(s^2+as+b)W=1$, 所以 $Y=WR$. 卷积定理遂得积分公式. $w$ 是因果单位脉冲响应; 其导数初值为1正是脉冲造成的单位跳跃.
\end{solution}

% MIT 18.03 Spring 2010, Mattuck Notes and Exercises 4–7.
% 题面与手写解答逐页核对；所有原题号均保留，不将相似题删为重复。
\originalchapter{原题题库四：线性系统}
\originalsection{4A：矩阵复习}
\begin{exercise}\label{cw:bank-4A-1}原题4A-1. 验证
\[\begin{pmatrix}2&0&1\\1&-1&-2\end{pmatrix}\begin{pmatrix}1&0&-1\\0&2&1\\-1&0&2\end{pmatrix}=\begin{pmatrix}1&0&0\\3&-2&-6\end{pmatrix}.\]
\end{exercise}
\begin{solution}（官方译审）矩阵乘积的第 $i,j$ 项为第 $i$ 行与第 $j$ 列的内积. 第一行给出 $2-1=1$, $0$, $-2+2=0$；第二行给出 $1+2=3$, $-2$, $-1-1-4=-6$. 六个元素与右端逐一相同.\end{solution}
\begin{exercise}\label{cw:bank-4A-2}原题4A-2. 设 $A=\begin{pmatrix}1&2\\3&-1\end{pmatrix}$, $B=\begin{pmatrix}0&-1\\2&1\end{pmatrix}$，证明 $AB\ne BA$.\end{exercise}
\begin{solution}（官方译审）直接按行列相乘得到 $AB=\begin{pmatrix}4&1\\-2&-4\end{pmatrix}$, $BA=\begin{pmatrix}-3&1\\5&3\end{pmatrix}$. 左上元素分别为 $4$ 与 $-3$，因此两乘积不同.\end{solution}
\begin{exercise}\label{cw:bank-4A-3}原题4A-3. 求 $A=\begin{pmatrix}2&2\\3&2\end{pmatrix}$ 的逆，并分别验证 $AA^{-1}=I$ 和 $A^{-1}A=I$.\end{exercise}
\begin{solution}（官方译审）$\det A=-2$，故 $A^{-1}=\begin{pmatrix}-1&1\\3/2&-1\end{pmatrix}$. 两个次序的乘积分别为
\[AA^{-1}=\begin{pmatrix}-2+3&2-2\\-3+3&3-2\end{pmatrix}=I,
\quad A^{-1}A=\begin{pmatrix}-2+3&-2+2\\3-3&3-2\end{pmatrix}=I.\]
\end{solution}
\begin{exercise}\label{cw:bank-4A-4}原题4A-4. 验证 Notes LS.1 给出的二阶矩阵求逆公式：当 $ad-bc\ne0$ 时，$\begin{pmatrix}a&b\\c&d\end{pmatrix}^{-1}=(ad-bc)^{-1}\begin{pmatrix}d&-b\\-c&a\end{pmatrix}$.\end{exercise}
\begin{solution}（官方译审）在两个次序下，未除以行列式的乘积均为 $\begin{pmatrix}ad-bc&0\\0&ad-bc\end{pmatrix}$. 因 $ad-bc\ne0$，除后均得 $I$，恰好满足逆矩阵定义. 行列式为零时不能使用该公式，矩阵也不可逆.\end{solution}
\begin{exercise}\label{cw:bank-4A-5}原题4A-5. 设 $A=\begin{pmatrix}0&1\\1&1\end{pmatrix}$，求 $A^3=A\cdot A\cdot A$.\end{exercise}
\begin{solution}（官方译审）先算 $A^2=\begin{pmatrix}1&1\\1&2\end{pmatrix}$，再乘一次得到 $A^3=\begin{pmatrix}1&2\\2&3\end{pmatrix}$.\end{solution}
\begin{exercise}\label{cw:bank-4A-6}原题4A-6. 什么 $c$ 值使 $\mathbf x_1=(1,2,c)$, $\mathbf x_2=(-1,0,1)$, $\mathbf x_3=(2,3,0)$ 线性相关？对该值用试探或其他方法找出 $c_1\mathbf x_1+c_2\mathbf x_2+c_3\mathbf x_3=0$ 的非平凡关系.\end{exercise}
\begin{solution}（官方译审）以三向量为列的行列式为 $1-3c$，故恰在 $c=1/3$ 时相关. 此时 $3\mathbf x_1-\mathbf x_2-2\mathbf x_3=(0,0,0)$，三系数并非全零.\end{solution}
\originalsection{4B：一阶系统、消元与矩阵表示}
\begin{exercise}\label{cw:bank-4B-1}原题4B-1. 把下列方程写成等价的一阶系统：(a) $x''+5x'+tx^2=0$；(b) $y''-x^2y'+(1-x^2)y=\sin x$.\end{exercise}
\begin{solution}（官方译审）(a) 令 $v=x'$，得 $x'=v$, $v'=-5v-tx^2$. (b) 此处自变量为 $x$，令 $z=y'$，得 $y'=z$, $z'=x^2z+(x^2-1)y+\sin x$. 从系统第一式再求导便恢复原式；反之每个原式解与其导数构成系统解.\end{solution}
\begin{exercise}\label{cw:bank-4B-2}原题4B-2. 把初值问题 $y^{(3)}+p(t)y''+q(t)y'+r(t)y=0$, $y(0)=y_0$, $y'(0)=y'_0$, $y''(0)=y''_0$ 写成三条一阶线性方程及矩阵初值问题.\end{exercise}
\begin{solution}（官方译审）令 $u_1=y$, $u_2=y'$, $u_3=y''$. 则
\[u_1'=u_2,\quad u_2'=u_3,\quad u_3'=-ru_1-qu_2-pu_3,
\qquad \mathbf u'=\begin{pmatrix}0&1&0\\0&0&1\\-r&-q&-p\end{pmatrix}\mathbf u,
\quad\mathbf u(0)=\begin{pmatrix}y_0\\y'_0\\y''_0\end{pmatrix}.\]
第一式、第二式保证三个变量确为同一函数的连续导数，故转换保持全部初值.\end{solution}
\begin{exercise}\label{cw:bank-4B-3}原题4B-3. 把 $\binom{x}{y}'=\begin{pmatrix}1&1\\4&1\end{pmatrix}\binom{x}{y}$ 展开为两条方程. (a) 消去 $y$ 得到 $x$ 的二阶方程；(b) 把二阶方程改为一阶系统，并找出连接新旧系统的变量变换.\end{exercise}
\begin{solution}（官方译审）原系统为 $x'=x+y$, $y'=4x+y$. (a) $y=x'-x$，代入第二式得 $x''-x'=4x+x'-x$，即 $x''-2x'-3x=0$. (b) 令 $u=x$, $v=x'$，得 $u'=v$, $v'=3u+2v$. 变换为 $u=x$, $v=x+y$，逆为 $x=u$, $y=v-u$；行列式为 $1$，故无丢解或增解.\end{solution}
\begin{exercise}\label{cw:bank-4B-4}原题4B-4. 对 $x'=4x-y$, $y'=2x+y$：(a) 用矩阵记号验证 $(e^{3t},e^{3t})$ 和 $(e^{2t},2e^{2t})$ 是解；(b) 验证两解构成基本解组；(c) 用两个任意常数写出向量形式及分量形式的通解.\end{exercise}
\begin{solution}（官方译审）令 $A=\begin{pmatrix}4&-1\\2&1\end{pmatrix}$. 有 $A\binom11=3\binom11$, $A\binom12=2\binom12$，这与两指数函数的导数相同. 两解为列的行列式为 $e^{5t}\ne0$，因此独立. 通解为
\[\mathbf x=c_1e^{3t}\binom11+c_2e^{2t}\binom12,
\quad x=c_1e^{3t}+c_2e^{2t},\quad y=c_1e^{3t}+2c_2e^{2t}.\]
\end{solution}
\begin{exercise}\label{cw:bank-4B-5}原题4B-5. 设 $A=\begin{pmatrix}1&3\\3&1\end{pmatrix}$. (a) 验证 $\binom11e^{4t}$ 和 $\binom1{-1}e^{-2t}$ 是独立的系统解；(b) 求 $\mathbf x'=A\mathbf x$, $\mathbf x(0)=\binom51$ 的解.\end{exercise}
\begin{solution}（官方译审）两个常向量分别满足 $Av_1=4v_1$, $Av_2=-2v_2$，且解矩阵行列式为 $-2e^{2t}\ne0$. 初值要求 $c_1+c_2=5$, $c_1-c_2=1$，所以
\[\mathbf x(t)=3e^{4t}\binom11+2e^{-2t}\binom1{-1}.\]
\end{solution}
\begin{exercise}\label{cw:bank-4B-6}原题4B-6. 用两种方法解 $\binom{x}{y}'=\begin{pmatrix}1&1\\0&1\end{pmatrix}\binom{x}{y}$：(a) 先解第二式，再代入第一式；(b) 由第一式消去 $y$ 得二阶方程，解后恢复 $y$. 两种结果是否相同？\end{exercise}
\begin{solution}（官方译审）(a) $y=C_2e^t$，第一式为 $x'-x=C_2e^t$，乘 $e^{-t}$ 得 $(e^{-t}x)'=C_2$，故 $x=(C_1+C_2t)e^t$. (b) $y=x'-x$，所以 $x''-2x'+x=0$，重根 $1$ 给出同一 $x$，再由 $y=x'-x$ 得 $C_2e^t$. 两方法完全一致.\end{solution}
\begin{exercise}\label{cw:bank-4B-7}原题4B-7. 放射性物质 $R$ 衰变为 $S$，$S$ 也继续衰变. 用 $x,y$ 表示时刻 $t$ 的两物质量. (a) 证明模型为 $\mathbf x'=\begin{pmatrix}-a&0\\a&-b\end{pmatrix}\mathbf x$，其中 $a,b$ 为常数；(b) 用4B-6的一种消元方法解系统；(c) 初始只有 $R$，数量为 $x_0$，求两物质量.\end{exercise}
\begin{solution}（官方思路展开）(a) $R$ 每单位时间减少 $ax$，同量转入 $S$；$S$ 再减少 $by$，故 $x'=-ax$, $y'=ax-by$，物理衰变率 $a,b\ge0$. (b) 先得 $x=Ce^{-at}$，继而 $(e^{bt}y)'=aCe^{(b-a)t}$. 当 $a\ne b$ 时
\[x=Ce^{-at},\qquad y=De^{-bt}+\frac{aC}{b-a}e^{-at}.\]
当 $a=b$ 时应单独积分，得 $y=(D+aCt)e^{-at}$. (c) 若 $x(0)=x_0$, $y(0)=0$，则
\[x=x_0e^{-at},\qquad y=\begin{cases}\dfrac{ax_0}{b-a}(e^{-at}-e^{-bt}),&a\ne b,\\ax_0te^{-at},&a=b.\end{cases}\]
这些式子包括 $a=0$ 的边界，且 $x_0\ge0$ 时两物质量非负. 消元若需对两式同时求导，可能产生额外解，须代回原系统；本题直接解第一式并代入第二式，没有该问题.\end{solution}
\originalsection{4C：特征值与特征向量}
\begin{exercise}\label{cw:bank-4C-1}原题4C-1. 分别求 $\mathbf x'=A\mathbf x$ 的通解，先求特征值、特征向量，再构造正规模态：
\[(a)\ A=\begin{pmatrix}-3&4\\-2&3\end{pmatrix},\quad
(b)\ A=\begin{pmatrix}4&-3\\8&-6\end{pmatrix},\quad
(c)\ A=\begin{pmatrix}1&-1&0\\1&2&1\\-2&1&-1\end{pmatrix}.\]
\end{exercise}
\begin{solution}（官方译审）(a) $\det(\lambda I-A)=\lambda^2-1$；$\lambda=1,-1$ 对应 $(1,1)^T,(2,1)^T$. 故 $\mathbf x=c_1e^t\binom11+c_2e^{-t}\binom21$. (b) 特征多项式为 $\lambda(\lambda+2)$，对应向量为 $(3,4)^T,(1,2)^T$，故 $\mathbf x=c_1\binom34+c_2e^{-2t}\binom12$. (c) 展开行列式得 $(\lambda-1)(\lambda-2)(\lambda+1)$，代入 $Av=\lambda v$ 分别可取
\[v_1=(1,0,-1)^T,\quad v_2=(1,-1,-1)^T,\quad v_3=(1,2,-7)^T.\]
于是 $\mathbf x=c_1v_1e^t+c_2v_2e^{2t}+c_3v_3e^{-t}$. 三个互异特征值的向量独立；或直接算其列行列式为 $6\ne0$，故模态给出全部解.\end{solution}
\begin{exercise}\label{cw:bank-4C-2}原题4C-2. 证明：$0$ 是 $n\times n$ 常矩阵 $A$ 的特征值，当且仅当 $A$ 奇异，即 $\det A=0$. 可用特征方程或定义.\end{exercise}
\begin{solution}（官方译审）$0$ 是特征值等价于存在 $v\ne0$ 满足 $Av=0$，即齐次线性方程有非零解，等价于 $A$ 不可逆，也等价于 $\det A=0$. 此论证不要求特征值彼此不同.\end{solution}
\begin{exercise}\label{cw:bank-4C-3}原题4C-3. 求上三角矩阵 $A=\begin{pmatrix}a&*&*\\0&b&*\\0&0&c\end{pmatrix}$ 的特征值（星号为任意数），并推广至 $n\times n$ 上三角矩阵.\end{exercise}
\begin{solution}（官方译审）$\det(\lambda I-A)=(\lambda-a)(\lambda-b)(\lambda-c)$，故特征值为 $a,b,c$，包括重复次数. 一般上三角矩阵的行列式为对角元乘积，因而特征多项式为 $\prod_{j=1}^n(\lambda-a_{jj})$，特征值恰为对角元. 重根不保证能找到相应数量的独立特征向量.\end{solution}
\begin{exercise}\label{cw:bank-4C-4}原题4C-4. 若 $v$ 是 $A$ 对应特征值 $m$ 的特征向量，证明它也是 $A^2$ 对应 $m^2$ 的特征向量. 提示：使用 $Av=mv$.\end{exercise}
\begin{solution}（官方译审）$A^2v=A(Av)=A(mv)=mAv=m^2v$，且原特征向量 $v\ne0$，因此符合定义. 即使 $m=0$ 也成立.\end{solution}
\begin{exercise}\label{cw:bank-4C-5}原题4C-5. 用特征值与特征向量解4B-7的放射性衰变问题：$x'=-ax$, $y'=ax-by$，初值 $x(0)=x_0$, $y(0)=0$.\end{exercise}
\begin{solution}（官方思路展开）若 $a\ne b$，特征值为 $-a,-b$，对应向量可取 $v_a=(b-a,a)^T$, $v_b=(0,1)^T$；直接相乘可验算. 因而 $\mathbf x=Cv_ae^{-at}+Dv_be^{-bt}$，初值给 $C=x_0/(b-a)$, $D=-ax_0/(b-a)$，即4B-7所列解. 若 $a=b>0$，矩阵为 $-aI+N$，$N=\begin{pmatrix}0&0\\a&0\end{pmatrix}$，$N^2=0$；于是 $\mathbf x=e^{-at}(I+tN)\binom{x_0}0=(x_0e^{-at},ax_0te^{-at})^T$. 若 $a=b=0$，系统恒定，解为 $(x_0,0)^T$. 重根情况不能强行套两个不同特征向量的公式.\end{solution}
\begin{exercise}\label{cw:bank-4C-6}原题4C-6. Smith和Jones各有一群兔子，每年Smith群的比例 $a$ 转移至Jones处，Jones群的比例 $b$ 转移至Smith处；假设每只兔子的自然增长率为每年1只. 用 $S,J$ 表示两群数量. (a) 建立微分方程；(b) $a=b=1/2$，初始 $S=20,J=10$，求数量随时间的变化；(c) 证明无论转移率及初值如何，兔群数量都不会振荡.\end{exercise}
\begin{solution}（官方译审）(a) $S'=(1-a)S+bJ$, $J'=aS+(1-b)J$. (b) $T=S+J$ 满足 $T'=T$，差 $D=S-J$ 在此参数下满足 $D'=0$. 所以 $T=30e^t$, $D=10$，即 $S=15e^t+5$, $J=15e^t-5$. (c) 特征多项式为 $(\lambda-1)(\lambda-(1-a-b))$，物理上 $a,b\ge0$；除 $a=b=0$ 的恒等矩阵情形外，它有两个实特征值. 每个分量为两个实指数的线性组合，其导数最多有一次变号（若某项系数或指数为零则更少），故不可能反复振荡. $a=b=0$ 时两个分量各自按 $e^t$ 增长，也无振荡.\end{solution}
\begin{exercise}\label{cw:bank-4C-7}原题4C-7. 反馈环中，黑箱 $B_1$ 输入 $x_1(t)$，输出 $(x'_1-x_1)/4$，此输出作为 $B_2$ 的输入 $x_2$；$B_2$ 输出 $x'_2-x_2$，回送为 $B_1$ 的输入 $x_1$. 若 $x_1(0)=1$, $x_2(0)=0$，求两信号随时间的变化.\end{exercise}
\begin{solution}（官方译审）接线给 $x'_1=x_1+4x_2$, $x'_2=x_1+x_2$. 特征多项式为 $(\lambda-1)^2-4$，根为 $3,-1$，对应 $(2,1)^T,(-2,1)^T$. 初值解出系数 $1/4,-1/4$，故
\[x_1=\frac12(e^{3t}+e^{-t}),\qquad x_2=\frac14(e^{3t}-e^{-t}).\]
\end{solution}
\originalsection{4D：复特征值与重特征值}
\begin{exercise}\label{cw:bank-4D-1}原题4D-1. 解 $\mathbf x'=\begin{pmatrix}1&-5\\1&-1\end{pmatrix}\mathbf x$.\end{exercise}
\begin{solution}（官方译审）特征多项式为 $\lambda^2+4$，根 $\pm2i$；对应 $2i$ 可取 $v=(1+2i,1)^T$. 取 $ve^{2it}$ 的实部、虚部得两实解，故
\[\mathbf x=c_1\binom{\cos2t-2\sin2t}{\cos2t}+c_2\binom{\sin2t+2\cos2t}{\sin2t}.\]
两列在 $t=0$ 的行列式为 $-2\ne0$，因此组成基本解组.\end{solution}
\begin{exercise}\label{cw:bank-4D-2}原题4D-2. 解 $\mathbf x'=\begin{pmatrix}3&-4\\4&3\end{pmatrix}\mathbf x$.\end{exercise}
\begin{solution}（官方译审）矩阵为 $3I+4J$，$J=\begin{pmatrix}0&-1\\1&0\end{pmatrix}$，特征值 $3\pm4i$. 由旋转解得到
\[\mathbf x=e^{3t}\left[c_1\binom{\cos4t}{\sin4t}+c_2\binom{-\sin4t}{\cos4t}\right].\]
旋转矩阵导数为 $4J$ 乘该矩阵，乘 $e^{3t}$ 后恰为原系统；其行列式为 $e^{6t}$，给出通解.\end{solution}
\begin{exercise}\label{cw:bank-4D-3}原题4D-3. 解 $\mathbf x'=\begin{pmatrix}2&3&3\\0&-1&-3\\0&0&2\end{pmatrix}\mathbf x$.\end{exercise}
\begin{solution}（官方译审）上三角矩阵特征值为 $-1,2,2$. 解 $(A+I)v=0$ 得 $v=(-1,1,0)^T$；解 $(A-2I)v=0$ 得条件 $v_2+v_3=0$，可取 $(1,0,0)^T,(0,1,-1)^T$. 因此
\[\mathbf x=c_1e^{-t}(-1,1,0)^T+c_2e^{2t}(1,0,0)^T+c_3e^{2t}(0,1,-1)^T.\]
三个向量独立，重根 $2$ 的特征空间维数为2，本题不需要 $te^{2t}$ 项.\end{solution}
\begin{exercise}\label{cw:bank-4D-4}原题4D-4. 三个完全相同的盐水室两两通过相同半透膜相隔. $A_i$ 为第 $i$ 室盐量，$x_i=A_i-A_0$ 为相对同一参考量的偏差. 膜间扩散速率与两侧浓度差成正比，比例常数取1. (a) 导出 $\mathbf x'=A\mathbf x$，$A=\begin{pmatrix}-2&1&1\\1&-2&1\\1&1&-2\end{pmatrix}$；(b) 找三个正规模态，解释其初始状态及物理合理性.\end{exercise}
\begin{solution}（官方译审）(a) 第一室净输入为 $(A_2-A_1)+(A_3-A_1)$，即 $x'_1=-2x_1+x_2+x_3$，其余循环置换，得到所给矩阵. (b) $v_0=(1,1,1)^T$ 对应特征值0；$v_1=(1,-1,0)^T$, $v_2=(0,1,-1)^T$ 对应 $-3$. 三模态是 $v_0$, $e^{-3t}v_1$, $e^{-3t}v_2$. 第一种三室盐量相同，不再扩散；后两种总偏差为零，一室高于参考量、一室低于参考量，差异指数消失. 模态系数需使实际初始盐量 $A_i\ge0$；偏差可正可负，无须把负偏差误当负盐量. 总盐量守恒，最终均匀.\end{solution}
\originalsection{4E：解耦}
\begin{exercise}\label{cw:bank-4E-1}原题4E-1. 对 $x'=4x+2y$, $y'=3x-y$，给出与 $x,y$ 线性相关的新变量 $u,v$ 使系统解耦，并用直接代入验证.\end{exercise}
\begin{solution}（官方译审）特征多项式 $\lambda^2-3\lambda-10=(\lambda-5)(\lambda+2)$，对应向量 $(2,1)^T,(1,-3)^T$. 令 $x=2u+v$, $y=u-3v$，逆为 $u=(3x+y)/7$, $v=(x-2y)/7$. 直接计算
\[u'=\frac{3(4x+2y)+(3x-y)}7=5u,\qquad
v'=\frac{4x+2y-2(3x-y)}7=-2v.\]
变换行列式为 $-7$，所以确为等价的解耦系统.\end{solution}
\begin{exercise}\label{cw:bank-4E-2}原题4E-2. 对4D-4的三室扩散系统回答4E-1同样的问题：使用正规模态给出解耦变量，最后用矩阵代入验证.\end{exercise}
\begin{solution}（官方译审）令 $P=\begin{pmatrix}1&1&0\\1&-1&1\\1&0&-1\end{pmatrix}$，写 $\mathbf x=P\mathbf u$. 三列是4D-4的特征向量，$\det P=3\ne0$，且 $AP=P\diag(0,-3,-3)$. 因而 $\mathbf u'=P^{-1}AP\mathbf u=\diag(0,-3,-3)\mathbf u$. 具体 $u_0=(x_1+x_2+x_3)/3$, $u_1=x_1-u_0$, $u_2=u_0-x_3$，分别满足 $u'_0=0$, $u'_1=-3u_1$, $u'_2=-3u_2$.\end{solution}
\originalsection{4F：线性系统理论}
\begin{exercise}\label{cw:bank-4F-1}原题4F-1. 对 $x''+p(t)x'+q(t)x=0$：(a) 按通常方法改为一阶系统；(b) 证明原方程两解的Wronskian与其对应系统解的Wronskian相同.\end{exercise}
\begin{solution}（官方译审）(a) 令 $\mathbf u=(x,x')^T$，得 $\mathbf u'=\begin{pmatrix}0&1\\-q&-p\end{pmatrix}\mathbf u$. (b) 两解对应 $\mathbf u_j=(x_j,x'_j)^T$，故 $\det[\mathbf u_1,\mathbf u_2]=x_1x'_2-x_2x'_1$，正是标量二阶方程的Wronskian定义.\end{solution}
\begin{exercise}\label{cw:bank-4F-2}原题4F-2. 给定 $\mathbf x_1=(t,1)^T$, $\mathbf x_2=(t^2,2t)^T$：(a) 用定义证明线性独立；(b) 计算Wronskian；(c) 与 Notes LS.5 定理5C所述连续线性系统的Wronskian非零判据如何相容？(d) 找出以这两函数为解的 $\mathbf x'=A(t)\mathbf x$，确认(c)的说明. 可将矩阵元素视作未知函数求解.\end{exercise}
\begin{solution}（官方译审）(a) 若 $c_1\mathbf x_1+c_2\mathbf x_2\equiv0$，第二分量为 $c_1+2c_2t\equiv0$，故 $c_1=c_2=0$. (b) $W=2t^2-t^2=t^2$. (c) 定理要求两函数是同一连续系数系统在一个区间上的解；一般向量函数并无“Wronskian要么恒零要么处处非零”的性质. (d) 设 $A=\begin{pmatrix}a&b\\c&d\end{pmatrix}$. 由两解分别得 $at+b=1$, $ct+d=0$, $at^2+2bt=2t$, $ct^2+2dt=2$. 在 $t\ne0$ 解为
\[A(t)=\begin{pmatrix}0&1\\-2/t^2&2/t\end{pmatrix}.\]
在不含0的区间上连续且 $W\ne0$，符合定理；无法延为包含0的连续系数矩阵. $t=0$ 时第二函数值为零而导数为 $(0,2)^T\ne0$，已直接排除了该延拓.\end{solution}
\begin{exercise}\label{cw:bank-4F-3}原题4F-3. 原题写“常二阶矩阵 $A$ 有两个不同的 eigenvectors $m_1,m_2$，对应特征向量 $v_1,v_2$”；此处 $m_1,m_2$ 应为两个不同的特征值. 证明 $\mathbf x_j=v_je^{m_jt}$ 独立：(a) 用行列式判据说明等价于 $v_1,v_2$ 独立；(b) 用 $Av_j=m_jv_j$ 证明 $c_1v_1+c_2v_2=0$ 蕴含两系数为零.\end{exercise}
\begin{solution}（官方思路展开）(a) $\det[\mathbf x_1,\mathbf x_2]=e^{(m_1+m_2)t}\det[v_1,v_2]$，指数非零，故两个判据等价. (b) 对关系乘 $A$ 得 $c_1m_1v_1+c_2m_2v_2=0$；减去原关系的 $m_1$ 倍得 $c_2(m_2-m_1)v_2=0$. 因 $m_2\ne m_1$, $v_2\ne0$，所以 $c_2=0$，再有 $c_1v_1=0$ 给 $c_1=0$.\end{solution}
\begin{exercise}\label{cw:bank-4F-4}原题4F-4. 设 $\mathbf x'=A\mathbf x$，其中 $A$ 是可逆常矩阵. 若某解在 $t_0$ 的速度向量为零，证明它恒为零. 此结论对 $A$ 的最低要求是什么？例如 $A$ 能否依赖于 $t$？\end{exercise}
\begin{solution}（官方译审）$0=\mathbf x'(t_0)=A\mathbf x(t_0)$ 与可逆性给 $\mathbf x(t_0)=0$. 齐次系统零解满足同一初值，唯一性得所论解恒零. 可以令 $A=A(t)$，只须所考虑连通区间上初值问题唯一，例如 $A$ 连续，并且 $A(t_0)$ 可逆；无需在其他时刻可逆. 对常矩阵，“对所有解均成立”要求 $\ker A=\{0\}$，否则非零常向量 $v\in\ker A$ 就是速度恒零的反例.\end{solution}
\originalsection{4G：基本矩阵}
\begin{exercise}\label{cw:bank-4G-1}原题4G-1. 已知两独立解 $v_1e^{3t}$, $v_2e^{2t}$，其中 $v_1=(1,1)^T$, $v_2=(1,2)^T$. (a) 分别求初值 $(0,1)^T$ 与 $(1,0)^T$ 的解；(b) 利用(a)方便地求初值 $(a,b)^T$ 的解.\end{exercise}
\begin{solution}（官方译审）通解为 $c_1v_1e^{3t}+c_2v_2e^{2t}$. (a) 解线性方程得 $\mathbf u=-v_1e^{3t}+v_2e^{2t}$，$\mathbf v=2v_1e^{3t}-v_2e^{2t}$，分别具有所指定初值. (b) 线性性给 $a\mathbf v+b\mathbf u=(2a-b)v_1e^{3t}+(b-a)v_2e^{2t}$.\end{solution}
\begin{exercise}\label{cw:bank-4G-2}原题4G-2. 对 $\mathbf x'=\begin{pmatrix}5&-1\\3&1\end{pmatrix}\mathbf x$：(a) 用正规模态构造基本矩阵，求 $x(0)=2$, $y(0)=-1$ 的解；(b) 求在0归一化的基本矩阵，再用它求同一初值问题.\end{exercise}
\begin{solution}（官方译审）特征多项式 $\lambda^2-6\lambda+8$ 的根为 $2,4$，向量可取 $(1,3)^T,(1,1)^T$. (a)
\[F(t)=\begin{pmatrix}e^{2t}&e^{4t}\\3e^{2t}&e^{4t}\end{pmatrix},\quad
F(0)^{-1}=\tfrac12\begin{pmatrix}-1&1\\3&-1\end{pmatrix},\quad
\mathbf x=-\tfrac32e^{2t}\binom13+\tfrac72e^{4t}\binom11.\]
(b) 归一化矩阵为
\[\Phi(t)=F(t)F(0)^{-1}=\tfrac12\begin{pmatrix}-e^{2t}+3e^{4t}&e^{2t}-e^{4t}\\-3e^{2t}+3e^{4t}&3e^{2t}-e^{4t}\end{pmatrix}.\]
$\Phi(0)=I$，$\Phi(t)(2,-1)^T$ 给出同一解.\end{solution}
\begin{exercise}\label{cw:bank-4G-3}原题4G-3（星号题）. 把4G-2的矩阵换为 $A=\begin{pmatrix}3&-2\\2&-2\end{pmatrix}$，同样求：(a) 正规模态基本矩阵及初值 $(2,-1)^T$ 的解；(b) 在0归一化的基本矩阵及同一初值解.\end{exercise}
\begin{solution}（编者补解）$\det(\lambda I-A)=\lambda^2-\lambda-2$，根为 $2,-1$，向量 $(2,1)^T,(1,2)^T$. (a) 取 $F=\begin{pmatrix}2e^{2t}&e^{-t}\\e^{2t}&2e^{-t}\end{pmatrix}$，$F(0)^{-1}=\frac13\begin{pmatrix}2&-1\\-1&2\end{pmatrix}$，给 $\mathbf x=\frac53e^{2t}(2,1)^T-\frac43e^{-t}(1,2)^T$. (b)
\[\Phi=\frac13\begin{pmatrix}4e^{2t}-e^{-t}&-2e^{2t}+2e^{-t}\\2e^{2t}-2e^{-t}&-e^{2t}+4e^{-t}\end{pmatrix}.\]
直接检查 $\Phi(0)=I$, $\Phi'=A\Phi$；因此乘初值即为(a)的解.\end{solution}
\originalsection{4H：矩阵指数}
\begin{exercise}\label{cw:bank-4H-1}原题4H-1. 当 $A=\diag(a,b)$ 时计算 $e^{At}$，并直接验证 $\mathbf x=e^{At}\mathbf x_0$ 解 $\mathbf x'=A\mathbf x$, $\mathbf x(0)=\mathbf x_0$.\end{exercise}
\begin{solution}（官方译审）$A^n=\diag(a^n,b^n)$，按绝对收敛的指数级数求和得 $e^{At}=\diag(e^{at},e^{bt})$. 若 $\mathbf x_0=(u,v)^T$，解为 $(ue^{at},ve^{bt})^T$；求导分别为 $a$ 与 $b$ 乘原分量，在0恢复初值.\end{solution}
\begin{exercise}\label{cw:bank-4H-2}原题4H-2. 当 $A=\begin{pmatrix}0&1\\-1&0\end{pmatrix}$ 时计算 $e^{At}$，并作4H-1所要求的初值及微分方程验证.\end{exercise}
\begin{solution}（官方译审）$A^2=-I$，分开偶数与奇数次项，得 $e^{At}=I\cos t+A\sin t=\begin{pmatrix}\cos t&\sin t\\-\sin t&\cos t\end{pmatrix}$. 其导数为 $A$ 乘它，且0时为 $I$，所以乘任意 $\mathbf x_0$ 满足所给初值问题.\end{solution}
\begin{exercise}\label{cw:bank-4H-3}原题4H-3. 当 $A=\begin{pmatrix}1&1\\0&1\end{pmatrix}$ 时直接由无穷级数计算 $e^{At}$，并作4H-1所要求的验证.\end{exercise}
\begin{solution}（编者补解）写 $A=I+N$，$N^2=0$，二项式给 $A^n=I+nN$. 故
\[e^{At}=\sum_{n\ge0}\frac{t^n}{n!}I+\sum_{n\ge1}\frac{nt^n}{n!}N=e^t(I+tN)=e^t\begin{pmatrix}1&t\\0&1\end{pmatrix}.\]
导数为 $e^t[I+(t+1)N]$，同样等于 $A e^{At}$；0时为 $I$，验证完成.\end{solution}
\begin{exercise}\label{cw:bank-4H-4}原题4H-4. 若 $A$ 的元素依赖 $t$，证明 $e^{At}$ 解 $\mathbf x'=A\mathbf x$ 的论证哪里失效？\end{exercise}
\begin{solution}（官方译审）对常矩阵有 $(A^n)'=0$；对变矩阵则 $(A^2)'=A'A+AA'$，两个乘积一般不能交换. 更直接地，指数中 $tA(t)$ 的导数包含 $tA'(t)$，甚至标量 $A(t)=t$ 也使候选 $e^{t^2}$ 的导数为 $2te^{t^2}$，并非 $te^{t^2}$. 因此不能把常系数级数求导公式照搬；一般须求满足 $\Phi'=A(t)\Phi$ 的基本矩阵.\end{solution}
\begin{exercise}\label{cw:bank-4H-5}原题4H-5. 证明：(a) $e^{kIt}=Ie^{kt}$；(b) $Ae^{At}=e^{At}A$. 其中 $k$ 为常数，$I$ 为单位矩阵，$A$ 为任意常方阵.\end{exercise}
\begin{solution}（编者补解）(a) $(kI)^n=k^nI$，逐项求和即得. (b) 在任何有限 $t$ 区间，指数矩阵级数在矩阵范数下绝对一致收敛，可把常矩阵乘法移入求和；$A A^n=A^n A=A^{n+1}$，故左右乘的级数相同.\end{solution}
\begin{exercise}\label{cw:bank-4H-6}原题4H-6. 用 Notes 中第三种方法（写成 $A=B+C$）再算4H-3的矩阵指数.\end{exercise}
\begin{solution}（编者补解）取 $B=I$, $C=\begin{pmatrix}0&1\\0&0\end{pmatrix}$，有 $BC=CB$, $C^2=0$. 由可交换矩阵的二项式及绝对收敛级数乘法，$e^{(B+C)t}=e^{Bt}e^{Ct}=e^t(I+tC)$，正是4H-3结果. 可交换条件是此乘法公式的依据.\end{solution}
\begin{exercise}\label{cw:bank-4H-7}原题4H-7. 设 $A=\begin{pmatrix}1&0\\2&1\end{pmatrix}$，用三种方法算 $e^{At}$：(a) 直接用无穷级数定义；(b) 把矩阵分成较简单矩阵的和；(c) 消元求基本矩阵，再归一化.\end{exercise}
\begin{solution}（官方译审）(a) 归纳得 $A^n=\begin{pmatrix}1&0\\2n&1\end{pmatrix}$，所以 $e^{At}=e^t\begin{pmatrix}1&0\\2t&1\end{pmatrix}$. (b) 令 $A=I+N$, $N=\begin{pmatrix}0&0\\2&0\end{pmatrix}$，可交换且 $N^2=0$，得 $e^{At}=e^t(I+tN)$. (c) 由 $x'=x$, $y'=2x+y$ 得 $x=C_1e^t$, $y=(2C_1t+C_2)e^t$，基本矩阵恰为上式，且0时为 $I$，已经归一化.\end{solution}
\originalsection{4I：非齐次系统}
\begin{exercise}\label{cw:bank-4I-1}原题4I-1. 用参数变易法解
\[\mathbf x'=\begin{pmatrix}1&1\\4&-2\end{pmatrix}\mathbf x+\binom2{-8}-t\binom58.\]
\end{exercise}
\begin{solution}（官方译审）齐次特征值为 $-3,2$，向量 $(1,-4)^T,(1,1)^T$，取
\[F=\begin{pmatrix}e^{-3t}&e^{2t}\\-4e^{-3t}&e^{2t}\end{pmatrix},\qquad
F^{-1}=\tfrac15\begin{pmatrix}e^{3t}&-e^{3t}\\4e^{-2t}&e^{-2t}\end{pmatrix}.\]
写 $\mathbf x_p=F\mathbf v$，则 $\mathbf v'=F^{-1}\mathbf f=((2+3t/5)e^{3t},-28te^{-2t}/5)^T$. 分别积分（取积分常数为零）得
\[v_1=e^{3t}(t/5+3/5),\qquad v_2=e^{-2t}(14t/5+7/5).\]
故 $\mathbf x_p=(3t+2,2t-1)^T$；将它求导得 $(3,2)^T$，代入右端也得同向量. 通解为该特解加 $c_1(1,-4)^Te^{-3t}+c_2(1,1)^Te^{2t}$.\end{solution}
\begin{exercise}\label{cw:bank-4I-2}原题4I-2. 对 $\mathbf x'=\begin{pmatrix}1&1\\4&-2\end{pmatrix}\mathbf x+\binom{e^{-2t}}{-2e^t}$：(a) 用参数变易法求解；(b) 也用待定系数，写强迫项为 $(1,0)^Te^{-2t}+(0,-2)^Te^t$，试取 $\mathbf x_p=ce^{-2t}+de^t$.\end{exercise}
\begin{solution}（官方译审）(a) 用4I-1的 $F$，有
\[\mathbf v'=\tfrac15\binom{e^t+2e^{4t}}{4e^{-4t}-2e^{-t}},\quad
\mathbf v=\tfrac15\binom{e^t+\frac12e^{4t}}{-e^{-4t}+2e^{-t}}.\]
乘 $F$ 得 $\mathbf x_p=(e^t/2,-e^{-2t})^T$. (b) 比较两指数的系数得 $(-2I-A)c=(1,0)^T$, $(I-A)d=(0,-2)^T$，解为 $c=(0,-1)^T$, $d=(1/2,0)^T$，与(a)相同. 通解为该特解加4I-1的齐次通解.\end{solution}
\begin{exercise}\label{cw:bank-4I-3}原题4I-3（星号题）. 用待定系数解 $\mathbf x'=\begin{pmatrix}-1&1\\-5&3\end{pmatrix}\mathbf x+(\sin t,0)^T$.\end{exercise}
\begin{solution}（编者补解）令 $\mathbf x_p=u\sin t+v\cos t$，则比较正余弦得 $-v=Au+(1,0)^T$, $u=Av$，即 $(I+A^2)v=-(1,0)^T$. $A^2=\begin{pmatrix}-4&2\\-10&4\end{pmatrix}$，解得 $v=(-1,-2)^T$, $u=(-1,-1)^T$. 因此 $\mathbf x_p=(-\sin t-\cos t,-\sin t-2\cos t)^T$. 齐次矩阵写为 $A=I+B$，$B^2=-I$，故全部解为
\[\mathbf x=e^t(I\cos t+B\sin t)c+\mathbf x_p,\qquad c\in\R^2.\]
这样既给出特解也包括两独立齐次解.\end{solution}
\begin{exercise}\label{cw:bank-4I-4}原题4I-4. 用待定系数解 $\mathbf x'=\begin{pmatrix}2&-1\\3&-2\end{pmatrix}\mathbf x+(1,-1)^Te^t$.\end{exercise}
\begin{solution}（官方思路展开，纠正手写中间行符号）$A$ 的特征值为 $1,-1$，向量 $(1,1)^T,(1,3)^T$. 强迫频率发生共振，试取 $\mathbf x_p=e^t(c+td)$. 比较 $t$ 与常数项得 $(I-A)d=0$, $(I-A)c=(1,-1)^T-d$. 写 $d=k(1,1)^T$；第二式的相容条件给 $k=2$，可取 $c=(1,0)^T$. 因此
\[\mathbf x=e^t\binom{1+2t}{2t}+c_1e^t\binom11+c_2e^{-t}\binom13.\]
按题面的 $(1,-1)^T$ 代入，特解导数为 $e^t(3+2t,2+2t)^T$，右端一致. 官方手写某一行误写 $(1,1)^T$，本题面未改.\end{solution}
\begin{exercise}\label{cw:bank-4I-5}原题4I-5. 设 $\mathbf x'=A\mathbf x+\mathbf x_0$，其中 $A$ 为可逆常 $n\times n$ 矩阵，$\mathbf x_0$ 为常向量，求一特解.\end{exercise}
\begin{solution}（官方译审）取常解 $\mathbf x_p=-A^{-1}\mathbf x_0$. 它的导数为零，而 $A\mathbf x_p+\mathbf x_0=0$，故满足方程. 可逆性保证代数方程可解；若 $A$ 奇异，只在 $-\mathbf x_0$ 属于 $A$ 的像时才有常特解.\end{solution}
\originalchapter{原题题库五：系统的相图}
\originalsection{5A：相平面与轨道}
\begin{exercise}\label{cw:bank-5A-1}原题5A-1. 求各非线性自治系统的平衡点：(a) $x'=x^2-y^2$, $y'=x-xy$；(b) $x'=1-x+y$, $y'=y+2x^2$.\end{exercise}
\begin{solution}（官方译审）(a) 第二式为 $x(1-y)=0$. 若 $x=0$，第一式给 $y=0$；若 $y=1$，第一式给 $x=\pm1$. 所以平衡点为 $(0,0),(1,1),(-1,1)$. (b) 第一式给 $y=x-1$，代第二式得 $2x^2+x-1=(2x-1)(x+1)=0$，故平衡点为 $(1/2,-1/2),(-1,-2)$.\end{solution}
\begin{exercise}\label{cw:bank-5A-2}原题5A-2. 写出等价一阶系统并找平衡点：(a) $x''+a(x^2-1)x'+x=0$；(b) $x''-x'+1-x^2=0$.\end{exercise}
\begin{solution}（官方译审）令 $y=x'$. (a) $x'=y$, $y'=-a(x^2-1)y-x$，平衡时 $y=0$, $x=0$，只有原点. (b) $x'=y$, $y'=y-1+x^2$，平衡时 $y=0$, $x^2=1$，得 $(\pm1,0)$. 两种转换均由 $x'=y$ 的求导恢复原式.\end{solution}
\begin{exercise}\label{cw:bank-5A-3}原题5A-3. 与 $x'=f(x,y)$, $y'=g(x,y)$ 相比，下列系统的轨道与平衡点一般有什么关系？(a) $x'=-f$, $y'=-g$；(b) $x'=g$, $y'=-f$.\end{exercise}
\begin{solution}（官方译审）(a) 平衡点不变；若原解为 $\mathbf x(t)$，则新解为 $\mathbf x(-t)$，几何轨道相同、箭头相反. (b) 平衡点仍由 $f=g=0$ 决定，不变；每个非平衡点的方向向量从 $(f,g)$ 顺时针转90度为 $(g,-f)$. 两方向点积为零，故相遇处的两族轨道正交；一般轨道形状并不相同，不能把(b)误认为只是时间反向.\end{solution}
\begin{exercise}\label{cw:bank-5A-4}原题5A-4. 对自治系统 $\mathbf x'=(f(x,y),g(x,y))^T$：(a) 若 $\mathbf x_1(t)$ 为解，证明 $\mathbf x_2(t)=\mathbf x_1(t-t_0)$ 也为解，并说明几何关系；(b) 设 $f,g$ 处处连续可微，利用初值唯一性与(a)证明两条不同轨道不能相交. 注意同一相点可在不同时间到达.\end{exercise}
\begin{solution}（官方译审并补证）(a) 链式法则给 $\mathbf x_2'(t)=F(\mathbf x_1(t-t_0))=F(\mathbf x_2(t))$. 它沿同一有向轨道运动，仅时间原点平移；定义区间也相应平移. (b) 若 $\mathbf x_1(t_1)=\mathbf x_2(t_2)=p$，令 $\widetilde{\mathbf x}_2(t)=\mathbf x_2(t-t_1+t_2)$，它在 $t_1$ 与第一解有同一初值. 局部唯一性使它们在公共区间相同；沿最大解延拓后，它们是同一轨道. 因而不同轨道不能穿交. 连续可微保证局部唯一性，不能据此额外声称所有解对所有实数时间存在.\end{solution}
\originalsection{5B：线性系统相图}
\begin{exercise}\label{cw:bank-5B-1}原题5B-1. 对 $x'=-x$, $y'=-2y$：(a) 消去 $t$ 得 $dy/dx=F(x,y)$，解并画曲线；(b) 求参数解并标运动方向，指出(a)漏掉哪些轨道；(c) 一条典型(a)曲线需几条轨道覆盖？在图中指出；(d) 若改为 $x'=x$, $y'=2y$，(a)(b)的图如何改变？\end{exercise}
\begin{solution}（官方译审）(a) 当 $x\ne0$，$dy/dx=2y/x$，积分得 $y=Cx^2$. (b) $x=Ae^{-t}$, $y=Be^{-2t}$，所以箭头指向原点. $x=0$ 被除法漏掉：正、负 $y$ 半轴及原点是另外三条轨道. (c) 一条完整抛物线由左支、右支、原点三条轨道覆盖；解不能在有限时间穿过原点. (d) 轨道集合保持不变，所有非平衡箭头反向，原点从渐近稳定结点变为不稳定结点. 图中同一抛物线的左右分支分别代表不同轨道.
\begin{center}
\begin{tikzpicture}\begin{axis}[width=.47\linewidth,height=5.2cm,axis lines=middle,xlabel={$x$},ylabel={$y$},xmin=-2,xmax=2,ymin=-2,ymax=2,title={抛物线向原点},tick label style={font=\footnotesize},clip=true]
\addplot[structurecolor,thin,no marks] coordinates {(-1.26491,-2.00000) (-1.25589,-1.97157) (-1.24693,-1.94354) (-1.23804,-1.91592) (-1.22921,-1.88868) (-1.22044,-1.86183) (-1.21173,-1.83537) (-1.20309,-1.80928) (-1.19451,-1.78356) (-1.18599,-1.75821) (-1.17753,-1.73321) (-1.16913,-1.70858) (-1.16079,-1.68429) (-1.15251,-1.66035) (-1.14429,-1.63674) (-1.13613,-1.61348) (-1.12802,-1.59054) (-1.11998,-1.56793) (-1.11199,-1.54564) (-1.10406,-1.52367) (-1.09618,-1.50201) (-1.08836,-1.48066) (-1.08060,-1.45961) (-1.07289,-1.43887) (-1.06524,-1.41841) (-1.05764,-1.39825) (-1.05009,-1.37837) (-1.04260,-1.35878) (-1.03517,-1.33946) (-1.02778,-1.32042) (-1.02045,-1.30165) (-1.01317,-1.28315) (-1.00595,-1.26491) (-0.99877,-1.24693) (-0.99165,-1.22921) (-0.98457,-1.21173) (-0.97755,-1.19451) (-0.97058,-1.17753) (-0.96365,-1.16079) (-0.95678,-1.14429) (-0.94996,-1.12802) (-0.94318,-1.11199) (-0.93645,-1.09618) (-0.92977,-1.08060) (-0.92314,-1.06524) (-0.91656,-1.05009) (-0.91002,-1.03517) (-0.90353,-1.02045) (-0.89708,-1.00595) (-0.89068,-0.99165) (-0.88433,-0.97755) (-0.87802,-0.96365) (-0.87176,-0.94996) (-0.86554,-0.93645) (-0.85937,-0.92314) (-0.85324,-0.91002) (-0.84715,-0.89708) (-0.84111,-0.88433) (-0.83511,-0.87176) (-0.82915,-0.85937) (-0.82324,-0.84715) (-0.81737,-0.83511) (-0.81154,-0.82324) (-0.80575,-0.81154) (-0.80000,-0.80000)};
\draw[-{Latex},structurecolor] (axis cs:-1.09618,-1.50201) -- (axis cs:-1.08836,-1.48066);
\draw[-{Latex},structurecolor] (axis cs:-0.92977,-1.08060) -- (axis cs:-0.92314,-1.06524);
\addplot[structurecolor,thin,no marks] coordinates {(-0.80000,-0.80000) (-0.76337,-0.72841) (-0.72841,-0.66322) (-0.69505,-0.60387) (-0.66322,-0.54983) (-0.63285,-0.50063) (-0.60387,-0.45583) (-0.57622,-0.41503) (-0.54983,-0.37789) (-0.52465,-0.34408) (-0.50063,-0.31328) (-0.47770,-0.28525) (-0.45583,-0.25972) (-0.43495,-0.23648) (-0.41503,-0.21532) (-0.39603,-0.19605) (-0.37789,-0.17850) (-0.36059,-0.16253) (-0.34408,-0.14799) (-0.32832,-0.13474) (-0.31328,-0.12268) (-0.29894,-0.11171) (-0.28525,-0.10171) (-0.27219,-0.09261) (-0.25972,-0.08432) (-0.24783,-0.07677) (-0.23648,-0.06990) (-0.22565,-0.06365) (-0.21532,-0.05795) (-0.20546,-0.05277) (-0.19605,-0.04804) (-0.18707,-0.04374) (-0.17850,-0.03983) (-0.17033,-0.03627) (-0.16253,-0.03302) (-0.15509,-0.03007) (-0.14799,-0.02737) (-0.14121,-0.02492) (-0.13474,-0.02269) (-0.12857,-0.02066) (-0.12268,-0.01881) (-0.11707,-0.01713) (-0.11171,-0.01560) (-0.10659,-0.01420) (-0.10171,-0.01293) (-0.09705,-0.01177) (-0.09261,-0.01072) (-0.08837,-0.00976) (-0.08432,-0.00889) (-0.08046,-0.00809) (-0.07677,-0.00737) (-0.07326,-0.00671) (-0.06990,-0.00611) (-0.06670,-0.00556) (-0.06365,-0.00506) (-0.06073,-0.00461) (-0.05795,-0.00420) (-0.05530,-0.00382) (-0.05277,-0.00348) (-0.05035,-0.00317) (-0.04804,-0.00289) (-0.04584,-0.00263) (-0.04374,-0.00239) (-0.04174,-0.00218) (-0.03983,-0.00198)};
\draw[-{Latex},structurecolor] (axis cs:-0.31328,-0.12268) -- (axis cs:-0.29894,-0.11171);
\draw[-{Latex},structurecolor] (axis cs:-0.10659,-0.01420) -- (axis cs:-0.10171,-0.01293);
\addplot[structurecolor,thin,no marks] coordinates {(-2.00000,-1.87500) (-1.97157,-1.82207) (-1.94354,-1.77064) (-1.91592,-1.72066) (-1.88868,-1.67209) (-1.86183,-1.62489) (-1.83537,-1.57902) (-1.80928,-1.53445) (-1.78356,-1.49113) (-1.75821,-1.44904) (-1.73321,-1.40814) (-1.70858,-1.36839) (-1.68429,-1.32976) (-1.66035,-1.29223) (-1.63674,-1.25575) (-1.61348,-1.22030) (-1.59054,-1.18585) (-1.56793,-1.15238) (-1.54564,-1.11985) (-1.52367,-1.08824) (-1.50201,-1.05752) (-1.48066,-1.02767) (-1.45961,-0.99866) (-1.43887,-0.97047) (-1.41841,-0.94308) (-1.39825,-0.91645) (-1.37837,-0.89058) (-1.35878,-0.86544) (-1.33946,-0.84102) (-1.32042,-0.81727) (-1.30165,-0.79420) (-1.28315,-0.77179) (-1.26491,-0.75000) (-1.24693,-0.72883) (-1.22921,-0.70826) (-1.21173,-0.68826) (-1.19451,-0.66883) (-1.17753,-0.64995) (-1.16079,-0.63161) (-1.14429,-0.61378) (-1.12802,-0.59645) (-1.11199,-0.57962) (-1.09618,-0.56325) (-1.08060,-0.54736) (-1.06524,-0.53190) (-1.05009,-0.51689) (-1.03517,-0.50230) (-1.02045,-0.48812) (-1.00595,-0.47434) (-0.99165,-0.46095) (-0.97755,-0.44794) (-0.96365,-0.43530) (-0.94996,-0.42301) (-0.93645,-0.41107) (-0.92314,-0.39946) (-0.91002,-0.38819) (-0.89708,-0.37723) (-0.88433,-0.36658) (-0.87176,-0.35623) (-0.85937,-0.34618) (-0.84715,-0.33641) (-0.83511,-0.32691) (-0.82324,-0.31768) (-0.81154,-0.30871) (-0.80000,-0.30000)};
\draw[-{Latex},structurecolor] (axis cs:-1.50201,-1.05752) -- (axis cs:-1.48066,-1.02767);
\draw[-{Latex},structurecolor] (axis cs:-1.08060,-0.54736) -- (axis cs:-1.06524,-0.53190);
\addplot[structurecolor,thin,no marks] coordinates {(-0.80000,-0.30000) (-0.76337,-0.27315) (-0.72841,-0.24871) (-0.69505,-0.22645) (-0.66322,-0.20619) (-0.63285,-0.18774) (-0.60387,-0.17093) (-0.57622,-0.15564) (-0.54983,-0.14171) (-0.52465,-0.12903) (-0.50063,-0.11748) (-0.47770,-0.10697) (-0.45583,-0.09740) (-0.43495,-0.08868) (-0.41503,-0.08074) (-0.39603,-0.07352) (-0.37789,-0.06694) (-0.36059,-0.06095) (-0.34408,-0.05549) (-0.32832,-0.05053) (-0.31328,-0.04601) (-0.29894,-0.04189) (-0.28525,-0.03814) (-0.27219,-0.03473) (-0.25972,-0.03162) (-0.24783,-0.02879) (-0.23648,-0.02621) (-0.22565,-0.02387) (-0.21532,-0.02173) (-0.20546,-0.01979) (-0.19605,-0.01802) (-0.18707,-0.01640) (-0.17850,-0.01494) (-0.17033,-0.01360) (-0.16253,-0.01238) (-0.15509,-0.01127) (-0.14799,-0.01027) (-0.14121,-0.00935) (-0.13474,-0.00851) (-0.12857,-0.00775) (-0.12268,-0.00706) (-0.11707,-0.00642) (-0.11171,-0.00585) (-0.10659,-0.00533) (-0.10171,-0.00485) (-0.09705,-0.00442) (-0.09261,-0.00402) (-0.08837,-0.00366) (-0.08432,-0.00333) (-0.08046,-0.00303) (-0.07677,-0.00276) (-0.07326,-0.00252) (-0.06990,-0.00229) (-0.06670,-0.00209) (-0.06365,-0.00190) (-0.06073,-0.00173) (-0.05795,-0.00157) (-0.05530,-0.00143) (-0.05277,-0.00131) (-0.05035,-0.00119) (-0.04804,-0.00108) (-0.04584,-0.00099) (-0.04374,-0.00090) (-0.04174,-0.00082) (-0.03983,-0.00074)};
\draw[-{Latex},structurecolor] (axis cs:-0.31328,-0.04601) -- (axis cs:-0.29894,-0.04189);
\draw[-{Latex},structurecolor] (axis cs:-0.10659,-0.00533) -- (axis cs:-0.10171,-0.00485);
\addplot[structurecolor,thin,no marks] coordinates {(-2.00000,1.87500) (-1.97157,1.82207) (-1.94354,1.77064) (-1.91592,1.72066) (-1.88868,1.67209) (-1.86183,1.62489) (-1.83537,1.57902) (-1.80928,1.53445) (-1.78356,1.49113) (-1.75821,1.44904) (-1.73321,1.40814) (-1.70858,1.36839) (-1.68429,1.32976) (-1.66035,1.29223) (-1.63674,1.25575) (-1.61348,1.22030) (-1.59054,1.18585) (-1.56793,1.15238) (-1.54564,1.11985) (-1.52367,1.08824) (-1.50201,1.05752) (-1.48066,1.02767) (-1.45961,0.99866) (-1.43887,0.97047) (-1.41841,0.94308) (-1.39825,0.91645) (-1.37837,0.89058) (-1.35878,0.86544) (-1.33946,0.84102) (-1.32042,0.81727) (-1.30165,0.79420) (-1.28315,0.77179) (-1.26491,0.75000) (-1.24693,0.72883) (-1.22921,0.70826) (-1.21173,0.68826) (-1.19451,0.66883) (-1.17753,0.64995) (-1.16079,0.63161) (-1.14429,0.61378) (-1.12802,0.59645) (-1.11199,0.57962) (-1.09618,0.56325) (-1.08060,0.54736) (-1.06524,0.53190) (-1.05009,0.51689) (-1.03517,0.50230) (-1.02045,0.48812) (-1.00595,0.47434) (-0.99165,0.46095) (-0.97755,0.44794) (-0.96365,0.43530) (-0.94996,0.42301) (-0.93645,0.41107) (-0.92314,0.39946) (-0.91002,0.38819) (-0.89708,0.37723) (-0.88433,0.36658) (-0.87176,0.35623) (-0.85937,0.34618) (-0.84715,0.33641) (-0.83511,0.32691) (-0.82324,0.31768) (-0.81154,0.30871) (-0.80000,0.30000)};
\draw[-{Latex},structurecolor] (axis cs:-1.50201,1.05752) -- (axis cs:-1.48066,1.02767);
\draw[-{Latex},structurecolor] (axis cs:-1.08060,0.54736) -- (axis cs:-1.06524,0.53190);
\addplot[structurecolor,thin,no marks] coordinates {(-0.80000,0.30000) (-0.76337,0.27315) (-0.72841,0.24871) (-0.69505,0.22645) (-0.66322,0.20619) (-0.63285,0.18774) (-0.60387,0.17093) (-0.57622,0.15564) (-0.54983,0.14171) (-0.52465,0.12903) (-0.50063,0.11748) (-0.47770,0.10697) (-0.45583,0.09740) (-0.43495,0.08868) (-0.41503,0.08074) (-0.39603,0.07352) (-0.37789,0.06694) (-0.36059,0.06095) (-0.34408,0.05549) (-0.32832,0.05053) (-0.31328,0.04601) (-0.29894,0.04189) (-0.28525,0.03814) (-0.27219,0.03473) (-0.25972,0.03162) (-0.24783,0.02879) (-0.23648,0.02621) (-0.22565,0.02387) (-0.21532,0.02173) (-0.20546,0.01979) (-0.19605,0.01802) (-0.18707,0.01640) (-0.17850,0.01494) (-0.17033,0.01360) (-0.16253,0.01238) (-0.15509,0.01127) (-0.14799,0.01027) (-0.14121,0.00935) (-0.13474,0.00851) (-0.12857,0.00775) (-0.12268,0.00706) (-0.11707,0.00642) (-0.11171,0.00585) (-0.10659,0.00533) (-0.10171,0.00485) (-0.09705,0.00442) (-0.09261,0.00402) (-0.08837,0.00366) (-0.08432,0.00333) (-0.08046,0.00303) (-0.07677,0.00276) (-0.07326,0.00252) (-0.06990,0.00229) (-0.06670,0.00209) (-0.06365,0.00190) (-0.06073,0.00173) (-0.05795,0.00157) (-0.05530,0.00143) (-0.05277,0.00131) (-0.05035,0.00119) (-0.04804,0.00108) (-0.04584,0.00099) (-0.04374,0.00090) (-0.04174,0.00082) (-0.03983,0.00074)};
\draw[-{Latex},structurecolor] (axis cs:-0.31328,0.04601) -- (axis cs:-0.29894,0.04189);
\draw[-{Latex},structurecolor] (axis cs:-0.10659,0.00533) -- (axis cs:-0.10171,0.00485);
\addplot[structurecolor,thin,no marks] coordinates {(-1.26491,2.00000) (-1.25589,1.97157) (-1.24693,1.94354) (-1.23804,1.91592) (-1.22921,1.88868) (-1.22044,1.86183) (-1.21173,1.83537) (-1.20309,1.80928) (-1.19451,1.78356) (-1.18599,1.75821) (-1.17753,1.73321) (-1.16913,1.70858) (-1.16079,1.68429) (-1.15251,1.66035) (-1.14429,1.63674) (-1.13613,1.61348) (-1.12802,1.59054) (-1.11998,1.56793) (-1.11199,1.54564) (-1.10406,1.52367) (-1.09618,1.50201) (-1.08836,1.48066) (-1.08060,1.45961) (-1.07289,1.43887) (-1.06524,1.41841) (-1.05764,1.39825) (-1.05009,1.37837) (-1.04260,1.35878) (-1.03517,1.33946) (-1.02778,1.32042) (-1.02045,1.30165) (-1.01317,1.28315) (-1.00595,1.26491) (-0.99877,1.24693) (-0.99165,1.22921) (-0.98457,1.21173) (-0.97755,1.19451) (-0.97058,1.17753) (-0.96365,1.16079) (-0.95678,1.14429) (-0.94996,1.12802) (-0.94318,1.11199) (-0.93645,1.09618) (-0.92977,1.08060) (-0.92314,1.06524) (-0.91656,1.05009) (-0.91002,1.03517) (-0.90353,1.02045) (-0.89708,1.00595) (-0.89068,0.99165) (-0.88433,0.97755) (-0.87802,0.96365) (-0.87176,0.94996) (-0.86554,0.93645) (-0.85937,0.92314) (-0.85324,0.91002) (-0.84715,0.89708) (-0.84111,0.88433) (-0.83511,0.87176) (-0.82915,0.85937) (-0.82324,0.84715) (-0.81737,0.83511) (-0.81154,0.82324) (-0.80575,0.81154) (-0.80000,0.80000)};
\draw[-{Latex},structurecolor] (axis cs:-1.09618,1.50201) -- (axis cs:-1.08836,1.48066);
\draw[-{Latex},structurecolor] (axis cs:-0.92977,1.08060) -- (axis cs:-0.92314,1.06524);
\addplot[structurecolor,thin,no marks] coordinates {(-0.80000,0.80000) (-0.76337,0.72841) (-0.72841,0.66322) (-0.69505,0.60387) (-0.66322,0.54983) (-0.63285,0.50063) (-0.60387,0.45583) (-0.57622,0.41503) (-0.54983,0.37789) (-0.52465,0.34408) (-0.50063,0.31328) (-0.47770,0.28525) (-0.45583,0.25972) (-0.43495,0.23648) (-0.41503,0.21532) (-0.39603,0.19605) (-0.37789,0.17850) (-0.36059,0.16253) (-0.34408,0.14799) (-0.32832,0.13474) (-0.31328,0.12268) (-0.29894,0.11171) (-0.28525,0.10171) (-0.27219,0.09261) (-0.25972,0.08432) (-0.24783,0.07677) (-0.23648,0.06990) (-0.22565,0.06365) (-0.21532,0.05795) (-0.20546,0.05277) (-0.19605,0.04804) (-0.18707,0.04374) (-0.17850,0.03983) (-0.17033,0.03627) (-0.16253,0.03302) (-0.15509,0.03007) (-0.14799,0.02737) (-0.14121,0.02492) (-0.13474,0.02269) (-0.12857,0.02066) (-0.12268,0.01881) (-0.11707,0.01713) (-0.11171,0.01560) (-0.10659,0.01420) (-0.10171,0.01293) (-0.09705,0.01177) (-0.09261,0.01072) (-0.08837,0.00976) (-0.08432,0.00889) (-0.08046,0.00809) (-0.07677,0.00737) (-0.07326,0.00671) (-0.06990,0.00611) (-0.06670,0.00556) (-0.06365,0.00506) (-0.06073,0.00461) (-0.05795,0.00420) (-0.05530,0.00382) (-0.05277,0.00348) (-0.05035,0.00317) (-0.04804,0.00289) (-0.04584,0.00263) (-0.04374,0.00239) (-0.04174,0.00218) (-0.03983,0.00198)};
\draw[-{Latex},structurecolor] (axis cs:-0.31328,0.12268) -- (axis cs:-0.29894,0.11171);
\draw[-{Latex},structurecolor] (axis cs:-0.10659,0.01420) -- (axis cs:-0.10171,0.01293);
\addplot[structurecolor,thin,no marks] coordinates {(1.26491,-2.00000) (1.25589,-1.97157) (1.24693,-1.94354) (1.23804,-1.91592) (1.22921,-1.88868) (1.22044,-1.86183) (1.21173,-1.83537) (1.20309,-1.80928) (1.19451,-1.78356) (1.18599,-1.75821) (1.17753,-1.73321) (1.16913,-1.70858) (1.16079,-1.68429) (1.15251,-1.66035) (1.14429,-1.63674) (1.13613,-1.61348) (1.12802,-1.59054) (1.11998,-1.56793) (1.11199,-1.54564) (1.10406,-1.52367) (1.09618,-1.50201) (1.08836,-1.48066) (1.08060,-1.45961) (1.07289,-1.43887) (1.06524,-1.41841) (1.05764,-1.39825) (1.05009,-1.37837) (1.04260,-1.35878) (1.03517,-1.33946) (1.02778,-1.32042) (1.02045,-1.30165) (1.01317,-1.28315) (1.00595,-1.26491) (0.99877,-1.24693) (0.99165,-1.22921) (0.98457,-1.21173) (0.97755,-1.19451) (0.97058,-1.17753) (0.96365,-1.16079) (0.95678,-1.14429) (0.94996,-1.12802) (0.94318,-1.11199) (0.93645,-1.09618) (0.92977,-1.08060) (0.92314,-1.06524) (0.91656,-1.05009) (0.91002,-1.03517) (0.90353,-1.02045) (0.89708,-1.00595) (0.89068,-0.99165) (0.88433,-0.97755) (0.87802,-0.96365) (0.87176,-0.94996) (0.86554,-0.93645) (0.85937,-0.92314) (0.85324,-0.91002) (0.84715,-0.89708) (0.84111,-0.88433) (0.83511,-0.87176) (0.82915,-0.85937) (0.82324,-0.84715) (0.81737,-0.83511) (0.81154,-0.82324) (0.80575,-0.81154) (0.80000,-0.80000)};
\draw[-{Latex},structurecolor] (axis cs:1.09618,-1.50201) -- (axis cs:1.08836,-1.48066);
\draw[-{Latex},structurecolor] (axis cs:0.92977,-1.08060) -- (axis cs:0.92314,-1.06524);
\addplot[structurecolor,thin,no marks] coordinates {(0.80000,-0.80000) (0.76337,-0.72841) (0.72841,-0.66322) (0.69505,-0.60387) (0.66322,-0.54983) (0.63285,-0.50063) (0.60387,-0.45583) (0.57622,-0.41503) (0.54983,-0.37789) (0.52465,-0.34408) (0.50063,-0.31328) (0.47770,-0.28525) (0.45583,-0.25972) (0.43495,-0.23648) (0.41503,-0.21532) (0.39603,-0.19605) (0.37789,-0.17850) (0.36059,-0.16253) (0.34408,-0.14799) (0.32832,-0.13474) (0.31328,-0.12268) (0.29894,-0.11171) (0.28525,-0.10171) (0.27219,-0.09261) (0.25972,-0.08432) (0.24783,-0.07677) (0.23648,-0.06990) (0.22565,-0.06365) (0.21532,-0.05795) (0.20546,-0.05277) (0.19605,-0.04804) (0.18707,-0.04374) (0.17850,-0.03983) (0.17033,-0.03627) (0.16253,-0.03302) (0.15509,-0.03007) (0.14799,-0.02737) (0.14121,-0.02492) (0.13474,-0.02269) (0.12857,-0.02066) (0.12268,-0.01881) (0.11707,-0.01713) (0.11171,-0.01560) (0.10659,-0.01420) (0.10171,-0.01293) (0.09705,-0.01177) (0.09261,-0.01072) (0.08837,-0.00976) (0.08432,-0.00889) (0.08046,-0.00809) (0.07677,-0.00737) (0.07326,-0.00671) (0.06990,-0.00611) (0.06670,-0.00556) (0.06365,-0.00506) (0.06073,-0.00461) (0.05795,-0.00420) (0.05530,-0.00382) (0.05277,-0.00348) (0.05035,-0.00317) (0.04804,-0.00289) (0.04584,-0.00263) (0.04374,-0.00239) (0.04174,-0.00218) (0.03983,-0.00198)};
\draw[-{Latex},structurecolor] (axis cs:0.31328,-0.12268) -- (axis cs:0.29894,-0.11171);
\draw[-{Latex},structurecolor] (axis cs:0.10659,-0.01420) -- (axis cs:0.10171,-0.01293);
\addplot[structurecolor,thin,no marks] coordinates {(2.00000,-1.87500) (1.97157,-1.82207) (1.94354,-1.77064) (1.91592,-1.72066) (1.88868,-1.67209) (1.86183,-1.62489) (1.83537,-1.57902) (1.80928,-1.53445) (1.78356,-1.49113) (1.75821,-1.44904) (1.73321,-1.40814) (1.70858,-1.36839) (1.68429,-1.32976) (1.66035,-1.29223) (1.63674,-1.25575) (1.61348,-1.22030) (1.59054,-1.18585) (1.56793,-1.15238) (1.54564,-1.11985) (1.52367,-1.08824) (1.50201,-1.05752) (1.48066,-1.02767) (1.45961,-0.99866) (1.43887,-0.97047) (1.41841,-0.94308) (1.39825,-0.91645) (1.37837,-0.89058) (1.35878,-0.86544) (1.33946,-0.84102) (1.32042,-0.81727) (1.30165,-0.79420) (1.28315,-0.77179) (1.26491,-0.75000) (1.24693,-0.72883) (1.22921,-0.70826) (1.21173,-0.68826) (1.19451,-0.66883) (1.17753,-0.64995) (1.16079,-0.63161) (1.14429,-0.61378) (1.12802,-0.59645) (1.11199,-0.57962) (1.09618,-0.56325) (1.08060,-0.54736) (1.06524,-0.53190) (1.05009,-0.51689) (1.03517,-0.50230) (1.02045,-0.48812) (1.00595,-0.47434) (0.99165,-0.46095) (0.97755,-0.44794) (0.96365,-0.43530) (0.94996,-0.42301) (0.93645,-0.41107) (0.92314,-0.39946) (0.91002,-0.38819) (0.89708,-0.37723) (0.88433,-0.36658) (0.87176,-0.35623) (0.85937,-0.34618) (0.84715,-0.33641) (0.83511,-0.32691) (0.82324,-0.31768) (0.81154,-0.30871) (0.80000,-0.30000)};
\draw[-{Latex},structurecolor] (axis cs:1.50201,-1.05752) -- (axis cs:1.48066,-1.02767);
\draw[-{Latex},structurecolor] (axis cs:1.08060,-0.54736) -- (axis cs:1.06524,-0.53190);
\addplot[structurecolor,thin,no marks] coordinates {(0.80000,-0.30000) (0.76337,-0.27315) (0.72841,-0.24871) (0.69505,-0.22645) (0.66322,-0.20619) (0.63285,-0.18774) (0.60387,-0.17093) (0.57622,-0.15564) (0.54983,-0.14171) (0.52465,-0.12903) (0.50063,-0.11748) (0.47770,-0.10697) (0.45583,-0.09740) (0.43495,-0.08868) (0.41503,-0.08074) (0.39603,-0.07352) (0.37789,-0.06694) (0.36059,-0.06095) (0.34408,-0.05549) (0.32832,-0.05053) (0.31328,-0.04601) (0.29894,-0.04189) (0.28525,-0.03814) (0.27219,-0.03473) (0.25972,-0.03162) (0.24783,-0.02879) (0.23648,-0.02621) (0.22565,-0.02387) (0.21532,-0.02173) (0.20546,-0.01979) (0.19605,-0.01802) (0.18707,-0.01640) (0.17850,-0.01494) (0.17033,-0.01360) (0.16253,-0.01238) (0.15509,-0.01127) (0.14799,-0.01027) (0.14121,-0.00935) (0.13474,-0.00851) (0.12857,-0.00775) (0.12268,-0.00706) (0.11707,-0.00642) (0.11171,-0.00585) (0.10659,-0.00533) (0.10171,-0.00485) (0.09705,-0.00442) (0.09261,-0.00402) (0.08837,-0.00366) (0.08432,-0.00333) (0.08046,-0.00303) (0.07677,-0.00276) (0.07326,-0.00252) (0.06990,-0.00229) (0.06670,-0.00209) (0.06365,-0.00190) (0.06073,-0.00173) (0.05795,-0.00157) (0.05530,-0.00143) (0.05277,-0.00131) (0.05035,-0.00119) (0.04804,-0.00108) (0.04584,-0.00099) (0.04374,-0.00090) (0.04174,-0.00082) (0.03983,-0.00074)};
\draw[-{Latex},structurecolor] (axis cs:0.31328,-0.04601) -- (axis cs:0.29894,-0.04189);
\draw[-{Latex},structurecolor] (axis cs:0.10659,-0.00533) -- (axis cs:0.10171,-0.00485);
\addplot[structurecolor,thin,no marks] coordinates {(2.00000,1.87500) (1.97157,1.82207) (1.94354,1.77064) (1.91592,1.72066) (1.88868,1.67209) (1.86183,1.62489) (1.83537,1.57902) (1.80928,1.53445) (1.78356,1.49113) (1.75821,1.44904) (1.73321,1.40814) (1.70858,1.36839) (1.68429,1.32976) (1.66035,1.29223) (1.63674,1.25575) (1.61348,1.22030) (1.59054,1.18585) (1.56793,1.15238) (1.54564,1.11985) (1.52367,1.08824) (1.50201,1.05752) (1.48066,1.02767) (1.45961,0.99866) (1.43887,0.97047) (1.41841,0.94308) (1.39825,0.91645) (1.37837,0.89058) (1.35878,0.86544) (1.33946,0.84102) (1.32042,0.81727) (1.30165,0.79420) (1.28315,0.77179) (1.26491,0.75000) (1.24693,0.72883) (1.22921,0.70826) (1.21173,0.68826) (1.19451,0.66883) (1.17753,0.64995) (1.16079,0.63161) (1.14429,0.61378) (1.12802,0.59645) (1.11199,0.57962) (1.09618,0.56325) (1.08060,0.54736) (1.06524,0.53190) (1.05009,0.51689) (1.03517,0.50230) (1.02045,0.48812) (1.00595,0.47434) (0.99165,0.46095) (0.97755,0.44794) (0.96365,0.43530) (0.94996,0.42301) (0.93645,0.41107) (0.92314,0.39946) (0.91002,0.38819) (0.89708,0.37723) (0.88433,0.36658) (0.87176,0.35623) (0.85937,0.34618) (0.84715,0.33641) (0.83511,0.32691) (0.82324,0.31768) (0.81154,0.30871) (0.80000,0.30000)};
\draw[-{Latex},structurecolor] (axis cs:1.50201,1.05752) -- (axis cs:1.48066,1.02767);
\draw[-{Latex},structurecolor] (axis cs:1.08060,0.54736) -- (axis cs:1.06524,0.53190);
\addplot[structurecolor,thin,no marks] coordinates {(0.80000,0.30000) (0.76337,0.27315) (0.72841,0.24871) (0.69505,0.22645) (0.66322,0.20619) (0.63285,0.18774) (0.60387,0.17093) (0.57622,0.15564) (0.54983,0.14171) (0.52465,0.12903) (0.50063,0.11748) (0.47770,0.10697) (0.45583,0.09740) (0.43495,0.08868) (0.41503,0.08074) (0.39603,0.07352) (0.37789,0.06694) (0.36059,0.06095) (0.34408,0.05549) (0.32832,0.05053) (0.31328,0.04601) (0.29894,0.04189) (0.28525,0.03814) (0.27219,0.03473) (0.25972,0.03162) (0.24783,0.02879) (0.23648,0.02621) (0.22565,0.02387) (0.21532,0.02173) (0.20546,0.01979) (0.19605,0.01802) (0.18707,0.01640) (0.17850,0.01494) (0.17033,0.01360) (0.16253,0.01238) (0.15509,0.01127) (0.14799,0.01027) (0.14121,0.00935) (0.13474,0.00851) (0.12857,0.00775) (0.12268,0.00706) (0.11707,0.00642) (0.11171,0.00585) (0.10659,0.00533) (0.10171,0.00485) (0.09705,0.00442) (0.09261,0.00402) (0.08837,0.00366) (0.08432,0.00333) (0.08046,0.00303) (0.07677,0.00276) (0.07326,0.00252) (0.06990,0.00229) (0.06670,0.00209) (0.06365,0.00190) (0.06073,0.00173) (0.05795,0.00157) (0.05530,0.00143) (0.05277,0.00131) (0.05035,0.00119) (0.04804,0.00108) (0.04584,0.00099) (0.04374,0.00090) (0.04174,0.00082) (0.03983,0.00074)};
\draw[-{Latex},structurecolor] (axis cs:0.31328,0.04601) -- (axis cs:0.29894,0.04189);
\draw[-{Latex},structurecolor] (axis cs:0.10659,0.00533) -- (axis cs:0.10171,0.00485);
\addplot[structurecolor,thin,no marks] coordinates {(1.26491,2.00000) (1.25589,1.97157) (1.24693,1.94354) (1.23804,1.91592) (1.22921,1.88868) (1.22044,1.86183) (1.21173,1.83537) (1.20309,1.80928) (1.19451,1.78356) (1.18599,1.75821) (1.17753,1.73321) (1.16913,1.70858) (1.16079,1.68429) (1.15251,1.66035) (1.14429,1.63674) (1.13613,1.61348) (1.12802,1.59054) (1.11998,1.56793) (1.11199,1.54564) (1.10406,1.52367) (1.09618,1.50201) (1.08836,1.48066) (1.08060,1.45961) (1.07289,1.43887) (1.06524,1.41841) (1.05764,1.39825) (1.05009,1.37837) (1.04260,1.35878) (1.03517,1.33946) (1.02778,1.32042) (1.02045,1.30165) (1.01317,1.28315) (1.00595,1.26491) (0.99877,1.24693) (0.99165,1.22921) (0.98457,1.21173) (0.97755,1.19451) (0.97058,1.17753) (0.96365,1.16079) (0.95678,1.14429) (0.94996,1.12802) (0.94318,1.11199) (0.93645,1.09618) (0.92977,1.08060) (0.92314,1.06524) (0.91656,1.05009) (0.91002,1.03517) (0.90353,1.02045) (0.89708,1.00595) (0.89068,0.99165) (0.88433,0.97755) (0.87802,0.96365) (0.87176,0.94996) (0.86554,0.93645) (0.85937,0.92314) (0.85324,0.91002) (0.84715,0.89708) (0.84111,0.88433) (0.83511,0.87176) (0.82915,0.85937) (0.82324,0.84715) (0.81737,0.83511) (0.81154,0.82324) (0.80575,0.81154) (0.80000,0.80000)};
\draw[-{Latex},structurecolor] (axis cs:1.09618,1.50201) -- (axis cs:1.08836,1.48066);
\draw[-{Latex},structurecolor] (axis cs:0.92977,1.08060) -- (axis cs:0.92314,1.06524);
\addplot[structurecolor,thin,no marks] coordinates {(0.80000,0.80000) (0.76337,0.72841) (0.72841,0.66322) (0.69505,0.60387) (0.66322,0.54983) (0.63285,0.50063) (0.60387,0.45583) (0.57622,0.41503) (0.54983,0.37789) (0.52465,0.34408) (0.50063,0.31328) (0.47770,0.28525) (0.45583,0.25972) (0.43495,0.23648) (0.41503,0.21532) (0.39603,0.19605) (0.37789,0.17850) (0.36059,0.16253) (0.34408,0.14799) (0.32832,0.13474) (0.31328,0.12268) (0.29894,0.11171) (0.28525,0.10171) (0.27219,0.09261) (0.25972,0.08432) (0.24783,0.07677) (0.23648,0.06990) (0.22565,0.06365) (0.21532,0.05795) (0.20546,0.05277) (0.19605,0.04804) (0.18707,0.04374) (0.17850,0.03983) (0.17033,0.03627) (0.16253,0.03302) (0.15509,0.03007) (0.14799,0.02737) (0.14121,0.02492) (0.13474,0.02269) (0.12857,0.02066) (0.12268,0.01881) (0.11707,0.01713) (0.11171,0.01560) (0.10659,0.01420) (0.10171,0.01293) (0.09705,0.01177) (0.09261,0.01072) (0.08837,0.00976) (0.08432,0.00889) (0.08046,0.00809) (0.07677,0.00737) (0.07326,0.00671) (0.06990,0.00611) (0.06670,0.00556) (0.06365,0.00506) (0.06073,0.00461) (0.05795,0.00420) (0.05530,0.00382) (0.05277,0.00348) (0.05035,0.00317) (0.04804,0.00289) (0.04584,0.00263) (0.04374,0.00239) (0.04174,0.00218) (0.03983,0.00198)};
\draw[-{Latex},structurecolor] (axis cs:0.31328,0.12268) -- (axis cs:0.29894,0.11171);
\draw[-{Latex},structurecolor] (axis cs:0.10659,0.01420) -- (axis cs:0.10171,0.01293);
\addplot[structurecolor,thin,no marks] coordinates {(0.00000,2.00000) (0.00000,1.99103) (0.00000,1.98210) (0.00000,1.97321) (0.00000,1.96436) (0.00000,1.95555) (0.00000,1.94678) (0.00000,1.93805) (0.00000,1.92936) (0.00000,1.92070) (0.00000,1.91209) (0.00000,1.90351) (0.00000,1.89498) (0.00000,1.88648) (0.00000,1.87802) (0.00000,1.86959) (0.00000,1.86121) (0.00000,1.85286) (0.00000,1.84455) (0.00000,1.83628) (0.00000,1.82804) (0.00000,1.81985) (0.00000,1.81168) (0.00000,1.80356) (0.00000,1.79547) (0.00000,1.78742) (0.00000,1.77940) (0.00000,1.77142) (0.00000,1.76347) (0.00000,1.75557) (0.00000,1.74769) (0.00000,1.73985) (0.00000,1.73205) (0.00000,1.72428) (0.00000,1.71655) (0.00000,1.70885) (0.00000,1.70119) (0.00000,1.69356) (0.00000,1.68596) (0.00000,1.67840) (0.00000,1.67087) (0.00000,1.66338) (0.00000,1.65592) (0.00000,1.64849) (0.00000,1.64110) (0.00000,1.63374) (0.00000,1.62641) (0.00000,1.61912) (0.00000,1.61185) (0.00000,1.60463) (0.00000,1.59743) (0.00000,1.59026) (0.00000,1.58313) (0.00000,1.57603) (0.00000,1.56896) (0.00000,1.56193) (0.00000,1.55492) (0.00000,1.54795) (0.00000,1.54101) (0.00000,1.53409) (0.00000,1.52721) (0.00000,1.52036) (0.00000,1.51355) (0.00000,1.50676) (0.00000,1.50000)};
\draw[-{Latex},structurecolor] (axis cs:0.00000,1.82804) -- (axis cs:0.00000,1.81985);
\draw[-{Latex},structurecolor] (axis cs:0.00000,1.64849) -- (axis cs:0.00000,1.64110);
\addplot[structurecolor,thin,no marks] coordinates {(0.00000,1.50000) (0.00000,1.36577) (0.00000,1.24354) (0.00000,1.13226) (0.00000,1.03093) (0.00000,0.93868) (0.00000,0.85467) (0.00000,0.77819) (0.00000,0.70855) (0.00000,0.64514) (0.00000,0.58741) (0.00000,0.53484) (0.00000,0.48698) (0.00000,0.44340) (0.00000,0.40372) (0.00000,0.36759) (0.00000,0.33470) (0.00000,0.30474) (0.00000,0.27747) (0.00000,0.25264) (0.00000,0.23003) (0.00000,0.20945) (0.00000,0.19070) (0.00000,0.17364) (0.00000,0.15810) (0.00000,0.14395) (0.00000,0.13107) (0.00000,0.11934) (0.00000,0.10866) (0.00000,0.09894) (0.00000,0.09008) (0.00000,0.08202) (0.00000,0.07468) (0.00000,0.06800) (0.00000,0.06191) (0.00000,0.05637) (0.00000,0.05133) (0.00000,0.04673) (0.00000,0.04255) (0.00000,0.03874) (0.00000,0.03528) (0.00000,0.03212) (0.00000,0.02925) (0.00000,0.02663) (0.00000,0.02425) (0.00000,0.02208) (0.00000,0.02010) (0.00000,0.01830) (0.00000,0.01666) (0.00000,0.01517) (0.00000,0.01381) (0.00000,0.01258) (0.00000,0.01145) (0.00000,0.01043) (0.00000,0.00949) (0.00000,0.00864) (0.00000,0.00787) (0.00000,0.00717) (0.00000,0.00653) (0.00000,0.00594) (0.00000,0.00541) (0.00000,0.00493) (0.00000,0.00448) (0.00000,0.00408) (0.00000,0.00372)};
\draw[-{Latex},structurecolor] (axis cs:0.00000,0.23003) -- (axis cs:0.00000,0.20945);
\draw[-{Latex},structurecolor] (axis cs:0.00000,0.02663) -- (axis cs:0.00000,0.02425);
\addplot[structurecolor,thin,no marks] coordinates {(0.00000,-2.00000) (0.00000,-1.99103) (0.00000,-1.98210) (0.00000,-1.97321) (0.00000,-1.96436) (0.00000,-1.95555) (0.00000,-1.94678) (0.00000,-1.93805) (0.00000,-1.92936) (0.00000,-1.92070) (0.00000,-1.91209) (0.00000,-1.90351) (0.00000,-1.89498) (0.00000,-1.88648) (0.00000,-1.87802) (0.00000,-1.86959) (0.00000,-1.86121) (0.00000,-1.85286) (0.00000,-1.84455) (0.00000,-1.83628) (0.00000,-1.82804) (0.00000,-1.81985) (0.00000,-1.81168) (0.00000,-1.80356) (0.00000,-1.79547) (0.00000,-1.78742) (0.00000,-1.77940) (0.00000,-1.77142) (0.00000,-1.76347) (0.00000,-1.75557) (0.00000,-1.74769) (0.00000,-1.73985) (0.00000,-1.73205) (0.00000,-1.72428) (0.00000,-1.71655) (0.00000,-1.70885) (0.00000,-1.70119) (0.00000,-1.69356) (0.00000,-1.68596) (0.00000,-1.67840) (0.00000,-1.67087) (0.00000,-1.66338) (0.00000,-1.65592) (0.00000,-1.64849) (0.00000,-1.64110) (0.00000,-1.63374) (0.00000,-1.62641) (0.00000,-1.61912) (0.00000,-1.61185) (0.00000,-1.60463) (0.00000,-1.59743) (0.00000,-1.59026) (0.00000,-1.58313) (0.00000,-1.57603) (0.00000,-1.56896) (0.00000,-1.56193) (0.00000,-1.55492) (0.00000,-1.54795) (0.00000,-1.54101) (0.00000,-1.53409) (0.00000,-1.52721) (0.00000,-1.52036) (0.00000,-1.51355) (0.00000,-1.50676) (0.00000,-1.50000)};
\draw[-{Latex},structurecolor] (axis cs:0.00000,-1.82804) -- (axis cs:0.00000,-1.81985);
\draw[-{Latex},structurecolor] (axis cs:0.00000,-1.64849) -- (axis cs:0.00000,-1.64110);
\addplot[structurecolor,thin,no marks] coordinates {(0.00000,-1.50000) (0.00000,-1.36577) (0.00000,-1.24354) (0.00000,-1.13226) (0.00000,-1.03093) (0.00000,-0.93868) (0.00000,-0.85467) (0.00000,-0.77819) (0.00000,-0.70855) (0.00000,-0.64514) (0.00000,-0.58741) (0.00000,-0.53484) (0.00000,-0.48698) (0.00000,-0.44340) (0.00000,-0.40372) (0.00000,-0.36759) (0.00000,-0.33470) (0.00000,-0.30474) (0.00000,-0.27747) (0.00000,-0.25264) (0.00000,-0.23003) (0.00000,-0.20945) (0.00000,-0.19070) (0.00000,-0.17364) (0.00000,-0.15810) (0.00000,-0.14395) (0.00000,-0.13107) (0.00000,-0.11934) (0.00000,-0.10866) (0.00000,-0.09894) (0.00000,-0.09008) (0.00000,-0.08202) (0.00000,-0.07468) (0.00000,-0.06800) (0.00000,-0.06191) (0.00000,-0.05637) (0.00000,-0.05133) (0.00000,-0.04673) (0.00000,-0.04255) (0.00000,-0.03874) (0.00000,-0.03528) (0.00000,-0.03212) (0.00000,-0.02925) (0.00000,-0.02663) (0.00000,-0.02425) (0.00000,-0.02208) (0.00000,-0.02010) (0.00000,-0.01830) (0.00000,-0.01666) (0.00000,-0.01517) (0.00000,-0.01381) (0.00000,-0.01258) (0.00000,-0.01145) (0.00000,-0.01043) (0.00000,-0.00949) (0.00000,-0.00864) (0.00000,-0.00787) (0.00000,-0.00717) (0.00000,-0.00653) (0.00000,-0.00594) (0.00000,-0.00541) (0.00000,-0.00493) (0.00000,-0.00448) (0.00000,-0.00408) (0.00000,-0.00372)};
\draw[-{Latex},structurecolor] (axis cs:0.00000,-0.23003) -- (axis cs:0.00000,-0.20945);
\draw[-{Latex},structurecolor] (axis cs:0.00000,-0.02663) -- (axis cs:0.00000,-0.02425);
\addplot[only marks,mark=*,mark size=1.4pt] coordinates {(0,0)};
\end{axis}\end{tikzpicture}
\end{center}
\end{solution}
\begin{exercise}\label{cw:bank-5B-2}原题5B-2. 对 $x'=y$, $y'=x$ 回答5B-1的(a)消元作图、(b)参数解与漏轨道、(c)曲线覆盖、(d)时间方向反转等四问；(d)改为 $x'=-y$, $y'=-x$.\end{exercise}
\begin{solution}（官方译审）(a) $y\,dy=x\,dx$，故 $y^2-x^2=C$. (b) 解为 $\mathbf x=Ae^t(1,1)^T+Be^{-t}(1,-1)^T$：沿 $y=x$ 离开原点，沿 $y=-x$ 进入原点，其他轨道为双曲线分支. 消元在 $y=0$ 的瞬间不能作除法，但曲线可连续穿过该点；原点平衡轨道须另外补上. (c) 非零 $C$ 的完整双曲线有两条分支，各由一条轨道覆盖；$C=0$ 是两条交叉直线，共四条半直线轨道及原点五条. 若把“曲线”只指一个连通分支，则非退化分支只需一条轨道. (d) 各非平衡箭头反向，稳定与不稳定特征直线互换，仍为鞍点.
\begin{center}
\begin{tikzpicture}\begin{axis}[width=.47\linewidth,height=5.2cm,axis lines=middle,xlabel={$x$},ylabel={$y$},xmin=-2,xmax=2,ymin=-2,ymax=2,title={双曲线与鞍点},tick label style={font=\footnotesize},clip=true]
\addplot[structurecolor,thin,no marks] coordinates {(-1.87617,2.00000) (-1.80636,1.93467) (-1.73883,1.87177) (-1.67349,1.81123) (-1.61025,1.75297) (-1.54905,1.69692) (-1.48979,1.64301) (-1.43241,1.59117) (-1.37684,1.54133) (-1.32300,1.49343) (-1.27082,1.44741) (-1.22025,1.40321) (-1.17121,1.36079) (-1.12365,1.32007) (-1.07750,1.28102) (-1.03271,1.24358) (-0.98922,1.20771) (-0.94698,1.17336) (-0.90593,1.14049) (-0.86602,1.10905) (-0.82720,1.07901) (-0.78942,1.05033) (-0.75264,1.02297) (-0.71680,0.99690) (-0.68187,0.97208) (-0.64780,0.94849) (-0.61454,0.92610) (-0.58206,0.90487) (-0.55030,0.88478) (-0.51925,0.86580) (-0.48884,0.84792) (-0.45906,0.83110) (-0.42985,0.81533) (-0.40118,0.80059) (-0.37301,0.78685) (-0.34532,0.77411) (-0.31806,0.76234) (-0.29120,0.75153) (-0.26471,0.74167) (-0.23855,0.73274) (-0.21270,0.72473) (-0.18711,0.71764) (-0.16175,0.71145) (-0.13660,0.70616) (-0.11162,0.70175) (-0.08678,0.69823) (-0.06206,0.69559) (-0.03741,0.69383) (-0.01280,0.69294) (0.01178,0.69292) (0.03639,0.69378) (0.06103,0.69550) (0.08576,0.69811) (0.11059,0.70159) (0.13556,0.70596) (0.16071,0.71121) (0.18605,0.71737) (0.21163,0.72442) (0.23748,0.73239) (0.26362,0.74128) (0.29010,0.75110) (0.31694,0.76187) (0.34418,0.77360) (0.37186,0.78631) (0.40000,0.80000)};
\draw[-{Latex},structurecolor] (axis cs:-0.82720,1.07901) -- (axis cs:-0.78942,1.05033);
\draw[-{Latex},structurecolor] (axis cs:-0.13660,0.70616) -- (axis cs:-0.11162,0.70175);
\addplot[structurecolor,thin,no marks] coordinates {(0.40000,0.80000) (0.41472,0.80746) (0.42959,0.81520) (0.44460,0.82320) (0.45975,0.83149) (0.47506,0.84005) (0.49054,0.84890) (0.50617,0.85803) (0.52198,0.86744) (0.53796,0.87715) (0.55412,0.88716) (0.57047,0.89746) (0.58700,0.90806) (0.60374,0.91897) (0.62068,0.93018) (0.63783,0.94171) (0.65519,0.95356) (0.67277,0.96572) (0.69057,0.97821) (0.70861,0.99103) (0.72689,1.00417) (0.74541,1.01766) (0.76418,1.03149) (0.78321,1.04566) (0.80250,1.06019) (0.82205,1.07507) (0.84189,1.09031) (0.86201,1.10592) (0.88241,1.12190) (0.90312,1.13825) (0.92412,1.15499) (0.94544,1.17212) (0.96707,1.18964) (0.98903,1.20755) (1.01132,1.22588) (1.03395,1.24461) (1.05693,1.26376) (1.08026,1.28334) (1.10396,1.30335) (1.12802,1.32379) (1.15246,1.34468) (1.17730,1.36602) (1.20252,1.38782) (1.22815,1.41009) (1.25419,1.43283) (1.28065,1.45605) (1.30755,1.47976) (1.33488,1.50396) (1.36266,1.52867) (1.39089,1.55389) (1.41960,1.57964) (1.44878,1.60591) (1.47844,1.63273) (1.50860,1.66009) (1.53927,1.68801) (1.57046,1.71649) (1.60217,1.74555) (1.63442,1.77520) (1.66722,1.80544) (1.70058,1.83629) (1.73451,1.86776) (1.76902,1.89985) (1.80413,1.93258) (1.83984,1.96596) (1.87617,2.00000)};
\draw[-{Latex},structurecolor] (axis cs:0.72689,1.00417) -- (axis cs:0.74541,1.01766);
\draw[-{Latex},structurecolor] (axis cs:1.22815,1.41009) -- (axis cs:1.25419,1.43283);
\addplot[structurecolor,thin,no marks] coordinates {(2.00000,-1.87617) (1.93467,-1.80636) (1.87177,-1.73883) (1.81123,-1.67349) (1.75297,-1.61025) (1.69692,-1.54905) (1.64301,-1.48979) (1.59117,-1.43241) (1.54133,-1.37684) (1.49343,-1.32300) (1.44741,-1.27082) (1.40321,-1.22025) (1.36079,-1.17121) (1.32007,-1.12365) (1.28102,-1.07750) (1.24358,-1.03271) (1.20771,-0.98922) (1.17336,-0.94698) (1.14049,-0.90593) (1.10905,-0.86602) (1.07901,-0.82720) (1.05033,-0.78942) (1.02297,-0.75264) (0.99690,-0.71680) (0.97208,-0.68187) (0.94849,-0.64780) (0.92610,-0.61454) (0.90487,-0.58206) (0.88478,-0.55030) (0.86580,-0.51925) (0.84792,-0.48884) (0.83110,-0.45906) (0.81533,-0.42985) (0.80059,-0.40118) (0.78685,-0.37301) (0.77411,-0.34532) (0.76234,-0.31806) (0.75153,-0.29120) (0.74167,-0.26471) (0.73274,-0.23855) (0.72473,-0.21270) (0.71764,-0.18711) (0.71145,-0.16175) (0.70616,-0.13660) (0.70175,-0.11162) (0.69823,-0.08678) (0.69559,-0.06206) (0.69383,-0.03741) (0.69294,-0.01280) (0.69292,0.01178) (0.69378,0.03639) (0.69550,0.06103) (0.69811,0.08576) (0.70159,0.11059) (0.70596,0.13556) (0.71121,0.16071) (0.71737,0.18605) (0.72442,0.21163) (0.73239,0.23748) (0.74128,0.26362) (0.75110,0.29010) (0.76187,0.31694) (0.77360,0.34418) (0.78631,0.37186) (0.80000,0.40000)};
\draw[-{Latex},structurecolor] (axis cs:1.07901,-0.82720) -- (axis cs:1.05033,-0.78942);
\draw[-{Latex},structurecolor] (axis cs:0.70616,-0.13660) -- (axis cs:0.70175,-0.11162);
\addplot[structurecolor,thin,no marks] coordinates {(0.80000,0.40000) (0.80746,0.41472) (0.81520,0.42959) (0.82320,0.44460) (0.83149,0.45975) (0.84005,0.47506) (0.84890,0.49054) (0.85803,0.50617) (0.86744,0.52198) (0.87715,0.53796) (0.88716,0.55412) (0.89746,0.57047) (0.90806,0.58700) (0.91897,0.60374) (0.93018,0.62068) (0.94171,0.63783) (0.95356,0.65519) (0.96572,0.67277) (0.97821,0.69057) (0.99103,0.70861) (1.00417,0.72689) (1.01766,0.74541) (1.03149,0.76418) (1.04566,0.78321) (1.06019,0.80250) (1.07507,0.82205) (1.09031,0.84189) (1.10592,0.86201) (1.12190,0.88241) (1.13825,0.90312) (1.15499,0.92412) (1.17212,0.94544) (1.18964,0.96707) (1.20755,0.98903) (1.22588,1.01132) (1.24461,1.03395) (1.26376,1.05693) (1.28334,1.08026) (1.30335,1.10396) (1.32379,1.12802) (1.34468,1.15246) (1.36602,1.17730) (1.38782,1.20252) (1.41009,1.22815) (1.43283,1.25419) (1.45605,1.28065) (1.47976,1.30755) (1.50396,1.33488) (1.52867,1.36266) (1.55389,1.39089) (1.57964,1.41960) (1.60591,1.44878) (1.63273,1.47844) (1.66009,1.50860) (1.68801,1.53927) (1.71649,1.57046) (1.74555,1.60217) (1.77520,1.63442) (1.80544,1.66722) (1.83629,1.70058) (1.86776,1.73451) (1.89985,1.76902) (1.93258,1.80413) (1.96596,1.83984) (2.00000,1.87617)};
\draw[-{Latex},structurecolor] (axis cs:1.00417,0.72689) -- (axis cs:1.01766,0.74541);
\draw[-{Latex},structurecolor] (axis cs:1.41009,1.22815) -- (axis cs:1.43283,1.25419);
\addplot[structurecolor,thin,no marks] coordinates {(-1.87617,2.00000) (-1.83984,1.96596) (-1.80413,1.93258) (-1.76902,1.89985) (-1.73451,1.86776) (-1.70058,1.83629) (-1.66722,1.80544) (-1.63442,1.77520) (-1.60217,1.74555) (-1.57046,1.71649) (-1.53927,1.68801) (-1.50860,1.66009) (-1.47844,1.63273) (-1.44878,1.60591) (-1.41960,1.57964) (-1.39089,1.55389) (-1.36266,1.52867) (-1.33488,1.50396) (-1.30755,1.47976) (-1.28065,1.45605) (-1.25419,1.43283) (-1.22815,1.41009) (-1.20252,1.38782) (-1.17730,1.36602) (-1.15246,1.34468) (-1.12802,1.32379) (-1.10396,1.30335) (-1.08026,1.28334) (-1.05693,1.26376) (-1.03395,1.24461) (-1.01132,1.22588) (-0.98903,1.20755) (-0.96707,1.18964) (-0.94544,1.17212) (-0.92412,1.15499) (-0.90312,1.13825) (-0.88241,1.12190) (-0.86201,1.10592) (-0.84189,1.09031) (-0.82205,1.07507) (-0.80250,1.06019) (-0.78321,1.04566) (-0.76418,1.03149) (-0.74541,1.01766) (-0.72689,1.00417) (-0.70861,0.99103) (-0.69057,0.97821) (-0.67277,0.96572) (-0.65519,0.95356) (-0.63783,0.94171) (-0.62068,0.93018) (-0.60374,0.91897) (-0.58700,0.90806) (-0.57047,0.89746) (-0.55412,0.88716) (-0.53796,0.87715) (-0.52198,0.86744) (-0.50617,0.85803) (-0.49054,0.84890) (-0.47506,0.84005) (-0.45975,0.83149) (-0.44460,0.82320) (-0.42959,0.81520) (-0.41472,0.80746) (-0.40000,0.80000)};
\draw[-{Latex},structurecolor] (axis cs:-1.25419,1.43283) -- (axis cs:-1.22815,1.41009);
\draw[-{Latex},structurecolor] (axis cs:-0.74541,1.01766) -- (axis cs:-0.72689,1.00417);
\addplot[structurecolor,thin,no marks] coordinates {(-0.40000,0.80000) (-0.37186,0.78631) (-0.34418,0.77360) (-0.31694,0.76187) (-0.29010,0.75110) (-0.26362,0.74128) (-0.23748,0.73239) (-0.21163,0.72442) (-0.18605,0.71737) (-0.16071,0.71121) (-0.13556,0.70596) (-0.11059,0.70159) (-0.08576,0.69811) (-0.06103,0.69550) (-0.03639,0.69378) (-0.01178,0.69292) (0.01280,0.69294) (0.03741,0.69383) (0.06206,0.69559) (0.08678,0.69823) (0.11162,0.70175) (0.13660,0.70616) (0.16175,0.71145) (0.18711,0.71764) (0.21270,0.72473) (0.23855,0.73274) (0.26471,0.74167) (0.29120,0.75153) (0.31806,0.76234) (0.34532,0.77411) (0.37301,0.78685) (0.40118,0.80059) (0.42985,0.81533) (0.45906,0.83110) (0.48884,0.84792) (0.51925,0.86580) (0.55030,0.88478) (0.58206,0.90487) (0.61454,0.92610) (0.64780,0.94849) (0.68187,0.97208) (0.71680,0.99690) (0.75264,1.02297) (0.78942,1.05033) (0.82720,1.07901) (0.86602,1.10905) (0.90593,1.14049) (0.94698,1.17336) (0.98922,1.20771) (1.03271,1.24358) (1.07750,1.28102) (1.12365,1.32007) (1.17121,1.36079) (1.22025,1.40321) (1.27082,1.44741) (1.32300,1.49343) (1.37684,1.54133) (1.43241,1.59117) (1.48979,1.64301) (1.54905,1.69692) (1.61025,1.75297) (1.67349,1.81123) (1.73883,1.87177) (1.80636,1.93467) (1.87617,2.00000)};
\draw[-{Latex},structurecolor] (axis cs:0.11162,0.70175) -- (axis cs:0.13660,0.70616);
\draw[-{Latex},structurecolor] (axis cs:0.78942,1.05033) -- (axis cs:0.82720,1.07901);
\addplot[structurecolor,thin,no marks] coordinates {(-2.00000,1.87617) (-1.96596,1.83984) (-1.93258,1.80413) (-1.89985,1.76902) (-1.86776,1.73451) (-1.83629,1.70058) (-1.80544,1.66722) (-1.77520,1.63442) (-1.74555,1.60217) (-1.71649,1.57046) (-1.68801,1.53927) (-1.66009,1.50860) (-1.63273,1.47844) (-1.60591,1.44878) (-1.57964,1.41960) (-1.55389,1.39089) (-1.52867,1.36266) (-1.50396,1.33488) (-1.47976,1.30755) (-1.45605,1.28065) (-1.43283,1.25419) (-1.41009,1.22815) (-1.38782,1.20252) (-1.36602,1.17730) (-1.34468,1.15246) (-1.32379,1.12802) (-1.30335,1.10396) (-1.28334,1.08026) (-1.26376,1.05693) (-1.24461,1.03395) (-1.22588,1.01132) (-1.20755,0.98903) (-1.18964,0.96707) (-1.17212,0.94544) (-1.15499,0.92412) (-1.13825,0.90312) (-1.12190,0.88241) (-1.10592,0.86201) (-1.09031,0.84189) (-1.07507,0.82205) (-1.06019,0.80250) (-1.04566,0.78321) (-1.03149,0.76418) (-1.01766,0.74541) (-1.00417,0.72689) (-0.99103,0.70861) (-0.97821,0.69057) (-0.96572,0.67277) (-0.95356,0.65519) (-0.94171,0.63783) (-0.93018,0.62068) (-0.91897,0.60374) (-0.90806,0.58700) (-0.89746,0.57047) (-0.88716,0.55412) (-0.87715,0.53796) (-0.86744,0.52198) (-0.85803,0.50617) (-0.84890,0.49054) (-0.84005,0.47506) (-0.83149,0.45975) (-0.82320,0.44460) (-0.81520,0.42959) (-0.80746,0.41472) (-0.80000,0.40000)};
\draw[-{Latex},structurecolor] (axis cs:-1.43283,1.25419) -- (axis cs:-1.41009,1.22815);
\draw[-{Latex},structurecolor] (axis cs:-1.01766,0.74541) -- (axis cs:-1.00417,0.72689);
\addplot[structurecolor,thin,no marks] coordinates {(-0.80000,0.40000) (-0.78631,0.37186) (-0.77360,0.34418) (-0.76187,0.31694) (-0.75110,0.29010) (-0.74128,0.26362) (-0.73239,0.23748) (-0.72442,0.21163) (-0.71737,0.18605) (-0.71121,0.16071) (-0.70596,0.13556) (-0.70159,0.11059) (-0.69811,0.08576) (-0.69550,0.06103) (-0.69378,0.03639) (-0.69292,0.01178) (-0.69294,-0.01280) (-0.69383,-0.03741) (-0.69559,-0.06206) (-0.69823,-0.08678) (-0.70175,-0.11162) (-0.70616,-0.13660) (-0.71145,-0.16175) (-0.71764,-0.18711) (-0.72473,-0.21270) (-0.73274,-0.23855) (-0.74167,-0.26471) (-0.75153,-0.29120) (-0.76234,-0.31806) (-0.77411,-0.34532) (-0.78685,-0.37301) (-0.80059,-0.40118) (-0.81533,-0.42985) (-0.83110,-0.45906) (-0.84792,-0.48884) (-0.86580,-0.51925) (-0.88478,-0.55030) (-0.90487,-0.58206) (-0.92610,-0.61454) (-0.94849,-0.64780) (-0.97208,-0.68187) (-0.99690,-0.71680) (-1.02297,-0.75264) (-1.05033,-0.78942) (-1.07901,-0.82720) (-1.10905,-0.86602) (-1.14049,-0.90593) (-1.17336,-0.94698) (-1.20771,-0.98922) (-1.24358,-1.03271) (-1.28102,-1.07750) (-1.32007,-1.12365) (-1.36079,-1.17121) (-1.40321,-1.22025) (-1.44741,-1.27082) (-1.49343,-1.32300) (-1.54133,-1.37684) (-1.59117,-1.43241) (-1.64301,-1.48979) (-1.69692,-1.54905) (-1.75297,-1.61025) (-1.81123,-1.67349) (-1.87177,-1.73883) (-1.93467,-1.80636) (-2.00000,-1.87617)};
\draw[-{Latex},structurecolor] (axis cs:-0.70175,-0.11162) -- (axis cs:-0.70616,-0.13660);
\draw[-{Latex},structurecolor] (axis cs:-1.05033,-0.78942) -- (axis cs:-1.07901,-0.82720);
\addplot[structurecolor,thin,no marks] coordinates {(1.87617,-2.00000) (1.83984,-1.96596) (1.80413,-1.93258) (1.76902,-1.89985) (1.73451,-1.86776) (1.70058,-1.83629) (1.66722,-1.80544) (1.63442,-1.77520) (1.60217,-1.74555) (1.57046,-1.71649) (1.53927,-1.68801) (1.50860,-1.66009) (1.47844,-1.63273) (1.44878,-1.60591) (1.41960,-1.57964) (1.39089,-1.55389) (1.36266,-1.52867) (1.33488,-1.50396) (1.30755,-1.47976) (1.28065,-1.45605) (1.25419,-1.43283) (1.22815,-1.41009) (1.20252,-1.38782) (1.17730,-1.36602) (1.15246,-1.34468) (1.12802,-1.32379) (1.10396,-1.30335) (1.08026,-1.28334) (1.05693,-1.26376) (1.03395,-1.24461) (1.01132,-1.22588) (0.98903,-1.20755) (0.96707,-1.18964) (0.94544,-1.17212) (0.92412,-1.15499) (0.90312,-1.13825) (0.88241,-1.12190) (0.86201,-1.10592) (0.84189,-1.09031) (0.82205,-1.07507) (0.80250,-1.06019) (0.78321,-1.04566) (0.76418,-1.03149) (0.74541,-1.01766) (0.72689,-1.00417) (0.70861,-0.99103) (0.69057,-0.97821) (0.67277,-0.96572) (0.65519,-0.95356) (0.63783,-0.94171) (0.62068,-0.93018) (0.60374,-0.91897) (0.58700,-0.90806) (0.57047,-0.89746) (0.55412,-0.88716) (0.53796,-0.87715) (0.52198,-0.86744) (0.50617,-0.85803) (0.49054,-0.84890) (0.47506,-0.84005) (0.45975,-0.83149) (0.44460,-0.82320) (0.42959,-0.81520) (0.41472,-0.80746) (0.40000,-0.80000)};
\draw[-{Latex},structurecolor] (axis cs:1.25419,-1.43283) -- (axis cs:1.22815,-1.41009);
\draw[-{Latex},structurecolor] (axis cs:0.74541,-1.01766) -- (axis cs:0.72689,-1.00417);
\addplot[structurecolor,thin,no marks] coordinates {(0.40000,-0.80000) (0.37186,-0.78631) (0.34418,-0.77360) (0.31694,-0.76187) (0.29010,-0.75110) (0.26362,-0.74128) (0.23748,-0.73239) (0.21163,-0.72442) (0.18605,-0.71737) (0.16071,-0.71121) (0.13556,-0.70596) (0.11059,-0.70159) (0.08576,-0.69811) (0.06103,-0.69550) (0.03639,-0.69378) (0.01178,-0.69292) (-0.01280,-0.69294) (-0.03741,-0.69383) (-0.06206,-0.69559) (-0.08678,-0.69823) (-0.11162,-0.70175) (-0.13660,-0.70616) (-0.16175,-0.71145) (-0.18711,-0.71764) (-0.21270,-0.72473) (-0.23855,-0.73274) (-0.26471,-0.74167) (-0.29120,-0.75153) (-0.31806,-0.76234) (-0.34532,-0.77411) (-0.37301,-0.78685) (-0.40118,-0.80059) (-0.42985,-0.81533) (-0.45906,-0.83110) (-0.48884,-0.84792) (-0.51925,-0.86580) (-0.55030,-0.88478) (-0.58206,-0.90487) (-0.61454,-0.92610) (-0.64780,-0.94849) (-0.68187,-0.97208) (-0.71680,-0.99690) (-0.75264,-1.02297) (-0.78942,-1.05033) (-0.82720,-1.07901) (-0.86602,-1.10905) (-0.90593,-1.14049) (-0.94698,-1.17336) (-0.98922,-1.20771) (-1.03271,-1.24358) (-1.07750,-1.28102) (-1.12365,-1.32007) (-1.17121,-1.36079) (-1.22025,-1.40321) (-1.27082,-1.44741) (-1.32300,-1.49343) (-1.37684,-1.54133) (-1.43241,-1.59117) (-1.48979,-1.64301) (-1.54905,-1.69692) (-1.61025,-1.75297) (-1.67349,-1.81123) (-1.73883,-1.87177) (-1.80636,-1.93467) (-1.87617,-2.00000)};
\draw[-{Latex},structurecolor] (axis cs:-0.11162,-0.70175) -- (axis cs:-0.13660,-0.70616);
\draw[-{Latex},structurecolor] (axis cs:-0.78942,-1.05033) -- (axis cs:-0.82720,-1.07901);
\addplot[structurecolor,thin,no marks] coordinates {(2.00000,-1.87617) (1.96596,-1.83984) (1.93258,-1.80413) (1.89985,-1.76902) (1.86776,-1.73451) (1.83629,-1.70058) (1.80544,-1.66722) (1.77520,-1.63442) (1.74555,-1.60217) (1.71649,-1.57046) (1.68801,-1.53927) (1.66009,-1.50860) (1.63273,-1.47844) (1.60591,-1.44878) (1.57964,-1.41960) (1.55389,-1.39089) (1.52867,-1.36266) (1.50396,-1.33488) (1.47976,-1.30755) (1.45605,-1.28065) (1.43283,-1.25419) (1.41009,-1.22815) (1.38782,-1.20252) (1.36602,-1.17730) (1.34468,-1.15246) (1.32379,-1.12802) (1.30335,-1.10396) (1.28334,-1.08026) (1.26376,-1.05693) (1.24461,-1.03395) (1.22588,-1.01132) (1.20755,-0.98903) (1.18964,-0.96707) (1.17212,-0.94544) (1.15499,-0.92412) (1.13825,-0.90312) (1.12190,-0.88241) (1.10592,-0.86201) (1.09031,-0.84189) (1.07507,-0.82205) (1.06019,-0.80250) (1.04566,-0.78321) (1.03149,-0.76418) (1.01766,-0.74541) (1.00417,-0.72689) (0.99103,-0.70861) (0.97821,-0.69057) (0.96572,-0.67277) (0.95356,-0.65519) (0.94171,-0.63783) (0.93018,-0.62068) (0.91897,-0.60374) (0.90806,-0.58700) (0.89746,-0.57047) (0.88716,-0.55412) (0.87715,-0.53796) (0.86744,-0.52198) (0.85803,-0.50617) (0.84890,-0.49054) (0.84005,-0.47506) (0.83149,-0.45975) (0.82320,-0.44460) (0.81520,-0.42959) (0.80746,-0.41472) (0.80000,-0.40000)};
\draw[-{Latex},structurecolor] (axis cs:1.43283,-1.25419) -- (axis cs:1.41009,-1.22815);
\draw[-{Latex},structurecolor] (axis cs:1.01766,-0.74541) -- (axis cs:1.00417,-0.72689);
\addplot[structurecolor,thin,no marks] coordinates {(0.80000,-0.40000) (0.78631,-0.37186) (0.77360,-0.34418) (0.76187,-0.31694) (0.75110,-0.29010) (0.74128,-0.26362) (0.73239,-0.23748) (0.72442,-0.21163) (0.71737,-0.18605) (0.71121,-0.16071) (0.70596,-0.13556) (0.70159,-0.11059) (0.69811,-0.08576) (0.69550,-0.06103) (0.69378,-0.03639) (0.69292,-0.01178) (0.69294,0.01280) (0.69383,0.03741) (0.69559,0.06206) (0.69823,0.08678) (0.70175,0.11162) (0.70616,0.13660) (0.71145,0.16175) (0.71764,0.18711) (0.72473,0.21270) (0.73274,0.23855) (0.74167,0.26471) (0.75153,0.29120) (0.76234,0.31806) (0.77411,0.34532) (0.78685,0.37301) (0.80059,0.40118) (0.81533,0.42985) (0.83110,0.45906) (0.84792,0.48884) (0.86580,0.51925) (0.88478,0.55030) (0.90487,0.58206) (0.92610,0.61454) (0.94849,0.64780) (0.97208,0.68187) (0.99690,0.71680) (1.02297,0.75264) (1.05033,0.78942) (1.07901,0.82720) (1.10905,0.86602) (1.14049,0.90593) (1.17336,0.94698) (1.20771,0.98922) (1.24358,1.03271) (1.28102,1.07750) (1.32007,1.12365) (1.36079,1.17121) (1.40321,1.22025) (1.44741,1.27082) (1.49343,1.32300) (1.54133,1.37684) (1.59117,1.43241) (1.64301,1.48979) (1.69692,1.54905) (1.75297,1.61025) (1.81123,1.67349) (1.87177,1.73883) (1.93467,1.80636) (2.00000,1.87617)};
\draw[-{Latex},structurecolor] (axis cs:0.70175,0.11162) -- (axis cs:0.70616,0.13660);
\draw[-{Latex},structurecolor] (axis cs:1.05033,0.78942) -- (axis cs:1.07901,0.82720);
\addplot[structurecolor,thin,no marks] coordinates {(1.87617,-2.00000) (1.80636,-1.93467) (1.73883,-1.87177) (1.67349,-1.81123) (1.61025,-1.75297) (1.54905,-1.69692) (1.48979,-1.64301) (1.43241,-1.59117) (1.37684,-1.54133) (1.32300,-1.49343) (1.27082,-1.44741) (1.22025,-1.40321) (1.17121,-1.36079) (1.12365,-1.32007) (1.07750,-1.28102) (1.03271,-1.24358) (0.98922,-1.20771) (0.94698,-1.17336) (0.90593,-1.14049) (0.86602,-1.10905) (0.82720,-1.07901) (0.78942,-1.05033) (0.75264,-1.02297) (0.71680,-0.99690) (0.68187,-0.97208) (0.64780,-0.94849) (0.61454,-0.92610) (0.58206,-0.90487) (0.55030,-0.88478) (0.51925,-0.86580) (0.48884,-0.84792) (0.45906,-0.83110) (0.42985,-0.81533) (0.40118,-0.80059) (0.37301,-0.78685) (0.34532,-0.77411) (0.31806,-0.76234) (0.29120,-0.75153) (0.26471,-0.74167) (0.23855,-0.73274) (0.21270,-0.72473) (0.18711,-0.71764) (0.16175,-0.71145) (0.13660,-0.70616) (0.11162,-0.70175) (0.08678,-0.69823) (0.06206,-0.69559) (0.03741,-0.69383) (0.01280,-0.69294) (-0.01178,-0.69292) (-0.03639,-0.69378) (-0.06103,-0.69550) (-0.08576,-0.69811) (-0.11059,-0.70159) (-0.13556,-0.70596) (-0.16071,-0.71121) (-0.18605,-0.71737) (-0.21163,-0.72442) (-0.23748,-0.73239) (-0.26362,-0.74128) (-0.29010,-0.75110) (-0.31694,-0.76187) (-0.34418,-0.77360) (-0.37186,-0.78631) (-0.40000,-0.80000)};
\draw[-{Latex},structurecolor] (axis cs:0.82720,-1.07901) -- (axis cs:0.78942,-1.05033);
\draw[-{Latex},structurecolor] (axis cs:0.13660,-0.70616) -- (axis cs:0.11162,-0.70175);
\addplot[structurecolor,thin,no marks] coordinates {(-0.40000,-0.80000) (-0.41472,-0.80746) (-0.42959,-0.81520) (-0.44460,-0.82320) (-0.45975,-0.83149) (-0.47506,-0.84005) (-0.49054,-0.84890) (-0.50617,-0.85803) (-0.52198,-0.86744) (-0.53796,-0.87715) (-0.55412,-0.88716) (-0.57047,-0.89746) (-0.58700,-0.90806) (-0.60374,-0.91897) (-0.62068,-0.93018) (-0.63783,-0.94171) (-0.65519,-0.95356) (-0.67277,-0.96572) (-0.69057,-0.97821) (-0.70861,-0.99103) (-0.72689,-1.00417) (-0.74541,-1.01766) (-0.76418,-1.03149) (-0.78321,-1.04566) (-0.80250,-1.06019) (-0.82205,-1.07507) (-0.84189,-1.09031) (-0.86201,-1.10592) (-0.88241,-1.12190) (-0.90312,-1.13825) (-0.92412,-1.15499) (-0.94544,-1.17212) (-0.96707,-1.18964) (-0.98903,-1.20755) (-1.01132,-1.22588) (-1.03395,-1.24461) (-1.05693,-1.26376) (-1.08026,-1.28334) (-1.10396,-1.30335) (-1.12802,-1.32379) (-1.15246,-1.34468) (-1.17730,-1.36602) (-1.20252,-1.38782) (-1.22815,-1.41009) (-1.25419,-1.43283) (-1.28065,-1.45605) (-1.30755,-1.47976) (-1.33488,-1.50396) (-1.36266,-1.52867) (-1.39089,-1.55389) (-1.41960,-1.57964) (-1.44878,-1.60591) (-1.47844,-1.63273) (-1.50860,-1.66009) (-1.53927,-1.68801) (-1.57046,-1.71649) (-1.60217,-1.74555) (-1.63442,-1.77520) (-1.66722,-1.80544) (-1.70058,-1.83629) (-1.73451,-1.86776) (-1.76902,-1.89985) (-1.80413,-1.93258) (-1.83984,-1.96596) (-1.87617,-2.00000)};
\draw[-{Latex},structurecolor] (axis cs:-0.72689,-1.00417) -- (axis cs:-0.74541,-1.01766);
\draw[-{Latex},structurecolor] (axis cs:-1.22815,-1.41009) -- (axis cs:-1.25419,-1.43283);
\addplot[structurecolor,thin,no marks] coordinates {(-2.00000,1.87617) (-1.93467,1.80636) (-1.87177,1.73883) (-1.81123,1.67349) (-1.75297,1.61025) (-1.69692,1.54905) (-1.64301,1.48979) (-1.59117,1.43241) (-1.54133,1.37684) (-1.49343,1.32300) (-1.44741,1.27082) (-1.40321,1.22025) (-1.36079,1.17121) (-1.32007,1.12365) (-1.28102,1.07750) (-1.24358,1.03271) (-1.20771,0.98922) (-1.17336,0.94698) (-1.14049,0.90593) (-1.10905,0.86602) (-1.07901,0.82720) (-1.05033,0.78942) (-1.02297,0.75264) (-0.99690,0.71680) (-0.97208,0.68187) (-0.94849,0.64780) (-0.92610,0.61454) (-0.90487,0.58206) (-0.88478,0.55030) (-0.86580,0.51925) (-0.84792,0.48884) (-0.83110,0.45906) (-0.81533,0.42985) (-0.80059,0.40118) (-0.78685,0.37301) (-0.77411,0.34532) (-0.76234,0.31806) (-0.75153,0.29120) (-0.74167,0.26471) (-0.73274,0.23855) (-0.72473,0.21270) (-0.71764,0.18711) (-0.71145,0.16175) (-0.70616,0.13660) (-0.70175,0.11162) (-0.69823,0.08678) (-0.69559,0.06206) (-0.69383,0.03741) (-0.69294,0.01280) (-0.69292,-0.01178) (-0.69378,-0.03639) (-0.69550,-0.06103) (-0.69811,-0.08576) (-0.70159,-0.11059) (-0.70596,-0.13556) (-0.71121,-0.16071) (-0.71737,-0.18605) (-0.72442,-0.21163) (-0.73239,-0.23748) (-0.74128,-0.26362) (-0.75110,-0.29010) (-0.76187,-0.31694) (-0.77360,-0.34418) (-0.78631,-0.37186) (-0.80000,-0.40000)};
\draw[-{Latex},structurecolor] (axis cs:-1.07901,0.82720) -- (axis cs:-1.05033,0.78942);
\draw[-{Latex},structurecolor] (axis cs:-0.70616,0.13660) -- (axis cs:-0.70175,0.11162);
\addplot[structurecolor,thin,no marks] coordinates {(-0.80000,-0.40000) (-0.80746,-0.41472) (-0.81520,-0.42959) (-0.82320,-0.44460) (-0.83149,-0.45975) (-0.84005,-0.47506) (-0.84890,-0.49054) (-0.85803,-0.50617) (-0.86744,-0.52198) (-0.87715,-0.53796) (-0.88716,-0.55412) (-0.89746,-0.57047) (-0.90806,-0.58700) (-0.91897,-0.60374) (-0.93018,-0.62068) (-0.94171,-0.63783) (-0.95356,-0.65519) (-0.96572,-0.67277) (-0.97821,-0.69057) (-0.99103,-0.70861) (-1.00417,-0.72689) (-1.01766,-0.74541) (-1.03149,-0.76418) (-1.04566,-0.78321) (-1.06019,-0.80250) (-1.07507,-0.82205) (-1.09031,-0.84189) (-1.10592,-0.86201) (-1.12190,-0.88241) (-1.13825,-0.90312) (-1.15499,-0.92412) (-1.17212,-0.94544) (-1.18964,-0.96707) (-1.20755,-0.98903) (-1.22588,-1.01132) (-1.24461,-1.03395) (-1.26376,-1.05693) (-1.28334,-1.08026) (-1.30335,-1.10396) (-1.32379,-1.12802) (-1.34468,-1.15246) (-1.36602,-1.17730) (-1.38782,-1.20252) (-1.41009,-1.22815) (-1.43283,-1.25419) (-1.45605,-1.28065) (-1.47976,-1.30755) (-1.50396,-1.33488) (-1.52867,-1.36266) (-1.55389,-1.39089) (-1.57964,-1.41960) (-1.60591,-1.44878) (-1.63273,-1.47844) (-1.66009,-1.50860) (-1.68801,-1.53927) (-1.71649,-1.57046) (-1.74555,-1.60217) (-1.77520,-1.63442) (-1.80544,-1.66722) (-1.83629,-1.70058) (-1.86776,-1.73451) (-1.89985,-1.76902) (-1.93258,-1.80413) (-1.96596,-1.83984) (-2.00000,-1.87617)};
\draw[-{Latex},structurecolor] (axis cs:-1.00417,-0.72689) -- (axis cs:-1.01766,-0.74541);
\draw[-{Latex},structurecolor] (axis cs:-1.41009,-1.22815) -- (axis cs:-1.43283,-1.25419);
\addplot[structurecolor,thin,no marks] coordinates {(0.00996,0.00996) (0.01044,0.01044) (0.01094,0.01094) (0.01146,0.01146) (0.01201,0.01201) (0.01259,0.01259) (0.01319,0.01319) (0.01382,0.01382) (0.01449,0.01449) (0.01518,0.01518) (0.01591,0.01591) (0.01668,0.01668) (0.01748,0.01748) (0.01831,0.01831) (0.01919,0.01919) (0.02011,0.02011) (0.02108,0.02108) (0.02209,0.02209) (0.02315,0.02315) (0.02426,0.02426) (0.02543,0.02543) (0.02665,0.02665) (0.02793,0.02793) (0.02927,0.02927) (0.03067,0.03067) (0.03214,0.03214) (0.03369,0.03369) (0.03530,0.03530) (0.03700,0.03700) (0.03877,0.03877) (0.04063,0.04063) (0.04258,0.04258) (0.04463,0.04463) (0.04677,0.04677) (0.04901,0.04901) (0.05136,0.05136) (0.05383,0.05383) (0.05641,0.05641) (0.05912,0.05912) (0.06196,0.06196) (0.06493,0.06493) (0.06805,0.06805) (0.07131,0.07131) (0.07473,0.07473) (0.07832,0.07832) (0.08208,0.08208) (0.08602,0.08602) (0.09015,0.09015) (0.09447,0.09447) (0.09901,0.09901) (0.10376,0.10376) (0.10874,0.10874) (0.11396,0.11396) (0.11943,0.11943) (0.12516,0.12516) (0.13116,0.13116) (0.13746,0.13746) (0.14405,0.14405) (0.15097,0.15097) (0.15821,0.15821) (0.16581,0.16581) (0.17376,0.17376) (0.18210,0.18210) (0.19084,0.19084) (0.20000,0.20000)};
\draw[-{Latex},structurecolor] (axis cs:0.02543,0.02543) -- (axis cs:0.02665,0.02665);
\draw[-{Latex},structurecolor] (axis cs:0.07473,0.07473) -- (axis cs:0.07832,0.07832);
\addplot[structurecolor,thin,no marks] coordinates {(0.20000,0.20000) (0.20733,0.20733) (0.21492,0.21492) (0.22279,0.22279) (0.23096,0.23096) (0.23942,0.23942) (0.24819,0.24819) (0.25728,0.25728) (0.26670,0.26670) (0.27647,0.27647) (0.28660,0.28660) (0.29710,0.29710) (0.30799,0.30799) (0.31927,0.31927) (0.33096,0.33096) (0.34309,0.34309) (0.35566,0.35566) (0.36868,0.36868) (0.38219,0.38219) (0.39619,0.39619) (0.41071,0.41071) (0.42575,0.42575) (0.44135,0.44135) (0.45751,0.45751) (0.47427,0.47427) (0.49165,0.49165) (0.50966,0.50966) (0.52833,0.52833) (0.54768,0.54768) (0.56775,0.56775) (0.58855,0.58855) (0.61011,0.61011) (0.63246,0.63246) (0.65562,0.65562) (0.67964,0.67964) (0.70454,0.70454) (0.73035,0.73035) (0.75710,0.75710) (0.78484,0.78484) (0.81359,0.81359) (0.84339,0.84339) (0.87429,0.87429) (0.90632,0.90632) (0.93952,0.93952) (0.97394,0.97394) (1.00961,1.00961) (1.04660,1.04660) (1.08494,1.08494) (1.12468,1.12468) (1.16588,1.16588) (1.20859,1.20859) (1.25287,1.25287) (1.29876,1.29876) (1.34634,1.34634) (1.39566,1.39566) (1.44679,1.44679) (1.49979,1.49979) (1.55473,1.55473) (1.61168,1.61168) (1.67073,1.67073) (1.73193,1.73193) (1.79537,1.79537) (1.86114,1.86114) (1.92932,1.92932) (2.00000,2.00000)};
\draw[-{Latex},structurecolor] (axis cs:0.41071,0.41071) -- (axis cs:0.42575,0.42575);
\draw[-{Latex},structurecolor] (axis cs:0.93952,0.93952) -- (axis cs:0.97394,0.97394);
\addplot[structurecolor,thin,no marks] coordinates {(2.00000,-2.00000) (1.92932,-1.92932) (1.86114,-1.86114) (1.79537,-1.79537) (1.73193,-1.73193) (1.67073,-1.67073) (1.61168,-1.61168) (1.55473,-1.55473) (1.49979,-1.49979) (1.44679,-1.44679) (1.39566,-1.39566) (1.34634,-1.34634) (1.29876,-1.29876) (1.25287,-1.25287) (1.20859,-1.20859) (1.16588,-1.16588) (1.12468,-1.12468) (1.08494,-1.08494) (1.04660,-1.04660) (1.00961,-1.00961) (0.97394,-0.97394) (0.93952,-0.93952) (0.90632,-0.90632) (0.87429,-0.87429) (0.84339,-0.84339) (0.81359,-0.81359) (0.78484,-0.78484) (0.75710,-0.75710) (0.73035,-0.73035) (0.70454,-0.70454) (0.67964,-0.67964) (0.65562,-0.65562) (0.63246,-0.63246) (0.61011,-0.61011) (0.58855,-0.58855) (0.56775,-0.56775) (0.54768,-0.54768) (0.52833,-0.52833) (0.50966,-0.50966) (0.49165,-0.49165) (0.47427,-0.47427) (0.45751,-0.45751) (0.44135,-0.44135) (0.42575,-0.42575) (0.41071,-0.41071) (0.39619,-0.39619) (0.38219,-0.38219) (0.36868,-0.36868) (0.35566,-0.35566) (0.34309,-0.34309) (0.33096,-0.33096) (0.31927,-0.31927) (0.30799,-0.30799) (0.29710,-0.29710) (0.28660,-0.28660) (0.27647,-0.27647) (0.26670,-0.26670) (0.25728,-0.25728) (0.24819,-0.24819) (0.23942,-0.23942) (0.23096,-0.23096) (0.22279,-0.22279) (0.21492,-0.21492) (0.20733,-0.20733) (0.20000,-0.20000)};
\draw[-{Latex},structurecolor] (axis cs:0.97394,-0.97394) -- (axis cs:0.93952,-0.93952);
\draw[-{Latex},structurecolor] (axis cs:0.42575,-0.42575) -- (axis cs:0.41071,-0.41071);
\addplot[structurecolor,thin,no marks] coordinates {(0.20000,-0.20000) (0.19084,-0.19084) (0.18210,-0.18210) (0.17376,-0.17376) (0.16581,-0.16581) (0.15821,-0.15821) (0.15097,-0.15097) (0.14405,-0.14405) (0.13746,-0.13746) (0.13116,-0.13116) (0.12516,-0.12516) (0.11943,-0.11943) (0.11396,-0.11396) (0.10874,-0.10874) (0.10376,-0.10376) (0.09901,-0.09901) (0.09447,-0.09447) (0.09015,-0.09015) (0.08602,-0.08602) (0.08208,-0.08208) (0.07832,-0.07832) (0.07473,-0.07473) (0.07131,-0.07131) (0.06805,-0.06805) (0.06493,-0.06493) (0.06196,-0.06196) (0.05912,-0.05912) (0.05641,-0.05641) (0.05383,-0.05383) (0.05136,-0.05136) (0.04901,-0.04901) (0.04677,-0.04677) (0.04463,-0.04463) (0.04258,-0.04258) (0.04063,-0.04063) (0.03877,-0.03877) (0.03700,-0.03700) (0.03530,-0.03530) (0.03369,-0.03369) (0.03214,-0.03214) (0.03067,-0.03067) (0.02927,-0.02927) (0.02793,-0.02793) (0.02665,-0.02665) (0.02543,-0.02543) (0.02426,-0.02426) (0.02315,-0.02315) (0.02209,-0.02209) (0.02108,-0.02108) (0.02011,-0.02011) (0.01919,-0.01919) (0.01831,-0.01831) (0.01748,-0.01748) (0.01668,-0.01668) (0.01591,-0.01591) (0.01518,-0.01518) (0.01449,-0.01449) (0.01382,-0.01382) (0.01319,-0.01319) (0.01259,-0.01259) (0.01201,-0.01201) (0.01146,-0.01146) (0.01094,-0.01094) (0.01044,-0.01044) (0.00996,-0.00996)};
\draw[-{Latex},structurecolor] (axis cs:0.07832,-0.07832) -- (axis cs:0.07473,-0.07473);
\draw[-{Latex},structurecolor] (axis cs:0.02665,-0.02665) -- (axis cs:0.02543,-0.02543);
\addplot[only marks,mark=*,mark size=1.4pt] coordinates {(0,0)};
\end{axis}\end{tikzpicture}
\end{center}
\end{solution}
\begin{exercise}\label{cw:bank-5B-3}原题5B-3. 对 $x'=y$, $y'=-2x$ 回答5B-1(a)(b)：消元作图；用系统方向场为曲线标运动方向.\end{exercise}
\begin{solution}（官方译审）(a) $dy/dx=-2x/y$，得 $2x^2+y^2=C$. 对 $C>0$ 是椭圆，$C=0$ 是原点；无实曲线对应 $C<0$. (b) 在 $(x,0)$ 且 $x>0$ 时向量为 $(0,-2x)$，故沿椭圆顺时针运动. 参数化 $x=A\cos\sqrt2t+B\sin\sqrt2t$, $y=\sqrt2(-A\sin\sqrt2t+B\cos\sqrt2t)$ 可验证每条椭圆为一个周期轨道，原点为平衡轨道.
\begin{center}
\begin{tikzpicture}\begin{axis}[width=.47\linewidth,height=5.2cm,axis lines=middle,xlabel={$x$},ylabel={$y$},xmin=-1.5,xmax=1.5,ymin=-2,ymax=2,title={顺时针椭圆},tick label style={font=\footnotesize},clip=true]
\addplot[structurecolor,thin,no marks] coordinates {(0.40000,0.00000) (0.39807,-0.05545) (0.39231,-0.11036) (0.38278,-0.16421) (0.36955,-0.21648) (0.35277,-0.26666) (0.33259,-0.31428) (0.30920,-0.35887) (0.28284,-0.40000) (0.25376,-0.43728) (0.22223,-0.47035) (0.18856,-0.49889) (0.15307,-0.52263) (0.11611,-0.54133) (0.07804,-0.55482) (0.03921,-0.56296) (-0.00000,-0.56569) (-0.03921,-0.56296) (-0.07804,-0.55482) (-0.11611,-0.54133) (-0.15307,-0.52263) (-0.18856,-0.49889) (-0.22223,-0.47035) (-0.25376,-0.43728) (-0.28284,-0.40000) (-0.30920,-0.35887) (-0.33259,-0.31428) (-0.35277,-0.26666) (-0.36955,-0.21648) (-0.38278,-0.16421) (-0.39231,-0.11036) (-0.39807,-0.05545) (-0.40000,0.00000) (-0.39807,0.05545) (-0.39231,0.11036) (-0.38278,0.16421) (-0.36955,0.21648) (-0.35277,0.26666) (-0.33259,0.31428) (-0.30920,0.35887) (-0.28284,0.40000) (-0.25376,0.43728) (-0.22223,0.47035) (-0.18856,0.49889) (-0.15307,0.52263) (-0.11611,0.54133) (-0.07804,0.55482) (-0.03921,0.56296) (0.00000,0.56569) (0.03921,0.56296) (0.07804,0.55482) (0.11611,0.54133) (0.15307,0.52263) (0.18856,0.49889) (0.22223,0.47035) (0.25376,0.43728) (0.28284,0.40000) (0.30920,0.35887) (0.33259,0.31428) (0.35277,0.26666) (0.36955,0.21648) (0.38278,0.16421) (0.39231,0.11036) (0.39807,0.05545) (0.40000,-0.00000)};
\draw[-{Latex},structurecolor] (axis cs:-0.15307,-0.52263) -- (axis cs:-0.18856,-0.49889);
\draw[-{Latex},structurecolor] (axis cs:-0.18856,0.49889) -- (axis cs:-0.15307,0.52263);
\addplot[structurecolor,thin,no marks] coordinates {(0.80000,0.00000) (0.79615,-0.11089) (0.78463,-0.22072) (0.76555,-0.32842) (0.73910,-0.43296) (0.70554,-0.53332) (0.66518,-0.62856) (0.61841,-0.71773) (0.56569,-0.80000) (0.50751,-0.87456) (0.44446,-0.94070) (0.37712,-0.99778) (0.30615,-1.04525) (0.23223,-1.08265) (0.15607,-1.10963) (0.07841,-1.12592) (-0.00000,-1.13137) (-0.07841,-1.12592) (-0.15607,-1.10963) (-0.23223,-1.08265) (-0.30615,-1.04525) (-0.37712,-0.99778) (-0.44446,-0.94070) (-0.50751,-0.87456) (-0.56569,-0.80000) (-0.61841,-0.71773) (-0.66518,-0.62856) (-0.70554,-0.53332) (-0.73910,-0.43296) (-0.76555,-0.32842) (-0.78463,-0.22072) (-0.79615,-0.11089) (-0.80000,0.00000) (-0.79615,0.11089) (-0.78463,0.22072) (-0.76555,0.32842) (-0.73910,0.43296) (-0.70554,0.53332) (-0.66518,0.62856) (-0.61841,0.71773) (-0.56569,0.80000) (-0.50751,0.87456) (-0.44446,0.94070) (-0.37712,0.99778) (-0.30615,1.04525) (-0.23223,1.08265) (-0.15607,1.10963) (-0.07841,1.12592) (0.00000,1.13137) (0.07841,1.12592) (0.15607,1.10963) (0.23223,1.08265) (0.30615,1.04525) (0.37712,0.99778) (0.44446,0.94070) (0.50751,0.87456) (0.56569,0.80000) (0.61841,0.71773) (0.66518,0.62856) (0.70554,0.53332) (0.73910,0.43296) (0.76555,0.32842) (0.78463,0.22072) (0.79615,0.11089) (0.80000,-0.00000)};
\draw[-{Latex},structurecolor] (axis cs:-0.30615,-1.04525) -- (axis cs:-0.37712,-0.99778);
\draw[-{Latex},structurecolor] (axis cs:-0.37712,0.99778) -- (axis cs:-0.30615,1.04525);
\addplot[structurecolor,thin,no marks] coordinates {(1.20000,0.00000) (1.19422,-0.16634) (1.17694,-0.33108) (1.14833,-0.49263) (1.10866,-0.64944) (1.05831,-0.79999) (0.99776,-0.94283) (0.92761,-1.07660) (0.84853,-1.20000) (0.76127,-1.31184) (0.66668,-1.41105) (0.56568,-1.49667) (0.45922,-1.56788) (0.34834,-1.62398) (0.23411,-1.66445) (0.11762,-1.68888) (-0.00000,-1.69706) (-0.11762,-1.68888) (-0.23411,-1.66445) (-0.34834,-1.62398) (-0.45922,-1.56788) (-0.56568,-1.49667) (-0.66668,-1.41105) (-0.76127,-1.31184) (-0.84853,-1.20000) (-0.92761,-1.07660) (-0.99776,-0.94283) (-1.05831,-0.79999) (-1.10866,-0.64944) (-1.14833,-0.49263) (-1.17694,-0.33108) (-1.19422,-0.16634) (-1.20000,0.00000) (-1.19422,0.16634) (-1.17694,0.33108) (-1.14833,0.49263) (-1.10866,0.64944) (-1.05831,0.79999) (-0.99776,0.94283) (-0.92761,1.07660) (-0.84853,1.20000) (-0.76127,1.31184) (-0.66668,1.41105) (-0.56568,1.49667) (-0.45922,1.56788) (-0.34834,1.62398) (-0.23411,1.66445) (-0.11762,1.68888) (0.00000,1.69706) (0.11762,1.68888) (0.23411,1.66445) (0.34834,1.62398) (0.45922,1.56788) (0.56568,1.49667) (0.66668,1.41105) (0.76127,1.31184) (0.84853,1.20000) (0.92761,1.07660) (0.99776,0.94283) (1.05831,0.79999) (1.10866,0.64944) (1.14833,0.49263) (1.17694,0.33108) (1.19422,0.16634) (1.20000,-0.00000)};
\draw[-{Latex},structurecolor] (axis cs:-0.45922,-1.56788) -- (axis cs:-0.56568,-1.49667);
\draw[-{Latex},structurecolor] (axis cs:-0.56568,1.49667) -- (axis cs:-0.45922,1.56788);
\addplot[only marks,mark=*,mark size=1.4pt] coordinates {(0,0)};
\end{axis}\end{tikzpicture}
\end{center}
\end{solution}
\begin{exercise}\label{cw:bank-5B-4}原题5B-4. 对以下五个系统各执行 Notes GS.4 的作图程序：(i) 求矩阵特征值，判定原点几何类型及稳定性；(ii) 若根为实数，求特征向量、画相应轨道及附近轨道并标正向箭头；(iii) 若根为复数，用若干方向场向量判断方向及曲线的大致形状.
\[\begin{array}{ll}
(a)\ x'=2x-3y,\ y'=x-2y;&(b)\ x'=2x,\ y'=3x+y;\\
(c)\ x'=-2x-2y,\ y'=-x-3y;&(d)\ x'=x-2y,\ y'=x+y;\\
(e)\ x'=x+y,\ y'=-2x-y.&
\end{array}\]
\end{exercise}
\begin{solution}（官方译审）(a) 特征多项式 $\lambda^2-1$，$\lambda=1,-1$ 对应 $(3,1)^T,(1,1)^T$，原点是不稳定鞍点；前者向外、后者向内，附近轨道为 $c_1e^t(3,1)^T+c_2e^{-t}(1,1)^T$. 此实根情形适用(i)(ii)，(iii)不适用.

(b) 特征值 $2,1$，向量 $(1,3)^T,(0,1)^T$，是向外的不稳定结点. 非直线轨道正向远离时趋近平行较快模态 $(1,3)^T$，逆向靠近原点时与较慢模态 $(0,1)^T$ 相切，参数式为 $c_1e^{2t}(1,3)^T+c_2e^t(0,1)^T$. 适用(i)(ii).

(c) 特征多项式 $(\lambda+1)(\lambda+4)$，$-1,-4$ 对应 $(2,-1)^T,(1,1)^T$，是渐近稳定结点. 所有箭头向内；一般轨道正向接近原点时与慢模态 $(2,-1)^T$ 相切，参数式为 $c_1e^{-t}(2,-1)^T+c_2e^{-4t}(1,1)^T$. 适用(i)(ii).

(d) 特征值 $1\pm i\sqrt2$，为不稳定螺旋点；$(1,0)$ 处向量 $(1,1)$，逆时针向外. 可用 $x=e^t(a\cos\sqrt2t+b\sin\sqrt2t)$, $y=e^t(a\sin\sqrt2t-b\cos\sqrt2t)/\sqrt2$ 重建轨道. 适用(i)(iii)，(ii)不适用.

(e) 特征值 $\pm i$，为稳定但不渐近稳定的中心；$(1,0)$ 处向量 $(1,-2)$，顺时针. 二次型 $H=2x^2+2xy+y^2$ 正定，直接求导得 $H'=0$，故轨道为倾斜椭圆 $H=C>0$. 适用(i)(iii). 原程序(iii)把复根图形统称为“spiral”；本小问纯虚根实际产生闭椭圆，应按计算修正. 下图按方程作出轨道；箭头给正向运动.
\begin{center}
\begin{tikzpicture}\begin{axis}[width=.47\linewidth,height=5.2cm,axis lines=middle,xlabel={$x$},ylabel={$y$},xmin=-2,xmax=2,ymin=-2,ymax=2,title={（a）鞍点},tick label style={font=\footnotesize},clip=true]
\addplot[structurecolor,thin,no marks] coordinates {(2.00000,2.00000) (1.94627,1.94627) (1.89398,1.89398) (1.84309,1.84309) (1.79358,1.79358) (1.74539,1.74539) (1.69850,1.69850) (1.65286,1.65286) (1.60846,1.60846) (1.56524,1.56524) (1.52319,1.52319) (1.48227,1.48227) (1.44245,1.44245) (1.40369,1.40369) (1.36598,1.36598) (1.32928,1.32928) (1.29357,1.29357) (1.25882,1.25882) (1.22500,1.22500) (1.19208,1.19208) (1.16006,1.16006) (1.12889,1.12889) (1.09856,1.09856) (1.06905,1.06905) (1.04033,1.04033) (1.01238,1.01238) (0.98518,0.98518) (0.95871,0.95871) (0.93295,0.93295) (0.90789,0.90789) (0.88349,0.88349) (0.85976,0.85976) (0.83666,0.83666) (0.81418,0.81418) (0.79231,0.79231) (0.77102,0.77102) (0.75031,0.75031) (0.73015,0.73015) (0.71053,0.71053) (0.69144,0.69144) (0.67287,0.67287) (0.65479,0.65479) (0.63720,0.63720) (0.62008,0.62008) (0.60342,0.60342) (0.58721,0.58721) (0.57143,0.57143) (0.55608,0.55608) (0.54114,0.54114) (0.52660,0.52660) (0.51245,0.51245) (0.49868,0.49868) (0.48529,0.48529) (0.47225,0.47225) (0.45956,0.45956) (0.44721,0.44721) (0.43520,0.43520) (0.42351,0.42351) (0.41213,0.41213) (0.40106,0.40106) (0.39028,0.39028) (0.37980,0.37980) (0.36959,0.36959) (0.35966,0.35966) (0.35000,0.35000)};
\draw[-{Latex},structurecolor] (axis cs:1.16006,1.16006) -- (axis cs:1.12889,1.12889);
\draw[-{Latex},structurecolor] (axis cs:0.62008,0.62008) -- (axis cs:0.60342,0.60342);
\addplot[structurecolor,thin,no marks] coordinates {(0.35000,0.35000) (0.32370,0.32370) (0.29937,0.29937) (0.27687,0.27687) (0.25607,0.25607) (0.23682,0.23682) (0.21902,0.21902) (0.20256,0.20256) (0.18734,0.18734) (0.17326,0.17326) (0.16024,0.16024) (0.14820,0.14820) (0.13706,0.13706) (0.12676,0.12676) (0.11724,0.11724) (0.10842,0.10842) (0.10028,0.10028) (0.09274,0.09274) (0.08577,0.08577) (0.07933,0.07933) (0.07336,0.07336) (0.06785,0.06785) (0.06275,0.06275) (0.05804,0.05804) (0.05367,0.05367) (0.04964,0.04964) (0.04591,0.04591) (0.04246,0.04246) (0.03927,0.03927) (0.03632,0.03632) (0.03359,0.03359) (0.03106,0.03106) (0.02873,0.02873) (0.02657,0.02657) (0.02457,0.02457) (0.02273,0.02273) (0.02102,0.02102) (0.01944,0.01944) (0.01798,0.01798) (0.01663,0.01663) (0.01538,0.01538) (0.01422,0.01422) (0.01315,0.01315) (0.01216,0.01216) (0.01125,0.01125) (0.01041,0.01041) (0.00962,0.00962) (0.00890,0.00890) (0.00823,0.00823) (0.00761,0.00761) (0.00704,0.00704) (0.00651,0.00651) (0.00602,0.00602) (0.00557,0.00557) (0.00515,0.00515) (0.00476,0.00476) (0.00441,0.00441) (0.00407,0.00407) (0.00377,0.00377) (0.00349,0.00349) (0.00322,0.00322) (0.00298,0.00298) (0.00276,0.00276) (0.00255,0.00255) (0.00236,0.00236)};
\draw[-{Latex},structurecolor] (axis cs:0.07336,0.07336) -- (axis cs:0.06785,0.06785);
\draw[-{Latex},structurecolor] (axis cs:0.01216,0.01216) -- (axis cs:0.01125,0.01125);
\addplot[structurecolor,thin,no marks] coordinates {(-1.76841,-2.00000) (-1.72610,-1.96173) (-1.68431,-1.92404) (-1.64302,-1.88693) (-1.60222,-1.85038) (-1.56190,-1.81439) (-1.52204,-1.77894) (-1.48264,-1.74402) (-1.44368,-1.70962) (-1.40516,-1.67573) (-1.36705,-1.64234) (-1.32936,-1.60944) (-1.29206,-1.57702) (-1.25514,-1.54507) (-1.21860,-1.51359) (-1.18243,-1.48256) (-1.14660,-1.45197) (-1.11112,-1.42181) (-1.07598,-1.39208) (-1.04115,-1.36276) (-1.00663,-1.33385) (-0.97242,-1.30534) (-0.93849,-1.27722) (-0.90485,-1.24948) (-0.87148,-1.22212) (-0.83836,-1.19512) (-0.80550,-1.16847) (-0.77288,-1.14218) (-0.74049,-1.11623) (-0.70832,-1.09061) (-0.67636,-1.06531) (-0.64460,-1.04034) (-0.61304,-1.01567) (-0.58166,-0.99131) (-0.55045,-0.96724) (-0.51940,-0.94347) (-0.48852,-0.91997) (-0.45778,-0.89675) (-0.42717,-0.87380) (-0.39669,-0.85111) (-0.36633,-0.82867) (-0.33608,-0.80648) (-0.30593,-0.78453) (-0.27588,-0.76282) (-0.24590,-0.74133) (-0.21600,-0.72006) (-0.18616,-0.69901) (-0.15638,-0.67817) (-0.12665,-0.65753) (-0.09695,-0.63709) (-0.06728,-0.61684) (-0.03763,-0.59677) (-0.00799,-0.57688) (0.02164,-0.55717) (0.05128,-0.53762) (0.08094,-0.51823) (0.11062,-0.49899) (0.14033,-0.47990) (0.17009,-0.46096) (0.19989,-0.44216) (0.22976,-0.42348) (0.25969,-0.40494) (0.28971,-0.38651) (0.31981,-0.36820) (0.35000,-0.35000)};
\draw[-{Latex},structurecolor] (axis cs:-1.00663,-1.33385) -- (axis cs:-0.97242,-1.30534);
\draw[-{Latex},structurecolor] (axis cs:-0.27588,-0.76282) -- (axis cs:-0.24590,-0.74133);
\addplot[structurecolor,thin,no marks] coordinates {(0.35000,-0.35000) (0.37167,-0.33704) (0.39340,-0.32413) (0.41519,-0.31127) (0.43704,-0.29846) (0.45896,-0.28569) (0.48095,-0.27297) (0.50301,-0.26029) (0.52515,-0.24765) (0.54737,-0.23505) (0.56968,-0.22248) (0.59207,-0.20995) (0.61455,-0.19745) (0.63713,-0.18498) (0.65980,-0.17254) (0.68257,-0.16012) (0.70545,-0.14773) (0.72844,-0.13536) (0.75154,-0.12301) (0.77475,-0.11069) (0.79808,-0.09837) (0.82153,-0.08608) (0.84511,-0.07379) (0.86882,-0.06152) (0.89266,-0.04926) (0.91664,-0.03700) (0.94076,-0.02475) (0.96502,-0.01250) (0.98943,-0.00026) (1.01399,0.01198) (1.03871,0.02423) (1.06358,0.03648) (1.08862,0.04873) (1.11382,0.06100) (1.13920,0.07327) (1.16475,0.08555) (1.19047,0.09785) (1.21638,0.11016) (1.24248,0.12249) (1.26876,0.13484) (1.29524,0.14720) (1.32192,0.15959) (1.34880,0.17201) (1.37588,0.18445) (1.40318,0.19692) (1.43069,0.20942) (1.45842,0.22195) (1.48638,0.23451) (1.51456,0.24711) (1.54297,0.25975) (1.57162,0.27243) (1.60051,0.28515) (1.62964,0.29792) (1.65902,0.31073) (1.68866,0.32358) (1.71855,0.33649) (1.74871,0.34945) (1.77914,0.36246) (1.80983,0.37553) (1.84081,0.38865) (1.87206,0.40183) (1.90361,0.41508) (1.93544,0.42839) (1.96757,0.44176) (2.00000,0.45520)};
\draw[-{Latex},structurecolor] (axis cs:0.79808,-0.09837) -- (axis cs:0.82153,-0.08608);
\draw[-{Latex},structurecolor] (axis cs:1.37588,0.18445) -- (axis cs:1.40318,0.19692);
\addplot[structurecolor,thin,no marks] coordinates {(1.76841,2.00000) (1.72610,1.96173) (1.68431,1.92404) (1.64302,1.88693) (1.60222,1.85038) (1.56190,1.81439) (1.52204,1.77894) (1.48264,1.74402) (1.44368,1.70962) (1.40516,1.67573) (1.36705,1.64234) (1.32936,1.60944) (1.29206,1.57702) (1.25514,1.54507) (1.21860,1.51359) (1.18243,1.48256) (1.14660,1.45197) (1.11112,1.42181) (1.07598,1.39208) (1.04115,1.36276) (1.00663,1.33385) (0.97242,1.30534) (0.93849,1.27722) (0.90485,1.24948) (0.87148,1.22212) (0.83836,1.19512) (0.80550,1.16847) (0.77288,1.14218) (0.74049,1.11623) (0.70832,1.09061) (0.67636,1.06531) (0.64460,1.04034) (0.61304,1.01567) (0.58166,0.99131) (0.55045,0.96724) (0.51940,0.94347) (0.48852,0.91997) (0.45778,0.89675) (0.42717,0.87380) (0.39669,0.85111) (0.36633,0.82867) (0.33608,0.80648) (0.30593,0.78453) (0.27588,0.76282) (0.24590,0.74133) (0.21600,0.72006) (0.18616,0.69901) (0.15638,0.67817) (0.12665,0.65753) (0.09695,0.63709) (0.06728,0.61684) (0.03763,0.59677) (0.00799,0.57688) (-0.02164,0.55717) (-0.05128,0.53762) (-0.08094,0.51823) (-0.11062,0.49899) (-0.14033,0.47990) (-0.17009,0.46096) (-0.19989,0.44216) (-0.22976,0.42348) (-0.25969,0.40494) (-0.28971,0.38651) (-0.31981,0.36820) (-0.35000,0.35000)};
\draw[-{Latex},structurecolor] (axis cs:1.00663,1.33385) -- (axis cs:0.97242,1.30534);
\draw[-{Latex},structurecolor] (axis cs:0.27588,0.76282) -- (axis cs:0.24590,0.74133);
\addplot[structurecolor,thin,no marks] coordinates {(-0.35000,0.35000) (-0.37167,0.33704) (-0.39340,0.32413) (-0.41519,0.31127) (-0.43704,0.29846) (-0.45896,0.28569) (-0.48095,0.27297) (-0.50301,0.26029) (-0.52515,0.24765) (-0.54737,0.23505) (-0.56968,0.22248) (-0.59207,0.20995) (-0.61455,0.19745) (-0.63713,0.18498) (-0.65980,0.17254) (-0.68257,0.16012) (-0.70545,0.14773) (-0.72844,0.13536) (-0.75154,0.12301) (-0.77475,0.11069) (-0.79808,0.09837) (-0.82153,0.08608) (-0.84511,0.07379) (-0.86882,0.06152) (-0.89266,0.04926) (-0.91664,0.03700) (-0.94076,0.02475) (-0.96502,0.01250) (-0.98943,0.00026) (-1.01399,-0.01198) (-1.03871,-0.02423) (-1.06358,-0.03648) (-1.08862,-0.04873) (-1.11382,-0.06100) (-1.13920,-0.07327) (-1.16475,-0.08555) (-1.19047,-0.09785) (-1.21638,-0.11016) (-1.24248,-0.12249) (-1.26876,-0.13484) (-1.29524,-0.14720) (-1.32192,-0.15959) (-1.34880,-0.17201) (-1.37588,-0.18445) (-1.40318,-0.19692) (-1.43069,-0.20942) (-1.45842,-0.22195) (-1.48638,-0.23451) (-1.51456,-0.24711) (-1.54297,-0.25975) (-1.57162,-0.27243) (-1.60051,-0.28515) (-1.62964,-0.29792) (-1.65902,-0.31073) (-1.68866,-0.32358) (-1.71855,-0.33649) (-1.74871,-0.34945) (-1.77914,-0.36246) (-1.80983,-0.37553) (-1.84081,-0.38865) (-1.87206,-0.40183) (-1.90361,-0.41508) (-1.93544,-0.42839) (-1.96757,-0.44176) (-2.00000,-0.45520)};
\draw[-{Latex},structurecolor] (axis cs:-0.79808,0.09837) -- (axis cs:-0.82153,0.08608);
\draw[-{Latex},structurecolor] (axis cs:-1.37588,-0.18445) -- (axis cs:-1.40318,-0.19692);
\addplot[structurecolor,thin,no marks] coordinates {(-2.00000,-2.00000) (-1.94627,-1.94627) (-1.89398,-1.89398) (-1.84309,-1.84309) (-1.79358,-1.79358) (-1.74539,-1.74539) (-1.69850,-1.69850) (-1.65286,-1.65286) (-1.60846,-1.60846) (-1.56524,-1.56524) (-1.52319,-1.52319) (-1.48227,-1.48227) (-1.44245,-1.44245) (-1.40369,-1.40369) (-1.36598,-1.36598) (-1.32928,-1.32928) (-1.29357,-1.29357) (-1.25882,-1.25882) (-1.22500,-1.22500) (-1.19208,-1.19208) (-1.16006,-1.16006) (-1.12889,-1.12889) (-1.09856,-1.09856) (-1.06905,-1.06905) (-1.04033,-1.04033) (-1.01238,-1.01238) (-0.98518,-0.98518) (-0.95871,-0.95871) (-0.93295,-0.93295) (-0.90789,-0.90789) (-0.88349,-0.88349) (-0.85976,-0.85976) (-0.83666,-0.83666) (-0.81418,-0.81418) (-0.79231,-0.79231) (-0.77102,-0.77102) (-0.75031,-0.75031) (-0.73015,-0.73015) (-0.71053,-0.71053) (-0.69144,-0.69144) (-0.67287,-0.67287) (-0.65479,-0.65479) (-0.63720,-0.63720) (-0.62008,-0.62008) (-0.60342,-0.60342) (-0.58721,-0.58721) (-0.57143,-0.57143) (-0.55608,-0.55608) (-0.54114,-0.54114) (-0.52660,-0.52660) (-0.51245,-0.51245) (-0.49868,-0.49868) (-0.48529,-0.48529) (-0.47225,-0.47225) (-0.45956,-0.45956) (-0.44721,-0.44721) (-0.43520,-0.43520) (-0.42351,-0.42351) (-0.41213,-0.41213) (-0.40106,-0.40106) (-0.39028,-0.39028) (-0.37980,-0.37980) (-0.36959,-0.36959) (-0.35966,-0.35966) (-0.35000,-0.35000)};
\draw[-{Latex},structurecolor] (axis cs:-1.16006,-1.16006) -- (axis cs:-1.12889,-1.12889);
\draw[-{Latex},structurecolor] (axis cs:-0.62008,-0.62008) -- (axis cs:-0.60342,-0.60342);
\addplot[structurecolor,thin,no marks] coordinates {(-0.35000,-0.35000) (-0.32370,-0.32370) (-0.29937,-0.29937) (-0.27687,-0.27687) (-0.25607,-0.25607) (-0.23682,-0.23682) (-0.21902,-0.21902) (-0.20256,-0.20256) (-0.18734,-0.18734) (-0.17326,-0.17326) (-0.16024,-0.16024) (-0.14820,-0.14820) (-0.13706,-0.13706) (-0.12676,-0.12676) (-0.11724,-0.11724) (-0.10842,-0.10842) (-0.10028,-0.10028) (-0.09274,-0.09274) (-0.08577,-0.08577) (-0.07933,-0.07933) (-0.07336,-0.07336) (-0.06785,-0.06785) (-0.06275,-0.06275) (-0.05804,-0.05804) (-0.05367,-0.05367) (-0.04964,-0.04964) (-0.04591,-0.04591) (-0.04246,-0.04246) (-0.03927,-0.03927) (-0.03632,-0.03632) (-0.03359,-0.03359) (-0.03106,-0.03106) (-0.02873,-0.02873) (-0.02657,-0.02657) (-0.02457,-0.02457) (-0.02273,-0.02273) (-0.02102,-0.02102) (-0.01944,-0.01944) (-0.01798,-0.01798) (-0.01663,-0.01663) (-0.01538,-0.01538) (-0.01422,-0.01422) (-0.01315,-0.01315) (-0.01216,-0.01216) (-0.01125,-0.01125) (-0.01041,-0.01041) (-0.00962,-0.00962) (-0.00890,-0.00890) (-0.00823,-0.00823) (-0.00761,-0.00761) (-0.00704,-0.00704) (-0.00651,-0.00651) (-0.00602,-0.00602) (-0.00557,-0.00557) (-0.00515,-0.00515) (-0.00476,-0.00476) (-0.00441,-0.00441) (-0.00407,-0.00407) (-0.00377,-0.00377) (-0.00349,-0.00349) (-0.00322,-0.00322) (-0.00298,-0.00298) (-0.00276,-0.00276) (-0.00255,-0.00255) (-0.00236,-0.00236)};
\draw[-{Latex},structurecolor] (axis cs:-0.07336,-0.07336) -- (axis cs:-0.06785,-0.06785);
\draw[-{Latex},structurecolor] (axis cs:-0.01216,-0.01216) -- (axis cs:-0.01125,-0.01125);
\addplot[structurecolor,thin,no marks] coordinates {(-1.88104,-2.00000) (-1.81978,-1.94208) (-1.75992,-1.88565) (-1.70140,-1.83067) (-1.64419,-1.77709) (-1.58825,-1.72487) (-1.53352,-1.67398) (-1.47996,-1.62437) (-1.42754,-1.57601) (-1.37622,-1.52885) (-1.32595,-1.48287) (-1.27670,-1.43802) (-1.22843,-1.39428) (-1.18109,-1.35161) (-1.13467,-1.30997) (-1.08911,-1.26934) (-1.04439,-1.22968) (-1.00048,-1.19096) (-0.95732,-1.15316) (-0.91491,-1.11624) (-0.87319,-1.08018) (-0.83215,-1.04494) (-0.79174,-1.01051) (-0.75194,-0.97685) (-0.71271,-0.94394) (-0.67404,-0.91176) (-0.63588,-0.88027) (-0.59820,-0.84946) (-0.56099,-0.81930) (-0.52420,-0.78977) (-0.48782,-0.76085) (-0.45181,-0.73250) (-0.41615,-0.70472) (-0.38081,-0.67749) (-0.34576,-0.65077) (-0.31098,-0.62454) (-0.27643,-0.59880) (-0.24209,-0.57352) (-0.20795,-0.54868) (-0.17396,-0.52425) (-0.14010,-0.50023) (-0.10635,-0.47660) (-0.07269,-0.45333) (-0.03907,-0.43040) (-0.00549,-0.40781) (0.02808,-0.38553) (0.06168,-0.36355) (0.09533,-0.34184) (0.12905,-0.32040) (0.16287,-0.29920) (0.19681,-0.27823) (0.23090,-0.25747) (0.26517,-0.23692) (0.29965,-0.21654) (0.33435,-0.19633) (0.36931,-0.17627) (0.40456,-0.15635) (0.44011,-0.13654) (0.47600,-0.11684) (0.51226,-0.09723) (0.54891,-0.07770) (0.58598,-0.05822) (0.62350,-0.03879) (0.66149,-0.01939) (0.70000,0.00000)};
\draw[-{Latex},structurecolor] (axis cs:-0.87319,-1.08018) -- (axis cs:-0.83215,-1.04494);
\draw[-{Latex},structurecolor] (axis cs:-0.03907,-0.43040) -- (axis cs:-0.00549,-0.40781);
\addplot[structurecolor,thin,no marks] coordinates {(0.70000,0.00000) (0.71592,0.00794) (0.73193,0.01588) (0.74804,0.02382) (0.76424,0.03176) (0.78054,0.03971) (0.79694,0.04766) (0.81344,0.05562) (0.83005,0.06358) (0.84676,0.07156) (0.86358,0.07954) (0.88052,0.08753) (0.89756,0.09554) (0.91473,0.10355) (0.93201,0.11158) (0.94941,0.11963) (0.96693,0.12769) (0.98457,0.13576) (1.00235,0.14386) (1.02025,0.15197) (1.03828,0.16010) (1.05645,0.16825) (1.07475,0.17643) (1.09319,0.18463) (1.11177,0.19285) (1.13049,0.20109) (1.14936,0.20936) (1.16838,0.21766) (1.18755,0.22599) (1.20687,0.23434) (1.22634,0.24273) (1.24597,0.25114) (1.26577,0.25959) (1.28572,0.26807) (1.30584,0.27659) (1.32613,0.28514) (1.34659,0.29373) (1.36722,0.30236) (1.38803,0.31102) (1.40902,0.31973) (1.43018,0.32847) (1.45153,0.33726) (1.47307,0.34609) (1.49480,0.35497) (1.51672,0.36389) (1.53883,0.37286) (1.56115,0.38188) (1.58366,0.39094) (1.60638,0.40006) (1.62930,0.40923) (1.65243,0.41845) (1.67578,0.42772) (1.69934,0.43705) (1.72312,0.44644) (1.74712,0.45588) (1.77134,0.46538) (1.79580,0.47494) (1.82048,0.48456) (1.84540,0.49425) (1.87055,0.50400) (1.89595,0.51381) (1.92159,0.52369) (1.94747,0.53363) (1.97361,0.54365) (2.00000,0.55373)};
\draw[-{Latex},structurecolor] (axis cs:1.03828,0.16010) -- (axis cs:1.05645,0.16825);
\draw[-{Latex},structurecolor] (axis cs:1.49480,0.35497) -- (axis cs:1.51672,0.36389);
\addplot[structurecolor,thin,no marks] coordinates {(1.88104,2.00000) (1.81978,1.94208) (1.75992,1.88565) (1.70140,1.83067) (1.64419,1.77709) (1.58825,1.72487) (1.53352,1.67398) (1.47996,1.62437) (1.42754,1.57601) (1.37622,1.52885) (1.32595,1.48287) (1.27670,1.43802) (1.22843,1.39428) (1.18109,1.35161) (1.13467,1.30997) (1.08911,1.26934) (1.04439,1.22968) (1.00048,1.19096) (0.95732,1.15316) (0.91491,1.11624) (0.87319,1.08018) (0.83215,1.04494) (0.79174,1.01051) (0.75194,0.97685) (0.71271,0.94394) (0.67404,0.91176) (0.63588,0.88027) (0.59820,0.84946) (0.56099,0.81930) (0.52420,0.78977) (0.48782,0.76085) (0.45181,0.73250) (0.41615,0.70472) (0.38081,0.67749) (0.34576,0.65077) (0.31098,0.62454) (0.27643,0.59880) (0.24209,0.57352) (0.20795,0.54868) (0.17396,0.52425) (0.14010,0.50023) (0.10635,0.47660) (0.07269,0.45333) (0.03907,0.43040) (0.00549,0.40781) (-0.02808,0.38553) (-0.06168,0.36355) (-0.09533,0.34184) (-0.12905,0.32040) (-0.16287,0.29920) (-0.19681,0.27823) (-0.23090,0.25747) (-0.26517,0.23692) (-0.29965,0.21654) (-0.33435,0.19633) (-0.36931,0.17627) (-0.40456,0.15635) (-0.44011,0.13654) (-0.47600,0.11684) (-0.51226,0.09723) (-0.54891,0.07770) (-0.58598,0.05822) (-0.62350,0.03879) (-0.66149,0.01939) (-0.70000,0.00000)};
\draw[-{Latex},structurecolor] (axis cs:0.87319,1.08018) -- (axis cs:0.83215,1.04494);
\draw[-{Latex},structurecolor] (axis cs:0.03907,0.43040) -- (axis cs:0.00549,0.40781);
\addplot[structurecolor,thin,no marks] coordinates {(-0.70000,0.00000) (-0.71592,-0.00794) (-0.73193,-0.01588) (-0.74804,-0.02382) (-0.76424,-0.03176) (-0.78054,-0.03971) (-0.79694,-0.04766) (-0.81344,-0.05562) (-0.83005,-0.06358) (-0.84676,-0.07156) (-0.86358,-0.07954) (-0.88052,-0.08753) (-0.89756,-0.09554) (-0.91473,-0.10355) (-0.93201,-0.11158) (-0.94941,-0.11963) (-0.96693,-0.12769) (-0.98457,-0.13576) (-1.00235,-0.14386) (-1.02025,-0.15197) (-1.03828,-0.16010) (-1.05645,-0.16825) (-1.07475,-0.17643) (-1.09319,-0.18463) (-1.11177,-0.19285) (-1.13049,-0.20109) (-1.14936,-0.20936) (-1.16838,-0.21766) (-1.18755,-0.22599) (-1.20687,-0.23434) (-1.22634,-0.24273) (-1.24597,-0.25114) (-1.26577,-0.25959) (-1.28572,-0.26807) (-1.30584,-0.27659) (-1.32613,-0.28514) (-1.34659,-0.29373) (-1.36722,-0.30236) (-1.38803,-0.31102) (-1.40902,-0.31973) (-1.43018,-0.32847) (-1.45153,-0.33726) (-1.47307,-0.34609) (-1.49480,-0.35497) (-1.51672,-0.36389) (-1.53883,-0.37286) (-1.56115,-0.38188) (-1.58366,-0.39094) (-1.60638,-0.40006) (-1.62930,-0.40923) (-1.65243,-0.41845) (-1.67578,-0.42772) (-1.69934,-0.43705) (-1.72312,-0.44644) (-1.74712,-0.45588) (-1.77134,-0.46538) (-1.79580,-0.47494) (-1.82048,-0.48456) (-1.84540,-0.49425) (-1.87055,-0.50400) (-1.89595,-0.51381) (-1.92159,-0.52369) (-1.94747,-0.53363) (-1.97361,-0.54365) (-2.00000,-0.55373)};
\draw[-{Latex},structurecolor] (axis cs:-1.03828,-0.16010) -- (axis cs:-1.05645,-0.16825);
\draw[-{Latex},structurecolor] (axis cs:-1.49480,-0.35497) -- (axis cs:-1.51672,-0.36389);
\addplot[structurecolor,thin,no marks] coordinates {(1.66120,2.00000) (1.63094,1.97361) (1.60090,1.94747) (1.57106,1.92159) (1.54142,1.89595) (1.51199,1.87055) (1.48274,1.84540) (1.45369,1.82048) (1.42482,1.79580) (1.39614,1.77134) (1.36764,1.74712) (1.33931,1.72312) (1.31115,1.69934) (1.28317,1.67578) (1.25534,1.65243) (1.22768,1.62930) (1.20018,1.60638) (1.17283,1.58366) (1.14563,1.56115) (1.11858,1.53883) (1.09168,1.51672) (1.06491,1.49480) (1.03828,1.47307) (1.01179,1.45153) (0.98542,1.43018) (0.95918,1.40902) (0.93306,1.38803) (0.90707,1.36722) (0.88119,1.34659) (0.85542,1.32613) (0.82977,1.30584) (0.80422,1.28572) (0.77877,1.26577) (0.75343,1.24597) (0.72818,1.22634) (0.70302,1.20687) (0.67796,1.18755) (0.65298,1.16838) (0.62809,1.14936) (0.60327,1.13049) (0.57854,1.11177) (0.55388,1.09319) (0.52929,1.07475) (0.50476,1.05645) (0.48031,1.03828) (0.45591,1.02025) (0.43158,1.00235) (0.40729,0.98457) (0.38307,0.96693) (0.35889,0.94941) (0.33475,0.93201) (0.31066,0.91473) (0.28661,0.89756) (0.26260,0.88052) (0.23862,0.86358) (0.21467,0.84676) (0.19075,0.83005) (0.16685,0.81344) (0.14297,0.79694) (0.11912,0.78054) (0.09528,0.76424) (0.07145,0.74804) (0.04763,0.73193) (0.02381,0.71592) (0.00000,0.70000)};
\draw[-{Latex},structurecolor] (axis cs:1.09168,1.51672) -- (axis cs:1.06491,1.49480);
\draw[-{Latex},structurecolor] (axis cs:0.50476,1.05645) -- (axis cs:0.48031,1.03828);
\addplot[structurecolor,thin,no marks] coordinates {(0.00000,0.70000) (-0.02780,0.68153) (-0.05561,0.66317) (-0.08343,0.64493) (-0.11126,0.62681) (-0.13911,0.60879) (-0.16699,0.59088) (-0.19489,0.57308) (-0.22283,0.55538) (-0.25081,0.53777) (-0.27883,0.52026) (-0.30690,0.50283) (-0.33503,0.48550) (-0.36321,0.46825) (-0.39146,0.45108) (-0.41978,0.43400) (-0.44817,0.41699) (-0.47663,0.40005) (-0.50519,0.38318) (-0.53383,0.36638) (-0.56256,0.34964) (-0.59139,0.33297) (-0.62033,0.31635) (-0.64937,0.29979) (-0.67853,0.28328) (-0.70781,0.26682) (-0.73721,0.25041) (-0.76674,0.23404) (-0.79641,0.21771) (-0.82621,0.20142) (-0.85616,0.18517) (-0.88626,0.16894) (-0.91652,0.15275) (-0.94693,0.13659) (-0.97751,0.12045) (-1.00827,0.10433) (-1.03919,0.08822) (-1.07031,0.07214) (-1.10161,0.05606) (-1.13310,0.04000) (-1.16479,0.02394) (-1.19669,0.00789) (-1.22879,-0.00816) (-1.26111,-0.02422) (-1.29365,-0.04027) (-1.32642,-0.05634) (-1.35942,-0.07241) (-1.39266,-0.08850) (-1.42614,-0.10460) (-1.45988,-0.12072) (-1.49386,-0.13686) (-1.52812,-0.15303) (-1.56263,-0.16922) (-1.59743,-0.18545) (-1.63250,-0.20170) (-1.66786,-0.21799) (-1.70351,-0.23432) (-1.73946,-0.25069) (-1.77571,-0.26710) (-1.81228,-0.28356) (-1.84916,-0.30007) (-1.88637,-0.31663) (-1.92391,-0.33325) (-1.96178,-0.34993) (-2.00000,-0.36667)};
\draw[-{Latex},structurecolor] (axis cs:-0.56256,0.34964) -- (axis cs:-0.59139,0.33297);
\draw[-{Latex},structurecolor] (axis cs:-1.26111,-0.02422) -- (axis cs:-1.29365,-0.04027);
\addplot[structurecolor,thin,no marks] coordinates {(-1.66120,-2.00000) (-1.63094,-1.97361) (-1.60090,-1.94747) (-1.57106,-1.92159) (-1.54142,-1.89595) (-1.51199,-1.87055) (-1.48274,-1.84540) (-1.45369,-1.82048) (-1.42482,-1.79580) (-1.39614,-1.77134) (-1.36764,-1.74712) (-1.33931,-1.72312) (-1.31115,-1.69934) (-1.28317,-1.67578) (-1.25534,-1.65243) (-1.22768,-1.62930) (-1.20018,-1.60638) (-1.17283,-1.58366) (-1.14563,-1.56115) (-1.11858,-1.53883) (-1.09168,-1.51672) (-1.06491,-1.49480) (-1.03828,-1.47307) (-1.01179,-1.45153) (-0.98542,-1.43018) (-0.95918,-1.40902) (-0.93306,-1.38803) (-0.90707,-1.36722) (-0.88119,-1.34659) (-0.85542,-1.32613) (-0.82977,-1.30584) (-0.80422,-1.28572) (-0.77877,-1.26577) (-0.75343,-1.24597) (-0.72818,-1.22634) (-0.70302,-1.20687) (-0.67796,-1.18755) (-0.65298,-1.16838) (-0.62809,-1.14936) (-0.60327,-1.13049) (-0.57854,-1.11177) (-0.55388,-1.09319) (-0.52929,-1.07475) (-0.50476,-1.05645) (-0.48031,-1.03828) (-0.45591,-1.02025) (-0.43158,-1.00235) (-0.40729,-0.98457) (-0.38307,-0.96693) (-0.35889,-0.94941) (-0.33475,-0.93201) (-0.31066,-0.91473) (-0.28661,-0.89756) (-0.26260,-0.88052) (-0.23862,-0.86358) (-0.21467,-0.84676) (-0.19075,-0.83005) (-0.16685,-0.81344) (-0.14297,-0.79694) (-0.11912,-0.78054) (-0.09528,-0.76424) (-0.07145,-0.74804) (-0.04763,-0.73193) (-0.02381,-0.71592) (0.00000,-0.70000)};
\draw[-{Latex},structurecolor] (axis cs:-1.09168,-1.51672) -- (axis cs:-1.06491,-1.49480);
\draw[-{Latex},structurecolor] (axis cs:-0.50476,-1.05645) -- (axis cs:-0.48031,-1.03828);
\addplot[structurecolor,thin,no marks] coordinates {(0.00000,-0.70000) (0.02780,-0.68153) (0.05561,-0.66317) (0.08343,-0.64493) (0.11126,-0.62681) (0.13911,-0.60879) (0.16699,-0.59088) (0.19489,-0.57308) (0.22283,-0.55538) (0.25081,-0.53777) (0.27883,-0.52026) (0.30690,-0.50283) (0.33503,-0.48550) (0.36321,-0.46825) (0.39146,-0.45108) (0.41978,-0.43400) (0.44817,-0.41699) (0.47663,-0.40005) (0.50519,-0.38318) (0.53383,-0.36638) (0.56256,-0.34964) (0.59139,-0.33297) (0.62033,-0.31635) (0.64937,-0.29979) (0.67853,-0.28328) (0.70781,-0.26682) (0.73721,-0.25041) (0.76674,-0.23404) (0.79641,-0.21771) (0.82621,-0.20142) (0.85616,-0.18517) (0.88626,-0.16894) (0.91652,-0.15275) (0.94693,-0.13659) (0.97751,-0.12045) (1.00827,-0.10433) (1.03919,-0.08822) (1.07031,-0.07214) (1.10161,-0.05606) (1.13310,-0.04000) (1.16479,-0.02394) (1.19669,-0.00789) (1.22879,0.00816) (1.26111,0.02422) (1.29365,0.04027) (1.32642,0.05634) (1.35942,0.07241) (1.39266,0.08850) (1.42614,0.10460) (1.45988,0.12072) (1.49386,0.13686) (1.52812,0.15303) (1.56263,0.16922) (1.59743,0.18545) (1.63250,0.20170) (1.66786,0.21799) (1.70351,0.23432) (1.73946,0.25069) (1.77571,0.26710) (1.81228,0.28356) (1.84916,0.30007) (1.88637,0.31663) (1.92391,0.33325) (1.96178,0.34993) (2.00000,0.36667)};
\draw[-{Latex},structurecolor] (axis cs:0.56256,-0.34964) -- (axis cs:0.59139,-0.33297);
\draw[-{Latex},structurecolor] (axis cs:1.26111,0.02422) -- (axis cs:1.29365,0.04027);
\addplot[only marks,mark=*,mark size=1.4pt] coordinates {(0,0)};
\end{axis}\end{tikzpicture}
\hfill
\begin{tikzpicture}\begin{axis}[width=.47\linewidth,height=5.2cm,axis lines=middle,xlabel={$x$},ylabel={$y$},xmin=-2,xmax=2,ymin=-2,ymax=2,title={（b）不稳定结点},tick label style={font=\footnotesize},clip=true]
\addplot[structurecolor,thin,no marks] coordinates {(0.00002,-0.00467) (0.00002,-0.00504) (0.00002,-0.00545) (0.00003,-0.00589) (0.00003,-0.00636) (0.00003,-0.00687) (0.00004,-0.00742) (0.00005,-0.00801) (0.00006,-0.00865) (0.00006,-0.00933) (0.00008,-0.01007) (0.00009,-0.01087) (0.00010,-0.01173) (0.00012,-0.01266) (0.00014,-0.01366) (0.00017,-0.01473) (0.00019,-0.01588) (0.00023,-0.01712) (0.00026,-0.01845) (0.00031,-0.01988) (0.00036,-0.02142) (0.00042,-0.02306) (0.00049,-0.02482) (0.00058,-0.02671) (0.00068,-0.02873) (0.00079,-0.03089) (0.00092,-0.03319) (0.00108,-0.03564) (0.00126,-0.03825) (0.00148,-0.04103) (0.00173,-0.04397) (0.00202,-0.04709) (0.00236,-0.05038) (0.00276,-0.05386) (0.00322,-0.05751) (0.00377,-0.06133) (0.00441,-0.06532) (0.00515,-0.06947) (0.00602,-0.07375) (0.00704,-0.07816) (0.00823,-0.08265) (0.00962,-0.08720) (0.01125,-0.09175) (0.01315,-0.09624) (0.01538,-0.10059) (0.01798,-0.10471) (0.02102,-0.10848) (0.02457,-0.11176) (0.02873,-0.11436) (0.03359,-0.11608) (0.03927,-0.11666) (0.04591,-0.11579) (0.05367,-0.11310) (0.06275,-0.10814) (0.07336,-0.10039) (0.08577,-0.08921) (0.10028,-0.07385) (0.11724,-0.05342) (0.13706,-0.02686) (0.16024,0.00708) (0.18734,0.04989) (0.21902,0.10333) (0.25607,0.16945) (0.29937,0.25072) (0.35000,0.35000)};
\draw[-{Latex},structurecolor] (axis cs:0.00036,-0.02142) -- (axis cs:0.00042,-0.02306);
\draw[-{Latex},structurecolor] (axis cs:0.01315,-0.09624) -- (axis cs:0.01538,-0.10059);
\addplot[structurecolor,thin,no marks] coordinates {(0.35000,0.35000) (0.35619,0.36242) (0.36250,0.37511) (0.36891,0.38808) (0.37544,0.40133) (0.38209,0.41488) (0.38885,0.42873) (0.39573,0.44287) (0.40274,0.45733) (0.40987,0.47209) (0.41712,0.48718) (0.42450,0.50260) (0.43202,0.51835) (0.43966,0.53443) (0.44744,0.55086) (0.45536,0.56765) (0.46342,0.58479) (0.47162,0.60230) (0.47997,0.62018) (0.48847,0.63845) (0.49711,0.65710) (0.50591,0.67614) (0.51486,0.69559) (0.52398,0.71545) (0.53325,0.73572) (0.54269,0.75642) (0.55229,0.77756) (0.56207,0.79914) (0.57202,0.82116) (0.58214,0.84365) (0.59244,0.86661) (0.60293,0.89004) (0.61360,0.91396) (0.62446,0.93837) (0.63551,0.96329) (0.64676,0.98873) (0.65821,1.01468) (0.66986,1.04117) (0.68171,1.06821) (0.69378,1.09580) (0.70606,1.12395) (0.71856,1.15268) (0.73127,1.18200) (0.74422,1.21191) (0.75739,1.24243) (0.77079,1.27358) (0.78444,1.30535) (0.79832,1.33777) (0.81245,1.37084) (0.82683,1.40458) (0.84146,1.43901) (0.85635,1.47412) (0.87151,1.50995) (0.88694,1.54649) (0.90263,1.58376) (0.91861,1.62179) (0.93487,1.66057) (0.95141,1.70013) (0.96825,1.74048) (0.98539,1.78163) (1.00283,1.82360) (1.02058,1.86641) (1.03864,1.91007) (1.05703,1.95459) (1.07573,2.00000)};
\draw[-{Latex},structurecolor] (axis cs:0.49711,0.65710) -- (axis cs:0.50591,0.67614);
\draw[-{Latex},structurecolor] (axis cs:0.74422,1.21191) -- (axis cs:0.75739,1.24243);
\addplot[structurecolor,thin,no marks] coordinates {(0.00002,-0.00939) (0.00002,-0.01014) (0.00002,-0.01096) (0.00003,-0.01185) (0.00003,-0.01280) (0.00003,-0.01384) (0.00004,-0.01495) (0.00005,-0.01616) (0.00006,-0.01746) (0.00006,-0.01886) (0.00008,-0.02038) (0.00009,-0.02201) (0.00010,-0.02378) (0.00012,-0.02568) (0.00014,-0.02774) (0.00017,-0.02995) (0.00019,-0.03234) (0.00023,-0.03492) (0.00026,-0.03770) (0.00031,-0.04069) (0.00036,-0.04392) (0.00042,-0.04739) (0.00049,-0.05113) (0.00058,-0.05516) (0.00068,-0.05948) (0.00079,-0.06414) (0.00092,-0.06914) (0.00108,-0.07452) (0.00126,-0.08029) (0.00148,-0.08648) (0.00173,-0.09312) (0.00202,-0.10023) (0.00236,-0.10784) (0.00276,-0.11599) (0.00322,-0.12468) (0.00377,-0.13397) (0.00441,-0.14386) (0.00515,-0.15439) (0.00602,-0.16557) (0.00704,-0.17744) (0.00823,-0.19000) (0.00962,-0.20327) (0.01125,-0.21725) (0.01315,-0.23194) (0.01538,-0.24732) (0.01798,-0.26337) (0.02102,-0.28003) (0.02457,-0.29724) (0.02873,-0.31492) (0.03359,-0.33293) (0.03927,-0.35113) (0.04591,-0.36932) (0.05367,-0.38723) (0.06275,-0.40454) (0.07336,-0.42087) (0.08577,-0.43574) (0.10028,-0.44854) (0.11724,-0.45855) (0.13706,-0.46491) (0.16024,-0.46656) (0.18734,-0.46224) (0.21902,-0.45042) (0.25607,-0.42929) (0.29937,-0.39668) (0.35000,-0.35000)};
\draw[-{Latex},structurecolor] (axis cs:0.00036,-0.04392) -- (axis cs:0.00042,-0.04739);
\draw[-{Latex},structurecolor] (axis cs:0.01315,-0.23194) -- (axis cs:0.01538,-0.24732);
\addplot[structurecolor,thin,no marks] coordinates {(0.35000,-0.35000) (0.35873,-0.34116) (0.36767,-0.33189) (0.37684,-0.32217) (0.38624,-0.31198) (0.39587,-0.30131) (0.40574,-0.29014) (0.41586,-0.27847) (0.42623,-0.26627) (0.43686,-0.25353) (0.44775,-0.24023) (0.45892,-0.22635) (0.47036,-0.21188) (0.48209,-0.19681) (0.49411,-0.18110) (0.50643,-0.16475) (0.51906,-0.14773) (0.53200,-0.13003) (0.54527,-0.11162) (0.55887,-0.09248) (0.57280,-0.07259) (0.58709,-0.05194) (0.60173,-0.03049) (0.61673,-0.00822) (0.63211,0.01489) (0.64787,0.03886) (0.66403,0.06373) (0.68059,0.08951) (0.69756,0.11623) (0.71495,0.14392) (0.73278,0.17262) (0.75105,0.20233) (0.76978,0.23311) (0.78898,0.26496) (0.80865,0.29794) (0.82882,0.33207) (0.84948,0.36737) (0.87067,0.40390) (0.89238,0.44167) (0.91463,0.48073) (0.93744,0.52110) (0.96082,0.56284) (0.98477,0.60598) (1.00933,0.65055) (1.03450,0.69659) (1.06030,0.74416) (1.08674,0.79328) (1.11384,0.84401) (1.14161,0.89639) (1.17008,0.95046) (1.19926,1.00628) (1.22916,1.06388) (1.25981,1.12332) (1.29123,1.18465) (1.32343,1.24793) (1.35643,1.31320) (1.39025,1.38052) (1.42492,1.44995) (1.46045,1.52154) (1.49687,1.59536) (1.53420,1.67146) (1.57245,1.74991) (1.61166,1.83077) (1.65185,1.91411) (1.69304,2.00000)};
\draw[-{Latex},structurecolor] (axis cs:0.57280,-0.07259) -- (axis cs:0.58709,-0.05194);
\draw[-{Latex},structurecolor] (axis cs:1.00933,0.65055) -- (axis cs:1.03450,0.69659);
\addplot[structurecolor,thin,no marks] coordinates {(-0.00002,0.00939) (-0.00002,0.01014) (-0.00002,0.01096) (-0.00003,0.01185) (-0.00003,0.01280) (-0.00003,0.01384) (-0.00004,0.01495) (-0.00005,0.01616) (-0.00006,0.01746) (-0.00006,0.01886) (-0.00008,0.02038) (-0.00009,0.02201) (-0.00010,0.02378) (-0.00012,0.02568) (-0.00014,0.02774) (-0.00017,0.02995) (-0.00019,0.03234) (-0.00023,0.03492) (-0.00026,0.03770) (-0.00031,0.04069) (-0.00036,0.04392) (-0.00042,0.04739) (-0.00049,0.05113) (-0.00058,0.05516) (-0.00068,0.05948) (-0.00079,0.06414) (-0.00092,0.06914) (-0.00108,0.07452) (-0.00126,0.08029) (-0.00148,0.08648) (-0.00173,0.09312) (-0.00202,0.10023) (-0.00236,0.10784) (-0.00276,0.11599) (-0.00322,0.12468) (-0.00377,0.13397) (-0.00441,0.14386) (-0.00515,0.15439) (-0.00602,0.16557) (-0.00704,0.17744) (-0.00823,0.19000) (-0.00962,0.20327) (-0.01125,0.21725) (-0.01315,0.23194) (-0.01538,0.24732) (-0.01798,0.26337) (-0.02102,0.28003) (-0.02457,0.29724) (-0.02873,0.31492) (-0.03359,0.33293) (-0.03927,0.35113) (-0.04591,0.36932) (-0.05367,0.38723) (-0.06275,0.40454) (-0.07336,0.42087) (-0.08577,0.43574) (-0.10028,0.44854) (-0.11724,0.45855) (-0.13706,0.46491) (-0.16024,0.46656) (-0.18734,0.46224) (-0.21902,0.45042) (-0.25607,0.42929) (-0.29937,0.39668) (-0.35000,0.35000)};
\draw[-{Latex},structurecolor] (axis cs:-0.00036,0.04392) -- (axis cs:-0.00042,0.04739);
\draw[-{Latex},structurecolor] (axis cs:-0.01315,0.23194) -- (axis cs:-0.01538,0.24732);
\addplot[structurecolor,thin,no marks] coordinates {(-0.35000,0.35000) (-0.35873,0.34116) (-0.36767,0.33189) (-0.37684,0.32217) (-0.38624,0.31198) (-0.39587,0.30131) (-0.40574,0.29014) (-0.41586,0.27847) (-0.42623,0.26627) (-0.43686,0.25353) (-0.44775,0.24023) (-0.45892,0.22635) (-0.47036,0.21188) (-0.48209,0.19681) (-0.49411,0.18110) (-0.50643,0.16475) (-0.51906,0.14773) (-0.53200,0.13003) (-0.54527,0.11162) (-0.55887,0.09248) (-0.57280,0.07259) (-0.58709,0.05194) (-0.60173,0.03049) (-0.61673,0.00822) (-0.63211,-0.01489) (-0.64787,-0.03886) (-0.66403,-0.06373) (-0.68059,-0.08951) (-0.69756,-0.11623) (-0.71495,-0.14392) (-0.73278,-0.17262) (-0.75105,-0.20233) (-0.76978,-0.23311) (-0.78898,-0.26496) (-0.80865,-0.29794) (-0.82882,-0.33207) (-0.84948,-0.36737) (-0.87067,-0.40390) (-0.89238,-0.44167) (-0.91463,-0.48073) (-0.93744,-0.52110) (-0.96082,-0.56284) (-0.98477,-0.60598) (-1.00933,-0.65055) (-1.03450,-0.69659) (-1.06030,-0.74416) (-1.08674,-0.79328) (-1.11384,-0.84401) (-1.14161,-0.89639) (-1.17008,-0.95046) (-1.19926,-1.00628) (-1.22916,-1.06388) (-1.25981,-1.12332) (-1.29123,-1.18465) (-1.32343,-1.24793) (-1.35643,-1.31320) (-1.39025,-1.38052) (-1.42492,-1.44995) (-1.46045,-1.52154) (-1.49687,-1.59536) (-1.53420,-1.67146) (-1.57245,-1.74991) (-1.61166,-1.83077) (-1.65185,-1.91411) (-1.69304,-2.00000)};
\draw[-{Latex},structurecolor] (axis cs:-0.57280,0.07259) -- (axis cs:-0.58709,0.05194);
\draw[-{Latex},structurecolor] (axis cs:-1.00933,-0.65055) -- (axis cs:-1.03450,-0.69659);
\addplot[structurecolor,thin,no marks] coordinates {(-0.00002,0.00467) (-0.00002,0.00504) (-0.00002,0.00545) (-0.00003,0.00589) (-0.00003,0.00636) (-0.00003,0.00687) (-0.00004,0.00742) (-0.00005,0.00801) (-0.00006,0.00865) (-0.00006,0.00933) (-0.00008,0.01007) (-0.00009,0.01087) (-0.00010,0.01173) (-0.00012,0.01266) (-0.00014,0.01366) (-0.00017,0.01473) (-0.00019,0.01588) (-0.00023,0.01712) (-0.00026,0.01845) (-0.00031,0.01988) (-0.00036,0.02142) (-0.00042,0.02306) (-0.00049,0.02482) (-0.00058,0.02671) (-0.00068,0.02873) (-0.00079,0.03089) (-0.00092,0.03319) (-0.00108,0.03564) (-0.00126,0.03825) (-0.00148,0.04103) (-0.00173,0.04397) (-0.00202,0.04709) (-0.00236,0.05038) (-0.00276,0.05386) (-0.00322,0.05751) (-0.00377,0.06133) (-0.00441,0.06532) (-0.00515,0.06947) (-0.00602,0.07375) (-0.00704,0.07816) (-0.00823,0.08265) (-0.00962,0.08720) (-0.01125,0.09175) (-0.01315,0.09624) (-0.01538,0.10059) (-0.01798,0.10471) (-0.02102,0.10848) (-0.02457,0.11176) (-0.02873,0.11436) (-0.03359,0.11608) (-0.03927,0.11666) (-0.04591,0.11579) (-0.05367,0.11310) (-0.06275,0.10814) (-0.07336,0.10039) (-0.08577,0.08921) (-0.10028,0.07385) (-0.11724,0.05342) (-0.13706,0.02686) (-0.16024,-0.00708) (-0.18734,-0.04989) (-0.21902,-0.10333) (-0.25607,-0.16945) (-0.29937,-0.25072) (-0.35000,-0.35000)};
\draw[-{Latex},structurecolor] (axis cs:-0.00036,0.02142) -- (axis cs:-0.00042,0.02306);
\draw[-{Latex},structurecolor] (axis cs:-0.01315,0.09624) -- (axis cs:-0.01538,0.10059);
\addplot[structurecolor,thin,no marks] coordinates {(-0.35000,-0.35000) (-0.35619,-0.36242) (-0.36250,-0.37511) (-0.36891,-0.38808) (-0.37544,-0.40133) (-0.38209,-0.41488) (-0.38885,-0.42873) (-0.39573,-0.44287) (-0.40274,-0.45733) (-0.40987,-0.47209) (-0.41712,-0.48718) (-0.42450,-0.50260) (-0.43202,-0.51835) (-0.43966,-0.53443) (-0.44744,-0.55086) (-0.45536,-0.56765) (-0.46342,-0.58479) (-0.47162,-0.60230) (-0.47997,-0.62018) (-0.48847,-0.63845) (-0.49711,-0.65710) (-0.50591,-0.67614) (-0.51486,-0.69559) (-0.52398,-0.71545) (-0.53325,-0.73572) (-0.54269,-0.75642) (-0.55229,-0.77756) (-0.56207,-0.79914) (-0.57202,-0.82116) (-0.58214,-0.84365) (-0.59244,-0.86661) (-0.60293,-0.89004) (-0.61360,-0.91396) (-0.62446,-0.93837) (-0.63551,-0.96329) (-0.64676,-0.98873) (-0.65821,-1.01468) (-0.66986,-1.04117) (-0.68171,-1.06821) (-0.69378,-1.09580) (-0.70606,-1.12395) (-0.71856,-1.15268) (-0.73127,-1.18200) (-0.74422,-1.21191) (-0.75739,-1.24243) (-0.77079,-1.27358) (-0.78444,-1.30535) (-0.79832,-1.33777) (-0.81245,-1.37084) (-0.82683,-1.40458) (-0.84146,-1.43901) (-0.85635,-1.47412) (-0.87151,-1.50995) (-0.88694,-1.54649) (-0.90263,-1.58376) (-0.91861,-1.62179) (-0.93487,-1.66057) (-0.95141,-1.70013) (-0.96825,-1.74048) (-0.98539,-1.78163) (-1.00283,-1.82360) (-1.02058,-1.86641) (-1.03864,-1.91007) (-1.05703,-1.95459) (-1.07573,-2.00000)};
\draw[-{Latex},structurecolor] (axis cs:-0.49711,-0.65710) -- (axis cs:-0.50591,-0.67614);
\draw[-{Latex},structurecolor] (axis cs:-0.74422,-1.21191) -- (axis cs:-0.75739,-1.24243);
\addplot[structurecolor,thin,no marks] coordinates {(0.00003,-0.01405) (0.00004,-0.01519) (0.00004,-0.01641) (0.00005,-0.01773) (0.00006,-0.01916) (0.00007,-0.02070) (0.00008,-0.02237) (0.00009,-0.02416) (0.00011,-0.02610) (0.00013,-0.02819) (0.00015,-0.03045) (0.00018,-0.03289) (0.00021,-0.03551) (0.00024,-0.03834) (0.00028,-0.04139) (0.00033,-0.04468) (0.00039,-0.04823) (0.00045,-0.05204) (0.00053,-0.05615) (0.00062,-0.06058) (0.00072,-0.06533) (0.00085,-0.07045) (0.00099,-0.07595) (0.00116,-0.08187) (0.00135,-0.08821) (0.00158,-0.09503) (0.00185,-0.10233) (0.00216,-0.11016) (0.00252,-0.11854) (0.00295,-0.12751) (0.00345,-0.13709) (0.00403,-0.14732) (0.00472,-0.15823) (0.00551,-0.16984) (0.00645,-0.18219) (0.00754,-0.19530) (0.00881,-0.20918) (0.01030,-0.22385) (0.01204,-0.23933) (0.01408,-0.25560) (0.01646,-0.27266) (0.01925,-0.29047) (0.02250,-0.30900) (0.02631,-0.32818) (0.03076,-0.34792) (0.03596,-0.36808) (0.04204,-0.38851) (0.04915,-0.40900) (0.05746,-0.42928) (0.06718,-0.44902) (0.07854,-0.46780) (0.09182,-0.48511) (0.10735,-0.50033) (0.12550,-0.51269) (0.14673,-0.52127) (0.17154,-0.52495) (0.20055,-0.52239) (0.23447,-0.51197) (0.27412,-0.49177) (0.32048,-0.45948) (0.37468,-0.41234) (0.43805,-0.34709) (0.51213,-0.25983) (0.59874,-0.14596) (0.70000,0.00000)};
\draw[-{Latex},structurecolor] (axis cs:0.00072,-0.06533) -- (axis cs:0.00085,-0.07045);
\draw[-{Latex},structurecolor] (axis cs:0.02631,-0.32818) -- (axis cs:0.03076,-0.34792);
\addplot[structurecolor,thin,no marks] coordinates {(0.70000,0.00000) (0.71031,0.01552) (0.72077,0.03138) (0.73138,0.04759) (0.74216,0.06416) (0.75309,0.08108) (0.76418,0.09838) (0.77543,0.11604) (0.78685,0.13408) (0.79844,0.15251) (0.81020,0.17134) (0.82213,0.19056) (0.83424,0.21018) (0.84652,0.23022) (0.85899,0.25068) (0.87164,0.27156) (0.88448,0.29288) (0.89750,0.31464) (0.91072,0.33685) (0.92413,0.35951) (0.93774,0.38264) (0.95155,0.40624) (0.96557,0.43031) (0.97979,0.45488) (0.99422,0.47994) (1.00886,0.50550) (1.02372,0.53158) (1.03879,0.55818) (1.05409,0.58530) (1.06962,0.61297) (1.08537,0.64118) (1.10135,0.66995) (1.11757,0.69928) (1.13403,0.72919) (1.15073,0.75969) (1.16768,0.79077) (1.18488,0.82246) (1.20233,0.85477) (1.22003,0.88770) (1.23800,0.92126) (1.25623,0.95547) (1.27473,0.99033) (1.29351,1.02586) (1.31256,1.06206) (1.33189,1.09896) (1.35150,1.13655) (1.37140,1.17485) (1.39160,1.21388) (1.41210,1.25364) (1.43289,1.29414) (1.45399,1.33541) (1.47541,1.37744) (1.49714,1.42026) (1.51919,1.46387) (1.54156,1.50830) (1.56426,1.55354) (1.58730,1.59962) (1.61067,1.64655) (1.63440,1.69434) (1.65847,1.74301) (1.68289,1.79257) (1.70767,1.84303) (1.73282,1.89441) (1.75834,1.94673) (1.78424,2.00000)};
\draw[-{Latex},structurecolor] (axis cs:0.93774,0.38264) -- (axis cs:0.95155,0.40624);
\draw[-{Latex},structurecolor] (axis cs:1.31256,1.06206) -- (axis cs:1.33189,1.09896);
\addplot[structurecolor,thin,no marks] coordinates {(-0.00003,0.01405) (-0.00004,0.01519) (-0.00004,0.01641) (-0.00005,0.01773) (-0.00006,0.01916) (-0.00007,0.02070) (-0.00008,0.02237) (-0.00009,0.02416) (-0.00011,0.02610) (-0.00013,0.02819) (-0.00015,0.03045) (-0.00018,0.03289) (-0.00021,0.03551) (-0.00024,0.03834) (-0.00028,0.04139) (-0.00033,0.04468) (-0.00039,0.04823) (-0.00045,0.05204) (-0.00053,0.05615) (-0.00062,0.06058) (-0.00072,0.06533) (-0.00085,0.07045) (-0.00099,0.07595) (-0.00116,0.08187) (-0.00135,0.08821) (-0.00158,0.09503) (-0.00185,0.10233) (-0.00216,0.11016) (-0.00252,0.11854) (-0.00295,0.12751) (-0.00345,0.13709) (-0.00403,0.14732) (-0.00472,0.15823) (-0.00551,0.16984) (-0.00645,0.18219) (-0.00754,0.19530) (-0.00881,0.20918) (-0.01030,0.22385) (-0.01204,0.23933) (-0.01408,0.25560) (-0.01646,0.27266) (-0.01925,0.29047) (-0.02250,0.30900) (-0.02631,0.32818) (-0.03076,0.34792) (-0.03596,0.36808) (-0.04204,0.38851) (-0.04915,0.40900) (-0.05746,0.42928) (-0.06718,0.44902) (-0.07854,0.46780) (-0.09182,0.48511) (-0.10735,0.50033) (-0.12550,0.51269) (-0.14673,0.52127) (-0.17154,0.52495) (-0.20055,0.52239) (-0.23447,0.51197) (-0.27412,0.49177) (-0.32048,0.45948) (-0.37468,0.41234) (-0.43805,0.34709) (-0.51213,0.25983) (-0.59874,0.14596) (-0.70000,0.00000)};
\draw[-{Latex},structurecolor] (axis cs:-0.00072,0.06533) -- (axis cs:-0.00085,0.07045);
\draw[-{Latex},structurecolor] (axis cs:-0.02631,0.32818) -- (axis cs:-0.03076,0.34792);
\addplot[structurecolor,thin,no marks] coordinates {(-0.70000,0.00000) (-0.71031,-0.01552) (-0.72077,-0.03138) (-0.73138,-0.04759) (-0.74216,-0.06416) (-0.75309,-0.08108) (-0.76418,-0.09838) (-0.77543,-0.11604) (-0.78685,-0.13408) (-0.79844,-0.15251) (-0.81020,-0.17134) (-0.82213,-0.19056) (-0.83424,-0.21018) (-0.84652,-0.23022) (-0.85899,-0.25068) (-0.87164,-0.27156) (-0.88448,-0.29288) (-0.89750,-0.31464) (-0.91072,-0.33685) (-0.92413,-0.35951) (-0.93774,-0.38264) (-0.95155,-0.40624) (-0.96557,-0.43031) (-0.97979,-0.45488) (-0.99422,-0.47994) (-1.00886,-0.50550) (-1.02372,-0.53158) (-1.03879,-0.55818) (-1.05409,-0.58530) (-1.06962,-0.61297) (-1.08537,-0.64118) (-1.10135,-0.66995) (-1.11757,-0.69928) (-1.13403,-0.72919) (-1.15073,-0.75969) (-1.16768,-0.79077) (-1.18488,-0.82246) (-1.20233,-0.85477) (-1.22003,-0.88770) (-1.23800,-0.92126) (-1.25623,-0.95547) (-1.27473,-0.99033) (-1.29351,-1.02586) (-1.31256,-1.06206) (-1.33189,-1.09896) (-1.35150,-1.13655) (-1.37140,-1.17485) (-1.39160,-1.21388) (-1.41210,-1.25364) (-1.43289,-1.29414) (-1.45399,-1.33541) (-1.47541,-1.37744) (-1.49714,-1.42026) (-1.51919,-1.46387) (-1.54156,-1.50830) (-1.56426,-1.55354) (-1.58730,-1.59962) (-1.61067,-1.64655) (-1.63440,-1.69434) (-1.65847,-1.74301) (-1.68289,-1.79257) (-1.70767,-1.84303) (-1.73282,-1.89441) (-1.75834,-1.94673) (-1.78424,-2.00000)};
\draw[-{Latex},structurecolor] (axis cs:-0.93774,-0.38264) -- (axis cs:-0.95155,-0.40624);
\draw[-{Latex},structurecolor] (axis cs:-1.31256,-1.06206) -- (axis cs:-1.33189,-1.09896);
\addplot[structurecolor,thin,no marks] coordinates {(0.00000,0.00472) (0.00000,0.00510) (0.00000,0.00551) (0.00000,0.00596) (0.00000,0.00645) (0.00000,0.00697) (0.00000,0.00754) (0.00000,0.00815) (0.00000,0.00881) (0.00000,0.00953) (0.00000,0.01030) (0.00000,0.01114) (0.00000,0.01204) (0.00000,0.01302) (0.00000,0.01408) (0.00000,0.01523) (0.00000,0.01646) (0.00000,0.01780) (0.00000,0.01925) (0.00000,0.02081) (0.00000,0.02250) (0.00000,0.02433) (0.00000,0.02631) (0.00000,0.02844) (0.00000,0.03076) (0.00000,0.03326) (0.00000,0.03596) (0.00000,0.03888) (0.00000,0.04204) (0.00000,0.04545) (0.00000,0.04915) (0.00000,0.05314) (0.00000,0.05746) (0.00000,0.06213) (0.00000,0.06718) (0.00000,0.07264) (0.00000,0.07854) (0.00000,0.08492) (0.00000,0.09182) (0.00000,0.09928) (0.00000,0.10735) (0.00000,0.11607) (0.00000,0.12550) (0.00000,0.13570) (0.00000,0.14673) (0.00000,0.15865) (0.00000,0.17154) (0.00000,0.18548) (0.00000,0.20055) (0.00000,0.21685) (0.00000,0.23447) (0.00000,0.25352) (0.00000,0.27412) (0.00000,0.29640) (0.00000,0.32048) (0.00000,0.34653) (0.00000,0.37468) (0.00000,0.40513) (0.00000,0.43805) (0.00000,0.47364) (0.00000,0.51213) (0.00000,0.55375) (0.00000,0.59874) (0.00000,0.64739) (0.00000,0.70000)};
\draw[-{Latex},structurecolor] (axis cs:0.00000,0.02250) -- (axis cs:0.00000,0.02433);
\draw[-{Latex},structurecolor] (axis cs:0.00000,0.13570) -- (axis cs:0.00000,0.14673);
\addplot[structurecolor,thin,no marks] coordinates {(0.00000,0.70000) (0.00000,0.71158) (0.00000,0.72335) (0.00000,0.73531) (0.00000,0.74747) (0.00000,0.75983) (0.00000,0.77240) (0.00000,0.78517) (0.00000,0.79816) (0.00000,0.81136) (0.00000,0.82478) (0.00000,0.83842) (0.00000,0.85229) (0.00000,0.86638) (0.00000,0.88071) (0.00000,0.89528) (0.00000,0.91008) (0.00000,0.92513) (0.00000,0.94044) (0.00000,0.95599) (0.00000,0.97180) (0.00000,0.98787) (0.00000,1.00421) (0.00000,1.02082) (0.00000,1.03770) (0.00000,1.05486) (0.00000,1.07231) (0.00000,1.09004) (0.00000,1.10807) (0.00000,1.12640) (0.00000,1.14503) (0.00000,1.16397) (0.00000,1.18322) (0.00000,1.20278) (0.00000,1.22268) (0.00000,1.24290) (0.00000,1.26345) (0.00000,1.28435) (0.00000,1.30559) (0.00000,1.32719) (0.00000,1.34914) (0.00000,1.37145) (0.00000,1.39413) (0.00000,1.41719) (0.00000,1.44063) (0.00000,1.46445) (0.00000,1.48867) (0.00000,1.51329) (0.00000,1.53832) (0.00000,1.56376) (0.00000,1.58963) (0.00000,1.61592) (0.00000,1.64264) (0.00000,1.66981) (0.00000,1.69743) (0.00000,1.72550) (0.00000,1.75404) (0.00000,1.78305) (0.00000,1.81253) (0.00000,1.84251) (0.00000,1.87298) (0.00000,1.90396) (0.00000,1.93545) (0.00000,1.96746) (0.00000,2.00000)};
\draw[-{Latex},structurecolor] (axis cs:0.00000,0.97180) -- (axis cs:0.00000,0.98787);
\draw[-{Latex},structurecolor] (axis cs:0.00000,1.41719) -- (axis cs:0.00000,1.44063);
\addplot[structurecolor,thin,no marks] coordinates {(0.00000,-0.00472) (0.00000,-0.00510) (0.00000,-0.00551) (0.00000,-0.00596) (0.00000,-0.00645) (0.00000,-0.00697) (0.00000,-0.00754) (0.00000,-0.00815) (0.00000,-0.00881) (0.00000,-0.00953) (0.00000,-0.01030) (0.00000,-0.01114) (0.00000,-0.01204) (0.00000,-0.01302) (0.00000,-0.01408) (0.00000,-0.01523) (0.00000,-0.01646) (0.00000,-0.01780) (0.00000,-0.01925) (0.00000,-0.02081) (0.00000,-0.02250) (0.00000,-0.02433) (0.00000,-0.02631) (0.00000,-0.02844) (0.00000,-0.03076) (0.00000,-0.03326) (0.00000,-0.03596) (0.00000,-0.03888) (0.00000,-0.04204) (0.00000,-0.04545) (0.00000,-0.04915) (0.00000,-0.05314) (0.00000,-0.05746) (0.00000,-0.06213) (0.00000,-0.06718) (0.00000,-0.07264) (0.00000,-0.07854) (0.00000,-0.08492) (0.00000,-0.09182) (0.00000,-0.09928) (0.00000,-0.10735) (0.00000,-0.11607) (0.00000,-0.12550) (0.00000,-0.13570) (0.00000,-0.14673) (0.00000,-0.15865) (0.00000,-0.17154) (0.00000,-0.18548) (0.00000,-0.20055) (0.00000,-0.21685) (0.00000,-0.23447) (0.00000,-0.25352) (0.00000,-0.27412) (0.00000,-0.29640) (0.00000,-0.32048) (0.00000,-0.34653) (0.00000,-0.37468) (0.00000,-0.40513) (0.00000,-0.43805) (0.00000,-0.47364) (0.00000,-0.51213) (0.00000,-0.55375) (0.00000,-0.59874) (0.00000,-0.64739) (0.00000,-0.70000)};
\draw[-{Latex},structurecolor] (axis cs:0.00000,-0.02250) -- (axis cs:0.00000,-0.02433);
\draw[-{Latex},structurecolor] (axis cs:0.00000,-0.13570) -- (axis cs:0.00000,-0.14673);
\addplot[structurecolor,thin,no marks] coordinates {(0.00000,-0.70000) (0.00000,-0.71158) (0.00000,-0.72335) (0.00000,-0.73531) (0.00000,-0.74747) (0.00000,-0.75983) (0.00000,-0.77240) (0.00000,-0.78517) (0.00000,-0.79816) (0.00000,-0.81136) (0.00000,-0.82478) (0.00000,-0.83842) (0.00000,-0.85229) (0.00000,-0.86638) (0.00000,-0.88071) (0.00000,-0.89528) (0.00000,-0.91008) (0.00000,-0.92513) (0.00000,-0.94044) (0.00000,-0.95599) (0.00000,-0.97180) (0.00000,-0.98787) (0.00000,-1.00421) (0.00000,-1.02082) (0.00000,-1.03770) (0.00000,-1.05486) (0.00000,-1.07231) (0.00000,-1.09004) (0.00000,-1.10807) (0.00000,-1.12640) (0.00000,-1.14503) (0.00000,-1.16397) (0.00000,-1.18322) (0.00000,-1.20278) (0.00000,-1.22268) (0.00000,-1.24290) (0.00000,-1.26345) (0.00000,-1.28435) (0.00000,-1.30559) (0.00000,-1.32719) (0.00000,-1.34914) (0.00000,-1.37145) (0.00000,-1.39413) (0.00000,-1.41719) (0.00000,-1.44063) (0.00000,-1.46445) (0.00000,-1.48867) (0.00000,-1.51329) (0.00000,-1.53832) (0.00000,-1.56376) (0.00000,-1.58963) (0.00000,-1.61592) (0.00000,-1.64264) (0.00000,-1.66981) (0.00000,-1.69743) (0.00000,-1.72550) (0.00000,-1.75404) (0.00000,-1.78305) (0.00000,-1.81253) (0.00000,-1.84251) (0.00000,-1.87298) (0.00000,-1.90396) (0.00000,-1.93545) (0.00000,-1.96746) (0.00000,-2.00000)};
\draw[-{Latex},structurecolor] (axis cs:0.00000,-0.97180) -- (axis cs:0.00000,-0.98787);
\draw[-{Latex},structurecolor] (axis cs:0.00000,-1.41719) -- (axis cs:0.00000,-1.44063);
\addplot[only marks,mark=*,mark size=1.4pt] coordinates {(0,0)};
\end{axis}\end{tikzpicture}
\hfill
\begin{tikzpicture}\begin{axis}[width=.47\linewidth,height=5.2cm,axis lines=middle,xlabel={$x$},ylabel={$y$},xmin=-2,xmax=2,ymin=-2,ymax=2,title={（c）稳定结点},tick label style={font=\footnotesize},clip=true]
\addplot[structurecolor,thin,no marks] coordinates {(2.00000,2.00000) (1.94627,1.94627) (1.89398,1.89398) (1.84309,1.84309) (1.79358,1.79358) (1.74539,1.74539) (1.69850,1.69850) (1.65286,1.65286) (1.60846,1.60846) (1.56524,1.56524) (1.52319,1.52319) (1.48227,1.48227) (1.44245,1.44245) (1.40369,1.40369) (1.36598,1.36598) (1.32928,1.32928) (1.29357,1.29357) (1.25882,1.25882) (1.22500,1.22500) (1.19208,1.19208) (1.16006,1.16006) (1.12889,1.12889) (1.09856,1.09856) (1.06905,1.06905) (1.04033,1.04033) (1.01238,1.01238) (0.98518,0.98518) (0.95871,0.95871) (0.93295,0.93295) (0.90789,0.90789) (0.88349,0.88349) (0.85976,0.85976) (0.83666,0.83666) (0.81418,0.81418) (0.79231,0.79231) (0.77102,0.77102) (0.75031,0.75031) (0.73015,0.73015) (0.71053,0.71053) (0.69144,0.69144) (0.67287,0.67287) (0.65479,0.65479) (0.63720,0.63720) (0.62008,0.62008) (0.60342,0.60342) (0.58721,0.58721) (0.57143,0.57143) (0.55608,0.55608) (0.54114,0.54114) (0.52660,0.52660) (0.51245,0.51245) (0.49868,0.49868) (0.48529,0.48529) (0.47225,0.47225) (0.45956,0.45956) (0.44721,0.44721) (0.43520,0.43520) (0.42351,0.42351) (0.41213,0.41213) (0.40106,0.40106) (0.39028,0.39028) (0.37980,0.37980) (0.36959,0.36959) (0.35966,0.35966) (0.35000,0.35000)};
\draw[-{Latex},structurecolor] (axis cs:1.16006,1.16006) -- (axis cs:1.12889,1.12889);
\draw[-{Latex},structurecolor] (axis cs:0.62008,0.62008) -- (axis cs:0.60342,0.60342);
\addplot[structurecolor,thin,no marks] coordinates {(0.35000,0.35000) (0.25607,0.25607) (0.18734,0.18734) (0.13706,0.13706) (0.10028,0.10028) (0.07336,0.07336) (0.05367,0.05367) (0.03927,0.03927) (0.02873,0.02873) (0.02102,0.02102) (0.01538,0.01538) (0.01125,0.01125) (0.00823,0.00823) (0.00602,0.00602) (0.00441,0.00441) (0.00322,0.00322) (0.00236,0.00236) (0.00173,0.00173) (0.00126,0.00126) (0.00092,0.00092) (0.00068,0.00068) (0.00049,0.00049) (0.00036,0.00036) (0.00026,0.00026) (0.00019,0.00019) (0.00014,0.00014) (0.00010,0.00010) (0.00008,0.00008) (0.00006,0.00006) (0.00004,0.00004) (0.00003,0.00003) (0.00002,0.00002) (0.00002,0.00002) (0.00001,0.00001) (0.00001,0.00001) (0.00001,0.00001) (0.00000,0.00000) (0.00000,0.00000) (0.00000,0.00000) (0.00000,0.00000) (0.00000,0.00000) (0.00000,0.00000) (0.00000,0.00000) (0.00000,0.00000) (0.00000,0.00000) (0.00000,0.00000) (0.00000,0.00000) (0.00000,0.00000) (0.00000,0.00000) (0.00000,0.00000) (0.00000,0.00000) (0.00000,0.00000) (0.00000,0.00000) (0.00000,0.00000) (0.00000,0.00000) (0.00000,0.00000) (0.00000,0.00000) (0.00000,0.00000) (0.00000,0.00000) (0.00000,0.00000) (0.00000,0.00000) (0.00000,0.00000) (0.00000,0.00000) (0.00000,0.00000) (0.00000,0.00000)};
\draw[-{Latex},structurecolor] (axis cs:0.00068,0.00068) -- (axis cs:0.00049,0.00049);
\addplot[structurecolor,thin,no marks] coordinates {(-0.66267,-2.00000) (-0.61001,-1.93389) (-0.55971,-1.87026) (-0.51166,-1.80903) (-0.46579,-1.75009) (-0.42198,-1.69336) (-0.38017,-1.63876) (-0.34027,-1.58619) (-0.30220,-1.53558) (-0.26589,-1.48686) (-0.23126,-1.43994) (-0.19825,-1.39477) (-0.16679,-1.35126) (-0.13681,-1.30936) (-0.10826,-1.26901) (-0.08107,-1.23014) (-0.05519,-1.19270) (-0.03056,-1.15662) (-0.00714,-1.12187) (0.01513,-1.08838) (0.03630,-1.05610) (0.05641,-1.02500) (0.07550,-0.99502) (0.09363,-0.96613) (0.11082,-0.93827) (0.12712,-0.91141) (0.14256,-0.88552) (0.15719,-0.86054) (0.17103,-0.83646) (0.18412,-0.81323) (0.19649,-0.79082) (0.20818,-0.76920) (0.21920,-0.74834) (0.22960,-0.72821) (0.23939,-0.70878) (0.24860,-0.69003) (0.25725,-0.67192) (0.26538,-0.65444) (0.27300,-0.63756) (0.28014,-0.62126) (0.28681,-0.60552) (0.29304,-0.59031) (0.29884,-0.57562) (0.30424,-0.56142) (0.30926,-0.54769) (0.31390,-0.53443) (0.31818,-0.52160) (0.32213,-0.50920) (0.32576,-0.49721) (0.32907,-0.48561) (0.33210,-0.47439) (0.33484,-0.46354) (0.33731,-0.45303) (0.33952,-0.44286) (0.34149,-0.43302) (0.34323,-0.42349) (0.34474,-0.41426) (0.34604,-0.40532) (0.34714,-0.39665) (0.34805,-0.38826) (0.34877,-0.38013) (0.34932,-0.37225) (0.34970,-0.36460) (0.34993,-0.35719) (0.35000,-0.35000)};
\draw[-{Latex},structurecolor] (axis cs:0.03630,-1.05610) -- (axis cs:0.05641,-1.02500);
\draw[-{Latex},structurecolor] (axis cs:0.30424,-0.56142) -- (axis cs:0.30926,-0.54769);
\addplot[structurecolor,thin,no marks] coordinates {(0.35000,-0.35000) (0.34624,-0.30115) (0.33671,-0.26203) (0.32348,-0.23027) (0.30800,-0.20414) (0.29131,-0.18234) (0.27414,-0.16391) (0.25700,-0.14813) (0.24021,-0.13447) (0.22401,-0.12251) (0.20853,-0.11195) (0.19385,-0.10255) (0.18001,-0.09412) (0.16701,-0.08652) (0.15485,-0.07963) (0.14349,-0.07336) (0.13292,-0.06764) (0.12308,-0.06240) (0.11394,-0.05760) (0.10546,-0.05319) (0.09759,-0.04913) (0.09030,-0.04540) (0.08355,-0.04195) (0.07729,-0.03878) (0.07150,-0.03585) (0.06614,-0.03314) (0.06118,-0.03064) (0.05659,-0.02833) (0.05234,-0.02620) (0.04841,-0.02423) (0.04477,-0.02240) (0.04141,-0.02072) (0.03830,-0.01916) (0.03542,-0.01772) (0.03276,-0.01639) (0.03030,-0.01515) (0.02802,-0.01401) (0.02592,-0.01296) (0.02397,-0.01199) (0.02217,-0.01109) (0.02050,-0.01025) (0.01896,-0.00948) (0.01754,-0.00877) (0.01622,-0.00811) (0.01500,-0.00750) (0.01387,-0.00694) (0.01283,-0.00642) (0.01187,-0.00593) (0.01097,-0.00549) (0.01015,-0.00508) (0.00939,-0.00469) (0.00868,-0.00434) (0.00803,-0.00401) (0.00743,-0.00371) (0.00687,-0.00343) (0.00635,-0.00318) (0.00587,-0.00294) (0.00543,-0.00272) (0.00502,-0.00251) (0.00465,-0.00232) (0.00430,-0.00215) (0.00397,-0.00199) (0.00368,-0.00184) (0.00340,-0.00170) (0.00314,-0.00157)};
\draw[-{Latex},structurecolor] (axis cs:0.09759,-0.04913) -- (axis cs:0.09030,-0.04540);
\draw[-{Latex},structurecolor] (axis cs:0.01622,-0.00811) -- (axis cs:0.01500,-0.00750);
\addplot[structurecolor,thin,no marks] coordinates {(0.66267,2.00000) (0.61001,1.93389) (0.55971,1.87026) (0.51166,1.80903) (0.46579,1.75009) (0.42198,1.69336) (0.38017,1.63876) (0.34027,1.58619) (0.30220,1.53558) (0.26589,1.48686) (0.23126,1.43994) (0.19825,1.39477) (0.16679,1.35126) (0.13681,1.30936) (0.10826,1.26901) (0.08107,1.23014) (0.05519,1.19270) (0.03056,1.15662) (0.00714,1.12187) (-0.01513,1.08838) (-0.03630,1.05610) (-0.05641,1.02500) (-0.07550,0.99502) (-0.09363,0.96613) (-0.11082,0.93827) (-0.12712,0.91141) (-0.14256,0.88552) (-0.15719,0.86054) (-0.17103,0.83646) (-0.18412,0.81323) (-0.19649,0.79082) (-0.20818,0.76920) (-0.21920,0.74834) (-0.22960,0.72821) (-0.23939,0.70878) (-0.24860,0.69003) (-0.25725,0.67192) (-0.26538,0.65444) (-0.27300,0.63756) (-0.28014,0.62126) (-0.28681,0.60552) (-0.29304,0.59031) (-0.29884,0.57562) (-0.30424,0.56142) (-0.30926,0.54769) (-0.31390,0.53443) (-0.31818,0.52160) (-0.32213,0.50920) (-0.32576,0.49721) (-0.32907,0.48561) (-0.33210,0.47439) (-0.33484,0.46354) (-0.33731,0.45303) (-0.33952,0.44286) (-0.34149,0.43302) (-0.34323,0.42349) (-0.34474,0.41426) (-0.34604,0.40532) (-0.34714,0.39665) (-0.34805,0.38826) (-0.34877,0.38013) (-0.34932,0.37225) (-0.34970,0.36460) (-0.34993,0.35719) (-0.35000,0.35000)};
\draw[-{Latex},structurecolor] (axis cs:-0.03630,1.05610) -- (axis cs:-0.05641,1.02500);
\draw[-{Latex},structurecolor] (axis cs:-0.30424,0.56142) -- (axis cs:-0.30926,0.54769);
\addplot[structurecolor,thin,no marks] coordinates {(-0.35000,0.35000) (-0.34624,0.30115) (-0.33671,0.26203) (-0.32348,0.23027) (-0.30800,0.20414) (-0.29131,0.18234) (-0.27414,0.16391) (-0.25700,0.14813) (-0.24021,0.13447) (-0.22401,0.12251) (-0.20853,0.11195) (-0.19385,0.10255) (-0.18001,0.09412) (-0.16701,0.08652) (-0.15485,0.07963) (-0.14349,0.07336) (-0.13292,0.06764) (-0.12308,0.06240) (-0.11394,0.05760) (-0.10546,0.05319) (-0.09759,0.04913) (-0.09030,0.04540) (-0.08355,0.04195) (-0.07729,0.03878) (-0.07150,0.03585) (-0.06614,0.03314) (-0.06118,0.03064) (-0.05659,0.02833) (-0.05234,0.02620) (-0.04841,0.02423) (-0.04477,0.02240) (-0.04141,0.02072) (-0.03830,0.01916) (-0.03542,0.01772) (-0.03276,0.01639) (-0.03030,0.01515) (-0.02802,0.01401) (-0.02592,0.01296) (-0.02397,0.01199) (-0.02217,0.01109) (-0.02050,0.01025) (-0.01896,0.00948) (-0.01754,0.00877) (-0.01622,0.00811) (-0.01500,0.00750) (-0.01387,0.00694) (-0.01283,0.00642) (-0.01187,0.00593) (-0.01097,0.00549) (-0.01015,0.00508) (-0.00939,0.00469) (-0.00868,0.00434) (-0.00803,0.00401) (-0.00743,0.00371) (-0.00687,0.00343) (-0.00635,0.00318) (-0.00587,0.00294) (-0.00543,0.00272) (-0.00502,0.00251) (-0.00465,0.00232) (-0.00430,0.00215) (-0.00397,0.00199) (-0.00368,0.00184) (-0.00340,0.00170) (-0.00314,0.00157)};
\draw[-{Latex},structurecolor] (axis cs:-0.09759,0.04913) -- (axis cs:-0.09030,0.04540);
\draw[-{Latex},structurecolor] (axis cs:-0.01622,0.00811) -- (axis cs:-0.01500,0.00750);
\addplot[structurecolor,thin,no marks] coordinates {(-2.00000,-2.00000) (-1.94627,-1.94627) (-1.89398,-1.89398) (-1.84309,-1.84309) (-1.79358,-1.79358) (-1.74539,-1.74539) (-1.69850,-1.69850) (-1.65286,-1.65286) (-1.60846,-1.60846) (-1.56524,-1.56524) (-1.52319,-1.52319) (-1.48227,-1.48227) (-1.44245,-1.44245) (-1.40369,-1.40369) (-1.36598,-1.36598) (-1.32928,-1.32928) (-1.29357,-1.29357) (-1.25882,-1.25882) (-1.22500,-1.22500) (-1.19208,-1.19208) (-1.16006,-1.16006) (-1.12889,-1.12889) (-1.09856,-1.09856) (-1.06905,-1.06905) (-1.04033,-1.04033) (-1.01238,-1.01238) (-0.98518,-0.98518) (-0.95871,-0.95871) (-0.93295,-0.93295) (-0.90789,-0.90789) (-0.88349,-0.88349) (-0.85976,-0.85976) (-0.83666,-0.83666) (-0.81418,-0.81418) (-0.79231,-0.79231) (-0.77102,-0.77102) (-0.75031,-0.75031) (-0.73015,-0.73015) (-0.71053,-0.71053) (-0.69144,-0.69144) (-0.67287,-0.67287) (-0.65479,-0.65479) (-0.63720,-0.63720) (-0.62008,-0.62008) (-0.60342,-0.60342) (-0.58721,-0.58721) (-0.57143,-0.57143) (-0.55608,-0.55608) (-0.54114,-0.54114) (-0.52660,-0.52660) (-0.51245,-0.51245) (-0.49868,-0.49868) (-0.48529,-0.48529) (-0.47225,-0.47225) (-0.45956,-0.45956) (-0.44721,-0.44721) (-0.43520,-0.43520) (-0.42351,-0.42351) (-0.41213,-0.41213) (-0.40106,-0.40106) (-0.39028,-0.39028) (-0.37980,-0.37980) (-0.36959,-0.36959) (-0.35966,-0.35966) (-0.35000,-0.35000)};
\draw[-{Latex},structurecolor] (axis cs:-1.16006,-1.16006) -- (axis cs:-1.12889,-1.12889);
\draw[-{Latex},structurecolor] (axis cs:-0.62008,-0.62008) -- (axis cs:-0.60342,-0.60342);
\addplot[structurecolor,thin,no marks] coordinates {(-0.35000,-0.35000) (-0.25607,-0.25607) (-0.18734,-0.18734) (-0.13706,-0.13706) (-0.10028,-0.10028) (-0.07336,-0.07336) (-0.05367,-0.05367) (-0.03927,-0.03927) (-0.02873,-0.02873) (-0.02102,-0.02102) (-0.01538,-0.01538) (-0.01125,-0.01125) (-0.00823,-0.00823) (-0.00602,-0.00602) (-0.00441,-0.00441) (-0.00322,-0.00322) (-0.00236,-0.00236) (-0.00173,-0.00173) (-0.00126,-0.00126) (-0.00092,-0.00092) (-0.00068,-0.00068) (-0.00049,-0.00049) (-0.00036,-0.00036) (-0.00026,-0.00026) (-0.00019,-0.00019) (-0.00014,-0.00014) (-0.00010,-0.00010) (-0.00008,-0.00008) (-0.00006,-0.00006) (-0.00004,-0.00004) (-0.00003,-0.00003) (-0.00002,-0.00002) (-0.00002,-0.00002) (-0.00001,-0.00001) (-0.00001,-0.00001) (-0.00001,-0.00001) (-0.00000,-0.00000) (-0.00000,-0.00000) (-0.00000,-0.00000) (-0.00000,-0.00000) (-0.00000,-0.00000) (-0.00000,-0.00000) (-0.00000,-0.00000) (-0.00000,-0.00000) (-0.00000,-0.00000) (-0.00000,-0.00000) (-0.00000,-0.00000) (-0.00000,-0.00000) (-0.00000,-0.00000) (-0.00000,-0.00000) (-0.00000,-0.00000) (-0.00000,-0.00000) (-0.00000,-0.00000) (-0.00000,-0.00000) (-0.00000,-0.00000) (-0.00000,-0.00000) (-0.00000,-0.00000) (-0.00000,-0.00000) (-0.00000,-0.00000) (-0.00000,-0.00000) (-0.00000,-0.00000) (-0.00000,-0.00000) (-0.00000,-0.00000) (-0.00000,-0.00000) (-0.00000,-0.00000)};
\draw[-{Latex},structurecolor] (axis cs:-0.00068,-0.00068) -- (axis cs:-0.00049,-0.00049);
\addplot[structurecolor,thin,no marks] coordinates {(2.00000,0.92765) (1.96145,0.89622) (1.92381,0.86566) (1.88708,0.83595) (1.85121,0.80707) (1.81620,0.77900) (1.78202,0.75170) (1.74865,0.72518) (1.71607,0.69939) (1.68425,0.67433) (1.65319,0.64997) (1.62285,0.62630) (1.59322,0.60329) (1.56429,0.58093) (1.53603,0.55920) (1.50843,0.53809) (1.48147,0.51758) (1.45513,0.49764) (1.42940,0.47828) (1.40427,0.45946) (1.37972,0.44118) (1.35572,0.42342) (1.33228,0.40617) (1.30938,0.38942) (1.28699,0.37315) (1.26512,0.35734) (1.24374,0.34199) (1.22284,0.32708) (1.20241,0.31261) (1.18245,0.29855) (1.16293,0.28490) (1.14385,0.27165) (1.12519,0.25879) (1.10695,0.24630) (1.08911,0.23418) (1.07167,0.22242) (1.05461,0.21100) (1.03793,0.19992) (1.02161,0.18917) (1.00565,0.17874) (0.99003,0.16862) (0.97476,0.15880) (0.95981,0.14927) (0.94519,0.14003) (0.93088,0.13107) (0.91688,0.12238) (0.90317,0.11395) (0.88976,0.10578) (0.87663,0.09786) (0.86378,0.09018) (0.85120,0.08274) (0.83888,0.07553) (0.82682,0.06854) (0.81501,0.06176) (0.80344,0.05520) (0.79212,0.04884) (0.78102,0.04269) (0.77016,0.03673) (0.75951,0.03095) (0.74908,0.02536) (0.73886,0.01995) (0.72885,0.01472) (0.71904,0.00965) (0.70942,0.00474) (0.70000,0.00000)};
\draw[-{Latex},structurecolor] (axis cs:1.37972,0.44118) -- (axis cs:1.35572,0.42342);
\draw[-{Latex},structurecolor] (axis cs:0.94519,0.14003) -- (axis cs:0.93088,0.13107);
\addplot[structurecolor,thin,no marks] coordinates {(0.70000,0.00000) (0.60231,-0.04509) (0.52406,-0.07469) (0.46054,-0.09321) (0.40827,-0.10386) (0.36467,-0.10897) (0.32782,-0.11023) (0.29627,-0.10886) (0.26894,-0.10574) (0.24503,-0.10150) (0.22391,-0.09658) (0.20510,-0.09130) (0.18824,-0.08589) (0.17303,-0.08049) (0.15925,-0.07522) (0.14672,-0.07013) (0.13527,-0.06528) (0.12480,-0.06068) (0.11520,-0.05634) (0.10638,-0.05227) (0.09827,-0.04846) (0.09080,-0.04490) (0.08391,-0.04159) (0.07756,-0.03851) (0.07169,-0.03565) (0.06628,-0.03300) (0.06128,-0.03054) (0.05666,-0.02826) (0.05240,-0.02614) (0.04845,-0.02418) (0.04480,-0.02237) (0.04143,-0.02070) (0.03832,-0.01914) (0.03544,-0.01771) (0.03277,-0.01638) (0.03031,-0.01515) (0.02803,-0.01401) (0.02592,-0.01296) (0.02397,-0.01198) (0.02217,-0.01108) (0.02050,-0.01025) (0.01896,-0.00948) (0.01754,-0.00877) (0.01622,-0.00811) (0.01500,-0.00750) (0.01387,-0.00694) (0.01283,-0.00642) (0.01187,-0.00593) (0.01098,-0.00549) (0.01015,-0.00508) (0.00939,-0.00469) (0.00868,-0.00434) (0.00803,-0.00401) (0.00743,-0.00371) (0.00687,-0.00343) (0.00635,-0.00318) (0.00587,-0.00294) (0.00543,-0.00272) (0.00502,-0.00251) (0.00465,-0.00232) (0.00430,-0.00215) (0.00397,-0.00199) (0.00368,-0.00184) (0.00340,-0.00170) (0.00314,-0.00157)};
\draw[-{Latex},structurecolor] (axis cs:0.09827,-0.04846) -- (axis cs:0.09080,-0.04490);
\draw[-{Latex},structurecolor] (axis cs:0.01622,-0.00811) -- (axis cs:0.01500,-0.00750);
\addplot[structurecolor,thin,no marks] coordinates {(-2.00000,-0.92765) (-1.96145,-0.89622) (-1.92381,-0.86566) (-1.88708,-0.83595) (-1.85121,-0.80707) (-1.81620,-0.77900) (-1.78202,-0.75170) (-1.74865,-0.72518) (-1.71607,-0.69939) (-1.68425,-0.67433) (-1.65319,-0.64997) (-1.62285,-0.62630) (-1.59322,-0.60329) (-1.56429,-0.58093) (-1.53603,-0.55920) (-1.50843,-0.53809) (-1.48147,-0.51758) (-1.45513,-0.49764) (-1.42940,-0.47828) (-1.40427,-0.45946) (-1.37972,-0.44118) (-1.35572,-0.42342) (-1.33228,-0.40617) (-1.30938,-0.38942) (-1.28699,-0.37315) (-1.26512,-0.35734) (-1.24374,-0.34199) (-1.22284,-0.32708) (-1.20241,-0.31261) (-1.18245,-0.29855) (-1.16293,-0.28490) (-1.14385,-0.27165) (-1.12519,-0.25879) (-1.10695,-0.24630) (-1.08911,-0.23418) (-1.07167,-0.22242) (-1.05461,-0.21100) (-1.03793,-0.19992) (-1.02161,-0.18917) (-1.00565,-0.17874) (-0.99003,-0.16862) (-0.97476,-0.15880) (-0.95981,-0.14927) (-0.94519,-0.14003) (-0.93088,-0.13107) (-0.91688,-0.12238) (-0.90317,-0.11395) (-0.88976,-0.10578) (-0.87663,-0.09786) (-0.86378,-0.09018) (-0.85120,-0.08274) (-0.83888,-0.07553) (-0.82682,-0.06854) (-0.81501,-0.06176) (-0.80344,-0.05520) (-0.79212,-0.04884) (-0.78102,-0.04269) (-0.77016,-0.03673) (-0.75951,-0.03095) (-0.74908,-0.02536) (-0.73886,-0.01995) (-0.72885,-0.01472) (-0.71904,-0.00965) (-0.70942,-0.00474) (-0.70000,0.00000)};
\draw[-{Latex},structurecolor] (axis cs:-1.37972,-0.44118) -- (axis cs:-1.35572,-0.42342);
\draw[-{Latex},structurecolor] (axis cs:-0.94519,-0.14003) -- (axis cs:-0.93088,-0.13107);
\addplot[structurecolor,thin,no marks] coordinates {(-0.70000,0.00000) (-0.60231,0.04509) (-0.52406,0.07469) (-0.46054,0.09321) (-0.40827,0.10386) (-0.36467,0.10897) (-0.32782,0.11023) (-0.29627,0.10886) (-0.26894,0.10574) (-0.24503,0.10150) (-0.22391,0.09658) (-0.20510,0.09130) (-0.18824,0.08589) (-0.17303,0.08049) (-0.15925,0.07522) (-0.14672,0.07013) (-0.13527,0.06528) (-0.12480,0.06068) (-0.11520,0.05634) (-0.10638,0.05227) (-0.09827,0.04846) (-0.09080,0.04490) (-0.08391,0.04159) (-0.07756,0.03851) (-0.07169,0.03565) (-0.06628,0.03300) (-0.06128,0.03054) (-0.05666,0.02826) (-0.05240,0.02614) (-0.04845,0.02418) (-0.04480,0.02237) (-0.04143,0.02070) (-0.03832,0.01914) (-0.03544,0.01771) (-0.03277,0.01638) (-0.03031,0.01515) (-0.02803,0.01401) (-0.02592,0.01296) (-0.02397,0.01198) (-0.02217,0.01108) (-0.02050,0.01025) (-0.01896,0.00948) (-0.01754,0.00877) (-0.01622,0.00811) (-0.01500,0.00750) (-0.01387,0.00694) (-0.01283,0.00642) (-0.01187,0.00593) (-0.01098,0.00549) (-0.01015,0.00508) (-0.00939,0.00469) (-0.00868,0.00434) (-0.00803,0.00401) (-0.00743,0.00371) (-0.00687,0.00343) (-0.00635,0.00318) (-0.00587,0.00294) (-0.00543,0.00272) (-0.00502,0.00251) (-0.00465,0.00232) (-0.00430,0.00215) (-0.00397,0.00199) (-0.00368,0.00184) (-0.00340,0.00170) (-0.00314,0.00157)};
\draw[-{Latex},structurecolor] (axis cs:-0.09827,0.04846) -- (axis cs:-0.09080,0.04490);
\draw[-{Latex},structurecolor] (axis cs:-0.01622,0.00811) -- (axis cs:-0.01500,0.00750);
\addplot[structurecolor,thin,no marks] coordinates {(1.03598,2.00000) (1.00594,1.96515) (0.97655,1.93097) (0.94778,1.89745) (0.91963,1.86456) (0.89209,1.83230) (0.86513,1.80066) (0.83876,1.76962) (0.81296,1.73917) (0.78771,1.70931) (0.76301,1.68001) (0.73884,1.65127) (0.71520,1.62307) (0.69207,1.59541) (0.66944,1.56828) (0.64730,1.54166) (0.62565,1.51555) (0.60447,1.48993) (0.58376,1.46480) (0.56349,1.44014) (0.54368,1.41595) (0.52429,1.39222) (0.50534,1.36893) (0.48680,1.34609) (0.46867,1.32367) (0.45095,1.30168) (0.43361,1.28011) (0.41667,1.25894) (0.40010,1.23816) (0.38390,1.21778) (0.36806,1.19779) (0.35257,1.17816) (0.33744,1.15891) (0.32264,1.14001) (0.30818,1.12148) (0.29404,1.10328) (0.28023,1.08543) (0.26672,1.06791) (0.25353,1.05072) (0.24063,1.03385) (0.22803,1.01729) (0.21572,1.00104) (0.20369,0.98510) (0.19194,0.96945) (0.18046,0.95409) (0.16925,0.93902) (0.15829,0.92422) (0.14759,0.90970) (0.13714,0.89545) (0.12694,0.88146) (0.11697,0.86773) (0.10724,0.85426) (0.09774,0.84103) (0.08846,0.82804) (0.07940,0.81530) (0.07056,0.80279) (0.06193,0.79050) (0.05351,0.77845) (0.04529,0.76661) (0.03727,0.75499) (0.02944,0.74359) (0.02181,0.73239) (0.01436,0.72139) (0.00709,0.71060) (0.00000,0.70000)};
\draw[-{Latex},structurecolor] (axis cs:0.54368,1.41595) -- (axis cs:0.52429,1.39222);
\draw[-{Latex},structurecolor] (axis cs:0.19194,0.96945) -- (axis cs:0.18046,0.95409);
\addplot[structurecolor,thin,no marks] coordinates {(0.00000,0.70000) (-0.09018,0.55722) (-0.14937,0.44937) (-0.18641,0.36733) (-0.20772,0.30441) (-0.21794,0.25570) (-0.22047,0.21758) (-0.21773,0.18740) (-0.21148,0.16320) (-0.20299,0.14353) (-0.19315,0.12733) (-0.18260,0.11380) (-0.17177,0.10235) (-0.16099,0.09254) (-0.15044,0.08403) (-0.14027,0.07658) (-0.13056,0.07000) (-0.12135,0.06413) (-0.11268,0.05886) (-0.10454,0.05411) (-0.09692,0.04981) (-0.08981,0.04589) (-0.08319,0.04232) (-0.07703,0.03904) (-0.07131,0.03604) (-0.06600,0.03328) (-0.06108,0.03074) (-0.05651,0.02841) (-0.05228,0.02625) (-0.04837,0.02427) (-0.04475,0.02243) (-0.04139,0.02074) (-0.03829,0.01917) (-0.03541,0.01773) (-0.03275,0.01639) (-0.03029,0.01516) (-0.02802,0.01402) (-0.02591,0.01296) (-0.02397,0.01199) (-0.02217,0.01109) (-0.02050,0.01025) (-0.01896,0.00948) (-0.01754,0.00877) (-0.01622,0.00811) (-0.01500,0.00750) (-0.01387,0.00694) (-0.01283,0.00642) (-0.01187,0.00593) (-0.01097,0.00549) (-0.01015,0.00508) (-0.00939,0.00469) (-0.00868,0.00434) (-0.00803,0.00401) (-0.00743,0.00371) (-0.00687,0.00343) (-0.00635,0.00318) (-0.00587,0.00294) (-0.00543,0.00272) (-0.00502,0.00251) (-0.00465,0.00232) (-0.00430,0.00215) (-0.00397,0.00199) (-0.00368,0.00184) (-0.00340,0.00170) (-0.00314,0.00157)};
\draw[-{Latex},structurecolor] (axis cs:-0.09692,0.04981) -- (axis cs:-0.08981,0.04589);
\draw[-{Latex},structurecolor] (axis cs:-0.01622,0.00811) -- (axis cs:-0.01500,0.00750);
\addplot[structurecolor,thin,no marks] coordinates {(-1.03598,-2.00000) (-1.00594,-1.96515) (-0.97655,-1.93097) (-0.94778,-1.89745) (-0.91963,-1.86456) (-0.89209,-1.83230) (-0.86513,-1.80066) (-0.83876,-1.76962) (-0.81296,-1.73917) (-0.78771,-1.70931) (-0.76301,-1.68001) (-0.73884,-1.65127) (-0.71520,-1.62307) (-0.69207,-1.59541) (-0.66944,-1.56828) (-0.64730,-1.54166) (-0.62565,-1.51555) (-0.60447,-1.48993) (-0.58376,-1.46480) (-0.56349,-1.44014) (-0.54368,-1.41595) (-0.52429,-1.39222) (-0.50534,-1.36893) (-0.48680,-1.34609) (-0.46867,-1.32367) (-0.45095,-1.30168) (-0.43361,-1.28011) (-0.41667,-1.25894) (-0.40010,-1.23816) (-0.38390,-1.21778) (-0.36806,-1.19779) (-0.35257,-1.17816) (-0.33744,-1.15891) (-0.32264,-1.14001) (-0.30818,-1.12148) (-0.29404,-1.10328) (-0.28023,-1.08543) (-0.26672,-1.06791) (-0.25353,-1.05072) (-0.24063,-1.03385) (-0.22803,-1.01729) (-0.21572,-1.00104) (-0.20369,-0.98510) (-0.19194,-0.96945) (-0.18046,-0.95409) (-0.16925,-0.93902) (-0.15829,-0.92422) (-0.14759,-0.90970) (-0.13714,-0.89545) (-0.12694,-0.88146) (-0.11697,-0.86773) (-0.10724,-0.85426) (-0.09774,-0.84103) (-0.08846,-0.82804) (-0.07940,-0.81530) (-0.07056,-0.80279) (-0.06193,-0.79050) (-0.05351,-0.77845) (-0.04529,-0.76661) (-0.03727,-0.75499) (-0.02944,-0.74359) (-0.02181,-0.73239) (-0.01436,-0.72139) (-0.00709,-0.71060) (0.00000,-0.70000)};
\draw[-{Latex},structurecolor] (axis cs:-0.54368,-1.41595) -- (axis cs:-0.52429,-1.39222);
\draw[-{Latex},structurecolor] (axis cs:-0.19194,-0.96945) -- (axis cs:-0.18046,-0.95409);
\addplot[structurecolor,thin,no marks] coordinates {(0.00000,-0.70000) (0.09018,-0.55722) (0.14937,-0.44937) (0.18641,-0.36733) (0.20772,-0.30441) (0.21794,-0.25570) (0.22047,-0.21758) (0.21773,-0.18740) (0.21148,-0.16320) (0.20299,-0.14353) (0.19315,-0.12733) (0.18260,-0.11380) (0.17177,-0.10235) (0.16099,-0.09254) (0.15044,-0.08403) (0.14027,-0.07658) (0.13056,-0.07000) (0.12135,-0.06413) (0.11268,-0.05886) (0.10454,-0.05411) (0.09692,-0.04981) (0.08981,-0.04589) (0.08319,-0.04232) (0.07703,-0.03904) (0.07131,-0.03604) (0.06600,-0.03328) (0.06108,-0.03074) (0.05651,-0.02841) (0.05228,-0.02625) (0.04837,-0.02427) (0.04475,-0.02243) (0.04139,-0.02074) (0.03829,-0.01917) (0.03541,-0.01773) (0.03275,-0.01639) (0.03029,-0.01516) (0.02802,-0.01402) (0.02591,-0.01296) (0.02397,-0.01199) (0.02217,-0.01109) (0.02050,-0.01025) (0.01896,-0.00948) (0.01754,-0.00877) (0.01622,-0.00811) (0.01500,-0.00750) (0.01387,-0.00694) (0.01283,-0.00642) (0.01187,-0.00593) (0.01097,-0.00549) (0.01015,-0.00508) (0.00939,-0.00469) (0.00868,-0.00434) (0.00803,-0.00401) (0.00743,-0.00371) (0.00687,-0.00343) (0.00635,-0.00318) (0.00587,-0.00294) (0.00543,-0.00272) (0.00502,-0.00251) (0.00465,-0.00232) (0.00430,-0.00215) (0.00397,-0.00199) (0.00368,-0.00184) (0.00340,-0.00170) (0.00314,-0.00157)};
\draw[-{Latex},structurecolor] (axis cs:0.09692,-0.04981) -- (axis cs:0.08981,-0.04589);
\draw[-{Latex},structurecolor] (axis cs:0.01622,-0.00811) -- (axis cs:0.01500,-0.00750);
\addplot[only marks,mark=*,mark size=1.4pt] coordinates {(0,0)};
\end{axis}\end{tikzpicture}
\hfill
\begin{tikzpicture}\begin{axis}[width=.47\linewidth,height=5.2cm,axis lines=middle,xlabel={$x$},ylabel={$y$},xmin=-2,xmax=2,ymin=-2,ymax=2,title={（d）向外逆时针},tick label style={font=\footnotesize},clip=true]
\addplot[structurecolor,thin,no marks] coordinates {(0.00403,0.00048) (0.00425,0.00086) (0.00442,0.00128) (0.00453,0.00175) (0.00458,0.00226) (0.00454,0.00281) (0.00440,0.00341) (0.00416,0.00403) (0.00379,0.00468) (0.00328,0.00535) (0.00262,0.00603) (0.00180,0.00670) (0.00081,0.00735) (-0.00037,0.00797) (-0.00174,0.00853) (-0.00331,0.00902) (-0.00508,0.00942) (-0.00705,0.00969) (-0.00921,0.00982) (-0.01155,0.00978) (-0.01406,0.00954) (-0.01672,0.00906) (-0.01950,0.00833) (-0.02236,0.00731) (-0.02526,0.00597) (-0.02815,0.00429) (-0.03098,0.00223) (-0.03367,-0.00021) (-0.03615,-0.00307) (-0.03833,-0.00634) (-0.04012,-0.01005) (-0.04142,-0.01418) (-0.04212,-0.01873) (-0.04211,-0.02368) (-0.04126,-0.02900) (-0.03945,-0.03464) (-0.03656,-0.04056) (-0.03245,-0.04667) (-0.02701,-0.05289) (-0.02011,-0.05911) (-0.01164,-0.06522) (-0.00151,-0.07107) (0.01036,-0.07650) (0.02403,-0.08134) (0.03954,-0.08539) (0.05689,-0.08843) (0.07605,-0.09024) (0.09694,-0.09057) (0.11945,-0.08916) (0.14340,-0.08575) (0.16856,-0.08006) (0.19465,-0.07183) (0.22129,-0.06078) (0.24806,-0.04667) (0.27445,-0.02924) (0.29987,-0.00828) (0.32366,0.01638) (0.34506,0.04488) (0.36326,0.07732) (0.37734,0.11372) (0.38635,0.15402) (0.38922,0.19809) (0.38489,0.24569) (0.37220,0.29649) (0.35000,0.35000)};
\draw[-{Latex},structurecolor] (axis cs:-0.01406,0.00954) -- (axis cs:-0.01672,0.00906);
\draw[-{Latex},structurecolor] (axis cs:0.02403,-0.08134) -- (axis cs:0.03954,-0.08539);
\addplot[structurecolor,thin,no marks] coordinates {(0.35000,0.35000) (0.34270,0.36397) (0.33469,0.37807) (0.32595,0.39228) (0.31647,0.40659) (0.30621,0.42100) (0.29518,0.43548) (0.28334,0.45002) (0.27067,0.46460) (0.25717,0.47922) (0.24281,0.49385) (0.22758,0.50848) (0.21145,0.52308) (0.19441,0.53765) (0.17645,0.55215) (0.15755,0.56658) (0.13769,0.58091) (0.11686,0.59512) (0.09505,0.60919) (0.07224,0.62309) (0.04841,0.63680) (0.02356,0.65030) (-0.00232,0.66356) (-0.02925,0.67656) (-0.05723,0.68926) (-0.08629,0.70165) (-0.11641,0.71370) (-0.14761,0.72537) (-0.17990,0.73664) (-0.21328,0.74748) (-0.24776,0.75785) (-0.28334,0.76773) (-0.32001,0.77708) (-0.35778,0.78587) (-0.39666,0.79408) (-0.43663,0.80165) (-0.47769,0.80856) (-0.51984,0.81478) (-0.56307,0.82027) (-0.60737,0.82499) (-0.65273,0.82890) (-0.69915,0.83197) (-0.74660,0.83416) (-0.79507,0.83544) (-0.84455,0.83575) (-0.89501,0.83507) (-0.94644,0.83336) (-0.99881,0.83057) (-1.05209,0.82666) (-1.10626,0.82160) (-1.16130,0.81534) (-1.21716,0.80785) (-1.27382,0.79907) (-1.33123,0.78897) (-1.38937,0.77751) (-1.44818,0.76465) (-1.50763,0.75035) (-1.56767,0.73456) (-1.62824,0.71724) (-1.68931,0.69836) (-1.75081,0.67786) (-1.81269,0.65572) (-1.87489,0.63189) (-1.93735,0.60633) (-2.00000,0.57901)};
\draw[-{Latex},structurecolor] (axis cs:0.04841,0.63680) -- (axis cs:0.02356,0.65030);
\draw[-{Latex},structurecolor] (axis cs:-0.79507,0.83544) -- (axis cs:-0.84455,0.83575);
\addplot[structurecolor,thin,no marks] coordinates {(-0.00070,-0.00285) (-0.00027,-0.00312) (0.00023,-0.00337) (0.00082,-0.00361) (0.00149,-0.00381) (0.00224,-0.00396) (0.00307,-0.00407) (0.00399,-0.00412) (0.00498,-0.00409) (0.00604,-0.00397) (0.00716,-0.00376) (0.00833,-0.00344) (0.00954,-0.00299) (0.01075,-0.00241) (0.01196,-0.00168) (0.01314,-0.00080) (0.01425,0.00025) (0.01527,0.00147) (0.01617,0.00286) (0.01689,0.00444) (0.01740,0.00620) (0.01766,0.00813) (0.01761,0.01022) (0.01720,0.01247) (0.01638,0.01485) (0.01510,0.01734) (0.01330,0.01991) (0.01094,0.02252) (0.00796,0.02512) (0.00432,0.02767) (-0.00002,0.03010) (-0.00510,0.03234) (-0.01093,0.03433) (-0.01754,0.03597) (-0.02491,0.03717) (-0.03304,0.03785) (-0.04189,0.03789) (-0.05140,0.03719) (-0.06151,0.03563) (-0.07211,0.03311) (-0.08308,0.02950) (-0.09426,0.02470) (-0.10546,0.01860) (-0.11647,0.01110) (-0.12704,0.00211) (-0.13688,-0.00845) (-0.14567,-0.02062) (-0.15307,-0.03444) (-0.15869,-0.04991) (-0.16213,-0.06702) (-0.16293,-0.08569) (-0.16065,-0.10582) (-0.15481,-0.12727) (-0.14491,-0.14982) (-0.13047,-0.17322) (-0.11100,-0.19715) (-0.08605,-0.22123) (-0.05518,-0.24500) (-0.01799,-0.26795) (0.02585,-0.28947) (0.07658,-0.30890) (0.13438,-0.32551) (0.19930,-0.33849) (0.27125,-0.34696) (0.35000,-0.35000)};
\draw[-{Latex},structurecolor] (axis cs:0.01740,0.00620) -- (axis cs:0.01766,0.00813);
\draw[-{Latex},structurecolor] (axis cs:-0.11647,0.01110) -- (axis cs:-0.12704,0.00211);
\addplot[structurecolor,thin,no marks] coordinates {(0.35000,-0.35000) (0.37628,-0.34967) (0.40318,-0.34867) (0.43069,-0.34697) (0.45879,-0.34452) (0.48744,-0.34131) (0.51664,-0.33729) (0.54634,-0.33243) (0.57652,-0.32670) (0.60715,-0.32006) (0.63818,-0.31249) (0.66960,-0.30394) (0.70134,-0.29440) (0.73338,-0.28381) (0.76566,-0.27215) (0.79814,-0.25940) (0.83076,-0.24550) (0.86348,-0.23045) (0.89623,-0.21419) (0.92896,-0.19671) (0.96160,-0.17797) (0.99409,-0.15795) (1.02636,-0.13662) (1.05833,-0.11394) (1.08994,-0.08991) (1.12109,-0.06449) (1.15172,-0.03765) (1.18174,-0.00939) (1.21106,0.02033) (1.23958,0.05151) (1.26723,0.08417) (1.29389,0.11834) (1.31947,0.15401) (1.34387,0.19120) (1.36698,0.22992) (1.38869,0.27017) (1.40889,0.31195) (1.42747,0.35527) (1.44431,0.40011) (1.45928,0.44647) (1.47227,0.49434) (1.48315,0.54371) (1.49180,0.59456) (1.49807,0.64687) (1.50186,0.70062) (1.50301,0.75577) (1.50139,0.81230) (1.49687,0.87017) (1.48930,0.92934) (1.47855,0.98976) (1.46448,1.05138) (1.44694,1.11414) (1.42580,1.17800) (1.40090,1.24288) (1.37210,1.30871) (1.33927,1.37542) (1.30226,1.44292) (1.26092,1.51113) (1.21511,1.57996) (1.16470,1.64932) (1.10954,1.71908) (1.04951,1.78915) (0.98445,1.85941) (0.91425,1.92974) (0.83877,2.00000)};
\draw[-{Latex},structurecolor] (axis cs:0.96160,-0.17797) -- (axis cs:0.99409,-0.15795);
\draw[-{Latex},structurecolor] (axis cs:1.49807,0.64687) -- (axis cs:1.50186,0.70062);
\addplot[structurecolor,thin,no marks] coordinates {(0.00070,0.00285) (0.00027,0.00312) (-0.00023,0.00337) (-0.00082,0.00361) (-0.00149,0.00381) (-0.00224,0.00396) (-0.00307,0.00407) (-0.00399,0.00412) (-0.00498,0.00409) (-0.00604,0.00397) (-0.00716,0.00376) (-0.00833,0.00344) (-0.00954,0.00299) (-0.01075,0.00241) (-0.01196,0.00168) (-0.01314,0.00080) (-0.01425,-0.00025) (-0.01527,-0.00147) (-0.01617,-0.00286) (-0.01689,-0.00444) (-0.01740,-0.00620) (-0.01766,-0.00813) (-0.01761,-0.01022) (-0.01720,-0.01247) (-0.01638,-0.01485) (-0.01510,-0.01734) (-0.01330,-0.01991) (-0.01094,-0.02252) (-0.00796,-0.02512) (-0.00432,-0.02767) (0.00002,-0.03010) (0.00510,-0.03234) (0.01093,-0.03433) (0.01754,-0.03597) (0.02491,-0.03717) (0.03304,-0.03785) (0.04189,-0.03789) (0.05140,-0.03719) (0.06151,-0.03563) (0.07211,-0.03311) (0.08308,-0.02950) (0.09426,-0.02470) (0.10546,-0.01860) (0.11647,-0.01110) (0.12704,-0.00211) (0.13688,0.00845) (0.14567,0.02062) (0.15307,0.03444) (0.15869,0.04991) (0.16213,0.06702) (0.16293,0.08569) (0.16065,0.10582) (0.15481,0.12727) (0.14491,0.14982) (0.13047,0.17322) (0.11100,0.19715) (0.08605,0.22123) (0.05518,0.24500) (0.01799,0.26795) (-0.02585,0.28947) (-0.07658,0.30890) (-0.13438,0.32551) (-0.19930,0.33849) (-0.27125,0.34696) (-0.35000,0.35000)};
\draw[-{Latex},structurecolor] (axis cs:-0.01740,-0.00620) -- (axis cs:-0.01766,-0.00813);
\draw[-{Latex},structurecolor] (axis cs:0.11647,-0.01110) -- (axis cs:0.12704,-0.00211);
\addplot[structurecolor,thin,no marks] coordinates {(-0.35000,0.35000) (-0.37628,0.34967) (-0.40318,0.34867) (-0.43069,0.34697) (-0.45879,0.34452) (-0.48744,0.34131) (-0.51664,0.33729) (-0.54634,0.33243) (-0.57652,0.32670) (-0.60715,0.32006) (-0.63818,0.31249) (-0.66960,0.30394) (-0.70134,0.29440) (-0.73338,0.28381) (-0.76566,0.27215) (-0.79814,0.25940) (-0.83076,0.24550) (-0.86348,0.23045) (-0.89623,0.21419) (-0.92896,0.19671) (-0.96160,0.17797) (-0.99409,0.15795) (-1.02636,0.13662) (-1.05833,0.11394) (-1.08994,0.08991) (-1.12109,0.06449) (-1.15172,0.03765) (-1.18174,0.00939) (-1.21106,-0.02033) (-1.23958,-0.05151) (-1.26723,-0.08417) (-1.29389,-0.11834) (-1.31947,-0.15401) (-1.34387,-0.19120) (-1.36698,-0.22992) (-1.38869,-0.27017) (-1.40889,-0.31195) (-1.42747,-0.35527) (-1.44431,-0.40011) (-1.45928,-0.44647) (-1.47227,-0.49434) (-1.48315,-0.54371) (-1.49180,-0.59456) (-1.49807,-0.64687) (-1.50186,-0.70062) (-1.50301,-0.75577) (-1.50139,-0.81230) (-1.49687,-0.87017) (-1.48930,-0.92934) (-1.47855,-0.98976) (-1.46448,-1.05138) (-1.44694,-1.11414) (-1.42580,-1.17800) (-1.40090,-1.24288) (-1.37210,-1.30871) (-1.33927,-1.37542) (-1.30226,-1.44292) (-1.26092,-1.51113) (-1.21511,-1.57996) (-1.16470,-1.64932) (-1.10954,-1.71908) (-1.04951,-1.78915) (-0.98445,-1.85941) (-0.91425,-1.92974) (-0.83877,-2.00000)};
\draw[-{Latex},structurecolor] (axis cs:-0.96160,0.17797) -- (axis cs:-0.99409,0.15795);
\draw[-{Latex},structurecolor] (axis cs:-1.49807,-0.64687) -- (axis cs:-1.50186,-0.70062);
\addplot[structurecolor,thin,no marks] coordinates {(-0.00403,-0.00048) (-0.00425,-0.00086) (-0.00442,-0.00128) (-0.00453,-0.00175) (-0.00458,-0.00226) (-0.00454,-0.00281) (-0.00440,-0.00341) (-0.00416,-0.00403) (-0.00379,-0.00468) (-0.00328,-0.00535) (-0.00262,-0.00603) (-0.00180,-0.00670) (-0.00081,-0.00735) (0.00037,-0.00797) (0.00174,-0.00853) (0.00331,-0.00902) (0.00508,-0.00942) (0.00705,-0.00969) (0.00921,-0.00982) (0.01155,-0.00978) (0.01406,-0.00954) (0.01672,-0.00906) (0.01950,-0.00833) (0.02236,-0.00731) (0.02526,-0.00597) (0.02815,-0.00429) (0.03098,-0.00223) (0.03367,0.00021) (0.03615,0.00307) (0.03833,0.00634) (0.04012,0.01005) (0.04142,0.01418) (0.04212,0.01873) (0.04211,0.02368) (0.04126,0.02900) (0.03945,0.03464) (0.03656,0.04056) (0.03245,0.04667) (0.02701,0.05289) (0.02011,0.05911) (0.01164,0.06522) (0.00151,0.07107) (-0.01036,0.07650) (-0.02403,0.08134) (-0.03954,0.08539) (-0.05689,0.08843) (-0.07605,0.09024) (-0.09694,0.09057) (-0.11945,0.08916) (-0.14340,0.08575) (-0.16856,0.08006) (-0.19465,0.07183) (-0.22129,0.06078) (-0.24806,0.04667) (-0.27445,0.02924) (-0.29987,0.00828) (-0.32366,-0.01638) (-0.34506,-0.04488) (-0.36326,-0.07732) (-0.37734,-0.11372) (-0.38635,-0.15402) (-0.38922,-0.19809) (-0.38489,-0.24569) (-0.37220,-0.29649) (-0.35000,-0.35000)};
\draw[-{Latex},structurecolor] (axis cs:0.01406,-0.00954) -- (axis cs:0.01672,-0.00906);
\draw[-{Latex},structurecolor] (axis cs:-0.02403,0.08134) -- (axis cs:-0.03954,0.08539);
\addplot[structurecolor,thin,no marks] coordinates {(-0.35000,-0.35000) (-0.34270,-0.36397) (-0.33469,-0.37807) (-0.32595,-0.39228) (-0.31647,-0.40659) (-0.30621,-0.42100) (-0.29518,-0.43548) (-0.28334,-0.45002) (-0.27067,-0.46460) (-0.25717,-0.47922) (-0.24281,-0.49385) (-0.22758,-0.50848) (-0.21145,-0.52308) (-0.19441,-0.53765) (-0.17645,-0.55215) (-0.15755,-0.56658) (-0.13769,-0.58091) (-0.11686,-0.59512) (-0.09505,-0.60919) (-0.07224,-0.62309) (-0.04841,-0.63680) (-0.02356,-0.65030) (0.00232,-0.66356) (0.02925,-0.67656) (0.05723,-0.68926) (0.08629,-0.70165) (0.11641,-0.71370) (0.14761,-0.72537) (0.17990,-0.73664) (0.21328,-0.74748) (0.24776,-0.75785) (0.28334,-0.76773) (0.32001,-0.77708) (0.35778,-0.78587) (0.39666,-0.79408) (0.43663,-0.80165) (0.47769,-0.80856) (0.51984,-0.81478) (0.56307,-0.82027) (0.60737,-0.82499) (0.65273,-0.82890) (0.69915,-0.83197) (0.74660,-0.83416) (0.79507,-0.83544) (0.84455,-0.83575) (0.89501,-0.83507) (0.94644,-0.83336) (0.99881,-0.83057) (1.05209,-0.82666) (1.10626,-0.82160) (1.16130,-0.81534) (1.21716,-0.80785) (1.27382,-0.79907) (1.33123,-0.78897) (1.38937,-0.77751) (1.44818,-0.76465) (1.50763,-0.75035) (1.56767,-0.73456) (1.62824,-0.71724) (1.68931,-0.69836) (1.75081,-0.67786) (1.81269,-0.65572) (1.87489,-0.63189) (1.93735,-0.60633) (2.00000,-0.57901)};
\draw[-{Latex},structurecolor] (axis cs:-0.04841,-0.63680) -- (axis cs:-0.02356,-0.65030);
\draw[-{Latex},structurecolor] (axis cs:0.79507,-0.83544) -- (axis cs:0.84455,-0.83575);
\addplot[structurecolor,thin,no marks] coordinates {(0.00333,-0.00236) (0.00397,-0.00226) (0.00465,-0.00209) (0.00535,-0.00186) (0.00606,-0.00155) (0.00678,-0.00115) (0.00748,-0.00066) (0.00815,-0.00008) (0.00877,0.00060) (0.00933,0.00138) (0.00979,0.00227) (0.01014,0.00327) (0.01034,0.00436) (0.01038,0.00556) (0.01022,0.00685) (0.00982,0.00822) (0.00917,0.00967) (0.00823,0.01116) (0.00696,0.01269) (0.00534,0.01422) (0.00334,0.01573) (0.00094,0.01719) (-0.00189,0.01855) (-0.00516,0.01978) (-0.00888,0.02082) (-0.01306,0.02163) (-0.01768,0.02214) (-0.02273,0.02230) (-0.02819,0.02205) (-0.03401,0.02132) (-0.04015,0.02005) (-0.04652,0.01816) (-0.05306,0.01560) (-0.05965,0.01229) (-0.06617,0.00818) (-0.07249,0.00321) (-0.07845,-0.00266) (-0.08386,-0.00948) (-0.08852,-0.01725) (-0.09222,-0.02600) (-0.09472,-0.03572) (-0.09577,-0.04637) (-0.09510,-0.05791) (-0.09244,-0.07025) (-0.08750,-0.08329) (-0.07999,-0.09688) (-0.06963,-0.11086) (-0.05613,-0.12500) (-0.03925,-0.13907) (-0.01873,-0.15276) (0.00563,-0.16575) (0.03399,-0.17765) (0.06648,-0.18805) (0.10315,-0.19648) (0.14398,-0.20246) (0.18887,-0.20544) (0.23761,-0.20485) (0.28989,-0.20012) (0.34527,-0.19062) (0.40319,-0.17575) (0.46293,-0.15488) (0.52360,-0.12742) (0.58418,-0.09279) (0.64345,-0.05047) (0.70000,0.00000)};
\draw[-{Latex},structurecolor] (axis cs:0.00334,0.01573) -- (axis cs:0.00094,0.01719);
\draw[-{Latex},structurecolor] (axis cs:-0.09244,-0.07025) -- (axis cs:-0.08750,-0.08329);
\addplot[structurecolor,thin,no marks] coordinates {(0.70000,0.00000) (0.71689,0.01753) (0.73331,0.03590) (0.74920,0.05512) (0.76452,0.07521) (0.77920,0.09617) (0.79318,0.11799) (0.80641,0.14070) (0.81882,0.16428) (0.83035,0.18875) (0.84093,0.21409) (0.85050,0.24031) (0.85898,0.26740) (0.86632,0.29536) (0.87243,0.32418) (0.87724,0.35384) (0.88069,0.38434) (0.88269,0.41566) (0.88318,0.44779) (0.88206,0.48071) (0.87927,0.51439) (0.87473,0.54881) (0.86835,0.58395) (0.86006,0.61978) (0.84978,0.65627) (0.83742,0.69338) (0.82290,0.73107) (0.80615,0.76931) (0.78707,0.80806) (0.76560,0.84726) (0.74165,0.88687) (0.71513,0.92683) (0.68598,0.96710) (0.65410,1.00761) (0.61943,1.04829) (0.58188,1.08909) (0.54138,1.12994) (0.49786,1.17075) (0.45125,1.21146) (0.40148,1.25199) (0.34847,1.29225) (0.29217,1.33215) (0.23252,1.37160) (0.16944,1.41052) (0.10290,1.44879) (0.03284,1.48632) (-0.04079,1.52300) (-0.11804,1.55873) (-0.19894,1.59338) (-0.28352,1.62684) (-0.37183,1.65899) (-0.46387,1.68971) (-0.55967,1.71886) (-0.65924,1.74632) (-0.76259,1.77194) (-0.86972,1.79560) (-0.98061,1.81715) (-1.09526,1.83644) (-1.21364,1.85332) (-1.33573,1.86765) (-1.46147,1.87927) (-1.59083,1.88802) (-1.72375,1.89374) (-1.86017,1.89627) (-2.00000,1.89545)};
\draw[-{Latex},structurecolor] (axis cs:0.87927,0.51439) -- (axis cs:0.87473,0.54881);
\draw[-{Latex},structurecolor] (axis cs:0.16944,1.41052) -- (axis cs:0.10290,1.44879);
\addplot[structurecolor,thin,no marks] coordinates {(-0.00333,0.00236) (-0.00397,0.00226) (-0.00465,0.00209) (-0.00535,0.00186) (-0.00606,0.00155) (-0.00678,0.00115) (-0.00748,0.00066) (-0.00815,0.00008) (-0.00877,-0.00060) (-0.00933,-0.00138) (-0.00979,-0.00227) (-0.01014,-0.00327) (-0.01034,-0.00436) (-0.01038,-0.00556) (-0.01022,-0.00685) (-0.00982,-0.00822) (-0.00917,-0.00967) (-0.00823,-0.01116) (-0.00696,-0.01269) (-0.00534,-0.01422) (-0.00334,-0.01573) (-0.00094,-0.01719) (0.00189,-0.01855) (0.00516,-0.01978) (0.00888,-0.02082) (0.01306,-0.02163) (0.01768,-0.02214) (0.02273,-0.02230) (0.02819,-0.02205) (0.03401,-0.02132) (0.04015,-0.02005) (0.04652,-0.01816) (0.05306,-0.01560) (0.05965,-0.01229) (0.06617,-0.00818) (0.07249,-0.00321) (0.07845,0.00266) (0.08386,0.00948) (0.08852,0.01725) (0.09222,0.02600) (0.09472,0.03572) (0.09577,0.04637) (0.09510,0.05791) (0.09244,0.07025) (0.08750,0.08329) (0.07999,0.09688) (0.06963,0.11086) (0.05613,0.12500) (0.03925,0.13907) (0.01873,0.15276) (-0.00563,0.16575) (-0.03399,0.17765) (-0.06648,0.18805) (-0.10315,0.19648) (-0.14398,0.20246) (-0.18887,0.20544) (-0.23761,0.20485) (-0.28989,0.20012) (-0.34527,0.19062) (-0.40319,0.17575) (-0.46293,0.15488) (-0.52360,0.12742) (-0.58418,0.09279) (-0.64345,0.05047) (-0.70000,0.00000)};
\draw[-{Latex},structurecolor] (axis cs:-0.00334,-0.01573) -- (axis cs:-0.00094,-0.01719);
\draw[-{Latex},structurecolor] (axis cs:0.09244,0.07025) -- (axis cs:0.08750,0.08329);
\addplot[structurecolor,thin,no marks] coordinates {(-0.70000,0.00000) (-0.71689,-0.01753) (-0.73331,-0.03590) (-0.74920,-0.05512) (-0.76452,-0.07521) (-0.77920,-0.09617) (-0.79318,-0.11799) (-0.80641,-0.14070) (-0.81882,-0.16428) (-0.83035,-0.18875) (-0.84093,-0.21409) (-0.85050,-0.24031) (-0.85898,-0.26740) (-0.86632,-0.29536) (-0.87243,-0.32418) (-0.87724,-0.35384) (-0.88069,-0.38434) (-0.88269,-0.41566) (-0.88318,-0.44779) (-0.88206,-0.48071) (-0.87927,-0.51439) (-0.87473,-0.54881) (-0.86835,-0.58395) (-0.86006,-0.61978) (-0.84978,-0.65627) (-0.83742,-0.69338) (-0.82290,-0.73107) (-0.80615,-0.76931) (-0.78707,-0.80806) (-0.76560,-0.84726) (-0.74165,-0.88687) (-0.71513,-0.92683) (-0.68598,-0.96710) (-0.65410,-1.00761) (-0.61943,-1.04829) (-0.58188,-1.08909) (-0.54138,-1.12994) (-0.49786,-1.17075) (-0.45125,-1.21146) (-0.40148,-1.25199) (-0.34847,-1.29225) (-0.29217,-1.33215) (-0.23252,-1.37160) (-0.16944,-1.41052) (-0.10290,-1.44879) (-0.03284,-1.48632) (0.04079,-1.52300) (0.11804,-1.55873) (0.19894,-1.59338) (0.28352,-1.62684) (0.37183,-1.65899) (0.46387,-1.68971) (0.55967,-1.71886) (0.65924,-1.74632) (0.76259,-1.77194) (0.86972,-1.79560) (0.98061,-1.81715) (1.09526,-1.83644) (1.21364,-1.85332) (1.33573,-1.86765) (1.46147,-1.87927) (1.59083,-1.88802) (1.72375,-1.89374) (1.86017,-1.89627) (2.00000,-1.89545)};
\draw[-{Latex},structurecolor] (axis cs:-0.87927,-0.51439) -- (axis cs:-0.87473,-0.54881);
\draw[-{Latex},structurecolor] (axis cs:-0.16944,-1.41052) -- (axis cs:-0.10290,-1.44879);
\addplot[structurecolor,thin,no marks] coordinates {(0.00473,0.00333) (0.00452,0.00397) (0.00419,0.00465) (0.00372,0.00535) (0.00309,0.00606) (0.00230,0.00678) (0.00133,0.00748) (0.00017,0.00815) (-0.00119,0.00877) (-0.00276,0.00933) (-0.00454,0.00979) (-0.00653,0.01014) (-0.00873,0.01034) (-0.01112,0.01038) (-0.01370,0.01022) (-0.01645,0.00982) (-0.01933,0.00917) (-0.02232,0.00823) (-0.02538,0.00696) (-0.02845,0.00534) (-0.03147,0.00334) (-0.03438,0.00094) (-0.03711,-0.00189) (-0.03956,-0.00516) (-0.04164,-0.00888) (-0.04325,-0.01306) (-0.04428,-0.01768) (-0.04461,-0.02273) (-0.04411,-0.02819) (-0.04265,-0.03401) (-0.04010,-0.04015) (-0.03632,-0.04652) (-0.03119,-0.05306) (-0.02457,-0.05965) (-0.01635,-0.06617) (-0.00641,-0.07249) (0.00533,-0.07845) (0.01895,-0.08386) (0.03451,-0.08852) (0.05201,-0.09222) (0.07144,-0.09472) (0.09274,-0.09577) (0.11582,-0.09510) (0.14050,-0.09244) (0.16657,-0.08750) (0.19376,-0.07999) (0.22172,-0.06963) (0.25001,-0.05613) (0.27814,-0.03925) (0.30553,-0.01873) (0.33150,0.00563) (0.35530,0.03399) (0.37610,0.06648) (0.39297,0.10315) (0.40492,0.14398) (0.41087,0.18887) (0.40971,0.23761) (0.40024,0.28989) (0.38125,0.34527) (0.35150,0.40319) (0.30976,0.46293) (0.25484,0.52360) (0.18559,0.58418) (0.10095,0.64345) (0.00000,0.70000)};
\draw[-{Latex},structurecolor] (axis cs:-0.03147,0.00334) -- (axis cs:-0.03438,0.00094);
\draw[-{Latex},structurecolor] (axis cs:0.14050,-0.09244) -- (axis cs:0.16657,-0.08750);
\addplot[structurecolor,thin,no marks] coordinates {(0.00000,0.70000) (-0.01776,0.70871) (-0.03596,0.71731) (-0.05461,0.72578) (-0.07371,0.73413) (-0.09325,0.74233) (-0.11325,0.75039) (-0.13369,0.75829) (-0.15460,0.76604) (-0.17596,0.77361) (-0.19778,0.78101) (-0.22006,0.78822) (-0.24279,0.79525) (-0.26599,0.80206) (-0.28966,0.80868) (-0.31378,0.81507) (-0.33837,0.82123) (-0.36341,0.82716) (-0.38893,0.83285) (-0.41490,0.83828) (-0.44133,0.84346) (-0.46823,0.84836) (-0.49558,0.85298) (-0.52339,0.85731) (-0.55166,0.86134) (-0.58038,0.86507) (-0.60956,0.86847) (-0.63918,0.87155) (-0.66925,0.87429) (-0.69976,0.87669) (-0.73072,0.87872) (-0.76211,0.88039) (-0.79394,0.88168) (-0.82619,0.88259) (-0.85887,0.88309) (-0.89197,0.88319) (-0.92548,0.88287) (-0.95940,0.88212) (-0.99372,0.88093) (-1.02845,0.87929) (-1.06356,0.87719) (-1.09906,0.87462) (-1.13493,0.87156) (-1.17117,0.86801) (-1.20777,0.86396) (-1.24473,0.85939) (-1.28203,0.85430) (-1.31967,0.84867) (-1.35763,0.84250) (-1.39590,0.83576) (-1.43449,0.82846) (-1.47336,0.82057) (-1.51252,0.81210) (-1.55196,0.80302) (-1.59165,0.79333) (-1.63160,0.78301) (-1.67177,0.77206) (-1.71217,0.76046) (-1.75278,0.74821) (-1.79359,0.73528) (-1.83457,0.72168) (-1.87572,0.70739) (-1.91702,0.69240) (-1.95845,0.67670) (-2.00000,0.66028)};
\draw[-{Latex},structurecolor] (axis cs:-0.44133,0.84346) -- (axis cs:-0.46823,0.84836);
\draw[-{Latex},structurecolor] (axis cs:-1.17117,0.86801) -- (axis cs:-1.20777,0.86396);
\addplot[structurecolor,thin,no marks] coordinates {(-0.00473,-0.00333) (-0.00452,-0.00397) (-0.00419,-0.00465) (-0.00372,-0.00535) (-0.00309,-0.00606) (-0.00230,-0.00678) (-0.00133,-0.00748) (-0.00017,-0.00815) (0.00119,-0.00877) (0.00276,-0.00933) (0.00454,-0.00979) (0.00653,-0.01014) (0.00873,-0.01034) (0.01112,-0.01038) (0.01370,-0.01022) (0.01645,-0.00982) (0.01933,-0.00917) (0.02232,-0.00823) (0.02538,-0.00696) (0.02845,-0.00534) (0.03147,-0.00334) (0.03438,-0.00094) (0.03711,0.00189) (0.03956,0.00516) (0.04164,0.00888) (0.04325,0.01306) (0.04428,0.01768) (0.04461,0.02273) (0.04411,0.02819) (0.04265,0.03401) (0.04010,0.04015) (0.03632,0.04652) (0.03119,0.05306) (0.02457,0.05965) (0.01635,0.06617) (0.00641,0.07249) (-0.00533,0.07845) (-0.01895,0.08386) (-0.03451,0.08852) (-0.05201,0.09222) (-0.07144,0.09472) (-0.09274,0.09577) (-0.11582,0.09510) (-0.14050,0.09244) (-0.16657,0.08750) (-0.19376,0.07999) (-0.22172,0.06963) (-0.25001,0.05613) (-0.27814,0.03925) (-0.30553,0.01873) (-0.33150,-0.00563) (-0.35530,-0.03399) (-0.37610,-0.06648) (-0.39297,-0.10315) (-0.40492,-0.14398) (-0.41087,-0.18887) (-0.40971,-0.23761) (-0.40024,-0.28989) (-0.38125,-0.34527) (-0.35150,-0.40319) (-0.30976,-0.46293) (-0.25484,-0.52360) (-0.18559,-0.58418) (-0.10095,-0.64345) (0.00000,-0.70000)};
\draw[-{Latex},structurecolor] (axis cs:0.03147,-0.00334) -- (axis cs:0.03438,-0.00094);
\draw[-{Latex},structurecolor] (axis cs:-0.14050,0.09244) -- (axis cs:-0.16657,0.08750);
\addplot[structurecolor,thin,no marks] coordinates {(0.00000,-0.70000) (0.01776,-0.70871) (0.03596,-0.71731) (0.05461,-0.72578) (0.07371,-0.73413) (0.09325,-0.74233) (0.11325,-0.75039) (0.13369,-0.75829) (0.15460,-0.76604) (0.17596,-0.77361) (0.19778,-0.78101) (0.22006,-0.78822) (0.24279,-0.79525) (0.26599,-0.80206) (0.28966,-0.80868) (0.31378,-0.81507) (0.33837,-0.82123) (0.36341,-0.82716) (0.38893,-0.83285) (0.41490,-0.83828) (0.44133,-0.84346) (0.46823,-0.84836) (0.49558,-0.85298) (0.52339,-0.85731) (0.55166,-0.86134) (0.58038,-0.86507) (0.60956,-0.86847) (0.63918,-0.87155) (0.66925,-0.87429) (0.69976,-0.87669) (0.73072,-0.87872) (0.76211,-0.88039) (0.79394,-0.88168) (0.82619,-0.88259) (0.85887,-0.88309) (0.89197,-0.88319) (0.92548,-0.88287) (0.95940,-0.88212) (0.99372,-0.88093) (1.02845,-0.87929) (1.06356,-0.87719) (1.09906,-0.87462) (1.13493,-0.87156) (1.17117,-0.86801) (1.20777,-0.86396) (1.24473,-0.85939) (1.28203,-0.85430) (1.31967,-0.84867) (1.35763,-0.84250) (1.39590,-0.83576) (1.43449,-0.82846) (1.47336,-0.82057) (1.51252,-0.81210) (1.55196,-0.80302) (1.59165,-0.79333) (1.63160,-0.78301) (1.67177,-0.77206) (1.71217,-0.76046) (1.75278,-0.74821) (1.79359,-0.73528) (1.83457,-0.72168) (1.87572,-0.70739) (1.91702,-0.69240) (1.95845,-0.67670) (2.00000,-0.66028)};
\draw[-{Latex},structurecolor] (axis cs:0.44133,-0.84346) -- (axis cs:0.46823,-0.84836);
\draw[-{Latex},structurecolor] (axis cs:1.17117,-0.86801) -- (axis cs:1.20777,-0.86396);
\addplot[only marks,mark=*,mark size=1.4pt] coordinates {(0,0)};
\end{axis}\end{tikzpicture}
\hfill
\begin{tikzpicture}\begin{axis}[width=.47\linewidth,height=5.2cm,axis lines=middle,xlabel={$x$},ylabel={$y$},xmin=-2,xmax=2,ymin=-2,ymax=2,title={（e）顺时针中心},tick label style={font=\footnotesize},clip=true]
\addplot[structurecolor,thin,no marks] coordinates {(0.35000,0.35000) (0.40356,0.26698) (0.45467,0.18234) (0.50300,0.09658) (0.54826,0.01024) (0.59017,-0.07617) (0.62849,-0.16211) (0.66297,-0.24707) (0.69341,-0.33052) (0.71961,-0.41195) (0.74143,-0.49086) (0.75872,-0.56679) (0.77139,-0.63925) (0.77935,-0.70782) (0.78255,-0.77207) (0.78098,-0.83161) (0.77465,-0.88607) (0.76359,-0.93513) (0.74788,-0.97849) (0.72760,-1.01587) (0.70288,-1.04706) (0.67387,-1.07186) (0.64076,-1.09012) (0.60373,-1.10174) (0.56302,-1.10663) (0.51888,-1.10477) (0.47157,-1.09617) (0.42139,-1.08088) (0.36863,-1.05900) (0.31363,-1.03066) (0.25671,-0.99603) (0.19822,-0.95533) (0.13853,-0.90880) (0.07799,-0.85672) (0.01698,-0.79942) (-0.04414,-0.73724) (-0.10499,-0.67056) (-0.16520,-0.59979) (-0.22440,-0.52537) (-0.28223,-0.44774) (-0.33834,-0.36737) (-0.39238,-0.28477) (-0.44403,-0.20043) (-0.49298,-0.11486) (-0.53891,-0.02860) (-0.58156,0.05784) (-0.62066,0.14393) (-0.65598,0.22913) (-0.68729,0.31294) (-0.71441,0.39484) (-0.73717,0.47434) (-0.75543,0.55094) (-0.76909,0.62417) (-0.77805,0.69360) (-0.78227,0.75880) (-0.78172,0.81937) (-0.77639,0.87494) (-0.76633,0.92518) (-0.75160,0.96977) (-0.73228,1.00844) (-0.70849,1.04096) (-0.68038,1.06714) (-0.64812,1.08680) (-0.61191,1.09983) (-0.57197,1.10615)};
\draw[-{Latex},structurecolor] (axis cs:0.70288,-1.04706) -- (axis cs:0.67387,-1.07186);
\draw[-{Latex},structurecolor] (axis cs:-0.49298,-0.11486) -- (axis cs:-0.53891,-0.02860);
\addplot[structurecolor,thin,no marks] coordinates {(0.35000,-0.35000) (0.34893,-0.37625) (0.34574,-0.40020) (0.34043,-0.42171) (0.33305,-0.44065) (0.32363,-0.45690) (0.31225,-0.47037) (0.29895,-0.48096) (0.28384,-0.48862) (0.26699,-0.49330) (0.24851,-0.49497) (0.22852,-0.49362) (0.20713,-0.48926) (0.18448,-0.48191) (0.16071,-0.47163) (0.13595,-0.45847) (0.11036,-0.44251) (0.08410,-0.42385) (0.05733,-0.40260) (0.03021,-0.37890) (0.00290,-0.35289) (-0.02442,-0.32473) (-0.05160,-0.29458) (-0.07846,-0.26264) (-0.10484,-0.22909) (-0.13058,-0.19415) (-0.15552,-0.15802) (-0.17952,-0.12093) (-0.20242,-0.08310) (-0.22409,-0.04477) (-0.24439,-0.00616) (-0.26320,0.03249) (-0.28040,0.07094) (-0.29589,0.10895) (-0.30958,0.14630) (-0.32138,0.18276) (-0.33122,0.21810) (-0.33904,0.25212) (-0.34478,0.28459) (-0.34843,0.31533) (-0.34995,0.34414) (-0.34934,0.37086) (-0.34659,0.39531) (-0.34173,0.41736) (-0.33479,0.43685) (-0.32580,0.45368) (-0.31483,0.46774) (-0.30193,0.47895) (-0.28720,0.48724) (-0.27071,0.49256) (-0.25257,0.49487) (-0.23289,0.49416) (-0.21178,0.49044) (-0.18939,0.48372) (-0.16584,0.47406) (-0.14128,0.46150) (-0.11586,0.44613) (-0.08973,0.42803) (-0.06305,0.40733) (-0.03599,0.38414) (-0.00871,0.35860) (0.01862,0.33088) (0.04584,0.30114) (0.07279,0.26956) (0.09928,0.23634)};
\draw[-{Latex},structurecolor] (axis cs:0.00290,-0.35289) -- (axis cs:-0.02442,-0.32473);
\draw[-{Latex},structurecolor] (axis cs:-0.34173,0.41736) -- (axis cs:-0.33479,0.43685);
\addplot[structurecolor,thin,no marks] coordinates {(-0.35000,0.35000) (-0.34893,0.37625) (-0.34574,0.40020) (-0.34043,0.42171) (-0.33305,0.44065) (-0.32363,0.45690) (-0.31225,0.47037) (-0.29895,0.48096) (-0.28384,0.48862) (-0.26699,0.49330) (-0.24851,0.49497) (-0.22852,0.49362) (-0.20713,0.48926) (-0.18448,0.48191) (-0.16071,0.47163) (-0.13595,0.45847) (-0.11036,0.44251) (-0.08410,0.42385) (-0.05733,0.40260) (-0.03021,0.37890) (-0.00290,0.35289) (0.02442,0.32473) (0.05160,0.29458) (0.07846,0.26264) (0.10484,0.22909) (0.13058,0.19415) (0.15552,0.15802) (0.17952,0.12093) (0.20242,0.08310) (0.22409,0.04477) (0.24439,0.00616) (0.26320,-0.03249) (0.28040,-0.07094) (0.29589,-0.10895) (0.30958,-0.14630) (0.32138,-0.18276) (0.33122,-0.21810) (0.33904,-0.25212) (0.34478,-0.28459) (0.34843,-0.31533) (0.34995,-0.34414) (0.34934,-0.37086) (0.34659,-0.39531) (0.34173,-0.41736) (0.33479,-0.43685) (0.32580,-0.45368) (0.31483,-0.46774) (0.30193,-0.47895) (0.28720,-0.48724) (0.27071,-0.49256) (0.25257,-0.49487) (0.23289,-0.49416) (0.21178,-0.49044) (0.18939,-0.48372) (0.16584,-0.47406) (0.14128,-0.46150) (0.11586,-0.44613) (0.08973,-0.42803) (0.06305,-0.40733) (0.03599,-0.38414) (0.00871,-0.35860) (-0.01862,-0.33088) (-0.04584,-0.30114) (-0.07279,-0.26956) (-0.09928,-0.23634)};
\draw[-{Latex},structurecolor] (axis cs:-0.00290,0.35289) -- (axis cs:0.02442,0.32473);
\draw[-{Latex},structurecolor] (axis cs:0.34173,-0.41736) -- (axis cs:0.33479,-0.43685);
\addplot[structurecolor,thin,no marks] coordinates {(-0.35000,-0.35000) (-0.40356,-0.26698) (-0.45467,-0.18234) (-0.50300,-0.09658) (-0.54826,-0.01024) (-0.59017,0.07617) (-0.62849,0.16211) (-0.66297,0.24707) (-0.69341,0.33052) (-0.71961,0.41195) (-0.74143,0.49086) (-0.75872,0.56679) (-0.77139,0.63925) (-0.77935,0.70782) (-0.78255,0.77207) (-0.78098,0.83161) (-0.77465,0.88607) (-0.76359,0.93513) (-0.74788,0.97849) (-0.72760,1.01587) (-0.70288,1.04706) (-0.67387,1.07186) (-0.64076,1.09012) (-0.60373,1.10174) (-0.56302,1.10663) (-0.51888,1.10477) (-0.47157,1.09617) (-0.42139,1.08088) (-0.36863,1.05900) (-0.31363,1.03066) (-0.25671,0.99603) (-0.19822,0.95533) (-0.13853,0.90880) (-0.07799,0.85672) (-0.01698,0.79942) (0.04414,0.73724) (0.10499,0.67056) (0.16520,0.59979) (0.22440,0.52537) (0.28223,0.44774) (0.33834,0.36737) (0.39238,0.28477) (0.44403,0.20043) (0.49298,0.11486) (0.53891,0.02860) (0.58156,-0.05784) (0.62066,-0.14393) (0.65598,-0.22913) (0.68729,-0.31294) (0.71441,-0.39484) (0.73717,-0.47434) (0.75543,-0.55094) (0.76909,-0.62417) (0.77805,-0.69360) (0.78227,-0.75880) (0.78172,-0.81937) (0.77639,-0.87494) (0.76633,-0.92518) (0.75160,-0.96977) (0.73228,-1.00844) (0.70849,-1.04096) (0.68038,-1.06714) (0.64812,-1.08680) (0.61191,-1.09983) (0.57197,-1.10615)};
\draw[-{Latex},structurecolor] (axis cs:-0.70288,1.04706) -- (axis cs:-0.67387,1.07186);
\draw[-{Latex},structurecolor] (axis cs:0.49298,0.11486) -- (axis cs:0.53891,0.02860);
\addplot[structurecolor,thin,no marks] coordinates {(0.70000,0.00000) (0.75250,-0.10926) (0.80040,-0.21786) (0.84343,-0.32513) (0.88130,-0.43041) (0.91381,-0.53307) (0.94073,-0.63248) (0.96192,-0.72803) (0.97724,-0.81914) (0.98660,-0.90525) (0.98994,-0.98583) (0.98724,-1.06041) (0.97852,-1.12851) (0.96383,-1.18973) (0.94326,-1.24370) (0.91693,-1.29007) (0.88501,-1.32858) (0.84770,-1.35898) (0.80521,-1.38109) (0.75781,-1.39478) (0.70578,-1.39995) (0.64945,-1.39659) (0.58916,-1.38470) (0.52528,-1.36437) (0.45819,-1.33572) (0.38830,-1.29892) (0.31605,-1.25419) (0.24187,-1.20181) (0.16621,-1.14211) (0.08954,-1.07543) (0.01232,-1.00219) (-0.06497,-0.92284) (-0.14187,-0.83786) (-0.21790,-0.74777) (-0.29260,-0.65312) (-0.36552,-0.55448) (-0.43621,-0.45246) (-0.50423,-0.34768) (-0.56918,-0.24078) (-0.63066,-0.13241) (-0.68829,-0.02323) (-0.74172,0.08609) (-0.79063,0.19489) (-0.83471,0.30249) (-0.87370,0.40825) (-0.90736,0.51152) (-0.93549,0.61167) (-0.95791,0.70809) (-0.97448,0.80019) (-0.98512,0.88740) (-0.98974,0.96921) (-0.98832,1.04510) (-0.98087,1.11461) (-0.96744,1.17733) (-0.94811,1.23286) (-0.92300,1.28087) (-0.89225,1.32107) (-0.85606,1.35321) (-0.81465,1.37709) (-0.76827,1.39258) (-0.71720,1.39957) (-0.66176,1.39802) (-0.60228,1.38794) (-0.53913,1.36939) (-0.47268,1.34249)};
\draw[-{Latex},structurecolor] (axis cs:0.70578,-1.39995) -- (axis cs:0.64945,-1.39659);
\draw[-{Latex},structurecolor] (axis cs:-0.83471,0.30249) -- (axis cs:-0.87370,0.40825);
\addplot[structurecolor,thin,no marks] coordinates {(-0.70000,0.00000) (-0.75250,0.10926) (-0.80040,0.21786) (-0.84343,0.32513) (-0.88130,0.43041) (-0.91381,0.53307) (-0.94073,0.63248) (-0.96192,0.72803) (-0.97724,0.81914) (-0.98660,0.90525) (-0.98994,0.98583) (-0.98724,1.06041) (-0.97852,1.12851) (-0.96383,1.18973) (-0.94326,1.24370) (-0.91693,1.29007) (-0.88501,1.32858) (-0.84770,1.35898) (-0.80521,1.38109) (-0.75781,1.39478) (-0.70578,1.39995) (-0.64945,1.39659) (-0.58916,1.38470) (-0.52528,1.36437) (-0.45819,1.33572) (-0.38830,1.29892) (-0.31605,1.25419) (-0.24187,1.20181) (-0.16621,1.14211) (-0.08954,1.07543) (-0.01232,1.00219) (0.06497,0.92284) (0.14187,0.83786) (0.21790,0.74777) (0.29260,0.65312) (0.36552,0.55448) (0.43621,0.45246) (0.50423,0.34768) (0.56918,0.24078) (0.63066,0.13241) (0.68829,0.02323) (0.74172,-0.08609) (0.79063,-0.19489) (0.83471,-0.30249) (0.87370,-0.40825) (0.90736,-0.51152) (0.93549,-0.61167) (0.95791,-0.70809) (0.97448,-0.80019) (0.98512,-0.88740) (0.98974,-0.96921) (0.98832,-1.04510) (0.98087,-1.11461) (0.96744,-1.17733) (0.94811,-1.23286) (0.92300,-1.28087) (0.89225,-1.32107) (0.85606,-1.35321) (0.81465,-1.37709) (0.76827,-1.39258) (0.71720,-1.39957) (0.66176,-1.39802) (0.60228,-1.38794) (0.53913,-1.36939) (0.47268,-1.34249)};
\draw[-{Latex},structurecolor] (axis cs:-0.70578,1.39995) -- (axis cs:-0.64945,1.39659);
\draw[-{Latex},structurecolor] (axis cs:0.83471,-0.30249) -- (axis cs:0.87370,-0.40825);
\addplot[structurecolor,thin,no marks] coordinates {(0.00000,0.70000) (0.05463,0.64323) (0.10893,0.58254) (0.16256,0.51830) (0.21521,0.45089) (0.26654,0.38073) (0.31624,0.30825) (0.36401,0.23389) (0.40957,0.15811) (0.45262,0.08135) (0.49292,0.00411) (0.53020,-0.07317) (0.56426,-0.14999) (0.59487,-0.22590) (0.62185,-0.30044) (0.64504,-0.37314) (0.66429,-0.44356) (0.67949,-0.51128) (0.69054,-0.57588) (0.69739,-0.63697) (0.69998,-0.69417) (0.69829,-0.74713) (0.69235,-0.79554) (0.68219,-0.83910) (0.66786,-0.87753) (0.64946,-0.91062) (0.62710,-0.93814) (0.60091,-0.95995) (0.57105,-0.97590) (0.53771,-0.98589) (0.50110,-0.98987) (0.46142,-0.98781) (0.41893,-0.97973) (0.37388,-0.96567) (0.32656,-0.94572) (0.27724,-0.92000) (0.22623,-0.88866) (0.17384,-0.85191) (0.12039,-0.80996) (0.06620,-0.76307) (0.01161,-0.71152) (-0.04305,-0.65563) (-0.09744,-0.59574) (-0.15125,-0.53222) (-0.20413,-0.46545) (-0.25576,-0.39584) (-0.30583,-0.32382) (-0.35404,-0.24982) (-0.40009,-0.17430) (-0.44370,-0.09771) (-0.48460,-0.02053) (-0.52255,0.05678) (-0.55731,0.13374) (-0.58866,0.20988) (-0.61643,0.28475) (-0.64044,0.35787) (-0.66054,0.42882) (-0.67661,0.49715) (-0.68855,0.56244) (-0.69629,0.62431) (-0.69978,0.68236) (-0.69901,0.73626) (-0.69397,0.78566) (-0.68470,0.83027) (-0.67125,0.86981)};
\draw[-{Latex},structurecolor] (axis cs:0.69998,-0.69417) -- (axis cs:0.69829,-0.74713);
\draw[-{Latex},structurecolor] (axis cs:-0.15125,-0.53222) -- (axis cs:-0.20413,-0.46545);
\addplot[structurecolor,thin,no marks] coordinates {(0.00000,-0.70000) (-0.05463,-0.64323) (-0.10893,-0.58254) (-0.16256,-0.51830) (-0.21521,-0.45089) (-0.26654,-0.38073) (-0.31624,-0.30825) (-0.36401,-0.23389) (-0.40957,-0.15811) (-0.45262,-0.08135) (-0.49292,-0.00411) (-0.53020,0.07317) (-0.56426,0.14999) (-0.59487,0.22590) (-0.62185,0.30044) (-0.64504,0.37314) (-0.66429,0.44356) (-0.67949,0.51128) (-0.69054,0.57588) (-0.69739,0.63697) (-0.69998,0.69417) (-0.69829,0.74713) (-0.69235,0.79554) (-0.68219,0.83910) (-0.66786,0.87753) (-0.64946,0.91062) (-0.62710,0.93814) (-0.60091,0.95995) (-0.57105,0.97590) (-0.53771,0.98589) (-0.50110,0.98987) (-0.46142,0.98781) (-0.41893,0.97973) (-0.37388,0.96567) (-0.32656,0.94572) (-0.27724,0.92000) (-0.22623,0.88866) (-0.17384,0.85191) (-0.12039,0.80996) (-0.06620,0.76307) (-0.01161,0.71152) (0.04305,0.65563) (0.09744,0.59574) (0.15125,0.53222) (0.20413,0.46545) (0.25576,0.39584) (0.30583,0.32382) (0.35404,0.24982) (0.40009,0.17430) (0.44370,0.09771) (0.48460,0.02053) (0.52255,-0.05678) (0.55731,-0.13374) (0.58866,-0.20988) (0.61643,-0.28475) (0.64044,-0.35787) (0.66054,-0.42882) (0.67661,-0.49715) (0.68855,-0.56244) (0.69629,-0.62431) (0.69978,-0.68236) (0.69901,-0.73626) (0.69397,-0.78566) (0.68470,-0.83027) (0.67125,-0.86981)};
\draw[-{Latex},structurecolor] (axis cs:-0.69998,0.69417) -- (axis cs:-0.69829,0.74713);
\draw[-{Latex},structurecolor] (axis cs:0.15125,0.53222) -- (axis cs:0.20413,0.46545);
\addplot[only marks,mark=*,mark size=1.4pt] coordinates {(0,0)};
\end{axis}\end{tikzpicture}
\end{center}
\end{solution}
\begin{exercise}\label{cw:bank-5B-5}原题5B-5. 弹簧阻尼振子满足 $mx''+cx'+kx=0$，题干设 $m,c,k>0$. (a) 写等价一阶系统；(b) 判平衡点类型与稳定性并用物理直觉检查：(i) $c=0$；(ii) $c\approx0$, $m,k\gg1$；(iii) 能否为鞍点？题干与(i)的 $c=0$ 相冲突，按(i)作为无阻尼边界情形处理.\end{exercise}
\begin{solution}（官方译审）(a) 令 $y=x'$，得 $x'=y$, $y'=-(k/m)x-(c/m)y$. (b) 特征根为 $(-c\pm\sqrt{c^2-4mk})/(2m)$. (i) $c=0$ 时为 $\pm i\sqrt{k/m}$，原点为中心，稳定而非渐近稳定；能量 $E=(my^2+kx^2)/2$ 守恒，轨道为椭圆. (ii) 小正阻尼满足 $c^2<4mk$ 时为渐近稳定螺旋点，实部 $-c/(2m)<0$；能量导数 $E'=-cy^2\le0$，振幅缓慢衰减，与结论相容. (iii) 不会为鞍点，因为矩阵行列式 $k/m>0$，不可能有异号实根. 若阻尼较大，则两个负实根形成稳定结点，临界阻尼为重负根，仍渐近稳定.\end{solution}
\originalsection{5C：非线性系统相图}
\begin{exercise}\label{cw:bank-5C-1}原题5C-1. 对 $x'=x-y+xy$, $y'=3x-2y-xy$，判原点几何类型和稳定性，并画附近若干轨道.\end{exercise}
\begin{solution}（官方译审）Jacobian为 $J_0=\begin{pmatrix}1&-1\\3&-2\end{pmatrix}$，特征多项式 $\lambda^2+\lambda+1$，根 $(-1\pm i\sqrt3)/2$. 原点为渐近稳定螺旋点；在正 $x$ 轴附近 $y'>0$，逆时针向内. 非线性余项二阶，实部严格负，线性化稳定性定理适用. 下图用原方程积分作局部相图；类型与稳定性由上述解析判据支持.
\begin{center}
\begin{tikzpicture}\begin{axis}[width=.47\linewidth,height=5.2cm,axis lines=middle,xlabel={$x$},ylabel={$y$},xmin=-0.4,xmax=0.4,ymin=-0.5,ymax=0.5,title={原点附近：向内逆时针},tick label style={font=\footnotesize},clip=true]
\addplot[structurecolor,thin,no marks] coordinates {(0.15000,0.00000) (0.17439,0.07563) (0.19257,0.13537) (0.20606,0.18184) (0.21603,0.21756) (0.22336,0.24477) (0.22870,0.26533) (0.23252,0.28075) (0.23519,0.29222) (0.23698,0.30067) (0.23809,0.30681) (0.23864,0.31117) (0.23876,0.31417) (0.23852,0.31611) (0.23795,0.31721) (0.23711,0.31764) (0.23600,0.31752) (0.23464,0.31693) (0.23303,0.31593) (0.23116,0.31455) (0.22903,0.31281) (0.22660,0.31073) (0.22388,0.30829) (0.22082,0.30549) (0.21742,0.30230) (0.21363,0.29870) (0.20944,0.29467) (0.20480,0.29016) (0.19970,0.28515) (0.19410,0.27960) (0.18799,0.27346) (0.18132,0.26670) (0.17410,0.25928) (0.16631,0.25116) (0.15795,0.24233) (0.14902,0.23275) (0.13957,0.22241) (0.12961,0.21131) (0.11922,0.19948) (0.10846,0.18693) (0.09743,0.17372) (0.08624,0.15993) (0.07502,0.14566) (0.06390,0.13103) (0.05304,0.11619) (0.04258,0.10130) (0.03268,0.08656) (0.02348,0.07215) (0.01510,0.05829) (0.00764,0.04516) (0.00118,0.03295) (-0.00425,0.02181) (-0.00863,0.01188) (-0.01199,0.00325) (-0.01438,-0.00402) (-0.01588,-0.00993) (-0.01659,-0.01451) (-0.01661,-0.01783) (-0.01606,-0.01998) (-0.01506,-0.02108) (-0.01372,-0.02127) (-0.01215,-0.02072) (-0.01044,-0.01955) (-0.00868,-0.01792) (-0.00694,-0.01598)};
\draw[-{Latex},structurecolor] (axis cs:0.22903,0.31281) -- (axis cs:0.22660,0.31073);
\draw[-{Latex},structurecolor] (axis cs:0.06390,0.13103) -- (axis cs:0.05304,0.11619);
\addplot[structurecolor,thin,no marks] coordinates {(0.00000,0.15000) (-0.02582,0.09673) (-0.04679,0.04935) (-0.06270,0.00804) (-0.07354,-0.02691) (-0.07956,-0.05518) (-0.08123,-0.07667) (-0.07917,-0.09148) (-0.07415,-0.10002) (-0.06695,-0.10292) (-0.05832,-0.10103) (-0.04893,-0.09529) (-0.03936,-0.08668) (-0.03003,-0.07613) (-0.02130,-0.06449) (-0.01338,-0.05246) (-0.00642,-0.04063) (-0.00047,-0.02943) (0.00445,-0.01917) (0.00838,-0.01006) (0.01136,-0.00220) (0.01348,0.00437) (0.01483,0.00968) (0.01550,0.01377) (0.01559,0.01674) (0.01519,0.01869) (0.01441,0.01975) (0.01333,0.02003) (0.01203,0.01966) (0.01058,0.01876) (0.00905,0.01745) (0.00751,0.01583) (0.00599,0.01401) (0.00454,0.01206) (0.00320,0.01007) (0.00198,0.00811) (0.00091,0.00623) (-0.00001,0.00447) (-0.00077,0.00288) (-0.00138,0.00146) (-0.00184,0.00024) (-0.00216,-0.00079) (-0.00236,-0.00161) (-0.00244,-0.00224) (-0.00244,-0.00269) (-0.00235,-0.00297) (-0.00221,-0.00311) (-0.00202,-0.00313) (-0.00179,-0.00304) (-0.00155,-0.00287) (-0.00130,-0.00264) (-0.00105,-0.00236) (-0.00081,-0.00205) (-0.00058,-0.00173) (-0.00038,-0.00141) (-0.00019,-0.00110) (-0.00004,-0.00081) (0.00010,-0.00054) (0.00020,-0.00030) (0.00028,-0.00009) (0.00034,0.00009) (0.00038,0.00023) (0.00040,0.00034) (0.00040,0.00042) (0.00039,0.00048)};
\draw[-{Latex},structurecolor] (axis cs:0.01136,-0.00220) -- (axis cs:0.01348,0.00437);
\draw[-{Latex},structurecolor] (axis cs:-0.00244,-0.00224) -- (axis cs:-0.00244,-0.00269);
\addplot[structurecolor,thin,no marks] coordinates {(-0.15000,0.00000) (-0.17164,-0.07735) (-0.18063,-0.13950) (-0.17820,-0.18506) (-0.16650,-0.21361) (-0.14809,-0.22597) (-0.12558,-0.22410) (-0.10129,-0.21075) (-0.07706,-0.18906) (-0.05418,-0.16209) (-0.03350,-0.13256) (-0.01547,-0.10266) (-0.00023,-0.07401) (0.01224,-0.04770) (0.02210,-0.02439) (0.02956,-0.00440) (0.03487,0.01222) (0.03830,0.02556) (0.04010,0.03585) (0.04051,0.04334) (0.03976,0.04836) (0.03808,0.05119) (0.03564,0.05217) (0.03263,0.05156) (0.02922,0.04966) (0.02556,0.04673) (0.02178,0.04300) (0.01799,0.03870) (0.01430,0.03402) (0.01080,0.02915) (0.00756,0.02424) (0.00463,0.01945) (0.00205,0.01487) (-0.00016,0.01062) (-0.00199,0.00675) (-0.00345,0.00333) (-0.00454,0.00039) (-0.00531,-0.00207) (-0.00576,-0.00404) (-0.00596,-0.00554) (-0.00592,-0.00661) (-0.00570,-0.00728) (-0.00533,-0.00759) (-0.00485,-0.00761) (-0.00429,-0.00737) (-0.00369,-0.00693) (-0.00307,-0.00634) (-0.00246,-0.00565) (-0.00187,-0.00489) (-0.00133,-0.00411) (-0.00083,-0.00332) (-0.00040,-0.00256) (-0.00002,-0.00185) (0.00029,-0.00120) (0.00054,-0.00063) (0.00073,-0.00013) (0.00086,0.00029) (0.00094,0.00062) (0.00098,0.00088) (0.00098,0.00107) (0.00095,0.00119) (0.00090,0.00125) (0.00082,0.00126) (0.00073,0.00123) (0.00064,0.00117)};
\draw[-{Latex},structurecolor] (axis cs:0.03976,0.04836) -- (axis cs:0.03808,0.05119);
\draw[-{Latex},structurecolor] (axis cs:-0.00485,-0.00761) -- (axis cs:-0.00429,-0.00737);
\addplot[structurecolor,thin,no marks] coordinates {(0.00000,-0.15000) (0.02518,-0.09635) (0.04488,-0.04891) (0.05976,-0.00847) (0.07049,0.02493) (0.07770,0.05165) (0.08193,0.07232) (0.08368,0.08761) (0.08334,0.09822) (0.08127,0.10483) (0.07779,0.10803) (0.07317,0.10834) (0.06765,0.10626) (0.06146,0.10219) (0.05480,0.09652) (0.04786,0.08958) (0.04082,0.08167) (0.03384,0.07309) (0.02708,0.06409) (0.02066,0.05491) (0.01469,0.04579) (0.00927,0.03691) (0.00448,0.02846) (0.00034,0.02060) (-0.00312,0.01343) (-0.00589,0.00706) (-0.00800,0.00154) (-0.00948,-0.00308) (-0.01040,-0.00683) (-0.01082,-0.00971) (-0.01081,-0.01178) (-0.01043,-0.01311) (-0.00978,-0.01377) (-0.00891,-0.01386) (-0.00790,-0.01348) (-0.00680,-0.01271) (-0.00566,-0.01165) (-0.00454,-0.01039) (-0.00346,-0.00901) (-0.00246,-0.00756) (-0.00154,-0.00612) (-0.00074,-0.00473) (-0.00005,-0.00342) (0.00052,-0.00223) (0.00098,-0.00116) (0.00133,-0.00025) (0.00158,0.00052) (0.00173,0.00114) (0.00181,0.00162) (0.00181,0.00196) (0.00175,0.00218) (0.00165,0.00230) (0.00152,0.00232) (0.00136,0.00227) (0.00118,0.00215) (0.00099,0.00199) (0.00081,0.00178) (0.00063,0.00156) (0.00046,0.00133) (0.00031,0.00109) (0.00017,0.00086) (0.00005,0.00064) (-0.00005,0.00044) (-0.00013,0.00026) (-0.00020,0.00010)};
\draw[-{Latex},structurecolor] (axis cs:0.01469,0.04579) -- (axis cs:0.00927,0.03691);
\draw[-{Latex},structurecolor] (axis cs:0.00052,-0.00223) -- (axis cs:0.00098,-0.00116);
\addplot[only marks,mark=*,mark size=1.4pt] coordinates {(0,0)};
\end{axis}\end{tikzpicture}
\end{center}
\end{solution}
\begin{exercise}\label{cw:bank-5C-2}原题5C-2. 对 $x'=x+2x^2-y^2$, $y'=x-2y+x^3$ 重做5C-1：判原点类型与稳定性，画附近轨道.\end{exercise}
\begin{solution}（官方译审）$J_0=\begin{pmatrix}1&0\\1&-2\end{pmatrix}$，根 $1,-2$，向量 $(3,1)^T,(0,1)^T$. 原点为不稳定鞍点，局部不稳定、稳定曲线分别与这两方向相切，箭头沿前者离开、后者进入. 鞍点邻近轨道在两条分离曲线之间转弯；非线性稳定流形弯曲，不能把整个直线当作原系统不变轨道.
\begin{center}
\begin{tikzpicture}\begin{axis}[width=.47\linewidth,height=5.2cm,axis lines=middle,xlabel={$x$},ylabel={$y$},xmin=-0.3,xmax=0.3,ymin=-0.3,ymax=0.3,title={原点附近：鞍点},tick label style={font=\footnotesize},clip=true]
\addplot[structurecolor,thin,no marks] coordinates {(-0.03069,-0.30000) (-0.03170,-0.29514) (-0.03270,-0.29037) (-0.03368,-0.28568) (-0.03465,-0.28109) (-0.03560,-0.27659) (-0.03654,-0.27216) (-0.03747,-0.26783) (-0.03838,-0.26357) (-0.03927,-0.25940) (-0.04016,-0.25530) (-0.04104,-0.25128) (-0.04190,-0.24734) (-0.04275,-0.24347) (-0.04360,-0.23968) (-0.04443,-0.23596) (-0.04526,-0.23231) (-0.04607,-0.22873) (-0.04688,-0.22521) (-0.04768,-0.22177) (-0.04847,-0.21839) (-0.04926,-0.21507) (-0.05004,-0.21182) (-0.05081,-0.20863) (-0.05157,-0.20550) (-0.05233,-0.20243) (-0.05309,-0.19942) (-0.05384,-0.19647) (-0.05458,-0.19358) (-0.05532,-0.19074) (-0.05606,-0.18795) (-0.05679,-0.18522) (-0.05752,-0.18255) (-0.05824,-0.17992) (-0.05897,-0.17735) (-0.05968,-0.17482) (-0.06040,-0.17235) (-0.06111,-0.16992) (-0.06182,-0.16754) (-0.06253,-0.16521) (-0.06324,-0.16293) (-0.06394,-0.16069) (-0.06465,-0.15849) (-0.06535,-0.15634) (-0.06605,-0.15423) (-0.06675,-0.15216) (-0.06744,-0.15014) (-0.06814,-0.14815) (-0.06884,-0.14621) (-0.06953,-0.14430) (-0.07023,-0.14243) (-0.07092,-0.14060) (-0.07162,-0.13881) (-0.07232,-0.13706) (-0.07301,-0.13534) (-0.07371,-0.13365) (-0.07440,-0.13201) (-0.07510,-0.13039) (-0.07580,-0.12881) (-0.07650,-0.12726) (-0.07719,-0.12575) (-0.07789,-0.12426) (-0.07859,-0.12281) (-0.07930,-0.12139) (-0.08000,-0.12000)};
\draw[-{Latex},structurecolor] (axis cs:-0.04847,-0.21839) -- (axis cs:-0.04926,-0.21507);
\draw[-{Latex},structurecolor] (axis cs:-0.06535,-0.15634) -- (axis cs:-0.06605,-0.15423);
\addplot[structurecolor,thin,no marks] coordinates {(-0.08000,-0.12000) (-0.08244,-0.11543) (-0.08489,-0.11120) (-0.08736,-0.10728) (-0.08986,-0.10366) (-0.09239,-0.10033) (-0.09494,-0.09727) (-0.09753,-0.09445) (-0.10015,-0.09188) (-0.10280,-0.08953) (-0.10548,-0.08740) (-0.10820,-0.08548) (-0.11096,-0.08374) (-0.11376,-0.08219) (-0.11659,-0.08082) (-0.11946,-0.07960) (-0.12237,-0.07855) (-0.12533,-0.07764) (-0.12832,-0.07688) (-0.13135,-0.07625) (-0.13443,-0.07576) (-0.13754,-0.07538) (-0.14070,-0.07512) (-0.14389,-0.07498) (-0.14713,-0.07494) (-0.15040,-0.07500) (-0.15372,-0.07516) (-0.15707,-0.07542) (-0.16047,-0.07576) (-0.16390,-0.07619) (-0.16736,-0.07671) (-0.17087,-0.07730) (-0.17441,-0.07797) (-0.17799,-0.07872) (-0.18160,-0.07953) (-0.18524,-0.08041) (-0.18892,-0.08136) (-0.19263,-0.08237) (-0.19637,-0.08344) (-0.20014,-0.08457) (-0.20393,-0.08576) (-0.20776,-0.08701) (-0.21161,-0.08830) (-0.21548,-0.08965) (-0.21938,-0.09105) (-0.22330,-0.09250) (-0.22724,-0.09399) (-0.23120,-0.09553) (-0.23518,-0.09711) (-0.23917,-0.09874) (-0.24318,-0.10041) (-0.24720,-0.10211) (-0.25123,-0.10386) (-0.25528,-0.10565) (-0.25933,-0.10747) (-0.26339,-0.10932) (-0.26745,-0.11122) (-0.27152,-0.11314) (-0.27560,-0.11510) (-0.27967,-0.11709) (-0.28374,-0.11910) (-0.28781,-0.12115) (-0.29188,-0.12323) (-0.29594,-0.12533) (-0.30000,-0.12746)};
\draw[-{Latex},structurecolor] (axis cs:-0.13443,-0.07576) -- (axis cs:-0.13754,-0.07538);
\draw[-{Latex},structurecolor] (axis cs:-0.21548,-0.08965) -- (axis cs:-0.21938,-0.09105);
\addplot[structurecolor,thin,no marks] coordinates {(0.00008,-0.30000) (-0.00205,-0.28565) (-0.00404,-0.27203) (-0.00589,-0.25911) (-0.00762,-0.24685) (-0.00925,-0.23522) (-0.01079,-0.22419) (-0.01224,-0.21372) (-0.01362,-0.20379) (-0.01494,-0.19436) (-0.01619,-0.18542) (-0.01740,-0.17693) (-0.01856,-0.16889) (-0.01968,-0.16125) (-0.02076,-0.15401) (-0.02182,-0.14713) (-0.02285,-0.14062) (-0.02387,-0.13444) (-0.02486,-0.12858) (-0.02584,-0.12302) (-0.02681,-0.11776) (-0.02777,-0.11277) (-0.02872,-0.10804) (-0.02967,-0.10356) (-0.03062,-0.09932) (-0.03157,-0.09531) (-0.03252,-0.09151) (-0.03347,-0.08792) (-0.03443,-0.08452) (-0.03540,-0.08130) (-0.03637,-0.07827) (-0.03736,-0.07540) (-0.03835,-0.07270) (-0.03936,-0.07014) (-0.04037,-0.06774) (-0.04141,-0.06547) (-0.04245,-0.06334) (-0.04351,-0.06134) (-0.04459,-0.05946) (-0.04568,-0.05769) (-0.04679,-0.05604) (-0.04792,-0.05449) (-0.04906,-0.05304) (-0.05023,-0.05169) (-0.05141,-0.05043) (-0.05262,-0.04927) (-0.05385,-0.04819) (-0.05509,-0.04719) (-0.05636,-0.04626) (-0.05765,-0.04542) (-0.05897,-0.04464) (-0.06031,-0.04394) (-0.06167,-0.04330) (-0.06306,-0.04272) (-0.06447,-0.04221) (-0.06590,-0.04176) (-0.06736,-0.04136) (-0.06885,-0.04102) (-0.07036,-0.04073) (-0.07190,-0.04049) (-0.07346,-0.04030) (-0.07506,-0.04016) (-0.07668,-0.04006) (-0.07832,-0.04001) (-0.08000,-0.04000)};
\draw[-{Latex},structurecolor] (axis cs:-0.02681,-0.11776) -- (axis cs:-0.02777,-0.11277);
\draw[-{Latex},structurecolor] (axis cs:-0.05023,-0.05169) -- (axis cs:-0.05141,-0.05043);
\addplot[structurecolor,thin,no marks] coordinates {(-0.08000,-0.04000) (-0.08213,-0.04005) (-0.08431,-0.04016) (-0.08654,-0.04033) (-0.08881,-0.04056) (-0.09113,-0.04084) (-0.09349,-0.04118) (-0.09590,-0.04158) (-0.09836,-0.04202) (-0.10086,-0.04251) (-0.10341,-0.04306) (-0.10601,-0.04364) (-0.10866,-0.04428) (-0.11136,-0.04496) (-0.11410,-0.04568) (-0.11689,-0.04644) (-0.11973,-0.04725) (-0.12262,-0.04810) (-0.12556,-0.04899) (-0.12854,-0.04992) (-0.13158,-0.05088) (-0.13466,-0.05189) (-0.13779,-0.05293) (-0.14096,-0.05401) (-0.14418,-0.05512) (-0.14745,-0.05628) (-0.15076,-0.05746) (-0.15412,-0.05869) (-0.15752,-0.05995) (-0.16097,-0.06124) (-0.16445,-0.06257) (-0.16798,-0.06393) (-0.17156,-0.06532) (-0.17517,-0.06675) (-0.17882,-0.06821) (-0.18251,-0.06970) (-0.18623,-0.07123) (-0.19000,-0.07279) (-0.19379,-0.07438) (-0.19763,-0.07600) (-0.20149,-0.07766) (-0.20538,-0.07934) (-0.20931,-0.08106) (-0.21326,-0.08280) (-0.21724,-0.08458) (-0.22124,-0.08638) (-0.22527,-0.08822) (-0.22932,-0.09009) (-0.23340,-0.09198) (-0.23749,-0.09390) (-0.24159,-0.09585) (-0.24572,-0.09783) (-0.24986,-0.09984) (-0.25401,-0.10187) (-0.25817,-0.10393) (-0.26234,-0.10602) (-0.26652,-0.10813) (-0.27070,-0.11027) (-0.27489,-0.11243) (-0.27908,-0.11462) (-0.28327,-0.11683) (-0.28746,-0.11906) (-0.29164,-0.12132) (-0.29582,-0.12359) (-0.30000,-0.12589)};
\draw[-{Latex},structurecolor] (axis cs:-0.13158,-0.05088) -- (axis cs:-0.13466,-0.05189);
\draw[-{Latex},structurecolor] (axis cs:-0.21326,-0.08280) -- (axis cs:-0.21724,-0.08458);
\addplot[structurecolor,thin,no marks] coordinates {(-0.02356,0.30000) (-0.02490,0.29255) (-0.02620,0.28527) (-0.02746,0.27814) (-0.02869,0.27117) (-0.02989,0.26434) (-0.03106,0.25767) (-0.03220,0.25115) (-0.03331,0.24476) (-0.03440,0.23852) (-0.03546,0.23241) (-0.03650,0.22643) (-0.03752,0.22059) (-0.03851,0.21487) (-0.03949,0.20927) (-0.04045,0.20380) (-0.04139,0.19845) (-0.04232,0.19321) (-0.04323,0.18809) (-0.04413,0.18308) (-0.04501,0.17818) (-0.04588,0.17338) (-0.04674,0.16869) (-0.04760,0.16410) (-0.04844,0.15961) (-0.04927,0.15522) (-0.05009,0.15092) (-0.05091,0.14672) (-0.05171,0.14260) (-0.05252,0.13857) (-0.05331,0.13463) (-0.05411,0.13078) (-0.05489,0.12701) (-0.05568,0.12331) (-0.05646,0.11970) (-0.05723,0.11616) (-0.05801,0.11270) (-0.05878,0.10931) (-0.05955,0.10600) (-0.06032,0.10275) (-0.06108,0.09957) (-0.06185,0.09646) (-0.06262,0.09341) (-0.06338,0.09043) (-0.06415,0.08750) (-0.06492,0.08464) (-0.06569,0.08184) (-0.06646,0.07910) (-0.06723,0.07641) (-0.06801,0.07377) (-0.06878,0.07119) (-0.06956,0.06866) (-0.07034,0.06619) (-0.07113,0.06376) (-0.07191,0.06138) (-0.07271,0.05905) (-0.07350,0.05676) (-0.07430,0.05452) (-0.07510,0.05232) (-0.07590,0.05017) (-0.07671,0.04806) (-0.07753,0.04598) (-0.07835,0.04395) (-0.07917,0.04196) (-0.08000,0.04000)};
\draw[-{Latex},structurecolor] (axis cs:-0.04501,0.17818) -- (axis cs:-0.04588,0.17338);
\draw[-{Latex},structurecolor] (axis cs:-0.06338,0.09043) -- (axis cs:-0.06415,0.08750);
\addplot[structurecolor,thin,no marks] coordinates {(-0.08000,0.04000) (-0.08215,0.03513) (-0.08434,0.03049) (-0.08657,0.02606) (-0.08883,0.02183) (-0.09113,0.01779) (-0.09348,0.01391) (-0.09587,0.01020) (-0.09830,0.00663) (-0.10078,0.00320) (-0.10330,-0.00009) (-0.10587,-0.00327) (-0.10848,-0.00634) (-0.11115,-0.00930) (-0.11386,-0.01218) (-0.11662,-0.01496) (-0.11943,-0.01766) (-0.12228,-0.02029) (-0.12519,-0.02285) (-0.12815,-0.02535) (-0.13115,-0.02780) (-0.13420,-0.03019) (-0.13731,-0.03253) (-0.14046,-0.03484) (-0.14365,-0.03710) (-0.14690,-0.03934) (-0.15019,-0.04154) (-0.15353,-0.04372) (-0.15692,-0.04587) (-0.16035,-0.04801) (-0.16383,-0.05012) (-0.16735,-0.05223) (-0.17092,-0.05432) (-0.17452,-0.05641) (-0.17817,-0.05849) (-0.18186,-0.06056) (-0.18558,-0.06263) (-0.18935,-0.06471) (-0.19315,-0.06678) (-0.19699,-0.06886) (-0.20086,-0.07094) (-0.20477,-0.07302) (-0.20870,-0.07512) (-0.21267,-0.07722) (-0.21666,-0.07933) (-0.22068,-0.08145) (-0.22473,-0.08359) (-0.22880,-0.08574) (-0.23290,-0.08789) (-0.23701,-0.09007) (-0.24114,-0.09225) (-0.24529,-0.09446) (-0.24946,-0.09667) (-0.25364,-0.09891) (-0.25783,-0.10115) (-0.26203,-0.10342) (-0.26624,-0.10570) (-0.27046,-0.10800) (-0.27468,-0.11031) (-0.27890,-0.11264) (-0.28312,-0.11499) (-0.28735,-0.11735) (-0.29157,-0.11973) (-0.29579,-0.12213) (-0.30000,-0.12454)};
\draw[-{Latex},structurecolor] (axis cs:-0.13115,-0.02780) -- (axis cs:-0.13420,-0.03019);
\draw[-{Latex},structurecolor] (axis cs:-0.21267,-0.07722) -- (axis cs:-0.21666,-0.07933);
\addplot[structurecolor,thin,no marks] coordinates {(-0.04220,0.30000) (-0.04298,0.29611) (-0.04375,0.29227) (-0.04451,0.28846) (-0.04526,0.28470) (-0.04600,0.28098) (-0.04673,0.27730) (-0.04745,0.27366) (-0.04816,0.27006) (-0.04886,0.26650) (-0.04956,0.26297) (-0.05025,0.25949) (-0.05093,0.25604) (-0.05161,0.25263) (-0.05227,0.24926) (-0.05293,0.24593) (-0.05358,0.24263) (-0.05423,0.23936) (-0.05487,0.23614) (-0.05550,0.23295) (-0.05613,0.22979) (-0.05675,0.22667) (-0.05737,0.22358) (-0.05798,0.22052) (-0.05859,0.21750) (-0.05919,0.21451) (-0.05978,0.21155) (-0.06037,0.20863) (-0.06096,0.20574) (-0.06154,0.20288) (-0.06211,0.20005) (-0.06269,0.19725) (-0.06326,0.19448) (-0.06382,0.19174) (-0.06438,0.18903) (-0.06494,0.18635) (-0.06549,0.18370) (-0.06604,0.18108) (-0.06659,0.17848) (-0.06713,0.17592) (-0.06767,0.17338) (-0.06821,0.17087) (-0.06875,0.16838) (-0.06928,0.16593) (-0.06981,0.16350) (-0.07033,0.16109) (-0.07086,0.15871) (-0.07138,0.15636) (-0.07190,0.15403) (-0.07242,0.15173) (-0.07294,0.14945) (-0.07345,0.14720) (-0.07396,0.14497) (-0.07447,0.14277) (-0.07498,0.14058) (-0.07549,0.13842) (-0.07600,0.13629) (-0.07650,0.13418) (-0.07700,0.13209) (-0.07751,0.13002) (-0.07801,0.12797) (-0.07851,0.12595) (-0.07901,0.12394) (-0.07950,0.12196) (-0.08000,0.12000)};
\draw[-{Latex},structurecolor] (axis cs:-0.05613,0.22979) -- (axis cs:-0.05675,0.22667);
\draw[-{Latex},structurecolor] (axis cs:-0.06928,0.16593) -- (axis cs:-0.06981,0.16350);
\addplot[structurecolor,thin,no marks] coordinates {(-0.08000,0.12000) (-0.08250,0.11039) (-0.08500,0.10128) (-0.08748,0.09264) (-0.08997,0.08443) (-0.09247,0.07664) (-0.09497,0.06924) (-0.09750,0.06221) (-0.10005,0.05551) (-0.10262,0.04914) (-0.10522,0.04307) (-0.10785,0.03728) (-0.11051,0.03175) (-0.11321,0.02647) (-0.11595,0.02143) (-0.11873,0.01659) (-0.12154,0.01196) (-0.12440,0.00752) (-0.12730,0.00325) (-0.13025,-0.00085) (-0.13323,-0.00480) (-0.13627,-0.00861) (-0.13934,-0.01229) (-0.14247,-0.01585) (-0.14563,-0.01929) (-0.14884,-0.02263) (-0.15210,-0.02588) (-0.15540,-0.02903) (-0.15875,-0.03211) (-0.16214,-0.03510) (-0.16557,-0.03803) (-0.16905,-0.04090) (-0.17256,-0.04370) (-0.17612,-0.04646) (-0.17972,-0.04916) (-0.18336,-0.05183) (-0.18704,-0.05445) (-0.19075,-0.05704) (-0.19450,-0.05960) (-0.19829,-0.06214) (-0.20211,-0.06464) (-0.20596,-0.06713) (-0.20984,-0.06960) (-0.21376,-0.07205) (-0.21770,-0.07449) (-0.22167,-0.07692) (-0.22566,-0.07934) (-0.22968,-0.08175) (-0.23372,-0.08416) (-0.23778,-0.08656) (-0.24186,-0.08897) (-0.24596,-0.09137) (-0.25007,-0.09378) (-0.25419,-0.09619) (-0.25833,-0.09860) (-0.26248,-0.10101) (-0.26664,-0.10343) (-0.27080,-0.10586) (-0.27497,-0.10829) (-0.27914,-0.11074) (-0.28332,-0.11319) (-0.28749,-0.11564) (-0.29167,-0.11811) (-0.29584,-0.12059) (-0.30000,-0.12307)};
\draw[-{Latex},structurecolor] (axis cs:-0.13323,-0.00480) -- (axis cs:-0.13627,-0.00861);
\draw[-{Latex},structurecolor] (axis cs:-0.21376,-0.07205) -- (axis cs:-0.21770,-0.07449);
\addplot[structurecolor,thin,no marks] coordinates {(0.06563,-0.30000) (0.06554,-0.29603) (0.06547,-0.29210) (0.06541,-0.28822) (0.06536,-0.28439) (0.06532,-0.28060) (0.06530,-0.27686) (0.06529,-0.27316) (0.06529,-0.26951) (0.06530,-0.26590) (0.06533,-0.26234) (0.06536,-0.25881) (0.06541,-0.25533) (0.06547,-0.25189) (0.06554,-0.24849) (0.06562,-0.24513) (0.06571,-0.24181) (0.06581,-0.23853) (0.06592,-0.23529) (0.06604,-0.23208) (0.06617,-0.22892) (0.06631,-0.22579) (0.06646,-0.22269) (0.06662,-0.21964) (0.06679,-0.21662) (0.06696,-0.21363) (0.06715,-0.21068) (0.06734,-0.20776) (0.06755,-0.20487) (0.06776,-0.20202) (0.06798,-0.19920) (0.06821,-0.19642) (0.06845,-0.19366) (0.06870,-0.19094) (0.06895,-0.18825) (0.06921,-0.18559) (0.06948,-0.18295) (0.06976,-0.18035) (0.07005,-0.17778) (0.07034,-0.17524) (0.07064,-0.17272) (0.07095,-0.17024) (0.07127,-0.16778) (0.07159,-0.16535) (0.07192,-0.16294) (0.07226,-0.16056) (0.07261,-0.15821) (0.07296,-0.15589) (0.07332,-0.15359) (0.07368,-0.15131) (0.07406,-0.14906) (0.07444,-0.14684) (0.07483,-0.14464) (0.07522,-0.14246) (0.07562,-0.14031) (0.07603,-0.13818) (0.07644,-0.13607) (0.07687,-0.13398) (0.07729,-0.13192) (0.07773,-0.12988) (0.07817,-0.12786) (0.07862,-0.12587) (0.07907,-0.12389) (0.07953,-0.12193) (0.08000,-0.12000)};
\draw[-{Latex},structurecolor] (axis cs:0.06617,-0.22892) -- (axis cs:0.06631,-0.22579);
\draw[-{Latex},structurecolor] (axis cs:0.07159,-0.16535) -- (axis cs:0.07192,-0.16294);
\addplot[structurecolor,thin,no marks] coordinates {(0.08000,-0.12000) (0.08128,-0.11495) (0.08260,-0.11004) (0.08396,-0.10526) (0.08538,-0.10061) (0.08683,-0.09608) (0.08834,-0.09167) (0.08989,-0.08737) (0.09148,-0.08319) (0.09313,-0.07911) (0.09481,-0.07513) (0.09655,-0.07125) (0.09833,-0.06746) (0.10016,-0.06376) (0.10205,-0.06015) (0.10398,-0.05662) (0.10596,-0.05317) (0.10799,-0.04980) (0.11007,-0.04650) (0.11221,-0.04327) (0.11441,-0.04010) (0.11665,-0.03700) (0.11896,-0.03396) (0.12132,-0.03097) (0.12375,-0.02804) (0.12623,-0.02516) (0.12878,-0.02233) (0.13139,-0.01955) (0.13407,-0.01681) (0.13682,-0.01411) (0.13963,-0.01146) (0.14252,-0.00883) (0.14548,-0.00624) (0.14852,-0.00369) (0.15164,-0.00116) (0.15484,0.00135) (0.15812,0.00382) (0.16149,0.00628) (0.16495,0.00872) (0.16850,0.01114) (0.17214,0.01355) (0.17589,0.01594) (0.17973,0.01832) (0.18369,0.02070) (0.18775,0.02307) (0.19192,0.02544) (0.19622,0.02780) (0.20063,0.03017) (0.20517,0.03254) (0.20985,0.03492) (0.21466,0.03731) (0.21961,0.03972) (0.22471,0.04213) (0.22996,0.04457) (0.23537,0.04702) (0.24095,0.04950) (0.24671,0.05201) (0.25264,0.05454) (0.25877,0.05711) (0.26509,0.05972) (0.27162,0.06236) (0.27837,0.06505) (0.28534,0.06779) (0.29254,0.07058) (0.30000,0.07343)};
\draw[-{Latex},structurecolor] (axis cs:0.11441,-0.04010) -- (axis cs:0.11665,-0.03700);
\draw[-{Latex},structurecolor] (axis cs:0.18369,0.02070) -- (axis cs:0.18775,0.02307);
\addplot[structurecolor,thin,no marks] coordinates {(0.04975,-0.30000) (0.04936,-0.29228) (0.04900,-0.28474) (0.04870,-0.27739) (0.04844,-0.27022) (0.04822,-0.26322) (0.04804,-0.25639) (0.04790,-0.24972) (0.04780,-0.24321) (0.04774,-0.23686) (0.04771,-0.23066) (0.04772,-0.22461) (0.04776,-0.21870) (0.04782,-0.21293) (0.04792,-0.20730) (0.04805,-0.20180) (0.04821,-0.19643) (0.04840,-0.19118) (0.04861,-0.18606) (0.04885,-0.18106) (0.04911,-0.17617) (0.04940,-0.17139) (0.04971,-0.16673) (0.05004,-0.16217) (0.05040,-0.15772) (0.05078,-0.15336) (0.05118,-0.14911) (0.05160,-0.14495) (0.05205,-0.14089) (0.05251,-0.13692) (0.05299,-0.13303) (0.05349,-0.12924) (0.05401,-0.12552) (0.05455,-0.12189) (0.05511,-0.11834) (0.05569,-0.11487) (0.05628,-0.11147) (0.05690,-0.10814) (0.05753,-0.10489) (0.05817,-0.10171) (0.05884,-0.09859) (0.05952,-0.09554) (0.06022,-0.09256) (0.06094,-0.08963) (0.06167,-0.08677) (0.06243,-0.08397) (0.06319,-0.08122) (0.06398,-0.07853) (0.06478,-0.07590) (0.06560,-0.07331) (0.06644,-0.07078) (0.06729,-0.06830) (0.06816,-0.06587) (0.06905,-0.06348) (0.06996,-0.06114) (0.07088,-0.05885) (0.07182,-0.05660) (0.07278,-0.05439) (0.07376,-0.05222) (0.07475,-0.05009) (0.07576,-0.04800) (0.07679,-0.04595) (0.07784,-0.04393) (0.07891,-0.04195) (0.08000,-0.04000)};
\draw[-{Latex},structurecolor] (axis cs:0.04911,-0.17617) -- (axis cs:0.04940,-0.17139);
\draw[-{Latex},structurecolor] (axis cs:0.06094,-0.08963) -- (axis cs:0.06167,-0.08677);
\addplot[structurecolor,thin,no marks] coordinates {(0.08000,-0.04000) (0.08145,-0.03750) (0.08293,-0.03506) (0.08445,-0.03267) (0.08600,-0.03033) (0.08759,-0.02804) (0.08922,-0.02579) (0.09088,-0.02359) (0.09258,-0.02143) (0.09431,-0.01931) (0.09609,-0.01722) (0.09790,-0.01518) (0.09976,-0.01316) (0.10166,-0.01118) (0.10360,-0.00923) (0.10559,-0.00731) (0.10762,-0.00541) (0.10970,-0.00354) (0.11183,-0.00170) (0.11401,0.00012) (0.11623,0.00193) (0.11851,0.00371) (0.12084,0.00547) (0.12323,0.00722) (0.12567,0.00895) (0.12817,0.01067) (0.13073,0.01238) (0.13335,0.01407) (0.13603,0.01576) (0.13878,0.01744) (0.14160,0.01911) (0.14449,0.02078) (0.14744,0.02244) (0.15047,0.02410) (0.15358,0.02576) (0.15676,0.02743) (0.16003,0.02909) (0.16338,0.03076) (0.16682,0.03243) (0.17034,0.03411) (0.17396,0.03580) (0.17768,0.03749) (0.18149,0.03920) (0.18541,0.04092) (0.18943,0.04266) (0.19356,0.04441) (0.19781,0.04618) (0.20218,0.04797) (0.20667,0.04978) (0.21128,0.05162) (0.21603,0.05348) (0.22092,0.05536) (0.22595,0.05728) (0.23114,0.05923) (0.23647,0.06121) (0.24197,0.06323) (0.24764,0.06529) (0.25348,0.06739) (0.25951,0.06953) (0.26572,0.07173) (0.27214,0.07397) (0.27877,0.07627) (0.28561,0.07862) (0.29269,0.08104) (0.30000,0.08352)};
\draw[-{Latex},structurecolor] (axis cs:0.11623,0.00193) -- (axis cs:0.11851,0.00371);
\draw[-{Latex},structurecolor] (axis cs:0.18541,0.04092) -- (axis cs:0.18943,0.04266);
\addplot[structurecolor,thin,no marks] coordinates {(0.03417,0.30000) (0.03297,0.28673) (0.03191,0.27405) (0.03099,0.26194) (0.03019,0.25038) (0.02951,0.23934) (0.02894,0.22881) (0.02847,0.21876) (0.02809,0.20917) (0.02779,0.20002) (0.02757,0.19129) (0.02743,0.18297) (0.02735,0.17504) (0.02734,0.16748) (0.02739,0.16027) (0.02749,0.15340) (0.02765,0.14686) (0.02786,0.14062) (0.02811,0.13469) (0.02841,0.12904) (0.02875,0.12366) (0.02914,0.11854) (0.02956,0.11368) (0.03002,0.10905) (0.03051,0.10464) (0.03105,0.10046) (0.03161,0.09649) (0.03221,0.09271) (0.03285,0.08913) (0.03351,0.08573) (0.03421,0.08251) (0.03494,0.07945) (0.03570,0.07656) (0.03649,0.07382) (0.03732,0.07122) (0.03817,0.06877) (0.03906,0.06645) (0.03998,0.06427) (0.04094,0.06220) (0.04193,0.06026) (0.04295,0.05843) (0.04400,0.05671) (0.04509,0.05510) (0.04621,0.05359) (0.04738,0.05218) (0.04857,0.05086) (0.04981,0.04963) (0.05109,0.04849) (0.05240,0.04743) (0.05376,0.04646) (0.05516,0.04556) (0.05660,0.04474) (0.05808,0.04399) (0.05962,0.04331) (0.06120,0.04270) (0.06283,0.04216) (0.06451,0.04168) (0.06624,0.04127) (0.06803,0.04091) (0.06987,0.04062) (0.07177,0.04038) (0.07373,0.04021) (0.07576,0.04008) (0.07785,0.04001) (0.08000,0.04000)};
\draw[-{Latex},structurecolor] (axis cs:0.02875,0.12366) -- (axis cs:0.02914,0.11854);
\draw[-{Latex},structurecolor] (axis cs:0.04621,0.05359) -- (axis cs:0.04738,0.05218);
\addplot[structurecolor,thin,no marks] coordinates {(0.08000,0.04000) (0.08146,0.04002) (0.08296,0.04006) (0.08448,0.04013) (0.08604,0.04022) (0.08763,0.04033) (0.08926,0.04046) (0.09092,0.04061) (0.09262,0.04079) (0.09435,0.04099) (0.09612,0.04121) (0.09793,0.04145) (0.09979,0.04172) (0.10168,0.04200) (0.10361,0.04231) (0.10559,0.04264) (0.10762,0.04299) (0.10968,0.04337) (0.11180,0.04376) (0.11397,0.04418) (0.11618,0.04462) (0.11845,0.04509) (0.12077,0.04557) (0.12314,0.04608) (0.12557,0.04662) (0.12806,0.04717) (0.13061,0.04775) (0.13322,0.04836) (0.13589,0.04899) (0.13863,0.04964) (0.14143,0.05032) (0.14430,0.05102) (0.14725,0.05176) (0.15027,0.05251) (0.15337,0.05330) (0.15654,0.05411) (0.15980,0.05496) (0.16314,0.05583) (0.16657,0.05673) (0.17009,0.05766) (0.17370,0.05862) (0.17741,0.05962) (0.18122,0.06065) (0.18513,0.06171) (0.18915,0.06281) (0.19328,0.06394) (0.19753,0.06511) (0.20189,0.06632) (0.20638,0.06757) (0.21100,0.06886) (0.21576,0.07020) (0.22065,0.07157) (0.22569,0.07300) (0.23088,0.07447) (0.23623,0.07599) (0.24174,0.07756) (0.24742,0.07919) (0.25328,0.08087) (0.25933,0.08261) (0.26557,0.08441) (0.27201,0.08628) (0.27867,0.08821) (0.28554,0.09021) (0.29265,0.09228) (0.30000,0.09444)};
\draw[-{Latex},structurecolor] (axis cs:0.11618,0.04462) -- (axis cs:0.11845,0.04509);
\draw[-{Latex},structurecolor] (axis cs:0.18513,0.06171) -- (axis cs:0.18915,0.06281);
\addplot[structurecolor,thin,no marks] coordinates {(0.06041,0.30000) (0.06023,0.29528) (0.06006,0.29065) (0.05992,0.28609) (0.05980,0.28161) (0.05971,0.27721) (0.05963,0.27289) (0.05957,0.26864) (0.05953,0.26447) (0.05951,0.26037) (0.05951,0.25633) (0.05953,0.25237) (0.05956,0.24848) (0.05961,0.24466) (0.05968,0.24091) (0.05977,0.23722) (0.05987,0.23360) (0.05999,0.23004) (0.06012,0.22654) (0.06027,0.22311) (0.06043,0.21974) (0.06061,0.21643) (0.06081,0.21317) (0.06101,0.20998) (0.06124,0.20685) (0.06147,0.20377) (0.06172,0.20074) (0.06198,0.19778) (0.06226,0.19486) (0.06255,0.19200) (0.06285,0.18920) (0.06317,0.18644) (0.06349,0.18374) (0.06383,0.18108) (0.06418,0.17848) (0.06455,0.17592) (0.06492,0.17342) (0.06531,0.17096) (0.06571,0.16854) (0.06612,0.16617) (0.06654,0.16385) (0.06698,0.16157) (0.06742,0.15934) (0.06788,0.15714) (0.06835,0.15499) (0.06883,0.15289) (0.06932,0.15082) (0.06982,0.14879) (0.07033,0.14681) (0.07085,0.14486) (0.07139,0.14295) (0.07193,0.14108) (0.07249,0.13925) (0.07305,0.13746) (0.07363,0.13570) (0.07422,0.13397) (0.07482,0.13229) (0.07543,0.13063) (0.07605,0.12902) (0.07668,0.12743) (0.07732,0.12588) (0.07797,0.12436) (0.07864,0.12288) (0.07931,0.12142) (0.08000,0.12000)};
\draw[-{Latex},structurecolor] (axis cs:0.06043,0.21974) -- (axis cs:0.06061,0.21643);
\draw[-{Latex},structurecolor] (axis cs:0.06788,0.15714) -- (axis cs:0.06835,0.15499);
\addplot[structurecolor,thin,no marks] coordinates {(0.08000,0.12000) (0.08131,0.11743) (0.08265,0.11497) (0.08404,0.11261) (0.08547,0.11034) (0.08693,0.10818) (0.08844,0.10611) (0.08998,0.10413) (0.09157,0.10224) (0.09320,0.10044) (0.09488,0.09873) (0.09659,0.09710) (0.09836,0.09555) (0.10016,0.09408) (0.10202,0.09269) (0.10392,0.09137) (0.10587,0.09013) (0.10787,0.08896) (0.10992,0.08787) (0.11203,0.08684) (0.11419,0.08589) (0.11640,0.08500) (0.11867,0.08418) (0.12099,0.08342) (0.12338,0.08273) (0.12583,0.08210) (0.12834,0.08153) (0.13091,0.08102) (0.13355,0.08058) (0.13626,0.08020) (0.13905,0.07987) (0.14190,0.07960) (0.14483,0.07939) (0.14783,0.07924) (0.15092,0.07915) (0.15409,0.07911) (0.15734,0.07913) (0.16069,0.07921) (0.16412,0.07935) (0.16765,0.07954) (0.17127,0.07978) (0.17500,0.08009) (0.17883,0.08045) (0.18277,0.08086) (0.18683,0.08134) (0.19100,0.08187) (0.19529,0.08246) (0.19970,0.08312) (0.20425,0.08383) (0.20893,0.08460) (0.21376,0.08543) (0.21873,0.08633) (0.22385,0.08729) (0.22914,0.08832) (0.23459,0.08942) (0.24021,0.09058) (0.24601,0.09181) (0.25200,0.09312) (0.25819,0.09450) (0.26459,0.09596) (0.27120,0.09750) (0.27803,0.09912) (0.28510,0.10083) (0.29242,0.10263) (0.30000,0.10452)};
\draw[-{Latex},structurecolor] (axis cs:0.11419,0.08589) -- (axis cs:0.11640,0.08500);
\draw[-{Latex},structurecolor] (axis cs:0.18277,0.08086) -- (axis cs:0.18683,0.08134);
\addplot[structurecolor,thin,no marks] coordinates {(0.01765,0.30000) (0.01703,0.29496) (0.01644,0.29001) (0.01586,0.28513) (0.01531,0.28034) (0.01477,0.27562) (0.01425,0.27098) (0.01374,0.26641) (0.01325,0.26192) (0.01278,0.25750) (0.01233,0.25316) (0.01188,0.24889) (0.01146,0.24468) (0.01104,0.24055) (0.01064,0.23648) (0.01026,0.23249) (0.00988,0.22855) (0.00952,0.22468) (0.00917,0.22088) (0.00883,0.21714) (0.00851,0.21346) (0.00819,0.20984) (0.00788,0.20628) (0.00759,0.20278) (0.00730,0.19934) (0.00702,0.19595) (0.00676,0.19262) (0.00650,0.18935) (0.00625,0.18613) (0.00600,0.18296) (0.00577,0.17985) (0.00554,0.17679) (0.00532,0.17378) (0.00511,0.17082) (0.00490,0.16791) (0.00470,0.16505) (0.00451,0.16223) (0.00432,0.15947) (0.00414,0.15674) (0.00397,0.15407) (0.00380,0.15144) (0.00363,0.14885) (0.00347,0.14631) (0.00332,0.14381) (0.00317,0.14135) (0.00303,0.13893) (0.00289,0.13655) (0.00275,0.13422) (0.00262,0.13192) (0.00249,0.12966) (0.00237,0.12744) (0.00225,0.12526) (0.00214,0.12311) (0.00203,0.12100) (0.00192,0.11892) (0.00181,0.11688) (0.00171,0.11488) (0.00161,0.11290) (0.00152,0.11096) (0.00142,0.10906) (0.00133,0.10718) (0.00125,0.10534) (0.00116,0.10353) (0.00108,0.10175) (0.00100,0.10000)};
\draw[-{Latex},structurecolor] (axis cs:0.00851,0.21346) -- (axis cs:0.00819,0.20984);
\draw[-{Latex},structurecolor] (axis cs:0.00332,0.14381) -- (axis cs:0.00317,0.14135);
\addplot[structurecolor,thin,no marks] coordinates {(0.00100,0.10000) (0.00038,0.08558) (-0.00010,0.07321) (-0.00048,0.06260) (-0.00080,0.05350) (-0.00106,0.04569) (-0.00129,0.03900) (-0.00150,0.03325) (-0.00170,0.02833) (-0.00189,0.02410) (-0.00209,0.02047) (-0.00229,0.01735) (-0.00249,0.01467) (-0.00271,0.01236) (-0.00294,0.01037) (-0.00318,0.00864) (-0.00344,0.00715) (-0.00372,0.00586) (-0.00403,0.00473) (-0.00435,0.00374) (-0.00470,0.00288) (-0.00508,0.00211) (-0.00549,0.00142) (-0.00593,0.00080) (-0.00641,0.00024) (-0.00692,-0.00028) (-0.00748,-0.00076) (-0.00807,-0.00121) (-0.00872,-0.00164) (-0.00941,-0.00206) (-0.01016,-0.00247) (-0.01097,-0.00288) (-0.01184,-0.00329) (-0.01278,-0.00370) (-0.01379,-0.00413) (-0.01488,-0.00457) (-0.01605,-0.00503) (-0.01732,-0.00551) (-0.01867,-0.00602) (-0.02013,-0.00655) (-0.02170,-0.00712) (-0.02339,-0.00772) (-0.02520,-0.00836) (-0.02714,-0.00905) (-0.02922,-0.00978) (-0.03146,-0.01056) (-0.03385,-0.01140) (-0.03641,-0.01230) (-0.03915,-0.01326) (-0.04208,-0.01428) (-0.04521,-0.01538) (-0.04855,-0.01656) (-0.05210,-0.01781) (-0.05589,-0.01915) (-0.05992,-0.02059) (-0.06420,-0.02212) (-0.06874,-0.02375) (-0.07356,-0.02549) (-0.07865,-0.02734) (-0.08403,-0.02931) (-0.08971,-0.03141) (-0.09569,-0.03363) (-0.10198,-0.03599) (-0.10858,-0.03849) (-0.11549,-0.04113)};
\draw[-{Latex},structurecolor] (axis cs:-0.00470,0.00288) -- (axis cs:-0.00508,0.00211);
\draw[-{Latex},structurecolor] (axis cs:-0.02714,-0.00905) -- (axis cs:-0.02922,-0.00978);
\addplot[structurecolor,thin,no marks] coordinates {(0.01709,-0.30000) (0.01649,-0.29483) (0.01591,-0.28975) (0.01535,-0.28477) (0.01480,-0.27987) (0.01428,-0.27505) (0.01377,-0.27033) (0.01328,-0.26568) (0.01281,-0.26112) (0.01235,-0.25664) (0.01191,-0.25223) (0.01148,-0.24791) (0.01107,-0.24366) (0.01067,-0.23948) (0.01028,-0.23538) (0.00991,-0.23135) (0.00955,-0.22739) (0.00920,-0.22350) (0.00886,-0.21968) (0.00854,-0.21593) (0.00822,-0.21224) (0.00792,-0.20861) (0.00762,-0.20505) (0.00734,-0.20155) (0.00706,-0.19811) (0.00680,-0.19473) (0.00654,-0.19141) (0.00629,-0.18815) (0.00605,-0.18495) (0.00581,-0.18180) (0.00559,-0.17870) (0.00537,-0.17566) (0.00516,-0.17267) (0.00495,-0.16974) (0.00475,-0.16685) (0.00456,-0.16401) (0.00438,-0.16123) (0.00420,-0.15849) (0.00402,-0.15580) (0.00386,-0.15315) (0.00369,-0.15056) (0.00354,-0.14800) (0.00338,-0.14549) (0.00324,-0.14303) (0.00309,-0.14060) (0.00295,-0.13822) (0.00282,-0.13588) (0.00269,-0.13358) (0.00256,-0.13132) (0.00244,-0.12910) (0.00232,-0.12691) (0.00221,-0.12477) (0.00210,-0.12266) (0.00199,-0.12058) (0.00189,-0.11855) (0.00179,-0.11654) (0.00169,-0.11458) (0.00159,-0.11264) (0.00150,-0.11074) (0.00141,-0.10887) (0.00132,-0.10704) (0.00124,-0.10523) (0.00116,-0.10346) (0.00108,-0.10171) (0.00100,-0.10000)};
\draw[-{Latex},structurecolor] (axis cs:0.00822,-0.21224) -- (axis cs:0.00792,-0.20861);
\draw[-{Latex},structurecolor] (axis cs:0.00324,-0.14303) -- (axis cs:0.00309,-0.14060);
\addplot[structurecolor,thin,no marks] coordinates {(0.00100,-0.10000) (0.00038,-0.08549) (-0.00010,-0.07311) (-0.00048,-0.06256) (-0.00079,-0.05355) (-0.00106,-0.04587) (-0.00129,-0.03932) (-0.00150,-0.03374) (-0.00170,-0.02897) (-0.00190,-0.02491) (-0.00210,-0.02145) (-0.00230,-0.01851) (-0.00251,-0.01601) (-0.00273,-0.01388) (-0.00297,-0.01208) (-0.00322,-0.01056) (-0.00348,-0.00927) (-0.00377,-0.00819) (-0.00408,-0.00729) (-0.00441,-0.00654) (-0.00477,-0.00593) (-0.00516,-0.00543) (-0.00557,-0.00503) (-0.00602,-0.00473) (-0.00651,-0.00450) (-0.00703,-0.00433) (-0.00760,-0.00424) (-0.00820,-0.00420) (-0.00886,-0.00421) (-0.00957,-0.00426) (-0.01033,-0.00437) (-0.01115,-0.00451) (-0.01204,-0.00470) (-0.01299,-0.00493) (-0.01402,-0.00519) (-0.01513,-0.00550) (-0.01632,-0.00584) (-0.01760,-0.00622) (-0.01898,-0.00665) (-0.02047,-0.00711) (-0.02206,-0.00762) (-0.02377,-0.00818) (-0.02561,-0.00878) (-0.02759,-0.00944) (-0.02970,-0.01015) (-0.03197,-0.01091) (-0.03440,-0.01174) (-0.03700,-0.01263) (-0.03978,-0.01358) (-0.04275,-0.01461) (-0.04593,-0.01571) (-0.04931,-0.01689) (-0.05292,-0.01816) (-0.05676,-0.01951) (-0.06085,-0.02096) (-0.06518,-0.02251) (-0.06978,-0.02416) (-0.07466,-0.02592) (-0.07982,-0.02779) (-0.08526,-0.02978) (-0.09101,-0.03190) (-0.09706,-0.03415) (-0.10341,-0.03654) (-0.11008,-0.03907) (-0.11707,-0.04174)};
\draw[-{Latex},structurecolor] (axis cs:-0.00477,-0.00593) -- (axis cs:-0.00516,-0.00543);
\draw[-{Latex},structurecolor] (axis cs:-0.02759,-0.00944) -- (axis cs:-0.02970,-0.01015);
\addplot[only marks,mark=*,mark size=1.4pt] coordinates {(0,0)};
\end{axis}\end{tikzpicture}
\end{center}
\end{solution}
\begin{exercise}\label{cw:bank-5C-3}原题5C-3. 对 $x'=2x+y+xy^3$, $y'=x-2y-xy$ 重做5C-1：判原点类型与稳定性，画附近轨道.\end{exercise}
\begin{solution}（官方译审）$J_0=\begin{pmatrix}2&1\\1&-2\end{pmatrix}$，根 $\pm\sqrt5$；可取向量 $(1,\sqrt5-2)^T$ 与 $(1,-\sqrt5-2)^T$. 原点为不稳定鞍点，出入方向分别与这两向量相切，双曲型判据不受高阶项改变. 下图保留原方程中的 $xy^3$，不误转录成 $xy$.
\begin{center}
\begin{tikzpicture}\begin{axis}[width=.47\linewidth,height=5.2cm,axis lines=middle,xlabel={$x$},ylabel={$y$},xmin=-0.3,xmax=0.3,ymin=-0.3,ymax=0.3,title={原点附近：鞍点},tick label style={font=\footnotesize},clip=true]
\addplot[structurecolor,thin,no marks] coordinates {(0.03577,-0.30000) (0.03397,-0.29496) (0.03218,-0.29002) (0.03041,-0.28517) (0.02864,-0.28042) (0.02688,-0.27576) (0.02514,-0.27119) (0.02340,-0.26672) (0.02166,-0.26233) (0.01994,-0.25802) (0.01822,-0.25380) (0.01651,-0.24967) (0.01481,-0.24562) (0.01311,-0.24164) (0.01141,-0.23775) (0.00972,-0.23393) (0.00803,-0.23019) (0.00635,-0.22653) (0.00466,-0.22293) (0.00298,-0.21942) (0.00130,-0.21597) (-0.00038,-0.21259) (-0.00206,-0.20928) (-0.00374,-0.20604) (-0.00542,-0.20286) (-0.00710,-0.19975) (-0.00879,-0.19670) (-0.01048,-0.19372) (-0.01217,-0.19080) (-0.01386,-0.18794) (-0.01556,-0.18514) (-0.01727,-0.18240) (-0.01898,-0.17971) (-0.02069,-0.17709) (-0.02242,-0.17452) (-0.02415,-0.17201) (-0.02588,-0.16955) (-0.02763,-0.16714) (-0.02938,-0.16479) (-0.03115,-0.16249) (-0.03292,-0.16024) (-0.03471,-0.15805) (-0.03650,-0.15590) (-0.03831,-0.15380) (-0.04013,-0.15175) (-0.04196,-0.14975) (-0.04380,-0.14779) (-0.04566,-0.14589) (-0.04754,-0.14402) (-0.04943,-0.14221) (-0.05133,-0.14043) (-0.05325,-0.13871) (-0.05519,-0.13702) (-0.05715,-0.13538) (-0.05912,-0.13378) (-0.06111,-0.13222) (-0.06313,-0.13070) (-0.06516,-0.12923) (-0.06721,-0.12779) (-0.06929,-0.12640) (-0.07138,-0.12504) (-0.07350,-0.12372) (-0.07564,-0.12244) (-0.07781,-0.12120) (-0.08000,-0.12000)};
\draw[-{Latex},structurecolor] (axis cs:0.00130,-0.21597) -- (axis cs:-0.00038,-0.21259);
\draw[-{Latex},structurecolor] (axis cs:-0.03831,-0.15380) -- (axis cs:-0.04013,-0.15175);
\addplot[structurecolor,thin,no marks] coordinates {(-0.08000,-0.12000) (-0.08215,-0.11887) (-0.08432,-0.11777) (-0.08652,-0.11671) (-0.08874,-0.11569) (-0.09099,-0.11469) (-0.09327,-0.11373) (-0.09557,-0.11281) (-0.09791,-0.11191) (-0.10027,-0.11105) (-0.10266,-0.11022) (-0.10508,-0.10943) (-0.10754,-0.10866) (-0.11002,-0.10793) (-0.11254,-0.10723) (-0.11509,-0.10656) (-0.11768,-0.10592) (-0.12030,-0.10531) (-0.12296,-0.10473) (-0.12565,-0.10419) (-0.12838,-0.10367) (-0.13115,-0.10318) (-0.13396,-0.10273) (-0.13681,-0.10230) (-0.13970,-0.10191) (-0.14263,-0.10154) (-0.14560,-0.10120) (-0.14861,-0.10090) (-0.15167,-0.10062) (-0.15478,-0.10037) (-0.15793,-0.10016) (-0.16113,-0.09997) (-0.16437,-0.09981) (-0.16767,-0.09968) (-0.17101,-0.09958) (-0.17440,-0.09951) (-0.17785,-0.09947) (-0.18135,-0.09946) (-0.18490,-0.09948) (-0.18851,-0.09953) (-0.19218,-0.09960) (-0.19590,-0.09971) (-0.19968,-0.09985) (-0.20351,-0.10001) (-0.20741,-0.10021) (-0.21137,-0.10044) (-0.21540,-0.10070) (-0.21948,-0.10098) (-0.22363,-0.10130) (-0.22785,-0.10165) (-0.23214,-0.10203) (-0.23649,-0.10244) (-0.24092,-0.10288) (-0.24542,-0.10336) (-0.24999,-0.10386) (-0.25463,-0.10440) (-0.25935,-0.10497) (-0.26415,-0.10557) (-0.26902,-0.10620) (-0.27398,-0.10687) (-0.27901,-0.10757) (-0.28413,-0.10830) (-0.28933,-0.10907) (-0.29462,-0.10987) (-0.30000,-0.11071)};
\draw[-{Latex},structurecolor] (axis cs:-0.12838,-0.10367) -- (axis cs:-0.13115,-0.10318);
\draw[-{Latex},structurecolor] (axis cs:-0.20351,-0.10001) -- (axis cs:-0.20741,-0.10021);
\addplot[structurecolor,thin,no marks] coordinates {(0.06451,-0.30000) (0.06129,-0.28730) (0.05819,-0.27516) (0.05520,-0.26355) (0.05233,-0.25247) (0.04955,-0.24187) (0.04687,-0.23174) (0.04428,-0.22205) (0.04178,-0.21280) (0.03936,-0.20395) (0.03701,-0.19549) (0.03473,-0.18740) (0.03251,-0.17967) (0.03036,-0.17228) (0.02827,-0.16522) (0.02623,-0.15847) (0.02423,-0.15202) (0.02229,-0.14585) (0.02038,-0.13996) (0.01851,-0.13433) (0.01668,-0.12895) (0.01488,-0.12381) (0.01310,-0.11890) (0.01135,-0.11421) (0.00962,-0.10973) (0.00791,-0.10546) (0.00622,-0.10138) (0.00453,-0.09749) (0.00285,-0.09378) (0.00118,-0.09024) (-0.00049,-0.08687) (-0.00216,-0.08366) (-0.00383,-0.08060) (-0.00551,-0.07769) (-0.00720,-0.07492) (-0.00891,-0.07229) (-0.01063,-0.06979) (-0.01237,-0.06741) (-0.01413,-0.06516) (-0.01592,-0.06303) (-0.01774,-0.06101) (-0.01959,-0.05911) (-0.02147,-0.05731) (-0.02340,-0.05561) (-0.02536,-0.05402) (-0.02738,-0.05252) (-0.02944,-0.05112) (-0.03156,-0.04981) (-0.03374,-0.04859) (-0.03597,-0.04746) (-0.03827,-0.04642) (-0.04065,-0.04546) (-0.04309,-0.04458) (-0.04562,-0.04378) (-0.04822,-0.04306) (-0.05092,-0.04242) (-0.05371,-0.04185) (-0.05659,-0.04136) (-0.05958,-0.04095) (-0.06268,-0.04061) (-0.06589,-0.04034) (-0.06922,-0.04015) (-0.07268,-0.04002) (-0.07627,-0.03998) (-0.08000,-0.04000)};
\draw[-{Latex},structurecolor] (axis cs:0.01668,-0.12895) -- (axis cs:0.01488,-0.12381);
\draw[-{Latex},structurecolor] (axis cs:-0.02340,-0.05561) -- (axis cs:-0.02536,-0.05402);
\addplot[structurecolor,thin,no marks] coordinates {(-0.08000,-0.04000) (-0.08178,-0.04004) (-0.08360,-0.04009) (-0.08545,-0.04016) (-0.08733,-0.04024) (-0.08925,-0.04034) (-0.09120,-0.04045) (-0.09319,-0.04059) (-0.09521,-0.04073) (-0.09728,-0.04090) (-0.09938,-0.04108) (-0.10152,-0.04127) (-0.10370,-0.04149) (-0.10592,-0.04172) (-0.10818,-0.04196) (-0.11049,-0.04223) (-0.11283,-0.04251) (-0.11523,-0.04280) (-0.11767,-0.04312) (-0.12015,-0.04345) (-0.12268,-0.04380) (-0.12526,-0.04416) (-0.12789,-0.04455) (-0.13057,-0.04495) (-0.13331,-0.04537) (-0.13609,-0.04581) (-0.13893,-0.04627) (-0.14182,-0.04674) (-0.14477,-0.04724) (-0.14778,-0.04775) (-0.15084,-0.04829) (-0.15396,-0.04884) (-0.15715,-0.04942) (-0.16039,-0.05001) (-0.16370,-0.05063) (-0.16708,-0.05126) (-0.17052,-0.05192) (-0.17403,-0.05260) (-0.17760,-0.05330) (-0.18125,-0.05402) (-0.18497,-0.05477) (-0.18876,-0.05553) (-0.19262,-0.05632) (-0.19657,-0.05714) (-0.20059,-0.05798) (-0.20468,-0.05884) (-0.20886,-0.05973) (-0.21313,-0.06064) (-0.21747,-0.06158) (-0.22190,-0.06255) (-0.22642,-0.06354) (-0.23103,-0.06456) (-0.23573,-0.06561) (-0.24053,-0.06668) (-0.24542,-0.06778) (-0.25041,-0.06892) (-0.25549,-0.07008) (-0.26068,-0.07127) (-0.26597,-0.07250) (-0.27136,-0.07375) (-0.27687,-0.07504) (-0.28248,-0.07636) (-0.28821,-0.07772) (-0.29404,-0.07911) (-0.30000,-0.08053)};
\draw[-{Latex},structurecolor] (axis cs:-0.12268,-0.04380) -- (axis cs:-0.12526,-0.04416);
\draw[-{Latex},structurecolor] (axis cs:-0.19657,-0.05714) -- (axis cs:-0.20059,-0.05798);
\addplot[structurecolor,thin,no marks] coordinates {(-0.08408,0.30000) (-0.08253,0.29210) (-0.08104,0.28439) (-0.07961,0.27688) (-0.07822,0.26955) (-0.07690,0.26240) (-0.07562,0.25543) (-0.07440,0.24863) (-0.07322,0.24199) (-0.07210,0.23552) (-0.07103,0.22921) (-0.07001,0.22306) (-0.06903,0.21706) (-0.06810,0.21121) (-0.06723,0.20550) (-0.06639,0.19993) (-0.06561,0.19450) (-0.06487,0.18920) (-0.06417,0.18403) (-0.06353,0.17899) (-0.06292,0.17407) (-0.06236,0.16927) (-0.06185,0.16459) (-0.06138,0.16002) (-0.06095,0.15556) (-0.06056,0.15122) (-0.06022,0.14697) (-0.05992,0.14283) (-0.05967,0.13879) (-0.05946,0.13484) (-0.05929,0.13099) (-0.05916,0.12723) (-0.05907,0.12356) (-0.05903,0.11997) (-0.05903,0.11647) (-0.05907,0.11305) (-0.05915,0.10970) (-0.05928,0.10644) (-0.05945,0.10325) (-0.05966,0.10013) (-0.05992,0.09708) (-0.06022,0.09410) (-0.06056,0.09118) (-0.06094,0.08833) (-0.06137,0.08555) (-0.06185,0.08282) (-0.06237,0.08015) (-0.06293,0.07753) (-0.06354,0.07497) (-0.06420,0.07246) (-0.06490,0.07000) (-0.06565,0.06760) (-0.06645,0.06523) (-0.06730,0.06292) (-0.06819,0.06065) (-0.06913,0.05842) (-0.07013,0.05623) (-0.07118,0.05408) (-0.07227,0.05197) (-0.07342,0.04989) (-0.07463,0.04785) (-0.07589,0.04584) (-0.07720,0.04386) (-0.07857,0.04192) (-0.08000,0.04000)};
\draw[-{Latex},structurecolor] (axis cs:-0.06292,0.17407) -- (axis cs:-0.06236,0.16927);
\draw[-{Latex},structurecolor] (axis cs:-0.06094,0.08833) -- (axis cs:-0.06137,0.08555);
\addplot[structurecolor,thin,no marks] coordinates {(-0.08000,0.04000) (-0.08127,0.03838) (-0.08259,0.03677) (-0.08395,0.03518) (-0.08536,0.03361) (-0.08681,0.03206) (-0.08831,0.03052) (-0.08986,0.02900) (-0.09145,0.02749) (-0.09310,0.02600) (-0.09479,0.02452) (-0.09654,0.02305) (-0.09834,0.02159) (-0.10019,0.02014) (-0.10210,0.01870) (-0.10406,0.01727) (-0.10608,0.01585) (-0.10815,0.01443) (-0.11029,0.01303) (-0.11248,0.01162) (-0.11474,0.01022) (-0.11705,0.00883) (-0.11943,0.00743) (-0.12188,0.00604) (-0.12439,0.00466) (-0.12697,0.00327) (-0.12962,0.00188) (-0.13234,0.00049) (-0.13513,-0.00090) (-0.13799,-0.00229) (-0.14093,-0.00369) (-0.14395,-0.00509) (-0.14704,-0.00649) (-0.15022,-0.00790) (-0.15347,-0.00932) (-0.15681,-0.01074) (-0.16024,-0.01218) (-0.16375,-0.01362) (-0.16735,-0.01507) (-0.17104,-0.01653) (-0.17483,-0.01801) (-0.17871,-0.01949) (-0.18269,-0.02099) (-0.18677,-0.02251) (-0.19095,-0.02404) (-0.19523,-0.02558) (-0.19962,-0.02715) (-0.20412,-0.02873) (-0.20874,-0.03033) (-0.21346,-0.03195) (-0.21831,-0.03359) (-0.22327,-0.03525) (-0.22835,-0.03694) (-0.23356,-0.03865) (-0.23890,-0.04038) (-0.24437,-0.04215) (-0.24997,-0.04394) (-0.25571,-0.04576) (-0.26159,-0.04761) (-0.26761,-0.04949) (-0.27378,-0.05140) (-0.28010,-0.05335) (-0.28658,-0.05533) (-0.29321,-0.05735) (-0.30000,-0.05940)};
\draw[-{Latex},structurecolor] (axis cs:-0.11474,0.01022) -- (axis cs:-0.11705,0.00883);
\draw[-{Latex},structurecolor] (axis cs:-0.18677,-0.02251) -- (axis cs:-0.19095,-0.02404);
\addplot[structurecolor,thin,no marks] coordinates {(-0.09330,0.30000) (-0.09264,0.29598) (-0.09199,0.29201) (-0.09135,0.28808) (-0.09073,0.28421) (-0.09013,0.28038) (-0.08954,0.27660) (-0.08896,0.27287) (-0.08841,0.26918) (-0.08786,0.26553) (-0.08733,0.26193) (-0.08682,0.25838) (-0.08632,0.25487) (-0.08584,0.25140) (-0.08537,0.24797) (-0.08491,0.24459) (-0.08447,0.24124) (-0.08405,0.23794) (-0.08364,0.23468) (-0.08324,0.23145) (-0.08286,0.22827) (-0.08249,0.22512) (-0.08214,0.22201) (-0.08180,0.21894) (-0.08148,0.21591) (-0.08117,0.21292) (-0.08087,0.20996) (-0.08059,0.20703) (-0.08032,0.20414) (-0.08007,0.20129) (-0.07983,0.19847) (-0.07961,0.19568) (-0.07940,0.19293) (-0.07920,0.19021) (-0.07902,0.18753) (-0.07885,0.18487) (-0.07869,0.18225) (-0.07855,0.17965) (-0.07842,0.17709) (-0.07831,0.17456) (-0.07821,0.17206) (-0.07812,0.16959) (-0.07805,0.16715) (-0.07799,0.16473) (-0.07795,0.16235) (-0.07792,0.15999) (-0.07790,0.15766) (-0.07790,0.15536) (-0.07791,0.15308) (-0.07794,0.15083) (-0.07798,0.14861) (-0.07803,0.14641) (-0.07810,0.14424) (-0.07818,0.14209) (-0.07827,0.13997) (-0.07838,0.13787) (-0.07850,0.13579) (-0.07864,0.13374) (-0.07879,0.13171) (-0.07896,0.12970) (-0.07914,0.12772) (-0.07933,0.12576) (-0.07954,0.12382) (-0.07976,0.12190) (-0.08000,0.12000)};
\draw[-{Latex},structurecolor] (axis cs:-0.08286,0.22827) -- (axis cs:-0.08249,0.22512);
\draw[-{Latex},structurecolor] (axis cs:-0.07799,0.16473) -- (axis cs:-0.07795,0.16235);
\addplot[structurecolor,thin,no marks] coordinates {(-0.08000,0.12000) (-0.08054,0.11613) (-0.08114,0.11235) (-0.08180,0.10865) (-0.08252,0.10503) (-0.08331,0.10149) (-0.08417,0.09803) (-0.08509,0.09464) (-0.08607,0.09132) (-0.08713,0.08806) (-0.08825,0.08487) (-0.08944,0.08175) (-0.09070,0.07868) (-0.09203,0.07567) (-0.09343,0.07272) (-0.09491,0.06982) (-0.09646,0.06698) (-0.09808,0.06418) (-0.09979,0.06143) (-0.10157,0.05872) (-0.10343,0.05606) (-0.10537,0.05343) (-0.10739,0.05085) (-0.10950,0.04830) (-0.11170,0.04579) (-0.11398,0.04330) (-0.11635,0.04085) (-0.11881,0.03843) (-0.12137,0.03604) (-0.12403,0.03366) (-0.12678,0.03132) (-0.12963,0.02899) (-0.13258,0.02668) (-0.13564,0.02439) (-0.13881,0.02212) (-0.14208,0.01985) (-0.14547,0.01760) (-0.14897,0.01536) (-0.15259,0.01313) (-0.15633,0.01091) (-0.16020,0.00868) (-0.16419,0.00646) (-0.16832,0.00425) (-0.17257,0.00203) (-0.17696,-0.00020) (-0.18150,-0.00243) (-0.18618,-0.00466) (-0.19100,-0.00691) (-0.19598,-0.00916) (-0.20111,-0.01143) (-0.20641,-0.01371) (-0.21187,-0.01601) (-0.21749,-0.01833) (-0.22330,-0.02067) (-0.22927,-0.02304) (-0.23544,-0.02542) (-0.24179,-0.02784) (-0.24833,-0.03029) (-0.25507,-0.03276) (-0.26201,-0.03527) (-0.26917,-0.03782) (-0.27654,-0.04041) (-0.28413,-0.04304) (-0.29195,-0.04571) (-0.30000,-0.04843)};
\draw[-{Latex},structurecolor] (axis cs:-0.10343,0.05606) -- (axis cs:-0.10537,0.05343);
\draw[-{Latex},structurecolor] (axis cs:-0.17257,0.00203) -- (axis cs:-0.17696,-0.00020);
\addplot[structurecolor,thin,no marks] coordinates {(0.09293,-0.30000) (0.09227,-0.29593) (0.09163,-0.29191) (0.09101,-0.28795) (0.09041,-0.28403) (0.08981,-0.28017) (0.08924,-0.27636) (0.08868,-0.27259) (0.08814,-0.26887) (0.08761,-0.26521) (0.08710,-0.26158) (0.08660,-0.25801) (0.08611,-0.25448) (0.08565,-0.25100) (0.08519,-0.24755) (0.08475,-0.24416) (0.08433,-0.24080) (0.08392,-0.23749) (0.08353,-0.23422) (0.08314,-0.23100) (0.08278,-0.22781) (0.08243,-0.22466) (0.08209,-0.22155) (0.08176,-0.21848) (0.08145,-0.21545) (0.08116,-0.21246) (0.08087,-0.20950) (0.08061,-0.20658) (0.08035,-0.20370) (0.08011,-0.20085) (0.07988,-0.19804) (0.07967,-0.19526) (0.07947,-0.19252) (0.07928,-0.18981) (0.07910,-0.18713) (0.07894,-0.18449) (0.07880,-0.18187) (0.07866,-0.17929) (0.07854,-0.17674) (0.07843,-0.17422) (0.07834,-0.17173) (0.07826,-0.16927) (0.07819,-0.16684) (0.07814,-0.16444) (0.07809,-0.16207) (0.07806,-0.15973) (0.07805,-0.15741) (0.07805,-0.15512) (0.07806,-0.15286) (0.07808,-0.15062) (0.07812,-0.14841) (0.07817,-0.14623) (0.07823,-0.14407) (0.07830,-0.14194) (0.07839,-0.13983) (0.07849,-0.13774) (0.07861,-0.13568) (0.07874,-0.13364) (0.07888,-0.13163) (0.07903,-0.12963) (0.07920,-0.12766) (0.07938,-0.12571) (0.07957,-0.12379) (0.07978,-0.12188) (0.08000,-0.12000)};
\draw[-{Latex},structurecolor] (axis cs:0.08278,-0.22781) -- (axis cs:0.08243,-0.22466);
\draw[-{Latex},structurecolor] (axis cs:0.07814,-0.16444) -- (axis cs:0.07809,-0.16207);
\addplot[structurecolor,thin,no marks] coordinates {(0.08000,-0.12000) (0.08053,-0.11593) (0.08112,-0.11196) (0.08178,-0.10808) (0.08251,-0.10429) (0.08330,-0.10059) (0.08416,-0.09697) (0.08508,-0.09343) (0.08607,-0.08997) (0.08713,-0.08659) (0.08826,-0.08328) (0.08946,-0.08004) (0.09073,-0.07686) (0.09208,-0.07375) (0.09349,-0.07070) (0.09498,-0.06771) (0.09655,-0.06478) (0.09819,-0.06191) (0.09991,-0.05908) (0.10171,-0.05631) (0.10359,-0.05358) (0.10555,-0.05090) (0.10759,-0.04826) (0.10972,-0.04567) (0.11193,-0.04311) (0.11424,-0.04059) (0.11663,-0.03811) (0.11911,-0.03566) (0.12169,-0.03325) (0.12437,-0.03086) (0.12714,-0.02850) (0.13001,-0.02616) (0.13298,-0.02386) (0.13606,-0.02157) (0.13925,-0.01930) (0.14254,-0.01706) (0.14595,-0.01483) (0.14947,-0.01261) (0.15310,-0.01041) (0.15686,-0.00822) (0.16074,-0.00605) (0.16474,-0.00388) (0.16888,-0.00172) (0.17314,0.00043) (0.17755,0.00258) (0.18209,0.00473) (0.18677,0.00688) (0.19159,0.00902) (0.19657,0.01117) (0.20170,0.01332) (0.20699,0.01548) (0.21244,0.01764) (0.21805,0.01982) (0.22383,0.02200) (0.22979,0.02419) (0.23592,0.02639) (0.24224,0.02861) (0.24875,0.03084) (0.25545,0.03309) (0.26235,0.03536) (0.26945,0.03765) (0.27676,0.03996) (0.28428,0.04229) (0.29203,0.04465) (0.30000,0.04704)};
\draw[-{Latex},structurecolor] (axis cs:0.10359,-0.05358) -- (axis cs:0.10555,-0.05090);
\draw[-{Latex},structurecolor] (axis cs:0.17314,0.00043) -- (axis cs:0.17755,0.00258);
\addplot[structurecolor,thin,no marks] coordinates {(0.08357,-0.30000) (0.08205,-0.29194) (0.08058,-0.28410) (0.07917,-0.27647) (0.07782,-0.26903) (0.07652,-0.26179) (0.07528,-0.25474) (0.07409,-0.24787) (0.07296,-0.24118) (0.07187,-0.23467) (0.07084,-0.22832) (0.06986,-0.22214) (0.06892,-0.21612) (0.06804,-0.21025) (0.06720,-0.20453) (0.06641,-0.19895) (0.06566,-0.19352) (0.06496,-0.18823) (0.06431,-0.18307) (0.06370,-0.17804) (0.06313,-0.17313) (0.06261,-0.16835) (0.06213,-0.16369) (0.06170,-0.15914) (0.06130,-0.15470) (0.06095,-0.15038) (0.06064,-0.14616) (0.06037,-0.14204) (0.06015,-0.13802) (0.05996,-0.13410) (0.05982,-0.13027) (0.05971,-0.12654) (0.05965,-0.12289) (0.05963,-0.11933) (0.05965,-0.11585) (0.05971,-0.11246) (0.05981,-0.10914) (0.05995,-0.10590) (0.06013,-0.10273) (0.06035,-0.09964) (0.06061,-0.09661) (0.06091,-0.09365) (0.06126,-0.09076) (0.06165,-0.08793) (0.06207,-0.08516) (0.06254,-0.08245) (0.06305,-0.07980) (0.06361,-0.07720) (0.06421,-0.07466) (0.06485,-0.07217) (0.06553,-0.06973) (0.06626,-0.06734) (0.06703,-0.06500) (0.06785,-0.06270) (0.06871,-0.06045) (0.06962,-0.05824) (0.07058,-0.05607) (0.07159,-0.05393) (0.07264,-0.05184) (0.07374,-0.04978) (0.07489,-0.04776) (0.07609,-0.04577) (0.07734,-0.04382) (0.07864,-0.04189) (0.08000,-0.04000)};
\draw[-{Latex},structurecolor] (axis cs:0.06313,-0.17313) -- (axis cs:0.06261,-0.16835);
\draw[-{Latex},structurecolor] (axis cs:0.06165,-0.08793) -- (axis cs:0.06207,-0.08516);
\addplot[structurecolor,thin,no marks] coordinates {(0.08000,-0.04000) (0.08127,-0.03831) (0.08259,-0.03664) (0.08395,-0.03500) (0.08536,-0.03337) (0.08681,-0.03177) (0.08832,-0.03018) (0.08987,-0.02861) (0.09147,-0.02706) (0.09312,-0.02553) (0.09482,-0.02401) (0.09657,-0.02250) (0.09837,-0.02101) (0.10023,-0.01953) (0.10214,-0.01807) (0.10411,-0.01661) (0.10614,-0.01517) (0.10822,-0.01374) (0.11036,-0.01231) (0.11256,-0.01089) (0.11483,-0.00948) (0.11715,-0.00808) (0.11954,-0.00669) (0.12200,-0.00529) (0.12452,-0.00391) (0.12711,-0.00253) (0.12976,-0.00115) (0.13249,0.00023) (0.13529,0.00160) (0.13817,0.00298) (0.14111,0.00435) (0.14414,0.00573) (0.14724,0.00710) (0.15043,0.00848) (0.15369,0.00986) (0.15704,0.01124) (0.16047,0.01262) (0.16399,0.01401) (0.16760,0.01541) (0.17130,0.01681) (0.17509,0.01822) (0.17898,0.01963) (0.18296,0.02106) (0.18704,0.02249) (0.19123,0.02393) (0.19551,0.02538) (0.19991,0.02684) (0.20441,0.02832) (0.20902,0.02980) (0.21374,0.03130) (0.21858,0.03281) (0.22354,0.03434) (0.22861,0.03588) (0.23382,0.03744) (0.23914,0.03901) (0.24460,0.04060) (0.25019,0.04221) (0.25591,0.04384) (0.26177,0.04548) (0.26777,0.04715) (0.27391,0.04884) (0.28021,0.05054) (0.28665,0.05227) (0.29325,0.05403) (0.30000,0.05580)};
\draw[-{Latex},structurecolor] (axis cs:0.11483,-0.00948) -- (axis cs:0.11715,-0.00808);
\draw[-{Latex},structurecolor] (axis cs:0.18704,0.02249) -- (axis cs:0.19123,0.02393);
\addplot[structurecolor,thin,no marks] coordinates {(-0.06517,0.30000) (-0.06207,0.28811) (-0.05906,0.27668) (-0.05615,0.26571) (-0.05334,0.25516) (-0.05061,0.24503) (-0.04796,0.23531) (-0.04539,0.22597) (-0.04290,0.21700) (-0.04048,0.20839) (-0.03812,0.20013) (-0.03583,0.19220) (-0.03360,0.18459) (-0.03143,0.17730) (-0.02931,0.17029) (-0.02724,0.16358) (-0.02521,0.15714) (-0.02323,0.15097) (-0.02129,0.14505) (-0.01938,0.13938) (-0.01750,0.13395) (-0.01566,0.12874) (-0.01384,0.12376) (-0.01205,0.11898) (-0.01028,0.11441) (-0.00852,0.11004) (-0.00678,0.10586) (-0.00506,0.10186) (-0.00334,0.09803) (-0.00162,0.09437) (0.00009,0.09088) (0.00180,0.08755) (0.00351,0.08436) (0.00524,0.08132) (0.00696,0.07843) (0.00871,0.07567) (0.01046,0.07304) (0.01224,0.07054) (0.01404,0.06817) (0.01586,0.06591) (0.01770,0.06377) (0.01958,0.06174) (0.02150,0.05982) (0.02345,0.05800) (0.02544,0.05629) (0.02748,0.05467) (0.02956,0.05316) (0.03170,0.05173) (0.03389,0.05040) (0.03614,0.04915) (0.03845,0.04799) (0.04083,0.04692) (0.04329,0.04593) (0.04581,0.04502) (0.04842,0.04419) (0.05112,0.04344) (0.05390,0.04276) (0.05677,0.04216) (0.05975,0.04164) (0.06283,0.04118) (0.06602,0.04080) (0.06933,0.04050) (0.07276,0.04026) (0.07631,0.04010) (0.08000,0.04000)};
\draw[-{Latex},structurecolor] (axis cs:-0.01750,0.13395) -- (axis cs:-0.01566,0.12874);
\draw[-{Latex},structurecolor] (axis cs:0.02345,0.05800) -- (axis cs:0.02544,0.05629);
\addplot[structurecolor,thin,no marks] coordinates {(0.08000,0.04000) (0.08179,0.03998) (0.08362,0.03997) (0.08548,0.03999) (0.08737,0.04001) (0.08930,0.04005) (0.09126,0.04011) (0.09326,0.04018) (0.09529,0.04027) (0.09736,0.04038) (0.09947,0.04050) (0.10162,0.04064) (0.10381,0.04079) (0.10604,0.04096) (0.10831,0.04114) (0.11062,0.04134) (0.11298,0.04156) (0.11538,0.04179) (0.11783,0.04204) (0.12032,0.04230) (0.12286,0.04258) (0.12545,0.04288) (0.12808,0.04319) (0.13077,0.04352) (0.13351,0.04387) (0.13630,0.04423) (0.13914,0.04461) (0.14204,0.04501) (0.14500,0.04542) (0.14801,0.04585) (0.15108,0.04630) (0.15420,0.04677) (0.15739,0.04725) (0.16064,0.04775) (0.16396,0.04827) (0.16733,0.04881) (0.17078,0.04937) (0.17429,0.04994) (0.17786,0.05054) (0.18151,0.05115) (0.18523,0.05178) (0.18902,0.05243) (0.19289,0.05310) (0.19683,0.05379) (0.20084,0.05450) (0.20494,0.05523) (0.20912,0.05598) (0.21337,0.05675) (0.21772,0.05754) (0.22214,0.05835) (0.22665,0.05918) (0.23126,0.06004) (0.23595,0.06091) (0.24073,0.06181) (0.24561,0.06273) (0.25059,0.06368) (0.25566,0.06464) (0.26083,0.06563) (0.26610,0.06664) (0.27148,0.06768) (0.27696,0.06873) (0.28256,0.06982) (0.28826,0.07093) (0.29407,0.07206) (0.30000,0.07322)};
\draw[-{Latex},structurecolor] (axis cs:0.12286,0.04258) -- (axis cs:0.12545,0.04288);
\draw[-{Latex},structurecolor] (axis cs:0.19683,0.05379) -- (axis cs:0.20084,0.05450);
\addplot[structurecolor,thin,no marks] coordinates {(-0.03480,0.30000) (-0.03304,0.29523) (-0.03128,0.29054) (-0.02953,0.28593) (-0.02779,0.28140) (-0.02606,0.27695) (-0.02434,0.27257) (-0.02262,0.26827) (-0.02091,0.26404) (-0.01920,0.25989) (-0.01750,0.25581) (-0.01580,0.25180) (-0.01411,0.24787) (-0.01242,0.24400) (-0.01073,0.24020) (-0.00905,0.23647) (-0.00737,0.23280) (-0.00569,0.22920) (-0.00401,0.22567) (-0.00233,0.22220) (-0.00066,0.21879) (0.00102,0.21544) (0.00270,0.21216) (0.00438,0.20893) (0.00605,0.20577) (0.00773,0.20266) (0.00942,0.19962) (0.01110,0.19663) (0.01279,0.19369) (0.01449,0.19081) (0.01618,0.18799) (0.01788,0.18522) (0.01959,0.18250) (0.02130,0.17984) (0.02302,0.17723) (0.02475,0.17467) (0.02648,0.17216) (0.02822,0.16970) (0.02997,0.16729) (0.03172,0.16492) (0.03349,0.16261) (0.03527,0.16034) (0.03705,0.15812) (0.03885,0.15595) (0.04066,0.15382) (0.04248,0.15174) (0.04431,0.14970) (0.04615,0.14770) (0.04801,0.14575) (0.04988,0.14384) (0.05177,0.14197) (0.05367,0.14015) (0.05558,0.13836) (0.05752,0.13662) (0.05947,0.13492) (0.06143,0.13325) (0.06342,0.13163) (0.06542,0.13004) (0.06744,0.12850) (0.06948,0.12699) (0.07154,0.12552) (0.07362,0.12408) (0.07573,0.12269) (0.07785,0.12133) (0.08000,0.12000)};
\draw[-{Latex},structurecolor] (axis cs:-0.00066,0.21879) -- (axis cs:0.00102,0.21544);
\draw[-{Latex},structurecolor] (axis cs:0.03885,0.15595) -- (axis cs:0.04066,0.15382);
\addplot[structurecolor,thin,no marks] coordinates {(0.08000,0.12000) (0.08217,0.11871) (0.08437,0.11746) (0.08659,0.11624) (0.08883,0.11505) (0.09110,0.11390) (0.09340,0.11278) (0.09572,0.11170) (0.09807,0.11065) (0.10045,0.10963) (0.10286,0.10865) (0.10530,0.10770) (0.10777,0.10678) (0.11027,0.10589) (0.11280,0.10503) (0.11537,0.10421) (0.11797,0.10341) (0.12060,0.10265) (0.12327,0.10192) (0.12598,0.10121) (0.12872,0.10054) (0.13150,0.09990) (0.13432,0.09929) (0.13717,0.09870) (0.14007,0.09815) (0.14301,0.09762) (0.14599,0.09713) (0.14901,0.09666) (0.15207,0.09622) (0.15518,0.09581) (0.15834,0.09543) (0.16154,0.09508) (0.16479,0.09475) (0.16808,0.09446) (0.17143,0.09419) (0.17482,0.09395) (0.17827,0.09373) (0.18176,0.09355) (0.18532,0.09339) (0.18892,0.09326) (0.19258,0.09315) (0.19630,0.09308) (0.20007,0.09303) (0.20390,0.09301) (0.20779,0.09301) (0.21175,0.09304) (0.21576,0.09310) (0.21984,0.09319) (0.22398,0.09330) (0.22818,0.09344) (0.23246,0.09361) (0.23680,0.09381) (0.24121,0.09403) (0.24569,0.09428) (0.25024,0.09455) (0.25486,0.09485) (0.25956,0.09518) (0.26434,0.09554) (0.26919,0.09592) (0.27412,0.09633) (0.27913,0.09677) (0.28422,0.09724) (0.28940,0.09773) (0.29466,0.09825) (0.30000,0.09880)};
\draw[-{Latex},structurecolor] (axis cs:0.12872,0.10054) -- (axis cs:0.13150,0.09990);
\draw[-{Latex},structurecolor] (axis cs:0.20390,0.09301) -- (axis cs:0.20779,0.09301);
\addplot[only marks,mark=*,mark size=1.4pt] coordinates {(0,0)};
\end{axis}\end{tikzpicture}
\end{center}
\end{solution}
\begin{exercise}\label{cw:bank-5C-4}原题5C-4. 对 $x'=1-y$, $y'=x^2-y^2$ 执行 Notes GS.6 作图程序：求全部平衡点，分析各点，画附近轨道，再补与之相容的轨道；需要时使用方向场向量.\end{exercise}
\begin{solution}（官方译审并补证）平衡条件给 $y=1$, $x=\pm1$. $J=\begin{pmatrix}0&-1\\2x&-2y\end{pmatrix}$. 在 $(1,1)$，根 $-1\pm i$，为渐近稳定螺旋点；在该点右侧向量竖直向上，所以逆时针. 在 $(-1,1)$，根 $-1\pm\sqrt3$，为鞍点；对应方向可取 $(1,-\lambda)^T$，正根方向离开、负根方向进入. 用零流线补图：$y<1$ 向右，$y>1$ 向左；$|x|>|y|$ 向上，$|x|<|y|$ 向下. 这些符号在各区域共同决定箭头，轨道不能相交. 下图为原系统的有界窗口示意，不由窗口外截断推断全局终点.
\begin{center}
\begin{tikzpicture}\begin{axis}[width=.47\linewidth,height=5.2cm,axis lines=middle,xlabel={$x$},ylabel={$y$},xmin=-2,xmax=2.5,ymin=-0.5,ymax=2.5,title={两平衡点及零流线},tick label style={font=\footnotesize},clip=true]
\addplot[structurecolor,thin,no marks] coordinates {(-1.81715,-0.50000) (-1.80627,-0.47776) (-1.79555,-0.45565) (-1.78500,-0.43368) (-1.77460,-0.41183) (-1.76436,-0.39013) (-1.75428,-0.36856) (-1.74436,-0.34714) (-1.73459,-0.32585) (-1.72498,-0.30471) (-1.71552,-0.28372) (-1.70622,-0.26288) (-1.69707,-0.24218) (-1.68806,-0.22164) (-1.67921,-0.20125) (-1.67051,-0.18102) (-1.66195,-0.16095) (-1.65354,-0.14103) (-1.64527,-0.12128) (-1.63715,-0.10169) (-1.62917,-0.08226) (-1.62133,-0.06299) (-1.61363,-0.04390) (-1.60607,-0.02497) (-1.59865,-0.00620) (-1.59137,0.01239) (-1.58422,0.03081) (-1.57720,0.04905) (-1.57032,0.06713) (-1.56357,0.08503) (-1.55695,0.10275) (-1.55045,0.12030) (-1.54409,0.13767) (-1.53785,0.15487) (-1.53174,0.17189) (-1.52575,0.18873) (-1.51988,0.20539) (-1.51413,0.22187) (-1.50850,0.23818) (-1.50300,0.25430) (-1.49761,0.27025) (-1.49233,0.28602) (-1.48717,0.30160) (-1.48212,0.31701) (-1.47719,0.33224) (-1.47236,0.34729) (-1.46765,0.36216) (-1.46304,0.37685) (-1.45854,0.39137) (-1.45414,0.40571) (-1.44985,0.41987) (-1.44566,0.43385) (-1.44158,0.44766) (-1.43759,0.46130) (-1.43370,0.47476) (-1.42991,0.48805) (-1.42622,0.50116) (-1.42262,0.51411) (-1.41912,0.52688) (-1.41571,0.53948) (-1.41239,0.55192) (-1.40916,0.56419) (-1.40602,0.57629) (-1.40297,0.58823) (-1.40000,0.60000)};
\draw[-{Latex},structurecolor] (axis cs:-1.62917,-0.08226) -- (axis cs:-1.62133,-0.06299);
\draw[-{Latex},structurecolor] (axis cs:-1.48212,0.31701) -- (axis cs:-1.47719,0.33224);
\addplot[structurecolor,thin,no marks] coordinates {(-1.40000,0.60000) (-1.38827,0.64853) (-1.37800,0.69420) (-1.36913,0.73713) (-1.36155,0.77744) (-1.35519,0.81526) (-1.34998,0.85072) (-1.34584,0.88397) (-1.34270,0.91513) (-1.34051,0.94435) (-1.33921,0.97175) (-1.33873,0.99746) (-1.33903,1.02161) (-1.34007,1.04433) (-1.34179,1.06572) (-1.34416,1.08590) (-1.34715,1.10498) (-1.35071,1.12305) (-1.35483,1.14021) (-1.35947,1.15654) (-1.36460,1.17214) (-1.37022,1.18708) (-1.37629,1.20143) (-1.38280,1.21527) (-1.38974,1.22866) (-1.39709,1.24165) (-1.40484,1.25431) (-1.41298,1.26668) (-1.42150,1.27882) (-1.43040,1.29078) (-1.43967,1.30258) (-1.44931,1.31428) (-1.45931,1.32592) (-1.46968,1.33751) (-1.48041,1.34911) (-1.49150,1.36074) (-1.50295,1.37243) (-1.51477,1.38421) (-1.52696,1.39609) (-1.53953,1.40812) (-1.55247,1.42031) (-1.56580,1.43268) (-1.57952,1.44526) (-1.59363,1.45807) (-1.60815,1.47112) (-1.62308,1.48445) (-1.63843,1.49805) (-1.65421,1.51197) (-1.67043,1.52621) (-1.68710,1.54079) (-1.70423,1.55573) (-1.72183,1.57105) (-1.73992,1.58676) (-1.75851,1.60289) (-1.77760,1.61946) (-1.79722,1.63648) (-1.81739,1.65396) (-1.83810,1.67194) (-1.85938,1.69042) (-1.88125,1.70943) (-1.90373,1.72898) (-1.92682,1.74911) (-1.95055,1.76982) (-1.97494,1.79114) (-2.00000,1.81309)};
\draw[-{Latex},structurecolor] (axis cs:-1.36460,1.17214) -- (axis cs:-1.37022,1.18708);
\draw[-{Latex},structurecolor] (axis cs:-1.59363,1.45807) -- (axis cs:-1.60815,1.47112);
\addplot[structurecolor,thin,no marks] coordinates {(-1.83204,-0.50000) (-1.80720,-0.44800) (-1.78323,-0.39670) (-1.76012,-0.34614) (-1.73786,-0.29634) (-1.71643,-0.24732) (-1.69582,-0.19911) (-1.67601,-0.15174) (-1.65699,-0.10523) (-1.63875,-0.05959) (-1.62128,-0.01484) (-1.60455,0.02899) (-1.58855,0.07190) (-1.57326,0.11386) (-1.55868,0.15487) (-1.54477,0.19493) (-1.53154,0.23402) (-1.51895,0.27214) (-1.50700,0.30928) (-1.49567,0.34545) (-1.48494,0.38065) (-1.47479,0.41489) (-1.46521,0.44816) (-1.45618,0.48047) (-1.44770,0.51185) (-1.43973,0.54228) (-1.43226,0.57180) (-1.42529,0.60041) (-1.41879,0.62812) (-1.41275,0.65496) (-1.40716,0.68093) (-1.40199,0.70607) (-1.39725,0.73038) (-1.39290,0.75390) (-1.38895,0.77663) (-1.38537,0.79859) (-1.38216,0.81982) (-1.37929,0.84033) (-1.37677,0.86015) (-1.37458,0.87929) (-1.37270,0.89778) (-1.37113,0.91564) (-1.36985,0.93289) (-1.36886,0.94956) (-1.36815,0.96566) (-1.36770,0.98122) (-1.36751,0.99626) (-1.36757,1.01081) (-1.36788,1.02487) (-1.36841,1.03847) (-1.36917,1.05164) (-1.37015,1.06439) (-1.37134,1.07673) (-1.37273,1.08870) (-1.37432,1.10030) (-1.37611,1.11156) (-1.37808,1.12249) (-1.38024,1.13310) (-1.38257,1.14343) (-1.38507,1.15347) (-1.38774,1.16325) (-1.39057,1.17277) (-1.39356,1.18207) (-1.39670,1.19114) (-1.40000,1.20000)};
\draw[-{Latex},structurecolor] (axis cs:-1.48494,0.38065) -- (axis cs:-1.47479,0.41489);
\draw[-{Latex},structurecolor] (axis cs:-1.36886,0.94956) -- (axis cs:-1.36815,0.96566);
\addplot[structurecolor,thin,no marks] coordinates {(-1.40000,1.20000) (-1.40399,1.20998) (-1.40816,1.21972) (-1.41253,1.22924) (-1.41707,1.23856) (-1.42180,1.24768) (-1.42670,1.25663) (-1.43178,1.26543) (-1.43702,1.27408) (-1.44243,1.28261) (-1.44801,1.29102) (-1.45375,1.29933) (-1.45965,1.30755) (-1.46571,1.31570) (-1.47192,1.32378) (-1.47829,1.33181) (-1.48482,1.33979) (-1.49151,1.34774) (-1.49834,1.35566) (-1.50533,1.36357) (-1.51248,1.37148) (-1.51978,1.37939) (-1.52723,1.38731) (-1.53484,1.39525) (-1.54260,1.40321) (-1.55051,1.41121) (-1.55859,1.41926) (-1.56682,1.42735) (-1.57520,1.43549) (-1.58375,1.44370) (-1.59245,1.45198) (-1.60132,1.46033) (-1.61035,1.46877) (-1.61955,1.47729) (-1.62891,1.48590) (-1.63844,1.49461) (-1.64814,1.50342) (-1.65802,1.51234) (-1.66806,1.52138) (-1.67829,1.53053) (-1.68869,1.53981) (-1.69928,1.54922) (-1.71005,1.55877) (-1.72100,1.56845) (-1.73215,1.57828) (-1.74349,1.58826) (-1.75502,1.59839) (-1.76676,1.60868) (-1.77869,1.61914) (-1.79083,1.62976) (-1.80318,1.64056) (-1.81574,1.65154) (-1.82851,1.66270) (-1.84150,1.67405) (-1.85472,1.68559) (-1.86816,1.69733) (-1.88183,1.70927) (-1.89574,1.72143) (-1.90989,1.73379) (-1.92427,1.74637) (-1.93891,1.75918) (-1.95379,1.77222) (-1.96893,1.78549) (-1.98433,1.79900) (-2.00000,1.81275)};
\draw[-{Latex},structurecolor] (axis cs:-1.51248,1.37148) -- (axis cs:-1.51978,1.37939);
\draw[-{Latex},structurecolor] (axis cs:-1.72100,1.56845) -- (axis cs:-1.73215,1.57828);
\addplot[structurecolor,thin,no marks] coordinates {(-1.57685,-0.50000) (-1.55391,-0.46570) (-1.53150,-0.43200) (-1.50960,-0.39890) (-1.48821,-0.36641) (-1.46732,-0.33453) (-1.44691,-0.30326) (-1.42699,-0.27260) (-1.40753,-0.24255) (-1.38854,-0.21312) (-1.36999,-0.18430) (-1.35189,-0.15609) (-1.33422,-0.12850) (-1.31697,-0.10152) (-1.30013,-0.07515) (-1.28370,-0.04939) (-1.26766,-0.02424) (-1.25200,0.00031) (-1.23672,0.02426) (-1.22181,0.04762) (-1.20725,0.07038) (-1.19304,0.09256) (-1.17917,0.11416) (-1.16563,0.13519) (-1.15241,0.15564) (-1.13950,0.17553) (-1.12690,0.19486) (-1.11459,0.21365) (-1.10256,0.23189) (-1.09082,0.24960) (-1.07934,0.26679) (-1.06813,0.28345) (-1.05717,0.29960) (-1.04645,0.31526) (-1.03598,0.33041) (-1.02573,0.34509) (-1.01571,0.35929) (-1.00591,0.37302) (-0.99631,0.38629) (-0.98691,0.39912) (-0.97771,0.41151) (-0.96870,0.42346) (-0.95987,0.43499) (-0.95122,0.44612) (-0.94273,0.45684) (-0.93441,0.46716) (-0.92624,0.47710) (-0.91823,0.48667) (-0.91036,0.49587) (-0.90263,0.50471) (-0.89503,0.51320) (-0.88756,0.52135) (-0.88022,0.52916) (-0.87299,0.53666) (-0.86588,0.54384) (-0.85888,0.55071) (-0.85198,0.55728) (-0.84518,0.56356) (-0.83847,0.56956) (-0.83186,0.57528) (-0.82533,0.58073) (-0.81888,0.58592) (-0.81252,0.59086) (-0.80622,0.59555) (-0.80000,0.60000)};
\draw[-{Latex},structurecolor] (axis cs:-1.20725,0.07038) -- (axis cs:-1.19304,0.09256);
\draw[-{Latex},structurecolor] (axis cs:-0.95122,0.44612) -- (axis cs:-0.94273,0.45684);
\addplot[structurecolor,thin,no marks] coordinates {(-0.80000,0.60000) (-0.74034,0.63290) (-0.68435,0.64785) (-0.62960,0.64939) (-0.57426,0.64089) (-0.51698,0.62479) (-0.45671,0.60295) (-0.39268,0.57680) (-0.32430,0.54755) (-0.25119,0.51635) (-0.17312,0.48436) (-0.09007,0.45286) (-0.00224,0.42333) (0.08994,0.39746) (0.18576,0.37716) (0.28418,0.36450) (0.38385,0.36156) (0.48309,0.37026) (0.57996,0.39201) (0.67238,0.42737) (0.75823,0.47581) (0.83560,0.53545) (0.90297,0.60320) (0.95938,0.67513) (1.00450,0.74700) (1.03865,0.81489) (1.06271,0.87577) (1.07794,0.92769) (1.08582,0.96984) (1.08787,1.00236) (1.08554,1.02604) (1.08013,1.04206) (1.07273,1.05174) (1.06422,1.05642) (1.05530,1.05729) (1.04646,1.05542) (1.03808,1.05167) (1.03037,1.04675) (1.02350,1.04120) (1.01751,1.03544) (1.01242,1.02978) (1.00819,1.02442) (1.00476,1.01950) (1.00206,1.01511) (1.00001,1.01129) (0.99851,1.00802) (0.99747,1.00531) (0.99682,1.00311) (0.99648,1.00137) (0.99637,1.00004) (0.99645,0.99906) (0.99665,0.99839) (0.99694,0.99796) (0.99728,0.99774) (0.99764,0.99766) (0.99800,0.99771) (0.99835,0.99784) (0.99868,0.99802) (0.99897,0.99824) (0.99923,0.99847) (0.99945,0.99871) (0.99963,0.99893) (0.99978,0.99914) (0.99990,0.99933) (0.99999,0.99950)};
\draw[-{Latex},structurecolor] (axis cs:0.75823,0.47581) -- (axis cs:0.83560,0.53545);
\draw[-{Latex},structurecolor] (axis cs:1.00206,1.01511) -- (axis cs:1.00001,1.01129);
\addplot[structurecolor,thin,no marks] coordinates {(-0.44550,2.50000) (-0.45825,2.44880) (-0.47057,2.39983) (-0.48248,2.35294) (-0.49399,2.30802) (-0.50512,2.26495) (-0.51588,2.22363) (-0.52630,2.18397) (-0.53637,2.14586) (-0.54613,2.10924) (-0.55557,2.07402) (-0.56471,2.04013) (-0.57357,2.00750) (-0.58215,1.97607) (-0.59046,1.94578) (-0.59852,1.91657) (-0.60633,1.88840) (-0.61389,1.86121) (-0.62123,1.83496) (-0.62834,1.80961) (-0.63524,1.78512) (-0.64193,1.76144) (-0.64842,1.73855) (-0.65471,1.71641) (-0.66082,1.69498) (-0.66674,1.67424) (-0.67249,1.65416) (-0.67806,1.63471) (-0.68347,1.61587) (-0.68872,1.59761) (-0.69381,1.57990) (-0.69876,1.56274) (-0.70355,1.54609) (-0.70821,1.52994) (-0.71272,1.51426) (-0.71711,1.49904) (-0.72136,1.48427) (-0.72549,1.46992) (-0.72949,1.45598) (-0.73338,1.44244) (-0.73715,1.42928) (-0.74081,1.41649) (-0.74436,1.40406) (-0.74780,1.39197) (-0.75114,1.38020) (-0.75438,1.36876) (-0.75752,1.35763) (-0.76057,1.34680) (-0.76352,1.33625) (-0.76639,1.32598) (-0.76917,1.31599) (-0.77186,1.30625) (-0.77447,1.29677) (-0.77699,1.28753) (-0.77944,1.27853) (-0.78181,1.26975) (-0.78411,1.26120) (-0.78633,1.25287) (-0.78849,1.24474) (-0.79057,1.23681) (-0.79258,1.22908) (-0.79453,1.22154) (-0.79642,1.21418) (-0.79824,1.20701) (-0.80000,1.20000)};
\draw[-{Latex},structurecolor] (axis cs:-0.63524,1.78512) -- (axis cs:-0.64193,1.76144);
\draw[-{Latex},structurecolor] (axis cs:-0.74780,1.39197) -- (axis cs:-0.75114,1.38020);
\addplot[structurecolor,thin,no marks] coordinates {(-0.80000,1.20000) (-0.82274,1.09804) (-0.83234,1.02916) (-0.83288,0.98029) (-0.82681,0.94364) (-0.81565,0.91439) (-0.80028,0.88944) (-0.78122,0.86677) (-0.75869,0.84501) (-0.73278,0.82325) (-0.70343,0.80084) (-0.67049,0.77734) (-0.63377,0.75242) (-0.59304,0.72589) (-0.54802,0.69764) (-0.49846,0.66766) (-0.44408,0.63603) (-0.38464,0.60296) (-0.31995,0.56879) (-0.24986,0.53405) (-0.17434,0.49945) (-0.09349,0.46596) (-0.00758,0.43479) (0.08294,0.40744) (0.17731,0.38568) (0.27452,0.37146) (0.37323,0.36678) (0.47180,0.37352) (0.56833,0.39311) (0.66075,0.42622) (0.74695,0.47243) (0.82502,0.53006) (0.89337,0.59619) (0.95096,0.66702) (0.99740,0.73837) (1.03290,0.80632) (1.05827,0.86772) (1.07469,0.92050) (1.08361,0.96371) (1.08653,0.99735) (1.08490,1.02211) (1.08003,1.03911) (1.07302,1.04965) (1.06478,1.05504) (1.05602,1.05649) (1.04728,1.05506) (1.03892,1.05164) (1.03121,1.04696) (1.02428,1.04158) (1.01823,1.03592) (1.01306,1.03031) (1.00874,1.02496) (1.00523,1.02004) (1.00245,1.01562) (1.00032,1.01174) (0.99875,1.00843) (0.99766,1.00567) (0.99696,1.00341) (0.99657,1.00162) (0.99643,1.00024) (0.99648,0.99922) (0.99666,0.99851) (0.99693,0.99805) (0.99726,0.99780) (0.99761,0.99770)};
\draw[-{Latex},structurecolor] (axis cs:-0.17434,0.49945) -- (axis cs:-0.09349,0.46596);
\draw[-{Latex},structurecolor] (axis cs:1.06478,1.05504) -- (axis cs:1.05602,1.05649);
\addplot[structurecolor,thin,no marks] coordinates {(0.15931,2.50000) (0.13043,2.38099) (0.10382,2.27269) (0.07929,2.17374) (0.05667,2.08298) (0.03579,1.99946) (0.01652,1.92235) (-0.00126,1.85096) (-0.01766,1.78468) (-0.03278,1.72300) (-0.04670,1.66546) (-0.05951,1.61166) (-0.07127,1.56126) (-0.08205,1.51395) (-0.09192,1.46947) (-0.10091,1.42756) (-0.10909,1.38803) (-0.11650,1.35067) (-0.12317,1.31531) (-0.12916,1.28180) (-0.13449,1.25000) (-0.13920,1.21979) (-0.14332,1.19104) (-0.14688,1.16365) (-0.14990,1.13754) (-0.15241,1.11260) (-0.15442,1.08877) (-0.15598,1.06597) (-0.15708,1.04413) (-0.15775,1.02319) (-0.15801,1.00310) (-0.15788,0.98380) (-0.15737,0.96524) (-0.15649,0.94739) (-0.15526,0.93019) (-0.15369,0.91362) (-0.15180,0.89764) (-0.14959,0.88220) (-0.14708,0.86729) (-0.14427,0.85288) (-0.14118,0.83893) (-0.13781,0.82543) (-0.13417,0.81234) (-0.13028,0.79966) (-0.12614,0.78736) (-0.12175,0.77542) (-0.11713,0.76382) (-0.11227,0.75255) (-0.10720,0.74159) (-0.10191,0.73094) (-0.09640,0.72057) (-0.09070,0.71047) (-0.08479,0.70064) (-0.07868,0.69106) (-0.07239,0.68172) (-0.06591,0.67262) (-0.05926,0.66374) (-0.05242,0.65508) (-0.04542,0.64662) (-0.03824,0.63838) (-0.03091,0.63033) (-0.02341,0.62247) (-0.01576,0.61480) (-0.00795,0.60731) (0.00000,0.60000)};
\draw[-{Latex},structurecolor] (axis cs:-0.13449,1.25000) -- (axis cs:-0.13920,1.21979);
\draw[-{Latex},structurecolor] (axis cs:-0.13028,0.79966) -- (axis cs:-0.12614,0.78736);
\addplot[structurecolor,thin,no marks] coordinates {(0.00000,0.60000) (0.06663,0.54879) (0.14051,0.50707) (0.22019,0.47467) (0.30417,0.45205) (0.39087,0.44004) (0.47855,0.43958) (0.56536,0.45140) (0.64934,0.47576) (0.72856,0.51211) (0.80125,0.55903) (0.86593,0.61415) (0.92156,0.67437) (0.96760,0.73627) (1.00407,0.79650) (1.03143,0.85224) (1.05058,0.90145) (1.06263,0.94297) (1.06882,0.97645) (1.07040,1.00219) (1.06850,1.02092) (1.06417,1.03361) (1.05826,1.04131) (1.05147,1.04506) (1.04434,1.04580) (1.03727,1.04435) (1.03055,1.04140) (1.02438,1.03750) (1.01886,1.03308) (1.01405,1.02848) (1.00996,1.02394) (1.00656,1.01964) (1.00380,1.01569) (1.00164,1.01216) (0.99998,1.00907) (0.99878,1.00645) (0.99795,1.00426) (0.99742,1.00248) (0.99715,1.00108) (0.99707,1.00001) (0.99713,0.99923) (0.99730,0.99869) (0.99753,0.99835) (0.99780,0.99817) (0.99810,0.99811) (0.99839,0.99815) (0.99867,0.99825) (0.99893,0.99840) (0.99917,0.99858) (0.99938,0.99877) (0.99955,0.99896) (0.99970,0.99914) (0.99982,0.99931) (0.99992,0.99946) (0.99999,0.99959) (1.00005,0.99971) (1.00009,0.99981) (1.00011,0.99988) (1.00012,0.99995) (1.00013,1.00000) (1.00012,1.00003) (1.00012,1.00005) (1.00011,1.00007) (1.00010,1.00008) (1.00008,1.00008)};
\draw[-{Latex},structurecolor] (axis cs:1.06850,1.02092) -- (axis cs:1.06417,1.03361);
\draw[-{Latex},structurecolor] (axis cs:0.99780,0.99817) -- (axis cs:0.99810,0.99811);
\addplot[structurecolor,thin,no marks] coordinates {(0.30263,2.50000) (0.29257,2.45880) (0.28279,2.41892) (0.27327,2.38029) (0.26401,2.34285) (0.25501,2.30656) (0.24624,2.27136) (0.23771,2.23720) (0.22941,2.20404) (0.22133,2.17184) (0.21347,2.14054) (0.20582,2.11013) (0.19837,2.08055) (0.19112,2.05178) (0.18406,2.02378) (0.17720,1.99652) (0.17051,1.96998) (0.16400,1.94413) (0.15767,1.91893) (0.15150,1.89437) (0.14550,1.87043) (0.13966,1.84707) (0.13398,1.82429) (0.12845,1.80205) (0.12307,1.78034) (0.11784,1.75914) (0.11274,1.73843) (0.10779,1.71820) (0.10297,1.69843) (0.09829,1.67911) (0.09374,1.66021) (0.08931,1.64173) (0.08501,1.62366) (0.08083,1.60597) (0.07677,1.58866) (0.07282,1.57172) (0.06899,1.55513) (0.06527,1.53889) (0.06166,1.52298) (0.05816,1.50739) (0.05476,1.49212) (0.05146,1.47715) (0.04827,1.46247) (0.04517,1.44809) (0.04217,1.43398) (0.03927,1.42014) (0.03646,1.40657) (0.03374,1.39325) (0.03111,1.38018) (0.02857,1.36735) (0.02611,1.35476) (0.02374,1.34240) (0.02145,1.33026) (0.01925,1.31833) (0.01712,1.30662) (0.01508,1.29512) (0.01311,1.28381) (0.01122,1.27270) (0.00940,1.26178) (0.00766,1.25105) (0.00598,1.24049) (0.00438,1.23011) (0.00285,1.21991) (0.00139,1.20987) (0.00000,1.20000)};
\draw[-{Latex},structurecolor] (axis cs:0.14550,1.87043) -- (axis cs:0.13966,1.84707);
\draw[-{Latex},structurecolor] (axis cs:0.04517,1.44809) -- (axis cs:0.04217,1.43398);
\addplot[structurecolor,thin,no marks] coordinates {(0.00000,1.20000) (-0.01560,1.01054) (-0.00596,0.87276) (0.02245,0.76804) (0.06535,0.68604) (0.11969,0.62087) (0.18312,0.56930) (0.25366,0.52974) (0.32948,0.50166) (0.40879,0.48512) (0.48977,0.48046) (0.57051,0.48800) (0.64913,0.50766) (0.72377,0.53880) (0.79273,0.58004) (0.85459,0.62921) (0.90833,0.68355) (0.95338,0.73999) (0.98963,0.79551) (1.01746,0.84751) (1.03757,0.89403) (1.05093,0.93388) (1.05861,0.96658) (1.06174,0.99225) (1.06137,1.01142) (1.05846,1.02489) (1.05384,1.03358) (1.04817,1.03841) (1.04199,1.04022) (1.03572,1.03980) (1.02964,1.03778) (1.02397,1.03470) (1.01883,1.03100) (1.01430,1.02700) (1.01039,1.02296) (1.00711,1.01906) (1.00442,1.01543) (1.00227,1.01214) (1.00061,1.00923) (0.99937,1.00672) (0.99849,1.00461) (0.99791,1.00288) (0.99757,1.00149) (0.99742,1.00042) (0.99742,0.99961) (0.99753,0.99904) (0.99772,0.99866) (0.99794,0.99844) (0.99820,0.99834) (0.99846,0.99834) (0.99871,0.99841) (0.99895,0.99852) (0.99917,0.99867) (0.99937,0.99883) (0.99954,0.99900) (0.99968,0.99917) (0.99980,0.99932) (0.99989,0.99946) (0.99997,0.99959) (1.00002,0.99970) (1.00006,0.99979) (1.00009,0.99987) (1.00010,0.99993) (1.00011,0.99998) (1.00011,1.00001)};
\draw[-{Latex},structurecolor] (axis cs:1.03757,0.89403) -- (axis cs:1.05093,0.93388);
\draw[-{Latex},structurecolor] (axis cs:0.99742,1.00042) -- (axis cs:0.99742,0.99961);
\addplot[structurecolor,thin,no marks] coordinates {(0.37294,2.50000) (0.31765,2.27752) (0.27050,2.09091) (0.23021,1.93218) (0.19580,1.79553) (0.16648,1.67670) (0.14160,1.57242) (0.12063,1.48022) (0.10314,1.39813) (0.08874,1.32459) (0.07714,1.25835) (0.06804,1.19838) (0.06123,1.14384) (0.05650,1.09404) (0.05367,1.04839) (0.05259,1.00639) (0.05312,0.96764) (0.05513,0.93178) (0.05853,0.89850) (0.06320,0.86754) (0.06907,0.83868) (0.07605,0.81172) (0.08407,0.78650) (0.09307,0.76287) (0.10298,0.74070) (0.11375,0.71989) (0.12532,0.70034) (0.13765,0.68197) (0.15069,0.66471) (0.16440,0.64850) (0.17873,0.63329) (0.19365,0.61903) (0.20913,0.60569) (0.22511,0.59323) (0.24158,0.58163) (0.25849,0.57088) (0.27582,0.56095) (0.29352,0.55183) (0.31158,0.54352) (0.32995,0.53600) (0.34860,0.52928) (0.36751,0.52334) (0.38663,0.51820) (0.40595,0.51384) (0.42543,0.51028) (0.44503,0.50751) (0.46472,0.50553) (0.48448,0.50435) (0.50427,0.50396) (0.52406,0.50437) (0.54382,0.50558) (0.56351,0.50757) (0.58311,0.51036) (0.60258,0.51392) (0.62190,0.51825) (0.64102,0.52335) (0.65993,0.52918) (0.67859,0.53575) (0.69698,0.54302) (0.71506,0.55098) (0.73280,0.55959) (0.75020,0.56883) (0.76721,0.57867) (0.78382,0.58907) (0.80000,0.60000)};
\draw[-{Latex},structurecolor] (axis cs:0.06907,0.83868) -- (axis cs:0.07605,0.81172);
\draw[-{Latex},structurecolor] (axis cs:0.40595,0.51384) -- (axis cs:0.42543,0.51028);
\addplot[structurecolor,thin,no marks] coordinates {(0.80000,0.60000) (0.85891,0.64703) (0.91006,0.69876) (0.95295,0.75229) (0.98752,0.80481) (1.01413,0.85392) (1.03344,0.89785) (1.04638,0.93551) (1.05395,0.96647) (1.05720,0.99085) (1.05712,1.00915) (1.05462,1.02210) (1.05045,1.03054) (1.04526,1.03533) (1.03956,1.03726) (1.03373,1.03704) (1.02806,1.03530) (1.02275,1.03253) (1.01793,1.02914) (1.01366,1.02545) (1.00998,1.02170) (1.00688,1.01806) (1.00432,1.01465) (1.00228,1.01156) (1.00069,1.00882) (0.99950,1.00646) (0.99866,1.00446) (0.99809,1.00282) (0.99776,1.00150) (0.99761,1.00047) (0.99760,0.99970) (0.99769,0.99915) (0.99786,0.99878) (0.99807,0.99856) (0.99830,0.99846) (0.99854,0.99845) (0.99878,0.99851) (0.99901,0.99861) (0.99921,0.99875) (0.99940,0.99890) (0.99956,0.99906) (0.99969,0.99921) (0.99980,0.99935) (0.99989,0.99949) (0.99996,0.99961) (1.00002,0.99971) (1.00005,0.99980) (1.00008,0.99987) (1.00010,0.99993) (1.00010,0.99998) (1.00010,1.00001) (1.00010,1.00003) (1.00009,1.00005) (1.00008,1.00006) (1.00007,1.00007) (1.00006,1.00007) (1.00005,1.00006) (1.00004,1.00006) (1.00003,1.00005) (1.00003,1.00005) (1.00002,1.00004) (1.00001,1.00003) (1.00001,1.00003) (1.00000,1.00002) (1.00000,1.00002)};
\draw[-{Latex},structurecolor] (axis cs:1.00998,1.02170) -- (axis cs:1.00688,1.01806);
\draw[-{Latex},structurecolor] (axis cs:0.99989,0.99949) -- (axis cs:0.99996,0.99961);
\addplot[structurecolor,thin,no marks] coordinates {(1.23710,2.50000) (1.22192,2.45251) (1.20723,2.40699) (1.19299,2.36332) (1.17919,2.32139) (1.16582,2.28109) (1.15285,2.24232) (1.14027,2.20501) (1.12807,2.16907) (1.11623,2.13442) (1.10474,2.10100) (1.09359,2.06874) (1.08277,2.03758) (1.07226,2.00746) (1.06205,1.97835) (1.05214,1.95017) (1.04251,1.92290) (1.03316,1.89648) (1.02408,1.87089) (1.01525,1.84607) (1.00668,1.82200) (0.99835,1.79864) (0.99026,1.77597) (0.98240,1.75395) (0.97476,1.73255) (0.96733,1.71176) (0.96012,1.69154) (0.95311,1.67188) (0.94630,1.65274) (0.93969,1.63412) (0.93327,1.61599) (0.92702,1.59833) (0.92096,1.58113) (0.91508,1.56437) (0.90936,1.54803) (0.90381,1.53210) (0.89842,1.51656) (0.89319,1.50140) (0.88811,1.48661) (0.88318,1.47218) (0.87840,1.45808) (0.87376,1.44433) (0.86926,1.43089) (0.86490,1.41776) (0.86067,1.40494) (0.85657,1.39241) (0.85260,1.38016) (0.84876,1.36819) (0.84503,1.35648) (0.84143,1.34504) (0.83794,1.33384) (0.83456,1.32289) (0.83130,1.31218) (0.82814,1.30170) (0.82510,1.29144) (0.82215,1.28140) (0.81931,1.27158) (0.81657,1.26196) (0.81392,1.25254) (0.81137,1.24332) (0.80892,1.23429) (0.80656,1.22545) (0.80428,1.21679) (0.80210,1.20831) (0.80000,1.20000)};
\draw[-{Latex},structurecolor] (axis cs:1.00668,1.82200) -- (axis cs:0.99835,1.79864);
\draw[-{Latex},structurecolor] (axis cs:0.86490,1.41776) -- (axis cs:0.86067,1.40494);
\addplot[structurecolor,thin,no marks] coordinates {(0.80000,1.20000) (0.77758,1.09243) (0.76965,1.01310) (0.77239,0.95494) (0.78288,0.91329) (0.79880,0.88489) (0.81829,0.86727) (0.83982,0.85849) (0.86214,0.85684) (0.88426,0.86086) (0.90540,0.86917) (0.92498,0.88056) (0.94262,0.89389) (0.95809,0.90822) (0.97129,0.92273) (0.98226,0.93678) (0.99110,0.94988) (0.99799,0.96170) (1.00315,0.97206) (1.00680,0.98088) (1.00920,0.98818) (1.01057,0.99404) (1.01113,0.99858) (1.01107,1.00197) (1.01057,1.00437) (1.00975,1.00593) (1.00875,1.00682) (1.00765,1.00719) (1.00652,1.00715) (1.00543,1.00683) (1.00440,1.00630) (1.00347,1.00565) (1.00264,1.00494) (1.00192,1.00422) (1.00132,1.00351) (1.00082,1.00285) (1.00042,1.00225) (1.00011,1.00171) (0.99988,1.00125) (0.99972,1.00086) (0.99961,1.00054) (0.99955,1.00028) (0.99952,1.00008) (0.99952,0.99993) (0.99954,0.99982) (0.99958,0.99975) (0.99962,0.99971) (0.99967,0.99969) (0.99971,0.99969) (0.99976,0.99970) (0.99981,0.99973) (0.99985,0.99975) (0.99988,0.99978) (0.99991,0.99981) (0.99994,0.99985) (0.99996,0.99987) (0.99998,0.99990) (0.99999,0.99992) (1.00000,0.99994) (1.00001,0.99996) (1.00002,0.99998) (1.00002,0.99999) (1.00002,1.00000) (1.00002,1.00000) (1.00002,1.00001)};
\draw[-{Latex},structurecolor] (axis cs:1.00920,0.98818) -- (axis cs:1.01057,0.99404);
\draw[-{Latex},structurecolor] (axis cs:0.99952,0.99993) -- (axis cs:0.99954,0.99982);
\addplot[structurecolor,thin,no marks] coordinates {(-1.12787,-0.50000) (-1.07745,-0.46661) (-1.02813,-0.43585) (-0.97981,-0.40763) (-0.93241,-0.38184) (-0.88584,-0.35840) (-0.84004,-0.33719) (-0.79491,-0.31812) (-0.75040,-0.30108) (-0.70644,-0.28597) (-0.66296,-0.27268) (-0.61990,-0.26111) (-0.57721,-0.25116) (-0.53482,-0.24271) (-0.49271,-0.23568) (-0.45080,-0.22995) (-0.40908,-0.22543) (-0.36748,-0.22200) (-0.32599,-0.21956) (-0.28456,-0.21801) (-0.24318,-0.21725) (-0.20180,-0.21717) (-0.16042,-0.21765) (-0.11902,-0.21860) (-0.07757,-0.21990) (-0.03608,-0.22144) (0.00546,-0.22311) (0.04707,-0.22479) (0.08873,-0.22635) (0.13044,-0.22769) (0.17219,-0.22868) (0.21396,-0.22919) (0.25574,-0.22910) (0.29751,-0.22827) (0.33924,-0.22658) (0.38089,-0.22389) (0.42243,-0.22008) (0.46382,-0.21501) (0.50501,-0.20855) (0.54596,-0.20059) (0.58661,-0.19099) (0.62691,-0.17963) (0.66679,-0.16643) (0.70618,-0.15126) (0.74503,-0.13405) (0.78326,-0.11472) (0.82079,-0.09322) (0.85755,-0.06951) (0.89347,-0.04355) (0.92847,-0.01537) (0.96247,0.01502) (0.99541,0.04757) (1.02720,0.08220) (1.05778,0.11881) (1.08709,0.15726) (1.11505,0.19741) (1.14163,0.23907) (1.16677,0.28206) (1.19043,0.32614) (1.21257,0.37109) (1.23318,0.41665) (1.25222,0.46258) (1.26971,0.50862) (1.28563,0.55451) (1.30000,0.60000)};
\draw[-{Latex},structurecolor] (axis cs:-0.24318,-0.21725) -- (axis cs:-0.20180,-0.21717);
\draw[-{Latex},structurecolor] (axis cs:0.70618,-0.15126) -- (axis cs:0.74503,-0.13405);
\addplot[structurecolor,thin,no marks] coordinates {(1.30000,0.60000) (1.34676,0.79750) (1.36495,0.96332) (1.36036,1.08838) (1.33946,1.17265) (1.30824,1.22180) (1.27157,1.24371) (1.23310,1.24611) (1.19533,1.23554) (1.15990,1.21700) (1.12774,1.19414) (1.09932,1.16952) (1.07477,1.14483) (1.05401,1.12121) (1.03680,1.09932) (1.02286,1.07956) (1.01182,1.06210) (1.00333,1.04697) (0.99702,1.03411) (0.99256,1.02338) (0.98962,1.01463) (0.98790,1.00765) (0.98714,1.00224) (0.98713,0.99819) (0.98765,0.99531) (0.98854,0.99339) (0.98967,0.99225) (0.99093,0.99174) (0.99223,0.99171) (0.99351,0.99203) (0.99471,0.99259) (0.99581,0.99331) (0.99679,0.99412) (0.99765,0.99496) (0.99837,0.99578) (0.99897,0.99656) (0.99945,0.99728) (0.99982,0.99791) (1.00010,0.99847) (1.00031,0.99894) (1.00044,0.99932) (1.00052,0.99964) (1.00056,0.99988) (1.00056,1.00006) (1.00054,1.00019) (1.00050,1.00028) (1.00045,1.00033) (1.00040,1.00036) (1.00034,1.00036) (1.00029,1.00035) (1.00023,1.00032) (1.00019,1.00029) (1.00014,1.00026) (1.00011,1.00022) (1.00007,1.00019) (1.00005,1.00015) (1.00003,1.00012) (1.00001,1.00009) (1.00000,1.00007) (0.99999,1.00005) (0.99998,1.00003) (0.99998,1.00002) (0.99998,1.00001) (0.99998,1.00000) (0.99998,0.99999)};
\draw[-{Latex},structurecolor] (axis cs:0.98962,1.01463) -- (axis cs:0.98790,1.00765);
\draw[-{Latex},structurecolor] (axis cs:1.00056,1.00006) -- (axis cs:1.00054,1.00019);
\addplot[structurecolor,thin,no marks] coordinates {(-1.14564,-0.50000) (-1.07718,-0.45324) (-1.01078,-0.41136) (-0.94622,-0.37410) (-0.88327,-0.34122) (-0.82176,-0.31246) (-0.76149,-0.28756) (-0.70229,-0.26626) (-0.64400,-0.24830) (-0.58647,-0.23342) (-0.52956,-0.22137) (-0.47316,-0.21187) (-0.41713,-0.20468) (-0.36140,-0.19954) (-0.30586,-0.19618) (-0.25043,-0.19434) (-0.19507,-0.19377) (-0.13970,-0.19421) (-0.08430,-0.19537) (-0.02883,-0.19700) (0.02672,-0.19880) (0.08235,-0.20050) (0.13805,-0.20180) (0.19380,-0.20241) (0.24955,-0.20201) (0.30526,-0.20031) (0.36085,-0.19699) (0.41624,-0.19173) (0.47134,-0.18423) (0.52603,-0.17418) (0.58020,-0.16128) (0.63369,-0.14528) (0.68636,-0.12592) (0.73806,-0.10300) (0.78860,-0.07634) (0.83782,-0.04583) (0.88554,-0.01142) (0.93157,0.02687) (0.97573,0.06895) (1.01786,0.11465) (1.05779,0.16372) (1.09537,0.21582) (1.13048,0.27053) (1.16299,0.32738) (1.19284,0.38585) (1.21994,0.44534) (1.24427,0.50528) (1.26582,0.56505) (1.28462,0.62408) (1.30071,0.68180) (1.31416,0.73772) (1.32507,0.79136) (1.33355,0.84236) (1.33973,0.89039) (1.34376,0.93523) (1.34579,0.97670) (1.34598,1.01471) (1.34448,1.04923) (1.34147,1.08028) (1.33709,1.10792) (1.33151,1.13228) (1.32487,1.15350) (1.31732,1.17173) (1.30899,1.18717) (1.30000,1.20000)};
\draw[-{Latex},structurecolor] (axis cs:0.02672,-0.19880) -- (axis cs:0.08235,-0.20050);
\draw[-{Latex},structurecolor] (axis cs:1.16299,0.32738) -- (axis cs:1.19284,0.38585);
\addplot[structurecolor,thin,no marks] coordinates {(1.30000,1.20000) (1.26636,1.22665) (1.23021,1.23326) (1.19418,1.22620) (1.15997,1.21048) (1.12865,1.18983) (1.10077,1.16689) (1.07653,1.14347) (1.05590,1.12076) (1.03871,1.09954) (1.02469,1.08023) (1.01353,1.06305) (1.00487,1.04809) (0.99839,1.03530) (0.99373,1.02459) (0.99060,1.01579) (0.98871,1.00873) (0.98779,1.00322) (0.98763,0.99907) (0.98803,0.99606) (0.98881,0.99402) (0.98985,0.99278) (0.99104,0.99216) (0.99228,0.99203) (0.99351,0.99227) (0.99468,0.99277) (0.99576,0.99343) (0.99673,0.99419) (0.99757,0.99499) (0.99829,0.99579) (0.99889,0.99654) (0.99937,0.99724) (0.99976,0.99787) (1.00004,0.99842) (1.00025,0.99889) (1.00040,0.99927) (1.00048,0.99959) (1.00053,0.99984) (1.00054,1.00002) (1.00052,1.00016) (1.00049,1.00025) (1.00045,1.00031) (1.00039,1.00034) (1.00034,1.00035) (1.00029,1.00034) (1.00024,1.00032) (1.00019,1.00029) (1.00015,1.00026) (1.00011,1.00022) (1.00008,1.00019) (1.00005,1.00015) (1.00003,1.00012) (1.00001,1.00009) (1.00000,1.00007) (0.99999,1.00005) (0.99998,1.00003) (0.99998,1.00002) (0.99998,1.00001) (0.99998,1.00000) (0.99998,0.99999) (0.99998,0.99999) (0.99998,0.99999) (0.99998,0.99999) (0.99999,0.99998) (0.99999,0.99999)};
\draw[-{Latex},structurecolor] (axis cs:0.98881,0.99402) -- (axis cs:0.98985,0.99278);
\addplot[gray,dashed,domain=-2:2.5,samples=60] {1};
\addplot[gray,dashed,domain=-2:2.5,samples=60] {x};
\addplot[gray,dashed,domain=-2:2.5,samples=60] {-x};
\addplot[only marks,mark=*,mark size=1.4pt] coordinates {(-1,1) (1,1)};
\end{axis}\end{tikzpicture}
\end{center}
\end{solution}
\begin{exercise}\label{cw:bank-5C-5}原题5C-5. 对 $x'=x-x^2-xy$, $y'=3y-xy-2y^2$ 重做5C-4的全套平衡点及相图程序.\end{exercise}
\begin{solution}（官方译审并补证）$x'=x(1-x-y)$, $y'=y(3-x-2y)$，得平衡点 $(0,0),(1,0),(0,3/2),(-1,2)$. Jacobian为
\[J=\begin{pmatrix}1-2x-y&-x\\-y&3-x-4y\end{pmatrix}.\]
在原点根 $1,3$，为不稳定结点；在 $(1,0)$ 根 $-1,2$，方向 $(1,0)^T,(-1,3)^T$，为鞍点. 在 $(0,3/2)$ 根 $-1/2,-3$，方向 $(5,-3)^T,(0,1)^T$，为渐近稳定结点. 在 $(-1,2)$，$J=\begin{pmatrix}1&1\\-2&-4\end{pmatrix}$，特征多项式 $\lambda^2+3\lambda-2$，根 $(-3\pm\sqrt{17})/2$，为鞍点，方向 $(1,\lambda-1)^T$. 正象限中 $x$ 在 $x+y<1$ 时增加，$y$ 在 $x+2y<3$ 时增加，超过各直线后相应减少；坐标轴是不变轨道. 下图依这些零流线、方向符号及原方程轨道作出，第四点虽在物理非负区域外，数学上仍必须登记.
\begin{center}
\begin{tikzpicture}\begin{axis}[width=.47\linewidth,height=5.2cm,axis lines=middle,xlabel={$x$},ylabel={$y$},xmin=-2,xmax=2,ymin=-0.3,ymax=3,title={四平衡点及竞争轨道},tick label style={font=\footnotesize},clip=true]
\addplot[structurecolor,thin,no marks] coordinates {(-0.00003,0.00000) (-0.00003,0.00000) (-0.00004,0.00000) (-0.00004,0.00000) (-0.00005,0.00000) (-0.00006,0.00000) (-0.00007,0.00000) (-0.00008,0.00000) (-0.00010,0.00000) (-0.00011,0.00000) (-0.00013,0.00000) (-0.00015,0.00000) (-0.00018,0.00000) (-0.00021,0.00000) (-0.00024,0.00000) (-0.00028,0.00000) (-0.00033,0.00000) (-0.00039,0.00000) (-0.00045,0.00000) (-0.00053,0.00000) (-0.00062,0.00000) (-0.00073,0.00000) (-0.00085,0.00000) (-0.00099,0.00000) (-0.00116,0.00000) (-0.00136,0.00000) (-0.00159,0.00000) (-0.00186,0.00000) (-0.00217,0.00000) (-0.00254,0.00000) (-0.00297,0.00000) (-0.00347,0.00000) (-0.00406,0.00000) (-0.00475,0.00000) (-0.00556,0.00000) (-0.00651,0.00000) (-0.00762,0.00000) (-0.00892,0.00000) (-0.01044,0.00000) (-0.01223,0.00000) (-0.01432,0.00000) (-0.01679,0.00000) (-0.01968,0.00000) (-0.02309,0.00000) (-0.02710,0.00001) (-0.03183,0.00001) (-0.03741,0.00002) (-0.04402,0.00003) (-0.05185,0.00005) (-0.06115,0.00008) (-0.07224,0.00014) (-0.08550,0.00022) (-0.10142,0.00036) (-0.12063,0.00058) (-0.14395,0.00095) (-0.17247,0.00156) (-0.20762,0.00257) (-0.25144,0.00425) (-0.30674,0.00708) (-0.37768,0.01189) (-0.47049,0.02020) (-0.59501,0.03477) (-0.76758,0.06086) (-1.01717,0.10882) (-1.40000,0.20000)};
\draw[-{Latex},structurecolor] (axis cs:-0.00062,0.00000) -- (axis cs:-0.00073,0.00000);
\draw[-{Latex},structurecolor] (axis cs:-0.02309,0.00000) -- (axis cs:-0.02710,0.00001);
\addplot[structurecolor,thin,no marks] coordinates {(-1.40000,0.20000) (-1.40718,0.20187) (-1.41442,0.20376) (-1.42172,0.20566) (-1.42907,0.20759) (-1.43647,0.20954) (-1.44394,0.21150) (-1.45146,0.21349) (-1.45903,0.21550) (-1.46667,0.21752) (-1.47437,0.21957) (-1.48212,0.22164) (-1.48994,0.22373) (-1.49781,0.22585) (-1.50575,0.22798) (-1.51376,0.23014) (-1.52182,0.23232) (-1.52995,0.23452) (-1.53814,0.23674) (-1.54640,0.23899) (-1.55473,0.24126) (-1.56312,0.24356) (-1.57158,0.24588) (-1.58011,0.24822) (-1.58871,0.25059) (-1.59737,0.25298) (-1.60611,0.25540) (-1.61492,0.25784) (-1.62380,0.26031) (-1.63276,0.26281) (-1.64179,0.26533) (-1.65089,0.26788) (-1.66007,0.27045) (-1.66933,0.27305) (-1.67867,0.27568) (-1.68808,0.27834) (-1.69757,0.28103) (-1.70715,0.28374) (-1.71680,0.28649) (-1.72654,0.28926) (-1.73636,0.29206) (-1.74627,0.29490) (-1.75626,0.29776) (-1.76633,0.30066) (-1.77650,0.30358) (-1.78675,0.30654) (-1.79710,0.30953) (-1.80753,0.31255) (-1.81806,0.31560) (-1.82868,0.31869) (-1.83939,0.32181) (-1.85020,0.32497) (-1.86111,0.32816) (-1.87211,0.33138) (-1.88321,0.33464) (-1.89442,0.33793) (-1.90572,0.34127) (-1.91713,0.34463) (-1.92865,0.34804) (-1.94027,0.35148) (-1.95199,0.35496) (-1.96383,0.35848) (-1.97577,0.36204) (-1.98783,0.36563) (-2.00000,0.36927)};
\draw[-{Latex},structurecolor] (axis cs:-1.55473,0.24126) -- (axis cs:-1.56312,0.24356);
\draw[-{Latex},structurecolor] (axis cs:-1.76633,0.30066) -- (axis cs:-1.77650,0.30358);
\addplot[structurecolor,thin,no marks] coordinates {(-0.00003,0.00000) (-0.00004,0.00000) (-0.00004,0.00000) (-0.00005,0.00000) (-0.00006,0.00000) (-0.00007,0.00000) (-0.00008,0.00000) (-0.00010,0.00000) (-0.00011,0.00000) (-0.00013,0.00000) (-0.00016,0.00000) (-0.00018,0.00000) (-0.00021,0.00000) (-0.00025,0.00000) (-0.00029,0.00000) (-0.00034,0.00000) (-0.00040,0.00000) (-0.00046,0.00000) (-0.00054,0.00000) (-0.00063,0.00000) (-0.00074,0.00000) (-0.00087,0.00000) (-0.00101,0.00000) (-0.00119,0.00000) (-0.00139,0.00000) (-0.00162,0.00000) (-0.00190,0.00000) (-0.00222,0.00000) (-0.00260,0.00000) (-0.00304,0.00000) (-0.00355,0.00000) (-0.00415,0.00000) (-0.00486,0.00000) (-0.00569,0.00000) (-0.00665,0.00000) (-0.00779,0.00000) (-0.00912,0.00000) (-0.01068,0.00000) (-0.01250,0.00000) (-0.01465,0.00001) (-0.01717,0.00001) (-0.02013,0.00002) (-0.02361,0.00003) (-0.02772,0.00005) (-0.03256,0.00007) (-0.03828,0.00012) (-0.04504,0.00019) (-0.05306,0.00030) (-0.06259,0.00049) (-0.07395,0.00079) (-0.08754,0.00128) (-0.10385,0.00208) (-0.12353,0.00338) (-0.14741,0.00552) (-0.17654,0.00902) (-0.21236,0.01480) (-0.25674,0.02439) (-0.31220,0.04033) (-0.38211,0.06692) (-0.47092,0.11121) (-0.58429,0.18437) (-0.72883,0.30283) (-0.91103,0.48691) (-1.13495,0.75338) (-1.40000,1.10000)};
\draw[-{Latex},structurecolor] (axis cs:-0.00074,0.00000) -- (axis cs:-0.00087,0.00000);
\draw[-{Latex},structurecolor] (axis cs:-0.02772,0.00005) -- (axis cs:-0.03256,0.00007);
\addplot[structurecolor,thin,no marks] coordinates {(-1.40000,1.10000) (-1.40833,1.11107) (-1.41669,1.12219) (-1.42508,1.13334) (-1.43350,1.14453) (-1.44196,1.15576) (-1.45045,1.16703) (-1.45896,1.17832) (-1.46751,1.18965) (-1.47609,1.20102) (-1.48471,1.21241) (-1.49335,1.22383) (-1.50203,1.23528) (-1.51073,1.24675) (-1.51947,1.25825) (-1.52824,1.26977) (-1.53704,1.28132) (-1.54587,1.29288) (-1.55473,1.30447) (-1.56362,1.31607) (-1.57255,1.32769) (-1.58151,1.33932) (-1.59050,1.35097) (-1.59952,1.36263) (-1.60857,1.37430) (-1.61765,1.38598) (-1.62677,1.39766) (-1.63592,1.40936) (-1.64510,1.42106) (-1.65431,1.43276) (-1.66356,1.44447) (-1.67284,1.45618) (-1.68215,1.46789) (-1.69150,1.47959) (-1.70088,1.49130) (-1.71029,1.50300) (-1.71974,1.51470) (-1.72922,1.52639) (-1.73874,1.53807) (-1.74829,1.54975) (-1.75788,1.56141) (-1.76750,1.57307) (-1.77716,1.58471) (-1.78686,1.59634) (-1.79660,1.60796) (-1.80637,1.61956) (-1.81618,1.63115) (-1.82603,1.64271) (-1.83592,1.65426) (-1.84585,1.66579) (-1.85582,1.67731) (-1.86583,1.68880) (-1.87588,1.70026) (-1.88597,1.71171) (-1.89611,1.72313) (-1.90629,1.73453) (-1.91651,1.74590) (-1.92678,1.75725) (-1.93710,1.76857) (-1.94746,1.77987) (-1.95787,1.79114) (-1.96833,1.80238) (-1.97884,1.81359) (-1.98939,1.82477) (-2.00000,1.83592)};
\draw[-{Latex},structurecolor] (axis cs:-1.57255,1.32769) -- (axis cs:-1.58151,1.33932);
\draw[-{Latex},structurecolor] (axis cs:-1.78686,1.59634) -- (axis cs:-1.79660,1.60796);
\addplot[structurecolor,thin,no marks] coordinates {(-0.00005,0.00000) (-0.00006,0.00000) (-0.00007,0.00000) (-0.00008,0.00000) (-0.00009,0.00000) (-0.00010,0.00000) (-0.00012,0.00000) (-0.00014,0.00000) (-0.00017,0.00000) (-0.00019,0.00000) (-0.00023,0.00000) (-0.00027,0.00000) (-0.00031,0.00000) (-0.00036,0.00000) (-0.00042,0.00000) (-0.00050,0.00000) (-0.00058,0.00000) (-0.00068,0.00000) (-0.00079,0.00000) (-0.00093,0.00000) (-0.00109,0.00000) (-0.00127,0.00000) (-0.00148,0.00000) (-0.00174,0.00000) (-0.00203,0.00000) (-0.00237,0.00000) (-0.00278,0.00000) (-0.00325,0.00000) (-0.00380,0.00000) (-0.00444,0.00000) (-0.00520,0.00000) (-0.00608,0.00000) (-0.00712,0.00000) (-0.00834,0.00000) (-0.00976,0.00000) (-0.01143,0.00001) (-0.01339,0.00001) (-0.01569,0.00001) (-0.01839,0.00002) (-0.02156,0.00004) (-0.02530,0.00006) (-0.02971,0.00010) (-0.03491,0.00016) (-0.04105,0.00025) (-0.04833,0.00041) (-0.05696,0.00066) (-0.06723,0.00106) (-0.07949,0.00172) (-0.09417,0.00278) (-0.11181,0.00451) (-0.13312,0.00733) (-0.15898,0.01194) (-0.19054,0.01951) (-0.22926,0.03197) (-0.27701,0.05245) (-0.33611,0.08607) (-0.40926,0.14078) (-0.49921,0.22817) (-0.60793,0.36292) (-0.73500,0.55855) (-0.87575,0.81754) (-1.02094,1.12009) (-1.15994,1.42536) (-1.28615,1.69208) (-1.40000,1.90000)};
\draw[-{Latex},structurecolor] (axis cs:-0.00109,0.00000) -- (axis cs:-0.00127,0.00000);
\draw[-{Latex},structurecolor] (axis cs:-0.04105,0.00025) -- (axis cs:-0.04833,0.00041);
\addplot[structurecolor,thin,no marks] coordinates {(-1.40000,1.90000) (-1.40791,1.91277) (-1.41580,1.92524) (-1.42365,1.93743) (-1.43148,1.94934) (-1.43929,1.96097) (-1.44708,1.97233) (-1.45486,1.98342) (-1.46263,1.99426) (-1.47038,2.00485) (-1.47814,2.01519) (-1.48589,2.02530) (-1.49365,2.03517) (-1.50141,2.04482) (-1.50918,2.05425) (-1.51697,2.06347) (-1.52477,2.07249) (-1.53259,2.08131) (-1.54044,2.08994) (-1.54832,2.09839) (-1.55623,2.10666) (-1.56417,2.11476) (-1.57216,2.12269) (-1.58019,2.13048) (-1.58826,2.13811) (-1.59639,2.14560) (-1.60457,2.15295) (-1.61281,2.16017) (-1.62111,2.16727) (-1.62948,2.17424) (-1.63792,2.18111) (-1.64644,2.18787) (-1.65503,2.19453) (-1.66371,2.20110) (-1.67247,2.20758) (-1.68133,2.21397) (-1.69028,2.22029) (-1.69933,2.22654) (-1.70848,2.23272) (-1.71775,2.23883) (-1.72713,2.24490) (-1.73662,2.25091) (-1.74624,2.25687) (-1.75599,2.26280) (-1.76586,2.26869) (-1.77588,2.27455) (-1.78604,2.28038) (-1.79635,2.28619) (-1.80681,2.29198) (-1.81743,2.29776) (-1.82821,2.30354) (-1.83916,2.30931) (-1.85029,2.31508) (-1.86160,2.32086) (-1.87310,2.32664) (-1.88480,2.33245) (-1.89670,2.33827) (-1.90880,2.34411) (-1.92112,2.34998) (-1.93367,2.35589) (-1.94644,2.36183) (-1.95945,2.36781) (-1.97271,2.37384) (-1.98622,2.37992) (-2.00000,2.38605)};
\draw[-{Latex},structurecolor] (axis cs:-1.55623,2.10666) -- (axis cs:-1.56417,2.11476);
\draw[-{Latex},structurecolor] (axis cs:-1.75599,2.26280) -- (axis cs:-1.76586,2.26869);
\addplot[structurecolor,thin,no marks] coordinates {(-0.00002,0.00000) (-0.00002,0.00000) (-0.00003,0.00000) (-0.00003,0.00000) (-0.00004,0.00000) (-0.00004,0.00000) (-0.00005,0.00000) (-0.00006,0.00000) (-0.00007,0.00000) (-0.00008,0.00000) (-0.00009,0.00000) (-0.00011,0.00000) (-0.00013,0.00000) (-0.00015,0.00000) (-0.00017,0.00000) (-0.00020,0.00000) (-0.00024,0.00000) (-0.00028,0.00000) (-0.00032,0.00000) (-0.00038,0.00000) (-0.00044,0.00000) (-0.00052,0.00000) (-0.00061,0.00000) (-0.00071,0.00000) (-0.00083,0.00000) (-0.00097,0.00000) (-0.00113,0.00000) (-0.00133,0.00000) (-0.00155,0.00000) (-0.00181,0.00000) (-0.00212,0.00000) (-0.00248,0.00000) (-0.00290,0.00000) (-0.00339,0.00000) (-0.00397,0.00000) (-0.00464,0.00000) (-0.00543,0.00000) (-0.00636,0.00000) (-0.00744,0.00000) (-0.00871,0.00000) (-0.01020,0.00000) (-0.01194,0.00000) (-0.01399,0.00000) (-0.01639,0.00001) (-0.01922,0.00001) (-0.02254,0.00002) (-0.02645,0.00003) (-0.03107,0.00005) (-0.03651,0.00007) (-0.04295,0.00012) (-0.05058,0.00019) (-0.05964,0.00031) (-0.07043,0.00050) (-0.08333,0.00081) (-0.09880,0.00132) (-0.11744,0.00214) (-0.14002,0.00349) (-0.16756,0.00570) (-0.20137,0.00935) (-0.24326,0.01541) (-0.29565,0.02551) (-0.36191,0.04247) (-0.44669,0.07103) (-0.55645,0.11919) (-0.70000,0.20000)};
\draw[-{Latex},structurecolor] (axis cs:-0.01639,0.00001) -- (axis cs:-0.01922,0.00001);
\addplot[structurecolor,thin,no marks] coordinates {(-0.70000,0.20000) (-0.71102,0.20698) (-0.72224,0.21419) (-0.73367,0.22165) (-0.74531,0.22937) (-0.75717,0.23733) (-0.76925,0.24557) (-0.78155,0.25408) (-0.79408,0.26287) (-0.80684,0.27195) (-0.81984,0.28133) (-0.83307,0.29101) (-0.84655,0.30101) (-0.86028,0.31132) (-0.87427,0.32197) (-0.88851,0.33296) (-0.90301,0.34429) (-0.91779,0.35598) (-0.93283,0.36803) (-0.94814,0.38045) (-0.96374,0.39325) (-0.97962,0.40644) (-0.99580,0.42002) (-1.01226,0.43400) (-1.02902,0.44840) (-1.04609,0.46322) (-1.06346,0.47846) (-1.08115,0.49413) (-1.09915,0.51024) (-1.11747,0.52680) (-1.13612,0.54381) (-1.15509,0.56128) (-1.17441,0.57922) (-1.19406,0.59762) (-1.21405,0.61649) (-1.23440,0.63584) (-1.25509,0.65567) (-1.27615,0.67598) (-1.29756,0.69678) (-1.31935,0.71806) (-1.34150,0.73984) (-1.36403,0.76210) (-1.38695,0.78485) (-1.41025,0.80808) (-1.43394,0.83180) (-1.45803,0.85601) (-1.48252,0.88069) (-1.50742,0.90585) (-1.53273,0.93147) (-1.55846,0.95756) (-1.58462,0.98411) (-1.61121,1.01111) (-1.63824,1.03855) (-1.66572,1.06643) (-1.69365,1.09472) (-1.72204,1.12344) (-1.75090,1.15255) (-1.78025,1.18206) (-1.81008,1.21194) (-1.84041,1.24220) (-1.87126,1.27281) (-1.90262,1.30376) (-1.93453,1.33503) (-1.96698,1.36662) (-2.00000,1.39852)};
\draw[-{Latex},structurecolor] (axis cs:-0.96374,0.39325) -- (axis cs:-0.97962,0.40644);
\draw[-{Latex},structurecolor] (axis cs:-1.41025,0.80808) -- (axis cs:-1.43394,0.83180);
\addplot[structurecolor,thin,no marks] coordinates {(-0.00003,0.00000) (-0.00003,0.00000) (-0.00004,0.00000) (-0.00004,0.00000) (-0.00005,0.00000) (-0.00006,0.00000) (-0.00007,0.00000) (-0.00008,0.00000) (-0.00009,0.00000) (-0.00010,0.00000) (-0.00012,0.00000) (-0.00014,0.00000) (-0.00017,0.00000) (-0.00020,0.00000) (-0.00023,0.00000) (-0.00027,0.00000) (-0.00031,0.00000) (-0.00037,0.00000) (-0.00043,0.00000) (-0.00050,0.00000) (-0.00058,0.00000) (-0.00068,0.00000) (-0.00080,0.00000) (-0.00093,0.00000) (-0.00109,0.00000) (-0.00128,0.00000) (-0.00149,0.00000) (-0.00175,0.00000) (-0.00204,0.00000) (-0.00239,0.00000) (-0.00279,0.00000) (-0.00327,0.00000) (-0.00382,0.00000) (-0.00447,0.00000) (-0.00523,0.00000) (-0.00612,0.00000) (-0.00717,0.00000) (-0.00839,0.00000) (-0.00982,0.00001) (-0.01150,0.00001) (-0.01347,0.00002) (-0.01579,0.00003) (-0.01851,0.00005) (-0.02170,0.00008) (-0.02547,0.00012) (-0.02990,0.00020) (-0.03514,0.00032) (-0.04132,0.00052) (-0.04865,0.00083) (-0.05733,0.00134) (-0.06767,0.00216) (-0.07999,0.00349) (-0.09474,0.00565) (-0.11244,0.00915) (-0.13375,0.01485) (-0.15951,0.02414) (-0.19071,0.03926) (-0.22855,0.06380) (-0.27434,0.10334) (-0.32931,0.16607) (-0.39411,0.26285) (-0.46806,0.40526) (-0.54803,0.59998) (-0.62796,0.84008) (-0.70000,1.10000)};
\draw[-{Latex},structurecolor] (axis cs:-0.02170,0.00008) -- (axis cs:-0.02547,0.00012);
\addplot[structurecolor,thin,no marks] coordinates {(-0.70000,1.10000) (-0.75753,1.34372) (-0.79780,1.54270) (-0.82207,1.68704) (-0.83377,1.78246) (-0.83657,1.84109) (-0.83341,1.87482) (-0.82632,1.89265) (-0.81661,1.90067) (-0.80507,1.90270) (-0.79217,1.90106) (-0.77818,1.89713) (-0.76328,1.89173) (-0.74757,1.88533) (-0.73112,1.87820) (-0.71399,1.87053) (-0.69621,1.86241) (-0.67784,1.85390) (-0.65893,1.84506) (-0.63954,1.83593) (-0.61970,1.82653) (-0.59950,1.81690) (-0.57900,1.80707) (-0.55826,1.79707) (-0.53736,1.78693) (-0.51638,1.77670) (-0.49538,1.76639) (-0.47445,1.75606) (-0.45365,1.74573) (-0.43306,1.73544) (-0.41275,1.72523) (-0.39278,1.71512) (-0.37320,1.70516) (-0.35408,1.69536) (-0.33546,1.68576) (-0.31737,1.67637) (-0.29986,1.66723) (-0.28296,1.65835) (-0.26667,1.64975) (-0.25104,1.64143) (-0.23605,1.63342) (-0.22172,1.62572) (-0.20806,1.61833) (-0.19505,1.61126) (-0.18269,1.60450) (-0.17096,1.59806) (-0.15986,1.59193) (-0.14937,1.58612) (-0.13947,1.58060) (-0.13014,1.57538) (-0.12136,1.57045) (-0.11310,1.56579) (-0.10535,1.56140) (-0.09808,1.55727) (-0.09127,1.55339) (-0.08489,1.54974) (-0.07893,1.54632) (-0.07336,1.54311) (-0.06815,1.54011) (-0.06330,1.53730) (-0.05877,1.53468) (-0.05455,1.53222) (-0.05062,1.52993) (-0.04696,1.52780) (-0.04356,1.52581)};
\draw[-{Latex},structurecolor] (axis cs:-0.61970,1.82653) -- (axis cs:-0.59950,1.81690);
\draw[-{Latex},structurecolor] (axis cs:-0.19505,1.61126) -- (axis cs:-0.18269,1.60450);
\addplot[structurecolor,thin,no marks] coordinates {(-0.96162,3.00000) (-0.94902,2.92358) (-0.93732,2.85410) (-0.92644,2.79072) (-0.91629,2.73271) (-0.90679,2.67947) (-0.89789,2.63047) (-0.88952,2.58526) (-0.88164,2.54346) (-0.87420,2.50473) (-0.86716,2.46877) (-0.86050,2.43532) (-0.85418,2.40416) (-0.84817,2.37508) (-0.84244,2.34790) (-0.83699,2.32246) (-0.83178,2.29863) (-0.82679,2.27626) (-0.82202,2.25525) (-0.81745,2.23549) (-0.81305,2.21689) (-0.80883,2.19936) (-0.80476,2.18283) (-0.80084,2.16722) (-0.79706,2.15247) (-0.79341,2.13852) (-0.78987,2.12532) (-0.78645,2.11281) (-0.78314,2.10096) (-0.77992,2.08972) (-0.77680,2.07904) (-0.77376,2.06891) (-0.77080,2.05927) (-0.76793,2.05011) (-0.76512,2.04139) (-0.76238,2.03309) (-0.75970,2.02518) (-0.75709,2.01764) (-0.75453,2.01045) (-0.75202,2.00358) (-0.74956,1.99703) (-0.74715,1.99077) (-0.74479,1.98478) (-0.74246,1.97906) (-0.74018,1.97358) (-0.73793,1.96834) (-0.73571,1.96331) (-0.73353,1.95850) (-0.73138,1.95388) (-0.72926,1.94945) (-0.72717,1.94520) (-0.72510,1.94112) (-0.72306,1.93720) (-0.72104,1.93343) (-0.71904,1.92980) (-0.71707,1.92631) (-0.71511,1.92295) (-0.71317,1.91971) (-0.71124,1.91659) (-0.70933,1.91358) (-0.70744,1.91067) (-0.70556,1.90786) (-0.70370,1.90515) (-0.70184,1.90253) (-0.70000,1.90000)};
\draw[-{Latex},structurecolor] (axis cs:-0.81305,2.21689) -- (axis cs:-0.80883,2.19936);
\draw[-{Latex},structurecolor] (axis cs:-0.74246,1.97906) -- (axis cs:-0.74018,1.97358);
\addplot[structurecolor,thin,no marks] coordinates {(-0.70000,1.90000) (-0.67872,1.87511) (-0.65813,1.85682) (-0.63775,1.84219) (-0.61733,1.82958) (-0.59678,1.81805) (-0.57607,1.80711) (-0.55520,1.79644) (-0.53423,1.78591) (-0.51320,1.77544) (-0.49219,1.76500) (-0.47127,1.75459) (-0.45049,1.74422) (-0.42994,1.73391) (-0.40967,1.72370) (-0.38976,1.71360) (-0.37025,1.70365) (-0.35120,1.69388) (-0.33265,1.68431) (-0.31465,1.67496) (-0.29723,1.66586) (-0.28042,1.65702) (-0.26424,1.64846) (-0.24870,1.64019) (-0.23381,1.63222) (-0.21959,1.62457) (-0.20602,1.61723) (-0.19311,1.61020) (-0.18085,1.60349) (-0.16922,1.59710) (-0.15822,1.59102) (-0.14782,1.58525) (-0.13801,1.57978) (-0.12876,1.57461) (-0.12006,1.56972) (-0.11188,1.56510) (-0.10420,1.56075) (-0.09701,1.55666) (-0.09026,1.55281) (-0.08395,1.54920) (-0.07805,1.54581) (-0.07254,1.54264) (-0.06739,1.53967) (-0.06258,1.53689) (-0.05810,1.53429) (-0.05393,1.53186) (-0.05004,1.52960) (-0.04643,1.52748) (-0.04306,1.52552) (-0.03993,1.52368) (-0.03702,1.52197) (-0.03431,1.52038) (-0.03180,1.51890) (-0.02947,1.51753) (-0.02730,1.51625) (-0.02529,1.51506) (-0.02343,1.51396) (-0.02170,1.51294) (-0.02009,1.51199) (-0.01861,1.51110) (-0.01723,1.51028) (-0.01595,1.50952) (-0.01477,1.50882) (-0.01367,1.50817) (-0.01265,1.50756)};
\draw[-{Latex},structurecolor] (axis cs:-0.29723,1.66586) -- (axis cs:-0.28042,1.65702);
\draw[-{Latex},structurecolor] (axis cs:-0.06258,1.53689) -- (axis cs:-0.05810,1.53429);
\addplot[structurecolor,thin,no marks] coordinates {(0.00001,0.00000) (0.00001,0.00000) (0.00001,0.00000) (0.00001,0.00000) (0.00002,0.00000) (0.00002,0.00000) (0.00002,0.00000) (0.00003,0.00000) (0.00003,0.00000) (0.00004,0.00000) (0.00004,0.00000) (0.00005,0.00000) (0.00006,0.00000) (0.00007,0.00000) (0.00008,0.00000) (0.00009,0.00000) (0.00011,0.00000) (0.00012,0.00000) (0.00015,0.00000) (0.00017,0.00000) (0.00020,0.00000) (0.00023,0.00000) (0.00027,0.00000) (0.00032,0.00000) (0.00037,0.00000) (0.00043,0.00000) (0.00051,0.00000) (0.00059,0.00000) (0.00069,0.00000) (0.00081,0.00000) (0.00095,0.00000) (0.00111,0.00000) (0.00129,0.00000) (0.00151,0.00000) (0.00177,0.00000) (0.00206,0.00000) (0.00241,0.00000) (0.00282,0.00000) (0.00329,0.00000) (0.00385,0.00000) (0.00450,0.00000) (0.00525,0.00001) (0.00614,0.00001) (0.00717,0.00001) (0.00837,0.00002) (0.00977,0.00004) (0.01140,0.00006) (0.01331,0.00009) (0.01552,0.00015) (0.01810,0.00024) (0.02109,0.00038) (0.02457,0.00061) (0.02861,0.00097) (0.03328,0.00154) (0.03867,0.00244) (0.04490,0.00387) (0.05206,0.00613) (0.06026,0.00968) (0.06960,0.01526) (0.08019,0.02395) (0.09206,0.03741) (0.10521,0.05801) (0.11950,0.08905) (0.13462,0.13477) (0.15000,0.20000)};
\draw[-{Latex},structurecolor] (axis cs:0.00717,0.00001) -- (axis cs:0.00837,0.00002);
\addplot[structurecolor,thin,no marks] coordinates {(0.15000,0.20000) (0.16475,0.28908) (0.17769,0.40386) (0.18757,0.54134) (0.19331,0.69255) (0.19442,0.84428) (0.19115,0.98339) (0.18434,1.10109) (0.17509,1.19440) (0.16445,1.26496) (0.15325,1.31672) (0.14209,1.35416) (0.13131,1.38120) (0.12110,1.40092) (0.11157,1.41556) (0.10273,1.42670) (0.09456,1.43541) (0.08705,1.44242) (0.08015,1.44821) (0.07380,1.45311) (0.06798,1.45734) (0.06263,1.46105) (0.05772,1.46434) (0.05320,1.46729) (0.04905,1.46996) (0.04524,1.47237) (0.04173,1.47457) (0.03850,1.47659) (0.03552,1.47843) (0.03278,1.48011) (0.03026,1.48166) (0.02794,1.48309) (0.02580,1.48440) (0.02382,1.48560) (0.02200,1.48671) (0.02032,1.48773) (0.01877,1.48867) (0.01734,1.48954) (0.01602,1.49034) (0.01480,1.49108) (0.01368,1.49176) (0.01264,1.49239) (0.01168,1.49297) (0.01079,1.49350) (0.00998,1.49400) (0.00922,1.49445) (0.00852,1.49487) (0.00788,1.49526) (0.00728,1.49562) (0.00673,1.49595) (0.00623,1.49626) (0.00576,1.49654) (0.00532,1.49680) (0.00492,1.49704) (0.00455,1.49727) (0.00421,1.49747) (0.00389,1.49766) (0.00360,1.49784) (0.00332,1.49800) (0.00307,1.49815) (0.00284,1.49829) (0.00263,1.49842) (0.00243,1.49854) (0.00225,1.49865) (0.00208,1.49875)};
\draw[-{Latex},structurecolor] (axis cs:0.06798,1.45734) -- (axis cs:0.06263,1.46105);
\draw[-{Latex},structurecolor] (axis cs:0.01079,1.49350) -- (axis cs:0.00998,1.49400);
\addplot[structurecolor,thin,no marks] coordinates {(0.00002,0.00000) (0.00002,0.00000) (0.00003,0.00000) (0.00003,0.00000) (0.00004,0.00000) (0.00004,0.00000) (0.00005,0.00000) (0.00006,0.00000) (0.00007,0.00000) (0.00008,0.00000) (0.00009,0.00000) (0.00011,0.00000) (0.00012,0.00000) (0.00014,0.00000) (0.00017,0.00000) (0.00020,0.00000) (0.00023,0.00000) (0.00027,0.00000) (0.00032,0.00000) (0.00037,0.00000) (0.00043,0.00000) (0.00051,0.00000) (0.00059,0.00000) (0.00069,0.00000) (0.00081,0.00000) (0.00094,0.00000) (0.00110,0.00000) (0.00129,0.00000) (0.00151,0.00000) (0.00176,0.00000) (0.00206,0.00000) (0.00240,0.00000) (0.00281,0.00000) (0.00328,0.00000) (0.00384,0.00000) (0.00448,0.00001) (0.00524,0.00001) (0.00612,0.00002) (0.00714,0.00003) (0.00834,0.00005) (0.00974,0.00008) (0.01137,0.00013) (0.01327,0.00021) (0.01547,0.00034) (0.01804,0.00053) (0.02103,0.00085) (0.02449,0.00135) (0.02851,0.00215) (0.03315,0.00342) (0.03852,0.00543) (0.04469,0.00861) (0.05177,0.01361) (0.05984,0.02144) (0.06897,0.03363) (0.07918,0.05242) (0.09042,0.08099) (0.10251,0.12353) (0.11505,0.18504) (0.12741,0.27039) (0.13868,0.38242) (0.14780,0.51932) (0.15378,0.67294) (0.15603,0.82984) (0.15458,0.97566) (0.15000,1.10000)};
\draw[-{Latex},structurecolor] (axis cs:0.01547,0.00034) -- (axis cs:0.01804,0.00053);
\addplot[structurecolor,thin,no marks] coordinates {(0.15000,1.10000) (0.14317,1.19875) (0.13498,1.27312) (0.12617,1.32717) (0.11725,1.36572) (0.10855,1.39305) (0.10027,1.41254) (0.09250,1.42668) (0.08527,1.43715) (0.07858,1.44514) (0.07240,1.45142) (0.06671,1.45651) (0.06148,1.46074) (0.05667,1.46435) (0.05224,1.46748) (0.04817,1.47024) (0.04443,1.47271) (0.04099,1.47493) (0.03781,1.47694) (0.03490,1.47877) (0.03221,1.48044) (0.02973,1.48197) (0.02745,1.48338) (0.02534,1.48467) (0.02340,1.48585) (0.02161,1.48694) (0.01996,1.48795) (0.01844,1.48887) (0.01704,1.48972) (0.01574,1.49051) (0.01454,1.49123) (0.01344,1.49190) (0.01242,1.49252) (0.01148,1.49309) (0.01061,1.49361) (0.00980,1.49410) (0.00906,1.49455) (0.00838,1.49496) (0.00774,1.49534) (0.00716,1.49570) (0.00662,1.49602) (0.00612,1.49632) (0.00566,1.49660) (0.00523,1.49686) (0.00483,1.49710) (0.00447,1.49731) (0.00413,1.49752) (0.00382,1.49770) (0.00353,1.49788) (0.00327,1.49804) (0.00302,1.49819) (0.00279,1.49832) (0.00258,1.49845) (0.00239,1.49857) (0.00221,1.49867) (0.00204,1.49877) (0.00189,1.49887) (0.00175,1.49895) (0.00162,1.49903) (0.00149,1.49910) (0.00138,1.49917) (0.00128,1.49923) (0.00118,1.49929) (0.00109,1.49934) (0.00101,1.49939)};
\draw[-{Latex},structurecolor] (axis cs:0.03221,1.48044) -- (axis cs:0.02973,1.48197);
\draw[-{Latex},structurecolor] (axis cs:0.00523,1.49686) -- (axis cs:0.00483,1.49710);
\addplot[structurecolor,thin,no marks] coordinates {(0.21926,3.00000) (0.21733,2.96198) (0.21546,2.92533) (0.21363,2.89000) (0.21185,2.85590) (0.21012,2.82299) (0.20842,2.79120) (0.20677,2.76048) (0.20516,2.73078) (0.20358,2.70205) (0.20205,2.67424) (0.20054,2.64731) (0.19907,2.62123) (0.19764,2.59596) (0.19623,2.57146) (0.19485,2.54769) (0.19350,2.52464) (0.19218,2.50226) (0.19089,2.48053) (0.18962,2.45943) (0.18838,2.43892) (0.18717,2.41899) (0.18597,2.39961) (0.18480,2.38077) (0.18365,2.36243) (0.18252,2.34459) (0.18141,2.32722) (0.18033,2.31031) (0.17926,2.29384) (0.17821,2.27779) (0.17717,2.26215) (0.17616,2.24691) (0.17516,2.23204) (0.17418,2.21755) (0.17322,2.20341) (0.17227,2.18962) (0.17134,2.17615) (0.17042,2.16301) (0.16951,2.15018) (0.16862,2.13765) (0.16775,2.12542) (0.16688,2.11346) (0.16603,2.10178) (0.16519,2.09037) (0.16437,2.07921) (0.16355,2.06830) (0.16275,2.05763) (0.16196,2.04720) (0.16118,2.03700) (0.16041,2.02701) (0.15965,2.01725) (0.15890,2.00769) (0.15816,1.99833) (0.15744,1.98917) (0.15672,1.98021) (0.15601,1.97142) (0.15531,1.96282) (0.15461,1.95440) (0.15393,1.94615) (0.15325,1.93807) (0.15259,1.93014) (0.15193,1.92238) (0.15128,1.91477) (0.15064,1.90731) (0.15000,1.90000)};
\draw[-{Latex},structurecolor] (axis cs:0.18838,2.43892) -- (axis cs:0.18717,2.41899);
\draw[-{Latex},structurecolor] (axis cs:0.16519,2.09037) -- (axis cs:0.16437,2.07921);
\addplot[structurecolor,thin,no marks] coordinates {(0.15000,1.90000) (0.12977,1.69913) (0.11517,1.59482) (0.10368,1.53715) (0.09414,1.50451) (0.08592,1.48610) (0.07869,1.47607) (0.07223,1.47104) (0.06639,1.46899) (0.06109,1.46872) (0.05626,1.46947) (0.05183,1.47080) (0.04778,1.47242) (0.04406,1.47416) (0.04063,1.47591) (0.03749,1.47762) (0.03459,1.47926) (0.03193,1.48080) (0.02947,1.48225) (0.02721,1.48360) (0.02512,1.48485) (0.02320,1.48600) (0.02143,1.48707) (0.01979,1.48806) (0.01828,1.48897) (0.01689,1.48982) (0.01560,1.49059) (0.01442,1.49131) (0.01332,1.49197) (0.01231,1.49259) (0.01138,1.49315) (0.01052,1.49367) (0.00972,1.49415) (0.00898,1.49459) (0.00830,1.49500) (0.00768,1.49538) (0.00710,1.49573) (0.00656,1.49606) (0.00607,1.49635) (0.00561,1.49663) (0.00518,1.49688) (0.00479,1.49712) (0.00443,1.49734) (0.00410,1.49754) (0.00379,1.49772) (0.00350,1.49790) (0.00324,1.49805) (0.00300,1.49820) (0.00277,1.49834) (0.00256,1.49846) (0.00237,1.49858) (0.00219,1.49869) (0.00203,1.49878) (0.00187,1.49888) (0.00173,1.49896) (0.00160,1.49904) (0.00148,1.49911) (0.00137,1.49918) (0.00127,1.49924) (0.00117,1.49930) (0.00108,1.49935) (0.00100,1.49940) (0.00093,1.49944) (0.00086,1.49949) (0.00079,1.49952)};
\draw[-{Latex},structurecolor] (axis cs:0.02512,1.48485) -- (axis cs:0.02320,1.48600);
\draw[-{Latex},structurecolor] (axis cs:0.00410,1.49754) -- (axis cs:0.00379,1.49772);
\addplot[structurecolor,thin,no marks] coordinates {(0.00011,0.00000) (0.00012,0.00000) (0.00014,0.00000) (0.00017,0.00000) (0.00020,0.00000) (0.00023,0.00000) (0.00027,0.00000) (0.00031,0.00000) (0.00037,0.00000) (0.00043,0.00000) (0.00050,0.00000) (0.00059,0.00000) (0.00069,0.00000) (0.00080,0.00000) (0.00094,0.00000) (0.00109,0.00000) (0.00128,0.00000) (0.00150,0.00000) (0.00175,0.00000) (0.00204,0.00000) (0.00239,0.00000) (0.00279,0.00000) (0.00326,0.00000) (0.00381,0.00000) (0.00445,0.00000) (0.00520,0.00000) (0.00607,0.00000) (0.00709,0.00000) (0.00828,0.00000) (0.00967,0.00000) (0.01129,0.00000) (0.01317,0.00000) (0.01537,0.00000) (0.01792,0.00000) (0.02089,0.00000) (0.02433,0.00000) (0.02833,0.00000) (0.03296,0.00000) (0.03832,0.00000) (0.04452,0.00001) (0.05166,0.00001) (0.05987,0.00002) (0.06929,0.00002) (0.08007,0.00004) (0.09236,0.00006) (0.10632,0.00009) (0.12210,0.00015) (0.13986,0.00023) (0.15973,0.00036) (0.18182,0.00056) (0.20620,0.00087) (0.23291,0.00135) (0.26192,0.00207) (0.29311,0.00317) (0.32629,0.00482) (0.36120,0.00729) (0.39743,0.01094) (0.43449,0.01632) (0.47177,0.02414) (0.50854,0.03541) (0.54394,0.05142) (0.57697,0.07383) (0.60650,0.10462) (0.63129,0.14598) (0.65000,0.20000)};
\draw[-{Latex},structurecolor] (axis cs:0.00239,0.00000) -- (axis cs:0.00279,0.00000);
\draw[-{Latex},structurecolor] (axis cs:0.08007,0.00004) -- (axis cs:0.09236,0.00006);
\addplot[structurecolor,thin,no marks] coordinates {(0.65000,0.20000) (0.66135,0.26820) (0.66426,0.35088) (0.65809,0.44652) (0.64288,0.55154) (0.61943,0.66074) (0.58925,0.76831) (0.55428,0.86912) (0.51657,0.95959) (0.47797,1.03806) (0.43991,1.10444) (0.40340,1.15973) (0.36906,1.20542) (0.33720,1.24316) (0.30791,1.27445) (0.28114,1.30058) (0.25675,1.32260) (0.23459,1.34135) (0.21447,1.35746) (0.19620,1.37144) (0.17961,1.38368) (0.16453,1.39446) (0.15082,1.40403) (0.13834,1.41256) (0.12697,1.42020) (0.11659,1.42708) (0.10713,1.43328) (0.09848,1.43889) (0.09056,1.44398) (0.08333,1.44860) (0.07669,1.45281) (0.07062,1.45664) (0.06504,1.46015) (0.05993,1.46335) (0.05523,1.46627) (0.05092,1.46895) (0.04695,1.47141) (0.04330,1.47366) (0.03994,1.47573) (0.03686,1.47763) (0.03401,1.47938) (0.03139,1.48098) (0.02898,1.48246) (0.02676,1.48381) (0.02471,1.48506) (0.02282,1.48621) (0.02107,1.48727) (0.01946,1.48825) (0.01798,1.48915) (0.01661,1.48998) (0.01535,1.49075) (0.01418,1.49145) (0.01310,1.49211) (0.01211,1.49271) (0.01119,1.49326) (0.01034,1.49377) (0.00956,1.49425) (0.00884,1.49468) (0.00817,1.49509) (0.00755,1.49546) (0.00698,1.49580) (0.00645,1.49612) (0.00597,1.49641) (0.00552,1.49669) (0.00510,1.49694)};
\draw[-{Latex},structurecolor] (axis cs:0.17961,1.38368) -- (axis cs:0.16453,1.39446);
\draw[-{Latex},structurecolor] (axis cs:0.02676,1.48381) -- (axis cs:0.02471,1.48506);
\addplot[structurecolor,thin,no marks] coordinates {(2.00000,1.51155) (1.92981,1.46937) (1.86501,1.43143) (1.80496,1.39721) (1.74914,1.36625) (1.69711,1.33817) (1.64848,1.31265) (1.60290,1.28942) (1.56010,1.26822) (1.51980,1.24887) (1.48178,1.23118) (1.44585,1.21499) (1.41183,1.20016) (1.37955,1.18658) (1.34889,1.17413) (1.31971,1.16272) (1.29191,1.15227) (1.26537,1.14269) (1.24002,1.13393) (1.21575,1.12591) (1.19251,1.11858) (1.17022,1.11190) (1.14882,1.10582) (1.12825,1.10029) (1.10846,1.09527) (1.08939,1.09074) (1.07102,1.08666) (1.05329,1.08300) (1.03617,1.07974) (1.01962,1.07685) (1.00361,1.07430) (0.98812,1.07208) (0.97311,1.07017) (0.95856,1.06854) (0.94445,1.06719) (0.93074,1.06609) (0.91743,1.06523) (0.90449,1.06460) (0.89191,1.06418) (0.87967,1.06397) (0.86774,1.06394) (0.85613,1.06410) (0.84481,1.06442) (0.83377,1.06491) (0.82299,1.06555) (0.81248,1.06633) (0.80221,1.06725) (0.79217,1.06830) (0.78236,1.06946) (0.77277,1.07074) (0.76338,1.07213) (0.75420,1.07362) (0.74521,1.07521) (0.73640,1.07689) (0.72777,1.07865) (0.71932,1.08050) (0.71102,1.08242) (0.70289,1.08441) (0.69491,1.08646) (0.68708,1.08858) (0.67940,1.09076) (0.67185,1.09300) (0.66444,1.09529) (0.65716,1.09762) (0.65000,1.10000)};
\draw[-{Latex},structurecolor] (axis cs:1.19251,1.11858) -- (axis cs:1.17022,1.11190);
\draw[-{Latex},structurecolor] (axis cs:0.83377,1.06491) -- (axis cs:0.82299,1.06555);
\addplot[structurecolor,thin,no marks] coordinates {(0.65000,1.10000) (0.58015,1.12764) (0.52069,1.15720) (0.46920,1.18678) (0.42407,1.21523) (0.38416,1.24188) (0.34867,1.26642) (0.31696,1.28876) (0.28853,1.30895) (0.26296,1.32709) (0.23993,1.34336) (0.21914,1.35794) (0.20034,1.37100) (0.18330,1.38270) (0.16785,1.39321) (0.15382,1.40265) (0.14106,1.41116) (0.12944,1.41884) (0.11885,1.42577) (0.10918,1.43205) (0.10035,1.43775) (0.09228,1.44292) (0.08490,1.44763) (0.07813,1.45192) (0.07194,1.45583) (0.06625,1.45939) (0.06104,1.46266) (0.05625,1.46564) (0.05185,1.46837) (0.04781,1.47088) (0.04409,1.47318) (0.04067,1.47528) (0.03753,1.47722) (0.03463,1.47900) (0.03196,1.48063) (0.02950,1.48214) (0.02724,1.48352) (0.02515,1.48479) (0.02323,1.48596) (0.02145,1.48704) (0.01981,1.48804) (0.01830,1.48896) (0.01691,1.48980) (0.01562,1.49058) (0.01444,1.49130) (0.01334,1.49196) (0.01233,1.49258) (0.01139,1.49314) (0.01053,1.49366) (0.00973,1.49414) (0.00899,1.49459) (0.00831,1.49500) (0.00769,1.49538) (0.00710,1.49573) (0.00657,1.49605) (0.00607,1.49635) (0.00561,1.49663) (0.00519,1.49688) (0.00480,1.49712) (0.00444,1.49733) (0.00410,1.49754) (0.00379,1.49772) (0.00351,1.49789) (0.00324,1.49805) (0.00300,1.49820)};
\draw[-{Latex},structurecolor] (axis cs:0.10035,1.43775) -- (axis cs:0.09228,1.44292);
\draw[-{Latex},structurecolor] (axis cs:0.01562,1.49058) -- (axis cs:0.01444,1.49130);
\addplot[structurecolor,thin,no marks] coordinates {(0.96636,3.00000) (0.95800,2.96543) (0.94983,2.93191) (0.94184,2.89941) (0.93403,2.86788) (0.92639,2.83728) (0.91892,2.80756) (0.91160,2.77870) (0.90444,2.75066) (0.89743,2.72340) (0.89055,2.69689) (0.88382,2.67110) (0.87722,2.64602) (0.87076,2.62160) (0.86441,2.59782) (0.85819,2.57467) (0.85209,2.55211) (0.84610,2.53013) (0.84022,2.50870) (0.83445,2.48780) (0.82878,2.46743) (0.82321,2.44755) (0.81775,2.42815) (0.81237,2.40922) (0.80709,2.39073) (0.80190,2.37268) (0.79680,2.35506) (0.79178,2.33783) (0.78685,2.32100) (0.78199,2.30455) (0.77721,2.28847) (0.77251,2.27275) (0.76789,2.25737) (0.76333,2.24233) (0.75885,2.22762) (0.75443,2.21322) (0.75009,2.19912) (0.74580,2.18533) (0.74159,2.17182) (0.73743,2.15859) (0.73333,2.14564) (0.72929,2.13295) (0.72531,2.12052) (0.72139,2.10833) (0.71752,2.09640) (0.71371,2.08469) (0.70995,2.07322) (0.70624,2.06197) (0.70258,2.05095) (0.69896,2.04013) (0.69540,2.02952) (0.69188,2.01911) (0.68841,2.00890) (0.68499,1.99888) (0.68160,1.98904) (0.67826,1.97939) (0.67497,1.96991) (0.67171,1.96061) (0.66849,1.95148) (0.66532,1.94251) (0.66218,1.93370) (0.65908,1.92505) (0.65602,1.91655) (0.65299,1.90820) (0.65000,1.90000)};
\draw[-{Latex},structurecolor] (axis cs:0.82878,2.46743) -- (axis cs:0.82321,2.44755);
\draw[-{Latex},structurecolor] (axis cs:0.72139,2.10833) -- (axis cs:0.71752,2.09640);
\addplot[structurecolor,thin,no marks] coordinates {(0.65000,1.90000) (0.52885,1.61101) (0.45072,1.47629) (0.39405,1.40897) (0.34985,1.37592) (0.31366,1.36170) (0.28304,1.35819) (0.25656,1.36079) (0.23330,1.36679) (0.21265,1.37452) (0.19416,1.38299) (0.17753,1.39159) (0.16252,1.39998) (0.14891,1.40796) (0.13655,1.41544) (0.12531,1.42237) (0.11507,1.42875) (0.10573,1.43462) (0.09719,1.43999) (0.08939,1.44490) (0.08224,1.44940) (0.07570,1.45350) (0.06971,1.45726) (0.06421,1.46070) (0.05916,1.46384) (0.05453,1.46672) (0.05027,1.46936) (0.04635,1.47178) (0.04275,1.47400) (0.03944,1.47604) (0.03639,1.47792) (0.03358,1.47964) (0.03100,1.48122) (0.02862,1.48268) (0.02642,1.48402) (0.02440,1.48525) (0.02253,1.48639) (0.02081,1.48743) (0.01922,1.48840) (0.01776,1.48929) (0.01641,1.49011) (0.01516,1.49086) (0.01401,1.49156) (0.01294,1.49220) (0.01196,1.49280) (0.01105,1.49335) (0.01022,1.49385) (0.00944,1.49432) (0.00873,1.49475) (0.00807,1.49515) (0.00746,1.49551) (0.00689,1.49585) (0.00637,1.49617) (0.00589,1.49646) (0.00545,1.49673) (0.00504,1.49697) (0.00466,1.49720) (0.00431,1.49741) (0.00398,1.49761) (0.00368,1.49779) (0.00340,1.49796) (0.00315,1.49811) (0.00291,1.49825) (0.00269,1.49838) (0.00249,1.49851)};
\draw[-{Latex},structurecolor] (axis cs:0.08224,1.44940) -- (axis cs:0.07570,1.45350);
\draw[-{Latex},structurecolor] (axis cs:0.01294,1.49220) -- (axis cs:0.01196,1.49280);
\addplot[structurecolor,thin,no marks] coordinates {(2.00000,0.03773) (1.96251,0.03839) (1.92700,0.03910) (1.89333,0.03984) (1.86135,0.04062) (1.83094,0.04143) (1.80199,0.04229) (1.77441,0.04318) (1.74809,0.04411) (1.72295,0.04509) (1.69892,0.04610) (1.67593,0.04716) (1.65390,0.04826) (1.63278,0.04940) (1.61252,0.05058) (1.59306,0.05181) (1.57435,0.05309) (1.55636,0.05441) (1.53904,0.05579) (1.52235,0.05721) (1.50626,0.05868) (1.49074,0.06021) (1.47575,0.06179) (1.46127,0.06343) (1.44727,0.06512) (1.43372,0.06687) (1.42061,0.06867) (1.40790,0.07054) (1.39559,0.07248) (1.38364,0.07447) (1.37204,0.07654) (1.36078,0.07867) (1.34983,0.08086) (1.33918,0.08314) (1.32882,0.08548) (1.31873,0.08790) (1.30890,0.09039) (1.29932,0.09297) (1.28997,0.09562) (1.28084,0.09836) (1.27193,0.10119) (1.26322,0.10410) (1.25470,0.10710) (1.24636,0.11019) (1.23819,0.11337) (1.23019,0.11665) (1.22234,0.12003) (1.21464,0.12351) (1.20709,0.12708) (1.19966,0.13077) (1.19236,0.13456) (1.18518,0.13846) (1.17812,0.14247) (1.17116,0.14659) (1.16429,0.15083) (1.15753,0.15518) (1.15085,0.15966) (1.14426,0.16425) (1.13774,0.16897) (1.13130,0.17382) (1.12492,0.17879) (1.11861,0.18390) (1.11235,0.18913) (1.10615,0.19450) (1.10000,0.20000)};
\draw[-{Latex},structurecolor] (axis cs:1.50626,0.05868) -- (axis cs:1.49074,0.06021);
\draw[-{Latex},structurecolor] (axis cs:1.24636,0.11019) -- (axis cs:1.23819,0.11337);
\addplot[structurecolor,thin,no marks] coordinates {(1.10000,0.20000) (1.04967,0.25184) (1.00049,0.31396) (0.95094,0.38622) (0.90008,0.46747) (0.84753,0.55548) (0.79348,0.64710) (0.73855,0.73876) (0.68365,0.82703) (0.62978,0.90917) (0.57786,0.98341) (0.52863,1.04899) (0.48257,1.10594) (0.43993,1.15485) (0.40079,1.19661) (0.36506,1.23221) (0.33258,1.26258) (0.30313,1.28860) (0.27645,1.31100) (0.25230,1.33039) (0.23043,1.34729) (0.21062,1.36209) (0.19267,1.37514) (0.17638,1.38669) (0.16158,1.39697) (0.14813,1.40616) (0.13588,1.41440) (0.12472,1.42181) (0.11454,1.42850) (0.10525,1.43454) (0.09676,1.44002) (0.08900,1.44500) (0.08189,1.44952) (0.07538,1.45365) (0.06941,1.45741) (0.06394,1.46084) (0.05891,1.46398) (0.05430,1.46686) (0.05006,1.46949) (0.04616,1.47190) (0.04258,1.47411) (0.03928,1.47614) (0.03624,1.47801) (0.03345,1.47972) (0.03087,1.48130) (0.02850,1.48275) (0.02631,1.48408) (0.02430,1.48531) (0.02244,1.48644) (0.02073,1.48749) (0.01914,1.48845) (0.01769,1.48933) (0.01634,1.49015) (0.01510,1.49090) (0.01395,1.49159) (0.01289,1.49224) (0.01191,1.49283) (0.01101,1.49337) (0.01018,1.49388) (0.00940,1.49434) (0.00869,1.49477) (0.00804,1.49517) (0.00743,1.49553) (0.00687,1.49587) (0.00635,1.49618)};
\draw[-{Latex},structurecolor] (axis cs:0.23043,1.34729) -- (axis cs:0.21062,1.36209);
\draw[-{Latex},structurecolor] (axis cs:0.03345,1.47972) -- (axis cs:0.03087,1.48130);
\addplot[structurecolor,thin,no marks] coordinates {(2.00000,1.52946) (1.97250,1.51243) (1.94585,1.49608) (1.92001,1.48039) (1.89494,1.46532) (1.87060,1.45084) (1.84697,1.43691) (1.82400,1.42352) (1.80168,1.41063) (1.77997,1.39823) (1.75885,1.38628) (1.73830,1.37477) (1.71828,1.36369) (1.69878,1.35300) (1.67978,1.34269) (1.66125,1.33275) (1.64319,1.32315) (1.62556,1.31390) (1.60836,1.30496) (1.59157,1.29633) (1.57517,1.28800) (1.55915,1.27996) (1.54350,1.27218) (1.52819,1.26467) (1.51323,1.25741) (1.49860,1.25039) (1.48428,1.24360) (1.47027,1.23704) (1.45655,1.23070) (1.44312,1.22456) (1.42997,1.21863) (1.41709,1.21289) (1.40446,1.20733) (1.39208,1.20196) (1.37995,1.19677) (1.36805,1.19174) (1.35638,1.18687) (1.34493,1.18216) (1.33369,1.17761) (1.32266,1.17320) (1.31184,1.16894) (1.30121,1.16481) (1.29077,1.16082) (1.28051,1.15696) (1.27044,1.15323) (1.26054,1.14961) (1.25080,1.14612) (1.24124,1.14274) (1.23183,1.13947) (1.22258,1.13631) (1.21348,1.13325) (1.20453,1.13030) (1.19572,1.12745) (1.18705,1.12469) (1.17852,1.12202) (1.17012,1.11945) (1.16186,1.11697) (1.15371,1.11457) (1.14569,1.11225) (1.13780,1.11002) (1.13001,1.10786) (1.12235,1.10579) (1.11479,1.10379) (1.10734,1.10186) (1.10000,1.10000)};
\draw[-{Latex},structurecolor] (axis cs:1.57517,1.28800) -- (axis cs:1.55915,1.27996);
\draw[-{Latex},structurecolor] (axis cs:1.28051,1.15696) -- (axis cs:1.27044,1.15323);
\addplot[structurecolor,thin,no marks] coordinates {(1.10000,1.10000) (0.92763,1.07069) (0.80206,1.07072) (0.70460,1.08643) (0.62553,1.11050) (0.55934,1.13858) (0.50267,1.16798) (0.45341,1.19705) (0.41011,1.22478) (0.37176,1.25063) (0.33759,1.27436) (0.30703,1.29591) (0.27960,1.31536) (0.25492,1.33283) (0.23268,1.34849) (0.21258,1.36253) (0.19440,1.37510) (0.17792,1.38638) (0.16297,1.39651) (0.14938,1.40563) (0.13701,1.41384) (0.12575,1.42126) (0.11548,1.42796) (0.10611,1.43404) (0.09755,1.43955) (0.08971,1.44456) (0.08255,1.44912) (0.07598,1.45328) (0.06996,1.45707) (0.06444,1.46053) (0.05938,1.46369) (0.05472,1.46659) (0.05045,1.46925) (0.04652,1.47168) (0.04291,1.47391) (0.03958,1.47596) (0.03652,1.47784) (0.03370,1.47956) (0.03111,1.48115) (0.02872,1.48262) (0.02652,1.48396) (0.02448,1.48520) (0.02261,1.48634) (0.02088,1.48739) (0.01929,1.48836) (0.01782,1.48925) (0.01646,1.49007) (0.01521,1.49083) (0.01406,1.49153) (0.01299,1.49218) (0.01200,1.49277) (0.01109,1.49332) (0.01025,1.49383) (0.00948,1.49430) (0.00876,1.49473) (0.00810,1.49513) (0.00748,1.49550) (0.00692,1.49584) (0.00640,1.49615) (0.00591,1.49645) (0.00547,1.49671) (0.00505,1.49696) (0.00467,1.49719) (0.00432,1.49740) (0.00399,1.49760)};
\draw[-{Latex},structurecolor] (axis cs:0.13701,1.41384) -- (axis cs:0.12575,1.42126);
\draw[-{Latex},structurecolor] (axis cs:0.02088,1.48739) -- (axis cs:0.01929,1.48836);
\addplot[structurecolor,thin,no marks] coordinates {(1.65180,3.00000) (1.63754,2.96712) (1.62359,2.93515) (1.60992,2.90406) (1.59654,2.87382) (1.58344,2.84438) (1.57060,2.81572) (1.55802,2.78782) (1.54568,2.76064) (1.53359,2.73415) (1.52173,2.70833) (1.51009,2.68316) (1.49868,2.65861) (1.48747,2.63466) (1.47648,2.61129) (1.46568,2.58849) (1.45508,2.56622) (1.44467,2.54448) (1.43444,2.52324) (1.42439,2.50249) (1.41451,2.48222) (1.40480,2.46240) (1.39526,2.44302) (1.38587,2.42408) (1.37664,2.40555) (1.36756,2.38742) (1.35863,2.36969) (1.34984,2.35233) (1.34120,2.33534) (1.33269,2.31871) (1.32431,2.30242) (1.31606,2.28647) (1.30794,2.27084) (1.29994,2.25554) (1.29206,2.24054) (1.28430,2.22584) (1.27665,2.21143) (1.26911,2.19730) (1.26169,2.18345) (1.25437,2.16987) (1.24715,2.15654) (1.24004,2.14347) (1.23302,2.13065) (1.22611,2.11807) (1.21928,2.10572) (1.21256,2.09360) (1.20592,2.08171) (1.19937,2.07003) (1.19291,2.05856) (1.18653,2.04730) (1.18024,2.03624) (1.17403,2.02537) (1.16790,2.01470) (1.16184,2.00421) (1.15587,1.99391) (1.14997,1.98378) (1.14414,1.97383) (1.13838,1.96404) (1.13270,1.95443) (1.12708,1.94497) (1.12154,1.93567) (1.11606,1.92653) (1.11064,1.91754) (1.10529,1.90870) (1.10000,1.90000)};
\draw[-{Latex},structurecolor] (axis cs:1.41451,2.48222) -- (axis cs:1.40480,2.46240);
\draw[-{Latex},structurecolor] (axis cs:1.22611,2.11807) -- (axis cs:1.21928,2.10572);
\addplot[structurecolor,thin,no marks] coordinates {(1.10000,1.90000) (0.84948,1.53956) (0.70174,1.38538) (0.60083,1.31396) (0.52548,1.28281) (0.46582,1.27322) (0.41663,1.27563) (0.37493,1.28468) (0.33887,1.29719) (0.30725,1.31125) (0.27925,1.32569) (0.25429,1.33982) (0.23191,1.35326) (0.21178,1.36581) (0.19360,1.37740) (0.17716,1.38803) (0.16225,1.39772) (0.14872,1.40653) (0.13640,1.41454) (0.12519,1.42181) (0.11497,1.42842) (0.10564,1.43442) (0.09711,1.43988) (0.08932,1.44484) (0.08218,1.44937) (0.07565,1.45350) (0.06966,1.45727) (0.06416,1.46071) (0.05912,1.46386) (0.05449,1.46674) (0.05023,1.46938) (0.04632,1.47180) (0.04272,1.47402) (0.03941,1.47606) (0.03637,1.47793) (0.03356,1.47965) (0.03098,1.48123) (0.02860,1.48269) (0.02640,1.48403) (0.02438,1.48526) (0.02252,1.48640) (0.02080,1.48744) (0.01921,1.48841) (0.01775,1.48930) (0.01639,1.49011) (0.01515,1.49087) (0.01400,1.49157) (0.01293,1.49221) (0.01195,1.49280) (0.01105,1.49335) (0.01021,1.49386) (0.00944,1.49432) (0.00872,1.49475) (0.00806,1.49515) (0.00745,1.49552) (0.00689,1.49586) (0.00637,1.49617) (0.00589,1.49646) (0.00544,1.49673) (0.00503,1.49698) (0.00465,1.49720) (0.00430,1.49741) (0.00398,1.49761) (0.00368,1.49779) (0.00340,1.49796)};
\draw[-{Latex},structurecolor] (axis cs:0.11497,1.42842) -- (axis cs:0.10564,1.43442);
\draw[-{Latex},structurecolor] (axis cs:0.01775,1.48930) -- (axis cs:0.01639,1.49011);
\addplot[gray,dashed,domain=-2:2,samples=60] {1-x};
\addplot[gray,dashed,domain=-2:2,samples=60] {(3-x)/2};
\addplot[only marks,mark=*,mark size=1.4pt] coordinates {(0,0) (1,0) (0,1.5) (-1,2)};
\end{axis}\end{tikzpicture}
\end{center}
\end{solution}
\originalsection{5D：极限环}
\begin{exercise}\label{cw:bank-5D-1}原题5D-1. Notes LC 例1为 $x'=-y+x(1-x^2-y^2)$, $y'=x+y(1-x^2-y^2)$. (a) 证明原点是唯一平衡点，可比较非零平衡点处 $y/x$ 与 $-x/y$；(b) 验证 $(\cos t,\sin t)$ 为周期解并指出轨道；(c) 证明所有其他非零解的轨道都不断接近(b)的轨道，故它是唯一周期轨道且为渐近稳定极限环.\end{exercise}
\begin{solution}（官方译审）(a) 平衡时，用第一式乘 $-y$、第二式乘 $x$ 后相加得 $x^2+y^2=0$，故只有原点；这也避免了除以零. (b) 单位圆上非线性径向项消失，求导 $(\cos t,\sin t)'=(-\sin t,\cos t)$ 正好满足系统，最小周期 $2\pi$. (c) 对 $r=\sqrt{x^2+y^2}>0$，$r'=r(1-r^2)$, $\theta'=1$. 若 $0<r<1$ 则径向向外，$r>1$ 则向内；令 $\rho=r^2$，可解为
\[\rho(t)=\frac1{1+(\rho(0)^{-1}-1)e^{-2t}},\qquad t\ge0.\]
所以非零解全正向存在且 $r\to1$. 初始半径接近1时始终夹在初值与1之间，因此轨道集合稳定且吸引附近轨道. 其他半径严格单调，不可能周期，故单位圆是唯一非恒定周期轨道. 不同相位是同一轨道的时间平移，不应误称唯一时间参数函数.\end{solution}
\begin{exercise}\label{cw:bank-5D-2}原题5D-2. 证明下列系统在所给区域 $R$ 中无闭轨道，除(c)外 $R=\R^2$：
\[\begin{array}{ll}
(a)\ x'=x+x^3+y^3,&y'=y+x^3+y^3;\\
(b)\ x'=x^2+y^2,&y'=1+x-y;\\
(c)\ x'=2x+x^2+y^2,&y'=x^2-y^2,\quad R=\{x<-1\};\\
(d)\ x'=ax+bx^2-2cxy+dy^2,&y'=ex+fx^2-2bxy+cy^2.
\end{array}\]
(d) 求六常数的一组条件以保证全平面无闭轨道.\end{exercise}
\begin{solution}（官方译审）(a) 散度为 $2+3x^2+3y^2>0$，Bendixson判据在单连通全平面排除闭轨道. (b) $x'=x^2+y^2\ge0$；非恒定周期解要求一个周期内 $x$ 的总增量为零，只能有 $x=y=0$ 恒成立，但该点 $y'=1$，矛盾. (c) 半平面内散度为 $2+2x-2y$ 未必定号，不能乱用(a)；若闭轨道存在，它内部也位于凸半平面 $x<-1$，须包含平衡点. 而平衡条件 $y=\pm x$, $2x+2x^2=0$ 只给 $x=0,-1$，均不在 $R$，所以闭轨道不存在. 这里调用平面闭轨道包围平衡点的标准判据. (d) 散度恒为 $a$，因此 $a\ne0$ 是对其余五参数不附限制的充分条件. $a=0$ 时此判据不作决定：例如 $b=c=d=f=0,e=-1$ 时 $x'=0$, $y'=-x$ 无闭轨道，故所给充分条件不是必要条件；不能据散度零声称一定有闭轨道.\end{solution}
\begin{exercise}\label{cw:bank-5D-3}原题5D-3. Liénard方程为 $x''+u(x)x'+v(x)=0$. 证明若下列任一条件成立便无非恒定周期解：(a) 对所有 $x$ 有 $u(x)>0$；(b) 对所有 $x$ 有 $v(x)>0$. 提示：考察相应系统.\end{exercise}
\begin{solution}（官方译审）相应系统为 $x'=y$, $y'=-v(x)-u(x)y$. (a) 散度为 $-u(x)<0$，在通常 $C^1$ 条件下由Bendixson判据排除. 也可只用连续系数及经典解：能量 $E=y^2/2+\int_0^xv(s)\dd s$ 满足 $E'=-u(x)y^2$；周期积分为零，迫使 $y\equiv0$，解恒定. (b) 若周期为 $T$，积分第二式得
\[0=-\int_0^T v(x(t))\dd t-\int_0^T u(x(t))x'(t)\dd t.
\]
后一项是 $U(x(T))-U(x(0))=0$，其中 $U'=u$；前一项严格负，矛盾. 此证法也落实了“无平衡点故无闭轨道”的结论，无需假设 $u$ 定号.\end{solution}
\begin{exercise}\label{cw:bank-5D-4}原题5D-4（星号题）. van der Pol方程为 $x''-a(1-x^2)x'+x=0$，$a>0$. (a) 验证它满足Levinson--Smith定理条件，从而有唯一极限环；(b) 证明相应系统仅有原点平衡，判几何类型及稳定性，注意类型随参数改变.\end{exercise}
\begin{solution}（编者补解）(a) 使用 Notes LC 第5页的引用版本：$u$ 偶且连续，$v$ 奇且正半轴为正，$V(x)=\int_0^xv\to\infty$；存在 $k>0$ 使 $U=\int_0^xu$ 在 $(0,k)$ 为负，在 $x>k$ 为正、递增且趋无穷. 本题 $u=a(x^2-1)$, $v=x$, $V=x^2/2$, $U=a(x^3/3-x)$，取 $k=\sqrt3$，全部条件逐项成立，故定理给唯一吸引极限环. 此全局定理作为工具引用，不以数值轨道代替其证明. (b) 系统 $x'=y$, $y'=a(1-x^2)y-x$，平衡时 $y=x=0$. $J_0=\begin{pmatrix}0&1\\-1&a\end{pmatrix}$，根 $(a\pm\sqrt{a^2-4})/2$. $0<a<2$ 为不稳定螺旋点；$a>2$ 为不稳定结点；$a=2$ 为重根1、仅一条特征方向 $(1,1)^T$ 的不稳定退化结点. 各情况都双曲，原点均不稳定. $a=0$ 不在定理参数内，其系统是中心，无唯一孤立极限环.\end{solution}
\begin{exercise}\label{cw:bank-5D-5}原题5D-5（星号题）. 设 $r=\sqrt{x^2+y^2}$，系统 $x'=-y+xf(r)$, $y'=x+yf(r)$. (a) 若 $f$ 有正零点 $a$，证明有圆形周期解；(b) 若 $f$ 在 $r\approx a$ 附近递减，证明该解为稳定极限环，可观察方向场.\end{exercise}
\begin{solution}（编者补解；官方仅回指5D-1）对于 $r>0$，内积及叉积给 $r'=rf(r)$, $\theta'=1$. (a) $f(a)=0$ 时 $r(t)=a$, $\theta=t+\theta_0$，即 $(a\cos(t+\theta_0),a\sin(t+\theta_0))$，周期 $2\pi$. (b) 在 $a$ 附近严格递减且连续时，$r<a$ 有 $f(r)>0$，$r>a$ 有 $f(r)<0$，邻近半径单调趋 $a$，始终夹在初值与 $a$ 之间，故轨道稳定且吸引. 邻域内无其他零点，其余半径单调，圆周孤立，确为极限环. 若原文“decreasing”只解释为不增，结论不充分：$f\equiv0$ 会产生一族圆周而无极限环. 因此须按严格递减或至少两侧严格异号解释；题面保留后明确补足此条件.\end{solution}
\originalsection{5E：线性化的边界与Volterra原理}
\begin{exercise}\label{cw:bank-5E-1}原题5E-1. 各系统原点为平衡点，判线性化类型及稳定性，再说明仅凭线性化对非线性原点还允许哪些可能：(a) $x'=x-4y-xy^2$, $y'=2x-y+x^2y$；(b) $x'=3x-y+x^2+y^2$, $y'=-6x+2y+3xy$.\end{exercise}
\begin{solution}（官方译审并补证）(a) $J_0=\begin{pmatrix}1&-4\\2&-1\end{pmatrix}$，根 $\pm i\sqrt7$，线性化为稳定中心. 纯虚根非双曲，仅凭一阶项不能区分非线性中心、向内螺旋或向外螺旋. 对本题本身还能算散度 $x^2-y^2$，并非定号，因而上述不确定性不能用Bendixson一句消除. (b) $J_0=\begin{pmatrix}3&-1\\-6&2\end{pmatrix}$，根 $0,5$，线性化有整条 $y=3x$ 的平衡点线，横向向外，因此不稳定. 对非线性系统，正根已保证不稳定，但零特征方向的高阶行为未定：可能是孤立不稳定结点、鞍点或带退化行为的平衡点，不能断言也有整条平衡线. 官方只列常见相图情形，这些并非零根时的穷尽分类.\end{solution}
\begin{exercise}\label{cw:bank-5E-2}原题5E-2. 各系统含一个线性化非结构稳定的平衡点. 先求全部平衡点并判线性化类型，再画若干示意显示仅凭线性化时非线性轨道可能怎样变化：(a) $x'=y$, $y'=x(1-x)$；(b) $x'=x^2-x+y$, $y'=-yx^2-y$.\end{exercise}
\begin{solution}（官方译审并补证）(a) 平衡点 $(0,0),(1,0)$；$J=\begin{pmatrix}0&1\\1-2x&0\end{pmatrix}$. 原点根 $\pm1$，鞍点；$(1,0)$ 根 $\pm i$，线性中心. 一阶判据允许中心、吸引螺旋、排斥螺旋三种常见图形. 对给定完整方程可进一步确定：$H=y^2/2-x^2/2+x^3/3$ 守恒，在 $(1,0)$ 有严格局部极小，附近非临界小等值线为闭曲线，所以实际为中心. 图中的螺旋仅演示同一线性化可能对应的高阶扰动，不把它们冒称本方程轨道.

(b) $y'=-y(x^2+1)$ 给平衡必有 $y=0$，第一式给 $x=0,1$. 原点矩阵 $\begin{pmatrix}-1&1\\0&-1\end{pmatrix}$ 为重负根退化结点，渐近稳定；$(1,0)$ 矩阵 $\begin{pmatrix}1&1\\0&-2\end{pmatrix}$，根 $1,-2$ 为鞍点，方向 $(1,0)^T,(1,-3)^T$. 重根结点的精细结点/螺旋形状对矩阵微扰不稳定，这是原教材的几何类别用语；它是双曲汇，拓扑意义的结构稳定性及渐近稳定性仍保持. 改变一阶系数可使重实特征根变为复根，因而教材可比较结点与螺旋两种稳定图形；仅保持当前Jacobian的高阶项不能据此声称原点成为真螺旋. 对给定系统，$y$ 的符号保持且 $|y|$ 正向严格减少，轨道不能反复穿过 $y=0$，故本系统不可能真正螺旋绕原点，实际为稳定结点. 下列示意分别给出教材几何类别的比较，随后原方程相图给实际行为.
\begin{center}
\begin{tikzpicture}\begin{axis}[width=.31\linewidth,height=4cm,axis lines=middle,xlabel={$\xi$},ylabel={$\eta$},xmin=-1.6,xmax=1.6,ymin=-1.6,ymax=1.6,title={中心},ticks=none]\addplot[structurecolor,domain=0:6.2832,samples=160] ({0.35*(1)*cos(deg(x))},{-0.35*(1)*sin(deg(x))});\addplot[structurecolor,domain=0:6.2832,samples=160] ({0.75*(1)*cos(deg(x))},{-0.75*(1)*sin(deg(x))});\end{axis}\end{tikzpicture}
\hfill
\begin{tikzpicture}\begin{axis}[width=.31\linewidth,height=4cm,axis lines=middle,xlabel={$\xi$},ylabel={$\eta$},xmin=-1.6,xmax=1.6,ymin=-1.6,ymax=1.6,title={吸引螺旋},ticks=none]\addplot[structurecolor,domain=0:15,samples=160] ({0.35*(exp(-.07*x))*cos(deg(x))},{-0.35*(exp(-.07*x))*sin(deg(x))});\addplot[structurecolor,domain=0:15,samples=160] ({0.75*(exp(-.07*x))*cos(deg(x))},{-0.75*(exp(-.07*x))*sin(deg(x))});\end{axis}\end{tikzpicture}
\hfill
\begin{tikzpicture}\begin{axis}[width=.31\linewidth,height=4cm,axis lines=middle,xlabel={$\xi$},ylabel={$\eta$},xmin=-1.6,xmax=1.6,ymin=-1.6,ymax=1.6,title={排斥螺旋},ticks=none]\addplot[structurecolor,domain=0:15,samples=160] ({0.35*(exp(.07*x))*cos(deg(x))},{-0.35*(exp(.07*x))*sin(deg(x))});\addplot[structurecolor,domain=0:15,samples=160] ({0.75*(exp(.07*x))*cos(deg(x))},{-0.75*(exp(.07*x))*sin(deg(x))});\end{axis}\end{tikzpicture}
\end{center}
\noindent 上排为(a)在中心线性化附近的三类示意；坐标已平移至该平衡点. 对(b)，结点与螺旋的比较指改变一阶系数、使重实特征根变为复根的微扰；两者均保持渐近稳定. 下图显示给定原方程的实际结点.
\begin{center}
\begin{tikzpicture}\begin{axis}[width=.47\linewidth,height=5.2cm,axis lines=middle,xlabel={$x$},ylabel={$y$},xmin=-0.4,xmax=1.8,ymin=-0.8,ymax=0.8,title={（a）实际中心与鞍点},tick label style={font=\footnotesize},clip=true]
\addplot[structurecolor,thin,no marks] coordinates {(0.94338,-0.36980) (0.89768,-0.36066) (0.85340,-0.34705) (0.81107,-0.32960) (0.77113,-0.30894) (0.73395,-0.28566) (0.69981,-0.26029) (0.66894,-0.23329) (0.64154,-0.20505) (0.61772,-0.17589) (0.59759,-0.14609) (0.58122,-0.11583) (0.56864,-0.08528) (0.55990,-0.05454) (0.55501,-0.02370) (0.55398,0.00718) (0.55681,0.03805) (0.56349,0.06886) (0.57402,0.09952) (0.58836,0.12995) (0.60649,0.16002) (0.62834,0.18955) (0.65384,0.21831) (0.68288,0.24602) (0.71529,0.27232) (0.75088,0.29678) (0.78939,0.31891) (0.83049,0.33815) (0.87378,0.35390) (0.91879,0.36551) (0.96496,0.37234) (1.01165,0.37376) (1.05815,0.36922) (1.10369,0.35825) (1.14743,0.34053) (1.18853,0.31595) (1.22613,0.28458) (1.25941,0.24679) (1.28759,0.20315) (1.30999,0.15454) (1.32606,0.10203) (1.33539,0.04689) (1.33773,-0.00946) (1.33303,-0.06556) (1.32141,-0.11997) (1.30317,-0.17131) (1.27876,-0.21836) (1.24880,-0.26013) (1.21399,-0.29582) (1.17512,-0.32494) (1.13303,-0.34723) (1.08860,-0.36266) (1.04265,-0.37143) (0.99600,-0.37392) (0.94941,-0.37062) (0.90356,-0.36212) (0.85907,-0.34905) (0.81646,-0.33207) (0.77619,-0.31179) (0.73863,-0.28883) (0.70407,-0.26370) (0.67277,-0.23688) (0.64490,-0.20878) (0.62061,-0.17973) (0.60000,-0.15000)};
\draw[-{Latex},structurecolor] (axis cs:0.60649,0.16002) -- (axis cs:0.62834,0.18955);
\draw[-{Latex},structurecolor] (axis cs:1.33303,-0.06556) -- (axis cs:1.32141,-0.11997);
\addplot[structurecolor,thin,no marks] coordinates {(0.60000,-0.15000) (0.58313,-0.11979) (0.57006,-0.08927) (0.56082,-0.05855) (0.55543,-0.02772) (0.55390,0.00316) (0.55622,0.03403) (0.56240,0.06485) (0.57243,0.09554) (0.58628,0.12600) (0.60391,0.15613) (0.62529,0.18574) (0.65032,0.21462) (0.67890,0.24249) (0.71088,0.26899) (0.74607,0.29372) (0.78422,0.31618) (0.82500,0.33583) (0.86803,0.35207) (0.91285,0.36425) (0.95890,0.37174) (1.00556,0.37390) (1.05213,0.37016) (1.09784,0.36005) (1.14186,0.34323) (1.18336,0.31954) (1.22146,0.28904) (1.25535,0.25206) (1.28423,0.20914) (1.30742,0.16112) (1.32434,0.10905) (1.33456,0.05418) (1.33783,-0.00210) (1.33404,-0.05832) (1.32331,-0.11303) (1.30591,-0.16484) (1.28227,-0.21251) (1.25299,-0.25501) (1.21877,-0.29154) (1.18038,-0.32153) (1.13867,-0.34471) (1.09449,-0.36103) (1.04869,-0.37066) (1.00209,-0.37394) (0.95545,-0.37136) (0.90947,-0.36350) (0.86477,-0.35099) (0.82189,-0.33448) (0.78129,-0.31460) (0.74335,-0.29195) (0.70839,-0.26708) (0.67666,-0.24046) (0.64834,-0.21251) (0.62357,-0.18356) (0.60247,-0.15390) (0.58512,-0.12375) (0.57155,-0.09326) (0.56181,-0.06256) (0.55592,-0.03174) (0.55388,-0.00087) (0.55570,0.03001) (0.56138,0.06084) (0.57090,0.09155) (0.58426,0.12205) (0.60140,0.15223)};
\draw[-{Latex},structurecolor] (axis cs:0.95890,0.37174) -- (axis cs:1.00556,0.37390);
\draw[-{Latex},structurecolor] (axis cs:1.00209,-0.37394) -- (axis cs:0.95545,-0.37136);
\addplot[structurecolor,thin,no marks] coordinates {(0.60140,-0.15223) (0.58426,-0.12205) (0.57090,-0.09155) (0.56138,-0.06084) (0.55570,-0.03001) (0.55388,0.00087) (0.55592,0.03174) (0.56181,0.06256) (0.57155,0.09326) (0.58512,0.12375) (0.60247,0.15390) (0.62357,0.18356) (0.64834,0.21251) (0.67666,0.24046) (0.70839,0.26708) (0.74335,0.29195) (0.78129,0.31460) (0.82189,0.33448) (0.86477,0.35099) (0.90947,0.36350) (0.95545,0.37136) (1.00209,0.37394) (1.04869,0.37066) (1.09449,0.36103) (1.13867,0.34471) (1.18038,0.32153) (1.21877,0.29154) (1.25299,0.25501) (1.28227,0.21251) (1.30591,0.16484) (1.32331,0.11303) (1.33404,0.05832) (1.33783,0.00210) (1.33456,-0.05418) (1.32434,-0.10905) (1.30742,-0.16112) (1.28423,-0.20914) (1.25535,-0.25206) (1.22146,-0.28904) (1.18336,-0.31954) (1.14186,-0.34323) (1.09784,-0.36005) (1.05213,-0.37016) (1.00556,-0.37390) (0.95890,-0.37174) (0.91285,-0.36425) (0.86803,-0.35207) (0.82500,-0.33583) (0.78422,-0.31618) (0.74607,-0.29372) (0.71088,-0.26899) (0.67890,-0.24249) (0.65032,-0.21462) (0.62529,-0.18574) (0.60391,-0.15613) (0.58628,-0.12600) (0.57243,-0.09554) (0.56240,-0.06485) (0.55622,-0.03403) (0.55390,-0.00316) (0.55543,0.02772) (0.56082,0.05855) (0.57006,0.08927) (0.58313,0.11979) (0.60000,0.15000)};
\draw[-{Latex},structurecolor] (axis cs:0.95545,0.37136) -- (axis cs:1.00209,0.37394);
\draw[-{Latex},structurecolor] (axis cs:1.00556,-0.37390) -- (axis cs:0.95890,-0.37174);
\addplot[structurecolor,thin,no marks] coordinates {(0.60000,0.15000) (0.62061,0.17973) (0.64490,0.20878) (0.67277,0.23688) (0.70407,0.26370) (0.73863,0.28883) (0.77619,0.31179) (0.81646,0.33207) (0.85907,0.34905) (0.90356,0.36212) (0.94941,0.37062) (0.99600,0.37392) (1.04265,0.37143) (1.08860,0.36266) (1.13303,0.34723) (1.17512,0.32494) (1.21399,0.29582) (1.24880,0.26013) (1.27876,0.21836) (1.30317,0.17131) (1.32141,0.11997) (1.33303,0.06556) (1.33773,0.00946) (1.33539,-0.04689) (1.32606,-0.10203) (1.30999,-0.15454) (1.28759,-0.20315) (1.25941,-0.24679) (1.22613,-0.28458) (1.18853,-0.31595) (1.14743,-0.34053) (1.10369,-0.35825) (1.05815,-0.36922) (1.01165,-0.37376) (0.96496,-0.37234) (0.91879,-0.36551) (0.87378,-0.35390) (0.83049,-0.33815) (0.78939,-0.31891) (0.75088,-0.29678) (0.71529,-0.27232) (0.68288,-0.24602) (0.65384,-0.21831) (0.62834,-0.18955) (0.60649,-0.16002) (0.58836,-0.12995) (0.57402,-0.09952) (0.56349,-0.06886) (0.55681,-0.03805) (0.55398,-0.00718) (0.55501,0.02370) (0.55990,0.05454) (0.56864,0.08528) (0.58122,0.11583) (0.59759,0.14609) (0.61772,0.17589) (0.64154,0.20505) (0.66894,0.23329) (0.69981,0.26029) (0.73395,0.28566) (0.77113,0.30894) (0.81107,0.32960) (0.85340,0.34705) (0.89768,0.36066) (0.94338,0.36980)};
\draw[-{Latex},structurecolor] (axis cs:1.32141,0.11997) -- (axis cs:1.33303,0.06556);
\draw[-{Latex},structurecolor] (axis cs:0.62834,-0.18955) -- (axis cs:0.60649,-0.16002);
\addplot[structurecolor,thin,no marks] coordinates {(1.15098,-0.08212) (1.13940,-0.10294) (1.12533,-0.12171) (1.10907,-0.13810) (1.09092,-0.15188) (1.07122,-0.16285) (1.05033,-0.17093) (1.02861,-0.17607) (1.00643,-0.17830) (0.98415,-0.17772) (0.96211,-0.17445) (0.94064,-0.16867) (0.92004,-0.16056) (0.90059,-0.15035) (0.88253,-0.13826) (0.86609,-0.12451) (0.85146,-0.10934) (0.83881,-0.09297) (0.82826,-0.07561) (0.81994,-0.05747) (0.81392,-0.03876) (0.81026,-0.01967) (0.80901,-0.00039) (0.81017,0.01889) (0.81372,0.03799) (0.81965,0.05672) (0.82788,0.07488) (0.83834,0.09228) (0.85091,0.10870) (0.86546,0.12392) (0.88183,0.13773) (0.89983,0.14990) (0.91923,0.16019) (0.93979,0.16838) (0.96123,0.17427) (0.98325,0.17764) (1.00552,0.17833) (1.02771,0.17621) (1.04946,0.17119) (1.07039,0.16324) (1.09014,0.15238) (1.10837,0.13871) (1.12471,0.12242) (1.13887,0.10374) (1.15056,0.08300) (1.15955,0.06057) (1.16565,0.03690) (1.16874,0.01246) (1.16876,-0.01223) (1.16570,-0.03668) (1.15962,-0.06036) (1.15066,-0.08280) (1.13899,-0.10356) (1.12486,-0.12226) (1.10853,-0.13857) (1.09032,-0.15226) (1.07058,-0.16315) (1.04966,-0.17113) (1.02792,-0.17618) (1.00573,-0.17833) (0.98345,-0.17766) (0.96143,-0.17431) (0.93998,-0.16845) (0.91941,-0.16027) (0.90000,-0.15000)};
\draw[-{Latex},structurecolor] (axis cs:0.81392,-0.03876) -- (axis cs:0.81026,-0.01967);
\draw[-{Latex},structurecolor] (axis cs:1.13887,0.10374) -- (axis cs:1.15056,0.08300);
\addplot[structurecolor,thin,no marks] coordinates {(0.90000,-0.15000) (0.88199,-0.13785) (0.86561,-0.12406) (0.85104,-0.10884) (0.83844,-0.09244) (0.82797,-0.07505) (0.81971,-0.05689) (0.81377,-0.03817) (0.81019,-0.01907) (0.80901,0.00021) (0.81024,0.01950) (0.81387,0.03859) (0.81987,0.05730) (0.82818,0.07544) (0.83870,0.09281) (0.85134,0.10919) (0.86595,0.12438) (0.88237,0.13814) (0.90041,0.15025) (0.91986,0.16048) (0.94045,0.16860) (0.96191,0.17441) (0.98394,0.17770) (1.00622,0.17831) (1.02840,0.17610) (1.05013,0.17099) (1.07103,0.16294) (1.09074,0.15199) (1.10891,0.13824) (1.12519,0.12187) (1.13928,0.10312) (1.15089,0.08232) (1.15979,0.05985) (1.16580,0.03614) (1.16879,0.01169) (1.16871,-0.01301) (1.16555,-0.03743) (1.15938,-0.06108) (1.15033,-0.08348) (1.13859,-0.10418) (1.12438,-0.12281) (1.10798,-0.13904) (1.08972,-0.15265) (1.06994,-0.16344) (1.04898,-0.17133) (1.02723,-0.17629) (1.00503,-0.17835) (0.98276,-0.17759) (0.96075,-0.17416) (0.93932,-0.16823) (0.91879,-0.15998) (0.89941,-0.14965) (0.88145,-0.13744) (0.86512,-0.12360) (0.85061,-0.10835) (0.83808,-0.09190) (0.82767,-0.07449) (0.81949,-0.05631) (0.81362,-0.03758) (0.81011,-0.01847) (0.80901,0.00082) (0.81032,0.02010) (0.81403,0.03918) (0.82010,0.05788) (0.82847,0.07600)};
\draw[-{Latex},structurecolor] (axis cs:0.94045,0.16860) -- (axis cs:0.96191,0.17441);
\draw[-{Latex},structurecolor] (axis cs:1.06994,-0.16344) -- (axis cs:1.04898,-0.17133);
\addplot[structurecolor,thin,no marks] coordinates {(0.82847,-0.07600) (0.82010,-0.05788) (0.81403,-0.03918) (0.81032,-0.02010) (0.80901,-0.00082) (0.81011,0.01847) (0.81362,0.03758) (0.81949,0.05631) (0.82767,0.07449) (0.83808,0.09190) (0.85061,0.10835) (0.86512,0.12360) (0.88145,0.13744) (0.89941,0.14965) (0.91879,0.15998) (0.93932,0.16823) (0.96075,0.17416) (0.98276,0.17759) (1.00503,0.17835) (1.02723,0.17629) (1.04898,0.17133) (1.06994,0.16344) (1.08972,0.15265) (1.10798,0.13904) (1.12438,0.12281) (1.13859,0.10418) (1.15033,0.08348) (1.15938,0.06108) (1.16555,0.03743) (1.16871,0.01301) (1.16879,-0.01169) (1.16580,-0.03614) (1.15979,-0.05985) (1.15089,-0.08232) (1.13928,-0.10312) (1.12519,-0.12187) (1.10891,-0.13824) (1.09074,-0.15199) (1.07103,-0.16294) (1.05013,-0.17099) (1.02840,-0.17610) (1.00622,-0.17831) (0.98394,-0.17770) (0.96191,-0.17441) (0.94045,-0.16860) (0.91986,-0.16048) (0.90041,-0.15025) (0.88237,-0.13814) (0.86595,-0.12438) (0.85134,-0.10919) (0.83870,-0.09281) (0.82818,-0.07544) (0.81987,-0.05730) (0.81387,-0.03859) (0.81024,-0.01950) (0.80901,-0.00021) (0.81019,0.01907) (0.81377,0.03817) (0.81971,0.05689) (0.82797,0.07505) (0.83844,0.09244) (0.85104,0.10884) (0.86561,0.12406) (0.88199,0.13785) (0.90000,0.15000)};
\draw[-{Latex},structurecolor] (axis cs:1.04898,0.17133) -- (axis cs:1.06994,0.16344);
\draw[-{Latex},structurecolor] (axis cs:0.96191,-0.17441) -- (axis cs:0.94045,-0.16860);
\addplot[structurecolor,thin,no marks] coordinates {(0.90000,0.15000) (0.91941,0.16027) (0.93998,0.16845) (0.96143,0.17431) (0.98345,0.17766) (1.00573,0.17833) (1.02792,0.17618) (1.04966,0.17113) (1.07058,0.16315) (1.09032,0.15226) (1.10853,0.13857) (1.12486,0.12226) (1.13899,0.10356) (1.15066,0.08280) (1.15962,0.06036) (1.16570,0.03668) (1.16876,0.01223) (1.16874,-0.01246) (1.16565,-0.03690) (1.15955,-0.06057) (1.15056,-0.08300) (1.13887,-0.10374) (1.12471,-0.12242) (1.10837,-0.13871) (1.09014,-0.15238) (1.07039,-0.16324) (1.04946,-0.17119) (1.02771,-0.17621) (1.00552,-0.17833) (0.98325,-0.17764) (0.96123,-0.17427) (0.93979,-0.16838) (0.91923,-0.16019) (0.89983,-0.14990) (0.88183,-0.13773) (0.86546,-0.12392) (0.85091,-0.10870) (0.83834,-0.09228) (0.82788,-0.07488) (0.81965,-0.05672) (0.81372,-0.03799) (0.81017,-0.01889) (0.80901,0.00039) (0.81026,0.01967) (0.81392,0.03876) (0.81994,0.05747) (0.82826,0.07561) (0.83881,0.09297) (0.85146,0.10934) (0.86609,0.12451) (0.88253,0.13826) (0.90059,0.15035) (0.92004,0.16056) (0.94064,0.16867) (0.96211,0.17445) (0.98415,0.17772) (1.00643,0.17830) (1.02861,0.17607) (1.05033,0.17093) (1.07122,0.16285) (1.09092,0.15188) (1.10907,0.13810) (1.12533,0.12171) (1.13940,0.10294) (1.15098,0.08212)};
\draw[-{Latex},structurecolor] (axis cs:1.15056,-0.08300) -- (axis cs:1.13887,-0.10374);
\draw[-{Latex},structurecolor] (axis cs:0.81026,0.01967) -- (axis cs:0.81392,0.03876);
\addplot[structurecolor,thin,no marks] coordinates {(1.12023,0.13249) (1.13569,0.11441) (1.14874,0.09407) (1.15912,0.07183) (1.16663,0.04811) (1.17111,0.02340) (1.17246,-0.00181) (1.17066,-0.02698) (1.16574,-0.05158) (1.15781,-0.07512) (1.14703,-0.09711) (1.13361,-0.11715) (1.11784,-0.13487) (1.10000,-0.15000) (1.08045,-0.16232) (1.05955,-0.17170) (1.03765,-0.17808) (1.01515,-0.18148) (0.99241,-0.18196) (0.96978,-0.17964) (0.94761,-0.17469) (0.92621,-0.16730) (0.90587,-0.15768) (0.88687,-0.14606) (0.86943,-0.13267) (0.85377,-0.11776) (0.84005,-0.10154) (0.82843,-0.08424) (0.81902,-0.06607) (0.81193,-0.04725) (0.80723,-0.02796) (0.80495,-0.00841) (0.80513,0.01122) (0.80776,0.03075) (0.81281,0.04998) (0.82023,0.06872) (0.82996,0.08678) (0.84189,0.10393) (0.85590,0.11998) (0.87183,0.13469) (0.88950,0.14784) (0.90871,0.15919) (0.92922,0.16850) (0.95074,0.17556) (0.97300,0.18014) (0.99567,0.18207) (1.01840,0.18117) (1.04084,0.17735) (1.06261,0.17054) (1.08335,0.16073) (1.10267,0.14800) (1.12023,0.13248) (1.13569,0.11441) (1.14874,0.09407) (1.15913,0.07183) (1.16664,0.04811) (1.17111,0.02339) (1.17246,-0.00181) (1.17066,-0.02698) (1.16574,-0.05159) (1.15781,-0.07512) (1.14702,-0.09712) (1.13361,-0.11715) (1.11783,-0.13488) (1.10000,-0.15000)};
\draw[-{Latex},structurecolor] (axis cs:0.94761,-0.17469) -- (axis cs:0.92621,-0.16730);
\draw[-{Latex},structurecolor] (axis cs:0.95074,0.17556) -- (axis cs:0.97300,0.18014);
\addplot[structurecolor,thin,no marks] coordinates {(1.10000,-0.15000) (1.08045,-0.16232) (1.05954,-0.17170) (1.03765,-0.17808) (1.01515,-0.18148) (0.99240,-0.18196) (0.96977,-0.17964) (0.94760,-0.17469) (0.92620,-0.16730) (0.90587,-0.15768) (0.88687,-0.14606) (0.86943,-0.13267) (0.85376,-0.11775) (0.84004,-0.10153) (0.82842,-0.08423) (0.81902,-0.06607) (0.81193,-0.04725) (0.80723,-0.02796) (0.80495,-0.00841) (0.80513,0.01123) (0.80776,0.03075) (0.81281,0.04999) (0.82023,0.06873) (0.82996,0.08678) (0.84189,0.10394) (0.85590,0.11998) (0.87183,0.13470) (0.88951,0.14784) (0.90872,0.15919) (0.92922,0.16850) (0.95075,0.17556) (0.97301,0.18014) (0.99567,0.18207) (1.01841,0.18117) (1.04084,0.17735) (1.06262,0.17053) (1.08335,0.16073) (1.10268,0.14799) (1.12024,0.13248) (1.13569,0.11441) (1.14874,0.09407) (1.15913,0.07182) (1.16664,0.04810) (1.17111,0.02339) (1.17246,-0.00182) (1.17066,-0.02699) (1.16574,-0.05159) (1.15781,-0.07513) (1.14702,-0.09712) (1.13361,-0.11716) (1.11783,-0.13488) (1.10000,-0.15000) (1.08045,-0.16232) (1.05954,-0.17170) (1.03765,-0.17808) (1.01514,-0.18148) (0.99240,-0.18196) (0.96977,-0.17964) (0.94760,-0.17469) (0.92620,-0.16730) (0.90587,-0.15768) (0.88686,-0.14605) (0.86943,-0.13267) (0.85376,-0.11775) (0.84004,-0.10153)};
\draw[-{Latex},structurecolor] (axis cs:0.80776,0.03075) -- (axis cs:0.81281,0.04999);
\draw[-{Latex},structurecolor] (axis cs:1.17111,0.02339) -- (axis cs:1.17246,-0.00182);
\addplot[structurecolor,thin,no marks] coordinates {(0.84004,0.10153) (0.85376,0.11775) (0.86943,0.13267) (0.88686,0.14605) (0.90587,0.15768) (0.92620,0.16730) (0.94760,0.17469) (0.96977,0.17964) (0.99240,0.18196) (1.01514,0.18148) (1.03765,0.17808) (1.05954,0.17170) (1.08045,0.16232) (1.10000,0.15000) (1.11783,0.13488) (1.13361,0.11716) (1.14702,0.09712) (1.15781,0.07513) (1.16574,0.05159) (1.17066,0.02699) (1.17246,0.00182) (1.17111,-0.02339) (1.16664,-0.04810) (1.15913,-0.07182) (1.14874,-0.09407) (1.13569,-0.11441) (1.12024,-0.13248) (1.10268,-0.14799) (1.08335,-0.16073) (1.06262,-0.17053) (1.04084,-0.17735) (1.01841,-0.18117) (0.99567,-0.18207) (0.97301,-0.18014) (0.95075,-0.17556) (0.92922,-0.16850) (0.90872,-0.15919) (0.88951,-0.14784) (0.87183,-0.13470) (0.85590,-0.11998) (0.84189,-0.10394) (0.82996,-0.08678) (0.82023,-0.06873) (0.81281,-0.04999) (0.80776,-0.03075) (0.80513,-0.01123) (0.80495,0.00841) (0.80723,0.02796) (0.81193,0.04725) (0.81902,0.06607) (0.82842,0.08423) (0.84004,0.10153) (0.85376,0.11775) (0.86943,0.13267) (0.88687,0.14606) (0.90587,0.15768) (0.92620,0.16730) (0.94760,0.17469) (0.96977,0.17964) (0.99240,0.18196) (1.01515,0.18148) (1.03765,0.17808) (1.05954,0.17170) (1.08045,0.16232) (1.10000,0.15000)};
\draw[-{Latex},structurecolor] (axis cs:1.17246,0.00182) -- (axis cs:1.17111,-0.02339);
\draw[-{Latex},structurecolor] (axis cs:0.81281,-0.04999) -- (axis cs:0.80776,-0.03075);
\addplot[structurecolor,thin,no marks] coordinates {(1.10000,0.15000) (1.11783,0.13488) (1.13361,0.11715) (1.14702,0.09712) (1.15781,0.07512) (1.16574,0.05159) (1.17066,0.02698) (1.17246,0.00181) (1.17111,-0.02339) (1.16664,-0.04811) (1.15913,-0.07183) (1.14874,-0.09407) (1.13569,-0.11441) (1.12023,-0.13248) (1.10267,-0.14800) (1.08335,-0.16073) (1.06261,-0.17054) (1.04084,-0.17735) (1.01840,-0.18117) (0.99567,-0.18207) (0.97300,-0.18014) (0.95074,-0.17556) (0.92922,-0.16850) (0.90871,-0.15919) (0.88950,-0.14784) (0.87183,-0.13469) (0.85590,-0.11998) (0.84189,-0.10393) (0.82996,-0.08678) (0.82023,-0.06872) (0.81281,-0.04998) (0.80776,-0.03075) (0.80513,-0.01122) (0.80495,0.00841) (0.80723,0.02796) (0.81193,0.04725) (0.81902,0.06607) (0.82843,0.08424) (0.84005,0.10154) (0.85377,0.11776) (0.86943,0.13267) (0.88687,0.14606) (0.90587,0.15768) (0.92621,0.16730) (0.94761,0.17469) (0.96978,0.17964) (0.99241,0.18196) (1.01515,0.18148) (1.03765,0.17808) (1.05955,0.17170) (1.08045,0.16232) (1.10000,0.15000) (1.11784,0.13487) (1.13361,0.11715) (1.14703,0.09711) (1.15781,0.07512) (1.16574,0.05158) (1.17066,0.02698) (1.17246,0.00181) (1.17111,-0.02340) (1.16663,-0.04811) (1.15912,-0.07183) (1.14874,-0.09407) (1.13569,-0.11441) (1.12023,-0.13249)};
\draw[-{Latex},structurecolor] (axis cs:0.97300,-0.18014) -- (axis cs:0.95074,-0.17556);
\draw[-{Latex},structurecolor] (axis cs:0.92621,0.16730) -- (axis cs:0.94761,0.17469);
\addplot[structurecolor,thin,no marks] coordinates {(1.36265,0.24871) (1.38975,0.18384) (1.40840,0.11388) (1.41808,0.04070) (1.41852,-0.03364) (1.40971,-0.10701) (1.39190,-0.17736) (1.36558,-0.24279) (1.33147,-0.30168) (1.29048,-0.35277) (1.24364,-0.39517) (1.19207,-0.42843) (1.13692,-0.45245) (1.07933,-0.46748) (1.02040,-0.47407) (0.96114,-0.47296) (0.90244,-0.46505) (0.84512,-0.45128) (0.78983,-0.43265) (0.73712,-0.41009) (0.68743,-0.38448) (0.64109,-0.35661) (0.59834,-0.32717) (0.55934,-0.29670) (0.52419,-0.26568) (0.49293,-0.23445) (0.46558,-0.20326) (0.44211,-0.17229) (0.42249,-0.14162) (0.40669,-0.11129) (0.39466,-0.08128) (0.38636,-0.05154) (0.38176,-0.02198) (0.38086,0.00750) (0.38364,0.03701) (0.39011,0.06665) (0.40031,0.09652) (0.41426,0.12668) (0.43199,0.15718) (0.45357,0.18801) (0.47901,0.21911) (0.50835,0.25035) (0.54159,0.28151) (0.57871,0.31230) (0.61964,0.34231) (0.66424,0.37103) (0.71231,0.39784) (0.76359,0.42200) (0.81767,0.44268) (0.87407,0.45896) (0.93218,0.46987) (0.99127,0.47444) (1.05049,0.47173) (1.10887,0.46093) (1.16536,0.44137) (1.21883,0.41267) (1.26814,0.37472) (1.31214,0.32781) (1.34975,0.27262) (1.37999,0.21022) (1.40206,0.14207) (1.41534,0.06993) (1.41946,-0.00420) (1.41430,-0.07821) (1.40000,-0.15000)};
\draw[-{Latex},structurecolor] (axis cs:0.68743,-0.38448) -- (axis cs:0.64109,-0.35661);
\draw[-{Latex},structurecolor] (axis cs:0.57871,0.31230) -- (axis cs:0.61964,0.34231);
\addplot[structurecolor,thin,no marks] coordinates {(1.40000,-0.15000) (1.37697,-0.21759) (1.34585,-0.27924) (1.30747,-0.33354) (1.26281,-0.37946) (1.21298,-0.41637) (1.15911,-0.44403) (1.10235,-0.46258) (1.04383,-0.47243) (0.98458,-0.47427) (0.92556,-0.46893) (0.86760,-0.45737) (0.81143,-0.44054) (0.75765,-0.41943) (0.70672,-0.39493) (0.65902,-0.36787) (0.61483,-0.33897) (0.57433,-0.30885) (0.53764,-0.27800) (0.50484,-0.24682) (0.47594,-0.21559) (0.45094,-0.18452) (0.42980,-0.15372) (0.41249,-0.12326) (0.39897,-0.09313) (0.38920,-0.06329) (0.38314,-0.03367) (0.38078,-0.00417) (0.38210,0.02531) (0.38711,0.05489) (0.39583,0.08466) (0.40828,0.11470) (0.42452,0.14507) (0.44456,0.17577) (0.46847,0.20678) (0.49627,0.23798) (0.52797,0.26920) (0.56356,0.30018) (0.60299,0.33055) (0.64615,0.35985) (0.69288,0.38750) (0.74293,0.41280) (0.79595,0.43496) (0.85150,0.45310) (0.90902,0.46625) (0.96782,0.47345) (1.02709,0.47373) (1.08592,0.46622) (1.14329,0.45019) (1.19810,0.42514) (1.24919,0.39084) (1.29542,0.34742) (1.33569,0.29540) (1.36895,0.23570) (1.39435,0.16963) (1.41117,0.09884) (1.41894,0.02525) (1.41745,-0.04906) (1.40673,-0.12198) (1.38710,-0.19145) (1.35909,-0.25564) (1.32347,-0.31300) (1.28117,-0.36234) (1.23326,-0.40286) (1.18084,-0.43419)};
\draw[-{Latex},structurecolor] (axis cs:0.47594,-0.21559) -- (axis cs:0.45094,-0.18452);
\draw[-{Latex},structurecolor] (axis cs:0.85150,0.45310) -- (axis cs:0.90902,0.46625);
\addplot[structurecolor,thin,no marks] coordinates {(1.18084,0.43419) (1.23326,0.40286) (1.28117,0.36234) (1.32347,0.31300) (1.35909,0.25564) (1.38710,0.19145) (1.40673,0.12198) (1.41745,0.04906) (1.41894,-0.02525) (1.41117,-0.09884) (1.39435,-0.16963) (1.36895,-0.23570) (1.33569,-0.29540) (1.29542,-0.34742) (1.24919,-0.39084) (1.19810,-0.42514) (1.14329,-0.45019) (1.08592,-0.46622) (1.02709,-0.47373) (0.96782,-0.47345) (0.90902,-0.46625) (0.85150,-0.45310) (0.79595,-0.43496) (0.74293,-0.41280) (0.69288,-0.38750) (0.64615,-0.35985) (0.60299,-0.33055) (0.56356,-0.30018) (0.52797,-0.26920) (0.49627,-0.23798) (0.46847,-0.20678) (0.44456,-0.17577) (0.42452,-0.14507) (0.40828,-0.11470) (0.39583,-0.08466) (0.38711,-0.05489) (0.38210,-0.02531) (0.38078,0.00417) (0.38314,0.03367) (0.38920,0.06329) (0.39897,0.09313) (0.41249,0.12326) (0.42980,0.15372) (0.45094,0.18452) (0.47594,0.21559) (0.50484,0.24682) (0.53764,0.27800) (0.57433,0.30885) (0.61483,0.33897) (0.65902,0.36787) (0.70672,0.39493) (0.75765,0.41943) (0.81143,0.44054) (0.86760,0.45737) (0.92556,0.46893) (0.98458,0.47427) (1.04383,0.47243) (1.10235,0.46258) (1.15911,0.44403) (1.21298,0.41637) (1.26281,0.37946) (1.30747,0.33354) (1.34585,0.27924) (1.37697,0.21759) (1.40000,0.15000)};
\draw[-{Latex},structurecolor] (axis cs:0.90902,-0.46625) -- (axis cs:0.85150,-0.45310);
\draw[-{Latex},structurecolor] (axis cs:0.45094,0.18452) -- (axis cs:0.47594,0.21559);
\addplot[structurecolor,thin,no marks] coordinates {(1.40000,0.15000) (1.41430,0.07821) (1.41946,0.00420) (1.41534,-0.06993) (1.40206,-0.14207) (1.37999,-0.21022) (1.34975,-0.27262) (1.31214,-0.32781) (1.26814,-0.37472) (1.21883,-0.41267) (1.16536,-0.44137) (1.10887,-0.46093) (1.05049,-0.47173) (0.99127,-0.47444) (0.93218,-0.46987) (0.87407,-0.45896) (0.81767,-0.44268) (0.76359,-0.42200) (0.71231,-0.39784) (0.66424,-0.37103) (0.61964,-0.34231) (0.57871,-0.31230) (0.54159,-0.28151) (0.50835,-0.25035) (0.47901,-0.21911) (0.45357,-0.18801) (0.43199,-0.15718) (0.41426,-0.12668) (0.40031,-0.09652) (0.39011,-0.06665) (0.38364,-0.03701) (0.38086,-0.00750) (0.38176,0.02198) (0.38636,0.05154) (0.39466,0.08128) (0.40669,0.11129) (0.42249,0.14162) (0.44211,0.17229) (0.46558,0.20326) (0.49293,0.23445) (0.52419,0.26568) (0.55934,0.29670) (0.59834,0.32717) (0.64109,0.35661) (0.68743,0.38448) (0.73712,0.41009) (0.78983,0.43265) (0.84512,0.45128) (0.90244,0.46505) (0.96114,0.47296) (1.02040,0.47407) (1.07933,0.46748) (1.13692,0.45245) (1.19207,0.42843) (1.24364,0.39517) (1.29048,0.35277) (1.33147,0.30168) (1.36558,0.24279) (1.39190,0.17736) (1.40971,0.10701) (1.41852,0.03364) (1.41808,-0.04070) (1.40840,-0.11388) (1.38975,-0.18384) (1.36265,-0.24871)};
\draw[-{Latex},structurecolor] (axis cs:0.61964,-0.34231) -- (axis cs:0.57871,-0.31230);
\draw[-{Latex},structurecolor] (axis cs:0.64109,0.35661) -- (axis cs:0.68743,0.38448);
\addplot[only marks,mark=*,mark size=1.4pt] coordinates {(0,0) (1,0)};
\end{axis}\end{tikzpicture}
\hfill
\begin{tikzpicture}\begin{axis}[width=.47\linewidth,height=5.2cm,axis lines=middle,xlabel={$x$},ylabel={$y$},xmin=-1,xmax=1.7,ymin=-0.8,ymax=0.8,title={（b）实际结点与鞍点},tick label style={font=\footnotesize},clip=true]
\addplot[structurecolor,thin,no marks] coordinates {(-1.00000,-0.68858) (-0.98255,-0.67043) (-0.96556,-0.65305) (-0.94899,-0.63639) (-0.93285,-0.62043) (-0.91711,-0.60510) (-0.90176,-0.59038) (-0.88679,-0.57623) (-0.87218,-0.56262) (-0.85793,-0.54952) (-0.84402,-0.53690) (-0.83045,-0.52473) (-0.81719,-0.51299) (-0.80424,-0.50166) (-0.79159,-0.49071) (-0.77924,-0.48013) (-0.76716,-0.46991) (-0.75535,-0.46001) (-0.74381,-0.45042) (-0.73253,-0.44114) (-0.72149,-0.43215) (-0.71069,-0.42342) (-0.70013,-0.41496) (-0.68979,-0.40675) (-0.67967,-0.39878) (-0.66976,-0.39103) (-0.66006,-0.38350) (-0.65056,-0.37619) (-0.64125,-0.36907) (-0.63214,-0.36214) (-0.62320,-0.35540) (-0.61445,-0.34884) (-0.60586,-0.34245) (-0.59745,-0.33622) (-0.58920,-0.33015) (-0.58111,-0.32423) (-0.57317,-0.31845) (-0.56538,-0.31282) (-0.55774,-0.30732) (-0.55024,-0.30196) (-0.54289,-0.29672) (-0.53566,-0.29160) (-0.52857,-0.28660) (-0.52161,-0.28171) (-0.51477,-0.27694) (-0.50805,-0.27227) (-0.50145,-0.26770) (-0.49497,-0.26324) (-0.48860,-0.25887) (-0.48234,-0.25459) (-0.47619,-0.25041) (-0.47014,-0.24631) (-0.46420,-0.24230) (-0.45835,-0.23837) (-0.45260,-0.23452) (-0.44695,-0.23075) (-0.44139,-0.22706) (-0.43592,-0.22344) (-0.43054,-0.21989) (-0.42524,-0.21641) (-0.42003,-0.21300) (-0.41490,-0.20966) (-0.40986,-0.20637) (-0.40489,-0.20316) (-0.40000,-0.20000)};
\draw[-{Latex},structurecolor] (axis cs:-0.72149,-0.43215) -- (axis cs:-0.71069,-0.42342);
\draw[-{Latex},structurecolor] (axis cs:-0.52161,-0.28171) -- (axis cs:-0.51477,-0.27694);
\addplot[structurecolor,thin,no marks] coordinates {(-0.40000,-0.20000) (-0.35804,-0.17336) (-0.32144,-0.15080) (-0.28929,-0.13154) (-0.26087,-0.11500) (-0.23564,-0.10071) (-0.21312,-0.08832) (-0.19297,-0.07754) (-0.17487,-0.06814) (-0.15857,-0.05992) (-0.14385,-0.05273) (-0.13056,-0.04643) (-0.11851,-0.04089) (-0.10760,-0.03603) (-0.09769,-0.03175) (-0.08870,-0.02799) (-0.08052,-0.02468) (-0.07309,-0.02176) (-0.06633,-0.01920) (-0.06019,-0.01693) (-0.05460,-0.01494) (-0.04951,-0.01318) (-0.04489,-0.01163) (-0.04068,-0.01026) (-0.03685,-0.00905) (-0.03338,-0.00799) (-0.03022,-0.00705) (-0.02735,-0.00622) (-0.02474,-0.00549) (-0.02237,-0.00484) (-0.02023,-0.00427) (-0.01828,-0.00377) (-0.01651,-0.00333) (-0.01491,-0.00294) (-0.01346,-0.00259) (-0.01214,-0.00229) (-0.01095,-0.00202) (-0.00987,-0.00178) (-0.00890,-0.00157) (-0.00802,-0.00139) (-0.00722,-0.00122) (-0.00650,-0.00108) (-0.00585,-0.00095) (-0.00527,-0.00084) (-0.00474,-0.00074) (-0.00426,-0.00066) (-0.00383,-0.00058) (-0.00344,-0.00051) (-0.00309,-0.00045) (-0.00278,-0.00040) (-0.00250,-0.00035) (-0.00224,-0.00031) (-0.00201,-0.00027) (-0.00180,-0.00024) (-0.00162,-0.00021) (-0.00145,-0.00019) (-0.00130,-0.00017) (-0.00117,-0.00015) (-0.00105,-0.00013) (-0.00094,-0.00011) (-0.00084,-0.00010) (-0.00075,-0.00009) (-0.00067,-0.00008) (-0.00060,-0.00007) (-0.00054,-0.00006)};
\draw[-{Latex},structurecolor] (axis cs:-0.05460,-0.01494) -- (axis cs:-0.04951,-0.01318);
\draw[-{Latex},structurecolor] (axis cs:-0.00527,-0.00084) -- (axis cs:-0.00474,-0.00074);
\addplot[structurecolor,thin,no marks] coordinates {(-1.00000,0.37828) (-0.98378,0.37313) (-0.96792,0.36814) (-0.95241,0.36329) (-0.93724,0.35858) (-0.92240,0.35399) (-0.90788,0.34954) (-0.89367,0.34519) (-0.87977,0.34097) (-0.86615,0.33685) (-0.85281,0.33283) (-0.83975,0.32891) (-0.82695,0.32509) (-0.81441,0.32136) (-0.80212,0.31771) (-0.79008,0.31415) (-0.77827,0.31067) (-0.76669,0.30727) (-0.75534,0.30394) (-0.74420,0.30068) (-0.73327,0.29749) (-0.72255,0.29437) (-0.71204,0.29131) (-0.70171,0.28831) (-0.69158,0.28537) (-0.68163,0.28248) (-0.67186,0.27966) (-0.66227,0.27688) (-0.65285,0.27416) (-0.64359,0.27149) (-0.63450,0.26886) (-0.62557,0.26628) (-0.61680,0.26375) (-0.60817,0.26126) (-0.59969,0.25881) (-0.59136,0.25640) (-0.58317,0.25403) (-0.57511,0.25170) (-0.56719,0.24941) (-0.55940,0.24716) (-0.55174,0.24493) (-0.54421,0.24275) (-0.53680,0.24059) (-0.52951,0.23847) (-0.52233,0.23638) (-0.51527,0.23432) (-0.50832,0.23229) (-0.50148,0.23028) (-0.49475,0.22831) (-0.48813,0.22636) (-0.48161,0.22444) (-0.47518,0.22255) (-0.46886,0.22068) (-0.46264,0.21883) (-0.45650,0.21701) (-0.45046,0.21521) (-0.44452,0.21344) (-0.43866,0.21168) (-0.43289,0.20995) (-0.42720,0.20824) (-0.42160,0.20656) (-0.41608,0.20489) (-0.41064,0.20324) (-0.40528,0.20161) (-0.40000,0.20000)};
\draw[-{Latex},structurecolor] (axis cs:-0.73327,0.29749) -- (axis cs:-0.72255,0.29437);
\draw[-{Latex},structurecolor] (axis cs:-0.52951,0.23847) -- (axis cs:-0.52233,0.23638);
\addplot[structurecolor,thin,no marks] coordinates {(-0.40000,0.20000) (-0.31621,0.17371) (-0.25064,0.15177) (-0.19863,0.13310) (-0.15696,0.11700) (-0.12334,0.10300) (-0.09609,0.09076) (-0.07395,0.08002) (-0.05594,0.07058) (-0.04131,0.06227) (-0.02944,0.05495) (-0.01985,0.04849) (-0.01214,0.04279) (-0.00598,0.03776) (-0.00111,0.03332) (0.00269,0.02941) (0.00562,0.02595) (0.00783,0.02290) (0.00945,0.02021) (0.01058,0.01784) (0.01132,0.01574) (0.01174,0.01389) (0.01191,0.01226) (0.01188,0.01082) (0.01169,0.00955) (0.01139,0.00842) (0.01099,0.00743) (0.01053,0.00656) (0.01003,0.00579) (0.00950,0.00511) (0.00896,0.00451) (0.00841,0.00398) (0.00787,0.00351) (0.00734,0.00310) (0.00683,0.00273) (0.00633,0.00241) (0.00586,0.00213) (0.00541,0.00188) (0.00498,0.00166) (0.00458,0.00146) (0.00421,0.00129) (0.00386,0.00114) (0.00353,0.00101) (0.00323,0.00089) (0.00295,0.00078) (0.00269,0.00069) (0.00245,0.00061) (0.00223,0.00054) (0.00203,0.00048) (0.00184,0.00042) (0.00167,0.00037) (0.00152,0.00033) (0.00138,0.00029) (0.00125,0.00025) (0.00113,0.00022) (0.00102,0.00020) (0.00092,0.00017) (0.00083,0.00015) (0.00075,0.00014) (0.00068,0.00012) (0.00061,0.00011) (0.00055,0.00009) (0.00050,0.00008) (0.00045,0.00007) (0.00040,0.00006)};
\draw[-{Latex},structurecolor] (axis cs:0.01132,0.01574) -- (axis cs:0.01174,0.01389);
\draw[-{Latex},structurecolor] (axis cs:0.00323,0.00089) -- (axis cs:0.00295,0.00078);
\addplot[structurecolor,thin,no marks] coordinates {(0.97426,-0.80000) (0.96211,-0.77735) (0.95013,-0.75560) (0.93830,-0.73470) (0.92662,-0.71462) (0.91509,-0.69530) (0.90370,-0.67672) (0.89244,-0.65884) (0.88132,-0.64162) (0.87032,-0.62503) (0.85944,-0.60904) (0.84867,-0.59362) (0.83802,-0.57875) (0.82749,-0.56440) (0.81706,-0.55054) (0.80673,-0.53717) (0.79650,-0.52424) (0.78638,-0.51175) (0.77635,-0.49968) (0.76641,-0.48800) (0.75657,-0.47670) (0.74681,-0.46577) (0.73715,-0.45518) (0.72757,-0.44493) (0.71808,-0.43500) (0.70867,-0.42538) (0.69934,-0.41605) (0.69010,-0.40701) (0.68093,-0.39823) (0.67185,-0.38972) (0.66284,-0.38146) (0.65391,-0.37343) (0.64505,-0.36565) (0.63628,-0.35808) (0.62757,-0.35073) (0.61895,-0.34358) (0.61039,-0.33663) (0.60191,-0.32988) (0.59350,-0.32331) (0.58517,-0.31691) (0.57691,-0.31069) (0.56872,-0.30463) (0.56060,-0.29873) (0.55256,-0.29299) (0.54458,-0.28739) (0.53668,-0.28194) (0.52885,-0.27662) (0.52109,-0.27144) (0.51340,-0.26638) (0.50579,-0.26146) (0.49824,-0.25665) (0.49076,-0.25196) (0.48336,-0.24738) (0.47602,-0.24290) (0.46876,-0.23854) (0.46156,-0.23428) (0.45444,-0.23011) (0.44739,-0.22604) (0.44041,-0.22207) (0.43350,-0.21818) (0.42666,-0.21438) (0.41989,-0.21067) (0.41319,-0.20703) (0.40656,-0.20348) (0.40000,-0.20000)};
\draw[-{Latex},structurecolor] (axis cs:0.75657,-0.47670) -- (axis cs:0.74681,-0.46577);
\draw[-{Latex},structurecolor] (axis cs:0.55256,-0.29299) -- (axis cs:0.54458,-0.28739);
\addplot[structurecolor,thin,no marks] coordinates {(0.40000,-0.20000) (0.34749,-0.17345) (0.29993,-0.15108) (0.25718,-0.13204) (0.21909,-0.11570) (0.18540,-0.10159) (0.15583,-0.08932) (0.13007,-0.07863) (0.10777,-0.06927) (0.08860,-0.06105) (0.07221,-0.05384) (0.05829,-0.04748) (0.04652,-0.04189) (0.03664,-0.03696) (0.02838,-0.03261) (0.02152,-0.02878) (0.01586,-0.02540) (0.01121,-0.02241) (0.00743,-0.01978) (0.00438,-0.01745) (0.00194,-0.01540) (0.00002,-0.01359) (-0.00149,-0.01200) (-0.00263,-0.01059) (-0.00349,-0.00934) (-0.00411,-0.00824) (-0.00453,-0.00728) (-0.00480,-0.00642) (-0.00494,-0.00567) (-0.00498,-0.00500) (-0.00495,-0.00441) (-0.00485,-0.00389) (-0.00471,-0.00344) (-0.00453,-0.00303) (-0.00433,-0.00268) (-0.00411,-0.00236) (-0.00389,-0.00208) (-0.00366,-0.00184) (-0.00343,-0.00162) (-0.00321,-0.00143) (-0.00299,-0.00126) (-0.00277,-0.00112) (-0.00257,-0.00098) (-0.00238,-0.00087) (-0.00219,-0.00077) (-0.00202,-0.00068) (-0.00186,-0.00060) (-0.00170,-0.00053) (-0.00156,-0.00047) (-0.00143,-0.00041) (-0.00131,-0.00036) (-0.00119,-0.00032) (-0.00109,-0.00028) (-0.00099,-0.00025) (-0.00090,-0.00022) (-0.00082,-0.00019) (-0.00074,-0.00017) (-0.00068,-0.00015) (-0.00061,-0.00013) (-0.00056,-0.00012) (-0.00050,-0.00010) (-0.00046,-0.00009) (-0.00041,-0.00008) (-0.00037,-0.00007) (-0.00034,-0.00006)};
\draw[-{Latex},structurecolor] (axis cs:0.00194,-0.01540) -- (axis cs:0.00002,-0.01359);
\draw[-{Latex},structurecolor] (axis cs:-0.00238,-0.00087) -- (axis cs:-0.00219,-0.00077);
\addplot[structurecolor,thin,no marks] coordinates {(0.12128,0.80000) (0.13448,0.78438) (0.14720,0.76902) (0.15945,0.75390) (0.17125,0.73903) (0.18262,0.72439) (0.19357,0.70999) (0.20411,0.69581) (0.21427,0.68187) (0.22404,0.66815) (0.23345,0.65465) (0.24250,0.64136) (0.25120,0.62830) (0.25958,0.61545) (0.26763,0.60282) (0.27536,0.59039) (0.28279,0.57818) (0.28993,0.56617) (0.29677,0.55436) (0.30334,0.54277) (0.30963,0.53137) (0.31566,0.52018) (0.32143,0.50918) (0.32696,0.49838) (0.33223,0.48778) (0.33727,0.47737) (0.34208,0.46715) (0.34667,0.45713) (0.35103,0.44729) (0.35518,0.43763) (0.35913,0.42817) (0.36287,0.41888) (0.36641,0.40978) (0.36976,0.40085) (0.37292,0.39210) (0.37590,0.38352) (0.37870,0.37512) (0.38133,0.36688) (0.38378,0.35881) (0.38607,0.35091) (0.38820,0.34317) (0.39017,0.33559) (0.39199,0.32817) (0.39366,0.32090) (0.39518,0.31379) (0.39656,0.30683) (0.39780,0.30002) (0.39891,0.29335) (0.39988,0.28682) (0.40073,0.28044) (0.40145,0.27420) (0.40204,0.26809) (0.40252,0.26211) (0.40288,0.25627) (0.40313,0.25056) (0.40327,0.24497) (0.40330,0.23951) (0.40323,0.23417) (0.40305,0.22895) (0.40278,0.22384) (0.40240,0.21885) (0.40194,0.21398) (0.40138,0.20921) (0.40073,0.20455) (0.40000,0.20000)};
\draw[-{Latex},structurecolor] (axis cs:0.30963,0.53137) -- (axis cs:0.31566,0.52018);
\draw[-{Latex},structurecolor] (axis cs:0.39366,0.32090) -- (axis cs:0.39518,0.31379);
\addplot[structurecolor,thin,no marks] coordinates {(0.40000,0.20000) (0.39335,0.17306) (0.38380,0.14986) (0.37186,0.12991) (0.35803,0.11275) (0.34272,0.09799) (0.32633,0.08527) (0.30919,0.07431) (0.29160,0.06484) (0.27384,0.05665) (0.25613,0.04956) (0.23866,0.04340) (0.22160,0.03805) (0.20509,0.03339) (0.18922,0.02932) (0.17408,0.02577) (0.15973,0.02266) (0.14620,0.01994) (0.13352,0.01756) (0.12167,0.01546) (0.11066,0.01362) (0.10046,0.01200) (0.09106,0.01058) (0.08241,0.00933) (0.07447,0.00823) (0.06722,0.00726) (0.06060,0.00640) (0.05457,0.00565) (0.04910,0.00498) (0.04413,0.00439) (0.03964,0.00388) (0.03557,0.00342) (0.03190,0.00302) (0.02859,0.00266) (0.02561,0.00235) (0.02293,0.00207) (0.02052,0.00183) (0.01836,0.00161) (0.01641,0.00143) (0.01467,0.00126) (0.01311,0.00111) (0.01171,0.00098) (0.01045,0.00086) (0.00933,0.00076) (0.00833,0.00067) (0.00743,0.00059) (0.00663,0.00052) (0.00591,0.00046) (0.00527,0.00041) (0.00470,0.00036) (0.00419,0.00032) (0.00374,0.00028) (0.00333,0.00025) (0.00297,0.00022) (0.00264,0.00019) (0.00235,0.00017) (0.00210,0.00015) (0.00187,0.00013) (0.00166,0.00012) (0.00148,0.00010) (0.00132,0.00009) (0.00117,0.00008) (0.00104,0.00007) (0.00093,0.00006) (0.00083,0.00006)};
\draw[-{Latex},structurecolor] (axis cs:0.11066,0.01362) -- (axis cs:0.10046,0.01200);
\draw[-{Latex},structurecolor] (axis cs:0.00933,0.00076) -- (axis cs:0.00833,0.00067);
\addplot[structurecolor,thin,no marks] coordinates {(1.16890,-0.80000) (1.16256,-0.78023) (1.15635,-0.76106) (1.15025,-0.74248) (1.14426,-0.72446) (1.13838,-0.70698) (1.13260,-0.69001) (1.12693,-0.67355) (1.12135,-0.65757) (1.11586,-0.64205) (1.11047,-0.62698) (1.10516,-0.61234) (1.09994,-0.59811) (1.09480,-0.58429) (1.08974,-0.57085) (1.08476,-0.55779) (1.07986,-0.54509) (1.07503,-0.53273) (1.07026,-0.52072) (1.06557,-0.50902) (1.06095,-0.49765) (1.05638,-0.48658) (1.05189,-0.47580) (1.04745,-0.46531) (1.04307,-0.45509) (1.03874,-0.44514) (1.03447,-0.43545) (1.03026,-0.42602) (1.02610,-0.41682) (1.02198,-0.40786) (1.01792,-0.39912) (1.01390,-0.39061) (1.00993,-0.38231) (1.00600,-0.37422) (1.00212,-0.36634) (0.99827,-0.35864) (0.99447,-0.35114) (0.99071,-0.34382) (0.98698,-0.33668) (0.98329,-0.32972) (0.97964,-0.32292) (0.97601,-0.31629) (0.97243,-0.30982) (0.96887,-0.30350) (0.96535,-0.29733) (0.96185,-0.29131) (0.95839,-0.28543) (0.95495,-0.27969) (0.95154,-0.27408) (0.94816,-0.26861) (0.94480,-0.26326) (0.94146,-0.25804) (0.93815,-0.25293) (0.93487,-0.24795) (0.93160,-0.24307) (0.92836,-0.23831) (0.92513,-0.23366) (0.92193,-0.22911) (0.91874,-0.22466) (0.91558,-0.22032) (0.91243,-0.21607) (0.90930,-0.21192) (0.90618,-0.20785) (0.90308,-0.20388) (0.90000,-0.20000)};
\draw[-{Latex},structurecolor] (axis cs:1.06095,-0.49765) -- (axis cs:1.05638,-0.48658);
\draw[-{Latex},structurecolor] (axis cs:0.96887,-0.30350) -- (axis cs:0.96535,-0.29733);
\addplot[structurecolor,thin,no marks] coordinates {(0.90000,-0.20000) (0.86462,-0.16014) (0.83047,-0.12918) (0.79696,-0.10495) (0.76366,-0.08583) (0.73031,-0.07064) (0.69673,-0.05849) (0.66285,-0.04872) (0.62869,-0.04081) (0.59432,-0.03437) (0.55987,-0.02910) (0.52551,-0.02475) (0.49142,-0.02115) (0.45781,-0.01814) (0.42490,-0.01563) (0.39288,-0.01350) (0.36195,-0.01171) (0.33227,-0.01018) (0.30398,-0.00887) (0.27719,-0.00774) (0.25197,-0.00677) (0.22838,-0.00594) (0.20644,-0.00521) (0.18612,-0.00457) (0.16741,-0.00402) (0.15025,-0.00354) (0.13458,-0.00311) (0.12033,-0.00274) (0.10741,-0.00242) (0.09573,-0.00213) (0.08520,-0.00188) (0.07574,-0.00166) (0.06726,-0.00146) (0.05967,-0.00129) (0.05288,-0.00114) (0.04684,-0.00100) (0.04145,-0.00088) (0.03666,-0.00078) (0.03241,-0.00069) (0.02863,-0.00061) (0.02529,-0.00054) (0.02232,-0.00047) (0.01970,-0.00042) (0.01738,-0.00037) (0.01533,-0.00033) (0.01351,-0.00029) (0.01191,-0.00025) (0.01050,-0.00022) (0.00925,-0.00020) (0.00815,-0.00017) (0.00718,-0.00015) (0.00633,-0.00014) (0.00557,-0.00012) (0.00491,-0.00011) (0.00432,-0.00009) (0.00381,-0.00008) (0.00335,-0.00007) (0.00295,-0.00006) (0.00260,-0.00006) (0.00229,-0.00005) (0.00201,-0.00004) (0.00177,-0.00004) (0.00156,-0.00003) (0.00137,-0.00003) (0.00121,-0.00003)};
\draw[-{Latex},structurecolor] (axis cs:0.25197,-0.00677) -- (axis cs:0.22838,-0.00594);
\draw[-{Latex},structurecolor] (axis cs:0.01738,-0.00037) -- (axis cs:0.01533,-0.00033);
\addplot[structurecolor,thin,no marks] coordinates {(0.65386,0.80000) (0.66127,0.78517) (0.66852,0.77052) (0.67561,0.75605) (0.68254,0.74176) (0.68932,0.72765) (0.69595,0.71373) (0.70244,0.69998) (0.70878,0.68642) (0.71498,0.67304) (0.72104,0.65985) (0.72696,0.64685) (0.73275,0.63403) (0.73841,0.62140) (0.74394,0.60895) (0.74934,0.59669) (0.75462,0.58461) (0.75977,0.57272) (0.76481,0.56102) (0.76973,0.54950) (0.77453,0.53816) (0.77922,0.52701) (0.78380,0.51604) (0.78828,0.50525) (0.79264,0.49464) (0.79690,0.48421) (0.80106,0.47396) (0.80512,0.46389) (0.80908,0.45399) (0.81295,0.44426) (0.81672,0.43471) (0.82039,0.42534) (0.82398,0.41613) (0.82747,0.40709) (0.83088,0.39821) (0.83421,0.38950) (0.83744,0.38096) (0.84060,0.37257) (0.84368,0.36435) (0.84667,0.35628) (0.84959,0.34837) (0.85244,0.34061) (0.85520,0.33301) (0.85790,0.32555) (0.86052,0.31825) (0.86308,0.31108) (0.86556,0.30407) (0.86798,0.29719) (0.87033,0.29046) (0.87262,0.28386) (0.87484,0.27739) (0.87701,0.27107) (0.87911,0.26487) (0.88115,0.25880) (0.88313,0.25286) (0.88506,0.24704) (0.88693,0.24135) (0.88874,0.23578) (0.89050,0.23033) (0.89221,0.22499) (0.89387,0.21977) (0.89548,0.21466) (0.89703,0.20967) (0.89854,0.20478) (0.90000,0.20000)};
\draw[-{Latex},structurecolor] (axis cs:0.77453,0.53816) -- (axis cs:0.77922,0.52701);
\draw[-{Latex},structurecolor] (axis cs:0.85790,0.32555) -- (axis cs:0.86052,0.31825);
\addplot[structurecolor,thin,no marks] coordinates {(0.90000,0.20000) (0.91174,0.15928) (0.91992,0.12657) (0.92514,0.10042) (0.92785,0.07960) (0.92840,0.06307) (0.92707,0.04998) (0.92405,0.03963) (0.91947,0.03145) (0.91342,0.02499) (0.90596,0.01988) (0.89710,0.01585) (0.88684,0.01266) (0.87517,0.01014) (0.86205,0.00814) (0.84746,0.00656) (0.83136,0.00530) (0.81373,0.00430) (0.79454,0.00350) (0.77379,0.00286) (0.75150,0.00235) (0.72771,0.00193) (0.70248,0.00160) (0.67589,0.00133) (0.64809,0.00111) (0.61921,0.00093) (0.58944,0.00079) (0.55898,0.00067) (0.52805,0.00057) (0.49690,0.00048) (0.46576,0.00041) (0.43487,0.00036) (0.40448,0.00031) (0.37480,0.00027) (0.34603,0.00023) (0.31834,0.00020) (0.29187,0.00018) (0.26675,0.00015) (0.24304,0.00013) (0.22080,0.00012) (0.20006,0.00010) (0.18081,0.00009) (0.16304,0.00008) (0.14670,0.00007) (0.13174,0.00006) (0.11810,0.00005) (0.10569,0.00005) (0.09445,0.00004) (0.08429,0.00004) (0.07514,0.00003) (0.06690,0.00003) (0.05951,0.00003) (0.05289,0.00002) (0.04697,0.00002) (0.04168,0.00002) (0.03697,0.00002) (0.03277,0.00001) (0.02903,0.00001) (0.02571,0.00001) (0.02276,0.00001) (0.02014,0.00001) (0.01782,0.00001) (0.01576,0.00001) (0.01393,0.00001) (0.01232,0.00001)};
\draw[-{Latex},structurecolor] (axis cs:0.75150,0.00235) -- (axis cs:0.72771,0.00193);
\draw[-{Latex},structurecolor] (axis cs:0.14670,0.00007) -- (axis cs:0.13174,0.00006);
\addplot[structurecolor,thin,no marks] coordinates {(1.24634,-0.80000) (1.24182,-0.78130) (1.23741,-0.76312) (1.23310,-0.74544) (1.22890,-0.72823) (1.22480,-0.71149) (1.22080,-0.69520) (1.21690,-0.67935) (1.21309,-0.66391) (1.20937,-0.64888) (1.20575,-0.63424) (1.20220,-0.61998) (1.19875,-0.60609) (1.19537,-0.59256) (1.19208,-0.57937) (1.18887,-0.56651) (1.18573,-0.55398) (1.18267,-0.54176) (1.17968,-0.52985) (1.17677,-0.51823) (1.17392,-0.50690) (1.17115,-0.49585) (1.16844,-0.48507) (1.16580,-0.47455) (1.16322,-0.46428) (1.16070,-0.45426) (1.15825,-0.44449) (1.15585,-0.43494) (1.15352,-0.42562) (1.15124,-0.41652) (1.14902,-0.40764) (1.14685,-0.39896) (1.14473,-0.39049) (1.14267,-0.38221) (1.14067,-0.37412) (1.13871,-0.36623) (1.13680,-0.35851) (1.13494,-0.35097) (1.13313,-0.34360) (1.13136,-0.33640) (1.12964,-0.32936) (1.12796,-0.32248) (1.12633,-0.31575) (1.12474,-0.30918) (1.12320,-0.30275) (1.12169,-0.29646) (1.12023,-0.29032) (1.11880,-0.28431) (1.11741,-0.27843) (1.11607,-0.27269) (1.11476,-0.26706) (1.11348,-0.26156) (1.11225,-0.25619) (1.11105,-0.25092) (1.10988,-0.24578) (1.10875,-0.24074) (1.10765,-0.23581) (1.10658,-0.23099) (1.10555,-0.22627) (1.10455,-0.22165) (1.10358,-0.21713) (1.10264,-0.21271) (1.10173,-0.20838) (1.10085,-0.20415) (1.10000,-0.20000)};
\draw[-{Latex},structurecolor] (axis cs:1.17392,-0.50690) -- (axis cs:1.17115,-0.49585);
\draw[-{Latex},structurecolor] (axis cs:1.12474,-0.30918) -- (axis cs:1.12320,-0.30275);
\addplot[structurecolor,thin,no marks] coordinates {(1.10000,-0.20000) (1.09690,-0.18435) (1.09423,-0.16996) (1.09195,-0.15673) (1.09006,-0.14455) (1.08852,-0.13334) (1.08732,-0.12301) (1.08644,-0.11349) (1.08587,-0.10471) (1.08559,-0.09662) (1.08559,-0.08915) (1.08587,-0.08226) (1.08640,-0.07589) (1.08719,-0.07002) (1.08823,-0.06460) (1.08951,-0.05959) (1.09103,-0.05496) (1.09279,-0.05069) (1.09477,-0.04674) (1.09700,-0.04309) (1.09945,-0.03972) (1.10214,-0.03660) (1.10506,-0.03372) (1.10822,-0.03107) (1.11163,-0.02861) (1.11528,-0.02634) (1.11918,-0.02424) (1.12334,-0.02230) (1.12777,-0.02051) (1.13247,-0.01885) (1.13746,-0.01733) (1.14275,-0.01592) (1.14834,-0.01461) (1.15425,-0.01341) (1.16050,-0.01230) (1.16711,-0.01128) (1.17408,-0.01033) (1.18145,-0.00946) (1.18923,-0.00866) (1.19744,-0.00791) (1.20611,-0.00723) (1.21527,-0.00660) (1.22496,-0.00602) (1.23519,-0.00549) (1.24602,-0.00500) (1.25748,-0.00455) (1.26962,-0.00413) (1.28248,-0.00375) (1.29611,-0.00340) (1.31059,-0.00307) (1.32598,-0.00278) (1.34234,-0.00251) (1.35976,-0.00226) (1.37833,-0.00203) (1.39816,-0.00182) (1.41936,-0.00163) (1.44205,-0.00146) (1.46638,-0.00130) (1.49252,-0.00116) (1.52066,-0.00102) (1.55100,-0.00091) (1.58380,-0.00080) (1.61934,-0.00070) (1.65794,-0.00061) (1.70000,-0.00053)};
\draw[-{Latex},structurecolor] (axis cs:1.09945,-0.03972) -- (axis cs:1.10214,-0.03660);
\draw[-{Latex},structurecolor] (axis cs:1.23519,-0.00549) -- (axis cs:1.24602,-0.00500);
\addplot[structurecolor,thin,no marks] coordinates {(0.79785,0.80000) (0.80490,0.78550) (0.81185,0.77116) (0.81868,0.75700) (0.82541,0.74300) (0.83202,0.72917) (0.83854,0.71551) (0.84495,0.70202) (0.85126,0.68870) (0.85747,0.67556) (0.86359,0.66259) (0.86961,0.64979) (0.87554,0.63716) (0.88138,0.62471) (0.88713,0.61243) (0.89279,0.60033) (0.89837,0.58840) (0.90387,0.57664) (0.90928,0.56506) (0.91462,0.55365) (0.91988,0.54241) (0.92506,0.53134) (0.93017,0.52044) (0.93521,0.50971) (0.94018,0.49915) (0.94507,0.48876) (0.94991,0.47854) (0.95468,0.46848) (0.95938,0.45859) (0.96403,0.44886) (0.96861,0.43930) (0.97313,0.42989) (0.97760,0.42065) (0.98202,0.41156) (0.98638,0.40264) (0.99069,0.39387) (0.99495,0.38525) (0.99916,0.37679) (1.00333,0.36848) (1.00745,0.36031) (1.01152,0.35230) (1.01555,0.34443) (1.01955,0.33671) (1.02350,0.32913) (1.02741,0.32170) (1.03129,0.31440) (1.03514,0.30724) (1.03895,0.30022) (1.04273,0.29334) (1.04648,0.28658) (1.05020,0.27996) (1.05389,0.27347) (1.05756,0.26710) (1.06120,0.26086) (1.06481,0.25474) (1.06841,0.24875) (1.07198,0.24287) (1.07554,0.23712) (1.07908,0.23148) (1.08260,0.22596) (1.08610,0.22055) (1.08960,0.21525) (1.09307,0.21006) (1.09654,0.20498) (1.10000,0.20000)};
\draw[-{Latex},structurecolor] (axis cs:0.91988,0.54241) -- (axis cs:0.92506,0.53134);
\draw[-{Latex},structurecolor] (axis cs:1.02350,0.32913) -- (axis cs:1.02741,0.32170);
\addplot[structurecolor,thin,no marks] coordinates {(1.10000,0.20000) (1.10549,0.19229) (1.11096,0.18484) (1.11643,0.17763) (1.12188,0.17067) (1.12734,0.16395) (1.13279,0.15746) (1.13826,0.15119) (1.14375,0.14514) (1.14925,0.13930) (1.15478,0.13367) (1.16034,0.12823) (1.16593,0.12299) (1.17157,0.11793) (1.17725,0.11306) (1.18298,0.10836) (1.18878,0.10383) (1.19463,0.09947) (1.20056,0.09526) (1.20656,0.09121) (1.21264,0.08731) (1.21881,0.08356) (1.22507,0.07994) (1.23143,0.07646) (1.23790,0.07311) (1.24448,0.06989) (1.25118,0.06679) (1.25801,0.06381) (1.26496,0.06094) (1.27206,0.05818) (1.27931,0.05553) (1.28672,0.05299) (1.29429,0.05054) (1.30203,0.04819) (1.30995,0.04593) (1.31807,0.04376) (1.32639,0.04168) (1.33491,0.03968) (1.34366,0.03776) (1.35264,0.03592) (1.36187,0.03415) (1.37135,0.03246) (1.38109,0.03083) (1.39112,0.02928) (1.40144,0.02778) (1.41208,0.02635) (1.42303,0.02498) (1.43433,0.02367) (1.44598,0.02241) (1.45801,0.02121) (1.47044,0.02006) (1.48328,0.01896) (1.49656,0.01791) (1.51029,0.01690) (1.52451,0.01594) (1.53924,0.01502) (1.55451,0.01414) (1.57034,0.01331) (1.58677,0.01251) (1.60384,0.01174) (1.62157,0.01102) (1.64001,0.01032) (1.65919,0.00966) (1.67918,0.00904) (1.70000,0.00844)};
\draw[-{Latex},structurecolor] (axis cs:1.21264,0.08731) -- (axis cs:1.21881,0.08356);
\draw[-{Latex},structurecolor] (axis cs:1.39112,0.02928) -- (axis cs:1.40144,0.02778);
\addplot[only marks,mark=*,mark size=1.4pt] coordinates {(0,0) (1,0)};
\end{axis}\end{tikzpicture}
\end{center}
\end{solution}
\begin{exercise}\label{cw:bank-5E-3}原题5E-3. Monet Gardens的Pristine Acres以野花吸引游客，Kandinsky蛀虫以花茎为食，数量按Volterra捕食模型周期变化. 园长拟施肥，使无蛀虫时野花增长率提高25\%，希望野花增长跑赢蛀虫. 你如何评价？用适当单位设野花 $x$、蛀虫 $y$，$x'=ax-pxy$, $y'=-by+qxy$，各参数为正常数.\end{exercise}
\begin{solution}（官方译审并补证）正平衡点为 $(b/q,a/p)$，施肥把 $a$ 改为 $5a/4$，只将蛀虫平衡数增为 $5a/(4p)$，野花平衡数不变. 题面说周期波动，评价应比较周期平均数而非只比较平衡点：正周期解上
\[0=\frac1T\int_0^T(\log x)'\dd t=a-p\bar y,\qquad
0=\frac1T\int_0^T(\log y)'\dd t=-b+q\bar x.\]
所以 $\bar x=b/q$, $\bar y=a/p$. 在维持同一捕食模型和周期假设下，施肥提高蛀虫平均数25\%，不提高野花平均数，不能实现预期目标. 这并不声称花量每个时刻都不变，也不把模型外的肥料效应纳入计算.\end{solution}
\originalchapter{原题题库六：幂级数}
\originalsection{6A：幂级数运算}
\begin{exercise}\label{cw:bank-6A-1}原题6A-1. 求各幂级数的收敛半径：
\[(a)\ \sum_{n=0}^\infty nx^n,\quad (b)\ \sum_{n=1}^\infty\frac{x^{2n}}{n2^n},\quad
(c)\ \sum_{n=1}^\infty n!x^n,\quad (d)\ \sum_{n=0}^\infty\frac{(2n)!}{(n!)^2}x^n.\]
原(b)写从0求和而分母含 $n$，零项未定义；这里把该零项排除，保留所有有定义的正阶项.
\end{exercise}
\begin{solution}（官方译审）(a) 相邻非零项绝对值比为 $(n+1)|x|/n\to|x|$，故 $R=1$. (b) 以非零项顺序取比得 $\frac{n}{2(n+1)}|x|^2\to|x|^2/2$，故 $R=\sqrt2$；不能忽略此处指数是 $2n$. (c) 比值 $(n+1)|x|\to\infty$ 对任意 $x\ne0$，所以只有 $x=0$ 收敛，$R=0$. (d) 比值 $\frac{(2n+2)(2n+1)}{(n+1)^2}|x|\to4|x|$，故 $R=1/4$. 收敛半径不包含端点结论：例如(b)在 $x=\pm\sqrt2$ 为调和级数，均发散.\end{solution}
\begin{exercise}\label{cw:bank-6A-2}原题6A-2. 从 $\sum_{n\ge0}x^n=(1-x)^{-1}$ 和 $\sum_{n\ge0}x^n/n!=e^x$ 出发，利用代入、加法、乘法、逐项微分或积分，求：(a) $(1-x)^{-2}$；(b) $xe^{-x^2}$；(c) $\arctan x$；(d) $\log(1+x)$ 的幂级数.\end{exercise}
\begin{solution}（官方译审）(a) 对几何级数微分并移项指标，$\frac1{(1-x)^2}=\sum_{n\ge0}(n+1)x^n$, $|x|<1$. (b) 将指数级数的变量换为 $-x^2$ 再乘 $x$，得 $xe^{-x^2}=\sum_{n\ge0}(-1)^nx^{2n+1}/n!$，对所有 $x$ 收敛. (c) $\frac1{1+x^2}=\sum_{n\ge0}(-1)^nx^{2n}$，从0积分得 $\arctan x=\sum_{n\ge0}(-1)^nx^{2n+1}/(2n+1)$，$|x|<1$，在 $x=\pm1$ 也按交错级数收敛至相应值. (d) 对 $\sum(-x)^n=(1+x)^{-1}$ 从0积分，得 $\log(1+x)=\sum_{n\ge1}(-1)^{n-1}x^n/n$, $|x|<1$；$x=1$ 也收敛，$x=-1$ 发散. 逐项积分、微分在收敛圆内部成立；积分常数由0处函数值决定.\end{solution}
\begin{exercise}\label{cw:bank-6A-3}原题6A-3. 设 $y=\sum_{n=0}^\infty x^{2n+1}/(2n+1)!$. 证明：(a) $y''-y=0$；(b) $y=\sinh x=(e^x-e^{-x})/2$.\end{exercise}
\begin{solution}（官方译审）该级数项比趋零，收敛半径为无穷，可在任何有界区间逐项求导. (a) $y'=\sum_{n\ge0}x^{2n}/(2n)!$，再次求导并把 $n\ge1$ 移为 $n-1\ge0$，得到原级数，故 $y''=y$. (b) $e^x$ 与 $e^{-x}$ 级数相减时偶项抵消、奇项加倍，除以2恰为 $y$，所以等于 $\sinh x$.\end{solution}
\begin{exercise}\label{cw:bank-6A-4}原题6A-4. 利用6A-2的运算求和：(a) $\sum_{n=0}^\infty x^{3n+2}$；(b) $\sum_{n=0}^\infty x^n/(n+1)$；(c) $\sum_{n=0}^\infty nx^n$.\end{exercise}
\begin{solution}（官方译审）(a) 把 $x^2$ 提出，剩下以 $x^3$ 为公比的几何级数，和为 $x^2/(1-x^3)$，$|x|<1$. (b) 几何级数从0积分得 $\sum x^{n+1}/(n+1)=-\log(1-x)$，故和为 $-\log(1-x)/x$，$0<|x|<1$，在0按连续延拓取值1. (c) 几何级数微分后乘 $x$，得 $x/(1-x)^2$，$|x|<1$. 每式都在该范围内由合法级数运算获得，不将有理函数的范围外延拓冒称原级数的和.\end{solution}
\originalsection{6B：一阶方程的级数解}
\begin{exercise}\label{cw:bank-6B-1}原题6B-1. 对非线性初值问题 $y'=x+y^2$, $y(0)=1$，用两种方法求级数解的前四个非零项：(a) 设 $y=\sum a_nx^n$，逐次代入求系数；(b) 对方程反复求导，再用Taylor公式.\end{exercise}
\begin{solution}（官方译审）(a) $a_0=1$；系数比较给
\[(n+1)a_{n+1}=\mathbf1_{n=1}+\sum_{j=0}^na_ja_{n-j}.
\]
所以 $a_1=1$, $2a_2=1+2a_1=3$, $3a_3=2a_2+a_1^2=4$. 因而 $y=1+x+\frac32x^2+\frac43x^3+O(x^4)$. (b) $y''=1+2yy'$, $y'''=2(y')^2+2yy''$；在0得 $y=1$, $y'=1$, $y''=3$, $y'''=8$. 按Taylor的阶乘系数得到同一四项. 方程右端解析，初值处存在解析局部解，这些系数对应实际局部解，而不只是形式运算.\end{solution}
\begin{exercise}\label{cw:bank-6B-2}原题6B-2. 假设 $y=\sum_{n\ge0}a_nx^n$，代入方程递推求系数，再尽可能求和为闭式：(a) $y'=x+y$, $y(0)=0$；(b) $y'=-xy$, $y(0)=1$；(c) $(1-x)y'-y=0$, $y(0)=1$.\end{exercise}
\begin{solution}（官方译审）(a) $(n+1)a_{n+1}=a_n+\mathbf1_{n=1}$，$a_0=a_1=0$, $a_n=1/n!$（$n\ge2$），故 $y=e^x-1-x$，$R=\infty$. (b) 常数项给 $a_1=0$；对 $n\ge1$，$(n+1)a_{n+1}=-a_{n-1}$，所以奇项为零，$a_{2k}=(-1)^k/(2^kk!)$. 故 $y=e^{-x^2/2}$，$R=\infty$. (c) 比较系数得 $(n+1)a_{n+1}-(n+1)a_n=0$，从 $a_0=1$ 得所有 $a_n=1$，故 $y=(1-x)^{-1}$，$R=1$，实际初值解的最大实区间为 $x<1$. 所有求和结果代回方程及初值均成立.\end{solution}
\originalsection{6C：二阶方程的级数解}
\begin{exercise}\label{cw:bank-6C-1}原题6C-1. 把以下各式改写为 $\sum_{n=N}^\infty b_nx^n$，说明 $N$ 并用原 $a_n$ 表示 $b_n$：(a) $\sum_{n=1}^\infty a_nx^{n+3}$；(b) $\sum_{n=0}^\infty n(n-1)a_nx^{n-2}$；(c) $\sum_{n=1}^\infty(n+1)a_nx^{n-1}$.\end{exercise}
\begin{solution}（官方译审）(a) 把旧指标设为 $n-3$，得 $N=4$, $b_n=a_{n-3}$. (b) 原 $n=0,1$ 两项因系数为零先删去，再把旧指标设为 $n+2$，得 $N=0$, $b_n=(n+2)(n+1)a_{n+2}$. (c) 把旧指标设为 $n+1$，得 $N=0$, $b_n=(n+2)a_{n+1}$. 重新索引必须同时改指数、系数和求和下限.\end{solution}
\begin{exercise}\label{cw:bank-6C-2}原题6C-2. 通过 $a_n$ 的递推关系，求 $y''-4y=0$ 的两个独立幂级数解 $\sum a_nx^n$.\end{exercise}
\begin{solution}（官方译审）代入得 $(n+2)(n+1)a_{n+2}=4a_n$，偶、奇系数各成独立链. 分别取初值 $(a_0,a_1)=(1,0),(0,1)$，得
\[y_1=\sum_{k\ge0}\frac{4^kx^{2k}}{(2k)!}=\cosh2x,
\qquad y_2=\sum_{k\ge0}\frac{4^kx^{2k+1}}{(2k+1)!}=\tfrac12\sinh2x.\]
两级数收敛半径无穷；0处初值列为单位矩阵，Wronskian在0为1，因此独立.\end{solution}
\begin{exercise}\label{cw:bank-6C-3}原题6C-3. 对 $y''+2xy'+2y=0$：(a) 求两个独立级数解 $y_1,y_2$；(b) 求各自收敛半径；(c) 用它们表示 $y(0)=1$, $y'(0)=-1$ 的解；(d) 原文要求“把级数表示为初等函数（即把级数求和为初等函数）”，提示一个易识别，另一个可用级数运算或降阶法得到. 官方解答写“Not solved”. 保留(d)原要求，下面说明其数学限制.\end{exercise}
\begin{solution}（编者补解）代入级数得 $(n+2)(n+1)a_{n+2}+2(n+1)a_n=0$，即 $a_{n+2}=-2a_n/(n+2)$. (a) 按初值 $(1,0),(0,1)$ 取
\[y_1=\sum_{k\ge0}\frac{(-1)^kx^{2k}}{k!}=e^{-x^2},\qquad
 y_2=\sum_{k\ge0}\frac{(-1)^k4^kk!}{(2k+1)!}x^{2k+1}.
\]
两解在0的初值不同且Wronskian为1，故独立. (b) 两级数相邻非零项的绝对值比分别为 $|x|^2/(k+1)$ 与 $2|x|^2/(2k+3)$，均趋零，因此各 $R=\infty$. (c) 所求解为 $y_1-y_2$，因为初值线性组合给系数 $1,-1$.

(d) 方程实际为 $(y'+2xy)'=0$. 对第二解有 $y'+2xy=1$，乘积分因子 $e^{x^2}$ 从0积分，得
\[y_2(x)=e^{-x^2}\int_0^x e^{s^2}\dd s.\]
这也是降阶公式所得的精确求积表示，且直接求导可恢复方程和初值. 但 $\int e^{s^2}\dd s$ 是非初等积分，因此两个独立解均化为初等函数的原要求不能实现. 一个标准的Liouville判据把此非初等性归结为：若它有初等原函数，须存在有理函数 $R$ 满足 $R'+2xR=1$. $R$ 若有有限极点，其导数最高极点阶比 $2xR$ 高一阶，不可能抵消，所以 $R$ 只能是多项式；非零多项式中 $2xR$ 的次数比 $R'$ 高，不能和为常数1，$R=0$ 也不满足. 这里将该非初等积分判据作为外部代数工具说明勘误，不把其一般理论算作本册的证明目标. 级数、全局收敛与精确求积已完整解决可成立的要求.\end{solution}
\begin{exercise}\label{cw:bank-6C-4}原题6C-4. Hermite方程 $y''-2xy'+ky=0$. 证明若 $k=2m$ 为正偶整数，则一个幂级数解是次数为 $m$ 的多项式.\end{exercise}
\begin{solution}（官方译审）代入得
\[a_{n+2}=\frac{2(n-m)}{(n+2)(n+1)}a_n.
\]
若 $m$ 偶，取 $a_0=1,a_1=0$；若 $m$ 奇，取 $a_0=0,a_1=1$. 异奇偶链恒零，与 $m$ 同奇偶的链在 $n<m$ 时乘子均非零，故 $a_m\ne0$；到 $n=m$ 乘子为零，使 $a_{m+2}=0$，随后同链全零. 因而解恰为次数 $m$ 的多项式，不只是次数不超过 $m$.\end{solution}
\begin{exercise}\label{cw:bank-6C-5}原题6C-5. 求Airy方程 $y''=xy$ 的两个独立 $x$ 幂级数解，确定其收敛半径，各写前三个非零项及一般项.\end{exercise}
\begin{solution}（官方译审并补证）比较常数项给 $a_2=0$，对 $n\ge1$ 有 $(n+2)(n+1)a_{n+2}=a_{n-1}$，等价于 $a_{n+3}=a_n/[(n+3)(n+2)]$（$n\ge0$）. 分别取 $(a_0,a_1)=(1,0),(0,1)$，得
\[\begin{aligned}
y_1&=\sum_{k\ge0}\frac{x^{3k}}{\prod_{j=1}^k(3j)(3j-1)}=1+\frac{x^3}{6}+\frac{x^6}{180}+\cdots,\\
y_2&=\sum_{k\ge0}\frac{x^{3k+1}}{\prod_{j=1}^k(3j+1)(3j)}=x+\frac{x^4}{12}+\frac{x^7}{504}+\cdots.
\end{aligned}\]
空乘积为1. 两级数项比为 $|x|^3/[(3k+3)(3k+2)]$ 与 $|x|^3/[(3k+4)(3k+3)]$，都趋零，故各 $R=\infty$；官方解答未写这一步，在此补齐. 初值矩阵为 $I$，保证独立.\end{solution}
\begin{exercise}\label{cw:bank-6C-6}原题6C-6. 求 $(1-x^2)y''-2xy'+6y=0$ 的两个独立幂级数解并确定各收敛半径. 从原方程能预料半径到什么程度？\end{exercise}
\begin{solution}（官方译审并补证）比较 $x^n$ 系数得
\[(n+2)(n+1)a_{n+2}-[n(n-1)+2n-6]a_n=0,
\quad a_{n+2}=\frac{(n+3)(n-2)}{(n+2)(n+1)}a_n.
\]
取 $(a_0,a_1)=(1,0)$，偶链在 $n=2$ 截断，$y_1=1-3x^2$，因此 $R_1=\infty$. 取 $(0,1)$，奇链为
\[y_2=x-\frac23x^3-\frac15x^5-\frac4{35}x^7-\cdots,
\quad a_{2k+1}=\prod_{j=0}^{k-1}\frac{(2j+4)(2j-1)}{(2j+3)(2j+2)}.
\]
奇链无零项，相邻非零项之比趋 $|x|^2$，故 $R_2=1$. 两解0处Wronskian为1. 标准形系数 $-2x/(1-x^2),6/(1-x^2)$ 在 $x=\pm1$ 奇异，解析存在定理保证围绕0至少在 $|x|<1$ 有收敛级数，因而预料 $R\ge1$；不能据系数奇点断言每个解恰好 $R=1$，因为多项式解跨越奇点仍解析. 官方只算奇链半径，这里分别补全两解半径.\end{solution}
\begin{exercise}\label{cw:bank-6C-7}原题6C-7. 三项系数递推常使一般项难求. 对 $y''+2y'+(x-1)y=0$，给递推关系以及两个独立幂级数解的前三个非零项.\end{exercise}
\begin{solution}（官方译审）置 $a_{-1}=0$，对 $n\ge0$ 有
\[(n+2)(n+1)a_{n+2}+2(n+1)a_{n+1}-a_n+a_{n-1}=0.
\]
取 $(a_0,a_1)=(1,0)$ 得 $a_2=1/2$, $a_3=-1/2$，所以 $y_1=1+x^2/2-x^3/2+O(x^4)$. 取 $(0,1)$ 得 $a_2=-1$, $a_3=5/6$，所以 $y_2=x-x^2+5x^3/6+O(x^4)$. 初值矩阵为 $I$，故独立. 系数为全平面解析的多项式，解析线性方程理论保证两解级数具有无穷收敛半径，虽然本题只要求前三项.\end{solution}
\originalchapter{原题题库七：Fourier级数}
原练习注明依据 Edwards 与 Penney 的 \textit{Elementary Differential Equations} 第8章练习；本题库所译来源为MIT官方发布的 Notes and Exercises 文件，保留其原题号与署名关系.
\originalsection{7A：Fourier级数与正交性}
\begin{exercise}\label{cw:bank-7A-1}原题7A-1. 求下列周期函数的最小正周期：(a) $\sin(\pi t/3)$；(b) $|\sin t|$；(c) $\cos^2(3t)$.\end{exercise}
\begin{solution}（官方译审）(a) 正弦角速度为 $\pi/3$，故周期 $2\pi/(\pi/3)=6$. (b) 绝对值把正负半波合并，$|\sin(t+\pi)|=|\sin t|$，连续相邻零点间距为 $\pi$，最小周期为 $\pi$. (c) $\cos^2(3t)=(1+\cos6t)/2$，最小周期为 $\pi/3$.\end{solution}
\begin{exercise}\label{cw:bank-7A-2}原题7A-2. 求周期 $2\pi$ 的函数的Fourier级数，定义于 $-\pi<t\le\pi$：(a) $f(t)=0$（$-\pi<t\le0$），$f(t)=1$（$0<t\le\pi$）；(b) $f(t)=-t$（$-\pi<t<0$），$f(t)=t$（$0\le t\le\pi$）.\end{exercise}
\begin{solution}（官方译审）采用 $f\sim a_0/2+\sum_{n\ge1}(a_n\cos nt+b_n\sin nt)$. (a) $a_0=1$, $a_n=\pi^{-1}\int_0^\pi\cos nt\dd t=0$, $b_n=(1-(-1)^n)/(n\pi)$，所以
\[f\sim\frac12+\frac2\pi\sum_{\substack{n\ge1\\n\text{奇}}}\frac{\sin nt}{n}.\]
在跳跃点级数取左右平均 $1/2$，不能把它写成该点所指定函数值. (b) $f=|t|$ 为偶函数，$b_n=0$, $a_0=\pi$；分部积分给 $a_n=2[(-1)^n-1]/(\pi n^2)$，故
\[f=\frac\pi2-\frac4\pi\sum_{\substack{n\ge1\\n\text{奇}}}\frac{\cos nt}{n^2}.\]
周期延拓连续，级数在所有点均取函数值.\end{solution}
\begin{exercise}\label{cw:bank-7A-3}原题7A-3. 用 $\cos A\cos B=[\cos(A+B)+\cos(A-B)]/2$ 再证 $\int_{-\pi}^{\pi}\cos mt\cos nt\dd t=0$（$m\ne n$），等于 $\pi$（$m=n$）的正交关系.\end{exercise}
\begin{solution}（官方译审）这里 $m,n$ 为正整数. 若不同，$m+n,m-n$ 均为非零整数，其余弦在对称整周期区间积分为零，积化和差给结果0. 若相同，积化和差为 $(1+\cos2nt)/2$，积分为 $(2\pi+0)/2=\pi$. 若把0频率也纳入，$m=n=0$ 的积分是 $2\pi$，原式的 $\pi$ 结论只适用于正频率.\end{solution}
\begin{exercise}\label{cw:bank-7A-4}原题7A-4. $f$ 为周期 $P$ 的函数，证明它在任何长度 $P$ 的区间上积分相同：(a) 用换元证明 $\int_P^{a+P}f(t)\dd t=\int_0^af(t)\dd t$；(b) 由(a)推出 $\int_a^{a+P}f=\int_0^Pf$.\end{exercise}
\begin{solution}（官方译审）假设 $f$ 在一个周期上可积，周期延拓可积. (a) 置 $t=u+P$，利用 $f(u+P)=f(u)$ 即得，任意实数 $a$ 均按有向积分成立. (b) 积分可加性给
\[\int_a^{a+P}f=\int_0^Pf+\int_P^{a+P}f-\int_0^af=\int_0^Pf.\]
公式不要求 $0\le a\le P$，亦无需先对 $a$ 求导.\end{solution}
\originalsection{7B：奇偶展开与边值问题}
\begin{exercise}\label{cw:bank-7B-1}原题7B-1. (a) 求 $1-t$ 在 $0<t<1$ 的Fourier余弦级数，并画级数和在 $[-2,2]$ 的图；(b) 对其Fourier正弦级数回答同样问题.\end{exercise}
\begin{solution}（官方译审）(a) 偶延拓、周期2，$a_0=2\int_0^1(1-t)\dd t=1$；分部积分给 $a_n=2(1-(-1)^n)/(n^2\pi^2)$. 因此
\[f_c(t)=\frac12+\frac4{\pi^2}\sum_{\substack{n\ge1\\n\text{奇}}}\frac{\cos(n\pi t)}{n^2}.\]
在 $[-1,1]$ 和为 $1-|t|$，再以周期2延拓，连续三角波. (b) 奇延拓、周期2，$b_n=2\int_0^1(1-t)\sin(n\pi t)\dd t=2/(n\pi)$，故
\[f_s(t)=\frac2\pi\sum_{n\ge1}\frac{\sin(n\pi t)}n.\]
在 $(0,2)$ 其和为 $1-t$，再周期延拓；在偶整数处左右极限为 $-1,1$，级数和为0，奇整数处连续等于0. 图中空心点给侧极限，实心点给跳跃处的级数和值.
\begin{center}
\begin{tikzpicture}\begin{axis}[width=.47\linewidth,height=4.5cm,axis lines=middle,xlabel={$t$},ylabel={$f_c$},xmin=-2.2,xmax=2.2,ymin=-.2,ymax=1.3,title={余弦级数的和},xtick={-2,-1,0,1,2}]\addplot[structurecolor,thick] coordinates {(-2,1)(-1,0)(0,1)(1,0)(2,1)};\end{axis}\end{tikzpicture}
\hfill
\begin{tikzpicture}\begin{axis}[width=.47\linewidth,height=4.5cm,axis lines=middle,xlabel={$t$},ylabel={$f_s$},xmin=-2.2,xmax=2.2,ymin=-1.3,ymax=1.3,title={正弦级数的和},xtick={-2,-1,0,1,2}]\addplot[structurecolor,thick] coordinates {(-2,1)(0,-1)};\addplot[structurecolor,thick] coordinates {(0,1)(2,-1)};\addplot[only marks,mark=o,structurecolor] coordinates {(-2,1)(0,-1)(0,1)(2,-1)};\addplot[only marks,mark=*,structurecolor] coordinates {(-2,0)(0,0)(2,0)};\end{axis}\end{tikzpicture}
\end{center}
\end{solution}
\begin{exercise}\label{cw:bank-7B-2}原题7B-2. 用Fourier级数求下列边值问题的形式解，可用已知展开：(a) $x''+2x=1$, $x(0)=x(\pi)=0$；(b) $x''+2x=t$, $x'(0)=x'(\pi)=0$，使用余弦级数.\end{exercise}
\begin{solution}（官方译审）(a) 在 $(0,\pi)$，$1=(4/\pi)\sum_{n\text{奇}}\sin nt/n$. 写 $x=\sum X_n\sin nt$，比较得 $(2-n^2)X_n=4/(\pi n)$ 对奇 $n$，偶 $n$ 为零. 故
\[x=\frac4\pi\sum_{\substack{n\ge1\\n\text{奇}}}\frac{\sin nt}{n(2-n^2)}.\]
(b) 余弦展开 $t=\pi/2-(4/\pi)\sum_{n\text{奇}}\cos nt/n^2$. 常数项除以2，其他项除以 $2-n^2$，得
\[x=\frac\pi4-\frac4\pi\sum_{\substack{n\ge1\\n\text{奇}}}\frac{\cos nt}{n^2(2-n^2)}.\]
$a_0/2$ 的记号不能使常数项多一倍. 因 $2\ne n^2$，两边值问题无整数模态共振. (a)函数级数和一阶导数绝对收敛，端值为零；(b)一阶导数绝对收敛且在两端为零. 二阶导数在开区间给相应强迫的Fourier展开，可按局部收敛解释原方程. 题目要求形式解，未把端点强迫的奇延拓值混同其开区间值.\end{solution}
\begin{exercise}\label{cw:bank-7B-3}原题7B-3. 设 $a>0$，证明 $\int_{-a}^0f(t)\dd t=\pm\int_0^af(t)\dd t$，其中偶函数取正号、奇函数取负号.\end{exercise}
\begin{solution}（官方译审）令 $t=-u$，得 $\int_{-a}^0f(t)\dd t=\int_0^af(-u)\dd u$. 若偶，$f(-u)=f(u)$；若奇，$f(-u)=-f(u)$，分别得到所给两式. 假设积分存在即可.\end{solution}
\begin{exercise}\label{cw:bank-7B-4}原题7B-4. 周期2且在 $0<t<2$ 等于 $t^2$ 的函数有Fourier级数
\[f(t)\sim\frac43+\frac4{\pi^2}\sum_{n\ge1}\frac{\cos(n\pi t)}{n^2}-\frac4\pi\sum_{n\ge1}\frac{\sin(n\pi t)}n.\]
逐项微分，并证明得到的级数不收敛到 $f'(t)$.\end{exercise}
\begin{solution}（官方译审并补证）形式微分给
\[-\frac4\pi\sum_{n\ge1}\frac{\sin(n\pi t)}n-4\sum_{n\ge1}\cos(n\pi t).
\]
不能合法逐项求导，因为原周期函数在偶整数处跳跃. 更强地，按同阶项结合后的第 $n$ 项为 $-4\sin(n\pi t)/(n\pi)-4\cos(n\pi t)$；第一部分趋零，而对任意固定实数 $t$，$\cos(n\pi t)$ 都不趋零. 若它趋零，则其子列 $\cos(2n\pi t)$ 也趋零，但恒等式 $\cos(2n\pi t)=2\cos^2(n\pi t)-1$ 使它趋 $-1$，矛盾. 所以形式导数级数在每一点都不满足项趋零条件，均不按通常意义收敛，更不能等于开区间上的 $2t$.\end{solution}
\originalsection{7C：共振与谐波响应}
\begin{exercise}\label{cw:bank-7C-1}原题7C-1. 不计算完整解，判断各系统是否出现纯共振：(a) $2x''+10x=F(t)$，$F=1$ 于 $(0,1)$，奇延拓且周期2；(b) $x''+4\pi^2x=F(t)$，$F=2t$ 于 $(0,1)$，奇延拓且周期2；(c) $x''+9x=F(t)$，$F=1$ 于 $(0,\pi)$，奇延拓且周期 $2\pi$.\end{exercise}
\begin{solution}（官方译审）纯共振要求强迫中有非零Fourier系数的频率恰等于固有频率. (a) 固有频率 $\sqrt5$，强迫只有奇数倍 $\pi$，没有匹配，不共振. (b) 固有频率 $2\pi$，而 $F$ 的正弦系数 $b_n=4(-1)^{n+1}/(n\pi)$，$n=2$ 系数非零，出现纯共振. (c) 固有频率3，方波强迫有 $b_n=4/(\pi n)$ 对奇 $n$；$n=3$ 非零，出现纯共振. 有匹配頻率且系数非零才产生世俗增长项，不能只看名义周期.\end{solution}
\begin{exercise}\label{cw:bank-7C-2}原题7C-2. 求 $x''+3x=F(t)$ 的一周期解，以Fourier级数表示；$F(t)=2t$ 于 $(0,\pi)$，为奇函数，周期 $2\pi$.\end{exercise}
\begin{solution}（官方译审）分部积分得 $b_n=(2/\pi)\int_0^\pi2t\sin nt\dd t=4(-1)^{n+1}/n$. 令 $x=\sum X_n\sin nt$，则 $(3-n^2)X_n=b_n$，$3$ 不是整数平方，无共振，故
\[x_p(t)=4\sum_{n\ge1}\frac{(-1)^{n+1}\sin nt}{n(3-n^2)}.
\]
系数为 $O(n^{-3})$，函数及一阶导数级数绝对一致收敛，确为 $2\pi$ 周期；在强迫连续的开区间，二阶导数Fourier级数给原方程. 加一般齐次项 $A\cos\sqrt3t+B\sin\sqrt3t$ 通常破坏这一周期，所以此处只给所求周期特解.\end{solution}
\begin{exercise}\label{cw:bank-7C-3}原题7C-3. 根据Fourier响应和无阻尼固有频率，判以下轻阻尼系统中哪个谐波振幅最大：(a) $2x''+0.1x'+18x=F(t)$，$F$ 与7C-2相同；(b) $3x''+x'+30x=F(t)$，$F=t-t^2$ 于 $(0,1)$，奇延拓且周期2.\end{exercise}
\begin{solution}（官方译审并补证）对于频率 $\omega$、强迫振幅 $|b|$，稳态响应振幅为
\[A=\frac{|b|}{\sqrt{(k-m\omega^2)^2+c^2\omega^2}}.
\]
(a) $\omega_0=3$，$b_n=4(-1)^{n+1}/n$，所以 $A_n=4/[n\sqrt{(18-2n^2)^2+0.01n^2}]$. $n=3$ 时 $A_3=40/9\approx4.44$. 对 $n=1,2$，分别小于 $1/4,1/5$；对 $n\ge4$，$n(2n^2-18)\ge56$，故 $A_n\le1/14$. 第3谐波明确最大，非仅凭固有频率猜测.

(b) $\omega_0=\sqrt{10}\approx\pi$，频率为 $n\pi$. 两次分部积分得
\[b_n=2\int_0^1(t-t^2)\sin(n\pi t)\dd t=\frac{4(1-(-1)^n)}{n^3\pi^3}.
\]
偶谐波消失，奇谐波不为零. $n=1$ 恰好很靠近固有频率，振幅
\[A_1=\frac8{\pi^3\sqrt{(30-3\pi^2)^2+\pi^2}}\approx0.0815.
\]
对奇 $n\ge3$，$|30-3n^2\pi^2|\ge27\pi^2-30>236$，且 $|b_n|\le8/(27\pi^3)$，故 $A_n<0.000041$；偶谐波为零. 因此第1谐波最大. 不计阻尼的共振频率仅提供候选，系数和响应分母共同完成比较.\end{solution}

% MIT 18.03 Spring 2010 PS1--PS5；源和逐叶覆盖见 research/ps01-05-coverage.json。
\originalchapter{作业一：一阶方程}
\label{cw:ps1}
本章保留官方作业的讲次编号，因而从第0题开始。第一部分中的 EP 指 Edwards--Penney 教材；官方作业只给出节号与题号，没有给出书题题面。本书对这些条目只列定位，未翻译、未解答，也不把它们计入已收录题面。Notes 指课程公开的习题集，下面的链接回指本书相应中文题目与解答；字母表示只指定该题的相应小问。第二部分的题面来自官方作业，解答依据官方答案审译，并在必要处补证或修正。涉及 Mathlet 的图均按方程独立重绘，图示不是实际运行交互程序的记录。
\originalsection{第一部分：指定练习}
\begin{exercise}\label{cw:ps1-I-0}
原题 I.0（自然增长与可分离方程）：Notes 1A-5(c)，见\hyperref[cw:bank-1A-5]{题1A-5(c)}；EP 第1.1节第32、33、35题，第1.4节第39、66题，第1.5节第1、9、20题。
\end{exercise}
\begin{solution}
公开 Notes 题的中文题面与完整解答见上述链接。EP 八道题只有书题定位，本章没有其合法可用的完整题面，因此不代拟题目、不声称已翻译或解出这些书题。
\end{solution}
\begin{exercise}\label{cw:ps1-I-1}
原题 I.1（方向场与等倾线）：Notes 1C-1(a)、(b)、(e)，见\hyperref[cw:bank-1C-1]{题1C-1}。
\end{exercise}
\begin{solution}按所列三个小问阅读链接中的题面与解答；(c)、(d)并非本次作业指定。\end{solution}
\begin{exercise}\label{cw:ps1-I-2}
原题 I.2（欧拉法）：Notes 1C-4，见\hyperref[cw:bank-1C-4]{题1C-4}。
\end{exercise}
\begin{solution}题面、欧拉迭代表与计算见链接。\end{solution}
\begin{exercise}\label{cw:ps1-I-3}
原题 I.3（线性模型）：EP 第1.5节第33、45题。两题均须明确写出微分方程。
\end{exercise}
\begin{solution}仅列教材定位；书题题面未公开收入本次材料，未翻译、未解答。\end{solution}
\originalsection{第二部分：题面与解答}
\begin{exercise}\label{cw:ps1-II-0}
原题 II.0。羚羊种群在短时间 $\Delta t$ 内的变化近似为 $k(t)x(t)\Delta t$。不考虑捕猎，时间以年计，种群以千只羚羊（ko）计。肯尼亚塔纳河地区出现一种病毒，使自然增长率变为
\[
 k(t)=\frac{k_0}{(a+t)^2},\qquad t\geq0,\quad a,k_0>0.
\]
(a) 常数 $a,k_0$ 的单位是什么？(b) 写出模型方程。(c) 求通解，考虑数学上所有允许的 $t,x$，正确处理 $\ln|x|$ 和分离变量时丢失的解。(d) 若 $x(0)=x_0>0$，种群是否在大时间趋于稳定？若是，求极限。
\end{exercise}
\begin{solution}
解答依据（编者译审或补解，AI 辅助）：官方 PS1 解答第1页，补充定义区间。由 $x'=kx$ 可知 $k$ 的单位是年$^{-1}$，所以 $a$ 的单位是年，而 $k_0$ 的单位是年，保证指数无量纲。短时间增量关系除以 $\Delta t$ 后取极限得
\[
 x'=\frac{k_0x}{(a+t)^2}.
\]
在 $x\ne0$ 的区间分离变量，积分得 $\ln|x|=-k_0/(a+t)+c$，于是 $x=Ce^{-k_0/(a+t)}$，其中 $C$ 可正可负；把 $C=0$ 纳入便补回零解。方程在 $t=-a$ 无定义，通解应分别理解为 $(-\infty,-a)$ 或 $(-a,\infty)$ 内的解，不能跨越奇点。在种群模型的 $t\geq0$ 上，由初值得
\[
 x(t)=x_0\exp\left(\frac{k_0}{a}-\frac{k_0}{a+t}\right),
 \qquad\lim_{t\to\infty}x(t)=x_0e^{k_0/a}.
\]
增长率虽始终为正，其从0到无穷的积分却有限，所以种群增长到有限极限。
\end{solution}
\begin{exercise}\label{cw:ps1-II-1}
原题 II.1。研究 $y'=y^2-x$。
(a) 在横、纵轴均标出 $-4$ 至 $4$ 的大图上，画斜率 $m=-1,0,1,2$ 的等倾线，并在同一图上依据它们描出若干解。原作业随后要求用 Isoclines Mathlet 选此方程，再画 $m=-2,0,1,2$ 的等倾线，比较手绘图并拖动初值观察解。
(b) 描出分界解：其上方的解向右增长无界，其下方的解最终下降。
(c) 解曲线在 $(a,b)$ 与 $m=-1$ 的等倾线相切，求 $(a,b)$。
(d) 若 $y(a)<b$，给出解在 $x\to\infty$ 时所渐近的显式函数 $f(x)$，判断解最终位于其上还是其下。原题要求说明：(i) 解不能在 $x>a$ 穿越 $m=-1$ 的等倾线；(ii) 若在 $x=a$ 低于 $m=1$ 的等倾线，最终必穿越它；(iii) 解最终必穿越零倾线。
(e) 若 $(c,d)$ 是临界点，$c,d$ 满足什么关系？(f) 所有临界点都是局部极大值吗？须验证。
\end{exercise}
\begin{solution}
解答依据（编者译审或补解，AI 辅助）：官方 PS1 解答第1--2页；(d)补充渐近收敛证明，并说明原题(iii)的条件缺漏。
等倾线为 $x=y^2-m$。沿线画斜率为 $m$ 的短线段便得到方向场；零倾线 $x=y^2$ 外侧 $y^2>x$ 的区域斜率为正，内侧斜率为负。下图直接由这些等式重绘，各条线段只表示局部斜率；灰色解曲线由 $y(0)=0.5,1.4$ 的数值积分叠加，供方向场描线的比较。
\begin{center}
\begin{tikzpicture}
\begin{axis}[width=.78\linewidth,height=6cm,xmin=-4,xmax=4,ymin=-4,ymax=4,domain=-4:4,axis lines=middle,xlabel={$x$},ylabel={$y$},samples=100,legend style={at={(1.02,1)},anchor=north west,font=\small}]
\addplot[domain=-2:2,blue] ({x*x+1},x);\addlegendentry{$m=-1$}
\addplot[domain=-2:2,black] ({x*x},x);\addlegendentry{$m=0$}
\addplot[domain=-2.2:2.2,red] ({x*x-1},x);\addlegendentry{$m=1$}
\addplot[domain=-2.4:2.4,orange] ({x*x-2},x);\addlegendentry{$m=2$}
\addplot[gray,thick,forget plot]coordinates{(0.000000,0.500000)(0.054054,0.512401)(0.108108,0.522503)(0.162162,0.530185)(0.216216,0.535313)(0.270270,0.537738)(0.324324,0.537300)(0.378378,0.533828)(0.432432,0.527141)(0.486486,0.517057)(0.540541,0.503389)(0.594595,0.485955)(0.648649,0.464583)(0.702703,0.439116)(0.756757,0.409422)(0.810811,0.375399)(0.864865,0.336990)(0.918919,0.294184)(0.972973,0.247030)(1.027027,0.195644)(1.081081,0.140213)(1.135135,0.080995)(1.189189,0.018328)(1.243243,-0.047383)(1.297297,-0.115667)(1.351351,-0.186003)(1.405405,-0.257827)(1.459459,-0.330555)(1.513514,-0.403598)(1.567568,-0.476380)(1.621622,-0.548356)(1.675676,-0.619026)(1.729730,-0.687947)(1.783784,-0.754746)(1.837838,-0.819117)(1.891892,-0.880831)(1.945946,-0.939727)(2.000000,-0.995714)(2.054054,-1.048758)(2.108108,-1.098881)(2.162162,-1.146145)(2.216216,-1.190650)(2.270270,-1.232522)(2.324324,-1.271907)(2.378378,-1.308963)(2.432432,-1.343854)(2.486486,-1.376747)(2.540541,-1.407806)(2.594595,-1.437191)(2.648649,-1.465053)(2.702703,-1.491536)(2.756757,-1.516774)(2.810811,-1.540889)(2.864865,-1.563996)(2.918919,-1.586196)(2.972973,-1.607582)(3.027027,-1.628239)(3.081081,-1.648240)(3.135135,-1.667653)(3.189189,-1.686536)(3.243243,-1.704943)(3.297297,-1.722920)(3.351351,-1.740507)(3.405405,-1.757740)(3.459459,-1.774652)(3.513514,-1.791270)(3.567568,-1.807618)(3.621622,-1.823718)(3.675676,-1.839589)(3.729730,-1.855246)(3.783784,-1.870706)(3.837838,-1.885980)(3.891892,-1.901079)(3.945946,-1.916014)(4.000000,-1.930794)};
\addplot[gray,thick,forget plot]coordinates{(0.000000,1.400000)(0.009970,1.419769)(0.019941,1.440004)(0.029911,1.460728)(0.039882,1.481964)(0.049852,1.503735)(0.059823,1.526069)(0.069793,1.548992)(0.079763,1.572533)(0.089734,1.596723)(0.099704,1.621595)(0.109675,1.647183)(0.119645,1.673524)(0.129616,1.700658)(0.139586,1.728628)(0.149556,1.757477)(0.159527,1.787254)(0.169497,1.818010)(0.179468,1.849801)(0.189438,1.882684)(0.199409,1.916725)(0.209379,1.951991)(0.219349,1.988556)(0.229320,2.026498)(0.239290,2.065903)(0.249261,2.106865)(0.259231,2.149483)(0.269202,2.193866)(0.279172,2.240132)(0.289143,2.288411)(0.299113,2.338842)(0.309083,2.391580)(0.319054,2.446792)(0.329024,2.504664)(0.338995,2.565398)(0.348965,2.629219)(0.358936,2.696373)(0.368906,2.767136)(0.378876,2.841812)(0.388847,2.920741)(0.398817,3.004302)(0.408788,3.092921)(0.418758,3.187077)(0.428729,3.287311)(0.438699,3.394234)(0.448669,3.508546)(0.458640,3.631041)(0.468610,3.762636)(0.478581,3.904386)};

\foreach \xx in {-3,-2,-1,0,1,2,3}{\foreach \yy in {-3,-2,-1,0,1,2,3}{\pgfmathsetmacro{\mm}{\yy*\yy-\xx}\pgfmathsetmacro{\dx}{.18/sqrt(1+\mm*\mm)}\pgfmathsetmacro{\dy}{\mm*\dx}\edef\fieldsegment{\noexpand\draw[gray] (axis cs:{\xx-\dx},{\yy-\dy})--(axis cs:{\xx+\dx},{\yy+\dy});}\fieldsegment}}
\end{axis}
\end{tikzpicture}
\end{center}
(b) 分界解的数值示意另见下图；它向右沿正支 $\sqrt x$ 上升，而较低的解转向负支。数值图用于定位分界，不作为所有解分类的证明。可用代换 $y=-u'/u$ 得到线性方程 $u''=xu$；取向右衰减的非零解 $u$ 所得到的曲线正是此分界解。因题目要求描图，下面采用从 $x=16$ 的近似渐近初值向左积分的曲线，并明确保留其数值性质。
% Numerical coordinates are produced from the equation, not from the Mathlet.
\begin{center}\begin{tikzpicture}\begin{axis}[width=.75\linewidth,height=5cm,axis lines=middle,xlabel={$x$},ylabel={$y$},xmin=-3,xmax=4,ymin=-3,ymax=4,domain=-3:4,samples=100]
\addplot[blue,thick]coordinates{(-2.026816,-2.990891)(-1.992179,-2.646832)(-1.957542,-2.361705)(-1.922905,-2.120841)(-1.888268,-1.914086)(-1.853631,-1.734187)(-1.818994,-1.575824)(-1.784358,-1.435004)(-1.749721,-1.308673)(-1.715084,-1.194455)(-1.680447,-1.090471)(-1.645810,-0.995220)(-1.611173,-0.907481)(-1.576536,-0.826258)(-1.541899,-0.750725)(-1.507263,-0.680193)(-1.472626,-0.614084)(-1.437989,-0.551905)(-1.403352,-0.493239)(-1.368715,-0.437726)(-1.334078,-0.385054)(-1.299441,-0.334954)(-1.264804,-0.287191)(-1.230168,-0.241558)(-1.195531,-0.197874)(-1.160894,-0.155977)(-1.126257,-0.115724)(-1.091620,-0.076989)(-1.056983,-0.039657)(-1.022346,-0.003627)(-0.987709,0.031195)(-0.953073,0.064890)(-0.918436,0.097533)(-0.883799,0.129194)(-0.849162,0.159933)(-0.814525,0.189808)(-0.779888,0.218871)(-0.745251,0.247168)(-0.710615,0.274743)(-0.675978,0.301636)(-0.641341,0.327885)(-0.606704,0.353523)(-0.572067,0.378582)(-0.537430,0.403090)(-0.502793,0.427076)(-0.468156,0.450564)(-0.433520,0.473577)(-0.398883,0.496139)(-0.364246,0.518268)(-0.329609,0.539985)(-0.294972,0.561306)(-0.260335,0.582249)(-0.225698,0.602830)(-0.191061,0.623064)(-0.156425,0.642963)(-0.121788,0.662542)(-0.087151,0.681812)(-0.052514,0.700786)(-0.017877,0.719474)(0.016760,0.737887)(0.051397,0.756034)(0.086034,0.773925)(0.120670,0.791569)(0.155307,0.808975)(0.189944,0.826149)(0.224581,0.843100)(0.259218,0.859835)(0.293855,0.876360)(0.328492,0.892683)(0.363128,0.908809)(0.397765,0.924745)(0.432402,0.940496)(0.467039,0.956068)(0.501676,0.971465)(0.536313,0.986693)(0.570950,1.001757)(0.605587,1.016660)(0.640223,1.031408)(0.674860,1.046004)(0.709497,1.060453)(0.744134,1.074759)(0.778771,1.088924)(0.813408,1.102953)(0.848045,1.116849)(0.882682,1.130616)(0.917318,1.144256)(0.951955,1.157772)(0.986592,1.171168)(1.021229,1.184446)(1.055866,1.197610)(1.090503,1.210661)(1.125140,1.223602)(1.159777,1.236436)(1.194413,1.249164)(1.229050,1.261790)(1.263687,1.274316)(1.298324,1.286743)(1.332961,1.299074)(1.367598,1.311311)(1.402235,1.323455)(1.436872,1.335508)(1.471508,1.347473)(1.506145,1.359351)(1.540782,1.371144)(1.575419,1.382853)(1.610056,1.394480)(1.644693,1.406027)(1.679330,1.417495)(1.713966,1.428885)(1.748603,1.440200)(1.783240,1.451439)(1.817877,1.462606)(1.852514,1.473700)(1.887151,1.484724)(1.921788,1.495678)(1.956425,1.506564)(1.991061,1.517382)(2.025698,1.528135)(2.060335,1.538823)(2.094972,1.549447)(2.129609,1.560008)(2.164246,1.570507)(2.198883,1.580946)(2.233520,1.591325)(2.268156,1.601645)(2.302793,1.611907)(2.337430,1.622112)(2.372067,1.632261)(2.406704,1.642355)(2.441341,1.652394)(2.475978,1.662380)(2.510615,1.672313)(2.545251,1.682194)(2.579888,1.692023)(2.614525,1.701802)(2.649162,1.711531)(2.683799,1.721211)(2.718436,1.730843)(2.753073,1.740426)(2.787709,1.749963)(2.822346,1.759454)(2.856983,1.768898)(2.891620,1.778297)(2.926257,1.787652)(2.960894,1.796963)(2.995531,1.806230)(3.030168,1.815455)(3.064804,1.824637)(3.099441,1.833778)(3.134078,1.842877)(3.168715,1.851936)(3.203352,1.860954)(3.237989,1.869933)(3.272626,1.878873)(3.307263,1.887774)(3.341899,1.896637)(3.376536,1.905462)(3.411173,1.914250)(3.445810,1.923001)(3.480447,1.931715)(3.515084,1.940394)(3.549721,1.949037)(3.584358,1.957645)(3.618994,1.966218)(3.653631,1.974757)(3.688268,1.983261)(3.722905,1.991733)(3.757542,2.000171)(3.792179,2.008576)(3.826816,2.016949)(3.861453,2.025290)(3.896089,2.033599)(3.930726,2.041877)(3.965363,2.050123)(4.000000,2.058339)};
\addplot[gray]coordinates{(0.000000,-1.000000)(0.036697,-0.965259)(0.073394,-0.934197)(0.110092,-0.906495)(0.146789,-0.881881)(0.183486,-0.860114)(0.220183,-0.840985)(0.256881,-0.824307)(0.293578,-0.809916)(0.330275,-0.797664)(0.366972,-0.787416)(0.403670,-0.779053)(0.440367,-0.772463)(0.477064,-0.767546)(0.513761,-0.764208)(0.550459,-0.762362)(0.587156,-0.761926)(0.623853,-0.762823)(0.660550,-0.764982)(0.697248,-0.768331)(0.733945,-0.772807)(0.770642,-0.778345)(0.807339,-0.784884)(0.844037,-0.792366)(0.880734,-0.800733)(0.917431,-0.809930)(0.954128,-0.819904)(0.990826,-0.830602)(1.027523,-0.841973)(1.064220,-0.853969)(1.100917,-0.866542)(1.137615,-0.879644)(1.174312,-0.893231)(1.211009,-0.907259)(1.247706,-0.921686)(1.284404,-0.936471)(1.321101,-0.951576)(1.357798,-0.966962)(1.394495,-0.982594)(1.431193,-0.998438)(1.467890,-1.014460)(1.504587,-1.030630)(1.541284,-1.046919)(1.577982,-1.063300)(1.614679,-1.079746)(1.651376,-1.096234)(1.688073,-1.112740)(1.724771,-1.129246)(1.761468,-1.145731)(1.798165,-1.162177)(1.834862,-1.178570)(1.871560,-1.194895)(1.908257,-1.211138)(1.944954,-1.227288)(1.981651,-1.243334)(2.018349,-1.259268)(2.055046,-1.275081)(2.091743,-1.290768)(2.128440,-1.306321)(2.165138,-1.321736)(2.201835,-1.337010)(2.238532,-1.352138)(2.275229,-1.367119)(2.311927,-1.381952)(2.348624,-1.396634)(2.385321,-1.411165)(2.422018,-1.425546)(2.458716,-1.439777)(2.495413,-1.453859)(2.532110,-1.467793)(2.568807,-1.481580)(2.605505,-1.495223)(2.642202,-1.508723)(2.678899,-1.522083)(2.715596,-1.535305)(2.752294,-1.548391)(2.788991,-1.561345)(2.825688,-1.574168)(2.862385,-1.586864)(2.899083,-1.599435)(2.935780,-1.611884)(2.972477,-1.624214)(3.009174,-1.636428)(3.045872,-1.648529)(3.082569,-1.660519)(3.119266,-1.672401)(3.155963,-1.684178)(3.192661,-1.695852)(3.229358,-1.707426)(3.266055,-1.718903)(3.302752,-1.730284)(3.339450,-1.741573)(3.376147,-1.752771)(3.412844,-1.763881)(3.449541,-1.774905)(3.486239,-1.785845)(3.522936,-1.796704)(3.559633,-1.807482)(3.596330,-1.818183)(3.633028,-1.828807)(3.669725,-1.839357)(3.706422,-1.849834)(3.743119,-1.860241)(3.779817,-1.870578)(3.816514,-1.880847)(3.853211,-1.891050)(3.889908,-1.901188)(3.926606,-1.911263)(3.963303,-1.921275)(4.000000,-1.931227)};
\addplot[gray]coordinates{(0.000000,0.000000)(0.036697,-0.000673)(0.073394,-0.002693)(0.110092,-0.006059)(0.146789,-0.010770)(0.183486,-0.016823)(0.220183,-0.024215)(0.256881,-0.032938)(0.293578,-0.042985)(0.330275,-0.054345)(0.366972,-0.067004)(0.403670,-0.080943)(0.440367,-0.096142)(0.477064,-0.112576)(0.513761,-0.130216)(0.550459,-0.149027)(0.587156,-0.168973)(0.623853,-0.190011)(0.660550,-0.212094)(0.697248,-0.235172)(0.733945,-0.259190)(0.770642,-0.284089)(0.807339,-0.309806)(0.844037,-0.336276)(0.880734,-0.363431)(0.917431,-0.391199)(0.954128,-0.419509)(0.990826,-0.448286)(1.027523,-0.477456)(1.064220,-0.506944)(1.100917,-0.536677)(1.137615,-0.566582)(1.174312,-0.596588)(1.211009,-0.626626)(1.247706,-0.656629)(1.284404,-0.686535)(1.321101,-0.716285)(1.357798,-0.745823)(1.394495,-0.775098)(1.431193,-0.804062)(1.467890,-0.832675)(1.504587,-0.860898)(1.541284,-0.888697)(1.577982,-0.916045)(1.614679,-0.942917)(1.651376,-0.969293)(1.688073,-0.995158)(1.724771,-1.020500)(1.761468,-1.045310)(1.798165,-1.069585)(1.834862,-1.093321)(1.871560,-1.116522)(1.908257,-1.139190)(1.944954,-1.161332)(1.981651,-1.182956)(2.018349,-1.204071)(2.055046,-1.224689)(2.091743,-1.244822)(2.128440,-1.264485)(2.165138,-1.283691)(2.201835,-1.302455)(2.238532,-1.320793)(2.275229,-1.338720)(2.311927,-1.356252)(2.348624,-1.373405)(2.385321,-1.390194)(2.422018,-1.406634)(2.458716,-1.422742)(2.495413,-1.438531)(2.532110,-1.454016)(2.568807,-1.469210)(2.605505,-1.484128)(2.642202,-1.498783)(2.678899,-1.513186)(2.715596,-1.527349)(2.752294,-1.541285)(2.788991,-1.555003)(2.825688,-1.568514)(2.862385,-1.581828)(2.899083,-1.594954)(2.935780,-1.607901)(2.972477,-1.620677)(3.009174,-1.633290)(3.045872,-1.645746)(3.082569,-1.658054)(3.119266,-1.670220)(3.155963,-1.682250)(3.192661,-1.694149)(3.229358,-1.705923)(3.266055,-1.717577)(3.302752,-1.729116)(3.339450,-1.740544)(3.376147,-1.751866)(3.412844,-1.763086)(3.449541,-1.774207)(3.486239,-1.785232)(3.522936,-1.796166)(3.559633,-1.807011)(3.596330,-1.817770)(3.633028,-1.828446)(3.669725,-1.839042)(3.706422,-1.849559)(3.743119,-1.860001)(3.779817,-1.870368)(3.816514,-1.880665)(3.853211,-1.890891)(3.889908,-1.901050)(3.926606,-1.911143)(3.963303,-1.921171)(4.000000,-1.931136)};
\addplot[gray]coordinates{(0.000000,0.500000)(0.036697,0.508664)(0.073394,0.516285)(0.110092,0.522828)(0.146789,0.528254)(0.183486,0.532521)(0.220183,0.535585)(0.256881,0.537398)(0.293578,0.537910)(0.330275,0.537070)(0.366972,0.534822)(0.403670,0.531110)(0.440367,0.525878)(0.477064,0.519068)(0.513761,0.510620)(0.550459,0.500478)(0.587156,0.488585)(0.623853,0.474886)(0.660550,0.459333)(0.697248,0.441876)(0.733945,0.422476)(0.770642,0.401098)(0.807339,0.377716)(0.844037,0.352311)(0.880734,0.324878)(0.917431,0.295420)(0.954128,0.263957)(0.990826,0.230519)(1.027523,0.195154)(1.064220,0.157923)(1.100917,0.118905)(1.137615,0.078194)(1.174312,0.035898)(1.211009,-0.007856)(1.247706,-0.052930)(1.284404,-0.099173)(1.321101,-0.146420)(1.357798,-0.194502)(1.394495,-0.243238)(1.431193,-0.292446)(1.467890,-0.341941)(1.504587,-0.391539)(1.541284,-0.441059)(1.577982,-0.490326)(1.614679,-0.539175)(1.651376,-0.587448)(1.688073,-0.635003)(1.724771,-0.681707)(1.761468,-0.727447)(1.798165,-0.772120)(1.834862,-0.815641)(1.871560,-0.857941)(1.908257,-0.898964)(1.944954,-0.938672)(1.981651,-0.977039)(2.018349,-1.014050)(2.055046,-1.049704)(2.091743,-1.084012)(2.128440,-1.116990)(2.165138,-1.148666)(2.201835,-1.179072)(2.238532,-1.208248)(2.275229,-1.236237)(2.311927,-1.263086)(2.348624,-1.288844)(2.385321,-1.313562)(2.422018,-1.337293)(2.458716,-1.360087)(2.495413,-1.381998)(2.532110,-1.403076)(2.568807,-1.423371)(2.605505,-1.442932)(2.642202,-1.461805)(2.678899,-1.480035)(2.715596,-1.497664)(2.752294,-1.514734)(2.788991,-1.531282)(2.825688,-1.547345)(2.862385,-1.562956)(2.899083,-1.578148)(2.935780,-1.592950)(2.972477,-1.607389)(3.009174,-1.621492)(3.045872,-1.635282)(3.082569,-1.648782)(3.119266,-1.662011)(3.155963,-1.674989)(3.192661,-1.687732)(3.229358,-1.700258)(3.266055,-1.712580)(3.302752,-1.724712)(3.339450,-1.736666)(3.376147,-1.748454)(3.412844,-1.760086)(3.449541,-1.771572)(3.486239,-1.782920)(3.522936,-1.794138)(3.559633,-1.805235)(3.596330,-1.816215)(3.633028,-1.827086)(3.669725,-1.837852)(3.706422,-1.848520)(3.743119,-1.859094)(3.779817,-1.869578)(3.816514,-1.879976)(3.853211,-1.890291)(3.889908,-1.900528)(3.926606,-1.910689)(3.963303,-1.920776)(4.000000,-1.930794)};
\addplot[gray]coordinates{(0.000000,1.000000)(0.010159,1.010211)(0.020318,1.020530)(0.030477,1.030961)(0.040636,1.041508)(0.050795,1.052176)(0.060954,1.062971)(0.071113,1.073897)(0.081272,1.084960)(0.091431,1.096165)(0.101590,1.107518)(0.111749,1.119025)(0.121908,1.130692)(0.132067,1.142526)(0.142226,1.154533)(0.152385,1.166721)(0.162544,1.179097)(0.172703,1.191669)(0.182862,1.204444)(0.193021,1.217431)(0.203180,1.230639)(0.213339,1.244077)(0.223498,1.257755)(0.233657,1.271682)(0.243816,1.285868)(0.253975,1.300326)(0.264134,1.315067)(0.274293,1.330102)(0.284452,1.345444)(0.294611,1.361107)(0.304770,1.377104)(0.314929,1.393451)(0.325088,1.410162)(0.335247,1.427255)(0.345406,1.444745)(0.355565,1.462652)(0.365724,1.480995)(0.375883,1.499793)(0.386042,1.519068)(0.396200,1.538842)(0.406359,1.559140)(0.416518,1.579986)(0.426677,1.601407)(0.436836,1.623432)(0.446995,1.646090)(0.457154,1.669415)(0.467313,1.693439)(0.477472,1.718199)(0.487631,1.743734)(0.497790,1.770084)(0.507949,1.797295)(0.518108,1.825413)(0.528267,1.854488)(0.538426,1.884574)(0.548585,1.915729)(0.558744,1.948017)(0.568903,1.981502)(0.579062,2.016258)(0.589221,2.052362)(0.599380,2.089898)(0.609539,2.128958)(0.619698,2.169638)(0.629857,2.212047)(0.640016,2.256300)(0.650175,2.302524)(0.660334,2.350856)(0.670493,2.401448)(0.680652,2.454464)(0.690811,2.510085)(0.700970,2.568511)(0.711129,2.629962)(0.721288,2.694681)(0.731447,2.762937)(0.741606,2.835028)(0.751765,2.911289)(0.761924,2.992092)(0.772083,3.077855)(0.782242,3.169048)(0.792401,3.266200)(0.802560,3.369914)(0.812719,3.480875)(0.822878,3.599863)(0.833037,3.727777)(0.843196,3.865653)};
\addplot[gray]coordinates{(0.000000,1.400000)(0.006769,1.413371)(0.013538,1.426954)(0.020307,1.440756)(0.027076,1.454783)(0.033845,1.469043)(0.040614,1.483543)(0.047382,1.498291)(0.054151,1.513295)(0.060920,1.528563)(0.067689,1.544105)(0.074458,1.559928)(0.081227,1.576042)(0.087996,1.592458)(0.094765,1.609185)(0.101534,1.626235)(0.108303,1.643617)(0.115072,1.661345)(0.121841,1.679429)(0.128609,1.697883)(0.135378,1.716719)(0.142147,1.735953)(0.148916,1.755597)(0.155685,1.775667)(0.162454,1.796179)(0.169223,1.817150)(0.175992,1.838597)(0.182761,1.860537)(0.189530,1.882991)(0.196299,1.905979)(0.203068,1.929521)(0.209836,1.953639)(0.216605,1.978358)(0.223374,2.003701)(0.230143,2.029695)(0.236912,2.056366)(0.243681,2.083744)(0.250450,2.111859)(0.257219,2.140742)(0.263988,2.170429)(0.270757,2.200954)(0.277526,2.232357)(0.284295,2.264676)(0.291063,2.297955)(0.297832,2.332239)(0.304601,2.367576)(0.311370,2.404018)(0.318139,2.441619)(0.324908,2.480437)(0.331677,2.520534)(0.338446,2.561977)(0.345215,2.604835)(0.351984,2.649186)(0.358753,2.695109)(0.365522,2.742693)(0.372290,2.792030)(0.379059,2.843221)(0.385828,2.896374)(0.392597,2.951607)(0.399366,3.009044)(0.406135,3.068824)(0.412904,3.131092)(0.419673,3.196011)(0.426442,3.263753)(0.433211,3.334510)(0.439980,3.408487)(0.446749,3.485912)(0.453517,3.567032)(0.460286,3.652119)(0.467055,3.741473)(0.473824,3.835423)(0.480593,3.934334)};
\addplot[black,dashed,domain=0:4]{-sqrt(x)};\end{axis}\end{tikzpicture}\end{center}

(c) 在等倾线 $x=y^2+1$ 上，其切线满足 $1=2yy'$。相切时两条曲线斜率都为 $-1$，故 $b=-1/2$，$a=b^2+1=5/4$。

(d) 这里 $m=-1$ 线须取负支 $U(x)=-\sqrt{x-1}$。在 $x>a$，$U'(x)=-1/(2\sqrt{x-1})>-1$，而解在接触这条线时斜率为 $-1$。从下方第一次穿越必须使 $y'-U'\geq0$，与此矛盾，故 $y<U$，这证明(i)。解在任意有限的 $x$ 区间也不会逃向负无穷：若 $y$ 足够负，则 $y^2-x>0$；配合上界 $U$，得到可延拓到所有 $x\geq a$ 的解。

令 $L(x)=-\sqrt{x+1}$，即 $m=1$ 线的负支。若始终 $y<L$，则 $y'>1$，于是 $y(x)>y(a)+(x-a)$，最终与 $y<L<0$ 矛盾。因此必穿越 $L$，这证明(ii)。同理，若解还低于零倾线负支 $N(x)=-\sqrt x$，则 $y'>0$，而 $N\to-\infty$，故必穿越 $N$；在接触 $N$ 时 $y'-N'=1/(2\sqrt x)>0$，所以穿越后永远在其上。

原题(iii)未写出“起初低于零倾线”的条件，作为无条件陈述并不正确。例如取 $y(a)=-1$，此时 $-\sqrt{5/4}<-1<-1/2$，解从一开始就在 $N$ 上方，并永远不穿越它。正确结论是：每个满足本问初值的解最终都处于 $N$ 上方；初始位于其下方的解必穿越一次。最终
\[
 -\sqrt x<y(x)<-\sqrt{x-1},\qquad
 0<y(x)+\sqrt x<\frac1{\sqrt x+\sqrt{x-1}}\longrightarrow0.
\]
所以 $f(x)=-\sqrt x$，且最终从上方渐近。这一夹逼才证明曲线间的纵向距离趋于0。

(e) 在临界点 $0=y'(c)=d^2-c$，故 $c=d^2$。(f) 对方程求导得 $y''=2yy'-1$。临界点处 $y''=-1<0$，故每个临界点都是严格局部极大值；此结论不依赖图像猜测。
\end{solution}
\begin{exercise}\label{cw:ps1-II-2}
原题 II.2。(a) 设 $y'=2x$、$y(0)=0$。求 $y(1)$，并分别用2、3、$n$个等长步的欧拉法估计它。可用 $1+\cdots+(n-1)=n(n-1)/2$。估计是否收敛？(b) 步长 $h=1/n$ 时的精确误差是多少？是否符合误差至多正比于 $h$ 的预测？
\end{exercise}
\begin{solution}
解答依据（编者译审或补解，AI 辅助）：官方 PS1 解答第2--3页，修正其表头把斜率误写为 $-y_k$ 的笔误。积分给 $y=x^2$，所以 $y(1)=1$。令 $x_k=kh$，欧拉公式为 $y_{k+1}=y_k+2kh^2$，从 $y_0=0$ 得
\[
 y_n=2h^2\sum_{k=0}^{n-1}k=n(n-1)h^2=1-\frac1n.
\]
2步得到 $1/2$，3步得到 $2/3$，一般估计趋于1。真值减去估计的误差恰为 $1/n=h$，因此不仅符合一阶误差预测，还给出本题的确切误差常数1。
\end{solution}
\begin{exercise}\label{cw:ps1-II-3}
原题 II.3。史高治为侄子设立信托，投资按恒定年利率 $I$ 连续增长，每年恒定支付 $q$ 美元；$I$ 的单位为年$^{-1}$，例如 $I=0.05$ 表示每年5\%。(a) 建模并解释。(b) 求通解。(c) 取 $I=0.05$，若每月支付1000美元并永远保持本金不变，应投入多少？(d) 若仍每月支付1000美元，但要求恰好20年后用完，应投入多少？精确到美分。
\end{exercise}
\begin{solution}
解答依据（编者译审或补解，AI 辅助）：官方 PS1 解答第3页。设 $x(t)$ 为账户美元余额、$t$ 为年。在 $\Delta t$ 内利息约为 $Ix(t)\Delta t$、支出为 $q\Delta t$，故 $x'=Ix-q$。若 $I\ne0$，令 $u=x-q/I$，则 $u'=Iu$，所以
\[
 x(t)=\frac qI+Ce^{It};
\]
其中 $C=0$ 包含恒定余额解。若 $I=0$，则 $x=C-qt$。
(c) $q=12000$ 美元/年，恒定余额为 $q/I=240000$ 美元。
(d) 令 $T=20$，终值条件 $0=q/I+Ce^{IT}$ 给 $C=-(q/I)e^{-IT}$，所以
\[
 x(0)=\frac qI(1-e^{-IT})=240000(1-e^{-1})
 \simeq151708.93\ \text{美元}.
\]
按同一公式，$0\leq t<20$ 的余额为正，且在 $t=20$ 为零，符合模型要求。
\end{solution}

\originalchapter{作业二：复数与线性响应}\label{cw:ps2}
\originalsection{第一部分：指定练习}
\begin{exercise}\label{cw:ps2-I-4}原题 I.4：EP 第1.5节第1、2、5题；应先检查是否可分离变量。\end{exercise}
\begin{solution}仅列教材定位，未取得书题题面，未翻译、未解答；第1题也在 PS1 I.0 指定过。\end{solution}
\begin{exercise}\label{cw:ps2-I-5}原题 I.5：Notes 2E-1、2、7，分别见\hyperref[cw:bank-2E-1]{题2E-1}、\hyperref[cw:bank-2E-2]{题2E-2}、\hyperref[cw:bank-2E-7]{题2E-7}。\end{exercise}
\begin{solution}题面与解答复用所链公开题。\end{solution}
\begin{exercise}\label{cw:ps2-I-6}原题 I.6：(a) Notes 2E-9、10，见\hyperref[cw:bank-2E-9]{题2E-9}、\hyperref[cw:bank-2E-10]{题2E-10}。(b) 将下列函数写为 $A\cos(\omega t-\phi)$，每题先画对应直角三角形：(i) $\cos2t+\sin2t$；(ii) $\cos\pi t-\sqrt3\sin\pi t$；(iii) $\cos(t-\pi/8)+\sin(t-\pi/8)$。\end{exercise}
\begin{solution}
解答依据（编者译审或补解，AI 辅助）：PS2题面第3页答案；(iii)修正其把 $3\pi/8$ 排成 $3t/8$ 的笔误。(a)见链接。(b) 若余弦、正弦系数为 $(a,b)$，取斜边 $A=\sqrt{a^2+b^2}$，角 $\phi$ 满足 $\cos\phi=a/A$、$\sin\phi=b/A$。因此三题分别为
\[
 \sqrt2\cos(2t-\pi/4),\quad 2\cos(\pi t+\pi/3),\quad
 \sqrt2\cos(t-3\pi/8).
\]
对应三角形的有向直角边为 $(1,1)$、$(1,-\sqrt3)$、$(1,1)$；第三题先对自变量 $s=t-\pi/8$ 作同一个三角形，再将相位加回。
\begin{center}\begin{tikzpicture}[scale=.7]
\foreach \offset/\aa/\bb/\texta/\textb in {0/1.7/1.7/1/1,4/1.4/-2.42/1/{-\sqrt3},8/1.7/1.7/1/1}{\begin{scope}[xshift=\offset cm]\draw[->](-.2,0)--(2.4,0);\draw[->](0,-2.6)--(0,2.3);\draw[thick](0,0)--(\aa,0)--(\aa,\bb)--cycle;\node[below]at(\aa/2,0){$\texta$};\node[right]at(\aa,\bb/2){$\textb$};\end{scope}}
\end{tikzpicture}\end{center}
\end{solution}
\begin{exercise}\label{cw:ps2-I-7}原题 I.7：(a) Notes 2E-15，见\hyperref[cw:bank-2E-15]{题2E-15}。(b) 分别对 (i) $x'+2x=e^{3t}$，(ii) 复值方程 $x'+2x=e^{3it}$，先找与右端同指数形式的特解，再写通解。\end{exercise}
\begin{solution}解答依据（编者译审或补解，AI 辅助）：PS2题面第3页。(a)见链接。(b)代入 $Be^{\lambda t}$ 得 $(\lambda+2)B=1$，齐次解为 $Ce^{-2t}$，故分别为 $e^{3t}/5+Ce^{-2t}$（$C\in\R$）与 $e^{3it}/(2+3i)+Ce^{-2t}$（$C\in\C$）。\end{solution}
\originalsection{第二部分：题面与解答}
\begin{exercise}\label{cw:ps2-II-4}
原题 II.4。虚构元素 St 的半衰期为 $t_S$，每次衰变以相同概率产生 Mi 或稳定元素 En；Mi 的半衰期为 $t_M$，衰变为 En。初始只有 St，选单位使 $x(0)=1$；$x,y,z$ 分别表示 St、Mi、En 的量，$y(0)=z(0)=0$。题面明确 $t_M\ne t_S$（原 PDF 的不等号在纯文本提取中可能丢失斜线，须以原页为准）。
(a) 描图并求三者的长期极限。(b) 用衰变常数 $\sigma,\mu$ 建立方程，先联系半衰期，再检验总量守恒。(c) 依次求解。(d) Mi 何时最多？(e) 若 $x(0)=2$ 有何变化？(f) 无关的问题：若 $x=e^t$ 是 $tx'+2x=q(t)$ 的解，求 $q$ 及通解。
\end{exercise}
\begin{solution}
解答依据（编者译审，AI 辅助）：官方 PS2 解答第1页，补充定义区间与相等衰变率的极限对照；实际题目要求两半衰期不同。
由 $e^{-\sigma t_S}=1/2$ 得 $\sigma=(\ln2)/t_S$，同理 $\mu=(\ln2)/t_M$；两者都为正且 $\mu\ne\sigma$。每单位时间 St 减少 $\sigma x$，其中一半进入 Mi，因此
\[
 x'=-\sigma x,\quad y'=\tfrac12\sigma x-\mu y,
 \quad z'=\tfrac12\sigma x+\mu y.
\]
相加恰为零，初始总量1，故 $x+y+z=1$。先得 $x=e^{-\sigma t}$；乘 $y$ 方程以 $e^{\mu t}$ 得
\[
 (e^{\mu t}y)'=\tfrac12\sigma e^{(\mu-\sigma)t}.
\]
从0积分，利用 $y(0)=0$，可得
\[
 y=\frac\sigma{2(\mu-\sigma)}(e^{-\sigma t}-e^{-\mu t}),\qquad
 z=1-e^{-\sigma t}-\frac\sigma{2(\mu-\sigma)}(e^{-\sigma t}-e^{-\mu t}).
\]
$x$ 单调下降到0；$z'=\sigma x/2+\mu y\geq0$，从0上升到1。对 $y$ 求导得
\[
 y'=\frac\sigma{2(\mu-\sigma)}(-\sigma e^{-\sigma t}+\mu e^{-\mu t}).
\]
它在0处为 $\sigma/2>0$，只有一个零点：由 $\sigma e^{-\sigma t}=\mu e^{-\mu t}$ 得
\[
 t_{\max}=\frac{\ln(\mu/\sigma)}{\mu-\sigma}>0.
\]
分子分母同号，且在此之后 $y'<0$，所以 Mi 先升后降并最终趋于0。下图选择 $\mu=2\sigma$，横轴为无量纲时间 $s=\sigma t$，只是符合原条件的一个参数例子。
\begin{center}\begin{tikzpicture}\begin{axis}[width=.75\linewidth,height=4.7cm,axis lines=left,xlabel={$s$},ylabel={相对量},domain=0:6,samples=100,legend pos=north east]
\addplot[blue]{exp(-x)};\addlegendentry{St}
\addplot[red]{(exp(-x)-exp(-2*x))/2};\addlegendentry{Mi}
\addplot[black]{1-1.5*exp(-x)+.5*exp(-2*x)};\addlegendentry{En}
\end{axis}\end{tikzpicture}\end{center}
若另行考虑相等半衰期的极限 $\mu\to\sigma$，积分式连续给出 $y=(\sigma t/2)e^{-\sigma t}$、$z=1-(1+\sigma t/2)e^{-\sigma t}$，峰值时间为 $1/\sigma$；这不是本题给定条件，而是可用于检验一般公式的补充。
(e) 方程组线性且齐次，把初始总量加倍便使 $x,y,z$ 全部加倍。
(f) 代入得 $q=(t+2)e^t$。在 $t\ne0$ 的区间，齐次方程 $th'+2h=0$ 给 $h=C/t^2$，故 $x=e^t+C/t^2$。若要求在包含0的区间上连续可微，则 $C=0$，且 $t=0$ 处的原方程也要求 $x(0)=1$。
\end{solution}
\begin{exercise}\label{cw:ps2-II-5}
原题 II.5。(a) 制作三列表格，分别给直角式、模与辐角以及复平面上的点：(i) $1-i$；(ii) 模为2、辐角为 $\pi/6$；(iii) 虚部为负的 $i$ 的平方根；(iv) 辐角在 $(0,\pi/2)$ 的1的六次根；(v) $((1+i)/\sqrt2)^{-13}$。(b) 求所有根：(i) $z^4+4=0$；(ii) $z^2+2z+2=0$。
\end{exercise}
\begin{solution}
解答依据（编者译审或补解，AI 辅助）：官方 PS2 解答第2页。辐角均按模 $2\pi$ 理解；主辐角选 $(-\pi,\pi]$。
\begin{center}\begin{tabular}{ccl}序号&直角式&模与主辐角\\\hline
(i)&$1-i$&$\sqrt2,-\pi/4$\\
(ii)&$\sqrt3+i$&$2,\pi/6$\\
(iii)&$(-1-i)/\sqrt2$&$1,-3\pi/4$\\
(iv)&$(1+\sqrt3 i)/2$&$1,\pi/3$\\
(v)&$(-1+i)/\sqrt2$&$1,3\pi/4$\\
\end{tabular}\end{center}
(iii) 因 $i=e^{i\pi/2}$，两个平方根为 $\pm e^{i\pi/4}$，其中负号给负虚部。(iv) 六次根辐角为 $k\pi/3$，指定范围只允许 $k=1$。(v) 底数是 $e^{i\pi/4}$，负13次方为 $e^{-13i\pi/4}=e^{3i\pi/4}$。相应点如下；它们的坐标就是表格直角式的实、虚部。
\begin{center}\begin{tikzpicture}[scale=1.3]
\draw[->](-1.2,0)--(2.2,0)node[right]{实部};\draw[->](0,-1.25)--(0,1.5)node[above]{虚部};
\foreach \xx/\yy/\lab in {1/-1/i,1.732/1/ii,-.707/-.707/iii,.5/.866/iv,-.707/.707/v}{\fill(\xx,\yy)circle(1.5pt)node[above right]{(\lab)};}
\end{tikzpicture}\end{center}
(b)(i) $z^4=-4=4e^{i(\pi+2k\pi)}$，故 $z=\sqrt2e^{i(\pi/4+k\pi/2)}$，$k=0,1,2,3$，即 $1+i,-1+i,-1-i,1-i$。(ii) 完成平方得 $(z+1)^2=-1$，根为 $-1\pm i$。
\end{solution}
\begin{exercise}\label{cw:ps2-II-6}
原题 II.6。(a) 在 II.5(a) 的表中加一列指数式 $Ae^{i\theta}$。(b) 求所有满足 $e^{a+bi}=-2$ 的复数。(c) 由 $(e^{it})^4=e^{4it}$ 导出以 $\sin t,\cos t$ 的幂表示 $\cos4t$ 的公式。(d) 对下列 $f$，画实函数图、找 $w=a+bi$ 使 $\Re(e^{wt})=f(t)$，并画 $e^{wt}$ 的轨迹：(i) $\cos2\pi t$；(ii) $e^{-t}$；(iii) $e^{-t}\cos2\pi t$；(iv) 恒等于1。原题建议借助 Complex Exponential Mathlet。
\end{exercise}
\begin{solution}
解答依据（编者译审或补解，AI 辅助）：官方 PS2 解答第2页。(a) 五个指数式依次为 $\sqrt2e^{-i\pi/4}$、$2e^{i\pi/6}$、$e^{-3i\pi/4}$、$e^{i\pi/3}$、$e^{3i\pi/4}$。
(b) 取模得 $e^a=2$，故 $a=\ln2$；取单位复数部分得 $e^{ib}=-1$，故 $b=(2k+1)\pi$、$k\in\mathbb Z$，这给出全部解。
(c) 展开并取实部：
\[
 \cos4t=\Re(\cos t+i\sin t)^4
 =\cos^4t-6\cos^2t\sin^2t+\sin^4t.
\]
(d) 利用 $e^{(a+bi)t}=e^{at}(\cos bt+i\sin bt)$，可分别取 $w=2\pi i,-1,-1+2\pi i,0$。轨迹分别是单位圆、正实轴、绕原点逆时针向内的螺线、单点1；这里“向内”指 $t$ 增大。若只取 $t\geq0$，第二项轨迹为 $(0,1]$，第三项只取螺线从1向内的部分。下图画 $0\leq t\leq2$ 的实函数与相应轨迹，恒值1与实轴上单点重合。
\begin{center}\begin{tikzpicture}
\begin{axis}[name=realplot,width=.46\linewidth,height=4.3cm,axis lines=middle,xlabel={$t$},domain=0:2,samples=140,ymin=-1.15,ymax=1.2]
\addplot[blue]{cos(deg(2*pi*x))};\addplot[red]{exp(-x)};\addplot[black]{exp(-x)*cos(deg(2*pi*x))};\addplot[orange]{1};
\end{axis}
\begin{axis}[at={(realplot.east)},anchor=west,xshift=1cm,width=.42\linewidth,height=4.3cm,axis equal image,axis lines=middle,xlabel={实部},ylabel={虚部},xlabel style={at={(axis description cs:1.03,.5)},anchor=west},ylabel style={at={(axis description cs:.5,1.03)},anchor=south},xmin=-1.15,xmax=1.15,ymin=-1.15,ymax=1.15,samples=150]
\addplot[blue,domain=0:360]({cos(x)},{sin(x)});
\addplot[red,domain=0:2]({exp(-x)},0);
\addplot[black,domain=0:2]({exp(-x)*cos(deg(2*pi*x))},{exp(-x)*sin(deg(2*pi*x))});\addplot[orange,only marks]coordinates{(1,0)};
\end{axis}\end{tikzpicture}\end{center}
\end{solution}
\begin{exercise}\label{cw:ps2-II-7}
原题 II.7。(a) 把 $\Re(e^{3it}/(\sqrt3+i))$ 写成 $a\cos3t+b\sin3t$ 和 $A\cos(3t-\phi)$，再通过分母的极式重算。原题第二条路线中出现了不一致的 $2+2i$ 与 $e^{2it}$；请辨明对应表达式。(b) 为 $z'+3z=e^{2it}$ 找形如 $we^{2it}$ 的特解并写通解。(c) 用(b)求 $x'+3x=\cos2t$ 的特解和通解。
\end{exercise}
\begin{solution}
解答依据（编者译审或补解，AI 辅助）：官方 PS2 解答第2页，修正题面(a)第二路线误植。(a) 有
\[
 \frac1{\sqrt3+i}=\frac{\sqrt3-i}{4},\qquad
 \Re\frac{e^{3it}}{\sqrt3+i}=\frac{\sqrt3}4\cos3t+\frac14\sin3t
 =\frac12\cos(3t-\pi/6).
\]
另一方面 $\sqrt3+i=2e^{i\pi/6}$，于是商为 $\frac12e^{i(3t-\pi/6)}$，取实部得到同一结果。题面的 $2+2i$ 是 $2\sqrt2e^{i\pi/4}$，若真的计算 $e^{2it}/(2+2i)$，实部为 $\frac1{2\sqrt2}\cos(2t-\pi/4)$，不会等于本题前一个表达式，因此不能把两式混用。
(b) 代入得 $(3+2i)w=1$，所以 $z=e^{2it}/(3+2i)+Ce^{-3t}$，$C\in\C$。(c) 取实部得
\[
 x=\frac{3\cos2t+2\sin2t}{13}+ce^{-3t},\qquad c\in\R.
\]
求导代回即可核验右端确为 $\cos2t$。
\end{solution}

\originalchapter{作业三：平衡与阻尼}\label{cw:ps3}
本章合并官方 PS3 的前后两半。原讲次10是考试，因此题号从9跳到11；不重编号。
\originalsection{第一部分：指定练习}
\begin{exercise}\label{cw:ps3-I-8}原题 I.8：Notes 1E-1，见\hyperref[cw:bank-1E-1]{题1E-1}。\end{exercise}
\begin{solution}题面与完整解答见链接。\end{solution}
\begin{exercise}\label{cw:ps3-I-9}
原题 I.9。考虑 $x'+2x=1$。(a) 分别用(i)分离变量，(ii)积分因子，(iii)把1写为 $e^{0t}$、试常数特解后加齐次解，求通解。(b) 画相线和包括平衡解在内的若干解，判断稳定性。(c) 对 $x(0)=0$，用三个等长欧拉步估计 $x(1)$。
\end{exercise}
\begin{solution}
解答依据（编者译审或补解，AI 辅助）：PS3前半第2页部分答案，补全三条推导。(a)(i) 非平衡时 $dx/(1-2x)=dt$，积分得 $-\frac12\ln|1-2x|=t+C$，恢复被除去的平衡解后，$x=\frac12+ce^{-2t}$。(ii) 乘 $e^{2t}$ 后有 $(e^{2t}x)'=e^{2t}$，积分得到同一通解。(iii) 常数特解 $A$ 满足 $2A=1$，加上齐次解 $ce^{-2t}$ 即得结果。(b) 在 $x<1/2$ 时 $x'>0$，在 $x>1/2$ 时 $x'<0$，相线箭头都指向 $1/2$；任意解与 $1/2$ 的差为 $ce^{-2t}\to0$，故平衡是渐近稳定的。
\begin{center}\begin{tikzpicture}[scale=1.1]
\draw[->](-2,0)--(2,0)node[right]{$x$};\fill(0,0)circle(2pt)node[below]{$1/2$};\draw[->,thick](-1.6,.15)--(-.5,.15);\draw[->,thick](1.6,.15)--(.5,.15);
\end{tikzpicture}\end{center}
下图另取 $C=-1,-1/2,0,1/2,1$ 描出 $x(t)=1/2+Ce^{-2t}$，其中水平虚线为平衡解，所有曲线沿时间正向趋于它。
\begin{center}\begin{tikzpicture}\begin{axis}[width=.72\linewidth,height=4.3cm,axis lines=middle,xlabel={$t$},ylabel={$x$},domain=0:3,samples=100,ymin=-.6,ymax=1.6]
\addplot[blue]{.5-exp(-2*x)};\addplot[red]{.5-exp(-2*x)/2};
\addplot[black,dashed]{.5};\addplot[red]{.5+exp(-2*x)/2};\addplot[blue]{.5+exp(-2*x)};
\end{axis}\end{tikzpicture}\end{center}
(c) $h=1/3$，$x_{k+1}=x_k+h(1-2x_k)=x_k/3+1/3$，依次得 $x_1=1/3,x_2=4/9,x_3=13/27$。真值是 $(1-e^{-2})/2$，欧拉估计 $13/27$ 为所求。
\end{solution}
\begin{exercise}\label{cw:ps3-I-11}原题 I.11：Notes 2C-1(a)、(b)，见\hyperref[cw:bank-2C-1]{题2C-1}；EP 第2.1节第43、44、47题，第2.3节第1、2、7、15题。\end{exercise}
\begin{solution}公开 Notes 指定小问复用链接；EP 七道题只列定位，未取得题面，未翻译、未解答。\end{solution}
\begin{exercise}\label{cw:ps3-I-12}原题 I.12：Notes 2C-1(c)、(d)、(e)以及2G-1，见\hyperref[cw:bank-2C-1]{题2C-1}、\hyperref[cw:bank-2G-1]{题2G-1}；EP 第2.3节第5、21、23题。\end{exercise}
\begin{solution}公开题按链接复用；EP 三道题仅列定位，未翻译、未解答。\end{solution}
\originalsection{第二部分：题面与解答}
\begin{exercise}\label{cw:ps3-II-8}
原题 II.8。没有捕猎时，保护区羚羊服从稳定种群为1千只的逻辑斯谛模型；政府希望研究恒定年捕猎率 $a$ 的影响。原题用 Phase Lines Mathlet 的首项
\[
 y'=(1-y)y-a
\]
和相线、分岔图研究此问题。
(a) 说明此模型的假设。(b) 对所有 $a$ 求平衡点数量、位置与稳定性。(c) 若每年允许捕杀187.5只，稳定种群是多少？低于什么临界种群会崩溃？(d) 对此 $a$ 描出五种解行为；相差时间平移的解视为同型。说明向前和向后时间的行为，并与相线各部分对应。(e) 给出分岔图中平衡点曲线的方程。
\end{exercise}
\begin{solution}
解答依据（编者译审或补解，AI 辅助）：官方 PS3 解答第1页，修正其(e)符号笔误，补充五类解的端点行为。
(a) 一般逻辑斯谛增长为 $y'=k_0(1-y/p)y$。这里 $p=1$ ko，另取 $k_0=1$ 年$^{-1}$，并假设捕猎率不随种群大小变化，因而减去常数 $a$（ko/年）。这一模型只在种群非负时有生物意义，抵达0以后应停止而不继续把负值解释为种群。
(b) 平衡条件 $y^2-y+a=0$ 给
\[
 y_\pm=\frac12\pm\sqrt{\frac14-a}.
\]
$a<1/4$ 时有两个平衡点。由于 $F(y)=y-y^2-a=-(y-y_-)(y-y_+)$，外侧 $F<0$、两根之间 $F>0$，故低根不稳定、高根渐近稳定。$a=1/4$ 时 $F=-(y-1/2)^2$，唯一平衡从上方吸引、从下方排斥，是半稳定。$a>1/4$ 时没有平衡，所有解均下降。若讨论捕猎，只允许 $a\geq0$；代数上的 $a<0$ 表示恒定引入种群。
(c) $a=187.5/1000=3/16$，所以 $y_-=1/4,y_+=3/4$，稳定值750只，崩溃阈值250只。
(d) 相线分为 $y<1/4$、$y=1/4$、$1/4<y<3/4$、$y=3/4$、$y>3/4$ 五部分。两个平衡分别给恒定解。其余三部分用分离变量统一描述：
\[
 \frac{y-3/4}{y-1/4}=Ke^{-t/2},\qquad
 y(t)=\frac{3/4-(K/4)e^{-t/2}}{1-Ke^{-t/2}}.
\]
取 $K<0$ 时，解在两根之间，向后趋于 $1/4$，向前趋于 $3/4$。取 $K>0$ 时，分母零点 $T=2\ln K$ 分开两个最大解区间：在 $t>T$，解从 $+\infty$ 下降到 $3/4$；在 $t<T$，解从 $1/4$ 的下方下降，$t\uparrow T$ 时趋于 $-\infty$，且在此前已抵达种群0。如下图分别选 $K=1,-1$，画数学模型三类非恒定解和两条平衡解。
\begin{center}\begin{tikzpicture}\begin{axis}[width=.76\linewidth,height=5.2cm,axis lines=middle,xlabel={$t$},ylabel={$y$},xmin=-5,xmax=5,ymin=-.5,ymax=1.5,samples=100]
\addplot[blue,domain=.45:5]{(.75-.25*exp(-x/2))/(1-exp(-x/2))};
\addplot[red,domain=-5:-.45]{(.75-.25*exp(-x/2))/(1-exp(-x/2))};
\addplot[black,domain=-5:5]{(.75+.25*exp(-x/2))/(1+exp(-x/2))};
\addplot[gray,dashed,domain=-5:5]{.25};\addplot[gray,dashed,domain=-5:5]{.75};
\end{axis}\end{tikzpicture}\end{center}
(e) 分岔曲线为 $a=y-y^2$，等价于 $y^2-y+a=0$。官方答案把最后一项误写为 $-a$；代回原方程可立即发现此误。两分支在 $(a,y)=(1/4,1/2)$ 合并。
\end{solution}
\begin{exercise}\label{cw:ps3-II-9}
原题 II.9。仍取 $y'=(1-y)y-3/16$，令稳定平衡为 $y_0$，$u=y-y_0$ 表示种群相对平衡的超额。(a) 写出 $u$ 的方程，检验自治性与零平衡。(b) 忽略 $u^2$ 等高阶项，求线性化方程及通解。(c) 写出近平衡近似；若 $y(10)-y_0=b$，估计 $y(11),y(12)$。(d) 若线性方程 $x'+p(t)x=q(t)$ 是自治方程，$p,q$ 应满足什么条件？
\end{exercise}
\begin{solution}
解答依据（编者译审或补解，AI 辅助）：官方 PS3 解答第1页，补上(c)的两个具体预测和(d)的理由。$y_0=3/4$，代入 $y=u+3/4$ 得
\[
 u'=(1/4-u)(u+3/4)-3/16=-\tfrac12u-u^2.
\]
右端无显含时间，且在 $u=0$ 时为0。(b) 线性化为 $u'=-u/2$，故 $u=Ce^{-t/2}$。(c) 当 $|u|\ll1$ 时二次项相对线性项很小，得到
\[
 y(t)\simeq\tfrac34+be^{-(t-10)/2},\quad
 y(11)\simeq\tfrac34+be^{-1/2},\quad y(12)\simeq\tfrac34+be^{-1}.
\]
这些是近平衡预测，不是原非线性方程的精确等式。精确超额也可写为
\[
 u(t)=\frac{be^{-(t-10)/2}}{1+2b(1-e^{-(t-10)/2})},
\]
在分母非零的区间成立；它显示被忽略的修正从 $b^2$ 开始。(d) 自治性要求 $q(t)-p(t)x$ 对每个 $x$ 都不依赖 $t$。取 $x=0$ 得 $q$ 常数，再取 $x=1$ 得 $p$ 也是常数。
\end{solution}
\begin{exercise}\label{cw:ps3-II-11}
原题 II.11。酒馆门可穿过门框，阻尼应使门缓慢回到闭合位置；若向门框推得足够快，它可能穿过后从另一侧回到闭合位置。用 Damped Vibration Mathlet 研究
\[
 mx''+bx'+kx=0,
\]
初值框的横坐标是 $x'(0)$、纵坐标是 $x(0)$。固定 $m=0.50,b=1.5,k=0.625=5/8$；Mathlet 可只能显示 $k=0.62$，计算仍用精确值。(a) 求特征多项式与根。(b) 用 $x(0)=x_0,x'(0)=v_0$ 写通解。(c) 哪些初值给出单一指数解？(d) 固定 $x_0$（示例 $x_0=0.25$），求使解在某个正时间等于0的 $v_0$ 的精确范围。
\end{exercise}
\begin{solution}
解答依据（编者译审或补解，AI 辅助）：官方 PS3 解答第2页，补充 $x_0=0$ 的边界。(a) 特征多项式
\[
 p(r)=\tfrac12r^2+\tfrac32r+\tfrac58
 =\tfrac12(r+\tfrac12)(r+\tfrac52),
\]
根为 $-1/2,-5/2$。(b) 设 $x=c_1e^{-t/2}+c_2e^{-5t/2}$，则 $c_1+c_2=x_0$，$-c_1/2-5c_2/2=v_0$。解这两个线性方程得
\[
 c_1=\tfrac14(5x_0+2v_0),\quad c_2=-\tfrac14(x_0+2v_0).
\]
(c) $c_1=0$ 或 $c_2=0$，即 $v_0=-5x_0/2$ 或 $v_0=-x_0/2$；零解位于两直线交点。
(d) 非零解有零点时，$e^{2t}=-c_2/c_1$。正时间零点要求 $-c_2/c_1>1$。若 $x_0>0$，这等价于 $c_1<0$，即 $v_0<-5x_0/2$；若 $x_0<0$，等价于 $v_0>-5x_0/2$。所以 $x_0=1/4$ 时须 $v_0<-5/8$。等号时是单一快指数解，只渐近到0，不在有限时间穿过。若 $x_0=0$ 且 $v_0\ne0$，$x=(v_0/2)(e^{-t/2}-e^{-5t/2})$，在 $t>0$ 没有零点；$x_0=v_0=0$ 时是恒零解。门的位置与速度可直接从公式重绘，无需把交互读数当作计算依据。
\end{solution}
\begin{exercise}\label{cw:ps3-II-12}
原题 II.12。继续使用同一模型，保持 $m=0.50,k=0.625$。(a) 降低 $b$，根在何值重合？代数验证并给此时通解。(b) 再取 $b=0.25$，任选非零初值，求根与通解。(c) 对 $x(0)=0,x'(0)=1$，求前五个零点的相邻间距和伪周期，判断振动是否随衰减而加快或减慢；再对不同初值说明相邻零点间距。
\end{exercise}
\begin{solution}
解答依据（编者译审或补解，AI 辅助）：官方 PS3 解答第2页，提供理论零点以替代未经运行的测量。(a) 归一化多项式为 $r^2+2br+5/4$，判别式为 $4b^2-5$。非负阻尼的重根条件是 $b=\sqrt5/2$，根为 $-\sqrt5/2$，通解 $(A+Bt)e^{-\sqrt5t/2}$。
(b) 当 $b=1/4$，多项式可写为 $(r+1/4)^2+19/16$，所以
\[
 r=-\tfrac14\pm i\tfrac{\sqrt{19}}4,\qquad
 x=e^{-t/4}\left(A\cos\tfrac{\sqrt{19}}4t+B\sin\tfrac{\sqrt{19}}4t\right).
\]
(c) 给定初值使 $A=0,B=4/\sqrt{19}$，零点为 $t_j=4j\pi/\sqrt{19}$，$j=0,1,2,\ldots$。含 $t=0$ 的前五个为
\[
 0,\quad\frac{4\pi}{\sqrt{19}},\quad\frac{8\pi}{\sqrt{19}},\quad\frac{12\pi}{\sqrt{19}},\quad\frac{16\pi}{\sqrt{19}}.
\]
若只数正零点，则从 $j=1$ 开始。每个相邻间距均为 $4\pi/\sqrt{19}\simeq2.8829$，伪周期为其两倍 $8\pi/\sqrt{19}\simeq5.7658$。非零一般解写成 $Re^{-t/4}\cos(\sqrt{19}t/4-\phi)$ 后，零点只整体平移，相邻间距保持相同。振幅衰减，振动的角频率和伪周期保持不变。
\begin{center}\begin{tikzpicture}\begin{axis}[width=.75\linewidth,height=4.6cm,axis lines=middle,xlabel={$t$},ylabel={$x$},domain=0:13,samples=180]
\addplot[blue]{4/sqrt(19)*exp(-x/4)*sin(deg(sqrt(19)*x/4))};
\addplot[gray,dashed]{4/sqrt(19)*exp(-x/4)};\addplot[gray,dashed]{-4/sqrt(19)*exp(-x/4)};
\end{axis}\end{tikzpicture}\end{center}
\end{solution}

\originalchapter{作业四：受迫响应与共振}\label{cw:ps4}
\originalsection{第一部分：指定练习}
\begin{exercise}\label{cw:ps4-I-13}原题 I.13：Notes 2F-6(b)、(c)、(d)，见\hyperref[cw:bank-2F-6]{题2F-6}。\end{exercise}
\begin{solution}仅复用指定三个小问的公开题面与解答。\end{solution}
\begin{exercise}\label{cw:ps4-I-14}原题 I.14：求 $x'''-x=e^{2t}$ 的通解。\end{exercise}
\begin{solution}
解答依据（编者译审或补解，AI 辅助）：PS4题面第3页答案。特征多项式 $r^3-1=(r-1)(r^2+r+1)$，根为 $1,-1/2\pm i\sqrt3/2$。试特解 $Be^{2t}$ 得 $7B=1$，因此
\[
 x=\tfrac17e^{2t}+c_1e^t+e^{-t/2}
 \left(c_2\cos\tfrac{\sqrt3}2t+c_3\sin\tfrac{\sqrt3}2t\right).
\]
\end{solution}
\begin{exercise}\label{cw:ps4-I-15}原题 I.15：Notes 2C-8(c)、(d)，见\hyperref[cw:bank-2C-8]{题2C-8}；EP 第2.5节第2、8、11、14题。\end{exercise}
\begin{solution}公开题按链接复用；EP 四道题仅列定位，题面未取得，未翻译、未解答。\end{solution}
\begin{exercise}\label{cw:ps4-I-16}原题 I.16：EP 第2.6节第17题。\end{exercise}
\begin{solution}仅列定位，题面未取得，未翻译、未解答。\end{solution}
\originalsection{第二部分：题面与解答}
\begin{exercise}\label{cw:ps4-II-13}
原题 II.13。(a) 求 $x'+2x=e^{3t}\cos4t$ 的实通解，建议用复数与指数响应公式。(b) 求 $x''+x'+2x=\cos t$ 的实通解，把稳态特解分别写成 $a\cos t+b\sin t$ 和 $A\cos(t-\phi)$。
原题随后使用 Forced Damped Vibrations Mathlet，力直接作用于质量，固定 $m=b=1,k=2$。(c) 在输入幅值 $A=0$、初值 $x(0)=0,x'(0)=1$ 下求伪周期。(d) 取输入 $A=1,\omega=1$，求稳态幅值、$x_p(0),x_p'(0)$，与图比较。(e) 对静止初值 $x(0)=x'(0)=0$，求瞬态项 $x_h$，使 $x=x_p+x_h$。
\end{exercise}
\begin{solution}
解答依据（编者译审或补解，AI 辅助）：官方 PS4 解答第1页。(a) 复特解 $z_p=e^{(3+4i)t}/(5+4i)$，分母有理化给
\[
 x=\frac{e^{3t}}{41}(5\cos4t+4\sin4t)+ce^{-2t}.
\]
(b) 特征多项式 $p(r)=r^2+r+2$，$p(i)=1+i$，故
\[
 x_p=\Re\frac{e^{it}}{1+i}=\tfrac12(\cos t+\sin t)
 =\frac1{\sqrt2}\cos(t-\pi/4).
\]
齐次根为 $-1/2\pm i\sqrt7/2$，实通解为
\[
 x=x_p+e^{-t/2}(C\cos\tfrac{\sqrt7}2t+D\sin\tfrac{\sqrt7}2t).
\]
(c) 无输入且给定初值时，$x=(2/\sqrt7)e^{-t/2}\sin(\sqrt7t/2)$，因此伪周期 $4\pi/\sqrt7\simeq4.7496$；这是理论值，不是交互测量值。(d) 幅值 $1/\sqrt2\simeq0.7071$，初值 $x_p(0)=x_p'(0)=1/2$。
(e) 需使 $x_h(0)=x_h'(0)=-1/2$。从 $C=-1/2$ 和 $-C/2+\sqrt7D/2=-1/2$ 得 $D=-3/(2\sqrt7)$，所以
\[
 x_h=e^{-t/2}\left(-\tfrac12\cos\tfrac{\sqrt7}2t
 -\frac3{2\sqrt7}\sin\tfrac{\sqrt7}2t\right).
\]
两项相加在 $t=0$ 的位移和速度均为0。图按所得公式复现三条曲线：
\begin{center}\begin{tikzpicture}\begin{axis}[width=.76\linewidth,height=5cm,axis lines=middle,xlabel={$t$},domain=0:12,samples=200,legend pos=north east]
\addplot[green!50!black]{(cos(deg(x))+sin(deg(x)))/2};\addlegendentry{稳态}
\addplot[red]{exp(-x/2)*(-cos(deg(sqrt(7)*x/2))/2-3*sin(deg(sqrt(7)*x/2))/(2*sqrt(7)))};\addlegendentry{瞬态}
\addplot[blue]{(cos(deg(x))+sin(deg(x)))/2+exp(-x/2)*(-cos(deg(sqrt(7)*x/2))/2-3*sin(deg(sqrt(7)*x/2))/(2*sqrt(7)))};\addlegendentry{总解}
\end{axis}\end{tikzpicture}\end{center}
\end{solution}
\begin{exercise}\label{cw:ps4-II-14}
原题 II.14。Amplitude and Phase: Second Order I Mathlet 研究经弹簧驱动的单位质量系统。输入为弹簧上端的位置 $y(t)=\cos\omega t$，输出为质量位置的稳态 $x(t)$，输入幅值为1。(a) 取 $b=0.5,k=4,\omega=2$，核算增益和时间滞后。(b) 求复增益、随 $\omega$ 变化的增益，以及相位滞后 $\phi$ 的正切。
\end{exercise}
\begin{solution}
解答依据（编者译审或补解，AI 辅助）：官方 PS4 解答第1页。弹簧力为 $k(y-x)$，阻尼力为 $-bx'$，故 $x''+bx'+kx=ky$，输入经弹簧而使右端带因子 $k$。
(a) $p(2i)=4-4+i=i$，于是 $x_p=\Re(4e^{2it}/i)=4\sin2t$，增益4，相位滞后 $\phi=\pi/2$，时间滞后为 $\phi/\omega=\pi/4\simeq0.7854$。
(b) 复输入 $e^{i\omega t}$ 对应稳态 $H(\omega)e^{i\omega t}$，其中
\[
 H(\omega)=\frac4{4-\omega^2+i\omega/2},\quad
 g(\omega)=\frac4{\sqrt{(4-\omega^2)^2+\omega^2/4}},\quad
 \tan\phi=\frac{\omega/2}{4-\omega^2}.
\]
$\omega>0$ 时应取 $\phi=\operatorname{Arg}(4-\omega^2+i\omega/2)\in(0,\pi)$；仅靠正切不能确定象限，$\omega=2$ 时正切无定义但相位仍为 $\pi/2$。
\end{solution}
\begin{exercise}\label{cw:ps4-II-15}
原题 II.15。求下列方程的实通解：(a) $x'+x=e^{-t}$，使用共振指数响应公式；(b) $x'''-x'=t^2+1$；(c) $2x''+x'=(t^2+1)e^{2t}$。
\end{exercise}
\begin{solution}
解答依据（编者译审或补解，AI 辅助）：(a)(b)官方 PS4 解答第1页；(c)官方文件未给解答，独立补算。
(a) $r+1$ 在 $r=-1$ 为零，因此试 $te^{-t}$。其导数加自身为 $e^{-t}$，通解 $x=(t+C)e^{-t}$。
(b) 令 $u=x'$，则 $u''-u=t^2+1$。试 $u=at^2+bt+c$，比较系数 $-a=1,-b=0,2a-c=1$，得 $u_p=-t^2-3$。积分得 $x_p=-t^3/3-3t$；齐次多项式 $r(r-1)(r+1)$ 给
\[
 x=-t^3/3-3t+C_1+C_2e^t+C_3e^{-t}.
\]
(c) 为避免反复求带指数的导数，令 $x=e^{2t}v$。则
\[
 2x''+x'=e^{2t}(2v''+9v'+10v).
\]
试 $v=At^2+Bt+C$，比较 $10A=1$、$18A+10B=0$、$4A+9B+10C=1$，得 $A=1/10$、$B=-9/50$、$C=111/500$。齐次根为0、$-1/2$，所以
\[
 x=e^{2t}\left(\frac{t^2}{10}-\frac{9t}{50}+\frac{111}{500}\right)
 +C_1+C_2e^{-t/2}.
\]
这里常数项经 $4A+9B+10C$ 核验为1。
\end{solution}
\begin{exercise}\label{cw:ps4-II-16}
原题 II.16。继续用 II.14 的经弹簧驱动系统，前三问固定 $k=4,b=0.5$。(a) 画增益随输入角频率的图，求使增益最大的实际共振频率 $\omega_r$ 和最大增益，判断是否恰为2。(b) 求相位滞后45度时的正频率。(c) 画 $\omega$ 从0到无穷时的复增益 Nyquist 轨迹，并标出相位接近 $\pi/4$ 的点。(d) 固定 $\omega=2,k=4$、改变 $b$，复增益在图中如何移动？对一般 $k$，调整 $\omega$ 使相位 $\pi/2$ 后再改变 $b$，相位为何保持？给出一般 $H$ 并解释。
\end{exercise}
\begin{solution}
解答依据（编者译审或补解，AI 辅助）：官方 PS4 答案文件没有本题解答，以下独立由 II.14 的公式求解。
(a) 最大化 $g$ 等价于最小化
\[
 D(\omega)=(4-\omega^2)^2+\omega^2/4,\qquad
 D'(\omega)=\omega(4\omega^2-31/2).
\]
正临界点 $\omega_r=\sqrt{31/8}\simeq1.9685$，在它之前 $D$ 下降、之后上升，所以给出唯一最大增益。此时 $4-\omega_r^2=1/8$，$D=63/64$，从而 $g_{\max}=32/\sqrt{63}\simeq4.0316$。它略高于 $g(2)=4$。
(b) 相位 $\pi/4$ 要求实部、虚部相等且为正，$4-\omega^2=\omega/2$，正解
\[
 \omega_{45}=\frac{\sqrt{65}-1}{4}\simeq1.7656.
\]
(c) 有
\[
 \Re H=\frac{4(4-\omega^2)}{D(\omega)},\qquad
 \Im H=-\frac{2\omega}{D(\omega)}.
\]
轨迹从 $H(0)=1$ 出发，经第四象限到 $H(2)=-4i$，随后进入第三象限，最终从下方趋于0。在 $\omega_{45}$ 处，$H=(4/\omega_{45})(1-i)$，下图标出该点，并同时给出增益图。
\begin{center}\begin{tikzpicture}
\begin{axis}[name=gainplot,width=.45\linewidth,height=4.7cm,axis lines=left,xlabel={$\omega$},ylabel={增益},domain=0:4,samples=220]
\addplot[blue]{4/sqrt((4-x*x)^2+x*x/4)};\addplot[only marks,red]coordinates{(1.968502,4.031621)};
\end{axis}
\begin{axis}[at={(gainplot.east)},anchor=west,xshift=1.1cm,width=.42\linewidth,height=4.7cm,axis equal image,axis lines=middle,xlabel={$\Re H$},ylabel={$\Im H$},xmin=-2.4,xmax=3.1,ymin=-4.3,ymax=.3,samples=250]
\addplot[blue,domain=0:10]({4*(4-x*x)/((4-x*x)^2+x*x/4)},{-2*x/((4-x*x)^2+x*x/4)});
\addplot[red,only marks]coordinates{(2.265564,-2.265564)};
\draw[red,->](axis cs:0,0)--(axis cs:2.265564,-2.265564);
\end{axis}\end{tikzpicture}\end{center}
(d) 一般的单位质量系统有
\[
 H(\omega)=\frac{k}{k-\omega^2+ib\omega}.
\]
对 $b>0,k>0,\omega>0$，$\phi=\pi/2$ 意味着 $H$ 是负纯虚数，等价于 $k-\omega^2=0$。所以 $\omega=\sqrt k$ 时无论正阻尼取何值，相位都保持 $\pi/2$，且 $H=-i\sqrt k/b$ 沿负虚轴移动。特别地 $k=4,\omega=2$ 时 $H=-2i/b$。$b=0$ 是无阻尼共振，稳态公式分母为零，不能纳入此论断。
\end{solution}

\originalchapter{作业五：电路与傅里叶级数}\label{cw:ps5}
本章合并官方 PS5 的两半；第19讲为考试，第一部分第18项是“无新题”，这两项保留为安排说明而不另造题目。
\originalsection{第一部分：指定练习}
\begin{exercise}\label{cw:ps5-I-17}
原题 I.17：EP 第2.5节第8题；另外，求 $x''+x=t^2+\cos(2t-1)$ 的一个解。
\end{exercise}
\begin{solution}
EP 题仅列定位，未翻译、未解答；它曾出现在 PS4 I.15。对于显式给出的方程，依据 PS5 后半第1页官方更正独立核验。设多项式特解 $At^2+Bt+C$，则 $x''+x=At^2+Bt+(2A+C)$，比较得 $A=1,B=0,C=-2$。对余弦项，二阶导加自身使系数乘 $1-2^2=-3$，故一个解为
\[
 x_p=t^2-2-\tfrac13\cos(2t-1).
\]
前半第2页曾把常数写成 $-1/2$，后半明确更正为 $-2$。直接代回，$(t^2-2)''+(t^2-2)=t^2$，确认更正；若需要通解，再加 $C_1\cos t+C_2\sin t$。
\end{solution}
原题 I.18：无新题。I.19 与 II.19：课堂考试，不是额外作业题。
\begin{exercise}\label{cw:ps5-I-20}
原题 I.20：Notes 7A-1、7A-2(a)，见\hyperref[cw:bank-7A-1]{题7A-1}、\hyperref[cw:bank-7A-2]{题7A-2(a)}。
\end{exercise}
\begin{solution}中文题面与解答复用上述公开题。\end{solution}
\begin{exercise}\label{cw:ps5-I-21}
原题 I.21：(a) 再做 Notes 7A-2(a)，这次先以方波 $\operatorname{sq}(t)$ 表示该函数；(b) 做 Notes 7A-2(b)，通过积分其导数的傅里叶级数求解。两题见\hyperref[cw:bank-7A-2]{题7A-2}。
\end{exercise}
\begin{solution}
解答依据（编者补解，AI 辅助）：公开 Notes 7A-2 的题面，按本次作业另行指定的方法重算。
(a) 链接中的函数在 $-\pi<t\leq0$ 为0、在 $0<t\leq\pi$ 为1，随后按 $2\pi$ 周期延拓。在非跳点处它等于 $(1+\operatorname{sq}(t))/2$，因此
\[
 f(t)\sim\tfrac12+\frac2\pi\sum_{j=0}^\infty\frac{\sin((2j+1)t)}{2j+1}.
\]
原题在跳点所定的0或1与方波表示的跳点值 $1/2$ 不同；单点值不改变系数，级数在跳点收敛到左右极限的平均 $1/2$。不能把这个表示写成处处等式。
(b) 链接中的另一函数在 $[-\pi,\pi]$ 为 $|t|$，按周期延拓。除角点外，$f'=\operatorname{sq}(t)$。从方波级数积分得
\[
 f(t)=C-\frac4\pi\sum_{j=0}^\infty\frac{\cos((2j+1)t)}{(2j+1)^2}.
\]
其平均值由一周期内三角形面积除以周期得 $C=\pi/2$。积分的合法性可由方波部分和在 $L^2$ 中收敛、柯西--施瓦茨不等式控制原函数的积分误差说明（详见 II.21(f) 的同一论证）；所得 $n^{-2}$ 系数级数绝对一致收敛，在角点也收敛到连续函数 $|t|$ 的周期延拓值。
\end{solution}
\originalsection{第二部分：题面与解答}
\begin{exercise}\label{cw:ps5-II-17}
原题 II.17。串联 RLC 电路的电流满足
\[
 LI''+RI'+\frac1C I=V',\qquad V(t)=V_0\sin\omega t.
\]
采用 SI：$R$ 为欧姆、$L$ 为亨利、$C$ 为法拉、$V$ 为伏特、$I$ 为安培。原 Mathlet 的滑块用毫亨（$10^{-3}$ H）、微法（$10^{-6}$ F）和毫秒（$10^{-3}$ s）。这里只研究稳态电流，不研究各元件的电压降。
(a) 找出电流与电压同相的参数；特别地 $L=500$ mH、$\omega=200$ rad/s 时，$R,C,V_0$ 应如何选？(b) 推导同相的参数关系。(c) 取 $R=100\,\Omega,L=1000$ mH、$C=100\,\mu$F、$V_0=500$ V，求最大电流幅值所对应的频率、幅值与相位滞后。(d) 对一般参数计算并验证(c)的三项结论。
\end{exercise}
\begin{solution}
解答依据（编者译审或补解，AI 辅助）：官方 PS5 解答第1页；补全省漏的 $V_0$ 和阻尼项记号。令 $R,L,C,V_0>0$、$\omega>0$。为使输入相位清楚，以复电压 $V_0e^{i\omega t}$ 的虚部表示原输入。复稳态电流与复电压的比为
\[
 H_I(\omega)=\frac{i\omega}{1/C-L\omega^2+iR\omega}
 =\frac1{R+i(L\omega-1/(C\omega))}.
\]
电流为 $I_p=\Im(V_0H_Ie^{i\omega t})$。同相要求 $H_I$ 为正实数，所以必须且只须 $L\omega=1/(C\omega)$，即 $LC\omega^2=1$。
(a)(b) 在 $L=0.5$ H、$\omega=200$ rad/s 时，$C=1/(0.5\cdot200^2)=50\,\mu$F；任意正 $R,V_0$ 均同相，电流幅值是 $V_0/R$。例如取 $R=100\,\Omega,V_0=500$ V 就得到 $5\sin200t$ A。
(c)(d) 对一般参数，电流幅值与相位滞后为
\[
 A_I(\omega)=\frac{V_0}{\sqrt{R^2+(L\omega-1/(C\omega))^2}},\qquad
 \phi_I=\arctan\frac{L\omega-1/(C\omega)}R.
\]
分母至少为 $R$，且正频率下恰在 $\omega_r=1/\sqrt{LC}$ 取等，因此此处幅值最大为 $V_0/R$、滞后为0。低频 $\phi_I<0$ 表示电流超前，高频 $\phi_I>0$ 表示滞后。给定 $L=1$ H、$C=10^{-4}$ F，得到 $\omega_r=100$ rad/s、幅值5 A、电流电压同相。下图画幅频曲线和共振时的同相波形；横轴波形采用无量纲 $s=100t$，电压除以100以便和电流对照。
\begin{center}\begin{tikzpicture}
\begin{axis}[name=rlcamp,width=.45\linewidth,height=4.4cm,axis lines=left,xlabel={$\omega$ (rad/s)},ylabel={幅值 (A)},domain=1:250,samples=180]
\addplot[blue]{500/sqrt(10000+(x-10000/x)^2)};
\end{axis}
\begin{axis}[at={(rlcamp.east)},anchor=west,xshift=1cm,width=.43\linewidth,height=4.4cm,axis lines=middle,xlabel={$s$},domain=0:13,samples=180]
\addplot[blue,thick]{5*sin(deg(x))};\addplot[green!60!black,dashed]{5*sin(deg(x))};
\end{axis}\end{tikzpicture}\end{center}
\end{solution}
\begin{exercise}\label{cw:ps5-II-18}
原题 II.18。对 $mx''+bx'+kx=0$，自然角频率为 $\omega_n=\sqrt{k/m}$，临界阻尼满足 $b/m=2\omega_n$。定义阻尼比 $\zeta$ 使 $b/m=2\omega_n\zeta$。假设 $0<\zeta<1$，观察到经过 $n$ 个周期后指数包络减为原来的一半。用 $n$ 和 $\alpha=(\ln2)/(2\pi)\simeq0.1103178$ 求 $\zeta$；算 $n=10$，并与工程近似 $\alpha/n$ 比较。
\end{exercise}
\begin{solution}
解答依据（编者译审或补解，AI 辅助）：官方 PS5 解答第1页。欠阻尼解为
\[
 x=Ae^{-\zeta\omega_n t}\cos(\omega_dt-\phi),
 \qquad\omega_d=\omega_n\sqrt{1-\zeta^2}.
\]
$n$ 个伪周期历时 $2\pi n/\omega_d$，包络比为 $\exp[-2\pi n\zeta/\sqrt{1-\zeta^2}]$。令其等于 $1/2$，得到
\[
 \frac{\zeta}{\sqrt{1-\zeta^2}}=\frac\alpha n,
 \qquad\zeta=\frac{\alpha/n}{\sqrt{1+(\alpha/n)^2}}
 =\frac\alpha{\sqrt{n^2+\alpha^2}}.
\]
$n=10$ 时，精确值约 $0.0110311087$，近似 $\alpha/10\simeq0.0110317800$；近似略大，相对误差约 $0.006085\%$。包络峰值递减与同类局部峰值的递减比例一致：$x'$ 仍是同一频率的衰减正弦，极值等间隔，同号峰值间隔一个伪周期。
\end{solution}
\begin{exercise}\label{cw:ps5-II-20}
原题 II.20。Fourier Coefficients Mathlet 可用滑块设置正弦或余弦的有限和。选择目标[B]，先判断奇偶性，决定使用正弦还是余弦、奇数项还是偶数项，并解释。
(a) 用目测调整滑块，记录近似最优系数。(b) 目标实际为周期 $2\pi$ 的函数
\[
 f(t)=\begin{cases}\pi/4,&-\pi/2<t<\pi/2,\\-\pi/4,&\pi/2<t<3\pi/2,\end{cases}
 \qquad f(\pm\pi/2)=0,
\]
按周期延拓。用积分公式计算系数，与(a)比较。(c) 从任意滑块值开始，逐个调系数，使所示的均方根距离最小，记录最佳系数并与(b)比较；解释为什么最佳拟合就是傅里叶系数，且每个系数的优化互不影响。
\end{exercise}
\begin{solution}
解答依据（编者译审或补解，AI 辅助）：官方 PS5 解答第1--2页；补充(c)正交最小化证明。本文没有运行 Mathlet 或取得个人目测记录；(a)提供由公式可重现的比较值和有限和图，不能把这些数值冒称交互实验结果。
目标 $f$ 为偶函数，所以所有 $b_n=0$。又 $f(t+\pi)=-f(t)$，故平均值与偶数频率余弦系数为0：平移半周期后，偶数频率余弦不变，而函数改变符号，积分只能为0。因此只需奇数余弦。
(b) 用约定 $f\sim a_0/2+\sum_{n\geq1}(a_n\cos nt+b_n\sin nt)$。平均值0；由偶性得
\begin{align*}
 a_n&=\frac2\pi\left(\int_0^{\pi/2}\frac\pi4\cos nt\,dt
 -\int_{\pi/2}^{\pi}\frac\pi4\cos nt\,dt\right)\\
 &=\frac1{2n}\left(2\sin(n\pi/2)-\sin n\pi\right)
 =\frac{\sin(n\pi/2)}n.
\end{align*}
所以
\[
 f(t)\sim\cos t-\frac13\cos3t+\frac15\cos5t-\frac17\cos7t+\cdots.
\]
它在连续点收敛到 $f$，在跳点收敛到左右值平均0，与题目的跳点值一致。前五个非零系数是 $1,-1/3,1/5,-1/7,1/9$，用这些值可重绘(a)要求的有限和图。
(c) 给定有限和 $S=c_0+\sum_{n=1}^N(c_n\cos nt+d_n\sin nt)$，定义
\[
 E(S)^2=\frac1{2\pi}\int_{-\pi}^{\pi}|f-S|^2\,dt.
\]
三角函数正交：不同频率的乘积积分为0，每个正弦、余弦的平方积分为 $\pi$。展开平方并利用(b)中系数定义，得到
\[
 E(S)^2=\|f\|_{\rm rms}^2-c_0^*{}^2
 -\tfrac12\sum_{n=1}^N(a_n^2+b_n^2)
 +(c_0-c_0^*)^2+\tfrac12\sum_{n=1}^N[(c_n-a_n)^2+(d_n-b_n)^2],
\]
其中 $c_0^*=a_0/2=0$、$\|f\|_{\rm rms}^2=\pi^2/16$。每个可变系数只出现在自己的一项平方中，故唯一最优值是 $c_0=0,c_n=a_n,d_n=0$，调整一个不影响另一个的最佳值。比如保留 $n=1,3,5,7,9$ 时，最小均方根距离为
\[
 \left[\frac{\pi^2}{16}-\frac12\left(1+\frac19+\frac1{25}+\frac1{49}+\frac1{81}\right)\right]^{1/2}
 \simeq0.1578537.
\]
此数值是积分公式的结果，不是软件画面的读取。有限和在跳点附近有振铃，不影响平方积分的最佳拟合性质。
\begin{center}\begin{tikzpicture}\begin{axis}[width=.8\linewidth,height=5cm,axis lines=middle,xlabel={$t$},xmin=-3.1416,xmax=3.1416,ymin=-1.05,ymax=1.05,domain=-pi:pi,samples=220]
\addplot[green!50!black,thick]coordinates{(-3.14159,-.785398)(-1.570796,-.785398)(-1.570796,.785398)(1.570796,.785398)(1.570796,-.785398)(3.14159,-.785398)};
\addplot[blue,domain=-pi:pi]{cos(deg(x))-cos(deg(3*x))/3+cos(deg(5*x))/5-cos(deg(7*x))/7+cos(deg(9*x))/9};
\addplot[only marks,mark=*,green!50!black]coordinates{(-1.570796,0)(1.570796,0)};
\end{axis}\end{tikzpicture}\end{center}
图中竖直线段仅标示跳跃位置，并非函数在该横坐标具有多个值；跳点的实际函数值由实心点0标出。
\end{solution}
\begin{exercise}\label{cw:ps5-II-21}
原题 II.21。(a) 求 $2\sin(t-\pi/3)$ 的傅里叶级数。
定义奇的 $2\pi$ 周期方波 $\operatorname{sq}(t)$：在 $0<t<\pi$ 为1，在 $-\pi<t<0$ 为 $-1$，在整数倍 $\pi$ 处为0。
(b) 把同一方波视为 $4\pi$ 周期函数，用积分公式求其傅里叶级数，比较两种周期约定。
以下各项不得重新计算傅里叶系数积分，应从方波级数通过代数和微积分求级数、描图并给出最小周期：(c) $\operatorname{sq}(t-\pi/4)$；(d) $1+2\operatorname{sq}(2\pi t)$；(e) II.20 的 $f(t)$；(f) $2\pi$ 周期函数 $g$，在 $-\pi/2\leq t\leq\pi/2$ 时 $g=t$，在 $\pi/2\leq t\leq3\pi/2$ 时 $g=\pi-t$。提示：先联系 $g'$ 与 $f$。
\end{exercise}
\begin{solution}
解答依据（编者译审或补解，AI 辅助）：官方 PS5 解答第2页；(a)补回官方答案遗漏的题面系数2，(f)说明积分级数的依据。
(a) 角差公式直接给
\[
 2\sin(t-\pi/3)=-\sqrt3\cos t+\sin t.
\]
这已是傅里叶表达式，只含第一频率；由正交性表示唯一。官方答案写的是没有因子2的 $\sin(t-\pi/3)$，不能直接作为本题答案。
(b) 对周期 $2L=4\pi$，$L=2\pi$，基函数为 $\sin(nt/2),\cos(nt/2)$。奇性给 $a_n=0$，而
\begin{align*}
 b_n&=\frac1\pi\left(\int_0^\pi\sin(nt/2)\,dt-\int_\pi^{2\pi}\sin(nt/2)\,dt\right)\\
 &=\frac2{\pi n}\left[1-2\cos(n\pi/2)+\cos n\pi\right].
\end{align*}
括号在 $n\equiv2\pmod4$ 时为4，其他整数 $n$ 时为0。因此
\[
 \operatorname{sq}(t)\sim\sum_{j=0}^\infty\frac8{\pi(4j+2)}\sin\frac{(4j+2)t}{2}
 =\frac4\pi\sum_{j=0}^\infty\frac{\sin((2j+1)t)}{2j+1}.
\]
改用较大的周期只使频率编号加倍，实际级数不变。
(c) 平移 $t$ 得
\[
 \operatorname{sq}(t-\pi/4)\sim\frac4\pi\sum_{\substack{n\geq1\\n\text{为奇数}}}
 \frac{\cos(n\pi/4)\sin nt-\sin(n\pi/4)\cos nt}{n}.
\]
最小周期仍为 $2\pi$，跳点为 $\pi/4+k\pi$，其间高低值交替。
(d) 伸缩及竖直平移给
\[
 1+2\operatorname{sq}(2\pi t)\sim1+\frac8\pi\sum_{j=0}^\infty
 \frac{\sin(2\pi(2j+1)t)}{2j+1}.
\]
最小周期1，每周期前半值3、后半值 $-1$，跳点值1。
(e) 由 $f(t)=\frac\pi4\operatorname{sq}(t+\pi/2)$，得
\[
 f(t)\sim\sum_{j=0}^\infty\frac{(-1)^j}{2j+1}\cos((2j+1)t),
\]
最小周期 $2\pi$，与 II.20 由积分所得结果一致。
(f) $g$ 是连续奇三角波，除角点外 $g'=(4/\pi)f$，故积分后候选为
\[
 g(t)=\frac4\pi\sum_{j=0}^\infty\frac{(-1)^j}{(2j+1)^2}\sin((2j+1)t).
\]
这里可严谨地积分级数：$f$ 的部分和在 $L^2[-\pi,\pi]$ 中收敛到 $f$；对任意 $t$，积分误差的绝对值由柯西--施瓦茨不等式控制为 $\sqrt{2\pi}\,\|f-S_N\|_2$，故积分一致收敛。每个积分后的系数为 $O(n^{-2})$，级数也绝对一致收敛；常数由 $g(0)=0$ 确定。所得积分在每个开区间的导数为 $\pm1$，且连续接合，正是题给 $g$，最小周期 $2\pi$。不能在角点声称 $g'$ 存在。

(c)、(d)、(e)、(f)的函数图如下，(e)也可参照 II.20 的较大图。竖线仅为跳跃位置，(c)跳点取0、(d)跳点取1、(e)跳点取0；(f)是连续折线。
\begin{center}\begin{tikzpicture}
\begin{axis}[name=cplot,width=.43\linewidth,height=3.7cm,axis lines=middle,title={(c)},xlabel={$t$},xmin=-3.1416,xmax=3.1416,ymin=-1.3,ymax=1.3]
\addplot[blue,thick]coordinates{(-3.14159,1)(-2.35619,1)(-2.35619,-1)(.785398,-1)(.785398,1)(3.14159,1)};
\addplot[blue,only marks]coordinates{(-2.35619,0)(.785398,0)};
\end{axis}
\begin{axis}[name=dplot,at={(cplot.east)},anchor=west,xshift=1.2cm,width=.43\linewidth,height=3.7cm,axis lines=middle,title={(d)},xlabel={$t$},xmin=0,xmax=2,ymin=-1.5,ymax=3.5]
\addplot[blue,thick]coordinates{(0,3)(.5,3)(.5,-1)(1,-1)(1,3)(1.5,3)(1.5,-1)(2,-1)};
\addplot[blue,only marks]coordinates{(0,1)(.5,1)(1,1)(1.5,1)(2,1)};
\end{axis}
\begin{axis}[name=eplot,at={(cplot.south)},anchor=north,yshift=-1.2cm,width=.43\linewidth,height=3.7cm,axis lines=middle,title={(e)},xlabel={$t$},xmin=-3.1416,xmax=3.1416,ymin=-1,ymax=1]
\addplot[blue,thick]coordinates{(-3.14159,-.785398)(-1.570796,-.785398)(-1.570796,.785398)(1.570796,.785398)(1.570796,-.785398)(3.14159,-.785398)};
\addplot[blue,only marks]coordinates{(-1.570796,0)(1.570796,0)};
\end{axis}
\begin{axis}[at={(dplot.south)},anchor=north,yshift=-1.2cm,width=.43\linewidth,height=3.7cm,axis lines=middle,title={(f)},xlabel={$t$},xmin=-3.1416,xmax=3.1416,ymin=-1.9,ymax=1.9]
\addplot[blue,thick]coordinates{(-3.14159,0)(-1.570796,-1.570796)(1.570796,1.570796)(3.14159,0)};
\end{axis}\end{tikzpicture}\end{center}
\end{solution}

% PS6--9 source ledger: research/ps06-09-coverage.json
\originalchapter{作业六：周期响应、阶跃与卷积}
本章保留原作业的讲次题号。$u$ 表示单位阶跃，$\delta$ 表示单位冲激；响应均取因果响应，卷积取 $\int_0^t f(t-\tau)g(\tau)\,d\tau$。第二部分以官方解为依据逐问核对，补充必要推导；明确指出的补解及修正不冒充原答案。
\originalsection{第一部分}
\begin{exercise}\label{cw:ps6-I-22}原题 I.22（3 月 31 日）：完成 Notes 7C-1、7C-2。\end{exercise}
\begin{solution}
\textit{编者补解及题库回指核对（AI 辅助）；教材题面缺项另标。}\par

完整题面和解答见习题 \ref{cw:bank-7C-1} 的 (a)、(b)、(c) 及习题 \ref{cw:bank-7C-2}。前者依次检查三种强迫的谐波与固有频率是否重合；后者求周期特解。此处保留原指定及子问映射，不重复排印同一题。
\end{solution}
原题 I.23、I.25 均写明无题。
\begin{exercise}\label{cw:ps6-I-24}原题 I.24（4 月 5 日）：分别求 (a) $D+kI$、(b) $D^2+\omega_n^2I$ 的单位冲激响应与单位阶跃响应。\end{exercise}
\begin{solution}
\textit{原题附答译审与补证（AI 辅助）。}\par

(a) 对 $t>0$，齐次方程为 $w'+kw=0$；跨过零点积分得 $w(0+)=1$，故
\[w=u(t)e^{-kt},\qquad v=\int_0^tw(\tau)\,d\tau=\begin{cases}u(t)(1-e^{-kt})/k,&k\ne0,\\u(t)t,&k=0.\end{cases}\]
(b) 在 $\omega_n>0$ 时，冲激不使位移跳跃，但使速度跳跃 $1$，所以 $w(0+)=0,w'(0+)=1$，得到
\[w=u(t)\frac{\sin\omega_nt}{\omega_n},\qquad v=u(t)\frac{1-\cos\omega_nt}{\omega_n^2}.\]
直接求导得到 $v'=w$；$v(0+)=v'(0+)=0$，且 $v''+\omega_n^2v=1$。若 $\omega_n=0$，连续极限分别为 $u(t)t$、$u(t)t^2/2$。
\end{solution}
\originalsection{第二部分}
\begin{exercise}\label{cw:ps6-II-22}原题 II.22（周期解）：$g$ 的周期为 $2\pi$，在 $-\pi/2\le t\le\pi/2$ 时 $g(t)=t$，在 $\pi/2\le t\le3\pi/2$ 时 $g(t)=\pi-t$（即前一作业 II.21(f) 的函数）。
(a) 求 $x''+\omega_n^2x=g(t)$ 的一个周期解，若存在。
(b) 哪些正 $\omega_n$ 使周期解不存在？
(c) 将 (b) 最小值记作 $\omega_r$，当 $\omega_n$ 略小于、略大于 $\omega_r$ 时，解大致如何？
(d) 哪些 $\omega_n$ 使周期解不唯一？
(e) 对 (d) 的值，是否所有解都周期？\end{exercise}
\begin{solution}
\textit{官方解译审与补证（AI 辅助）；其中另标编者补解或原文修正。}\par

由奇性与分段积分，$g$ 只有奇数正弦谐波，且
\[g(t)=\frac4\pi\sum_{j=0}^{\infty}\frac{(-1)^j}{(2j+1)^2}\sin((2j+1)t).\]
(a) 若 $\omega_n$ 不是正奇整数，可逐谐波除以 $\omega_n^2-(2j+1)^2$：
\[x_p(t)=\frac4\pi\sum_{j=0}^{\infty}\frac{(-1)^j\sin((2j+1)t)}{(2j+1)^2[\omega_n^2-(2j+1)^2]}.\]
系数为 $O(j^{-4})$，两次导数的系数为 $O(j^{-2})$，因此上述级数及两次求导均一致收敛，所得为实际 $2\pi$ 周期解。
(b) 若 $\omega_n=N$ 为正奇整数，强迫的 $\sin Nt$ 系数非零。相应特解有非零的 $t\cos Nt$ 项，齐次解只有有界正弦、余弦，无法消除此无界项，故无周期解。
(c) $\omega_r=1$。靠近 $1$ 时第一项支配：
\[x_p(t)\approx\frac{4\sin t}{\pi(\omega_n^2-1)}.\]
左侧系数负，近似反相且振幅很大；右侧系数正，近似同相且振幅很大。其余谐波在邻域内一致有界。
(d) 一般解为 $x_p+A\cos\omega_nt+B\sin\omega_nt$。若非零齐次项与 $x_p$ 共存，则其频率与强迫的基本频率须可公度。严格地说，若和的周期为 $T$，则其方程右端也以 $T$ 为周期，三角波的周期必须是 $2\pi$ 的整数倍；再由齐次项周期得到 $\omega_nT\in2\pi\mathbb Z$。所以恰在 $\omega_n\in\mathbb Q_{>0}$ 且不是正奇整数时，有多个周期解。
(e) 是。写 $\omega_n=p/q$（互素正整数），所有解均以 $2\pi q$ 为周期。若 $\omega_n$ 为正偶整数，$2\pi$ 已是所有解的周期；若 $\omega_n$ 无理，只有 $x_p$ 这一周期解。
\end{solution}
\begin{exercise}\label{cw:ps6-II-23}原题 II.23（阶跃与冲激）：对下列每个 $f$，(i) 画函数图；(ii) 画广义导数图；(iii) 用 $u(t-a)$、$\delta(t-a)$ 及普通函数写出 $f$、$f'$（少数点可不定义）。
(a) $t<0$ 时 $f=0$，$t>0$ 时 $f=-t$。
(b) $t<0$ 时 $f=0$，$t>0$ 时 $f=1-t$。
(c) $t<0$ 时 $f=0$，$0<t<1$ 时 $f=2t-1$，$t>1$ 时 $f=0$。
(d) $t<0$ 时 $f=0$，$t>0$ 时 $f=t-\lfloor t\rfloor$，$\lfloor t\rfloor$ 是不超过 $t$ 的最大整数。\end{exercise}
\begin{solution}
\textit{官方解译审与补证（AI 辅助）；其中另标编者补解或原文修正。}\par

广义导数由各光滑区间的斜率与各跳点的跳跃量组成：$f(a+)-f(a-)$ 乘 $\delta(t-a)$。于是
\begin{align*}
\text{(a)}\quad&f=-tu(t),& f'&=-u(t),\\
\text{(b)}\quad&f=(1-t)u(t),& f'&=-u(t)+\delta(t),\\
\text{(c)}\quad&f=(2t-1)[u(t)-u(t-1)],&f'&=2[u(t)-u(t-1)]-\delta(t)-\delta(t-1),\\
\text{(d)}\quad&f=tu(t)-\sum_{n=1}^{\infty}u(t-n),&f'&=u(t)-\sum_{n=1}^{\infty}\delta(t-n).
\end{align*}
(d) 的和在任何有界区间都只有有限非零项；$0$ 处连续，不产生冲激，每个正整数处从 $1$ 跳到 $0$，故冲激权重为 $-1$。下图将冲激画为带权箭头，箭头不是普通函数值。
\begin{center}
\begin{tikzpicture}[x=1.1cm,y=.65cm,>=Stealth,font=\small]
\foreach \j/\title in {0/{(a)},1/{(b)},2/{(c)},3/{(d)}}{
\begin{scope}[shift={(3.1*\j,0)}]
\draw[->](-.3,0)--(2.4,0) node[right]{$t$};\draw[->](0,-1.3)--(0,1.5);\node at (1,1.8){\title：$f$};
\end{scope}}
\draw[main,thick](0,0)--(1.2,-1.2);
\draw[main,thick](3.1,1)--(5.1,-1);
\draw[main,thick](6.2,-1)--(7.2,1);\draw[main,thick](7.2,0)--(8.4,0);
\foreach \n in {0,1}{\draw[main,thick](9.3+\n,0)--(10.3+\n,1);}\node at (11.5,.5){$\cdots$};
\foreach \j/\title in {0/{(a)},1/{(b)},2/{(c)},3/{(d)}}{
\begin{scope}[shift={(3.1*\j,-4)}]\draw[->](-.3,0)--(2.4,0) node[right]{$t$};\draw[->](0,-1.3)--(0,1.5);\node at (1,1.8){\title：$f'$};\end{scope}}
\draw[main,thick](0,-5)--(2.2,-5);
\draw[main,thick](3.1,-5)--(5.3,-5);\draw[->,thick](3.1,-4)--(3.1,-3) node[above right]{$+1$};
\draw[main,thick](6.2,-2.7)--(7.2,-2.7);\draw[main,thick](7.2,-4)--(8.4,-4);\foreach \x in {6.2,7.2}{\draw[->,thick](\x,-4)--(\x,-5) node[below]{$-1$};}\node at (6.7,-2.4){$2$};
\draw[main,thick](9.3,-3.3)--(11.5,-3.3);\foreach \x in {10.3,11.3}{\draw[->,thick](\x,-4)--(\x,-5) node[below]{$-1$};}\node at (9.7,-3){$1$};
\end{tikzpicture}
\end{center}
\end{solution}
\begin{exercise}\label{cw:ps6-II-24}原题 II.24：
(a) 求 $2D^2+4D+4I$ 的单位冲激响应 $w$。
(b) 求其单位阶跃响应 $v$。
(c) 验证 $v'=w$。
(d) 分别求以 (i) $2u(t)$、(ii) $u(t)t$、(iii) $u(t)t^2$ 为单位冲激响应的 LTI 微分算子 $p(D)$。\end{exercise}
\begin{solution}
\textit{官方解译审与补证（AI 辅助）；其中另标编者补解或原文修正。}\par

(a) 特征根为 $-1\pm i$，故 $t>0$ 时 $w=e^{-t}(A\cos t+B\sin t)$。冲激要求 $w(0+)=0$、$2w'(0+)=1$，于是
\[w=\tfrac12u(t)e^{-t}\sin t.\]
(b) 常数特解为 $1/4$；将齐次项加上后，以静止条件 $v(0+)=v'(0+)=0$ 确定系数：
\[v=\tfrac14u(t)[1-e^{-t}(\cos t+\sin t)].\]
(c) 因括号在零点为零，求广义导数不产生冲激，且
\[v'=\tfrac14u(t)e^{-t}[(\cos t+\sin t)-(\cos t-\sin t)]=w.\]
(d) 直接应用广义导数：(i) $D(2u)=2\delta$，故 $p(D)=D/2$；(ii) $D(ut)=u$、$D^2(ut)=\delta$，故 $p(D)=D^2$；(iii) $D(ut^2)=2ut$、$D^2(ut^2)=2u$、$D^3(ut^2)=2\delta$，故 $p(D)=D^3/2$。这些算子的低阶系数必须为零，因为在 $t>0$ 上作用所得多项式恒为零；更高阶非零项会产生额外的冲激导数，亦不允许。
\end{solution}
\begin{exercise}\label{cw:ps6-II-25}原题 II.25（卷积）：
(a) 令 $q(t)=\cos\omega t$，用 I.24 的 $D+kI$ 冲激响应求 $w*q$，验证它是 $x'+kx=q$ 的静止初值解。
(b) 令 $q(t)=1$，用 I.24 的 $D^2+\omega_n^2I$ 冲激响应求 $w*q$，验证它是 $x''+\omega_n^2x=q$ 的静止初值解。
(c) 计算 $t^2*t$、$t*t^2$，比较。
(d) 计算 $(t*t)*t$、$t*(t*t)$，比较。\end{exercise}
\begin{solution}
\textit{官方解译审与补证（AI 辅助）；其中另标编者补解或原文修正。}\par

(a) 当 $k^2+\omega^2>0$ 时，
\begin{align*}
x(t)&=e^{-kt}\operatorname{Re}\int_0^t e^{(k+i\omega)\tau}\,d\tau\\
&=\frac{k\cos\omega t+\omega\sin\omega t-ke^{-kt}}{k^2+\omega^2}.
\end{align*}
其导数为 $[-k\omega\sin\omega t+\omega^2\cos\omega t+k^2e^{-kt}]/(k^2+\omega^2)$，加 $kx$ 恰为 $\cos\omega t$，而 $x(0)=0$。$k=\omega=0$ 时直接积分得 $x=t$。
(b) $x=\omega_n^{-1}\int_0^t\sin\omega_n(t-\tau)\,d\tau=(1-\cos\omega_nt)/\omega_n^2$。于是 $x'=\sin\omega_nt/\omega_n$、$x''=\cos\omega_nt$，方程与两个零初值均成立。
(c) 展开被积式：
\[t^2*t=\int_0^t(t-\tau)^2\tau\,d\tau=(\tfrac12-\tfrac23+\tfrac14)t^4=\frac{t^4}{12},\]
\[t*t^2=\int_0^t(t-\tau)\tau^2\,d\tau=(\tfrac13-\tfrac14)t^4=\frac{t^4}{12}.\]
(d) $t*t=\int_0^t(t-\tau)\tau\,d\tau=t^3/6$。因此
\[(t*t)*t=\tfrac16\int_0^t(t-\tau)^3\tau\,d\tau=\frac{t^5}{120},\quad t*(t*t)=\tfrac16\int_0^t(t-\tau)\tau^3\,d\tau=\frac{t^5}{120}.\]
这些计算分别展示交换律与结合律，所有输入在此均按 $t\ge0$ 的因果函数理解。
\end{solution}
\originalchapter{作业七：拉普拉斯变换}
\originalsection{第一部分}
\begin{exercise}\label{cw:ps7-I-26}原题 I.26（4 月 9 日）：令 $z$ 为复数。从定义求 $\mathcal L[e^{zt}]$，并求积分的收敛域。\end{exercise}
\begin{solution}
\textit{编者补解及题库回指核对（AI 辅助）；教材题面缺项另标。}\par

\[\int_0^R e^{-(s-z)t}\,dt=\frac{1-e^{-(s-z)R}}{s-z}\quad(s\ne z).\]
极限存在恰需 $\operatorname{Re}(s-z)>0$，此时 $F(s)=1/(s-z)$。若实部小于零，边界项增长；若实部等于零而 $s\ne z$，边界项振荡无极限；$s=z$ 时积分为 $R$。因此收敛域为 $\operatorname{Re}s>\operatorname{Re}z$。
\end{solution}
\begin{exercise}\label{cw:ps7-I-27}原题 I.27（4 月 12 日）：完成 Notes 3A-9（(a) 先按建议用 $s$ 平移规则，再将 $e^{-t}\sin3t$ 写成复指数的线性组合重做）、Notes 3B-1；EP 第 4.3 节第 1、5 题。\end{exercise}
\begin{solution}
\textit{编者补解及题库回指核对（AI 辅助）；教材题面缺项另标。}\par

Notes 3A-9 的 (a)、(b) 见习题 \ref{cw:bank-3A-9}，Notes 3B-1 的 (a)--(e) 见习题 \ref{cw:bank-3B-1}。本作业额外要求的重做为
\[e^{-t}\sin3t=\frac{e^{(-1+3i)t}-e^{(-1-3i)t}}{2i},\]
\[\mathcal L[e^{-t}\sin3t]=\frac1{2i}\left(\frac1{s+1-3i}-\frac1{s+1+3i}\right)=\frac3{(s+1)^2+9},\quad\operatorname{Re}s>-1.\]
EP 4.3 第 1、5 题仅有教材定位，未取得可核实的公开题面（\texttt{external-textbook-statement-unavailable}），故此处不猜写题面或解答。
\end{solution}
\begin{exercise}\label{cw:ps7-I-28}原题 I.28（4 月 14 日）：完成 EP 第 4.3 节第 7、8、12、13、28、36 题。\end{exercise}
\begin{solution}
\textit{编者补解及题库回指核对（AI 辅助）；教材题面缺项另标。}\par

以上六项保留完整指定；公开作业未包含它们的题面，当前公开题源亦未取得相应教材正文，状态均为 \texttt{external-textbook-statement-unavailable}。不能用题号推测或借另一版教材同号题代替。
\end{solution}
\originalsection{第二部分}
\begin{exercise}\label{cw:ps7-II-26}原题 II.26：
(a) $F(s)=\mathcal L[f](s)$，$a>0$。用积分定义及换元求 $g(t)=f(at)$ 的变换，并用 $f(t)=t^n$ 的两边变换验证。
(b) 用二重积分及换元 $x=t-\tau,y=\tau$ 证明 $h=f*g$ 时 $H=FG$；从 $F=\int_0^\infty f(x)e^{-sx}\,dx$、$G=\int_0^\infty g(y)e^{-sy}\,dy$ 出发。
(c) 原题写作「$0<t<1$ 时 $f(t)=1$，$t>0$ 时 $f(t)=0$」，要求从定义求变换及收敛域。\end{exercise}
\begin{solution}
\textit{官方解译审与补证（AI 辅助）；其中另标编者补解或原文修正。}\par

(c) 原文两个条件在 $(0,1)$ 相互矛盾；官方解将后一条件按 $t>1$ 处理。以下明确采用该修正，即单位矩形脉冲，保留这一差异记录。
(a) 令 $v=at$，则
\[G(s)=\int_0^\infty f(at)e^{-st}\,dt=\frac1a\int_0^\infty f(v)e^{-(s/a)v}\,dv=\frac1aF(s/a).\]
若 $F$ 在 $\operatorname{Re}s>\sigma$ 收敛，$G$ 在 $\operatorname{Re}s>a\sigma$ 收敛。$n$ 为非负整数时，$F=n!/s^{n+1}$、$G=a^nn!/s^{n+1}$，而 $a^{-1}F(s/a)=a^nn!/s^{n+1}$，一致。
(b) 在 $\operatorname{Re}s$ 足够大、两积分绝对收敛时，Fubini 定理允许写
\[FG=\int_0^\infty\int_0^\infty f(x)g(y)e^{-s(x+y)}\,dx\,dy.\]
新区域为 $0\le\tau\le t<\infty$，Jacobi 行列式 $\det\begin{pmatrix}1&-1\\0&1\end{pmatrix}=1$。故
\[FG=\int_0^\infty e^{-st}\left[\int_0^t f(t-\tau)g(\tau)\,d\tau\right]dt=H(s).\]
绝对收敛的假设确保这里的乘积、换序与换元合法。
(c) 修正后的变换为
\[F(s)=\int_0^1 e^{-st}\,dt=\begin{cases}(1-e^{-s})/s,&s\ne0,\\1,&s=0.\end{cases}\]
这是有限区间积分，对全部复 $s$ 收敛；$s=0$ 是表达式的可去奇点，并非极点。
\end{solution}
\begin{exercise}\label{cw:ps7-II-27}原题 II.27：$a,b$ 为实数，$a\ne0$。
(a) 对 $aD+bI$，通过拉普拉斯变换到频域求解再变回，求单位冲激响应及阶跃响应。宜使用因果广义导数规则 $\mathcal L[x']=sX$。
(b) 用三种方法求 $ax'+bx=t$ 的静止初值解：(i) 待定系数加适当瞬态；(ii) 用 (a) 的 $w$ 求 $w*t$；(iii) 拉普拉斯变换。\end{exercise}
\begin{solution}
\textit{官方解译审与补证（AI 辅助）；其中另标编者补解或原文修正。}\par

(a) $(as+b)W=1$，故
\[W=\frac1{as+b},\quad w=\frac1a u(t)e^{-bt/a},\quad V=\frac1{s(as+b)}.\]
当 $b\ne0$ 时，$V=b^{-1}(s^{-1}-(s+b/a)^{-1})$，所以 $v=u(t)(1-e^{-bt/a})/b$；当 $b=0$ 时 $v=u(t)t/a$。
(b)(i) 设 $x_p=ct+d$，得到 $bc=1$、$ac+bd=0$。$b\ne0$ 时 $x_p=t/b-a/b^2$，加 $Ce^{-bt/a}$ 并令 $x(0)=0$ 得
\[x=u(t)\left[\frac tb-\frac a{b^2}(1-e^{-bt/a})\right].\]
$b=0$ 时直接积分 $x'=t/a$，得 $x=u(t)t^2/(2a)$。
(ii) $b\ne0$ 时，分部积分给出
\begin{align*}
(w*t)(t)&=\frac1a e^{-bt/a}\int_0^t\tau e^{b\tau/a}\,d\tau\\
&=\frac1a e^{-bt/a}\left[\frac{a\tau}{b}e^{b\tau/a}-\frac{a^2}{b^2}e^{b\tau/a}\right]_0^t\\
&=\frac tb-\frac a{b^2}(1-e^{-bt/a}).
\end{align*}
$b=0$ 时积分为 $a^{-1}\int_0^t\tau\,d\tau=t^2/(2a)$。
(iii) $X=1/[s^2(as+b)]$。$b\ne0$ 时分解为
\[X=-\frac a{b^2s}+\frac1{bs^2}+\frac a{b^2(s+b/a)},\]
反变换得同一公式；$b=0$ 时 $X=1/(as^3)$。在 $t>0$，$x'=(1-e^{-bt/a})/b$，代入 $ax'+bx=t$，同时 $x(0+)=0$；三种构造的解确实一致。
\end{solution}
\begin{exercise}\label{cw:ps7-II-28}原题 II.28：用拉普拉斯变换求以下算子的单位冲激响应：(a) $3D^2+6D+6I$；(b) $D^4+I$。\end{exercise}
\begin{solution}
\textit{官方解译审与补证（AI 辅助）；其中另标编者补解或原文修正。}\par

(a) $W=[3((s+1)^2+1)]^{-1}$，故 $w=u(t)e^{-t}\sin t/3$。
(b) 原题确为 $D^4+I$。官方解只列出四个根及部分分式框架，并说明原本想问 $D^4-I$；下面是对实际题面的编者补解。令 $\alpha=1/\sqrt2$，则
\[s^4+1=(s^2-\sqrt2s+1)(s^2+\sqrt2s+1),\]
\[\frac1{s^4+1}=\frac{-s/(2\sqrt2)+1/2}{s^2-\sqrt2s+1}+\frac{s/(2\sqrt2)+1/2}{s^2+\sqrt2s+1}.\]
将两分母写成 $(s\mp\alpha)^2+\alpha^2$，反变换并合并指数项，得
\[w(t)=\frac{u(t)}{\sqrt2}\left[\cosh(\alpha t)\sin(\alpha t)-\sinh(\alpha t)\cos(\alpha t)\right].\]
四个指数根为 $\pm\alpha\pm i\alpha$，每项在 $t>0$ 满足 $w^{(4)}+w=0$。更直接地，括号的展开从 $(2/3)(\alpha t)^3$ 开始，故 $w=t^3/6+O(t^7)$，于是 $w(0+)=w'(0+)=w''(0+)=0$、$w'''(0+)=1$。因果延拓的四阶导数恰产生单位冲激，验证了所求响应。若改问官方所述的 $D^4-I$，答案才是 $u(t)(\sinh t-\sin t)/2$；二者不可混用。
\end{solution}
\originalchapter{作业八：极点与线性系统}
原作业分两半发布，以下合并但保留讲次 29、32、33；讲次 30 无题、31 为第三次考试，不补造题目。第一半将 L32 日期写成 4 月 28 日，第二半实际题目写 4 月 26 日，以下依第二半。
\originalsection{第一部分}
\begin{exercise}\label{cw:ps8-I-29}原题 I.29：
(a) 用 $t$ 平移规则求 $f(t)=[u(t)-u(t-2\pi)]\sin t$ 的拉普拉斯变换。
(b) 画下列函数变换的极点图：(i) $f=1$；(ii) $f=e^{-t}+3e^{-3t}$；(iii) $f=\cos2t+e^{-t}\sin t$。\end{exercise}
\begin{solution}
\textit{原题附答译审与补证（AI 辅助）。}\par

(a) 令 $g=u(t)\sin t$，则 $g(t-2\pi)=u(t-2\pi)\sin t$，所以
\[F(s)=\frac{1-e^{-2\pi s}}{s^2+1}.\]
虽然分母在 $\pm i$ 为零，分子亦为零且抵消；$f$ 有紧支集，$F$ 在整个复平面解析，故这两个点不是极点。
(b) 三个变换分别为
\[\frac1s,\quad \frac1{s+1}+\frac3{s+3},\quad\frac{s}{s^2+4}+\frac1{(s+1)^2+1}.\]
无抵消，其极点依次为 $\{0\}$、$\{-1,-3\}$、$\{\pm2i,-1\pm i\}$。极点图如下，横轴为实部、纵轴为虚部，叉号为极点。
\begin{center}\begin{tikzpicture}[x=.8cm,y=.65cm,font=\small]
\foreach \a/\ttl in {0/{(i)},5/{(ii)},11/{(iii)}}{\begin{scope}[shift={(\a,0)}]\draw[->](-3.3,0)--(1,0) node[right]{$\operatorname{Re}s$};\draw[->](0,-2.3)--(0,2.5) node[above]{$\operatorname{Im}s$};\node at (-1,3){\ttl};\end{scope}}
\node[main] at (0,0){$\times$};\node[main] at (2,0){$\times$};\node[main] at (4,0){$\times$};\node[below] at (2,0){$-3$};\node[below] at (4,0){$-1$};
\foreach \x/\y in {11/2,11/-2,10/1,10/-1}{\node[main] at (\x,\y){$\times$};}\node[right] at (11,2){$2i$};\node[right] at (11,-2){$-2i$};\node[above left] at (10,1){$-1+i$};\node[below left] at (10,-1){$-1-i$};
\end{tikzpicture}\end{center}
\end{solution}
\begin{exercise}\label{cw:ps8-I-32}原题 I.32：
(a) 计算五个矩阵乘积
\[\begin{pmatrix}1&2\end{pmatrix}\binom xy,\quad\binom12\begin{pmatrix}x&y\end{pmatrix},\quad\begin{pmatrix}a&b\\c&d\end{pmatrix}\binom xy,\quad\begin{pmatrix}1&2\end{pmatrix}\begin{pmatrix}a&b\\c&d\end{pmatrix},\quad\begin{pmatrix}a&b\\c&d\end{pmatrix}\begin{pmatrix}x&u\\y&v\end{pmatrix}.\]
(b) 完成 Notes 4A-2、4B-1、4B-6。\end{exercise}
\begin{solution}
\textit{编者补解及题库回指核对（AI 辅助）；教材题面缺项另标。}\par

(a) 每一项用左矩阵的一行与右矩阵的一列作内积，依次得
\[x+2y,\quad \begin{pmatrix}x&y\\2x&2y\end{pmatrix},\quad\binom{ax+by}{cx+dy},\quad\begin{pmatrix}a+2c&b+2d\end{pmatrix},\quad\begin{pmatrix}ax+by&au+bv\\cx+dy&cu+dv\end{pmatrix}.\]
(b) 见习题 \ref{cw:bank-4A-2}、\ref{cw:bank-4B-1} 的 (a)、(b)、\ref{cw:bank-4B-6} 的 (a)、(b)，后一题两种消元方法须比较。
\end{solution}
\begin{exercise}\label{cw:ps8-I-33}原题 I.33：完成 Notes 4C-1；可参考 EP 第 344 页的 $3\times3$ 行列式计算说明。\end{exercise}
\begin{solution}
\textit{编者补解及题库回指核对（AI 辅助）；教材题面缺项另标。}\par
完整的三个矩阵题及解答见习题 \ref{cw:bank-4C-1} 的 (a)、(b)、(c)。EP 页码在此只作为计算方法的参考指定，不作为另一个未给题面的作业题。\end{solution}
\originalsection{第二部分}
\begin{exercise}\label{cw:ps8-II-29}原题 II.29（极点）：
(a) 对四个极点图 (1) $\{1,i,-i\}$；(2) $\{-1+4i,-1-4i\}$；(3) $\{-1\}$；(4) 空图，分别 (i) 描述具有该极点图的所有 $f$ 的共同特征；(ii) 写出两个 $f,F$ 的例子。
(b) 埃塞俄比亚的一次考古发掘发现一套机械系统，为避免拆开，研究者施加单位冲激，测得 $w(t)=u(t)e^{-t/2}\sin(3t/2)$。
(i) 求系统的特征多项式与传递函数 $W(s)$。
(ii) 画极点图。
(iii) 运送卡车的振动会激起系统振动，应避开哪个使响应振幅最大的圆频率？原题此处再次标成 (ii)，此册按第三小问记。
(iv) 在 Mathlet「Amplitude Response and Pole Diagram」中设定系统参数，检查峰值并转动三维图，解释黄色、绿色盒状图，黄色盒底曲线，红箭头及黄菱形。第二半纠正：Mathlet 的 $b,k$ 应输入 (i) 的 $b/m,k/m$。\end{exercise}
\begin{solution}
\textit{官方解译审与补证（AI 辅助）；其中另标编者补解或原文修正。}\par

(a) 需要先说明推断范围。若 $F$ 为严格真有理函数、图中各极点均为单极点、$f$ 为实函数，则部分分式反演给出以下共同结构。仅列极点位置、不列阶数或允许非有理变换时，不能声称这些式子描述「所有函数」：重极点引入 $t$ 的多项式因子，整函数项还可加入无新极点的时域成分。以下结构与官方解采用的简单指数模式相同，同时说明其适用范围。
(1) $f=Ae^t+B\cos t+C\sin t$，$A\ne0$ 且 $(B,C)\ne(0,0)$。长期增长由 $e^t$ 支配，另有圆频率 $1$ 的恒振幅振荡。两个例子为
\[f_1=e^t+\sin t,\quad F_1=\frac1{s-1}+\frac1{s^2+1};\qquad f_2=2e^t+\cos t,\quad F_2=\frac2{s-1}+\frac{s}{s^2+1}.\]
(2) $f=e^{-t}(B\cos4t+C\sin4t)$，$(B,C)\ne(0,0)$，是包络按 $e^{-t}$ 衰减的圆频率 $4$ 振荡。例子：
\[f_1=e^{-t}\sin4t,\ F_1=\frac4{(s+1)^2+16};\qquad f_2=e^{-t}\cos4t,\ F_2=\frac{s+1}{(s+1)^2+16}.\]
(3) $f=Ae^{-t}$、$A\ne0$，无振荡并指数衰减；可取 $A=1,2$，相应 $F=1/(s+1),2/(s+1)$。
(4) 空图只表明变换无极点，不能仅据此断定每种 $f$ 的长期衰减速度。若限于上述严格真有理类，只能有零函数；允许一般因果函数或广义函数时，可取
\[f_1=\delta(t-1),\quad F_1=e^{-s};\qquad f_2=u(t)-u(t-1),\quad F_2=\frac{1-e^{-s}}s,\quad F_2(0)=1.\]
两者均在有限时间后为零，但这只是这些例子的性质。官方解将空极点图概括成「比任意指数衰减更快」，此结论需要额外函数类假设，不从极点位置单独推出。
(b)(i) 变换给出
\[W(s)=\frac{3/2}{(s+1/2)^2+9/4}=\frac{3/2}{s^2+s+5/2},\quad p(s)=\frac23s^2+\frac23s+\frac53.\]
因此机械参数为 $m=2/3,b=2/3,k=5/3$。也可从 $w'(0+)=3/2=1/m$ 确定首项系数，不能只写首一特征多项式而丢失传递函数的增益。
(ii) 极点为 $-1/2\pm3i/2$。
(iii) 对单位幅度谐波，增益
\[|W(i\omega)|=\frac{3/2}{\sqrt{(5/2-\omega^2)^2+\omega^2}}.\]
令 $z=\omega^2\ge0$，分母平方为 $(5/2-z)^2+z=(z-2)^2+9/4$，故唯一正峰值频率为 $\omega_r=\sqrt2$，最大增益为 $1$。
(iv) 以下依据函数和官方图例作解析解释，未声称实际操作了 Mathlet。输入其归一化参数 $b/m=1,k/m=5/2$；其首一系统增益与实际 $W$ 相差常数 $2/3$，峰值频率一致。黄色盒状图处于虚轴上方的竖直平面，按官方图例，其底边界曲线是振幅响应 $(i\omega,|W(i\omega)|)$。绿色盒状图处于实轴上方的平面，顶部是三维曲面 $|W(s)|$ 在实轴上的截线、底部是实轴。黄色盒的底曲线将输入频率与相应增益配对，并不是虚轴本身，红箭头位于 $W$ 的极点上方，黄菱形在 $(\pm i\omega,|W(i\omega)|)$ 处标记所选输入频率与相应增益。
\begin{center}\begin{tikzpicture}
\begin{axis}[width=6cm,height=4.8cm,xmin=-1.6,xmax=.5,ymin=-2,ymax=2,axis lines=middle,xlabel={$\operatorname{Re}s$},ylabel={$\operatorname{Im}s$},title={极点图},xtick={-.5},ytick={-1.5,1.5}]
\addplot[only marks,mark=x,mark size=4pt,main] coordinates{(-.5,-1.5)(-.5,1.5)};
\end{axis}
\begin{axis}[at={(6.2cm,0)},width=6cm,height=4.8cm,xmin=0,xmax=4,ymin=0,ymax=1.15,axis lines=left,xlabel={$\omega$},ylabel={增益},title={实际系统振幅响应},xtick={0,1.4142,3},xticklabels={$0$,$\sqrt2$,$3$}]
\addplot[main,thick,domain=0:4,samples=120]{1.5/sqrt((2.5-x*x)^2+x*x)};\addplot[only marks] coordinates{(1.41421356,1)};
\end{axis}\end{tikzpicture}\end{center}
\end{solution}
\begin{exercise}\label{cw:ps8-II-32}原题 II.32（线性系统与矩阵）：研究 $x''+4x'+3x=0$、$x''+x'+\tfrac52x=0$。
(a) 每个方程求两个线性无关实解 $x_1,x_2$（取指数或 $e^{rt}\cos\omega t,e^{rt}\sin\omega t$），写一般实解、阻尼类型，并算 $x_1',x_2'$。
(b) 写两个伴随矩阵。
(c) 在 Mathlet「Linear Phase Portraits: Matrix Entry」选择伴随矩阵，将 $c,d$ 设为第一方程的对应值；可在左上 $(d,-c)$ 平面调整，点击相平面增添轨道，再清除并恢复初始轨道。画相平面，标时间箭头，指出 (a) 的两个基本解所对应的轨道。
(d) 若图示最大 $x=2$，一条钩形轨道在 $t=0$ 经过 $x$ 轴的 $x=1$ 点，求该解并画原方程的 $x(t)$。
(e) 写同一轨道上的一般解；可设在 $t=a$ 时 $x(a)=1$。
(f) 将 $c,d$ 改成第二方程值，画带时间箭头的相图。同一尺度下，求 $t=0$ 经过 $(0,1)$ 的解；何时再次过 $y$ 轴？粗画 $x(t),y(t)$。
(g) 解释这些伴随矩阵的轨道过 $x$ 轴时切向量为何竖直。\end{exercise}
\begin{solution}
\textit{官方解译审与补证（AI 辅助）；其中另标编者补解或原文修正。}\par

(a) 第一式特征多项式 $(r+1)(r+3)$，可取
\[x_1=e^{-t},\ x_2=e^{-3t},\quad x_1'=-e^{-t},\ x_2'=-3e^{-3t},\quad x=C_1e^{-t}+C_2e^{-3t};\]
是过阻尼。第二式根为 $-1/2\pm3i/2$，可取
\[x_1=e^{-t/2}\cos\frac{3t}2,\quad x_2=e^{-t/2}\sin\frac{3t}2,\quad x=C_1x_1+C_2x_2;\]
\[x_1'=e^{-t/2}\left(-\tfrac12\cos\tfrac{3t}2-\tfrac32\sin\tfrac{3t}2\right),\quad x_2'=e^{-t/2}\left(-\tfrac12\sin\tfrac{3t}2+\tfrac32\cos\tfrac{3t}2\right).\]
两根有非零虚部，是欠阻尼。两对解的初值向量均线性无关，故构成全部实解。
(b) 置 $y=x'$，伴随矩阵分别为
\[A_1=\begin{pmatrix}0&1\\-3&-4\end{pmatrix},\qquad A_2=\begin{pmatrix}0&1\\-5/2&-1\end{pmatrix}.\]
(c) 这里用解析轨道替代实际 Mathlet 操作的陈述。第一系统稳定结点，基本解向量为 $e^{-t}(1,-1)^T$、$e^{-3t}(1,-3)^T$；其轨道在 $y=-x$、$y=-3x$，箭头均指向原点。非特征线轨道也趋于原点，通常以慢衰减的 $y=-x$ 为到达切线。
(d) $x(0)=1,y(0)=0$ 给出 $C_1+C_2=1$、$-C_1-3C_2=0$，所以
\[x=\tfrac12(3e^{-t}-e^{-3t}),\qquad y=\tfrac32(-e^{-t}+e^{-3t}).\]
$x$ 在 $t=-\tfrac12\log3$ 过零，在 $t=0$ 达最大值 $1$，随后单调下降到 $0$；$x''=\tfrac32(e^{-t}-3e^{-3t})$ 在 $t=\tfrac12\log3$ 变号。过去 $t\to-\infty$ 时 $x\to-\infty$，故画出钩形转向。
(e) 自治方程的同轨道解是时间平移：
\[\binom{x(t)}{y(t)}=\frac12\binom{3e^{-(t-a)}-e^{-3(t-a)}}{-3e^{-(t-a)}+3e^{-3(t-a)}},\quad a\in\mathbb R.\]
(f) 第二系统为顺时针稳定螺旋：在正 $x$ 轴上速度向下。由 $x(0)=0,x'(0)=1$ 得
\[x=\frac23e^{-t/2}\sin\frac{3t}2,\qquad y=e^{-t/2}\left(\cos\frac{3t}2-\frac13\sin\frac{3t}2\right).\]
$x=0$ 恰在 $t=2n\pi/3$，$n\in\mathbb Z$；下一次为 $2\pi/3$。二者均为衰减振荡，$x(0)=0,x'(0)=1$，$y(0)=1,y'(0)=-1$。
(g) $x$ 轴上 $y=x'=0$，故 $\mathbf u'=(0,x'')^T=(0,cx)^T$，非原点处因本题 $c\ne0$ 而为非零竖直向量。原点是平衡点，不谈非零切向量。
\begin{center}\begin{tikzpicture}
\begin{axis}[width=6cm,height=5cm,xmin=-1.2,xmax=1.5,ymin=-2,ymax=2,axis lines=middle,xlabel={$x$},ylabel={$y=x'$},title={第一系统：稳定结点},ticks=none]
\addplot[dashed,domain=-1.2:1.2]{-x};\addplot[dashed,domain=-.66:.66]{-3*x};
\foreach \c/\d in {1/.2,1.5/-.5,-1/-.2,-1.5/.5}{\addplot[main,domain=-.25:3,samples=70,->]({\c*exp(-x)+\d*exp(-3*x)},{-\c*exp(-x)-3*\d*exp(-3*x)});}
\node at (axis cs:.8,-1.6){$y=-3x$};\node at (axis cs:1.1,-.8){$y=-x$};
\end{axis}
\begin{axis}[at={(6.2cm,0)},width=6cm,height=5cm,xmin=-1.2,xmax=1.2,ymin=-1.5,ymax=1.5,axis lines=middle,xlabel={$x$},ylabel={$y$},title={第二系统：顺时针螺旋},ticks=none]
\addplot[main,thick,domain=-.6:9,samples=160,->]({2/3*exp(-x/2)*sin(deg(1.5*x))},{exp(-x/2)*(cos(deg(1.5*x))-sin(deg(1.5*x))/3)});\addplot[only marks] coordinates{(0,1)};
\end{axis}\end{tikzpicture}
\begin{tikzpicture}
\begin{axis}[width=6cm,height=4.5cm,xmin=-.8,xmax=4,ymin=-1,ymax=1.2,axis lines=middle,xlabel={$t$},ylabel={$x$},title={(d) 位移曲线}]
\addplot[main,domain=-.8:4,samples=100]{(3*exp(-x)-exp(-3*x))/2};\addplot[only marks]coordinates{(0,1)};
\end{axis}
\begin{axis}[at={(6.2cm,0)},width=6cm,height=4.5cm,xmin=0,xmax=7,ymin=-.7,ymax=1.2,axis lines=middle,xlabel={$t$},title={(f) 时间曲线},legend style={font=\small}]
\addplot[main,domain=0:7,samples=120]{2/3*exp(-x/2)*sin(deg(1.5*x))};\addlegendentry{$x(t)$}
\addplot[structurecolor,dashed,domain=0:7,samples=120]{exp(-x/2)*(cos(deg(1.5*x))-sin(deg(1.5*x))/3)};\addlegendentry{$y(t)$}
\end{axis}\end{tikzpicture}\end{center}
\end{solution}
\begin{exercise}\label{cw:ps8-II-33}原题 II.33：
(a) 求 $x''+4x'+3x=0$ 的伴随矩阵特征值、特征向量，画特征线，各写一个沿线运动的解，与 II.32 特别是 (c) 比较。
(b) 对 $A=\begin{pmatrix}2&-4\\1&-3\end{pmatrix}$ 求特征值、特征向量，画特征线，各写出轨道处于该线的所有解。
(c) 在上述 Mathlet 中设该矩阵，画相平面，至少在每个特征线夹角区域各选一条轨道，并标时间箭头。\end{exercise}
\begin{solution}
\textit{官方解译审与补证（AI 辅助）；其中另标编者补解或原文修正。}\par

(a) $\det(\lambda I-A_1)=\lambda^2+4\lambda+3$，特征值 $-1,-3$。由第一行 $y=\lambda x$，特征向量可取 $(1,-1)^T,(1,-3)^T$；特征线及解正是 II.32(c) 的两条线和 $e^{-t}(1,-1)^T,e^{-3t}(1,-3)^T$。
(b) $\det(\lambda I-A)=\lambda^2+\lambda-2=(\lambda+2)(\lambda-1)$。$\lambda=-2$ 时向量可取 $(1,1)^T$，$\lambda=1$ 时可取 $(4,1)^T$。所有沿线解分别为 $Ce^{-2t}(1,1)^T$、$Ce^t(4,1)^T$，$C\in\mathbb R$。第一线向内，第二线向外。
(c) 全解为 $\mathbf u=C_1e^{-2t}(1,1)^T+C_2e^t(4,1)^T$。取 $(C_1,C_2)$ 的四组符号即得四个夹角的代表轨道；它们过去趋近稳定特征线方向、将来沿不稳定特征线逃离，形成鞍点。下图由此式构造，并未宣称运行 Mathlet。
\begin{center}\begin{tikzpicture}\begin{axis}[width=8cm,height=6cm,xmin=-4,xmax=4,ymin=-3,ymax=3,axis lines=middle,xlabel={$x$},ylabel={$y$},title={鞍点与四个夹角的轨道},ticks=none]
\addplot[dashed,domain=-3:3]{x};\addplot[dashed,domain=-4:4]{x/4};
\foreach \c/\d in {1/1,1/-1,-1/1,-1/-1}{\addplot[main,thick,domain=-.8:1.4,samples=90,->]({\c*exp(-2*x)+.4*\d*exp(x)},{\c*exp(-2*x)+.1*\d*exp(x)});}
\draw[->,thick](axis cs:2,2)--(axis cs:1,1);\draw[->,thick](axis cs:-2,-2)--(axis cs:-1,-1);\draw[->,thick](axis cs:1,.25)--(axis cs:2,.5);\draw[->,thick](axis cs:-1,-.25)--(axis cs:-2,-.5);
\node at (axis cs:1.7,2.5){$\lambda=-2$};\node at (axis cs:2.8,.4){$\lambda=1$};
\end{axis}\end{tikzpicture}\end{center}
\end{solution}
\originalchapter{作业九：相图分类与矩阵指数}
\originalsection{第一部分}
\begin{exercise}\label{cw:ps9-I-34}原题 I.34：完成 Notes 4D-2；EP 第 5.6 节第 1 题，可用「Linear Phase Portraits: Matrix Entry」辅助画相图。\end{exercise}
\begin{solution}
\textit{编者补解及题库回指核对（AI 辅助）；教材题面缺项另标。}\par
Notes 4D-2 的完整题面与解答见习题 \ref{cw:bank-4D-2}。EP 5.6 第 1 题只取得教材定位，未获可核实公开题面，标为 \texttt{external-textbook-statement-unavailable}，不以其他矩阵题代替。\end{solution}
\begin{exercise}\label{cw:ps9-I-35}原题 I.35：完成 Notes 5B-4，可用同一 Mathlet 检查。\end{exercise}
\begin{solution}
\textit{编者补解及题库回指核对（AI 辅助）；教材题面缺项另标。}\par
见习题 \ref{cw:bank-5B-4} 的 (a)--(e)，分别检查特征值及稳定性、实特征线或复根的旋转方向，再画带时间方向的轨道。\end{solution}
\begin{exercise}\label{cw:ps9-I-36}原题 I.36：完成 EP 第 5.7 节第 1、3 题。\end{exercise}
\begin{solution}
\textit{编者补解及题库回指核对（AI 辅助）；教材题面缺项另标。}\par
两项均只有教材题号指定，未获得相应公开题面，均标为 \texttt{external-textbook-statement-unavailable}。\end{solution}
\originalsection{第二部分}
\begin{exercise}\label{cw:ps9-II-34}原题 II.34：
(a) 生物学家研究楠塔基特岛的长尾鼬与草地田鼠。$x,y$ 是相对基线的数量，单位为百只：如 $x(2)=5$ 表示 $t=2$ 时多 500 只鼬，$y(2)=-5$ 表示少 500 只田鼠。模型为
\[x'=0.5x+y,\qquad y'=-2.25x+0.5y.\]
求特征值；取一根求复指数解，再取实、虚部产生两个实解。初始多 100 只鼬、田鼠恰为基线，求初值解；判断初始两种数量增减及田鼠何时再次等于基线。画约 $-2\le t\le2$ 的相图和两时间曲线，标明轨道上 $t=k\pi/3$、$k=-2,-1,0,1,2$ 五点，以说明三图对应。
(b) 使用上述 Mathlet，取消伴随矩阵选项、显示复平面特征值，将 $A=\begin{pmatrix}1&b\\0&1\end{pmatrix}$ 中 $b$ 从 $-4$ 调到 $4$。对每个 $b$ 核实特征值、描述全部特征向量、写全部正规模态 $e^{\lambda t}v$；若缺损则补一个独立解。各选一个可完全对角化与缺损的 $b$ 画相图。\end{exercise}
\begin{solution}
\textit{官方解译审与补证（AI 辅助）；其中另标编者补解或原文修正。}\par

(a) $\det(\lambda I-A)=(\lambda-1/2)^2+9/4$，根为 $\lambda=1/2\pm3i/2$。对正虚部根取 $v=(1,3i/2)^T$，则复解 $e^{(1/2+3i/2)t}v$ 的实、虚部为
\[\mathbf u_1=e^{t/2}\binom{\cos(3t/2)}{-(3/2)\sin(3t/2)},\qquad\mathbf u_2=e^{t/2}\binom{\sin(3t/2)}{(3/2)\cos(3t/2)}.\]
零点初值分别为 $(1,0)^T,(0,3/2)^T$，线性无关。给定初值 $\mathbf u(0)=(1,0)^T$ 恰选 $\mathbf u_1$。初始速度 $A(1,0)^T=(1/2,-9/4)^T$，因此鼬增加、田鼠减少，数量变化率分别为 $50$ 与 $-225$ 只每题面时间单位（原题未另指定时间单位）。$y(t)=0$ 在 $t=2n\pi/3$，故下一次为 $2\pi/3$。
相图为顺时针不稳定螺旋。在 $X=x,Y=2y/3$ 坐标中，$X^2+Y^2=e^t$ 且转角为 $-3t/2$，因此画图时形状、扩张及旋转方向均可直接检查。五个标记点为
\[\begin{array}{c|ccccc}
t&-2\pi/3&-\pi/3&0&\pi/3&2\pi/3\\\hline
(x,y)&(-e^{-\pi/3},0)&(0,\tfrac32e^{-\pi/6})&(1,0)&(0,-\tfrac32e^{\pi/6})&(-e^{\pi/3},0)
\end{array}\]
两端标记略超出 $[-2,2]$，所以图将时间范围延至 $[-2.1,2.1]$ 以保留原要求全部五点。
\begin{center}\begin{tikzpicture}
\begin{axis}[width=4.5cm,height=5cm,xmin=-3,xmax=2,ymin=-3,ymax=1.5,axis lines=middle,xlabel={$x$},ylabel={$y$},title={初值解的相轨道},ticks=none]
\addplot[main,thick,domain=-2.1:2.1,samples=160,->]({exp(x/2)*cos(deg(1.5*x))},{-1.5*exp(x/2)*sin(deg(1.5*x))});
\addplot[only marks,mark=*]coordinates{(-.35092,0)(0,.88858)(1,0)(0,-2.5321)(-2.84965,0)};
\node[below] at(axis cs:-.35092,0){$-2$};\node[above left]at(axis cs:0,.88858){$-1$};\node[above]at(axis cs:1,0){$0$};\node[right]at(axis cs:0,-2.5321){$1$};\node[above]at(axis cs:-2.84965,0){$2$};
\end{axis}
\begin{axis}[at={(4.6cm,0)},width=4.5cm,height=5cm,xmin=-2.1,xmax=2.1,ymin=-3,ymax=1.5,axis lines=middle,xlabel={$t$},title={$x(t)$},xtick={-2.0944,0,2.0944},xticklabels={$-2\pi/3$,$0$,$2\pi/3$},tick label style={font=\scriptsize}]
\addplot[main,thick,domain=-2.1:2.1,samples=150]{exp(x/2)*cos(deg(1.5*x))};
\end{axis}
\begin{axis}[at={(9.2cm,0)},width=4.5cm,height=5cm,xmin=-2.1,xmax=2.1,ymin=-3,ymax=1.5,axis lines=middle,xlabel={$t$},title={$y(t)$},xtick={-2.0944,0,2.0944},xticklabels={$-2\pi/3$,$0$,$2\pi/3$},tick label style={font=\scriptsize}]
\addplot[structurecolor,dashed,domain=-2.1:2.1,samples=150]{-1.5*exp(x/2)*sin(deg(1.5*x))};
\end{axis}\end{tikzpicture}\end{center}
轨道点旁的整数是 $k$，相应横坐标时刻由右图刻度给出；一个时刻的两曲线纵坐标正好组成左图一点。
(b) 这是按矩阵计算所得的核实，不是实际运行 Mathlet 的观察记录。特征多项式恒为 $(\lambda-1)^2$。若 $b=0$，矩阵为 $I$，所有非零向量均为特征向量，全部解都是 $e^tv$，形成径向向外的星形结点。若 $b\ne0$，$\ker(A-I)=\{(c,0)^T\}$，非零特征向量只在 $x$ 轴，正规模态为 $ce^t(1,0)^T$。取广义特征向量 $w=(0,1/b)^T$，满足 $(A-I)w=v=(1,0)^T$，补解为 $e^t(tv+w)$。等价的全部解为
\[\mathbf u(t)=e^t\binom{C_1+bC_2t}{C_2}.\]
其初值矩阵可逆，故已给齐全部解。$b=1$ 是缺损不稳定结点；在上半平面一般轨道向右剪切，角速度满足 $xy'-yx'=-by^2$，故 $b>0$ 为顺时针剪切，$b<0$ 为逆时针剪切。
\begin{center}\begin{tikzpicture}
\begin{axis}[width=5.8cm,height=4.8cm,xmin=-3,xmax=3,ymin=-3,ymax=3,axis lines=middle,xlabel={$x$},ylabel={$y$},title={$b=0$：不稳定星形结点},ticks=none]
\draw[->,main,thick](axis cs:0.30000,0.00000)--(axis cs:2.30000,0.00000);
\draw[->,main,thick](axis cs:0.21213,0.21213)--(axis cs:1.62635,1.62635);
\draw[->,main,thick](axis cs:0.00000,0.30000)--(axis cs:0.00000,2.30000);
\draw[->,main,thick](axis cs:-0.21213,0.21213)--(axis cs:-1.62635,1.62635);
\draw[->,main,thick](axis cs:-0.30000,0.00000)--(axis cs:-2.30000,0.00000);
\draw[->,main,thick](axis cs:-0.21213,-0.21213)--(axis cs:-1.62635,-1.62635);
\draw[->,main,thick](axis cs:-0.00000,-0.30000)--(axis cs:-0.00000,-2.30000);
\draw[->,main,thick](axis cs:0.21213,-0.21213)--(axis cs:1.62635,-1.62635);
\end{axis}
\begin{axis}[at={(6.1cm,0)},width=5.8cm,height=4.8cm,xmin=-3,xmax=3,ymin=-3,ymax=3,axis lines=middle,xlabel={$x$},ylabel={$y$},title={$b=1$：缺损不稳定结点},ticks=none]
\foreach \c/\d in {-1/1,1/1,-1/-1,1/-1}{\addplot[main,thick,domain=-3:.9,samples=90,->]({exp(x)*(\c+\d*x)},{exp(x)*\d});}
\draw[->,thick](axis cs:.2,0)--(axis cs:2.4,0);\draw[->,thick](axis cs:-.2,0)--(axis cs:-2.4,0);\node at(axis cs:1.3,-.45){特征线 $y=0$};
\end{axis}\end{tikzpicture}\end{center}
\end{solution}
\begin{exercise}\label{cw:ps9-II-35}原题 II.35：研究 $\mathbf u'=A\mathbf u$，$A=\begin{pmatrix}a&-3\\1&-1\end{pmatrix}$。在上述 Mathlet 中从 $a=-4$ 调到 $4$，关注迹--行列式平面的标记、复平面特征值及相图变化。
(a) 求迹、行列式及标记轨迹的方程。
(b) 求红边界的四个参数值，即 $\det A=0$、$\det A=(\operatorname{tr}A/2)^2$ 的两个值（一个超出 Mathlet 原滑条范围）及 $\operatorname{tr}A=0$。
(c) 在 $[-5,4]$ 上标四个值，分类五区间与四边界的九种相图；注明稳定性及有关旋转方向，分类可用螺旋、结点、鞍点、中心、星形、缺损结点、退化梳形等。
(d) 四个特殊值及五个区间各选一值画相图，注明所有特征线与时间方向。\end{exercise}
\begin{solution}
\textit{官方解译审与补证（AI 辅助）；其中另标编者补解或原文修正。}\par

(a) $T=\operatorname{tr}A=a-1$、$\Delta=\det A=3-a$，所以参数在 $(T,\Delta)$ 平面沿直线 $T+\Delta=2$ 运动。特征根为
\[\lambda_\pm=\frac{a-1\pm\sqrt{a^2+2a-11}}2.\]
(b) 四个分界为
\[a_-= -1-2\sqrt3\approx-4.4641,\quad1,\quad a_+=-1+2\sqrt3\approx2.4641,\quad3.\]
前后二次曲线交点来自 $4(3-a)=(a-1)^2$。$a_-$ 确在原滑条 $[-4,4]$ 外，须解析补足。
(c) 在复根区间，角速度的分子为
\[xy'-yx'=x^2-(a+1)xy+3y^2.\]
当 $a_-<a<a_+$ 时这是正定二次型，故均为逆时针。边界时退化为 $(x+\sqrt3y)^2$ 或 $(x-\sqrt3y)^2$，非特征线上的剪切亦为逆时针。分类如下；官方解页 2 有将 $a_+$ 的 $-1$ 写成 $+1$ 的笔误，此表按特征判别式核算。
\begin{center}\begin{tabular}{ll}\toprule
参数&类型\\\midrule
$a<a_-$&稳定结点\\
$a=a_-$&稳定缺损结点，逆时针剪切\\
$a_-<a<1$&逆时针稳定螺旋\\
$a=1$&逆时针中心\\
$1<a<a_+$&逆时针不稳定螺旋\\
$a=a_+$&不稳定缺损结点，逆时针剪切\\
$a_+<a<3$&不稳定结点\\
$a=3$&不稳定退化梳形\\
$a>3$&鞍点\\\bottomrule\end{tabular}\end{center}
不存在星形边界：在重根处矩阵不是 $\lambda I$，只有一条特征线。$a=3$ 时根为 $0,2$，$y=x$ 上全为平衡点，其他轨道沿 $(3,1)^T$ 的平行方向离开平衡线。
(d) 下列九图依次取 $a=-5,a_-,0,1,2,a_+,11/4,3,4$。虚线为特征线，实曲线为解析矩阵指数给出的代表轨道，曲线箭头或图内速度箭头标时间增加。任一实根的特征线为 $x=(\lambda+1)y$，这给出图中线的精确定位；中心图的轨道为闭椭圆。为可复算，作图使用
\[A=\frac T2I+B,\quad B^2=\frac{T^2-4\Delta}{4}I.\]
若右端系数为 $\rho^2>0$，则 $e^{At}=e^{Tt/2}[I\cosh\rho t+B\sinh(\rho t)/\rho]$；系数为 $-\beta^2<0$ 时改为 $I\cos\beta t+B\sin(\beta t)/\beta$；零时为 $e^{Tt/2}(I+tB)$。这由幂级数将奇、偶次幂分别求和得到。
\begin{center}
\begin{tikzpicture}
\begin{axis}[at={(0.0cm,0)},width=4.3cm,height=4cm,xmin=-2,xmax=2,ymin=-2,ymax=2,axis lines=middle,xlabel={$x$},ylabel={$y$},title={$a=-5$},ticks=none,clip=true]
\addplot[dashed] coordinates{(-2.2000,0.7333)(2.2000,-0.7333)};
\addplot[dashed] coordinates{(-2.2000,2.2000)(2.2000,-2.2000)};
\addplot[main,thin] coordinates{(2.2375,-0.4288)(1.9091,-0.3412)(1.6265,-0.2674)(1.3835,-0.2053)(1.1747,-0.1534)(0.9954,-0.1100)(0.8416,-0.0740)(0.7096,-0.0443)(0.5967,-0.0199)(0.5000,0.0000)(0.4174,0.0160)(0.3469,0.0289)(0.2868,0.0389)(0.2357,0.0467)(0.1922,0.0525)(0.1554,0.0568)(0.1243,0.0597)(0.0980,0.0615)(0.0759,0.0623)(0.0574,0.0625)(0.0419,0.0620)(0.0290,0.0611)(0.0183,0.0597)(0.0095,0.0581)(0.0024,0.0563)(-0.0035,0.0543)(-0.0081,0.0522)(-0.0118,0.0500)(-0.0147,0.0478)(-0.0168,0.0456)(-0.0184,0.0433)(-0.0196,0.0411)(-0.0203,0.0390)(-0.0207,0.0369)(-0.0208,0.0349)(-0.0207,0.0329)(-0.0205,0.0311)(-0.0201,0.0293)(-0.0196,0.0275)(-0.0190,0.0259)(-0.0184,0.0243)(-0.0177,0.0228)(-0.0169,0.0214)(-0.0162,0.0201)(-0.0155,0.0188)(-0.0147,0.0176)(-0.0140,0.0165)(-0.0133,0.0155)(-0.0126,0.0145)(-0.0119,0.0135)(-0.0112,0.0127)(-0.0106,0.0118)(-0.0100,0.0111)(-0.0094,0.0103)(-0.0088,0.0096)(-0.0083,0.0090)(-0.0078,0.0084)(-0.0073,0.0078)(-0.0069,0.0073)(-0.0064,0.0068)(-0.0060,0.0064)(-0.0056,0.0059)(-0.0053,0.0055)(-0.0049,0.0052)(-0.0046,0.0048)(-0.0043,0.0045)};
\addplot[main,thin] coordinates{(2.3972,0.3791)(1.9647,0.4421)(1.5975,0.4889)(1.2865,0.5221)(1.0237,0.5442)(0.8022,0.5570)(0.6160,0.5622)(0.4601,0.5612)(0.3301,0.5553)(0.2221,0.5455)(0.1329,0.5325)(0.0596,0.5171)(0.0000,0.5000)(-0.0481,0.4816)(-0.0866,0.4623)(-0.1168,0.4425)(-0.1401,0.4224)(-0.1576,0.4024)(-0.1703,0.3825)(-0.1790,0.3629)(-0.1844,0.3438)(-0.1870,0.3252)(-0.1874,0.3072)(-0.1860,0.2899)(-0.1832,0.2733)(-0.1792,0.2573)(-0.1744,0.2421)(-0.1689,0.2276)(-0.1629,0.2138)(-0.1566,0.2007)(-0.1500,0.1882)(-0.1434,0.1765)(-0.1367,0.1654)(-0.1300,0.1549)(-0.1234,0.1450)(-0.1170,0.1357)(-0.1107,0.1270)(-0.1047,0.1187)(-0.0988,0.1110)(-0.0932,0.1037)(-0.0878,0.0969)(-0.0826,0.0905)(-0.0777,0.0845)(-0.0730,0.0789)(-0.0685,0.0737)(-0.0643,0.0688)(-0.0603,0.0642)(-0.0565,0.0599)(-0.0529,0.0559)(-0.0496,0.0521)(-0.0464,0.0486)(-0.0434,0.0453)(-0.0406,0.0422)(-0.0380,0.0394)(-0.0355,0.0367)(-0.0332,0.0342)(-0.0310,0.0319)(-0.0289,0.0297)(-0.0270,0.0277)(-0.0252,0.0258)(-0.0235,0.0241)(-0.0220,0.0224)(-0.0205,0.0209)(-0.0191,0.0195)(-0.0178,0.0181)(-0.0166,0.0169)(-0.0155,0.0157)(-0.0145,0.0147)(-0.0135,0.0137)};
\addplot[main,thin] coordinates{(-2.2375,0.4288)(-1.9091,0.3412)(-1.6265,0.2674)(-1.3835,0.2053)(-1.1747,0.1534)(-0.9954,0.1100)(-0.8416,0.0740)(-0.7096,0.0443)(-0.5967,0.0199)(-0.5000,0.0000)(-0.4174,-0.0160)(-0.3469,-0.0289)(-0.2868,-0.0389)(-0.2357,-0.0467)(-0.1922,-0.0525)(-0.1554,-0.0568)(-0.1243,-0.0597)(-0.0980,-0.0615)(-0.0759,-0.0623)(-0.0574,-0.0625)(-0.0419,-0.0620)(-0.0290,-0.0611)(-0.0183,-0.0597)(-0.0095,-0.0581)(-0.0024,-0.0563)(0.0035,-0.0543)(0.0081,-0.0522)(0.0118,-0.0500)(0.0147,-0.0478)(0.0168,-0.0456)(0.0184,-0.0433)(0.0196,-0.0411)(0.0203,-0.0390)(0.0207,-0.0369)(0.0208,-0.0349)(0.0207,-0.0329)(0.0205,-0.0311)(0.0201,-0.0293)(0.0196,-0.0275)(0.0190,-0.0259)(0.0184,-0.0243)(0.0177,-0.0228)(0.0169,-0.0214)(0.0162,-0.0201)(0.0155,-0.0188)(0.0147,-0.0176)(0.0140,-0.0165)(0.0133,-0.0155)(0.0126,-0.0145)(0.0119,-0.0135)(0.0112,-0.0127)(0.0106,-0.0118)(0.0100,-0.0111)(0.0094,-0.0103)(0.0088,-0.0096)(0.0083,-0.0090)(0.0078,-0.0084)(0.0073,-0.0078)(0.0069,-0.0073)(0.0064,-0.0068)(0.0060,-0.0064)(0.0056,-0.0059)(0.0053,-0.0055)(0.0049,-0.0052)(0.0046,-0.0048)(0.0043,-0.0045)};
\addplot[main,thin] coordinates{(-2.3972,-0.3791)(-1.9647,-0.4421)(-1.5975,-0.4889)(-1.2865,-0.5221)(-1.0237,-0.5442)(-0.8022,-0.5570)(-0.6160,-0.5622)(-0.4601,-0.5612)(-0.3301,-0.5553)(-0.2221,-0.5455)(-0.1329,-0.5325)(-0.0596,-0.5171)(0.0000,-0.5000)(0.0481,-0.4816)(0.0866,-0.4623)(0.1168,-0.4425)(0.1401,-0.4224)(0.1576,-0.4024)(0.1703,-0.3825)(0.1790,-0.3629)(0.1844,-0.3438)(0.1870,-0.3252)(0.1874,-0.3072)(0.1860,-0.2899)(0.1832,-0.2733)(0.1792,-0.2573)(0.1744,-0.2421)(0.1689,-0.2276)(0.1629,-0.2138)(0.1566,-0.2007)(0.1500,-0.1882)(0.1434,-0.1765)(0.1367,-0.1654)(0.1300,-0.1549)(0.1234,-0.1450)(0.1170,-0.1357)(0.1107,-0.1270)(0.1047,-0.1187)(0.0988,-0.1110)(0.0932,-0.1037)(0.0878,-0.0969)(0.0826,-0.0905)(0.0777,-0.0845)(0.0730,-0.0789)(0.0685,-0.0737)(0.0643,-0.0688)(0.0603,-0.0642)(0.0565,-0.0599)(0.0529,-0.0559)(0.0496,-0.0521)(0.0464,-0.0486)(0.0434,-0.0453)(0.0406,-0.0422)(0.0380,-0.0394)(0.0355,-0.0367)(0.0332,-0.0342)(0.0310,-0.0319)(0.0289,-0.0297)(0.0270,-0.0277)(0.0252,-0.0258)(0.0235,-0.0241)(0.0220,-0.0224)(0.0205,-0.0209)(0.0191,-0.0195)(0.0178,-0.0181)(0.0166,-0.0169)(0.0155,-0.0157)(0.0145,-0.0147)(0.0135,-0.0137)};
\draw[->,thick](axis cs:1.0,0.0)--(axis cs:0.6862,0.0628);
\draw[->,thick](axis cs:0.0,1.0)--(axis cs:-0.3036,0.8988);
\draw[->,thick](axis cs:-1.0,0.0)--(axis cs:-0.6862,-0.0628);
\draw[->,thick](axis cs:0.0,-1.0)--(axis cs:0.3036,-0.8988);
\end{axis}
\begin{axis}[at={(4.35cm,0)},width=4.3cm,height=4cm,xmin=-2,xmax=2,ymin=-2,ymax=2,axis lines=middle,xlabel={$x$},ylabel={$y$},title={$a=a_-$},ticks=none,clip=true]
\addplot[dashed] coordinates{(-2.2000,1.2702)(2.2000,-1.2702)};
\addplot[dashed] coordinates{(-2.2000,1.2702)(2.2000,-1.2702)};
\addplot[main,thin] coordinates{(2.4577,-0.5746)(2.1471,-0.4738)(1.8731,-0.3868)(1.6315,-0.3118)(1.4186,-0.2475)(1.2311,-0.1924)(1.0663,-0.1454)(0.9215,-0.1055)(0.7944,-0.0718)(0.6829,-0.0434)(0.5853,-0.0197)(0.5000,0.0000)(0.4255,0.0162)(0.3605,0.0294)(0.3039,0.0400)(0.2547,0.0483)(0.2120,0.0548)(0.1751,0.0597)(0.1432,0.0631)(0.1157,0.0654)(0.0921,0.0668)(0.0719,0.0673)(0.0546,0.0672)(0.0400,0.0664)(0.0275,0.0653)(0.0171,0.0638)(0.0083,0.0620)(0.0011,0.0600)(-0.0049,0.0578)(-0.0098,0.0555)(-0.0137,0.0531)(-0.0168,0.0507)(-0.0193,0.0483)(-0.0211,0.0459)(-0.0224,0.0435)(-0.0233,0.0412)(-0.0238,0.0389)(-0.0241,0.0367)(-0.0240,0.0346)(-0.0238,0.0325)(-0.0234,0.0306)(-0.0229,0.0287)(-0.0223,0.0269)(-0.0216,0.0252)(-0.0208,0.0235)(-0.0200,0.0220)(-0.0192,0.0205)(-0.0183,0.0192)(-0.0174,0.0179)(-0.0166,0.0166)(-0.0157,0.0155)(-0.0149,0.0144)(-0.0141,0.0134)(-0.0133,0.0125)(-0.0125,0.0116)(-0.0118,0.0107)(-0.0111,0.0100)(-0.0104,0.0092)(-0.0097,0.0086)(-0.0091,0.0079)(-0.0085,0.0073)(-0.0080,0.0068)(-0.0074,0.0063)(-0.0069,0.0058)(-0.0065,0.0054)(-0.0060,0.0050)(-0.0056,0.0046)(-0.0052,0.0042)};
\addplot[main,thin] coordinates{(2.4760,0.3481)(2.0731,0.4155)(1.7237,0.4673)(1.4213,0.5060)(1.1603,0.5334)(0.9355,0.5513)(0.7424,0.5613)(0.5772,0.5646)(0.4363,0.5625)(0.3166,0.5559)(0.2154,0.5457)(0.1302,0.5326)(0.0591,0.5171)(0.0000,0.5000)(-0.0486,0.4816)(-0.0881,0.4622)(-0.1199,0.4424)(-0.1450,0.4222)(-0.1644,0.4019)(-0.1790,0.3818)(-0.1894,0.3619)(-0.1963,0.3424)(-0.2004,0.3234)(-0.2019,0.3050)(-0.2015,0.2873)(-0.1993,0.2701)(-0.1959,0.2537)(-0.1913,0.2380)(-0.1859,0.2231)(-0.1799,0.2088)(-0.1734,0.1953)(-0.1665,0.1825)(-0.1594,0.1704)(-0.1522,0.1589)(-0.1450,0.1481)(-0.1378,0.1380)(-0.1306,0.1284)(-0.1236,0.1195)(-0.1168,0.1111)(-0.1102,0.1032)(-0.1038,0.0958)(-0.0976,0.0889)(-0.0917,0.0825)(-0.0861,0.0765)(-0.0807,0.0709)(-0.0755,0.0656)(-0.0706,0.0608)(-0.0660,0.0562)(-0.0616,0.0520)(-0.0575,0.0481)(-0.0536,0.0445)(-0.0499,0.0411)(-0.0465,0.0380)(-0.0432,0.0351)(-0.0402,0.0324)(-0.0374,0.0299)(-0.0347,0.0276)(-0.0322,0.0254)(-0.0299,0.0234)(-0.0277,0.0216)(-0.0257,0.0199)(-0.0238,0.0184)(-0.0220,0.0169)(-0.0204,0.0156)(-0.0189,0.0143)(-0.0174,0.0132)(-0.0161,0.0121)(-0.0149,0.0112)(-0.0138,0.0103)(-0.0127,0.0095)};
\addplot[main,thin] coordinates{(-2.4577,0.5746)(-2.1471,0.4738)(-1.8731,0.3868)(-1.6315,0.3118)(-1.4186,0.2475)(-1.2311,0.1924)(-1.0663,0.1454)(-0.9215,0.1055)(-0.7944,0.0718)(-0.6829,0.0434)(-0.5853,0.0197)(-0.5000,0.0000)(-0.4255,-0.0162)(-0.3605,-0.0294)(-0.3039,-0.0400)(-0.2547,-0.0483)(-0.2120,-0.0548)(-0.1751,-0.0597)(-0.1432,-0.0631)(-0.1157,-0.0654)(-0.0921,-0.0668)(-0.0719,-0.0673)(-0.0546,-0.0672)(-0.0400,-0.0664)(-0.0275,-0.0653)(-0.0171,-0.0638)(-0.0083,-0.0620)(-0.0011,-0.0600)(0.0049,-0.0578)(0.0098,-0.0555)(0.0137,-0.0531)(0.0168,-0.0507)(0.0193,-0.0483)(0.0211,-0.0459)(0.0224,-0.0435)(0.0233,-0.0412)(0.0238,-0.0389)(0.0241,-0.0367)(0.0240,-0.0346)(0.0238,-0.0325)(0.0234,-0.0306)(0.0229,-0.0287)(0.0223,-0.0269)(0.0216,-0.0252)(0.0208,-0.0235)(0.0200,-0.0220)(0.0192,-0.0205)(0.0183,-0.0192)(0.0174,-0.0179)(0.0166,-0.0166)(0.0157,-0.0155)(0.0149,-0.0144)(0.0141,-0.0134)(0.0133,-0.0125)(0.0125,-0.0116)(0.0118,-0.0107)(0.0111,-0.0100)(0.0104,-0.0092)(0.0097,-0.0086)(0.0091,-0.0079)(0.0085,-0.0073)(0.0080,-0.0068)(0.0074,-0.0063)(0.0069,-0.0058)(0.0065,-0.0054)(0.0060,-0.0050)(0.0056,-0.0046)(0.0052,-0.0042)};
\addplot[main,thin] coordinates{(-2.4760,-0.3481)(-2.0731,-0.4155)(-1.7237,-0.4673)(-1.4213,-0.5060)(-1.1603,-0.5334)(-0.9355,-0.5513)(-0.7424,-0.5613)(-0.5772,-0.5646)(-0.4363,-0.5625)(-0.3166,-0.5559)(-0.2154,-0.5457)(-0.1302,-0.5326)(-0.0591,-0.5171)(0.0000,-0.5000)(0.0486,-0.4816)(0.0881,-0.4622)(0.1199,-0.4424)(0.1450,-0.4222)(0.1644,-0.4019)(0.1790,-0.3818)(0.1894,-0.3619)(0.1963,-0.3424)(0.2004,-0.3234)(0.2019,-0.3050)(0.2015,-0.2873)(0.1993,-0.2701)(0.1959,-0.2537)(0.1913,-0.2380)(0.1859,-0.2231)(0.1799,-0.2088)(0.1734,-0.1953)(0.1665,-0.1825)(0.1594,-0.1704)(0.1522,-0.1589)(0.1450,-0.1481)(0.1378,-0.1380)(0.1306,-0.1284)(0.1236,-0.1195)(0.1168,-0.1111)(0.1102,-0.1032)(0.1038,-0.0958)(0.0976,-0.0889)(0.0917,-0.0825)(0.0861,-0.0765)(0.0807,-0.0709)(0.0755,-0.0656)(0.0706,-0.0608)(0.0660,-0.0562)(0.0616,-0.0520)(0.0575,-0.0481)(0.0536,-0.0445)(0.0499,-0.0411)(0.0465,-0.0380)(0.0432,-0.0351)(0.0402,-0.0324)(0.0374,-0.0299)(0.0347,-0.0276)(0.0322,-0.0254)(0.0299,-0.0234)(0.0277,-0.0216)(0.0257,-0.0199)(0.0238,-0.0184)(0.0220,-0.0169)(0.0204,-0.0156)(0.0189,-0.0143)(0.0174,-0.0132)(0.0161,-0.0121)(0.0149,-0.0112)(0.0138,-0.0103)(0.0127,-0.0095)};
\draw[->,thick](axis cs:1.0,0.0)--(axis cs:0.6877,0.0699);
\draw[->,thick](axis cs:0.0,1.0)--(axis cs:-0.3036,0.8988);
\draw[->,thick](axis cs:-1.0,0.0)--(axis cs:-0.6877,-0.0699);
\draw[->,thick](axis cs:0.0,-1.0)--(axis cs:0.3036,-0.8988);
\end{axis}
\begin{axis}[at={(8.7cm,0)},width=4.3cm,height=4cm,xmin=-2,xmax=2,ymin=-2,ymax=2,axis lines=middle,xlabel={$x$},ylabel={$y$},title={$a=0$},ticks=none,clip=true]
\addplot[main,thin] coordinates{(-0.3197,-0.4952)(-0.2670,-0.4881)(-0.2151,-0.4794)(-0.1643,-0.4693)(-0.1146,-0.4577)(-0.0663,-0.4448)(-0.0193,-0.4307)(0.0260,-0.4154)(0.0696,-0.3992)(0.1115,-0.3820)(0.1515,-0.3640)(0.1895,-0.3452)(0.2254,-0.3258)(0.2593,-0.3059)(0.2909,-0.2855)(0.3204,-0.2647)(0.3477,-0.2437)(0.3726,-0.2225)(0.3953,-0.2012)(0.4158,-0.1799)(0.4339,-0.1587)(0.4498,-0.1376)(0.4634,-0.1168)(0.4748,-0.0962)(0.4840,-0.0760)(0.4911,-0.0562)(0.4961,-0.0369)(0.4990,-0.0182)(0.5000,0.0000)(0.4991,0.0175)(0.4963,0.0344)(0.4917,0.0505)(0.4855,0.0659)(0.4776,0.0805)(0.4683,0.0942)(0.4575,0.1072)(0.4453,0.1193)(0.4319,0.1305)(0.4174,0.1408)(0.4018,0.1502)(0.3852,0.1588)(0.3678,0.1664)(0.3496,0.1731)(0.3307,0.1790)(0.3113,0.1840)(0.2913,0.1881)(0.2710,0.1914)(0.2504,0.1938)(0.2295,0.1954)(0.2085,0.1962)(0.1875,0.1963)(0.1665,0.1956)(0.1456,0.1942)(0.1249,0.1922)(0.1044,0.1894)(0.0843,0.1861)(0.0646,0.1822)(0.0453,0.1777)(0.0265,0.1727)(0.0083,0.1673)(-0.0093,0.1614)(-0.0263,0.1551)(-0.0425,0.1484)(-0.0581,0.1415)(-0.0728,0.1342)(-0.0868,0.1267)(-0.1000,0.1190)(-0.1123,0.1111)(-0.1238,0.1030)(-0.1344,0.0949)(-0.1441,0.0867)(-0.1529,0.0784)(-0.1609,0.0701)(-0.1680,0.0619)(-0.1742,0.0537)(-0.1795,0.0456)(-0.1840,0.0377)(-0.1876,0.0298)(-0.1904,0.0221)(-0.1923,0.0147)(-0.1935,0.0074)(-0.1939,0.0003)(-0.1936,-0.0065)(-0.1925,-0.0130)(-0.1908,-0.0193)};
\addplot[main,thin] coordinates{(1.4856,0.1756)(1.4643,0.2212)(1.4383,0.2643)(1.4078,0.3050)(1.3731,0.3431)(1.3344,0.3785)(1.2920,0.4113)(1.2463,0.4414)(1.1975,0.4688)(1.1460,0.4935)(1.0919,0.5154)(1.0356,0.5347)(0.9774,0.5512)(0.9176,0.5651)(0.8564,0.5764)(0.7942,0.5851)(0.7311,0.5914)(0.6675,0.5951)(0.6037,0.5966)(0.5398,0.5957)(0.4761,0.5926)(0.4129,0.5874)(0.3503,0.5802)(0.2886,0.5710)(0.2280,0.5600)(0.1687,0.5473)(0.1108,0.5330)(0.0545,0.5172)(0.0000,0.5000)(-0.0526,0.4815)(-0.1031,0.4619)(-0.1515,0.4412)(-0.1977,0.4196)(-0.2414,0.3972)(-0.2827,0.3740)(-0.3215,0.3503)(-0.3578,0.3261)(-0.3914,0.3015)(-0.4224,0.2766)(-0.4507,0.2516)(-0.4763,0.2265)(-0.4992,0.2014)(-0.5194,0.1765)(-0.5370,0.1517)(-0.5520,0.1273)(-0.5643,0.1032)(-0.5741,0.0796)(-0.5814,0.0566)(-0.5862,0.0341)(-0.5887,0.0123)(-0.5889,-0.0088)(-0.5869,-0.0291)(-0.5827,-0.0486)(-0.5765,-0.0673)(-0.5683,-0.0850)(-0.5583,-0.1018)(-0.5465,-0.1176)(-0.5331,-0.1324)(-0.5182,-0.1462)(-0.5018,-0.1590)(-0.4842,-0.1707)(-0.4653,-0.1814)(-0.4453,-0.1910)(-0.4244,-0.1995)(-0.4026,-0.2070)(-0.3801,-0.2135)(-0.3569,-0.2189)(-0.3332,-0.2234)(-0.3091,-0.2268)(-0.2846,-0.2293)(-0.2600,-0.2308)(-0.2352,-0.2313)(-0.2104,-0.2310)(-0.1857,-0.2299)(-0.1612,-0.2279)(-0.1369,-0.2251)(-0.1130,-0.2216)(-0.0895,-0.2174)(-0.0664,-0.2125)(-0.0440,-0.2070)(-0.0221,-0.2009)(-0.0009,-0.1942)(0.0195,-0.1871)(0.0391,-0.1795)(0.0579,-0.1715)};
\addplot[main,thin] coordinates{(0.3197,0.4952)(0.2670,0.4881)(0.2151,0.4794)(0.1643,0.4693)(0.1146,0.4577)(0.0663,0.4448)(0.0193,0.4307)(-0.0260,0.4154)(-0.0696,0.3992)(-0.1115,0.3820)(-0.1515,0.3640)(-0.1895,0.3452)(-0.2254,0.3258)(-0.2593,0.3059)(-0.2909,0.2855)(-0.3204,0.2647)(-0.3477,0.2437)(-0.3726,0.2225)(-0.3953,0.2012)(-0.4158,0.1799)(-0.4339,0.1587)(-0.4498,0.1376)(-0.4634,0.1168)(-0.4748,0.0962)(-0.4840,0.0760)(-0.4911,0.0562)(-0.4961,0.0369)(-0.4990,0.0182)(-0.5000,0.0000)(-0.4991,-0.0175)(-0.4963,-0.0344)(-0.4917,-0.0505)(-0.4855,-0.0659)(-0.4776,-0.0805)(-0.4683,-0.0942)(-0.4575,-0.1072)(-0.4453,-0.1193)(-0.4319,-0.1305)(-0.4174,-0.1408)(-0.4018,-0.1502)(-0.3852,-0.1588)(-0.3678,-0.1664)(-0.3496,-0.1731)(-0.3307,-0.1790)(-0.3113,-0.1840)(-0.2913,-0.1881)(-0.2710,-0.1914)(-0.2504,-0.1938)(-0.2295,-0.1954)(-0.2085,-0.1962)(-0.1875,-0.1963)(-0.1665,-0.1956)(-0.1456,-0.1942)(-0.1249,-0.1922)(-0.1044,-0.1894)(-0.0843,-0.1861)(-0.0646,-0.1822)(-0.0453,-0.1777)(-0.0265,-0.1727)(-0.0083,-0.1673)(0.0093,-0.1614)(0.0263,-0.1551)(0.0425,-0.1484)(0.0581,-0.1415)(0.0728,-0.1342)(0.0868,-0.1267)(0.1000,-0.1190)(0.1123,-0.1111)(0.1238,-0.1030)(0.1344,-0.0949)(0.1441,-0.0867)(0.1529,-0.0784)(0.1609,-0.0701)(0.1680,-0.0619)(0.1742,-0.0537)(0.1795,-0.0456)(0.1840,-0.0377)(0.1876,-0.0298)(0.1904,-0.0221)(0.1923,-0.0147)(0.1935,-0.0074)(0.1939,-0.0003)(0.1936,0.0065)(0.1925,0.0130)(0.1908,0.0193)};
\addplot[main,thin] coordinates{(-1.4856,-0.1756)(-1.4643,-0.2212)(-1.4383,-0.2643)(-1.4078,-0.3050)(-1.3731,-0.3431)(-1.3344,-0.3785)(-1.2920,-0.4113)(-1.2463,-0.4414)(-1.1975,-0.4688)(-1.1460,-0.4935)(-1.0919,-0.5154)(-1.0356,-0.5347)(-0.9774,-0.5512)(-0.9176,-0.5651)(-0.8564,-0.5764)(-0.7942,-0.5851)(-0.7311,-0.5914)(-0.6675,-0.5951)(-0.6037,-0.5966)(-0.5398,-0.5957)(-0.4761,-0.5926)(-0.4129,-0.5874)(-0.3503,-0.5802)(-0.2886,-0.5710)(-0.2280,-0.5600)(-0.1687,-0.5473)(-0.1108,-0.5330)(-0.0545,-0.5172)(0.0000,-0.5000)(0.0526,-0.4815)(0.1031,-0.4619)(0.1515,-0.4412)(0.1977,-0.4196)(0.2414,-0.3972)(0.2827,-0.3740)(0.3215,-0.3503)(0.3578,-0.3261)(0.3914,-0.3015)(0.4224,-0.2766)(0.4507,-0.2516)(0.4763,-0.2265)(0.4992,-0.2014)(0.5194,-0.1765)(0.5370,-0.1517)(0.5520,-0.1273)(0.5643,-0.1032)(0.5741,-0.0796)(0.5814,-0.0566)(0.5862,-0.0341)(0.5887,-0.0123)(0.5889,0.0088)(0.5869,0.0291)(0.5827,0.0486)(0.5765,0.0673)(0.5683,0.0850)(0.5583,0.1018)(0.5465,0.1176)(0.5331,0.1324)(0.5182,0.1462)(0.5018,0.1590)(0.4842,0.1707)(0.4653,0.1814)(0.4453,0.1910)(0.4244,0.1995)(0.4026,0.2070)(0.3801,0.2135)(0.3569,0.2189)(0.3332,0.2234)(0.3091,0.2268)(0.2846,0.2293)(0.2600,0.2308)(0.2352,0.2313)(0.2104,0.2310)(0.1857,0.2299)(0.1612,0.2279)(0.1369,0.2251)(0.1130,0.2216)(0.0895,0.2174)(0.0664,0.2125)(0.0440,0.2070)(0.0221,0.2009)(0.0009,0.1942)(-0.0195,0.1871)(-0.0391,0.1795)(-0.0579,0.1715)};
\draw[->,thick](axis cs:1.0,0.0)--(axis cs:1.0000,0.3200);
\draw[->,thick](axis cs:0.0,1.0)--(axis cs:-0.3036,0.8988);
\draw[->,thick](axis cs:-1.0,0.0)--(axis cs:-1.0000,-0.3200);
\draw[->,thick](axis cs:0.0,-1.0)--(axis cs:0.3036,-0.8988);
\end{axis}
\end{tikzpicture}
\par
\begin{tikzpicture}
\begin{axis}[at={(0.0cm,0)},width=4.3cm,height=4cm,xmin=-2,xmax=2,ymin=-2,ymax=2,axis lines=middle,xlabel={$x$},ylabel={$y$},title={$a=1$},ticks=none,clip=true]
\addplot[main,thin] coordinates{(0.5000,0.0000)(0.5214,0.0224)(0.5407,0.0448)(0.5579,0.0669)(0.5728,0.0888)(0.5854,0.1103)(0.5956,0.1314)(0.6034,0.1520)(0.6088,0.1719)(0.6118,0.1911)(0.6123,0.2096)(0.6103,0.2273)(0.6058,0.2440)(0.5990,0.2597)(0.5897,0.2744)(0.5780,0.2880)(0.5640,0.3004)(0.5478,0.3116)(0.5293,0.3216)(0.5087,0.3303)(0.4861,0.3376)(0.4615,0.3436)(0.4350,0.3482)(0.4068,0.3514)(0.3769,0.3532)(0.3456,0.3535)(0.3128,0.3524)(0.2788,0.3500)(0.2437,0.3461)(0.2075,0.3408)(0.1706,0.3341)(0.1329,0.3261)(0.0947,0.3168)(0.0562,0.3062)(0.0174,0.2944)(-0.0215,0.2813)(-0.0602,0.2672)(-0.0988,0.2520)(-0.1369,0.2357)(-0.1745,0.2186)(-0.2114,0.2005)(-0.2474,0.1816)(-0.2824,0.1620)(-0.3163,0.1418)(-0.3489,0.1209)(-0.3801,0.0996)(-0.4098,0.0779)(-0.4379,0.0559)(-0.4641,0.0336)(-0.4885,0.0112)(-0.5110,-0.0112)(-0.5313,-0.0336)(-0.5496,-0.0559)(-0.5656,-0.0779)(-0.5794,-0.0996)(-0.5908,-0.1209)(-0.5998,-0.1418)(-0.6064,-0.1620)(-0.6106,-0.1816)(-0.6123,-0.2005)(-0.6116,-0.2186)(-0.6084,-0.2357)(-0.6027,-0.2520)(-0.5946,-0.2672)(-0.5841,-0.2813)(-0.5713,-0.2944)(-0.5562,-0.3062)(-0.5388,-0.3168)(-0.5193,-0.3261)(-0.4976,-0.3341)(-0.4740,-0.3408)(-0.4485,-0.3461)(-0.4211,-0.3500)(-0.3921,-0.3524)(-0.3614,-0.3535)(-0.3294,-0.3532)(-0.2960,-0.3514)(-0.2614,-0.3482)(-0.2257,-0.3436)(-0.1891,-0.3376)(-0.1518,-0.3303)(-0.1139,-0.3216)(-0.0755,-0.3116)(-0.0368,-0.3004)(0.0020,-0.2880)(0.0409,-0.2744)(0.0795,-0.2597)(0.1179,-0.2440)(0.1558,-0.2273)(0.1930,-0.2096)(0.2295,-0.1911)(0.2650,-0.1719)(0.2995,-0.1520)(0.3328,-0.1314)(0.3647,-0.1103)(0.3952,-0.0888)(0.4241,-0.0669)(0.4512,-0.0448)(0.4766,-0.0224)(0.5000,-0.0000)};
\addplot[main,thin] coordinates{(0.0000,0.5000)(-0.0673,0.4766)(-0.1343,0.4512)(-0.2007,0.4241)(-0.2664,0.3952)(-0.3310,0.3647)(-0.3942,0.3328)(-0.4559,0.2995)(-0.5157,0.2650)(-0.5734,0.2295)(-0.6289,0.1930)(-0.6818,0.1558)(-0.7319,0.1179)(-0.7792,0.0795)(-0.8232,0.0409)(-0.8640,0.0020)(-0.9013,-0.0368)(-0.9349,-0.0755)(-0.9648,-0.1139)(-0.9908,-0.1518)(-1.0128,-0.1891)(-1.0308,-0.2257)(-1.0445,-0.2614)(-1.0541,-0.2960)(-1.0595,-0.3294)(-1.0605,-0.3614)(-1.0573,-0.3921)(-1.0499,-0.4211)(-1.0382,-0.4485)(-1.0223,-0.4740)(-1.0023,-0.4976)(-0.9783,-0.5193)(-0.9503,-0.5388)(-0.9186,-0.5562)(-0.8831,-0.5713)(-0.8440,-0.5841)(-0.8016,-0.5946)(-0.7559,-0.6027)(-0.7072,-0.6084)(-0.6557,-0.6116)(-0.6015,-0.6123)(-0.5448,-0.6106)(-0.4860,-0.6064)(-0.4253,-0.5998)(-0.3628,-0.5908)(-0.2988,-0.5794)(-0.2337,-0.5656)(-0.1676,-0.5496)(-0.1008,-0.5313)(-0.0337,-0.5110)(0.0337,-0.4885)(0.1008,-0.4641)(0.1676,-0.4379)(0.2337,-0.4098)(0.2988,-0.3801)(0.3628,-0.3489)(0.4253,-0.3163)(0.4860,-0.2824)(0.5448,-0.2474)(0.6015,-0.2114)(0.6557,-0.1745)(0.7072,-0.1369)(0.7559,-0.0988)(0.8016,-0.0602)(0.8440,-0.0215)(0.8831,0.0174)(0.9186,0.0562)(0.9503,0.0947)(0.9783,0.1329)(1.0023,0.1706)(1.0223,0.2075)(1.0382,0.2437)(1.0499,0.2788)(1.0573,0.3128)(1.0605,0.3456)(1.0595,0.3769)(1.0541,0.4068)(1.0445,0.4350)(1.0308,0.4615)(1.0128,0.4861)(0.9908,0.5087)(0.9648,0.5293)(0.9349,0.5478)(0.9013,0.5640)(0.8640,0.5780)(0.8232,0.5897)(0.7792,0.5990)(0.7319,0.6058)(0.6818,0.6103)(0.6289,0.6123)(0.5734,0.6118)(0.5157,0.6088)(0.4559,0.6034)(0.3942,0.5956)(0.3310,0.5854)(0.2664,0.5728)(0.2007,0.5579)(0.1343,0.5407)(0.0673,0.5214)(-0.0000,0.5000)};
\addplot[main,thin] coordinates{(-0.5000,0.0000)(-0.5214,-0.0224)(-0.5407,-0.0448)(-0.5579,-0.0669)(-0.5728,-0.0888)(-0.5854,-0.1103)(-0.5956,-0.1314)(-0.6034,-0.1520)(-0.6088,-0.1719)(-0.6118,-0.1911)(-0.6123,-0.2096)(-0.6103,-0.2273)(-0.6058,-0.2440)(-0.5990,-0.2597)(-0.5897,-0.2744)(-0.5780,-0.2880)(-0.5640,-0.3004)(-0.5478,-0.3116)(-0.5293,-0.3216)(-0.5087,-0.3303)(-0.4861,-0.3376)(-0.4615,-0.3436)(-0.4350,-0.3482)(-0.4068,-0.3514)(-0.3769,-0.3532)(-0.3456,-0.3535)(-0.3128,-0.3524)(-0.2788,-0.3500)(-0.2437,-0.3461)(-0.2075,-0.3408)(-0.1706,-0.3341)(-0.1329,-0.3261)(-0.0947,-0.3168)(-0.0562,-0.3062)(-0.0174,-0.2944)(0.0215,-0.2813)(0.0602,-0.2672)(0.0988,-0.2520)(0.1369,-0.2357)(0.1745,-0.2186)(0.2114,-0.2005)(0.2474,-0.1816)(0.2824,-0.1620)(0.3163,-0.1418)(0.3489,-0.1209)(0.3801,-0.0996)(0.4098,-0.0779)(0.4379,-0.0559)(0.4641,-0.0336)(0.4885,-0.0112)(0.5110,0.0112)(0.5313,0.0336)(0.5496,0.0559)(0.5656,0.0779)(0.5794,0.0996)(0.5908,0.1209)(0.5998,0.1418)(0.6064,0.1620)(0.6106,0.1816)(0.6123,0.2005)(0.6116,0.2186)(0.6084,0.2357)(0.6027,0.2520)(0.5946,0.2672)(0.5841,0.2813)(0.5713,0.2944)(0.5562,0.3062)(0.5388,0.3168)(0.5193,0.3261)(0.4976,0.3341)(0.4740,0.3408)(0.4485,0.3461)(0.4211,0.3500)(0.3921,0.3524)(0.3614,0.3535)(0.3294,0.3532)(0.2960,0.3514)(0.2614,0.3482)(0.2257,0.3436)(0.1891,0.3376)(0.1518,0.3303)(0.1139,0.3216)(0.0755,0.3116)(0.0368,0.3004)(-0.0020,0.2880)(-0.0409,0.2744)(-0.0795,0.2597)(-0.1179,0.2440)(-0.1558,0.2273)(-0.1930,0.2096)(-0.2295,0.1911)(-0.2650,0.1719)(-0.2995,0.1520)(-0.3328,0.1314)(-0.3647,0.1103)(-0.3952,0.0888)(-0.4241,0.0669)(-0.4512,0.0448)(-0.4766,0.0224)(-0.5000,0.0000)};
\addplot[main,thin] coordinates{(0.0000,-0.5000)(0.0673,-0.4766)(0.1343,-0.4512)(0.2007,-0.4241)(0.2664,-0.3952)(0.3310,-0.3647)(0.3942,-0.3328)(0.4559,-0.2995)(0.5157,-0.2650)(0.5734,-0.2295)(0.6289,-0.1930)(0.6818,-0.1558)(0.7319,-0.1179)(0.7792,-0.0795)(0.8232,-0.0409)(0.8640,-0.0020)(0.9013,0.0368)(0.9349,0.0755)(0.9648,0.1139)(0.9908,0.1518)(1.0128,0.1891)(1.0308,0.2257)(1.0445,0.2614)(1.0541,0.2960)(1.0595,0.3294)(1.0605,0.3614)(1.0573,0.3921)(1.0499,0.4211)(1.0382,0.4485)(1.0223,0.4740)(1.0023,0.4976)(0.9783,0.5193)(0.9503,0.5388)(0.9186,0.5562)(0.8831,0.5713)(0.8440,0.5841)(0.8016,0.5946)(0.7559,0.6027)(0.7072,0.6084)(0.6557,0.6116)(0.6015,0.6123)(0.5448,0.6106)(0.4860,0.6064)(0.4253,0.5998)(0.3628,0.5908)(0.2988,0.5794)(0.2337,0.5656)(0.1676,0.5496)(0.1008,0.5313)(0.0337,0.5110)(-0.0337,0.4885)(-0.1008,0.4641)(-0.1676,0.4379)(-0.2337,0.4098)(-0.2988,0.3801)(-0.3628,0.3489)(-0.4253,0.3163)(-0.4860,0.2824)(-0.5448,0.2474)(-0.6015,0.2114)(-0.6557,0.1745)(-0.7072,0.1369)(-0.7559,0.0988)(-0.8016,0.0602)(-0.8440,0.0215)(-0.8831,-0.0174)(-0.9186,-0.0562)(-0.9503,-0.0947)(-0.9783,-0.1329)(-1.0023,-0.1706)(-1.0223,-0.2075)(-1.0382,-0.2437)(-1.0499,-0.2788)(-1.0573,-0.3128)(-1.0605,-0.3456)(-1.0595,-0.3769)(-1.0541,-0.4068)(-1.0445,-0.4350)(-1.0308,-0.4615)(-1.0128,-0.4861)(-0.9908,-0.5087)(-0.9648,-0.5293)(-0.9349,-0.5478)(-0.9013,-0.5640)(-0.8640,-0.5780)(-0.8232,-0.5897)(-0.7792,-0.5990)(-0.7319,-0.6058)(-0.6818,-0.6103)(-0.6289,-0.6123)(-0.5734,-0.6118)(-0.5157,-0.6088)(-0.4559,-0.6034)(-0.3942,-0.5956)(-0.3310,-0.5854)(-0.2664,-0.5728)(-0.2007,-0.5579)(-0.1343,-0.5407)(-0.0673,-0.5214)(0.0000,-0.5000)};
\draw[->,thick](axis cs:1.0,0.0)--(axis cs:1.2263,0.2263);
\draw[->,thick](axis cs:0.0,1.0)--(axis cs:-0.3036,0.8988);
\draw[->,thick](axis cs:-1.0,0.0)--(axis cs:-1.2263,-0.2263);
\draw[->,thick](axis cs:0.0,-1.0)--(axis cs:0.3036,-0.8988);
\end{axis}
\begin{axis}[at={(4.35cm,0)},width=4.3cm,height=4cm,xmin=-2,xmax=2,ymin=-2,ymax=2,axis lines=middle,xlabel={$x$},ylabel={$y$},title={$a=2$},ticks=none,clip=true]
\addplot[main,thin] coordinates{(-0.2037,-0.2668)(-0.1892,-0.2643)(-0.1741,-0.2614)(-0.1581,-0.2580)(-0.1413,-0.2542)(-0.1238,-0.2500)(-0.1055,-0.2452)(-0.0863,-0.2400)(-0.0664,-0.2342)(-0.0456,-0.2280)(-0.0240,-0.2212)(-0.0016,-0.2139)(0.0216,-0.2061)(0.0456,-0.1976)(0.0705,-0.1887)(0.0961,-0.1791)(0.1226,-0.1690)(0.1498,-0.1583)(0.1779,-0.1470)(0.2067,-0.1351)(0.2364,-0.1226)(0.2668,-0.1095)(0.2979,-0.0957)(0.3298,-0.0813)(0.3625,-0.0663)(0.3958,-0.0507)(0.4299,-0.0344)(0.4646,-0.0175)(0.5000,0.0000)(0.5360,0.0182)(0.5727,0.0370)(0.6099,0.0564)(0.6477,0.0765)(0.6860,0.0972)(0.7249,0.1186)(0.7642,0.1405)(0.8039,0.1631)(0.8440,0.1863)(0.8845,0.2101)(0.9254,0.2345)(0.9665,0.2594)(1.0078,0.2850)(1.0494,0.3111)(1.0911,0.3377)(1.1328,0.3649)(1.1747,0.3926)(1.2165,0.4207)(1.2583,0.4494)(1.3000,0.4785)(1.3414,0.5081)(1.3827,0.5380)(1.4237,0.5684)(1.4643,0.5991)(1.5044,0.6302)(1.5441,0.6615)(1.5832,0.6932)(1.6217,0.7251)(1.6595,0.7572)(1.6965,0.7895)(1.7327,0.8220)(1.7679,0.8546)(1.8021,0.8872)(1.8352,0.9199)(1.8671,0.9526)(1.8977,0.9852)(1.9270,1.0178)(1.9549,1.0502)(1.9812,1.0824)(2.0060,1.1143)(2.0290,1.1460)(2.0502,1.1774)(2.0696,1.2084)(2.0869,1.2389)(2.1022,1.2689)(2.1153,1.2984)(2.1261,1.3272)(2.1346,1.3554)(2.1406,1.3829)(2.1440,1.4095)(2.1448,1.4353)(2.1428,1.4602)(2.1379,1.4841)(2.1302,1.5069)(2.1193,1.5286)(2.1053,1.5490)};
\addplot[main,thin] coordinates{(0.8003,0.5966)(0.7929,0.6036)(0.7842,0.6101)(0.7741,0.6160)(0.7627,0.6214)(0.7499,0.6261)(0.7357,0.6302)(0.7200,0.6337)(0.7027,0.6364)(0.6840,0.6384)(0.6637,0.6396)(0.6417,0.6401)(0.6182,0.6397)(0.5929,0.6385)(0.5660,0.6365)(0.5374,0.6335)(0.5070,0.6296)(0.4749,0.6247)(0.4410,0.6189)(0.4053,0.6120)(0.3677,0.6041)(0.3284,0.5951)(0.2871,0.5850)(0.2440,0.5738)(0.1990,0.5615)(0.1521,0.5479)(0.1033,0.5332)(0.0526,0.5172)(0.0000,0.5000)(-0.0545,0.4815)(-0.1110,0.4617)(-0.1693,0.4406)(-0.2296,0.4181)(-0.2917,0.3943)(-0.3557,0.3691)(-0.4216,0.3425)(-0.4894,0.3145)(-0.5589,0.2851)(-0.6303,0.2542)(-0.7034,0.2219)(-0.7783,0.1881)(-0.8549,0.1529)(-0.9332,0.1162)(-1.0131,0.0779)(-1.0946,0.0382)(-1.1777,-0.0030)(-1.2622,-0.0457)(-1.3482,-0.0899)(-1.4355,-0.1356)(-1.5242,-0.1827)(-1.6141,-0.2314)(-1.7052,-0.2815)(-1.7973,-0.3331)(-1.8905,-0.3861)(-1.9846,-0.4405)(-2.0796,-0.4964)(-2.1753,-0.5536)(-2.2717,-0.6122)(-2.3686,-0.6721)(-2.4660,-0.7334)};
\addplot[main,thin] coordinates{(0.2037,0.2668)(0.1892,0.2643)(0.1741,0.2614)(0.1581,0.2580)(0.1413,0.2542)(0.1238,0.2500)(0.1055,0.2452)(0.0863,0.2400)(0.0664,0.2342)(0.0456,0.2280)(0.0240,0.2212)(0.0016,0.2139)(-0.0216,0.2061)(-0.0456,0.1976)(-0.0705,0.1887)(-0.0961,0.1791)(-0.1226,0.1690)(-0.1498,0.1583)(-0.1779,0.1470)(-0.2067,0.1351)(-0.2364,0.1226)(-0.2668,0.1095)(-0.2979,0.0957)(-0.3298,0.0813)(-0.3625,0.0663)(-0.3958,0.0507)(-0.4299,0.0344)(-0.4646,0.0175)(-0.5000,0.0000)(-0.5360,-0.0182)(-0.5727,-0.0370)(-0.6099,-0.0564)(-0.6477,-0.0765)(-0.6860,-0.0972)(-0.7249,-0.1186)(-0.7642,-0.1405)(-0.8039,-0.1631)(-0.8440,-0.1863)(-0.8845,-0.2101)(-0.9254,-0.2345)(-0.9665,-0.2594)(-1.0078,-0.2850)(-1.0494,-0.3111)(-1.0911,-0.3377)(-1.1328,-0.3649)(-1.1747,-0.3926)(-1.2165,-0.4207)(-1.2583,-0.4494)(-1.3000,-0.4785)(-1.3414,-0.5081)(-1.3827,-0.5380)(-1.4237,-0.5684)(-1.4643,-0.5991)(-1.5044,-0.6302)(-1.5441,-0.6615)(-1.5832,-0.6932)(-1.6217,-0.7251)(-1.6595,-0.7572)(-1.6965,-0.7895)(-1.7327,-0.8220)(-1.7679,-0.8546)(-1.8021,-0.8872)(-1.8352,-0.9199)(-1.8671,-0.9526)(-1.8977,-0.9852)(-1.9270,-1.0178)(-1.9549,-1.0502)(-1.9812,-1.0824)(-2.0060,-1.1143)(-2.0290,-1.1460)(-2.0502,-1.1774)(-2.0696,-1.2084)(-2.0869,-1.2389)(-2.1022,-1.2689)(-2.1153,-1.2984)(-2.1261,-1.3272)(-2.1346,-1.3554)(-2.1406,-1.3829)(-2.1440,-1.4095)(-2.1448,-1.4353)(-2.1428,-1.4602)(-2.1379,-1.4841)(-2.1302,-1.5069)(-2.1193,-1.5286)(-2.1053,-1.5490)};
\addplot[main,thin] coordinates{(-0.8003,-0.5966)(-0.7929,-0.6036)(-0.7842,-0.6101)(-0.7741,-0.6160)(-0.7627,-0.6214)(-0.7499,-0.6261)(-0.7357,-0.6302)(-0.7200,-0.6337)(-0.7027,-0.6364)(-0.6840,-0.6384)(-0.6637,-0.6396)(-0.6417,-0.6401)(-0.6182,-0.6397)(-0.5929,-0.6385)(-0.5660,-0.6365)(-0.5374,-0.6335)(-0.5070,-0.6296)(-0.4749,-0.6247)(-0.4410,-0.6189)(-0.4053,-0.6120)(-0.3677,-0.6041)(-0.3284,-0.5951)(-0.2871,-0.5850)(-0.2440,-0.5738)(-0.1990,-0.5615)(-0.1521,-0.5479)(-0.1033,-0.5332)(-0.0526,-0.5172)(0.0000,-0.5000)(0.0545,-0.4815)(0.1110,-0.4617)(0.1693,-0.4406)(0.2296,-0.4181)(0.2917,-0.3943)(0.3557,-0.3691)(0.4216,-0.3425)(0.4894,-0.3145)(0.5589,-0.2851)(0.6303,-0.2542)(0.7034,-0.2219)(0.7783,-0.1881)(0.8549,-0.1529)(0.9332,-0.1162)(1.0131,-0.0779)(1.0946,-0.0382)(1.1777,0.0030)(1.2622,0.0457)(1.3482,0.0899)(1.4355,0.1356)(1.5242,0.1827)(1.6141,0.2314)(1.7052,0.2815)(1.7973,0.3331)(1.8905,0.3861)(1.9846,0.4405)(2.0796,0.4964)(2.1753,0.5536)(2.2717,0.6122)(2.3686,0.6721)(2.4660,0.7334)};
\draw[->,thick](axis cs:1.0,0.0)--(axis cs:1.2862,0.1431);
\draw[->,thick](axis cs:0.0,1.0)--(axis cs:-0.3036,0.8988);
\draw[->,thick](axis cs:-1.0,0.0)--(axis cs:-1.2862,-0.1431);
\draw[->,thick](axis cs:0.0,-1.0)--(axis cs:0.3036,-0.8988);
\end{axis}
\begin{axis}[at={(8.7cm,0)},width=4.3cm,height=4cm,xmin=-2,xmax=2,ymin=-2,ymax=2,axis lines=middle,xlabel={$x$},ylabel={$y$},title={$a=a_+$},ticks=none,clip=true]
\addplot[dashed] coordinates{(2.2000,1.2702)(-2.2000,-1.2702)};
\addplot[dashed] coordinates{(2.2000,1.2702)(-2.2000,-1.2702)};
\addplot[main,thin] coordinates{(-0.1760,-0.2405)(-0.1654,-0.2380)(-0.1541,-0.2353)(-0.1421,-0.2322)(-0.1294,-0.2288)(-0.1159,-0.2251)(-0.1015,-0.2210)(-0.0863,-0.2166)(-0.0703,-0.2117)(-0.0533,-0.2065)(-0.0354,-0.2008)(-0.0165,-0.1946)(0.0034,-0.1880)(0.0244,-0.1810)(0.0465,-0.1734)(0.0697,-0.1653)(0.0941,-0.1566)(0.1198,-0.1473)(0.1468,-0.1375)(0.1752,-0.1270)(0.2049,-0.1159)(0.2361,-0.1041)(0.2688,-0.0916)(0.3030,-0.0783)(0.3389,-0.0643)(0.3765,-0.0495)(0.4158,-0.0339)(0.4570,-0.0174)(0.5000,0.0000)(0.5450,0.0183)(0.5920,0.0376)(0.6412,0.0579)(0.6925,0.0793)(0.7461,0.1018)(0.8020,0.1253)(0.8604,0.1501)(0.9213,0.1761)(0.9849,0.2034)(1.0511,0.2319)(1.1202,0.2619)(1.1922,0.2933)(1.2672,0.3261)(1.3454,0.3605)(1.4268,0.3965)(1.5116,0.4341)(1.5999,0.4735)(1.6918,0.5146)(1.7874,0.5576)(1.8870,0.6025)(1.9905,0.6493)(2.0982,0.6983)(2.2102,0.7494)(2.3267,0.8027)(2.4478,0.8583)};
\addplot[main,thin] coordinates{(0.7214,0.6570)(0.7140,0.6591)(0.7058,0.6609)(0.6966,0.6623)(0.6865,0.6633)(0.6753,0.6639)(0.6631,0.6641)(0.6497,0.6639)(0.6351,0.6631)(0.6194,0.6618)(0.6023,0.6601)(0.5839,0.6577)(0.5641,0.6548)(0.5429,0.6512)(0.5201,0.6470)(0.4958,0.6422)(0.4697,0.6366)(0.4420,0.6302)(0.4125,0.6231)(0.3811,0.6152)(0.3477,0.6064)(0.3123,0.5967)(0.2748,0.5860)(0.2350,0.5744)(0.1930,0.5618)(0.1486,0.5481)(0.1017,0.5332)(0.0522,0.5172)(0.0000,0.5000)(-0.0550,0.4815)(-0.1129,0.4617)(-0.1738,0.4404)(-0.2379,0.4178)(-0.3053,0.3936)(-0.3760,0.3678)(-0.4503,0.3404)(-0.5283,0.3113)(-0.6101,0.2804)(-0.6958,0.2477)(-0.7856,0.2130)(-0.8798,0.1763)(-0.9783,0.1376)(-1.0815,0.0966)(-1.1894,0.0534)(-1.3023,0.0078)(-1.4204,-0.0402)(-1.5438,-0.0908)(-1.6727,-0.1441)(-1.8074,-0.2000)(-1.9480,-0.2589)(-2.0949,-0.3207)(-2.2481,-0.3857)(-2.4080,-0.4538)};
\addplot[main,thin] coordinates{(0.1760,0.2405)(0.1654,0.2380)(0.1541,0.2353)(0.1421,0.2322)(0.1294,0.2288)(0.1159,0.2251)(0.1015,0.2210)(0.0863,0.2166)(0.0703,0.2117)(0.0533,0.2065)(0.0354,0.2008)(0.0165,0.1946)(-0.0034,0.1880)(-0.0244,0.1810)(-0.0465,0.1734)(-0.0697,0.1653)(-0.0941,0.1566)(-0.1198,0.1473)(-0.1468,0.1375)(-0.1752,0.1270)(-0.2049,0.1159)(-0.2361,0.1041)(-0.2688,0.0916)(-0.3030,0.0783)(-0.3389,0.0643)(-0.3765,0.0495)(-0.4158,0.0339)(-0.4570,0.0174)(-0.5000,0.0000)(-0.5450,-0.0183)(-0.5920,-0.0376)(-0.6412,-0.0579)(-0.6925,-0.0793)(-0.7461,-0.1018)(-0.8020,-0.1253)(-0.8604,-0.1501)(-0.9213,-0.1761)(-0.9849,-0.2034)(-1.0511,-0.2319)(-1.1202,-0.2619)(-1.1922,-0.2933)(-1.2672,-0.3261)(-1.3454,-0.3605)(-1.4268,-0.3965)(-1.5116,-0.4341)(-1.5999,-0.4735)(-1.6918,-0.5146)(-1.7874,-0.5576)(-1.8870,-0.6025)(-1.9905,-0.6493)(-2.0982,-0.6983)(-2.2102,-0.7494)(-2.3267,-0.8027)(-2.4478,-0.8583)};
\addplot[main,thin] coordinates{(-0.7214,-0.6570)(-0.7140,-0.6591)(-0.7058,-0.6609)(-0.6966,-0.6623)(-0.6865,-0.6633)(-0.6753,-0.6639)(-0.6631,-0.6641)(-0.6497,-0.6639)(-0.6351,-0.6631)(-0.6194,-0.6618)(-0.6023,-0.6601)(-0.5839,-0.6577)(-0.5641,-0.6548)(-0.5429,-0.6512)(-0.5201,-0.6470)(-0.4958,-0.6422)(-0.4697,-0.6366)(-0.4420,-0.6302)(-0.4125,-0.6231)(-0.3811,-0.6152)(-0.3477,-0.6064)(-0.3123,-0.5967)(-0.2748,-0.5860)(-0.2350,-0.5744)(-0.1930,-0.5618)(-0.1486,-0.5481)(-0.1017,-0.5332)(-0.0522,-0.5172)(0.0000,-0.5000)(0.0550,-0.4815)(0.1129,-0.4617)(0.1738,-0.4404)(0.2379,-0.4178)(0.3053,-0.3936)(0.3760,-0.3678)(0.4503,-0.3404)(0.5283,-0.3113)(0.6101,-0.2804)(0.6958,-0.2477)(0.7856,-0.2130)(0.8798,-0.1763)(0.9783,-0.1376)(1.0815,-0.0966)(1.1894,-0.0534)(1.3023,-0.0078)(1.4204,0.0402)(1.5438,0.0908)(1.6727,0.1441)(1.8074,0.2000)(1.9480,0.2589)(2.0949,0.3207)(2.2481,0.3857)(2.4080,0.4538)};
\draw[->,thick](axis cs:1.0,0.0)--(axis cs:1.2965,0.1203);
\draw[->,thick](axis cs:0.0,1.0)--(axis cs:-0.3036,0.8988);
\draw[->,thick](axis cs:-1.0,0.0)--(axis cs:-1.2965,-0.1203);
\draw[->,thick](axis cs:0.0,-1.0)--(axis cs:0.3036,-0.8988);
\end{axis}
\end{tikzpicture}
\par
\begin{tikzpicture}
\begin{axis}[at={(0.0cm,0)},width=4.3cm,height=4cm,xmin=-2,xmax=2,ymin=-2,ymax=2,axis lines=middle,xlabel={$x$},ylabel={$y$},title={$a=11/4$},ticks=none,clip=true]
\addplot[dashed] coordinates{(2.2000,0.8484)(-2.2000,-0.8484)};
\addplot[dashed] coordinates{(2.2000,1.9016)(-2.2000,-1.9016)};
\addplot[main,thin] coordinates{(-0.1608,-0.2268)(-0.1519,-0.2243)(-0.1425,-0.2216)(-0.1324,-0.2187)(-0.1217,-0.2155)(-0.1101,-0.2120)(-0.0979,-0.2082)(-0.0847,-0.2041)(-0.0707,-0.1997)(-0.0558,-0.1949)(-0.0399,-0.1897)(-0.0230,-0.1842)(-0.0050,-0.1782)(0.0142,-0.1718)(0.0347,-0.1649)(0.0564,-0.1575)(0.0795,-0.1496)(0.1041,-0.1411)(0.1302,-0.1321)(0.1580,-0.1224)(0.1876,-0.1120)(0.2189,-0.1010)(0.2522,-0.0892)(0.2876,-0.0766)(0.3252,-0.0631)(0.3651,-0.0488)(0.4074,-0.0336)(0.4523,-0.0173)(0.5000,0.0000)(0.5506,0.0184)(0.6043,0.0380)(0.6612,0.0589)(0.7216,0.0811)(0.7856,0.1047)(0.8535,0.1297)(0.9256,0.1564)(1.0019,0.1847)(1.0829,0.2148)(1.1687,0.2468)(1.2597,0.2807)(1.3561,0.3167)(1.4583,0.3550)(1.5666,0.3956)(1.6814,0.4387)(1.8030,0.4844)(1.9318,0.5329)(2.0684,0.5844)(2.2131,0.6390)(2.3663,0.6969)};
\addplot[main,thin] coordinates{(0.6804,0.6898)(0.6730,0.6893)(0.6649,0.6886)(0.6560,0.6876)(0.6464,0.6864)(0.6359,0.6848)(0.6246,0.6829)(0.6123,0.6806)(0.5990,0.6780)(0.5846,0.6749)(0.5692,0.6715)(0.5525,0.6676)(0.5346,0.6633)(0.5154,0.6584)(0.4947,0.6530)(0.4726,0.6471)(0.4488,0.6406)(0.4234,0.6334)(0.3962,0.6256)(0.3672,0.6170)(0.3361,0.6077)(0.3029,0.5976)(0.2675,0.5866)(0.2297,0.5748)(0.1894,0.5620)(0.1465,0.5482)(0.1007,0.5333)(0.0519,0.5172)(0.0000,0.5000)(-0.0553,0.4815)(-0.1141,0.4616)(-0.1767,0.4403)(-0.2432,0.4175)(-0.3140,0.3931)(-0.3892,0.3670)(-0.4692,0.3391)(-0.5542,0.3092)(-0.6444,0.2774)(-0.7403,0.2433)(-0.8421,0.2070)(-0.9502,0.1684)(-1.0649,0.1271)(-1.1867,0.0832)(-1.3160,0.0363)(-1.4532,-0.0135)(-1.5987,-0.0666)(-1.7531,-0.1230)(-1.9169,-0.1831)(-2.0906,-0.2469)(-2.2748,-0.3148)(-2.4702,-0.3870)};
\addplot[main,thin] coordinates{(0.1608,0.2268)(0.1519,0.2243)(0.1425,0.2216)(0.1324,0.2187)(0.1217,0.2155)(0.1101,0.2120)(0.0979,0.2082)(0.0847,0.2041)(0.0707,0.1997)(0.0558,0.1949)(0.0399,0.1897)(0.0230,0.1842)(0.0050,0.1782)(-0.0142,0.1718)(-0.0347,0.1649)(-0.0564,0.1575)(-0.0795,0.1496)(-0.1041,0.1411)(-0.1302,0.1321)(-0.1580,0.1224)(-0.1876,0.1120)(-0.2189,0.1010)(-0.2522,0.0892)(-0.2876,0.0766)(-0.3252,0.0631)(-0.3651,0.0488)(-0.4074,0.0336)(-0.4523,0.0173)(-0.5000,0.0000)(-0.5506,-0.0184)(-0.6043,-0.0380)(-0.6612,-0.0589)(-0.7216,-0.0811)(-0.7856,-0.1047)(-0.8535,-0.1297)(-0.9256,-0.1564)(-1.0019,-0.1847)(-1.0829,-0.2148)(-1.1687,-0.2468)(-1.2597,-0.2807)(-1.3561,-0.3167)(-1.4583,-0.3550)(-1.5666,-0.3956)(-1.6814,-0.4387)(-1.8030,-0.4844)(-1.9318,-0.5329)(-2.0684,-0.5844)(-2.2131,-0.6390)(-2.3663,-0.6969)};
\addplot[main,thin] coordinates{(-0.6804,-0.6898)(-0.6730,-0.6893)(-0.6649,-0.6886)(-0.6560,-0.6876)(-0.6464,-0.6864)(-0.6359,-0.6848)(-0.6246,-0.6829)(-0.6123,-0.6806)(-0.5990,-0.6780)(-0.5846,-0.6749)(-0.5692,-0.6715)(-0.5525,-0.6676)(-0.5346,-0.6633)(-0.5154,-0.6584)(-0.4947,-0.6530)(-0.4726,-0.6471)(-0.4488,-0.6406)(-0.4234,-0.6334)(-0.3962,-0.6256)(-0.3672,-0.6170)(-0.3361,-0.6077)(-0.3029,-0.5976)(-0.2675,-0.5866)(-0.2297,-0.5748)(-0.1894,-0.5620)(-0.1465,-0.5482)(-0.1007,-0.5333)(-0.0519,-0.5172)(0.0000,-0.5000)(0.0553,-0.4815)(0.1141,-0.4616)(0.1767,-0.4403)(0.2432,-0.4175)(0.3140,-0.3931)(0.3892,-0.3670)(0.4692,-0.3391)(0.5542,-0.3092)(0.6444,-0.2774)(0.7403,-0.2433)(0.8421,-0.2070)(0.9502,-0.1684)(1.0649,-0.1271)(1.1867,-0.0832)(1.3160,-0.0363)(1.4532,0.0135)(1.5987,0.0666)(1.7531,0.1230)(1.9169,0.1831)(2.0906,0.2469)(2.2748,0.3148)(2.4702,0.3870)};
\draw[->,thick](axis cs:1.0,0.0)--(axis cs:1.3007,0.1094);
\draw[->,thick](axis cs:0.0,1.0)--(axis cs:-0.3036,0.8988);
\draw[->,thick](axis cs:-1.0,0.0)--(axis cs:-1.3007,-0.1094);
\draw[->,thick](axis cs:0.0,-1.0)--(axis cs:0.3036,-0.8988);
\end{axis}
\begin{axis}[at={(4.35cm,0)},width=4.3cm,height=4cm,xmin=-2,xmax=2,ymin=-2,ymax=2,axis lines=middle,xlabel={$x$},ylabel={$y$},title={$a=3$},ticks=none,clip=true]
\addplot[dashed] coordinates{(2.2000,0.7333)(-2.2000,-0.7333)};
\addplot[dashed] coordinates{(2.2000,2.2000)(-2.2000,-2.2000)};
\addplot[main,thin] coordinates{(-0.1485,-0.2162)(-0.1410,-0.2137)(-0.1329,-0.2110)(-0.1242,-0.2081)(-0.1149,-0.2050)(-0.1049,-0.2016)(-0.0942,-0.1981)(-0.0827,-0.1942)(-0.0703,-0.1901)(-0.0570,-0.1857)(-0.0427,-0.1809)(-0.0273,-0.1758)(-0.0108,-0.1703)(0.0069,-0.1644)(0.0259,-0.1580)(0.0463,-0.1512)(0.0683,-0.1439)(0.0918,-0.1361)(0.1172,-0.1276)(0.1443,-0.1186)(0.1735,-0.1088)(0.2049,-0.0984)(0.2386,-0.0871)(0.2748,-0.0751)(0.3136,-0.0621)(0.3553,-0.0482)(0.4002,-0.0333)(0.4483,-0.0172)(0.5000,0.0000)(0.5555,0.0185)(0.6152,0.0384)(0.6792,0.0597)(0.7480,0.0827)(0.8219,0.1073)(0.9013,0.1338)(0.9865,0.1622)(1.0781,0.1927)(1.1764,0.2255)(1.2820,0.2607)(1.3955,0.2985)(1.5173,0.3391)(1.6482,0.3827)(1.7887,0.4296)(1.9397,0.4799)(2.1018,0.5339)(2.2759,0.5920)(2.4629,0.6543)};
\addplot[main,thin] coordinates{(0.6485,0.7162)(0.6410,0.7137)(0.6329,0.7110)(0.6242,0.7081)(0.6149,0.7050)(0.6049,0.7016)(0.5942,0.6981)(0.5827,0.6942)(0.5703,0.6901)(0.5570,0.6857)(0.5427,0.6809)(0.5273,0.6758)(0.5108,0.6703)(0.4931,0.6644)(0.4741,0.6580)(0.4537,0.6512)(0.4317,0.6439)(0.4082,0.6361)(0.3828,0.6276)(0.3557,0.6186)(0.3265,0.6088)(0.2951,0.5984)(0.2614,0.5871)(0.2252,0.5751)(0.1864,0.5621)(0.1447,0.5482)(0.0998,0.5333)(0.0517,0.5172)(0.0000,0.5000)(-0.0555,0.4815)(-0.1152,0.4616)(-0.1792,0.4403)(-0.2480,0.4173)(-0.3219,0.3927)(-0.4013,0.3662)(-0.4865,0.3378)(-0.5781,0.3073)(-0.6764,0.2745)(-0.7820,0.2393)(-0.8955,0.2015)(-1.0173,0.1609)(-1.1482,0.1173)(-1.2887,0.0704)(-1.4397,0.0201)(-1.6018,-0.0339)(-1.7759,-0.0920)(-1.9629,-0.1543)(-2.1638,-0.2213)(-2.3796,-0.2932)};
\addplot[main,thin] coordinates{(0.1485,0.2162)(0.1410,0.2137)(0.1329,0.2110)(0.1242,0.2081)(0.1149,0.2050)(0.1049,0.2016)(0.0942,0.1981)(0.0827,0.1942)(0.0703,0.1901)(0.0570,0.1857)(0.0427,0.1809)(0.0273,0.1758)(0.0108,0.1703)(-0.0069,0.1644)(-0.0259,0.1580)(-0.0463,0.1512)(-0.0683,0.1439)(-0.0918,0.1361)(-0.1172,0.1276)(-0.1443,0.1186)(-0.1735,0.1088)(-0.2049,0.0984)(-0.2386,0.0871)(-0.2748,0.0751)(-0.3136,0.0621)(-0.3553,0.0482)(-0.4002,0.0333)(-0.4483,0.0172)(-0.5000,0.0000)(-0.5555,-0.0185)(-0.6152,-0.0384)(-0.6792,-0.0597)(-0.7480,-0.0827)(-0.8219,-0.1073)(-0.9013,-0.1338)(-0.9865,-0.1622)(-1.0781,-0.1927)(-1.1764,-0.2255)(-1.2820,-0.2607)(-1.3955,-0.2985)(-1.5173,-0.3391)(-1.6482,-0.3827)(-1.7887,-0.4296)(-1.9397,-0.4799)(-2.1018,-0.5339)(-2.2759,-0.5920)(-2.4629,-0.6543)};
\addplot[main,thin] coordinates{(-0.6485,-0.7162)(-0.6410,-0.7137)(-0.6329,-0.7110)(-0.6242,-0.7081)(-0.6149,-0.7050)(-0.6049,-0.7016)(-0.5942,-0.6981)(-0.5827,-0.6942)(-0.5703,-0.6901)(-0.5570,-0.6857)(-0.5427,-0.6809)(-0.5273,-0.6758)(-0.5108,-0.6703)(-0.4931,-0.6644)(-0.4741,-0.6580)(-0.4537,-0.6512)(-0.4317,-0.6439)(-0.4082,-0.6361)(-0.3828,-0.6276)(-0.3557,-0.6186)(-0.3265,-0.6088)(-0.2951,-0.5984)(-0.2614,-0.5871)(-0.2252,-0.5751)(-0.1864,-0.5621)(-0.1447,-0.5482)(-0.0998,-0.5333)(-0.0517,-0.5172)(0.0000,-0.5000)(0.0555,-0.4815)(0.1152,-0.4616)(0.1792,-0.4403)(0.2480,-0.4173)(0.3219,-0.3927)(0.4013,-0.3662)(0.4865,-0.3378)(0.5781,-0.3073)(0.6764,-0.2745)(0.7820,-0.2393)(0.8955,-0.2015)(1.0173,-0.1609)(1.1482,-0.1173)(1.2887,-0.0704)(1.4397,-0.0201)(1.6018,0.0339)(1.7759,0.0920)(1.9629,0.1543)(2.1638,0.2213)(2.3796,0.2932)};
\draw[->,thick](axis cs:1.0,0.0)--(axis cs:1.3036,0.1012);
\draw[->,thick](axis cs:0.0,1.0)--(axis cs:-0.3036,0.8988);
\draw[->,thick](axis cs:-1.0,0.0)--(axis cs:-1.3036,-0.1012);
\draw[->,thick](axis cs:0.0,-1.0)--(axis cs:0.3036,-0.8988);
\end{axis}
\begin{axis}[at={(8.7cm,0)},width=4.3cm,height=4cm,xmin=-2,xmax=2,ymin=-2,ymax=2,axis lines=middle,xlabel={$x$},ylabel={$y$},title={$a=4$},ticks=none,clip=true]
\addplot[dashed] coordinates{(2.2000,0.5113)(-2.2000,-0.5113)};
\addplot[dashed] coordinates{(1.5339,2.2000)(-1.5339,-2.2000)};
\addplot[main,thin] coordinates{(-0.1089,-0.1826)(-0.1048,-0.1800)(-0.1003,-0.1772)(-0.0954,-0.1745)(-0.0902,-0.1716)(-0.0844,-0.1686)(-0.0781,-0.1656)(-0.0712,-0.1624)(-0.0636,-0.1591)(-0.0553,-0.1556)(-0.0461,-0.1519)(-0.0359,-0.1480)(-0.0246,-0.1439)(-0.0120,-0.1395)(0.0019,-0.1347)(0.0175,-0.1297)(0.0348,-0.1242)(0.0541,-0.1183)(0.0757,-0.1119)(0.0998,-0.1049)(0.1268,-0.0972)(0.1570,-0.0888)(0.1909,-0.0796)(0.2288,-0.0695)(0.2713,-0.0583)(0.3190,-0.0459)(0.3725,-0.0322)(0.4326,-0.0169)(0.5000,0.0000)(0.5757,0.0189)(0.6608,0.0399)(0.7564,0.0633)(0.8638,0.0895)(0.9846,0.1187)(1.1203,0.1515)(1.2729,0.1881)(1.4444,0.2291)(1.6373,0.2751)(1.8542,0.3266)(2.0982,0.3845)(2.3725,0.4493)};
\addplot[main,thin] coordinates{(0.5478,0.8041)(0.5399,0.7950)(0.5317,0.7859)(0.5234,0.7768)(0.5148,0.7678)(0.5059,0.7588)(0.4967,0.7497)(0.4871,0.7407)(0.4772,0.7316)(0.4667,0.7225)(0.4556,0.7133)(0.4440,0.7041)(0.4316,0.6947)(0.4184,0.6853)(0.4042,0.6757)(0.3890,0.6659)(0.3727,0.6559)(0.3549,0.6456)(0.3356,0.6351)(0.3146,0.6242)(0.2917,0.6130)(0.2665,0.6013)(0.2389,0.5890)(0.2085,0.5762)(0.1749,0.5627)(0.1377,0.5485)(0.0965,0.5334)(0.0508,0.5172)(0.0000,0.5000)(-0.0566,0.4815)(-0.1196,0.4615)(-0.1899,0.4399)(-0.2684,0.4164)(-0.3562,0.3909)(-0.4544,0.3630)(-0.5643,0.3324)(-0.6874,0.2988)(-0.8253,0.2618)(-0.9799,0.2210)(-1.1534,0.1759)(-1.3480,0.1259)(-1.5664,0.0704)(-1.8117,0.0087)(-2.0871,-0.0599)(-2.3965,-0.1364)};
\addplot[main,thin] coordinates{(0.1089,0.1826)(0.1048,0.1800)(0.1003,0.1772)(0.0954,0.1745)(0.0902,0.1716)(0.0844,0.1686)(0.0781,0.1656)(0.0712,0.1624)(0.0636,0.1591)(0.0553,0.1556)(0.0461,0.1519)(0.0359,0.1480)(0.0246,0.1439)(0.0120,0.1395)(-0.0019,0.1347)(-0.0175,0.1297)(-0.0348,0.1242)(-0.0541,0.1183)(-0.0757,0.1119)(-0.0998,0.1049)(-0.1268,0.0972)(-0.1570,0.0888)(-0.1909,0.0796)(-0.2288,0.0695)(-0.2713,0.0583)(-0.3190,0.0459)(-0.3725,0.0322)(-0.4326,0.0169)(-0.5000,0.0000)(-0.5757,-0.0189)(-0.6608,-0.0399)(-0.7564,-0.0633)(-0.8638,-0.0895)(-0.9846,-0.1187)(-1.1203,-0.1515)(-1.2729,-0.1881)(-1.4444,-0.2291)(-1.6373,-0.2751)(-1.8542,-0.3266)(-2.0982,-0.3845)(-2.3725,-0.4493)};
\addplot[main,thin] coordinates{(-0.5478,-0.8041)(-0.5399,-0.7950)(-0.5317,-0.7859)(-0.5234,-0.7768)(-0.5148,-0.7678)(-0.5059,-0.7588)(-0.4967,-0.7497)(-0.4871,-0.7407)(-0.4772,-0.7316)(-0.4667,-0.7225)(-0.4556,-0.7133)(-0.4440,-0.7041)(-0.4316,-0.6947)(-0.4184,-0.6853)(-0.4042,-0.6757)(-0.3890,-0.6659)(-0.3727,-0.6559)(-0.3549,-0.6456)(-0.3356,-0.6351)(-0.3146,-0.6242)(-0.2917,-0.6130)(-0.2665,-0.6013)(-0.2389,-0.5890)(-0.2085,-0.5762)(-0.1749,-0.5627)(-0.1377,-0.5485)(-0.0965,-0.5334)(-0.0508,-0.5172)(0.0000,-0.5000)(0.0566,-0.4815)(0.1196,-0.4615)(0.1899,-0.4399)(0.2684,-0.4164)(0.3562,-0.3909)(0.4544,-0.3630)(0.5643,-0.3324)(0.6874,-0.2988)(0.8253,-0.2618)(0.9799,-0.2210)(1.1534,-0.1759)(1.3480,-0.1259)(1.5664,-0.0704)(1.8117,-0.0087)(2.0871,0.0599)(2.3965,0.1364)};
\draw[->,thick](axis cs:1.0,0.0)--(axis cs:1.3104,0.0776);
\draw[->,thick](axis cs:0.0,1.0)--(axis cs:-0.3036,0.8988);
\draw[->,thick](axis cs:-1.0,0.0)--(axis cs:-1.3104,-0.0776);
\draw[->,thick](axis cs:0.0,-1.0)--(axis cs:0.3036,-0.8988);
\end{axis}
\end{tikzpicture}
\par
\end{center}
\end{solution}
\begin{exercise}\label{cw:ps9-II-36}原题 II.36（矩阵指数）：
(a) 如习题课 21 第 4 题，复数 $z=a+bi$ 对应 $A(z)=\begin{pmatrix}a&-b\\b&a\end{pmatrix}$，它表示乘法：$z(x+yi)=v+wi$ 时 $A(z)(x,y)^T=(v,w)^T$。求 $e^{A(z)t}$、$A(e^{zt})$。
(b) 若 $mx''+bx'+kx=0$ 的一对解满足 $x_1(0)=1,x_1'(0)=0,x_2(0)=0,x_2'(0)=1$，称在 $0$ 归一化。求 $x''+2x'+2x=0$ 的归一化解对，再求算子 $D^2+2D+2I$ 伴随矩阵的 $e^{At}$，作评论。
(c) 已知 $e^{3t}(1,1)^T$ 与 $e^{2t}(1,2)^T$ 满足 $\mathbf u'=A\mathbf u$：(i) 求 $\mathbf u_1(0)=(1,0)^T,\mathbf u_2(0)=(0,1)^T$ 的两个解；(ii) 求 $e^{At}$；(iii) 求 $A$。\end{exercise}
\begin{solution}
\textit{官方解译审与补证（AI 辅助）；其中另标编者补解或原文修正。}\par

(a) 记 $J=\begin{pmatrix}0&-1\\1&0\end{pmatrix}$，则 $J^2=-I$、$A=aI+bJ$。两部分可交换，故
\[e^{A(z)t}=e^{at}(I\cos bt+J\sin bt)=e^{at}\begin{pmatrix}\cos bt&-\sin bt\\\sin bt&\cos bt\end{pmatrix}.\]
另一方面 $e^{zt}=e^{at}(\cos bt+i\sin bt)$，代入映射 $A$ 得到同一矩阵。因此 $A(e^{zt})=e^{A(z)t}$，包括 $b=0$ 的情况。此等式反映复乘法与其二维实矩阵表示相容。
(b) 基本解 $y_1=e^{-t}\cos t,y_2=e^{-t}\sin t$ 的初值为 $(1,-1)$、$(0,1)$，所以
\[x_1=y_1+y_2=e^{-t}(\cos t+\sin t),\qquad x_2=y_2=e^{-t}\sin t.\]
$A=\begin{pmatrix}0&1\\-2&-2\end{pmatrix}$，将 $x_j$ 及 $x_j'$ 作第 $j$ 列，得
\[e^{At}=\begin{pmatrix}x_1&x_2\\x_1'&x_2'\end{pmatrix}=e^{-t}\begin{pmatrix}\cos t+\sin t&\sin t\\-2\sin t&\cos t-\sin t\end{pmatrix}.\]
此矩阵在 $t=0$ 为 $I$，各列解伴随系统，故由初值唯一性确为矩阵指数。一般基本解矩阵若不在 $0$ 等于 $I$，不能直接当 $e^{At}$；归一化正是将两列分别对应于两个坐标初值。
(c)(i) 令 $\mathbf u=c_1e^{3t}(1,1)^T+c_2e^{2t}(1,2)^T$，零点值为 $(c_1+c_2,c_1+2c_2)^T$。第一个初值得 $c_1=2,c_2=-1$，第二个得 $c_1=-1,c_2=1$，于是
\[\mathbf u_1=\binom{2e^{3t}-e^{2t}}{2e^{3t}-2e^{2t}},\qquad\mathbf u_2=\binom{-e^{3t}+e^{2t}}{-e^{3t}+2e^{2t}}.\]
(ii) 两列构成在 $0$ 归一化的基本矩阵：
\[e^{At}=\begin{pmatrix}2e^{3t}-e^{2t}&-e^{3t}+e^{2t}\\2e^{3t}-2e^{2t}&-e^{3t}+2e^{2t}\end{pmatrix}.\]
(iii) 在 $0$ 求导得 $A=\begin{pmatrix}4&-1\\2&1\end{pmatrix}$。亦可令 $S=\begin{pmatrix}1&1\\1&2\end{pmatrix}$，由 $AS=S\operatorname{diag}(3,2)$ 得 $A=S\operatorname{diag}(3,2)S^{-1}$，并核实两个给定模态确实满足系统。
\end{solution}

% MIT 18.03 S2010：实际22份习题课；原题编号保留于标签。
% 题面和解答分别登记，来源/页码/哈希见research/recitations-coverage.json。
\originalchapter{习题课：一阶模型与复数响应}
本部分保留原课程习题课编号。公开文件为第1--5、7--11、13--19、21--25次，共22次；第6、12、20次为考试复习，第26次为课程复习，均没有本次清单中的公开题册。每题的中文题面与解答分列；“官方译审”表示核对过配套解答后重新组织推导，“编者补解（AI辅助）”表示补上官方没有提供的解答；“译审并补证（AI辅助）”表示在官方路线基础上补入论证。此稿由AI辅助撰写，独立第二读者为同模型另次审阅，不是MIT官方答案或外部专家鉴定。源文件所载日期均为2010年；文件作者未独立署名时沿用课程署名，不推测著作年份。原题提到的 Mathlet 在此用方程分析和自绘图落实，不声称运行过交互程序。

\originalsection{第1次习题课：自然增长、捕猎与对数（2月2日）}
\begin{exercise}\label{cw:rec01-1}
某非洲政府为剑羚捕猎制定政策。设种群自然增长率为 $k$，恒定捕猎量为每年 $a$ 只。选择变量与单位，建立种群模型。
\end{exercise}
\begin{solution}
以年为时间单位，$x(t)$ 为剑羚只数，$k$ 的单位为年$^{-1}$，$a$ 为只/年。短时间内增加 $kx\Delta t$ 只、捕猎 $a\Delta t$ 只，除以 $\Delta t$ 并取极限得 $\dot x=kx-a$。两端均为只/年。通常模型取 $k>0,a\ge0$，只在 $x\ge0$ 时有种群意义。
\par\noindent{\small 依据：官方译审并补证（AI辅助）。}
\end{solution}

\begin{exercise}\label{cw:rec01-2}
上述模型中若 $a=0$，种群的倍增时间是多少？
\end{exercise}
\begin{solution}
由 $\dot x=kx$ 得 $x(t)=x(0)\ee^{kt}$。正初值及 $k>0$ 时，$\ee^{kT}=2$，故 $T=\log2/k$ 年。零种群没有倍增时间；$k\le0$ 时没有正向倍增。
\par\noindent{\small 依据：官方译审并补证（AI辅助）。}
\end{solution}

\begin{exercise}\label{cw:rec01-3}
求 $\dot x=kx-a$ 的通解。
\end{exercise}
\begin{solution}
$k\ne0$ 时令 $z=x-a/k$，则 $\dot z=kz$，故 $x=a/k+C\ee^{kt}$。以初值表示为 $x=a/k+(x_0-a/k)\ee^{kt}$。若 $k=0$，直接积分得 $x=x_0-at$。
\par\noindent{\small 依据：官方译审并补证（AI辅助）。}
\end{solution}

\begin{exercise}\label{cw:rec01-4}
直接核验上一题提出的解满足原微分方程。
\end{exercise}
\begin{solution}
$k\ne0$ 时 $\dot x=kC\ee^{kt}=k(a/k+C\ee^{kt})-a$；$k=0$ 时 $\dot x=-a$。此外代入 $t=0$ 得指定初值，故方程与初值同时满足。
\par\noindent{\small 依据：官方译审并补证（AI辅助）。}
\end{solution}

\begin{exercise}\label{cw:rec01-5}
找出模型的稳态（常值、平衡）解；检查单位，以及对 $k,a$ 的大小依赖是否符合种群含义。
\end{exercise}
\begin{solution}
令 $\dot x=0$ 得 $x_*=a/k$（$k\ne0$）。单位为只/年除以年$^{-1}$，即只；捕猎量越大，需要越大的自然种群才能抵消损失，增长率越大则所需平衡种群越小。若 $k=0,a>0$ 无平衡；若 $k=a=0$ 所有常值均为平衡。
\par\noindent{\small 依据：官方译审并补证（AI辅助）。}
\end{solution}

\begin{exercise}\label{cw:rec01-6}
画出平衡解及其他解，说明不稳定性。低于平衡的初值何时灭绝？特别地，初值为平衡的一半时，求灭绝时间。
\end{exercise}
\begin{solution}
取 $k>0,a>0$，有 $x-x_*=(x_0-x_*)\ee^{kt}$：高于平衡则增长，低于平衡则下降，故平衡不稳定。$0\le x_0<x_*$ 时解 $x(t_0)=0$ 得 $t_0=k^{-1}\log[a/(a-kx_0)]$；$x_0=x_*/2$ 时为 $\log2/k$。负值应截断为灭绝，而过大的正值亦超出无限资源模型的有效范围。图用无量纲量 $s=kt,X=kx/a$，曲线族 $X=1+(X_0-1)\ee^s$。\begin{center}\begin{tikzpicture}\begin{axis}[width=.68\textwidth,height=4.5cm,axis lines=middle,xlabel={$s$},ylabel={$X$},samples=130,xmin=0,xmax=1.5,ymin=0,ymax=4]
\addplot[black,domain=0:1.5]{1};\foreach \c in {-1,-.5,.5,1}{\addplot[blue,domain=0:1.5]{1+\c*exp(x)};}\end{axis}\end{tikzpicture}\end{center}
\par\noindent{\small 依据：官方译审并补证（AI辅助）。}
\end{solution}

\originalsection{第2次习题课：方向场、等斜线、分界线与漏斗（2月4日）}
\begin{exercise}\label{cw:rec02-1}
对 $y'=x-2y$ 画若干等斜线，特别是零斜线，绘制方向场并画几条积分曲线。等斜线指 $F(x,y)=m$ 的点集。
\end{exercise}
\begin{solution}
等斜线为 $y=x/2-m/2$；零斜线为 $y=x/2$。每条等斜线上放斜率 $m$ 的短线段。由积分因子计算的曲线 $y=x/2-1/4+C\ee^{-2x}$ 可用于核对图。蓝线为解，虚线为零斜线与斜率 $1/2$ 的直线解。\begin{center}\begin{tikzpicture}\begin{axis}[width=.68\textwidth,height=4.5cm,axis lines=middle,xlabel={$x$},ylabel={$y$},samples=130,xmin=-.5,xmax=3,ymin=-1.5,ymax=2]
\addplot[dashed,domain=-.5:3]{x/2};\addplot[dashed,domain=-.5:3]{x/2-.25};\foreach \c in {-.5,0,.5,1}{\addplot[blue,domain=-.5:3]{x/2-.25+\c*exp(-2*x)};}\draw[gray](axis cs:0,-1)--(axis cs:0.05367,-0.89267);\draw[gray](axis cs:0,-0.5)--(axis cs:0.08485,-0.41515);\draw[gray](axis cs:0,0)--(axis cs:0.12000,0.00000);\draw[gray](axis cs:0,0.5)--(axis cs:0.08485,0.41515);\draw[gray](axis cs:0,1)--(axis cs:0.05367,0.89267);\draw[gray](axis cs:0,1.5)--(axis cs:0.03795,1.38616);\draw[gray](axis cs:0.5,-1)--(axis cs:0.54457,-0.88858);\draw[gray](axis cs:0.5,-0.5)--(axis cs:0.56656,-0.40015);\draw[gray](axis cs:0.5,0)--(axis cs:0.60733,0.05367);\draw[gray](axis cs:0.5,0.5)--(axis cs:0.60733,0.44633);\draw[gray](axis cs:0.5,1)--(axis cs:0.56656,0.90015);\draw[gray](axis cs:0.5,1.5)--(axis cs:0.54457,1.38858);\draw[gray](axis cs:1,-1)--(axis cs:1.03795,-0.88616);\draw[gray](axis cs:1,-0.5)--(axis cs:1.05367,-0.39267);\draw[gray](axis cs:1,0)--(axis cs:1.08485,0.08485);\draw[gray](axis cs:1,0.5)--(axis cs:1.12000,0.50000);\draw[gray](axis cs:1,1)--(axis cs:1.08485,0.91515);\draw[gray](axis cs:1,1.5)--(axis cs:1.05367,1.39267);\draw[gray](axis cs:1.5,-1)--(axis cs:1.53297,-0.88462);\draw[gray](axis cs:1.5,-0.5)--(axis cs:1.54457,-0.38858);\draw[gray](axis cs:1.5,0)--(axis cs:1.56656,0.09985);\draw[gray](axis cs:1.5,0.5)--(axis cs:1.60733,0.55367);\draw[gray](axis cs:1.5,1)--(axis cs:1.60733,0.94633);\draw[gray](axis cs:1.5,1.5)--(axis cs:1.56656,1.40015);\draw[gray](axis cs:2,-1)--(axis cs:2.02910,-0.88358);\draw[gray](axis cs:2,-0.5)--(axis cs:2.03795,-0.38616);\draw[gray](axis cs:2,0)--(axis cs:2.05367,0.10733);\draw[gray](axis cs:2,0.5)--(axis cs:2.08485,0.58485);\draw[gray](axis cs:2,1)--(axis cs:2.12000,1.00000);\draw[gray](axis cs:2,1.5)--(axis cs:2.08485,1.41515);\draw[gray](axis cs:2.5,-1)--(axis cs:2.52603,-0.88286);\draw[gray](axis cs:2.5,-0.5)--(axis cs:2.53297,-0.38462);\draw[gray](axis cs:2.5,0)--(axis cs:2.54457,0.11142);\draw[gray](axis cs:2.5,0.5)--(axis cs:2.56656,0.59985);\draw[gray](axis cs:2.5,1)--(axis cs:2.60733,1.05367);\draw[gray](axis cs:2.5,1.5)--(axis cs:2.60733,1.44633);\end{axis}\end{tikzpicture}\end{center}
\par\noindent{\small 依据：官方译审并补证（AI辅助）。}
\end{solution}

\begin{exercise}\label{cw:rec02-2}
图中似乎有一条直线积分曲线。判定 $y=mx+b$ 是否可能为解，并求 $m,b$。
\end{exercise}
\begin{solution}
代入得 $m=(1-2m)x-2b$。逐项比较：$m=1/2$、$b=-1/4$，所以 $y=x/2-1/4$ 确为解。
\par\noindent{\small 依据：官方译审并补证（AI辅助）。}
\end{solution}

\begin{exercise}\label{cw:rec02-3}
对一般方程 $y'=F(x,y)$，若一条直线是积分曲线，它与等斜线有何关系？说明本题中的情况。
\end{exercise}
\begin{solution}
直线 $y=mx+b$ 是解意味着在线上处处 $F(x,mx+b)=m$，故这条线属于斜率 $m$ 的等斜线；等斜线的其他分支未必也是解。本例直线解恰是整个 $m=1/2$ 等斜线。
\par\noindent{\small 依据：官方译审并补证（AI辅助）。}
\end{solution}

\begin{exercise}\label{cw:rec02-4}
图中各解在 $x\to\infty$ 时似乎互相渐近。先用方向场说明，为什么它们会被固定斜率的平行线夹住。
\end{exercise}
\begin{solution}
取 $y=x/2-m/2$。解与边界相遇时，相对边界的导数为 $m-1/2$。下边界取 $m>1/2$，解只能向内向上穿越；上边界取 $m<1/2$，解只能向内向下穿越。因此包含初值的这两条线围成正向不变带。任意窄带最终都会捕获解，因 $y-(x/2-1/4)=C\ee^{-2x}\to0$。
\par\noindent{\small 依据：官方译审并补证（AI辅助）。}
\end{solution}

\begin{exercise}\label{cw:rec02-5}
$y'=x-2y$ 的解在哪里有驻点？一条解最多有几个？哪些初值 $y(0)=y_0$ 会产生驻点？它是极小还是极大？用二阶导数检验。
\end{exercise}
\begin{solution}
驻点位于 $y=x/2$。通解中 $C=y_0+1/4$，$y'=1/2-2C\ee^{-2x}$，所以在全实轴上有驻点当且仅当 $C>0$，且唯一，位置为 $x_c=\frac12\log(4C)$。若只考察 $x\ge0$，要求 $y_0\ge0$（严格内部要求 $y_0>0$）。在驻点 $y''=1-2y'=1>0$，故是极小；驻点前后导数由负转正。
\par\noindent{\small 依据：官方译审并补证（AI辅助）。}
\end{solution}

\begin{exercise}\label{cw:rec02-6}
对另一方程 $y'=y^2-x^2$，画若干等斜线、方向场与几条解；原题建议用 Isoclines Mathlet。
\end{exercise}
\begin{solution}
等斜线为 $y^2-x^2=m$，$m=0$ 为 $y=\pm x$；在 $|y|>|x|$ 区域斜率为正，内部为负。图中灰色短线据 $F=y^2-x^2$ 绘制，虚线为零斜线。蓝线据数值积分只作示意；下两题用解析论证确定分界与漏斗，数值图不作为证明。\begin{center}\begin{tikzpicture}\begin{axis}[width=.68\textwidth,height=4.5cm,axis lines=middle,xlabel={$x$},ylabel={$y$},samples=130,xmin=0,xmax=3,ymin=-3,ymax=3]
\addplot[dashed,domain=0:3]{x};\addplot[dashed,domain=0:3]{-x};\draw[gray](axis cs:0,-3)--(axis cs:0.01325,-2.88073);\draw[gray](axis cs:0,-2.5)--(axis cs:0.01896,-2.38151);\draw[gray](axis cs:0,-2)--(axis cs:0.02910,-1.88358);\draw[gray](axis cs:0,-1.5)--(axis cs:0.04874,-1.39034);\draw[gray](axis cs:0,-1)--(axis cs:0.08485,-0.91515);\draw[gray](axis cs:0,-0.5)--(axis cs:0.11642,-0.47090);\draw[gray](axis cs:0,0)--(axis cs:0.12000,0.00000);\draw[gray](axis cs:0,0.5)--(axis cs:0.11642,0.52910);\draw[gray](axis cs:0,1)--(axis cs:0.08485,1.08485);\draw[gray](axis cs:0,1.5)--(axis cs:0.04874,1.60966);\draw[gray](axis cs:0,2)--(axis cs:0.02910,2.11642);\draw[gray](axis cs:0,2.5)--(axis cs:0.01896,2.61849);\draw[gray](axis cs:0,3)--(axis cs:0.01325,3.11927);\draw[gray](axis cs:0.4,-3)--(axis cs:0.41349,-2.88076);\draw[gray](axis cs:0.4,-2.5)--(axis cs:0.41944,-2.38159);\draw[gray](axis cs:0.4,-2)--(axis cs:0.43024,-1.88387);\draw[gray](axis cs:0.4,-1.5)--(axis cs:0.45179,-1.39175);\draw[gray](axis cs:0.4,-1)--(axis cs:0.49188,-0.92282);\draw[gray](axis cs:0.4,-0.5)--(axis cs:0.51952,-0.48924);\draw[gray](axis cs:0.4,0)--(axis cs:0.51849,-0.01896);\draw[gray](axis cs:0.4,0.5)--(axis cs:0.51952,0.51076);\draw[gray](axis cs:0.4,1)--(axis cs:0.49188,1.07718);\draw[gray](axis cs:0.4,1.5)--(axis cs:0.45179,1.60825);\draw[gray](axis cs:0.4,2)--(axis cs:0.43024,2.11613);\draw[gray](axis cs:0.4,2.5)--(axis cs:0.41944,2.61841);\draw[gray](axis cs:0.4,3)--(axis cs:0.41349,3.11924);\draw[gray](axis cs:0.8,-3)--(axis cs:0.81425,-2.88085);\draw[gray](axis cs:0.8,-2.5)--(axis cs:0.82106,-2.38186);\draw[gray](axis cs:0.8,-2)--(axis cs:0.83423,-1.88499);\draw[gray](axis cs:0.8,-1.5)--(axis cs:0.86332,-1.39806);\draw[gray](axis cs:0.8,-1)--(axis cs:0.91291,-0.95935);\draw[gray](axis cs:0.8,-0.5)--(axis cs:0.91180,-0.54360);\draw[gray](axis cs:0.8,0)--(axis cs:0.90107,-0.06469);\draw[gray](axis cs:0.8,0.5)--(axis cs:0.91180,0.45640);\draw[gray](axis cs:0.8,1)--(axis cs:0.91291,1.04065);\draw[gray](axis cs:0.8,1.5)--(axis cs:0.86332,1.60194);\draw[gray](axis cs:0.8,2)--(axis cs:0.83423,2.11501);\draw[gray](axis cs:0.8,2.5)--(axis cs:0.82106,2.61814);\draw[gray](axis cs:0.8,3)--(axis cs:0.81425,3.11915);\draw[gray](axis cs:1.2,-3)--(axis cs:1.21574,-2.88104);\draw[gray](axis cs:1.2,-2.5)--(axis cs:1.22443,-2.38251);\draw[gray](axis cs:1.2,-2)--(axis cs:1.24366,-1.88823);\draw[gray](axis cs:1.2,-1.5)--(axis cs:1.29325,-1.42447);\draw[gray](axis cs:1.2,-1)--(axis cs:1.30984,-1.04833);\draw[gray](axis cs:1.2,-0.5)--(axis cs:1.27720,-0.59187);\draw[gray](axis cs:1.2,0)--(axis cs:1.26845,-0.09856);\draw[gray](axis cs:1.2,0.5)--(axis cs:1.27720,0.40813);\draw[gray](axis cs:1.2,1)--(axis cs:1.30984,0.95167);\draw[gray](axis cs:1.2,1.5)--(axis cs:1.29325,1.57553);\draw[gray](axis cs:1.2,2)--(axis cs:1.24366,2.11177);\draw[gray](axis cs:1.2,2.5)--(axis cs:1.22443,2.61749);\draw[gray](axis cs:1.2,3)--(axis cs:1.21574,3.11896);\draw[gray](axis cs:1.6,-3)--(axis cs:1.61841,-2.88142);\draw[gray](axis cs:1.6,-2.5)--(axis cs:1.63139,-2.38418);\draw[gray](axis cs:1.6,-2)--(axis cs:1.66845,-1.90144);\draw[gray](axis cs:1.6,-1.5)--(axis cs:1.71462,-1.53553);\draw[gray](axis cs:1.6,-1)--(axis cs:1.66476,-1.10103);\draw[gray](axis cs:1.6,-0.5)--(axis cs:1.64767,-0.61012);\draw[gray](axis cs:1.6,0)--(axis cs:1.64366,-0.11177);\draw[gray](axis cs:1.6,0.5)--(axis cs:1.64767,0.38988);\draw[gray](axis cs:1.6,1)--(axis cs:1.66476,0.89897);\draw[gray](axis cs:1.6,1.5)--(axis cs:1.71462,1.46447);\draw[gray](axis cs:1.6,2)--(axis cs:1.66845,2.09856);\draw[gray](axis cs:1.6,2.5)--(axis cs:1.63139,2.61582);\draw[gray](axis cs:1.6,3)--(axis cs:1.61841,3.11858);\draw[gray](axis cs:2,-3)--(axis cs:2.02353,-2.88233);\draw[gray](axis cs:2,-2.5)--(axis cs:2.04874,-2.39034);\draw[gray](axis cs:2,-2)--(axis cs:2.12000,-2.00000);\draw[gray](axis cs:2,-1.5)--(axis cs:2.05954,-1.60419);\draw[gray](axis cs:2,-1)--(axis cs:2.03795,-1.11384);\draw[gray](axis cs:2,-0.5)--(axis cs:2.03092,-0.61595);\draw[gray](axis cs:2,0)--(axis cs:2.02910,-0.11642);\draw[gray](axis cs:2,0.5)--(axis cs:2.03092,0.38405);\draw[gray](axis cs:2,1)--(axis cs:2.03795,0.88616);\draw[gray](axis cs:2,1.5)--(axis cs:2.05954,1.39581);\draw[gray](axis cs:2,2)--(axis cs:2.12000,2.00000);\draw[gray](axis cs:2,2.5)--(axis cs:2.04874,2.60966);\draw[gray](axis cs:2,3)--(axis cs:2.02353,3.11767);\draw[gray](axis cs:2.4,-3)--(axis cs:2.43539,-2.88534);\draw[gray](axis cs:2.4,-2.5)--(axis cs:2.50776,-2.44720);\draw[gray](axis cs:2.4,-2)--(axis cs:2.45928,-2.10433);\draw[gray](axis cs:2.4,-1.5)--(axis cs:2.43288,-1.61541);\draw[gray](axis cs:2.4,-1)--(axis cs:2.42467,-1.11744);\draw[gray](axis cs:2.4,-0.5)--(axis cs:2.42143,-0.61807);\draw[gray](axis cs:2.4,0)--(axis cs:2.42053,-0.11823);\draw[gray](axis cs:2.4,0.5)--(axis cs:2.42143,0.38193);\draw[gray](axis cs:2.4,1)--(axis cs:2.42467,0.88256);\draw[gray](axis cs:2.4,1.5)--(axis cs:2.43288,1.38459);\draw[gray](axis cs:2.4,2)--(axis cs:2.45928,1.89567);\draw[gray](axis cs:2.4,2.5)--(axis cs:2.50776,2.55280);\draw[gray](axis cs:2.4,3)--(axis cs:2.43539,3.11466);\draw[gray](axis cs:2.8,-3)--(axis cs:2.87835,-2.90911);\draw[gray](axis cs:2.8,-2.5)--(axis cs:2.86389,-2.60158);\draw[gray](axis cs:2.8,-2)--(axis cs:2.83024,-2.11613);\draw[gray](axis cs:2.8,-1.5)--(axis cs:2.82113,-1.61812);\draw[gray](axis cs:2.8,-1)--(axis cs:2.81736,-1.11874);\draw[gray](axis cs:2.8,-0.5)--(axis cs:2.81567,-0.61897);\draw[gray](axis cs:2.8,0)--(axis cs:2.81518,-0.11904);\draw[gray](axis cs:2.8,0.5)--(axis cs:2.81567,0.38103);\draw[gray](axis cs:2.8,1)--(axis cs:2.81736,0.88126);\draw[gray](axis cs:2.8,1.5)--(axis cs:2.82113,1.38188);\draw[gray](axis cs:2.8,2)--(axis cs:2.83024,1.88387);\draw[gray](axis cs:2.8,2.5)--(axis cs:2.86389,2.39842);\draw[gray](axis cs:2.8,3)--(axis cs:2.87835,3.09089);\end{axis}\end{tikzpicture}\end{center}
下图补入蓝色解曲线及红色分界线，均由原方程构造；红线严格的定义和证明见下一题。\begin{center}\begin{tikzpicture}\begin{axis}[width=.68\textwidth,height=5cm,axis lines=middle,xlabel={$x$},ylabel={$y$},xmin=0,xmax=3,ymin=-3,ymax=3,clip=true]\addplot[blue]coordinates{(0.00000,-1.00000)(0.04348,-0.95836)(0.08696,-0.92021)(0.13043,-0.88531)(0.17391,-0.85348)(0.21739,-0.82456)(0.26087,-0.79843)(0.30435,-0.77502)(0.34783,-0.75424)(0.39130,-0.73606)(0.43478,-0.72043)(0.47826,-0.70736)(0.52174,-0.69681)(0.56522,-0.68881)(0.60870,-0.68334)(0.65217,-0.68042)(0.69565,-0.68007)(0.73913,-0.68229)(0.78261,-0.68710)(0.82609,-0.69450)(0.86957,-0.70450)(0.91304,-0.71709)(0.95652,-0.73227)(1.00000,-0.75002)(1.04348,-0.77030)(1.08696,-0.79309)(1.13043,-0.81832)(1.17391,-0.84595)(1.21739,-0.87590)(1.26087,-0.90808)(1.30435,-0.94240)(1.34783,-0.97876)(1.39130,-1.01703)(1.43478,-1.05710)(1.47826,-1.09883)(1.52174,-1.14208)(1.56522,-1.18672)(1.60870,-1.23261)(1.65217,-1.27959)(1.69565,-1.32754)(1.73913,-1.37632)(1.78261,-1.42579)(1.82609,-1.47583)(1.86957,-1.52632)(1.91304,-1.57715)(1.95652,-1.62823)(2.00000,-1.67946)(2.04348,-1.73076)(2.08696,-1.78207)(2.13043,-1.83332)(2.17391,-1.88446)(2.21739,-1.93545)(2.26087,-1.98626)(2.30435,-2.03686)(2.34783,-2.08723)(2.39130,-2.13736)(2.43478,-2.18723)(2.47826,-2.23685)(2.52174,-2.28621)(2.56522,-2.33532)(2.60870,-2.38418)(2.65217,-2.43280)(2.69565,-2.48119)(2.73913,-2.52935)(2.78261,-2.57730)(2.82609,-2.62504)(2.86957,-2.67260)(2.91304,-2.71997)(2.95652,-2.76717)(3.00000,-2.81422)};\addplot[blue]coordinates{(0.00000,0.00000)(0.04348,-0.00003)(0.08696,-0.00022)(0.13043,-0.00074)(0.17391,-0.00175)(0.21739,-0.00342)(0.26087,-0.00592)(0.30435,-0.00939)(0.34783,-0.01402)(0.39130,-0.01995)(0.43478,-0.02735)(0.47826,-0.03637)(0.52174,-0.04717)(0.56522,-0.05990)(0.60870,-0.07469)(0.65217,-0.09168)(0.69565,-0.11098)(0.73913,-0.13272)(0.78261,-0.15699)(0.82609,-0.18386)(0.86957,-0.21340)(0.91304,-0.24566)(0.95652,-0.28065)(1.00000,-0.31837)(1.04348,-0.35878)(1.08696,-0.40183)(1.13043,-0.44744)(1.17391,-0.49550)(1.21739,-0.54588)(1.26087,-0.59841)(1.30435,-0.65292)(1.34783,-0.70921)(1.39130,-0.76708)(1.43478,-0.82630)(1.47826,-0.88664)(1.52174,-0.94788)(1.56522,-1.00980)(1.60870,-1.07218)(1.65217,-1.13481)(1.69565,-1.19750)(1.73913,-1.26008)(1.78261,-1.32239)(1.82609,-1.38430)(1.86957,-1.44569)(1.91304,-1.50646)(1.95652,-1.56656)(2.00000,-1.62592)(2.04348,-1.68450)(2.08696,-1.74228)(2.13043,-1.79927)(2.17391,-1.85545)(2.21739,-1.91086)(2.26087,-1.96550)(2.30435,-2.01942)(2.34783,-2.07264)(2.39130,-2.12521)(2.43478,-2.17716)(2.47826,-2.22854)(2.52174,-2.27938)(2.56522,-2.32973)(2.60870,-2.37963)(2.65217,-2.42911)(2.69565,-2.47820)(2.73913,-2.52695)(2.78261,-2.57537)(2.82609,-2.62351)(2.86957,-2.67138)(2.91304,-2.71901)(2.95652,-2.76641)(3.00000,-2.81362)};\addplot[blue]coordinates{(0.00000,0.50000)(0.04348,0.51108)(0.08696,0.52250)(0.13043,0.53412)(0.17391,0.54578)(0.21739,0.55734)(0.26087,0.56863)(0.30435,0.57948)(0.34783,0.58971)(0.39130,0.59913)(0.43478,0.60754)(0.47826,0.61471)(0.52174,0.62043)(0.56522,0.62443)(0.60870,0.62646)(0.65217,0.62624)(0.69565,0.62348)(0.73913,0.61786)(0.78261,0.60906)(0.82609,0.59674)(0.86957,0.58057)(0.91304,0.56019)(0.95652,0.53525)(1.00000,0.50543)(1.04348,0.47041)(1.08696,0.42990)(1.13043,0.38367)(1.17391,0.33153)(1.21739,0.27337)(1.26087,0.20916)(1.30435,0.13897)(1.34783,0.06298)(1.39130,-0.01853)(1.43478,-0.10516)(1.47826,-0.19639)(1.52174,-0.29161)(1.56522,-0.39012)(1.60870,-0.49115)(1.65217,-0.59391)(1.69565,-0.69757)(1.73913,-0.80135)(1.78261,-0.90450)(1.82609,-1.00632)(1.86957,-1.10622)(1.91304,-1.20369)(1.95652,-1.29835)(2.00000,-1.38990)(2.04348,-1.47815)(2.08696,-1.56300)(2.13043,-1.64445)(2.17391,-1.72256)(2.21739,-1.79743)(2.26087,-1.86923)(2.30435,-1.93815)(2.34783,-2.00440)(2.39130,-2.06820)(2.43478,-2.12977)(2.47826,-2.18933)(2.52174,-2.24711)(2.56522,-2.30328)(2.60870,-2.35805)(2.65217,-2.41159)(2.69565,-2.46404)(2.73913,-2.51555)(2.78261,-2.56624)(2.82609,-2.61622)(2.86957,-2.66559)(2.91304,-2.71442)(2.95652,-2.76280)(3.00000,-2.81079)};\addplot[blue]coordinates{(0.00000,1.00000)(0.01041,1.01052)(0.02082,1.02126)(0.03123,1.03222)(0.04164,1.04342)(0.05205,1.05485)(0.06245,1.06653)(0.07286,1.07846)(0.08327,1.09063)(0.09368,1.10308)(0.10409,1.11579)(0.11450,1.12877)(0.12491,1.14204)(0.13532,1.15560)(0.14573,1.16946)(0.15614,1.18363)(0.16654,1.19812)(0.17695,1.21294)(0.18736,1.22810)(0.19777,1.24362)(0.20818,1.25949)(0.21859,1.27574)(0.22900,1.29238)(0.23941,1.30943)(0.24982,1.32689)(0.26023,1.34479)(0.27063,1.36313)(0.28104,1.38195)(0.29145,1.40125)(0.30186,1.42106)(0.31227,1.44140)(0.32268,1.46229)(0.33309,1.48376)(0.34350,1.50583)(0.35391,1.52852)(0.36432,1.55187)(0.37472,1.57590)(0.38513,1.60066)(0.39554,1.62616)(0.40595,1.65246)(0.41636,1.67959)(0.42677,1.70760)(0.43718,1.73652)(0.44759,1.76641)(0.45800,1.79732)(0.46841,1.82931)(0.47881,1.86244)(0.48922,1.89678)(0.49963,1.93238)(0.51004,1.96934)(0.52045,2.00773)(0.53086,2.04765)(0.54127,2.08919)(0.55168,2.13245)(0.56209,2.17756)(0.57250,2.22463)(0.58291,2.27381)(0.59331,2.32525)(0.60372,2.37910)(0.61413,2.43556)(0.62454,2.49481)(0.63495,2.55709)(0.64536,2.62262)(0.65577,2.69170)(0.66618,2.76461)(0.67659,2.84169)(0.68700,2.92332)(0.69740,3.00992)(0.70781,3.10197)(0.71822,3.20000)};\addplot[red,thick]coordinates{(0.00000,0.67598)(0.04348,0.69642)(0.08696,0.71797)(0.13043,0.74057)(0.17391,0.76416)(0.21739,0.78870)(0.26087,0.81413)(0.30435,0.84040)(0.34783,0.86747)(0.39130,0.89530)(0.43478,0.92384)(0.47826,0.95307)(0.52174,0.98293)(0.56522,1.01340)(0.60870,1.04444)(0.65217,1.07603)(0.69565,1.10813)(0.73913,1.14071)(0.78261,1.17376)(0.82609,1.20725)(0.86957,1.24115)(0.91304,1.27544)(0.95652,1.31011)(1.00000,1.34513)(1.04348,1.38048)(1.08696,1.41616)(1.13043,1.45213)(1.17391,1.48840)(1.21739,1.52493)(1.26087,1.56173)(1.30435,1.59877)(1.34783,1.63605)(1.39130,1.67356)(1.43478,1.71127)(1.47826,1.74919)(1.52174,1.78731)(1.56522,1.82561)(1.60870,1.86408)(1.65217,1.90272)(1.69565,1.94153)(1.73913,1.98049)(1.78261,2.01959)(1.82609,2.05884)(1.86957,2.09822)(1.91304,2.13773)(1.95652,2.17736)(2.00000,2.21711)(2.04348,2.25697)(2.08696,2.29694)(2.13043,2.33702)(2.17391,2.37719)(2.21739,2.41746)(2.26087,2.45782)(2.30435,2.49827)(2.34783,2.53881)(2.39130,2.57943)(2.43478,2.62012)(2.47826,2.66089)(2.52174,2.70174)(2.56522,2.74265)(2.60870,2.78363)(2.65217,2.82468)(2.69565,2.86579)(2.73913,2.90696)(2.78261,2.94818)(2.82609,2.98947)(2.86957,3.03080)(2.91304,3.07219)(2.95652,3.11363)(3.00000,3.15512)};\addplot[dashed,domain=0:3]{x};\addplot[dashed,domain=0:3]{-x};\end{axis}\end{tikzpicture}\end{center}
\par\noindent{\small 依据：官方译审并补证（AI辅助）。}
\end{solution}

\begin{exercise}\label{cw:rec02-7}
分界线指一条解曲线：其上、下方的解在 $x$ 增加时有完全不同的命运。画出 $y'=y^2-x^2$ 的分界线，说明两侧的行为。
\end{exercise}
\begin{solution}
代换 $y=-u'/u$ 把方程化为 $u''=x^2u$。为明确分界线，取正函数
\[
 U(x)=\ee^{-x^2/2}\int_0^\infty t^{-1/2}\ee^{-t^2/2-\sqrt2xt}\dd t,\qquad x\ge0.
\]
积分可逐次求导，因为被积函数及其导数有高斯控制。分部积分 $\int_0^\infty (t^{1/2}\ee^{-t^2/2-\sqrt2xt})'\dd t=0$ 给出 $U''=x^2U$。因此 $s=-U'/U$ 是解，且 $s=x+\sqrt2\,\bigl(\int t^{1/2}\ee^{-t^2/2-\sqrt2xt}\dd t\bigr)/\bigl(\int t^{-1/2}\ee^{-t^2/2-\sqrt2xt}\dd t\bigr)>x$。
任意其他解可由降阶写成 $u=U(1+cJ)$，其中 $J(x)=\int_0^xU(v)^{-2}\dd v$。于是
\[
 y=s-\frac{cJ'}{1+cJ},\qquad y(0)-s(0)=-c/U(0)^2.
\]
$J$ 严格增加且无界。上方初值对应 $c<0$，$1+cJ$ 在有限时刻第一次为零，$y\to+\infty$；下方对应 $c>0$，分母始终正，解存在于全部 $x\ge0$，并趋向 $-x$，如下题所证。这正是分界线。示意图用红色表示分界线在右端靠近 $y=x$ 上方；须注意它并不是零斜线 $y=x$。
分界线的红色图在习题\ref{cw:rec02-6}中；在$x=0$，$s(0)=2\Gamma(3/4)/\Gamma(1/4)\approx0.67598$仅用作作图定位，证明不依赖此数值。
\par\noindent{\small 依据：官方译审并补证（AI辅助）；分界存在与两侧行为给出独立解析论证。}
\end{solution}

\begin{exercise}\label{cw:rec02-8}
方程 $y'=y^2-x^2$ 还存在漏斗，多条解在其中互相渐近。用两条等斜线解释，并找出这些解逼近的简单函数（该函数本身不是解）。
\end{exercise}
\begin{solution}
对于 $x>\sqrt{8/3}$，取下界 $L=-x$（等斜率0）和上界 $B=-\sqrt{x^2-2}$（等斜率$-2$）。边界上 $F(x,L)-L'=1>0$，$F(x,B)-B'=-2+x/\sqrt{x^2-2}<0$，故二者夹成不变漏斗，宽度 $x-\sqrt{x^2-2}\to0$，任何在里面的解满足 $y+x\to0$。
还可确认全部分界线以下的解最终进入此带：上一题的积分作 $t=v/(\sqrt2x)$，展开 $\ee^{-v^2/(4x^2)}$ 并以指数衰减控制余项，得
\[
 U=C\ee^{-x^2/2}x^{-1/2}\bigl(1-3/(16x^2)+O(x^{-4})\bigr),\quad
 J=\frac{\ee^{x^2}}{2C^2}\bigl(1+3/(8x^2)+O(x^{-4})\bigr).
\]
同样的展开可逐次求导，故 $s=x+1/(2x)+O(x^{-3})$，$J'/J=2x+O(x^{-3})$。当 $c>0$ 时 $y=-x+1/(2x)+O(x^{-3})$，最终在漏斗内。两解差 $d$ 满足 $d'=(y_1+y_2)d$，因此也趋于0。简单渐近函数是 $-x$；它不是解，因为其导数为$-1$而方程右端为0。官方末句的 $-x+1$ 是笔误，双曲线 $y^2=x^2-1$ 的下支渐近线同样为 $-x$。
\par\noindent{\small 依据：官方译审并补证（AI辅助）；修正官方渐近线并补不变带与渐近推导。}
\end{solution}

\originalsection{第3次习题课：Euler 法与线性模型（2月9日）}
\begin{exercise}\label{cw:rec03-1}
用 Euler 法估计 $y'=y^2-x^2,y(0)=-1$ 的 $y(1.5)$，步长 $h=0.5$。令 $x_{n+1}=x_n+h,y_{n+1}=y_n+m_nh,m_n=F(x_n,y_n)$，列出 $n,x_n,y_n,m_n,m_nh$ 并画折线。
\end{exercise}
\begin{solution}
逐步在旧节点取斜率：
\[\begin{array}{c|rrrr}n&x_n&y_n&m_n&m_nh\\\hline0&0&-1&1&1/2\\1&1/2&-1/2&0&0\\2&1&-1/2&-3/4&-3/8\\3&3/2&-7/8&-&-\end{array}\]
故估计为 $-7/8$。折线连接这四个点。\begin{center}\begin{tikzpicture}\begin{axis}[width=.68\textwidth,height=4.5cm,axis lines=middle,xlabel={$x$},ylabel={$y$},samples=130,xmin=0,xmax=1.6,ymin=-1.1,ymax=0]
\addplot[blue,mark=*] coordinates {(0,-1)(.5,-.5)(1,-.5)(1.5,-.875)};\end{axis}\end{tikzpicture}\end{center}
\par\noindent{\small 依据：官方译审并补证（AI辅助）。}
\end{solution}

\begin{exercise}\label{cw:rec03-2}
上一题的估计偏大还是偏小？
\end{exercise}
\begin{solution}
Euler 折线 $P$ 在三段分别为 $x-1,-1/2,-3x/4+1/4$。其斜率减 $F(x,P)$ 分别为 $2x,x^2-1/4$，以及 $-3/4+7x^2/16+3x/8-1/16$；在各区间均非负，第三式从 $x=1$ 的0递增。所以 $P$ 是上解。从 $d=P-y$ 的方程 $d'=(P+y)d+[P'-F(x,P)]$，积分因子直接给 $d\ge0$（首段内部严格正）。因而 $-7/8$ 偏大；这比仅检查三个点的凹凸性更可靠。
\par\noindent{\small 依据：官方译审并补证（AI辅助）。}
\end{solution}

\begin{exercise}\label{cw:rec03-3}
盐水槽体积为 $V$ 升，浓度 $c$ 克/升的盐水以 $r$ 升/分流入，同速流出，搅拌均匀。以 $x(t)$ 表示槽中盐的克数，写出方程并检查单位、化为标准线性形式。
\end{exercise}
\begin{solution}
流入盐量每分钟 $rc$ 克，流出浓度为 $x/V$，故 $\dot x=rc-rx/V$，即 $\dot x+(r/V)x=rc$。两端单位均为克/分；$V$ 恒定是因为进出体积速率相同。
\par\noindent{\small 依据：官方译审并补证（AI辅助）。}
\end{solution}

\begin{exercise}\label{cw:rec03-4}
设 $c,r$ 恒定，$V=1$ 升、$r=2$ 升/分，$x(0)=0$。求解、长期盐量及达到其一半的时间，并说明物理意义。
\end{exercise}
\begin{solution}
$\dot x+2x=2c$，积分因子 $\ee^{2t}$ 给 $x=c(1-\ee^{-2t})$ 克，其中体积1升已代入。极限为 $c$ 克，相应浓度与输入相同。若 $c>0$，半量时 $\ee^{-2t}=1/2$，即 $t=\log2/2$ 分。
\par\noindent{\small 依据：官方译审并补证（AI辅助）。}
\end{solution}

\begin{exercise}\label{cw:rec03-5}
第一槽的流出液进入另一体积1升且同速流出的均匀搅拌槽；第二槽在 $t=1$ 时无盐。建立第二槽盐量随时间的微分方程。
\end{exercise}
\begin{solution}
以 $z(t)$ 表示第二槽盐量，输入浓度为 $x(t)/1$，所以 $\dot z+2z=2c(1-\ee^{-2t})$，$z(1)=0$，$t\ge1$。保留原题初始时刻1；若进一步积分，$z=c-2ct\ee^{-2t}+c(2-\ee^2)\ee^{-2t}$，代入 $t=1$ 确为0。
\par\noindent{\small 依据：官方译审并补证（AI辅助）。}
\end{solution}

\begin{exercise}\label{cw:rec03-6}
画出电压源、电阻、电容串联电路，取顺时针电流 $I(t)$，源电压 $V(t)$，电阻 $R>0$、电容 $C>0$，满足 $R\dot I+C^{-1}I=\dot V$。当 $V(t)=V_0$ 恒定时，按给定 $I(0)$ 求解，并写作 $c\ee^{-t/\tau}$，求 $c,\tau,I(t+\tau)$。
\end{exercise}
\begin{solution}
$\dot V=0$，故 $\dot I=-I/(RC)$，积分得 $I(t)=I(0)\ee^{-t/(RC)}$。因此 $c=I(0),\tau=RC$（时间单位），且 $I(t+\tau)=I(t)/\ee$。以下示意图中框为电阻、双短线为电容、圆形为电压源。\begin{center}\begin{tikzpicture}[scale=.8]\draw(0,0)--(0,1.3) (0,2.1)--(0,3)--(1.3,3) (2.3,3)--(4,3)--(4,1.9) (4,1.6)--(4,0)--(0,0);\draw (0,1.7) circle(.4);\node at(0,1.7){$V$};\draw(1.3,2.8) rectangle(2.3,3.2);\node at(1.8,3.5){电阻 $R$};\draw(3.6,1.9)--(4.4,1.9) (3.6,1.6)--(4.4,1.6);\node[right]at(4.4,1.75){电容 $C$};\draw[->](2.4,.6)--(1.4,.6) node[midway,above]{$I$};\end{tikzpicture}\end{center}
\par\noindent{\small 依据：官方译审并补证（AI辅助）。}
\end{solution}

\originalsection{第4次习题课：一阶线性方程与积分因子（2月11日）}
\begin{exercise}\label{cw:rec04-1}
Cape Cod 盐池经窄水道连海，池面在短时间的变化与海池水位差及时间长度成正比。同一零点测海面 $y(t)$、池面 $x(t)$。
\begin{enumerate}[label=(\alph*)]\item 建立方程；说明系统、代表系统的算子、输入和输出。\item 高潮恰每 $4\pi$ 小时一次，$y=\cos(\omega t)$，长度用米、时间用小时，求 $\omega$。\item 用积分因子求解，可用 $\int\ee^{kt}\cos\omega t\dd t=\ee^{kt}(k\cos\omega t+\omega\sin\omega t)/(k^2+\omega^2)+C$。\item 另设 $x_p=a\cos\omega t+b\sin\omega t$ 求同一周期解。\end{enumerate}
\end{exercise}
\begin{solution}
(a) $\dot x=k(y-x)$，$k>0$，故 $(D+k)x=ky$。系统是盐池与水道，固定算子为 $D+k$，海面 $y$ 为输入，池面 $x$ 为输出，输入系数为$k$。
(b) $2\pi/\omega=4\pi$，故 $\omega=1/2$ 小时$^{-1}$。
(c) 乘 $\ee^{kt}$ 并积分得
\[x=\frac{k^2\cos\omega t+k\omega\sin\omega t}{k^2+\omega^2}+C\ee^{-kt}.
\]
(d) 比较余弦、正弦系数得到 $ka+\omega b=k$、$kb-\omega a=0$，解得 $a=k^2/(k^2+\omega^2),b=k\omega/(k^2+\omega^2)$，与(c)一致。$k=0$ 时池面恒定，不能从周期强迫确定其常值。
\par\noindent{\small 依据：官方译审并补证（AI辅助）。}
\end{solution}

\begin{exercise}\label{cw:rec04-2}
用解析方法求 $y'=x-2y$ 通解；核对第2次的直线解，画几条解，判断漏斗，并求 $y(0)=1$ 的解。
\end{exercise}
\begin{solution}
乘 $\ee^{2x}$ 得 $(\ee^{2x}y)'=x\ee^{2x}$；积分得 $y=x/2-1/4+C\ee^{-2x}$。$C=0$ 给直线解，任意窄的两条平行线最终夹住所有有限初值解。$y(0)=1$ 给 $C=5/4$，故 $y=x/2-1/4+5\ee^{-2x}/4$。图为同一方程，沿用习题\ref{cw:rec02-1} 的同册图；漏斗的边界方向论证见习题\ref{cw:rec02-4}。
\par\noindent{\small 依据：官方译审并补证（AI辅助）；同一方程方向图回指第2次，积分推导独立列出。}
\end{solution}

\begin{exercise}\label{cw:rec04-3}
把左侧识别为乘积导数，求 $x^2y'+2xy=\sin(2x)$ 通解。
\end{exercise}
\begin{solution}
左侧为 $(x^2y)'$，故 $x^2y=-\cos(2x)/2+C$。在不含0的区间 $y=[C-\cos(2x)/2]/x^2$。若要连续延至0，必须 $C=1/2$；这时 $y=\sin^2x/x^2$，可取 $y(0)=1$，且光滑延拓在0亦满足原式。任意常数公式不能直接在0相除。
\par\noindent{\small 依据：官方译审并补证（AI辅助）。}
\end{solution}

\originalsection{第5次习题课：复数、复指数与相位（2月18日）}
\begin{exercise}\label{cw:rec05-1}
在复平面标出 $z=1+\sqrt3\ii$，求极坐标；再标 $z^n,n=1,2,3,4$，分别写为 $a+b\ii$ 与 $A\ee^{\ii\theta}$，并标 $n=0,-1,-2,-3,-4$ 的点。
\end{exercise}
\begin{solution}
$z=2\ee^{\ii\pi/3}$，所以 $z^n=2^n\ee^{\ii n\pi/3}$。直角坐标依次为
\[\begin{array}{c|ccccc}n&0&1&2&3&4\\\hline z^n&1&1+\sqrt3\ii&-2+2\sqrt3\ii&-8&-8-8\sqrt3\ii\end{array}\]
\[z^{-1}=(1-\sqrt3\ii)/4,\quad z^{-2}=(-1-\sqrt3\ii)/8,\quad z^{-3}=-1/8,\quad z^{-4}=(-1+\sqrt3\ii)/32.\]
特别是 $z^{-4}$ 在上半平面，不能把负幂一概说成在下半平面。下图以两种尺度分别标正幂和非正幂。\begin{center}\begin{tikzpicture}\begin{axis}[width=.46\textwidth,height=4.5cm,axis lines=middle,xlabel={$\Re z$},ylabel={$\Im z$},xmin=-10,xmax=3,ymin=-15,ymax=5]\addplot[only marks,blue]coordinates{(1,1.732)(-2,3.464)(-8,0)(-8,-13.856)};\node at(axis cs:1,1.732)[above]{$z$};\node at(axis cs:-2,3.464)[above]{$z^2$};\node at(axis cs:-8,0)[above]{$z^3$};\node at(axis cs:-8,-13.856)[above]{$z^4$};\end{axis}\end{tikzpicture}\begin{tikzpicture}\begin{axis}[width=.46\textwidth,height=4.5cm,axis lines=middle,xlabel={$\Re z$},ylabel={$\Im z$},xmin=-.45,xmax=1.2,ymin=-.55,ymax=.25]\addplot[only marks,blue]coordinates{(1,0)(.25,-.433)(-.125,-.2165)(-.125,0)(-.03125,.0541)};\node at(axis cs:1,0)[above]{$z^0$};\node at(axis cs:.25,-.433)[right]{$z^{-1}$};\node at(axis cs:-.125,-.2165)[below]{$z^{-2}$};\node at(axis cs:-.18,.15)[above left]{$z^{-3}$};\draw[gray](axis cs:-.18,.13)--(axis cs:-.125,0);\node at(axis cs:.12,.15)[right]{$z^{-4}$};\draw[gray](axis cs:.12,.13)--(axis cs:-.03125,.0541);\end{axis}\end{tikzpicture}\end{center}
\par\noindent{\small 依据：官方译审并补证（AI辅助）。}
\end{solution}

\begin{exercise}\label{cw:rec05-2}
找出所有满足 $\ee^{a+b\ii}=1+\sqrt3\ii$ 的复数；取最小正 $b$ 后求 $\ee^{n(a+b\ii)}$，$n=1,2,3,4,0,-1,-2,-3,-4$。
\end{exercise}
\begin{solution}
比较模与辐角得 $a=\log2,b=\pi/3+2j\pi,j\in\mathbb Z$。最小正 $b$ 是 $\pi/3$。所有整数 $n$ 均有 $\ee^{n(a+b\ii)}=(\ee^{a+b\ii})^n=z^n$，所以各值、极形式及点位正是习题\ref{cw:rec05-1} 完整列出的九个值。
\par\noindent{\small 依据：官方译审并补证（AI辅助）；九个完全相同的数值回指上题，复对数分支另行推导。}
\end{solution}

\begin{exercise}\label{cw:rec05-3}
各式化为 $A\cos(\omega t-\phi)$，先画以余弦、正弦系数 $a,b$ 为直角边的三角形：
(a) $\cos2t+\sin2t$；(b) $\cos\pi t-\sqrt3\sin\pi t$；(c) $\Re[\ee^{\ii t}/(2+2\ii)]$。
\end{exercise}
\begin{solution}
展开 $A\cos(\omega t-\phi)=A\cos\phi\cos\omega t+A\sin\phi\sin\omega t$，故 $A=\sqrt{a^2+b^2}$，$\phi=\arg(a+b\ii)$。
(a) $(a,b)=(1,1)$，得 $\sqrt2\cos(2t-\pi/4)$；(b) $(1,-\sqrt3)$，得 $2\cos(\pi t+\pi/3)$；(c) $(2+2\ii)^{-1}=(1-\ii)/4$，实部为 $(\cos t+\sin t)/4$，得 $\sqrt2\cos(t-\pi/4)/4$。
三角形中的有向竖边决定相位正负，长度仍取绝对值。\begin{center}\begin{tikzpicture}[scale=.7]\foreach \p/\a/\b/\lab in {0/1/1/(a),4/1/-1.732/(b),8/1/1/(c)}{\begin{scope}[xshift=\p cm]\draw[->](-.3,0)--(2,0);\draw(0,0)--(1.4,0)--(1.4,\b)--cycle;\node[below]at(.7,0){$a$};\node[right]at(1.4,\b/2){$b$};\node[left]at(.7,\b/2){$A$};\node at(.7,2){\lab};\end{scope}}\end{tikzpicture}\end{center}
\par\noindent{\small 依据：官方译审并补证（AI辅助）。}
\end{solution}

\begin{exercise}\label{cw:rec05-4}
以 $w\ee^t$ 形式求 $\dot x+2x=\ee^t$ 的特解；同样求 $\dot z+2z=\ee^{2\ii t}$，并给两个方程的通解。
\end{exercise}
\begin{solution}
第一式代入得 $3w=1$，故 $x=\ee^t/3+C\ee^{-2t}$。第二式设 $z=w\ee^{2\ii t}$ 得 $(2+2\ii)w=1$，故 $z=\ee^{2\ii t}/(2+2\ii)+C\ee^{-2t}$；此处 $C$ 可以是复数。每个齐次差满足 $\dot h+2h=0$，所以没有遗漏。
\par\noindent{\small 依据：官方译审并补证（AI辅助）。}
\end{solution}

\begin{exercise}\label{cw:rec05-5}
用复替代求 $\dot x+2x=\cos2t$ 的解，再由同一计算得到正弦输入的解。
\end{exercise}
\begin{solution}
求解 $\dot z+2z=\ee^{2\ii t}$ 得 $z_p=(1-\ii)\ee^{2\ii t}/4$。由于算子实系数，实部、虚部分别满足余弦、正弦强迫。因此 $x_{p,\cos}=(\cos2t+\sin2t)/4$，$x_{p,\sin}=(\sin2t-\cos2t)/4$；两者分别加 $C\ee^{-2t}$ 得通解。
\par\noindent{\small 依据：官方译审并补证（AI辅助）。}
\end{solution}

\originalchapter{习题课：二阶方程、共振与 Fourier 响应}
\originalsection{第7次习题课：二阶方程的解（2月25日）}
\begin{exercise}\label{cw:rec07-1}
核验 $x=\cos\omega t$ 与 $x=\sin\omega t$ 都满足谐振子方程 $\ddot x+\omega^2x=0$。
\end{exercise}
\begin{solution}
余弦的二阶导数为 $-\omega^2\cos\omega t$，正弦的二阶导数为 $-\omega^2\sin\omega t$，代入各自抵消。
\par\noindent{\small 依据：官方译审并补证（AI辅助）。}
\end{solution}

\begin{exercise}\label{cw:rec07-2}
进一步核验一般正弦波 $x=A\cos(\omega t-\phi)$ 满足同一方程。
\end{exercise}
\begin{solution}
两次链式求导得 $\ddot x=-A\omega^2\cos(\omega t-\phi)=-\omega^2x$。振幅、相位任意不改变结论。
\par\noindent{\small 依据：官方译审并补证（AI辅助）。}
\end{solution}

\begin{exercise}\label{cw:rec07-3}
哪些 $x=A\cos(\omega t-\phi)$ 满足 $x(0)=0$？这是否违反微分方程的唯一性定理？
\end{exercise}
\begin{solution}
$A\cos\phi=0$。除零解外，$\phi=\pi/2+j\pi$，相应解为任意倍数的 $\sin\omega t$。二阶方程的唯一性需要同时指定位置和速度；只给位置不能唯一决定解。化为 $(x,v)'=(v,-\omega^2x)$ 后，一阶系统的两个初值也表明此点。
\par\noindent{\small 依据：官方译审并补证（AI辅助）。}
\end{solution}

\begin{exercise}\label{cw:rec07-4}
给定 $x_0,\dot x_0$，求 $\ddot x+\omega^2x=0$ 满足 $x(0)=x_0,\dot x(0)=\dot x_0$ 的解，有多少个？
\end{exercise}
\begin{solution}
$\omega\ne0$ 时 $x=x_0\cos\omega t+(\dot x_0/\omega)\sin\omega t$，其两个初值直接成立。差解的能量 $v^2+\omega^2x^2$ 恒定，零初值使能量恒为0，故解唯一。$\omega=0$ 时唯一解为 $x=x_0+\dot x_0t$。
\par\noindent{\small 依据：官方译审并补证（AI辅助）。}
\end{solution}

\begin{exercise}\label{cw:rec07-5}
若复数常数 $r$ 使 $\ee^{rt}$ 满足 $\ddot x+kx=0$，$r$ 必须为何值？
\end{exercise}
\begin{solution}
代入得 $(r^2+k)\ee^{rt}=0$，指数不为0，故 $r^2=-k$。$k>0$ 时 $r=\pm\ii\sqrt k$，$k=0$ 时 $r=0$，$k<0$ 时 $r=\pm\sqrt{-k}$。
\par\noindent{\small 依据：官方译审并补证（AI辅助）。}
\end{solution}

\begin{exercise}\label{cw:rec07-6}
对 $\ddot x-a^2x=0$ 求规范化解 $x_1(0)=1,\dot x_1(0)=0$ 及 $x_2(0)=0,\dot x_2(0)=1$。
\end{exercise}
\begin{solution}
$a\ne0$ 时从 $x=c_1\ee^{at}+c_2\ee^{-at}$ 比较初值得 $x_1=\cosh at$，$x_2=\sinh(at)/a$。二阶导数与初值逐一满足要求。$a=0$ 时分别为1、$t$，也是这两个公式的连续极限。
\par\noindent{\small 依据：官方译审并补证（AI辅助）。}
\end{solution}

\originalsection{第8次习题课：二阶常系数齐次方程（3月2日）}
\begin{exercise}\label{cw:rec08-1}
对 $\ddot x+\omega^2x=0$ 求特征多项式、根、指数解及其实部与虚部。
\end{exercise}
\begin{solution}
$p(s)=s^2+\omega^2$，根为 $\pm\ii\omega$，指数解为 $\ee^{\pm\ii\omega t}$，实部均为 $\cos\omega t$，虚部为 $\pm\sin\omega t$。$\omega\ne0$ 时后二者构成实基本解；$\omega=0$ 时须补重根解$t$。
\par\noindent{\small 依据：官方译审并补证（AI辅助）。}
\end{solution}

\begin{exercise}\label{cw:rec08-2}
已知 $\ee^{-t/2}\cos3t$ 是 $m\ddot x+b\dot x+kx=0$ 的解，$m,b,k$ 为实数、$m\ne0$。
(a) 确定它们的关系；(b) 写一个指数解；(c) 画指数解在复平面的轨迹及实部时间图，说明关系；(d) 给通解。
\end{exercise}
\begin{solution}
(a) 实部求导并比较独立的正余弦项，等价于特征根 $-1/2\pm3\ii$；因此 $p(s)=m[(s+1/2)^2+9]$，$b=m,k=37m/4$。
(b) $z=\ee^{(-1/2+3\ii)t}$。
(c) $|z|=\ee^{-t/2},\arg z=3t$，故逆时针向内螺旋；实部是这个轨迹在实轴上的投影，并按时间展开。
(d) $x=\ee^{-t/2}(C_1\cos3t+C_2\sin3t)$，两解初值矩阵行列式为3，保证独立。\begin{center}\begin{tikzpicture}\begin{axis}[width=.44\textwidth,height=4.2cm,axis lines=middle,xlabel={实部},ylabel={虚部},xmin=-1,xmax=1,ymin=-1,ymax=1]\addplot[blue,domain=0:10,samples=230]({exp(-x/2)*cos(deg(3*x))},{exp(-x/2)*sin(deg(3*x))});\end{axis}\end{tikzpicture}\begin{tikzpicture}\begin{axis}[width=.48\textwidth,height=4.2cm,axis lines=middle,xlabel={$t$},ylabel={$\Re z$},xmin=0,xmax=10,ymin=-1,ymax=1]\addplot[blue,domain=0:10,samples=230]{exp(-x/2)*cos(deg(3*x))};\addplot[dashed,domain=0:10]{exp(-x/2)};\addplot[dashed,domain=0:10]{-exp(-x/2)};\end{axis}\end{tikzpicture}\end{center}
\par\noindent{\small 依据：官方译审并补证（AI辅助）。}
\end{solution}

\begin{exercise}\label{cw:rec08-3}
阻尼正弦 $x=A\ee^{-at}\cos\omega t$ 的伪周期为 $2\pi/\omega$，相邻零点距离为其一半。相邻极大值点间距是否固定？是多少？
\end{exercise}
\begin{solution}
取 $A\ne0,\omega>0$。$\dot x=-A\ee^{-at}(a\cos\omega t+\omega\sin\omega t)$，驻点满足 $\tan\omega t=-a/\omega$，间距 $\pi/\omega$。驻点处 $\ddot x=-(a^2+\omega^2)x$，正峰为极大、负峰为极小，交替出现，因此相邻极大值相隔 $2\pi/\omega$，恒定。零函数或 $\omega=0$ 没有这类相邻波峰。
\par\noindent{\small 依据：官方译审并补证（AI辅助）。}
\end{solution}

\begin{exercise}\label{cw:rec08-4}
相邻极大值时刻为 $t_0,t_1$，求 $x(t_1)/x(t_0)$；能否由图估计 $a$？
\end{exercise}
\begin{solution}
上一题给 $\omega(t_1-t_0)=2\pi$，余弦因子相同，故比值为 $\ee^{-a(t_1-t_0)}$。于是 $a=\log[x(t_0)/x(t_1)]/(t_1-t_0)$，两个正峰的高度与间距即可确定衰减率。
\par\noindent{\small 依据：官方译审并补证（AI辅助）。}
\end{solution}

\begin{exercise}\label{cw:rec08-5}
何值的 $b$ 使 $\ddot x+b\dot x+x=0$ 临界阻尼？求规范化解 $x_1(0)=1,\dot x_1(0)=0$ 与 $x_2(0)=0,\dot x_2(0)=1$；再求 $x(0)=2,\dot x(0)=3$ 的解。
\end{exercise}
\begin{solution}
判别式 $b^2-4=0$ 给 $b=\pm2$；阻尼要求 $b>0$，故取2。通解 $(c_0+c_1t)\ee^{-t}$，初值为 $(c_0,c_1-c_0)$。因此 $x_1=(1+t)\ee^{-t}$、$x_2=t\ee^{-t}$，最后一解 $2x_1+3x_2=(2+5t)\ee^{-t}$。$b=-2$ 为重正根导致增长，不是临界阻尼。
\par\noindent{\small 依据：官方译审并补证（AI辅助）。}
\end{solution}

\originalsection{第9次习题课：指数与正弦输入（3月4日）}
\begin{exercise}\label{cw:rec09-1}
求 $A$ 使 $A\sin3t$ 满足 $\ddot x+4x=\sin3t$，并求通解。
\end{exercise}
\begin{solution}
代入左侧为 $(4-9)A\sin3t$，所以 $A=-1/5$。齐次根为 $\pm2\ii$，故 $x=-\sin3t/5+C_1\cos2t+C_2\sin2t$。
\par\noindent{\small 依据：官方译审并补证（AI辅助）。}
\end{solution}

\begin{exercise}\label{cw:rec09-2}
对 $\omega\ge0$ 求 $A$ 使 $A\cos\omega t$ 满足 $\ddot x+4x=\cos\omega t$；分别画 $\omega=0,1,3$ 的输入与响应，再画 $0\le\omega\le5$ 时 $A(\omega)$，说明共振点及正负号的意义。
\end{exercise}
\begin{solution}
$A(4-\omega^2)=1$，所以 $\omega\ne2$ 时 $A=1/(4-\omega^2)$。$\omega=0,1,3$ 时分别为 $1/4,1/3,-1/5$；正号为同相，负号为反相。$\omega=2$ 无此形式特解而发生共振，不能将无限值代作解。下图黑线为输入，蓝线为响应。\begin{center}\begin{tikzpicture}\begin{axis}[width=.3\textwidth,height=3.3cm,axis lines=middle,title={$\omega=0$},xmin=0,xmax=7,ymin=-1.2,ymax=1.2]\addplot[black,domain=0:7]{1};\addplot[blue,domain=0:7]{.25};\end{axis}\end{tikzpicture}\begin{tikzpicture}\begin{axis}[width=.3\textwidth,height=3.3cm,axis lines=middle,title={$\omega=1$},xmin=0,xmax=7,ymin=-1.2,ymax=1.2]\addplot[black,domain=0:7]{cos(deg(x))};\addplot[blue,domain=0:7]{cos(deg(x))/3};\end{axis}\end{tikzpicture}\begin{tikzpicture}\begin{axis}[width=.3\textwidth,height=3.3cm,axis lines=middle,title={$\omega=3$},xmin=0,xmax=7,ymin=-1.2,ymax=1.2]\addplot[black,domain=0:7,samples=150]{cos(deg(3*x))};\addplot[blue,domain=0:7,samples=150]{-cos(deg(3*x))/5};\end{axis}\end{tikzpicture}\end{center}\begin{center}\begin{tikzpicture}\begin{axis}[width=.68\textwidth,height=4.5cm,axis lines=middle,xlabel={$\omega$},ylabel={$A$},samples=130,xmin=0,xmax=5,ymin=-4,ymax=4]
\addplot[blue,domain=0:1.94]{1/(4-x*x)};\addplot[blue,domain=2.06:5]{1/(4-x*x)};\draw[dashed] (axis cs:2,-4)--(axis cs:2,4);\end{axis}\end{tikzpicture}\end{center}
\par\noindent{\small 依据：官方译审并补证（AI辅助）。}
\end{solution}

\begin{exercise}\label{cw:rec09-3}
求 $x^{(4)}-x=\ee^{-2t}$ 的指数型特解。
\end{exercise}
\begin{solution}
设 $x_p=B\ee^{-2t}$，四阶导数为 $16B\ee^{-2t}$，左端为 $15B\ee^{-2t}$，故 $x_p=\ee^{-2t}/15$。
\par\noindent{\small 依据：官方译审并补证（AI辅助）。}
\end{solution}

\begin{exercise}\label{cw:rec09-4}
求 $x^{(4)}-x=\cos2t$ 的正弦型特解。
\end{exercise}
\begin{solution}
余弦的四阶导数为 $16\cos2t$，因此 $x_p=\cos2t/15$ 满足方程。等价地 $p(2\ii)=(2\ii)^4-1=15$，取复指数响应的实部。
\par\noindent{\small 依据：官方译审并补证（AI辅助）。}
\end{solution}

\begin{exercise}\label{cw:rec09-5}
求第3、4题两个四阶非齐次方程的通解。
\end{exercise}
\begin{solution}
$p(s)=s^4-1=(s-1)(s+1)(s^2+1)$，四个互异根给 $x_h=C_1\ee^t+C_2\ee^{-t}+C_3\cos t+C_4\sin t$。分别加第3题 $\ee^{-2t}/15$ 或第4题 $\cos2t/15$ 就是所需通解。相同齐次方程的完整四维解空间在此只写一次。
\par\noindent{\small 依据：官方译审并补证（AI辅助）。}
\end{solution}

\originalsection{第10次习题课：增益、相位滞后与待定系数（3月9日）}
\begin{exercise}\label{cw:rec10-1}
解释指数响应公式（ERF）的记号：$p(D)x=A\ee^{rt}$ 在 $p(r)\ne0$ 时有 $x_p=A\ee^{rt}/p(r)$；若 $p(r)=0,p'(r)\ne0$，则有 $x_p=At\ee^{rt}/p'(r)$。
\end{exercise}
\begin{solution}
$D=\dd/\dd t$ 是微分算子，$p(s)=\sum_{j=0}^na_js^j$ 是常系数算子的多项式，$p(D)=\sum a_jD^j$；$r$ 是可取复值的输入指数，$A$ 是输入系数，$x_p$ 表示特解，$p'(r)$ 是多项式对其变量求导后在$r$的值。因 $D^j\ee^{rt}=r^j\ee^{rt}$，有 $p(D)\ee^{rt}=p(r)\ee^{rt}$，立即得到非共振公式。
对 $r$ 求导（有限和允许交换）得 $p(D)(t\ee^{rt})=p'(r)\ee^{rt}+p(r)t\ee^{rt}$，故在单根共振时得到第二式。多重根需要更高次$t$因子，不能继续除以零。实系数算子可取实部、虚部得到实正弦响应。
\par\noindent{\small 依据：编者补解（AI辅助）；配套官方解答缺第1题。}
\end{solution}

\begin{exercise}\label{cw:rec10-2}
说明由弹簧端驱动的质量--弹簧--阻尼器为何满足 $m\ddot x+b\dot x+kx=ky$。$x$ 为质量块位置，$y$ 为弹簧另一端的位置标记，且 $x=y$ 时弹簧自然放松。
\end{exercise}
\begin{solution}
原题以自然长度调整位置基准，使伸长为 $x-y$。弹簧力 $-k(x-y)$，固定阻尼器力 $-b\dot x$；Newton 方程为 $m\ddot x=-b\dot x-k(x-y)$，移项得到原式。关键是阻尼器相对固定支座，若阻尼器另一端也移动则其力需改写。
\par\noindent{\small 依据：官方译审并补证（AI辅助）。}
\end{solution}

\begin{exercise}\label{cw:rec10-3}
取 $m=1,b=3,k=4,y=A\cos t$，把输入复替代，求稳态指数响应 $z_p=Hy_{\rm cx}$、复增益 $H=|H|\ee^{-\ii\phi}$、增益与相位滞后、稳态实解；质量振幅比驱动端振幅 $A$ 大还是小？
\end{exercise}
\begin{solution}
$y_{\rm cx}=A\ee^{\ii t}$，$p(s)=s^2+3s+4$，$p(\ii)=3+3\ii$。故 $z_p=4A\ee^{\ii t}/(3+3\ii)$，$H=(2/3)(1-\ii)=(2\sqrt2/3)\ee^{-\ii\pi/4}$。增益 $2\sqrt2/3<1$，相位滞后 $\pi/4$，$x_p=(2\sqrt2A/3)\cos(t-\pi/4)$，振幅小于 $A$（按 $A\ge0$ 的振幅约定）。齐次根实部为$-3/2$，所以其他响应的瞬态衰减。
\par\noindent{\small 依据：官方译审并补证（AI辅助）。}
\end{solution}

\begin{exercise}\label{cw:rec10-4}
求 $\ddot x-x=t^2+t+1$ 的多项式解。
\end{exercise}
\begin{solution}
设 $x=at^2+bt+c$，左侧 $-at^2-bt+2a-c$。比较得 $a=-1,b=-1,c=-3$，所以 $x=-t^2-t-3$。若多项式次数更高，最高项在$-x$中无法被二阶导数抵消，故没有另一个多项式解。
\par\noindent{\small 依据：官方译审并补证（AI辅助）。}
\end{solution}

\begin{exercise}\label{cw:rec10-5}
用复替代并尝试 ERF 求 $\ddot x+4x=\cos2t$ 的一个解。
\end{exercise}
\begin{solution}
$p(s)=s^2+4$，$p(2\ii)=0,p'(2\ii)=4\ii$，使用单根共振公式：$z_p=t\ee^{2\ii t}/(4\ii)$。取实部得 $x_p=t\sin2t/4$。直接求导给 $\ddot x_p+4x_p=\cos2t$，体现振幅线性增长。
\par\noindent{\small 依据：官方译审并补证（AI辅助）。}
\end{solution}

\originalsection{第11次习题课：频率响应（3月11日）}
本次系统为直接外力驱动：$\ddot x+\frac14\dot x+2x=F_{\rm ext}$，输入 $F_{\rm ext}=\cos\omega t$（振幅1），稳态响应 $x_p=g(\omega)\cos(\omega t-\phi)$。原题建议用 Amplitude and Phase: Second order, IV；以下均为解析计算。
\begin{exercise}\label{cw:rec11-1}
以 $F_{\rm cx}=\ee^{\ii\omega t}$ 求复增益 $H(\omega)$，使 $z_p=H(\omega)F_{\rm cx}$。
\end{exercise}
\begin{solution}
$p(\ii\omega)=2-\omega^2+\ii\omega/4$，故 $H=(2-\omega^2+\ii\omega/4)^{-1}$。乘共轭可写为 $(2-\omega^2-\ii\omega/4)/[(2-\omega^2)^2+\omega^2/16]$。$\omega\ge0$ 时分母不为0。
\par\noindent{\small 依据：官方译审并补证（AI辅助）。}
\end{solution}

\begin{exercise}\label{cw:rec11-2}
求 $g(\omega)=|H(\omega)|$，以及 $\omega=1$ 的振幅。
\end{exercise}
\begin{solution}
$g=[(2-\omega^2)^2+\omega^2/16]^{-1/2}=[\omega^4-63\omega^2/16+4]^{-1/2}$。所以 $g(1)=4/\sqrt{17}$。因输入振幅为1，响应振幅正是这个数。
\par\noindent{\small 依据：官方译审并补证（AI辅助）。}
\end{solution}

\begin{exercise}\label{cw:rec11-3}
求使振幅最大的共振圆频率 $\omega_r$；可最小化分母的平方。
\end{exercise}
\begin{solution}
令 $z=\omega^2\ge0$，$g^{-2}=z^2-63z/16+4$ 是开口向上抛物线，唯一最小点为 $z=63/32$。故 $\omega_r=\sqrt{63/32}=3\sqrt{14}/8$；这种全局最小论证同时排除端点$\omega=0$。
\par\noindent{\small 依据：官方译审并补证（AI辅助）。}
\end{solution}

\begin{exercise}\label{cw:rec11-4}
在共振频率相位似乎约为 $\pi/2$，它是否恰好等于 $\pi/2$？究竟哪个频率使滞后为四分之一周期？
\end{exercise}
\begin{solution}
相位 $\phi=\arg(2-\omega^2+\ii\omega/4)$ 取 $[0,\pi]$。恰为 $\pi/2$ 要求实部0且虚部正，故 $\omega=\sqrt2$。共振点实部为 $2-63/32=1/32>0$，所以相位略小于$\pi/2$；其正切为 $3\sqrt{14}$，约为$84.91^\circ$。
\par\noindent{\small 依据：官方译审并补证（AI辅助）；补官方未明确给出的精确频率。}
\end{solution}

\begin{exercise}\label{cw:rec11-5}
哪个圆频率使相位滞后为 $\pi/4$？哪个使其为 $3\pi/4$？
\end{exercise}
\begin{solution}
第一种要求 $2-\omega^2=\omega/4>0$，即 $4\omega^2+\omega-8=0$，正根 $\omega=(\sqrt{129}-1)/8$。第二种要求 $2-\omega^2=-\omega/4<0$，正根 $(\sqrt{129}+1)/8$。用象限约束消除仅取正切导致的相位歧义。
\par\noindent{\small 依据：官方译审并补证（AI辅助）。}
\end{solution}

\begin{exercise}\label{cw:rec11-6}
另求 $\ddot x+3\dot x+2x=t\ee^{-t}$ 的一个解。
\end{exercise}
\begin{solution}
令 $x=u\ee^{-t}$，算子平移得左端 $\ee^{-t}(\ddot u+\dot u)$，故需 $\ddot u+\dot u=t$。取 $\dot u=t-1$，积分取 $u=t^2/2-t$，得到 $x_p=(t^2/2-t)\ee^{-t}$。
\par\noindent{\small 依据：官方译审并补证（AI辅助）。}
\end{solution}

\originalsection{第13次习题课：Fourier 级数入门与周期（3月18日）}
\begin{exercise}\label{cw:rec13-1}
迅速写出 $\ddot x+\omega_n^2x=0$ 的通解。
\end{exercise}
\begin{solution}
$\omega_n>0$ 时，特征根为 $\pm\ii\omega_n$，通解 $C_1\cos\omega_nt+C_2\sin\omega_nt$；若允许 $\omega_n=0$，通解为 $C_1+C_2t$。
\par\noindent{\small 依据：官方译审并补证（AI辅助）。}
\end{solution}

\begin{exercise}\label{cw:rec13-2}
解释为什么在 $\omega\ne\pm\omega_n$ 时，$\ddot x+\omega_n^2x=a\cos\omega t$ 和 $\ddot x+\omega_n^2x=b\sin\omega t$ 分别有特解 $a\cos\omega t/(\omega_n^2-\omega^2)$ 与 $b\sin\omega t/(\omega_n^2-\omega^2)$。
\end{exercise}
\begin{solution}
两次微分对正余弦都乘 $-\omega^2$，加 $\omega_n^2$ 后乘 $\omega_n^2-\omega^2$，故除以此非零因子恰还原输入。亦可用 $p(\ii\omega)=\omega_n^2-\omega^2$ 的指数响应公式。
\par\noindent{\small 依据：官方译审并补证（AI辅助）。}
\end{solution}

\begin{exercise}\label{cw:rec13-3}
对 $\ddot x+\omega_n^2x=\cos\omega_nt$ 求特解、通解。是否存在对所有 $t$ 有 $|x(t)|<10^6$ 的解？是否存在周期解？
\end{exercise}
\begin{solution}
$\omega_n>0$ 时共振给 $x_p=t\sin\omega_nt/(2\omega_n)$，通解再加 $C_1\cos\omega_nt+C_2\sin\omega_nt$。取 $t_j=(2j\pi+\pi/2)/\omega_n$，则 $x(t_j)=t_j/(2\omega_n)+C_2\to\infty$。所以没有有界解，也没有连续周期解。$\omega_n=0$ 时方程为 $\ddot x=1$，通解 $t^2/2+C_1+C_2t$，结论仍然相同。
\par\noindent{\small 依据：官方译审并补证（AI辅助）。}
\end{solution}

\begin{exercise}\label{cw:rec13-4}
在同一坐标系画 $\sin t,\sin2t$ 及和 $f(t)=\sin t+\sin2t$。在两项同为0的点求导数，求三个函数的最小周期。
\end{exercise}
\begin{solution}
$f(k\pi)=0$，$f'(k\pi)=(-1)^k+2$，偶数$k$斜率3，奇数$k$斜率1。三个最小周期为 $2\pi,\pi,2\pi$；和式中频率1项非零，不能缩成$\pi$。图中黑虚线为$\sin t$、灰线为$\sin2t$、蓝线为和。\begin{center}\begin{tikzpicture}\begin{axis}[width=.68\textwidth,height=4.5cm,axis lines=middle,xlabel={$t$},ylabel={$f$},samples=130,xmin=0,xmax=6.283,ymin=-2,ymax=2]
\addplot[black,dashed,domain=0:6.283]{sin(deg(x))};\addplot[gray,domain=0:6.283]{sin(deg(2*x))};\addplot[blue,domain=0:6.283]{sin(deg(x))+sin(deg(2*x))};\end{axis}\end{tikzpicture}\end{center}
\par\noindent{\small 依据：官方译审并补证（AI辅助）。}
\end{solution}

\begin{exercise}\label{cw:rec13-5}
设 $b_1,b_2\ne0$，哪些 $\omega_n$ 使 $\ddot x+\omega_n^2x=b_1\sin t+b_2\sin2t$ 有周期解？若有，举出一个。
\end{exercise}
\begin{solution}
$\omega_n^2\ne1,4$ 时可取 $x_p=b_1\sin t/(\omega_n^2-1)+b_2\sin2t/(\omega_n^2-4)$，周期$2\pi$。若 $\omega_n^2$ 等于1或4，对应非零输入产生 $t\cos\omega_nt$，有界的其他频率项和齐次项不能消除此增长，故无周期解。若约定$\omega_n\ge0$，只排除1、2；允许负值则同时排除$-1,-2$。
\par\noindent{\small 依据：官方译审并补证（AI辅助）。}
\end{solution}

\begin{exercise}\label{cw:rec13-6}
哪些实数 $\omega$ 使 $\sin t+\sin\omega t$ 周期？周期是多少？
\end{exercise}
\begin{solution}
若 $\omega=-1$，和恒为0，每个 $P>0$ 都是周期，没有最小正周期。$\omega=0$ 或1时最小周期为$2\pi$。
当 $\omega\notin\{0,\pm1\}$，函数由四个不同频率 $\pm1,\pm\omega$ 的指数项组成。若有周期$P$，相减 $f(t+P)-f(t)=0$，不同指数函数线性独立（在0求前四阶导数得到非零 Vandermonde 行列式），所以 $\ee^{\ii P}=\ee^{\ii\omega P}=1$。因此 $P=2\pi q$、$\omega P=2\pi p$，$\omega=p/q$ 必须有理。反之，若 $\omega=p/q$ 为既约分数，$q>0$，则 $2\pi q$ 是最小周期，全部周期为其正整数倍。负频率改变符号但不改变可公度条件。
\par\noindent{\small 依据：官方译审并补证（AI辅助）；补频率独立论证与负频率抵消例外。}
\end{solution}

\originalsection{第14次习题课：Fourier 系数与变换（3月30日）}
约定标准方波 $\operatorname{sq}(t)$ 为奇函数，周期$2\pi$，在$0<t<\pi$取1，跳点取左右平均。其级数为 $\frac4\pi\sum_{j\ge0}\sin((2j+1)t)/(2j+1)$。本节函数在跳点同样按中点值理解。
\begin{exercise}\label{cw:rec14-1}
画偶函数 $f$：周期$2\pi$，$0<t<\pi/2$ 时为2，$\pi/2<t<\pi$ 时为0。
(a) 用系数积分求级数，先利用奇偶性；(b) 用标准方波及三角恒等式另求；(c) 求$f-1$的级数。
\end{exercise}
\begin{solution}
(a) 偶性给 $b_n=0$；$a_0=\pi^{-1}\int_{-\pi}^\pi f=2$，$a_n=(4/\pi)\int_0^{\pi/2}\cos nt\dd t=4\sin(n\pi/2)/(n\pi)$。所以
\[f=1+\frac4\pi\sum_{j\ge0}\frac{(-1)^j\cos((2j+1)t)}{2j+1}.\]
(b) $f=1+\operatorname{sq}(t+\pi/2)$，因为平移后的方波正段恰为中心半周期；用 $\sin((2j+1)(t+\pi/2))=(-1)^j\cos((2j+1)t)$ 得同式。
(c) 去常数项1即可。以下图上跳点圆点取中点值1。\begin{center}\begin{tikzpicture}\begin{axis}[width=.68\textwidth,height=4.5cm,axis lines=middle,xlabel={$t$},ylabel={$f$},samples=130,xmin=-3.2,xmax=3.2,ymin=-.2,ymax=2.4]
\addplot[blue]coordinates{(-3.142,0)(-1.571,0)};\addplot[blue]coordinates{(-1.571,2)(1.571,2)};\addplot[blue]coordinates{(1.571,0)(3.142,0)};\addplot[only marks]coordinates{(-1.571,1)(1.571,1)};\end{axis}\end{tikzpicture}\end{center}
\par\noindent{\small 依据：官方译审并补证（AI辅助）。}
\end{solution}

\begin{exercise}\label{cw:rec14-2}
求 $\sin^2t$ 的 Fourier 级数。
\end{exercise}
\begin{solution}
倍角公式 $\sin^2t=(1-\cos2t)/2$ 已是有限 Fourier 级数。按 $a_0/2$ 的 convention，$a_0=1,a_2=-1/2$，其他$a_n,b_n$均为0。
\par\noindent{\small 依据：官方译审并补证（AI辅助）。}
\end{solution}

\begin{exercise}\label{cw:rec14-3}
画奇函数 $g$，周期$\pi$，在 $0<x<\pi/2$ 时为1。$2\pi$亦是其周期；用标准方波求其级数。
\end{exercise}
\begin{solution}
$g(x)=\operatorname{sq}(2x)$，所以 $g=\frac4\pi\sum_{j\ge0}\sin(2(2j+1)x)/(2j+1)$。频率只出现$2,6,10,\ldots$，系数分母仍是$1,3,5,\ldots$。\begin{center}\begin{tikzpicture}\begin{axis}[width=.68\textwidth,height=4.5cm,axis lines=middle,xlabel={$x$},ylabel={$g$},samples=130,xmin=-3.3,xmax=3.3,ymin=-1.3,ymax=1.3]
\addplot[blue] coordinates{(-3.142,1)(-1.571,1)};\addplot[blue]coordinates{(-1.571,-1)(0,-1)};\addplot[blue]coordinates{(0,1)(1.571,1)};\addplot[blue]coordinates{(1.571,-1)(3.142,-1)};\addplot[only marks]coordinates{(-3.142,0)(-1.571,0)(0,0)(1.571,0)(3.142,0)};\end{axis}\end{tikzpicture}\end{center}
\par\noindent{\small 依据：官方译审并补证（AI辅助）。}
\end{solution}

\begin{exercise}\label{cw:rec14-4}
画奇函数 $h$，周期$2\pi$；$0<t<\pi/2$ 时 $h=t$，$\pi/2<t<\pi$ 时 $h=\pi-t$。由第1(c)题求级数。
\end{exercise}
\begin{solution}
$h$ 连续、分段线性，除转角外 $h'=f-1$。对第1(c)题积分，并用$h(0)=0$，得
\[h(t)=\frac4\pi\sum_{j\ge0}\frac{(-1)^j\sin((2j+1)t)}{(2j+1)^2}.\]
积分合法：方波是$L^2$函数，其 Fourier 部分和在$L^2$中收敛，积分误差由 Cauchy--Schwarz 至多为$\sqrt{2\pi}$倍$L^2$误差；所得$1/n^2$级数一致收敛，故处处给出连续的$h$。\begin{center}\begin{tikzpicture}\begin{axis}[width=.68\textwidth,height=4.5cm,axis lines=middle,xlabel={$t$},ylabel={$h$},samples=130,xmin=-3.3,xmax=3.3,ymin=-2,ymax=2]
\addplot[blue] coordinates{(-3.142,0)(-1.571,-1.571)(0,0)(1.571,1.571)(3.142,0)};\end{axis}\end{tikzpicture}\end{center}
\par\noindent{\small 依据：官方译审并补证（AI辅助）。}
\end{solution}

\begin{exercise}\label{cw:rec14-5}
解释为什么任意定义在对称区间的函数 $g(x)$ 都唯一分解成偶函数与奇函数之和；求 $\ee^x$ 的两部分。
\end{exercise}
\begin{solution}
$E=[g(x)+g(-x)]/2$ 为偶，$O=[g(x)-g(-x)]/2$ 为奇，且$E+O=g$。若$g=E_1+O_1$，反射给$g(-x)=E_1-O_1$，相加相减强制得到上述两个公式，故唯一。指数函数的两部分是$\cosh x$与$\sinh x$。
\par\noindent{\small 依据：官方译审并补证（AI辅助）。}
\end{solution}

\originalsection{第15次习题课：Fourier 级数与谐波响应（4月1日）}
\begin{exercise}\label{cw:rec15-1}
偶函数 $f$ 周期$2\pi$，在 $-\pi<t<\pi$ 等于 $|t|$，画图。已知
\[f(t)=\frac\pi2-\frac4\pi\sum_{k\ge1,\ k\text{奇}}\frac{\cos kt}{k^2}.\]
把其图形水平缩放至圆频率$\omega>0$，求 $g$ 与$f$的关系及$g$的级数。
\end{exercise}
\begin{solution}
$g(t)=f(\omega t)$，周期变为$2\pi/\omega$，振幅不变。代入得 $g=\pi/2-(4/\pi)\sum_{k\text{奇}}\cos(k\omega t)/k^2$。注意只改变余弦自变量，不能在系数额外引入$\omega^2$。\begin{center}\begin{tikzpicture}\begin{axis}[width=.68\textwidth,height=4.5cm,axis lines=middle,xlabel={$t$},ylabel={$f$},samples=130,xmin=-6.5,xmax=6.5,ymin=-.2,ymax=3.5]
\addplot[blue]coordinates{(-6.283,0)(-3.142,3.142)(0,0)(3.142,3.142)(6.283,0)};\end{axis}\end{tikzpicture}\end{center}
\par\noindent{\small 依据：官方译审并补证（AI辅助）。}
\end{solution}

\begin{exercise}\label{cw:rec15-2}
用上一题 $f$ 驱动谐振子 $\ddot x+\omega_n^2x=f(t)$。何时存在周期解？若存在，写成 Fourier 级数。
\end{exercise}
\begin{solution}
当 $\omega_n\ge0$ 且 $\omega_n\notin\{0,1,3,5,\ldots\}$ 时，取
\[x_p=\frac\pi{2\omega_n^2}-\frac4\pi\sum_{k\text{奇}}\frac{\cos kt}{k^2(\omega_n^2-k^2)}.\]
固定非共振$\omega_n$时，尾项为$O(k^{-4})$，二阶导数项为$O(k^{-2})$，两者一致收敛，故可逐项二次求导并得到$f$，这确为$C^2$周期解。若$\omega_n$为奇整数，共振输入系数非零产生线性增长；若$\omega_n=0$，$x'(t+2\pi)-x'(t)=\int_t^{t+2\pi}f=\pi^2\ne0$，亦不可能周期。
\par\noindent{\small 依据：官方译审并补证（AI辅助）。}
\end{solution}

\begin{exercise}\label{cw:rec15-3}
改用圆频率$\omega>0$的$g(t)$驱动 $\ddot x+\omega_n^2x=g(t)$，求周期解及存在条件。
\end{exercise}
\begin{solution}
把每个频率改为$k\omega$，得到
\[x_p=\frac\pi{2\omega_n^2}-\frac4\pi\sum_{k\text{奇}}\frac{\cos(k\omega t)}{k^2[\omega_n^2-(k\omega)^2]}.
\]
存在条件是 $\omega_n>0$ 且不等于任何正奇数倍$\omega$。同上一题，尾项与二阶导数分别为$O(k^{-4})$、$O(k^{-2})$，可逐项核验；$\omega_n=0$ 时非零平均输入排除周期解。
\par\noindent{\small 依据：官方译审并补证（AI辅助）。}
\end{solution}

\begin{exercise}\label{cw:rec15-4}
固定$\omega$而调节自然圆频率$\omega_n$（例如调谐无线电电路）。在哪些值没有周期响应？当$\omega_n$略大或略小时响应如何变化？
\end{exercise}
\begin{solution}
非周期值为 $\omega_n=0$ 或 $k\omega$（$k$正奇数）。靠近非零共振时，频率$k\omega$的系数 $-4/[\pi k^2(\omega_n^2-k^2\omega^2)]$ 绝对值很大；略高于共振时此系数负，与该负余弦输入同相，略低时系数正，与该输入反相。靠近0时平均响应$\pi/(2\omega_n^2)$变得很大。精确共振的$t$增长项不能用大的有限稳态系数替代。
\par\noindent{\small 依据：官方译审并补证（AI辅助）。}
\end{solution}

\begin{exercise}\label{cw:rec15-5}
是否存在具有不止一个周期解的频率？给出条件。
\end{exercise}
\begin{solution}
一般解为 $x_p+C_1\cos\omega_nt+C_2\sin\omega_nt$。在非共振条件下，若 $\omega_n/\omega\in\mathbb Q$，齐次周期与输入周期可公度，任意$C_1,C_2$都给周期解，所以有无穷多个。例如$\omega_n=2\omega$。
若此比无理，则仅$C_1=C_2=0$给周期解：$x_p$含非零频率$\omega$项（连续周期函数的 Fourier 系数唯一），增加非零自然频率项后，一旦存在周期，两个非零频率都须是其基本频率的整数倍，矛盾。因此不能不加可公度条件便断言加任意齐次解仍周期。
\par\noindent{\small 依据：官方译审并补证（AI辅助）；修正任意齐次周期项可相加的缺失条件。}
\end{solution}

\originalchapter{习题课：阶跃、冲量、卷积与 Laplace 变换}
\originalsection{第16次习题课：广义导数与响应（4月6日）}
\begin{exercise}\label{cw:rec16-1}
令 $Q(t)=0$（$t<1$），$Q=2t-2$（$1<t<2$），$Q=2t-1$（$2<t<3$），$Q=5$（$t>3$）。
(a) 画图，判断是否分段光滑；(b) 求广义导数$q=Q'$并画图；(c) 举一个给出 $\dot x+kx=q$、静止初态的情境，可自选$k$甚至取负；(d) 举一个给出 $2\ddot x+4\dot x+4x=q$、静止初态的情境。
\end{exercise}
\begin{solution}
(a) 每段为一次函数或常数，端点有限，故分段光滑；$t=2$从2跳到3，跳幅1。
(b) $q=2u(t-1)-2u(t-3)+\delta(t-2)$；在$(1,3)$普通部分为2，外部为0，另有$t=2$的单位冲量。图的箭头高度仅标强度1，不把$\delta$当普通有限函数。
(c) 取$k=-r,r>0$，银行余额按年利率$r$连续增长；$t=1$到3以每年2单位存款，$t=2$另一次性存入1单位，之前无余额，即$\dot x-rx=q$。必须用$k=-r$才能表示正利息。
(d) 质量2、阻尼4、刚度4的静止弹簧系统受此外力：时间1到3恒力2，时间2施加总冲量1。冲量使速度跳$1/2$，位置连续。\begin{center}\begin{tikzpicture}\begin{axis}[width=.44\textwidth,height=3.8cm,axis lines=middle,xlabel={$t$},ylabel={$Q$},xmin=0,xmax=4,ymin=0,ymax=5.5]\addplot[blue]coordinates{(0,0)(1,0)(2,2)};\addplot[blue]coordinates{(2,3)(3,5)(4,5)};\draw[dashed](axis cs:2,2)--(axis cs:2,3);\end{axis}\end{tikzpicture}\begin{tikzpicture}\begin{axis}[width=.44\textwidth,height=3.8cm,axis lines=middle,xlabel={$t$},ylabel={$q$},xmin=0,xmax=4,ymin=0,ymax=5.5]\addplot[blue]coordinates{(0,0)(1,0)};\addplot[blue]coordinates{(1,2)(3,2)};\addplot[blue]coordinates{(3,0)(4,0)};\draw[red,->,thick](axis cs:2,2)--(axis cs:2,4)node[above]{强度1};\end{axis}\end{tikzpicture}\end{center}
\par\noindent{\small 依据：官方译审并补证（AI辅助）。}
\end{solution}

\begin{exercise}\label{cw:rec16-2}
求算子 $2D+I$ 的单位阶跃响应与单位冲量响应，并画图。
\end{exercise}
\begin{solution}
阶跃响应在$t>0$满足 $2\dot v+v=1,v(0+)=0$，解为 $v=u(t)(1-\ee^{-t/2})$。其广义导数 $w=\frac12u(t)\ee^{-t/2}$，原点无额外$\delta$，因为$v$连续。$w$的跳幅$1/2$使$(2D+I)w=\delta$。黑线为$v$、蓝线为$w$，$t<0$均为0。\begin{center}\begin{tikzpicture}\begin{axis}[width=.68\textwidth,height=4.5cm,axis lines=middle,xlabel={$t$},ylabel={响应},samples=130,xmin=0,xmax=7,ymin=0,ymax=1.1]
\addplot[black,domain=0:7]{1-exp(-x/2)};\addplot[blue,domain=0:7]{.5*exp(-x/2)};\end{axis}\end{tikzpicture}\end{center}
\par\noindent{\small 依据：官方译审并补证（AI辅助）。}
\end{solution}

\begin{exercise}\label{cw:rec16-3}
求 $D^2+2D$ 的单位冲量响应并画图。
\end{exercise}
\begin{solution}
冲量积分给 $w(0+)=0,w'(0+)=1$，$t>0$满足 $w''+2w'=0$。故 $w'=\ee^{-2t}$，$w=\frac12u(t)(1-\ee^{-2t})$。分布求导第二次在0产生系数1的$\delta$，普通部分抵消。\begin{center}\begin{tikzpicture}\begin{axis}[width=.68\textwidth,height=4.5cm,axis lines=middle,xlabel={$t$},ylabel={$w$},samples=130,xmin=0,xmax=4,ymin=0,ymax=.6]
\addplot[blue,domain=0:4]{(1-exp(-2*x))/2};\end{axis}\end{tikzpicture}\end{center}
\par\noindent{\small 依据：官方译审并补证（AI辅助）。}
\end{solution}

\begin{exercise}\label{cw:rec16-4}
利用上一题求 $\ddot x+2\dot x=3\delta(t-1)$ 的静止初态解。
\end{exercise}
\begin{solution}
时间不变性与线性性给 $x(t)=3w(t-1)=\frac32u(t-1)[1-\ee^{-2(t-1)}]$。时间1前为0，位置在1连续，速度从0跳到3；积分原方程穿过1可检验冲量强度恰为3。
\par\noindent{\small 依据：官方译审并补证（AI辅助）。}
\end{solution}

\originalsection{第17次习题课：卷积（4月8日）}
采用因果约定：所有函数在$t<0$为0，$f*g(t)=\int_0^tf(t-\tau)g(\tau)\dd\tau$。若$w$是$p(D)$的单位冲量响应，则静止初态响应为$w*q$。
\begin{exercise}\label{cw:rec17-1}
(a) 求 $t*u(t)$，并一般求$q*u$；(b) 求$u*t$，并一般求$u*q$。两者关系反映什么性质？
\end{exercise}
\begin{solution}
(a) $t*u=\int_0^t(t-\tau)\dd\tau=t^2/2$，$q*u=\int_0^tq(t-\tau)\dd\tau$。(b) $u*t=\int_0^t\tau\dd\tau=t^2/2$，$u*q=\int_0^tq(\tau)\dd\tau$。用$v=t-\tau$换元，两一般式相同；同样换元证明$f*g=g*f$，即卷积交换性。所有结果在$t<0$延为0。
\par\noindent{\small 依据：官方译审并补证（AI辅助）。}
\end{solution}

\begin{exercise}\label{cw:rec17-2}
哪一个微分算子$p(D)$的单位冲量响应是$u(t)$？上一题计算$u*q$是否验证了卷积响应结论？
\end{exercise}
\begin{solution}
$Du=\delta$，故$p(D)=D$。$x=u*q=\int_0^tq(\tau)\dd\tau$给原点前0，且$Dx=q$；普通连续$q$时由微积分基本定理成立，含冲量时由因果积分的广义导数成立。因此此算子确满足卷积响应结论。
\par\noindent{\small 依据：官方译审并补证（AI辅助）。}
\end{solution}

\begin{exercise}\label{cw:rec17-3}
(a) 若$f$在$a$连续，应如何理解$f(t)\delta(t-a)$？(b) 若$f$连续且$t<0$为0，解释$a\ge0$时$f*\delta(t-a)=f(t-a)$；说明$a=0$的含义。
\end{exercise}
\begin{solution}
(a) 作为测度乘连续函数，任意测试函数$\psi$满足 $\int\psi f\dd\delta_a=\psi(a)f(a)$，故乘积为$f(a)\delta(t-a)$。若坚持一般分布框架可先对光滑$f$定义，再用本题的点质量测度解释连续情形。(b) $t\ge a$ 时卷积抽样为$f(t-a)$，$t<a$时没有支持点，结果为0，与因果延拓$f(t-a)=0$一致。包括原点冲量时积分应从$0^-$开始，不能误取半个冲量。$a=0$给$f*\delta=f$，即卷积单位元。
\par\noindent{\small 依据：官方译审并补证（AI辅助）。}
\end{solution}

\begin{exercise}\label{cw:rec17-4}
(a) 核验 $u(t)\sin(\omega_nt)/\omega_n$ 是 $D^2+\omega_n^2I$ 的单位冲量响应；(b) 用共振 ERF 求 $\ddot x+x=\sin t,x(0)=\dot x(0)=0$；(c) 在 $t=2\pi n,n\in\mathbb N_{>0}$ 计算$\sin t*\sin t$，并与(b)核对，可用$\sin^2t=(1-\cos2t)/2$。
\end{exercise}
\begin{solution}
(a) 记普通部分$v=\sin\omega_nt/\omega_n$，$v(0)=0,v'(0)=1$，故$D^2(uv)=u v''+\delta$；加入$\omega_n^2uv$后普通项抵消，留下$\delta$。$\omega_n=0$时连续极限为$u(t)t$。
(b) 复替代$z''+z=\ee^{\ii t}$给$z_p=t\ee^{\ii t}/(2\ii)$，虚部$-t\cos t/2$；补齐初值后 $x=(\sin t-t\cos t)/2$。
(c) $\sin(2\pi n-\tau)=-\sin\tau$，所以卷积值为$-\int_0^{2\pi n}\sin^2\tau\dd\tau=-\pi n$。代入(b)也为$-\pi n$。两者不仅在这些点相同，卷积一般计算用积化和差得$(\sin t-t\cos t)/2$。
\par\noindent{\small 依据：官方译审并补证（AI辅助）。}
\end{solution}

\originalsection{第18次习题课：Laplace 变换（4月13日）}
本次变换按因果延拓从$0^-$积分，$\Lap[Df]=sF(s)$中的$D$为广义导数；对普通右导数则$\Lap[f_r']=sF-f(0+)$。所有公式在相应右半平面成立。
\begin{exercise}\label{cw:rec18-1}
从规则与表求 $u(t)\ee^{-t}(t^2+1)$ 的 Laplace 变换。
\end{exercise}
\begin{solution}
$\Lap[t^2+1]=2/s^3+1/s$；乘$\ee^{-t}$使$s$换成$s+1$，故 $F=2/(s+1)^3+1/(s+1)=(s^2+2s+3)/(s+1)^3$，$\Re s>-1$。
\par\noindent{\small 依据：官方译审并补证（AI辅助）。}
\end{solution}

\begin{exercise}\label{cw:rec18-2}
令 $f(t)=\ee^{-t}\cos3t$（隐含因果延拓）。求$F(s)$、广义导数$f'$及其变换，核验$t$导数规则。
\end{exercise}
\begin{solution}
$F=(s+1)/[(s+1)^2+9]$。原点跳幅为1，故 $Df=\delta-u(t)\ee^{-t}(\cos3t+3\sin3t)$。变换为 $1-(s+1)/[(s+1)^2+9]-9/[(s+1)^2+9]=s(s+1)/[(s+1)^2+9]=sF$，正好符合广义导数规则。
\par\noindent{\small 依据：官方译审并补证（AI辅助）。}
\end{solution}

\begin{exercise}\label{cw:rec18-3}
分别求 $(2s+1)/(s^2+9)$、$(s^2+2)/(s^3-s)$、$2/[s^2(s-1)]$ 的逆 Laplace 变换。
\end{exercise}
\begin{solution}
第一式分成$2s/(s^2+9)+(1/3)3/(s^2+9)$，故为$u(t)(2\cos3t+\sin3t/3)$。
第二式的部分分式为 $-2/s+(3/2)/(s-1)+(3/2)/(s+1)$，故为$u(t)[-2+\frac32(\ee^t+\ee^{-t})]$。
第三式为 $-2/s-2/s^2+2/(s-1)$，故为$2u(t)(\ee^t-1-t)$。每式分母通分即可核对系数，因果条件固定其右半平面解释。
\par\noindent{\small 依据：官方译审并补证（AI辅助）。}
\end{solution}

\begin{exercise}\label{cw:rec18-4}
用 Laplace 变换求$D+2I$的单位阶跃与单位冲量响应。
\end{exercise}
\begin{solution}
阶跃输入变换为$1/s$，$V=1/[s(s+2)]=\frac12(1/s-1/(s+2))$，故$v=\frac12u(t)(1-\ee^{-2t})$。冲量输入变换为1，$W=1/(s+2)$，故$w=u(t)\ee^{-2t}$。$v'=w$，且$w$原点跳幅1，均可检验。
\par\noindent{\small 依据：官方译审并补证（AI辅助）。}
\end{solution}

\begin{exercise}\label{cw:rec18-5}
用 Laplace 变换求 $\dot x+2x=t^2,x(0+)=1$。
\end{exercise}
\begin{solution}
对普通导数取变换：$(s+2)X-1=2/s^3$，故 $X=(s^3+2)/[s^3(s+2)]=1/s^3-1/(2s^2)+1/(4s)+3/[4(s+2)]$。逆变换给 $x=t^2/2-t/2+1/4+3\ee^{-2t}/4$，$t\ge0$。在0值为1，求导加$2x$后指数项抵消、多项式项为$t^2$。若使用因果广义导数，则方程右侧另加$\delta(t)$，结果相同。
\par\noindent{\small 依据：官方译审并补证（AI辅助）。}
\end{solution}

\originalsection{第19次习题课：Laplace 变换与初值（4月15日）}
\begin{exercise}\label{cw:rec19-1}
求$D+3I$的单位冲量响应。
\end{exercise}
\begin{solution}
$(s+3)W=1$，故$w=u(t)\ee^{-3t}$；跳幅1使广义导数产生单位冲量，普通部分满足$w'+3w=0$。
\par\noindent{\small 依据：官方译审并补证（AI辅助）。}
\end{solution}

\begin{exercise}\label{cw:rec19-2}
用 Laplace 变换求 $\dot x+3x=\ee^{-t}$ 的静止初值解$x(0)=0$。
\end{exercise}
\begin{solution}
$X=1/[(s+1)(s+3)]=\frac12[1/(s+1)-1/(s+3)]$，所以 $x=\frac12(\ee^{-t}-\ee^{-3t})$，$t\ge0$；在0值为0，代入方程得到输入$\ee^{-t}$。
\par\noindent{\small 依据：官方译审并补证（AI辅助）。}
\end{solution}

\begin{exercise}\label{cw:rec19-3}
用 Laplace 变换求 $D^3+D$ 的单位冲量响应。
\end{exercise}
\begin{solution}
$W=1/[s(s^2+1)]=1/s-s/(s^2+1)$，故$w=u(t)(1-\cos t)$。普通部分满足$w^{(3)}+w'=0$；$w(0)=w'(0)=0,w''(0)=1$，所以三阶广义导数在0恰产生$\delta$。
\par\noindent{\small 依据：官方译审并补证（AI辅助）。}
\end{solution}

\begin{exercise}\label{cw:rec19-4}
(a) 用普通导数规则求 $x_r'+3x=1,x(0+)=2$；(b) 以 $x(0-)=0$，把原点跳跃写为 $Dx+3x=u(t)+2\delta(t)$，再求解。
\end{exercise}
\begin{solution}
(a) $(s+3)X-2=1/s$，$X=1/[s(s+3)]+2/(s+3)$，故$x=1/3+5\ee^{-3t}/3$，$t\ge0$。(b) 广义规则给$(s+3)X=1/s+2$，同一变换；因果解为$u(t)(1/3+5\ee^{-3t}/3)$。原点跳幅2验证了输入冲量系数，原点后普通方程一致。
\par\noindent{\small 依据：官方译审并补证（AI辅助）。}
\end{solution}

\originalchapter{习题课：线性系统与非线性相图}
\originalsection{第21次习题课：矩阵与一阶线性系统（4月27日）}
\begin{exercise}\label{cw:rec21-1}
计算四个矩阵乘积：
\[[1\quad2]\binom xy,\qquad\binom12[x\quad y],\qquad\begin{pmatrix}a&b\\c&d\end{pmatrix}\binom xy,\qquad\begin{pmatrix}1&2\\3&4\end{pmatrix}\begin{pmatrix}x&u\\y&v\end{pmatrix}.\]
\end{exercise}
\begin{solution}
按照行乘列，结果依次为标量$x+2y$、外积矩阵$\begin{pmatrix}x&y\\2x&2y\end{pmatrix}$、列向量$\binom{ax+by}{cx+dy}$、矩阵$\begin{pmatrix}x+2y&u+2v\\3x+4y&3u+4v\end{pmatrix}$。第二个乘积次序与第一个不同，输出尺寸亦不同。
\par\noindent{\small 依据：官方译审并补证（AI辅助）。}
\end{solution}

\begin{exercise}\label{cw:rec21-2}
矩阵$A$把单位方形顶点 $0,e_1,e_1+e_2,e_2$ 送至$0,Ae_1,A(e_1+e_2),Ae_2$。对以下五个$A$依序连接这些像点，画图并描述变换：
\[\begin{pmatrix}1&0\\0&2\end{pmatrix},\quad\begin{pmatrix}1&1\\0&1\end{pmatrix},\quad\begin{pmatrix}1&0\\0&-1\end{pmatrix},\quad\begin{pmatrix}0&1\\1&0\end{pmatrix},\quad\begin{pmatrix}1&1\\1&-1\end{pmatrix}.\]
\end{exercise}
\begin{solution}
像点依次为 $(0,0),(1,0),(1,2),(0,2)$；$(0,0),(1,0),(2,1),(1,1)$；\par $(0,0),(1,0),(1,-1),(0,-1)$；$(0,0),(0,1),(1,1),(1,0)$；\par $(0,0),(1,1),(2,0),(1,-1)$。各组均再闭合回原点。
对应为竖直放大2倍、水平剪切、关于$x$轴反射、关于$y=x$反射、以及先关于$x$轴反射再逆时针旋转$\pi/4$并放大$\sqrt2$。最后一个次序可由 $A=\sqrt2 R_{\pi/4}\diag(1,-1)$核验。\begin{center}\begin{tikzpicture}\begin{axis}[width=.28\textwidth,height=3.5cm,axis lines=middle,title={竖直伸长},xmin=-1.5,xmax=2.5,ymin=-1.5,ymax=2.5,xlabel={$x$},ylabel={$y$}]\addplot[blue,mark=*]coordinates {(0,0) (1,0) (1,2) (0,2) (0,0)};\end{axis}\end{tikzpicture}\begin{tikzpicture}\begin{axis}[width=.28\textwidth,height=3.5cm,axis lines=middle,title={剪切},xmin=-1.5,xmax=2.5,ymin=-1.5,ymax=2.5,xlabel={$x$},ylabel={$y$}]\addplot[blue,mark=*]coordinates {(0,0) (1,0) (2,1) (1,1) (0,0)};\end{axis}\end{tikzpicture}\begin{tikzpicture}\begin{axis}[width=.28\textwidth,height=3.5cm,axis lines=middle,title={关于横轴反射},xmin=-1.5,xmax=2.5,ymin=-1.5,ymax=2.5,xlabel={$x$},ylabel={$y$}]\addplot[blue,mark=*]coordinates {(0,0) (1,0) (1,-1) (0,-1) (0,0)};\end{axis}\end{tikzpicture}\par\begin{tikzpicture}\begin{axis}[width=.28\textwidth,height=3.5cm,axis lines=middle,title={交换坐标},xmin=-1.5,xmax=2.5,ymin=-1.5,ymax=2.5,xlabel={$x$},ylabel={$y$}]\addplot[blue,mark=*]coordinates {(0,0) (0,1) (1,1) (1,0) (0,0)};\end{axis}\end{tikzpicture}\begin{tikzpicture}\begin{axis}[width=.28\textwidth,height=3.5cm,axis lines=middle,title={反射、旋转与伸长},xmin=-1.5,xmax=2.5,ymin=-1.5,ymax=2.5,xlabel={$x$},ylabel={$y$}]\addplot[blue,mark=*]coordinates {(0,0) (1,1) (2,0) (1,-1) (0,0)};\end{axis}\end{tikzpicture}\end{center}
\par\noindent{\small 依据：官方译审并补证（AI辅助）。}
\end{solution}

\begin{exercise}\label{cw:rec21-3}
求$\ddot x+2\dot x+2x=0$的伴随矩阵$A$，写两个实独立解。令$x_1(0)=0,\dot x_1(0)=1$，求$x_1$及$u_1=(x_1,\dot x_1)^T$、$u_1(0)$；画$x_1,\dot x_1$时间图与$u_1$轨迹，比较。再画几条轨迹，特别是$u_2(0)=e_1$的轨迹。轨迹穿越$x$轴时角度为何？
\end{exercise}
\begin{solution}
设$v=\dot x$，$\dot u=Au$，$A=\begin{pmatrix}0&1\\-2&-2\end{pmatrix}$。特征根$-1\pm\ii$给$\ee^{-t}\cos t,\ee^{-t}\sin t$。
$x_1=\ee^{-t}\sin t$，$u_1=\ee^{-t}(\sin t,\cos t-\sin t)^T$，初值$e_2$。另一规范解$x_2=\ee^{-t}(\cos t+\sin t)$，$u_2=\ee^{-t}(\cos t+\sin t,-2\sin t)^T$，初值$e_1$。它们的线性组合给整个相图。
时间图中黑线为$x_1$、蓝线为$\dot x_1$；相图把同一时刻的二者作为坐标，向原点顺时针螺旋。在$x$轴上$v=0$，速度$(\dot x,\dot v)=(0,-2x)$竖直，故除平衡点外都垂直穿越，夹角$\pi/2$。\begin{center}\begin{tikzpicture}\begin{axis}[width=.68\textwidth,height=4.5cm,axis lines=middle,xlabel={$t$},ylabel={$x_1,\dot x_1$},samples=130,xmin=0,xmax=8,ymin=-.5,ymax=1.1]
\addplot[black,domain=0:8]{exp(-x)*sin(deg(x))};\addplot[blue,domain=0:8]{exp(-x)*(cos(deg(x))-sin(deg(x)))};\end{axis}\end{tikzpicture}\end{center}\begin{center}\begin{tikzpicture}\begin{axis}[width=.5\textwidth,height=3.8cm,axis lines=middle,title={伴随系统相图},xlabel={$x$},ylabel={$y$},xmin=-2.4,xmax=2.4,ymin=-2.4,ymax=2.4,clip=true]\addplot[blue,no marks] coordinates {(0.02346,3.05467) (0.32628,2.40884) (0.56155,1.83847) (0.73754,1.34114) (0.86216,0.91327) (0.94290,0.55049) (0.98671,0.24783) (1.00000,0.00000) (0.98854,-0.19844) (0.95753,-0.35296) (0.91154,-0.46889) (0.85457,-0.55136) (0.79007,-0.60521) (0.72097,-0.63496) (0.64971,-0.64476) (0.57829,-0.63835) (0.50833,-0.61912) (0.44107,-0.59004) (0.37747,-0.55372) (0.31821,-0.51240) (0.26372,-0.46799) (0.21426,-0.42210) (0.16993,-0.37602) (0.13067,-0.33081) (0.09635,-0.28731) (0.06674,-0.24612) (0.04156,-0.20770) (0.02047,-0.17235) (0.00313,-0.14024) (-0.01082,-0.11142) (-0.02175,-0.08589) (-0.03002,-0.06355) (-0.03598,-0.04425) (-0.03996,-0.02782) (-0.04226,-0.01405) (-0.04317,-0.00272) (-0.04295,0.00642) (-0.04182,0.01360) (-0.03999,0.01904) (-0.03764,0.02298) (-0.03493,0.02563) (-0.03199,0.02718) (-0.02893,0.02782) (-0.02583,0.02772) (-0.02279,0.02703) (-0.01984,0.02588) (-0.01705,0.02438) (-0.01443,0.02265) (-0.01202,0.02076) (-0.00982,0.01879) (-0.00785,0.01679) (-0.00609,0.01483) (-0.00455,0.01292)};\draw[blue,->] (axis cs:0.94727,-0.38536)--(axis cs:0.89377,-0.49956);\addplot[blue,no marks] coordinates {(-0.02346,-3.05467) (-0.32628,-2.40884) (-0.56155,-1.83847) (-0.73754,-1.34114) (-0.86216,-0.91327) (-0.94290,-0.55049) (-0.98671,-0.24783) (-1.00000,0.00000) (-0.98854,0.19844) (-0.95753,0.35296) (-0.91154,0.46889) (-0.85457,0.55136) (-0.79007,0.60521) (-0.72097,0.63496) (-0.64971,0.64476) (-0.57829,0.63835) (-0.50833,0.61912) (-0.44107,0.59004) (-0.37747,0.55372) (-0.31821,0.51240) (-0.26372,0.46799) (-0.21426,0.42210) (-0.16993,0.37602) (-0.13067,0.33081) (-0.09635,0.28731) (-0.06674,0.24612) (-0.04156,0.20770) (-0.02047,0.17235) (-0.00313,0.14024) (0.01082,0.11142) (0.02175,0.08589) (0.03002,0.06355) (0.03598,0.04425) (0.03996,0.02782) (0.04226,0.01405) (0.04317,0.00272) (0.04295,-0.00642) (0.04182,-0.01360) (0.03999,-0.01904) (0.03764,-0.02298) (0.03493,-0.02563) (0.03199,-0.02718) (0.02893,-0.02782) (0.02583,-0.02772) (0.02279,-0.02703) (0.01984,-0.02588) (0.01705,-0.02438) (0.01443,-0.02265) (0.01202,-0.02076) (0.00982,-0.01879) (0.00785,-0.01679) (0.00609,-0.01483) (0.00455,-0.01292)};\draw[blue,->] (axis cs:-0.94727,0.38536)--(axis cs:-0.89377,0.49956);\addplot[blue,no marks] coordinates {(-1.88847,3.42156) (-1.52734,3.07813) (-1.20442,2.73512) (-0.91924,2.40003) (-0.67057,2.07868) (-0.45664,1.77543) (-0.27524,1.49338) (-0.12391,1.23454) (0.00000,1.00000) (0.09922,0.79010) (0.17648,0.60457) (0.23445,0.44265) (0.27568,0.30321) (0.30261,0.18486) (0.31748,0.08601) (0.32238,0.00495) (0.31918,-0.06006) (0.30956,-0.11079) (0.29502,-0.14897) (0.27686,-0.17624) (0.25620,-0.19419) (0.23400,-0.20427) (0.21105,-0.20783) (0.18801,-0.20609) (0.16541,-0.20014) (0.14365,-0.19095) (0.12306,-0.17938) (0.10385,-0.16615) (0.08618,-0.15188) (0.07012,-0.13710) (0.05571,-0.12224) (0.04294,-0.10764) (0.03177,-0.09357) (0.02213,-0.08023) (0.01391,-0.06778) (0.00703,-0.05631) (0.00136,-0.04589) (-0.00321,-0.03653) (-0.00680,-0.02822) (-0.00952,-0.02095) (-0.01149,-0.01466) (-0.01281,-0.00930) (-0.01359,-0.00481) (-0.01391,-0.00110) (-0.01386,0.00189) (-0.01352,0.00424) (-0.01294,0.00603) (-0.01219,0.00734) (-0.01132,0.00822) (-0.01038,0.00874) (-0.00939,0.00896) (-0.00840,0.00895) (-0.00741,0.00873) (-0.00646,0.00837)};\draw[blue,->] (axis cs:0.19268,0.56191)--(axis cs:0.24978,0.39421);\addplot[blue,no marks] coordinates {(1.88847,-3.42156) (1.52734,-3.07813) (1.20442,-2.73512) (0.91924,-2.40003) (0.67057,-2.07868) (0.45664,-1.77543) (0.27524,-1.49338) (0.12391,-1.23454) (0.00000,-1.00000) (-0.09922,-0.79010) (-0.17648,-0.60457) (-0.23445,-0.44265) (-0.27568,-0.30321) (-0.30261,-0.18486) (-0.31748,-0.08601) (-0.32238,-0.00495) (-0.31918,0.06006) (-0.30956,0.11079) (-0.29502,0.14897) (-0.27686,0.17624) (-0.25620,0.19419) (-0.23400,0.20427) (-0.21105,0.20783) (-0.18801,0.20609) (-0.16541,0.20014) (-0.14365,0.19095) (-0.12306,0.17938) (-0.10385,0.16615) (-0.08618,0.15188) (-0.07012,0.13710) (-0.05571,0.12224) (-0.04294,0.10764) (-0.03177,0.09357) (-0.02213,0.08023) (-0.01391,0.06778) (-0.00703,0.05631) (-0.00136,0.04589) (0.00321,0.03653) (0.00680,0.02822) (0.00952,0.02095) (0.01149,0.01466) (0.01281,0.00930) (0.01359,0.00481) (0.01391,0.00110) (0.01386,-0.00189) (0.01352,-0.00424) (0.01294,-0.00603) (0.01219,-0.00734) (0.01132,-0.00822) (0.01038,-0.00874) (0.00939,-0.00896) (0.00840,-0.00895) (0.00741,-0.00873) (0.00646,-0.00837)};\draw[blue,->] (axis cs:-0.19268,-0.56191)--(axis cs:-0.24978,-0.39421);\addplot[blue,no marks] coordinates {(0.06697,3.41982) (0.40552,2.68870) (0.66765,2.04387) (0.86280,1.48237) (1.00000,1.00000) (1.08777,0.59165) (1.13401,0.25161) (1.14598,-0.02624) (1.13025,-0.24815) (1.09268,-0.42035) (1.03845,-0.54896) (0.97209,-0.63981) (0.89747,-0.69841) (0.81789,-0.72991) (0.73609,-0.73901) (0.65433,-0.72996) (0.57441,-0.70659) (0.49772,-0.67227) (0.42531,-0.62993) (0.35794,-0.58210) (0.29608,-0.53095) (0.24000,-0.47826) (0.18980,-0.42550) (0.14541,-0.37385) (0.10665,-0.32424) (0.07325,-0.27734) (0.04490,-0.23366) (0.02120,-0.19353) (0.00175,-0.15711) (-0.01386,-0.12448) (-0.02605,-0.09560) (-0.03524,-0.07037) (-0.04182,-0.04860) (-0.04616,-0.03010) (-0.04862,-0.01463) (-0.04951,-0.00191) (-0.04913,0.00832) (-0.04774,0.01632) (-0.04558,0.02237) (-0.04284,0.02672) (-0.03969,0.02961) (-0.03630,0.03127) (-0.03278,0.03191) (-0.02924,0.03172) (-0.02576,0.03086) (-0.02240,0.02950) (-0.01922,0.02775) (-0.01624,0.02574) (-0.01350,0.02356) (-0.01101,0.02129)};\draw[blue,->] (axis cs:1.13995,0.17655)--(axis cs:1.14355,-0.10535);\addplot[blue,no marks] coordinates {(-0.06697,-3.41982) (-0.40552,-2.68870) (-0.66765,-2.04387) (-0.86280,-1.48237) (-1.00000,-1.00000) (-1.08777,-0.59165) (-1.13401,-0.25161) (-1.14598,0.02624) (-1.13025,0.24815) (-1.09268,0.42035) (-1.03845,0.54896) (-0.97209,0.63981) (-0.89747,0.69841) (-0.81789,0.72991) (-0.73609,0.73901) (-0.65433,0.72996) (-0.57441,0.70659) (-0.49772,0.67227) (-0.42531,0.62993) (-0.35794,0.58210) (-0.29608,0.53095) (-0.24000,0.47826) (-0.18980,0.42550) (-0.14541,0.37385) (-0.10665,0.32424) (-0.07325,0.27734) (-0.04490,0.23366) (-0.02120,0.19353) (-0.00175,0.15711) (0.01386,0.12448) (0.02605,0.09560) (0.03524,0.07037) (0.04182,0.04860) (0.04616,0.03010) (0.04862,0.01463) (0.04951,0.00191) (0.04913,-0.00832) (0.04774,-0.01632) (0.04558,-0.02237) (0.04284,-0.02672) (0.03969,-0.02961) (0.03630,-0.03127) (0.03278,-0.03191) (0.02924,-0.03172) (0.02576,-0.03086) (0.02240,-0.02950) (0.01922,-0.02775) (0.01624,-0.02574) (0.01350,-0.02356) (0.01101,-0.02129)};\draw[blue,->] (axis cs:-1.13995,-0.17655)--(axis cs:-1.14355,0.10535);\end{axis}\end{tikzpicture}\end{center}
\par\noindent{\small 依据：官方译审并补证（AI辅助）。}
\end{solution}

\begin{exercise}\label{cw:rec21-4}
若$(a+b\ii)(x+y\ii)=v+w\ii$，求把$(x,y)^T$送到$(v,w)^T$的实矩阵；再取$a+b\ii=2,\ii,1+\ii$，分别写矩阵并画单位方形的像。
\end{exercise}
\begin{solution}
展开得$v=ax-by,w=bx+ay$，所以$A=\begin{pmatrix}a&-b\\b&a\end{pmatrix}$。三例分别为$2I,\begin{pmatrix}0&-1\\1&0\end{pmatrix},\begin{pmatrix}1&-1\\1&1\end{pmatrix}$，即放大2倍、逆时针旋转$\pi/2$、逆时针旋转$\pi/4$并放大$\sqrt2$。\begin{center}\begin{tikzpicture}\begin{axis}[width=.28\textwidth,height=3.5cm,axis lines=middle,title={乘以2},xmin=-1.5,xmax=2.5,ymin=-1.5,ymax=2.5,xlabel={$x$},ylabel={$y$}]\addplot[blue,mark=*]coordinates {(0,0) (2,0) (2,2) (0,2) (0,0)};\end{axis}\end{tikzpicture}\begin{tikzpicture}\begin{axis}[width=.28\textwidth,height=3.5cm,axis lines=middle,title={乘以$\ii$},xmin=-1.5,xmax=2.5,ymin=-1.5,ymax=2.5,xlabel={$x$},ylabel={$y$}]\addplot[blue,mark=*]coordinates {(0,0) (0,1) (-1,1) (-1,0) (0,0)};\end{axis}\end{tikzpicture}\begin{tikzpicture}\begin{axis}[width=.28\textwidth,height=3.5cm,axis lines=middle,title={乘以$1+\ii$},xmin=-1.5,xmax=2.5,ymin=-1.5,ymax=2.5,xlabel={$x$},ylabel={$y$}]\addplot[blue,mark=*]coordinates {(0,0) (1,1) (0,2) (-1,1) (0,0)};\end{axis}\end{tikzpicture}\end{center}
\par\noindent{\small 依据：官方译审并补证（AI辅助）。}
\end{solution}

\originalsection{第22次习题课：特征值与特征向量（4月29日）}
\begin{exercise}\label{cw:rec22-1}
解系统$\dot x=-5x-3y,\dot y=6x+4y$。
(a) 写系数矩阵、迹、行列式、$p_A(\lambda)=\det(A-\lambda I)$及其系数关系；(b) 求特征值和非零特征向量；(c) 画特征线及时间方向，解释每线为何分成三条轨迹；取线上非零点，写所有经过它的解；(d) 把初值$u(0)=(1,0)^T$分解，写解、画轨迹，再写具有同一轨迹的所有解；(e) 补完整相图。
\end{exercise}
\begin{solution}
(a) $A=\begin{pmatrix}-5&-3\\6&4\end{pmatrix}$，$\tau=-1,\Delta=-2$，$p_A=\lambda^2-\tau\lambda+\Delta=\lambda^2+\lambda-2$。
(b) $\lambda_+=1,v_+=(1,-2)^T$；$\lambda_-=-2,v_-=(1,-1)^T$，直接乘矩阵核验。
(c) 特征线$y=-2x$向外，$y=-x$向内；线上解为$ce^{\lambda t}v$。因指数始终正，正半线、原点、负半线是三条不同轨迹。若固定非零点$p=cv$，所有在某时刻$t_0$经过$p$的解为$u(t)=\ee^{\lambda(t-t_0)}p$，$t_0\in\mathbb R$。
(d) $(1,0)^T=-v_++2v_-$，故 $u_0(t)=(-\ee^t+2\ee^{-2t},2\ee^t-2\ee^{-2t})^T$；同一轨迹的所有参数化解为$u_0(t-t_0)$，相当于$c_+=-\ee^{-t_0},c_-=2\ee^{2t_0}$。
(e) 一般解$u=c_+\ee^tv_++c_-\ee^{-2t}v_-$，相图为鞍点；非特征轨迹前向沿不稳定线远离，反向沿稳定线远离。图的箭头为时间正向。\begin{center}\begin{tikzpicture}\begin{axis}[width=.55\textwidth,height=3.8cm,axis lines=middle,title={稳定线$y=-x$},xlabel={$x$},ylabel={$y$},xmin=-2.4,xmax=2.4,ymin=-2.4,ymax=2.4,clip=true]\addplot[blue,no marks] coordinates {(3.17894,-2.46241) (2.72835,-1.97089) (2.31851,-1.51777) (1.94474,-1.09826) (1.60286,-0.70802) (1.28908,-0.34312) (1.00000,0.00000) (0.73255,0.32458) (0.48396,0.63356) (0.25170,0.92966) (0.03351,1.21534) (-0.17269,1.49288) (-0.36878,1.76439) (-0.55649,2.03183) (-0.73740,2.29702) (-0.91296,2.56168) (-1.08452,2.82743) (-1.25333,3.09581) (-1.42054,3.36827)};\draw[blue,->] (axis cs:-0.07096,1.35499)--(axis cs:-0.49351,1.94124);\addplot[blue,no marks] coordinates {(3.39472,-3.39472) (3.03773,-3.03773) (2.71828,-2.71828) (2.43243,-2.43243) (2.17663,-2.17663) (1.94773,-1.94773) (1.74291,-1.74291) (1.55962,-1.55962) (1.39561,-1.39561) (1.24885,-1.24885) (1.11752,-1.11752) (1.00000,-1.00000) (0.89484,-0.89484) (0.80074,-0.80074) (0.71653,-0.71653) (0.64118,-0.64118) (0.57375,-0.57375) (0.51342,-0.51342) (0.45943,-0.45943) (0.41111,-0.41111) (0.36788,-0.36788) (0.32919,-0.32919) (0.29457,-0.29457) (0.26360,-0.26360) (0.23588,-0.23588) (0.21107,-0.21107) (0.18888,-0.18888) (0.16901,-0.16901) (0.15124,-0.15124) (0.13534,-0.13534) (0.12110,-0.12110) (0.10837,-0.10837) (0.09697,-0.09697) (0.08677,-0.08677) (0.07765,-0.07765) (0.06948,-0.06948) (0.06218,-0.06218) (0.05564,-0.05564) (0.04979,-0.04979) (0.04455,-0.04455) (0.03987,-0.03987) (0.03567,-0.03567) (0.03192,-0.03192) (0.02857,-0.02857) (0.02556,-0.02556) (0.02287,-0.02287) (0.02047,-0.02047) (0.01832,-0.01832)};\draw[blue,->] (axis cs:0.60653,-0.60653)--(axis cs:0.47711,-0.47711);\addplot[blue,no marks] coordinates {(-3.39472,3.39472) (-3.03773,3.03773) (-2.71828,2.71828) (-2.43243,2.43243) (-2.17663,2.17663) (-1.94773,1.94773) (-1.74291,1.74291) (-1.55962,1.55962) (-1.39561,1.39561) (-1.24885,1.24885) (-1.11752,1.11752) (-1.00000,1.00000) (-0.89484,0.89484) (-0.80074,0.80074) (-0.71653,0.71653) (-0.64118,0.64118) (-0.57375,0.57375) (-0.51342,0.51342) (-0.45943,0.45943) (-0.41111,0.41111) (-0.36788,0.36788) (-0.32919,0.32919) (-0.29457,0.29457) (-0.26360,0.26360) (-0.23588,0.23588) (-0.21107,0.21107) (-0.18888,0.18888) (-0.16901,0.16901) (-0.15124,0.15124) (-0.13534,0.13534) (-0.12110,0.12110) (-0.10837,0.10837) (-0.09697,0.09697) (-0.08677,0.08677) (-0.07765,0.07765) (-0.06948,0.06948) (-0.06218,0.06218) (-0.05564,0.05564) (-0.04979,0.04979) (-0.04455,0.04455) (-0.03987,0.03987) (-0.03567,0.03567) (-0.03192,0.03192) (-0.02857,0.02857) (-0.02556,0.02556) (-0.02287,0.02287) (-0.02047,0.02047) (-0.01832,0.01832)};\draw[blue,->] (axis cs:-0.60653,0.60653)--(axis cs:-0.47711,0.47711);\addplot[blue,no marks] coordinates {(0.36788,-0.73576) (0.38890,-0.77779) (0.41111,-0.82222) (0.43460,-0.86920) (0.45943,-0.91885) (0.48567,-0.97134) (0.51342,-1.02683) (0.54275,-1.08549) (0.57375,-1.14751) (0.60653,-1.21306) (0.64118,-1.28236) (0.67781,-1.35562) (0.71653,-1.43306) (0.75747,-1.51493) (0.80074,-1.60147) (0.84648,-1.69296) (0.89484,-1.78968) (0.94596,-1.89192) (1.00000,-2.00000) (1.05713,-2.11426) (1.11752,-2.23504) (1.18136,-2.36272) (1.24885,-2.49770) (1.32019,-2.64039) (1.39561,-2.79122) (1.47534,-2.95068) (1.55962,-3.11925) (1.64872,-3.29744) (1.74291,-3.48582)};\addplot[blue,no marks] coordinates {(-0.36788,0.73576) (-0.38890,0.77779) (-0.41111,0.82222) (-0.43460,0.86920) (-0.45943,0.91885) (-0.48567,0.97134) (-0.51342,1.02683) (-0.54275,1.08549) (-0.57375,1.14751) (-0.60653,1.21306) (-0.64118,1.28236) (-0.67781,1.35562) (-0.71653,1.43306) (-0.75747,1.51493) (-0.80074,1.60147) (-0.84648,1.69296) (-0.89484,1.78968) (-0.94596,1.89192) (-1.00000,2.00000) (-1.05713,2.11426) (-1.11752,2.23504) (-1.18136,2.36272) (-1.24885,2.49770) (-1.32019,2.64039) (-1.39561,2.79122) (-1.47534,2.95068) (-1.55962,3.11925) (-1.64872,3.29744) (-1.74291,3.48582)};\addplot[blue,no marks] coordinates {(3.28025,-2.76683) (2.85198,-2.30923) (2.46398,-1.89022) (2.11175,-1.50522) (1.79125,-1.15006) (1.49882,-0.82101) (1.23120,-0.51467) (0.98544,-0.22798) (0.75889,0.04185) (0.54913,0.29735) (0.35401,0.54083) (0.17156,0.77440) (0.00000,1.00000) (-0.16229,1.21942) (-0.31678,1.43430) (-0.46483,1.64619) (-0.60767,1.85652) (-0.74644,2.06663) (-0.88220,2.27781) (-1.01591,2.49126) (-1.14851,2.70813) (-1.28084,2.92956) (-1.41372,3.15663) (-1.54790,3.39038)};\addplot[blue,no marks] coordinates {(-3.28025,2.76683) (-2.85198,2.30923) (-2.46398,1.89022) (-2.11175,1.50522) (-1.79125,1.15006) (-1.49882,0.82101) (-1.23120,0.51467) (-0.98544,0.22798) (-0.75889,-0.04185) (-0.54913,-0.29735) (-0.35401,-0.54083) (-0.17156,-0.77440) (0.00000,-1.00000) (0.16229,-1.21942) (0.31678,-1.43430) (0.46483,-1.64619) (0.60767,-1.85652) (0.74644,-2.06663) (0.88220,-2.27781) (1.01591,-2.49126) (1.14851,-2.70813) (1.28084,-2.92956) (1.41372,-3.15663) (1.54790,-3.39038)};\end{axis}\end{tikzpicture}\end{center}
\par\noindent{\small 依据：官方译审并补证（AI辅助）。}
\end{solution}

\begin{exercise}\label{cw:rec22-2}
对$\dot x=4x+3y,\dot y=-6x-5y$，完成第1题同一(a)--(e)序列，包括初值$(1,0)^T$及所有同轨迹解。
\end{exercise}
\begin{solution}
(a) $B=\begin{pmatrix}4&3\\-6&-5\end{pmatrix}$，迹$-1$、行列式$-2$，仍有$p_B=\lambda^2+\lambda-2$。
(b) $\lambda_+=1,w_+=(1,-1)^T$；$\lambda_-=-2,w_-=(1,-2)^T$。
(c) 此时$y=-x$向外，$y=-2x$向内；每线仍为正半线、零点、负半线三轨迹。经过$p=cw_\pm\ne0$的所有解为$\ee^{\lambda_\pm(t-t_0)}p$。
(d) $(1,0)^T=2w_+-w_-$，故 $v_0(t)=(2\ee^t-\ee^{-2t},-2\ee^t+2\ee^{-2t})^T$；所有同轨迹解为$v_0(t-t_0)$。
(e) 通解$c_+\ee^tw_++c_-\ee^{-2t}w_-$，仍是鞍点，但两条特征线的稳定性交换；不能把整个向量场简单说成$-A$，因为两矩阵的迹相同。\begin{center}\begin{tikzpicture}\begin{axis}[width=.55\textwidth,height=3.8cm,axis lines=middle,title={稳定线$y=-2x$},xlabel={$x$},ylabel={$y$},xmin=-2.4,xmax=2.4,ymin=-2.4,ymax=2.4,clip=true]\addplot[blue,no marks] coordinates {(-0.82101,2.99764) (-0.51467,2.46241) (-0.22798,1.97089) (0.04185,1.51777) (0.29735,1.09826) (0.54083,0.70802) (0.77440,0.34312) (1.00000,0.00000) (1.21942,-0.32458) (1.43430,-0.63356) (1.64619,-0.92966) (1.85652,-1.21534) (2.06663,-1.49288) (2.27781,-1.76439) (2.49126,-2.03183) (2.70813,-2.29702) (2.92956,-2.56168) (3.15663,-2.82743) (3.39038,-3.09581)};\addplot[blue,no marks] coordinates {(0.36788,-0.36788) (0.38890,-0.38890) (0.41111,-0.41111) (0.43460,-0.43460) (0.45943,-0.45943) (0.48567,-0.48567) (0.51342,-0.51342) (0.54275,-0.54275) (0.57375,-0.57375) (0.60653,-0.60653) (0.64118,-0.64118) (0.67781,-0.67781) (0.71653,-0.71653) (0.75747,-0.75747) (0.80074,-0.80074) (0.84648,-0.84648) (0.89484,-0.89484) (0.94596,-0.94596) (1.00000,-1.00000) (1.05713,-1.05713) (1.11752,-1.11752) (1.18136,-1.18136) (1.24885,-1.24885) (1.32019,-1.32019) (1.39561,-1.39561) (1.47534,-1.47534) (1.55962,-1.55962) (1.64872,-1.64872) (1.74291,-1.74291) (1.84248,-1.84248) (1.94773,-1.94773) (2.05900,-2.05900) (2.17663,-2.17663) (2.30098,-2.30098) (2.43243,-2.43243) (2.57138,-2.57138) (2.71828,-2.71828) (2.87357,-2.87357) (3.03773,-3.03773) (3.21127,-3.21127) (3.39472,-3.39472)};\draw[blue,->] (axis cs:1.28403,-1.28403)--(axis cs:1.44773,-1.44773);\addplot[blue,no marks] coordinates {(-0.36788,0.36788) (-0.38890,0.38890) (-0.41111,0.41111) (-0.43460,0.43460) (-0.45943,0.45943) (-0.48567,0.48567) (-0.51342,0.51342) (-0.54275,0.54275) (-0.57375,0.57375) (-0.60653,0.60653) (-0.64118,0.64118) (-0.67781,0.67781) (-0.71653,0.71653) (-0.75747,0.75747) (-0.80074,0.80074) (-0.84648,0.84648) (-0.89484,0.89484) (-0.94596,0.94596) (-1.00000,1.00000) (-1.05713,1.05713) (-1.11752,1.11752) (-1.18136,1.18136) (-1.24885,1.24885) (-1.32019,1.32019) (-1.39561,1.39561) (-1.47534,1.47534) (-1.55962,1.55962) (-1.64872,1.64872) (-1.74291,1.74291) (-1.84248,1.84248) (-1.94773,1.94773) (-2.05900,2.05900) (-2.17663,2.17663) (-2.30098,2.30098) (-2.43243,2.43243) (-2.57138,2.57138) (-2.71828,2.71828) (-2.87357,2.87357) (-3.03773,3.03773) (-3.21127,3.21127) (-3.39472,3.39472)};\draw[blue,->] (axis cs:-1.28403,1.28403)--(axis cs:-1.44773,1.44773);\addplot[blue,no marks] coordinates {(1.74291,-3.48582) (1.55962,-3.11925) (1.39561,-2.79122) (1.24885,-2.49770) (1.11752,-2.23504) (1.00000,-2.00000) (0.89484,-1.78968) (0.80074,-1.60147) (0.71653,-1.43306) (0.64118,-1.28236) (0.57375,-1.14751) (0.51342,-1.02683) (0.45943,-0.91885) (0.41111,-0.82222) (0.36788,-0.73576) (0.32919,-0.65839) (0.29457,-0.58915) (0.26360,-0.52719) (0.23588,-0.47175) (0.21107,-0.42214) (0.18888,-0.37775) (0.16901,-0.33803) (0.15124,-0.30248) (0.13534,-0.27067) (0.12110,-0.24221) (0.10837,-0.21674) (0.09697,-0.19394) (0.08677,-0.17355) (0.07765,-0.15530) (0.06948,-0.13897) (0.06218,-0.12435) (0.05564,-0.11128) (0.04979,-0.09957) (0.04455,-0.08910) (0.03987,-0.07973) (0.03567,-0.07135) (0.03192,-0.06384) (0.02857,-0.05713) (0.02556,-0.05112) (0.02287,-0.04575) (0.02047,-0.04094) (0.01832,-0.03663)};\draw[blue,->] (axis cs:0.60653,-1.21306)--(axis cs:0.47711,-0.95423);\addplot[blue,no marks] coordinates {(-1.74291,3.48582) (-1.55962,3.11925) (-1.39561,2.79122) (-1.24885,2.49770) (-1.11752,2.23504) (-1.00000,2.00000) (-0.89484,1.78968) (-0.80074,1.60147) (-0.71653,1.43306) (-0.64118,1.28236) (-0.57375,1.14751) (-0.51342,1.02683) (-0.45943,0.91885) (-0.41111,0.82222) (-0.36788,0.73576) (-0.32919,0.65839) (-0.29457,0.58915) (-0.26360,0.52719) (-0.23588,0.47175) (-0.21107,0.42214) (-0.18888,0.37775) (-0.16901,0.33803) (-0.15124,0.30248) (-0.13534,0.27067) (-0.12110,0.24221) (-0.10837,0.21674) (-0.09697,0.19394) (-0.08677,0.17355) (-0.07765,0.15530) (-0.06948,0.13897) (-0.06218,0.12435) (-0.05564,0.11128) (-0.04979,0.09957) (-0.04455,0.08910) (-0.03987,0.07973) (-0.03567,0.07135) (-0.03192,0.06384) (-0.02857,0.05713) (-0.02556,0.05112) (-0.02287,0.04575) (-0.02047,0.04094) (-0.01832,0.03663)};\draw[blue,->] (axis cs:-0.60653,1.21306)--(axis cs:-0.47711,0.95423);\addplot[blue,no marks] coordinates {(-1.23120,3.17894) (-0.98544,2.72835) (-0.75889,2.31851) (-0.54913,1.94474) (-0.35401,1.60286) (-0.17156,1.28908) (0.00000,1.00000) (0.16229,0.73255) (0.31678,0.48396) (0.46483,0.25170) (0.60767,0.03351) (0.74644,-0.17269) (0.88220,-0.36878) (1.01591,-0.55649) (1.14851,-0.73740) (1.28084,-0.91296) (1.41372,-1.08452) (1.54790,-1.25333) (1.68414,-1.42054) (1.82313,-1.58725) (1.96556,-1.75449) (2.11210,-1.92322) (2.26341,-2.09440) (2.42014,-2.26890) (2.58295,-2.44761) (2.75247,-2.63136) (2.92936,-2.82100) (3.11430,-3.01733) (3.30795,-3.22117)};\draw[blue,->] (axis cs:0.67749,-0.07096)--(axis cs:0.97062,-0.49351);\addplot[blue,no marks] coordinates {(1.23120,-3.17894) (0.98544,-2.72835) (0.75889,-2.31851) (0.54913,-1.94474) (0.35401,-1.60286) (0.17156,-1.28908) (0.00000,-1.00000) (-0.16229,-0.73255) (-0.31678,-0.48396) (-0.46483,-0.25170) (-0.60767,-0.03351) (-0.74644,0.17269) (-0.88220,0.36878) (-1.01591,0.55649) (-1.14851,0.73740) (-1.28084,0.91296) (-1.41372,1.08452) (-1.54790,1.25333) (-1.68414,1.42054) (-1.82313,1.58725) (-1.96556,1.75449) (-2.11210,1.92322) (-2.26341,2.09440) (-2.42014,2.26890) (-2.58295,2.44761) (-2.75247,2.63136) (-2.92936,2.82100) (-3.11430,3.01733) (-3.30795,3.22117)};\draw[blue,->] (axis cs:-0.67749,0.07096)--(axis cs:-0.97062,0.49351);\end{axis}\end{tikzpicture}\end{center}
\par\noindent{\small 依据：官方译审并补证（AI辅助）。}
\end{solution}

\originalsection{第23次习题课：线性相图与迹--行列式平面（5月4日）}
以下固定参数族 $A(a)=\begin{pmatrix}a&2\\-2&-1\end{pmatrix}$。
\begin{exercise}\label{cw:rec23-1}
以$a$表示迹、行列式、特征多项式与特征值。
\end{exercise}
\begin{solution}
$\tau=a-1,\Delta=4-a$，$p_A=\lambda^2-(a-1)\lambda+4-a$，所以 $\lambda_\pm=[a-1\pm\sqrt{(a+5)(a-3)}]/2$。复根时平方根按共轭两支取值。
\par\noindent{\small 依据：官方译审并补证（AI辅助）。}
\end{solution}

\begin{exercise}\label{cw:rec23-2}
将行列式写成迹的函数，画迹--行列式平面、临界抛物线$\Delta=\tau^2/4$及本族曲线。标出$a=-6,-5,-2,1,2,3,4,5$的点。
\end{exercise}
\begin{solution}
$a=\tau+1$，故$\Delta=3-\tau$。点依次为$(-7,10),(-6,9),(-3,6),(0,3),(1,2),(2,1),(3,0),(4,-1)$。交点由$\tau^2/4=3-\tau$得$\tau=-6,2$，对应$a=-5,3$。\begin{center}\begin{tikzpicture}\begin{axis}[width=.68\textwidth,height=4.5cm,axis lines=middle,xlabel={$\tau$},ylabel={$\Delta$},samples=130,xmin=-8,xmax=5,ymin=-2,ymax=16]
\addplot[black,domain=-8:5]{x*x/4};\addplot[blue,domain=-8:5]{3-x};\addplot[only marks,red] coordinates{(-7,10)(-6,9)(-3,6)(0,3)(1,2)(2,1)(3,0)(4,-1)};\node[above left] at(axis cs:-7,10){$-6$};\node[below left] at(axis cs:-6,9){$-5$};\node[above] at(axis cs:-3,6){$-2$};\node[above right] at(axis cs:0,3){$1$};\node[above] at(axis cs:1,2){$2$};\node[below] at(axis cs:2,1){$3$};\node[above right] at(axis cs:3,0){$4$};\node[above right] at(axis cs:4,-1){$5$};\end{axis}\end{tikzpicture}\end{center}
\par\noindent{\small 依据：官方译审并补证（AI辅助）。}
\end{solution}

\begin{exercise}\label{cw:rec23-3}
对上述八个$a$列出：(1)特征值；(2)由特征值确定的相图种类及稳定类型：所有实部负称稳定，至少一个实部正称不稳定，否则不指定；(3)额外信息：螺旋的旋转方向，重根是否亏损。每例画小相图，传达分类即可，不要求猜准特征方向。
\end{exercise}
\begin{solution}
计算结果如下；“不指定”沿用本题措辞，中心实际为中性稳定。
\[\begin{array}{c|c|l}
a&\text{特征值}&\text{相图与稳定类型}\\\hline
-6&-5,-2&\text{稳定结点}\\
-5&-3,-3&\text{亏损稳定结点}\\
-2&(-3\pm\ii\sqrt{15})/2&\text{顺时针稳定螺旋}\\
1&\pm\ii\sqrt3&\text{顺时针中心；不指定}\\
2&(1\pm\ii\sqrt7)/2&\text{顺时针不稳定螺旋}\\
3&1,1&\text{亏损不稳定结点}\\
4&0,3&\text{一条平衡线与不稳定横向运动}\\
5&2\pm\sqrt5&\text{鞍点；不稳定}
\end{array}\]
非实根情形，在正$x$轴速度的$y$分量为$-2x<0$，故顺时针。重根处$A-\lambda I\ne0$且平方为0，所以只有一条特征方向，均亏损。$a=4$不是孤立结点：零特征空间由平衡点构成。以下用实际矩阵指数计算坐标重绘，相图类别由上述解析结果确定。\begin{center}\begin{tikzpicture}\begin{axis}[width=.28\textwidth,height=3.8cm,axis lines=middle,title={$a=-6$},xlabel={$x$},ylabel={$y$},xmin=-2.4,xmax=2.4,ymin=-2.4,ymax=2.4,clip=true]\addplot[blue,no marks] coordinates {(2.47250,0.55125) (1.71722,0.25689) (1.17962,0.06125) (0.79840,-0.06510) (0.52935,-0.14318) (0.34062,-0.18796) (0.20926,-0.21008) (0.11877,-0.21706) (0.05729,-0.21419) (0.01630,-0.20516) (-0.01029,-0.19252) (-0.02683,-0.17801) (-0.03644,-0.16280) (-0.04131,-0.14766) (-0.04302,-0.13307) (-0.04266,-0.11933) (-0.04099,-0.10658) (-0.03855,-0.09490) (-0.03571,-0.08428) (-0.03270,-0.07470) (-0.02968,-0.06610) (-0.02677,-0.05841) (-0.02402,-0.05156) (-0.02146,-0.04547) (-0.01912,-0.04008) (-0.01698,-0.03530) (-0.01506,-0.03108) (-0.01333,-0.02735) (-0.01178,-0.02406) (-0.01040,-0.02116) (-0.00917,-0.01861) (-0.00808,-0.01636) (-0.00712,-0.01438) (-0.00627,-0.01264) (-0.00552,-0.01111) (-0.00485,-0.00976) (-0.00427,-0.00858) (-0.00375,-0.00754) (-0.00330,-0.00662) (-0.00290,-0.00582) (-0.00255,-0.00511) (-0.00224,-0.00449)};\draw[blue,->] (axis cs:0.17983,-0.21335)--(axis cs:0.05061,-0.21325);\addplot[blue,no marks] coordinates {(-2.47250,-0.55125) (-1.71722,-0.25689) (-1.17962,-0.06125) (-0.79840,0.06510) (-0.52935,0.14318) (-0.34062,0.18796) (-0.20926,0.21008) (-0.11877,0.21706) (-0.05729,0.21419) (-0.01630,0.20516) (0.01029,0.19252) (0.02683,0.17801) (0.03644,0.16280) (0.04131,0.14766) (0.04302,0.13307) (0.04266,0.11933) (0.04099,0.10658) (0.03855,0.09490) (0.03571,0.08428) (0.03270,0.07470) (0.02968,0.06610) (0.02677,0.05841) (0.02402,0.05156) (0.02146,0.04547) (0.01912,0.04008) (0.01698,0.03530) (0.01506,0.03108) (0.01333,0.02735) (0.01178,0.02406) (0.01040,0.02116) (0.00917,0.01861) (0.00808,0.01636) (0.00712,0.01438) (0.00627,0.01264) (0.00552,0.01111) (0.00485,0.00976) (0.00427,0.00858) (0.00375,0.00754) (0.00330,0.00662) (0.00290,0.00582) (0.00255,0.00511) (0.00224,0.00449)};\draw[blue,->] (axis cs:-0.17983,0.21335)--(axis cs:-0.05061,0.21325);\addplot[blue,no marks] coordinates {(-2.52461,0.75892) (-1.61662,0.96717) (-0.98541,1.06692) (-0.55125,1.09438) (-0.25689,1.07500) (-0.06125,1.02650) (0.06510,0.96115) (0.14318,0.88729) (0.18796,0.81051) (0.21008,0.73445) (0.21706,0.66142) (0.21419,0.59277) (0.20516,0.52921) (0.19252,0.47102) (0.17801,0.41820) (0.16280,0.37057) (0.14766,0.32784) (0.13307,0.28967) (0.11933,0.25567) (0.10658,0.22546) (0.09490,0.19869) (0.08428,0.17499) (0.07470,0.15405) (0.06610,0.13556) (0.05841,0.11925) (0.05156,0.10488) (0.04547,0.09222) (0.04008,0.08107) (0.03530,0.07126) (0.03108,0.06263) (0.02735,0.05504) (0.02406,0.04837) (0.02116,0.04250) (0.01861,0.03734) (0.01636,0.03281) (0.01438,0.02883) (0.01264,0.02533) (0.01111,0.02225) (0.00976,0.01955) (0.00858,0.01717) (0.00754,0.01508) (0.00662,0.01325) (0.00582,0.01164) (0.00511,0.01023) (0.00449,0.00898)};\draw[blue,->] (axis cs:0.21335,0.71321)--(axis cs:0.21325,0.58374);\addplot[blue,no marks] coordinates {(2.52461,-0.75892) (1.61662,-0.96717) (0.98541,-1.06692) (0.55125,-1.09438) (0.25689,-1.07500) (0.06125,-1.02650) (-0.06510,-0.96115) (-0.14318,-0.88729) (-0.18796,-0.81051) (-0.21008,-0.73445) (-0.21706,-0.66142) (-0.21419,-0.59277) (-0.20516,-0.52921) (-0.19252,-0.47102) (-0.17801,-0.41820) (-0.16280,-0.37057) (-0.14766,-0.32784) (-0.13307,-0.28967) (-0.11933,-0.25567) (-0.10658,-0.22546) (-0.09490,-0.19869) (-0.08428,-0.17499) (-0.07470,-0.15405) (-0.06610,-0.13556) (-0.05841,-0.11925) (-0.05156,-0.10488) (-0.04547,-0.09222) (-0.04008,-0.08107) (-0.03530,-0.07126) (-0.03108,-0.06263) (-0.02735,-0.05504) (-0.02406,-0.04837) (-0.02116,-0.04250) (-0.01861,-0.03734) (-0.01636,-0.03281) (-0.01438,-0.02883) (-0.01264,-0.02533) (-0.01111,-0.02225) (-0.00976,-0.01955) (-0.00858,-0.01717) (-0.00754,-0.01508) (-0.00662,-0.01325) (-0.00582,-0.01164) (-0.00511,-0.01023) (-0.00449,-0.00898)};\draw[blue,->] (axis cs:-0.21335,-0.71321)--(axis cs:-0.21325,-0.58374);\addplot[blue,no marks] coordinates {(3.39211,2.58380) (2.54503,2.05233) (1.92125,1.64563) (1.46033,1.33189) (1.11837,1.08775) (0.86350,0.89605) (0.67253,0.74412) (0.52858,0.62255) (0.41934,0.52438) (0.33583,0.44436) (0.27148,0.37858) (0.22146,0.32404) (0.18223,0.27850) (0.15118,0.24019) (0.12637,0.20777) (0.10635,0.18018) (0.09006,0.15659) (0.07667,0.13634) (0.06559,0.11888) (0.05634,0.10379) (0.04857,0.09071) (0.04200,0.07935) (0.03641,0.06946) (0.03164,0.06084) (0.02754,0.05332) (0.02401,0.04675) (0.02096,0.04100) (0.01832,0.03597) (0.01602,0.03156) (0.01402,0.02770) (0.01228,0.02431) (0.01076,0.02134) (0.00943,0.01874) (0.00827,0.01645) (0.00726,0.01445) (0.00637,0.01269) (0.00559,0.01114) (0.00491,0.00979) (0.00431,0.00859) (0.00378,0.00755) (0.00332,0.00663) (0.00292,0.00582) (0.00256,0.00512) (0.00225,0.00449)};\draw[blue,->] (axis cs:0.39318,0.49986)--(axis cs:0.26386,0.37049);\addplot[blue,no marks] coordinates {(-3.39211,-2.58380) (-2.54503,-2.05233) (-1.92125,-1.64563) (-1.46033,-1.33189) (-1.11837,-1.08775) (-0.86350,-0.89605) (-0.67253,-0.74412) (-0.52858,-0.62255) (-0.41934,-0.52438) (-0.33583,-0.44436) (-0.27148,-0.37858) (-0.22146,-0.32404) (-0.18223,-0.27850) (-0.15118,-0.24019) (-0.12637,-0.20777) (-0.10635,-0.18018) (-0.09006,-0.15659) (-0.07667,-0.13634) (-0.06559,-0.11888) (-0.05634,-0.10379) (-0.04857,-0.09071) (-0.04200,-0.07935) (-0.03641,-0.06946) (-0.03164,-0.06084) (-0.02754,-0.05332) (-0.02401,-0.04675) (-0.02096,-0.04100) (-0.01832,-0.03597) (-0.01602,-0.03156) (-0.01402,-0.02770) (-0.01228,-0.02431) (-0.01076,-0.02134) (-0.00943,-0.01874) (-0.00827,-0.01645) (-0.00726,-0.01445) (-0.00637,-0.01269) (-0.00559,-0.01114) (-0.00491,-0.00979) (-0.00431,-0.00859) (-0.00378,-0.00755) (-0.00332,-0.00663) (-0.00292,-0.00582) (-0.00256,-0.00512) (-0.00225,-0.00449)};\draw[blue,->] (axis cs:-0.39318,-0.49986)--(axis cs:-0.26386,-0.37049);\end{axis}\end{tikzpicture}\hspace{.035\textwidth}\begin{tikzpicture}\begin{axis}[width=.28\textwidth,height=3.8cm,axis lines=middle,title={$a=-5$},xlabel={$x$},ylabel={$y$},xmin=-2.4,xmax=2.4,ymin=-2.4,ymax=2.4,clip=true]\addplot[blue,no marks] coordinates {(2.81339,0.86566) (2.10838,0.50482) (1.56467,0.24448) (1.14729,0.06038) (0.82855,-0.06628) (0.58664,-0.15007) (0.40435,-0.20218) (0.26817,-0.23118) (0.16749,-0.24362) (0.09402,-0.24445) (0.04128,-0.23737) (0.00425,-0.22517) (-0.02099,-0.20986) (-0.03744,-0.19293) (-0.04742,-0.17544) (-0.05270,-0.15810) (-0.05464,-0.14141) (-0.05424,-0.12568) (-0.05228,-0.11110) (-0.04932,-0.09774) (-0.04577,-0.08564) (-0.04194,-0.07476) (-0.03803,-0.06505) (-0.03419,-0.05644) (-0.03053,-0.04884) (-0.02709,-0.04217) (-0.02391,-0.03632) (-0.02101,-0.03123) (-0.01839,-0.02680) (-0.01604,-0.02296) (-0.01394,-0.01965) (-0.01209,-0.01678) (-0.01045,-0.01432) (-0.00902,-0.01220) (-0.00776,-0.01038) (-0.00667,-0.00883) (-0.00572,-0.00750) (-0.00490,-0.00636) (-0.00419,-0.00540) (-0.00358,-0.00457) (-0.00305,-0.00387) (-0.00260,-0.00327) (-0.00221,-0.00277)};\draw[blue,->] (axis cs:0.23618,-0.23618)--(axis cs:0.08569,-0.24387);\addplot[blue,no marks] coordinates {(-2.81339,-0.86566) (-2.10838,-0.50482) (-1.56467,-0.24448) (-1.14729,-0.06038) (-0.82855,0.06628) (-0.58664,0.15007) (-0.40435,0.20218) (-0.26817,0.23118) (-0.16749,0.24362) (-0.09402,0.24445) (-0.04128,0.23737) (-0.00425,0.22517) (0.02099,0.20986) (0.03744,0.19293) (0.04742,0.17544) (0.05270,0.15810) (0.05464,0.14141) (0.05424,0.12568) (0.05228,0.11110) (0.04932,0.09774) (0.04577,0.08564) (0.04194,0.07476) (0.03803,0.06505) (0.03419,0.05644) (0.03053,0.04884) (0.02709,0.04217) (0.02391,0.03632) (0.02101,0.03123) (0.01839,0.02680) (0.01604,0.02296) (0.01394,0.01965) (0.01209,0.01678) (0.01045,0.01432) (0.00902,0.01220) (0.00776,0.01038) (0.00667,0.00883) (0.00572,0.00750) (0.00490,0.00636) (0.00419,0.00540) (0.00358,0.00457) (0.00305,0.00387) (0.00260,0.00327) (0.00221,0.00277)};\draw[blue,->] (axis cs:-0.23618,0.23618)--(axis cs:-0.08569,0.24387);\addplot[blue,no marks] coordinates {(-2.90862,0.58172) (-2.02214,0.85143) (-1.35814,1.00765) (-0.86566,1.08207) (-0.50482,1.09873) (-0.24448,1.07571) (-0.06038,1.02652) (0.06628,0.96112) (0.15007,0.88679) (0.20218,0.80871) (0.23118,0.73053) (0.24362,0.65473) (0.24445,0.58291) (0.23737,0.51603) (0.22517,0.45458) (0.20986,0.39874) (0.19293,0.34843) (0.17544,0.30346) (0.15810,0.26350) (0.14141,0.22818) (0.12568,0.19712) (0.11110,0.16991) (0.09774,0.14617) (0.08564,0.12551) (0.07476,0.10758) (0.06505,0.09207) (0.05644,0.07869) (0.04884,0.06716) (0.04217,0.05724) (0.03632,0.04874) (0.03123,0.04145) (0.02680,0.03522) (0.02296,0.02989) (0.01965,0.02535) (0.01678,0.02148) (0.01432,0.01818) (0.01220,0.01538) (0.01038,0.01300) (0.00883,0.01099) (0.00750,0.00928) (0.00636,0.00783) (0.00540,0.00660) (0.00457,0.00556) (0.00387,0.00468) (0.00327,0.00394) (0.00277,0.00332)};\draw[blue,->] (axis cs:0.23618,0.70855)--(axis cs:0.24387,0.57343);\addplot[blue,no marks] coordinates {(2.90862,-0.58172) (2.02214,-0.85143) (1.35814,-1.00765) (0.86566,-1.08207) (0.50482,-1.09873) (0.24448,-1.07571) (0.06038,-1.02652) (-0.06628,-0.96112) (-0.15007,-0.88679) (-0.20218,-0.80871) (-0.23118,-0.73053) (-0.24362,-0.65473) (-0.24445,-0.58291) (-0.23737,-0.51603) (-0.22517,-0.45458) (-0.20986,-0.39874) (-0.19293,-0.34843) (-0.17544,-0.30346) (-0.15810,-0.26350) (-0.14141,-0.22818) (-0.12568,-0.19712) (-0.11110,-0.16991) (-0.09774,-0.14617) (-0.08564,-0.12551) (-0.07476,-0.10758) (-0.06505,-0.09207) (-0.05644,-0.07869) (-0.04884,-0.06716) (-0.04217,-0.05724) (-0.03632,-0.04874) (-0.03123,-0.04145) (-0.02680,-0.03522) (-0.02296,-0.02989) (-0.01965,-0.02535) (-0.01678,-0.02148) (-0.01432,-0.01818) (-0.01220,-0.01538) (-0.01038,-0.01300) (-0.00883,-0.01099) (-0.00750,-0.00928) (-0.00636,-0.00783) (-0.00540,-0.00660) (-0.00457,-0.00556) (-0.00387,-0.00468) (-0.00327,-0.00394) (-0.00277,-0.00332)};\draw[blue,->] (axis cs:-0.23618,-0.70855)--(axis cs:-0.24387,-0.57343);\addplot[blue,no marks] coordinates {(3.49034,3.49034) (2.87357,2.87357) (2.36579,2.36579) (1.94773,1.94773) (1.60355,1.60355) (1.32019,1.32019) (1.08690,1.08690) (0.89484,0.89484) (0.73671,0.73671) (0.60653,0.60653) (0.49935,0.49935) (0.41111,0.41111) (0.33847,0.33847) (0.27866,0.27866) (0.22942,0.22942) (0.18888,0.18888) (0.15550,0.15550) (0.12802,0.12802) (0.10540,0.10540) (0.08677,0.08677) (0.07144,0.07144) (0.05882,0.05882) (0.04842,0.04842) (0.03987,0.03987) (0.03282,0.03282) (0.02702,0.02702) (0.02225,0.02225) (0.01832,0.01832) (0.01508,0.01508) (0.01241,0.01241) (0.01022,0.01022) (0.00841,0.00841) (0.00693,0.00693) (0.00570,0.00570) (0.00470,0.00470) (0.00387,0.00387) (0.00318,0.00318) (0.00262,0.00262) (0.00216,0.00216) (0.00178,0.00178) (0.00146,0.00146) (0.00120,0.00120) (0.00099,0.00099) (0.00082,0.00082) (0.00067,0.00067) (0.00055,0.00055)};\draw[blue,->] (axis cs:0.47237,0.47237)--(axis cs:0.32956,0.32956);\addplot[blue,no marks] coordinates {(-3.49034,-3.49034) (-2.87357,-2.87357) (-2.36579,-2.36579) (-1.94773,-1.94773) (-1.60355,-1.60355) (-1.32019,-1.32019) (-1.08690,-1.08690) (-0.89484,-0.89484) (-0.73671,-0.73671) (-0.60653,-0.60653) (-0.49935,-0.49935) (-0.41111,-0.41111) (-0.33847,-0.33847) (-0.27866,-0.27866) (-0.22942,-0.22942) (-0.18888,-0.18888) (-0.15550,-0.15550) (-0.12802,-0.12802) (-0.10540,-0.10540) (-0.08677,-0.08677) (-0.07144,-0.07144) (-0.05882,-0.05882) (-0.04842,-0.04842) (-0.03987,-0.03987) (-0.03282,-0.03282) (-0.02702,-0.02702) (-0.02225,-0.02225) (-0.01832,-0.01832) (-0.01508,-0.01508) (-0.01241,-0.01241) (-0.01022,-0.01022) (-0.00841,-0.00841) (-0.00693,-0.00693) (-0.00570,-0.00570) (-0.00470,-0.00470) (-0.00387,-0.00387) (-0.00318,-0.00318) (-0.00262,-0.00262) (-0.00216,-0.00216) (-0.00178,-0.00178) (-0.00146,-0.00146) (-0.00120,-0.00120) (-0.00099,-0.00099) (-0.00082,-0.00082) (-0.00067,-0.00067) (-0.00055,-0.00055)};\draw[blue,->] (axis cs:-0.47237,-0.47237)--(axis cs:-0.32956,-0.32956);\end{axis}\end{tikzpicture}\hspace{.035\textwidth}\begin{tikzpicture}\begin{axis}[width=.28\textwidth,height=3.8cm,axis lines=middle,title={$a=-2$},xlabel={$x$},ylabel={$y$},xmin=-2.4,xmax=2.4,ymin=-2.4,ymax=2.4,clip=true]\addplot[blue,no marks] coordinates {(0.90068,3.45749) (1.19000,3.10823) (1.40098,2.74964) (1.54267,2.39156) (1.62405,2.04209) (1.65393,1.70770) (1.64072,1.39335) (1.59234,1.10271) (1.51611,0.83822) (1.41870,0.60130) (1.30607,0.39251) (1.18349,0.21164) (1.05551,0.05789) (0.92602,-0.07001) (0.79825,-0.17371) (0.67481,-0.25512) (0.55776,-0.31631) (0.44864,-0.35945) (0.34855,-0.38676) (0.25817,-0.40040) (0.17787,-0.40247) (0.10769,-0.39498) (0.04745,-0.37979) (-0.00323,-0.35859) (-0.04488,-0.33295) (-0.07814,-0.30423) (-0.10374,-0.27363) (-0.12245,-0.24218) (-0.13505,-0.21075) (-0.14234,-0.18006) (-0.14508,-0.15067) (-0.14403,-0.12304) (-0.13987,-0.09747) (-0.13324,-0.07420) (-0.12474,-0.05334) (-0.11489,-0.03495) (-0.10416,-0.01901) (-0.09294,-0.00545) (-0.08158,0.00584) (-0.07036,0.01500) (-0.05952,0.02219) (-0.04923,0.02761) (-0.03964,0.03144) (-0.03083,0.03388) (-0.02287,0.03511) (-0.01580,0.03532) (-0.00962,0.03468) (-0.00431,0.03337) (0.00016,0.03152) (0.00384,0.02928) (0.00678,0.02677) (0.00904,0.02409)};\draw[blue,->] (axis cs:0.52571,-0.33038)--(axis cs:0.33557,-0.38939);\addplot[blue,no marks] coordinates {(-0.90068,-3.45749) (-1.19000,-3.10823) (-1.40098,-2.74964) (-1.54267,-2.39156) (-1.62405,-2.04209) (-1.65393,-1.70770) (-1.64072,-1.39335) (-1.59234,-1.10271) (-1.51611,-0.83822) (-1.41870,-0.60130) (-1.30607,-0.39251) (-1.18349,-0.21164) (-1.05551,-0.05789) (-0.92602,0.07001) (-0.79825,0.17371) (-0.67481,0.25512) (-0.55776,0.31631) (-0.44864,0.35945) (-0.34855,0.38676) (-0.25817,0.40040) (-0.17787,0.40247) (-0.10769,0.39498) (-0.04745,0.37979) (0.00323,0.35859) (0.04488,0.33295) (0.07814,0.30423) (0.10374,0.27363) (0.12245,0.24218) (0.13505,0.21075) (0.14234,0.18006) (0.14508,0.15067) (0.14403,0.12304) (0.13987,0.09747) (0.13324,0.07420) (0.12474,0.05334) (0.11489,0.03495) (0.10416,0.01901) (0.09294,0.00545) (0.08158,-0.00584) (0.07036,-0.01500) (0.05952,-0.02219) (0.04923,-0.02761) (0.03964,-0.03144) (0.03083,-0.03388) (0.02287,-0.03511) (0.01580,-0.03532) (0.00962,-0.03468) (0.00431,-0.03337) (-0.00016,-0.03152) (-0.00384,-0.02928) (-0.00678,-0.02677) (-0.00904,-0.02409)};\draw[blue,->] (axis cs:-0.52571,0.33038)--(axis cs:-0.33557,0.38939);\addplot[blue,no marks] coordinates {(-3.45749,-0.82807) (-3.10823,-0.36412) (-2.74964,0.02616) (-2.39156,0.34689) (-2.04209,0.60300) (-1.70770,0.80008) (-1.39335,0.94404) (-1.10271,1.04098) (-0.83822,1.09700) (-0.60130,1.11804) (-0.39251,1.10981) (-0.21164,1.07767) (-0.05789,1.02657) (0.07001,0.96103) (0.17371,0.88511) (0.25512,0.80237) (0.31631,0.71591) (0.35945,0.62836) (0.38676,0.54192) (0.40040,0.45837) (0.40247,0.37911) (0.39498,0.30518) (0.37979,0.23734) (0.35859,0.17607) (0.33295,0.12160) (0.30423,0.07397) (0.27363,0.03307) (0.24218,-0.00136) (0.21075,-0.02967) (0.18006,-0.05231) (0.15067,-0.06975) (0.12304,-0.08251) (0.09747,-0.09113) (0.07420,-0.09614) (0.05334,-0.09807) (0.03495,-0.09742) (0.01901,-0.09466) (0.00545,-0.09022) (-0.00584,-0.08450) (-0.01500,-0.07786) (-0.02219,-0.07062) (-0.02761,-0.06304) (-0.03144,-0.05536) (-0.03388,-0.04777) (-0.03511,-0.04043) (-0.03532,-0.03346) (-0.03468,-0.02696) (-0.03337,-0.02099) (-0.03152,-0.01560) (-0.02928,-0.01080) (-0.02677,-0.00661) (-0.02409,-0.00300)};\draw[blue,->] (axis cs:0.33038,0.69090)--(axis cs:0.38939,0.53026);\addplot[blue,no marks] coordinates {(3.45749,0.82807) (3.10823,0.36412) (2.74964,-0.02616) (2.39156,-0.34689) (2.04209,-0.60300) (1.70770,-0.80008) (1.39335,-0.94404) (1.10271,-1.04098) (0.83822,-1.09700) (0.60130,-1.11804) (0.39251,-1.10981) (0.21164,-1.07767) (0.05789,-1.02657) (-0.07001,-0.96103) (-0.17371,-0.88511) (-0.25512,-0.80237) (-0.31631,-0.71591) (-0.35945,-0.62836) (-0.38676,-0.54192) (-0.40040,-0.45837) (-0.40247,-0.37911) (-0.39498,-0.30518) (-0.37979,-0.23734) (-0.35859,-0.17607) (-0.33295,-0.12160) (-0.30423,-0.07397) (-0.27363,-0.03307) (-0.24218,0.00136) (-0.21075,0.02967) (-0.18006,0.05231) (-0.15067,0.06975) (-0.12304,0.08251) (-0.09747,0.09113) (-0.07420,0.09614) (-0.05334,0.09807) (-0.03495,0.09742) (-0.01901,0.09466) (-0.00545,0.09022) (0.00584,0.08450) (0.01500,0.07786) (0.02219,0.07062) (0.02761,0.06304) (0.03144,0.05536) (0.03388,0.04777) (0.03511,0.04043) (0.03532,0.03346) (0.03468,0.02696) (0.03337,0.02099) (0.03152,0.01560) (0.02928,0.01080) (0.02677,0.00661) (0.02409,0.00300)};\draw[blue,->] (axis cs:-0.33038,-0.69090)--(axis cs:-0.38939,-0.53026);\addplot[blue,no marks] coordinates {(-3.26146,2.41712) (-2.55681,2.62943) (-1.91824,2.74411) (-1.34866,2.77581) (-0.84890,2.73845) (-0.41804,2.64509) (-0.05377,2.50777) (0.24736,2.33740) (0.48963,2.14369) (0.67789,1.93522) (0.81739,1.71935) (0.91356,1.50232) (0.97185,1.28930) (0.99762,1.08446) (0.99604,0.89102) (0.97196,0.71140) (0.92993,0.54725) (0.87406,0.39960) (0.80809,0.26891) (0.73530,0.15517) (0.65857,0.05798) (0.58034,-0.02337) (0.50267,-0.08980) (0.42724,-0.14244) (0.35537,-0.18253) (0.28807,-0.21136) (0.22608,-0.23026) (0.16988,-0.24056) (0.11973,-0.24354) (0.07570,-0.24042) (0.03772,-0.23237) (0.00559,-0.22042) (-0.02099,-0.20555) (-0.04239,-0.18860) (-0.05904,-0.17034) (-0.07140,-0.15141) (-0.07995,-0.13237) (-0.08515,-0.11366) (-0.08749,-0.09567) (-0.08742,-0.07866) (-0.08536,-0.06287) (-0.08171,-0.04842) (-0.07684,-0.03542) (-0.07108,-0.02391) (-0.06471,-0.01389) (-0.05798,-0.00532) (-0.05112,0.00186) (-0.04430,0.00772) (-0.03768,0.01238) (-0.03136,0.01592) (-0.02544,0.01848) (-0.01999,0.02016) (-0.01504,0.02109)};\draw[blue,->] (axis cs:0.85609,0.36053)--(axis cs:0.72495,0.14087);\addplot[blue,no marks] coordinates {(3.26146,-2.41712) (2.55681,-2.62943) (1.91824,-2.74411) (1.34866,-2.77581) (0.84890,-2.73845) (0.41804,-2.64509) (0.05377,-2.50777) (-0.24736,-2.33740) (-0.48963,-2.14369) (-0.67789,-1.93522) (-0.81739,-1.71935) (-0.91356,-1.50232) (-0.97185,-1.28930) (-0.99762,-1.08446) (-0.99604,-0.89102) (-0.97196,-0.71140) (-0.92993,-0.54725) (-0.87406,-0.39960) (-0.80809,-0.26891) (-0.73530,-0.15517) (-0.65857,-0.05798) (-0.58034,0.02337) (-0.50267,0.08980) (-0.42724,0.14244) (-0.35537,0.18253) (-0.28807,0.21136) (-0.22608,0.23026) (-0.16988,0.24056) (-0.11973,0.24354) (-0.07570,0.24042) (-0.03772,0.23237) (-0.00559,0.22042) (0.02099,0.20555) (0.04239,0.18860) (0.05904,0.17034) (0.07140,0.15141) (0.07995,0.13237) (0.08515,0.11366) (0.08749,0.09567) (0.08742,0.07866) (0.08536,0.06287) (0.08171,0.04842) (0.07684,0.03542) (0.07108,0.02391) (0.06471,0.01389) (0.05798,0.00532) (0.05112,-0.00186) (0.04430,-0.00772) (0.03768,-0.01238) (0.03136,-0.01592) (0.02544,-0.01848) (0.01999,-0.02016) (0.01504,-0.02109)};\draw[blue,->] (axis cs:-0.85609,-0.36053)--(axis cs:-0.72495,-0.14087);\end{axis}\end{tikzpicture}\par\begin{tikzpicture}\begin{axis}[width=.28\textwidth,height=3.8cm,axis lines=middle,title={$a=1$},xlabel={$x$},ylabel={$y$},xmin=-2.4,xmax=2.4,ymin=-2.4,ymax=2.4,clip=true]\addplot[blue,no marks] coordinates {(-0.73042,1.13972) (-0.62563,1.15332) (-0.51297,1.15239) (-0.39385,1.13696) (-0.26977,1.10721) (-0.14229,1.06352) (-0.01302,1.00645) (0.11641,0.93670) (0.24437,0.85516) (0.36926,0.76286) (0.48950,0.66095) (0.60358,0.55072) (0.71006,0.43355) (0.80760,0.31093) (0.89497,0.18439) (0.97108,0.05553) (1.03495,-0.07402) (1.08580,-0.20265) (1.12298,-0.32872) (1.14602,-0.45066) (1.15464,-0.56692) (1.14871,-0.67605) (1.12833,-0.77666) (1.09374,-0.86749) (1.04538,-0.94741) (0.98386,-1.01540) (0.90995,-1.07060) (0.82459,-1.11232) (0.72884,-1.14004) (0.62392,-1.15341) (0.51115,-1.15226) (0.39194,-1.13660) (0.26779,-1.10663) (0.14027,-1.06273) (0.01099,-1.00545) (-0.11843,-0.93551) (-0.24636,-0.85380) (-0.37119,-0.76133) (-0.49134,-0.65928) (-0.60531,-0.54893) (-0.71166,-0.43167) (-0.80905,-0.30897) (-0.89625,-0.18239) (-0.97217,-0.05351) (-1.03585,0.07605) (-1.08649,0.20465) (-1.12345,0.33067) (-1.14627,0.45253) (-1.15466,0.56869) (-1.14851,0.67769) (-1.12790,0.77816) (-1.09309,0.86883) (-1.04452,0.94857) (-0.98279,1.01636) (-0.90870,1.07136)};\draw[blue,->] (axis cs:1.14997,-0.48452)--(axis cs:1.14677,-0.69038);\addplot[blue,no marks] coordinates {(0.73042,-1.13972) (0.62563,-1.15332) (0.51297,-1.15239) (0.39385,-1.13696) (0.26977,-1.10721) (0.14229,-1.06352) (0.01302,-1.00645) (-0.11641,-0.93670) (-0.24437,-0.85516) (-0.36926,-0.76286) (-0.48950,-0.66095) (-0.60358,-0.55072) (-0.71006,-0.43355) (-0.80760,-0.31093) (-0.89497,-0.18439) (-0.97108,-0.05553) (-1.03495,0.07402) (-1.08580,0.20265) (-1.12298,0.32872) (-1.14602,0.45066) (-1.15464,0.56692) (-1.14871,0.67605) (-1.12833,0.77666) (-1.09374,0.86749) (-1.04538,0.94741) (-0.98386,1.01540) (-0.90995,1.07060) (-0.82459,1.11232) (-0.72884,1.14004) (-0.62392,1.15341) (-0.51115,1.15226) (-0.39194,1.13660) (-0.26779,1.10663) (-0.14027,1.06273) (-0.01099,1.00545) (0.11843,0.93551) (0.24636,0.85380) (0.37119,0.76133) (0.49134,0.65928) (0.60531,0.54893) (0.71166,0.43167) (0.80905,0.30897) (0.89625,0.18239) (0.97217,0.05351) (1.03585,-0.07605) (1.08649,-0.20465) (1.12345,-0.33067) (1.14627,-0.45253) (1.15466,-0.56869) (1.14851,-0.67769) (1.12790,-0.77816) (1.09309,-0.86883) (1.04452,-0.94857) (0.98279,-1.01636) (0.90870,-1.07136)};\draw[blue,->] (axis cs:-1.14997,0.48452)--(axis cs:-1.14677,0.69038);\addplot[blue,no marks] coordinates {(-1.13972,0.40930) (-1.15332,0.52769) (-1.15239,0.63942) (-1.13696,0.74311) (-1.10721,0.83744) (-1.06352,0.92123) (-1.00645,0.99343) (-0.93670,1.05311) (-0.85516,1.09954) (-0.76286,1.13212) (-0.66095,1.15045) (-0.55072,1.15430) (-0.43355,1.14361) (-0.31093,1.11853) (-0.18439,1.07936) (-0.05553,1.02661) (0.07402,0.96093) (0.20265,0.88316) (0.32872,0.79426) (0.45066,0.69537) (0.56692,0.58772) (0.67605,0.47267) (0.77666,0.35167) (0.86749,0.22624) (0.94741,0.09797) (1.01540,-0.03154) (1.07060,-0.16065) (1.11232,-0.28774) (1.14004,-0.41120) (1.15341,-0.52949) (1.15226,-0.64111) (1.13660,-0.74466) (1.10663,-0.83884) (1.06273,-0.92246) (1.00545,-0.99446) (0.93551,-1.05394) (0.85380,-1.10015) (0.76133,-1.13252) (0.65928,-1.15062) (0.54893,-1.15424) (0.43167,-1.14333) (0.30897,-1.11802) (0.18239,-1.07864) (0.05351,-1.02568) (-0.07605,-0.95980) (-0.20465,-0.88185) (-0.33067,-0.79279) (-0.45253,-0.69374) (-0.56869,-0.58597) (-0.67769,-0.47082) (-0.77816,-0.34974) (-0.86883,-0.22425) (-0.94857,-0.09595) (-1.01636,0.03357) (-1.07136,0.16266)};\draw[blue,->] (axis cs:0.48452,0.66545)--(axis cs:0.69038,0.45639);\addplot[blue,no marks] coordinates {(1.13972,-0.40930) (1.15332,-0.52769) (1.15239,-0.63942) (1.13696,-0.74311) (1.10721,-0.83744) (1.06352,-0.92123) (1.00645,-0.99343) (0.93670,-1.05311) (0.85516,-1.09954) (0.76286,-1.13212) (0.66095,-1.15045) (0.55072,-1.15430) (0.43355,-1.14361) (0.31093,-1.11853) (0.18439,-1.07936) (0.05553,-1.02661) (-0.07402,-0.96093) (-0.20265,-0.88316) (-0.32872,-0.79426) (-0.45066,-0.69537) (-0.56692,-0.58772) (-0.67605,-0.47267) (-0.77666,-0.35167) (-0.86749,-0.22624) (-0.94741,-0.09797) (-1.01540,0.03154) (-1.07060,0.16065) (-1.11232,0.28774) (-1.14004,0.41120) (-1.15341,0.52949) (-1.15226,0.64111) (-1.13660,0.74466) (-1.10663,0.83884) (-1.06273,0.92246) (-1.00545,0.99446) (-0.93551,1.05394) (-0.85380,1.10015) (-0.76133,1.13252) (-0.65928,1.15062) (-0.54893,1.15424) (-0.43167,1.14333) (-0.30897,1.11802) (-0.18239,1.07864) (-0.05351,1.02568) (0.07605,0.95980) (0.20465,0.88185) (0.33067,0.79279) (0.45253,0.69374) (0.56869,0.58597) (0.67769,0.47082) (0.77816,0.34974) (0.86883,0.22425) (0.94857,0.09595) (1.01636,-0.03357) (1.07136,-0.16266)};\draw[blue,->] (axis cs:-0.48452,-0.66545)--(axis cs:-0.69038,-0.45639);\addplot[blue,no marks] coordinates {(-1.87014,1.54902) (-1.77894,1.68100) (-1.66536,1.79181) (-1.53080,1.88007) (-1.37698,1.94465) (-1.20581,1.98476) (-1.01947,1.99987) (-0.82029,1.98981) (-0.61079,1.95470) (-0.39359,1.89498) (-0.17144,1.81140) (0.05287,1.70501) (0.27651,1.57716) (0.49667,1.42946) (0.71058,1.26376) (0.91554,1.08214) (1.10898,0.88691) (1.28845,0.68051) (1.45171,0.46554) (1.59668,0.24471) (1.72156,0.02080) (1.82476,-0.20338) (1.90499,-0.42499) (1.96123,-0.64125) (1.99279,-0.84944) (1.99925,-1.04693) (1.98055,-1.23125) (1.93691,-1.40006) (1.86889,-1.55125) (1.77733,-1.68290) (1.66341,-1.79337) (1.52854,-1.88127) (1.37442,-1.94547) (1.20300,-1.98519) (1.01644,-1.99991) (0.81708,-1.98945) (0.60744,-1.95395) (0.39014,-1.89385) (0.16794,-1.80990) (-0.05638,-1.70317) (-0.27999,-1.57500) (-0.50008,-1.42700) (-0.71387,-1.26103) (-0.91867,-1.07918) (-1.11190,-0.88375) (-1.29114,-0.67720) (-1.45412,-0.46212) (-1.59880,-0.24122) (-1.72335,-0.01728) (-1.82620,0.20687) (-1.90606,0.42842) (-1.96192,0.64458) (-1.99308,0.85262) (-1.99915,1.04993) (-1.98006,1.23402)};\draw[blue,->] (axis cs:1.63449,0.18092)--(axis cs:1.83715,-0.23398);\addplot[blue,no marks] coordinates {(1.87014,-1.54902) (1.77894,-1.68100) (1.66536,-1.79181) (1.53080,-1.88007) (1.37698,-1.94465) (1.20581,-1.98476) (1.01947,-1.99987) (0.82029,-1.98981) (0.61079,-1.95470) (0.39359,-1.89498) (0.17144,-1.81140) (-0.05287,-1.70501) (-0.27651,-1.57716) (-0.49667,-1.42946) (-0.71058,-1.26376) (-0.91554,-1.08214) (-1.10898,-0.88691) (-1.28845,-0.68051) (-1.45171,-0.46554) (-1.59668,-0.24471) (-1.72156,-0.02080) (-1.82476,0.20338) (-1.90499,0.42499) (-1.96123,0.64125) (-1.99279,0.84944) (-1.99925,1.04693) (-1.98055,1.23125) (-1.93691,1.40006) (-1.86889,1.55125) (-1.77733,1.68290) (-1.66341,1.79337) (-1.52854,1.88127) (-1.37442,1.94547) (-1.20300,1.98519) (-1.01644,1.99991) (-0.81708,1.98945) (-0.60744,1.95395) (-0.39014,1.89385) (-0.16794,1.80990) (0.05638,1.70317) (0.27999,1.57500) (0.50008,1.42700) (0.71387,1.26103) (0.91867,1.07918) (1.11190,0.88375) (1.29114,0.67720) (1.45412,0.46212) (1.59880,0.24122) (1.72335,0.01728) (1.82620,-0.20687) (1.90606,-0.42842) (1.96192,-0.64458) (1.99308,-0.85262) (1.99915,-1.04993) (1.98006,-1.23402)};\draw[blue,->] (axis cs:-1.63449,-0.18092)--(axis cs:-1.83715,0.23398);\end{axis}\end{tikzpicture}\hspace{.035\textwidth}\begin{tikzpicture}\begin{axis}[width=.28\textwidth,height=3.8cm,axis lines=middle,title={$a=2$},xlabel={$x$},ylabel={$y$},xmin=-2.4,xmax=2.4,ymin=-2.4,ymax=2.4,clip=true]\addplot[blue,no marks] coordinates {(-0.51788,0.88895) (-0.46603,0.89495) (-0.40667,0.89360) (-0.33980,0.88439) (-0.26552,0.86691) (-0.18395,0.84073) (-0.09533,0.80550) (0.00006,0.76094) (0.10183,0.70678) (0.20953,0.64287) (0.32263,0.56910) (0.44050,0.48546) (0.56245,0.39200) (0.68768,0.28889) (0.81533,0.17637) (0.94446,0.05478) (1.07404,-0.07543) (1.20298,-0.21370) (1.33012,-0.35937) (1.45424,-0.51169) (1.57405,-0.66976) (1.68824,-0.83261) (1.79542,-0.99914) (1.89421,-1.16815) (1.98318,-1.33835) (2.06091,-1.50834) (2.12597,-1.67663) (2.17694,-1.84164) (2.21245,-2.00173) (2.23113,-2.15517) (2.23172,-2.30019) (2.21299,-2.43498) (2.17381,-2.55766) (2.11315,-2.66638) (2.03009,-2.75925) (1.92387,-2.83440) (1.79385,-2.89000) (1.63956,-2.92426) (1.46071,-2.93544) (1.25722,-2.92191) (1.02920,-2.88213) (0.77698,-2.81468) (0.50114,-2.71830) (0.20247,-2.59189) (-0.11795,-2.43454) (-0.45880,-2.24553) (-0.81849,-2.02438) (-1.19518,-1.77087) (-1.58671,-1.48502) (-1.99071,-1.16715) (-2.40448,-0.81787) (-2.82511,-0.43811) (-3.24940,-0.02912)};\draw[blue,->] (axis cs:1.48897,-0.55630)--(axis cs:1.70338,-0.85525);\addplot[blue,no marks] coordinates {(0.51788,-0.88895) (0.46603,-0.89495) (0.40667,-0.89360) (0.33980,-0.88439) (0.26552,-0.86691) (0.18395,-0.84073) (0.09533,-0.80550) (-0.00006,-0.76094) (-0.10183,-0.70678) (-0.20953,-0.64287) (-0.32263,-0.56910) (-0.44050,-0.48546) (-0.56245,-0.39200) (-0.68768,-0.28889) (-0.81533,-0.17637) (-0.94446,-0.05478) (-1.07404,0.07543) (-1.20298,0.21370) (-1.33012,0.35937) (-1.45424,0.51169) (-1.57405,0.66976) (-1.68824,0.83261) (-1.79542,0.99914) (-1.89421,1.16815) (-1.98318,1.33835) (-2.06091,1.50834) (-2.12597,1.67663) (-2.17694,1.84164) (-2.21245,2.00173) (-2.23113,2.15517) (-2.23172,2.30019) (-2.21299,2.43498) (-2.17381,2.55766) (-2.11315,2.66638) (-2.03009,2.75925) (-1.92387,2.83440) (-1.79385,2.89000) (-1.63956,2.92426) (-1.46071,2.93544) (-1.25722,2.92191) (-1.02920,2.88213) (-0.77698,2.81468) (-0.50114,2.71830) (-0.20247,2.59189) (0.11795,2.43454) (0.45880,2.24553) (0.81849,2.02438) (1.19518,1.77087) (1.58671,1.48502) (1.99071,1.16715) (2.40448,0.81787) (2.82511,0.43811) (3.24940,0.02912)};\draw[blue,->] (axis cs:-1.48897,0.55630)--(axis cs:-1.70338,0.85525);\addplot[blue,no marks] coordinates {(-0.88895,0.81555) (-0.89495,0.87640) (-0.89360,0.93372) (-0.88439,0.98679) (-0.86691,1.03485) (-0.84073,1.07714) (-0.80550,1.11293) (-0.76094,1.14146) (-0.70678,1.16200) (-0.64287,1.17384) (-0.56910,1.17628) (-0.48546,1.16869) (-0.39200,1.15045) (-0.28889,1.12101) (-0.17637,1.07988) (-0.05478,1.02662) (0.07543,0.96090) (0.21370,0.88243) (0.35937,0.79106) (0.51169,0.68671) (0.66976,0.56942) (0.83261,0.43933) (0.99914,0.29672) (1.16815,0.14198) (1.33835,-0.02434) (1.50834,-0.20160) (1.67663,-0.38897) (1.84164,-0.58552) (2.00173,-0.79014) (2.15517,-1.00162) (2.30019,-1.21857) (2.43498,-1.43948) (2.55766,-1.66269) (2.66638,-1.88643) (2.75925,-2.10878) (2.83440,-2.32773) (2.89000,-2.54116) (2.92426,-2.74683) (2.93544,-2.94245) (2.92191,-3.12564) (2.88213,-3.29399) (2.81468,-3.44504) (0.43811,-3.48227) (0.02912,-3.29307) (-0.40751,-3.06264) (-0.86985,-2.79021)};\draw[blue,->] (axis cs:0.55630,0.65452)--(axis cs:0.85525,0.42050);\addplot[blue,no marks] coordinates {(0.88895,-0.81555) (0.89495,-0.87640) (0.89360,-0.93372) (0.88439,-0.98679) (0.86691,-1.03485) (0.84073,-1.07714) (0.80550,-1.11293) (0.76094,-1.14146) (0.70678,-1.16200) (0.64287,-1.17384) (0.56910,-1.17628) (0.48546,-1.16869) (0.39200,-1.15045) (0.28889,-1.12101) (0.17637,-1.07988) (0.05478,-1.02662) (-0.07543,-0.96090) (-0.21370,-0.88243) (-0.35937,-0.79106) (-0.51169,-0.68671) (-0.66976,-0.56942) (-0.83261,-0.43933) (-0.99914,-0.29672) (-1.16815,-0.14198) (-1.33835,0.02434) (-1.50834,0.20160) (-1.67663,0.38897) (-1.84164,0.58552) (-2.00173,0.79014) (-2.15517,1.00162) (-2.30019,1.21857) (-2.43498,1.43948) (-2.55766,1.66269) (-2.66638,1.88643) (-2.75925,2.10878) (-2.83440,2.32773) (-2.89000,2.54116) (-2.92426,2.74683) (-2.93544,2.94245) (-2.92191,3.12564) (-2.88213,3.29399) (-2.81468,3.44504) (-0.43811,3.48227) (-0.02912,3.29307) (0.40751,3.06264) (0.86985,2.79021)};\draw[blue,->] (axis cs:-0.55630,-0.65452)--(axis cs:-0.85525,-0.42050);\addplot[blue,no marks] coordinates {(-1.40683,1.70450) (-1.36098,1.77135) (-1.30026,1.82732) (-1.22420,1.87118) (-1.13242,1.90175) (-1.02468,1.91787) (-0.90083,1.91844) (-0.76087,1.90240) (-0.60495,1.86878) (-0.43333,1.81671) (-0.24647,1.74539) (-0.04495,1.65415) (0.17045,1.54246) (0.39879,1.40990) (0.63897,1.25625) (0.88968,1.08140) (1.14947,0.88547) (1.41668,0.66874) (1.68950,0.43169) (1.96593,0.17503) (2.24381,-0.10034) (2.52084,-0.39328) (2.79456,-0.70242) (3.06236,-1.02617) (3.32154,-1.36270) (-3.22028,-3.32219)};\addplot[blue,no marks] coordinates {(1.40683,-1.70450) (1.36098,-1.77135) (1.30026,-1.82732) (1.22420,-1.87118) (1.13242,-1.90175) (1.02468,-1.91787) (0.90083,-1.91844) (0.76087,-1.90240) (0.60495,-1.86878) (0.43333,-1.81671) (0.24647,-1.74539) (0.04495,-1.65415) (-0.17045,-1.54246) (-0.39879,-1.40990) (-0.63897,-1.25625) (-0.88968,-1.08140) (-1.14947,-0.88547) (-1.41668,-0.66874) (-1.68950,-0.43169) (-1.96593,-0.17503) (-2.24381,0.10034) (-2.52084,0.39328) (-2.79456,0.70242) (-3.06236,1.02617) (-3.32154,1.36270) (3.22028,3.32219)};\end{axis}\end{tikzpicture}\hspace{.035\textwidth}\begin{tikzpicture}\begin{axis}[width=.28\textwidth,height=3.8cm,axis lines=middle,title={$a=3$},xlabel={$x$},ylabel={$y$},xmin=-2.4,xmax=2.4,ymin=-2.4,ymax=2.4,clip=true]\addplot[blue,no marks] coordinates {(-0.36788,0.73576) (-0.34163,0.73415) (-0.31022,0.72902) (-0.27307,0.71991) (-0.22955,0.70631) (-0.17898,0.68767) (-0.12061,0.66336) (-0.05362,0.63271) (0.02288,0.59498) (0.10987,0.54937) (0.20841,0.49497) (0.31965,0.43083) (0.44485,0.35588) (0.58539,0.26896) (0.74276,0.16881) (0.91857,0.05403) (1.11460,-0.07687) (1.33276,-0.22554) (1.57515,-0.39379) (1.84402,-0.58355) (2.14183,-0.79696) (2.47126,-1.03633) (2.83520,-1.30419) (3.23680,-1.60328)};\addplot[blue,no marks] coordinates {(0.36788,-0.73576) (0.34163,-0.73415) (0.31022,-0.72902) (0.27307,-0.71991) (0.22955,-0.70631) (0.17898,-0.68767) (0.12061,-0.66336) (0.05362,-0.63271) (-0.02288,-0.59498) (-0.10987,-0.54937) (-0.20841,-0.49497) (-0.31965,-0.43083) (-0.44485,-0.35588) (-0.58539,-0.26896) (-0.74276,-0.16881) (-0.91857,-0.05403) (-1.11460,0.07687) (-1.33276,0.22554) (-1.57515,0.39379) (-1.84402,0.58355) (-2.14183,0.79696) (-2.47126,1.03633) (-2.83520,1.30419) (-3.23680,1.60328)};\addplot[blue,no marks] coordinates {(-0.73576,1.10364) (-0.73415,1.12666) (-0.72902,1.14781) (-0.71991,1.16675) (-0.70631,1.18307) (-0.68767,1.19635) (-0.66336,1.20611) (-0.63271,1.21180) (-0.59498,1.21285) (-0.54937,1.20861) (-0.49497,1.19836) (-0.43083,1.18132) (-0.35588,1.15662) (-0.26896,1.12332) (-0.16881,1.08037) (-0.05403,1.02664) (0.07687,0.96086) (0.22554,0.88167) (0.39379,0.78757) (0.58355,0.67692) (0.79696,0.54791) (1.03633,0.39859) (1.30419,0.22682) (1.60328,0.03025) (1.93657,-0.19366) (2.30730,-0.44769) (2.71901,-0.73487) (3.17550,-1.05850)};\draw[blue,->] (axis cs:0.64201,0.64201)--(axis cs:1.07132,0.37641);\addplot[blue,no marks] coordinates {(0.73576,-1.10364) (0.73415,-1.12666) (0.72902,-1.14781) (0.71991,-1.16675) (0.70631,-1.18307) (0.68767,-1.19635) (0.66336,-1.20611) (0.63271,-1.21180) (0.59498,-1.21285) (0.54937,-1.20861) (0.49497,-1.19836) (0.43083,-1.18132) (0.35588,-1.15662) (0.26896,-1.12332) (0.16881,-1.08037) (0.05403,-1.02664) (-0.07687,-0.96086) (-0.22554,-0.88167) (-0.39379,-0.78757) (-0.58355,-0.67692) (-0.79696,-0.54791) (-1.03633,-0.39859) (-1.30419,-0.22682) (-1.60328,-0.03025) (-1.93657,0.19366) (-2.30730,0.44769) (-2.71901,0.73487) (-3.17550,1.05850)};\draw[blue,->] (axis cs:-0.64201,-0.64201)--(axis cs:-1.07132,-0.37641);\addplot[blue,no marks] coordinates {(-1.10364,1.83940) (-1.07578,1.86080) (-1.03924,1.87683) (-0.99298,1.88666) (-0.93586,1.88938) (-0.86665,1.88402) (-0.78397,1.86946) (-0.68633,1.84451) (-0.57210,1.80783) (-0.43949,1.75798) (-0.28656,1.69333) (-0.11118,1.61215) (0.08897,1.51250) (0.31643,1.39228) (0.57395,1.24918) (0.86454,1.08067) (1.19147,0.88399) (1.55831,0.65613) (1.96893,0.39379) (2.42756,0.09337) (2.93879,-0.24905)};\addplot[blue,no marks] coordinates {(1.10364,-1.83940) (1.07578,-1.86080) (1.03924,-1.87683) (0.99298,-1.88666) (0.93586,-1.88938) (0.86665,-1.88402) (0.78397,-1.86946) (0.68633,-1.84451) (0.57210,-1.80783) (0.43949,-1.75798) (0.28656,-1.69333) (0.11118,-1.61215) (-0.08897,-1.51250) (-0.31643,-1.39228) (-0.57395,-1.24918) (-0.86454,-1.08067) (-1.19147,-0.88399) (-1.55831,-0.65613) (-1.96893,-0.39379) (-2.42756,-0.09337) (-2.93879,0.24905)};\end{axis}\end{tikzpicture}\par\begin{tikzpicture}\begin{axis}[width=.28\textwidth,height=3.8cm,axis lines=middle,title={$a=4$},xlabel={$x$},ylabel={$y$},xmin=-2.4,xmax=2.4,ymin=-2.4,ymax=2.4,clip=true]\addplot[blue,no marks] coordinates {(-0.26695,0.63348) (-0.25270,0.62635) (-0.23540,0.61770) (-0.21438,0.60719) (-0.18884,0.59442) (-0.15783,0.57891) (-0.12016,0.56008) (-0.07441,0.53720) (-0.01883,0.50942) (0.04867,0.47566) (0.13067,0.43467) (0.23026,0.38487) (0.35122,0.32439) (0.49815,0.25092) (0.67662,0.16169) (0.89339,0.05330) (1.15669,-0.07835) (1.47651,-0.23825) (1.86496,-0.43248) (2.33679,-0.66840) (2.90990,-0.95495)};\addplot[blue,no marks] coordinates {(0.26695,-0.63348) (0.25270,-0.62635) (0.23540,-0.61770) (0.21438,-0.60719) (0.18884,-0.59442) (0.15783,-0.57891) (0.12016,-0.56008) (0.07441,-0.53720) (0.01883,-0.50942) (-0.04867,-0.47566) (-0.13067,-0.43467) (-0.23026,-0.38487) (-0.35122,-0.32439) (-0.49815,-0.25092) (-0.67662,-0.16169) (-0.89339,-0.05330) (-1.15669,0.07835) (-1.47651,0.23825) (-1.86496,0.43248) (-2.33679,0.66840) (-2.90990,0.95495)};\addplot[blue,no marks] coordinates {(-0.63348,1.31674) (-0.62635,1.31318) (-0.61770,1.30885) (-0.60719,1.30359) (-0.59442,1.29721) (-0.57891,1.28946) (-0.56008,1.28004) (-0.53720,1.26860) (-0.50942,1.25471) (-0.47566,1.23783) (-0.43467,1.21733) (-0.38487,1.19244) (-0.32439,1.16219) (-0.25092,1.12546) (-0.16169,1.08084) (-0.05330,1.02665) (0.07835,0.96083) (0.23825,0.88087) (0.43248,0.78376) (0.66840,0.66580) (0.95495,0.52252) (1.30301,0.34850) (1.72577,0.13711) (2.23927,-0.11964) (2.86299,-0.43150)};\draw[blue,->] (axis cs:0.74467,0.62767)--(axis cs:1.35624,0.32188);\addplot[blue,no marks] coordinates {(0.63348,-1.31674) (0.62635,-1.31318) (0.61770,-1.30885) (0.60719,-1.30359) (0.59442,-1.29721) (0.57891,-1.28946) (0.56008,-1.28004) (0.53720,-1.26860) (0.50942,-1.25471) (0.47566,-1.23783) (0.43467,-1.21733) (0.38487,-1.19244) (0.32439,-1.16219) (0.25092,-1.12546) (0.16169,-1.08084) (0.05330,-1.02665) (-0.07835,-0.96083) (-0.23825,-0.88087) (-0.43248,-0.78376) (-0.66840,-0.66580) (-0.95495,-0.52252) (-1.30301,-0.34850) (-1.72577,-0.13711) (-2.23927,0.11964) (-2.86299,0.43150)};\draw[blue,->] (axis cs:-0.74467,-0.62767)--(axis cs:-1.35624,-0.32188);\addplot[blue,no marks] coordinates {(-0.90043,1.95021) (-0.87905,1.93953) (-0.85309,1.92655) (-0.82156,1.91078) (-0.78326,1.89163) (-0.73674,1.86837) (-0.68024,1.84012) (-0.61161,1.80580) (-0.52825,1.76412) (-0.42699,1.71350) (-0.30400,1.65200) (-0.15462,1.57731) (0.02683,1.48658) (0.24723,1.37639) (0.51493,1.24253) (0.84009,1.07996) (1.23504,0.88248) (1.71476,0.64262) (2.29744,0.35128) (3.00519,-0.00260)};\addplot[blue,no marks] coordinates {(0.90043,-1.95021) (0.87905,-1.93953) (0.85309,-1.92655) (0.82156,-1.91078) (0.78326,-1.89163) (0.73674,-1.86837) (0.68024,-1.84012) (0.61161,-1.80580) (0.52825,-1.76412) (0.42699,-1.71350) (0.30400,-1.65200) (0.15462,-1.57731) (-0.02683,-1.48658) (-0.24723,-1.37639) (-0.51493,-1.24253) (-0.84009,-1.07996) (-1.23504,-0.88248) (-1.71476,-0.64262) (-2.29744,-0.35128) (-3.00519,0.00260)};\end{axis}\end{tikzpicture}\hspace{.035\textwidth}\begin{tikzpicture}\begin{axis}[width=.28\textwidth,height=3.8cm,axis lines=middle,title={$a=5$},xlabel={$x$},ylabel={$y$},xmin=-2.4,xmax=2.4,ymin=-2.4,ymax=2.4,clip=true]\addplot[blue,no marks] coordinates {(-0.19937,0.55982) (-0.19073,0.54918) (-0.18046,0.53802) (-0.16801,0.52614) (-0.15268,0.51327) (-0.13354,0.49905) (-0.10938,0.48302) (-0.07860,0.46456) (-0.03908,0.44287) (0.01195,0.41688) (0.07813,0.38520) (0.16428,0.34599) (0.27672,0.29684) (0.42376,0.23457) (0.61635,0.15498) (0.86890,0.05259) (1.20037,-0.07986) (1.63570,-0.25189) (2.20773,-0.47605) (2.95965,-0.76883)};\addplot[blue,no marks] coordinates {(0.19937,-0.55982) (0.19073,-0.54918) (0.18046,-0.53802) (0.16801,-0.52614) (0.15268,-0.51327) (0.13354,-0.49905) (0.10938,-0.48302) (0.07860,-0.46456) (0.03908,-0.44287) (-0.01195,-0.41688) (-0.07813,-0.38520) (-0.16428,-0.34599) (-0.27672,-0.29684) (-0.42376,-0.23457) (-0.61635,-0.15498) (-0.86890,-0.05259) (-1.20037,0.07986) (-1.63570,0.25189) (-2.20773,0.47605) (-2.95965,0.76883)};\addplot[blue,no marks] coordinates {(-0.55982,1.48009) (-0.54918,1.45680) (-0.53802,1.43360) (-0.52614,1.41042) (-0.51327,1.38714) (-0.49905,1.36362) (-0.48302,1.33969) (-0.46456,1.31510) (-0.44287,1.28954) (-0.41688,1.26260) (-0.38520,1.23374) (-0.34599,1.20227) (-0.29684,1.16724) (-0.23457,1.12745) (-0.15498,1.08130) (-0.05259,1.02667) (0.07986,0.96079) (0.25189,0.88003) (0.47605,0.77959) (0.76883,0.65315) (1.15196,0.49244) (1.65400,0.28651) (2.31258,0.02102) (3.17716,-0.32296)};\draw[blue,->] (axis cs:0.86797,0.61115)--(axis cs:1.73409,0.25400);\addplot[blue,no marks] coordinates {(0.55982,-1.48009) (0.54918,-1.45680) (0.53802,-1.43360) (0.52614,-1.41042) (0.51327,-1.38714) (0.49905,-1.36362) (0.48302,-1.33969) (0.46456,-1.31510) (0.44287,-1.28954) (0.41688,-1.26260) (0.38520,-1.23374) (0.34599,-1.20227) (0.29684,-1.16724) (0.23457,-1.12745) (0.15498,-1.08130) (0.05259,-1.02667) (-0.07986,-0.96079) (-0.25189,-0.88003) (-0.47605,-0.77959) (-0.76883,-0.65315) (-1.15196,-0.49244) (-1.65400,-0.28651) (-2.31258,-0.02102) (-3.17716,0.32296)};\draw[blue,->] (axis cs:-0.86797,-0.61115)--(axis cs:-1.73409,-0.25400);\addplot[blue,no marks] coordinates {(-0.75919,2.03991) (-0.73991,2.00598) (-0.71848,1.97162) (-0.69415,1.93656) (-0.66595,1.90041) (-0.63259,1.86268) (-0.59240,1.82271) (-0.54316,1.77966) (-0.48195,1.73241) (-0.40494,1.67948) (-0.30707,1.61895) (-0.18171,1.54826) (-0.02012,1.46409) (0.18919,1.36202) (0.46137,1.23628) (0.81632,1.07925) (1.28023,0.88093) (1.88760,0.62814) (2.68378,0.30354)};\addplot[blue,no marks] coordinates {(0.75919,-2.03991) (0.73991,-2.00598) (0.71848,-1.97162) (0.69415,-1.93656) (0.66595,-1.90041) (0.63259,-1.86268) (0.59240,-1.82271) (0.54316,-1.77966) (0.48195,-1.73241) (0.40494,-1.67948) (0.30707,-1.61895) (0.18171,-1.54826) (0.02012,-1.46409) (-0.18919,-1.36202) (-0.46137,-1.23628) (-0.81632,-1.07925) (-1.28023,-0.88093) (-1.88760,-0.62814) (-2.68378,-0.30354)};\end{axis}\end{tikzpicture}\end{center}
\par\noindent{\small 依据：官方译审并补证（AI辅助）。}
\end{solution}

\begin{exercise}\label{cw:rec23-4}
给出这八种情形的通解方法，特别处理亏损结点。
\end{exercise}
\begin{solution}
以下公式直接给出每例的所有解，不仅列方法名。令 $h=(a-1)/2,N=A-hI$，Cayley--Hamilton 给 $N^2=d^2I$，$d^2=(a+5)(a-3)/4$。
\[
 u(t)=\begin{cases}
 \ee^{ht}[\cosh(dt)I+\sinh(dt)N/d]c,&d^2>0,\\
 \ee^{ht}[\cos(\beta t)I+\sin(\beta t)N/\beta]c,&d^2=-\beta^2<0,\\
 \ee^{ht}(I+tN)c,&d=0,
 \end{cases}\qquad c\in\mathbb R^2.
\]
对每行求导并用$N^2$即得到$\dot u=Au$；$t=0$矩阵为$I$，故覆盖全部初值。各参数代入为
\[\begin{array}{c|rrrrrrrr}a&-6&-5&-2&1&2&3&4&5\\\hline h&-7/2&-3&-3/2&0&1/2&1&3/2&2\\ d^2&9/4&0&-15/4&-3&-7/4&0&9/4&5\end{array}\]
特别地，$a=-5$时$N=\begin{pmatrix}-2&2\\-2&2\end{pmatrix}$，$a=3$时$N=\begin{pmatrix}2&2\\-2&-2\end{pmatrix}$；两者均$N^2=0,N\ne0$，故$tN$项不可删。这明确补上官方此题未给出的推导。
\par\noindent{\small 依据：官方分类译审；通解与亏损矩阵指数为编者补解（AI辅助）。}
\end{solution}

\originalsection{第24次习题课：矩阵指数与非齐次系统（5月6日）}
本次沿用原题的长期定义：“稳定”指所有解趋于0；“不稳定”指一般初值的解无界；其余称中性稳定。该定义讨论长期极限，不与 Lyapunov 稳定的术语混淆。
\begin{exercise}\label{cw:rec24-1}
在迹--行列式平面，何处能保证所有具有同样$(\tau,\Delta)$的实二阶矩阵稳定、不稳定、中性稳定？有无相同迹与行列式而长期行为不同的矩阵？
\end{exercise}
\begin{solution}
$\tau<0,\Delta>0$ 恰好对应两个根实部负，指数衰减（重根的$t$因子亦被负指数压过），因此稳定。$\Delta<0$有一正一负根，或$\tau>0$有正实部根，均不稳定。
$\tau=0,\Delta>0$时为非零纯虚根，矩阵指数有界周期而不趋于0，中性稳定；$\Delta=0,\tau<0$时根为0与负数，解趋于零特征方向上的常量，亦中性稳定。
剩下原点$(0,0)$：Cayley--Hamilton使$A^2=0$。若$A=0$，全部解常值，中性稳定；若$A\ne0$，$\ee^{At}=I+tA$，一般$c$满足$Ac\ne0$，解线性无界，按本题定义为不稳定。例如$\begin{pmatrix}0&1\\0&0\end{pmatrix}$与零矩阵同迹同列式却不同。以上区域与原点情形覆盖整个平面。
\par\noindent{\small 依据：官方译审并补证（AI辅助）；补幂零边界与本题稳定定义。}
\end{solution}

\begin{exercise}\label{cw:rec24-2}
设$A=\begin{pmatrix}1&1\\-4&1\end{pmatrix}$。
(a) 求基本矩阵$\Phi(t)$；(b) 求$\ee^{At}$；(c) 求初值$u(0)=(1,2)^T$的解；(d) 求$\dot u=Au+(5,10)^T$的一个特解、通解与$u(0)=0$的解。
\end{exercise}
\begin{solution}
(a) $1+2\ii$的特征向量$(1,2\ii)^T$，复解实虚部给
\[\Phi=\ee^t\begin{pmatrix}\cos2t&\sin2t\\-2\sin2t&2\cos2t\end{pmatrix},\quad\det\Phi=2\ee^{2t}\ne0.\]
(b) $\Phi(0)=\diag(1,2)$，因此
\[E(t)=\ee^{At}=\ee^t\begin{pmatrix}\cos2t&\frac12\sin2t\\-2\sin2t&\cos2t\end{pmatrix}.\]
(c) $E(t)(1,2)^T=\ee^t(\cos2t+\sin2t,2\cos2t-2\sin2t)^T$。
(d) 常值特解$p=-A^{-1}(5,10)^T=(1,-6)^T$，直接有$Ap+(5,10)^T=0$。通解$p+E(t)c$；初值零时$c=-p=(-1,6)^T$，故
\[u(t)=\binom{1-\ee^t\cos2t+3\ee^t\sin2t}{-6+2\ee^t\sin2t+6\ee^t\cos2t}.\]
$E(0)=I,E'=AE$检验了全部初值与方程。
\par\noindent{\small 依据：官方译审并补证（AI辅助）。}
\end{solution}

\begin{exercise}\label{cw:rec24-3}
已知常值$u_1=(1,1)^T$和$u_2=(\ee^t,-\ee^t)^T$是$\dot u=Bu$的解。
(a) 求通解，以及初值$(2,2)^T$、$(1,0)^T$的两个解；(b) 求基本矩阵、$\ee^{Bt}$、$\ee^B$；(c) 求$B$的特征值与向量；(d) 求$B$。
\end{exercise}
\begin{solution}
(a) 两初值向量独立，故$u=c_1(1,1)^T+c_2\ee^t(1,-1)^T$。第一初值给常值$(2,2)^T$；第二初值给$u=((1+\ee^t)/2,(1-\ee^t)/2)^T$。
(b) $\Phi=\begin{pmatrix}1&\ee^t\\1&-\ee^t\end{pmatrix}$，$\Phi(0)^{-1}=\frac12\begin{pmatrix}1&1\\1&-1\end{pmatrix}$，故
\[\ee^{Bt}=\frac12\begin{pmatrix}1+\ee^t&1-\ee^t\\1-\ee^t&1+\ee^t\end{pmatrix},\qquad\ee^B=\frac12\begin{pmatrix}1+\ee&1-\ee\\1-\ee&1+\ee\end{pmatrix}.\]
(c) 直接从$Bu_1=u_1'=0$、$Bu_2=u_2'=u_2$得特征值0、1，向量$(1,1)^T,(1,-1)^T$；无需对指数特征值取存在分支歧义的复对数。
(d) 令$P=\begin{pmatrix}1&1\\1&-1\end{pmatrix}$，$B=P\diag(0,1)P^{-1}=\frac12\begin{pmatrix}1&-1\\-1&1\end{pmatrix}$。
\par\noindent{\small 依据：官方译审并补证（AI辅助）。}
\end{solution}

\originalsection{第25次习题课：自治系统与生态模型（5月11日）}
令$x$为昆虫种群、$y$为鸟类种群，使用题设的方便单位。模型 $\dot x=(2-x-y)x,\dot y=(x-1)y$ 表示鸟捕食虫：无鸟时虫 logistic 增长，无虫时鸟指数消亡。
\begin{exercise}\label{cw:rec25-1}
求系统所有临界点（平衡点）。
\end{exercise}
\begin{solution}
同时解$x(2-x-y)=0,y(x-1)=0$。若$x=0$则$y=0$；若$y=0$且$x\ne0$则$x=2$；若两者非零则$x=1,y=1$。所以全部为$(0,0),(2,0),(1,1)$，没有遗漏其他实平衡。
\par\noindent{\small 依据：官方译审并补证（AI辅助）。}
\end{solution}

\begin{exercise}\label{cw:rec25-2}
在每个平衡点求线性化，画附近轨迹。
\end{exercise}
\begin{solution}
Jacobian为$J=\begin{pmatrix}2-2x-y&-x\\y&x-1\end{pmatrix}$。
在$(0,0)$为$\diag(2,-1)$：鞍点，$x$轴向外、$y$轴向内。
在$(2,0)$为$\begin{pmatrix}-2&-2\\0&1\end{pmatrix}$：鞍点，稳定方向$(1,0)$，不稳定方向$(2,-3)$（斜率$-3/2$）。
在$(1,1)$为$\begin{pmatrix}-1&-1\\1&0\end{pmatrix}$：特征值$(-1\pm\ii\sqrt3)/2$，逆时针稳定螺旋。三个均双曲，线性化决定局部相图类型；图中坐标为相对平衡点的扰动$(\xi,\eta)$（小图的$x,y$轴分别指这两个扰动）。\begin{center}\begin{tikzpicture}\begin{axis}[width=.28\textwidth,height=3.8cm,axis lines=middle,title={原点附近},xlabel={$x$},ylabel={$y$},xmin=-2.4,xmax=2.4,ymin=-2.4,ymax=2.4,clip=true]\addplot[blue,no marks] coordinates {(0.13534,0.00000) (0.15695,0.00000) (0.18201,0.00000) (0.21107,0.00000) (0.24478,0.00000) (0.28386,0.00000) (0.32919,0.00000) (0.38176,0.00000) (0.44272,0.00000) (0.51342,0.00000) (0.59540,0.00000) (0.69048,0.00000) (0.80074,0.00000) (0.92860,0.00000) (1.07689,0.00000) (1.24885,0.00000) (1.44827,0.00000) (1.67954,0.00000) (1.94773,0.00000) (2.25876,0.00000) (2.61945,0.00000) (3.03773,0.00000)};\draw[blue,->] (axis cs:1.64872,0.00000)--(axis cs:2.09594,0.00000);\addplot[blue,no marks] coordinates {(-0.13534,0.00000) (-0.15695,0.00000) (-0.18201,0.00000) (-0.21107,0.00000) (-0.24478,0.00000) (-0.28386,0.00000) (-0.32919,0.00000) (-0.38176,0.00000) (-0.44272,0.00000) (-0.51342,0.00000) (-0.59540,0.00000) (-0.69048,0.00000) (-0.80074,0.00000) (-0.92860,0.00000) (-1.07689,0.00000) (-1.24885,0.00000) (-1.44827,0.00000) (-1.67954,0.00000) (-1.94773,0.00000) (-2.25876,0.00000) (-2.61945,0.00000) (-3.03773,0.00000)};\draw[blue,->] (axis cs:-1.64872,0.00000)--(axis cs:-2.09594,0.00000);\addplot[blue,no marks] coordinates {(0.00000,2.71828) (0.00000,2.52420) (0.00000,2.34398) (0.00000,2.17663) (0.00000,2.02122) (0.00000,1.87692) (0.00000,1.74291) (0.00000,1.61847) (0.00000,1.50292) (0.00000,1.39561) (0.00000,1.29597) (0.00000,1.20344) (0.00000,1.11752) (0.00000,1.03773) (0.00000,0.96364) (0.00000,0.89484) (0.00000,0.83095) (0.00000,0.77162) (0.00000,0.71653) (0.00000,0.66537) (0.00000,0.61787) (0.00000,0.57375) (0.00000,0.53279) (0.00000,0.49475) (0.00000,0.45943) (0.00000,0.42662) (0.00000,0.39616) (0.00000,0.36788) (0.00000,0.34161) (0.00000,0.31722) (0.00000,0.29457) (0.00000,0.27354) (0.00000,0.25401) (0.00000,0.23588) (0.00000,0.21904) (0.00000,0.20340) (0.00000,0.18888) (0.00000,0.17539) (0.00000,0.16287) (0.00000,0.15124) (0.00000,0.14044) (0.00000,0.13041) (0.00000,0.12110) (0.00000,0.11246) (0.00000,0.10443) (0.00000,0.09697) (0.00000,0.09005) (0.00000,0.08362) (0.00000,0.07765) (0.00000,0.07211) (0.00000,0.06696) (0.00000,0.06218) (0.00000,0.05774) (0.00000,0.05362) (0.00000,0.04979)};\draw[blue,->] (axis cs:0.00000,0.77880)--(axis cs:0.00000,0.69073);\addplot[blue,no marks] coordinates {(0.00000,-2.71828) (0.00000,-2.52420) (0.00000,-2.34398) (0.00000,-2.17663) (0.00000,-2.02122) (0.00000,-1.87692) (0.00000,-1.74291) (0.00000,-1.61847) (0.00000,-1.50292) (0.00000,-1.39561) (0.00000,-1.29597) (0.00000,-1.20344) (0.00000,-1.11752) (0.00000,-1.03773) (0.00000,-0.96364) (0.00000,-0.89484) (0.00000,-0.83095) (0.00000,-0.77162) (0.00000,-0.71653) (0.00000,-0.66537) (0.00000,-0.61787) (0.00000,-0.57375) (0.00000,-0.53279) (0.00000,-0.49475) (0.00000,-0.45943) (0.00000,-0.42662) (0.00000,-0.39616) (0.00000,-0.36788) (0.00000,-0.34161) (0.00000,-0.31722) (0.00000,-0.29457) (0.00000,-0.27354) (0.00000,-0.25401) (0.00000,-0.23588) (0.00000,-0.21904) (0.00000,-0.20340) (0.00000,-0.18888) (0.00000,-0.17539) (0.00000,-0.16287) (0.00000,-0.15124) (0.00000,-0.14044) (0.00000,-0.13041) (0.00000,-0.12110) (0.00000,-0.11246) (0.00000,-0.10443) (0.00000,-0.09697) (0.00000,-0.09005) (0.00000,-0.08362) (0.00000,-0.07765) (0.00000,-0.07211) (0.00000,-0.06696) (0.00000,-0.06218) (0.00000,-0.05774) (0.00000,-0.05362) (0.00000,-0.04979)};\draw[blue,->] (axis cs:0.00000,-0.77880)--(axis cs:0.00000,-0.69073);\addplot[blue,no marks] coordinates {(0.13534,2.71828) (0.15695,2.52420) (0.18201,2.34398) (0.21107,2.17663) (0.24478,2.02122) (0.28386,1.87692) (0.32919,1.74291) (0.38176,1.61847) (0.44272,1.50292) (0.51342,1.39561) (0.59540,1.29597) (0.69048,1.20344) (0.80074,1.11752) (0.92860,1.03773) (1.07689,0.96364) (1.24885,0.89484) (1.44827,0.83095) (1.67954,0.77162) (1.94773,0.71653) (2.25876,0.66537) (2.61945,0.61787) (3.03773,0.57375)};\draw[blue,->] (axis cs:1.64872,0.77880)--(axis cs:2.09594,0.69073);\addplot[blue,no marks] coordinates {(-0.13534,-2.71828) (-0.15695,-2.52420) (-0.18201,-2.34398) (-0.21107,-2.17663) (-0.24478,-2.02122) (-0.28386,-1.87692) (-0.32919,-1.74291) (-0.38176,-1.61847) (-0.44272,-1.50292) (-0.51342,-1.39561) (-0.59540,-1.29597) (-0.69048,-1.20344) (-0.80074,-1.11752) (-0.92860,-1.03773) (-1.07689,-0.96364) (-1.24885,-0.89484) (-1.44827,-0.83095) (-1.67954,-0.77162) (-1.94773,-0.71653) (-2.25876,-0.66537) (-2.61945,-0.61787) (-3.03773,-0.57375)};\draw[blue,->] (axis cs:-1.64872,-0.77880)--(axis cs:-2.09594,-0.69073);\end{axis}\end{tikzpicture}\hspace{.035\textwidth}\begin{tikzpicture}\begin{axis}[width=.28\textwidth,height=3.8cm,axis lines=middle,title={$(2,0)$附近},xlabel={$x$},ylabel={$y$},xmin=-2.4,xmax=2.4,ymin=-2.4,ymax=2.4,clip=true]\addplot[blue,no marks] coordinates {(3.03773,0.00000) (2.61945,0.00000) (2.25876,0.00000) (1.94773,0.00000) (1.67954,0.00000) (1.44827,0.00000) (1.24885,0.00000) (1.07689,0.00000) (0.92860,0.00000) (0.80074,0.00000) (0.69048,0.00000) (0.59540,0.00000) (0.51342,0.00000) (0.44272,0.00000) (0.38176,0.00000) (0.32919,0.00000) (0.28386,0.00000) (0.24478,0.00000) (0.21107,0.00000) (0.18201,0.00000) (0.15695,0.00000) (0.13534,0.00000) (0.11670,0.00000) (0.10063,0.00000) (0.08677,0.00000) (0.07483,0.00000) (0.06452,0.00000) (0.05564,0.00000) (0.04798,0.00000) (0.04137,0.00000) (0.03567,0.00000) (0.03076,0.00000) (0.02653,0.00000) (0.02287,0.00000) (0.01972,0.00000) (0.01701,0.00000) (0.01467,0.00000) (0.01265,0.00000) (0.01091,0.00000) (0.00940,0.00000) (0.00811,0.00000) (0.00699,0.00000) (0.00603,0.00000) (0.00520,0.00000) (0.00448,0.00000) (0.00387,0.00000) (0.00333,0.00000) (0.00287,0.00000) (0.00248,0.00000)};\draw[blue,->] (axis cs:0.60653,0.00000)--(axis cs:0.47711,0.00000);\addplot[blue,no marks] coordinates {(-3.03773,0.00000) (-2.61945,0.00000) (-2.25876,0.00000) (-1.94773,0.00000) (-1.67954,0.00000) (-1.44827,0.00000) (-1.24885,0.00000) (-1.07689,0.00000) (-0.92860,0.00000) (-0.80074,0.00000) (-0.69048,0.00000) (-0.59540,0.00000) (-0.51342,0.00000) (-0.44272,0.00000) (-0.38176,0.00000) (-0.32919,0.00000) (-0.28386,0.00000) (-0.24478,0.00000) (-0.21107,0.00000) (-0.18201,0.00000) (-0.15695,0.00000) (-0.13534,0.00000) (-0.11670,0.00000) (-0.10063,0.00000) (-0.08677,0.00000) (-0.07483,0.00000) (-0.06452,0.00000) (-0.05564,0.00000) (-0.04798,0.00000) (-0.04137,0.00000) (-0.03567,0.00000) (-0.03076,0.00000) (-0.02653,0.00000) (-0.02287,0.00000) (-0.01972,0.00000) (-0.01701,0.00000) (-0.01467,0.00000) (-0.01265,0.00000) (-0.01091,0.00000) (-0.00940,0.00000) (-0.00811,0.00000) (-0.00699,0.00000) (-0.00603,0.00000) (-0.00520,0.00000) (-0.00448,0.00000) (-0.00387,0.00000) (-0.00333,0.00000) (-0.00287,0.00000) (-0.00248,0.00000)};\draw[blue,->] (axis cs:-0.60653,0.00000)--(axis cs:-0.47711,0.00000);\addplot[blue,no marks] coordinates {(3.37842,0.42662) (2.85219,0.45943) (2.39373,0.49475) (1.99335,0.53279) (1.64265,0.57375) (1.33439,0.61787) (1.06226,0.66537) (0.82080,0.71653) (0.60528,0.77162) (0.41155,0.83095) (0.23601,0.89484) (0.07550,0.96364) (-0.07275,1.03773) (-0.21119,1.11752) (-0.34198,1.20344) (-0.46705,1.29597) (-0.58813,1.39561) (-0.70680,1.50292) (-0.82447,1.61847) (-0.94248,1.74291) (-1.06203,1.87692) (-1.18430,2.02122) (-1.31037,2.17663) (-1.44132,2.34398) (-1.57817,2.52420) (-1.72196,2.71828) (-1.87372,2.92728) (-2.03448,3.15235) (-2.20530,3.39472)};\draw[blue,->] (axis cs:-0.45166,1.28403)--(axis cs:-0.64708,1.44773);\addplot[blue,no marks] coordinates {(-3.37842,-0.42662) (-2.85219,-0.45943) (-2.39373,-0.49475) (-1.99335,-0.53279) (-1.64265,-0.57375) (-1.33439,-0.61787) (-1.06226,-0.66537) (-0.82080,-0.71653) (-0.60528,-0.77162) (-0.41155,-0.83095) (-0.23601,-0.89484) (-0.07550,-0.96364) (0.07275,-1.03773) (0.21119,-1.11752) (0.34198,-1.20344) (0.46705,-1.29597) (0.58813,-1.39561) (0.70680,-1.50292) (0.82447,-1.61847) (0.94248,-1.74291) (1.06203,-1.87692) (1.18430,-2.02122) (1.31037,-2.17663) (1.44132,-2.34398) (1.57817,-2.52420) (1.72196,-2.71828) (1.87372,-2.92728) (2.03448,-3.15235) (2.20530,-3.39472)};\draw[blue,->] (axis cs:0.45166,-1.28403)--(axis cs:0.64708,-1.44773);\addplot[blue,no marks] coordinates {(3.32101,0.66537) (2.76854,0.71653) (2.28481,0.77162) (1.85982,0.83095) (1.48486,0.89484) (1.15238,0.96364) (0.85585,1.03773) (0.58955,1.11752) (0.34850,1.20344) (0.12836,1.29597) (-0.07471,1.39561) (-0.26408,1.50292) (-0.44271,1.61847) (-0.61328,1.74291) (-0.77817,1.87692) (-0.93952,2.02122) (-1.09930,2.17663) (-1.25931,2.34398) (-1.42123,2.52420) (-1.58663,2.71828) (-1.75702,2.92728) (-1.93385,3.15235) (-2.11852,3.39472)};\draw[blue,->] (axis cs:0.15487,1.28403)--(axis cs:-0.16997,1.44773);\addplot[blue,no marks] coordinates {(-3.32101,-0.66537) (-2.76854,-0.71653) (-2.28481,-0.77162) (-1.85982,-0.83095) (-1.48486,-0.89484) (-1.15238,-0.96364) (-0.85585,-1.03773) (-0.58955,-1.11752) (-0.34850,-1.20344) (-0.12836,-1.29597) (0.07471,-1.39561) (0.26408,-1.50292) (0.44271,-1.61847) (0.61328,-1.74291) (0.77817,-1.87692) (0.93952,-2.02122) (1.09930,-2.17663) (1.25931,-2.34398) (1.42123,-2.52420) (1.58663,-2.71828) (1.75702,-2.92728) (1.93385,-3.15235) (2.11852,-3.39472)};\draw[blue,->] (axis cs:-0.15487,-1.28403)--(axis cs:0.16997,-1.44773);\end{axis}\end{tikzpicture}\hspace{.035\textwidth}\begin{tikzpicture}\begin{axis}[width=.28\textwidth,height=3.8cm,axis lines=middle,title={$(1,1)$附近},xlabel={$x$},ylabel={$y$},xmin=-2.4,xmax=2.4,ymin=-2.4,ymax=2.4,clip=true]\addplot[blue,no marks] coordinates {(1.79325,-1.45022) (1.73725,-1.22115) (1.66071,-1.00071) (1.56734,-0.79132) (1.46066,-0.59494) (1.34396,-0.41307) (1.22029,-0.24681) (1.09246,-0.09688) (0.96297,0.03635) (0.83407,0.15281) (0.70772,0.25270) (0.58559,0.33647) (0.46909,0.40477) (0.35937,0.45838) (0.25734,0.49827) (0.16366,0.52546) (0.07880,0.54108) (0.00301,0.54628) (-0.06361,0.54226) (-0.12113,0.53019) (-0.16975,0.51124) (-0.20979,0.48655) (-0.24165,0.45720) (-0.26584,0.42423) (-0.28291,0.38859) (-0.29345,0.35116) (-0.29809,0.31276) (-0.29747,0.27411) (-0.29225,0.23584) (-0.28306,0.19851) (-0.27053,0.16260) (-0.25527,0.12849) (-0.23785,0.09651) (-0.21880,0.06690) (-0.19863,0.03983) (-0.17778,0.01543) (-0.15667,-0.00624) (-0.13566,-0.02519) (-0.11506,-0.04143) (-0.09517,-0.05505) (-0.07619,-0.06615) (-0.05832,-0.07485) (-0.04171,-0.08132) (-0.02645,-0.08572) (-0.01264,-0.08824) (-0.00031,-0.08906) (0.01053,-0.08839) (0.01988,-0.08640) (0.02779,-0.08329) (0.03429,-0.07925) (0.03947,-0.07446) (0.04339,-0.06908) (0.04616,-0.06326) (0.04786,-0.05715) (0.04860,-0.05089)};\draw[blue,->] (axis cs:0.75244,0.21890)--(axis cs:0.63766,0.30227);\addplot[blue,no marks] coordinates {(-1.79325,1.45022) (-1.73725,1.22115) (-1.66071,1.00071) (-1.56734,0.79132) (-1.46066,0.59494) (-1.34396,0.41307) (-1.22029,0.24681) (-1.09246,0.09688) (-0.96297,-0.03635) (-0.83407,-0.15281) (-0.70772,-0.25270) (-0.58559,-0.33647) (-0.46909,-0.40477) (-0.35937,-0.45838) (-0.25734,-0.49827) (-0.16366,-0.52546) (-0.07880,-0.54108) (-0.00301,-0.54628) (0.06361,-0.54226) (0.12113,-0.53019) (0.16975,-0.51124) (0.20979,-0.48655) (0.24165,-0.45720) (0.26584,-0.42423) (0.28291,-0.38859) (0.29345,-0.35116) (0.29809,-0.31276) (0.29747,-0.27411) (0.29225,-0.23584) (0.28306,-0.19851) (0.27053,-0.16260) (0.25527,-0.12849) (0.23785,-0.09651) (0.21880,-0.06690) (0.19863,-0.03983) (0.17778,-0.01543) (0.15667,0.00624) (0.13566,0.02519) (0.11506,0.04143) (0.09517,0.05505) (0.07619,0.06615) (0.05832,0.07485) (0.04171,0.08132) (0.02645,0.08572) (0.01264,0.08824) (0.00031,0.08906) (-0.01053,0.08839) (-0.01988,0.08640) (-0.02779,0.08329) (-0.03429,0.07925) (-0.03947,0.07446) (-0.04339,0.06908) (-0.04616,0.06326) (-0.04786,0.05715) (-0.04860,0.05089)};\draw[blue,->] (axis cs:-0.75244,-0.21890)--(axis cs:-0.63766,-0.30227);\addplot[blue,no marks] coordinates {(1.45022,0.34303) (1.22115,0.51609) (1.00071,0.66000) (0.79132,0.77602) (0.59494,0.86572) (0.41307,0.93089) (0.24681,0.97348) (0.09688,0.99558) (-0.03635,0.99932) (-0.15281,0.98688) (-0.25270,0.96042) (-0.33647,0.92206) (-0.40477,0.87386) (-0.45838,0.81776) (-0.49827,0.75561) (-0.52546,0.68913) (-0.54108,0.61988) (-0.54628,0.54930) (-0.54226,0.47865) (-0.53019,0.40906) (-0.51124,0.34149) (-0.48655,0.27676) (-0.45720,0.21555) (-0.42423,0.15838) (-0.38859,0.10568) (-0.35116,0.05772) (-0.31276,0.01468) (-0.27411,-0.02336) (-0.23584,-0.05640) (-0.19851,-0.08454) (-0.16260,-0.10793) (-0.12849,-0.12678) (-0.09651,-0.14134) (-0.06690,-0.15190) (-0.03983,-0.15879) (-0.01543,-0.16234) (0.00624,-0.16291) (0.02519,-0.16084) (0.04143,-0.15650) (0.05505,-0.15022) (0.06615,-0.14233) (0.07485,-0.13317) (0.08132,-0.12302) (0.08572,-0.11218) (0.08824,-0.10088) (0.08906,-0.08937) (0.08839,-0.07786) (0.08640,-0.06651) (0.08329,-0.05551) (0.07925,-0.04496) (0.07446,-0.03499) (0.06908,-0.02568) (0.06326,-0.01710) (0.05715,-0.00929) (0.05089,-0.00229)};\draw[blue,->] (axis cs:-0.21890,0.97135)--(axis cs:-0.30227,0.93994);\addplot[blue,no marks] coordinates {(-1.45022,-0.34303) (-1.22115,-0.51609) (-1.00071,-0.66000) (-0.79132,-0.77602) (-0.59494,-0.86572) (-0.41307,-0.93089) (-0.24681,-0.97348) (-0.09688,-0.99558) (0.03635,-0.99932) (0.15281,-0.98688) (0.25270,-0.96042) (0.33647,-0.92206) (0.40477,-0.87386) (0.45838,-0.81776) (0.49827,-0.75561) (0.52546,-0.68913) (0.54108,-0.61988) (0.54628,-0.54930) (0.54226,-0.47865) (0.53019,-0.40906) (0.51124,-0.34149) (0.48655,-0.27676) (0.45720,-0.21555) (0.42423,-0.15838) (0.38859,-0.10568) (0.35116,-0.05772) (0.31276,-0.01468) (0.27411,0.02336) (0.23584,0.05640) (0.19851,0.08454) (0.16260,0.10793) (0.12849,0.12678) (0.09651,0.14134) (0.06690,0.15190) (0.03983,0.15879) (0.01543,0.16234) (-0.00624,0.16291) (-0.02519,0.16084) (-0.04143,0.15650) (-0.05505,0.15022) (-0.06615,0.14233) (-0.07485,0.13317) (-0.08132,0.12302) (-0.08572,0.11218) (-0.08824,0.10088) (-0.08906,0.08937) (-0.08839,0.07786) (-0.08640,0.06651) (-0.08329,0.05551) (-0.07925,0.04496) (-0.07446,0.03499) (-0.06908,0.02568) (-0.06326,0.01710) (-0.05715,0.00929) (-0.05089,0.00229)};\draw[blue,->] (axis cs:0.21890,-0.97135)--(axis cs:0.30227,-0.93994);\addplot[blue,no marks] coordinates {(3.24347,-1.10719) (2.95840,-0.70506) (2.66142,-0.34072) (2.35867,-0.01531) (2.05560,0.27078) (1.75702,0.51782) (1.46710,0.72668) (1.18933,0.89870) (0.92662,1.03567) (0.68126,1.13969) (0.45502,1.21312) (0.24911,1.25854) (0.06433,1.27862) (-0.09901,1.27614) (-0.24093,1.25388) (-0.36180,1.21459) (-0.46228,1.16096) (-0.54327,1.09558) (-0.60587,1.02091) (-0.65132,0.93924) (-0.68099,0.85273) (-0.69634,0.76331) (-0.69886,0.67275) (-0.69007,0.58261) (-0.67150,0.49427) (-0.64461,0.40888) (-0.61085,0.32744) (-0.57158,0.25075) (-0.52809,0.17944) (-0.48157,0.11397) (-0.43313,0.05467) (-0.38377,0.00172) (-0.33436,-0.04483) (-0.28570,-0.08500) (-0.23846,-0.11896) (-0.19321,-0.14691) (-0.15042,-0.16916) (-0.11047,-0.18603) (-0.07363,-0.19793) (-0.04012,-0.20527) (-0.01004,-0.20848) (0.01653,-0.20802) (0.03961,-0.20434) (0.05927,-0.19790) (0.07560,-0.18912) (0.08875,-0.17844) (0.09891,-0.16624) (0.10628,-0.15291) (0.11108,-0.13880) (0.11355,-0.12421) (0.11393,-0.10945) (0.11247,-0.09476) (0.10942,-0.08036) (0.10502,-0.06645) (0.09950,-0.05318)};\draw[blue,->] (axis cs:0.53354,1.19025)--(axis cs:0.33539,1.24221);\addplot[blue,no marks] coordinates {(-3.24347,1.10719) (-2.95840,0.70506) (-2.66142,0.34072) (-2.35867,0.01531) (-2.05560,-0.27078) (-1.75702,-0.51782) (-1.46710,-0.72668) (-1.18933,-0.89870) (-0.92662,-1.03567) (-0.68126,-1.13969) (-0.45502,-1.21312) (-0.24911,-1.25854) (-0.06433,-1.27862) (0.09901,-1.27614) (0.24093,-1.25388) (0.36180,-1.21459) (0.46228,-1.16096) (0.54327,-1.09558) (0.60587,-1.02091) (0.65132,-0.93924) (0.68099,-0.85273) (0.69634,-0.76331) (0.69886,-0.67275) (0.69007,-0.58261) (0.67150,-0.49427) (0.64461,-0.40888) (0.61085,-0.32744) (0.57158,-0.25075) (0.52809,-0.17944) (0.48157,-0.11397) (0.43313,-0.05467) (0.38377,-0.00172) (0.33436,0.04483) (0.28570,0.08500) (0.23846,0.11896) (0.19321,0.14691) (0.15042,0.16916) (0.11047,0.18603) (0.07363,0.19793) (0.04012,0.20527) (0.01004,0.20848) (-0.01653,0.20802) (-0.03961,0.20434) (-0.05927,0.19790) (-0.07560,0.18912) (-0.08875,0.17844) (-0.09891,0.16624) (-0.10628,0.15291) (-0.11108,0.13880) (-0.11355,0.12421) (-0.11393,0.10945) (-0.11247,0.09476) (-0.10942,0.08036) (-0.10502,0.06645) (-0.09950,0.05318)};\draw[blue,->] (axis cs:-0.53354,-1.19025)--(axis cs:-0.33539,-1.24221);\end{axis}\end{tikzpicture}\end{center}
\par\noindent{\small 依据：官方译审并补证（AI辅助）。}
\end{solution}

\begin{exercise}\label{cw:rec25-3}
在虫--鸟平面标出仅有昆虫和仅有鸟时的相线；结合线性化画完整种群相图。
\end{exercise}
\begin{solution}
坐标轴不变：$y=0$时$\dot x=x(2-x)$，$0<x<2$向右，$x>2$向左；$x=0$时$\dot y=-y$，向原点。第一象限内，$x+y<2$时$x$增加，$x+y>2$时减少；$x>1$时$y$增加，$x<1$时减少。
不仅局部可以判断，所有正初值都趋于$(1,1)$。取
\[V=x-1-\log x+y-1-\log y\ge0,\qquad\dot V=-(x-1)^2\le0.\]
正初值的子水平集在第一象限内紧致，远离坐标轴，所以解全局存在并有非空紧的极限集。该极限集必须位于$\dot V=0$的$x=1$上：否则负导数邻域的反复进入会使$V$下降超过其有限下界。极限集又随流不变，而$x=1$上$\dot x=1-y$，仅$y=1$能留在线上。因此唯一极限点为$(1,1)$。图中蓝线为方程数值积分所得示意轨迹，虚线为$x=1$及$x+y=2$；收敛由上述论证支持。
\begin{center}\begin{tikzpicture}\begin{axis}[width=.65\textwidth,height=5cm,axis lines=middle,xlabel={昆虫$x$},ylabel={鸟类$y$},xmin=0,xmax=2.7,ymin=0,ymax=2.8,clip=true]\addplot[blue]coordinates{(0.80000,0.40000)(0.86462,0.39330)(0.92892,0.38924)(0.99188,0.38769)(1.05257,0.38857)(1.11010,0.39178)(1.16373,0.39723)(1.21285,0.40485)(1.25701,0.41456)(1.29594,0.42629)(1.32950,0.43995)(1.35771,0.45546)(1.38066,0.47274)(1.39856,0.49168)(1.41168,0.51218)(1.42032,0.53411)(1.42483,0.55735)(1.42555,0.58176)(1.42285,0.60717)(1.41707,0.63342)(1.40855,0.66032)(1.39763,0.68769)(1.38463,0.71534)(1.36985,0.74305)(1.35357,0.77063)(1.33607,0.79787)(1.31761,0.82458)(1.29841,0.85056)(1.27871,0.87565)(1.25872,0.89967)(1.23861,0.92249)(1.21857,0.94398)(1.19876,0.96402)(1.17931,0.98255)(1.16035,0.99950)(1.14198,1.01483)(1.12430,1.02852)(1.10739,1.04058)(1.09131,1.05104)(1.07611,1.05993)(1.06183,1.06731)(1.04850,1.07325)(1.03612,1.07783)(1.02472,1.08113)(1.01427,1.08324)(1.00479,1.08427)(0.99624,1.08432)(0.98861,1.08349)(0.98187,1.08187)(0.97599,1.07957)(0.97094,1.07668)(0.96667,1.07329)(0.96314,1.06950)(0.96032,1.06538)(0.95816,1.06101)(0.95661,1.05646)(0.95562,1.05180)(0.95516,1.04708)(0.95518,1.04235)(0.95562,1.03768)(0.95645,1.03309)(0.95762,1.02863)(0.95909,1.02432)(0.96081,1.02019)(0.96275,1.01627)(0.96487,1.01257)(0.96712,1.00911)(0.96949,1.00589)(0.97192,1.00293)(0.97440,1.00022)(0.97689,0.99777)(0.97938,0.99558)(0.98183,0.99364)(0.98422,0.99194)(0.98655,0.99048)(0.98878,0.98925)(0.99092,0.98824)(0.99294,0.98744)(0.99484,0.98683)(0.99661,0.98641)(0.99825,0.98616)(0.99975,0.98606)(1.00111,0.98610)(1.00233,0.98627)(1.00341,0.98656)(1.00436,0.98695)(1.00517,0.98742)(1.00586,0.98797)(1.00642,0.98858)(1.00687,0.98924)(1.00721,0.98995)(1.00745,0.99068)(1.00759,0.99143)(1.00765,0.99219)(1.00763,0.99296)(1.00753,0.99372)(1.00738,0.99446)(1.00716,0.99519)(1.00690,0.99590)(1.00660,0.99658)(1.00626,0.99722)(1.00589,0.99783)(1.00551,0.99841)(1.00510,0.99894)(1.00469,0.99943)(1.00427,0.99988)(1.00385,1.00029)(1.00343,1.00066)(1.00302,1.00099)(1.00263,1.00127)(1.00224,1.00152)(1.00187,1.00172)(1.00151,1.00189)(1.00118,1.00203)(1.00087,1.00213)(1.00058,1.00220)(1.00031,1.00225)(1.00006,1.00227)(0.99984,1.00226)(0.99964,1.00224)(0.99946,1.00219)(0.99931,1.00213)(0.99917,1.00205)(0.99906,1.00196)(0.99897,1.00186)(0.99889,1.00175)(0.99884,1.00164)(0.99880,1.00152)(0.99877,1.00140)(0.99876,1.00127)(0.99877,1.00115)(0.99878,1.00102)(0.99880,1.00090)(0.99884,1.00078)(0.99888,1.00067)(0.99893,1.00056)(0.99898,1.00045)(0.99904,1.00035)(0.99910,1.00026)(0.99917,1.00017)};\draw[blue,->](axis cs:0.99055,0.38770)--(axis cs:1.10811,0.39162);\addplot[blue]coordinates{(2.00000,0.80000)(1.85199,0.87802)(1.72557,0.95047)(1.61565,1.01676)(1.51885,1.07644)(1.43289,1.12919)(1.35616,1.17488)(1.28748,1.21351)(1.22596,1.24523)(1.17091,1.27030)(1.12172,1.28910)(1.07792,1.30207)(1.03908,1.30971)(1.00480,1.31256)(0.97473,1.31116)(0.94856,1.30606)(0.92598,1.29780)(0.90673,1.28688)(0.89054,1.27378)(0.87718,1.25893)(0.86642,1.24276)(0.85806,1.22561)(0.85189,1.20781)(0.84775,1.18966)(0.84545,1.17141)(0.84483,1.15326)(0.84572,1.13541)(0.84799,1.11802)(0.85148,1.10122)(0.85606,1.08511)(0.86158,1.06978)(0.86793,1.05530)(0.87496,1.04172)(0.88257,1.02907)(0.89063,1.01738)(0.89904,1.00666)(0.90768,0.99691)(0.91645,0.98812)(0.92526,0.98027)(0.93403,0.97335)(0.94266,0.96733)(0.95110,0.96217)(0.95926,0.95784)(0.96710,0.95429)(0.97456,0.95150)(0.98160,0.94940)(0.98819,0.94796)(0.99430,0.94713)(0.99992,0.94686)(1.00502,0.94710)(1.00961,0.94780)(1.01368,0.94892)(1.01725,0.95040)(1.02032,0.95221)(1.02290,0.95428)(1.02502,0.95659)(1.02671,0.95909)(1.02797,0.96174)(1.02885,0.96450)(1.02936,0.96733)(1.02954,0.97021)(1.02942,0.97310)(1.02903,0.97597)(1.02840,0.97880)(1.02755,0.98156)(1.02652,0.98424)(1.02533,0.98681)(1.02402,0.98927)(1.02260,0.99159)(1.02110,0.99378)(1.01955,0.99582)(1.01795,0.99770)(1.01634,0.99942)(1.01473,1.00099)(1.01314,1.00239)(1.01157,1.00364)(1.01004,1.00473)(1.00856,1.00568)(1.00715,1.00647)(1.00580,1.00713)(1.00453,1.00765)(1.00334,1.00805)(1.00223,1.00833)(1.00121,1.00851)(1.00027,1.00858)(0.99942,1.00856)(0.99866,1.00847)(0.99799,1.00830)(0.99740,1.00806)(0.99689,1.00777)(0.99646,1.00743)(0.99611,1.00705)(0.99583,1.00665)(0.99561,1.00621)(0.99546,1.00576)(0.99537,1.00529)(0.99533,1.00482)(0.99533,1.00435)(0.99539,1.00388)(0.99548,1.00342)(0.99561,1.00297)(0.99576,1.00253)(0.99595,1.00211)(0.99615,1.00171)(0.99637,1.00134)(0.99661,1.00098)(0.99686,1.00065)(0.99711,1.00035)(0.99737,1.00007)(0.99762,0.99982)(0.99788,0.99959)(0.99813,0.99939)(0.99838,0.99922)(0.99861,0.99906)(0.99884,0.99894)(0.99906,0.99883)(0.99927,0.99875)(0.99946,0.99868)(0.99964,0.99864)(0.99981,0.99861)(0.99996,0.99860)(1.00010,0.99860)(1.00022,0.99862)(1.00033,0.99865)(1.00043,0.99868)(1.00051,0.99873)(1.00058,0.99879)(1.00064,0.99885)(1.00068,0.99891)(1.00072,0.99899)(1.00074,0.99906)(1.00076,0.99913)(1.00077,0.99921)(1.00076,0.99929)(1.00076,0.99937)(1.00074,0.99944)(1.00072,0.99951)(1.00069,0.99959)(1.00066,0.99965)(1.00063,0.99972)};\draw[blue,->](axis cs:1.61785,1.01541)--(axis cs:1.43579,1.12743);\addplot[blue]coordinates{(0.30000,1.80000)(0.29888,1.67732)(0.30132,1.56310)(0.30704,1.45727)(0.31585,1.35960)(0.32766,1.26979)(0.34241,1.18750)(0.36008,1.11236)(0.38068,1.04399)(0.40420,0.98199)(0.43065,0.92601)(0.46000,0.87567)(0.49218,0.83064)(0.52709,0.79060)(0.56456,0.75523)(0.60436,0.72426)(0.64618,0.69742)(0.68966,0.67448)(0.73437,0.65519)(0.77980,0.63935)(0.82542,0.62676)(0.87070,0.61725)(0.91505,0.61063)(0.95795,0.60675)(0.99889,0.60544)(1.03743,0.60656)(1.07318,0.60997)(1.10584,0.61551)(1.13519,0.62304)(1.16109,0.63242)(1.18347,0.64351)(1.20235,0.65616)(1.21777,0.67021)(1.22987,0.68551)(1.23880,0.70190)(1.24474,0.71922)(1.24791,0.73730)(1.24854,0.75598)(1.24685,0.77509)(1.24310,0.79447)(1.23752,0.81394)(1.23035,0.83336)(1.22181,0.85256)(1.21212,0.87140)(1.20148,0.88975)(1.19010,0.90748)(1.17815,0.92447)(1.16580,0.94062)(1.15321,0.95586)(1.14051,0.97010)(1.12784,0.98330)(1.11530,0.99541)(1.10299,1.00641)(1.09102,1.01629)(1.07944,1.02505)(1.06832,1.03270)(1.05771,1.03927)(1.04767,1.04479)(1.03821,1.04931)(1.02937,1.05289)(1.02116,1.05556)(1.01359,1.05741)(1.00666,1.05848)(1.00038,1.05885)(0.99472,1.05858)(0.98969,1.05775)(0.98526,1.05641)(0.98141,1.05463)(0.97813,1.05248)(0.97537,1.05001)(0.97313,1.04729)(0.97136,1.04436)(0.97004,1.04128)(0.96914,1.03809)(0.96862,1.03484)(0.96845,1.03156)(0.96861,1.02830)(0.96905,1.02507)(0.96974,1.02191)(0.97067,1.01885)(0.97178,1.01590)(0.97307,1.01308)(0.97449,1.01041)(0.97602,1.00789)(0.97764,1.00554)(0.97933,1.00337)(0.98105,1.00137)(0.98280,0.99954)(0.98454,0.99790)(0.98627,0.99644)(0.98797,0.99514)(0.98963,0.99402)(0.99123,0.99307)(0.99276,0.99227)(0.99422,0.99162)(0.99560,0.99111)(0.99688,0.99073)(0.99808,0.99048)(0.99918,0.99035)(1.00019,0.99032)(1.00110,0.99038)(1.00191,0.99053)(1.00262,0.99076)(1.00324,0.99105)(1.00378,0.99140)(1.00422,0.99180)(1.00458,0.99224)(1.00486,0.99272)(1.00507,0.99321)(1.00521,0.99373)(1.00529,0.99426)(1.00531,0.99479)(1.00527,0.99532)(1.00519,0.99584)(1.00506,0.99636)(1.00490,0.99686)(1.00471,0.99734)(1.00449,0.99780)(1.00425,0.99824)(1.00399,0.99866)(1.00371,0.99904)(1.00343,0.99940)(1.00314,0.99973)(1.00285,1.00004)(1.00256,1.00031)(1.00227,1.00055)(1.00199,1.00077)(1.00172,1.00095)(1.00145,1.00111)(1.00120,1.00125)(1.00096,1.00136)(1.00073,1.00144)(1.00052,1.00150)(1.00032,1.00155)(1.00014,1.00157)(0.99998,1.00158)(0.99983,1.00157)(0.99970,1.00154)(0.99958,1.00151)(0.99948,1.00146)};\draw[blue,->](axis cs:0.30688,1.45945)--(axis cs:0.32719,1.27286);\addplot[blue]coordinates{(2.50000,2.50000)(1.87883,2.81240)(1.45044,3.00332)(1.14532,3.09220)(0.92406,3.10121)(0.76165,3.05125)(0.64134,2.95999)(0.55159,2.84148)(0.48432,2.70638)(0.43384,2.56257)(0.39611,2.41571)(0.36826,2.26980)(0.34820,2.12760)(0.33447,1.99094)(0.32598,1.86099)(0.32198,1.73843)(0.32191,1.62362)(0.32540,1.51666)(0.33219,1.41748)(0.34211,1.32590)(0.35504,1.24167)(0.37094,1.16448)(0.38975,1.09400)(0.41146,1.02988)(0.43604,0.97178)(0.46346,0.91937)(0.49363,0.87231)(0.52646,0.83029)(0.56179,0.79301)(0.59940,0.76020)(0.63901,0.73158)(0.68029,0.70692)(0.72284,0.68598)(0.76621,0.66855)(0.80991,0.65443)(0.85343,0.64343)(0.89624,0.63538)(0.93782,0.63010)(0.97771,0.62743)(1.01545,0.62723)(1.05068,0.62934)(1.08309,0.63361)(1.11243,0.63989)(1.13855,0.64805)(1.16135,0.65793)(1.18083,0.66938)(1.19700,0.68226)(1.20998,0.69641)(1.21987,0.71166)(1.22686,0.72787)(1.23113,0.74487)(1.23288,0.76250)(1.23235,0.78059)(1.22975,0.79898)(1.22531,0.81751)(1.21925,0.83603)(1.21180,0.85439)(1.20316,0.87244)(1.19354,0.89005)(1.18313,0.90710)(1.17210,0.92348)(1.16061,0.93908)(1.14884,0.95383)(1.13690,0.96766)(1.12493,0.98050)(1.11304,0.99232)(1.10134,1.00309)(1.08990,1.01280)(1.07881,1.02144)(1.06814,1.02902)(1.05793,1.03557)(1.04823,1.04111)(1.03907,1.04570)(1.03049,1.04936)(1.02250,1.05216)(1.01512,1.05415)(1.00834,1.05539)(1.00217,1.05594)(0.99660,1.05587)(0.99163,1.05524)(0.98724,1.05411)(0.98340,1.05255)(0.98011,1.05061)(0.97733,1.04836)(0.97504,1.04584)(0.97322,1.04312)(0.97183,1.04023)(0.97085,1.03723)(0.97024,1.03415)(0.96998,1.03104)(0.97003,1.02793)(0.97037,1.02485)(0.97096,1.02182)(0.97177,1.01888)(0.97278,1.01603)(0.97395,1.01331)(0.97527,1.01072)(0.97670,1.00828)(0.97822,1.00599)(0.97980,1.00386)(0.98143,1.00191)(0.98309,1.00012)(0.98475,0.99850)(0.98641,0.99705)(0.98803,0.99577)(0.98963,0.99465)(0.99117,0.99369)(0.99265,0.99288)(0.99406,0.99221)(0.99539,0.99169)(0.99665,0.99129)(0.99781,0.99102)(0.99889,0.99085)(0.99988,0.99079)(1.00077,0.99082)(1.00157,0.99094)(1.00228,0.99113)(1.00290,0.99139)(1.00343,0.99171)(1.00388,0.99208)(1.00425,0.99249)(1.00455,0.99293)(1.00477,0.99339)(1.00492,0.99388)(1.00501,0.99437)(1.00505,0.99488)(1.00503,0.99538)(1.00496,0.99589)(1.00486,0.99638)(1.00472,0.99686)(1.00454,0.99732)(1.00434,0.99777)(1.00412,0.99820)(1.00387,0.99860)(1.00362,0.99897)(1.00335,0.99933)(1.00308,0.99965)(1.00280,0.99994)(1.00252,1.00021)(1.00225,1.00045)};\draw[blue,->](axis cs:1.15087,3.09121)--(axis cs:0.76664,3.05384);\addplot[dashed,domain=0:2]{2-x};\draw[dashed](axis cs:1,0)--(axis cs:1,2.8);\addplot[only marks]coordinates{(0,0)(2,0)(1,1)};\draw[->,thick](axis cs:.3,0)--(axis cs:1.2,0);\draw[->,thick](axis cs:2.7,0)--(axis cs:2.1,0);\draw[->,thick](axis cs:0,2.3)--(axis cs:0,1.1);\end{axis}\end{tikzpicture}\end{center}
\par\noindent{\small 依据：官方相线译审；全局收敛为编者补解（AI辅助）。}
\end{solution}

\begin{exercise}\label{cw:rec25-4}
选择正初值，由相图画$x(t),y(t)$。它们的长期极限为何？以何种指数率靠近？是否越过极限并振荡？若是，伪周期为何？
\end{exercise}
\begin{solution}
选$x(0)=0.8,y(0)=0.4$。由上一题知两者均趋于1。近$(1,1)$，扰动$z$满足$\dot z=Jz+O(|z|^2)$，$J$的根为$-1/2\pm\ii\sqrt3/2$。因此主导衰减率为$\ee^{-t/2}$，渐近振荡圆频率$\sqrt3/2$，伪周期$4\pi/\sqrt3$。除恰在平衡点的初值外，近点轨迹旋转使两个坐标反复过冲和欠冲；早期振荡周期不必完全相等。
为说明线性化指数不只是图上猜测，先作可逆实线性变换，将$J$化为$-I/2+\beta\begin{pmatrix}0&-1\\1&0\end{pmatrix}$，其中$\beta=\sqrt3/2$。变换后的极坐标满足$\dot r=-r/2+O(r^2),\dot\theta=\beta+O(r)$。小扰动下$r=O(\ee^{-t/4})$可积，故$\log r+t/2$的导数$O(r)$可积，非零解有$r\ee^{t/2}\to C>0$，而$\theta-\beta t$趋于常数。可逆线性变换将圆转为椭圆，两原坐标因而都有所述主导衰减和渐近频率；平衡解本身没有振荡。
\begin{center}\begin{tikzpicture}\begin{axis}[width=.65\textwidth,height=4cm,axis lines=middle,xlabel={$t$},ylabel={种群},xmin=0,xmax=18,ymin=0,ymax=1.7]\addplot[black]coordinates{(0.00000,0.80000)(0.12081,0.87750)(0.24161,0.95424)(0.36242,1.02852)(0.48322,1.09876)(0.60403,1.16358)(0.72483,1.22193)(0.84564,1.27306)(0.96644,1.31657)(1.08725,1.35235)(1.20805,1.38054)(1.32886,1.40146)(1.44966,1.41559)(1.57047,1.42347)(1.69128,1.42570)(1.81208,1.42289)(1.93289,1.41564)(2.05369,1.40456)(2.17450,1.39020)(2.29530,1.37310)(2.41611,1.35377)(2.53691,1.33267)(2.65772,1.31024)(2.77852,1.28690)(2.89933,1.26301)(3.02013,1.23890)(3.14094,1.21489)(3.26174,1.19123)(3.38255,1.16817)(3.50336,1.14591)(3.62416,1.12460)(3.74497,1.10440)(3.86577,1.08541)(3.98658,1.06770)(4.10738,1.05134)(4.22819,1.03636)(4.34899,1.02277)(4.46980,1.01057)(4.59060,0.99973)(4.71141,0.99023)(4.83221,0.98202)(4.95302,0.97504)(5.07383,0.96924)(5.19463,0.96455)(5.31544,0.96090)(5.43624,0.95820)(5.55705,0.95640)(5.67785,0.95539)(5.79866,0.95512)(5.91946,0.95548)(6.04027,0.95642)(6.16107,0.95785)(6.28188,0.95970)(6.40268,0.96189)(6.52349,0.96437)(6.64430,0.96705)(6.76510,0.96989)(6.88591,0.97283)(7.00671,0.97581)(7.12752,0.97880)(7.24832,0.98174)(7.36913,0.98461)(7.48993,0.98737)(7.61074,0.99000)(7.73154,0.99247)(7.85235,0.99477)(7.97315,0.99689)(8.09396,0.99881)(8.21477,1.00053)(8.33557,1.00205)(8.45638,1.00337)(8.57718,1.00449)(8.69799,1.00543)(8.81879,1.00619)(8.93960,1.00677)(9.06040,1.00720)(9.18121,1.00748)(9.30201,1.00762)(9.42282,1.00765)(9.54362,1.00756)(9.66443,1.00739)(9.78523,1.00713)(9.90604,1.00680)(10.02685,1.00642)(10.14765,1.00599)(10.26846,1.00553)(10.38926,1.00504)(10.51007,1.00454)(10.63087,1.00404)(10.75168,1.00354)(10.87248,1.00305)(10.99329,1.00257)(11.11409,1.00211)(11.23490,1.00167)(11.35570,1.00126)(11.47651,1.00088)(11.59732,1.00054)(11.71812,1.00022)(11.83893,0.99994)(11.95973,0.99969)(12.08054,0.99947)(12.20134,0.99929)(12.32215,0.99913)(12.44295,0.99901)(12.56376,0.99891)(12.68456,0.99884)(12.80537,0.99879)(12.92617,0.99877)(13.04698,0.99876)(13.16779,0.99877)(13.28859,0.99880)(13.40940,0.99884)(13.53020,0.99890)(13.65101,0.99896)(13.77181,0.99903)(13.89262,0.99910)(14.01342,0.99918)(14.13423,0.99926)(14.25503,0.99934)(14.37584,0.99942)(14.49664,0.99950)(14.61745,0.99958)(14.73826,0.99966)(14.85906,0.99973)(14.97987,0.99979)(15.10067,0.99985)(15.22148,0.99991)(15.34228,0.99996)(15.46309,1.00001)(15.58389,1.00005)(15.70470,1.00009)(15.82550,1.00012)(15.94631,1.00014)(16.06711,1.00016)(16.18792,1.00018)(16.30872,1.00019)(16.42953,1.00020)(16.55034,1.00020)(16.67114,1.00020)(16.79195,1.00020)(16.91275,1.00020)(17.03356,1.00019)(17.15436,1.00018)(17.27517,1.00017)(17.39597,1.00016)(17.51678,1.00015)(17.63758,1.00013)(17.75839,1.00012)(17.87919,1.00011)(18.00000,1.00009)};\addplot[blue]coordinates{(0.00000,0.40000)(0.12081,0.39228)(0.24161,0.38832)(0.36242,0.38793)(0.48322,0.39094)(0.60403,0.39721)(0.72483,0.40660)(0.84564,0.41897)(0.96644,0.43420)(1.08725,0.45213)(1.20805,0.47263)(1.32886,0.49553)(1.44966,0.52063)(1.57047,0.54773)(1.69128,0.57659)(1.81208,0.60694)(1.93289,0.63850)(2.05369,0.67096)(2.17450,0.70397)(2.29530,0.73721)(2.41611,0.77031)(2.53691,0.80294)(2.65772,0.83474)(2.77852,0.86540)(2.89933,0.89463)(3.02013,0.92217)(3.14094,0.94780)(3.26174,0.97133)(3.38255,0.99264)(3.50336,1.01164)(3.62416,1.02830)(3.74497,1.04261)(3.86577,1.05462)(3.98658,1.06440)(4.10738,1.07207)(4.22819,1.07775)(4.34899,1.08159)(4.46980,1.08375)(4.59060,1.08441)(4.71141,1.08374)(4.83221,1.08191)(4.95302,1.07910)(5.07383,1.07546)(5.19463,1.07116)(5.31544,1.06633)(5.43624,1.06112)(5.55705,1.05566)(5.67785,1.05004)(5.79866,1.04437)(5.91946,1.03874)(6.04027,1.03322)(6.16107,1.02788)(6.28188,1.02277)(6.40268,1.01794)(6.52349,1.01341)(6.64430,1.00922)(6.76510,1.00538)(6.88591,1.00191)(7.00671,0.99880)(7.12752,0.99607)(7.24832,0.99370)(7.36913,0.99168)(7.48993,0.99001)(7.61074,0.98866)(7.73154,0.98761)(7.85235,0.98685)(7.97315,0.98636)(8.09396,0.98610)(8.21477,0.98606)(8.33557,0.98622)(8.45638,0.98654)(8.57718,0.98702)(8.69799,0.98761)(8.81879,0.98830)(8.93960,0.98908)(9.06040,0.98992)(9.18121,0.99080)(9.30201,0.99170)(9.42282,0.99262)(9.54362,0.99353)(9.66443,0.99443)(9.78523,0.99530)(9.90604,0.99614)(10.02685,0.99694)(10.14765,0.99768)(10.26846,0.99838)(10.38926,0.99902)(10.51007,0.99960)(10.63087,1.00011)(10.75168,1.00057)(10.87248,1.00097)(10.99329,1.00131)(11.11409,1.00159)(11.23490,1.00182)(11.35570,1.00200)(11.47651,1.00213)(11.59732,1.00221)(11.71812,1.00226)(11.83893,1.00227)(11.95973,1.00224)(12.08054,1.00219)(12.20134,1.00212)(12.32215,1.00202)(12.44295,1.00191)(12.56376,1.00178)(12.68456,1.00165)(12.80537,1.00150)(12.92617,1.00135)(13.04698,1.00120)(13.16779,1.00106)(13.28859,1.00091)(13.40940,1.00077)(13.53020,1.00063)(13.65101,1.00050)(13.77181,1.00038)(13.89262,1.00026)(14.01342,1.00016)(14.13423,1.00007)(14.25503,0.99998)(14.37584,0.99991)(14.49664,0.99984)(14.61745,0.99979)(14.73826,0.99974)(14.85906,0.99970)(14.97987,0.99967)(15.10067,0.99965)(15.22148,0.99964)(15.34228,0.99963)(15.46309,0.99963)(15.58389,0.99963)(15.70470,0.99964)(15.82550,0.99965)(15.94631,0.99967)(16.06711,0.99969)(16.18792,0.99971)(16.30872,0.99973)(16.42953,0.99975)(16.55034,0.99978)(16.67114,0.99980)(16.79195,0.99983)(16.91275,0.99985)(17.03356,0.99987)(17.15436,0.99990)(17.27517,0.99992)(17.39597,0.99994)(17.51678,0.99996)(17.63758,0.99997)(17.75839,0.99999)(17.87919,1.00000)(18.00000,1.00001)};\addplot[dashed,domain=0:18]{1};\end{axis}\end{tikzpicture}\end{center}黑线为$x$，蓝线为$y$。
\par\noindent{\small 依据：官方译审并补证（AI辅助）；区分渐近伪周期与全程固定周期。}
\end{solution}

\begin{exercise}\label{cw:rec25-5}
引入马拉硫磷试图减少虫口，两物种繁殖率均下降，系统变为$\dot x=(2-x-y-a)x,\dot y=(x-1-b)y$，$a,b$为小正数。正共存平衡如何变化？是否成功减少昆虫？
\end{exercise}
\begin{solution}
正平衡必须满足$x=1+b$、$y=2-a-x=1-a-b$。共存条件$a+b<1$；在此条件下昆虫从1增加为$1+b$，鸟从1降为$1-a-b$，所以按长期共存平衡标准，措施反而增加昆虫。若$a+b\ge1$则正共存平衡消失，不能继续沿用这个结论，已超出题设“小正数”的共存情形。
\par\noindent{\small 依据：官方译审并补证（AI辅助）。}
\end{solution}


% Figures independently redrawn from equations or actually inspected source graphs.
\newcommand{\ODEExamRiccatiField}{%
\begin{center}\begin{tikzpicture}[x=.62cm,y=.53cm]
\clip (-1.6,-4.6) rectangle (6.7,4.8);
\draw[->](-1,0)--(6.35,0)node[right]{$x$};\draw[->](0,-4)--(0,4.35)node[above]{$y$};
\draw[gray!55](-1.027,-3.883)--(-0.973,-4.117);
\draw[gray!55](-1.035,-3.385)--(-0.965,-3.615);
\draw[gray!55](-1.045,-2.889)--(-0.955,-3.111);
\draw[gray!55](-1.058,-2.395)--(-0.942,-2.605);
\draw[gray!55](-1.075,-1.906)--(-0.925,-2.094);
\draw[gray!55](-1.093,-1.424)--(-0.907,-1.576);
\draw[gray!55](-1.107,-0.946)--(-0.893,-1.054);
\draw[gray!55](-1.115,-0.464)--(-0.885,-0.536);
\draw[gray!55](-1.116,0.029)--(-0.884,-0.029);
\draw[gray!55](-1.115,0.536)--(-0.885,0.464);
\draw[gray!55](-1.107,1.054)--(-0.893,0.946);
\draw[gray!55](-1.093,1.576)--(-0.907,1.424);
\draw[gray!55](-1.075,2.094)--(-0.925,1.906);
\draw[gray!55](-1.058,2.605)--(-0.942,2.395);
\draw[gray!55](-1.045,3.111)--(-0.955,2.889);
\draw[gray!55](-1.035,3.615)--(-0.965,3.385);
\draw[gray!55](-1.027,4.117)--(-0.973,3.883);
\draw[gray!55](-0.528,-3.883)--(-0.472,-4.117);
\draw[gray!55](-0.536,-3.386)--(-0.464,-3.614);
\draw[gray!55](-0.547,-2.889)--(-0.453,-3.111);
\draw[gray!55](-0.561,-2.397)--(-0.439,-2.603);
\draw[gray!55](-0.580,-1.910)--(-0.420,-2.090);
\draw[gray!55](-0.599,-1.432)--(-0.401,-1.568);
\draw[gray!55](-0.612,-0.958)--(-0.388,-1.042);
\draw[gray!55](-0.618,-0.478)--(-0.382,-0.522);
\draw[gray!55](-0.619,0.015)--(-0.381,-0.015);
\draw[gray!55](-0.618,0.522)--(-0.382,0.478);
\draw[gray!55](-0.612,1.042)--(-0.388,0.958);
\draw[gray!55](-0.599,1.568)--(-0.401,1.432);
\draw[gray!55](-0.580,2.090)--(-0.420,1.910);
\draw[gray!55](-0.561,2.603)--(-0.439,2.397);
\draw[gray!55](-0.547,3.111)--(-0.453,2.889);
\draw[gray!55](-0.536,3.614)--(-0.464,3.386);
\draw[gray!55](-0.528,4.117)--(-0.472,3.883);
\draw[gray!55](-0.029,-3.884)--(0.029,-4.116);
\draw[gray!55](-0.037,-3.386)--(0.037,-3.614);
\draw[gray!55](-0.049,-2.890)--(0.049,-3.110);
\draw[gray!55](-0.065,-2.399)--(0.065,-2.601);
\draw[gray!55](-0.085,-1.915)--(0.085,-2.085);
\draw[gray!55](-0.105,-1.441)--(0.105,-1.559);
\draw[gray!55](-0.116,-0.971)--(0.116,-1.029);
\draw[gray!55](-0.120,-0.493)--(0.120,-0.507);
\draw[gray!55](-0.120,0.000)--(0.120,0.000);
\draw[gray!55](-0.120,0.507)--(0.120,0.493);
\draw[gray!55](-0.116,1.029)--(0.116,0.971);
\draw[gray!55](-0.105,1.559)--(0.105,1.441);
\draw[gray!55](-0.085,2.085)--(0.085,1.915);
\draw[gray!55](-0.065,2.601)--(0.065,2.399);
\draw[gray!55](-0.049,3.110)--(0.049,2.890);
\draw[gray!55](-0.037,3.614)--(0.037,3.386);
\draw[gray!55](-0.029,4.116)--(0.029,3.884);
\draw[gray!55](0.470,-3.884)--(0.530,-4.116);
\draw[gray!55](0.461,-3.386)--(0.539,-3.614);
\draw[gray!55](0.449,-2.891)--(0.551,-3.109);
\draw[gray!55](0.431,-2.401)--(0.569,-2.599);
\draw[gray!55](0.410,-1.921)--(0.590,-2.079);
\draw[gray!55](0.390,-1.452)--(0.610,-1.548);
\draw[gray!55](0.381,-0.985)--(0.619,-1.015);
\draw[gray!55](0.380,-0.507)--(0.620,-0.493);
\draw[gray!55](0.381,-0.015)--(0.619,0.015);
\draw[gray!55](0.380,0.493)--(0.620,0.507);
\draw[gray!55](0.381,1.015)--(0.619,0.985);
\draw[gray!55](0.390,1.548)--(0.610,1.452);
\draw[gray!55](0.410,2.079)--(0.590,1.921);
\draw[gray!55](0.431,2.599)--(0.569,2.401);
\draw[gray!55](0.449,3.109)--(0.551,2.891);
\draw[gray!55](0.461,3.614)--(0.539,3.386);
\draw[gray!55](0.470,4.116)--(0.530,3.884);
\draw[gray!55](0.969,-3.884)--(1.031,-4.116);
\draw[gray!55](0.960,-3.387)--(1.040,-3.613);
\draw[gray!55](0.946,-2.893)--(1.054,-3.107);
\draw[gray!55](0.927,-2.405)--(1.073,-2.595);
\draw[gray!55](0.904,-1.928)--(1.096,-2.072);
\draw[gray!55](0.885,-1.464)--(1.115,-1.536);
\draw[gray!55](0.880,-1.000)--(1.120,-1.000);
\draw[gray!55](0.882,-0.522)--(1.118,-0.478);
\draw[gray!55](0.884,-0.029)--(1.116,0.029);
\draw[gray!55](0.882,0.478)--(1.118,0.522);
\draw[gray!55](0.880,1.000)--(1.120,1.000);
\draw[gray!55](0.885,1.536)--(1.115,1.464);
\draw[gray!55](0.904,2.072)--(1.096,1.928);
\draw[gray!55](0.927,2.595)--(1.073,2.405);
\draw[gray!55](0.946,3.107)--(1.054,2.893);
\draw[gray!55](0.960,3.613)--(1.040,3.387);
\draw[gray!55](0.969,4.116)--(1.031,3.884);
\draw[gray!55](1.468,-3.884)--(1.532,-4.116);
\draw[gray!55](1.458,-3.388)--(1.542,-3.612);
\draw[gray!55](1.444,-2.894)--(1.556,-3.106);
\draw[gray!55](1.423,-2.408)--(1.577,-2.592);
\draw[gray!55](1.398,-1.936)--(1.602,-2.064);
\draw[gray!55](1.382,-1.478)--(1.618,-1.522);
\draw[gray!55](1.381,-1.015)--(1.619,-0.985);
\draw[gray!55](1.385,-0.536)--(1.615,-0.464);
\draw[gray!55](1.388,-0.042)--(1.612,0.042);
\draw[gray!55](1.385,0.464)--(1.615,0.536);
\draw[gray!55](1.381,0.985)--(1.619,1.015);
\draw[gray!55](1.382,1.522)--(1.618,1.478);
\draw[gray!55](1.398,2.064)--(1.602,1.936);
\draw[gray!55](1.423,2.592)--(1.577,2.408);
\draw[gray!55](1.444,3.106)--(1.556,2.894);
\draw[gray!55](1.458,3.612)--(1.542,3.388);
\draw[gray!55](1.468,4.116)--(1.532,3.884);
\draw[gray!55](1.967,-3.885)--(2.033,-4.115);
\draw[gray!55](1.956,-3.388)--(2.044,-3.612);
\draw[gray!55](1.940,-2.896)--(2.060,-3.104);
\draw[gray!55](1.918,-2.413)--(2.082,-2.587);
\draw[gray!55](1.893,-1.946)--(2.107,-2.054);
\draw[gray!55](1.880,-1.493)--(2.120,-1.507);
\draw[gray!55](1.884,-1.029)--(2.116,-0.971);
\draw[gray!55](1.890,-0.548)--(2.110,-0.452);
\draw[gray!55](1.893,-0.054)--(2.107,0.054);
\draw[gray!55](1.890,0.452)--(2.110,0.548);
\draw[gray!55](1.884,0.971)--(2.116,1.029);
\draw[gray!55](1.880,1.507)--(2.120,1.493);
\draw[gray!55](1.893,2.054)--(2.107,1.946);
\draw[gray!55](1.918,2.587)--(2.082,2.413);
\draw[gray!55](1.940,3.104)--(2.060,2.896);
\draw[gray!55](1.956,3.612)--(2.044,3.388);
\draw[gray!55](1.967,4.115)--(2.033,3.885);
\draw[gray!55](2.466,-3.885)--(2.534,-4.115);
\draw[gray!55](2.454,-3.389)--(2.546,-3.611);
\draw[gray!55](2.437,-2.898)--(2.563,-3.102);
\draw[gray!55](2.412,-2.418)--(2.588,-2.582);
\draw[gray!55](2.388,-1.958)--(2.612,-2.042);
\draw[gray!55](2.380,-1.507)--(2.620,-1.493);
\draw[gray!55](2.388,-1.042)--(2.612,-0.958);
\draw[gray!55](2.395,-0.559)--(2.605,-0.441);
\draw[gray!55](2.398,-0.064)--(2.602,0.064);
\draw[gray!55](2.395,0.441)--(2.605,0.559);
\draw[gray!55](2.388,0.958)--(2.612,1.042);
\draw[gray!55](2.380,1.493)--(2.620,1.507);
\draw[gray!55](2.388,2.042)--(2.612,1.958);
\draw[gray!55](2.412,2.582)--(2.588,2.418);
\draw[gray!55](2.437,3.102)--(2.563,2.898);
\draw[gray!55](2.454,3.611)--(2.546,3.389);
\draw[gray!55](2.466,4.115)--(2.534,3.885);
\draw[gray!55](2.965,-3.885)--(3.035,-4.115);
\draw[gray!55](2.952,-3.390)--(3.048,-3.610);
\draw[gray!55](2.933,-2.900)--(3.067,-3.100);
\draw[gray!55](2.907,-2.424)--(3.093,-2.576);
\draw[gray!55](2.884,-1.971)--(3.116,-2.029);
\draw[gray!55](2.882,-1.522)--(3.118,-1.478);
\draw[gray!55](2.893,-1.054)--(3.107,-0.946);
\draw[gray!55](2.901,-0.568)--(3.099,-0.432);
\draw[gray!55](2.904,-0.072)--(3.096,0.072);
\draw[gray!55](2.901,0.432)--(3.099,0.568);
\draw[gray!55](2.893,0.946)--(3.107,1.054);
\draw[gray!55](2.882,1.478)--(3.118,1.522);
\draw[gray!55](2.884,2.029)--(3.116,1.971);
\draw[gray!55](2.907,2.576)--(3.093,2.424);
\draw[gray!55](2.933,3.100)--(3.067,2.900);
\draw[gray!55](2.952,3.610)--(3.048,3.390);
\draw[gray!55](2.965,4.115)--(3.035,3.885);
\draw[gray!55](3.463,-3.886)--(3.537,-4.114);
\draw[gray!55](3.450,-3.391)--(3.550,-3.609);
\draw[gray!55](3.429,-2.903)--(3.571,-3.097);
\draw[gray!55](3.401,-2.432)--(3.599,-2.568);
\draw[gray!55](3.381,-1.985)--(3.619,-2.015);
\draw[gray!55](3.385,-1.536)--(3.615,-1.464);
\draw[gray!55](3.398,-1.064)--(3.602,-0.936);
\draw[gray!55](3.407,-0.576)--(3.593,-0.424);
\draw[gray!55](3.410,-0.079)--(3.590,0.079);
\draw[gray!55](3.407,0.424)--(3.593,0.576);
\draw[gray!55](3.398,0.936)--(3.602,1.064);
\draw[gray!55](3.385,1.464)--(3.615,1.536);
\draw[gray!55](3.381,2.015)--(3.619,1.985);
\draw[gray!55](3.401,2.568)--(3.599,2.432);
\draw[gray!55](3.429,3.097)--(3.571,2.903);
\draw[gray!55](3.450,3.609)--(3.550,3.391);
\draw[gray!55](3.463,4.114)--(3.537,3.886);
\draw[gray!55](3.962,-3.886)--(4.038,-4.114);
\draw[gray!55](3.948,-3.392)--(4.052,-3.608);
\draw[gray!55](3.925,-2.906)--(4.075,-3.094);
\draw[gray!55](3.895,-2.441)--(4.105,-2.559);
\draw[gray!55](3.880,-2.000)--(4.120,-2.000);
\draw[gray!55](3.890,-1.548)--(4.110,-1.452);
\draw[gray!55](3.904,-1.072)--(4.096,-0.928);
\draw[gray!55](3.912,-0.582)--(4.088,-0.418);
\draw[gray!55](3.915,-0.085)--(4.085,0.085);
\draw[gray!55](3.912,0.418)--(4.088,0.582);
\draw[gray!55](3.904,0.928)--(4.096,1.072);
\draw[gray!55](3.890,1.452)--(4.110,1.548);
\draw[gray!55](3.880,2.000)--(4.120,2.000);
\draw[gray!55](3.895,2.559)--(4.105,2.441);
\draw[gray!55](3.925,3.094)--(4.075,2.906);
\draw[gray!55](3.948,3.608)--(4.052,3.392);
\draw[gray!55](3.962,4.114)--(4.038,3.886);
\draw[gray!55](4.461,-3.887)--(4.539,-4.113);
\draw[gray!55](4.445,-3.393)--(4.555,-3.607);
\draw[gray!55](4.420,-2.910)--(4.580,-3.090);
\draw[gray!55](4.390,-2.452)--(4.610,-2.548);
\draw[gray!55](4.381,-2.015)--(4.619,-1.985);
\draw[gray!55](4.395,-1.559)--(4.605,-1.441);
\draw[gray!55](4.410,-1.079)--(4.590,-0.921);
\draw[gray!55](4.418,-0.587)--(4.582,-0.413);
\draw[gray!55](4.420,-0.090)--(4.580,0.090);
\draw[gray!55](4.418,0.413)--(4.582,0.587);
\draw[gray!55](4.410,0.921)--(4.590,1.079);
\draw[gray!55](4.395,1.441)--(4.605,1.559);
\draw[gray!55](4.381,1.985)--(4.619,2.015);
\draw[gray!55](4.390,2.548)--(4.610,2.452);
\draw[gray!55](4.420,3.090)--(4.580,2.910);
\draw[gray!55](4.445,3.607)--(4.555,3.393);
\draw[gray!55](4.461,4.113)--(4.539,3.887);
\draw[gray!55](4.959,-3.887)--(5.041,-4.113);
\draw[gray!55](4.942,-3.395)--(5.058,-3.605);
\draw[gray!55](4.915,-2.915)--(5.085,-3.085);
\draw[gray!55](4.885,-2.464)--(5.115,-2.536);
\draw[gray!55](4.884,-2.029)--(5.116,-1.971);
\draw[gray!55](4.901,-1.568)--(5.099,-1.432);
\draw[gray!55](4.915,-1.085)--(5.085,-0.915);
\draw[gray!55](4.923,-0.592)--(5.077,-0.408);
\draw[gray!55](4.925,-0.094)--(5.075,0.094);
\draw[gray!55](4.923,0.408)--(5.077,0.592);
\draw[gray!55](4.915,0.915)--(5.085,1.085);
\draw[gray!55](4.901,1.432)--(5.099,1.568);
\draw[gray!55](4.884,1.971)--(5.116,2.029);
\draw[gray!55](4.885,2.536)--(5.115,2.464);
\draw[gray!55](4.915,3.085)--(5.085,2.915);
\draw[gray!55](4.942,3.605)--(5.058,3.395);
\draw[gray!55](4.959,4.113)--(5.041,3.887);
\draw[gray!55](5.457,-3.888)--(5.543,-4.112);
\draw[gray!55](5.439,-3.397)--(5.561,-3.603);
\draw[gray!55](5.410,-2.921)--(5.590,-3.079);
\draw[gray!55](5.382,-2.478)--(5.618,-2.522);
\draw[gray!55](5.388,-2.042)--(5.612,-1.958);
\draw[gray!55](5.407,-1.576)--(5.593,-1.424);
\draw[gray!55](5.420,-1.090)--(5.580,-0.910);
\draw[gray!55](5.427,-0.595)--(5.573,-0.405);
\draw[gray!55](5.429,-0.097)--(5.571,0.097);
\draw[gray!55](5.427,0.405)--(5.573,0.595);
\draw[gray!55](5.420,0.910)--(5.580,1.090);
\draw[gray!55](5.407,1.424)--(5.593,1.576);
\draw[gray!55](5.388,1.958)--(5.612,2.042);
\draw[gray!55](5.382,2.522)--(5.618,2.478);
\draw[gray!55](5.410,3.079)--(5.590,2.921);
\draw[gray!55](5.439,3.603)--(5.561,3.397);
\draw[gray!55](5.457,4.112)--(5.543,3.888);
\draw[gray!55](5.955,-3.889)--(6.045,-4.111);
\draw[gray!55](5.935,-3.399)--(6.065,-3.601);
\draw[gray!55](5.904,-2.928)--(6.096,-3.072);
\draw[gray!55](5.880,-2.493)--(6.120,-2.507);
\draw[gray!55](5.893,-2.054)--(6.107,-1.946);
\draw[gray!55](5.912,-1.582)--(6.088,-1.418);
\draw[gray!55](5.925,-1.094)--(6.075,-0.906);
\draw[gray!55](5.931,-0.599)--(6.069,-0.401);
\draw[gray!55](5.933,-0.100)--(6.067,0.100);
\draw[gray!55](5.931,0.401)--(6.069,0.599);
\draw[gray!55](5.925,0.906)--(6.075,1.094);
\draw[gray!55](5.912,1.418)--(6.088,1.582);
\draw[gray!55](5.893,1.946)--(6.107,2.054);
\draw[gray!55](5.880,2.507)--(6.120,2.493);
\draw[gray!55](5.904,3.072)--(6.096,2.928);
\draw[gray!55](5.935,3.601)--(6.065,3.399);
\draw[gray!55](5.955,4.111)--(6.045,3.889);
\draw[structurecolor] plot coordinates {(5.9548,-2.4403) (5.7118,-2.3899) (5.4738,-2.3396) (5.2409,-2.2893) (5.0131,-2.2390) (4.7903,-2.1887) (4.5726,-2.1384) (4.3600,-2.0881) (4.1524,-2.0377) (3.9498,-1.9874) (3.7524,-1.9371) (3.5600,-1.8868) (3.3727,-1.8365) (3.1904,-1.7862) (3.0132,-1.7358) (2.8410,-1.6855) (2.6739,-1.6352) (2.5119,-1.5849) (2.3550,-1.5346) (2.2031,-1.4843) (2.0562,-1.4340) (1.9145,-1.3836) (1.7778,-1.3333) (1.6461,-1.2830) (1.5196,-1.2327) (1.3980,-1.1824) (1.2816,-1.1321) (1.1702,-1.0818) (1.0639,-1.0314) (0.9626,-0.9811) (0.8664,-0.9308) (0.7753,-0.8805) (0.6892,-0.8302) (0.6082,-0.7799) (0.5323,-0.7296) (0.4614,-0.6792) (0.3956,-0.6289) (0.3348,-0.5786) (0.2791,-0.5283) (0.2285,-0.4780) (0.1829,-0.4277) (0.1424,-0.3774) (0.1070,-0.3270) (0.0766,-0.2767) (0.0513,-0.2264) (0.0310,-0.1761) (0.0158,-0.1258) (0.0057,-0.0755) (0.0006,-0.0252) (0.0006,0.0252) (0.0057,0.0755) (0.0158,0.1258) (0.0310,0.1761) (0.0513,0.2264) (0.0766,0.2767) (0.1070,0.3270) (0.1424,0.3774) (0.1829,0.4277) (0.2285,0.4780) (0.2791,0.5283) (0.3348,0.5786) (0.3956,0.6289) (0.4614,0.6792) (0.5323,0.7296) (0.6082,0.7799) (0.6892,0.8302) (0.7753,0.8805) (0.8664,0.9308) (0.9626,0.9811) (1.0639,1.0314) (1.1702,1.0818) (1.2816,1.1321) (1.3980,1.1824) (1.5196,1.2327) (1.6461,1.2830) (1.7778,1.3333) (1.9145,1.3836) (2.0562,1.4340) (2.2031,1.4843) (2.3550,1.5346) (2.5119,1.5849) (2.6739,1.6352) (2.8410,1.6855) (3.0132,1.7358) (3.1904,1.7862) (3.3727,1.8365) (3.5600,1.8868) (3.7524,1.9371) (3.9498,1.9874) (4.1524,2.0377) (4.3600,2.0881) (4.5726,2.1384) (4.7903,2.1887) (5.0131,2.2390) (5.2409,2.2893) (5.4738,2.3396) (5.7118,2.3899) (5.9548,2.4403)};
\draw[main,thick] plot coordinates {(0.0000,0.0000) (-0.0042,0.0000) (-0.0084,0.0000) (-0.0126,0.0000) (-0.0167,0.0000) (-0.0209,0.0001) (-0.0251,0.0001) (-0.0293,0.0001) (-0.0335,0.0001) (-0.0377,0.0002) (-0.0418,0.0002) (-0.0460,0.0003) (-0.0502,0.0003) (-0.0544,0.0004) (-0.0586,0.0004) (-0.0628,0.0005) (-0.0669,0.0006) (-0.0711,0.0006) (-0.0753,0.0007) (-0.0795,0.0008) (-0.0837,0.0009) (-0.0879,0.0010) (-0.0921,0.0011) (-0.0962,0.0012) (-0.1004,0.0013) (-0.1046,0.0014) (-0.1088,0.0015) (-0.1130,0.0016) (-0.1172,0.0017) (-0.1213,0.0018) (-0.1255,0.0020) (-0.1297,0.0021) (-0.1339,0.0022) (-0.1381,0.0024) (-0.1423,0.0025) (-0.1464,0.0027) (-0.1506,0.0028) (-0.1548,0.0030) (-0.1590,0.0032) (-0.1632,0.0033) (-0.1674,0.0035) (-0.1715,0.0037) (-0.1757,0.0039) (-0.1799,0.0040) (-0.1841,0.0042) (-0.1883,0.0044) (-0.1925,0.0046) (-0.1967,0.0048) (-0.2008,0.0050) (-0.2050,0.0053) (-0.2092,0.0055) (-0.2134,0.0057) (-0.2176,0.0059) (-0.2218,0.0061) (-0.2259,0.0064) (-0.2301,0.0066) (-0.2343,0.0069) (-0.2385,0.0071) (-0.2427,0.0074) (-0.2469,0.0076) (-0.2510,0.0079) (-0.2552,0.0081) (-0.2594,0.0084) (-0.2636,0.0087) (-0.2678,0.0090) (-0.2720,0.0092) (-0.2762,0.0095) (-0.2803,0.0098) (-0.2845,0.0101) (-0.2887,0.0104) (-0.2929,0.0107) (-0.2971,0.0110) (-0.3013,0.0113) (-0.3054,0.0117) (-0.3096,0.0120) (-0.3138,0.0123) (-0.3180,0.0126) (-0.3222,0.0130) (-0.3264,0.0133) (-0.3305,0.0137) (-0.3347,0.0140) (-0.3389,0.0144) (-0.3431,0.0147) (-0.3473,0.0151) (-0.3515,0.0154) (-0.3556,0.0158) (-0.3598,0.0162) (-0.3640,0.0166) (-0.3682,0.0170) (-0.3724,0.0173) (-0.3766,0.0177) (-0.3808,0.0181) (-0.3849,0.0185) (-0.3891,0.0189) (-0.3933,0.0193) (-0.3975,0.0198) (-0.4017,0.0202) (-0.4059,0.0206) (-0.4100,0.0210) (-0.4142,0.0215) (-0.4184,0.0219) (-0.4226,0.0223) (-0.4268,0.0228) (-0.4310,0.0232) (-0.4351,0.0237) (-0.4393,0.0241) (-0.4435,0.0246) (-0.4477,0.0251) (-0.4519,0.0255) (-0.4561,0.0260) (-0.4603,0.0265) (-0.4644,0.0270) (-0.4686,0.0275) (-0.4728,0.0280) (-0.4770,0.0285) (-0.4812,0.0290) (-0.4854,0.0295) (-0.4895,0.0300) (-0.4937,0.0305) (-0.4979,0.0310) (-0.5021,0.0315) (-0.5063,0.0321) (-0.5105,0.0326) (-0.5146,0.0331) (-0.5188,0.0337) (-0.5230,0.0342) (-0.5272,0.0348) (-0.5314,0.0353) (-0.5356,0.0359) (-0.5397,0.0365) (-0.5439,0.0370) (-0.5481,0.0376) (-0.5523,0.0382) (-0.5565,0.0388) (-0.5607,0.0393) (-0.5649,0.0399) (-0.5690,0.0405) (-0.5732,0.0411) (-0.5774,0.0417) (-0.5816,0.0423) (-0.5858,0.0429) (-0.5900,0.0436) (-0.5941,0.0442) (-0.5983,0.0448) (-0.6025,0.0454) (-0.6067,0.0461) (-0.6109,0.0467) (-0.6151,0.0474) (-0.6192,0.0480) (-0.6234,0.0487) (-0.6276,0.0493) (-0.6318,0.0500) (-0.6360,0.0506) (-0.6402,0.0513) (-0.6444,0.0520) (-0.6485,0.0527) (-0.6527,0.0533) (-0.6569,0.0540) (-0.6611,0.0547) (-0.6653,0.0554) (-0.6695,0.0561) (-0.6736,0.0568) (-0.6778,0.0575) (-0.6820,0.0583) (-0.6862,0.0590) (-0.6904,0.0597) (-0.6946,0.0604) (-0.6987,0.0612) (-0.7029,0.0619) (-0.7071,0.0626) (-0.7113,0.0634) (-0.7155,0.0641) (-0.7197,0.0649) (-0.7238,0.0657) (-0.7280,0.0664) (-0.7322,0.0672) (-0.7364,0.0680) (-0.7406,0.0687) (-0.7448,0.0695) (-0.7490,0.0703) (-0.7531,0.0711) (-0.7573,0.0719) (-0.7615,0.0727) (-0.7657,0.0735) (-0.7699,0.0743) (-0.7741,0.0751) (-0.7782,0.0759) (-0.7824,0.0768) (-0.7866,0.0776) (-0.7908,0.0784) (-0.7950,0.0792) (-0.7992,0.0801) (-0.8033,0.0809) (-0.8075,0.0818) (-0.8117,0.0826) (-0.8159,0.0835) (-0.8201,0.0844) (-0.8243,0.0852) (-0.8285,0.0861) (-0.8326,0.0870) (-0.8368,0.0879) (-0.8410,0.0887) (-0.8452,0.0896) (-0.8494,0.0905) (-0.8536,0.0914) (-0.8577,0.0923) (-0.8619,0.0932) (-0.8661,0.0942) (-0.8703,0.0951) (-0.8745,0.0960) (-0.8787,0.0969) (-0.8828,0.0978) (-0.8870,0.0988) (-0.8912,0.0997) (-0.8954,0.1007) (-0.8996,0.1016) (-0.9038,0.1026) (-0.9079,0.1035) (-0.9121,0.1045) (-0.9163,0.1055) (-0.9205,0.1064) (-0.9247,0.1074) (-0.9289,0.1084) (-0.9331,0.1094) (-0.9372,0.1104) (-0.9414,0.1114) (-0.9456,0.1124) (-0.9498,0.1134) (-0.9540,0.1144) (-0.9582,0.1154) (-0.9623,0.1164) (-0.9665,0.1174) (-0.9707,0.1185) (-0.9749,0.1195) (-0.9791,0.1205) (-0.9833,0.1216) (-0.9874,0.1226) (-0.9916,0.1237) (-0.9958,0.1247) (-1.0000,0.1258)};
\draw[main,thick] plot coordinates {(0.0000,0.0000) (0.0251,0.0001) (0.0502,0.0003) (0.0753,0.0007) (0.1004,0.0013) (0.1255,0.0020) (0.1506,0.0028) (0.1757,0.0039) (0.2008,0.0050) (0.2259,0.0064) (0.2510,0.0079) (0.2762,0.0095) (0.3013,0.0113) (0.3264,0.0133) (0.3515,0.0154) (0.3766,0.0177) (0.4017,0.0202) (0.4268,0.0228) (0.4519,0.0255) (0.4770,0.0284) (0.5021,0.0315) (0.5272,0.0347) (0.5523,0.0381) (0.5774,0.0416) (0.6025,0.0453) (0.6276,0.0492) (0.6527,0.0532) (0.6778,0.0573) (0.7029,0.0616) (0.7280,0.0661) (0.7531,0.0707) (0.7782,0.0755) (0.8033,0.0804) (0.8285,0.0855) (0.8536,0.0907) (0.8787,0.0961) (0.9038,0.1016) (0.9289,0.1073) (0.9540,0.1131) (0.9791,0.1191) (1.0042,0.1253) (1.0293,0.1315) (1.0544,0.1380) (1.0795,0.1445) (1.1046,0.1512) (1.1297,0.1581) (1.1548,0.1651) (1.1799,0.1723) (1.2050,0.1796) (1.2301,0.1870) (1.2552,0.1946) (1.2803,0.2023) (1.3054,0.2101) (1.3305,0.2181) (1.3556,0.2262) (1.3808,0.2345) (1.4059,0.2429) (1.4310,0.2514) (1.4561,0.2600) (1.4812,0.2688) (1.5063,0.2777) (1.5314,0.2867) (1.5565,0.2959) (1.5816,0.3052) (1.6067,0.3146) (1.6318,0.3241) (1.6569,0.3337) (1.6820,0.3435) (1.7071,0.3534) (1.7322,0.3634) (1.7573,0.3735) (1.7824,0.3837) (1.8075,0.3940) (1.8326,0.4044) (1.8577,0.4149) (1.8828,0.4256) (1.9079,0.4363) (1.9331,0.4471) (1.9582,0.4581) (1.9833,0.4691) (2.0084,0.4802) (2.0335,0.4914) (2.0586,0.5027) (2.0837,0.5141) (2.1088,0.5255) (2.1339,0.5371) (2.1590,0.5487) (2.1841,0.5604) (2.2092,0.5722) (2.2343,0.5840) (2.2594,0.5959) (2.2845,0.6079) (2.3096,0.6200) (2.3347,0.6321) (2.3598,0.6442) (2.3849,0.6565) (2.4100,0.6688) (2.4351,0.6811) (2.4603,0.6935) (2.4854,0.7060) (2.5105,0.7185) (2.5356,0.7310) (2.5607,0.7436) (2.5858,0.7562) (2.6109,0.7688) (2.6360,0.7815) (2.6611,0.7943) (2.6862,0.8070) (2.7113,0.8198) (2.7364,0.8326) (2.7615,0.8455) (2.7866,0.8583) (2.8117,0.8712) (2.8368,0.8841) (2.8619,0.8970) (2.8870,0.9099) (2.9121,0.9228) (2.9372,0.9358) (2.9623,0.9487) (2.9874,0.9617) (3.0126,0.9746) (3.0377,0.9875) (3.0628,1.0005) (3.0879,1.0134) (3.1130,1.0264) (3.1381,1.0393) (3.1632,1.0522) (3.1883,1.0651) (3.2134,1.0780) (3.2385,1.0908) (3.2636,1.1037) (3.2887,1.1165) (3.3138,1.1293) (3.3389,1.1421) (3.3640,1.1548) (3.3891,1.1676) (3.4142,1.1803) (3.4393,1.1929) (3.4644,1.2056) (3.4895,1.2182) (3.5146,1.2308) (3.5397,1.2433) (3.5649,1.2558) (3.5900,1.2682) (3.6151,1.2807) (3.6402,1.2930) (3.6653,1.3054) (3.6904,1.3177) (3.7155,1.3299) (3.7406,1.3421) (3.7657,1.3542) (3.7908,1.3663) (3.8159,1.3784) (3.8410,1.3904) (3.8661,1.4023) (3.8912,1.4142) (3.9163,1.4261) (3.9414,1.4379) (3.9665,1.4496) (3.9916,1.4613) (4.0167,1.4729) (4.0418,1.4845) (4.0669,1.4960) (4.0921,1.5074) (4.1172,1.5188) (4.1423,1.5301) (4.1674,1.5414) (4.1925,1.5526) (4.2176,1.5638) (4.2427,1.5749) (4.2678,1.5859) (4.2929,1.5969) (4.3180,1.6078) (4.3431,1.6186) (4.3682,1.6294) (4.3933,1.6401) (4.4184,1.6508) (4.4435,1.6614) (4.4686,1.6719) (4.4937,1.6824) (4.5188,1.6928) (4.5439,1.7031) (4.5690,1.7134) (4.5941,1.7236) (4.6192,1.7338) (4.6444,1.7439) (4.6695,1.7539) (4.6946,1.7639) (4.7197,1.7738) (4.7448,1.7836) (4.7699,1.7934) (4.7950,1.8031) (4.8201,1.8128) (4.8452,1.8224) (4.8703,1.8319) (4.8954,1.8414) (4.9205,1.8508) (4.9456,1.8602) (4.9707,1.8695) (4.9958,1.8787) (5.0209,1.8879) (5.0460,1.8970) (5.0711,1.9060) (5.0962,1.9150) (5.1213,1.9240) (5.1464,1.9329) (5.1715,1.9417) (5.1967,1.9504) (5.2218,1.9592) (5.2469,1.9678) (5.2720,1.9764) (5.2971,1.9850) (5.3222,1.9934) (5.3473,2.0019) (5.3724,2.0103) (5.3975,2.0186) (5.4226,2.0269) (5.4477,2.0351) (5.4728,2.0433) (5.4979,2.0514) (5.5230,2.0594) (5.5481,2.0675) (5.5732,2.0754) (5.5983,2.0834) (5.6234,2.0912) (5.6485,2.0990) (5.6736,2.1068) (5.6987,2.1146) (5.7238,2.1222) (5.7490,2.1299) (5.7741,2.1375) (5.7992,2.1450) (5.8243,2.1525) (5.8494,2.1599) (5.8745,2.1674) (5.8996,2.1747) (5.9247,2.1820) (5.9498,2.1893) (5.9749,2.1966) (6.0000,2.2038)};
\fill[main](0,0)circle(2pt);
\draw[main,thick] plot coordinates {(0.0000,-2.0000) (-0.0042,-1.9958) (-0.0084,-1.9917) (-0.0126,-1.9875) (-0.0167,-1.9834) (-0.0209,-1.9792) (-0.0251,-1.9751) (-0.0293,-1.9710) (-0.0335,-1.9669) (-0.0377,-1.9629) (-0.0418,-1.9588) (-0.0460,-1.9547) (-0.0502,-1.9507) (-0.0544,-1.9467) (-0.0586,-1.9427) (-0.0628,-1.9387) (-0.0669,-1.9347) (-0.0711,-1.9307) (-0.0753,-1.9267) (-0.0795,-1.9228) (-0.0837,-1.9188) (-0.0879,-1.9149) (-0.0921,-1.9110) (-0.0962,-1.9071) (-0.1004,-1.9032) (-0.1046,-1.8993) (-0.1088,-1.8954) (-0.1130,-1.8915) (-0.1172,-1.8877) (-0.1213,-1.8838) (-0.1255,-1.8800) (-0.1297,-1.8762) (-0.1339,-1.8724) (-0.1381,-1.8686) (-0.1423,-1.8648) (-0.1464,-1.8610) (-0.1506,-1.8572) (-0.1548,-1.8535) (-0.1590,-1.8497) (-0.1632,-1.8460) (-0.1674,-1.8422) (-0.1715,-1.8385) (-0.1757,-1.8348) (-0.1799,-1.8311) (-0.1841,-1.8274) (-0.1883,-1.8237) (-0.1925,-1.8201) (-0.1967,-1.8164) (-0.2008,-1.8127) (-0.2050,-1.8091) (-0.2092,-1.8055) (-0.2134,-1.8018) (-0.2176,-1.7982) (-0.2218,-1.7946) (-0.2259,-1.7910) (-0.2301,-1.7874) (-0.2343,-1.7839) (-0.2385,-1.7803) (-0.2427,-1.7767) (-0.2469,-1.7732) (-0.2510,-1.7696) (-0.2552,-1.7661) (-0.2594,-1.7626) (-0.2636,-1.7591) (-0.2678,-1.7556) (-0.2720,-1.7521) (-0.2762,-1.7486) (-0.2803,-1.7451) (-0.2845,-1.7416) (-0.2887,-1.7381) (-0.2929,-1.7347) (-0.2971,-1.7312) (-0.3013,-1.7278) (-0.3054,-1.7244) (-0.3096,-1.7209) (-0.3138,-1.7175) (-0.3180,-1.7141) (-0.3222,-1.7107) (-0.3264,-1.7073) (-0.3305,-1.7039) (-0.3347,-1.7005) (-0.3389,-1.6972) (-0.3431,-1.6938) (-0.3473,-1.6905) (-0.3515,-1.6871) (-0.3556,-1.6838) (-0.3598,-1.6804) (-0.3640,-1.6771) (-0.3682,-1.6738) (-0.3724,-1.6705) (-0.3766,-1.6672) (-0.3808,-1.6639) (-0.3849,-1.6606) (-0.3891,-1.6573) (-0.3933,-1.6540) (-0.3975,-1.6507) (-0.4017,-1.6475) (-0.4059,-1.6442) (-0.4100,-1.6410) (-0.4142,-1.6377) (-0.4184,-1.6345) (-0.4226,-1.6313) (-0.4268,-1.6281) (-0.4310,-1.6248) (-0.4351,-1.6216) (-0.4393,-1.6184) (-0.4435,-1.6152) (-0.4477,-1.6120) (-0.4519,-1.6089) (-0.4561,-1.6057) (-0.4603,-1.6025) (-0.4644,-1.5993) (-0.4686,-1.5962) (-0.4728,-1.5930) (-0.4770,-1.5899) (-0.4812,-1.5867) (-0.4854,-1.5836) (-0.4895,-1.5805) (-0.4937,-1.5774) (-0.4979,-1.5742) (-0.5021,-1.5711) (-0.5063,-1.5680) (-0.5105,-1.5649) (-0.5146,-1.5618) (-0.5188,-1.5588) (-0.5230,-1.5557) (-0.5272,-1.5526) (-0.5314,-1.5495) (-0.5356,-1.5465) (-0.5397,-1.5434) (-0.5439,-1.5403) (-0.5481,-1.5373) (-0.5523,-1.5343) (-0.5565,-1.5312) (-0.5607,-1.5282) (-0.5649,-1.5252) (-0.5690,-1.5221) (-0.5732,-1.5191) (-0.5774,-1.5161) (-0.5816,-1.5131) (-0.5858,-1.5101) (-0.5900,-1.5071) (-0.5941,-1.5041) (-0.5983,-1.5011) (-0.6025,-1.4982) (-0.6067,-1.4952) (-0.6109,-1.4922) (-0.6151,-1.4892) (-0.6192,-1.4863) (-0.6234,-1.4833) (-0.6276,-1.4804) (-0.6318,-1.4774) (-0.6360,-1.4745) (-0.6402,-1.4715) (-0.6444,-1.4686) (-0.6485,-1.4657) (-0.6527,-1.4628) (-0.6569,-1.4598) (-0.6611,-1.4569) (-0.6653,-1.4540) (-0.6695,-1.4511) (-0.6736,-1.4482) (-0.6778,-1.4453) (-0.6820,-1.4424) (-0.6862,-1.4395) (-0.6904,-1.4367) (-0.6946,-1.4338) (-0.6987,-1.4309) (-0.7029,-1.4280) (-0.7071,-1.4252) (-0.7113,-1.4223) (-0.7155,-1.4194) (-0.7197,-1.4166) (-0.7238,-1.4137) (-0.7280,-1.4109) (-0.7322,-1.4081) (-0.7364,-1.4052) (-0.7406,-1.4024) (-0.7448,-1.3996) (-0.7490,-1.3967) (-0.7531,-1.3939) (-0.7573,-1.3911) (-0.7615,-1.3883) (-0.7657,-1.3855) (-0.7699,-1.3827) (-0.7741,-1.3799) (-0.7782,-1.3771) (-0.7824,-1.3743) (-0.7866,-1.3715) (-0.7908,-1.3687) (-0.7950,-1.3659) (-0.7992,-1.3631) (-0.8033,-1.3603) (-0.8075,-1.3576) (-0.8117,-1.3548) (-0.8159,-1.3520) (-0.8201,-1.3493) (-0.8243,-1.3465) (-0.8285,-1.3437) (-0.8326,-1.3410) (-0.8368,-1.3382) (-0.8410,-1.3355) (-0.8452,-1.3327) (-0.8494,-1.3300) (-0.8536,-1.3273) (-0.8577,-1.3245) (-0.8619,-1.3218) (-0.8661,-1.3191) (-0.8703,-1.3164) (-0.8745,-1.3136) (-0.8787,-1.3109) (-0.8828,-1.3082) (-0.8870,-1.3055) (-0.8912,-1.3028) (-0.8954,-1.3001) (-0.8996,-1.2974) (-0.9038,-1.2947) (-0.9079,-1.2920) (-0.9121,-1.2893) (-0.9163,-1.2866) (-0.9205,-1.2839) (-0.9247,-1.2812) (-0.9289,-1.2785) (-0.9331,-1.2758) (-0.9372,-1.2732) (-0.9414,-1.2705) (-0.9456,-1.2678) (-0.9498,-1.2652) (-0.9540,-1.2625) (-0.9582,-1.2598) (-0.9623,-1.2572) (-0.9665,-1.2545) (-0.9707,-1.2518) (-0.9749,-1.2492) (-0.9791,-1.2465) (-0.9833,-1.2439) (-0.9874,-1.2412) (-0.9916,-1.2386) (-0.9958,-1.2360) (-1.0000,-1.2333)};
\draw[main,thick] plot coordinates {(0.0000,-2.0000) (0.0088,-2.0088) (0.0175,-2.0176) (0.0263,-2.0265) (0.0350,-2.0355) (0.0438,-2.0445) (0.0526,-2.0536) (0.0613,-2.0628) (0.0701,-2.0720) (0.0788,-2.0813) (0.0876,-2.0906) (0.0963,-2.1000) (0.1051,-2.1095) (0.1139,-2.1191) (0.1226,-2.1287) (0.1314,-2.1384) (0.1401,-2.1481) (0.1489,-2.1580) (0.1577,-2.1679) (0.1664,-2.1779) (0.1752,-2.1879) (0.1839,-2.1980) (0.1927,-2.2083) (0.2015,-2.2186) (0.2102,-2.2289) (0.2190,-2.2394) (0.2277,-2.2499) (0.2365,-2.2606) (0.2453,-2.2713) (0.2540,-2.2821) (0.2628,-2.2930) (0.2715,-2.3040) (0.2803,-2.3150) (0.2890,-2.3262) (0.2978,-2.3375) (0.3066,-2.3488) (0.3153,-2.3603) (0.3241,-2.3719) (0.3328,-2.3835) (0.3416,-2.3953) (0.3504,-2.4071) (0.3591,-2.4191) (0.3679,-2.4312) (0.3766,-2.4434) (0.3854,-2.4557) (0.3942,-2.4681) (0.4029,-2.4807) (0.4117,-2.4933) (0.4204,-2.5061) (0.4292,-2.5190) (0.4380,-2.5320) (0.4467,-2.5451) (0.4555,-2.5584) (0.4642,-2.5718) (0.4730,-2.5853) (0.4817,-2.5990) (0.4905,-2.6128) (0.4993,-2.6268) (0.5080,-2.6408) (0.5168,-2.6551) (0.5255,-2.6695) (0.5343,-2.6840) (0.5431,-2.6987) (0.5518,-2.7135) (0.5606,-2.7285) (0.5693,-2.7436) (0.5781,-2.7590) (0.5869,-2.7745) (0.5956,-2.7901) (0.6044,-2.8059) (0.6131,-2.8219) (0.6219,-2.8381) (0.6307,-2.8545) (0.6394,-2.8711) (0.6482,-2.8878) (0.6569,-2.9047) (0.6657,-2.9219) (0.6744,-2.9392) (0.6832,-2.9568) (0.6920,-2.9745) (0.7007,-2.9925) (0.7095,-3.0107) (0.7182,-3.0291) (0.7270,-3.0477) (0.7358,-3.0666) (0.7445,-3.0857) (0.7533,-3.1050) (0.7620,-3.1246) (0.7708,-3.1444) (0.7796,-3.1645) (0.7883,-3.1849) (0.7971,-3.2055) (0.8058,-3.2264) (0.8146,-3.2476) (0.8233,-3.2690) (0.8321,-3.2908) (0.8409,-3.3128) (0.8496,-3.3351) (0.8584,-3.3578) (0.8671,-3.3808) (0.8759,-3.4040) (0.8847,-3.4277) (0.8934,-3.4516) (0.9022,-3.4759) (0.9109,-3.5006) (0.9197,-3.5256) (0.9285,-3.5510) (0.9372,-3.5768) (0.9460,-3.6029) (0.9547,-3.6295) (0.9635,-3.6565) (0.9723,-3.6838) (0.9810,-3.7116) (0.9898,-3.7399) (0.9985,-3.7686) (1.0073,-3.7977) (1.0160,-3.8273) (1.0248,-3.8574) (1.0336,-3.8880) (1.0423,-3.9191) (1.0511,-3.9507) (1.0598,-3.9828)};
\fill[main](0,-2)circle(2pt);
\draw[structurecolor,dashed,thick]plot coordinates{(-1.0000,-0.7897) (-0.9560,-0.8074) (-0.9119,-0.8250) (-0.8679,-0.8425) (-0.8239,-0.8598) (-0.7799,-0.8769) (-0.7358,-0.8939) (-0.6918,-0.9107) (-0.6478,-0.9274) (-0.6038,-0.9439) (-0.5597,-0.9603) (-0.5157,-0.9765) (-0.4717,-0.9926) (-0.4277,-1.0086) (-0.3836,-1.0244) (-0.3396,-1.0401) (-0.2956,-1.0557) (-0.2516,-1.0712) (-0.2075,-1.0865) (-0.1635,-1.1017) (-0.1195,-1.1168) (-0.0755,-1.1318) (-0.0314,-1.1467) (0.0126,-1.1614) (0.0566,-1.1761) (0.1006,-1.1906) (0.1447,-1.2051) (0.1887,-1.2194) (0.2327,-1.2337) (0.2767,-1.2478) (0.3208,-1.2618) (0.3648,-1.2758) (0.4088,-1.2896) (0.4528,-1.3034) (0.4969,-1.3171) (0.5409,-1.3307) (0.5849,-1.3441) (0.6289,-1.3576) (0.6730,-1.3709) (0.7170,-1.3841) (0.7610,-1.3973) (0.8050,-1.4103) (0.8491,-1.4233) (0.8931,-1.4362) (0.9371,-1.4491) (0.9811,-1.4618) (1.0252,-1.4745) (1.0692,-1.4871) (1.1132,-1.4997) (1.1572,-1.5121) (1.2013,-1.5245) (1.2453,-1.5368) (1.2893,-1.5491) (1.3333,-1.5613) (1.3774,-1.5734) (1.4214,-1.5855) (1.4654,-1.5975) (1.5094,-1.6094) (1.5535,-1.6212) (1.5975,-1.6330) (1.6415,-1.6448) (1.6855,-1.6565) (1.7296,-1.6681) (1.7736,-1.6796) (1.8176,-1.6911) (1.8616,-1.7026) (1.9057,-1.7140) (1.9497,-1.7253) (1.9937,-1.7366) (2.0377,-1.7478) (2.0818,-1.7590) (2.1258,-1.7701) (2.1698,-1.7811) (2.2138,-1.7922) (2.2579,-1.8031) (2.3019,-1.8140) (2.3459,-1.8249) (2.3899,-1.8357) (2.4340,-1.8464) (2.4780,-1.8572) (2.5220,-1.8678) (2.5660,-1.8784) (2.6101,-1.8890) (2.6541,-1.8995) (2.6981,-1.9100) (2.7421,-1.9204) (2.7862,-1.9308) (2.8302,-1.9412) (2.8742,-1.9515) (2.9182,-1.9617) (2.9623,-1.9720) (3.0063,-1.9821) (3.0503,-1.9923) (3.0943,-2.0024) (3.1384,-2.0124) (3.1824,-2.0224) (3.2264,-2.0324) (3.2704,-2.0423) (3.3145,-2.0522) (3.3585,-2.0621) (3.4025,-2.0719) (3.4465,-2.0817) (3.4906,-2.0914) (3.5346,-2.1011) (3.5786,-2.1108) (3.6226,-2.1204) (3.6667,-2.1300) (3.7107,-2.1396) (3.7547,-2.1491) (3.7987,-2.1586) (3.8428,-2.1680) (3.8868,-2.1775) (3.9308,-2.1869) (3.9748,-2.1962) (4.0189,-2.2055) (4.0629,-2.2148) (4.1069,-2.2241) (4.1509,-2.2333) (4.1950,-2.2425) (4.2390,-2.2517) (4.2830,-2.2608) (4.3270,-2.2699) (4.3711,-2.2790) (4.4151,-2.2880) (4.4591,-2.2970) (4.5031,-2.3060) (4.5472,-2.3149) (4.5912,-2.3239) (4.6352,-2.3327) (4.6792,-2.3416) (4.7233,-2.3504) (4.7673,-2.3592) (4.8113,-2.3680) (4.8553,-2.3768) (4.8994,-2.3855) (4.9434,-2.3942) (4.9874,-2.4029) (5.0314,-2.4115) (5.0755,-2.4201) (5.1195,-2.4287) (5.1635,-2.4373) (5.2075,-2.4458) (5.2516,-2.4543) (5.2956,-2.4628) (5.3396,-2.4712) (5.3836,-2.4797) (5.4277,-2.4881) (5.4717,-2.4965) (5.5157,-2.5048) (5.5597,-2.5132) (5.6038,-2.5215) (5.6478,-2.5298) (5.6918,-2.5380) (5.7358,-2.5463) (5.7799,-2.5545) (5.8239,-2.5627) (5.8679,-2.5709) (5.9119,-2.5790) (5.9560,-2.5871) (6.0000,-2.5952)};\node[fill=white,inner sep=1pt]at(4,-3.6){分界解};\draw(-1,-.08)--(-1,.08);\draw(1,-.08)--(1,.08);\draw(2,-.08)--(2,.08);\node[below,font=\scriptsize,fill=white,inner sep=.2pt]at(2,-.08){$ 2 $};\draw(3,-.08)--(3,.08);\draw(4,-.08)--(4,.08);\node[below,font=\scriptsize,fill=white,inner sep=.2pt]at(4,-.08){$ 4 $};\draw(5,-.08)--(5,.08);\draw(6,-.08)--(6,.08);\node[below,font=\scriptsize,fill=white,inner sep=.2pt]at(6,-.08){$ 6 $};\draw(-.08,-4)--(.08,-4);\node[left,font=\scriptsize,fill=white,inner sep=.2pt]at(-.08,-4){$ -4 $};\draw(-.08,-3)--(.08,-3);\draw(-.08,-2)--(.08,-2);\node[left,font=\scriptsize,fill=white,inner sep=.2pt]at(-.08,-2){$ -2 $};\draw(-.08,-1)--(.08,-1);\draw(-.08,1)--(.08,1);\draw(-.08,2)--(.08,2);\node[left,font=\scriptsize,fill=white,inner sep=.2pt]at(-.08,2){$ 2 $};\draw(-.08,3)--(.08,3);\draw(-.08,4)--(.08,4);\node[left,font=\scriptsize,fill=white,inner sep=.2pt]at(-.08,4){$ 4 $};\end{tikzpicture}\end{center}
}

\newcommand{\ODEExamFinalField}{%
\begin{center}\begin{tikzpicture}[x=.62cm,y=.53cm]
\clip (-1.6,-4.6) rectangle (5.7,5.8);
\draw[->](-1,0)--(5.35,0)node[right]{$x$};\draw[->](0,-4)--(0,5.35)node[above]{$y$};
\draw[gray!55](-1.024,-3.882)--(-0.976,-4.118);
\draw[gray!55](-1.029,-3.383)--(-0.971,-3.617);
\draw[gray!55](-1.035,-2.885)--(-0.965,-3.115);
\draw[gray!55](-1.044,-2.388)--(-0.956,-2.612);
\draw[gray!55](-1.054,-1.893)--(-0.946,-2.107);
\draw[gray!55](-1.065,-1.399)--(-0.935,-1.601);
\draw[gray!55](-1.075,-0.906)--(-0.925,-1.094);
\draw[gray!55](-1.082,-0.413)--(-0.918,-0.587);
\draw[gray!55](-1.085,0.085)--(-0.915,-0.085);
\draw[gray!55](-1.082,0.587)--(-0.918,0.413);
\draw[gray!55](-1.075,1.094)--(-0.925,0.906);
\draw[gray!55](-1.065,1.601)--(-0.935,1.399);
\draw[gray!55](-1.054,2.107)--(-0.946,1.893);
\draw[gray!55](-1.044,2.612)--(-0.956,2.388);
\draw[gray!55](-1.035,3.115)--(-0.965,2.885);
\draw[gray!55](-1.029,3.617)--(-0.971,3.383);
\draw[gray!55](-1.024,4.118)--(-0.976,3.882);
\draw[gray!55](-1.020,4.618)--(-0.980,4.382);
\draw[gray!55](-1.016,5.119)--(-0.984,4.881);
\draw[gray!55](-0.526,-3.883)--(-0.474,-4.117);
\draw[gray!55](-0.532,-3.384)--(-0.468,-3.616);
\draw[gray!55](-0.541,-2.887)--(-0.459,-3.113);
\draw[gray!55](-0.552,-2.392)--(-0.448,-2.608);
\draw[gray!55](-0.567,-1.900)--(-0.433,-2.100);
\draw[gray!55](-0.582,-1.413)--(-0.418,-1.587);
\draw[gray!55](-0.596,-0.928)--(-0.404,-1.072);
\draw[gray!55](-0.605,-0.441)--(-0.395,-0.559);
\draw[gray!55](-0.607,0.054)--(-0.393,-0.054);
\draw[gray!55](-0.605,0.559)--(-0.395,0.441);
\draw[gray!55](-0.596,1.072)--(-0.404,0.928);
\draw[gray!55](-0.582,1.587)--(-0.418,1.413);
\draw[gray!55](-0.567,2.100)--(-0.433,1.900);
\draw[gray!55](-0.552,2.608)--(-0.448,2.392);
\draw[gray!55](-0.541,3.113)--(-0.459,2.887);
\draw[gray!55](-0.532,3.616)--(-0.468,3.384);
\draw[gray!55](-0.526,4.117)--(-0.474,3.883);
\draw[gray!55](-0.521,4.618)--(-0.479,4.382);
\draw[gray!55](-0.518,5.119)--(-0.482,4.881);
\draw[gray!55](-0.029,-3.884)--(0.029,-4.116);
\draw[gray!55](-0.037,-3.386)--(0.037,-3.614);
\draw[gray!55](-0.049,-2.890)--(0.049,-3.110);
\draw[gray!55](-0.065,-2.399)--(0.065,-2.601);
\draw[gray!55](-0.085,-1.915)--(0.085,-2.085);
\draw[gray!55](-0.105,-1.441)--(0.105,-1.559);
\draw[gray!55](-0.116,-0.971)--(0.116,-1.029);
\draw[gray!55](-0.120,-0.493)--(0.120,-0.507);
\draw[gray!55](-0.120,0.000)--(0.120,0.000);
\draw[gray!55](-0.120,0.507)--(0.120,0.493);
\draw[gray!55](-0.116,1.029)--(0.116,0.971);
\draw[gray!55](-0.105,1.559)--(0.105,1.441);
\draw[gray!55](-0.085,2.085)--(0.085,1.915);
\draw[gray!55](-0.065,2.601)--(0.065,2.399);
\draw[gray!55](-0.049,3.110)--(0.049,2.890);
\draw[gray!55](-0.037,3.614)--(0.037,3.386);
\draw[gray!55](-0.029,4.116)--(0.029,3.884);
\draw[gray!55](-0.023,4.618)--(0.023,4.382);
\draw[gray!55](-0.019,5.118)--(0.019,4.882);
\draw[gray!55](0.467,-3.885)--(0.533,-4.115);
\draw[gray!55](0.456,-3.388)--(0.544,-3.612);
\draw[gray!55](0.440,-2.896)--(0.560,-3.104);
\draw[gray!55](0.418,-2.413)--(0.582,-2.587);
\draw[gray!55](0.393,-1.946)--(0.607,-2.054);
\draw[gray!55](0.380,-1.493)--(0.620,-1.507);
\draw[gray!55](0.384,-1.029)--(0.616,-0.971);
\draw[gray!55](0.390,-0.548)--(0.610,-0.452);
\draw[gray!55](0.393,-0.054)--(0.607,0.054);
\draw[gray!55](0.390,0.452)--(0.610,0.548);
\draw[gray!55](0.384,0.971)--(0.616,1.029);
\draw[gray!55](0.380,1.507)--(0.620,1.493);
\draw[gray!55](0.393,2.054)--(0.607,1.946);
\draw[gray!55](0.418,2.587)--(0.582,2.413);
\draw[gray!55](0.440,3.104)--(0.560,2.896);
\draw[gray!55](0.456,3.612)--(0.544,3.388);
\draw[gray!55](0.467,4.115)--(0.533,3.885);
\draw[gray!55](0.474,4.617)--(0.526,4.383);
\draw[gray!55](0.479,5.118)--(0.521,4.882);
\draw[gray!55](0.962,-3.886)--(1.038,-4.114);
\draw[gray!55](0.948,-3.392)--(1.052,-3.608);
\draw[gray!55](0.925,-2.906)--(1.075,-3.094);
\draw[gray!55](0.895,-2.441)--(1.105,-2.559);
\draw[gray!55](0.880,-2.000)--(1.120,-2.000);
\draw[gray!55](0.890,-1.548)--(1.110,-1.452);
\draw[gray!55](0.904,-1.072)--(1.096,-0.928);
\draw[gray!55](0.912,-0.582)--(1.088,-0.418);
\draw[gray!55](0.915,-0.085)--(1.085,0.085);
\draw[gray!55](0.912,0.418)--(1.088,0.582);
\draw[gray!55](0.904,0.928)--(1.096,1.072);
\draw[gray!55](0.890,1.452)--(1.110,1.548);
\draw[gray!55](0.880,2.000)--(1.120,2.000);
\draw[gray!55](0.895,2.559)--(1.105,2.441);
\draw[gray!55](0.925,3.094)--(1.075,2.906);
\draw[gray!55](0.948,3.608)--(1.052,3.392);
\draw[gray!55](0.962,4.114)--(1.038,3.886);
\draw[gray!55](0.971,4.617)--(1.029,4.383);
\draw[gray!55](0.978,5.118)--(1.022,4.882);
\draw[gray!55](1.455,-3.889)--(1.545,-4.111);
\draw[gray!55](1.435,-3.399)--(1.565,-3.601);
\draw[gray!55](1.404,-2.928)--(1.596,-3.072);
\draw[gray!55](1.380,-2.493)--(1.620,-2.507);
\draw[gray!55](1.393,-2.054)--(1.607,-1.946);
\draw[gray!55](1.412,-1.582)--(1.588,-1.418);
\draw[gray!55](1.425,-1.094)--(1.575,-0.906);
\draw[gray!55](1.431,-0.599)--(1.569,-0.401);
\draw[gray!55](1.433,-0.100)--(1.567,0.100);
\draw[gray!55](1.431,0.401)--(1.569,0.599);
\draw[gray!55](1.425,0.906)--(1.575,1.094);
\draw[gray!55](1.412,1.418)--(1.588,1.582);
\draw[gray!55](1.393,1.946)--(1.607,2.054);
\draw[gray!55](1.380,2.507)--(1.620,2.493);
\draw[gray!55](1.404,3.072)--(1.596,2.928);
\draw[gray!55](1.435,3.601)--(1.565,3.399);
\draw[gray!55](1.455,4.111)--(1.545,3.889);
\draw[gray!55](1.468,4.616)--(1.532,4.384);
\draw[gray!55](1.475,5.117)--(1.525,4.883);
\draw[gray!55](1.946,-3.893)--(2.054,-4.107);
\draw[gray!55](1.918,-3.413)--(2.082,-3.587);
\draw[gray!55](1.884,-2.971)--(2.116,-3.029);
\draw[gray!55](1.890,-2.548)--(2.110,-2.452);
\draw[gray!55](1.915,-2.085)--(2.085,-1.915);
\draw[gray!55](1.931,-1.599)--(2.069,-1.401);
\draw[gray!55](1.940,-1.104)--(2.060,-0.896);
\draw[gray!55](1.945,-0.607)--(2.055,-0.393);
\draw[gray!55](1.946,-0.107)--(2.054,0.107);
\draw[gray!55](1.945,0.393)--(2.055,0.607);
\draw[gray!55](1.940,0.896)--(2.060,1.104);
\draw[gray!55](1.931,1.401)--(2.069,1.599);
\draw[gray!55](1.915,1.915)--(2.085,2.085);
\draw[gray!55](1.890,2.452)--(2.110,2.548);
\draw[gray!55](1.884,3.029)--(2.116,2.971);
\draw[gray!55](1.918,3.587)--(2.082,3.413);
\draw[gray!55](1.946,4.107)--(2.054,3.893);
\draw[gray!55](1.963,4.614)--(2.037,4.386);
\draw[gray!55](1.973,5.117)--(2.027,4.883);
\draw[gray!55](2.433,-3.900)--(2.567,-4.100);
\draw[gray!55](2.395,-3.441)--(2.605,-3.559);
\draw[gray!55](2.384,-3.029)--(2.616,-2.971);
\draw[gray!55](2.412,-2.582)--(2.588,-2.418);
\draw[gray!55](2.433,-2.100)--(2.567,-1.900);
\draw[gray!55](2.445,-1.607)--(2.555,-1.393);
\draw[gray!55](2.451,-1.110)--(2.549,-0.890);
\draw[gray!55](2.454,-0.611)--(2.546,-0.389);
\draw[gray!55](2.455,-0.111)--(2.545,0.111);
\draw[gray!55](2.454,0.389)--(2.546,0.611);
\draw[gray!55](2.451,0.890)--(2.549,1.110);
\draw[gray!55](2.445,1.393)--(2.555,1.607);
\draw[gray!55](2.433,1.900)--(2.567,2.100);
\draw[gray!55](2.412,2.418)--(2.588,2.582);
\draw[gray!55](2.384,2.971)--(2.616,3.029);
\draw[gray!55](2.395,3.559)--(2.605,3.441);
\draw[gray!55](2.433,4.100)--(2.567,3.900);
\draw[gray!55](2.456,4.612)--(2.544,4.388);
\draw[gray!55](2.469,5.116)--(2.531,4.884);
\draw[gray!55](2.915,-3.915)--(3.085,-4.085);
\draw[gray!55](2.880,-3.493)--(3.120,-3.507);
\draw[gray!55](2.904,-3.072)--(3.096,-2.928);
\draw[gray!55](2.931,-2.599)--(3.069,-2.401);
\draw[gray!55](2.946,-2.107)--(3.054,-1.893);
\draw[gray!55](2.954,-1.611)--(3.046,-1.389);
\draw[gray!55](2.959,-1.113)--(3.041,-0.887);
\draw[gray!55](2.961,-0.614)--(3.039,-0.386);
\draw[gray!55](2.962,-0.114)--(3.038,0.114);
\draw[gray!55](2.961,0.386)--(3.039,0.614);
\draw[gray!55](2.959,0.887)--(3.041,1.113);
\draw[gray!55](2.954,1.389)--(3.046,1.611);
\draw[gray!55](2.946,1.893)--(3.054,2.107);
\draw[gray!55](2.931,2.401)--(3.069,2.599);
\draw[gray!55](2.904,2.928)--(3.096,3.072);
\draw[gray!55](2.880,3.507)--(3.120,3.493);
\draw[gray!55](2.915,4.085)--(3.085,3.915);
\draw[gray!55](2.948,4.608)--(3.052,4.392);
\draw[gray!55](2.965,5.115)--(3.035,4.885);
\draw[gray!55](3.393,-3.946)--(3.607,-4.054);
\draw[gray!55](3.390,-3.548)--(3.610,-3.452);
\draw[gray!55](3.425,-3.094)--(3.575,-2.906);
\draw[gray!55](3.445,-2.607)--(3.555,-2.393);
\draw[gray!55](3.455,-2.111)--(3.545,-1.889);
\draw[gray!55](3.461,-1.614)--(3.539,-1.386);
\draw[gray!55](3.465,-1.115)--(3.535,-0.885);
\draw[gray!55](3.466,-0.615)--(3.534,-0.385);
\draw[gray!55](3.467,-0.115)--(3.533,0.115);
\draw[gray!55](3.466,0.385)--(3.534,0.615);
\draw[gray!55](3.465,0.885)--(3.535,1.115);
\draw[gray!55](3.461,1.386)--(3.539,1.614);
\draw[gray!55](3.455,1.889)--(3.545,2.111);
\draw[gray!55](3.445,2.393)--(3.555,2.607);
\draw[gray!55](3.425,2.906)--(3.575,3.094);
\draw[gray!55](3.390,3.452)--(3.610,3.548);
\draw[gray!55](3.393,4.054)--(3.607,3.946);
\draw[gray!55](3.435,4.601)--(3.565,4.399);
\draw[gray!55](3.459,5.113)--(3.541,4.887);
\draw[gray!55](3.880,-4.000)--(4.120,-4.000);
\draw[gray!55](3.912,-3.582)--(4.088,-3.418);
\draw[gray!55](3.940,-3.104)--(4.060,-2.896);
\draw[gray!55](3.954,-2.611)--(4.046,-2.389);
\draw[gray!55](3.962,-2.114)--(4.038,-1.886);
\draw[gray!55](3.966,-1.615)--(4.034,-1.385);
\draw[gray!55](3.969,-1.116)--(4.031,-0.884);
\draw[gray!55](3.970,-0.616)--(4.030,-0.384);
\draw[gray!55](3.971,-0.116)--(4.029,0.116);
\draw[gray!55](3.970,0.384)--(4.030,0.616);
\draw[gray!55](3.969,0.884)--(4.031,1.116);
\draw[gray!55](3.966,1.385)--(4.034,1.615);
\draw[gray!55](3.962,1.886)--(4.038,2.114);
\draw[gray!55](3.954,2.389)--(4.046,2.611);
\draw[gray!55](3.940,2.896)--(4.060,3.104);
\draw[gray!55](3.912,3.418)--(4.088,3.582);
\draw[gray!55](3.880,4.000)--(4.120,4.000);
\draw[gray!55](3.918,4.587)--(4.082,4.413);
\draw[gray!55](3.951,5.110)--(4.049,4.890);
\draw[gray!55](4.393,-4.054)--(4.607,-3.946);
\draw[gray!55](4.431,-3.599)--(4.569,-3.401);
\draw[gray!55](4.451,-3.110)--(4.549,-2.890);
\draw[gray!55](4.461,-2.614)--(4.539,-2.386);
\draw[gray!55](4.467,-2.115)--(4.533,-1.885);
\draw[gray!55](4.470,-1.616)--(4.530,-1.384);
\draw[gray!55](4.473,-1.117)--(4.527,-0.883);
\draw[gray!55](4.474,-0.617)--(4.526,-0.383);
\draw[gray!55](4.474,-0.117)--(4.526,0.117);
\draw[gray!55](4.474,0.383)--(4.526,0.617);
\draw[gray!55](4.473,0.883)--(4.527,1.117);
\draw[gray!55](4.470,1.384)--(4.530,1.616);
\draw[gray!55](4.467,1.885)--(4.533,2.115);
\draw[gray!55](4.461,2.386)--(4.539,2.614);
\draw[gray!55](4.451,2.890)--(4.549,3.110);
\draw[gray!55](4.431,3.401)--(4.569,3.599);
\draw[gray!55](4.393,3.946)--(4.607,4.054);
\draw[gray!55](4.395,4.559)--(4.605,4.441);
\draw[gray!55](4.440,5.104)--(4.560,4.896);
\draw[gray!55](4.915,-4.085)--(5.085,-3.915);
\draw[gray!55](4.945,-3.607)--(5.055,-3.393);
\draw[gray!55](4.959,-3.113)--(5.041,-2.887);
\draw[gray!55](4.966,-2.615)--(5.034,-2.385);
\draw[gray!55](4.971,-2.116)--(5.029,-1.884);
\draw[gray!55](4.974,-1.617)--(5.026,-1.383);
\draw[gray!55](4.975,-1.117)--(5.025,-0.883);
\draw[gray!55](4.976,-0.618)--(5.024,-0.382);
\draw[gray!55](4.976,-0.118)--(5.024,0.118);
\draw[gray!55](4.976,0.382)--(5.024,0.618);
\draw[gray!55](4.975,0.883)--(5.025,1.117);
\draw[gray!55](4.974,1.383)--(5.026,1.617);
\draw[gray!55](4.971,1.884)--(5.029,2.116);
\draw[gray!55](4.966,2.385)--(5.034,2.615);
\draw[gray!55](4.959,2.887)--(5.041,3.113);
\draw[gray!55](4.945,3.393)--(5.055,3.607);
\draw[gray!55](4.915,3.915)--(5.085,4.085);
\draw[gray!55](4.880,4.507)--(5.120,4.493);
\draw[gray!55](4.925,5.094)--(5.075,4.906);
\draw[structurecolor,dashed] plot coordinates {(3.0000,-4.0000) (2.8876,-3.9434) (2.7768,-3.8868) (2.6676,-3.8302) (2.5600,-3.7736) (2.4540,-3.7170) (2.3496,-3.6604) (2.2468,-3.6038) (2.1456,-3.5472) (2.0460,-3.4906) (1.9480,-3.4340) (1.8516,-3.3774) (1.7569,-3.3208) (1.6637,-3.2642) (1.5721,-3.2075) (1.4821,-3.1509) (1.3937,-3.0943) (1.3070,-3.0377) (1.2218,-2.9811) (1.1382,-2.9245) (1.0562,-2.8679) (0.9759,-2.8113) (0.8971,-2.7547) (0.8200,-2.6981) (0.7444,-2.6415) (0.6704,-2.5849) (0.5981,-2.5283) (0.5273,-2.4717) (0.4582,-2.4151) (0.3906,-2.3585) (0.3247,-2.3019) (0.2603,-2.2453) (0.1976,-2.1887) (0.1364,-2.1321) (0.0769,-2.0755) (0.0190,-2.0189) (-0.0374,-1.9623) (-0.0921,-1.9057) (-0.1452,-1.8491) (-0.1968,-1.7925) (-0.2467,-1.7358) (-0.2950,-1.6792) (-0.3418,-1.6226) (-0.3869,-1.5660) (-0.4304,-1.5094) (-0.4723,-1.4528) (-0.5126,-1.3962) (-0.5514,-1.3396) (-0.5885,-1.2830) (-0.6240,-1.2264) (-0.6579,-1.1698) (-0.6902,-1.1132) (-0.7209,-1.0566) (-0.7500,-1.0000) (-0.7775,-0.9434) (-0.8034,-0.8868) (-0.8277,-0.8302) (-0.8504,-0.7736) (-0.8715,-0.7170) (-0.8910,-0.6604) (-0.9089,-0.6038) (-0.9252,-0.5472) (-0.9398,-0.4906) (-0.9529,-0.4340) (-0.9644,-0.3774) (-0.9743,-0.3208) (-0.9826,-0.2642) (-0.9892,-0.2075) (-0.9943,-0.1509) (-0.9978,-0.0943) (-0.9996,-0.0377) (-0.9999,0.0189) (-0.9986,0.0755) (-0.9956,0.1321) (-0.9911,0.1887) (-0.9850,0.2453) (-0.9772,0.3019) (-0.9679,0.3585) (-0.9569,0.4151) (-0.9444,0.4717) (-0.9302,0.5283) (-0.9145,0.5849) (-0.8971,0.6415) (-0.8782,0.6981) (-0.8576,0.7547) (-0.8354,0.8113) (-0.8117,0.8679) (-0.7863,0.9245) (-0.7593,0.9811) (-0.7308,1.0377) (-0.7006,1.0943) (-0.6688,1.1509) (-0.6355,1.2075) (-0.6005,1.2642) (-0.5639,1.3208) (-0.5257,1.3774) (-0.4859,1.4340) (-0.4446,1.4906) (-0.4016,1.5472) (-0.3570,1.6038) (-0.3108,1.6604) (-0.2630,1.7170) (-0.2136,1.7736) (-0.1626,1.8302) (-0.1100,1.8868) (-0.0558,1.9434) (0.0000,2.0000) (0.0574,2.0566) (0.1164,2.1132) (0.1770,2.1698) (0.2392,2.2264) (0.3030,2.2830) (0.3685,2.3396) (0.4355,2.3962) (0.5041,2.4528) (0.5743,2.5094) (0.6461,2.5660) (0.7196,2.6226) (0.7946,2.6792) (0.8712,2.7358) (0.9494,2.7925) (1.0293,2.8491) (1.1107,2.9057) (1.1938,2.9623) (1.2784,3.0189) (1.3646,3.0755) (1.4525,3.1321) (1.5419,3.1887) (1.6330,3.2453) (1.7256,3.3019) (1.8199,3.3585) (1.9157,3.4151) (2.0132,3.4717) (2.1122,3.5283) (2.2129,3.5849) (2.3151,3.6415) (2.4190,3.6981) (2.5245,3.7547) (2.6315,3.8113) (2.7402,3.8679) (2.8505,3.9245) (2.9624,3.9811) (3.0758,4.0377) (3.1909,4.0943) (3.3076,4.1509) (3.4259,4.2075) (3.5457,4.2642) (3.6672,4.3208) (3.7903,4.3774) (3.9150,4.4340) (4.0413,4.4906) (4.1692,4.5472) (4.2987,4.6038) (4.4298,4.6604) (4.5625,4.7170) (4.6968,4.7736) (4.8327,4.8302) (4.9702,4.8868)};
\draw[structurecolor] plot coordinates {(4.0000,-4.0000) (3.8876,-3.9434) (3.7768,-3.8868) (3.6676,-3.8302) (3.5600,-3.7736) (3.4540,-3.7170) (3.3496,-3.6604) (3.2468,-3.6038) (3.1456,-3.5472) (3.0460,-3.4906) (2.9480,-3.4340) (2.8516,-3.3774) (2.7569,-3.3208) (2.6637,-3.2642) (2.5721,-3.2075) (2.4821,-3.1509) (2.3937,-3.0943) (2.3070,-3.0377) (2.2218,-2.9811) (2.1382,-2.9245) (2.0562,-2.8679) (1.9759,-2.8113) (1.8971,-2.7547) (1.8200,-2.6981) (1.7444,-2.6415) (1.6704,-2.5849) (1.5981,-2.5283) (1.5273,-2.4717) (1.4582,-2.4151) (1.3906,-2.3585) (1.3247,-2.3019) (1.2603,-2.2453) (1.1976,-2.1887) (1.1364,-2.1321) (1.0769,-2.0755) (1.0190,-2.0189) (0.9626,-1.9623) (0.9079,-1.9057) (0.8548,-1.8491) (0.8032,-1.7925) (0.7533,-1.7358) (0.7050,-1.6792) (0.6582,-1.6226) (0.6131,-1.5660) (0.5696,-1.5094) (0.5277,-1.4528) (0.4874,-1.3962) (0.4486,-1.3396) (0.4115,-1.2830) (0.3760,-1.2264) (0.3421,-1.1698) (0.3098,-1.1132) (0.2791,-1.0566) (0.2500,-1.0000) (0.2225,-0.9434) (0.1966,-0.8868) (0.1723,-0.8302) (0.1496,-0.7736) (0.1285,-0.7170) (0.1090,-0.6604) (0.0911,-0.6038) (0.0748,-0.5472) (0.0602,-0.4906) (0.0471,-0.4340) (0.0356,-0.3774) (0.0257,-0.3208) (0.0174,-0.2642) (0.0108,-0.2075) (0.0057,-0.1509) (0.0022,-0.0943) (0.0004,-0.0377) (0.0001,0.0189) (0.0014,0.0755) (0.0044,0.1321) (0.0089,0.1887) (0.0150,0.2453) (0.0228,0.3019) (0.0321,0.3585) (0.0431,0.4151) (0.0556,0.4717) (0.0698,0.5283) (0.0855,0.5849) (0.1029,0.6415) (0.1218,0.6981) (0.1424,0.7547) (0.1646,0.8113) (0.1883,0.8679) (0.2137,0.9245) (0.2407,0.9811) (0.2692,1.0377) (0.2994,1.0943) (0.3312,1.1509) (0.3645,1.2075) (0.3995,1.2642) (0.4361,1.3208) (0.4743,1.3774) (0.5141,1.4340) (0.5554,1.4906) (0.5984,1.5472) (0.6430,1.6038) (0.6892,1.6604) (0.7370,1.7170) (0.7864,1.7736) (0.8374,1.8302) (0.8900,1.8868) (0.9442,1.9434) (1.0000,2.0000) (1.0574,2.0566) (1.1164,2.1132) (1.1770,2.1698) (1.2392,2.2264) (1.3030,2.2830) (1.3685,2.3396) (1.4355,2.3962) (1.5041,2.4528) (1.5743,2.5094) (1.6461,2.5660) (1.7196,2.6226) (1.7946,2.6792) (1.8712,2.7358) (1.9494,2.7925) (2.0293,2.8491) (2.1107,2.9057) (2.1938,2.9623) (2.2784,3.0189) (2.3646,3.0755) (2.4525,3.1321) (2.5419,3.1887) (2.6330,3.2453) (2.7256,3.3019) (2.8199,3.3585) (2.9157,3.4151) (3.0132,3.4717) (3.1122,3.5283) (3.2129,3.5849) (3.3151,3.6415) (3.4190,3.6981) (3.5245,3.7547) (3.6315,3.8113) (3.7402,3.8679) (3.8505,3.9245) (3.9624,3.9811) (4.0758,4.0377) (4.1909,4.0943) (4.3076,4.1509) (4.4259,4.2075) (4.5457,4.2642) (4.6672,4.3208) (4.7903,4.3774) (4.9150,4.4340)};
\draw[structurecolor,dashed] plot coordinates {(5.0000,-4.0000) (4.8876,-3.9434) (4.7768,-3.8868) (4.6676,-3.8302) (4.5600,-3.7736) (4.4540,-3.7170) (4.3496,-3.6604) (4.2468,-3.6038) (4.1456,-3.5472) (4.0460,-3.4906) (3.9480,-3.4340) (3.8516,-3.3774) (3.7569,-3.3208) (3.6637,-3.2642) (3.5721,-3.2075) (3.4821,-3.1509) (3.3937,-3.0943) (3.3070,-3.0377) (3.2218,-2.9811) (3.1382,-2.9245) (3.0562,-2.8679) (2.9759,-2.8113) (2.8971,-2.7547) (2.8200,-2.6981) (2.7444,-2.6415) (2.6704,-2.5849) (2.5981,-2.5283) (2.5273,-2.4717) (2.4582,-2.4151) (2.3906,-2.3585) (2.3247,-2.3019) (2.2603,-2.2453) (2.1976,-2.1887) (2.1364,-2.1321) (2.0769,-2.0755) (2.0190,-2.0189) (1.9626,-1.9623) (1.9079,-1.9057) (1.8548,-1.8491) (1.8032,-1.7925) (1.7533,-1.7358) (1.7050,-1.6792) (1.6582,-1.6226) (1.6131,-1.5660) (1.5696,-1.5094) (1.5277,-1.4528) (1.4874,-1.3962) (1.4486,-1.3396) (1.4115,-1.2830) (1.3760,-1.2264) (1.3421,-1.1698) (1.3098,-1.1132) (1.2791,-1.0566) (1.2500,-1.0000) (1.2225,-0.9434) (1.1966,-0.8868) (1.1723,-0.8302) (1.1496,-0.7736) (1.1285,-0.7170) (1.1090,-0.6604) (1.0911,-0.6038) (1.0748,-0.5472) (1.0602,-0.4906) (1.0471,-0.4340) (1.0356,-0.3774) (1.0257,-0.3208) (1.0174,-0.2642) (1.0108,-0.2075) (1.0057,-0.1509) (1.0022,-0.0943) (1.0004,-0.0377) (1.0001,0.0189) (1.0014,0.0755) (1.0044,0.1321) (1.0089,0.1887) (1.0150,0.2453) (1.0228,0.3019) (1.0321,0.3585) (1.0431,0.4151) (1.0556,0.4717) (1.0698,0.5283) (1.0855,0.5849) (1.1029,0.6415) (1.1218,0.6981) (1.1424,0.7547) (1.1646,0.8113) (1.1883,0.8679) (1.2137,0.9245) (1.2407,0.9811) (1.2692,1.0377) (1.2994,1.0943) (1.3312,1.1509) (1.3645,1.2075) (1.3995,1.2642) (1.4361,1.3208) (1.4743,1.3774) (1.5141,1.4340) (1.5554,1.4906) (1.5984,1.5472) (1.6430,1.6038) (1.6892,1.6604) (1.7370,1.7170) (1.7864,1.7736) (1.8374,1.8302) (1.8900,1.8868) (1.9442,1.9434) (2.0000,2.0000) (2.0574,2.0566) (2.1164,2.1132) (2.1770,2.1698) (2.2392,2.2264) (2.3030,2.2830) (2.3685,2.3396) (2.4355,2.3962) (2.5041,2.4528) (2.5743,2.5094) (2.6461,2.5660) (2.7196,2.6226) (2.7946,2.6792) (2.8712,2.7358) (2.9494,2.7925) (3.0293,2.8491) (3.1107,2.9057) (3.1938,2.9623) (3.2784,3.0189) (3.3646,3.0755) (3.4525,3.1321) (3.5419,3.1887) (3.6330,3.2453) (3.7256,3.3019) (3.8199,3.3585) (3.9157,3.4151) (4.0132,3.4717) (4.1122,3.5283) (4.2129,3.5849) (4.3151,3.6415) (4.4190,3.6981) (4.5245,3.7547) (4.6315,3.8113) (4.7402,3.8679) (4.8505,3.9245) (4.9624,3.9811)};
\draw[main,thick] plot coordinates {(1.0000,0.0000) (0.9916,-0.0083) (0.9833,-0.0166) (0.9749,-0.0248) (0.9665,-0.0329) (0.9582,-0.0410) (0.9498,-0.0489) (0.9414,-0.0568) (0.9331,-0.0647) (0.9247,-0.0724) (0.9163,-0.0801) (0.9079,-0.0878) (0.8996,-0.0953) (0.8912,-0.1028) (0.8828,-0.1102) (0.8745,-0.1175) (0.8661,-0.1247) (0.8577,-0.1319) (0.8494,-0.1390) (0.8410,-0.1461) (0.8326,-0.1530) (0.8243,-0.1599) (0.8159,-0.1667) (0.8075,-0.1734) (0.7992,-0.1801) (0.7908,-0.1867) (0.7824,-0.1932) (0.7741,-0.1996) (0.7657,-0.2060) (0.7573,-0.2123) (0.7490,-0.2185) (0.7406,-0.2246) (0.7322,-0.2306) (0.7238,-0.2366) (0.7155,-0.2425) (0.7071,-0.2483) (0.6987,-0.2541) (0.6904,-0.2598) (0.6820,-0.2654) (0.6736,-0.2709) (0.6653,-0.2763) (0.6569,-0.2817) (0.6485,-0.2870) (0.6402,-0.2922) (0.6318,-0.2974) (0.6234,-0.3024) (0.6151,-0.3074) (0.6067,-0.3123) (0.5983,-0.3172) (0.5900,-0.3219) (0.5816,-0.3266) (0.5732,-0.3312) (0.5649,-0.3357) (0.5565,-0.3402) (0.5481,-0.3446) (0.5397,-0.3489) (0.5314,-0.3531) (0.5230,-0.3572) (0.5146,-0.3613) (0.5063,-0.3653) (0.4979,-0.3692) (0.4895,-0.3731) (0.4812,-0.3768) (0.4728,-0.3805) (0.4644,-0.3841) (0.4561,-0.3877) (0.4477,-0.3911) (0.4393,-0.3945) (0.4310,-0.3978) (0.4226,-0.4011) (0.4142,-0.4042) (0.4059,-0.4073) (0.3975,-0.4103) (0.3891,-0.4133) (0.3808,-0.4161) (0.3724,-0.4189) (0.3640,-0.4216) (0.3556,-0.4243) (0.3473,-0.4268) (0.3389,-0.4293) (0.3305,-0.4317) (0.3222,-0.4341) (0.3138,-0.4363) (0.3054,-0.4385) (0.2971,-0.4406) (0.2887,-0.4427) (0.2803,-0.4447) (0.2720,-0.4466) (0.2636,-0.4484) (0.2552,-0.4501) (0.2469,-0.4518) (0.2385,-0.4534) (0.2301,-0.4549) (0.2218,-0.4564) (0.2134,-0.4578) (0.2050,-0.4591) (0.1967,-0.4603) (0.1883,-0.4615) (0.1799,-0.4626) (0.1715,-0.4636) (0.1632,-0.4646) (0.1548,-0.4654) (0.1464,-0.4662) (0.1381,-0.4670) (0.1297,-0.4676) (0.1213,-0.4682) (0.1130,-0.4687) (0.1046,-0.4692) (0.0962,-0.4696) (0.0879,-0.4699) (0.0795,-0.4701) (0.0711,-0.4703) (0.0628,-0.4704) (0.0544,-0.4704) (0.0460,-0.4704) (0.0377,-0.4703) (0.0293,-0.4701) (0.0209,-0.4698) (0.0126,-0.4695) (0.0042,-0.4691) (-0.0042,-0.4687) (-0.0126,-0.4681) (-0.0209,-0.4675) (-0.0293,-0.4669) (-0.0377,-0.4661) (-0.0460,-0.4653) (-0.0544,-0.4644) (-0.0628,-0.4635) (-0.0711,-0.4625) (-0.0795,-0.4614) (-0.0879,-0.4603) (-0.0962,-0.4591) (-0.1046,-0.4578) (-0.1130,-0.4564) (-0.1213,-0.4550) (-0.1297,-0.4535) (-0.1381,-0.4520) (-0.1464,-0.4504) (-0.1548,-0.4487) (-0.1632,-0.4469) (-0.1715,-0.4451) (-0.1799,-0.4432) (-0.1883,-0.4413) (-0.1967,-0.4393) (-0.2050,-0.4372) (-0.2134,-0.4350) (-0.2218,-0.4328) (-0.2301,-0.4306) (-0.2385,-0.4282) (-0.2469,-0.4258) (-0.2552,-0.4233) (-0.2636,-0.4208) (-0.2720,-0.4182) (-0.2803,-0.4155) (-0.2887,-0.4127) (-0.2971,-0.4099) (-0.3054,-0.4071) (-0.3138,-0.4041) (-0.3222,-0.4011) (-0.3305,-0.3981) (-0.3389,-0.3949) (-0.3473,-0.3917) (-0.3556,-0.3885) (-0.3640,-0.3852) (-0.3724,-0.3818) (-0.3808,-0.3783) (-0.3891,-0.3748) (-0.3975,-0.3712) (-0.4059,-0.3676) (-0.4142,-0.3639) (-0.4226,-0.3601) (-0.4310,-0.3562) (-0.4393,-0.3523) (-0.4477,-0.3484) (-0.4561,-0.3443) (-0.4644,-0.3402) (-0.4728,-0.3361) (-0.4812,-0.3319) (-0.4895,-0.3276) (-0.4979,-0.3232) (-0.5063,-0.3188) (-0.5146,-0.3143) (-0.5230,-0.3098) (-0.5314,-0.3052) (-0.5397,-0.3005) (-0.5481,-0.2958) (-0.5565,-0.2910) (-0.5649,-0.2861) (-0.5732,-0.2812) (-0.5816,-0.2762) (-0.5900,-0.2711) (-0.5983,-0.2660) (-0.6067,-0.2608) (-0.6151,-0.2555) (-0.6234,-0.2502) (-0.6318,-0.2448) (-0.6402,-0.2394) (-0.6485,-0.2339) (-0.6569,-0.2283) (-0.6653,-0.2227) (-0.6736,-0.2170) (-0.6820,-0.2112) (-0.6904,-0.2054) (-0.6987,-0.1995) (-0.7071,-0.1935) (-0.7155,-0.1875) (-0.7238,-0.1814) (-0.7322,-0.1752) (-0.7406,-0.1690) (-0.7490,-0.1627) (-0.7573,-0.1564) (-0.7657,-0.1499) (-0.7741,-0.1435) (-0.7824,-0.1369) (-0.7908,-0.1303) (-0.7992,-0.1236) (-0.8075,-0.1168) (-0.8159,-0.1100) (-0.8243,-0.1031) (-0.8326,-0.0962) (-0.8410,-0.0892) (-0.8494,-0.0821) (-0.8577,-0.0749) (-0.8661,-0.0677) (-0.8745,-0.0604) (-0.8828,-0.0530) (-0.8912,-0.0456) (-0.8996,-0.0381) (-0.9079,-0.0306) (-0.9163,-0.0229) (-0.9247,-0.0152) (-0.9331,-0.0074) (-0.9414,0.0004) (-0.9498,0.0083) (-0.9582,0.0163) (-0.9665,0.0243) (-0.9749,0.0325) (-0.9833,0.0407) (-0.9916,0.0489) (-1.0000,0.0573)};
\draw[main,thick] plot coordinates {(1.0000,0.0000) (1.0167,0.0169) (1.0335,0.0340) (1.0502,0.0515) (1.0669,0.0692) (1.0837,0.0871) (1.1004,0.1054) (1.1172,0.1239) (1.1339,0.1426) (1.1506,0.1617) (1.1674,0.1809) (1.1841,0.2005) (1.2008,0.2202) (1.2176,0.2402) (1.2343,0.2605) (1.2510,0.2810) (1.2678,0.3017) (1.2845,0.3227) (1.3013,0.3438) (1.3180,0.3652) (1.3347,0.3868) (1.3515,0.4087) (1.3682,0.4307) (1.3849,0.4529) (1.4017,0.4753) (1.4184,0.4979) (1.4351,0.5207) (1.4519,0.5437) (1.4686,0.5668) (1.4854,0.5902) (1.5021,0.6136) (1.5188,0.6373) (1.5356,0.6611) (1.5523,0.6850) (1.5690,0.7091) (1.5858,0.7333) (1.6025,0.7577) (1.6192,0.7822) (1.6360,0.8068) (1.6527,0.8315) (1.6695,0.8563) (1.6862,0.8812) (1.7029,0.9062) (1.7197,0.9314) (1.7364,0.9565) (1.7531,0.9818) (1.7699,1.0072) (1.7866,1.0326) (1.8033,1.0580) (1.8201,1.0836) (1.8368,1.1091) (1.8536,1.1347) (1.8703,1.1604) (1.8870,1.1861) (1.9038,1.2118) (1.9205,1.2375) (1.9372,1.2633) (1.9540,1.2890) (1.9707,1.3148) (1.9874,1.3405) (2.0042,1.3662) (2.0209,1.3920) (2.0377,1.4177) (2.0544,1.4434) (2.0711,1.4690) (2.0879,1.4946) (2.1046,1.5202) (2.1213,1.5457) (2.1381,1.5712) (2.1548,1.5966) (2.1715,1.6220) (2.1883,1.6473) (2.2050,1.6725) (2.2218,1.6977) (2.2385,1.7228) (2.2552,1.7478) (2.2720,1.7727) (2.2887,1.7975) (2.3054,1.8223) (2.3222,1.8469) (2.3389,1.8715) (2.3556,1.8959) (2.3724,1.9202) (2.3891,1.9445) (2.4059,1.9686) (2.4226,1.9926) (2.4393,2.0164) (2.4561,2.0402) (2.4728,2.0638) (2.4895,2.0873) (2.5063,2.1107) (2.5230,2.1339) (2.5397,2.1570) (2.5565,2.1800) (2.5732,2.2028) (2.5900,2.2255) (2.6067,2.2481) (2.6234,2.2705) (2.6402,2.2928) (2.6569,2.3149) (2.6736,2.3368) (2.6904,2.3587) (2.7071,2.3803) (2.7238,2.4019) (2.7406,2.4232) (2.7573,2.4445) (2.7741,2.4655) (2.7908,2.4864) (2.8075,2.5072) (2.8243,2.5278) (2.8410,2.5483) (2.8577,2.5686) (2.8745,2.5887) (2.8912,2.6087) (2.9079,2.6286) (2.9247,2.6482) (2.9414,2.6678) (2.9582,2.6871) (2.9749,2.7064) (2.9916,2.7254) (3.0084,2.7443) (3.0251,2.7631) (3.0418,2.7817) (3.0586,2.8002) (3.0753,2.8185) (3.0921,2.8366) (3.1088,2.8546) (3.1255,2.8725) (3.1423,2.8902) (3.1590,2.9078) (3.1757,2.9252) (3.1925,2.9425) (3.2092,2.9596) (3.2259,2.9766) (3.2427,2.9934) (3.2594,3.0101) (3.2762,3.0267) (3.2929,3.0432) (3.3096,3.0594) (3.3264,3.0756) (3.3431,3.0916) (3.3598,3.1075) (3.3766,3.1233) (3.3933,3.1389) (3.4100,3.1544) (3.4268,3.1698) (3.4435,3.1850) (3.4603,3.2002) (3.4770,3.2152) (3.4937,3.2300) (3.5105,3.2448) (3.5272,3.2594) (3.5439,3.2740) (3.5607,3.2884) (3.5774,3.3027) (3.5941,3.3168) (3.6109,3.3309) (3.6276,3.3449) (3.6444,3.3587) (3.6611,3.3725) (3.6778,3.3861) (3.6946,3.3996) (3.7113,3.4130) (3.7280,3.4264) (3.7448,3.4396) (3.7615,3.4527) (3.7782,3.4657) (3.7950,3.4787) (3.8117,3.4915) (3.8285,3.5042) (3.8452,3.5169) (3.8619,3.5294) (3.8787,3.5419) (3.8954,3.5543) (3.9121,3.5666) (3.9289,3.5788) (3.9456,3.5909) (3.9623,3.6030) (3.9791,3.6149) (3.9958,3.6268) (4.0126,3.6386) (4.0293,3.6503) (4.0460,3.6620) (4.0628,3.6735) (4.0795,3.6850) (4.0962,3.6965) (4.1130,3.7078) (4.1297,3.7191) (4.1464,3.7303) (4.1632,3.7414) (4.1799,3.7525) (4.1967,3.7635) (4.2134,3.7745) (4.2301,3.7853) (4.2469,3.7961) (4.2636,3.8069) (4.2803,3.8176) (4.2971,3.8282) (4.3138,3.8388) (4.3305,3.8493) (4.3473,3.8597) (4.3640,3.8701) (4.3808,3.8805) (4.3975,3.8908) (4.4142,3.9010) (4.4310,3.9112) (4.4477,3.9213) (4.4644,3.9314) (4.4812,3.9414) (4.4979,3.9514) (4.5146,3.9613) (4.5314,3.9712) (4.5481,3.9810) (4.5649,3.9908) (4.5816,4.0005) (4.5983,4.0102) (4.6151,4.0199) (4.6318,4.0295) (4.6485,4.0390) (4.6653,4.0486) (4.6820,4.0580) (4.6987,4.0675) (4.7155,4.0769) (4.7322,4.0862) (4.7490,4.0956) (4.7657,4.1048) (4.7824,4.1141) (4.7992,4.1233) (4.8159,4.1324) (4.8326,4.1416) (4.8494,4.1507) (4.8661,4.1597) (4.8828,4.1688) (4.8996,4.1777) (4.9163,4.1867) (4.9331,4.1956) (4.9498,4.2045) (4.9665,4.2134) (4.9833,4.2222) (5.0000,4.2310)};
\fill[main](1,0)circle(2pt);
\draw(-1,-.08)--(-1,.08);\draw(1,-.08)--(1,.08);\draw(2,-.08)--(2,.08);\node[below,font=\scriptsize,fill=white,inner sep=.2pt]at(2,-.08){$ 2 $};\draw(3,-.08)--(3,.08);\draw(4,-.08)--(4,.08);\node[below,font=\scriptsize,fill=white,inner sep=.2pt]at(4,-.08){$ 4 $};\draw(5,-.08)--(5,.08);\node[below,font=\scriptsize,fill=white,inner sep=.2pt]at(5,-.08){$ 5 $};\draw(-.08,-4)--(.08,-4);\node[left,font=\scriptsize,fill=white,inner sep=.2pt]at(-.08,-4){$ -4 $};\draw(-.08,-3)--(.08,-3);\draw(-.08,-2)--(.08,-2);\node[left,font=\scriptsize,fill=white,inner sep=.2pt]at(-.08,-2){$ -2 $};\draw(-.08,-1)--(.08,-1);\draw(-.08,1)--(.08,1);\draw(-.08,2)--(.08,2);\node[left,font=\scriptsize,fill=white,inner sep=.2pt]at(-.08,2){$ 2 $};\draw(-.08,3)--(.08,3);\draw(-.08,4)--(.08,4);\node[left,font=\scriptsize,fill=white,inner sep=.2pt]at(-.08,4){$ 4 $};\draw(-.08,5)--(.08,5);\end{tikzpicture}\end{center}
}

\newcommand{\ODEExamPracticeField}{%
\begin{center}\begin{tikzpicture}[x=.62cm,y=.53cm]
\clip (-6.6,-6.6) rectangle (6.7,6.8);
\draw[->](-6,0)--(6.35,0)node[right]{$x$};\draw[->](0,-6)--(0,6.35)node[above]{$y$};
\draw[gray!55](-6.120,-6.000)--(-5.880,-6.000);
\draw[gray!55](-6.021,-5.382)--(-5.979,-5.618);
\draw[gray!55](-6.011,-4.880)--(-5.989,-5.120);
\draw[gray!55](-6.008,-4.380)--(-5.992,-4.620);
\draw[gray!55](-6.006,-3.880)--(-5.994,-4.120);
\draw[gray!55](-6.005,-3.380)--(-5.995,-3.620);
\draw[gray!55](-6.004,-2.880)--(-5.996,-3.120);
\draw[gray!55](-6.004,-2.380)--(-5.996,-2.620);
\draw[gray!55](-6.004,-1.880)--(-5.996,-2.120);
\draw[gray!55](-6.004,-1.380)--(-5.996,-1.620);
\draw[gray!55](-6.003,-0.880)--(-5.997,-1.120);
\draw[gray!55](-6.003,-0.380)--(-5.997,-0.620);
\draw[gray!55](-6.003,0.120)--(-5.997,-0.120);
\draw[gray!55](-6.003,0.620)--(-5.997,0.380);
\draw[gray!55](-6.003,1.120)--(-5.997,0.880);
\draw[gray!55](-6.004,1.620)--(-5.996,1.380);
\draw[gray!55](-6.004,2.120)--(-5.996,1.880);
\draw[gray!55](-6.004,2.620)--(-5.996,2.380);
\draw[gray!55](-6.004,3.120)--(-5.996,2.880);
\draw[gray!55](-6.005,3.620)--(-5.995,3.380);
\draw[gray!55](-6.006,4.120)--(-5.994,3.880);
\draw[gray!55](-6.008,4.620)--(-5.992,4.380);
\draw[gray!55](-6.011,5.120)--(-5.989,4.880);
\draw[gray!55](-6.021,5.618)--(-5.979,5.382);
\draw[gray!55](-6.120,6.000)--(-5.880,6.000);
\draw[gray!55](-5.521,-6.118)--(-5.479,-5.882);
\draw[gray!55](-5.620,-5.500)--(-5.380,-5.500);
\draw[gray!55](-5.522,-4.882)--(-5.478,-5.118);
\draw[gray!55](-5.512,-4.381)--(-5.488,-4.619);
\draw[gray!55](-5.508,-3.880)--(-5.492,-4.120);
\draw[gray!55](-5.507,-3.380)--(-5.493,-3.620);
\draw[gray!55](-5.506,-2.880)--(-5.494,-3.120);
\draw[gray!55](-5.505,-2.380)--(-5.495,-2.620);
\draw[gray!55](-5.505,-1.880)--(-5.495,-2.120);
\draw[gray!55](-5.504,-1.380)--(-5.496,-1.620);
\draw[gray!55](-5.504,-0.880)--(-5.496,-1.120);
\draw[gray!55](-5.504,-0.380)--(-5.496,-0.620);
\draw[gray!55](-5.504,0.120)--(-5.496,-0.120);
\draw[gray!55](-5.504,0.620)--(-5.496,0.380);
\draw[gray!55](-5.504,1.120)--(-5.496,0.880);
\draw[gray!55](-5.504,1.620)--(-5.496,1.380);
\draw[gray!55](-5.505,2.120)--(-5.495,1.880);
\draw[gray!55](-5.505,2.620)--(-5.495,2.380);
\draw[gray!55](-5.506,3.120)--(-5.494,2.880);
\draw[gray!55](-5.507,3.620)--(-5.493,3.380);
\draw[gray!55](-5.508,4.120)--(-5.492,3.880);
\draw[gray!55](-5.512,4.619)--(-5.488,4.381);
\draw[gray!55](-5.522,5.118)--(-5.478,4.882);
\draw[gray!55](-5.620,5.500)--(-5.380,5.500);
\draw[gray!55](-5.521,5.882)--(-5.479,6.118);
\draw[gray!55](-5.011,-6.120)--(-4.989,-5.880);
\draw[gray!55](-5.022,-5.618)--(-4.978,-5.382);
\draw[gray!55](-5.120,-5.000)--(-4.880,-5.000);
\draw[gray!55](-5.025,-4.383)--(-4.975,-4.617);
\draw[gray!55](-5.013,-3.881)--(-4.987,-4.119);
\draw[gray!55](-5.009,-3.380)--(-4.991,-3.620);
\draw[gray!55](-5.007,-2.880)--(-4.993,-3.120);
\draw[gray!55](-5.006,-2.380)--(-4.994,-2.620);
\draw[gray!55](-5.006,-1.880)--(-4.994,-2.120);
\draw[gray!55](-5.005,-1.380)--(-4.995,-1.620);
\draw[gray!55](-5.005,-0.880)--(-4.995,-1.120);
\draw[gray!55](-5.005,-0.380)--(-4.995,-0.620);
\draw[gray!55](-5.005,0.120)--(-4.995,-0.120);
\draw[gray!55](-5.005,0.620)--(-4.995,0.380);
\draw[gray!55](-5.005,1.120)--(-4.995,0.880);
\draw[gray!55](-5.005,1.620)--(-4.995,1.380);
\draw[gray!55](-5.006,2.120)--(-4.994,1.880);
\draw[gray!55](-5.006,2.620)--(-4.994,2.380);
\draw[gray!55](-5.007,3.120)--(-4.993,2.880);
\draw[gray!55](-5.009,3.620)--(-4.991,3.380);
\draw[gray!55](-5.013,4.119)--(-4.987,3.881);
\draw[gray!55](-5.025,4.617)--(-4.975,4.383);
\draw[gray!55](-5.120,5.000)--(-4.880,5.000);
\draw[gray!55](-5.022,5.382)--(-4.978,5.618);
\draw[gray!55](-5.011,5.880)--(-4.989,6.120);
\draw[gray!55](-4.508,-6.120)--(-4.492,-5.880);
\draw[gray!55](-4.512,-5.619)--(-4.488,-5.381);
\draw[gray!55](-4.525,-5.117)--(-4.475,-4.883);
\draw[gray!55](-4.620,-4.500)--(-4.380,-4.500);
\draw[gray!55](-4.527,-3.883)--(-4.473,-4.117);
\draw[gray!55](-4.515,-3.381)--(-4.485,-3.619);
\draw[gray!55](-4.511,-2.880)--(-4.489,-3.120);
\draw[gray!55](-4.509,-2.380)--(-4.491,-2.620);
\draw[gray!55](-4.507,-1.880)--(-4.493,-2.120);
\draw[gray!55](-4.507,-1.380)--(-4.493,-1.620);
\draw[gray!55](-4.506,-0.880)--(-4.494,-1.120);
\draw[gray!55](-4.506,-0.380)--(-4.494,-0.620);
\draw[gray!55](-4.506,0.120)--(-4.494,-0.120);
\draw[gray!55](-4.506,0.620)--(-4.494,0.380);
\draw[gray!55](-4.506,1.120)--(-4.494,0.880);
\draw[gray!55](-4.507,1.620)--(-4.493,1.380);
\draw[gray!55](-4.507,2.120)--(-4.493,1.880);
\draw[gray!55](-4.509,2.620)--(-4.491,2.380);
\draw[gray!55](-4.511,3.120)--(-4.489,2.880);
\draw[gray!55](-4.515,3.619)--(-4.485,3.381);
\draw[gray!55](-4.527,4.117)--(-4.473,3.883);
\draw[gray!55](-4.620,4.500)--(-4.380,4.500);
\draw[gray!55](-4.525,4.883)--(-4.475,5.117);
\draw[gray!55](-4.512,5.381)--(-4.488,5.619);
\draw[gray!55](-4.508,5.880)--(-4.492,6.120);
\draw[gray!55](-4.006,-6.120)--(-3.994,-5.880);
\draw[gray!55](-4.008,-5.620)--(-3.992,-5.380);
\draw[gray!55](-4.013,-5.119)--(-3.987,-4.881);
\draw[gray!55](-4.027,-4.617)--(-3.973,-4.383);
\draw[gray!55](-4.120,-4.000)--(-3.880,-4.000);
\draw[gray!55](-4.031,-3.384)--(-3.969,-3.616);
\draw[gray!55](-4.017,-2.881)--(-3.983,-3.119);
\draw[gray!55](-4.012,-2.381)--(-3.988,-2.619);
\draw[gray!55](-4.010,-1.880)--(-3.990,-2.120);
\draw[gray!55](-4.009,-1.380)--(-3.991,-1.620);
\draw[gray!55](-4.008,-0.880)--(-3.992,-1.120);
\draw[gray!55](-4.008,-0.380)--(-3.992,-0.620);
\draw[gray!55](-4.007,0.120)--(-3.993,-0.120);
\draw[gray!55](-4.008,0.620)--(-3.992,0.380);
\draw[gray!55](-4.008,1.120)--(-3.992,0.880);
\draw[gray!55](-4.009,1.620)--(-3.991,1.380);
\draw[gray!55](-4.010,2.120)--(-3.990,1.880);
\draw[gray!55](-4.012,2.619)--(-3.988,2.381);
\draw[gray!55](-4.017,3.119)--(-3.983,2.881);
\draw[gray!55](-4.031,3.616)--(-3.969,3.384);
\draw[gray!55](-4.120,4.000)--(-3.880,4.000);
\draw[gray!55](-4.027,4.383)--(-3.973,4.617);
\draw[gray!55](-4.013,4.881)--(-3.987,5.119);
\draw[gray!55](-4.008,5.380)--(-3.992,5.620);
\draw[gray!55](-4.006,5.880)--(-3.994,6.120);
\draw[gray!55](-3.505,-6.120)--(-3.495,-5.880);
\draw[gray!55](-3.507,-5.620)--(-3.493,-5.380);
\draw[gray!55](-3.509,-5.120)--(-3.491,-4.880);
\draw[gray!55](-3.515,-4.619)--(-3.485,-4.381);
\draw[gray!55](-3.531,-4.116)--(-3.469,-3.884);
\draw[gray!55](-3.620,-3.500)--(-3.380,-3.500);
\draw[gray!55](-3.535,-2.885)--(-3.465,-3.115);
\draw[gray!55](-3.520,-2.382)--(-3.480,-2.618);
\draw[gray!55](-3.514,-1.881)--(-3.486,-2.119);
\draw[gray!55](-3.512,-1.381)--(-3.488,-1.619);
\draw[gray!55](-3.511,-0.880)--(-3.489,-1.120);
\draw[gray!55](-3.510,-0.380)--(-3.490,-0.620);
\draw[gray!55](-3.510,0.120)--(-3.490,-0.120);
\draw[gray!55](-3.510,0.620)--(-3.490,0.380);
\draw[gray!55](-3.511,1.120)--(-3.489,0.880);
\draw[gray!55](-3.512,1.619)--(-3.488,1.381);
\draw[gray!55](-3.514,2.119)--(-3.486,1.881);
\draw[gray!55](-3.520,2.618)--(-3.480,2.382);
\draw[gray!55](-3.535,3.115)--(-3.465,2.885);
\draw[gray!55](-3.620,3.500)--(-3.380,3.500);
\draw[gray!55](-3.531,3.884)--(-3.469,4.116);
\draw[gray!55](-3.515,4.381)--(-3.485,4.619);
\draw[gray!55](-3.509,4.880)--(-3.491,5.120);
\draw[gray!55](-3.507,5.380)--(-3.493,5.620);
\draw[gray!55](-3.505,5.880)--(-3.495,6.120);
\draw[gray!55](-3.004,-6.120)--(-2.996,-5.880);
\draw[gray!55](-3.006,-5.620)--(-2.994,-5.380);
\draw[gray!55](-3.007,-5.120)--(-2.993,-4.880);
\draw[gray!55](-3.011,-4.620)--(-2.989,-4.380);
\draw[gray!55](-3.017,-4.119)--(-2.983,-3.881);
\draw[gray!55](-3.035,-3.615)--(-2.965,-3.385);
\draw[gray!55](-3.120,-3.000)--(-2.880,-3.000);
\draw[gray!55](-3.041,-2.387)--(-2.959,-2.613);
\draw[gray!55](-3.024,-1.882)--(-2.976,-2.118);
\draw[gray!55](-3.018,-1.381)--(-2.982,-1.619);
\draw[gray!55](-3.015,-0.881)--(-2.985,-1.119);
\draw[gray!55](-3.014,-0.381)--(-2.986,-0.619);
\draw[gray!55](-3.013,0.119)--(-2.987,-0.119);
\draw[gray!55](-3.014,0.619)--(-2.986,0.381);
\draw[gray!55](-3.015,1.119)--(-2.985,0.881);
\draw[gray!55](-3.018,1.619)--(-2.982,1.381);
\draw[gray!55](-3.024,2.118)--(-2.976,1.882);
\draw[gray!55](-3.041,2.613)--(-2.959,2.387);
\draw[gray!55](-3.120,3.000)--(-2.880,3.000);
\draw[gray!55](-3.035,3.385)--(-2.965,3.615);
\draw[gray!55](-3.017,3.881)--(-2.983,4.119);
\draw[gray!55](-3.011,4.380)--(-2.989,4.620);
\draw[gray!55](-3.007,4.880)--(-2.993,5.120);
\draw[gray!55](-3.006,5.380)--(-2.994,5.620);
\draw[gray!55](-3.004,5.880)--(-2.996,6.120);
\draw[gray!55](-2.504,-6.120)--(-2.496,-5.880);
\draw[gray!55](-2.505,-5.620)--(-2.495,-5.380);
\draw[gray!55](-2.506,-5.120)--(-2.494,-4.880);
\draw[gray!55](-2.509,-4.620)--(-2.491,-4.380);
\draw[gray!55](-2.512,-4.119)--(-2.488,-3.881);
\draw[gray!55](-2.520,-3.618)--(-2.480,-3.382);
\draw[gray!55](-2.541,-3.113)--(-2.459,-2.887);
\draw[gray!55](-2.620,-2.500)--(-2.380,-2.500);
\draw[gray!55](-2.549,-1.890)--(-2.451,-2.110);
\draw[gray!55](-2.529,-1.384)--(-2.471,-1.616);
\draw[gray!55](-2.522,-0.882)--(-2.478,-1.118);
\draw[gray!55](-2.520,-0.382)--(-2.480,-0.618);
\draw[gray!55](-2.519,0.118)--(-2.481,-0.118);
\draw[gray!55](-2.520,0.618)--(-2.480,0.382);
\draw[gray!55](-2.522,1.118)--(-2.478,0.882);
\draw[gray!55](-2.529,1.616)--(-2.471,1.384);
\draw[gray!55](-2.549,2.110)--(-2.451,1.890);
\draw[gray!55](-2.620,2.500)--(-2.380,2.500);
\draw[gray!55](-2.541,2.887)--(-2.459,3.113);
\draw[gray!55](-2.520,3.382)--(-2.480,3.618);
\draw[gray!55](-2.512,3.881)--(-2.488,4.119);
\draw[gray!55](-2.509,4.380)--(-2.491,4.620);
\draw[gray!55](-2.506,4.880)--(-2.494,5.120);
\draw[gray!55](-2.505,5.380)--(-2.495,5.620);
\draw[gray!55](-2.504,5.880)--(-2.496,6.120);
\draw[gray!55](-2.004,-6.120)--(-1.996,-5.880);
\draw[gray!55](-2.005,-5.620)--(-1.995,-5.380);
\draw[gray!55](-2.006,-5.120)--(-1.994,-4.880);
\draw[gray!55](-2.007,-4.620)--(-1.993,-4.380);
\draw[gray!55](-2.010,-4.120)--(-1.990,-3.880);
\draw[gray!55](-2.014,-3.619)--(-1.986,-3.381);
\draw[gray!55](-2.024,-3.118)--(-1.976,-2.882);
\draw[gray!55](-2.049,-2.610)--(-1.951,-2.390);
\draw[gray!55](-2.120,-2.000)--(-1.880,-2.000);
\draw[gray!55](-2.060,-1.396)--(-1.940,-1.604);
\draw[gray!55](-2.038,-0.886)--(-1.962,-1.114);
\draw[gray!55](-2.031,-0.384)--(-1.969,-0.616);
\draw[gray!55](-2.029,0.116)--(-1.971,-0.116);
\draw[gray!55](-2.031,0.616)--(-1.969,0.384);
\draw[gray!55](-2.038,1.114)--(-1.962,0.886);
\draw[gray!55](-2.060,1.604)--(-1.940,1.396);
\draw[gray!55](-2.120,2.000)--(-1.880,2.000);
\draw[gray!55](-2.049,2.390)--(-1.951,2.610);
\draw[gray!55](-2.024,2.882)--(-1.976,3.118);
\draw[gray!55](-2.014,3.381)--(-1.986,3.619);
\draw[gray!55](-2.010,3.880)--(-1.990,4.120);
\draw[gray!55](-2.007,4.380)--(-1.993,4.620);
\draw[gray!55](-2.006,4.880)--(-1.994,5.120);
\draw[gray!55](-2.005,5.380)--(-1.995,5.620);
\draw[gray!55](-2.004,5.880)--(-1.996,6.120);
\draw[gray!55](-1.504,-6.120)--(-1.496,-5.880);
\draw[gray!55](-1.504,-5.620)--(-1.496,-5.380);
\draw[gray!55](-1.505,-5.120)--(-1.495,-4.880);
\draw[gray!55](-1.507,-4.620)--(-1.493,-4.380);
\draw[gray!55](-1.509,-4.120)--(-1.491,-3.880);
\draw[gray!55](-1.512,-3.619)--(-1.488,-3.381);
\draw[gray!55](-1.518,-3.119)--(-1.482,-2.881);
\draw[gray!55](-1.529,-2.616)--(-1.471,-2.384);
\draw[gray!55](-1.560,-2.104)--(-1.440,-1.896);
\draw[gray!55](-1.620,-1.500)--(-1.380,-1.500);
\draw[gray!55](-1.575,-0.906)--(-1.425,-1.094);
\draw[gray!55](-1.554,-0.393)--(-1.446,-0.607);
\draw[gray!55](-1.549,0.110)--(-1.451,-0.110);
\draw[gray!55](-1.554,0.607)--(-1.446,0.393);
\draw[gray!55](-1.575,1.094)--(-1.425,0.906);
\draw[gray!55](-1.620,1.500)--(-1.380,1.500);
\draw[gray!55](-1.560,1.896)--(-1.440,2.104);
\draw[gray!55](-1.529,2.384)--(-1.471,2.616);
\draw[gray!55](-1.518,2.881)--(-1.482,3.119);
\draw[gray!55](-1.512,3.381)--(-1.488,3.619);
\draw[gray!55](-1.509,3.880)--(-1.491,4.120);
\draw[gray!55](-1.507,4.380)--(-1.493,4.620);
\draw[gray!55](-1.505,4.880)--(-1.495,5.120);
\draw[gray!55](-1.504,5.380)--(-1.496,5.620);
\draw[gray!55](-1.504,5.880)--(-1.496,6.120);
\draw[gray!55](-1.003,-6.120)--(-0.997,-5.880);
\draw[gray!55](-1.004,-5.620)--(-0.996,-5.380);
\draw[gray!55](-1.005,-5.120)--(-0.995,-4.880);
\draw[gray!55](-1.006,-4.620)--(-0.994,-4.380);
\draw[gray!55](-1.008,-4.120)--(-0.992,-3.880);
\draw[gray!55](-1.011,-3.620)--(-0.989,-3.380);
\draw[gray!55](-1.015,-3.119)--(-0.985,-2.881);
\draw[gray!55](-1.022,-2.618)--(-0.978,-2.382);
\draw[gray!55](-1.038,-2.114)--(-0.962,-1.886);
\draw[gray!55](-1.075,-1.594)--(-0.925,-1.406);
\draw[gray!55](-1.120,-1.000)--(-0.880,-1.000);
\draw[gray!55](-1.096,-0.428)--(-0.904,-0.572);
\draw[gray!55](-1.085,0.085)--(-0.915,-0.085);
\draw[gray!55](-1.096,0.572)--(-0.904,0.428);
\draw[gray!55](-1.120,1.000)--(-0.880,1.000);
\draw[gray!55](-1.075,1.406)--(-0.925,1.594);
\draw[gray!55](-1.038,1.886)--(-0.962,2.114);
\draw[gray!55](-1.022,2.382)--(-0.978,2.618);
\draw[gray!55](-1.015,2.881)--(-0.985,3.119);
\draw[gray!55](-1.011,3.380)--(-0.989,3.620);
\draw[gray!55](-1.008,3.880)--(-0.992,4.120);
\draw[gray!55](-1.006,4.380)--(-0.994,4.620);
\draw[gray!55](-1.005,4.880)--(-0.995,5.120);
\draw[gray!55](-1.004,5.380)--(-0.996,5.620);
\draw[gray!55](-1.003,5.880)--(-0.997,6.120);
\draw[gray!55](-0.503,-6.120)--(-0.497,-5.880);
\draw[gray!55](-0.504,-5.620)--(-0.496,-5.380);
\draw[gray!55](-0.505,-5.120)--(-0.495,-4.880);
\draw[gray!55](-0.506,-4.620)--(-0.494,-4.380);
\draw[gray!55](-0.508,-4.120)--(-0.492,-3.880);
\draw[gray!55](-0.510,-3.620)--(-0.490,-3.380);
\draw[gray!55](-0.514,-3.119)--(-0.486,-2.881);
\draw[gray!55](-0.520,-2.618)--(-0.480,-2.382);
\draw[gray!55](-0.531,-2.116)--(-0.469,-1.884);
\draw[gray!55](-0.554,-1.607)--(-0.446,-1.393);
\draw[gray!55](-0.596,-1.072)--(-0.404,-0.928);
\draw[gray!55](-0.620,-0.500)--(-0.380,-0.500);
\draw[gray!55](-0.616,0.029)--(-0.384,-0.029);
\draw[gray!55](-0.620,0.500)--(-0.380,0.500);
\draw[gray!55](-0.596,0.928)--(-0.404,1.072);
\draw[gray!55](-0.554,1.393)--(-0.446,1.607);
\draw[gray!55](-0.531,1.884)--(-0.469,2.116);
\draw[gray!55](-0.520,2.382)--(-0.480,2.618);
\draw[gray!55](-0.514,2.881)--(-0.486,3.119);
\draw[gray!55](-0.510,3.380)--(-0.490,3.620);
\draw[gray!55](-0.508,3.880)--(-0.492,4.120);
\draw[gray!55](-0.506,4.380)--(-0.494,4.620);
\draw[gray!55](-0.505,4.880)--(-0.495,5.120);
\draw[gray!55](-0.504,5.380)--(-0.496,5.620);
\draw[gray!55](-0.503,5.880)--(-0.497,6.120);
\draw[gray!55](-0.003,-6.120)--(0.003,-5.880);
\draw[gray!55](-0.004,-5.620)--(0.004,-5.380);
\draw[gray!55](-0.005,-5.120)--(0.005,-4.880);
\draw[gray!55](-0.006,-4.620)--(0.006,-4.380);
\draw[gray!55](-0.007,-4.120)--(0.007,-3.880);
\draw[gray!55](-0.010,-3.620)--(0.010,-3.380);
\draw[gray!55](-0.013,-3.119)--(0.013,-2.881);
\draw[gray!55](-0.019,-2.618)--(0.019,-2.382);
\draw[gray!55](-0.029,-2.116)--(0.029,-1.884);
\draw[gray!55](-0.049,-1.610)--(0.049,-1.390);
\draw[gray!55](-0.085,-1.085)--(0.085,-0.915);
\draw[gray!55](-0.116,-0.529)--(0.116,-0.471);
\draw[gray!55](-0.120,0.000)--(0.120,0.000);
\draw[gray!55](-0.116,0.471)--(0.116,0.529);
\draw[gray!55](-0.085,0.915)--(0.085,1.085);
\draw[gray!55](-0.049,1.390)--(0.049,1.610);
\draw[gray!55](-0.029,1.884)--(0.029,2.116);
\draw[gray!55](-0.019,2.382)--(0.019,2.618);
\draw[gray!55](-0.013,2.881)--(0.013,3.119);
\draw[gray!55](-0.010,3.380)--(0.010,3.620);
\draw[gray!55](-0.007,3.880)--(0.007,4.120);
\draw[gray!55](-0.006,4.380)--(0.006,4.620);
\draw[gray!55](-0.005,4.880)--(0.005,5.120);
\draw[gray!55](-0.004,5.380)--(0.004,5.620);
\draw[gray!55](-0.003,5.880)--(0.003,6.120);
\draw[gray!55](0.497,-6.120)--(0.503,-5.880);
\draw[gray!55](0.496,-5.620)--(0.504,-5.380);
\draw[gray!55](0.495,-5.120)--(0.505,-4.880);
\draw[gray!55](0.494,-4.620)--(0.506,-4.380);
\draw[gray!55](0.492,-4.120)--(0.508,-3.880);
\draw[gray!55](0.490,-3.620)--(0.510,-3.380);
\draw[gray!55](0.486,-3.119)--(0.514,-2.881);
\draw[gray!55](0.480,-2.618)--(0.520,-2.382);
\draw[gray!55](0.469,-2.116)--(0.531,-1.884);
\draw[gray!55](0.446,-1.607)--(0.554,-1.393);
\draw[gray!55](0.404,-1.072)--(0.596,-0.928);
\draw[gray!55](0.380,-0.500)--(0.620,-0.500);
\draw[gray!55](0.384,0.029)--(0.616,-0.029);
\draw[gray!55](0.380,0.500)--(0.620,0.500);
\draw[gray!55](0.404,0.928)--(0.596,1.072);
\draw[gray!55](0.446,1.393)--(0.554,1.607);
\draw[gray!55](0.469,1.884)--(0.531,2.116);
\draw[gray!55](0.480,2.382)--(0.520,2.618);
\draw[gray!55](0.486,2.881)--(0.514,3.119);
\draw[gray!55](0.490,3.380)--(0.510,3.620);
\draw[gray!55](0.492,3.880)--(0.508,4.120);
\draw[gray!55](0.494,4.380)--(0.506,4.620);
\draw[gray!55](0.495,4.880)--(0.505,5.120);
\draw[gray!55](0.496,5.380)--(0.504,5.620);
\draw[gray!55](0.497,5.880)--(0.503,6.120);
\draw[gray!55](0.997,-6.120)--(1.003,-5.880);
\draw[gray!55](0.996,-5.620)--(1.004,-5.380);
\draw[gray!55](0.995,-5.120)--(1.005,-4.880);
\draw[gray!55](0.994,-4.620)--(1.006,-4.380);
\draw[gray!55](0.992,-4.120)--(1.008,-3.880);
\draw[gray!55](0.989,-3.620)--(1.011,-3.380);
\draw[gray!55](0.985,-3.119)--(1.015,-2.881);
\draw[gray!55](0.978,-2.618)--(1.022,-2.382);
\draw[gray!55](0.962,-2.114)--(1.038,-1.886);
\draw[gray!55](0.925,-1.594)--(1.075,-1.406);
\draw[gray!55](0.880,-1.000)--(1.120,-1.000);
\draw[gray!55](0.904,-0.428)--(1.096,-0.572);
\draw[gray!55](0.915,0.085)--(1.085,-0.085);
\draw[gray!55](0.904,0.572)--(1.096,0.428);
\draw[gray!55](0.880,1.000)--(1.120,1.000);
\draw[gray!55](0.925,1.406)--(1.075,1.594);
\draw[gray!55](0.962,1.886)--(1.038,2.114);
\draw[gray!55](0.978,2.382)--(1.022,2.618);
\draw[gray!55](0.985,2.881)--(1.015,3.119);
\draw[gray!55](0.989,3.380)--(1.011,3.620);
\draw[gray!55](0.992,3.880)--(1.008,4.120);
\draw[gray!55](0.994,4.380)--(1.006,4.620);
\draw[gray!55](0.995,4.880)--(1.005,5.120);
\draw[gray!55](0.996,5.380)--(1.004,5.620);
\draw[gray!55](0.997,5.880)--(1.003,6.120);
\draw[gray!55](1.496,-6.120)--(1.504,-5.880);
\draw[gray!55](1.496,-5.620)--(1.504,-5.380);
\draw[gray!55](1.495,-5.120)--(1.505,-4.880);
\draw[gray!55](1.493,-4.620)--(1.507,-4.380);
\draw[gray!55](1.491,-4.120)--(1.509,-3.880);
\draw[gray!55](1.488,-3.619)--(1.512,-3.381);
\draw[gray!55](1.482,-3.119)--(1.518,-2.881);
\draw[gray!55](1.471,-2.616)--(1.529,-2.384);
\draw[gray!55](1.440,-2.104)--(1.560,-1.896);
\draw[gray!55](1.380,-1.500)--(1.620,-1.500);
\draw[gray!55](1.425,-0.906)--(1.575,-1.094);
\draw[gray!55](1.446,-0.393)--(1.554,-0.607);
\draw[gray!55](1.451,0.110)--(1.549,-0.110);
\draw[gray!55](1.446,0.607)--(1.554,0.393);
\draw[gray!55](1.425,1.094)--(1.575,0.906);
\draw[gray!55](1.380,1.500)--(1.620,1.500);
\draw[gray!55](1.440,1.896)--(1.560,2.104);
\draw[gray!55](1.471,2.384)--(1.529,2.616);
\draw[gray!55](1.482,2.881)--(1.518,3.119);
\draw[gray!55](1.488,3.381)--(1.512,3.619);
\draw[gray!55](1.491,3.880)--(1.509,4.120);
\draw[gray!55](1.493,4.380)--(1.507,4.620);
\draw[gray!55](1.495,4.880)--(1.505,5.120);
\draw[gray!55](1.496,5.380)--(1.504,5.620);
\draw[gray!55](1.496,5.880)--(1.504,6.120);
\draw[gray!55](1.996,-6.120)--(2.004,-5.880);
\draw[gray!55](1.995,-5.620)--(2.005,-5.380);
\draw[gray!55](1.994,-5.120)--(2.006,-4.880);
\draw[gray!55](1.993,-4.620)--(2.007,-4.380);
\draw[gray!55](1.990,-4.120)--(2.010,-3.880);
\draw[gray!55](1.986,-3.619)--(2.014,-3.381);
\draw[gray!55](1.976,-3.118)--(2.024,-2.882);
\draw[gray!55](1.951,-2.610)--(2.049,-2.390);
\draw[gray!55](1.880,-2.000)--(2.120,-2.000);
\draw[gray!55](1.940,-1.396)--(2.060,-1.604);
\draw[gray!55](1.962,-0.886)--(2.038,-1.114);
\draw[gray!55](1.969,-0.384)--(2.031,-0.616);
\draw[gray!55](1.971,0.116)--(2.029,-0.116);
\draw[gray!55](1.969,0.616)--(2.031,0.384);
\draw[gray!55](1.962,1.114)--(2.038,0.886);
\draw[gray!55](1.940,1.604)--(2.060,1.396);
\draw[gray!55](1.880,2.000)--(2.120,2.000);
\draw[gray!55](1.951,2.390)--(2.049,2.610);
\draw[gray!55](1.976,2.882)--(2.024,3.118);
\draw[gray!55](1.986,3.381)--(2.014,3.619);
\draw[gray!55](1.990,3.880)--(2.010,4.120);
\draw[gray!55](1.993,4.380)--(2.007,4.620);
\draw[gray!55](1.994,4.880)--(2.006,5.120);
\draw[gray!55](1.995,5.380)--(2.005,5.620);
\draw[gray!55](1.996,5.880)--(2.004,6.120);
\draw[gray!55](2.496,-6.120)--(2.504,-5.880);
\draw[gray!55](2.495,-5.620)--(2.505,-5.380);
\draw[gray!55](2.494,-5.120)--(2.506,-4.880);
\draw[gray!55](2.491,-4.620)--(2.509,-4.380);
\draw[gray!55](2.488,-4.119)--(2.512,-3.881);
\draw[gray!55](2.480,-3.618)--(2.520,-3.382);
\draw[gray!55](2.459,-3.113)--(2.541,-2.887);
\draw[gray!55](2.380,-2.500)--(2.620,-2.500);
\draw[gray!55](2.451,-1.890)--(2.549,-2.110);
\draw[gray!55](2.471,-1.384)--(2.529,-1.616);
\draw[gray!55](2.478,-0.882)--(2.522,-1.118);
\draw[gray!55](2.480,-0.382)--(2.520,-0.618);
\draw[gray!55](2.481,0.118)--(2.519,-0.118);
\draw[gray!55](2.480,0.618)--(2.520,0.382);
\draw[gray!55](2.478,1.118)--(2.522,0.882);
\draw[gray!55](2.471,1.616)--(2.529,1.384);
\draw[gray!55](2.451,2.110)--(2.549,1.890);
\draw[gray!55](2.380,2.500)--(2.620,2.500);
\draw[gray!55](2.459,2.887)--(2.541,3.113);
\draw[gray!55](2.480,3.382)--(2.520,3.618);
\draw[gray!55](2.488,3.881)--(2.512,4.119);
\draw[gray!55](2.491,4.380)--(2.509,4.620);
\draw[gray!55](2.494,4.880)--(2.506,5.120);
\draw[gray!55](2.495,5.380)--(2.505,5.620);
\draw[gray!55](2.496,5.880)--(2.504,6.120);
\draw[gray!55](2.996,-6.120)--(3.004,-5.880);
\draw[gray!55](2.994,-5.620)--(3.006,-5.380);
\draw[gray!55](2.993,-5.120)--(3.007,-4.880);
\draw[gray!55](2.989,-4.620)--(3.011,-4.380);
\draw[gray!55](2.983,-4.119)--(3.017,-3.881);
\draw[gray!55](2.965,-3.615)--(3.035,-3.385);
\draw[gray!55](2.880,-3.000)--(3.120,-3.000);
\draw[gray!55](2.959,-2.387)--(3.041,-2.613);
\draw[gray!55](2.976,-1.882)--(3.024,-2.118);
\draw[gray!55](2.982,-1.381)--(3.018,-1.619);
\draw[gray!55](2.985,-0.881)--(3.015,-1.119);
\draw[gray!55](2.986,-0.381)--(3.014,-0.619);
\draw[gray!55](2.987,0.119)--(3.013,-0.119);
\draw[gray!55](2.986,0.619)--(3.014,0.381);
\draw[gray!55](2.985,1.119)--(3.015,0.881);
\draw[gray!55](2.982,1.619)--(3.018,1.381);
\draw[gray!55](2.976,2.118)--(3.024,1.882);
\draw[gray!55](2.959,2.613)--(3.041,2.387);
\draw[gray!55](2.880,3.000)--(3.120,3.000);
\draw[gray!55](2.965,3.385)--(3.035,3.615);
\draw[gray!55](2.983,3.881)--(3.017,4.119);
\draw[gray!55](2.989,4.380)--(3.011,4.620);
\draw[gray!55](2.993,4.880)--(3.007,5.120);
\draw[gray!55](2.994,5.380)--(3.006,5.620);
\draw[gray!55](2.996,5.880)--(3.004,6.120);
\draw[gray!55](3.495,-6.120)--(3.505,-5.880);
\draw[gray!55](3.493,-5.620)--(3.507,-5.380);
\draw[gray!55](3.491,-5.120)--(3.509,-4.880);
\draw[gray!55](3.485,-4.619)--(3.515,-4.381);
\draw[gray!55](3.469,-4.116)--(3.531,-3.884);
\draw[gray!55](3.380,-3.500)--(3.620,-3.500);
\draw[gray!55](3.465,-2.885)--(3.535,-3.115);
\draw[gray!55](3.480,-2.382)--(3.520,-2.618);
\draw[gray!55](3.486,-1.881)--(3.514,-2.119);
\draw[gray!55](3.488,-1.381)--(3.512,-1.619);
\draw[gray!55](3.489,-0.880)--(3.511,-1.120);
\draw[gray!55](3.490,-0.380)--(3.510,-0.620);
\draw[gray!55](3.490,0.120)--(3.510,-0.120);
\draw[gray!55](3.490,0.620)--(3.510,0.380);
\draw[gray!55](3.489,1.120)--(3.511,0.880);
\draw[gray!55](3.488,1.619)--(3.512,1.381);
\draw[gray!55](3.486,2.119)--(3.514,1.881);
\draw[gray!55](3.480,2.618)--(3.520,2.382);
\draw[gray!55](3.465,3.115)--(3.535,2.885);
\draw[gray!55](3.380,3.500)--(3.620,3.500);
\draw[gray!55](3.469,3.884)--(3.531,4.116);
\draw[gray!55](3.485,4.381)--(3.515,4.619);
\draw[gray!55](3.491,4.880)--(3.509,5.120);
\draw[gray!55](3.493,5.380)--(3.507,5.620);
\draw[gray!55](3.495,5.880)--(3.505,6.120);
\draw[gray!55](3.994,-6.120)--(4.006,-5.880);
\draw[gray!55](3.992,-5.620)--(4.008,-5.380);
\draw[gray!55](3.987,-5.119)--(4.013,-4.881);
\draw[gray!55](3.973,-4.617)--(4.027,-4.383);
\draw[gray!55](3.880,-4.000)--(4.120,-4.000);
\draw[gray!55](3.969,-3.384)--(4.031,-3.616);
\draw[gray!55](3.983,-2.881)--(4.017,-3.119);
\draw[gray!55](3.988,-2.381)--(4.012,-2.619);
\draw[gray!55](3.990,-1.880)--(4.010,-2.120);
\draw[gray!55](3.991,-1.380)--(4.009,-1.620);
\draw[gray!55](3.992,-0.880)--(4.008,-1.120);
\draw[gray!55](3.992,-0.380)--(4.008,-0.620);
\draw[gray!55](3.993,0.120)--(4.007,-0.120);
\draw[gray!55](3.992,0.620)--(4.008,0.380);
\draw[gray!55](3.992,1.120)--(4.008,0.880);
\draw[gray!55](3.991,1.620)--(4.009,1.380);
\draw[gray!55](3.990,2.120)--(4.010,1.880);
\draw[gray!55](3.988,2.619)--(4.012,2.381);
\draw[gray!55](3.983,3.119)--(4.017,2.881);
\draw[gray!55](3.969,3.616)--(4.031,3.384);
\draw[gray!55](3.880,4.000)--(4.120,4.000);
\draw[gray!55](3.973,4.383)--(4.027,4.617);
\draw[gray!55](3.987,4.881)--(4.013,5.119);
\draw[gray!55](3.992,5.380)--(4.008,5.620);
\draw[gray!55](3.994,5.880)--(4.006,6.120);
\draw[gray!55](4.492,-6.120)--(4.508,-5.880);
\draw[gray!55](4.488,-5.619)--(4.512,-5.381);
\draw[gray!55](4.475,-5.117)--(4.525,-4.883);
\draw[gray!55](4.380,-4.500)--(4.620,-4.500);
\draw[gray!55](4.473,-3.883)--(4.527,-4.117);
\draw[gray!55](4.485,-3.381)--(4.515,-3.619);
\draw[gray!55](4.489,-2.880)--(4.511,-3.120);
\draw[gray!55](4.491,-2.380)--(4.509,-2.620);
\draw[gray!55](4.493,-1.880)--(4.507,-2.120);
\draw[gray!55](4.493,-1.380)--(4.507,-1.620);
\draw[gray!55](4.494,-0.880)--(4.506,-1.120);
\draw[gray!55](4.494,-0.380)--(4.506,-0.620);
\draw[gray!55](4.494,0.120)--(4.506,-0.120);
\draw[gray!55](4.494,0.620)--(4.506,0.380);
\draw[gray!55](4.494,1.120)--(4.506,0.880);
\draw[gray!55](4.493,1.620)--(4.507,1.380);
\draw[gray!55](4.493,2.120)--(4.507,1.880);
\draw[gray!55](4.491,2.620)--(4.509,2.380);
\draw[gray!55](4.489,3.120)--(4.511,2.880);
\draw[gray!55](4.485,3.619)--(4.515,3.381);
\draw[gray!55](4.473,4.117)--(4.527,3.883);
\draw[gray!55](4.380,4.500)--(4.620,4.500);
\draw[gray!55](4.475,4.883)--(4.525,5.117);
\draw[gray!55](4.488,5.381)--(4.512,5.619);
\draw[gray!55](4.492,5.880)--(4.508,6.120);
\draw[gray!55](4.989,-6.120)--(5.011,-5.880);
\draw[gray!55](4.978,-5.618)--(5.022,-5.382);
\draw[gray!55](4.880,-5.000)--(5.120,-5.000);
\draw[gray!55](4.975,-4.383)--(5.025,-4.617);
\draw[gray!55](4.987,-3.881)--(5.013,-4.119);
\draw[gray!55](4.991,-3.380)--(5.009,-3.620);
\draw[gray!55](4.993,-2.880)--(5.007,-3.120);
\draw[gray!55](4.994,-2.380)--(5.006,-2.620);
\draw[gray!55](4.994,-1.880)--(5.006,-2.120);
\draw[gray!55](4.995,-1.380)--(5.005,-1.620);
\draw[gray!55](4.995,-0.880)--(5.005,-1.120);
\draw[gray!55](4.995,-0.380)--(5.005,-0.620);
\draw[gray!55](4.995,0.120)--(5.005,-0.120);
\draw[gray!55](4.995,0.620)--(5.005,0.380);
\draw[gray!55](4.995,1.120)--(5.005,0.880);
\draw[gray!55](4.995,1.620)--(5.005,1.380);
\draw[gray!55](4.994,2.120)--(5.006,1.880);
\draw[gray!55](4.994,2.620)--(5.006,2.380);
\draw[gray!55](4.993,3.120)--(5.007,2.880);
\draw[gray!55](4.991,3.620)--(5.009,3.380);
\draw[gray!55](4.987,4.119)--(5.013,3.881);
\draw[gray!55](4.975,4.617)--(5.025,4.383);
\draw[gray!55](4.880,5.000)--(5.120,5.000);
\draw[gray!55](4.978,5.382)--(5.022,5.618);
\draw[gray!55](4.989,5.880)--(5.011,6.120);
\draw[gray!55](5.479,-6.118)--(5.521,-5.882);
\draw[gray!55](5.380,-5.500)--(5.620,-5.500);
\draw[gray!55](5.478,-4.882)--(5.522,-5.118);
\draw[gray!55](5.488,-4.381)--(5.512,-4.619);
\draw[gray!55](5.492,-3.880)--(5.508,-4.120);
\draw[gray!55](5.493,-3.380)--(5.507,-3.620);
\draw[gray!55](5.494,-2.880)--(5.506,-3.120);
\draw[gray!55](5.495,-2.380)--(5.505,-2.620);
\draw[gray!55](5.495,-1.880)--(5.505,-2.120);
\draw[gray!55](5.496,-1.380)--(5.504,-1.620);
\draw[gray!55](5.496,-0.880)--(5.504,-1.120);
\draw[gray!55](5.496,-0.380)--(5.504,-0.620);
\draw[gray!55](5.496,0.120)--(5.504,-0.120);
\draw[gray!55](5.496,0.620)--(5.504,0.380);
\draw[gray!55](5.496,1.120)--(5.504,0.880);
\draw[gray!55](5.496,1.620)--(5.504,1.380);
\draw[gray!55](5.495,2.120)--(5.505,1.880);
\draw[gray!55](5.495,2.620)--(5.505,2.380);
\draw[gray!55](5.494,3.120)--(5.506,2.880);
\draw[gray!55](5.493,3.620)--(5.507,3.380);
\draw[gray!55](5.492,4.120)--(5.508,3.880);
\draw[gray!55](5.488,4.619)--(5.512,4.381);
\draw[gray!55](5.478,5.118)--(5.522,4.882);
\draw[gray!55](5.380,5.500)--(5.620,5.500);
\draw[gray!55](5.479,5.882)--(5.521,6.118);
\draw[gray!55](5.880,-6.000)--(6.120,-6.000);
\draw[gray!55](5.979,-5.382)--(6.021,-5.618);
\draw[gray!55](5.989,-4.880)--(6.011,-5.120);
\draw[gray!55](5.992,-4.380)--(6.008,-4.620);
\draw[gray!55](5.994,-3.880)--(6.006,-4.120);
\draw[gray!55](5.995,-3.380)--(6.005,-3.620);
\draw[gray!55](5.996,-2.880)--(6.004,-3.120);
\draw[gray!55](5.996,-2.380)--(6.004,-2.620);
\draw[gray!55](5.996,-1.880)--(6.004,-2.120);
\draw[gray!55](5.996,-1.380)--(6.004,-1.620);
\draw[gray!55](5.997,-0.880)--(6.003,-1.120);
\draw[gray!55](5.997,-0.380)--(6.003,-0.620);
\draw[gray!55](5.997,0.120)--(6.003,-0.120);
\draw[gray!55](5.997,0.620)--(6.003,0.380);
\draw[gray!55](5.997,1.120)--(6.003,0.880);
\draw[gray!55](5.996,1.620)--(6.004,1.380);
\draw[gray!55](5.996,2.120)--(6.004,1.880);
\draw[gray!55](5.996,2.620)--(6.004,2.380);
\draw[gray!55](5.996,3.120)--(6.004,2.880);
\draw[gray!55](5.995,3.620)--(6.005,3.380);
\draw[gray!55](5.994,4.120)--(6.006,3.880);
\draw[gray!55](5.992,4.620)--(6.008,4.380);
\draw[gray!55](5.989,5.120)--(6.011,4.880);
\draw[gray!55](5.979,5.618)--(6.021,5.382);
\draw[gray!55](5.880,6.000)--(6.120,6.000);
\draw[main,thick] plot coordinates {(2.0000,0.0000) (1.9665,0.1315) (1.9331,0.2574) (1.8996,0.3769) (1.8661,0.4892) (1.8326,0.5939) (1.7992,0.6904) (1.7657,0.7786) (1.7322,0.8585) (1.6987,0.9302) (1.6653,0.9939) (1.6318,1.0498) (1.5983,1.0985) (1.5649,1.1402) (1.5314,1.1755) (1.4979,1.2049) (1.4644,1.2287) (1.4310,1.2475) (1.3975,1.2618) (1.3640,1.2718) (1.3305,1.2782) (1.2971,1.2811) (1.2636,1.2810) (1.2301,1.2782) (1.1967,1.2730) (1.1632,1.2657) (1.1297,1.2564) (1.0962,1.2455) (1.0628,1.2331) (1.0293,1.2194) (0.9958,1.2045) (0.9623,1.1887) (0.9289,1.1720) (0.8954,1.1545) (0.8619,1.1365) (0.8285,1.1178) (0.7950,1.0988) (0.7615,1.0794) (0.7280,1.0596) (0.6946,1.0397) (0.6611,1.0196) (0.6276,0.9994) (0.5941,0.9791) (0.5607,0.9588) (0.5272,0.9386) (0.4937,0.9185) (0.4603,0.8985) (0.4268,0.8786) (0.3933,0.8590) (0.3598,0.8396) (0.3264,0.8205) (0.2929,0.8017) (0.2594,0.7832) (0.2259,0.7651) (0.1925,0.7475) (0.1590,0.7302) (0.1255,0.7135) (0.0921,0.6972) (0.0586,0.6815) (0.0251,0.6664) (-0.0084,0.6518) (-0.0418,0.6379) (-0.0753,0.6247) (-0.1088,0.6122) (-0.1423,0.6004) (-0.1757,0.5894) (-0.2092,0.5792) (-0.2427,0.5699) (-0.2762,0.5615) (-0.3096,0.5539) (-0.3431,0.5473) (-0.3766,0.5418) (-0.4100,0.5372) (-0.4435,0.5337) (-0.4770,0.5313) (-0.5105,0.5300) (-0.5439,0.5300) (-0.5774,0.5311) (-0.6109,0.5334) (-0.6444,0.5370) (-0.6778,0.5419) (-0.7113,0.5481) (-0.7448,0.5557) (-0.7782,0.5646) (-0.8117,0.5749) (-0.8452,0.5866) (-0.8787,0.5997) (-0.9121,0.6142) (-0.9456,0.6301) (-0.9791,0.6475) (-1.0126,0.6662) (-1.0460,0.6864) (-1.0795,0.7079) (-1.1130,0.7308) (-1.1464,0.7551) (-1.1799,0.7806) (-1.2134,0.8075) (-1.2469,0.8355) (-1.2803,0.8648) (-1.3138,0.8952) (-1.3473,0.9267) (-1.3808,0.9592) (-1.4142,0.9927) (-1.4477,1.0271) (-1.4812,1.0624) (-1.5146,1.0984) (-1.5481,1.1352) (-1.5816,1.1726) (-1.6151,1.2106) (-1.6485,1.2491) (-1.6820,1.2880) (-1.7155,1.3274) (-1.7490,1.3671) (-1.7824,1.4070) (-1.8159,1.4472) (-1.8494,1.4875) (-1.8828,1.5280) (-1.9163,1.5686) (-1.9498,1.6091) (-1.9833,1.6497) (-2.0167,1.6902) (-2.0502,1.7307) (-2.0837,1.7711) (-2.1172,1.8114) (-2.1506,1.8515) (-2.1841,1.8915) (-2.2176,1.9313) (-2.2510,1.9710) (-2.2845,2.0105) (-2.3180,2.0498) (-2.3515,2.0889) (-2.3849,2.1278) (-2.4184,2.1666) (-2.4519,2.2051) (-2.4854,2.2435) (-2.5188,2.2817) (-2.5523,2.3197) (-2.5858,2.3576) (-2.6192,2.3952) (-2.6527,2.4328) (-2.6862,2.4701) (-2.7197,2.5073) (-2.7531,2.5444) (-2.7866,2.5814) (-2.8201,2.6182) (-2.8536,2.6549) (-2.8870,2.6915) (-2.9205,2.7279) (-2.9540,2.7643) (-2.9874,2.8005) (-3.0209,2.8367) (-3.0544,2.8728) (-3.0879,2.9088) (-3.1213,2.9447) (-3.1548,2.9805) (-3.1883,3.0163) (-3.2218,3.0520) (-3.2552,3.0876) (-3.2887,3.1231) (-3.3222,3.1587) (-3.3556,3.1941) (-3.3891,3.2295) (-3.4226,3.2648) (-3.4561,3.3001) (-3.4895,3.3354) (-3.5230,3.3706) (-3.5565,3.4057) (-3.5900,3.4408) (-3.6234,3.4759) (-3.6569,3.5109) (-3.6904,3.5459) (-3.7238,3.5809) (-3.7573,3.6158) (-3.7908,3.6507) (-3.8243,3.6856) (-3.8577,3.7205) (-3.8912,3.7553) (-3.9247,3.7900) (-3.9582,3.8248) (-3.9916,3.8595) (-4.0251,3.8942) (-4.0586,3.9289) (-4.0921,3.9636) (-4.1255,3.9982) (-4.1590,4.0328) (-4.1925,4.0674) (-4.2259,4.1020) (-4.2594,4.1365) (-4.2929,4.1710) (-4.3264,4.2056) (-4.3598,4.2400) (-4.3933,4.2745) (-4.4268,4.3090) (-4.4603,4.3434) (-4.4937,4.3778) (-4.5272,4.4122) (-4.5607,4.4466) (-4.5941,4.4810) (-4.6276,4.5154) (-4.6611,4.5497) (-4.6946,4.5840) (-4.7280,4.6184) (-4.7615,4.6527) (-4.7950,4.6870) (-4.8285,4.7212) (-4.8619,4.7555) (-4.8954,4.7898) (-4.9289,4.8240) (-4.9623,4.8582) (-4.9958,4.8924) (-5.0293,4.9267) (-5.0628,4.9609) (-5.0962,4.9950) (-5.1297,5.0292) (-5.1632,5.0634) (-5.1967,5.0975) (-5.2301,5.1317) (-5.2636,5.1658) (-5.2971,5.2000) (-5.3305,5.2341) (-5.3640,5.2682) (-5.3975,5.3023) (-5.4310,5.3364) (-5.4644,5.3705) (-5.4979,5.4045) (-5.5314,5.4386) (-5.5649,5.4727) (-5.5983,5.5067) (-5.6318,5.5408) (-5.6653,5.5748) (-5.6987,5.6088) (-5.7322,5.6429) (-5.7657,5.6769) (-5.7992,5.7109) (-5.8326,5.7449) (-5.8661,5.7789) (-5.8996,5.8129) (-5.9331,5.8469) (-5.9665,5.8809) (-6.0000,5.9148)};
\draw[main,thick] plot coordinates {(2.0000,0.0000) (2.0167,-0.0675) (2.0335,-0.1359) (2.0502,-0.2052) (2.0669,-0.2752) (2.0837,-0.3456) (2.1004,-0.4165) (2.1172,-0.4875) (2.1339,-0.5585) (2.1506,-0.6294) (2.1674,-0.7000) (2.1841,-0.7702) (2.2008,-0.8398) (2.2176,-0.9086) (2.2343,-0.9767) (2.2510,-1.0438) (2.2678,-1.1098) (2.2845,-1.1747) (2.3013,-1.2383) (2.3180,-1.3006) (2.3347,-1.3615) (2.3515,-1.4210) (2.3682,-1.4790) (2.3849,-1.5355) (2.4017,-1.5905) (2.4184,-1.6439) (2.4351,-1.6958) (2.4519,-1.7461) (2.4686,-1.7949) (2.4854,-1.8423) (2.5021,-1.8881) (2.5188,-1.9325) (2.5356,-1.9755) (2.5523,-2.0171) (2.5690,-2.0574) (2.5858,-2.0963) (2.6025,-2.1341) (2.6192,-2.1706) (2.6360,-2.2060) (2.6527,-2.2403) (2.6695,-2.2736) (2.6862,-2.3058) (2.7029,-2.3372) (2.7197,-2.3676) (2.7364,-2.3971) (2.7531,-2.4259) (2.7699,-2.4539) (2.7866,-2.4812) (2.8033,-2.5078) (2.8201,-2.5337) (2.8368,-2.5591) (2.8536,-2.5839) (2.8703,-2.6082) (2.8870,-2.6320) (2.9038,-2.6553) (2.9205,-2.6782) (2.9372,-2.7007) (2.9540,-2.7229) (2.9707,-2.7446) (2.9874,-2.7661) (3.0042,-2.7873) (3.0209,-2.8082) (3.0377,-2.8288) (3.0544,-2.8492) (3.0711,-2.8694) (3.0879,-2.8893) (3.1046,-2.9091) (3.1213,-2.9287) (3.1381,-2.9481) (3.1548,-2.9674) (3.1715,-2.9865) (3.1883,-3.0055) (3.2050,-3.0244) (3.2218,-3.0432) (3.2385,-3.0619) (3.2552,-3.0804) (3.2720,-3.0989) (3.2887,-3.1173) (3.3054,-3.1357) (3.3222,-3.1539) (3.3389,-3.1721) (3.3556,-3.1903) (3.3724,-3.2084) (3.3891,-3.2264) (3.4059,-3.2444) (3.4226,-3.2624) (3.4393,-3.2803) (3.4561,-3.2981) (3.4728,-3.3160) (3.4895,-3.3338) (3.5063,-3.3515) (3.5230,-3.3693) (3.5397,-3.3870) (3.5565,-3.4047) (3.5732,-3.4224) (3.5900,-3.4400) (3.6067,-3.4577) (3.6234,-3.4753) (3.6402,-3.4929) (3.6569,-3.5104) (3.6736,-3.5280) (3.6904,-3.5455) (3.7071,-3.5631) (3.7238,-3.5806) (3.7406,-3.5981) (3.7573,-3.6156) (3.7741,-3.6331) (3.7908,-3.6506) (3.8075,-3.6680) (3.8243,-3.6855) (3.8410,-3.7029) (3.8577,-3.7203) (3.8745,-3.7378) (3.8912,-3.7552) (3.9079,-3.7726) (3.9247,-3.7900) (3.9414,-3.8074) (3.9582,-3.8247) (3.9749,-3.8421) (3.9916,-3.8595) (4.0084,-3.8768) (4.0251,-3.8942) (4.0418,-3.9115) (4.0586,-3.9289) (4.0753,-3.9462) (4.0921,-3.9636) (4.1088,-3.9809) (4.1255,-3.9982) (4.1423,-4.0155) (4.1590,-4.0328) (4.1757,-4.0501) (4.1925,-4.0674) (4.2092,-4.0847) (4.2259,-4.1020) (4.2427,-4.1192) (4.2594,-4.1365) (4.2762,-4.1538) (4.2929,-4.1710) (4.3096,-4.1883) (4.3264,-4.2056) (4.3431,-4.2228) (4.3598,-4.2400) (4.3766,-4.2573) (4.3933,-4.2745) (4.4100,-4.2917) (4.4268,-4.3090) (4.4435,-4.3262) (4.4603,-4.3434) (4.4770,-4.3606) (4.4937,-4.3778) (4.5105,-4.3950) (4.5272,-4.4122) (4.5439,-4.4294) (4.5607,-4.4466) (4.5774,-4.4638) (4.5941,-4.4810) (4.6109,-4.4982) (4.6276,-4.5154) (4.6444,-4.5325) (4.6611,-4.5497) (4.6778,-4.5669) (4.6946,-4.5840) (4.7113,-4.6012) (4.7280,-4.6184) (4.7448,-4.6355) (4.7615,-4.6527) (4.7782,-4.6698) (4.7950,-4.6870) (4.8117,-4.7041) (4.8285,-4.7212) (4.8452,-4.7384) (4.8619,-4.7555) (4.8787,-4.7726) (4.8954,-4.7898) (4.9121,-4.8069) (4.9289,-4.8240) (4.9456,-4.8411) (4.9623,-4.8582) (4.9791,-4.8753) (4.9958,-4.8924) (5.0126,-4.9096) (5.0293,-4.9267) (5.0460,-4.9438) (5.0628,-4.9609) (5.0795,-4.9779) (5.0962,-4.9950) (5.1130,-5.0121) (5.1297,-5.0292) (5.1464,-5.0463) (5.1632,-5.0634) (5.1799,-5.0805) (5.1967,-5.0975) (5.2134,-5.1146) (5.2301,-5.1317) (5.2469,-5.1488) (5.2636,-5.1658) (5.2803,-5.1829) (5.2971,-5.2000) (5.3138,-5.2170) (5.3305,-5.2341) (5.3473,-5.2511) (5.3640,-5.2682) (5.3808,-5.2852) (5.3975,-5.3023) (5.4142,-5.3193) (5.4310,-5.3364) (5.4477,-5.3534) (5.4644,-5.3705) (5.4812,-5.3875) (5.4979,-5.4045) (5.5146,-5.4216) (5.5314,-5.4386) (5.5481,-5.4556) (5.5649,-5.4727) (5.5816,-5.4897) (5.5983,-5.5067) (5.6151,-5.5237) (5.6318,-5.5408) (5.6485,-5.5578) (5.6653,-5.5748) (5.6820,-5.5918) (5.6987,-5.6088) (5.7155,-5.6259) (5.7322,-5.6429) (5.7490,-5.6599) (5.7657,-5.6769) (5.7824,-5.6939) (5.7992,-5.7109) (5.8159,-5.7279) (5.8326,-5.7449) (5.8494,-5.7619) (5.8661,-5.7789) (5.8828,-5.7959) (5.8996,-5.8129) (5.9163,-5.8299) (5.9331,-5.8469) (5.9498,-5.8639) (5.9665,-5.8809) (5.9833,-5.8978) (6.0000,-5.9148)};
\fill[main](2,0)circle(2pt);
\draw[structurecolor,dashed]plot coordinates{(-5.9830,-5.6388) (-5.9452,-5.5987) (-5.9074,-5.5585) (-5.8696,-5.5184) (-5.8319,-5.4783) (-5.7942,-5.4381) (-5.7566,-5.3980) (-5.7190,-5.3579) (-5.6814,-5.3177) (-5.6438,-5.2776) (-5.6063,-5.2375) (-5.5689,-5.1973) (-5.5314,-5.1572) (-5.4940,-5.1171) (-5.4567,-5.0769) (-5.4193,-5.0368) (-5.3821,-4.9967) (-5.3448,-4.9565) (-5.3076,-4.9164) (-5.2705,-4.8763) (-5.2334,-4.8361) (-5.1963,-4.7960) (-5.1593,-4.7559) (-5.1223,-4.7157) (-5.0854,-4.6756) (-5.0485,-4.6355) (-5.0117,-4.5953) (-4.9749,-4.5552) (-4.9382,-4.5151) (-4.9015,-4.4749) (-4.8649,-4.4348) (-4.8283,-4.3946) (-4.7918,-4.3545) (-4.7554,-4.3144) (-4.7190,-4.2742) (-4.6827,-4.2341) (-4.6464,-4.1940) (-4.6103,-4.1538) (-4.5741,-4.1137) (-4.5381,-4.0736) (-4.5021,-4.0334) (-4.4662,-3.9933) (-4.4303,-3.9532) (-4.3945,-3.9130) (-4.3588,-3.8729) (-4.3232,-3.8328) (-4.2877,-3.7926) (-4.2522,-3.7525) (-4.2168,-3.7124) (-4.1815,-3.6722) (-4.1463,-3.6321) (-4.1112,-3.5920) (-4.0762,-3.5518) (-4.0413,-3.5117) (-4.0065,-3.4716) (-3.9717,-3.4314) (-3.9371,-3.3913) (-3.9026,-3.3512) (-3.8682,-3.3110) (-3.8339,-3.2709) (-3.7997,-3.2308) (-3.7657,-3.1906) (-3.7317,-3.1505) (-3.6979,-3.1104) (-3.6642,-3.0702) (-3.6306,-3.0301) (-3.5972,-2.9900) (-3.5639,-2.9498) (-3.5308,-2.9097) (-3.4978,-2.8696) (-3.4649,-2.8294) (-3.4322,-2.7893) (-3.3997,-2.7492) (-3.3673,-2.7090) (-3.3351,-2.6689) (-3.3031,-2.6288) (-3.2712,-2.5886) (-3.2396,-2.5485) (-3.2081,-2.5084) (-3.1768,-2.4682) (-3.1457,-2.4281) (-3.1149,-2.3880) (-3.0842,-2.3478) (-3.0538,-2.3077) (-3.0235,-2.2676) (-2.9936,-2.2274) (-2.9638,-2.1873) (-2.9343,-2.1472) (-2.9051,-2.1070) (-2.8761,-2.0669) (-2.8474,-2.0268) (-2.8190,-1.9866) (-2.7908,-1.9465) (-2.7630,-1.9064) (-2.7355,-1.8662) (-2.7082,-1.8261) (-2.6813,-1.7860) (-2.6548,-1.7458) (-2.6286,-1.7057) (-2.6027,-1.6656) (-2.5772,-1.6254) (-2.5521,-1.5853) (-2.5273,-1.5452) (-2.5030,-1.5050) (-2.4791,-1.4649) (-2.4556,-1.4247) (-2.4325,-1.3846) (-2.4099,-1.3445) (-2.3877,-1.3043) (-2.3661,-1.2642) (-2.3449,-1.2241) (-2.3242,-1.1839) (-2.3040,-1.1438) (-2.2843,-1.1037) (-2.2652,-1.0635) (-2.2466,-1.0234) (-2.2286,-0.9833) (-2.2112,-0.9431) (-2.1944,-0.9030) (-2.1782,-0.8629) (-2.1626,-0.8227) (-2.1477,-0.7826) (-2.1334,-0.7425) (-2.1197,-0.7023) (-2.1068,-0.6622) (-2.0945,-0.6221) (-2.0829,-0.5819) (-2.0721,-0.5418) (-2.0620,-0.5017) (-2.0526,-0.4615) (-2.0439,-0.4214) (-2.0360,-0.3813) (-2.0289,-0.3411) (-2.0225,-0.3010) (-2.0169,-0.2609) (-2.0121,-0.2207) (-2.0081,-0.1806) (-2.0049,-0.1405) (-2.0025,-0.1003) (-2.0009,-0.0602) (-2.0001,-0.0201) (-2.0001,0.0201) (-2.0009,0.0602) (-2.0025,0.1003) (-2.0049,0.1405) (-2.0081,0.1806) (-2.0121,0.2207) (-2.0169,0.2609) (-2.0225,0.3010) (-2.0289,0.3411) (-2.0360,0.3813) (-2.0439,0.4214) (-2.0526,0.4615) (-2.0620,0.5017) (-2.0721,0.5418) (-2.0829,0.5819) (-2.0945,0.6221) (-2.1068,0.6622) (-2.1197,0.7023) (-2.1334,0.7425) (-2.1477,0.7826) (-2.1626,0.8227) (-2.1782,0.8629) (-2.1944,0.9030) (-2.2112,0.9431) (-2.2286,0.9833) (-2.2466,1.0234) (-2.2652,1.0635) (-2.2843,1.1037) (-2.3040,1.1438) (-2.3242,1.1839) (-2.3449,1.2241) (-2.3661,1.2642) (-2.3877,1.3043) (-2.4099,1.3445) (-2.4325,1.3846) (-2.4556,1.4247) (-2.4791,1.4649) (-2.5030,1.5050) (-2.5273,1.5452) (-2.5521,1.5853) (-2.5772,1.6254) (-2.6027,1.6656) (-2.6286,1.7057) (-2.6548,1.7458) (-2.6813,1.7860) (-2.7082,1.8261) (-2.7355,1.8662) (-2.7630,1.9064) (-2.7908,1.9465) (-2.8190,1.9866) (-2.8474,2.0268) (-2.8761,2.0669) (-2.9051,2.1070) (-2.9343,2.1472) (-2.9638,2.1873) (-2.9936,2.2274) (-3.0235,2.2676) (-3.0538,2.3077) (-3.0842,2.3478) (-3.1149,2.3880) (-3.1457,2.4281) (-3.1768,2.4682) (-3.2081,2.5084) (-3.2396,2.5485) (-3.2712,2.5886) (-3.3031,2.6288) (-3.3351,2.6689) (-3.3673,2.7090) (-3.3997,2.7492) (-3.4322,2.7893) (-3.4649,2.8294) (-3.4978,2.8696) (-3.5308,2.9097) (-3.5639,2.9498) (-3.5972,2.9900) (-3.6306,3.0301) (-3.6642,3.0702) (-3.6979,3.1104) (-3.7317,3.1505) (-3.7657,3.1906) (-3.7997,3.2308) (-3.8339,3.2709) (-3.8682,3.3110) (-3.9026,3.3512) (-3.9371,3.3913) (-3.9717,3.4314) (-4.0065,3.4716) (-4.0413,3.5117) (-4.0762,3.5518) (-4.1112,3.5920) (-4.1463,3.6321) (-4.1815,3.6722) (-4.2168,3.7124) (-4.2522,3.7525) (-4.2877,3.7926) (-4.3232,3.8328) (-4.3588,3.8729) (-4.3945,3.9130) (-4.4303,3.9532) (-4.4662,3.9933) (-4.5021,4.0334) (-4.5381,4.0736) (-4.5741,4.1137) (-4.6103,4.1538) (-4.6464,4.1940) (-4.6827,4.2341) (-4.7190,4.2742) (-4.7554,4.3144) (-4.7918,4.3545) (-4.8283,4.3946) (-4.8649,4.4348) (-4.9015,4.4749) (-4.9382,4.5151) (-4.9749,4.5552) (-5.0117,4.5953) (-5.0485,4.6355) (-5.0854,4.6756) (-5.1223,4.7157) (-5.1593,4.7559) (-5.1963,4.7960) (-5.2334,4.8361) (-5.2705,4.8763) (-5.3076,4.9164) (-5.3448,4.9565) (-5.3821,4.9967) (-5.4193,5.0368) (-5.4567,5.0769) (-5.4940,5.1171) (-5.5314,5.1572) (-5.5689,5.1973) (-5.6063,5.2375) (-5.6438,5.2776) (-5.6814,5.3177) (-5.7190,5.3579) (-5.7566,5.3980) (-5.7942,5.4381) (-5.8319,5.4783) (-5.8696,5.5184) (-5.9074,5.5585) (-5.9452,5.5987) (-5.9830,5.6388)};\draw[structurecolor,dashed]plot coordinates{(5.9830,-5.6388) (5.9452,-5.5987) (5.9074,-5.5585) (5.8696,-5.5184) (5.8319,-5.4783) (5.7942,-5.4381) (5.7566,-5.3980) (5.7190,-5.3579) (5.6814,-5.3177) (5.6438,-5.2776) (5.6063,-5.2375) (5.5689,-5.1973) (5.5314,-5.1572) (5.4940,-5.1171) (5.4567,-5.0769) (5.4193,-5.0368) (5.3821,-4.9967) (5.3448,-4.9565) (5.3076,-4.9164) (5.2705,-4.8763) (5.2334,-4.8361) (5.1963,-4.7960) (5.1593,-4.7559) (5.1223,-4.7157) (5.0854,-4.6756) (5.0485,-4.6355) (5.0117,-4.5953) (4.9749,-4.5552) (4.9382,-4.5151) (4.9015,-4.4749) (4.8649,-4.4348) (4.8283,-4.3946) (4.7918,-4.3545) (4.7554,-4.3144) (4.7190,-4.2742) (4.6827,-4.2341) (4.6464,-4.1940) (4.6103,-4.1538) (4.5741,-4.1137) (4.5381,-4.0736) (4.5021,-4.0334) (4.4662,-3.9933) (4.4303,-3.9532) (4.3945,-3.9130) (4.3588,-3.8729) (4.3232,-3.8328) (4.2877,-3.7926) (4.2522,-3.7525) (4.2168,-3.7124) (4.1815,-3.6722) (4.1463,-3.6321) (4.1112,-3.5920) (4.0762,-3.5518) (4.0413,-3.5117) (4.0065,-3.4716) (3.9717,-3.4314) (3.9371,-3.3913) (3.9026,-3.3512) (3.8682,-3.3110) (3.8339,-3.2709) (3.7997,-3.2308) (3.7657,-3.1906) (3.7317,-3.1505) (3.6979,-3.1104) (3.6642,-3.0702) (3.6306,-3.0301) (3.5972,-2.9900) (3.5639,-2.9498) (3.5308,-2.9097) (3.4978,-2.8696) (3.4649,-2.8294) (3.4322,-2.7893) (3.3997,-2.7492) (3.3673,-2.7090) (3.3351,-2.6689) (3.3031,-2.6288) (3.2712,-2.5886) (3.2396,-2.5485) (3.2081,-2.5084) (3.1768,-2.4682) (3.1457,-2.4281) (3.1149,-2.3880) (3.0842,-2.3478) (3.0538,-2.3077) (3.0235,-2.2676) (2.9936,-2.2274) (2.9638,-2.1873) (2.9343,-2.1472) (2.9051,-2.1070) (2.8761,-2.0669) (2.8474,-2.0268) (2.8190,-1.9866) (2.7908,-1.9465) (2.7630,-1.9064) (2.7355,-1.8662) (2.7082,-1.8261) (2.6813,-1.7860) (2.6548,-1.7458) (2.6286,-1.7057) (2.6027,-1.6656) (2.5772,-1.6254) (2.5521,-1.5853) (2.5273,-1.5452) (2.5030,-1.5050) (2.4791,-1.4649) (2.4556,-1.4247) (2.4325,-1.3846) (2.4099,-1.3445) (2.3877,-1.3043) (2.3661,-1.2642) (2.3449,-1.2241) (2.3242,-1.1839) (2.3040,-1.1438) (2.2843,-1.1037) (2.2652,-1.0635) (2.2466,-1.0234) (2.2286,-0.9833) (2.2112,-0.9431) (2.1944,-0.9030) (2.1782,-0.8629) (2.1626,-0.8227) (2.1477,-0.7826) (2.1334,-0.7425) (2.1197,-0.7023) (2.1068,-0.6622) (2.0945,-0.6221) (2.0829,-0.5819) (2.0721,-0.5418) (2.0620,-0.5017) (2.0526,-0.4615) (2.0439,-0.4214) (2.0360,-0.3813) (2.0289,-0.3411) (2.0225,-0.3010) (2.0169,-0.2609) (2.0121,-0.2207) (2.0081,-0.1806) (2.0049,-0.1405) (2.0025,-0.1003) (2.0009,-0.0602) (2.0001,-0.0201) (2.0001,0.0201) (2.0009,0.0602) (2.0025,0.1003) (2.0049,0.1405) (2.0081,0.1806) (2.0121,0.2207) (2.0169,0.2609) (2.0225,0.3010) (2.0289,0.3411) (2.0360,0.3813) (2.0439,0.4214) (2.0526,0.4615) (2.0620,0.5017) (2.0721,0.5418) (2.0829,0.5819) (2.0945,0.6221) (2.1068,0.6622) (2.1197,0.7023) (2.1334,0.7425) (2.1477,0.7826) (2.1626,0.8227) (2.1782,0.8629) (2.1944,0.9030) (2.2112,0.9431) (2.2286,0.9833) (2.2466,1.0234) (2.2652,1.0635) (2.2843,1.1037) (2.3040,1.1438) (2.3242,1.1839) (2.3449,1.2241) (2.3661,1.2642) (2.3877,1.3043) (2.4099,1.3445) (2.4325,1.3846) (2.4556,1.4247) (2.4791,1.4649) (2.5030,1.5050) (2.5273,1.5452) (2.5521,1.5853) (2.5772,1.6254) (2.6027,1.6656) (2.6286,1.7057) (2.6548,1.7458) (2.6813,1.7860) (2.7082,1.8261) (2.7355,1.8662) (2.7630,1.9064) (2.7908,1.9465) (2.8190,1.9866) (2.8474,2.0268) (2.8761,2.0669) (2.9051,2.1070) (2.9343,2.1472) (2.9638,2.1873) (2.9936,2.2274) (3.0235,2.2676) (3.0538,2.3077) (3.0842,2.3478) (3.1149,2.3880) (3.1457,2.4281) (3.1768,2.4682) (3.2081,2.5084) (3.2396,2.5485) (3.2712,2.5886) (3.3031,2.6288) (3.3351,2.6689) (3.3673,2.7090) (3.3997,2.7492) (3.4322,2.7893) (3.4649,2.8294) (3.4978,2.8696) (3.5308,2.9097) (3.5639,2.9498) (3.5972,2.9900) (3.6306,3.0301) (3.6642,3.0702) (3.6979,3.1104) (3.7317,3.1505) (3.7657,3.1906) (3.7997,3.2308) (3.8339,3.2709) (3.8682,3.3110) (3.9026,3.3512) (3.9371,3.3913) (3.9717,3.4314) (4.0065,3.4716) (4.0413,3.5117) (4.0762,3.5518) (4.1112,3.5920) (4.1463,3.6321) (4.1815,3.6722) (4.2168,3.7124) (4.2522,3.7525) (4.2877,3.7926) (4.3232,3.8328) (4.3588,3.8729) (4.3945,3.9130) (4.4303,3.9532) (4.4662,3.9933) (4.5021,4.0334) (4.5381,4.0736) (4.5741,4.1137) (4.6103,4.1538) (4.6464,4.1940) (4.6827,4.2341) (4.7190,4.2742) (4.7554,4.3144) (4.7918,4.3545) (4.8283,4.3946) (4.8649,4.4348) (4.9015,4.4749) (4.9382,4.5151) (4.9749,4.5552) (5.0117,4.5953) (5.0485,4.6355) (5.0854,4.6756) (5.1223,4.7157) (5.1593,4.7559) (5.1963,4.7960) (5.2334,4.8361) (5.2705,4.8763) (5.3076,4.9164) (5.3448,4.9565) (5.3821,4.9967) (5.4193,5.0368) (5.4567,5.0769) (5.4940,5.1171) (5.5314,5.1572) (5.5689,5.1973) (5.6063,5.2375) (5.6438,5.2776) (5.6814,5.3177) (5.7190,5.3579) (5.7566,5.3980) (5.7942,5.4381) (5.8319,5.4783) (5.8696,5.5184) (5.9074,5.5585) (5.9452,5.5987) (5.9830,5.6388)};\draw[structurecolor]plot coordinates{(-6.0000,-6.0000) (-5.9599,-5.9599) (-5.9197,-5.9197) (-5.8796,-5.8796) (-5.8395,-5.8395) (-5.7993,-5.7993) (-5.7592,-5.7592) (-5.7191,-5.7191) (-5.6789,-5.6789) (-5.6388,-5.6388) (-5.5987,-5.5987) (-5.5585,-5.5585) (-5.5184,-5.5184) (-5.4783,-5.4783) (-5.4381,-5.4381) (-5.3980,-5.3980) (-5.3579,-5.3579) (-5.3177,-5.3177) (-5.2776,-5.2776) (-5.2375,-5.2375) (-5.1973,-5.1973) (-5.1572,-5.1572) (-5.1171,-5.1171) (-5.0769,-5.0769) (-5.0368,-5.0368) (-4.9967,-4.9967) (-4.9565,-4.9565) (-4.9164,-4.9164) (-4.8763,-4.8763) (-4.8361,-4.8361) (-4.7960,-4.7960) (-4.7559,-4.7559) (-4.7157,-4.7157) (-4.6756,-4.6756) (-4.6355,-4.6355) (-4.5953,-4.5953) (-4.5552,-4.5552) (-4.5151,-4.5151) (-4.4749,-4.4749) (-4.4348,-4.4348) (-4.3946,-4.3946) (-4.3545,-4.3545) (-4.3144,-4.3144) (-4.2742,-4.2742) (-4.2341,-4.2341) (-4.1940,-4.1940) (-4.1538,-4.1538) (-4.1137,-4.1137) (-4.0736,-4.0736) (-4.0334,-4.0334) (-3.9933,-3.9933) (-3.9532,-3.9532) (-3.9130,-3.9130) (-3.8729,-3.8729) (-3.8328,-3.8328) (-3.7926,-3.7926) (-3.7525,-3.7525) (-3.7124,-3.7124) (-3.6722,-3.6722) (-3.6321,-3.6321) (-3.5920,-3.5920) (-3.5518,-3.5518) (-3.5117,-3.5117) (-3.4716,-3.4716) (-3.4314,-3.4314) (-3.3913,-3.3913) (-3.3512,-3.3512) (-3.3110,-3.3110) (-3.2709,-3.2709) (-3.2308,-3.2308) (-3.1906,-3.1906) (-3.1505,-3.1505) (-3.1104,-3.1104) (-3.0702,-3.0702) (-3.0301,-3.0301) (-2.9900,-2.9900) (-2.9498,-2.9498) (-2.9097,-2.9097) (-2.8696,-2.8696) (-2.8294,-2.8294) (-2.7893,-2.7893) (-2.7492,-2.7492) (-2.7090,-2.7090) (-2.6689,-2.6689) (-2.6288,-2.6288) (-2.5886,-2.5886) (-2.5485,-2.5485) (-2.5084,-2.5084) (-2.4682,-2.4682) (-2.4281,-2.4281) (-2.3880,-2.3880) (-2.3478,-2.3478) (-2.3077,-2.3077) (-2.2676,-2.2676) (-2.2274,-2.2274) (-2.1873,-2.1873) (-2.1472,-2.1472) (-2.1070,-2.1070) (-2.0669,-2.0669) (-2.0268,-2.0268) (-1.9866,-1.9866) (-1.9465,-1.9465) (-1.9064,-1.9064) (-1.8662,-1.8662) (-1.8261,-1.8261) (-1.7860,-1.7860) (-1.7458,-1.7458) (-1.7057,-1.7057) (-1.6656,-1.6656) (-1.6254,-1.6254) (-1.5853,-1.5853) (-1.5452,-1.5452) (-1.5050,-1.5050) (-1.4649,-1.4649) (-1.4247,-1.4247) (-1.3846,-1.3846) (-1.3445,-1.3445) (-1.3043,-1.3043) (-1.2642,-1.2642) (-1.2241,-1.2241) (-1.1839,-1.1839) (-1.1438,-1.1438) (-1.1037,-1.1037) (-1.0635,-1.0635) (-1.0234,-1.0234) (-0.9833,-0.9833) (-0.9431,-0.9431) (-0.9030,-0.9030) (-0.8629,-0.8629) (-0.8227,-0.8227) (-0.7826,-0.7826) (-0.7425,-0.7425) (-0.7023,-0.7023) (-0.6622,-0.6622) (-0.6221,-0.6221) (-0.5819,-0.5819) (-0.5418,-0.5418) (-0.5017,-0.5017) (-0.4615,-0.4615) (-0.4214,-0.4214) (-0.3813,-0.3813) (-0.3411,-0.3411) (-0.3010,-0.3010) (-0.2609,-0.2609) (-0.2207,-0.2207) (-0.1806,-0.1806) (-0.1405,-0.1405) (-0.1003,-0.1003) (-0.0602,-0.0602) (-0.0201,-0.0201) (-0.0201,0.0201) (-0.0602,0.0602) (-0.1003,0.1003) (-0.1405,0.1405) (-0.1806,0.1806) (-0.2207,0.2207) (-0.2609,0.2609) (-0.3010,0.3010) (-0.3411,0.3411) (-0.3813,0.3813) (-0.4214,0.4214) (-0.4615,0.4615) (-0.5017,0.5017) (-0.5418,0.5418) (-0.5819,0.5819) (-0.6221,0.6221) (-0.6622,0.6622) (-0.7023,0.7023) (-0.7425,0.7425) (-0.7826,0.7826) (-0.8227,0.8227) (-0.8629,0.8629) (-0.9030,0.9030) (-0.9431,0.9431) (-0.9833,0.9833) (-1.0234,1.0234) (-1.0635,1.0635) (-1.1037,1.1037) (-1.1438,1.1438) (-1.1839,1.1839) (-1.2241,1.2241) (-1.2642,1.2642) (-1.3043,1.3043) (-1.3445,1.3445) (-1.3846,1.3846) (-1.4247,1.4247) (-1.4649,1.4649) (-1.5050,1.5050) (-1.5452,1.5452) (-1.5853,1.5853) (-1.6254,1.6254) (-1.6656,1.6656) (-1.7057,1.7057) (-1.7458,1.7458) (-1.7860,1.7860) (-1.8261,1.8261) (-1.8662,1.8662) (-1.9064,1.9064) (-1.9465,1.9465) (-1.9866,1.9866) (-2.0268,2.0268) (-2.0669,2.0669) (-2.1070,2.1070) (-2.1472,2.1472) (-2.1873,2.1873) (-2.2274,2.2274) (-2.2676,2.2676) (-2.3077,2.3077) (-2.3478,2.3478) (-2.3880,2.3880) (-2.4281,2.4281) (-2.4682,2.4682) (-2.5084,2.5084) (-2.5485,2.5485) (-2.5886,2.5886) (-2.6288,2.6288) (-2.6689,2.6689) (-2.7090,2.7090) (-2.7492,2.7492) (-2.7893,2.7893) (-2.8294,2.8294) (-2.8696,2.8696) (-2.9097,2.9097) (-2.9498,2.9498) (-2.9900,2.9900) (-3.0301,3.0301) (-3.0702,3.0702) (-3.1104,3.1104) (-3.1505,3.1505) (-3.1906,3.1906) (-3.2308,3.2308) (-3.2709,3.2709) (-3.3110,3.3110) (-3.3512,3.3512) (-3.3913,3.3913) (-3.4314,3.4314) (-3.4716,3.4716) (-3.5117,3.5117) (-3.5518,3.5518) (-3.5920,3.5920) (-3.6321,3.6321) (-3.6722,3.6722) (-3.7124,3.7124) (-3.7525,3.7525) (-3.7926,3.7926) (-3.8328,3.8328) (-3.8729,3.8729) (-3.9130,3.9130) (-3.9532,3.9532) (-3.9933,3.9933) (-4.0334,4.0334) (-4.0736,4.0736) (-4.1137,4.1137) (-4.1538,4.1538) (-4.1940,4.1940) (-4.2341,4.2341) (-4.2742,4.2742) (-4.3144,4.3144) (-4.3545,4.3545) (-4.3946,4.3946) (-4.4348,4.4348) (-4.4749,4.4749) (-4.5151,4.5151) (-4.5552,4.5552) (-4.5953,4.5953) (-4.6355,4.6355) (-4.6756,4.6756) (-4.7157,4.7157) (-4.7559,4.7559) (-4.7960,4.7960) (-4.8361,4.8361) (-4.8763,4.8763) (-4.9164,4.9164) (-4.9565,4.9565) (-4.9967,4.9967) (-5.0368,5.0368) (-5.0769,5.0769) (-5.1171,5.1171) (-5.1572,5.1572) (-5.1973,5.1973) (-5.2375,5.2375) (-5.2776,5.2776) (-5.3177,5.3177) (-5.3579,5.3579) (-5.3980,5.3980) (-5.4381,5.4381) (-5.4783,5.4783) (-5.5184,5.5184) (-5.5585,5.5585) (-5.5987,5.5987) (-5.6388,5.6388) (-5.6789,5.6789) (-5.7191,5.7191) (-5.7592,5.7592) (-5.7993,5.7993) (-5.8395,5.8395) (-5.8796,5.8796) (-5.9197,5.9197) (-5.9599,5.9599) (-6.0000,6.0000)};\draw[structurecolor]plot coordinates{(6.0000,-6.0000) (5.9599,-5.9599) (5.9197,-5.9197) (5.8796,-5.8796) (5.8395,-5.8395) (5.7993,-5.7993) (5.7592,-5.7592) (5.7191,-5.7191) (5.6789,-5.6789) (5.6388,-5.6388) (5.5987,-5.5987) (5.5585,-5.5585) (5.5184,-5.5184) (5.4783,-5.4783) (5.4381,-5.4381) (5.3980,-5.3980) (5.3579,-5.3579) (5.3177,-5.3177) (5.2776,-5.2776) (5.2375,-5.2375) (5.1973,-5.1973) (5.1572,-5.1572) (5.1171,-5.1171) (5.0769,-5.0769) (5.0368,-5.0368) (4.9967,-4.9967) (4.9565,-4.9565) (4.9164,-4.9164) (4.8763,-4.8763) (4.8361,-4.8361) (4.7960,-4.7960) (4.7559,-4.7559) (4.7157,-4.7157) (4.6756,-4.6756) (4.6355,-4.6355) (4.5953,-4.5953) (4.5552,-4.5552) (4.5151,-4.5151) (4.4749,-4.4749) (4.4348,-4.4348) (4.3946,-4.3946) (4.3545,-4.3545) (4.3144,-4.3144) (4.2742,-4.2742) (4.2341,-4.2341) (4.1940,-4.1940) (4.1538,-4.1538) (4.1137,-4.1137) (4.0736,-4.0736) (4.0334,-4.0334) (3.9933,-3.9933) (3.9532,-3.9532) (3.9130,-3.9130) (3.8729,-3.8729) (3.8328,-3.8328) (3.7926,-3.7926) (3.7525,-3.7525) (3.7124,-3.7124) (3.6722,-3.6722) (3.6321,-3.6321) (3.5920,-3.5920) (3.5518,-3.5518) (3.5117,-3.5117) (3.4716,-3.4716) (3.4314,-3.4314) (3.3913,-3.3913) (3.3512,-3.3512) (3.3110,-3.3110) (3.2709,-3.2709) (3.2308,-3.2308) (3.1906,-3.1906) (3.1505,-3.1505) (3.1104,-3.1104) (3.0702,-3.0702) (3.0301,-3.0301) (2.9900,-2.9900) (2.9498,-2.9498) (2.9097,-2.9097) (2.8696,-2.8696) (2.8294,-2.8294) (2.7893,-2.7893) (2.7492,-2.7492) (2.7090,-2.7090) (2.6689,-2.6689) (2.6288,-2.6288) (2.5886,-2.5886) (2.5485,-2.5485) (2.5084,-2.5084) (2.4682,-2.4682) (2.4281,-2.4281) (2.3880,-2.3880) (2.3478,-2.3478) (2.3077,-2.3077) (2.2676,-2.2676) (2.2274,-2.2274) (2.1873,-2.1873) (2.1472,-2.1472) (2.1070,-2.1070) (2.0669,-2.0669) (2.0268,-2.0268) (1.9866,-1.9866) (1.9465,-1.9465) (1.9064,-1.9064) (1.8662,-1.8662) (1.8261,-1.8261) (1.7860,-1.7860) (1.7458,-1.7458) (1.7057,-1.7057) (1.6656,-1.6656) (1.6254,-1.6254) (1.5853,-1.5853) (1.5452,-1.5452) (1.5050,-1.5050) (1.4649,-1.4649) (1.4247,-1.4247) (1.3846,-1.3846) (1.3445,-1.3445) (1.3043,-1.3043) (1.2642,-1.2642) (1.2241,-1.2241) (1.1839,-1.1839) (1.1438,-1.1438) (1.1037,-1.1037) (1.0635,-1.0635) (1.0234,-1.0234) (0.9833,-0.9833) (0.9431,-0.9431) (0.9030,-0.9030) (0.8629,-0.8629) (0.8227,-0.8227) (0.7826,-0.7826) (0.7425,-0.7425) (0.7023,-0.7023) (0.6622,-0.6622) (0.6221,-0.6221) (0.5819,-0.5819) (0.5418,-0.5418) (0.5017,-0.5017) (0.4615,-0.4615) (0.4214,-0.4214) (0.3813,-0.3813) (0.3411,-0.3411) (0.3010,-0.3010) (0.2609,-0.2609) (0.2207,-0.2207) (0.1806,-0.1806) (0.1405,-0.1405) (0.1003,-0.1003) (0.0602,-0.0602) (0.0201,-0.0201) (0.0201,0.0201) (0.0602,0.0602) (0.1003,0.1003) (0.1405,0.1405) (0.1806,0.1806) (0.2207,0.2207) (0.2609,0.2609) (0.3010,0.3010) (0.3411,0.3411) (0.3813,0.3813) (0.4214,0.4214) (0.4615,0.4615) (0.5017,0.5017) (0.5418,0.5418) (0.5819,0.5819) (0.6221,0.6221) (0.6622,0.6622) (0.7023,0.7023) (0.7425,0.7425) (0.7826,0.7826) (0.8227,0.8227) (0.8629,0.8629) (0.9030,0.9030) (0.9431,0.9431) (0.9833,0.9833) (1.0234,1.0234) (1.0635,1.0635) (1.1037,1.1037) (1.1438,1.1438) (1.1839,1.1839) (1.2241,1.2241) (1.2642,1.2642) (1.3043,1.3043) (1.3445,1.3445) (1.3846,1.3846) (1.4247,1.4247) (1.4649,1.4649) (1.5050,1.5050) (1.5452,1.5452) (1.5853,1.5853) (1.6254,1.6254) (1.6656,1.6656) (1.7057,1.7057) (1.7458,1.7458) (1.7860,1.7860) (1.8261,1.8261) (1.8662,1.8662) (1.9064,1.9064) (1.9465,1.9465) (1.9866,1.9866) (2.0268,2.0268) (2.0669,2.0669) (2.1070,2.1070) (2.1472,2.1472) (2.1873,2.1873) (2.2274,2.2274) (2.2676,2.2676) (2.3077,2.3077) (2.3478,2.3478) (2.3880,2.3880) (2.4281,2.4281) (2.4682,2.4682) (2.5084,2.5084) (2.5485,2.5485) (2.5886,2.5886) (2.6288,2.6288) (2.6689,2.6689) (2.7090,2.7090) (2.7492,2.7492) (2.7893,2.7893) (2.8294,2.8294) (2.8696,2.8696) (2.9097,2.9097) (2.9498,2.9498) (2.9900,2.9900) (3.0301,3.0301) (3.0702,3.0702) (3.1104,3.1104) (3.1505,3.1505) (3.1906,3.1906) (3.2308,3.2308) (3.2709,3.2709) (3.3110,3.3110) (3.3512,3.3512) (3.3913,3.3913) (3.4314,3.4314) (3.4716,3.4716) (3.5117,3.5117) (3.5518,3.5518) (3.5920,3.5920) (3.6321,3.6321) (3.6722,3.6722) (3.7124,3.7124) (3.7525,3.7525) (3.7926,3.7926) (3.8328,3.8328) (3.8729,3.8729) (3.9130,3.9130) (3.9532,3.9532) (3.9933,3.9933) (4.0334,4.0334) (4.0736,4.0736) (4.1137,4.1137) (4.1538,4.1538) (4.1940,4.1940) (4.2341,4.2341) (4.2742,4.2742) (4.3144,4.3144) (4.3545,4.3545) (4.3946,4.3946) (4.4348,4.4348) (4.4749,4.4749) (4.5151,4.5151) (4.5552,4.5552) (4.5953,4.5953) (4.6355,4.6355) (4.6756,4.6756) (4.7157,4.7157) (4.7559,4.7559) (4.7960,4.7960) (4.8361,4.8361) (4.8763,4.8763) (4.9164,4.9164) (4.9565,4.9565) (4.9967,4.9967) (5.0368,5.0368) (5.0769,5.0769) (5.1171,5.1171) (5.1572,5.1572) (5.1973,5.1973) (5.2375,5.2375) (5.2776,5.2776) (5.3177,5.3177) (5.3579,5.3579) (5.3980,5.3980) (5.4381,5.4381) (5.4783,5.4783) (5.5184,5.5184) (5.5585,5.5585) (5.5987,5.5987) (5.6388,5.6388) (5.6789,5.6789) (5.7191,5.7191) (5.7592,5.7592) (5.7993,5.7993) (5.8395,5.8395) (5.8796,5.8796) (5.9197,5.9197) (5.9599,5.9599) (6.0000,6.0000)};\draw[structurecolor,dashed]plot coordinates{(-5.6569,-6.0000) (-5.6143,-5.9599) (-5.5716,-5.9197) (-5.5290,-5.8796) (-5.4863,-5.8395) (-5.4436,-5.7993) (-5.4008,-5.7592) (-5.3580,-5.7191) (-5.3151,-5.6789) (-5.2722,-5.6388) (-5.2292,-5.5987) (-5.1863,-5.5585) (-5.1432,-5.5184) (-5.1001,-5.4783) (-5.0570,-5.4381) (-5.0138,-5.3980) (-4.9706,-5.3579) (-4.9273,-5.3177) (-4.8840,-5.2776) (-4.8406,-5.2375) (-4.7971,-5.1973) (-4.7536,-5.1572) (-4.7100,-5.1171) (-4.6664,-5.0769) (-4.6227,-5.0368) (-4.5789,-4.9967) (-4.5351,-4.9565) (-4.4912,-4.9164) (-4.4472,-4.8763) (-4.4032,-4.8361) (-4.3591,-4.7960) (-4.3149,-4.7559) (-4.2706,-4.7157) (-4.2262,-4.6756) (-4.1818,-4.6355) (-4.1373,-4.5953) (-4.0926,-4.5552) (-4.0479,-4.5151) (-4.0031,-4.4749) (-3.9582,-4.4348) (-3.9132,-4.3946) (-3.8680,-4.3545) (-3.8228,-4.3144) (-3.7775,-4.2742) (-3.7320,-4.2341) (-3.6864,-4.1940) (-3.6407,-4.1538) (-3.5948,-4.1137) (-3.5488,-4.0736) (-3.5027,-4.0334) (-3.4564,-3.9933) (-3.4099,-3.9532) (-3.3633,-3.9130) (-3.3165,-3.8729) (-3.2696,-3.8328) (-3.2224,-3.7926) (-3.1751,-3.7525) (-3.1276,-3.7124) (-3.0798,-3.6722) (-3.0319,-3.6321) (-2.9837,-3.5920) (-2.9352,-3.5518) (-2.8865,-3.5117) (-2.8376,-3.4716) (-2.7883,-3.4314) (-2.7388,-3.3913) (-2.6889,-3.3512) (-2.6387,-3.3110) (-2.5882,-3.2709) (-2.5373,-3.2308) (-2.4860,-3.1906) (-2.4343,-3.1505) (-2.3821,-3.1104) (-2.3295,-3.0702) (-2.2763,-3.0301) (-2.2226,-2.9900) (-2.1683,-2.9498) (-2.1134,-2.9097) (-2.0578,-2.8696) (-2.0014,-2.8294) (-1.9443,-2.7893) (-1.8862,-2.7492) (-1.8273,-2.7090) (-1.7672,-2.6689) (-1.7060,-2.6288) (-1.6435,-2.5886) (-1.5795,-2.5485) (-1.5139,-2.5084) (-1.4464,-2.4682) (-1.3768,-2.4281) (-1.3047,-2.3880) (-1.2298,-2.3478) (-1.1513,-2.3077) (-1.0686,-2.2676) (-0.9805,-2.2274) (-0.8856,-2.1873) (-0.7812,-2.1472) (-0.6630,-2.1070) (-0.5216,-2.0669) (-0.3282,-2.0268)};\draw[structurecolor,dashed]plot coordinates{(-0.3282,2.0268) (-0.5216,2.0669) (-0.6630,2.1070) (-0.7812,2.1472) (-0.8856,2.1873) (-0.9805,2.2274) (-1.0686,2.2676) (-1.1513,2.3077) (-1.2298,2.3478) (-1.3047,2.3880) (-1.3768,2.4281) (-1.4464,2.4682) (-1.5139,2.5084) (-1.5795,2.5485) (-1.6435,2.5886) (-1.7060,2.6288) (-1.7672,2.6689) (-1.8273,2.7090) (-1.8862,2.7492) (-1.9443,2.7893) (-2.0014,2.8294) (-2.0578,2.8696) (-2.1134,2.9097) (-2.1683,2.9498) (-2.2226,2.9900) (-2.2763,3.0301) (-2.3295,3.0702) (-2.3821,3.1104) (-2.4343,3.1505) (-2.4860,3.1906) (-2.5373,3.2308) (-2.5882,3.2709) (-2.6387,3.3110) (-2.6889,3.3512) (-2.7388,3.3913) (-2.7883,3.4314) (-2.8376,3.4716) (-2.8865,3.5117) (-2.9352,3.5518) (-2.9837,3.5920) (-3.0319,3.6321) (-3.0798,3.6722) (-3.1276,3.7124) (-3.1751,3.7525) (-3.2224,3.7926) (-3.2696,3.8328) (-3.3165,3.8729) (-3.3633,3.9130) (-3.4099,3.9532) (-3.4564,3.9933) (-3.5027,4.0334) (-3.5488,4.0736) (-3.5948,4.1137) (-3.6407,4.1538) (-3.6864,4.1940) (-3.7320,4.2341) (-3.7775,4.2742) (-3.8228,4.3144) (-3.8680,4.3545) (-3.9132,4.3946) (-3.9582,4.4348) (-4.0031,4.4749) (-4.0479,4.5151) (-4.0926,4.5552) (-4.1373,4.5953) (-4.1818,4.6355) (-4.2262,4.6756) (-4.2706,4.7157) (-4.3149,4.7559) (-4.3591,4.7960) (-4.4032,4.8361) (-4.4472,4.8763) (-4.4912,4.9164) (-4.5351,4.9565) (-4.5789,4.9967) (-4.6227,5.0368) (-4.6664,5.0769) (-4.7100,5.1171) (-4.7536,5.1572) (-4.7971,5.1973) (-4.8406,5.2375) (-4.8840,5.2776) (-4.9273,5.3177) (-4.9706,5.3579) (-5.0138,5.3980) (-5.0570,5.4381) (-5.1001,5.4783) (-5.1432,5.5184) (-5.1863,5.5585) (-5.2292,5.5987) (-5.2722,5.6388) (-5.3151,5.6789) (-5.3580,5.7191) (-5.4008,5.7592) (-5.4436,5.7993) (-5.4863,5.8395) (-5.5290,5.8796) (-5.5716,5.9197) (-5.6143,5.9599) (-5.6569,6.0000)};\draw[structurecolor,dashed]plot coordinates{(5.6569,-6.0000) (5.6143,-5.9599) (5.5716,-5.9197) (5.5290,-5.8796) (5.4863,-5.8395) (5.4436,-5.7993) (5.4008,-5.7592) (5.3580,-5.7191) (5.3151,-5.6789) (5.2722,-5.6388) (5.2292,-5.5987) (5.1863,-5.5585) (5.1432,-5.5184) (5.1001,-5.4783) (5.0570,-5.4381) (5.0138,-5.3980) (4.9706,-5.3579) (4.9273,-5.3177) (4.8840,-5.2776) (4.8406,-5.2375) (4.7971,-5.1973) (4.7536,-5.1572) (4.7100,-5.1171) (4.6664,-5.0769) (4.6227,-5.0368) (4.5789,-4.9967) (4.5351,-4.9565) (4.4912,-4.9164) (4.4472,-4.8763) (4.4032,-4.8361) (4.3591,-4.7960) (4.3149,-4.7559) (4.2706,-4.7157) (4.2262,-4.6756) (4.1818,-4.6355) (4.1373,-4.5953) (4.0926,-4.5552) (4.0479,-4.5151) (4.0031,-4.4749) (3.9582,-4.4348) (3.9132,-4.3946) (3.8680,-4.3545) (3.8228,-4.3144) (3.7775,-4.2742) (3.7320,-4.2341) (3.6864,-4.1940) (3.6407,-4.1538) (3.5948,-4.1137) (3.5488,-4.0736) (3.5027,-4.0334) (3.4564,-3.9933) (3.4099,-3.9532) (3.3633,-3.9130) (3.3165,-3.8729) (3.2696,-3.8328) (3.2224,-3.7926) (3.1751,-3.7525) (3.1276,-3.7124) (3.0798,-3.6722) (3.0319,-3.6321) (2.9837,-3.5920) (2.9352,-3.5518) (2.8865,-3.5117) (2.8376,-3.4716) (2.7883,-3.4314) (2.7388,-3.3913) (2.6889,-3.3512) (2.6387,-3.3110) (2.5882,-3.2709) (2.5373,-3.2308) (2.4860,-3.1906) (2.4343,-3.1505) (2.3821,-3.1104) (2.3295,-3.0702) (2.2763,-3.0301) (2.2226,-2.9900) (2.1683,-2.9498) (2.1134,-2.9097) (2.0578,-2.8696) (2.0014,-2.8294) (1.9443,-2.7893) (1.8862,-2.7492) (1.8273,-2.7090) (1.7672,-2.6689) (1.7060,-2.6288) (1.6435,-2.5886) (1.5795,-2.5485) (1.5139,-2.5084) (1.4464,-2.4682) (1.3768,-2.4281) (1.3047,-2.3880) (1.2298,-2.3478) (1.1513,-2.3077) (1.0686,-2.2676) (0.9805,-2.2274) (0.8856,-2.1873) (0.7812,-2.1472) (0.6630,-2.1070) (0.5216,-2.0669) (0.3282,-2.0268)};\draw[structurecolor,dashed]plot coordinates{(0.3282,2.0268) (0.5216,2.0669) (0.6630,2.1070) (0.7812,2.1472) (0.8856,2.1873) (0.9805,2.2274) (1.0686,2.2676) (1.1513,2.3077) (1.2298,2.3478) (1.3047,2.3880) (1.3768,2.4281) (1.4464,2.4682) (1.5139,2.5084) (1.5795,2.5485) (1.6435,2.5886) (1.7060,2.6288) (1.7672,2.6689) (1.8273,2.7090) (1.8862,2.7492) (1.9443,2.7893) (2.0014,2.8294) (2.0578,2.8696) (2.1134,2.9097) (2.1683,2.9498) (2.2226,2.9900) (2.2763,3.0301) (2.3295,3.0702) (2.3821,3.1104) (2.4343,3.1505) (2.4860,3.1906) (2.5373,3.2308) (2.5882,3.2709) (2.6387,3.3110) (2.6889,3.3512) (2.7388,3.3913) (2.7883,3.4314) (2.8376,3.4716) (2.8865,3.5117) (2.9352,3.5518) (2.9837,3.5920) (3.0319,3.6321) (3.0798,3.6722) (3.1276,3.7124) (3.1751,3.7525) (3.2224,3.7926) (3.2696,3.8328) (3.3165,3.8729) (3.3633,3.9130) (3.4099,3.9532) (3.4564,3.9933) (3.5027,4.0334) (3.5488,4.0736) (3.5948,4.1137) (3.6407,4.1538) (3.6864,4.1940) (3.7320,4.2341) (3.7775,4.2742) (3.8228,4.3144) (3.8680,4.3545) (3.9132,4.3946) (3.9582,4.4348) (4.0031,4.4749) (4.0479,4.5151) (4.0926,4.5552) (4.1373,4.5953) (4.1818,4.6355) (4.2262,4.6756) (4.2706,4.7157) (4.3149,4.7559) (4.3591,4.7960) (4.4032,4.8361) (4.4472,4.8763) (4.4912,4.9164) (4.5351,4.9565) (4.5789,4.9967) (4.6227,5.0368) (4.6664,5.0769) (4.7100,5.1171) (4.7536,5.1572) (4.7971,5.1973) (4.8406,5.2375) (4.8840,5.2776) (4.9273,5.3177) (4.9706,5.3579) (5.0138,5.3980) (5.0570,5.4381) (5.1001,5.4783) (5.1432,5.5184) (5.1863,5.5585) (5.2292,5.5987) (5.2722,5.6388) (5.3151,5.6789) (5.3580,5.7191) (5.4008,5.7592) (5.4436,5.7993) (5.4863,5.8395) (5.5290,5.8796) (5.5716,5.9197) (5.6143,5.9599) (5.6569,6.0000)};\draw(-6,-.08)--(-6,.08);\node[below,font=\scriptsize,fill=white,inner sep=.2pt]at(-6,-.08){$ -6 $};\draw(-5,-.08)--(-5,.08);\draw(-4,-.08)--(-4,.08);\node[below,font=\scriptsize,fill=white,inner sep=.2pt]at(-4,-.08){$ -4 $};\draw(-3,-.08)--(-3,.08);\draw(-2,-.08)--(-2,.08);\node[below,font=\scriptsize,fill=white,inner sep=.2pt]at(-2,-.08){$ -2 $};\draw(-1,-.08)--(-1,.08);\draw(1,-.08)--(1,.08);\draw(2,-.08)--(2,.08);\node[below,font=\scriptsize,fill=white,inner sep=.2pt]at(2,-.08){$ 2 $};\draw(3,-.08)--(3,.08);\draw(4,-.08)--(4,.08);\node[below,font=\scriptsize,fill=white,inner sep=.2pt]at(4,-.08){$ 4 $};\draw(5,-.08)--(5,.08);\draw(6,-.08)--(6,.08);\node[below,font=\scriptsize,fill=white,inner sep=.2pt]at(6,-.08){$ 6 $};\draw(-.08,-6)--(.08,-6);\node[left,font=\scriptsize,fill=white,inner sep=.2pt]at(-.08,-6){$ -6 $};\draw(-.08,-5)--(.08,-5);\draw(-.08,-4)--(.08,-4);\node[left,font=\scriptsize,fill=white,inner sep=.2pt]at(-.08,-4){$ -4 $};\draw(-.08,-3)--(.08,-3);\draw(-.08,-2)--(.08,-2);\node[left,font=\scriptsize,fill=white,inner sep=.2pt]at(-.08,-2){$ -2 $};\draw(-.08,-1)--(.08,-1);\draw(-.08,1)--(.08,1);\draw(-.08,2)--(.08,2);\node[left,font=\scriptsize,fill=white,inner sep=.2pt]at(-.08,2){$ 2 $};\draw(-.08,3)--(.08,3);\draw(-.08,4)--(.08,4);\node[left,font=\scriptsize,fill=white,inner sep=.2pt]at(-.08,4){$ 4 $};\draw(-.08,5)--(.08,5);\draw(-.08,6)--(.08,6);\node[left,font=\scriptsize,fill=white,inner sep=.2pt]at(-.08,6){$ 6 $};\end{tikzpicture}\end{center}
}

\newcommand{\ODEExamPracticeFinalField}{%
\begin{center}\begin{tikzpicture}[x=.62cm,y=.53cm]
\clip (-5.6,-5.6) rectangle (5.7,5.8);
\draw[->](-5,0)--(5.35,0)node[right]{$x$};\draw[->](0,-5)--(0,5.35)node[above]{$y$};
\draw[gray!55](-5.120,-5.000)--(-4.880,-5.000);
\draw[gray!55](-5.025,-4.617)--(-4.975,-4.383);
\draw[gray!55](-5.013,-4.119)--(-4.987,-3.881);
\draw[gray!55](-5.009,-3.620)--(-4.991,-3.380);
\draw[gray!55](-5.007,-3.120)--(-4.993,-2.880);
\draw[gray!55](-5.006,-2.620)--(-4.994,-2.380);
\draw[gray!55](-5.006,-2.120)--(-4.994,-1.880);
\draw[gray!55](-5.005,-1.620)--(-4.995,-1.380);
\draw[gray!55](-5.005,-1.120)--(-4.995,-0.880);
\draw[gray!55](-5.005,-0.620)--(-4.995,-0.380);
\draw[gray!55](-5.005,-0.120)--(-4.995,0.120);
\draw[gray!55](-5.005,0.380)--(-4.995,0.620);
\draw[gray!55](-5.005,0.880)--(-4.995,1.120);
\draw[gray!55](-5.005,1.380)--(-4.995,1.620);
\draw[gray!55](-5.006,1.880)--(-4.994,2.120);
\draw[gray!55](-5.006,2.380)--(-4.994,2.620);
\draw[gray!55](-5.007,2.880)--(-4.993,3.120);
\draw[gray!55](-5.009,3.380)--(-4.991,3.620);
\draw[gray!55](-5.013,3.881)--(-4.987,4.119);
\draw[gray!55](-5.025,4.383)--(-4.975,4.617);
\draw[gray!55](-5.120,5.000)--(-4.880,5.000);
\draw[gray!55](-4.525,-4.883)--(-4.475,-5.117);
\draw[gray!55](-4.620,-4.500)--(-4.380,-4.500);
\draw[gray!55](-4.527,-4.117)--(-4.473,-3.883);
\draw[gray!55](-4.515,-3.619)--(-4.485,-3.381);
\draw[gray!55](-4.511,-3.120)--(-4.489,-2.880);
\draw[gray!55](-4.509,-2.620)--(-4.491,-2.380);
\draw[gray!55](-4.507,-2.120)--(-4.493,-1.880);
\draw[gray!55](-4.507,-1.620)--(-4.493,-1.380);
\draw[gray!55](-4.506,-1.120)--(-4.494,-0.880);
\draw[gray!55](-4.506,-0.620)--(-4.494,-0.380);
\draw[gray!55](-4.506,-0.120)--(-4.494,0.120);
\draw[gray!55](-4.506,0.380)--(-4.494,0.620);
\draw[gray!55](-4.506,0.880)--(-4.494,1.120);
\draw[gray!55](-4.507,1.380)--(-4.493,1.620);
\draw[gray!55](-4.507,1.880)--(-4.493,2.120);
\draw[gray!55](-4.509,2.380)--(-4.491,2.620);
\draw[gray!55](-4.511,2.880)--(-4.489,3.120);
\draw[gray!55](-4.515,3.381)--(-4.485,3.619);
\draw[gray!55](-4.527,3.883)--(-4.473,4.117);
\draw[gray!55](-4.620,4.500)--(-4.380,4.500);
\draw[gray!55](-4.525,5.117)--(-4.475,4.883);
\draw[gray!55](-4.013,-4.881)--(-3.987,-5.119);
\draw[gray!55](-4.027,-4.383)--(-3.973,-4.617);
\draw[gray!55](-4.120,-4.000)--(-3.880,-4.000);
\draw[gray!55](-4.031,-3.616)--(-3.969,-3.384);
\draw[gray!55](-4.017,-3.119)--(-3.983,-2.881);
\draw[gray!55](-4.012,-2.619)--(-3.988,-2.381);
\draw[gray!55](-4.010,-2.120)--(-3.990,-1.880);
\draw[gray!55](-4.009,-1.620)--(-3.991,-1.380);
\draw[gray!55](-4.008,-1.120)--(-3.992,-0.880);
\draw[gray!55](-4.008,-0.620)--(-3.992,-0.380);
\draw[gray!55](-4.007,-0.120)--(-3.993,0.120);
\draw[gray!55](-4.008,0.380)--(-3.992,0.620);
\draw[gray!55](-4.008,0.880)--(-3.992,1.120);
\draw[gray!55](-4.009,1.380)--(-3.991,1.620);
\draw[gray!55](-4.010,1.880)--(-3.990,2.120);
\draw[gray!55](-4.012,2.381)--(-3.988,2.619);
\draw[gray!55](-4.017,2.881)--(-3.983,3.119);
\draw[gray!55](-4.031,3.384)--(-3.969,3.616);
\draw[gray!55](-4.120,4.000)--(-3.880,4.000);
\draw[gray!55](-4.027,4.617)--(-3.973,4.383);
\draw[gray!55](-4.013,5.119)--(-3.987,4.881);
\draw[gray!55](-3.509,-4.880)--(-3.491,-5.120);
\draw[gray!55](-3.515,-4.381)--(-3.485,-4.619);
\draw[gray!55](-3.531,-3.884)--(-3.469,-4.116);
\draw[gray!55](-3.620,-3.500)--(-3.380,-3.500);
\draw[gray!55](-3.535,-3.115)--(-3.465,-2.885);
\draw[gray!55](-3.520,-2.618)--(-3.480,-2.382);
\draw[gray!55](-3.514,-2.119)--(-3.486,-1.881);
\draw[gray!55](-3.512,-1.619)--(-3.488,-1.381);
\draw[gray!55](-3.511,-1.120)--(-3.489,-0.880);
\draw[gray!55](-3.510,-0.620)--(-3.490,-0.380);
\draw[gray!55](-3.510,-0.120)--(-3.490,0.120);
\draw[gray!55](-3.510,0.380)--(-3.490,0.620);
\draw[gray!55](-3.511,0.880)--(-3.489,1.120);
\draw[gray!55](-3.512,1.381)--(-3.488,1.619);
\draw[gray!55](-3.514,1.881)--(-3.486,2.119);
\draw[gray!55](-3.520,2.382)--(-3.480,2.618);
\draw[gray!55](-3.535,2.885)--(-3.465,3.115);
\draw[gray!55](-3.620,3.500)--(-3.380,3.500);
\draw[gray!55](-3.531,4.116)--(-3.469,3.884);
\draw[gray!55](-3.515,4.619)--(-3.485,4.381);
\draw[gray!55](-3.509,5.120)--(-3.491,4.880);
\draw[gray!55](-3.007,-4.880)--(-2.993,-5.120);
\draw[gray!55](-3.011,-4.380)--(-2.989,-4.620);
\draw[gray!55](-3.017,-3.881)--(-2.983,-4.119);
\draw[gray!55](-3.035,-3.385)--(-2.965,-3.615);
\draw[gray!55](-3.120,-3.000)--(-2.880,-3.000);
\draw[gray!55](-3.041,-2.613)--(-2.959,-2.387);
\draw[gray!55](-3.024,-2.118)--(-2.976,-1.882);
\draw[gray!55](-3.018,-1.619)--(-2.982,-1.381);
\draw[gray!55](-3.015,-1.119)--(-2.985,-0.881);
\draw[gray!55](-3.014,-0.619)--(-2.986,-0.381);
\draw[gray!55](-3.013,-0.119)--(-2.987,0.119);
\draw[gray!55](-3.014,0.381)--(-2.986,0.619);
\draw[gray!55](-3.015,0.881)--(-2.985,1.119);
\draw[gray!55](-3.018,1.381)--(-2.982,1.619);
\draw[gray!55](-3.024,1.882)--(-2.976,2.118);
\draw[gray!55](-3.041,2.387)--(-2.959,2.613);
\draw[gray!55](-3.120,3.000)--(-2.880,3.000);
\draw[gray!55](-3.035,3.615)--(-2.965,3.385);
\draw[gray!55](-3.017,4.119)--(-2.983,3.881);
\draw[gray!55](-3.011,4.620)--(-2.989,4.380);
\draw[gray!55](-3.007,5.120)--(-2.993,4.880);
\draw[gray!55](-2.506,-4.880)--(-2.494,-5.120);
\draw[gray!55](-2.509,-4.380)--(-2.491,-4.620);
\draw[gray!55](-2.512,-3.881)--(-2.488,-4.119);
\draw[gray!55](-2.520,-3.382)--(-2.480,-3.618);
\draw[gray!55](-2.541,-2.887)--(-2.459,-3.113);
\draw[gray!55](-2.620,-2.500)--(-2.380,-2.500);
\draw[gray!55](-2.549,-2.110)--(-2.451,-1.890);
\draw[gray!55](-2.529,-1.616)--(-2.471,-1.384);
\draw[gray!55](-2.522,-1.118)--(-2.478,-0.882);
\draw[gray!55](-2.520,-0.618)--(-2.480,-0.382);
\draw[gray!55](-2.519,-0.118)--(-2.481,0.118);
\draw[gray!55](-2.520,0.382)--(-2.480,0.618);
\draw[gray!55](-2.522,0.882)--(-2.478,1.118);
\draw[gray!55](-2.529,1.384)--(-2.471,1.616);
\draw[gray!55](-2.549,1.890)--(-2.451,2.110);
\draw[gray!55](-2.620,2.500)--(-2.380,2.500);
\draw[gray!55](-2.541,3.113)--(-2.459,2.887);
\draw[gray!55](-2.520,3.618)--(-2.480,3.382);
\draw[gray!55](-2.512,4.119)--(-2.488,3.881);
\draw[gray!55](-2.509,4.620)--(-2.491,4.380);
\draw[gray!55](-2.506,5.120)--(-2.494,4.880);
\draw[gray!55](-2.006,-4.880)--(-1.994,-5.120);
\draw[gray!55](-2.007,-4.380)--(-1.993,-4.620);
\draw[gray!55](-2.010,-3.880)--(-1.990,-4.120);
\draw[gray!55](-2.014,-3.381)--(-1.986,-3.619);
\draw[gray!55](-2.024,-2.882)--(-1.976,-3.118);
\draw[gray!55](-2.049,-2.390)--(-1.951,-2.610);
\draw[gray!55](-2.120,-2.000)--(-1.880,-2.000);
\draw[gray!55](-2.060,-1.604)--(-1.940,-1.396);
\draw[gray!55](-2.038,-1.114)--(-1.962,-0.886);
\draw[gray!55](-2.031,-0.616)--(-1.969,-0.384);
\draw[gray!55](-2.029,-0.116)--(-1.971,0.116);
\draw[gray!55](-2.031,0.384)--(-1.969,0.616);
\draw[gray!55](-2.038,0.886)--(-1.962,1.114);
\draw[gray!55](-2.060,1.396)--(-1.940,1.604);
\draw[gray!55](-2.120,2.000)--(-1.880,2.000);
\draw[gray!55](-2.049,2.610)--(-1.951,2.390);
\draw[gray!55](-2.024,3.118)--(-1.976,2.882);
\draw[gray!55](-2.014,3.619)--(-1.986,3.381);
\draw[gray!55](-2.010,4.120)--(-1.990,3.880);
\draw[gray!55](-2.007,4.620)--(-1.993,4.380);
\draw[gray!55](-2.006,5.120)--(-1.994,4.880);
\draw[gray!55](-1.505,-4.880)--(-1.495,-5.120);
\draw[gray!55](-1.507,-4.380)--(-1.493,-4.620);
\draw[gray!55](-1.509,-3.880)--(-1.491,-4.120);
\draw[gray!55](-1.512,-3.381)--(-1.488,-3.619);
\draw[gray!55](-1.518,-2.881)--(-1.482,-3.119);
\draw[gray!55](-1.529,-2.384)--(-1.471,-2.616);
\draw[gray!55](-1.560,-1.896)--(-1.440,-2.104);
\draw[gray!55](-1.620,-1.500)--(-1.380,-1.500);
\draw[gray!55](-1.575,-1.094)--(-1.425,-0.906);
\draw[gray!55](-1.554,-0.607)--(-1.446,-0.393);
\draw[gray!55](-1.549,-0.110)--(-1.451,0.110);
\draw[gray!55](-1.554,0.393)--(-1.446,0.607);
\draw[gray!55](-1.575,0.906)--(-1.425,1.094);
\draw[gray!55](-1.620,1.500)--(-1.380,1.500);
\draw[gray!55](-1.560,2.104)--(-1.440,1.896);
\draw[gray!55](-1.529,2.616)--(-1.471,2.384);
\draw[gray!55](-1.518,3.119)--(-1.482,2.881);
\draw[gray!55](-1.512,3.619)--(-1.488,3.381);
\draw[gray!55](-1.509,4.120)--(-1.491,3.880);
\draw[gray!55](-1.507,4.620)--(-1.493,4.380);
\draw[gray!55](-1.505,5.120)--(-1.495,4.880);
\draw[gray!55](-1.005,-4.880)--(-0.995,-5.120);
\draw[gray!55](-1.006,-4.380)--(-0.994,-4.620);
\draw[gray!55](-1.008,-3.880)--(-0.992,-4.120);
\draw[gray!55](-1.011,-3.380)--(-0.989,-3.620);
\draw[gray!55](-1.015,-2.881)--(-0.985,-3.119);
\draw[gray!55](-1.022,-2.382)--(-0.978,-2.618);
\draw[gray!55](-1.038,-1.886)--(-0.962,-2.114);
\draw[gray!55](-1.075,-1.406)--(-0.925,-1.594);
\draw[gray!55](-1.120,-1.000)--(-0.880,-1.000);
\draw[gray!55](-1.096,-0.572)--(-0.904,-0.428);
\draw[gray!55](-1.085,-0.085)--(-0.915,0.085);
\draw[gray!55](-1.096,0.428)--(-0.904,0.572);
\draw[gray!55](-1.120,1.000)--(-0.880,1.000);
\draw[gray!55](-1.075,1.594)--(-0.925,1.406);
\draw[gray!55](-1.038,2.114)--(-0.962,1.886);
\draw[gray!55](-1.022,2.618)--(-0.978,2.382);
\draw[gray!55](-1.015,3.119)--(-0.985,2.881);
\draw[gray!55](-1.011,3.620)--(-0.989,3.380);
\draw[gray!55](-1.008,4.120)--(-0.992,3.880);
\draw[gray!55](-1.006,4.620)--(-0.994,4.380);
\draw[gray!55](-1.005,5.120)--(-0.995,4.880);
\draw[gray!55](-0.505,-4.880)--(-0.495,-5.120);
\draw[gray!55](-0.506,-4.380)--(-0.494,-4.620);
\draw[gray!55](-0.508,-3.880)--(-0.492,-4.120);
\draw[gray!55](-0.510,-3.380)--(-0.490,-3.620);
\draw[gray!55](-0.514,-2.881)--(-0.486,-3.119);
\draw[gray!55](-0.520,-2.382)--(-0.480,-2.618);
\draw[gray!55](-0.531,-1.884)--(-0.469,-2.116);
\draw[gray!55](-0.554,-1.393)--(-0.446,-1.607);
\draw[gray!55](-0.596,-0.928)--(-0.404,-1.072);
\draw[gray!55](-0.620,-0.500)--(-0.380,-0.500);
\draw[gray!55](-0.616,-0.029)--(-0.384,0.029);
\draw[gray!55](-0.620,0.500)--(-0.380,0.500);
\draw[gray!55](-0.596,1.072)--(-0.404,0.928);
\draw[gray!55](-0.554,1.607)--(-0.446,1.393);
\draw[gray!55](-0.531,2.116)--(-0.469,1.884);
\draw[gray!55](-0.520,2.618)--(-0.480,2.382);
\draw[gray!55](-0.514,3.119)--(-0.486,2.881);
\draw[gray!55](-0.510,3.620)--(-0.490,3.380);
\draw[gray!55](-0.508,4.120)--(-0.492,3.880);
\draw[gray!55](-0.506,4.620)--(-0.494,4.380);
\draw[gray!55](-0.505,5.120)--(-0.495,4.880);
\draw[gray!55](-0.005,-4.880)--(0.005,-5.120);
\draw[gray!55](-0.006,-4.380)--(0.006,-4.620);
\draw[gray!55](-0.007,-3.880)--(0.007,-4.120);
\draw[gray!55](-0.010,-3.380)--(0.010,-3.620);
\draw[gray!55](-0.013,-2.881)--(0.013,-3.119);
\draw[gray!55](-0.019,-2.382)--(0.019,-2.618);
\draw[gray!55](-0.029,-1.884)--(0.029,-2.116);
\draw[gray!55](-0.049,-1.390)--(0.049,-1.610);
\draw[gray!55](-0.085,-0.915)--(0.085,-1.085);
\draw[gray!55](-0.116,-0.471)--(0.116,-0.529);
\draw[gray!55](-0.120,0.000)--(0.120,0.000);
\draw[gray!55](-0.116,0.529)--(0.116,0.471);
\draw[gray!55](-0.085,1.085)--(0.085,0.915);
\draw[gray!55](-0.049,1.610)--(0.049,1.390);
\draw[gray!55](-0.029,2.116)--(0.029,1.884);
\draw[gray!55](-0.019,2.618)--(0.019,2.382);
\draw[gray!55](-0.013,3.119)--(0.013,2.881);
\draw[gray!55](-0.010,3.620)--(0.010,3.380);
\draw[gray!55](-0.007,4.120)--(0.007,3.880);
\draw[gray!55](-0.006,4.620)--(0.006,4.380);
\draw[gray!55](-0.005,5.120)--(0.005,4.880);
\draw[gray!55](0.495,-4.880)--(0.505,-5.120);
\draw[gray!55](0.494,-4.380)--(0.506,-4.620);
\draw[gray!55](0.492,-3.880)--(0.508,-4.120);
\draw[gray!55](0.490,-3.380)--(0.510,-3.620);
\draw[gray!55](0.486,-2.881)--(0.514,-3.119);
\draw[gray!55](0.480,-2.382)--(0.520,-2.618);
\draw[gray!55](0.469,-1.884)--(0.531,-2.116);
\draw[gray!55](0.446,-1.393)--(0.554,-1.607);
\draw[gray!55](0.404,-0.928)--(0.596,-1.072);
\draw[gray!55](0.380,-0.500)--(0.620,-0.500);
\draw[gray!55](0.384,-0.029)--(0.616,0.029);
\draw[gray!55](0.380,0.500)--(0.620,0.500);
\draw[gray!55](0.404,1.072)--(0.596,0.928);
\draw[gray!55](0.446,1.607)--(0.554,1.393);
\draw[gray!55](0.469,2.116)--(0.531,1.884);
\draw[gray!55](0.480,2.618)--(0.520,2.382);
\draw[gray!55](0.486,3.119)--(0.514,2.881);
\draw[gray!55](0.490,3.620)--(0.510,3.380);
\draw[gray!55](0.492,4.120)--(0.508,3.880);
\draw[gray!55](0.494,4.620)--(0.506,4.380);
\draw[gray!55](0.495,5.120)--(0.505,4.880);
\draw[gray!55](0.995,-4.880)--(1.005,-5.120);
\draw[gray!55](0.994,-4.380)--(1.006,-4.620);
\draw[gray!55](0.992,-3.880)--(1.008,-4.120);
\draw[gray!55](0.989,-3.380)--(1.011,-3.620);
\draw[gray!55](0.985,-2.881)--(1.015,-3.119);
\draw[gray!55](0.978,-2.382)--(1.022,-2.618);
\draw[gray!55](0.962,-1.886)--(1.038,-2.114);
\draw[gray!55](0.925,-1.406)--(1.075,-1.594);
\draw[gray!55](0.880,-1.000)--(1.120,-1.000);
\draw[gray!55](0.904,-0.572)--(1.096,-0.428);
\draw[gray!55](0.915,-0.085)--(1.085,0.085);
\draw[gray!55](0.904,0.428)--(1.096,0.572);
\draw[gray!55](0.880,1.000)--(1.120,1.000);
\draw[gray!55](0.925,1.594)--(1.075,1.406);
\draw[gray!55](0.962,2.114)--(1.038,1.886);
\draw[gray!55](0.978,2.618)--(1.022,2.382);
\draw[gray!55](0.985,3.119)--(1.015,2.881);
\draw[gray!55](0.989,3.620)--(1.011,3.380);
\draw[gray!55](0.992,4.120)--(1.008,3.880);
\draw[gray!55](0.994,4.620)--(1.006,4.380);
\draw[gray!55](0.995,5.120)--(1.005,4.880);
\draw[gray!55](1.495,-4.880)--(1.505,-5.120);
\draw[gray!55](1.493,-4.380)--(1.507,-4.620);
\draw[gray!55](1.491,-3.880)--(1.509,-4.120);
\draw[gray!55](1.488,-3.381)--(1.512,-3.619);
\draw[gray!55](1.482,-2.881)--(1.518,-3.119);
\draw[gray!55](1.471,-2.384)--(1.529,-2.616);
\draw[gray!55](1.440,-1.896)--(1.560,-2.104);
\draw[gray!55](1.380,-1.500)--(1.620,-1.500);
\draw[gray!55](1.425,-1.094)--(1.575,-0.906);
\draw[gray!55](1.446,-0.607)--(1.554,-0.393);
\draw[gray!55](1.451,-0.110)--(1.549,0.110);
\draw[gray!55](1.446,0.393)--(1.554,0.607);
\draw[gray!55](1.425,0.906)--(1.575,1.094);
\draw[gray!55](1.380,1.500)--(1.620,1.500);
\draw[gray!55](1.440,2.104)--(1.560,1.896);
\draw[gray!55](1.471,2.616)--(1.529,2.384);
\draw[gray!55](1.482,3.119)--(1.518,2.881);
\draw[gray!55](1.488,3.619)--(1.512,3.381);
\draw[gray!55](1.491,4.120)--(1.509,3.880);
\draw[gray!55](1.493,4.620)--(1.507,4.380);
\draw[gray!55](1.495,5.120)--(1.505,4.880);
\draw[gray!55](1.994,-4.880)--(2.006,-5.120);
\draw[gray!55](1.993,-4.380)--(2.007,-4.620);
\draw[gray!55](1.990,-3.880)--(2.010,-4.120);
\draw[gray!55](1.986,-3.381)--(2.014,-3.619);
\draw[gray!55](1.976,-2.882)--(2.024,-3.118);
\draw[gray!55](1.951,-2.390)--(2.049,-2.610);
\draw[gray!55](1.880,-2.000)--(2.120,-2.000);
\draw[gray!55](1.940,-1.604)--(2.060,-1.396);
\draw[gray!55](1.962,-1.114)--(2.038,-0.886);
\draw[gray!55](1.969,-0.616)--(2.031,-0.384);
\draw[gray!55](1.971,-0.116)--(2.029,0.116);
\draw[gray!55](1.969,0.384)--(2.031,0.616);
\draw[gray!55](1.962,0.886)--(2.038,1.114);
\draw[gray!55](1.940,1.396)--(2.060,1.604);
\draw[gray!55](1.880,2.000)--(2.120,2.000);
\draw[gray!55](1.951,2.610)--(2.049,2.390);
\draw[gray!55](1.976,3.118)--(2.024,2.882);
\draw[gray!55](1.986,3.619)--(2.014,3.381);
\draw[gray!55](1.990,4.120)--(2.010,3.880);
\draw[gray!55](1.993,4.620)--(2.007,4.380);
\draw[gray!55](1.994,5.120)--(2.006,4.880);
\draw[gray!55](2.494,-4.880)--(2.506,-5.120);
\draw[gray!55](2.491,-4.380)--(2.509,-4.620);
\draw[gray!55](2.488,-3.881)--(2.512,-4.119);
\draw[gray!55](2.480,-3.382)--(2.520,-3.618);
\draw[gray!55](2.459,-2.887)--(2.541,-3.113);
\draw[gray!55](2.380,-2.500)--(2.620,-2.500);
\draw[gray!55](2.451,-2.110)--(2.549,-1.890);
\draw[gray!55](2.471,-1.616)--(2.529,-1.384);
\draw[gray!55](2.478,-1.118)--(2.522,-0.882);
\draw[gray!55](2.480,-0.618)--(2.520,-0.382);
\draw[gray!55](2.481,-0.118)--(2.519,0.118);
\draw[gray!55](2.480,0.382)--(2.520,0.618);
\draw[gray!55](2.478,0.882)--(2.522,1.118);
\draw[gray!55](2.471,1.384)--(2.529,1.616);
\draw[gray!55](2.451,1.890)--(2.549,2.110);
\draw[gray!55](2.380,2.500)--(2.620,2.500);
\draw[gray!55](2.459,3.113)--(2.541,2.887);
\draw[gray!55](2.480,3.618)--(2.520,3.382);
\draw[gray!55](2.488,4.119)--(2.512,3.881);
\draw[gray!55](2.491,4.620)--(2.509,4.380);
\draw[gray!55](2.494,5.120)--(2.506,4.880);
\draw[gray!55](2.993,-4.880)--(3.007,-5.120);
\draw[gray!55](2.989,-4.380)--(3.011,-4.620);
\draw[gray!55](2.983,-3.881)--(3.017,-4.119);
\draw[gray!55](2.965,-3.385)--(3.035,-3.615);
\draw[gray!55](2.880,-3.000)--(3.120,-3.000);
\draw[gray!55](2.959,-2.613)--(3.041,-2.387);
\draw[gray!55](2.976,-2.118)--(3.024,-1.882);
\draw[gray!55](2.982,-1.619)--(3.018,-1.381);
\draw[gray!55](2.985,-1.119)--(3.015,-0.881);
\draw[gray!55](2.986,-0.619)--(3.014,-0.381);
\draw[gray!55](2.987,-0.119)--(3.013,0.119);
\draw[gray!55](2.986,0.381)--(3.014,0.619);
\draw[gray!55](2.985,0.881)--(3.015,1.119);
\draw[gray!55](2.982,1.381)--(3.018,1.619);
\draw[gray!55](2.976,1.882)--(3.024,2.118);
\draw[gray!55](2.959,2.387)--(3.041,2.613);
\draw[gray!55](2.880,3.000)--(3.120,3.000);
\draw[gray!55](2.965,3.615)--(3.035,3.385);
\draw[gray!55](2.983,4.119)--(3.017,3.881);
\draw[gray!55](2.989,4.620)--(3.011,4.380);
\draw[gray!55](2.993,5.120)--(3.007,4.880);
\draw[gray!55](3.491,-4.880)--(3.509,-5.120);
\draw[gray!55](3.485,-4.381)--(3.515,-4.619);
\draw[gray!55](3.469,-3.884)--(3.531,-4.116);
\draw[gray!55](3.380,-3.500)--(3.620,-3.500);
\draw[gray!55](3.465,-3.115)--(3.535,-2.885);
\draw[gray!55](3.480,-2.618)--(3.520,-2.382);
\draw[gray!55](3.486,-2.119)--(3.514,-1.881);
\draw[gray!55](3.488,-1.619)--(3.512,-1.381);
\draw[gray!55](3.489,-1.120)--(3.511,-0.880);
\draw[gray!55](3.490,-0.620)--(3.510,-0.380);
\draw[gray!55](3.490,-0.120)--(3.510,0.120);
\draw[gray!55](3.490,0.380)--(3.510,0.620);
\draw[gray!55](3.489,0.880)--(3.511,1.120);
\draw[gray!55](3.488,1.381)--(3.512,1.619);
\draw[gray!55](3.486,1.881)--(3.514,2.119);
\draw[gray!55](3.480,2.382)--(3.520,2.618);
\draw[gray!55](3.465,2.885)--(3.535,3.115);
\draw[gray!55](3.380,3.500)--(3.620,3.500);
\draw[gray!55](3.469,4.116)--(3.531,3.884);
\draw[gray!55](3.485,4.619)--(3.515,4.381);
\draw[gray!55](3.491,5.120)--(3.509,4.880);
\draw[gray!55](3.987,-4.881)--(4.013,-5.119);
\draw[gray!55](3.973,-4.383)--(4.027,-4.617);
\draw[gray!55](3.880,-4.000)--(4.120,-4.000);
\draw[gray!55](3.969,-3.616)--(4.031,-3.384);
\draw[gray!55](3.983,-3.119)--(4.017,-2.881);
\draw[gray!55](3.988,-2.619)--(4.012,-2.381);
\draw[gray!55](3.990,-2.120)--(4.010,-1.880);
\draw[gray!55](3.991,-1.620)--(4.009,-1.380);
\draw[gray!55](3.992,-1.120)--(4.008,-0.880);
\draw[gray!55](3.992,-0.620)--(4.008,-0.380);
\draw[gray!55](3.993,-0.120)--(4.007,0.120);
\draw[gray!55](3.992,0.380)--(4.008,0.620);
\draw[gray!55](3.992,0.880)--(4.008,1.120);
\draw[gray!55](3.991,1.380)--(4.009,1.620);
\draw[gray!55](3.990,1.880)--(4.010,2.120);
\draw[gray!55](3.988,2.381)--(4.012,2.619);
\draw[gray!55](3.983,2.881)--(4.017,3.119);
\draw[gray!55](3.969,3.384)--(4.031,3.616);
\draw[gray!55](3.880,4.000)--(4.120,4.000);
\draw[gray!55](3.973,4.617)--(4.027,4.383);
\draw[gray!55](3.987,5.119)--(4.013,4.881);
\draw[gray!55](4.475,-4.883)--(4.525,-5.117);
\draw[gray!55](4.380,-4.500)--(4.620,-4.500);
\draw[gray!55](4.473,-4.117)--(4.527,-3.883);
\draw[gray!55](4.485,-3.619)--(4.515,-3.381);
\draw[gray!55](4.489,-3.120)--(4.511,-2.880);
\draw[gray!55](4.491,-2.620)--(4.509,-2.380);
\draw[gray!55](4.493,-2.120)--(4.507,-1.880);
\draw[gray!55](4.493,-1.620)--(4.507,-1.380);
\draw[gray!55](4.494,-1.120)--(4.506,-0.880);
\draw[gray!55](4.494,-0.620)--(4.506,-0.380);
\draw[gray!55](4.494,-0.120)--(4.506,0.120);
\draw[gray!55](4.494,0.380)--(4.506,0.620);
\draw[gray!55](4.494,0.880)--(4.506,1.120);
\draw[gray!55](4.493,1.380)--(4.507,1.620);
\draw[gray!55](4.493,1.880)--(4.507,2.120);
\draw[gray!55](4.491,2.380)--(4.509,2.620);
\draw[gray!55](4.489,2.880)--(4.511,3.120);
\draw[gray!55](4.485,3.381)--(4.515,3.619);
\draw[gray!55](4.473,3.883)--(4.527,4.117);
\draw[gray!55](4.380,4.500)--(4.620,4.500);
\draw[gray!55](4.475,5.117)--(4.525,4.883);
\draw[gray!55](4.880,-5.000)--(5.120,-5.000);
\draw[gray!55](4.975,-4.617)--(5.025,-4.383);
\draw[gray!55](4.987,-4.119)--(5.013,-3.881);
\draw[gray!55](4.991,-3.620)--(5.009,-3.380);
\draw[gray!55](4.993,-3.120)--(5.007,-2.880);
\draw[gray!55](4.994,-2.620)--(5.006,-2.380);
\draw[gray!55](4.994,-2.120)--(5.006,-1.880);
\draw[gray!55](4.995,-1.620)--(5.005,-1.380);
\draw[gray!55](4.995,-1.120)--(5.005,-0.880);
\draw[gray!55](4.995,-0.620)--(5.005,-0.380);
\draw[gray!55](4.995,-0.120)--(5.005,0.120);
\draw[gray!55](4.995,0.380)--(5.005,0.620);
\draw[gray!55](4.995,0.880)--(5.005,1.120);
\draw[gray!55](4.995,1.380)--(5.005,1.620);
\draw[gray!55](4.994,1.880)--(5.006,2.120);
\draw[gray!55](4.994,2.380)--(5.006,2.620);
\draw[gray!55](4.993,2.880)--(5.007,3.120);
\draw[gray!55](4.991,3.380)--(5.009,3.620);
\draw[gray!55](4.987,3.881)--(5.013,4.119);
\draw[gray!55](4.975,4.383)--(5.025,4.617);
\draw[gray!55](4.880,5.000)--(5.120,5.000);
\draw[main,thick] plot coordinates {(-2.0000,0.0000) (-2.0126,-0.0505) (-2.0251,-0.1016) (-2.0377,-0.1532) (-2.0502,-0.2052) (-2.0628,-0.2576) (-2.0753,-0.3104) (-2.0879,-0.3633) (-2.1004,-0.4165) (-2.1130,-0.4697) (-2.1255,-0.5230) (-2.1381,-0.5762) (-2.1506,-0.6294) (-2.1632,-0.6824) (-2.1757,-0.7351) (-2.1883,-0.7876) (-2.2008,-0.8398) (-2.2134,-0.8915) (-2.2259,-0.9428) (-2.2385,-0.9936) (-2.2510,-1.0438) (-2.2636,-1.0934) (-2.2762,-1.1424) (-2.2887,-1.1907) (-2.3013,-1.2383) (-2.3138,-1.2851) (-2.3264,-1.3312) (-2.3389,-1.3765) (-2.3515,-1.4210) (-2.3640,-1.4646) (-2.3766,-1.5074) (-2.3891,-1.5494) (-2.4017,-1.5905) (-2.4142,-1.6307) (-2.4268,-1.6700) (-2.4393,-1.7085) (-2.4519,-1.7461) (-2.4644,-1.7829) (-2.4770,-1.8188) (-2.4895,-1.8539) (-2.5021,-1.8881) (-2.5146,-1.9215) (-2.5272,-1.9542) (-2.5397,-1.9860) (-2.5523,-2.0171) (-2.5649,-2.0474) (-2.5774,-2.0770) (-2.5900,-2.1059) (-2.6025,-2.1341) (-2.6151,-2.1616) (-2.6276,-2.1885) (-2.6402,-2.2147) (-2.6527,-2.2403) (-2.6653,-2.2654) (-2.6778,-2.2898) (-2.6904,-2.3138) (-2.7029,-2.3372) (-2.7155,-2.3600) (-2.7280,-2.3825) (-2.7406,-2.4044) (-2.7531,-2.4259) (-2.7657,-2.4469) (-2.7782,-2.4676) (-2.7908,-2.4879) (-2.8033,-2.5078) (-2.8159,-2.5273) (-2.8285,-2.5465) (-2.8410,-2.5653) (-2.8536,-2.5839) (-2.8661,-2.6022) (-2.8787,-2.6201) (-2.8912,-2.6378) (-2.9038,-2.6553) (-2.9163,-2.6725) (-2.9289,-2.6895) (-2.9414,-2.7063) (-2.9540,-2.7229) (-2.9665,-2.7392) (-2.9791,-2.7554) (-2.9916,-2.7714) (-3.0042,-2.7873) (-3.0167,-2.8030) (-3.0293,-2.8185) (-3.0418,-2.8339) (-3.0544,-2.8492) (-3.0669,-2.8643) (-3.0795,-2.8794) (-3.0921,-2.8943) (-3.1046,-2.9091) (-3.1172,-2.9238) (-3.1297,-2.9384) (-3.1423,-2.9529) (-3.1548,-2.9674) (-3.1674,-2.9817) (-3.1799,-2.9960) (-3.1925,-3.0103) (-3.2050,-3.0244) (-3.2176,-3.0385) (-3.2301,-3.0525) (-3.2427,-3.0665) (-3.2552,-3.0804) (-3.2678,-3.0943) (-3.2803,-3.1081) (-3.2929,-3.1219) (-3.3054,-3.1357) (-3.3180,-3.1494) (-3.3305,-3.1631) (-3.3431,-3.1767) (-3.3556,-3.1903) (-3.3682,-3.2039) (-3.3808,-3.2174) (-3.3933,-3.2309) (-3.4059,-3.2444) (-3.4184,-3.2579) (-3.4310,-3.2713) (-3.4435,-3.2847) (-3.4561,-3.2981) (-3.4686,-3.3115) (-3.4812,-3.3249) (-3.4937,-3.3382) (-3.5063,-3.3515) (-3.5188,-3.3649) (-3.5314,-3.3782) (-3.5439,-3.3914) (-3.5565,-3.4047) (-3.5690,-3.4180) (-3.5816,-3.4312) (-3.5941,-3.4444) (-3.6067,-3.4577) (-3.6192,-3.4709) (-3.6318,-3.4841) (-3.6444,-3.4973) (-3.6569,-3.5104) (-3.6695,-3.5236) (-3.6820,-3.5368) (-3.6946,-3.5499) (-3.7071,-3.5631) (-3.7197,-3.5762) (-3.7322,-3.5894) (-3.7448,-3.6025) (-3.7573,-3.6156) (-3.7699,-3.6287) (-3.7824,-3.6418) (-3.7950,-3.6549) (-3.8075,-3.6680) (-3.8201,-3.6811) (-3.8326,-3.6942) (-3.8452,-3.7073) (-3.8577,-3.7203) (-3.8703,-3.7334) (-3.8828,-3.7465) (-3.8954,-3.7595) (-3.9079,-3.7726) (-3.9205,-3.7856) (-3.9331,-3.7987) (-3.9456,-3.8117) (-3.9582,-3.8247) (-3.9707,-3.8378) (-3.9833,-3.8508) (-3.9958,-3.8638) (-4.0084,-3.8768) (-4.0209,-3.8899) (-4.0335,-3.9029) (-4.0460,-3.9159) (-4.0586,-3.9289) (-4.0711,-3.9419) (-4.0837,-3.9549) (-4.0962,-3.9679) (-4.1088,-3.9809) (-4.1213,-3.9939) (-4.1339,-4.0068) (-4.1464,-4.0198) (-4.1590,-4.0328) (-4.1715,-4.0458) (-4.1841,-4.0587) (-4.1967,-4.0717) (-4.2092,-4.0847) (-4.2218,-4.0976) (-4.2343,-4.1106) (-4.2469,-4.1236) (-4.2594,-4.1365) (-4.2720,-4.1495) (-4.2845,-4.1624) (-4.2971,-4.1754) (-4.3096,-4.1883) (-4.3222,-4.2012) (-4.3347,-4.2142) (-4.3473,-4.2271) (-4.3598,-4.2400) (-4.3724,-4.2530) (-4.3849,-4.2659) (-4.3975,-4.2788) (-4.4100,-4.2917) (-4.4226,-4.3047) (-4.4351,-4.3176) (-4.4477,-4.3305) (-4.4603,-4.3434) (-4.4728,-4.3563) (-4.4854,-4.3692) (-4.4979,-4.3821) (-4.5105,-4.3950) (-4.5230,-4.4079) (-4.5356,-4.4208) (-4.5481,-4.4337) (-4.5607,-4.4466) (-4.5732,-4.4595) (-4.5858,-4.4724) (-4.5983,-4.4853) (-4.6109,-4.4982) (-4.6234,-4.5111) (-4.6360,-4.5240) (-4.6485,-4.5368) (-4.6611,-4.5497) (-4.6736,-4.5626) (-4.6862,-4.5755) (-4.6987,-4.5883) (-4.7113,-4.6012) (-4.7238,-4.6141) (-4.7364,-4.6269) (-4.7490,-4.6398) (-4.7615,-4.6527) (-4.7741,-4.6655) (-4.7866,-4.6784) (-4.7992,-4.6912) (-4.8117,-4.7041) (-4.8243,-4.7170) (-4.8368,-4.7298) (-4.8494,-4.7427) (-4.8619,-4.7555) (-4.8745,-4.7683) (-4.8870,-4.7812) (-4.8996,-4.7940) (-4.9121,-4.8069) (-4.9247,-4.8197) (-4.9372,-4.8326) (-4.9498,-4.8454) (-4.9623,-4.8582) (-4.9749,-4.8711) (-4.9874,-4.8839) (-5.0000,-4.8967)};
\draw[main,thick] plot coordinates {(-2.0000,0.0000) (-1.9707,0.1153) (-1.9414,0.2265) (-1.9121,0.3329) (-1.8828,0.4340) (-1.8536,0.5294) (-1.8243,0.6188) (-1.7950,0.7019) (-1.7657,0.7786) (-1.7364,0.8490) (-1.7071,0.9130) (-1.6778,0.9709) (-1.6485,1.0228) (-1.6192,1.0689) (-1.5900,1.1096) (-1.5607,1.1450) (-1.5314,1.1755) (-1.5021,1.2015) (-1.4728,1.2232) (-1.4435,1.2410) (-1.4142,1.2552) (-1.3849,1.2660) (-1.3556,1.2737) (-1.3264,1.2787) (-1.2971,1.2811) (-1.2678,1.2812) (-1.2385,1.2792) (-1.2092,1.2753) (-1.1799,1.2696) (-1.1506,1.2624) (-1.1213,1.2539) (-1.0921,1.2440) (-1.0628,1.2331) (-1.0335,1.2212) (-1.0042,1.2083) (-0.9749,1.1947) (-0.9456,1.1804) (-0.9163,1.1655) (-0.8870,1.1501) (-0.8577,1.1342) (-0.8285,1.1178) (-0.7992,1.1012) (-0.7699,1.0842) (-0.7406,1.0671) (-0.7113,1.0497) (-0.6820,1.0322) (-0.6527,1.0145) (-0.6234,0.9968) (-0.5941,0.9791) (-0.5649,0.9614) (-0.5356,0.9437) (-0.5063,0.9260) (-0.4770,0.9085) (-0.4477,0.8910) (-0.4184,0.8737) (-0.3891,0.8566) (-0.3598,0.8396) (-0.3305,0.8229) (-0.3013,0.8063) (-0.2720,0.7901) (-0.2427,0.7741) (-0.2134,0.7584) (-0.1841,0.7431) (-0.1548,0.7281) (-0.1255,0.7135) (-0.0962,0.6992) (-0.0669,0.6854) (-0.0377,0.6720) (-0.0084,0.6590) (0.0209,0.6465) (0.0502,0.6346) (0.0795,0.6231) (0.1088,0.6122) (0.1381,0.6019) (0.1674,0.5921) (0.1967,0.5830) (0.2259,0.5745) (0.2552,0.5666) (0.2845,0.5595) (0.3138,0.5530) (0.3431,0.5473) (0.3724,0.5424) (0.4017,0.5382) (0.4310,0.5349) (0.4603,0.5324) (0.4895,0.5307) (0.5188,0.5299) (0.5481,0.5300) (0.5774,0.5311) (0.6067,0.5330) (0.6360,0.5360) (0.6653,0.5399) (0.6946,0.5448) (0.7238,0.5508) (0.7531,0.5578) (0.7824,0.5658) (0.8117,0.5749) (0.8410,0.5850) (0.8703,0.5963) (0.8996,0.6086) (0.9289,0.6220) (0.9582,0.6365) (0.9874,0.6520) (1.0167,0.6687) (1.0460,0.6864) (1.0753,0.7052) (1.1046,0.7250) (1.1339,0.7458) (1.1632,0.7677) (1.1925,0.7906) (1.2218,0.8144) (1.2510,0.8391) (1.2803,0.8648) (1.3096,0.8913) (1.3389,0.9187) (1.3682,0.9469) (1.3975,0.9758) (1.4268,1.0055) (1.4561,1.0358) (1.4854,1.0668) (1.5146,1.0984) (1.5439,1.1305) (1.5732,1.1632) (1.6025,1.1962) (1.6318,1.2297) (1.6611,1.2636) (1.6904,1.2978) (1.7197,1.3323) (1.7490,1.3671) (1.7782,1.4020) (1.8075,1.4371) (1.8368,1.4724) (1.8661,1.5078) (1.8954,1.5432) (1.9247,1.5787) (1.9540,1.6142) (1.9833,1.6497) (2.0126,1.6852) (2.0418,1.7206) (2.0711,1.7560) (2.1004,1.7913) (2.1297,1.8264) (2.1590,1.8615) (2.1883,1.8965) (2.2176,1.9313) (2.2469,1.9661) (2.2762,2.0006) (2.3054,2.0351) (2.3347,2.0694) (2.3640,2.1035) (2.3933,2.1375) (2.4226,2.1714) (2.4519,2.2051) (2.4812,2.2387) (2.5105,2.2722) (2.5397,2.3055) (2.5690,2.3387) (2.5983,2.3717) (2.6276,2.4046) (2.6569,2.4374) (2.6862,2.4701) (2.7155,2.5027) (2.7448,2.5352) (2.7741,2.5675) (2.8033,2.5998) (2.8326,2.6320) (2.8619,2.6640) (2.8912,2.6960) (2.9205,2.7279) (2.9498,2.7597) (2.9791,2.7915) (3.0084,2.8232) (3.0377,2.8548) (3.0669,2.8863) (3.0962,2.9178) (3.1255,2.9492) (3.1548,2.9805) (3.1841,3.0118) (3.2134,3.0430) (3.2427,3.0742) (3.2720,3.1054) (3.3013,3.1365) (3.3305,3.1675) (3.3598,3.1985) (3.3891,3.2295) (3.4184,3.2604) (3.4477,3.2913) (3.4770,3.3222) (3.5063,3.3530) (3.5356,3.3838) (3.5649,3.4145) (3.5941,3.4452) (3.6234,3.4759) (3.6527,3.5066) (3.6820,3.5372) (3.7113,3.5678) (3.7406,3.5984) (3.7699,3.6289) (3.7992,3.6595) (3.8285,3.6900) (3.8577,3.7205) (3.8870,3.7509) (3.9163,3.7814) (3.9456,3.8118) (3.9749,3.8422) (4.0042,3.8725) (4.0335,3.9029) (4.0628,3.9332) (4.0921,3.9636) (4.1213,3.9939) (4.1506,4.0242) (4.1799,4.0544) (4.2092,4.0847) (4.2385,4.1149) (4.2678,4.1452) (4.2971,4.1754) (4.3264,4.2056) (4.3556,4.2357) (4.3849,4.2659) (4.4142,4.2961) (4.4435,4.3262) (4.4728,4.3563) (4.5021,4.3864) (4.5314,4.4165) (4.5607,4.4466) (4.5900,4.4767) (4.6192,4.5068) (4.6485,4.5368) (4.6778,4.5669) (4.7071,4.5969) (4.7364,4.6269) (4.7657,4.6570) (4.7950,4.6870) (4.8243,4.7170) (4.8536,4.7469) (4.8828,4.7769) (4.9121,4.8069) (4.9414,4.8368) (4.9707,4.8668) (5.0000,4.8967)};
\fill[main](-2,0)circle(2pt);
\draw[structurecolor,dashed]plot coordinates{(-4.7958,-5.0000) (-4.7610,-4.9666) (-4.7261,-4.9331) (-4.6911,-4.8997) (-4.6562,-4.8662) (-4.6212,-4.8328) (-4.5862,-4.7993) (-4.5512,-4.7659) (-4.5162,-4.7324) (-4.4811,-4.6990) (-4.4461,-4.6656) (-4.4109,-4.6321) (-4.3758,-4.5987) (-4.3406,-4.5652) (-4.3055,-4.5318) (-4.2702,-4.4983) (-4.2350,-4.4649) (-4.1997,-4.4314) (-4.1644,-4.3980) (-4.1291,-4.3645) (-4.0937,-4.3311) (-4.0583,-4.2977) (-4.0229,-4.2642) (-3.9874,-4.2308) (-3.9519,-4.1973) (-3.9164,-4.1639) (-3.8808,-4.1304) (-3.8452,-4.0970) (-3.8095,-4.0635) (-3.7738,-4.0301) (-3.7381,-3.9967) (-3.7023,-3.9632) (-3.6665,-3.9298) (-3.6306,-3.8963) (-3.5947,-3.8629) (-3.5587,-3.8294) (-3.5227,-3.7960) (-3.4866,-3.7625) (-3.4505,-3.7291) (-3.4144,-3.6957) (-3.3781,-3.6622) (-3.3418,-3.6288) (-3.3055,-3.5953) (-3.2691,-3.5619) (-3.2326,-3.5284) (-3.1961,-3.4950) (-3.1595,-3.4615) (-3.1228,-3.4281) (-3.0860,-3.3946) (-3.0492,-3.3612) (-3.0123,-3.3278) (-2.9753,-3.2943) (-2.9382,-3.2609) (-2.9011,-3.2274) (-2.8638,-3.1940) (-2.8265,-3.1605) (-2.7890,-3.1271) (-2.7515,-3.0936) (-2.7138,-3.0602) (-2.6761,-3.0268) (-2.6382,-2.9933) (-2.6002,-2.9599) (-2.5620,-2.9264) (-2.5237,-2.8930) (-2.4853,-2.8595) (-2.4468,-2.8261) (-2.4081,-2.7926) (-2.3692,-2.7592) (-2.3302,-2.7258) (-2.2910,-2.6923) (-2.2516,-2.6589) (-2.2120,-2.6254) (-2.1722,-2.5920) (-2.1322,-2.5585) (-2.0919,-2.5251) (-2.0514,-2.4916) (-2.0107,-2.4582) (-1.9696,-2.4247) (-1.9283,-2.3913) (-1.8867,-2.3579) (-1.8447,-2.3244) (-1.8024,-2.2910) (-1.7597,-2.2575) (-1.7165,-2.2241) (-1.6730,-2.1906) (-1.6289,-2.1572) (-1.5844,-2.1237) (-1.5393,-2.0903) (-1.4935,-2.0569) (-1.4471,-2.0234) (-1.4000,-1.9900) (-1.3520,-1.9565) (-1.3032,-1.9231) (-1.2533,-1.8896) (-1.2023,-1.8562) (-1.1500,-1.8227) (-1.0962,-1.7893) (-1.0407,-1.7559) (-0.9832,-1.7224) (-0.9234,-1.6890) (-0.8607,-1.6555) (-0.7944,-1.6221) (-0.7237,-1.5886) (-0.6470,-1.5552) (-0.5619,-1.5217) (-0.4637,-1.4883) (-0.3414,-1.4548) (-0.1428,-1.4214)};\draw[structurecolor,dashed]plot coordinates{(-0.1428,1.4214) (-0.3414,1.4548) (-0.4637,1.4883) (-0.5619,1.5217) (-0.6470,1.5552) (-0.7237,1.5886) (-0.7944,1.6221) (-0.8607,1.6555) (-0.9234,1.6890) (-0.9832,1.7224) (-1.0407,1.7559) (-1.0962,1.7893) (-1.1500,1.8227) (-1.2023,1.8562) (-1.2533,1.8896) (-1.3032,1.9231) (-1.3520,1.9565) (-1.4000,1.9900) (-1.4471,2.0234) (-1.4935,2.0569) (-1.5393,2.0903) (-1.5844,2.1237) (-1.6289,2.1572) (-1.6730,2.1906) (-1.7165,2.2241) (-1.7597,2.2575) (-1.8024,2.2910) (-1.8447,2.3244) (-1.8867,2.3579) (-1.9283,2.3913) (-1.9696,2.4247) (-2.0107,2.4582) (-2.0514,2.4916) (-2.0919,2.5251) (-2.1322,2.5585) (-2.1722,2.5920) (-2.2120,2.6254) (-2.2516,2.6589) (-2.2910,2.6923) (-2.3302,2.7258) (-2.3692,2.7592) (-2.4081,2.7926) (-2.4468,2.8261) (-2.4853,2.8595) (-2.5237,2.8930) (-2.5620,2.9264) (-2.6002,2.9599) (-2.6382,2.9933) (-2.6761,3.0268) (-2.7138,3.0602) (-2.7515,3.0936) (-2.7890,3.1271) (-2.8265,3.1605) (-2.8638,3.1940) (-2.9011,3.2274) (-2.9382,3.2609) (-2.9753,3.2943) (-3.0123,3.3278) (-3.0492,3.3612) (-3.0860,3.3946) (-3.1228,3.4281) (-3.1595,3.4615) (-3.1961,3.4950) (-3.2326,3.5284) (-3.2691,3.5619) (-3.3055,3.5953) (-3.3418,3.6288) (-3.3781,3.6622) (-3.4144,3.6957) (-3.4505,3.7291) (-3.4866,3.7625) (-3.5227,3.7960) (-3.5587,3.8294) (-3.5947,3.8629) (-3.6306,3.8963) (-3.6665,3.9298) (-3.7023,3.9632) (-3.7381,3.9967) (-3.7738,4.0301) (-3.8095,4.0635) (-3.8452,4.0970) (-3.8808,4.1304) (-3.9164,4.1639) (-3.9519,4.1973) (-3.9874,4.2308) (-4.0229,4.2642) (-4.0583,4.2977) (-4.0937,4.3311) (-4.1291,4.3645) (-4.1644,4.3980) (-4.1997,4.4314) (-4.2350,4.4649) (-4.2702,4.4983) (-4.3055,4.5318) (-4.3406,4.5652) (-4.3758,4.5987) (-4.4109,4.6321) (-4.4461,4.6656) (-4.4811,4.6990) (-4.5162,4.7324) (-4.5512,4.7659) (-4.5862,4.7993) (-4.6212,4.8328) (-4.6562,4.8662) (-4.6911,4.8997) (-4.7261,4.9331) (-4.7610,4.9666) (-4.7958,5.0000)};\draw[structurecolor,dashed]plot coordinates{(4.7958,-5.0000) (4.7610,-4.9666) (4.7261,-4.9331) (4.6911,-4.8997) (4.6562,-4.8662) (4.6212,-4.8328) (4.5862,-4.7993) (4.5512,-4.7659) (4.5162,-4.7324) (4.4811,-4.6990) (4.4461,-4.6656) (4.4109,-4.6321) (4.3758,-4.5987) (4.3406,-4.5652) (4.3055,-4.5318) (4.2702,-4.4983) (4.2350,-4.4649) (4.1997,-4.4314) (4.1644,-4.3980) (4.1291,-4.3645) (4.0937,-4.3311) (4.0583,-4.2977) (4.0229,-4.2642) (3.9874,-4.2308) (3.9519,-4.1973) (3.9164,-4.1639) (3.8808,-4.1304) (3.8452,-4.0970) (3.8095,-4.0635) (3.7738,-4.0301) (3.7381,-3.9967) (3.7023,-3.9632) (3.6665,-3.9298) (3.6306,-3.8963) (3.5947,-3.8629) (3.5587,-3.8294) (3.5227,-3.7960) (3.4866,-3.7625) (3.4505,-3.7291) (3.4144,-3.6957) (3.3781,-3.6622) (3.3418,-3.6288) (3.3055,-3.5953) (3.2691,-3.5619) (3.2326,-3.5284) (3.1961,-3.4950) (3.1595,-3.4615) (3.1228,-3.4281) (3.0860,-3.3946) (3.0492,-3.3612) (3.0123,-3.3278) (2.9753,-3.2943) (2.9382,-3.2609) (2.9011,-3.2274) (2.8638,-3.1940) (2.8265,-3.1605) (2.7890,-3.1271) (2.7515,-3.0936) (2.7138,-3.0602) (2.6761,-3.0268) (2.6382,-2.9933) (2.6002,-2.9599) (2.5620,-2.9264) (2.5237,-2.8930) (2.4853,-2.8595) (2.4468,-2.8261) (2.4081,-2.7926) (2.3692,-2.7592) (2.3302,-2.7258) (2.2910,-2.6923) (2.2516,-2.6589) (2.2120,-2.6254) (2.1722,-2.5920) (2.1322,-2.5585) (2.0919,-2.5251) (2.0514,-2.4916) (2.0107,-2.4582) (1.9696,-2.4247) (1.9283,-2.3913) (1.8867,-2.3579) (1.8447,-2.3244) (1.8024,-2.2910) (1.7597,-2.2575) (1.7165,-2.2241) (1.6730,-2.1906) (1.6289,-2.1572) (1.5844,-2.1237) (1.5393,-2.0903) (1.4935,-2.0569) (1.4471,-2.0234) (1.4000,-1.9900) (1.3520,-1.9565) (1.3032,-1.9231) (1.2533,-1.8896) (1.2023,-1.8562) (1.1500,-1.8227) (1.0962,-1.7893) (1.0407,-1.7559) (0.9832,-1.7224) (0.9234,-1.6890) (0.8607,-1.6555) (0.7944,-1.6221) (0.7237,-1.5886) (0.6470,-1.5552) (0.5619,-1.5217) (0.4637,-1.4883) (0.3414,-1.4548) (0.1428,-1.4214)};\draw[structurecolor,dashed]plot coordinates{(0.1428,1.4214) (0.3414,1.4548) (0.4637,1.4883) (0.5619,1.5217) (0.6470,1.5552) (0.7237,1.5886) (0.7944,1.6221) (0.8607,1.6555) (0.9234,1.6890) (0.9832,1.7224) (1.0407,1.7559) (1.0962,1.7893) (1.1500,1.8227) (1.2023,1.8562) (1.2533,1.8896) (1.3032,1.9231) (1.3520,1.9565) (1.4000,1.9900) (1.4471,2.0234) (1.4935,2.0569) (1.5393,2.0903) (1.5844,2.1237) (1.6289,2.1572) (1.6730,2.1906) (1.7165,2.2241) (1.7597,2.2575) (1.8024,2.2910) (1.8447,2.3244) (1.8867,2.3579) (1.9283,2.3913) (1.9696,2.4247) (2.0107,2.4582) (2.0514,2.4916) (2.0919,2.5251) (2.1322,2.5585) (2.1722,2.5920) (2.2120,2.6254) (2.2516,2.6589) (2.2910,2.6923) (2.3302,2.7258) (2.3692,2.7592) (2.4081,2.7926) (2.4468,2.8261) (2.4853,2.8595) (2.5237,2.8930) (2.5620,2.9264) (2.6002,2.9599) (2.6382,2.9933) (2.6761,3.0268) (2.7138,3.0602) (2.7515,3.0936) (2.7890,3.1271) (2.8265,3.1605) (2.8638,3.1940) (2.9011,3.2274) (2.9382,3.2609) (2.9753,3.2943) (3.0123,3.3278) (3.0492,3.3612) (3.0860,3.3946) (3.1228,3.4281) (3.1595,3.4615) (3.1961,3.4950) (3.2326,3.5284) (3.2691,3.5619) (3.3055,3.5953) (3.3418,3.6288) (3.3781,3.6622) (3.4144,3.6957) (3.4505,3.7291) (3.4866,3.7625) (3.5227,3.7960) (3.5587,3.8294) (3.5947,3.8629) (3.6306,3.8963) (3.6665,3.9298) (3.7023,3.9632) (3.7381,3.9967) (3.7738,4.0301) (3.8095,4.0635) (3.8452,4.0970) (3.8808,4.1304) (3.9164,4.1639) (3.9519,4.1973) (3.9874,4.2308) (4.0229,4.2642) (4.0583,4.2977) (4.0937,4.3311) (4.1291,4.3645) (4.1644,4.3980) (4.1997,4.4314) (4.2350,4.4649) (4.2702,4.4983) (4.3055,4.5318) (4.3406,4.5652) (4.3758,4.5987) (4.4109,4.6321) (4.4461,4.6656) (4.4811,4.6990) (4.5162,4.7324) (4.5512,4.7659) (4.5862,4.7993) (4.6212,4.8328) (4.6562,4.8662) (4.6911,4.8997) (4.7261,4.9331) (4.7610,4.9666) (4.7958,5.0000)};\draw[structurecolor]plot coordinates{(-5.0000,-5.0000) (-4.9666,-4.9666) (-4.9331,-4.9331) (-4.8997,-4.8997) (-4.8662,-4.8662) (-4.8328,-4.8328) (-4.7993,-4.7993) (-4.7659,-4.7659) (-4.7324,-4.7324) (-4.6990,-4.6990) (-4.6656,-4.6656) (-4.6321,-4.6321) (-4.5987,-4.5987) (-4.5652,-4.5652) (-4.5318,-4.5318) (-4.4983,-4.4983) (-4.4649,-4.4649) (-4.4314,-4.4314) (-4.3980,-4.3980) (-4.3645,-4.3645) (-4.3311,-4.3311) (-4.2977,-4.2977) (-4.2642,-4.2642) (-4.2308,-4.2308) (-4.1973,-4.1973) (-4.1639,-4.1639) (-4.1304,-4.1304) (-4.0970,-4.0970) (-4.0635,-4.0635) (-4.0301,-4.0301) (-3.9967,-3.9967) (-3.9632,-3.9632) (-3.9298,-3.9298) (-3.8963,-3.8963) (-3.8629,-3.8629) (-3.8294,-3.8294) (-3.7960,-3.7960) (-3.7625,-3.7625) (-3.7291,-3.7291) (-3.6957,-3.6957) (-3.6622,-3.6622) (-3.6288,-3.6288) (-3.5953,-3.5953) (-3.5619,-3.5619) (-3.5284,-3.5284) (-3.4950,-3.4950) (-3.4615,-3.4615) (-3.4281,-3.4281) (-3.3946,-3.3946) (-3.3612,-3.3612) (-3.3278,-3.3278) (-3.2943,-3.2943) (-3.2609,-3.2609) (-3.2274,-3.2274) (-3.1940,-3.1940) (-3.1605,-3.1605) (-3.1271,-3.1271) (-3.0936,-3.0936) (-3.0602,-3.0602) (-3.0268,-3.0268) (-2.9933,-2.9933) (-2.9599,-2.9599) (-2.9264,-2.9264) (-2.8930,-2.8930) (-2.8595,-2.8595) (-2.8261,-2.8261) (-2.7926,-2.7926) (-2.7592,-2.7592) (-2.7258,-2.7258) (-2.6923,-2.6923) (-2.6589,-2.6589) (-2.6254,-2.6254) (-2.5920,-2.5920) (-2.5585,-2.5585) (-2.5251,-2.5251) (-2.4916,-2.4916) (-2.4582,-2.4582) (-2.4247,-2.4247) (-2.3913,-2.3913) (-2.3579,-2.3579) (-2.3244,-2.3244) (-2.2910,-2.2910) (-2.2575,-2.2575) (-2.2241,-2.2241) (-2.1906,-2.1906) (-2.1572,-2.1572) (-2.1237,-2.1237) (-2.0903,-2.0903) (-2.0569,-2.0569) (-2.0234,-2.0234) (-1.9900,-1.9900) (-1.9565,-1.9565) (-1.9231,-1.9231) (-1.8896,-1.8896) (-1.8562,-1.8562) (-1.8227,-1.8227) (-1.7893,-1.7893) (-1.7559,-1.7559) (-1.7224,-1.7224) (-1.6890,-1.6890) (-1.6555,-1.6555) (-1.6221,-1.6221) (-1.5886,-1.5886) (-1.5552,-1.5552) (-1.5217,-1.5217) (-1.4883,-1.4883) (-1.4548,-1.4548) (-1.4214,-1.4214) (-1.3880,-1.3880) (-1.3545,-1.3545) (-1.3211,-1.3211) (-1.2876,-1.2876) (-1.2542,-1.2542) (-1.2207,-1.2207) (-1.1873,-1.1873) (-1.1538,-1.1538) (-1.1204,-1.1204) (-1.0870,-1.0870) (-1.0535,-1.0535) (-1.0201,-1.0201) (-0.9866,-0.9866) (-0.9532,-0.9532) (-0.9197,-0.9197) (-0.8863,-0.8863) (-0.8528,-0.8528) (-0.8194,-0.8194) (-0.7860,-0.7860) (-0.7525,-0.7525) (-0.7191,-0.7191) (-0.6856,-0.6856) (-0.6522,-0.6522) (-0.6187,-0.6187) (-0.5853,-0.5853) (-0.5518,-0.5518) (-0.5184,-0.5184) (-0.4849,-0.4849) (-0.4515,-0.4515) (-0.4181,-0.4181) (-0.3846,-0.3846) (-0.3512,-0.3512) (-0.3177,-0.3177) (-0.2843,-0.2843) (-0.2508,-0.2508) (-0.2174,-0.2174) (-0.1839,-0.1839) (-0.1505,-0.1505) (-0.1171,-0.1171) (-0.0836,-0.0836) (-0.0502,-0.0502) (-0.0167,-0.0167) (-0.0167,0.0167) (-0.0502,0.0502) (-0.0836,0.0836) (-0.1171,0.1171) (-0.1505,0.1505) (-0.1839,0.1839) (-0.2174,0.2174) (-0.2508,0.2508) (-0.2843,0.2843) (-0.3177,0.3177) (-0.3512,0.3512) (-0.3846,0.3846) (-0.4181,0.4181) (-0.4515,0.4515) (-0.4849,0.4849) (-0.5184,0.5184) (-0.5518,0.5518) (-0.5853,0.5853) (-0.6187,0.6187) (-0.6522,0.6522) (-0.6856,0.6856) (-0.7191,0.7191) (-0.7525,0.7525) (-0.7860,0.7860) (-0.8194,0.8194) (-0.8528,0.8528) (-0.8863,0.8863) (-0.9197,0.9197) (-0.9532,0.9532) (-0.9866,0.9866) (-1.0201,1.0201) (-1.0535,1.0535) (-1.0870,1.0870) (-1.1204,1.1204) (-1.1538,1.1538) (-1.1873,1.1873) (-1.2207,1.2207) (-1.2542,1.2542) (-1.2876,1.2876) (-1.3211,1.3211) (-1.3545,1.3545) (-1.3880,1.3880) (-1.4214,1.4214) (-1.4548,1.4548) (-1.4883,1.4883) (-1.5217,1.5217) (-1.5552,1.5552) (-1.5886,1.5886) (-1.6221,1.6221) (-1.6555,1.6555) (-1.6890,1.6890) (-1.7224,1.7224) (-1.7559,1.7559) (-1.7893,1.7893) (-1.8227,1.8227) (-1.8562,1.8562) (-1.8896,1.8896) (-1.9231,1.9231) (-1.9565,1.9565) (-1.9900,1.9900) (-2.0234,2.0234) (-2.0569,2.0569) (-2.0903,2.0903) (-2.1237,2.1237) (-2.1572,2.1572) (-2.1906,2.1906) (-2.2241,2.2241) (-2.2575,2.2575) (-2.2910,2.2910) (-2.3244,2.3244) (-2.3579,2.3579) (-2.3913,2.3913) (-2.4247,2.4247) (-2.4582,2.4582) (-2.4916,2.4916) (-2.5251,2.5251) (-2.5585,2.5585) (-2.5920,2.5920) (-2.6254,2.6254) (-2.6589,2.6589) (-2.6923,2.6923) (-2.7258,2.7258) (-2.7592,2.7592) (-2.7926,2.7926) (-2.8261,2.8261) (-2.8595,2.8595) (-2.8930,2.8930) (-2.9264,2.9264) (-2.9599,2.9599) (-2.9933,2.9933) (-3.0268,3.0268) (-3.0602,3.0602) (-3.0936,3.0936) (-3.1271,3.1271) (-3.1605,3.1605) (-3.1940,3.1940) (-3.2274,3.2274) (-3.2609,3.2609) (-3.2943,3.2943) (-3.3278,3.3278) (-3.3612,3.3612) (-3.3946,3.3946) (-3.4281,3.4281) (-3.4615,3.4615) (-3.4950,3.4950) (-3.5284,3.5284) (-3.5619,3.5619) (-3.5953,3.5953) (-3.6288,3.6288) (-3.6622,3.6622) (-3.6957,3.6957) (-3.7291,3.7291) (-3.7625,3.7625) (-3.7960,3.7960) (-3.8294,3.8294) (-3.8629,3.8629) (-3.8963,3.8963) (-3.9298,3.9298) (-3.9632,3.9632) (-3.9967,3.9967) (-4.0301,4.0301) (-4.0635,4.0635) (-4.0970,4.0970) (-4.1304,4.1304) (-4.1639,4.1639) (-4.1973,4.1973) (-4.2308,4.2308) (-4.2642,4.2642) (-4.2977,4.2977) (-4.3311,4.3311) (-4.3645,4.3645) (-4.3980,4.3980) (-4.4314,4.4314) (-4.4649,4.4649) (-4.4983,4.4983) (-4.5318,4.5318) (-4.5652,4.5652) (-4.5987,4.5987) (-4.6321,4.6321) (-4.6656,4.6656) (-4.6990,4.6990) (-4.7324,4.7324) (-4.7659,4.7659) (-4.7993,4.7993) (-4.8328,4.8328) (-4.8662,4.8662) (-4.8997,4.8997) (-4.9331,4.9331) (-4.9666,4.9666) (-5.0000,5.0000)};\draw[structurecolor]plot coordinates{(5.0000,-5.0000) (4.9666,-4.9666) (4.9331,-4.9331) (4.8997,-4.8997) (4.8662,-4.8662) (4.8328,-4.8328) (4.7993,-4.7993) (4.7659,-4.7659) (4.7324,-4.7324) (4.6990,-4.6990) (4.6656,-4.6656) (4.6321,-4.6321) (4.5987,-4.5987) (4.5652,-4.5652) (4.5318,-4.5318) (4.4983,-4.4983) (4.4649,-4.4649) (4.4314,-4.4314) (4.3980,-4.3980) (4.3645,-4.3645) (4.3311,-4.3311) (4.2977,-4.2977) (4.2642,-4.2642) (4.2308,-4.2308) (4.1973,-4.1973) (4.1639,-4.1639) (4.1304,-4.1304) (4.0970,-4.0970) (4.0635,-4.0635) (4.0301,-4.0301) (3.9967,-3.9967) (3.9632,-3.9632) (3.9298,-3.9298) (3.8963,-3.8963) (3.8629,-3.8629) (3.8294,-3.8294) (3.7960,-3.7960) (3.7625,-3.7625) (3.7291,-3.7291) (3.6957,-3.6957) (3.6622,-3.6622) (3.6288,-3.6288) (3.5953,-3.5953) (3.5619,-3.5619) (3.5284,-3.5284) (3.4950,-3.4950) (3.4615,-3.4615) (3.4281,-3.4281) (3.3946,-3.3946) (3.3612,-3.3612) (3.3278,-3.3278) (3.2943,-3.2943) (3.2609,-3.2609) (3.2274,-3.2274) (3.1940,-3.1940) (3.1605,-3.1605) (3.1271,-3.1271) (3.0936,-3.0936) (3.0602,-3.0602) (3.0268,-3.0268) (2.9933,-2.9933) (2.9599,-2.9599) (2.9264,-2.9264) (2.8930,-2.8930) (2.8595,-2.8595) (2.8261,-2.8261) (2.7926,-2.7926) (2.7592,-2.7592) (2.7258,-2.7258) (2.6923,-2.6923) (2.6589,-2.6589) (2.6254,-2.6254) (2.5920,-2.5920) (2.5585,-2.5585) (2.5251,-2.5251) (2.4916,-2.4916) (2.4582,-2.4582) (2.4247,-2.4247) (2.3913,-2.3913) (2.3579,-2.3579) (2.3244,-2.3244) (2.2910,-2.2910) (2.2575,-2.2575) (2.2241,-2.2241) (2.1906,-2.1906) (2.1572,-2.1572) (2.1237,-2.1237) (2.0903,-2.0903) (2.0569,-2.0569) (2.0234,-2.0234) (1.9900,-1.9900) (1.9565,-1.9565) (1.9231,-1.9231) (1.8896,-1.8896) (1.8562,-1.8562) (1.8227,-1.8227) (1.7893,-1.7893) (1.7559,-1.7559) (1.7224,-1.7224) (1.6890,-1.6890) (1.6555,-1.6555) (1.6221,-1.6221) (1.5886,-1.5886) (1.5552,-1.5552) (1.5217,-1.5217) (1.4883,-1.4883) (1.4548,-1.4548) (1.4214,-1.4214) (1.3880,-1.3880) (1.3545,-1.3545) (1.3211,-1.3211) (1.2876,-1.2876) (1.2542,-1.2542) (1.2207,-1.2207) (1.1873,-1.1873) (1.1538,-1.1538) (1.1204,-1.1204) (1.0870,-1.0870) (1.0535,-1.0535) (1.0201,-1.0201) (0.9866,-0.9866) (0.9532,-0.9532) (0.9197,-0.9197) (0.8863,-0.8863) (0.8528,-0.8528) (0.8194,-0.8194) (0.7860,-0.7860) (0.7525,-0.7525) (0.7191,-0.7191) (0.6856,-0.6856) (0.6522,-0.6522) (0.6187,-0.6187) (0.5853,-0.5853) (0.5518,-0.5518) (0.5184,-0.5184) (0.4849,-0.4849) (0.4515,-0.4515) (0.4181,-0.4181) (0.3846,-0.3846) (0.3512,-0.3512) (0.3177,-0.3177) (0.2843,-0.2843) (0.2508,-0.2508) (0.2174,-0.2174) (0.1839,-0.1839) (0.1505,-0.1505) (0.1171,-0.1171) (0.0836,-0.0836) (0.0502,-0.0502) (0.0167,-0.0167) (0.0167,0.0167) (0.0502,0.0502) (0.0836,0.0836) (0.1171,0.1171) (0.1505,0.1505) (0.1839,0.1839) (0.2174,0.2174) (0.2508,0.2508) (0.2843,0.2843) (0.3177,0.3177) (0.3512,0.3512) (0.3846,0.3846) (0.4181,0.4181) (0.4515,0.4515) (0.4849,0.4849) (0.5184,0.5184) (0.5518,0.5518) (0.5853,0.5853) (0.6187,0.6187) (0.6522,0.6522) (0.6856,0.6856) (0.7191,0.7191) (0.7525,0.7525) (0.7860,0.7860) (0.8194,0.8194) (0.8528,0.8528) (0.8863,0.8863) (0.9197,0.9197) (0.9532,0.9532) (0.9866,0.9866) (1.0201,1.0201) (1.0535,1.0535) (1.0870,1.0870) (1.1204,1.1204) (1.1538,1.1538) (1.1873,1.1873) (1.2207,1.2207) (1.2542,1.2542) (1.2876,1.2876) (1.3211,1.3211) (1.3545,1.3545) (1.3880,1.3880) (1.4214,1.4214) (1.4548,1.4548) (1.4883,1.4883) (1.5217,1.5217) (1.5552,1.5552) (1.5886,1.5886) (1.6221,1.6221) (1.6555,1.6555) (1.6890,1.6890) (1.7224,1.7224) (1.7559,1.7559) (1.7893,1.7893) (1.8227,1.8227) (1.8562,1.8562) (1.8896,1.8896) (1.9231,1.9231) (1.9565,1.9565) (1.9900,1.9900) (2.0234,2.0234) (2.0569,2.0569) (2.0903,2.0903) (2.1237,2.1237) (2.1572,2.1572) (2.1906,2.1906) (2.2241,2.2241) (2.2575,2.2575) (2.2910,2.2910) (2.3244,2.3244) (2.3579,2.3579) (2.3913,2.3913) (2.4247,2.4247) (2.4582,2.4582) (2.4916,2.4916) (2.5251,2.5251) (2.5585,2.5585) (2.5920,2.5920) (2.6254,2.6254) (2.6589,2.6589) (2.6923,2.6923) (2.7258,2.7258) (2.7592,2.7592) (2.7926,2.7926) (2.8261,2.8261) (2.8595,2.8595) (2.8930,2.8930) (2.9264,2.9264) (2.9599,2.9599) (2.9933,2.9933) (3.0268,3.0268) (3.0602,3.0602) (3.0936,3.0936) (3.1271,3.1271) (3.1605,3.1605) (3.1940,3.1940) (3.2274,3.2274) (3.2609,3.2609) (3.2943,3.2943) (3.3278,3.3278) (3.3612,3.3612) (3.3946,3.3946) (3.4281,3.4281) (3.4615,3.4615) (3.4950,3.4950) (3.5284,3.5284) (3.5619,3.5619) (3.5953,3.5953) (3.6288,3.6288) (3.6622,3.6622) (3.6957,3.6957) (3.7291,3.7291) (3.7625,3.7625) (3.7960,3.7960) (3.8294,3.8294) (3.8629,3.8629) (3.8963,3.8963) (3.9298,3.9298) (3.9632,3.9632) (3.9967,3.9967) (4.0301,4.0301) (4.0635,4.0635) (4.0970,4.0970) (4.1304,4.1304) (4.1639,4.1639) (4.1973,4.1973) (4.2308,4.2308) (4.2642,4.2642) (4.2977,4.2977) (4.3311,4.3311) (4.3645,4.3645) (4.3980,4.3980) (4.4314,4.4314) (4.4649,4.4649) (4.4983,4.4983) (4.5318,4.5318) (4.5652,4.5652) (4.5987,4.5987) (4.6321,4.6321) (4.6656,4.6656) (4.6990,4.6990) (4.7324,4.7324) (4.7659,4.7659) (4.7993,4.7993) (4.8328,4.8328) (4.8662,4.8662) (4.8997,4.8997) (4.9331,4.9331) (4.9666,4.9666) (5.0000,5.0000)};\draw[structurecolor,dashed]plot coordinates{(-4.9713,-4.7659) (-4.9392,-4.7324) (-4.9072,-4.6990) (-4.8752,-4.6656) (-4.8432,-4.6321) (-4.8112,-4.5987) (-4.7792,-4.5652) (-4.7473,-4.5318) (-4.7154,-4.4983) (-4.6835,-4.4649) (-4.6516,-4.4314) (-4.6198,-4.3980) (-4.5879,-4.3645) (-4.5561,-4.3311) (-4.5244,-4.2977) (-4.4926,-4.2642) (-4.4609,-4.2308) (-4.4292,-4.1973) (-4.3975,-4.1639) (-4.3658,-4.1304) (-4.3342,-4.0970) (-4.3026,-4.0635) (-4.2710,-4.0301) (-4.2395,-3.9967) (-4.2080,-3.9632) (-4.1765,-3.9298) (-4.1450,-3.8963) (-4.1136,-3.8629) (-4.0822,-3.8294) (-4.0509,-3.7960) (-4.0195,-3.7625) (-3.9883,-3.7291) (-3.9570,-3.6957) (-3.9258,-3.6622) (-3.8946,-3.6288) (-3.8635,-3.5953) (-3.8324,-3.5619) (-3.8013,-3.5284) (-3.7703,-3.4950) (-3.7393,-3.4615) (-3.7083,-3.4281) (-3.6775,-3.3946) (-3.6466,-3.3612) (-3.6158,-3.3278) (-3.5850,-3.2943) (-3.5543,-3.2609) (-3.5237,-3.2274) (-3.4931,-3.1940) (-3.4625,-3.1605) (-3.4320,-3.1271) (-3.4016,-3.0936) (-3.3712,-3.0602) (-3.3408,-3.0268) (-3.3106,-2.9933) (-3.2804,-2.9599) (-3.2502,-2.9264) (-3.2201,-2.8930) (-3.1901,-2.8595) (-3.1602,-2.8261) (-3.1303,-2.7926) (-3.1005,-2.7592) (-3.0708,-2.7258) (-3.0411,-2.6923) (-3.0116,-2.6589) (-2.9821,-2.6254) (-2.9527,-2.5920) (-2.9234,-2.5585) (-2.8941,-2.5251) (-2.8650,-2.4916) (-2.8360,-2.4582) (-2.8070,-2.4247) (-2.7782,-2.3913) (-2.7495,-2.3579) (-2.7208,-2.3244) (-2.6923,-2.2910) (-2.6639,-2.2575) (-2.6356,-2.2241) (-2.6075,-2.1906) (-2.5794,-2.1572) (-2.5515,-2.1237) (-2.5238,-2.0903) (-2.4961,-2.0569) (-2.4686,-2.0234) (-2.4413,-1.9900) (-2.4141,-1.9565) (-2.3871,-1.9231) (-2.3602,-1.8896) (-2.3335,-1.8562) (-2.3070,-1.8227) (-2.2807,-1.7893) (-2.2546,-1.7559) (-2.2286,-1.7224) (-2.2029,-1.6890) (-2.1773,-1.6555) (-2.1520,-1.6221) (-2.1269,-1.5886) (-2.1020,-1.5552) (-2.0774,-1.5217) (-2.0531,-1.4883) (-2.0289,-1.4548) (-2.0051,-1.4214) (-1.9815,-1.3880) (-1.9582,-1.3545) (-1.9353,-1.3211) (-1.9126,-1.2876) (-1.8902,-1.2542) (-1.8682,-1.2207) (-1.8465,-1.1873) (-1.8252,-1.1538) (-1.8042,-1.1204) (-1.7837,-1.0870) (-1.7635,-1.0535) (-1.7437,-1.0201) (-1.7244,-0.9866) (-1.7054,-0.9532) (-1.6870,-0.9197) (-1.6690,-0.8863) (-1.6515,-0.8528) (-1.6344,-0.8194) (-1.6179,-0.7860) (-1.6020,-0.7525) (-1.5865,-0.7191) (-1.5716,-0.6856) (-1.5573,-0.6522) (-1.5436,-0.6187) (-1.5305,-0.5853) (-1.5181,-0.5518) (-1.5062,-0.5184) (-1.4951,-0.4849) (-1.4845,-0.4515) (-1.4747,-0.4181) (-1.4656,-0.3846) (-1.4572,-0.3512) (-1.4495,-0.3177) (-1.4425,-0.2843) (-1.4363,-0.2508) (-1.4308,-0.2174) (-1.4261,-0.1839) (-1.4222,-0.1505) (-1.4190,-0.1171) (-1.4167,-0.0836) (-1.4151,-0.0502) (-1.4143,-0.0167) (-1.4143,0.0167) (-1.4151,0.0502) (-1.4167,0.0836) (-1.4190,0.1171) (-1.4222,0.1505) (-1.4261,0.1839) (-1.4308,0.2174) (-1.4363,0.2508) (-1.4425,0.2843) (-1.4495,0.3177) (-1.4572,0.3512) (-1.4656,0.3846) (-1.4747,0.4181) (-1.4845,0.4515) (-1.4951,0.4849) (-1.5062,0.5184) (-1.5181,0.5518) (-1.5305,0.5853) (-1.5436,0.6187) (-1.5573,0.6522) (-1.5716,0.6856) (-1.5865,0.7191) (-1.6020,0.7525) (-1.6179,0.7860) (-1.6344,0.8194) (-1.6515,0.8528) (-1.6690,0.8863) (-1.6870,0.9197) (-1.7054,0.9532) (-1.7244,0.9866) (-1.7437,1.0201) (-1.7635,1.0535) (-1.7837,1.0870) (-1.8042,1.1204) (-1.8252,1.1538) (-1.8465,1.1873) (-1.8682,1.2207) (-1.8902,1.2542) (-1.9126,1.2876) (-1.9353,1.3211) (-1.9582,1.3545) (-1.9815,1.3880) (-2.0051,1.4214) (-2.0289,1.4548) (-2.0531,1.4883) (-2.0774,1.5217) (-2.1020,1.5552) (-2.1269,1.5886) (-2.1520,1.6221) (-2.1773,1.6555) (-2.2029,1.6890) (-2.2286,1.7224) (-2.2546,1.7559) (-2.2807,1.7893) (-2.3070,1.8227) (-2.3335,1.8562) (-2.3602,1.8896) (-2.3871,1.9231) (-2.4141,1.9565) (-2.4413,1.9900) (-2.4686,2.0234) (-2.4961,2.0569) (-2.5238,2.0903) (-2.5515,2.1237) (-2.5794,2.1572) (-2.6075,2.1906) (-2.6356,2.2241) (-2.6639,2.2575) (-2.6923,2.2910) (-2.7208,2.3244) (-2.7495,2.3579) (-2.7782,2.3913) (-2.8070,2.4247) (-2.8360,2.4582) (-2.8650,2.4916) (-2.8941,2.5251) (-2.9234,2.5585) (-2.9527,2.5920) (-2.9821,2.6254) (-3.0116,2.6589) (-3.0411,2.6923) (-3.0708,2.7258) (-3.1005,2.7592) (-3.1303,2.7926) (-3.1602,2.8261) (-3.1901,2.8595) (-3.2201,2.8930) (-3.2502,2.9264) (-3.2804,2.9599) (-3.3106,2.9933) (-3.3408,3.0268) (-3.3712,3.0602) (-3.4016,3.0936) (-3.4320,3.1271) (-3.4625,3.1605) (-3.4931,3.1940) (-3.5237,3.2274) (-3.5543,3.2609) (-3.5850,3.2943) (-3.6158,3.3278) (-3.6466,3.3612) (-3.6775,3.3946) (-3.7083,3.4281) (-3.7393,3.4615) (-3.7703,3.4950) (-3.8013,3.5284) (-3.8324,3.5619) (-3.8635,3.5953) (-3.8946,3.6288) (-3.9258,3.6622) (-3.9570,3.6957) (-3.9883,3.7291) (-4.0195,3.7625) (-4.0509,3.7960) (-4.0822,3.8294) (-4.1136,3.8629) (-4.1450,3.8963) (-4.1765,3.9298) (-4.2080,3.9632) (-4.2395,3.9967) (-4.2710,4.0301) (-4.3026,4.0635) (-4.3342,4.0970) (-4.3658,4.1304) (-4.3975,4.1639) (-4.4292,4.1973) (-4.4609,4.2308) (-4.4926,4.2642) (-4.5244,4.2977) (-4.5561,4.3311) (-4.5879,4.3645) (-4.6198,4.3980) (-4.6516,4.4314) (-4.6835,4.4649) (-4.7154,4.4983) (-4.7473,4.5318) (-4.7792,4.5652) (-4.8112,4.5987) (-4.8432,4.6321) (-4.8752,4.6656) (-4.9072,4.6990) (-4.9392,4.7324) (-4.9713,4.7659)};\draw[structurecolor,dashed]plot coordinates{(4.9713,-4.7659) (4.9392,-4.7324) (4.9072,-4.6990) (4.8752,-4.6656) (4.8432,-4.6321) (4.8112,-4.5987) (4.7792,-4.5652) (4.7473,-4.5318) (4.7154,-4.4983) (4.6835,-4.4649) (4.6516,-4.4314) (4.6198,-4.3980) (4.5879,-4.3645) (4.5561,-4.3311) (4.5244,-4.2977) (4.4926,-4.2642) (4.4609,-4.2308) (4.4292,-4.1973) (4.3975,-4.1639) (4.3658,-4.1304) (4.3342,-4.0970) (4.3026,-4.0635) (4.2710,-4.0301) (4.2395,-3.9967) (4.2080,-3.9632) (4.1765,-3.9298) (4.1450,-3.8963) (4.1136,-3.8629) (4.0822,-3.8294) (4.0509,-3.7960) (4.0195,-3.7625) (3.9883,-3.7291) (3.9570,-3.6957) (3.9258,-3.6622) (3.8946,-3.6288) (3.8635,-3.5953) (3.8324,-3.5619) (3.8013,-3.5284) (3.7703,-3.4950) (3.7393,-3.4615) (3.7083,-3.4281) (3.6775,-3.3946) (3.6466,-3.3612) (3.6158,-3.3278) (3.5850,-3.2943) (3.5543,-3.2609) (3.5237,-3.2274) (3.4931,-3.1940) (3.4625,-3.1605) (3.4320,-3.1271) (3.4016,-3.0936) (3.3712,-3.0602) (3.3408,-3.0268) (3.3106,-2.9933) (3.2804,-2.9599) (3.2502,-2.9264) (3.2201,-2.8930) (3.1901,-2.8595) (3.1602,-2.8261) (3.1303,-2.7926) (3.1005,-2.7592) (3.0708,-2.7258) (3.0411,-2.6923) (3.0116,-2.6589) (2.9821,-2.6254) (2.9527,-2.5920) (2.9234,-2.5585) (2.8941,-2.5251) (2.8650,-2.4916) (2.8360,-2.4582) (2.8070,-2.4247) (2.7782,-2.3913) (2.7495,-2.3579) (2.7208,-2.3244) (2.6923,-2.2910) (2.6639,-2.2575) (2.6356,-2.2241) (2.6075,-2.1906) (2.5794,-2.1572) (2.5515,-2.1237) (2.5238,-2.0903) (2.4961,-2.0569) (2.4686,-2.0234) (2.4413,-1.9900) (2.4141,-1.9565) (2.3871,-1.9231) (2.3602,-1.8896) (2.3335,-1.8562) (2.3070,-1.8227) (2.2807,-1.7893) (2.2546,-1.7559) (2.2286,-1.7224) (2.2029,-1.6890) (2.1773,-1.6555) (2.1520,-1.6221) (2.1269,-1.5886) (2.1020,-1.5552) (2.0774,-1.5217) (2.0531,-1.4883) (2.0289,-1.4548) (2.0051,-1.4214) (1.9815,-1.3880) (1.9582,-1.3545) (1.9353,-1.3211) (1.9126,-1.2876) (1.8902,-1.2542) (1.8682,-1.2207) (1.8465,-1.1873) (1.8252,-1.1538) (1.8042,-1.1204) (1.7837,-1.0870) (1.7635,-1.0535) (1.7437,-1.0201) (1.7244,-0.9866) (1.7054,-0.9532) (1.6870,-0.9197) (1.6690,-0.8863) (1.6515,-0.8528) (1.6344,-0.8194) (1.6179,-0.7860) (1.6020,-0.7525) (1.5865,-0.7191) (1.5716,-0.6856) (1.5573,-0.6522) (1.5436,-0.6187) (1.5305,-0.5853) (1.5181,-0.5518) (1.5062,-0.5184) (1.4951,-0.4849) (1.4845,-0.4515) (1.4747,-0.4181) (1.4656,-0.3846) (1.4572,-0.3512) (1.4495,-0.3177) (1.4425,-0.2843) (1.4363,-0.2508) (1.4308,-0.2174) (1.4261,-0.1839) (1.4222,-0.1505) (1.4190,-0.1171) (1.4167,-0.0836) (1.4151,-0.0502) (1.4143,-0.0167) (1.4143,0.0167) (1.4151,0.0502) (1.4167,0.0836) (1.4190,0.1171) (1.4222,0.1505) (1.4261,0.1839) (1.4308,0.2174) (1.4363,0.2508) (1.4425,0.2843) (1.4495,0.3177) (1.4572,0.3512) (1.4656,0.3846) (1.4747,0.4181) (1.4845,0.4515) (1.4951,0.4849) (1.5062,0.5184) (1.5181,0.5518) (1.5305,0.5853) (1.5436,0.6187) (1.5573,0.6522) (1.5716,0.6856) (1.5865,0.7191) (1.6020,0.7525) (1.6179,0.7860) (1.6344,0.8194) (1.6515,0.8528) (1.6690,0.8863) (1.6870,0.9197) (1.7054,0.9532) (1.7244,0.9866) (1.7437,1.0201) (1.7635,1.0535) (1.7837,1.0870) (1.8042,1.1204) (1.8252,1.1538) (1.8465,1.1873) (1.8682,1.2207) (1.8902,1.2542) (1.9126,1.2876) (1.9353,1.3211) (1.9582,1.3545) (1.9815,1.3880) (2.0051,1.4214) (2.0289,1.4548) (2.0531,1.4883) (2.0774,1.5217) (2.1020,1.5552) (2.1269,1.5886) (2.1520,1.6221) (2.1773,1.6555) (2.2029,1.6890) (2.2286,1.7224) (2.2546,1.7559) (2.2807,1.7893) (2.3070,1.8227) (2.3335,1.8562) (2.3602,1.8896) (2.3871,1.9231) (2.4141,1.9565) (2.4413,1.9900) (2.4686,2.0234) (2.4961,2.0569) (2.5238,2.0903) (2.5515,2.1237) (2.5794,2.1572) (2.6075,2.1906) (2.6356,2.2241) (2.6639,2.2575) (2.6923,2.2910) (2.7208,2.3244) (2.7495,2.3579) (2.7782,2.3913) (2.8070,2.4247) (2.8360,2.4582) (2.8650,2.4916) (2.8941,2.5251) (2.9234,2.5585) (2.9527,2.5920) (2.9821,2.6254) (3.0116,2.6589) (3.0411,2.6923) (3.0708,2.7258) (3.1005,2.7592) (3.1303,2.7926) (3.1602,2.8261) (3.1901,2.8595) (3.2201,2.8930) (3.2502,2.9264) (3.2804,2.9599) (3.3106,2.9933) (3.3408,3.0268) (3.3712,3.0602) (3.4016,3.0936) (3.4320,3.1271) (3.4625,3.1605) (3.4931,3.1940) (3.5237,3.2274) (3.5543,3.2609) (3.5850,3.2943) (3.6158,3.3278) (3.6466,3.3612) (3.6775,3.3946) (3.7083,3.4281) (3.7393,3.4615) (3.7703,3.4950) (3.8013,3.5284) (3.8324,3.5619) (3.8635,3.5953) (3.8946,3.6288) (3.9258,3.6622) (3.9570,3.6957) (3.9883,3.7291) (4.0195,3.7625) (4.0509,3.7960) (4.0822,3.8294) (4.1136,3.8629) (4.1450,3.8963) (4.1765,3.9298) (4.2080,3.9632) (4.2395,3.9967) (4.2710,4.0301) (4.3026,4.0635) (4.3342,4.0970) (4.3658,4.1304) (4.3975,4.1639) (4.4292,4.1973) (4.4609,4.2308) (4.4926,4.2642) (4.5244,4.2977) (4.5561,4.3311) (4.5879,4.3645) (4.6198,4.3980) (4.6516,4.4314) (4.6835,4.4649) (4.7154,4.4983) (4.7473,4.5318) (4.7792,4.5652) (4.8112,4.5987) (4.8432,4.6321) (4.8752,4.6656) (4.9072,4.6990) (4.9392,4.7324) (4.9713,4.7659)};\draw(-5,-.08)--(-5,.08);\draw(-4,-.08)--(-4,.08);\node[below,font=\scriptsize,fill=white,inner sep=.2pt]at(-4,-.08){$ -4 $};\draw(-3,-.08)--(-3,.08);\draw(-2,-.08)--(-2,.08);\node[below,font=\scriptsize,fill=white,inner sep=.2pt]at(-2,-.08){$ -2 $};\draw(-1,-.08)--(-1,.08);\draw(1,-.08)--(1,.08);\draw(2,-.08)--(2,.08);\node[below,font=\scriptsize,fill=white,inner sep=.2pt]at(2,-.08){$ 2 $};\draw(3,-.08)--(3,.08);\draw(4,-.08)--(4,.08);\node[below,font=\scriptsize,fill=white,inner sep=.2pt]at(4,-.08){$ 4 $};\draw(5,-.08)--(5,.08);\node[below,font=\scriptsize,fill=white,inner sep=.2pt]at(5,-.08){$ 5 $};\draw(-.08,-5)--(.08,-5);\draw(-.08,-4)--(.08,-4);\node[left,font=\scriptsize,fill=white,inner sep=.2pt]at(-.08,-4){$ -4 $};\draw(-.08,-3)--(.08,-3);\draw(-.08,-2)--(.08,-2);\node[left,font=\scriptsize,fill=white,inner sep=.2pt]at(-.08,-2){$ -2 $};\draw(-.08,-1)--(.08,-1);\draw(-.08,1)--(.08,1);\draw(-.08,2)--(.08,2);\node[left,font=\scriptsize,fill=white,inner sep=.2pt]at(-.08,2){$ 2 $};\draw(-.08,3)--(.08,3);\draw(-.08,4)--(.08,4);\node[left,font=\scriptsize,fill=white,inner sep=.2pt]at(-.08,4){$ 4 $};\draw(-.08,5)--(.08,5);\end{tikzpicture}\end{center}
}

\newcommand{\ODEExamFinalPhase}{%
\begin{center}\begin{tikzpicture}[x=.8cm,y=.7cm]\draw[->](-1,-2)--(-1,3.5)node[above]{$x$};\draw[->](-1,-1.35)--(-1,-1.65);\draw[->](-1,-0.65)--(-1,-0.35);\draw[->](-1,1.15)--(-1,0.85);\draw[->](-1,2.35)--(-1,2.65);\fill(-1,-1)circle(1.8pt);\node[left]at(-1,-1){$ -1 $};\draw[dashed](0,-1)--(5,-1);\fill(-1,0)circle(1.8pt);\node[left]at(-1,0){$ 0 $};\draw[dashed](0,0)--(5,0);\fill(-1,2)circle(1.8pt);\node[left]at(-1,2){$ 2 $};\draw[dashed](0,2)--(5,2);\draw[->](0,-2)--(5.4,-2)node[right]{$t$};\draw[->](0,-2)--(0,3.3)node[above]{$x$};\draw[main,thick]plot coordinates{(2.0000,-1.5000) (1.9888,-1.4719) (1.9777,-1.4460) (1.9665,-1.4221) (1.9553,-1.4000) (1.9441,-1.3795) (1.9330,-1.3605) (1.9218,-1.3427) (1.9106,-1.3260) (1.8994,-1.3105) (1.8883,-1.2959) (1.8771,-1.2822) (1.8659,-1.2694) (1.8547,-1.2572) (1.8436,-1.2458) (1.8324,-1.2350) (1.8212,-1.2248) (1.8101,-1.2151) (1.7989,-1.2060) (1.7877,-1.1973) (1.7765,-1.1891) (1.7654,-1.1813) (1.7542,-1.1739) (1.7430,-1.1668) (1.7318,-1.1601) (1.7207,-1.1537) (1.7095,-1.1476) (1.6983,-1.1418) (1.6872,-1.1362) (1.6760,-1.1309) (1.6648,-1.1259) (1.6536,-1.1210) (1.6425,-1.1164) (1.6313,-1.1120) (1.6201,-1.1077) (1.6089,-1.1037) (1.5978,-1.0998) (1.5866,-1.0961) (1.5754,-1.0925) (1.5642,-1.0891) (1.5531,-1.0858) (1.5419,-1.0827) (1.5307,-1.0796) (1.5196,-1.0767) (1.5084,-1.0740) (1.4972,-1.0713) (1.4860,-1.0687) (1.4749,-1.0662) (1.4637,-1.0639) (1.4525,-1.0616) (1.4413,-1.0594) (1.4302,-1.0573) (1.4190,-1.0553) (1.4078,-1.0533) (1.3966,-1.0514) (1.3855,-1.0496) (1.3743,-1.0479) (1.3631,-1.0462) (1.3520,-1.0446) (1.3408,-1.0430) (1.3296,-1.0415) (1.3184,-1.0401) (1.3073,-1.0387) (1.2961,-1.0374) (1.2849,-1.0361) (1.2737,-1.0348) (1.2626,-1.0336) (1.2514,-1.0325) (1.2402,-1.0314) (1.2291,-1.0303) (1.2179,-1.0292) (1.2067,-1.0282) (1.1955,-1.0273) (1.1844,-1.0263) (1.1732,-1.0254) (1.1620,-1.0246) (1.1508,-1.0237) (1.1397,-1.0229) (1.1285,-1.0222) (1.1173,-1.0214) (1.1061,-1.0207) (1.0950,-1.0200) (1.0838,-1.0193) (1.0726,-1.0187) (1.0615,-1.0180) (1.0503,-1.0174) (1.0391,-1.0168) (1.0279,-1.0163) (1.0168,-1.0157) (1.0056,-1.0152) (0.9944,-1.0147) (0.9832,-1.0142) (0.9721,-1.0137) (0.9609,-1.0132) (0.9497,-1.0128) (0.9385,-1.0124) (0.9274,-1.0120) (0.9162,-1.0116) (0.9050,-1.0112) (0.8939,-1.0108) (0.8827,-1.0104) (0.8715,-1.0101) (0.8603,-1.0098) (0.8492,-1.0094) (0.8380,-1.0091) (0.8268,-1.0088) (0.8156,-1.0085) (0.8045,-1.0082) (0.7933,-1.0080) (0.7821,-1.0077) (0.7709,-1.0074) (0.7598,-1.0072) (0.7486,-1.0069) (0.7374,-1.0067) (0.7263,-1.0065) (0.7151,-1.0063) (0.7039,-1.0061) (0.6927,-1.0059) (0.6816,-1.0057) (0.6704,-1.0055) (0.6592,-1.0053) (0.6480,-1.0051) (0.6369,-1.0050) (0.6257,-1.0048) (0.6145,-1.0046) (0.6034,-1.0045) (0.5922,-1.0043) (0.5810,-1.0042) (0.5698,-1.0040) (0.5587,-1.0039) (0.5475,-1.0038) (0.5363,-1.0037) (0.5251,-1.0035) (0.5140,-1.0034) (0.5028,-1.0033) (0.4916,-1.0032) (0.4804,-1.0031) (0.4693,-1.0030) (0.4581,-1.0029) (0.4469,-1.0028) (0.4358,-1.0027) (0.4246,-1.0026) (0.4134,-1.0025) (0.4022,-1.0024) (0.3911,-1.0024) (0.3799,-1.0023) (0.3687,-1.0022) (0.3575,-1.0021) (0.3464,-1.0021) (0.3352,-1.0020) (0.3240,-1.0019) (0.3128,-1.0019) (0.3017,-1.0018) (0.2905,-1.0017) (0.2793,-1.0017) (0.2682,-1.0016) (0.2570,-1.0016) (0.2458,-1.0015) (0.2346,-1.0015) (0.2235,-1.0014) (0.2123,-1.0014) (0.2011,-1.0013) (0.1899,-1.0013) (0.1788,-1.0012) (0.1676,-1.0012) (0.1564,-1.0012) (0.1453,-1.0011) (0.1341,-1.0011) (0.1229,-1.0011) (0.1117,-1.0010) (0.1006,-1.0010) (0.0894,-1.0010) (0.0782,-1.0009) (0.0670,-1.0009) (0.0559,-1.0009) (0.0447,-1.0008) (0.0335,-1.0008) (0.0223,-1.0008) (0.0112,-1.0008) (0.0000,-1.0007)};\draw[main,thick]plot coordinates{(2.0000,-1.5000) (2.0013,-1.5033) (2.0025,-1.5067) (2.0038,-1.5100) (2.0050,-1.5135) (2.0063,-1.5169) (2.0075,-1.5204) (2.0088,-1.5239) (2.0101,-1.5275) (2.0113,-1.5311) (2.0126,-1.5347) (2.0138,-1.5383) (2.0151,-1.5420) (2.0163,-1.5458) (2.0176,-1.5496) (2.0189,-1.5534) (2.0201,-1.5572) (2.0214,-1.5611) (2.0226,-1.5651) (2.0239,-1.5691) (2.0251,-1.5731) (2.0264,-1.5772) (2.0277,-1.5813) (2.0289,-1.5854) (2.0302,-1.5897) (2.0314,-1.5939) (2.0327,-1.5982) (2.0339,-1.6026) (2.0352,-1.6070) (2.0365,-1.6114) (2.0377,-1.6159) (2.0390,-1.6204) (2.0402,-1.6250) (2.0415,-1.6297) (2.0427,-1.6344) (2.0440,-1.6392) (2.0452,-1.6440) (2.0465,-1.6489) (2.0478,-1.6538) (2.0490,-1.6588) (2.0503,-1.6639) (2.0515,-1.6690) (2.0528,-1.6742) (2.0540,-1.6794) (2.0553,-1.6847) (2.0566,-1.6901) (2.0578,-1.6955) (2.0591,-1.7010) (2.0603,-1.7066) (2.0616,-1.7123) (2.0628,-1.7180) (2.0641,-1.7238) (2.0654,-1.7297) (2.0666,-1.7356) (2.0679,-1.7417) (2.0691,-1.7478) (2.0704,-1.7540) (2.0716,-1.7603) (2.0729,-1.7666) (2.0742,-1.7731) (2.0754,-1.7797) (2.0767,-1.7863) (2.0779,-1.7930) (2.0792,-1.7998) (2.0804,-1.8068) (2.0817,-1.8138) (2.0830,-1.8209) (2.0842,-1.8282) (2.0855,-1.8355) (2.0867,-1.8429) (2.0880,-1.8505) (2.0892,-1.8582) (2.0905,-1.8660) (2.0918,-1.8739) (2.0930,-1.8819) (2.0943,-1.8901) (2.0955,-1.8984) (2.0968,-1.9068) (2.0980,-1.9154) (2.0993,-1.9241) (2.1006,-1.9329) (2.1018,-1.9419) (2.1031,-1.9510) (2.1043,-1.9603) (2.1056,-1.9698) (2.1068,-1.9794) (2.1081,-1.9892) (2.1094,-1.9991)};\draw[main,thick]plot coordinates{(2.0000,-0.5000) (1.9888,-0.5070) (1.9777,-0.5140) (1.9665,-0.5210) (1.9553,-0.5281) (1.9441,-0.5351) (1.9330,-0.5422) (1.9218,-0.5492) (1.9106,-0.5563) (1.8994,-0.5633) (1.8883,-0.5703) (1.8771,-0.5774) (1.8659,-0.5844) (1.8547,-0.5914) (1.8436,-0.5984) (1.8324,-0.6054) (1.8212,-0.6123) (1.8101,-0.6192) (1.7989,-0.6261) (1.7877,-0.6329) (1.7765,-0.6398) (1.7654,-0.6465) (1.7542,-0.6533) (1.7430,-0.6600) (1.7318,-0.6666) (1.7207,-0.6732) (1.7095,-0.6798) (1.6983,-0.6862) (1.6872,-0.6927) (1.6760,-0.6991) (1.6648,-0.7054) (1.6536,-0.7116) (1.6425,-0.7178) (1.6313,-0.7239) (1.6201,-0.7300) (1.6089,-0.7359) (1.5978,-0.7418) (1.5866,-0.7477) (1.5754,-0.7534) (1.5642,-0.7591) (1.5531,-0.7647) (1.5419,-0.7702) (1.5307,-0.7757) (1.5196,-0.7810) (1.5084,-0.7863) (1.4972,-0.7915) (1.4860,-0.7966) (1.4749,-0.8016) (1.4637,-0.8065) (1.4525,-0.8114) (1.4413,-0.8162) (1.4302,-0.8208) (1.4190,-0.8254) (1.4078,-0.8299) (1.3966,-0.8344) (1.3855,-0.8387) (1.3743,-0.8429) (1.3631,-0.8471) (1.3520,-0.8512) (1.3408,-0.8552) (1.3296,-0.8591) (1.3184,-0.8629) (1.3073,-0.8666) (1.2961,-0.8703) (1.2849,-0.8739) (1.2737,-0.8774) (1.2626,-0.8808) (1.2514,-0.8841) (1.2402,-0.8874) (1.2291,-0.8906) (1.2179,-0.8937) (1.2067,-0.8967) (1.1955,-0.8997) (1.1844,-0.9026) (1.1732,-0.9054) (1.1620,-0.9081) (1.1508,-0.9108) (1.1397,-0.9134) (1.1285,-0.9160) (1.1173,-0.9184) (1.1061,-0.9209) (1.0950,-0.9232) (1.0838,-0.9255) (1.0726,-0.9277) (1.0615,-0.9299) (1.0503,-0.9320) (1.0391,-0.9340) (1.0279,-0.9360) (1.0168,-0.9380) (1.0056,-0.9398) (0.9944,-0.9417) (0.9832,-0.9435) (0.9721,-0.9452) (0.9609,-0.9469) (0.9497,-0.9485) (0.9385,-0.9501) (0.9274,-0.9516) (0.9162,-0.9531) (0.9050,-0.9546) (0.8939,-0.9560) (0.8827,-0.9574) (0.8715,-0.9587) (0.8603,-0.9600) (0.8492,-0.9612) (0.8380,-0.9624) (0.8268,-0.9636) (0.8156,-0.9648) (0.8045,-0.9659) (0.7933,-0.9669) (0.7821,-0.9680) (0.7709,-0.9690) (0.7598,-0.9700) (0.7486,-0.9709) (0.7374,-0.9719) (0.7263,-0.9727) (0.7151,-0.9736) (0.7039,-0.9745) (0.6927,-0.9753) (0.6816,-0.9761) (0.6704,-0.9768) (0.6592,-0.9776) (0.6480,-0.9783) (0.6369,-0.9790) (0.6257,-0.9797) (0.6145,-0.9803) (0.6034,-0.9809) (0.5922,-0.9816) (0.5810,-0.9821) (0.5698,-0.9827) (0.5587,-0.9833) (0.5475,-0.9838) (0.5363,-0.9843) (0.5251,-0.9848) (0.5140,-0.9853) (0.5028,-0.9858) (0.4916,-0.9863) (0.4804,-0.9867) (0.4693,-0.9871) (0.4581,-0.9876) (0.4469,-0.9880) (0.4358,-0.9884) (0.4246,-0.9887) (0.4134,-0.9891) (0.4022,-0.9895) (0.3911,-0.9898) (0.3799,-0.9901) (0.3687,-0.9905) (0.3575,-0.9908) (0.3464,-0.9911) (0.3352,-0.9914) (0.3240,-0.9916) (0.3128,-0.9919) (0.3017,-0.9922) (0.2905,-0.9924) (0.2793,-0.9927) (0.2682,-0.9929) (0.2570,-0.9931) (0.2458,-0.9934) (0.2346,-0.9936) (0.2235,-0.9938) (0.2123,-0.9940) (0.2011,-0.9942) (0.1899,-0.9944) (0.1788,-0.9946) (0.1676,-0.9947) (0.1564,-0.9949) (0.1453,-0.9951) (0.1341,-0.9952) (0.1229,-0.9954) (0.1117,-0.9956) (0.1006,-0.9957) (0.0894,-0.9958) (0.0782,-0.9960) (0.0670,-0.9961) (0.0559,-0.9962) (0.0447,-0.9964) (0.0335,-0.9965) (0.0223,-0.9966) (0.0112,-0.9967) (0.0000,-0.9968)};\draw[main,thick]plot coordinates{(2.0000,-0.5000) (2.0168,-0.4895) (2.0335,-0.4791) (2.0503,-0.4688) (2.0670,-0.4585) (2.0838,-0.4484) (2.1006,-0.4382) (2.1173,-0.4282) (2.1341,-0.4183) (2.1508,-0.4085) (2.1676,-0.3988) (2.1844,-0.3892) (2.2011,-0.3798) (2.2179,-0.3704) (2.2346,-0.3613) (2.2514,-0.3522) (2.2682,-0.3433) (2.2849,-0.3345) (2.3017,-0.3258) (2.3184,-0.3174) (2.3352,-0.3090) (2.3520,-0.3008) (2.3687,-0.2928) (2.3855,-0.2849) (2.4022,-0.2772) (2.4190,-0.2696) (2.4358,-0.2622) (2.4525,-0.2550) (2.4693,-0.2479) (2.4860,-0.2409) (2.5028,-0.2341) (2.5196,-0.2275) (2.5363,-0.2210) (2.5531,-0.2147) (2.5698,-0.2085) (2.5866,-0.2024) (2.6034,-0.1966) (2.6201,-0.1908) (2.6369,-0.1852) (2.6536,-0.1798) (2.6704,-0.1744) (2.6872,-0.1693) (2.7039,-0.1642) (2.7207,-0.1593) (2.7374,-0.1545) (2.7542,-0.1499) (2.7709,-0.1453) (2.7877,-0.1409) (2.8045,-0.1367) (2.8212,-0.1325) (2.8380,-0.1284) (2.8547,-0.1245) (2.8715,-0.1207) (2.8883,-0.1170) (2.9050,-0.1133) (2.9218,-0.1098) (2.9385,-0.1064) (2.9553,-0.1031) (2.9721,-0.0999) (2.9888,-0.0968) (3.0056,-0.0938) (3.0223,-0.0908) (3.0391,-0.0880) (3.0559,-0.0852) (3.0726,-0.0825) (3.0894,-0.0799) (3.1061,-0.0774) (3.1229,-0.0750) (3.1397,-0.0726) (3.1564,-0.0703) (3.1732,-0.0680) (3.1899,-0.0659) (3.2067,-0.0638) (3.2235,-0.0617) (3.2402,-0.0598) (3.2570,-0.0579) (3.2737,-0.0560) (3.2905,-0.0542) (3.3073,-0.0525) (3.3240,-0.0508) (3.3408,-0.0492) (3.3575,-0.0476) (3.3743,-0.0461) (3.3911,-0.0446) (3.4078,-0.0431) (3.4246,-0.0417) (3.4413,-0.0404) (3.4581,-0.0391) (3.4749,-0.0378) (3.4916,-0.0366) (3.5084,-0.0354) (3.5251,-0.0343) (3.5419,-0.0332) (3.5587,-0.0321) (3.5754,-0.0310) (3.5922,-0.0300) (3.6089,-0.0291) (3.6257,-0.0281) (3.6425,-0.0272) (3.6592,-0.0263) (3.6760,-0.0255) (3.6927,-0.0246) (3.7095,-0.0238) (3.7263,-0.0231) (3.7430,-0.0223) (3.7598,-0.0216) (3.7765,-0.0209) (3.7933,-0.0202) (3.8101,-0.0195) (3.8268,-0.0189) (3.8436,-0.0183) (3.8603,-0.0177) (3.8771,-0.0171) (3.8939,-0.0165) (3.9106,-0.0160) (3.9274,-0.0155) (3.9441,-0.0150) (3.9609,-0.0145) (3.9777,-0.0140) (3.9944,-0.0136) (4.0112,-0.0131) (4.0279,-0.0127) (4.0447,-0.0123) (4.0615,-0.0119) (4.0782,-0.0115) (4.0950,-0.0111) (4.1117,-0.0107) (4.1285,-0.0104) (4.1453,-0.0100) (4.1620,-0.0097) (4.1788,-0.0094) (4.1955,-0.0091) (4.2123,-0.0088) (4.2291,-0.0085) (4.2458,-0.0082) (4.2626,-0.0080) (4.2793,-0.0077) (4.2961,-0.0074) (4.3128,-0.0072) (4.3296,-0.0070) (4.3464,-0.0067) (4.3631,-0.0065) (4.3799,-0.0063) (4.3966,-0.0061) (4.4134,-0.0059) (4.4302,-0.0057) (4.4469,-0.0055) (4.4637,-0.0053) (4.4804,-0.0051) (4.4972,-0.0050) (4.5140,-0.0048) (4.5307,-0.0047) (4.5475,-0.0045) (4.5642,-0.0044) (4.5810,-0.0042) (4.5978,-0.0041) (4.6145,-0.0039) (4.6313,-0.0038) (4.6480,-0.0037) (4.6648,-0.0036) (4.6816,-0.0034) (4.6983,-0.0033) (4.7151,-0.0032) (4.7318,-0.0031) (4.7486,-0.0030) (4.7654,-0.0029) (4.7821,-0.0028) (4.7989,-0.0027) (4.8156,-0.0026) (4.8324,-0.0026) (4.8492,-0.0025) (4.8659,-0.0024) (4.8827,-0.0023) (4.8994,-0.0022) (4.9162,-0.0022) (4.9330,-0.0021) (4.9497,-0.0020) (4.9665,-0.0020) (4.9832,-0.0019) (5.0000,-0.0018)};\draw[main,thick]plot coordinates{(2.0000,1.0000) (1.9888,1.0225) (1.9777,1.0452) (1.9665,1.0681) (1.9553,1.0912) (1.9441,1.1144) (1.9330,1.1378) (1.9218,1.1613) (1.9106,1.1848) (1.8994,1.2084) (1.8883,1.2320) (1.8771,1.2556) (1.8659,1.2791) (1.8547,1.3025) (1.8436,1.3258) (1.8324,1.3490) (1.8212,1.3719) (1.8101,1.3946) (1.7989,1.4171) (1.7877,1.4392) (1.7765,1.4611) (1.7654,1.4825) (1.7542,1.5036) (1.7430,1.5243) (1.7318,1.5445) (1.7207,1.5643) (1.7095,1.5836) (1.6983,1.6023) (1.6872,1.6206) (1.6760,1.6383) (1.6648,1.6555) (1.6536,1.6722) (1.6425,1.6883) (1.6313,1.7038) (1.6201,1.7188) (1.6089,1.7332) (1.5978,1.7470) (1.5866,1.7603) (1.5754,1.7730) (1.5642,1.7852) (1.5531,1.7969) (1.5419,1.8081) (1.5307,1.8187) (1.5196,1.8288) (1.5084,1.8385) (1.4972,1.8477) (1.4860,1.8564) (1.4749,1.8647) (1.4637,1.8726) (1.4525,1.8800) (1.4413,1.8871) (1.4302,1.8938) (1.4190,1.9001) (1.4078,1.9061) (1.3966,1.9117) (1.3855,1.9171) (1.3743,1.9221) (1.3631,1.9268) (1.3520,1.9313) (1.3408,1.9355) (1.3296,1.9395) (1.3184,1.9432) (1.3073,1.9468) (1.2961,1.9501) (1.2849,1.9532) (1.2737,1.9561) (1.2626,1.9589) (1.2514,1.9614) (1.2402,1.9639) (1.2291,1.9661) (1.2179,1.9683) (1.2067,1.9703) (1.1955,1.9722) (1.1844,1.9739) (1.1732,1.9756) (1.1620,1.9771) (1.1508,1.9786) (1.1397,1.9800) (1.1285,1.9812) (1.1173,1.9824) (1.1061,1.9836) (1.0950,1.9846) (1.0838,1.9856) (1.0726,1.9865) (1.0615,1.9874) (1.0503,1.9882) (1.0391,1.9890) (1.0279,1.9897) (1.0168,1.9903) (1.0056,1.9910) (0.9944,1.9915) (0.9832,1.9921) (0.9721,1.9926) (0.9609,1.9931) (0.9497,1.9935) (0.9385,1.9939) (0.9274,1.9943) (0.9162,1.9947) (0.9050,1.9950) (0.8939,1.9954) (0.8827,1.9957) (0.8715,1.9959) (0.8603,1.9962) (0.8492,1.9964) (0.8380,1.9967) (0.8268,1.9969) (0.8156,1.9971) (0.8045,1.9973) (0.7933,1.9975) (0.7821,1.9976) (0.7709,1.9978) (0.7598,1.9979) (0.7486,1.9981) (0.7374,1.9982) (0.7263,1.9983) (0.7151,1.9984) (0.7039,1.9985) (0.6927,1.9986) (0.6816,1.9987) (0.6704,1.9988) (0.6592,1.9989) (0.6480,1.9989) (0.6369,1.9990) (0.6257,1.9991) (0.6145,1.9991) (0.6034,1.9992) (0.5922,1.9992) (0.5810,1.9993) (0.5698,1.9993) (0.5587,1.9994) (0.5475,1.9994) (0.5363,1.9995) (0.5251,1.9995) (0.5140,1.9995) (0.5028,1.9996) (0.4916,1.9996) (0.4804,1.9996) (0.4693,1.9996) (0.4581,1.9997) (0.4469,1.9997) (0.4358,1.9997) (0.4246,1.9997) (0.4134,1.9997) (0.4022,1.9998) (0.3911,1.9998) (0.3799,1.9998) (0.3687,1.9998) (0.3575,1.9998) (0.3464,1.9998) (0.3352,1.9998) (0.3240,1.9998) (0.3128,1.9999) (0.3017,1.9999) (0.2905,1.9999) (0.2793,1.9999) (0.2682,1.9999) (0.2570,1.9999) (0.2458,1.9999) (0.2346,1.9999) (0.2235,1.9999) (0.2123,1.9999) (0.2011,1.9999) (0.1899,1.9999) (0.1788,1.9999) (0.1676,1.9999) (0.1564,1.9999) (0.1453,1.9999) (0.1341,2.0000) (0.1229,2.0000) (0.1117,2.0000) (0.1006,2.0000) (0.0894,2.0000) (0.0782,2.0000) (0.0670,2.0000) (0.0559,2.0000) (0.0447,2.0000) (0.0335,2.0000) (0.0223,2.0000) (0.0112,2.0000) (0.0000,2.0000)};\draw[main,thick]plot coordinates{(2.0000,1.0000) (2.0168,0.9668) (2.0335,0.9342) (2.0503,0.9022) (2.0670,0.8710) (2.0838,0.8406) (2.1006,0.8109) (2.1173,0.7821) (2.1341,0.7540) (2.1508,0.7268) (2.1676,0.7005) (2.1844,0.6749) (2.2011,0.6502) (2.2179,0.6264) (2.2346,0.6033) (2.2514,0.5811) (2.2682,0.5596) (2.2849,0.5389) (2.3017,0.5190) (2.3184,0.4998) (2.3352,0.4813) (2.3520,0.4635) (2.3687,0.4464) (2.3855,0.4299) (2.4022,0.4140) (2.4190,0.3987) (2.4358,0.3841) (2.4525,0.3699) (2.4693,0.3564) (2.4860,0.3433) (2.5028,0.3307) (2.5196,0.3187) (2.5363,0.3071) (2.5531,0.2959) (2.5698,0.2852) (2.5866,0.2748) (2.6034,0.2649) (2.6201,0.2553) (2.6369,0.2461) (2.6536,0.2373) (2.6704,0.2288) (2.6872,0.2206) (2.7039,0.2127) (2.7207,0.2051) (2.7374,0.1978) (2.7542,0.1908) (2.7709,0.1841) (2.7877,0.1776) (2.8045,0.1713) (2.8212,0.1652) (2.8380,0.1594) (2.8547,0.1538) (2.8715,0.1485) (2.8883,0.1433) (2.9050,0.1383) (2.9218,0.1334) (2.9385,0.1288) (2.9553,0.1243) (2.9721,0.1200) (2.9888,0.1158) (3.0056,0.1118) (3.0223,0.1080) (3.0391,0.1042) (3.0559,0.1007) (3.0726,0.0972) (3.0894,0.0939) (3.1061,0.0906) (3.1229,0.0875) (3.1397,0.0845) (3.1564,0.0816) (3.1732,0.0788) (3.1899,0.0762) (3.2067,0.0736) (3.2235,0.0711) (3.2402,0.0686) (3.2570,0.0663) (3.2737,0.0641) (3.2905,0.0619) (3.3073,0.0598) (3.3240,0.0578) (3.3408,0.0558) (3.3575,0.0539) (3.3743,0.0521) (3.3911,0.0503) (3.4078,0.0486) (3.4246,0.0470) (3.4413,0.0454) (3.4581,0.0439) (3.4749,0.0424) (3.4916,0.0410) (3.5084,0.0396) (3.5251,0.0383) (3.5419,0.0370) (3.5587,0.0358) (3.5754,0.0346) (3.5922,0.0334) (3.6089,0.0323) (3.6257,0.0312) (3.6425,0.0302) (3.6592,0.0292) (3.6760,0.0282) (3.6927,0.0272) (3.7095,0.0263) (3.7263,0.0254) (3.7430,0.0246) (3.7598,0.0238) (3.7765,0.0230) (3.7933,0.0222) (3.8101,0.0215) (3.8268,0.0208) (3.8436,0.0201) (3.8603,0.0194) (3.8771,0.0188) (3.8939,0.0181) (3.9106,0.0175) (3.9274,0.0170) (3.9441,0.0164) (3.9609,0.0158) (3.9777,0.0153) (3.9944,0.0148) (4.0112,0.0143) (4.0279,0.0138) (4.0447,0.0134) (4.0615,0.0129) (4.0782,0.0125) (4.0950,0.0121) (4.1117,0.0117) (4.1285,0.0113) (4.1453,0.0109) (4.1620,0.0106) (4.1788,0.0102) (4.1955,0.0099) (4.2123,0.0096) (4.2291,0.0092) (4.2458,0.0089) (4.2626,0.0086) (4.2793,0.0083) (4.2961,0.0081) (4.3128,0.0078) (4.3296,0.0075) (4.3464,0.0073) (4.3631,0.0071) (4.3799,0.0068) (4.3966,0.0066) (4.4134,0.0064) (4.4302,0.0062) (4.4469,0.0060) (4.4637,0.0058) (4.4804,0.0056) (4.4972,0.0054) (4.5140,0.0052) (4.5307,0.0050) (4.5475,0.0049) (4.5642,0.0047) (4.5810,0.0046) (4.5978,0.0044) (4.6145,0.0043) (4.6313,0.0041) (4.6480,0.0040) (4.6648,0.0039) (4.6816,0.0037) (4.6983,0.0036) (4.7151,0.0035) (4.7318,0.0034) (4.7486,0.0033) (4.7654,0.0032) (4.7821,0.0030) (4.7989,0.0029) (4.8156,0.0028) (4.8324,0.0028) (4.8492,0.0027) (4.8659,0.0026) (4.8827,0.0025) (4.8994,0.0024) (4.9162,0.0023) (4.9330,0.0023) (4.9497,0.0022) (4.9665,0.0021) (4.9832,0.0020) (5.0000,0.0020)};\draw[main,thick]plot coordinates{(2.0000,2.5000) (1.9888,2.4541) (1.9777,2.4137) (1.9665,2.3777) (1.9553,2.3457) (1.9441,2.3169) (1.9330,2.2911) (1.9218,2.2678) (1.9106,2.2467) (1.8994,2.2275) (1.8883,2.2100) (1.8771,2.1941) (1.8659,2.1796) (1.8547,2.1663) (1.8436,2.1541) (1.8324,2.1429) (1.8212,2.1325) (1.8101,2.1230) (1.7989,2.1143) (1.7877,2.1062) (1.7765,2.0987) (1.7654,2.0918) (1.7542,2.0855) (1.7430,2.0795) (1.7318,2.0741) (1.7207,2.0690) (1.7095,2.0643) (1.6983,2.0599) (1.6872,2.0558) (1.6760,2.0520) (1.6648,2.0485) (1.6536,2.0453) (1.6425,2.0422) (1.6313,2.0394) (1.6201,2.0368) (1.6089,2.0343) (1.5978,2.0320) (1.5866,2.0299) (1.5754,2.0279) (1.5642,2.0261) (1.5531,2.0243) (1.5419,2.0227) (1.5307,2.0212) (1.5196,2.0198) (1.5084,2.0185) (1.4972,2.0173) (1.4860,2.0162) (1.4749,2.0151) (1.4637,2.0141) (1.4525,2.0132) (1.4413,2.0123) (1.4302,2.0115) (1.4190,2.0108) (1.4078,2.0101) (1.3966,2.0094) (1.3855,2.0088) (1.3743,2.0082) (1.3631,2.0077) (1.3520,2.0072) (1.3408,2.0067) (1.3296,2.0063) (1.3184,2.0059) (1.3073,2.0055) (1.2961,2.0051) (1.2849,2.0048) (1.2737,2.0045) (1.2626,2.0042) (1.2514,2.0039) (1.2402,2.0037) (1.2291,2.0034) (1.2179,2.0032) (1.2067,2.0030) (1.1955,2.0028) (1.1844,2.0026) (1.1732,2.0024) (1.1620,2.0023) (1.1508,2.0021) (1.1397,2.0020) (1.1285,2.0019) (1.1173,2.0017) (1.1061,2.0016) (1.0950,2.0015) (1.0838,2.0014) (1.0726,2.0013) (1.0615,2.0013) (1.0503,2.0012) (1.0391,2.0011) (1.0279,2.0010) (1.0168,2.0010) (1.0056,2.0009) (0.9944,2.0008) (0.9832,2.0008) (0.9721,2.0007) (0.9609,2.0007) (0.9497,2.0006) (0.9385,2.0006) (0.9274,2.0006) (0.9162,2.0005) (0.9050,2.0005) (0.8939,2.0005) (0.8827,2.0004) (0.8715,2.0004) (0.8603,2.0004) (0.8492,2.0003) (0.8380,2.0003) (0.8268,2.0003) (0.8156,2.0003) (0.8045,2.0003) (0.7933,2.0002) (0.7821,2.0002) (0.7709,2.0002) (0.7598,2.0002) (0.7486,2.0002) (0.7374,2.0002) (0.7263,2.0002) (0.7151,2.0002) (0.7039,2.0001) (0.6927,2.0001) (0.6816,2.0001) (0.6704,2.0001) (0.6592,2.0001) (0.6480,2.0001) (0.6369,2.0001) (0.6257,2.0001) (0.6145,2.0001) (0.6034,2.0001) (0.5922,2.0001) (0.5810,2.0001) (0.5698,2.0001) (0.5587,2.0001) (0.5475,2.0001) (0.5363,2.0001) (0.5251,2.0001) (0.5140,2.0000) (0.5028,2.0000) (0.4916,2.0000) (0.4804,2.0000) (0.4693,2.0000) (0.4581,2.0000) (0.4469,2.0000) (0.4358,2.0000) (0.4246,2.0000) (0.4134,2.0000) (0.4022,2.0000) (0.3911,2.0000) (0.3799,2.0000) (0.3687,2.0000) (0.3575,2.0000) (0.3464,2.0000) (0.3352,2.0000) (0.3240,2.0000) (0.3128,2.0000) (0.3017,2.0000) (0.2905,2.0000) (0.2793,2.0000) (0.2682,2.0000) (0.2570,2.0000) (0.2458,2.0000) (0.2346,2.0000) (0.2235,2.0000) (0.2123,2.0000) (0.2011,2.0000) (0.1899,2.0000) (0.1788,2.0000) (0.1676,2.0000) (0.1564,2.0000) (0.1453,2.0000) (0.1341,2.0000) (0.1229,2.0000) (0.1117,2.0000) (0.1006,2.0000) (0.0894,2.0000) (0.0782,2.0000) (0.0670,2.0000) (0.0559,2.0000) (0.0447,2.0000) (0.0335,2.0000) (0.0223,2.0000) (0.0112,2.0000) (0.0000,2.0000)};\draw[main,thick]plot coordinates{(2.0000,2.5000) (2.0009,2.5038) (2.0017,2.5077) (2.0026,2.5116) (2.0035,2.5156) (2.0044,2.5196) (2.0052,2.5236) (2.0061,2.5277) (2.0070,2.5318) (2.0078,2.5360) (2.0087,2.5402) (2.0096,2.5445) (2.0105,2.5488) (2.0113,2.5531) (2.0122,2.5575) (2.0131,2.5620) (2.0140,2.5665) (2.0148,2.5710) (2.0157,2.5756) (2.0166,2.5803) (2.0174,2.5850) (2.0183,2.5897) (2.0192,2.5945) (2.0201,2.5994) (2.0209,2.6043) (2.0218,2.6093) (2.0227,2.6143) (2.0235,2.6194) (2.0244,2.6245) (2.0253,2.6298) (2.0262,2.6350) (2.0270,2.6404) (2.0279,2.6458) (2.0288,2.6512) (2.0296,2.6568) (2.0305,2.6623) (2.0314,2.6680) (2.0323,2.6737) (2.0331,2.6796) (2.0340,2.6854) (2.0349,2.6914) (2.0358,2.6974) (2.0366,2.7035) (2.0375,2.7097) (2.0384,2.7160) (2.0392,2.7223) (2.0401,2.7287) (2.0410,2.7352) (2.0419,2.7418) (2.0427,2.7485) (2.0436,2.7553) (2.0445,2.7621) (2.0453,2.7691) (2.0462,2.7761) (2.0471,2.7833) (2.0480,2.7905) (2.0488,2.7979) (2.0497,2.8053) (2.0506,2.8129) (2.0514,2.8205) (2.0523,2.8283) (2.0532,2.8361) (2.0541,2.8441) (2.0549,2.8522) (2.0558,2.8605) (2.0567,2.8688) (2.0575,2.8773) (2.0584,2.8859) (2.0593,2.8946) (2.0602,2.9035) (2.0610,2.9125) (2.0619,2.9216) (2.0628,2.9309) (2.0637,2.9403) (2.0645,2.9499) (2.0654,2.9596) (2.0663,2.9695) (2.0671,2.9795) (2.0680,2.9898)};\end{tikzpicture}\end{center}
}

\newcommand{\ODEExamCubicPhase}{%
\begin{center}\begin{tikzpicture}[x=.8cm,y=.7cm]\draw[->](-1,-2)--(-1,2.5)node[above]{$x$};\draw[->](-1,-1.35)--(-1,-1.65);\draw[->](-1,-0.65)--(-1,-0.35);\draw[->](-1,0.65)--(-1,0.35);\draw[->](-1,1.35)--(-1,1.65);\fill(-1,-1)circle(1.8pt);\node[left]at(-1,-1){$ -1 $};\draw[dashed](0,-1)--(5,-1);\fill(-1,0)circle(1.8pt);\node[left]at(-1,0){$ 0 $};\draw[dashed](0,0)--(5,0);\fill(-1,1)circle(1.8pt);\node[left]at(-1,1){$ 1 $};\draw[dashed](0,1)--(5,1);\draw[->](0,-2)--(5.4,-2)node[right]{$t$};\draw[->](0,-2)--(0,2.3)node[above]{$x$};\draw[main,thick]plot coordinates{(2.0000,-1.5000) (1.9888,-1.4797) (1.9777,-1.4606) (1.9665,-1.4427) (1.9553,-1.4257) (1.9441,-1.4097) (1.9330,-1.3946) (1.9218,-1.3803) (1.9106,-1.3667) (1.8994,-1.3538) (1.8883,-1.3415) (1.8771,-1.3298) (1.8659,-1.3186) (1.8547,-1.3080) (1.8436,-1.2978) (1.8324,-1.2881) (1.8212,-1.2789) (1.8101,-1.2700) (1.7989,-1.2615) (1.7877,-1.2533) (1.7765,-1.2455) (1.7654,-1.2379) (1.7542,-1.2307) (1.7430,-1.2238) (1.7318,-1.2171) (1.7207,-1.2107) (1.7095,-1.2045) (1.6983,-1.1986) (1.6872,-1.1928) (1.6760,-1.1873) (1.6648,-1.1819) (1.6536,-1.1768) (1.6425,-1.1718) (1.6313,-1.1670) (1.6201,-1.1624) (1.6089,-1.1579) (1.5978,-1.1536) (1.5866,-1.1494) (1.5754,-1.1453) (1.5642,-1.1414) (1.5531,-1.1376) (1.5419,-1.1339) (1.5307,-1.1303) (1.5196,-1.1269) (1.5084,-1.1235) (1.4972,-1.1203) (1.4860,-1.1172) (1.4749,-1.1141) (1.4637,-1.1112) (1.4525,-1.1083) (1.4413,-1.1055) (1.4302,-1.1028) (1.4190,-1.1002) (1.4078,-1.0976) (1.3966,-1.0951) (1.3855,-1.0927) (1.3743,-1.0904) (1.3631,-1.0881) (1.3520,-1.0859) (1.3408,-1.0838) (1.3296,-1.0817) (1.3184,-1.0797) (1.3073,-1.0777) (1.2961,-1.0758) (1.2849,-1.0739) (1.2737,-1.0721) (1.2626,-1.0703) (1.2514,-1.0686) (1.2402,-1.0670) (1.2291,-1.0653) (1.2179,-1.0637) (1.2067,-1.0622) (1.1955,-1.0607) (1.1844,-1.0592) (1.1732,-1.0578) (1.1620,-1.0564) (1.1508,-1.0551) (1.1397,-1.0537) (1.1285,-1.0525) (1.1173,-1.0512) (1.1061,-1.0500) (1.0950,-1.0488) (1.0838,-1.0477) (1.0726,-1.0465) (1.0615,-1.0454) (1.0503,-1.0444) (1.0391,-1.0433) (1.0279,-1.0423) (1.0168,-1.0413) (1.0056,-1.0403) (0.9944,-1.0394) (0.9832,-1.0385) (0.9721,-1.0376) (0.9609,-1.0367) (0.9497,-1.0358) (0.9385,-1.0350) (0.9274,-1.0342) (0.9162,-1.0334) (0.9050,-1.0326) (0.8939,-1.0319) (0.8827,-1.0311) (0.8715,-1.0304) (0.8603,-1.0297) (0.8492,-1.0290) (0.8380,-1.0284) (0.8268,-1.0277) (0.8156,-1.0271) (0.8045,-1.0264) (0.7933,-1.0258) (0.7821,-1.0252) (0.7709,-1.0247) (0.7598,-1.0241) (0.7486,-1.0235) (0.7374,-1.0230) (0.7263,-1.0225) (0.7151,-1.0220) (0.7039,-1.0215) (0.6927,-1.0210) (0.6816,-1.0205) (0.6704,-1.0200) (0.6592,-1.0196) (0.6480,-1.0191) (0.6369,-1.0187) (0.6257,-1.0183) (0.6145,-1.0179) (0.6034,-1.0175) (0.5922,-1.0171) (0.5810,-1.0167) (0.5698,-1.0163) (0.5587,-1.0159) (0.5475,-1.0156) (0.5363,-1.0152) (0.5251,-1.0149) (0.5140,-1.0145) (0.5028,-1.0142) (0.4916,-1.0139) (0.4804,-1.0136) (0.4693,-1.0133) (0.4581,-1.0130) (0.4469,-1.0127) (0.4358,-1.0124) (0.4246,-1.0121) (0.4134,-1.0118) (0.4022,-1.0116) (0.3911,-1.0113) (0.3799,-1.0111) (0.3687,-1.0108) (0.3575,-1.0106) (0.3464,-1.0103) (0.3352,-1.0101) (0.3240,-1.0099) (0.3128,-1.0096) (0.3017,-1.0094) (0.2905,-1.0092) (0.2793,-1.0090) (0.2682,-1.0088) (0.2570,-1.0086) (0.2458,-1.0084) (0.2346,-1.0082) (0.2235,-1.0081) (0.2123,-1.0079) (0.2011,-1.0077) (0.1899,-1.0075) (0.1788,-1.0074) (0.1676,-1.0072) (0.1564,-1.0070) (0.1453,-1.0069) (0.1341,-1.0067) (0.1229,-1.0066) (0.1117,-1.0064) (0.1006,-1.0063) (0.0894,-1.0061) (0.0782,-1.0060) (0.0670,-1.0059) (0.0559,-1.0057) (0.0447,-1.0056) (0.0335,-1.0055) (0.0223,-1.0054) (0.0112,-1.0052) (0.0000,-1.0051)};\draw[main,thick]plot coordinates{(2.0000,-1.5000) (2.0016,-1.5031) (2.0033,-1.5062) (2.0049,-1.5094) (2.0066,-1.5126) (2.0082,-1.5158) (2.0099,-1.5190) (2.0115,-1.5223) (2.0131,-1.5256) (2.0148,-1.5289) (2.0164,-1.5323) (2.0181,-1.5357) (2.0197,-1.5392) (2.0213,-1.5426) (2.0230,-1.5462) (2.0246,-1.5497) (2.0263,-1.5533) (2.0279,-1.5569) (2.0296,-1.5606) (2.0312,-1.5643) (2.0328,-1.5680) (2.0345,-1.5718) (2.0361,-1.5756) (2.0378,-1.5794) (2.0394,-1.5833) (2.0410,-1.5873) (2.0427,-1.5913) (2.0443,-1.5953) (2.0460,-1.5994) (2.0476,-1.6035) (2.0493,-1.6076) (2.0509,-1.6118) (2.0525,-1.6161) (2.0542,-1.6204) (2.0558,-1.6247) (2.0575,-1.6291) (2.0591,-1.6336) (2.0607,-1.6381) (2.0624,-1.6426) (2.0640,-1.6472) (2.0657,-1.6519) (2.0673,-1.6566) (2.0690,-1.6614) (2.0706,-1.6662) (2.0722,-1.6711) (2.0739,-1.6761) (2.0755,-1.6811) (2.0772,-1.6861) (2.0788,-1.6913) (2.0805,-1.6965) (2.0821,-1.7017) (2.0837,-1.7071) (2.0854,-1.7125) (2.0870,-1.7179) (2.0887,-1.7235) (2.0903,-1.7291) (2.0919,-1.7348) (2.0936,-1.7405) (2.0952,-1.7464) (2.0969,-1.7523) (2.0985,-1.7583) (2.1002,-1.7644) (2.1018,-1.7705) (2.1034,-1.7768) (2.1051,-1.7831) (2.1067,-1.7895) (2.1084,-1.7961) (2.1100,-1.8027) (2.1116,-1.8094) (2.1133,-1.8162) (2.1149,-1.8231) (2.1166,-1.8301) (2.1182,-1.8372) (2.1199,-1.8444) (2.1215,-1.8518) (2.1231,-1.8592) (2.1248,-1.8667) (2.1264,-1.8744) (2.1281,-1.8822) (2.1297,-1.8901) (2.1313,-1.8982) (2.1330,-1.9064) (2.1346,-1.9147) (2.1363,-1.9231) (2.1379,-1.9317) (2.1396,-1.9405) (2.1412,-1.9493) (2.1428,-1.9584) (2.1445,-1.9676) (2.1461,-1.9769) (2.1478,-1.9865) (2.1494,-1.9962)};\draw[main,thick]plot coordinates{(2.0000,-0.5000) (1.9888,-0.5042) (1.9777,-0.5084) (1.9665,-0.5126) (1.9553,-0.5168) (1.9441,-0.5211) (1.9330,-0.5253) (1.9218,-0.5296) (1.9106,-0.5338) (1.8994,-0.5381) (1.8883,-0.5424) (1.8771,-0.5467) (1.8659,-0.5510) (1.8547,-0.5552) (1.8436,-0.5595) (1.8324,-0.5638) (1.8212,-0.5681) (1.8101,-0.5724) (1.7989,-0.5767) (1.7877,-0.5810) (1.7765,-0.5853) (1.7654,-0.5896) (1.7542,-0.5939) (1.7430,-0.5982) (1.7318,-0.6025) (1.7207,-0.6068) (1.7095,-0.6111) (1.6983,-0.6153) (1.6872,-0.6196) (1.6760,-0.6239) (1.6648,-0.6281) (1.6536,-0.6324) (1.6425,-0.6366) (1.6313,-0.6408) (1.6201,-0.6451) (1.6089,-0.6493) (1.5978,-0.6534) (1.5866,-0.6576) (1.5754,-0.6618) (1.5642,-0.6659) (1.5531,-0.6701) (1.5419,-0.6742) (1.5307,-0.6783) (1.5196,-0.6824) (1.5084,-0.6864) (1.4972,-0.6905) (1.4860,-0.6945) (1.4749,-0.6985) (1.4637,-0.7025) (1.4525,-0.7065) (1.4413,-0.7104) (1.4302,-0.7143) (1.4190,-0.7182) (1.4078,-0.7221) (1.3966,-0.7259) (1.3855,-0.7298) (1.3743,-0.7336) (1.3631,-0.7373) (1.3520,-0.7411) (1.3408,-0.7448) (1.3296,-0.7485) (1.3184,-0.7522) (1.3073,-0.7558) (1.2961,-0.7594) (1.2849,-0.7630) (1.2737,-0.7665) (1.2626,-0.7700) (1.2514,-0.7735) (1.2402,-0.7770) (1.2291,-0.7804) (1.2179,-0.7838) (1.2067,-0.7872) (1.1955,-0.7905) (1.1844,-0.7938) (1.1732,-0.7971) (1.1620,-0.8003) (1.1508,-0.8035) (1.1397,-0.8067) (1.1285,-0.8098) (1.1173,-0.8129) (1.1061,-0.8159) (1.0950,-0.8190) (1.0838,-0.8220) (1.0726,-0.8249) (1.0615,-0.8279) (1.0503,-0.8308) (1.0391,-0.8336) (1.0279,-0.8364) (1.0168,-0.8392) (1.0056,-0.8420) (0.9944,-0.8447) (0.9832,-0.8474) (0.9721,-0.8500) (0.9609,-0.8527) (0.9497,-0.8552) (0.9385,-0.8578) (0.9274,-0.8603) (0.9162,-0.8628) (0.9050,-0.8652) (0.8939,-0.8676) (0.8827,-0.8700) (0.8715,-0.8724) (0.8603,-0.8747) (0.8492,-0.8770) (0.8380,-0.8792) (0.8268,-0.8814) (0.8156,-0.8836) (0.8045,-0.8858) (0.7933,-0.8879) (0.7821,-0.8900) (0.7709,-0.8920) (0.7598,-0.8940) (0.7486,-0.8960) (0.7374,-0.8980) (0.7263,-0.8999) (0.7151,-0.9018) (0.7039,-0.9037) (0.6927,-0.9055) (0.6816,-0.9073) (0.6704,-0.9091) (0.6592,-0.9108) (0.6480,-0.9126) (0.6369,-0.9142) (0.6257,-0.9159) (0.6145,-0.9175) (0.6034,-0.9192) (0.5922,-0.9207) (0.5810,-0.9223) (0.5698,-0.9238) (0.5587,-0.9253) (0.5475,-0.9268) (0.5363,-0.9282) (0.5251,-0.9297) (0.5140,-0.9311) (0.5028,-0.9324) (0.4916,-0.9338) (0.4804,-0.9351) (0.4693,-0.9364) (0.4581,-0.9377) (0.4469,-0.9389) (0.4358,-0.9402) (0.4246,-0.9414) (0.4134,-0.9426) (0.4022,-0.9437) (0.3911,-0.9449) (0.3799,-0.9460) (0.3687,-0.9471) (0.3575,-0.9482) (0.3464,-0.9492) (0.3352,-0.9503) (0.3240,-0.9513) (0.3128,-0.9523) (0.3017,-0.9533) (0.2905,-0.9542) (0.2793,-0.9552) (0.2682,-0.9561) (0.2570,-0.9570) (0.2458,-0.9579) (0.2346,-0.9588) (0.2235,-0.9596) (0.2123,-0.9605) (0.2011,-0.9613) (0.1899,-0.9621) (0.1788,-0.9629) (0.1676,-0.9637) (0.1564,-0.9644) (0.1453,-0.9652) (0.1341,-0.9659) (0.1229,-0.9666) (0.1117,-0.9673) (0.1006,-0.9680) (0.0894,-0.9687) (0.0782,-0.9693) (0.0670,-0.9700) (0.0559,-0.9706) (0.0447,-0.9712) (0.0335,-0.9719) (0.0223,-0.9725) (0.0112,-0.9730) (0.0000,-0.9736)};\draw[main,thick]plot coordinates{(2.0000,-0.5000) (2.0168,-0.4937) (2.0335,-0.4875) (2.0503,-0.4813) (2.0670,-0.4751) (2.0838,-0.4689) (2.1006,-0.4628) (2.1173,-0.4568) (2.1341,-0.4507) (2.1508,-0.4447) (2.1676,-0.4388) (2.1844,-0.4328) (2.2011,-0.4270) (2.2179,-0.4211) (2.2346,-0.4154) (2.2514,-0.4096) (2.2682,-0.4039) (2.2849,-0.3983) (2.3017,-0.3927) (2.3184,-0.3872) (2.3352,-0.3817) (2.3520,-0.3762) (2.3687,-0.3708) (2.3855,-0.3655) (2.4022,-0.3602) (2.4190,-0.3550) (2.4358,-0.3498) (2.4525,-0.3447) (2.4693,-0.3396) (2.4860,-0.3346) (2.5028,-0.3297) (2.5196,-0.3248) (2.5363,-0.3199) (2.5531,-0.3152) (2.5698,-0.3104) (2.5866,-0.3058) (2.6034,-0.3011) (2.6201,-0.2966) (2.6369,-0.2921) (2.6536,-0.2876) (2.6704,-0.2832) (2.6872,-0.2789) (2.7039,-0.2746) (2.7207,-0.2704) (2.7374,-0.2662) (2.7542,-0.2621) (2.7709,-0.2580) (2.7877,-0.2540) (2.8045,-0.2501) (2.8212,-0.2462) (2.8380,-0.2423) (2.8547,-0.2385) (2.8715,-0.2348) (2.8883,-0.2311) (2.9050,-0.2274) (2.9218,-0.2238) (2.9385,-0.2203) (2.9553,-0.2168) (2.9721,-0.2134) (2.9888,-0.2100) (3.0056,-0.2067) (3.0223,-0.2034) (3.0391,-0.2001) (3.0559,-0.1969) (3.0726,-0.1938) (3.0894,-0.1907) (3.1061,-0.1876) (3.1229,-0.1846) (3.1397,-0.1816) (3.1564,-0.1787) (3.1732,-0.1758) (3.1899,-0.1730) (3.2067,-0.1702) (3.2235,-0.1675) (3.2402,-0.1648) (3.2570,-0.1621) (3.2737,-0.1595) (3.2905,-0.1569) (3.3073,-0.1543) (3.3240,-0.1518) (3.3408,-0.1494) (3.3575,-0.1469) (3.3743,-0.1445) (3.3911,-0.1422) (3.4078,-0.1399) (3.4246,-0.1376) (3.4413,-0.1353) (3.4581,-0.1331) (3.4749,-0.1310) (3.4916,-0.1288) (3.5084,-0.1267) (3.5251,-0.1246) (3.5419,-0.1226) (3.5587,-0.1206) (3.5754,-0.1186) (3.5922,-0.1167) (3.6089,-0.1148) (3.6257,-0.1129) (3.6425,-0.1110) (3.6592,-0.1092) (3.6760,-0.1074) (3.6927,-0.1056) (3.7095,-0.1039) (3.7263,-0.1022) (3.7430,-0.1005) (3.7598,-0.0989) (3.7765,-0.0972) (3.7933,-0.0956) (3.8101,-0.0941) (3.8268,-0.0925) (3.8436,-0.0910) (3.8603,-0.0895) (3.8771,-0.0880) (3.8939,-0.0866) (3.9106,-0.0851) (3.9274,-0.0837) (3.9441,-0.0823) (3.9609,-0.0810) (3.9777,-0.0796) (3.9944,-0.0783) (4.0112,-0.0770) (4.0279,-0.0758) (4.0447,-0.0745) (4.0615,-0.0733) (4.0782,-0.0721) (4.0950,-0.0709) (4.1117,-0.0697) (4.1285,-0.0686) (4.1453,-0.0674) (4.1620,-0.0663) (4.1788,-0.0652) (4.1955,-0.0641) (4.2123,-0.0631) (4.2291,-0.0620) (4.2458,-0.0610) (4.2626,-0.0600) (4.2793,-0.0590) (4.2961,-0.0580) (4.3128,-0.0571) (4.3296,-0.0561) (4.3464,-0.0552) (4.3631,-0.0543) (4.3799,-0.0534) (4.3966,-0.0525) (4.4134,-0.0516) (4.4302,-0.0508) (4.4469,-0.0499) (4.4637,-0.0491) (4.4804,-0.0483) (4.4972,-0.0475) (4.5140,-0.0467) (4.5307,-0.0459) (4.5475,-0.0451) (4.5642,-0.0444) (4.5810,-0.0437) (4.5978,-0.0429) (4.6145,-0.0422) (4.6313,-0.0415) (4.6480,-0.0408) (4.6648,-0.0402) (4.6816,-0.0395) (4.6983,-0.0388) (4.7151,-0.0382) (4.7318,-0.0376) (4.7486,-0.0369) (4.7654,-0.0363) (4.7821,-0.0357) (4.7989,-0.0351) (4.8156,-0.0345) (4.8324,-0.0340) (4.8492,-0.0334) (4.8659,-0.0329) (4.8827,-0.0323) (4.8994,-0.0318) (4.9162,-0.0312) (4.9330,-0.0307) (4.9497,-0.0302) (4.9665,-0.0297) (4.9832,-0.0292) (5.0000,-0.0287)};\draw[main,thick]plot coordinates{(2.0000,0.5000) (1.9888,0.5042) (1.9777,0.5084) (1.9665,0.5126) (1.9553,0.5168) (1.9441,0.5211) (1.9330,0.5253) (1.9218,0.5296) (1.9106,0.5338) (1.8994,0.5381) (1.8883,0.5424) (1.8771,0.5467) (1.8659,0.5510) (1.8547,0.5552) (1.8436,0.5595) (1.8324,0.5638) (1.8212,0.5681) (1.8101,0.5724) (1.7989,0.5767) (1.7877,0.5810) (1.7765,0.5853) (1.7654,0.5896) (1.7542,0.5939) (1.7430,0.5982) (1.7318,0.6025) (1.7207,0.6068) (1.7095,0.6111) (1.6983,0.6153) (1.6872,0.6196) (1.6760,0.6239) (1.6648,0.6281) (1.6536,0.6324) (1.6425,0.6366) (1.6313,0.6408) (1.6201,0.6451) (1.6089,0.6493) (1.5978,0.6534) (1.5866,0.6576) (1.5754,0.6618) (1.5642,0.6659) (1.5531,0.6701) (1.5419,0.6742) (1.5307,0.6783) (1.5196,0.6824) (1.5084,0.6864) (1.4972,0.6905) (1.4860,0.6945) (1.4749,0.6985) (1.4637,0.7025) (1.4525,0.7065) (1.4413,0.7104) (1.4302,0.7143) (1.4190,0.7182) (1.4078,0.7221) (1.3966,0.7259) (1.3855,0.7298) (1.3743,0.7336) (1.3631,0.7373) (1.3520,0.7411) (1.3408,0.7448) (1.3296,0.7485) (1.3184,0.7522) (1.3073,0.7558) (1.2961,0.7594) (1.2849,0.7630) (1.2737,0.7665) (1.2626,0.7700) (1.2514,0.7735) (1.2402,0.7770) (1.2291,0.7804) (1.2179,0.7838) (1.2067,0.7872) (1.1955,0.7905) (1.1844,0.7938) (1.1732,0.7971) (1.1620,0.8003) (1.1508,0.8035) (1.1397,0.8067) (1.1285,0.8098) (1.1173,0.8129) (1.1061,0.8159) (1.0950,0.8190) (1.0838,0.8220) (1.0726,0.8249) (1.0615,0.8279) (1.0503,0.8308) (1.0391,0.8336) (1.0279,0.8364) (1.0168,0.8392) (1.0056,0.8420) (0.9944,0.8447) (0.9832,0.8474) (0.9721,0.8500) (0.9609,0.8527) (0.9497,0.8552) (0.9385,0.8578) (0.9274,0.8603) (0.9162,0.8628) (0.9050,0.8652) (0.8939,0.8676) (0.8827,0.8700) (0.8715,0.8724) (0.8603,0.8747) (0.8492,0.8770) (0.8380,0.8792) (0.8268,0.8814) (0.8156,0.8836) (0.8045,0.8858) (0.7933,0.8879) (0.7821,0.8900) (0.7709,0.8920) (0.7598,0.8940) (0.7486,0.8960) (0.7374,0.8980) (0.7263,0.8999) (0.7151,0.9018) (0.7039,0.9037) (0.6927,0.9055) (0.6816,0.9073) (0.6704,0.9091) (0.6592,0.9108) (0.6480,0.9126) (0.6369,0.9142) (0.6257,0.9159) (0.6145,0.9175) (0.6034,0.9192) (0.5922,0.9207) (0.5810,0.9223) (0.5698,0.9238) (0.5587,0.9253) (0.5475,0.9268) (0.5363,0.9282) (0.5251,0.9297) (0.5140,0.9311) (0.5028,0.9324) (0.4916,0.9338) (0.4804,0.9351) (0.4693,0.9364) (0.4581,0.9377) (0.4469,0.9389) (0.4358,0.9402) (0.4246,0.9414) (0.4134,0.9426) (0.4022,0.9437) (0.3911,0.9449) (0.3799,0.9460) (0.3687,0.9471) (0.3575,0.9482) (0.3464,0.9492) (0.3352,0.9503) (0.3240,0.9513) (0.3128,0.9523) (0.3017,0.9533) (0.2905,0.9542) (0.2793,0.9552) (0.2682,0.9561) (0.2570,0.9570) (0.2458,0.9579) (0.2346,0.9588) (0.2235,0.9596) (0.2123,0.9605) (0.2011,0.9613) (0.1899,0.9621) (0.1788,0.9629) (0.1676,0.9637) (0.1564,0.9644) (0.1453,0.9652) (0.1341,0.9659) (0.1229,0.9666) (0.1117,0.9673) (0.1006,0.9680) (0.0894,0.9687) (0.0782,0.9693) (0.0670,0.9700) (0.0559,0.9706) (0.0447,0.9712) (0.0335,0.9719) (0.0223,0.9725) (0.0112,0.9730) (0.0000,0.9736)};\draw[main,thick]plot coordinates{(2.0000,0.5000) (2.0168,0.4937) (2.0335,0.4875) (2.0503,0.4813) (2.0670,0.4751) (2.0838,0.4689) (2.1006,0.4628) (2.1173,0.4568) (2.1341,0.4507) (2.1508,0.4447) (2.1676,0.4388) (2.1844,0.4328) (2.2011,0.4270) (2.2179,0.4211) (2.2346,0.4154) (2.2514,0.4096) (2.2682,0.4039) (2.2849,0.3983) (2.3017,0.3927) (2.3184,0.3872) (2.3352,0.3817) (2.3520,0.3762) (2.3687,0.3708) (2.3855,0.3655) (2.4022,0.3602) (2.4190,0.3550) (2.4358,0.3498) (2.4525,0.3447) (2.4693,0.3396) (2.4860,0.3346) (2.5028,0.3297) (2.5196,0.3248) (2.5363,0.3199) (2.5531,0.3152) (2.5698,0.3104) (2.5866,0.3058) (2.6034,0.3011) (2.6201,0.2966) (2.6369,0.2921) (2.6536,0.2876) (2.6704,0.2832) (2.6872,0.2789) (2.7039,0.2746) (2.7207,0.2704) (2.7374,0.2662) (2.7542,0.2621) (2.7709,0.2580) (2.7877,0.2540) (2.8045,0.2501) (2.8212,0.2462) (2.8380,0.2423) (2.8547,0.2385) (2.8715,0.2348) (2.8883,0.2311) (2.9050,0.2274) (2.9218,0.2238) (2.9385,0.2203) (2.9553,0.2168) (2.9721,0.2134) (2.9888,0.2100) (3.0056,0.2067) (3.0223,0.2034) (3.0391,0.2001) (3.0559,0.1969) (3.0726,0.1938) (3.0894,0.1907) (3.1061,0.1876) (3.1229,0.1846) (3.1397,0.1816) (3.1564,0.1787) (3.1732,0.1758) (3.1899,0.1730) (3.2067,0.1702) (3.2235,0.1675) (3.2402,0.1648) (3.2570,0.1621) (3.2737,0.1595) (3.2905,0.1569) (3.3073,0.1543) (3.3240,0.1518) (3.3408,0.1494) (3.3575,0.1469) (3.3743,0.1445) (3.3911,0.1422) (3.4078,0.1399) (3.4246,0.1376) (3.4413,0.1353) (3.4581,0.1331) (3.4749,0.1310) (3.4916,0.1288) (3.5084,0.1267) (3.5251,0.1246) (3.5419,0.1226) (3.5587,0.1206) (3.5754,0.1186) (3.5922,0.1167) (3.6089,0.1148) (3.6257,0.1129) (3.6425,0.1110) (3.6592,0.1092) (3.6760,0.1074) (3.6927,0.1056) (3.7095,0.1039) (3.7263,0.1022) (3.7430,0.1005) (3.7598,0.0989) (3.7765,0.0972) (3.7933,0.0956) (3.8101,0.0941) (3.8268,0.0925) (3.8436,0.0910) (3.8603,0.0895) (3.8771,0.0880) (3.8939,0.0866) (3.9106,0.0851) (3.9274,0.0837) (3.9441,0.0823) (3.9609,0.0810) (3.9777,0.0796) (3.9944,0.0783) (4.0112,0.0770) (4.0279,0.0758) (4.0447,0.0745) (4.0615,0.0733) (4.0782,0.0721) (4.0950,0.0709) (4.1117,0.0697) (4.1285,0.0686) (4.1453,0.0674) (4.1620,0.0663) (4.1788,0.0652) (4.1955,0.0641) (4.2123,0.0631) (4.2291,0.0620) (4.2458,0.0610) (4.2626,0.0600) (4.2793,0.0590) (4.2961,0.0580) (4.3128,0.0571) (4.3296,0.0561) (4.3464,0.0552) (4.3631,0.0543) (4.3799,0.0534) (4.3966,0.0525) (4.4134,0.0516) (4.4302,0.0508) (4.4469,0.0499) (4.4637,0.0491) (4.4804,0.0483) (4.4972,0.0475) (4.5140,0.0467) (4.5307,0.0459) (4.5475,0.0451) (4.5642,0.0444) (4.5810,0.0437) (4.5978,0.0429) (4.6145,0.0422) (4.6313,0.0415) (4.6480,0.0408) (4.6648,0.0402) (4.6816,0.0395) (4.6983,0.0388) (4.7151,0.0382) (4.7318,0.0376) (4.7486,0.0369) (4.7654,0.0363) (4.7821,0.0357) (4.7989,0.0351) (4.8156,0.0345) (4.8324,0.0340) (4.8492,0.0334) (4.8659,0.0329) (4.8827,0.0323) (4.8994,0.0318) (4.9162,0.0312) (4.9330,0.0307) (4.9497,0.0302) (4.9665,0.0297) (4.9832,0.0292) (5.0000,0.0287)};\draw[main,thick]plot coordinates{(2.0000,1.5000) (1.9888,1.4797) (1.9777,1.4606) (1.9665,1.4427) (1.9553,1.4257) (1.9441,1.4097) (1.9330,1.3946) (1.9218,1.3803) (1.9106,1.3667) (1.8994,1.3538) (1.8883,1.3415) (1.8771,1.3298) (1.8659,1.3186) (1.8547,1.3080) (1.8436,1.2978) (1.8324,1.2881) (1.8212,1.2789) (1.8101,1.2700) (1.7989,1.2615) (1.7877,1.2533) (1.7765,1.2455) (1.7654,1.2379) (1.7542,1.2307) (1.7430,1.2238) (1.7318,1.2171) (1.7207,1.2107) (1.7095,1.2045) (1.6983,1.1986) (1.6872,1.1928) (1.6760,1.1873) (1.6648,1.1819) (1.6536,1.1768) (1.6425,1.1718) (1.6313,1.1670) (1.6201,1.1624) (1.6089,1.1579) (1.5978,1.1536) (1.5866,1.1494) (1.5754,1.1453) (1.5642,1.1414) (1.5531,1.1376) (1.5419,1.1339) (1.5307,1.1303) (1.5196,1.1269) (1.5084,1.1235) (1.4972,1.1203) (1.4860,1.1172) (1.4749,1.1141) (1.4637,1.1112) (1.4525,1.1083) (1.4413,1.1055) (1.4302,1.1028) (1.4190,1.1002) (1.4078,1.0976) (1.3966,1.0951) (1.3855,1.0927) (1.3743,1.0904) (1.3631,1.0881) (1.3520,1.0859) (1.3408,1.0838) (1.3296,1.0817) (1.3184,1.0797) (1.3073,1.0777) (1.2961,1.0758) (1.2849,1.0739) (1.2737,1.0721) (1.2626,1.0703) (1.2514,1.0686) (1.2402,1.0670) (1.2291,1.0653) (1.2179,1.0637) (1.2067,1.0622) (1.1955,1.0607) (1.1844,1.0592) (1.1732,1.0578) (1.1620,1.0564) (1.1508,1.0551) (1.1397,1.0537) (1.1285,1.0525) (1.1173,1.0512) (1.1061,1.0500) (1.0950,1.0488) (1.0838,1.0477) (1.0726,1.0465) (1.0615,1.0454) (1.0503,1.0444) (1.0391,1.0433) (1.0279,1.0423) (1.0168,1.0413) (1.0056,1.0403) (0.9944,1.0394) (0.9832,1.0385) (0.9721,1.0376) (0.9609,1.0367) (0.9497,1.0358) (0.9385,1.0350) (0.9274,1.0342) (0.9162,1.0334) (0.9050,1.0326) (0.8939,1.0319) (0.8827,1.0311) (0.8715,1.0304) (0.8603,1.0297) (0.8492,1.0290) (0.8380,1.0284) (0.8268,1.0277) (0.8156,1.0271) (0.8045,1.0264) (0.7933,1.0258) (0.7821,1.0252) (0.7709,1.0247) (0.7598,1.0241) (0.7486,1.0235) (0.7374,1.0230) (0.7263,1.0225) (0.7151,1.0220) (0.7039,1.0215) (0.6927,1.0210) (0.6816,1.0205) (0.6704,1.0200) (0.6592,1.0196) (0.6480,1.0191) (0.6369,1.0187) (0.6257,1.0183) (0.6145,1.0179) (0.6034,1.0175) (0.5922,1.0171) (0.5810,1.0167) (0.5698,1.0163) (0.5587,1.0159) (0.5475,1.0156) (0.5363,1.0152) (0.5251,1.0149) (0.5140,1.0145) (0.5028,1.0142) (0.4916,1.0139) (0.4804,1.0136) (0.4693,1.0133) (0.4581,1.0130) (0.4469,1.0127) (0.4358,1.0124) (0.4246,1.0121) (0.4134,1.0118) (0.4022,1.0116) (0.3911,1.0113) (0.3799,1.0111) (0.3687,1.0108) (0.3575,1.0106) (0.3464,1.0103) (0.3352,1.0101) (0.3240,1.0099) (0.3128,1.0096) (0.3017,1.0094) (0.2905,1.0092) (0.2793,1.0090) (0.2682,1.0088) (0.2570,1.0086) (0.2458,1.0084) (0.2346,1.0082) (0.2235,1.0081) (0.2123,1.0079) (0.2011,1.0077) (0.1899,1.0075) (0.1788,1.0074) (0.1676,1.0072) (0.1564,1.0070) (0.1453,1.0069) (0.1341,1.0067) (0.1229,1.0066) (0.1117,1.0064) (0.1006,1.0063) (0.0894,1.0061) (0.0782,1.0060) (0.0670,1.0059) (0.0559,1.0057) (0.0447,1.0056) (0.0335,1.0055) (0.0223,1.0054) (0.0112,1.0052) (0.0000,1.0051)};\draw[main,thick]plot coordinates{(2.0000,1.5000) (2.0016,1.5031) (2.0033,1.5062) (2.0049,1.5094) (2.0066,1.5126) (2.0082,1.5158) (2.0099,1.5190) (2.0115,1.5223) (2.0131,1.5256) (2.0148,1.5289) (2.0164,1.5323) (2.0181,1.5357) (2.0197,1.5392) (2.0213,1.5426) (2.0230,1.5462) (2.0246,1.5497) (2.0263,1.5533) (2.0279,1.5569) (2.0296,1.5606) (2.0312,1.5643) (2.0328,1.5680) (2.0345,1.5718) (2.0361,1.5756) (2.0378,1.5794) (2.0394,1.5833) (2.0410,1.5873) (2.0427,1.5913) (2.0443,1.5953) (2.0460,1.5994) (2.0476,1.6035) (2.0493,1.6076) (2.0509,1.6118) (2.0525,1.6161) (2.0542,1.6204) (2.0558,1.6247) (2.0575,1.6291) (2.0591,1.6336) (2.0607,1.6381) (2.0624,1.6426) (2.0640,1.6472) (2.0657,1.6519) (2.0673,1.6566) (2.0690,1.6614) (2.0706,1.6662) (2.0722,1.6711) (2.0739,1.6761) (2.0755,1.6811) (2.0772,1.6861) (2.0788,1.6913) (2.0805,1.6965) (2.0821,1.7017) (2.0837,1.7071) (2.0854,1.7125) (2.0870,1.7179) (2.0887,1.7235) (2.0903,1.7291) (2.0919,1.7348) (2.0936,1.7405) (2.0952,1.7464) (2.0969,1.7523) (2.0985,1.7583) (2.1002,1.7644) (2.1018,1.7705) (2.1034,1.7768) (2.1051,1.7831) (2.1067,1.7895) (2.1084,1.7961) (2.1100,1.8027) (2.1116,1.8094) (2.1133,1.8162) (2.1149,1.8231) (2.1166,1.8301) (2.1182,1.8372) (2.1199,1.8444) (2.1215,1.8518) (2.1231,1.8592) (2.1248,1.8667) (2.1264,1.8744) (2.1281,1.8822) (2.1297,1.8901) (2.1313,1.8982) (2.1330,1.9064) (2.1346,1.9147) (2.1363,1.9231) (2.1379,1.9317) (2.1396,1.9405) (2.1412,1.9493) (2.1428,1.9584) (2.1445,1.9676) (2.1461,1.9769) (2.1478,1.9865) (2.1494,1.9962)};\end{tikzpicture}\end{center}
}

\newcommand{\ODEExamPracticeFinalPhase}{%
\begin{center}\begin{tikzpicture}[x=.8cm,y=.7cm]\draw[->](-1,-1)--(-1,3.5)node[above]{$x$};\draw[->](-1,-0.35)--(-1,-0.65);\draw[->](-1,0.35)--(-1,0.65);\draw[->](-1,1.65)--(-1,1.35);\draw[->](-1,2.35)--(-1,2.65);\fill(-1,0)circle(1.8pt);\node[left]at(-1,0){$ 0 $};\draw[dashed](0,0)--(5,0);\fill(-1,1)circle(1.8pt);\node[left]at(-1,1){$ 1 $};\draw[dashed](0,1)--(5,1);\fill(-1,2)circle(1.8pt);\node[left]at(-1,2){$ 2 $};\draw[dashed](0,2)--(5,2);\draw[->](0,-1)--(5.4,-1)node[right]{$t$};\draw[->](0,-1)--(0,3.3)node[above]{$x$};\draw[main,thick]plot coordinates{(2.0000,-0.5000) (1.9888,-0.4797) (1.9777,-0.4606) (1.9665,-0.4427) (1.9553,-0.4257) (1.9441,-0.4097) (1.9330,-0.3946) (1.9218,-0.3803) (1.9106,-0.3667) (1.8994,-0.3538) (1.8883,-0.3415) (1.8771,-0.3298) (1.8659,-0.3186) (1.8547,-0.3080) (1.8436,-0.2978) (1.8324,-0.2881) (1.8212,-0.2789) (1.8101,-0.2700) (1.7989,-0.2615) (1.7877,-0.2533) (1.7765,-0.2455) (1.7654,-0.2379) (1.7542,-0.2307) (1.7430,-0.2238) (1.7318,-0.2171) (1.7207,-0.2107) (1.7095,-0.2045) (1.6983,-0.1986) (1.6872,-0.1928) (1.6760,-0.1873) (1.6648,-0.1819) (1.6536,-0.1768) (1.6425,-0.1718) (1.6313,-0.1670) (1.6201,-0.1624) (1.6089,-0.1579) (1.5978,-0.1536) (1.5866,-0.1494) (1.5754,-0.1453) (1.5642,-0.1414) (1.5531,-0.1376) (1.5419,-0.1339) (1.5307,-0.1303) (1.5196,-0.1269) (1.5084,-0.1235) (1.4972,-0.1203) (1.4860,-0.1172) (1.4749,-0.1141) (1.4637,-0.1112) (1.4525,-0.1083) (1.4413,-0.1055) (1.4302,-0.1028) (1.4190,-0.1002) (1.4078,-0.0976) (1.3966,-0.0951) (1.3855,-0.0927) (1.3743,-0.0904) (1.3631,-0.0881) (1.3520,-0.0859) (1.3408,-0.0838) (1.3296,-0.0817) (1.3184,-0.0797) (1.3073,-0.0777) (1.2961,-0.0758) (1.2849,-0.0739) (1.2737,-0.0721) (1.2626,-0.0703) (1.2514,-0.0686) (1.2402,-0.0670) (1.2291,-0.0653) (1.2179,-0.0637) (1.2067,-0.0622) (1.1955,-0.0607) (1.1844,-0.0592) (1.1732,-0.0578) (1.1620,-0.0564) (1.1508,-0.0551) (1.1397,-0.0537) (1.1285,-0.0525) (1.1173,-0.0512) (1.1061,-0.0500) (1.0950,-0.0488) (1.0838,-0.0477) (1.0726,-0.0465) (1.0615,-0.0454) (1.0503,-0.0444) (1.0391,-0.0433) (1.0279,-0.0423) (1.0168,-0.0413) (1.0056,-0.0403) (0.9944,-0.0394) (0.9832,-0.0385) (0.9721,-0.0376) (0.9609,-0.0367) (0.9497,-0.0358) (0.9385,-0.0350) (0.9274,-0.0342) (0.9162,-0.0334) (0.9050,-0.0326) (0.8939,-0.0319) (0.8827,-0.0311) (0.8715,-0.0304) (0.8603,-0.0297) (0.8492,-0.0290) (0.8380,-0.0284) (0.8268,-0.0277) (0.8156,-0.0271) (0.8045,-0.0264) (0.7933,-0.0258) (0.7821,-0.0252) (0.7709,-0.0247) (0.7598,-0.0241) (0.7486,-0.0235) (0.7374,-0.0230) (0.7263,-0.0225) (0.7151,-0.0220) (0.7039,-0.0215) (0.6927,-0.0210) (0.6816,-0.0205) (0.6704,-0.0200) (0.6592,-0.0196) (0.6480,-0.0191) (0.6369,-0.0187) (0.6257,-0.0183) (0.6145,-0.0179) (0.6034,-0.0175) (0.5922,-0.0171) (0.5810,-0.0167) (0.5698,-0.0163) (0.5587,-0.0159) (0.5475,-0.0156) (0.5363,-0.0152) (0.5251,-0.0149) (0.5140,-0.0145) (0.5028,-0.0142) (0.4916,-0.0139) (0.4804,-0.0136) (0.4693,-0.0133) (0.4581,-0.0130) (0.4469,-0.0127) (0.4358,-0.0124) (0.4246,-0.0121) (0.4134,-0.0118) (0.4022,-0.0116) (0.3911,-0.0113) (0.3799,-0.0111) (0.3687,-0.0108) (0.3575,-0.0106) (0.3464,-0.0103) (0.3352,-0.0101) (0.3240,-0.0099) (0.3128,-0.0096) (0.3017,-0.0094) (0.2905,-0.0092) (0.2793,-0.0090) (0.2682,-0.0088) (0.2570,-0.0086) (0.2458,-0.0084) (0.2346,-0.0082) (0.2235,-0.0081) (0.2123,-0.0079) (0.2011,-0.0077) (0.1899,-0.0075) (0.1788,-0.0074) (0.1676,-0.0072) (0.1564,-0.0070) (0.1453,-0.0069) (0.1341,-0.0067) (0.1229,-0.0066) (0.1117,-0.0064) (0.1006,-0.0063) (0.0894,-0.0061) (0.0782,-0.0060) (0.0670,-0.0059) (0.0559,-0.0057) (0.0447,-0.0056) (0.0335,-0.0055) (0.0223,-0.0054) (0.0112,-0.0052) (0.0000,-0.0051)};\draw[main,thick]plot coordinates{(2.0000,-0.5000) (2.0016,-0.5031) (2.0033,-0.5062) (2.0049,-0.5094) (2.0066,-0.5126) (2.0082,-0.5158) (2.0099,-0.5190) (2.0115,-0.5223) (2.0131,-0.5256) (2.0148,-0.5289) (2.0164,-0.5323) (2.0181,-0.5357) (2.0197,-0.5392) (2.0213,-0.5426) (2.0230,-0.5462) (2.0246,-0.5497) (2.0263,-0.5533) (2.0279,-0.5569) (2.0296,-0.5606) (2.0312,-0.5643) (2.0328,-0.5680) (2.0345,-0.5718) (2.0361,-0.5756) (2.0378,-0.5794) (2.0394,-0.5833) (2.0410,-0.5873) (2.0427,-0.5913) (2.0443,-0.5953) (2.0460,-0.5994) (2.0476,-0.6035) (2.0493,-0.6076) (2.0509,-0.6118) (2.0525,-0.6161) (2.0542,-0.6204) (2.0558,-0.6247) (2.0575,-0.6291) (2.0591,-0.6336) (2.0607,-0.6381) (2.0624,-0.6426) (2.0640,-0.6472) (2.0657,-0.6519) (2.0673,-0.6566) (2.0690,-0.6614) (2.0706,-0.6662) (2.0722,-0.6711) (2.0739,-0.6761) (2.0755,-0.6811) (2.0772,-0.6861) (2.0788,-0.6913) (2.0805,-0.6965) (2.0821,-0.7017) (2.0837,-0.7071) (2.0854,-0.7125) (2.0870,-0.7179) (2.0887,-0.7235) (2.0903,-0.7291) (2.0919,-0.7348) (2.0936,-0.7405) (2.0952,-0.7464) (2.0969,-0.7523) (2.0985,-0.7583) (2.1002,-0.7644) (2.1018,-0.7705) (2.1034,-0.7768) (2.1051,-0.7831) (2.1067,-0.7895) (2.1084,-0.7961) (2.1100,-0.8027) (2.1116,-0.8094) (2.1133,-0.8162) (2.1149,-0.8231) (2.1166,-0.8301) (2.1182,-0.8372) (2.1199,-0.8444) (2.1215,-0.8518) (2.1231,-0.8592) (2.1248,-0.8668) (2.1264,-0.8744) (2.1281,-0.8822) (2.1297,-0.8901) (2.1313,-0.8982) (2.1330,-0.9064) (2.1346,-0.9147) (2.1363,-0.9231) (2.1379,-0.9317) (2.1396,-0.9405) (2.1412,-0.9493) (2.1428,-0.9584) (2.1445,-0.9676) (2.1461,-0.9769) (2.1478,-0.9865) (2.1494,-0.9962)};\draw[main,thick]plot coordinates{(2.0000,0.5000) (1.9888,0.4958) (1.9777,0.4916) (1.9665,0.4874) (1.9553,0.4832) (1.9441,0.4789) (1.9330,0.4747) (1.9218,0.4704) (1.9106,0.4662) (1.8994,0.4619) (1.8883,0.4576) (1.8771,0.4533) (1.8659,0.4490) (1.8547,0.4448) (1.8436,0.4405) (1.8324,0.4362) (1.8212,0.4319) (1.8101,0.4276) (1.7989,0.4233) (1.7877,0.4190) (1.7765,0.4147) (1.7654,0.4104) (1.7542,0.4061) (1.7430,0.4018) (1.7318,0.3975) (1.7207,0.3932) (1.7095,0.3889) (1.6983,0.3847) (1.6872,0.3804) (1.6760,0.3761) (1.6648,0.3719) (1.6536,0.3676) (1.6425,0.3634) (1.6313,0.3592) (1.6201,0.3549) (1.6089,0.3507) (1.5978,0.3466) (1.5866,0.3424) (1.5754,0.3382) (1.5642,0.3341) (1.5531,0.3299) (1.5419,0.3258) (1.5307,0.3217) (1.5196,0.3176) (1.5084,0.3136) (1.4972,0.3095) (1.4860,0.3055) (1.4749,0.3015) (1.4637,0.2975) (1.4525,0.2935) (1.4413,0.2896) (1.4302,0.2857) (1.4190,0.2818) (1.4078,0.2779) (1.3966,0.2741) (1.3855,0.2702) (1.3743,0.2664) (1.3631,0.2627) (1.3520,0.2589) (1.3408,0.2552) (1.3296,0.2515) (1.3184,0.2478) (1.3073,0.2442) (1.2961,0.2406) (1.2849,0.2370) (1.2737,0.2335) (1.2626,0.2300) (1.2514,0.2265) (1.2402,0.2230) (1.2291,0.2196) (1.2179,0.2162) (1.2067,0.2128) (1.1955,0.2095) (1.1844,0.2062) (1.1732,0.2029) (1.1620,0.1997) (1.1508,0.1965) (1.1397,0.1933) (1.1285,0.1902) (1.1173,0.1871) (1.1061,0.1841) (1.0950,0.1810) (1.0838,0.1780) (1.0726,0.1751) (1.0615,0.1721) (1.0503,0.1692) (1.0391,0.1664) (1.0279,0.1636) (1.0168,0.1608) (1.0056,0.1580) (0.9944,0.1553) (0.9832,0.1526) (0.9721,0.1500) (0.9609,0.1473) (0.9497,0.1448) (0.9385,0.1422) (0.9274,0.1397) (0.9162,0.1372) (0.9050,0.1348) (0.8939,0.1324) (0.8827,0.1300) (0.8715,0.1276) (0.8603,0.1253) (0.8492,0.1230) (0.8380,0.1208) (0.8268,0.1186) (0.8156,0.1164) (0.8045,0.1142) (0.7933,0.1121) (0.7821,0.1100) (0.7709,0.1080) (0.7598,0.1060) (0.7486,0.1040) (0.7374,0.1020) (0.7263,0.1001) (0.7151,0.0982) (0.7039,0.0963) (0.6927,0.0945) (0.6816,0.0927) (0.6704,0.0909) (0.6592,0.0892) (0.6480,0.0874) (0.6369,0.0858) (0.6257,0.0841) (0.6145,0.0825) (0.6034,0.0808) (0.5922,0.0793) (0.5810,0.0777) (0.5698,0.0762) (0.5587,0.0747) (0.5475,0.0732) (0.5363,0.0718) (0.5251,0.0703) (0.5140,0.0689) (0.5028,0.0676) (0.4916,0.0662) (0.4804,0.0649) (0.4693,0.0636) (0.4581,0.0623) (0.4469,0.0611) (0.4358,0.0598) (0.4246,0.0586) (0.4134,0.0574) (0.4022,0.0563) (0.3911,0.0551) (0.3799,0.0540) (0.3687,0.0529) (0.3575,0.0518) (0.3464,0.0508) (0.3352,0.0497) (0.3240,0.0487) (0.3128,0.0477) (0.3017,0.0467) (0.2905,0.0458) (0.2793,0.0448) (0.2682,0.0439) (0.2570,0.0430) (0.2458,0.0421) (0.2346,0.0412) (0.2235,0.0404) (0.2123,0.0395) (0.2011,0.0387) (0.1899,0.0379) (0.1788,0.0371) (0.1676,0.0363) (0.1564,0.0356) (0.1453,0.0348) (0.1341,0.0341) (0.1229,0.0334) (0.1117,0.0327) (0.1006,0.0320) (0.0894,0.0313) (0.0782,0.0307) (0.0670,0.0300) (0.0559,0.0294) (0.0447,0.0288) (0.0335,0.0281) (0.0223,0.0275) (0.0112,0.0270) (0.0000,0.0264)};\draw[main,thick]plot coordinates{(2.0000,0.5000) (2.0168,0.5063) (2.0335,0.5125) (2.0503,0.5187) (2.0670,0.5249) (2.0838,0.5311) (2.1006,0.5372) (2.1173,0.5432) (2.1341,0.5493) (2.1508,0.5553) (2.1676,0.5612) (2.1844,0.5672) (2.2011,0.5730) (2.2179,0.5789) (2.2346,0.5846) (2.2514,0.5904) (2.2682,0.5961) (2.2849,0.6017) (2.3017,0.6073) (2.3184,0.6128) (2.3352,0.6183) (2.3520,0.6238) (2.3687,0.6292) (2.3855,0.6345) (2.4022,0.6398) (2.4190,0.6450) (2.4358,0.6502) (2.4525,0.6553) (2.4693,0.6604) (2.4860,0.6654) (2.5028,0.6703) (2.5196,0.6752) (2.5363,0.6801) (2.5531,0.6848) (2.5698,0.6896) (2.5866,0.6942) (2.6034,0.6989) (2.6201,0.7034) (2.6369,0.7079) (2.6536,0.7124) (2.6704,0.7168) (2.6872,0.7211) (2.7039,0.7254) (2.7207,0.7296) (2.7374,0.7338) (2.7542,0.7379) (2.7709,0.7420) (2.7877,0.7460) (2.8045,0.7499) (2.8212,0.7538) (2.8380,0.7577) (2.8547,0.7615) (2.8715,0.7652) (2.8883,0.7689) (2.9050,0.7726) (2.9218,0.7762) (2.9385,0.7797) (2.9553,0.7832) (2.9721,0.7866) (2.9888,0.7900) (3.0056,0.7933) (3.0223,0.7966) (3.0391,0.7999) (3.0559,0.8031) (3.0726,0.8062) (3.0894,0.8093) (3.1061,0.8124) (3.1229,0.8154) (3.1397,0.8184) (3.1564,0.8213) (3.1732,0.8242) (3.1899,0.8270) (3.2067,0.8298) (3.2235,0.8325) (3.2402,0.8352) (3.2570,0.8379) (3.2737,0.8405) (3.2905,0.8431) (3.3073,0.8457) (3.3240,0.8482) (3.3408,0.8506) (3.3575,0.8531) (3.3743,0.8555) (3.3911,0.8578) (3.4078,0.8601) (3.4246,0.8624) (3.4413,0.8647) (3.4581,0.8669) (3.4749,0.8690) (3.4916,0.8712) (3.5084,0.8733) (3.5251,0.8754) (3.5419,0.8774) (3.5587,0.8794) (3.5754,0.8814) (3.5922,0.8833) (3.6089,0.8852) (3.6257,0.8871) (3.6425,0.8890) (3.6592,0.8908) (3.6760,0.8926) (3.6927,0.8944) (3.7095,0.8961) (3.7263,0.8978) (3.7430,0.8995) (3.7598,0.9011) (3.7765,0.9028) (3.7933,0.9044) (3.8101,0.9059) (3.8268,0.9075) (3.8436,0.9090) (3.8603,0.9105) (3.8771,0.9120) (3.8939,0.9134) (3.9106,0.9149) (3.9274,0.9163) (3.9441,0.9177) (3.9609,0.9190) (3.9777,0.9204) (3.9944,0.9217) (4.0112,0.9230) (4.0279,0.9242) (4.0447,0.9255) (4.0615,0.9267) (4.0782,0.9279) (4.0950,0.9291) (4.1117,0.9303) (4.1285,0.9314) (4.1453,0.9326) (4.1620,0.9337) (4.1788,0.9348) (4.1955,0.9359) (4.2123,0.9369) (4.2291,0.9380) (4.2458,0.9390) (4.2626,0.9400) (4.2793,0.9410) (4.2961,0.9420) (4.3128,0.9429) (4.3296,0.9439) (4.3464,0.9448) (4.3631,0.9457) (4.3799,0.9466) (4.3966,0.9475) (4.4134,0.9484) (4.4302,0.9492) (4.4469,0.9501) (4.4637,0.9509) (4.4804,0.9517) (4.4972,0.9525) (4.5140,0.9533) (4.5307,0.9541) (4.5475,0.9549) (4.5642,0.9556) (4.5810,0.9563) (4.5978,0.9571) (4.6145,0.9578) (4.6313,0.9585) (4.6480,0.9592) (4.6648,0.9598) (4.6816,0.9605) (4.6983,0.9612) (4.7151,0.9618) (4.7318,0.9624) (4.7486,0.9631) (4.7654,0.9637) (4.7821,0.9643) (4.7989,0.9649) (4.8156,0.9655) (4.8324,0.9660) (4.8492,0.9666) (4.8659,0.9671) (4.8827,0.9677) (4.8994,0.9682) (4.9162,0.9688) (4.9330,0.9693) (4.9497,0.9698) (4.9665,0.9703) (4.9832,0.9708) (5.0000,0.9713)};\draw[main,thick]plot coordinates{(2.0000,1.5000) (1.9888,1.5042) (1.9777,1.5084) (1.9665,1.5126) (1.9553,1.5168) (1.9441,1.5211) (1.9330,1.5253) (1.9218,1.5296) (1.9106,1.5338) (1.8994,1.5381) (1.8883,1.5424) (1.8771,1.5467) (1.8659,1.5510) (1.8547,1.5552) (1.8436,1.5595) (1.8324,1.5638) (1.8212,1.5681) (1.8101,1.5724) (1.7989,1.5767) (1.7877,1.5810) (1.7765,1.5853) (1.7654,1.5896) (1.7542,1.5939) (1.7430,1.5982) (1.7318,1.6025) (1.7207,1.6068) (1.7095,1.6111) (1.6983,1.6153) (1.6872,1.6196) (1.6760,1.6239) (1.6648,1.6281) (1.6536,1.6324) (1.6425,1.6366) (1.6313,1.6408) (1.6201,1.6451) (1.6089,1.6493) (1.5978,1.6534) (1.5866,1.6576) (1.5754,1.6618) (1.5642,1.6659) (1.5531,1.6701) (1.5419,1.6742) (1.5307,1.6783) (1.5196,1.6824) (1.5084,1.6864) (1.4972,1.6905) (1.4860,1.6945) (1.4749,1.6985) (1.4637,1.7025) (1.4525,1.7065) (1.4413,1.7104) (1.4302,1.7143) (1.4190,1.7182) (1.4078,1.7221) (1.3966,1.7259) (1.3855,1.7298) (1.3743,1.7336) (1.3631,1.7373) (1.3520,1.7411) (1.3408,1.7448) (1.3296,1.7485) (1.3184,1.7522) (1.3073,1.7558) (1.2961,1.7594) (1.2849,1.7630) (1.2737,1.7665) (1.2626,1.7700) (1.2514,1.7735) (1.2402,1.7770) (1.2291,1.7804) (1.2179,1.7838) (1.2067,1.7872) (1.1955,1.7905) (1.1844,1.7938) (1.1732,1.7971) (1.1620,1.8003) (1.1508,1.8035) (1.1397,1.8067) (1.1285,1.8098) (1.1173,1.8129) (1.1061,1.8159) (1.0950,1.8190) (1.0838,1.8220) (1.0726,1.8249) (1.0615,1.8279) (1.0503,1.8308) (1.0391,1.8336) (1.0279,1.8364) (1.0168,1.8392) (1.0056,1.8420) (0.9944,1.8447) (0.9832,1.8474) (0.9721,1.8500) (0.9609,1.8527) (0.9497,1.8552) (0.9385,1.8578) (0.9274,1.8603) (0.9162,1.8628) (0.9050,1.8652) (0.8939,1.8676) (0.8827,1.8700) (0.8715,1.8724) (0.8603,1.8747) (0.8492,1.8770) (0.8380,1.8792) (0.8268,1.8814) (0.8156,1.8836) (0.8045,1.8858) (0.7933,1.8879) (0.7821,1.8900) (0.7709,1.8920) (0.7598,1.8940) (0.7486,1.8960) (0.7374,1.8980) (0.7263,1.8999) (0.7151,1.9018) (0.7039,1.9037) (0.6927,1.9055) (0.6816,1.9073) (0.6704,1.9091) (0.6592,1.9108) (0.6480,1.9126) (0.6369,1.9142) (0.6257,1.9159) (0.6145,1.9175) (0.6034,1.9192) (0.5922,1.9207) (0.5810,1.9223) (0.5698,1.9238) (0.5587,1.9253) (0.5475,1.9268) (0.5363,1.9282) (0.5251,1.9297) (0.5140,1.9311) (0.5028,1.9324) (0.4916,1.9338) (0.4804,1.9351) (0.4693,1.9364) (0.4581,1.9377) (0.4469,1.9389) (0.4358,1.9402) (0.4246,1.9414) (0.4134,1.9426) (0.4022,1.9437) (0.3911,1.9449) (0.3799,1.9460) (0.3687,1.9471) (0.3575,1.9482) (0.3464,1.9492) (0.3352,1.9503) (0.3240,1.9513) (0.3128,1.9523) (0.3017,1.9533) (0.2905,1.9542) (0.2793,1.9552) (0.2682,1.9561) (0.2570,1.9570) (0.2458,1.9579) (0.2346,1.9588) (0.2235,1.9596) (0.2123,1.9605) (0.2011,1.9613) (0.1899,1.9621) (0.1788,1.9629) (0.1676,1.9637) (0.1564,1.9644) (0.1453,1.9652) (0.1341,1.9659) (0.1229,1.9666) (0.1117,1.9673) (0.1006,1.9680) (0.0894,1.9687) (0.0782,1.9693) (0.0670,1.9700) (0.0559,1.9706) (0.0447,1.9712) (0.0335,1.9719) (0.0223,1.9725) (0.0112,1.9730) (0.0000,1.9736)};\draw[main,thick]plot coordinates{(2.0000,1.5000) (2.0168,1.4937) (2.0335,1.4875) (2.0503,1.4813) (2.0670,1.4751) (2.0838,1.4689) (2.1006,1.4628) (2.1173,1.4568) (2.1341,1.4507) (2.1508,1.4447) (2.1676,1.4388) (2.1844,1.4328) (2.2011,1.4270) (2.2179,1.4211) (2.2346,1.4154) (2.2514,1.4096) (2.2682,1.4039) (2.2849,1.3983) (2.3017,1.3927) (2.3184,1.3872) (2.3352,1.3817) (2.3520,1.3762) (2.3687,1.3708) (2.3855,1.3655) (2.4022,1.3602) (2.4190,1.3550) (2.4358,1.3498) (2.4525,1.3447) (2.4693,1.3396) (2.4860,1.3346) (2.5028,1.3297) (2.5196,1.3248) (2.5363,1.3199) (2.5531,1.3152) (2.5698,1.3104) (2.5866,1.3058) (2.6034,1.3011) (2.6201,1.2966) (2.6369,1.2921) (2.6536,1.2876) (2.6704,1.2832) (2.6872,1.2789) (2.7039,1.2746) (2.7207,1.2704) (2.7374,1.2662) (2.7542,1.2621) (2.7709,1.2580) (2.7877,1.2540) (2.8045,1.2501) (2.8212,1.2462) (2.8380,1.2423) (2.8547,1.2385) (2.8715,1.2348) (2.8883,1.2311) (2.9050,1.2274) (2.9218,1.2238) (2.9385,1.2203) (2.9553,1.2168) (2.9721,1.2134) (2.9888,1.2100) (3.0056,1.2067) (3.0223,1.2034) (3.0391,1.2001) (3.0559,1.1969) (3.0726,1.1938) (3.0894,1.1907) (3.1061,1.1876) (3.1229,1.1846) (3.1397,1.1816) (3.1564,1.1787) (3.1732,1.1758) (3.1899,1.1730) (3.2067,1.1702) (3.2235,1.1675) (3.2402,1.1648) (3.2570,1.1621) (3.2737,1.1595) (3.2905,1.1569) (3.3073,1.1543) (3.3240,1.1518) (3.3408,1.1494) (3.3575,1.1469) (3.3743,1.1445) (3.3911,1.1422) (3.4078,1.1399) (3.4246,1.1376) (3.4413,1.1353) (3.4581,1.1331) (3.4749,1.1310) (3.4916,1.1288) (3.5084,1.1267) (3.5251,1.1246) (3.5419,1.1226) (3.5587,1.1206) (3.5754,1.1186) (3.5922,1.1167) (3.6089,1.1148) (3.6257,1.1129) (3.6425,1.1110) (3.6592,1.1092) (3.6760,1.1074) (3.6927,1.1056) (3.7095,1.1039) (3.7263,1.1022) (3.7430,1.1005) (3.7598,1.0989) (3.7765,1.0972) (3.7933,1.0956) (3.8101,1.0941) (3.8268,1.0925) (3.8436,1.0910) (3.8603,1.0895) (3.8771,1.0880) (3.8939,1.0866) (3.9106,1.0851) (3.9274,1.0837) (3.9441,1.0823) (3.9609,1.0810) (3.9777,1.0796) (3.9944,1.0783) (4.0112,1.0770) (4.0279,1.0758) (4.0447,1.0745) (4.0615,1.0733) (4.0782,1.0721) (4.0950,1.0709) (4.1117,1.0697) (4.1285,1.0686) (4.1453,1.0674) (4.1620,1.0663) (4.1788,1.0652) (4.1955,1.0641) (4.2123,1.0631) (4.2291,1.0620) (4.2458,1.0610) (4.2626,1.0600) (4.2793,1.0590) (4.2961,1.0580) (4.3128,1.0571) (4.3296,1.0561) (4.3464,1.0552) (4.3631,1.0543) (4.3799,1.0534) (4.3966,1.0525) (4.4134,1.0516) (4.4302,1.0508) (4.4469,1.0499) (4.4637,1.0491) (4.4804,1.0483) (4.4972,1.0475) (4.5140,1.0467) (4.5307,1.0459) (4.5475,1.0451) (4.5642,1.0444) (4.5810,1.0437) (4.5978,1.0429) (4.6145,1.0422) (4.6313,1.0415) (4.6480,1.0408) (4.6648,1.0402) (4.6816,1.0395) (4.6983,1.0388) (4.7151,1.0382) (4.7318,1.0376) (4.7486,1.0369) (4.7654,1.0363) (4.7821,1.0357) (4.7989,1.0351) (4.8156,1.0345) (4.8324,1.0340) (4.8492,1.0334) (4.8659,1.0329) (4.8827,1.0323) (4.8994,1.0318) (4.9162,1.0312) (4.9330,1.0307) (4.9497,1.0302) (4.9665,1.0297) (4.9832,1.0292) (5.0000,1.0287)};\draw[main,thick]plot coordinates{(2.0000,2.5000) (1.9888,2.4797) (1.9777,2.4606) (1.9665,2.4427) (1.9553,2.4257) (1.9441,2.4097) (1.9330,2.3946) (1.9218,2.3803) (1.9106,2.3667) (1.8994,2.3538) (1.8883,2.3415) (1.8771,2.3298) (1.8659,2.3186) (1.8547,2.3080) (1.8436,2.2978) (1.8324,2.2881) (1.8212,2.2789) (1.8101,2.2700) (1.7989,2.2615) (1.7877,2.2533) (1.7765,2.2455) (1.7654,2.2379) (1.7542,2.2307) (1.7430,2.2238) (1.7318,2.2171) (1.7207,2.2107) (1.7095,2.2045) (1.6983,2.1986) (1.6872,2.1928) (1.6760,2.1873) (1.6648,2.1819) (1.6536,2.1768) (1.6425,2.1718) (1.6313,2.1670) (1.6201,2.1624) (1.6089,2.1579) (1.5978,2.1536) (1.5866,2.1494) (1.5754,2.1453) (1.5642,2.1414) (1.5531,2.1376) (1.5419,2.1339) (1.5307,2.1303) (1.5196,2.1269) (1.5084,2.1235) (1.4972,2.1203) (1.4860,2.1172) (1.4749,2.1141) (1.4637,2.1112) (1.4525,2.1083) (1.4413,2.1055) (1.4302,2.1028) (1.4190,2.1002) (1.4078,2.0976) (1.3966,2.0951) (1.3855,2.0927) (1.3743,2.0904) (1.3631,2.0881) (1.3520,2.0859) (1.3408,2.0838) (1.3296,2.0817) (1.3184,2.0797) (1.3073,2.0777) (1.2961,2.0758) (1.2849,2.0739) (1.2737,2.0721) (1.2626,2.0703) (1.2514,2.0686) (1.2402,2.0670) (1.2291,2.0653) (1.2179,2.0637) (1.2067,2.0622) (1.1955,2.0607) (1.1844,2.0592) (1.1732,2.0578) (1.1620,2.0564) (1.1508,2.0551) (1.1397,2.0537) (1.1285,2.0525) (1.1173,2.0512) (1.1061,2.0500) (1.0950,2.0488) (1.0838,2.0477) (1.0726,2.0465) (1.0615,2.0454) (1.0503,2.0444) (1.0391,2.0433) (1.0279,2.0423) (1.0168,2.0413) (1.0056,2.0403) (0.9944,2.0394) (0.9832,2.0385) (0.9721,2.0376) (0.9609,2.0367) (0.9497,2.0358) (0.9385,2.0350) (0.9274,2.0342) (0.9162,2.0334) (0.9050,2.0326) (0.8939,2.0319) (0.8827,2.0311) (0.8715,2.0304) (0.8603,2.0297) (0.8492,2.0290) (0.8380,2.0284) (0.8268,2.0277) (0.8156,2.0271) (0.8045,2.0264) (0.7933,2.0258) (0.7821,2.0252) (0.7709,2.0247) (0.7598,2.0241) (0.7486,2.0235) (0.7374,2.0230) (0.7263,2.0225) (0.7151,2.0220) (0.7039,2.0215) (0.6927,2.0210) (0.6816,2.0205) (0.6704,2.0200) (0.6592,2.0196) (0.6480,2.0191) (0.6369,2.0187) (0.6257,2.0183) (0.6145,2.0179) (0.6034,2.0175) (0.5922,2.0171) (0.5810,2.0167) (0.5698,2.0163) (0.5587,2.0159) (0.5475,2.0156) (0.5363,2.0152) (0.5251,2.0149) (0.5140,2.0145) (0.5028,2.0142) (0.4916,2.0139) (0.4804,2.0136) (0.4693,2.0133) (0.4581,2.0130) (0.4469,2.0127) (0.4358,2.0124) (0.4246,2.0121) (0.4134,2.0118) (0.4022,2.0116) (0.3911,2.0113) (0.3799,2.0111) (0.3687,2.0108) (0.3575,2.0106) (0.3464,2.0103) (0.3352,2.0101) (0.3240,2.0099) (0.3128,2.0096) (0.3017,2.0094) (0.2905,2.0092) (0.2793,2.0090) (0.2682,2.0088) (0.2570,2.0086) (0.2458,2.0084) (0.2346,2.0082) (0.2235,2.0081) (0.2123,2.0079) (0.2011,2.0077) (0.1899,2.0075) (0.1788,2.0074) (0.1676,2.0072) (0.1564,2.0070) (0.1453,2.0069) (0.1341,2.0067) (0.1229,2.0066) (0.1117,2.0064) (0.1006,2.0063) (0.0894,2.0061) (0.0782,2.0060) (0.0670,2.0059) (0.0559,2.0057) (0.0447,2.0056) (0.0335,2.0055) (0.0223,2.0054) (0.0112,2.0052) (0.0000,2.0051)};\draw[main,thick]plot coordinates{(2.0000,2.5000) (2.0016,2.5031) (2.0033,2.5062) (2.0049,2.5094) (2.0066,2.5126) (2.0082,2.5158) (2.0099,2.5190) (2.0115,2.5223) (2.0131,2.5256) (2.0148,2.5289) (2.0164,2.5323) (2.0181,2.5357) (2.0197,2.5392) (2.0213,2.5426) (2.0230,2.5462) (2.0246,2.5497) (2.0263,2.5533) (2.0279,2.5569) (2.0296,2.5606) (2.0312,2.5643) (2.0328,2.5680) (2.0345,2.5718) (2.0361,2.5756) (2.0378,2.5794) (2.0394,2.5833) (2.0410,2.5873) (2.0427,2.5913) (2.0443,2.5953) (2.0460,2.5994) (2.0476,2.6035) (2.0493,2.6076) (2.0509,2.6118) (2.0525,2.6161) (2.0542,2.6204) (2.0558,2.6247) (2.0575,2.6291) (2.0591,2.6336) (2.0607,2.6381) (2.0624,2.6426) (2.0640,2.6472) (2.0657,2.6519) (2.0673,2.6566) (2.0690,2.6614) (2.0706,2.6662) (2.0722,2.6711) (2.0739,2.6761) (2.0755,2.6811) (2.0772,2.6861) (2.0788,2.6913) (2.0805,2.6965) (2.0821,2.7017) (2.0837,2.7071) (2.0854,2.7125) (2.0870,2.7179) (2.0887,2.7235) (2.0903,2.7291) (2.0919,2.7348) (2.0936,2.7405) (2.0952,2.7464) (2.0969,2.7523) (2.0985,2.7583) (2.1002,2.7644) (2.1018,2.7705) (2.1034,2.7768) (2.1051,2.7831) (2.1067,2.7895) (2.1084,2.7961) (2.1100,2.8027) (2.1116,2.8094) (2.1133,2.8162) (2.1149,2.8231) (2.1166,2.8301) (2.1182,2.8372) (2.1199,2.8444) (2.1215,2.8518) (2.1231,2.8592) (2.1248,2.8668) (2.1264,2.8744) (2.1281,2.8822) (2.1297,2.8901) (2.1313,2.8982) (2.1330,2.9064) (2.1346,2.9147) (2.1363,2.9231) (2.1379,2.9317) (2.1396,2.9405) (2.1412,2.9493) (2.1428,2.9584) (2.1445,2.9676) (2.1461,2.9769) (2.1478,2.9865) (2.1494,2.9962)};\end{tikzpicture}\end{center}
}

\newcommand{\ODEExamRampDerivative}{%
\begin{center}\begin{tikzpicture}[x=1.1cm,y=.65cm]\begin{scope}[xshift=0cm]\draw[->](-1.6,0)--(2.3,0)node[right]{$t$};\draw[->](0,-1.5)--(0,1.7)node[above]{$f$};\draw[main,thick](-1.5,0)--(-1,0);\draw[main,thick](-1,-1)--(1,1);\draw[main,thick](1,0)--(2,0);\node[below]at(-1,0){$ -1 $};\node[below]at(1,0){$ 1 $};\end{scope}\begin{scope}[xshift=5.2cm]\draw[->](-1.6,0)--(2.3,0)node[right]{$t$};\draw[->](0,-1.5)--(0,1.7)node[above]{$f'$};\draw[main,thick](-1.5,0)--(-1,0);\draw[main,thick](-1,1)--(1,1);\draw[main,thick](1,0)--(2,0);\draw[structurecolor,very thick,->](-1,0)--(-1,-1);\node[above]at(-1,-1){$ -1 $};\draw[structurecolor,very thick,->](1,0)--(1,-1);\node[above]at(1,-1){$ -1 $};\node[below]at(-1,0){$ -1 $};\node[below]at(1,0){$ 1 $};\end{scope}\end{tikzpicture}\end{center}
}

\newcommand{\ODEExamVDerivative}{%
\begin{center}\begin{tikzpicture}[x=1.1cm,y=.65cm]\begin{scope}[xshift=0cm]\draw[->](-1.6,0)--(2.3,0)node[right]{$t$};\draw[->](0,-1.5)--(0,1.7)node[above]{$f'$};\draw[main,thick](-1.5,0)--(-1,0);\draw[main,thick](-1,-1)--(0,-1);\draw[main,thick](0,1)--(1,1);\draw[main,thick](1,0)--(2,0);\draw[structurecolor,very thick,->](-1,0)--(-1,1);\node[above]at(-1,1){$ 1 $};\draw[structurecolor,very thick,->](1,0)--(1,-1);\node[above]at(1,-1){$ -1 $};\node[below]at(-1,0){$ -1 $};\node[below]at(1,0){$ 1 $};\end{scope}\end{tikzpicture}\end{center}
}

\newcommand{\ODEExamPracticeRamp}{%
\begin{center}\begin{tikzpicture}[x=1.1cm,y=.65cm]\begin{scope}[xshift=0cm]\draw[->](-1.6,0)--(2.3,0)node[right]{$t$};\draw[->](0,-1.5)--(0,1.7)node[above]{$f$};\draw[main,thick](-1.5,0)--(0,0);\draw[main,thick](0,1)--(1,0);\draw[main,thick](1,1)--(2,1);\node[below]at(-1,0){$ -1 $};\node[below]at(1,0){$ 1 $};\end{scope}\begin{scope}[xshift=5.2cm]\draw[->](-1.6,0)--(2.3,0)node[right]{$t$};\draw[->](0,-1.5)--(0,1.7)node[above]{$f'$};\draw[main,thick](-1.5,0)--(0,0);\draw[main,thick](0,-1)--(1,-1);\draw[main,thick](1,0)--(2,0);\draw[structurecolor,very thick,->](0,0)--(0,1);\node[above]at(0,1){$ 1 $};\draw[structurecolor,very thick,->](1,0)--(1,1);\node[above]at(1,1){$ 1 $};\node[below]at(-1,0){$ -1 $};\node[below]at(1,0){$ 1 $};\end{scope}\end{tikzpicture}\end{center}
}

\newcommand{\ODEExamMotionDerivatives}{%
\begin{center}\begin{tikzpicture}[x=.8cm,y=.6cm]\draw[->](-.5,0)--(4.8,0)node[right]{$t$};\node at(2,2.6){速度};\draw[main,thick](-.4,0)--(0,0);\draw[main,thick](0,1)--(1,1);\draw[main,thick](1,-1)--(3,-1);\draw[main,thick](3,1)--(4,1);\draw[main,thick](4,0)--(4.6,0);\begin{scope}[xshift=5.4cm]\draw[->](-.5,0)--(4.8,0)node[right]{$t$};\node at(2,2.6){加速度};\foreach\x/\a in{0/1,1/-2,3/2,4/-1}{\draw[->,structurecolor,very thick](\x,0)--(\x,\a);\node[right]at(\x,\a){$\a$};\node[above]at(\x,0){$\x$};}\end{scope}\end{tikzpicture}\end{center}
}

\newcommand{\ODEExamNodePortrait}{%
\begin{center}\begin{tikzpicture}\begin{axis}[width=7cm,height=6cm,axis lines=middle,xmin=-3,xmax=3,ymin=-3,ymax=3,xlabel={$x$},ylabel={$y$},samples=70]\addplot[gray,domain=-3:3]{x};\addplot[gray,domain=-3:3]{-x};\addplot[main,domain=0:-1.5,->]({x-(-1)*x^2},{x+(-1)*x^2});\addplot[main,domain=0:1.5,->]({x-(-1)*x^2},{x+(-1)*x^2});\addplot[main,domain=0:-1.5,->]({x-(-0.4)*x^2},{x+(-0.4)*x^2});\addplot[main,domain=0:1.5,->]({x-(-0.4)*x^2},{x+(-0.4)*x^2});\addplot[main,domain=0:-1.5,->]({x-(0.4)*x^2},{x+(0.4)*x^2});\addplot[main,domain=0:1.5,->]({x-(0.4)*x^2},{x+(0.4)*x^2});\addplot[main,domain=0:-1.5,->]({x-(1)*x^2},{x+(1)*x^2});\addplot[main,domain=0:1.5,->]({x-(1)*x^2},{x+(1)*x^2});\end{axis}\end{tikzpicture}\end{center}
}

\newcommand{\ODEExamPoleDiagram}[2]{\begin{center}\begin{tikzpicture}[scale=.55]\draw[->](-3,0)--(1.3,0)node[right]{$\operatorname{Re}s$};\draw[->](0,-3.8)--(0,3.8)node[above]{$\operatorname{Im}s$};\node[main]at(0,0){$\times$};\node[below right]at(0,0){$0$};\node[main]at(-1,#1){$\times$};\node[main]at(-1,-#1){$\times$};\node[left]at(-1,#1){$-1+#1\ii$};\node[left]at(-1,-#1){$-1-#1\ii$};\end{tikzpicture}\end{center}}

% MIT 18.03 Spring 2010 eight examinations; Chinese translation and checked solutions.
\originalchapter{期中考试 I（2010 年 2 月 24 日）}
八卷解答的“译审”指核对原题与相应答案后中文重述并复算；“编者补解（AI 辅助）”是本册新增解答，非 MIT 官方答案。独立第二读者的复核状态见书外审校记录，不把未完成复核宣称为已完成。
本卷五题，满分 100 分；原卷不允许计算器和笔记。题号和小问沿用原卷。解答依据为当期官方答案的中文译审，并补足计算；原答案中的错误另行指出。
\begin{exercise}\label{cw:ex1-1}\textbf{原题 1.} 理想环境中，MIT 校园挪威鼠的数量每年增长为原来的 $e\approx2.718281828459045$ 倍。
(a) 用微分方程描述自然增长，求增长率 $k$。
(b) 校园可持续容纳的上限为 $R=1000$ 只，写出相应 Logistic 方程；若未求出 $k$，可以仍用 $k$ 表示。
(c) 防治部门每年以恒定速率 $a$ 杀鼠；若希望数量限制在可持续上限的 $75\%$，应取多少只/年？
\end{exercise}
\begin{solution}
\textit{依据：当期官方答案译审并补证（AI 辅助）；涉及原稿错误时另按题面重算。}\par
(a) 设 $x(t)$ 为鼠数，$t$ 以年计，则 $\dot x=kx$，$x(1)/x(0)=\ee^k=\ee$，故 $k=1\ \text{年}^{-1}$。
(b) $\dot x=x(1-x/1000)$。
(c) 加入捕杀后 $\dot x=x(1-x/1000)-a$。在 $x=750$ 平衡要求 $a=750(1-750/1000)=187.5$ 只/年。此平衡附近的导数为 $1-2(750)/1000=-1/2<0$，因此是稳定平衡。官方写出的 $3R/16$ 正确，但其数值 $375$ 错误；$R=1000$ 时应为 $187.5$。恒定捕杀还产生较低的不稳定平衡 $250$，因此该模型不能保证所有初始鼠数都趋于 $750$。
\end{solution}
\begin{exercise}\label{cw:ex1-2}\textbf{原题 2.}
对于 $\dot x=x(x-1)(x+2)$，(a) 画相线，标出临界点及其稳定、不稳定或半稳定类型；(b) 每一基本解型至少画一个代表，使每个解都是所画解之一的时间平移。
另考虑方向场 $y'=\tfrac14(x-y^2)$：(c) 标出零斜线；(d) 画至少两个很不同的解；(e) 若有分界解，画出它；(f) 判断：若一个解有极小值，则充分大的 $x$ 下总有 $y(x)<\sqrt x$（原题无需解释）。
\end{exercise}
\begin{solution}
\textit{依据：当期官方答案译审并补证（AI 辅助）；涉及原稿错误时另按题面重算。}\par
(a) 依次在 $(-\infty,-2),(-2,0),(0,1),(1,\infty)$ 上，右端符号为 $-,+,-,+$。因此 $-2,1$ 不稳定，$0$ 稳定。
(b) 基本解共有七型：三个常值解和四个区间中的非恒定解。内侧两型随时间趋于 $0$，外侧两型分别向 $-\infty,+\infty$ 逃逸；外侧解在有限时间爆破，因为 $\int^\infty\dd x/[x(x-1)(x+2)]$ 收敛。图中只截取爆破前的解段。
\begin{center}\begin{tikzpicture}[x=.72cm,y=.65cm]
\draw[->] (-.8,-3)--(-.8,2) node[above]{$x$};
\foreach \v in {-2,0,1}{\fill (-.8,\v) circle(1.5pt);\node[left] at(-.8,\v){$\v$};}
\draw[->](-.8,-2.4)--(-.8,-2.9);\draw[->](-.8,-1.7)--(-.8,-.5);\draw[->](-.8,.8)--(-.8,.2);\draw[->](-.8,1.3)--(-.8,1.9);
\draw[->](1,-3)--(6,-3)node[right]{$t$};\draw[->](1,-3)--(1,2)node[above]{$x$};
\foreach \v in {-2,0,1}{\draw[dashed](1,\v)--(5.7,\v);}
\draw[main] (1,-1.9)..controls(2,-1.9)and(2.6,-.05)..(5.6,-.02);
\draw[main] (1,.95)..controls(2,.95)and(2.6,.05)..(5.6,.02);
\draw[main] (1,1.08)..controls(4,1.12)and(4.6,1.4)..(4.9,2);
\draw[main] (1,-2.08)..controls(3.4,-2.1)and(4,-2.5)..(4.3,-3);
\end{tikzpicture}\end{center}
(c) 零斜线为 $x=y^2$，即向右开的抛物线。抛物线内侧 $x>y^2$ 时向上，外侧向下。
(d)--(e) 下图给出方向场及代表曲线；实线初值为 $y(0)=0$ 和 $y(0)=-2$，虚线为分界解。它在 $x\to+\infty$ 时接近下支 $-\sqrt x$ 的下侧，把趋向上支的解和向下爆破的解分开。为精确定义该线，可取满足 $A''(z)=zA(z)$、在正无穷衰减的正解 $A$，令 $c=16^{-1/3}$，则 $y_*(x)=4c A'(cx)/A(cx)$（$A$ 无零点的区间内）。直接求导得到 $y_*'=(x-y_*^2)/4$；这只是分界线的精确表述，原题作图不要求特殊函数。
\ODEExamRiccatiField
(f) 真。极值处 $y'=0$，且 $y''=(1-2yy')/4=1/4>0$；越过极小值后，曲线处于零斜线内侧。若试图从下方穿越上支 $\sqrt x$，交点处解的斜率为零，而上支斜率 $1/(2\sqrt x)>0$，不可能发生向上的穿越。因此它保持在上支下方。
\end{solution}
\begin{exercise}\label{cw:ex1-3}\textbf{原题 3.}
(a) 对 $y'=x+y$、$y(0)=1$，用步长 $h=1/2$ 的 Euler 法估计 $y(3/2)$。
(b) 求 $t\dot x+x=\cos t$、$x(\pi)=1$ 的解。
\end{exercise}
\begin{solution}
\textit{依据：当期官方答案译审并补证（AI 辅助）；涉及原稿错误时另按题面重算。}\par
(a) 递推 $y_{j+1}=y_j+\tfrac12(x_j+y_j)$，得到
\[ (x_j,y_j)=(0,1),(1/2,3/2),(1,5/2),(3/2,17/4). \]
故估计为 $17/4$。
(b) 左端为 $(tx)'$，积分得 $tx=\sin t+C$；初值给 $C=\pi$，故 $x=(\pi+\sin t)/t$，初值所在最大普通解区间为 $t>0$。
\end{solution}
\begin{exercise}\label{cw:ex1-4}\textbf{原题 4.}
(a) 求实数 $a,b$ 使 $1/(3+2\ii)=a+b\ii$。
(b) 求实数 $r,\theta$ 使 $1-\ii=r\ee^{\ii\theta}$。
(c) 求 $(1-\ii)^8=a+b\ii$ 中的实数 $a,b$。
(d) 求 $-1$ 的一个立方根 $a+b\ii$，要求 $b>0$。
(e) 求 $\ee^{\ln2+\ii\pi}=a+b\ii$ 中的实数 $a,b$。
(f) 将 $2\cos4t-2\sin4t$ 写成 $A\cos(\omega t-\phi)$。
\end{exercise}
\begin{solution}
\textit{依据：当期官方答案译审并补证（AI 辅助）；涉及原稿错误时另按题面重算。}\par
(a) 乘以共轭得 $(3-2\ii)/13$，故 $a=3/13,b=-2/13$。
(b) 模长 $\sqrt2$、辐角 $-\pi/4$，可取 $r=\sqrt2,\theta=-\pi/4$。
(c) $(\sqrt2)^8\ee^{-2\pi\ii}=16$，故 $(a,b)=(16,0)$。
(d) 三个辐角为 $\pi/3,\pi,5\pi/3$，正虚部的根是 $1/2+\ii\sqrt3/2$。
(e) $2\ee^{\ii\pi}=-2$，故 $(a,b)=(-2,0)$。
(f) 比较 $A\cos\phi=2,A\sin\phi=-2$，取 $A=2\sqrt2,\omega=4,\phi=-\pi/4$（相位可加 $2\pi$ 的整数倍）。
\end{solution}
\begin{exercise}\label{cw:ex1-5}\textbf{原题 5.}
(a) 求 $\dot x+3x=\ee^{2t}$ 的一个特解。
(b) 求该方程满足 $x(0)=1$ 的解。
(c) 写出一个指数右端的复线性方程，使 $\dot x+3x=\cos2t$ 是它的实部。
(d) 求 $\dot x+3x=\cos2t$ 的一个特解。
\end{exercise}
\begin{solution}
\textit{依据：当期官方答案译审并补证（AI 辅助）；涉及原稿错误时另按题面重算。}\par
(a) 令 $x=A\ee^{2t}$，得到 $5A=1$，故 $x_p=\ee^{2t}/5$。
(b) 通解为 $\ee^{2t}/5+C\ee^{-3t}$，初值给 $C=4/5$。
(c) $\dot z+3z=\ee^{2\ii t}$。
(d) $z_p=\ee^{2\ii t}/(3+2\ii)$，取实部得 $x_p=(3\cos2t+2\sin2t)/13$。
\end{solution}

\originalchapter{期中考试 II（2010 年 3 月 17 日）}
本卷五题，满分 100 分，不允许计算器和笔记。解答依据为当期官方答案的中文译审与计算补全。
\begin{exercise}\label{cw:ex2-1}\textbf{原题 1.}
(a) $\ddot x+\dot x+kx=0$ 在何值 $k$ 下临界阻尼？(b) $k$ 大于此值时为过阻尼还是欠阻尼？(c) 一个解在 $t=1,2$ 为零，中间不为零，求 $k$。
\end{exercise}
\begin{solution}
\textit{依据：当期官方答案译审并补证（AI 辅助）；涉及原稿错误时另按题面重算。}\par
(a) 特征方程为 $(s+1/2)^2+k-1/4=0$，重根时 $k=1/4$。
(b) 此时两根非实，故欠阻尼。
(c) 非零解两个相邻零点相距 $\pi/\omega_d=1$，故 $\omega_d=\pi$、$k=\pi^2+1/4$。实根或重根解不可能具有两个不同零点而中间不为零：除去一个正指数因子，剩余为单调指数的仿射式或关于 $t$ 的一次式。
\end{solution}
\begin{exercise}\label{cw:ex2-2}\textbf{原题 2.}
(a) 求 $\ddot x+x=5t\ee^{2t}$ 的一个解。
(b) 已知同方程的解 $y$ 满足 $y(0)=1,\dot y(0)=2$，用 $y$ 及其他函数写出满足 $x(0)=3,\dot x(0)=5$ 的解。
\end{exercise}
\begin{solution}
\textit{依据：当期官方答案译审并补证（AI 辅助）；涉及原稿错误时另按题面重算。}\par
(a) 令 $x=\ee^{2t}u$，方程化为 $u''+4u'+5u=5t$。取 $u=at+b$，比较系数得 $a=1,b=-4/5$，故 $x=(t-4/5)\ee^{2t}$。
(b) 两解之差满足 $h''+h=0$，初值为 $h(0)=2,h'(0)=3$，故 $x=y+2\cos t+3\sin t$。
\end{solution}
\begin{exercise}\label{cw:ex2-3}\textbf{原题 3.}
(a) 在 $\ddot x+b\dot x+kx=\cos\omega t$ 中固定 $b$、变化弹簧常数 $k$。哪个 $k$ 使正弦特解的振幅最大？答案可依赖 $\omega,b$。
(b) 另一问题：求 $x'''-x'=0$ 的通解。
\end{exercise}
\begin{solution}
\textit{依据：当期官方答案译审并补证（AI 辅助）；涉及原稿错误时另按题面重算。}\par
(a) 振幅为 $[(k-\omega^2)^2+b^2\omega^2]^{-1/2}$，故分母在 $k=\omega^2$ 最小。若 $b\omega\ne0$，最大振幅为 $1/|b\omega|$；若 $b\omega=0$ 且分母为零，则共振点没有有限振幅的正弦特解，须区分此边界情形。
(b) 特征多项式 $s(s-1)(s+1)$ 的根为 $0,1,-1$，故 $x=C_0+C_1\ee^t+C_2\ee^{-t}$。
\end{solution}
\begin{exercise}\label{cw:ex2-4}\textbf{原题 4.} 输入 $y$ 与响应 $x$ 满足 $\ddot x+\dot x+6x=6y$，输入为正弦信号。
(a) 求复增益 $H(\omega)$；(b) 求 $\omega=2$ 的增益；(c) 求该频率的相位滞后。
\end{exercise}
\begin{solution}
\textit{依据：当期官方答案译审并补证（AI 辅助）；涉及原稿错误时另按题面重算。}\par
(a) 输入取 $\ee^{\ii\omega t}$，响应取 $H\ee^{\ii\omega t}$，得到 $H=6/(6-\omega^2+\ii\omega)$。
(b) $H(2)=3/(1+\ii)=\tfrac32(1-\ii)$，故增益 $|H(2)|=3/\sqrt2$。
(c) 相位滞后为 $-\arg H(2)=\pi/4$。
\end{solution}
\begin{exercise}\label{cw:ex2-5}\textbf{原题 5.} 已知 $\tfrac12t\sin2t$ 为 $m\ddot x+b\dot x+kx=4\cos2t$ 的解。
(a) 写出右端改为 $4\cos(2t-1)$ 的一个解；(b) 写出右端改为 $8\cos2t$ 的一个解；(c) 求 $m,b,k$。
\end{exercise}
\begin{solution}
\textit{依据：当期官方答案译审并补证（AI 辅助）；涉及原稿错误时另按题面重算。}\par
(a) 时间平移 $1/2$ 得 $\tfrac12(t-1/2)\sin(2t-1)$。
(b) 线性缩放得 $t\sin2t$。
(c) 已知解的导数为 $x'=\tfrac12\sin2t+t\cos2t$、$x''=2\cos2t-2t\sin2t$。代入得到
\[2m\cos2t+bt\cos2t+\tfrac b2\sin2t+(k/2-2m)t\sin2t=4\cos2t.\]
比较独立函数的系数得 $b=0,m=2,k=8$。
\end{solution}

\originalchapter{期中考试 III（2010 年 4 月 23 日）}
本卷五题，满分 100 分，不允许计算器和笔记；原卷附有 Laplace 和 Fourier 公式页。本册统一列在考试附用公式节。解答依据为当期官方答案译审与必要推导。原附表将 $t\sin(\omega t)$ 与 $t\cos(\omega t)$ 的变换互换；本册按 $\mathcal L(tf)=-F'(s)$ 核算，分别为 $2\omega s/(s^2+\omega^2)^2$ 与 $(s^2-\omega^2)/(s^2+\omega^2)^2$，并明确补正这一笔误。
\begin{exercise}\label{cw:ex3-1}\textbf{原题 1.} 周期函数具有级数
\[f(t)=1+\sum_{j=1}^\infty 2^{-j}\cos(j\pi t).\]
(a) 求最小周期；(b) 判断为偶、奇、两者皆非或两者皆是；(c) 若存在，给出 $\ddot x+\omega_n^2x=f(t)$ 的一个周期解的 Fourier 级数；(d) 哪些 $\omega_n$ 下没有周期解？
\end{exercise}
\begin{solution}
\textit{依据：当期官方答案译审并补证（AI 辅助）；涉及原稿错误时另按题面重算。}\par
(a) 最小周期为 $2$；首谐波系数非零，故周期不可能小于首谐波的周期。
(b) 偶函数，且不是奇函数，因为平均值为 $1$。
(c) 在不共振时，逐个谐波响应给
\[x_p=\frac1{\omega_n^2}+\sum_{j=1}^\infty\frac{2^{-j}\cos(j\pi t)}{\omega_n^2-j^2\pi^2}.\]
其系数指数衰减，所以二次逐项微分也一致收敛，可代入验证。
(d) 自然频率按非负数约定时为 $0,\pi,2\pi,\ldots$。$\omega_n=0$ 时常数驱动产生二次增长；$\omega_n=j\pi$ 时第 $j$ 个非零谐波产生 $t\sin(j\pi t)$ 项，无法由齐次解抵消。若允许负 $\omega_n$，应同时包括负整数倍 $\pi$。
\end{solution}
\begin{exercise}\label{cw:ex3-2}\textbf{原题 2.} $f(t)=[u(t+1)-u(t-1)]t$。
(a) 画 $f$；(b) 画广义导数 $f'$；(c) 写出广义导数，分别指出常规部分 $f'_r$ 与奇异部分 $f'_s$。
\end{exercise}
\begin{solution}
\textit{依据：当期官方答案译审并补证（AI 辅助）；涉及原稿错误时另按题面重算。}\par
(a) $(-1,1)$ 上为 $t$，其余为零，在 $-1,1$ 处均有大小 $-1$ 的跳跃。
(b)--(c) 区间内部斜率为 $1$，两个端点各有一个权重 $-1$ 的向下冲激：
\[f'_r=u(t+1)-u(t-1),\qquad f'_s=-\delta(t+1)-\delta(t-1).\]
单点取值不影响这个广义导数。
\ODEExamRampDerivative
\end{solution}
\begin{exercise}\label{cw:ex3-3}\textbf{原题 3.} 算子 $p(D)$ 的单位冲激响应为 $w(t)=\ee^{-t}-\ee^{-3t}$（因果延拓）。
(a) 用卷积求静止初值下 $p(D)v=u(t)$ 的单位阶跃响应；(b) 求传递函数 $W(s)$；(c) 求特征多项式 $p(s)$。
\end{exercise}
\begin{solution}
\textit{依据：当期官方答案译审并补证（AI 辅助）；涉及原稿错误时另按题面重算。}\par
(a) $t\ge0$ 时 $v(t)=\int_0^t[\ee^{-\tau}-\ee^{-3\tau}]\dd\tau=2/3-\ee^{-t}+\ee^{-3t}/3$，负时间为零。
(b) $W(s)=1/(s+1)-1/(s+3)=2/[(s+1)(s+3)]$。
(c) $p(s)=1/W(s)=\tfrac12(s^2+4s+3)$。系数 $1/2$ 不能漏掉；$w'(0+)=2$ 与最高阶系数的倒数相符。
\end{solution}
\begin{exercise}\label{cw:ex3-4}\textbf{原题 4.}
(a) 求 Laplace 变换为 $\ee^{-s}(s-1)/s$ 的广义函数；(b) 求 Laplace 变换为 $(s+10)/(s^3+2s^2+10s)$ 的函数。
\end{exercise}
\begin{solution}
\textit{依据：当期官方答案译审并补证（AI 辅助）；涉及原稿错误时另按题面重算。}\par
(a) $(s-1)/s=1-1/s$ 对应 $\delta(t)-u(t)$，时间平移后为 $\delta(t-1)-u(t-1)$。
(b) 分母为 $s[(s+1)^2+9]$，且
\[\frac{s+10}{s[(s+1)^2+9]}=\frac1s-\frac{s+1}{(s+1)^2+9}.\]
因此 $f(t)=u(t)[1-\ee^{-t}\cos3t]$。
\end{solution}
\begin{exercise}\label{cw:ex3-5}\textbf{原题 5.} $W(s)=(s+10)/(s^3+2s^2+10s)$。
(a) 画极点图；(b) 若 $W$ 是算子 $p(D)$ 的传递函数，求静止初值下 $p(D)x=\sin2t$ 的解的 Laplace 变换。
\end{exercise}
\begin{solution}
\textit{依据：当期官方答案译审并补证（AI 辅助）；涉及原稿错误时另按题面重算。}\par
(a) 极点为 $0,-1\pm3\ii$，均无分子消去，见图。
\ODEExamPoleDiagram{3}{0}
(b) 静止初值保证 $X=WF$，故
\[X(s)=\frac{2(s+10)}{s(s^2+2s+10)(s^2+4)}.\]
这里的 $W$ 并非一个常系数多项式微分算子的倒数：$1/W=s(s^2+2s+10)/(s+10)$ 为有理函数。原题把 $p(D)$ 作为一般线性系统记号；这一点不影响题设指定的输入输出乘法关系。
\end{solution}

\originalchapter{期末考试（2010 年 5 月 18 日）}
原卷考试时间为 9:00--12:00，十题，共 360 分；禁止计算器和笔记，可使用卷末公式页。若题目指定方法，必须使用该方法。本卷未提供官方答案，以下全部为独立补解，不能与题头为 2008 年的练习期末答案混配。
\begin{exercise}\label{cw:final-1}\textbf{原题 1.}
(a)--(b) 新泽西州 Oyster Creek 核电站的氚以恒定 $1$ 千克/年漏入含水层，氚半衰期为 $12$ 年。
(a) 忽略其他来源和移除机制，建立含水层氚量的微分方程并确定所有常数；(b) 若持续泄漏很久，氚量是多少千克？
(c)--(g) 考虑 $y'=x-y^2/4$，初值 $y(1)=0$。
(c) 用步长 $1/2$ 的 Euler 法估计 $y(2)$；(d) 在同一平面画斜率 $m=-1,0,1$ 的等斜线；(e) 画该初值解；(f) 假设此解在 $x=a$ 处达到极小值，用 $a$ 表示 $y(a)$；(g) 估计 $y(100)$。
\end{exercise}
\begin{solution}
\textit{依据：编者补解（AI 辅助），无 MIT 官方答案可供核对。}\par
(a) 设 $M(t)$ 以千克计，衰减率满足 $\ee^{-12k}=1/2$，故 $k=\ln2/12\ \text{年}^{-1}$，方程为 $M'=1-(\ln2/12)M$。
(b) 通解 $M=12/\ln2+[M(0)-12/\ln2]\ee^{-(\ln2)t/12}$，长期量为 $12/\ln2\approx17.31$ 千克。
(c) 从 $(1,0)$ 出发，第一步得 $(3/2,1/2)$；第二步斜率为 $3/2-(1/2)^2/4=23/16$，故 $y(2)\approx1/2+(1/2)(23/16)=39/32$。
(d) 等斜线为 $x=y^2/4+m$，是顶点分别在 $(-1,0),(0,0),(1,0)$ 的右开抛物线。
(e) 初值处斜率为 $1$，向右解上升并渐近接近 $2\sqrt x$ 下侧；向左先下降，在负支零斜线上出现极小值。图中曲线由所给方程数值取点重绘，等斜线为精确曲线。
\ODEExamFinalField
(f) 极值条件给 $y(a)=\pm2\sqrt a$。又 $y''=1-(y/2)y'$，极值处为 $1$，均是严格极小值。该初值解在极值之后上升至 $y(1)=0$，因而最小值负，故 $y(a)=-2\sqrt a$。它不能在 $1$ 右侧首次向下穿过正支零斜线，因为该支斜率为正而解在那里斜率为零。
(g) 零斜线上支给估计 $y(100)\approx2\sqrt{100}=20$。进一步令 $y=2\sqrt x+\eta$，由
\[x^{-1/2}+\eta'=-\sqrt x\,\eta-\eta^2/4\]
主阶平衡得 $\eta\sim-1/x$，所以更精细的近似是 $19.99$。这是大 $x$ 渐近估计，不能视为 Euler 两步的继续或精确值。
\end{solution}
\begin{exercise}\label{cw:final-2}\textbf{原题 2.}
(a)--(c) 考虑 $\dot x=x^3-x^2-2x$。
(a) 画相线；(b) 每个临界点上、下及两临界点之间均画至少一个代表解；(c) 若 $x(0)=0.1$ 很小，则 $t>0$ 时 $x(t)$ 最宜用 $0.1\ee^{at}$ 近似，求 $a$。
(d) 解 $x y'+2y=-\sin x/x$、$y(\pi)=0$。
\end{exercise}
\begin{solution}
\textit{依据：编者补解（AI 辅助），无 MIT 官方答案可供核对。}\par
(a) 因式分解为 $x(x-2)(x+1)$，临界点为 $-1,0,2$。四个区间中的符号为 $-,+,-,+$，故 $-1,2$ 不稳定，$0$ 稳定。
(b) 包括三条常值线、$(-1,0)$ 内向上趋 $0$ 的解、$(0,2)$ 内向下趋 $0$ 的解，以及两条外侧向无穷逃逸的解。\ODEExamFinalPhase
(c) $f'(0)=-2$，所以线性近似为 $x(t)\approx0.1\ee^{-2t}$，$a=-2$；近似忽略了二次及三次项。
(d) 先除以 $x$，积分因子为 $x^2$，因此 $(x^2y)'=-\sin x$。得 $x^2y=\cos x+C$，初值给 $C=1$，故 $y=(1+\cos x)/x^2$，初值所在区间为 $x>0$。
\end{solution}
\begin{exercise}\label{cw:final-3}\textbf{原题 3.}
(a) 求复数 $r=a+b\ii$ 和正数 $\omega$，使 $\operatorname{Re}(\ee^{\ii\omega t}/r)=2\cos(2t-\pi/4)$。
(b) 将 $-8\ii$ 的所有立方根写成 $a+b\ii$。
(c) 求 $\ddot x+2\dot x+10x=0$ 的非零解的伪周期。
(d) $\ddot x+2\dot x+10x=\cos\omega t$ 的正弦特解在何角频率 $\omega_r$ 下振幅最大？
(e) 哪个角频率 $\omega_p$ 下相位滞后为 $\pi/2$？
\end{exercise}
\begin{solution}
\textit{依据：编者补解（AI 辅助），无 MIT 官方答案可供核对。}\par
(a) 频率必须为 $2$，且 $1/r=2\ee^{-\ii\pi/4}$，故 $r=\tfrac12\ee^{\ii\pi/4}=(1+\ii)/(2\sqrt2)$。
(b) 模长为 $2$，辐角可取 $-\pi/6,\pi/2,7\pi/6$，故根为 $\sqrt3-\ii,2\ii,-\sqrt3-\ii$。
(c) 特征根 $-1\pm3\ii$，解为 $\ee^{-t}(C\cos3t+D\sin3t)$，伪周期为 $2\pi/3$。
(d) 振幅分母的平方为 $(10-\omega^2)^2+4\omega^2$。令 $z=\omega^2\ge0$，此式为 $(z-8)^2+36$，最小处 $\omega_r=\sqrt8=2\sqrt2$。
(e) 相位滞后为 $\arg(10-\omega^2+2\ii\omega)$；正频率下其为 $\pi/2$ 当且仅当实部为零，故 $\omega_p=\sqrt{10}$。
\end{solution}
\begin{exercise}\label{cw:final-4}\textbf{原题 4.} $p(D)=D^2+4D+8I$。
(a) 求 $p(D)x=8t^2$ 的一个解；(b) 求 $p(D)x=\ee^{-t}$ 的一个解。
(c)--(d) 已知 $t^3$ 是 $p(D)x=q(t)$ 的解：(c) 求 $q$；(d) 求满足 $x(0)=0,\dot x(0)=2$ 的解。
\end{exercise}
\begin{solution}
\textit{依据：编者补解（AI 辅助），无 MIT 官方答案可供核对。}\par
(a) 取 $x=at^2+bt+c$，左端为 $8at^2+(8a+8b)t+2a+4b+8c$，故 $a=1,b=-1,c=1/4$，$x=t^2-t+1/4$。
(b) $p(-1)=1-4+8=5$，故 $x=\ee^{-t}/5$。
(c) 直接求导，$q=6t+12t^2+8t^3$。
(d) 齐次特征根 $-2\pm2\ii$，所以 $x=t^3+\ee^{-2t}(C\cos2t+D\sin2t)$。初值要求 $C=0,2D=2$，故 $x=t^3+\ee^{-2t}\sin2t$。
\end{solution}
\begin{exercise}\label{cw:final-5}\textbf{原题 5.}
(a) 若存在，求 $\ddot x+\omega_n^2x=\operatorname{sq}(2t)$ 的一个周期解。
(b)--(c) 设 $f$ 为 $2\pi$ 周期函数，$|t|\le\pi$ 时 $f(t)=|t|-\pi/2$。
(b) 判断奇偶性，求平均值并画图；(c) 求 Fourier 级数。
\end{exercise}
\begin{solution}
\textit{依据：编者补解（AI 辅助），无 MIT 官方答案可供核对。}\par
(a) 方波展开为 $\tfrac4\pi\sum_{j\text{ 正奇数}}\sin(2jt)/j$，故不共振时
\[x_p=\frac4\pi\sum_{j\text{ 正奇数}}\frac{\sin(2jt)}{j(\omega_n^2-4j^2)}.\]
$\omega_n=2,6,10,\ldots$ 时存在非零共振谐波，因而无周期解。$\omega_n=0$ 时平均驱动为零，上式仍给周期解；这与期中 III 第 1 题含非零平均值的情形不同。若允许负自然频率，排除相应负值。级数二次微分恢复方波的 Fourier 级数，跳点按分段方程解释。
(b) $f$ 是偶函数，$\int_{-\pi}^{\pi}f=0$，平均值为零。一个周期是连接 $(-\pi,\pi/2),(0,-\pi/2),(\pi,\pi/2)$ 的折线，作周期延拓。
\begin{center}\begin{tikzpicture}[x=.7cm,y=.65cm]\draw[->](-6.7,0)--(6.7,0)node[right]{$t$};\draw[->](0,-2)--(0,2)node[above]{$f$};\draw[main](-6.28,-1.57)--(-3.14,1.57)--(0,-1.57)--(3.14,1.57)--(6.28,-1.57);\foreach\x/\l in {-3.14/-\pi,3.14/\pi}{\node[below]at(\x,0){$\l$};}\node[left]at(0,1.57){$\pi/2$};\node[left]at(0,-1.57){$-\pi/2$};\end{tikzpicture}\end{center}
(c) $b_j=0,a_0=0$。分部积分给
\[a_j=\frac2\pi\int_0^\pi(t-\pi/2)\cos jt\dd t=\frac{2[(-1)^j-1]}{\pi j^2}.\]
故 $f(t)=-\tfrac4\pi\sum_{j\text{ 正奇数}}\cos jt/j^2$。函数连续，系数绝对可和，级数在每一点给出函数本身。
\end{solution}
\begin{exercise}\label{cw:final-6}\textbf{原题 6.}
(a)--(b) 函数 $f$ 由下图给定，刻度间隔均为 $1$；端点的单点值不指定。
(a) 画广义导数；(b) 用阶跃和冲激写出 $f'$ 并分出常规与奇异部分。
\begin{center}\begin{tikzpicture}[scale=1.05]\draw[->](-2.5,0)--(2.5,0)node[right]{$t$};\draw[->](0,-.4)--(0,1.5)node[above]{$f$};\foreach\x in {-2,-1,1,2}{\draw(\x,-.05)--(\x,.05);\node[below]at(\x,0){$\x$};}\draw(-.05,1)--(.05,1)node[right]{$1$};\draw[very thick](-2.4,0)--(-1,0);\draw[main,thick](-1,1)--(0,0)--(1,1);\draw[very thick](1,0)--(2.4,0);\end{tikzpicture}\end{center}
(c)--(e) 某算子单位冲激响应为 $w(t)=u(t)\ee^{-t}\sin2t$。
(c) 求静止初值下 $p(D)x=u(t)\ee^{-t}$ 的解；(d) 求静止初值下 $p(D)x=\delta(t-1)$ 的解；(e) 求特征多项式 $p(s)$。
\end{exercise}
\begin{solution}
\textit{依据：编者补解（AI 辅助），无 MIT 官方答案可供核对。}\par
(a)--(b) 图示限定 $f=|t|$ 于 $-1<t<1$，其余为零。在 $-1$ 跳升 $1$，在 $1$ 跳降 $1$；原点只有拐角，没有跳跃，所以一阶导数在原点没有冲激。
\[f'_r=-u(t+1)+2u(t)-u(t-1),\qquad f'_s=\delta(t+1)-\delta(t-1).\]
导数图是 $(-1,0)$ 高度 $-1$、$(0,1)$ 高度 $1$、外侧为零，加两端的冲激。\ODEExamVDerivative
(c) 卷积为
\[x(t)=u(t)\int_0^t\ee^{-(t-\tau)}\sin[2(t-\tau)]\ee^{-\tau}\dd\tau
=\tfrac12u(t)\ee^{-t}(1-\cos2t).\]
(d) 冲激平移响应，$x=u(t-1)\ee^{-(t-1)}\sin[2(t-1)]$。
(e) $W(s)=2/[(s+1)^2+4]$，故 $p(s)=\tfrac12(s^2+2s+5)$。注意最高阶系数是 $1/2$，与 $w'(0+)=2$ 相符。
\end{solution}
\begin{exercise}\label{cw:final-7}\textbf{原题 7.}
(a)--(c) 考虑 $p(D)=D^3+D$：(a) 求传递函数；(b) 求单位冲激响应；(c) 求静止初值下 $p(D)x=\cos2t$ 的解的 Laplace 变换。
(d) 画 $(s+1)\ee^{-s}/(s^3+2s^2+2s)$ 的极点图。
\end{exercise}
\begin{solution}
\textit{依据：编者补解（AI 辅助），无 MIT 官方答案可供核对。}\par
(a) $W=1/[s(s^2+1)]$。
(b) 拆分为 $1/s-s/(s^2+1)$，故 $w=u(t)(1-\cos t)$。其 $w(0+)=w'(0+)=0,w''(0+)=1$，符合三阶算子的冲激条件。
(c) $X=W\,s/(s^2+4)=1/[(s^2+1)(s^2+4)]$。
(d) 分母为 $s[(s+1)^2+1]$，极点为 $0,-1\pm\ii$；分子在这些点不为零，指数因子无有限零点，故均保留。\ODEExamPoleDiagram{1}{0}
\end{solution}
\begin{exercise}\label{cw:final-8}\textbf{原题 8.}
(a)--(b) $A=\begin{pmatrix}5&-1\\-1&5\end{pmatrix}$：(a) 求特征值；(b) 每个特征值求一个非零特征向量。
(c)--(e) $B$ 的特征值为 $2,-2$：(c) 求迹和行列式；(d) 若相应特征向量为 $(1,2)^T,(3,4)^T$，求 $\ee^{Bt}$；(e) 求 $\dot u=Bu,u(0)=(1,1)^T$ 的解。
\end{exercise}
\begin{solution}
\textit{依据：编者补解（AI 辅助），无 MIT 官方答案可供核对。}\par
(a)--(b) 特征多项式为 $(5-\lambda)^2-1$，特征值 $4,6$，可取特征向量分别为 $(1,1)^T,(1,-1)^T$。
(c) $\tr B=2+(-2)=0,\det B=-4$。
(d) 令 $P=\begin{pmatrix}1&3\\2&4\end{pmatrix}$，则 $P^{-1}=\begin{pmatrix}-2&3/2\\1&-1/2\end{pmatrix}$。因此，记 $E=\ee^{2t},F=\ee^{-2t}$，有
\[\ee^{Bt}=P\begin{pmatrix}E&0\\0&F\end{pmatrix}P^{-1}
=\begin{pmatrix}-2E+3F&\tfrac32(E-F)\\-4E+4F&3E-2F\end{pmatrix}.\]
在 $t=0$ 得单位矩阵，导数满足 $B\ee^{Bt}$。
(e) $P^{-1}(1,1)^T=(-1/2,1/2)^T$，故
\[u(t)=\tfrac12\begin{pmatrix}-\ee^{2t}+3\ee^{-2t}\\-2\ee^{2t}+4\ee^{-2t}\end{pmatrix}.\]
\end{solution}
\begin{exercise}\label{cw:final-9}\textbf{原题 9.} $A=\begin{pmatrix}1&2\\-2&a\end{pmatrix}$，$\dot u=Au$。
(a) 在迹--行列式平面画重特征值临界抛物线及参数 $a$ 对应的直线。
(b) 分别求下列类型对应的全部 $a$：
(i) 鞍点；(ii) 稳定结点；(iii) 不稳定结点；(iv) 稳定螺旋；(v) 不稳定螺旋；(vi) 稳定亏损结点；(vii) 不稳定亏损结点；(viii) 星形结点；(ix) 中心；(x) 退化情形。
\end{exercise}
\begin{solution}
\textit{依据：编者补解（AI 辅助），无 MIT 官方答案可供核对。}\par
(a) $T=a+1,\Delta=a+4$，参数直线为 $\Delta=T+3$；临界抛物线为 $\Delta=T^2/4$。交点为 $(T,\Delta)=(-2,1),(6,9)$。
\begin{center}\begin{tikzpicture}[x=.48cm,y=.24cm]
\draw[->](-5,0)--(8,0)node[right]{$T$};\draw[->](0,-2)--(0,13)node[above]{$\Delta$};
\foreach\t in{-4,-2,2,4,6,8}{\draw(\t,.15)--(\t,-.15)node[below]{$\t$};}
\foreach\d in{5,10}{\draw(.1,\d)--(-.1,\d)node[left]{$\d$};}
\draw[main,domain=-5:7,samples=80]plot(\x,{\x*\x/4});
\draw[structurecolor](-5,-2)--(8,11);\fill(-2,1)circle(2pt);\fill(6,9)circle(2pt);
\end{tikzpicture}\end{center}
(b) 判别式 $T^2-4\Delta=(a-5)(a+3)$。按原编号，答案为
\[
\begin{array}{c|l}
\text{(i)}&a<-4\\
\text{(ii)}&-4<a<-3\\
\text{(iii)}&a>5\\
\text{(iv)}&-3<a<-1\\
\text{(v)}&-1<a<5\\
\text{(vi)}&a=-3\\
\text{(vii)}&a=5\\
\text{(viii)}&\text{无}\\
\text{(ix)}&a=-1\\
\text{(x)}&a=-4
\end{array}
\]
其理由分别是：$\Delta<0$ 给鞍点；$\Delta>0$、判别式正给异根结点，迹的符号决定稳定性；判别式负给复根螺旋，迹为零则为中心。重根在 $a=-3,5$，矩阵均不是标量矩阵，故只有一个独立特征向量而亏损，不是星形结点。在 $a=-4$ 时一个根为零，平衡点不是孤立的；这就是本题的退化情形。表中普通结点排除了另列的亏损结点。
\end{solution}
\begin{exercise}\label{cw:final-10}\textbf{原题 10.}
\[\dot x=(y+1)(y-x+1),\qquad \dot y=(x-3)(x+y-1).\]
(a) 求并在平面上标出全部平衡点；(b) 求平衡点 $(1,0)$ 的 Jacobian；(c) 此点是稳定螺旋，趋于它的解在大 $t$ 时，其 $y$ 坐标可用 $A\ee^{at}\cos(\omega t-\phi)$ 近似；求所有解共有的常数，区别于依赖具体解的常数。
\end{exercise}
\begin{solution}
\textit{依据：编者补解（AI 辅助），无 MIT 官方答案可供核对。}\par
(a) $\dot x=0$ 给 $y=-1$ 或 $y=x-1$；$\dot y=0$ 给 $x=3$ 或 $y=1-x$。两族直线的四个交点为 $(3,-1),(2,-1),(3,2),(1,0)$。
\begin{center}\begin{tikzpicture}[scale=.7]\draw[->](-.5,0)--(4,0)node[right]{$x$};\draw[->](0,-1.6)--(0,2.7)node[above]{$y$};\foreach\x/\y in{3/-1,2/-1,3/2,1/0}{\fill[main](\x,\y)circle(2pt);}\node[below left]at(2,-1){$(2,-1)$};\node[below right]at(3,-1){$(3,-1)$};\node[above right]at(3,2){$(3,2)$};\node[above right]at(1,0){$(1,0)$};\end{tikzpicture}\end{center}
(b) $J(x,y)=\begin{pmatrix}-(y+1)&2y-x+2\\2x+y-4&x-3\end{pmatrix}$，故 $J(1,0)=\begin{pmatrix}-1&1\\-2&-2\end{pmatrix}$。
(c) 特征方程 $\lambda^2+3\lambda+4=0$ 给 $\lambda=(-3\pm\ii\sqrt7)/2$。因此共有的指数率 $a=-3/2$ 和角频率 $\omega=\sqrt7/2$；$A,\phi$ 依赖具体解。这里是平衡点附近的线性化主项；非线性系统的坐标并不在所有时刻都等于这个正弦表达式。
\end{solution}

\originalchapter{练习期中考试 I（2010）}
本文件的第 1--2 页为试题，第 2--3 页兼有官方解答。以下解答是官方译审与补证，发现的原稿错误在相应小问说明。
\begin{exercise}\label{cw:prex1-1}\textbf{原题 1.} 计算机芯片以正比于其温度与环境温度之差的速率散热。
(a) 环境恒为 $20^\circ\mathrm C$，时间以分钟计，建立芯片温度的微分方程，其中有一个尚不能确定的常数；(b) 求通解；(c) 芯片在 $t=0$ 断电时为 $70^\circ\mathrm C$，$t=10$ 分钟时为 $60^\circ\mathrm C$，据此确定方程。
\end{exercise}
\begin{solution}
\textit{依据：同文件官方答案译审并补证（AI 辅助）；官方遗漏或错误另标。}\par
(a) 设 $T(t)$ 为温度，$k>0$ 为散热系数，$T'=-k(T-20)$，即 $T'+kT=20k$。官方解答漏写了右端的 $k$，这里依物理模型补正。
(b) 分离变量或积分因子给 $T=20+C\ee^{-kt}$。
(c) $C=50$，$40=50\ee^{-10k}$，所以 $k=\ln(5/4)/10\ \text{分钟}^{-1}$，方程为 $T'=-(\ln(5/4)/10)(T-20)$。
\end{solution}
\begin{exercise}\label{cw:prex1-2}\textbf{原题 2.} 对 $y'=y^2-x^2,y(2)=0$，用步长 $0.1$ 的 Euler 法估计 $y(2.2)$。
\end{exercise}
\begin{solution}
\textit{依据：同文件官方答案译审并补证（AI 辅助）；官方遗漏或错误另标。}\par
第一步斜率为 $-4$，$y(2.1)\approx-0.4$；第二步斜率 $0.16-4.41=-4.25$，故 $y(2.2)\approx-0.4-0.425=-0.825$。
\end{solution}
\begin{exercise}\label{cw:prex1-3}\textbf{原题 3.} 考虑 $y'=y^2-x^2$，原方向场绘于 $-6\le x,y\le6$。
(a) 画并标出 $m=-4,0,4$ 的等斜线；(b) 画初值为 $y(2)=0$ 的解；(c) 估计此解的 $y(100)$，判断估计偏大还是偏小；(d) 某解在 $x=-1$ 有局部极值，求可能的 $y(-1)$，并计算说明是极大还是极小，不能只依赖图。
\end{exercise}
\begin{solution}
\textit{依据：同文件官方答案译审并补证（AI 辅助）；官方遗漏或错误另标。}\par
(a) 等斜线为 $y^2-x^2=m$，零斜线为 $y=\pm x$，$m=4$ 是上下双曲线，$m=-4$ 是左右双曲线。官方解答把零斜线误印为 $y=\pm1$，并改用了 $\pm2$ 的等斜线；这里均按题面 $\pm4$ 重绘。
(b) 初值处斜率 $-4$，曲线向右下降，后来在 $y=-x$ 上方逐渐贴近该直线。\ODEExamPracticeField
(c) 估计 $y(100)\approx-100$，偏小。证明曲线在 $y=-x$ 上方：令 $z=y+x$，在 $z=0$ 处 $z'=y'+1=1>0$，故不能向下穿越，且 $z(2)=2>0$。令 $y=-x+\eta$，主阶平衡 $-1\approx-2x\eta$ 给 $\eta\sim1/(2x)$，较精细估计为 $-99.995$。
(d) 极值条件 $y'=0$ 给 $y(-1)=\pm1$。又 $y''=2yy'-2x$，在该点为 $2>0$，两种初值均为严格极小值。
\end{solution}
\begin{exercise}\label{cw:prex1-4}\textbf{原题 4.}
(a) 求 $t\dot x+2x=t^2$ 的通解；(b) 求 $\dot x+2x=\cos2t$ 的正弦解，写成正弦与余弦之和，可用任何方法。
\end{exercise}
\begin{solution}
\textit{依据：同文件官方答案译审并补证（AI 辅助）；官方遗漏或错误另标。}\par
(a) 乘以 $t$，得 $(t^2x)'=t^3$，故 $x=t^2/4+C/t^2$，分别定义在 $t>0$ 或 $t<0$ 的区间。
(b) 复解为 $z=\ee^{2\ii t}/(2+2\ii)$，实部为 $x=(\cos2t+\sin2t)/4$。代入可见余弦系数为 $1$、正弦系数为 $0$。
\end{solution}
\begin{exercise}\label{cw:prex1-5}\textbf{原题 5.}
(a) 将 $-8\ii$ 的每个立方根先写成 $A\ee^{\ii\theta}$，再写成 $a+b\ii$。
(b)--(e) 考虑 $f(t)=-\cos(\pi t/2)-\sin(\pi t/2)$。
(b) 求正数 $A,\phi$ 使 $f=A\cos(\pi t/2-\phi)$；(c) 求周期 $P$；(d) 求时间滞后 $t_0$；(e) 画图并标出 $A,P,t_0$。
\end{exercise}
\begin{solution}
\textit{依据：同文件官方答案译审并补证（AI 辅助）；官方遗漏或错误另标。}\par
(a) 根为 $2\ee^{\ii\pi/2}=2\ii$、$2\ee^{7\ii\pi/6}=-\sqrt3-\ii$、$2\ee^{11\ii\pi/6}=\sqrt3-\ii$。
(b) 比较系数 $A\cos\phi=A\sin\phi=-1$，取 $A=\sqrt2,\phi=5\pi/4$。
(c) $P=2\pi/(\pi/2)=4$。
(d) $t_0=\phi/(\pi/2)=5/2$，允许按周期同余的滞后。
(e) 图在 $t=1/2$ 为谷，在 $t=5/2$ 为峰。
\begin{center}\begin{tikzpicture}\begin{axis}[width=10cm,height=4.8cm,axis lines=middle,xmin=-.2,xmax=7,ymin=-1.9,ymax=1.9,xlabel={$t$},ylabel={$f$},samples=120,xtick={.5,2.5,6.5},xticklabels={$1/2$,$t_0=5/2$,$13/2$},ytick={-1.414,1.414},yticklabels={$-\sqrt2$,$A=\sqrt2$}]\addplot[main,domain=0:7]{-cos(deg(pi*x/2))-sin(deg(pi*x/2))};\draw[<->](axis cs:2.5,1.65)--(axis cs:6.5,1.65)node[midway,below]{$P=4$};\end{axis}\end{tikzpicture}\end{center}
\end{solution}
\begin{exercise}\label{cw:prex1-6}\textbf{原题 6.} $\dot y=y^3-y$。
(a) 画相线，标出全部临界点并判定稳定性；(b) 画足够多的解，使每个 $-2<b<2$ 都被某条所画解经过；(c) 求解的拐点位置。原题将拐点定义为 $f(a)=b,f''(a)=0$ 的点。
\end{exercise}
\begin{solution}
\textit{依据：同文件官方答案译审并补证（AI 辅助）；官方遗漏或错误另标。}\par
(a) 临界点为 $-1,0,1$，四个区间上的右端符号为 $-,+,-,+$，故 $\pm1$ 不稳定，$0$ 稳定。
(b) 三条常值线加四个开区间中的一条非恒定解即可。\ODEExamCubicPhase
(c) $y''=(3y^2-1)(y^3-y)$，故非恒定解的拐点高度为 $y=\pm1/\sqrt3$，且在相应内侧区间确实改变凹凸性。按原题仅要求二阶导数为零的定义，还须列出三个常值解上的每一点，高度为 $0,\pm1$；按通常要求凹凸性改变的定义，这些常值解不算拐点。此处保留原定义并说明两种用法的差别。
\end{solution}

\originalchapter{练习期中考试 II 与复习指南（2010 年 3 月）}
\originalsection{复习指南的内容与适用条件}
原指南的九个主题依次为线性模型、齐次方程、线性叠加、指数响应、复数替换、待定系数、变参数、时间不变性及频率响应。其公式中的排印错误在此按模型补正。
线性方程 $\sum_{j=0}^n a_j(t)x^{(j)}=q(t)$ 的左侧代表系统，右侧由输入决定。本单元讨论常系数系统，记 $p(D)x=q$。弹簧系统的 $p(s)=ms^2+bs+k$：直接外力为 $q=F_{\rm ext}$，弹簧端点位移 $y$ 驱动时为 $q=ky$，阻尼器端点位移驱动时为 $q=b\dot y$。原指南最后一个模型把阻尼系数误印为质量 $m$。
齐次指数模式为 $\ee^{rt}$，其中 $p(r)=0$；重数 $\ell$ 的根需要 $t^j\ee^{rt}$，$0\le j<\ell$。实系数的共轭复根 $a\pm\ii\beta$ 对应实模式 $\ee^{at}\cos\beta t,\ee^{at}\sin\beta t$。$m>0,b,k\ge0$ 时按 $b^2-4mk$ 正、零、负区分过阻尼、临界阻尼、欠阻尼。欠阻尼的角频率为 $\sqrt{k/m-b^2/(4m^2)}$，指数包络为 $\ee^{-bt/(2m)}$，并允许振幅和相位常数。
线性性给 $p(D)(c_1x_1+c_2x_2)=c_1q_1+c_2q_2$，所以非齐次通解为一个特解加齐次通解。$p(r)\ne0$ 时指数特解为 $A\ee^{rt}/p(r)$；$r$ 为单根时改用 $At\ee^{rt}/p'(r)$。实系数允许取复解的实部或虚部求正弦响应。
若 $p(0)\ne0$，多项式驱动有唯一多项式特解，次数不超过驱动次数；若最低非零项为 $a_jD^j$，先令 $u=x^{(j)}$ 降阶，再积分得到特解。变参数替换 $x=\ee^{rt}u$ 利用 $p(D)(\ee^{rt}u)=\ee^{rt}p(D+r)u$。常系数还给时间平移不变性：$x(t-a)$ 响应于 $q(t-a)$。
输入 $y=\ee^{\ii\omega t}$ 时响应 $H(\omega)\ee^{\ii\omega t}$ 定义复增益，增益为 $|H|$，相位滞后为 $-\arg H$。输入振幅为 $A$ 时输出振幅是 $A|H|$，不能漏掉 $A$。共振使分母为零时不存在有限的纯正弦特解。
\originalsection{练习试卷与解答}
\begin{exercise}\label{cw:prex2-1}\textbf{原题 1.} 质量 $m=1$、弹簧常数 $k=25$，阻尼常数 $b$ 未知。某解满足 $\ddot x+b\dot x+25x=0$、$x(\pi/6)=x(\pi/2)=0$，且两零点之间 $x>0$。
(a) 判定阻尼类型；(b) 求 $b$。
\end{exercise}
\begin{solution}
\textit{依据：同文件官方答案译审并补证（AI 辅助）；官方遗漏或错误另标。}\par
(a) 两个不同零点迫使根为复根，因此欠阻尼。
(b) 相邻零点相距 $\pi/3$，$\omega_d=3$，所以 $25-b^2/4=9$。物理阻尼常数非负，故 $b=8$；若取消阻尼器的非负约定，方程数据也容许 $b=-8$。
\end{solution}
\begin{exercise}\label{cw:prex2-2}\textbf{原题 2.} 求 $3\ddot x+2\dot x+x=t^2$ 的一个解。
\end{exercise}
\begin{solution}
\textit{依据：同文件官方答案译审并补证（AI 辅助）；官方遗漏或错误另标。}\par
设 $x=at^2+bt+c$，左端为 $at^2+(4a+b)t+6a+2b+c$。比较系数得 $a=1,b=-4,c=2$，故 $x=t^2-4t+2$。
\end{solution}
\begin{exercise}\label{cw:prex2-3}\textbf{原题 3.} 求 $\ddot x+3\dot x+2x=\ee^{-t}$ 的一个解。
\end{exercise}
\begin{solution}
\textit{依据：同文件官方答案译审并补证（AI 辅助）；官方遗漏或错误另标。}\par
$p(s)=(s+1)(s+2)$，$p(-1)=0,p'(-1)=1$，故单根共振响应为 $x=t\ee^{-t}$。直接求导代入也得到 $\ee^{-t}$。
\end{solution}
\begin{exercise}\label{cw:prex2-4}\textbf{原题 4.} 考虑 $\ddot x+4\dot x+9x=\cos\omega t$ 的正弦解。
(a) 哪个 $\omega$ 下振幅最大？(b) 哪个 $\omega$ 下相位滞后恰为 $\pi/4$？
\end{exercise}
\begin{solution}
\textit{依据：同文件官方答案译审并补证（AI 辅助）；官方遗漏或错误另标。}\par
(a) 分母平方 $(9-\omega^2)^2+16\omega^2=(\omega^2-1)^2+80$，故非负频率取 $\omega=1$。
(b) $\arg(9-\omega^2+4\ii\omega)=\pi/4$ 要求实部、虚部相等且正，所以 $9-\omega^2=4\omega>0$，得 $\omega=\sqrt{13}-2$。
\end{solution}
\begin{exercise}\label{cw:prex2-5}\textbf{原题 5.} 系统输入 $y$、响应 $x$ 满足 $2\ddot x+\dot x+x=\dot y$，以角频率 $\omega$ 的正弦信号驱动。求复增益函数，以及 $\omega=1$ 时的增益和相位滞后。
\end{exercise}
\begin{solution}
\textit{依据：同文件官方答案译审并补证（AI 辅助）；官方遗漏或错误另标。}\par
复输入的导数引入因子 $\ii\omega$，所以 $H(\omega)=\ii\omega/(1-2\omega^2+\ii\omega)$。$H(1)=\ii/(-1+\ii)=(1-\ii)/2$，增益为 $1/\sqrt2$、相位滞后为 $\pi/4$。
\end{solution}
\begin{exercise}\label{cw:prex2-6}\textbf{原题 6.} 求 $x'''+x=\ee^{-t}\cos t$ 的一个解。
\end{exercise}
\begin{solution}
\textit{依据：同文件官方答案译审并补证（AI 辅助）；官方遗漏或错误另标。}\par
取复右端 $\ee^{(-1+\ii)t}$，$p(-1+\ii)=(-1+\ii)^3+1=3+2\ii$，故 $z=\ee^{(-1+\ii)t}/(3+2\ii)$。实部为 $x=\ee^{-t}(3\cos t+2\sin t)/13$。
\end{solution}
\begin{exercise}\label{cw:prex2-7}\textbf{原题 7.} 已知 $\cos t$ 与 $t$ 都是某常系数方程 $p(D)x=q(t)$ 的解。
(a) 写一个非零齐次解；(b) 写一个非齐次解使 $x(0)=2$；(c) 写 $p(D)x=q(t-1)$ 的一个解。
\end{exercise}
\begin{solution}
\textit{依据：同文件官方答案译审并补证（AI 辅助）；官方遗漏或错误另标。}\par
(a) 相减给 $h=\cos t-t$，它不恒为零且 $p(D)h=0$。
(b) $x=t+2h=2\cos t-t$，系数之和仍为 $1$，故右端仍为 $q$，且初值为 $2$。
(c) 时间平移给 $x=t-1$，也可取 $\cos(t-1)$；常系数条件保证 $p(D)x(t-1)=q(t-1)$。
\end{solution}

\originalchapter{练习期中考试 III 与复习指南（2010 年 4 月）}
\originalsection{阶跃、冲激、卷积与 Laplace 的复习指南}
原文件第 1--2 页的学习主题为阶跃与冲激的光滑近似、分段光滑信号及广义导数、静止初值、冲激响应、阶跃响应、卷积和串联系统、Laplace 规则、部分分式、极点与频率响应。理想阶跃 $u$ 在负时间为零、正时间为一，$u'=\delta$。可以用短时间尺度的光滑过渡帮助理解，但 $u,\delta$ 的精确运算使用广义函数，而不是把一个有限宽度近似直接当作等式。
对分段光滑 $f$，广义导数为各段普通导数加每个断点的跳跃权重冲激：$f'=f'_r+\sum_a[f(a+)-f(a-)]\delta(t-a)$；相应积分满足 $\int_{a-}^{c+}f'=f(c+)-f(a-)$。后续因果信号在 $t<0$ 为零；静止初值要求没有先前运动，并在无奇异输入时有零初始导数。
若 $p(s)=a_ns^n+\cdots+a_0$，$a_n\ne0$，单位冲激响应 $w=u\,z$ 满足正时间齐次方程及 $z(0+)=\cdots=z^{(n-2)}(0+)=0,z^{(n-1)}(0+)=1/a_n$。单位阶跃响应 $v$ 满足 $v'=w$。一般静止响应为 $x=w*q$，其中 $(f*g)(t)=\int_0^tf(t-\tau)g(\tau)\dd\tau$；卷积满足交换、结合及分配律，$f*\delta=f$，串联系统的响应核为两核卷积。
单边 Laplace 变换按 $0-$ 起积分，广义导数满足 $\Lap[f']=sF$，所以 $W=\Lap[w]=1/p(s)$，$X=WF$。常规导数则需扣除 $f(0+)$，不能混用。因果分段连续函数在通常指数增长条件下，同一变换确定所有连续点的值及断点两侧极限，单点值不受确定。
原文的 ``continued fractions'' 在上下文中应是部分分式分解，而不是连分数。对有理传递函数，其未被分子消去的分母零点给极点；因果有限指数多项式信号的收敛半平面在最右极点右侧。右侧极点决定长期的指数率和振荡频率，重极点还引入多项式因子；这不能不加条件推广到任意有非有理变换的函数。指数响应 $W(r)\ee^{rt}$ 在 $r$ 不为极点时成立；$W(\ii\omega)$ 给以右端为输入的复增益。
\originalsection{练习试卷与解答}
\begin{exercise}\label{cw:prex3-1}\textbf{原题 1.} $\omega>0$，用角频率 $\omega$ 的方波驱动谐振子：$\ddot x+4x=\operatorname{sq}(\omega t)$。
(a) 若存在，写一个周期解；(b) 哪些 $\omega$ 下没有周期解？
\end{exercise}
\begin{solution}
\textit{依据：同文件官方答案译审并补证（AI 辅助）；官方遗漏或错误另标。}\par
(a) 不共振时 $x=\tfrac4\pi\sum_{j\text{ 正奇数}}\sin(j\omega t)/[j(4-j^2\omega^2)]$。
(b) 频率 $j\omega=2$ 的非零谐波共振，故无周期解当且仅当 $\omega=2/j$，$j$ 为正奇数。此时强制产生线性增长的共振项，齐次解无法抵消；其余情形上述周期级数给出所需解。
\end{solution}
\begin{exercise}\label{cw:prex3-2}\textbf{原题 2.} $f(t)=0$ 于 $t<0$，$f(t)=1-t$ 于 $0<t<1$，$f(t)=1$ 于 $t>1$。
(a) 画图；(b) 画广义导数；(c) 用阶跃与冲激写广义导数。
\end{exercise}
\begin{solution}
\textit{依据：同文件官方答案译审并补证（AI 辅助）；官方遗漏或错误另标。}\par
(a) 负时间为零，$0$ 处跳升到 $1$，线性下降到 $(1,0)$，在 $1$ 处再次跳升到 $1$ 并保持。
(b)--(c) 常规部分在 $(0,1)$ 为 $-1$，其他地方为零；两个跳跃均为 $+1$，故 $f'=-u(t)+u(t-1)+\delta(t)+\delta(t-1)$。\ODEExamPracticeRamp
\end{solution}
\begin{exercise}\label{cw:prex3-3}\textbf{原题 3.}
(a) 求卷积 $t*t^6$；(b) 某算子的单位冲激响应为 $w=2u(t)t\ee^{-t}$，求静止初值下 $p(D)x=\ee^{-t}$ 的解（右端按因果信号理解）。
\end{exercise}
\begin{solution}
\textit{依据：同文件官方答案译审并补证（AI 辅助）；官方遗漏或错误另标。}\par
(a) $\int_0^t(t-\tau)\tau^6\dd\tau=t^8(1/7-1/8)=t^8/56$。
(b) $x=u(t)\int_0^t2(t-\tau)\ee^{-(t-\tau)}\ee^{-\tau}\dd\tau=u(t)t^2\ee^{-t}$。
\end{solution}
\begin{exercise}\label{cw:prex3-4}\textbf{原题 4.}
(a) 求静止初值下 $\ddot x+2\dot x+2x=1$ 的解的 Laplace 变换；(b) 求变换为 $2s/[(s+1)(s^2+2s+5)]$ 的函数。
\end{exercise}
\begin{solution}
\textit{依据：同文件官方答案译审并补证（AI 辅助）；官方遗漏或错误另标。}\par
(a) $(s^2+2s+2)X=1/s$，故 $X=1/[s(s^2+2s+2)]$。
(b) 部分分式为
\[\frac{2s}{(s+1)[(s+1)^2+4]}=-\frac{1}{2(s+1)}+\frac{(s+1)/2+2}{(s+1)^2+4}.\]
因此 $f=u(t)\ee^{-t}[-1/2+\tfrac12\cos2t+\sin2t]$。
\end{solution}
\begin{exercise}\label{cw:prex3-5}\textbf{原题 5.}
(a) 画 $2/[(s+1)(s^2+2s+5)]$ 的极点图；(b) 举一个函数，其 Laplace 变换恰有极点 $2,-3\pm4\ii$，无其他极点。
\end{exercise}
\begin{solution}
\textit{依据：同文件官方答案译审并补证（AI 辅助）；官方遗漏或错误另标。}\par
(a) 三个极点为 $-1,-1\pm2\ii$，见图。
\begin{center}\begin{tikzpicture}[scale=.7]\draw[->](-2.5,0)--(1,0)node[right]{$\operatorname{Re}s$};\draw[->](0,-2.7)--(0,2.7)node[above]{$\operatorname{Im}s$};\foreach\y in{-2,0,2}{\node[main]at(-1,\y){$\times$};}\node[below]at(-1,0){$-1$};\node[left]at(-1,2){$-1+2\ii$};\node[left]at(-1,-2){$-1-2\ii$};\end{tikzpicture}\end{center}
(b) 官方示例为 $\ee^{2t}+\ee^{-3t}\sin4t$。编者另给等价示例 $f=u(t)[\ee^{2t}+\ee^{-3t}\cos4t]$，变换为 $1/(s-2)+(s+3)/[(s+3)^2+16]$。两项在各自极点处的非零主部不被另一项取消，故恰有所需三个极点；官方的正弦示例以 $4/[(s+3)^2+16]$ 代替第二项，也满足同一要求。
\end{solution}

\originalchapter{练习期末考试（2010）}
题卷题头为 ``Practice Final Exam, 2010''，对应答案的题头实际为 ``Solutions of Spring 2008 Final Exam''。下面仅在具体小问题面与条件一致时采用“2008 原题答案与 2010 练习题核同后的译审”；题面不一致者独立补解，不能把整份 2008 答案宣称为 2010 官方解答。原第 1 题没有 (b)，以下保留其 (a),(c),(d),(e),(f)。第 2(d)、7(c) 有实质差异；另外修正答案中的算术错误和线性化坐标平移。
\begin{exercise}\label{cw:prfinal-1}\textbf{原题 1.} $y'=x^2-y^2$，设 $y=f(x)$ 是满足 $f(-2)=0$ 的解。
(a) 画斜率 $-2,0,2$ 的等斜线及线上的方向场；(c) 在同图画 $f$，求 $x=-2$ 时的斜率；(d) 估计 $f(100)$；(e) 若 $f$ 在 $x=a$ 达到极大值，求 $f(a)$；(f) 用两步 Euler 法估计 $f(-1)$。
\end{exercise}
\begin{solution}
\textit{依据：本题 (a,c,d,e,f) 分别与 2008 答案 (a,b,c,d,e) 核同；译审并补证。}
(a) 等斜线为 $x^2-y^2=m$，零斜线为 $y=\pm x$；$m=2$ 左右开、$m=-2$ 上下开，沿线画相应斜率的短线。
(c) 初始斜率为 $(-2)^2-0=4$。解从 $(-2,0)$ 上升，在左半平面上支 $y=-x$ 处转为下降；随后在右半平面 $y=x$ 附近取得极小值，再向上接近该线下侧。\ODEExamPracticeFinalField
(d) $f(100)\approx100$；代入 $f=x+\eta$ 的主阶平衡给 $\eta\sim-1/(2x)$，可细化为 $99.995$。
(e) 极大值处 $f'=0$，$f=\pm a$；又 $f''=2a<0$，故 $a<0$，该解在此处为正，得 $f(a)=-a$。
(f) 步长 $h=1/2$：第一步 $y_1=0+(1/2)4=2$，$x_1=-3/2$；第二步 $y_2=2+\tfrac12(9/4-4)=9/8$，因此 $f(-1)\approx1.125$。
\end{solution}
\begin{exercise}\label{cw:prfinal-2}\textbf{原题 2.}
(a)--(c) 考虑 $\dot x=2x-3x^2+x^3$。
(a) 画相线；(b) 各临界点之间及外侧均画至少一条代表解；(c) 非恒定解在 $t=a$ 有拐点时，$x(a)$ 可以是多少？
(d) 放射性 Bostonium（Bo）同位素半衰期为两年，衰变为 Cantabrigium（Ct）。MIT 反应堆装有这种材料；$t=0$ 时没有 Ct，从 $t=0$ 开始加入 Ct，至时刻 $t$ 累计加入量为 $t$ 摩尔。建立反应堆中 Ct 摩尔数的方程，写出初值。
(e) 解 $xy'+3y=x^2,y(1)=1$。
\end{exercise}
\begin{solution}
\textit{依据：(a)--(c)、(e) 与 2008 答案核同后译审并补证；(d) 题面不一致，编者补解（AI 辅助）。}
(a) 右端 $x(x-1)(x-2)$，四区间符号 $-,+,-,+$，$0,2$ 不稳定、$1$ 稳定。
(b) 三个常值解，两个内侧趋于 $1$ 的解和两个外侧向无穷逃逸的解构成全部七型。\ODEExamPracticeFinalPhase
(c) 令 $F(x)=x^3-3x^2+2x$，则 $x''=F'(x)F(x)$。非恒定解不会经过平衡点，所以 $F'(x)=3x^2-6x+2=0$，得 $x=1\pm1/\sqrt3$，两处确有凹凸性改变。
(d) 此处不能采用 2008 答案：该答案把所加材料改为 Bo，并取每年 $2$ 摩尔，两个条件都与 2010 题面不同。按现题，设 $B(t),C(t)$ 分别为 Bo、Ct 的摩尔数，一摩尔 Bo 衰变生成一摩尔 Ct，$k=\ln2/2$。若初始 Bo 量为 $B_0$，且无其他移除机制，则
\[B'=-kB,\quad B(0)=B_0;\qquad C'=1+kB,\quad C(0)=0.\]
因而 $B=B_0\ee^{-kt}$，$C'=1+kB_0\ee^{-kt}$，$C=t+B_0(1-\ee^{-kt})$。题面未给初始 Bo 装载量，因此不能唯一确定这一常数；若反应堆起初也没有 Bo，则特例为 $C'=1,C(0)=0,C=t$。这里如实保留材料和累计量，没有替原题暗中改写为另一道题。
(e) 除以 $x$ 后积分因子为 $x^3$，$(x^3y)'=x^4$，故 $y=x^2/5+C/x^3$。初值给 $C=4/5$，所以 $y=x^2/5+4/(5x^3)$，$x>0$。高清原页确认 2008 答案在除以 $x$ 后写为 $y'+3y/x=x$，与 2010 题面一致；不能把文本提取时错位的幂次当作题面差异。
\end{solution}
\begin{exercise}\label{cw:prfinal-3}\textbf{原题 3.}
(a) 求非负实数 $A,\omega,\phi$，使 $\operatorname{Re}[\ii\ee^{2\ii t}/(1+\ii)]=A\cos(\omega t-\phi)$；(b) 画 $\ee^{(1-\pi\ii)t}$ 的轨迹；(c) 将 $8\ii$ 的立方根写成 $a+b\ii$。
\end{exercise}
\begin{solution}
\textit{依据：三个小问均与 2008 原题答案核同，译审并补证。}
(a) $\ii/(1+\ii)=(1+\ii)/2=\ee^{\ii\pi/4}/\sqrt2$，故实部为 $\cos(2t+\pi/4)/\sqrt2$。取 $A=1/\sqrt2,\omega=2,\phi=7\pi/4$，满足非负要求。
(b) 轨迹为 $x=\ee^t\cos\pi t,y=-\ee^t\sin\pi t$，经过 $(1,0)$，随 $t$ 顺时针向外螺旋；每转一圈模长乘以 $\ee^2$。
\begin{center}\begin{tikzpicture}\begin{axis}[width=6cm,height=5cm,axis lines=middle,xmin=-4,xmax=5,ymin=-4,ymax=4,xlabel={实部},ylabel={虚部},samples=160]\addplot[main,domain=-1.5:1.4,->]({exp(x)*cos(deg(pi*x))},{-exp(x)*sin(deg(pi*x))});\addplot[only marks]coordinates{(1,0)};\end{axis}\end{tikzpicture}\end{center}
(c) $8\ii=8\ee^{\ii\pi/2}$，根为 $2\ee^{\ii\pi/6}=\sqrt3+\ii$、$2\ee^{5\ii\pi/6}=-\sqrt3+\ii$、$2\ee^{3\ii\pi/2}=-2\ii$。
\end{solution}
\begin{exercise}\label{cw:prfinal-4}\textbf{原题 4.}
(a)--(c) 对 $\ddot x+2\dot x+2x=q(t)$ 分别求一个解：
(a) $q=t^2+1$；(b) $q=\ee^{-2t}+1$；(c) $q=\sin t$，并求正弦解振幅。
(d)--(e) 已知 $t^3$ 是该方程的一个解：(d) 求 $q$；(e) 求通解。
\end{exercise}
\begin{solution}
\textit{依据：各小问题面核同；2008 答案译审并补证，其中 (d) 算术错误由编者重算。}
(a) 取 $x=at^2+bt+c$，系数条件 $2a=1,4a+2b=0,2a+2b+2c=1$ 给 $x=t^2/2-t+1$。
(b) $p(-2)=p(0)=2$，故 $x=(\ee^{-2t}+1)/2$。
(c) 取 $z=\ee^{\ii t}/(1+2\ii)$ 的虚部，得 $x=(\sin t-2\cos t)/5$，振幅为 $1/\sqrt5$。
(d) $q=(t^3)''+2(t^3)'+2t^3=6t+6t^2+2t^3$；2008 答案漏了最后的系数 $2$。
(e) 特征根 $-1\pm\ii$，通解为 $x=t^3+\ee^{-t}(C\cos t+D\sin t)$。
\end{solution}
\begin{exercise}\label{cw:prfinal-5}\textbf{原题 5.}
(a)--(b) $f(t)=\operatorname{sq}(t+\pi/2)$：(a) 画图；(b) 求 Fourier 级数并化简三角函数。
(c) 求 $\ddot x+x=\operatorname{sq}(t)$ 的一个解。
\end{exercise}
\begin{solution}
\textit{依据：与 2008 原题答案核同后的译审与补证。}
(a) $f=\operatorname{sgn}(\cos t)$，$-\pi/2<t<\pi/2$ 为 $1$、$\pi/2<t<3\pi/2$ 为 $-1$，以 $2\pi$ 周期延拓。
\begin{center}\begin{tikzpicture}[x=.65cm,y=.65cm]\draw[->](-5.3,0)--(5.3,0)node[right]{$t$};\draw[->](0,-1.4)--(0,1.5)node[above]{$f$};\draw[main,thick](-4.71,-1)--(-1.57,-1);\draw[main,thick](-1.57,1)--(1.57,1);\draw[main,thick](1.57,-1)--(4.71,-1);\node[below]at(-1.57,0){$-\pi/2$};\node[below]at(1.57,0){$\pi/2$};\node[below]at(4.71,0){$3\pi/2$};\node[left]at(0,1){$1$};\node[left]at(0,-1){$-1$};\end{tikzpicture}\end{center}
(b) 平移每个奇谐波得
\[f(t)=\frac4\pi\left(\cos t-\frac{\cos3t}{3}+\frac{\cos5t}{5}-\cdots\right).\]
(c) 首谐波共振，$(D^2+1)(-t\cos t/2)=\sin t$；其余谐波用非共振响应。因此
\[x_p=\frac4\pi\left[-\frac12t\cos t+\sum_{j\ge3,\ j\text{ 奇数}}\frac{\sin jt}{j(1-j^2)}\right].\]
本小问仅要求一个解，解含增长项，不能错误声称它是周期解。
\end{solution}
\begin{exercise}\label{cw:prfinal-6}\textbf{原题 6.}
(a)--(d) 一学生在夺旗游戏中的位移由下图给出，刻度间隔为 $1$。
(a) 画广义速度 $v$；(b) 用阶跃、必要时冲激写 $v$；(c) 画广义加速度 $v'$；(d) 用阶跃、必要时冲激写 $v'$。
\begin{center}\begin{tikzpicture}[scale=.95]\draw[->](-1,0)--(5.3,0)node[right]{$t$};\draw[->](0,-1.5)--(0,1.5)node[above]{位移};\foreach\x in{1,2,3,4,5}{\draw(\x,-.06)--(\x,.06);\node[below]at(\x,0){$\x$};}\node[left]at(0,1){$1$};\node[left]at(0,-1){$-1$};\draw[main,thick](-.8,0)--(0,0)--(1,1)--(3,-1)--(4,0)--(5,0);\end{tikzpicture}\end{center}
(e) 某系统单位冲激响应为 $w(t)$，输入 $q=1$ 于 $0<t<1$，其他时刻为零。求积分界限 $a(t),b(t)$，使静止初值解可写为 $x(t)=\int_{a(t)}^{b(t)}w(\tau)\dd\tau$。
\end{exercise}
\begin{solution}
\textit{依据：原图和全部小问与 2008 原题答案核同，译审并补证。}
(a)--(b) 位移连续，所以速度没有冲激；斜率在 $(0,1),(1,3),(3,4)$ 上依次为 $1,-1,1$，外侧为零。
\[v=u(t)-2u(t-1)+2u(t-3)-u(t-4).\]
(c)--(d) 速度在 $0,1,3,4$ 的跳跃为 $1,-2,2,-1$，故
\[v'=\delta(t)-2\delta(t-1)+2\delta(t-3)-\delta(t-4).\]
常规加速度为零，图为相应权重的四个冲激。
\ODEExamMotionDerivatives
(e) $t\ge0$ 时卷积为 $\int_0^tw(\tau)q(t-\tau)\dd\tau$，有效区间要求 $0<t-\tau<1$，故 $a(t)=\max(0,t-1),b(t)=t$。因果系统在 $t<0$ 响应为零，此时可另取 $a=b=0$。
\end{solution}
\begin{exercise}\label{cw:prfinal-7}\textbf{原题 7.} $p(D)=2D^2+8D+16I$。
(a) 求传递函数；(b) 求单位冲激响应；(c) 求静止初值下 $p(D)x=\sin t$ 的解的 Laplace 变换。
\end{exercise}
\begin{solution}
\textit{依据：(a)--(b) 与 2008 原题答案核同后译审；(c) 的初值不匹配，编者补解（AI 辅助）。}
(a) $W=1/(2s^2+8s+16)=1/[2((s+2)^2+4)]$。
(b) $w=\tfrac14u(t)\ee^{-2t}\sin2t$，其 $w'(0+)=1/2$ 对应最高阶系数 $2$。
(c) 静止初值给 $X=W/(s^2+1)=1/[(2s^2+8s+16)(s^2+1)]$。2008 答案采用 $x(0)=1,x'(0)=2$，额外出现 $2s+12$，与本题静止条件不符，不能沿用。
\end{solution}
\begin{exercise}\label{cw:prfinal-8}\textbf{原题 8.}
(a)--(b) $A=\begin{pmatrix}2&12\\3&2\end{pmatrix}$：(a) 求特征值；(b) 每个特征值求一个非零特征向量。
(c) $B$ 的特征值为 $1,2$，相应特征向量为 $(1,1)^T,(-1,1)^T$，求 $\ee^{Bt}$。
(d) 求 $\dot u=Bu,u(0)=(2,1)^T$ 的解。
\end{exercise}
\begin{solution}
\textit{依据：四个小问均与 2008 原题答案核同，译审并补证。}
(a)--(b) $\det(A-\lambda I)=(2-\lambda)^2-36$，根为 $8,-4$；相应向量可取 $(2,1)^T,(2,-1)^T$。
(c) $P=\begin{pmatrix}1&-1\\1&1\end{pmatrix}$，$P^{-1}=\tfrac12\begin{pmatrix}1&1\\-1&1\end{pmatrix}$，因此
\[\ee^{Bt}=\frac12\begin{pmatrix}\ee^t+\ee^{2t}&\ee^t-\ee^{2t}\\\ee^t-\ee^{2t}&\ee^t+\ee^{2t}\end{pmatrix}.\]
(d) 初值的特征向量系数为 $3/2,-1/2$，故 $u=\tfrac12(3\ee^t+\ee^{2t},3\ee^t-\ee^{2t})^T$。
\end{solution}
\begin{exercise}\label{cw:prfinal-9}\textbf{原题 9.}
(a) $B$ 的特征值为 $1,2$，相应向量为 $(1,1)^T,(-1,1)^T$，画 $\dot u=Bu$ 的相图。
(b) $A=\begin{pmatrix}a&-2\\2&1\end{pmatrix}$，对 $\dot u=Au$，分别求满足下列类型的所有 $a$：
(i) 鞍点；(ii) 星形结点；(iii) 稳定结点；(iv) 稳定螺旋并判转向；(v) 不稳定螺旋；(vi) 不稳定亏损结点。
\end{exercise}
\begin{solution}
\textit{依据：原小问与 2008 答案核同，译审并补证。}
(a) 两个正特征值给不稳定结点。特征坐标 $u=\xi(1,1)^T+\eta(-1,1)^T$ 中 $\xi'=\xi,\eta'=2\eta$，所以非直线轨道为 $\eta=C\xi^2$，在原点切于 $y=x$，向外流动；$y=\pm x$ 是直线轨道。图由这一精确参数式重绘。\ODEExamNodePortrait
(b) $T=a+1,\Delta=a+4,T^2-4\Delta=(a-5)(a+3)$。因此
\[
\begin{array}{c|l}
\text{(i)}&a<-4\\
\text{(ii)}&\text{无}\\
\text{(iii)}&-4<a<-3\\
\text{(iv)}&-3<a<-1\quad\text{逆时针}\\
\text{(v)}&-1<a<5\\
\text{(vi)}&a=5
\end{array}
\]
星形要求重根并有两个独立特征向量，即 $A$ 为标量矩阵；非零非对角元排除了它。螺旋区间在 $(x,y)=(1,0)$ 处 $\dot y=2>0$，故逆时针。结点的迹、行列式及判别式条件与期末第 9 题同，尽管两题矩阵不同，分类参数恰相同。
\end{solution}
\begin{exercise}\label{cw:prfinal-10}\textbf{原题 10.}
(a)--(c) 非线性系统 $\dot x=x^2-y^2,\dot y=x^2+y^2-8$。
(a) 求全部平衡点；(b) 求西南象限平衡点的 Jacobian；(c) 此点是稳定螺旋，趋于它的解的大 $t$ 横坐标，原题写作 $A\ee^{at}\cos(\omega t-\phi)$ 的近似，求所有此类解共有的常数。
(d) 回到第 2 题的 $\dot x=2x-3x^2+x^3$，写出趋于稳定平衡点的解在大 $t$ 时的近似公式。
\end{exercise}
\begin{solution}
\textit{依据：与 2008 原题答案核同后的译审与补证；(c) 原题省略平衡坐标平移，此处明确补正。}
(a) 联立 $x^2=y^2$、$x^2+y^2=8$，得 $(\pm2,\pm2)$ 四点。
(b) 西南点为 $(-2,-2)$，$J=\begin{pmatrix}2x&-2y\\2x&2y\end{pmatrix}$，故 $J(-2,-2)=\begin{pmatrix}-4&4\\-4&-4\end{pmatrix}$。
(c) 特征值为 $-4\pm4\ii$，共有参数为 $a=-4,\omega=4$。近似式应写为 $x(t)\approx-2+A\ee^{-4t}\cos(4t-\phi)$；若原式用于位移 $x+2$ 则成立，用于横坐标本身则漏了 $-2$。振幅和相位由具体解决定。
(d) 稳定点 $x=1$，$F'(1)=2-6+3=-1$。令 $z=x-1$，有 $z'=-z+z^3$，线性化给 $x(t)\approx1+C\ee^{-t}$，$C$ 随解变化。平衡解对应 $C=0$。
\end{solution}

\originalsection{考试附用公式与符号}
各原公式页的数学信息在此统一列出；$D=\mathrm d/\mathrm dt$，$u$ 是单位阶跃，$\delta_a=\delta(t-a)$，$u_a=u(t-a)$。$\operatorname{sq}$ 是 $2\pi$ 周期奇方波，在 $0<t<\pi$ 为 $1$、$-\pi<t<0$ 为 $-1$，跳点值不影响分段方程和广义函数。
若 $p(r)\ne0$，$p(D)x=A\ee^{rt}$ 有特解 $A\ee^{rt}/p(r)$；若 $p(r)=0,p'(r)\ne0$，特解为 $At\ee^{rt}/p'(r)$。亏损二阶矩阵 $A$ 的重根为 $\lambda$，$(A-\lambda I)v=0$，若 $(A-\lambda I)w=v$，则 $\ee^{\lambda t}(tv+w)$ 是系统解。
因果信号的 Laplace 变换为 $F(s)=\int_{0-}^\infty f(t)\ee^{-st}\dd t$，在适当右半平面定义。线性、$s$ 平移、时间平移、$s$ 微分、时间微分及卷积规则为
\[
\begin{aligned}
\Lap[af+bg]&=aF+bG,&\Lap[\ee^{at}f(t)]&=F(s-a),\\
\Lap[u(t-a)f(t-a)]&=\ee^{-as}F(s)\quad(a\ge0),&\Lap[tf]&=-F'(s),\\
\Lap[f']&=sF(s)\quad\text{（广义导数）},&\Lap[f'_r]&=sF(s)-f(0+)\quad\text{（正时间连续）},\\
\Lap[f*g]&=FG,&(f*g)(t)&=\int_0^tf(t-\tau)g(\tau)\dd\tau.
\end{aligned}
\]
单位冲激响应满足 $W=\Lap[w]=1/p(s)$；静止响应 $X=WF$。常用表为
\[
\begin{array}{c|c@{\qquad}c|c}
f&\Lap[f]&f&\Lap[f]\\\hline
u(t)&1/s&u(t)\ee^{at}&1/(s-a)\\
u(t)t^n&n!/s^{n+1}&\delta(t-a)&\ee^{-as}\\
u(t)\cos\omega t&s/(s^2+\omega^2)&u(t)\sin\omega t&\omega/(s^2+\omega^2)\\
u(t)t\sin\omega t&2\omega s/(s^2+\omega^2)^2&u(t)t\cos\omega t&(s^2-\omega^2)/(s^2+\omega^2)^2\\
u(t-a)&\ee^{-as}/s&&
\end{array}
\]
$2\pi$ 周期 Fourier 系数为
\[f\sim\frac{a_0}{2}+\sum_{j=1}^\infty(a_j\cos jt+b_j\sin jt),\qquad
 a_j=\frac1\pi\int_{-\pi}^\pi f\cos jt\dd t,\quad b_j=\frac1\pi\int_{-\pi}^\pi f\sin jt\dd t.\]
异频正弦或余弦的积分内积为零，同频平方积分为 $\pi$（频率正整数）；方波展开为 $\operatorname{sq}(t)\sim\tfrac4\pi\sum_{j\text{ 正奇数}}\sin jt/j$。这些规则原指南已给，本节把分段、因果和非共振条件显式保留。

\originalchapter{来源覆盖与补充范围}
\originalsection{18.03 的课堂主线}
本册按数学依赖顺序重组, 不是逐讲逐页翻译. 官方目录列出36份课堂讲义, 第10、19、31次为考试位置, 目录无相应课堂讲义. 下表的数字是原课 session 编号, 不是本册页码. 所有原件逐文件页数、官方链接、SHA-256 和持有状态见 \texttt{source-manifest.json}.
\begin{longtable}{p{.16\textwidth}p{.47\textwidth}p{.25\textwidth}}
\toprule 原课讲次 & 数学内容 & 本册对应 \\\midrule\endhead
1 & 方向场、分离变量、存在唯一性 & 第1、2章 \\
2 & 数值方法 & 第1章误差衔接; 算法见数值分析册 \\
3--4 & 线性建模、积分因子 & 第2章 \\
5--7 & 复数、复指数、正弦响应 & 第4、5章; 复数代数作先修 \\
8--9 & 相线、线性与非线性 & 第3、9、10章 \\
11--12 & 特征根、阻尼振动 & 第4、5章 \\
13--16 & 指数响应、拍、共振与增益 & 第5章 \\
17--18 & LTI、RLC 与工程响应 & 第5、6章; 原工程实例不逐一重述 \\
20--22 & Fourier 级数、周期响应与共振 & 第7章 \\
23--25 & 阶跃、脉冲、卷积 & 第6章 \\
26--30 & Laplace、极点、传递函数 & 第6章; 差分补充见第13章 \\
32--34 & 线性系统、特征值、重根 & 第8章 \\
35--36 & 相图、模态、矩阵指数 & 第8、9章 \\
37--39 & 非线性、单摆、极限环与混沌入门 & 第9--11、13章 \\
\bottomrule
\end{longtable}
Miller 的 \textit{Supplementary Notes} 整本为155个 PDF 页面, 封面署名年份含2004、2006、2008和2010; 本次课程版本是2010. 核对 PDF 正文发现第1章与前言的分件含不同版本的文字, 因而另行保留这两份, 不按文件名强行去重. 其余25章与3个附录的数学内容对应整本, 单独的 OCW 末页和版式可能不同, 不把它们视为字节相同文件. 这28个重复分件登记官方链接并归并到整本持有. 官方资源网页的个别描述有误: 标记 chapter\_8 的原件实际为 \textit{RLC circuits}, 以 PDF 正文为准.

补充讲义第1--6节的记号、线性建模和复指数纳入第1、2、5章; 第7--15节的拍、RLC、响应、阻尼和 Wronskian 纳入第4、5章; 第16--24节的 Fourier、脉冲、卷积和极点纳入第6、7章; 第25--26节的系统与相图纳入第8、9章. 附录A的流行病主题在第3章给出 SIR 动力学与最终规模推导; 附录B的桥梁颤振历史、附录C的 phugoid 工程详例留作延伸, 未声称本册完整译出. 这三篇随整本归档, 其来源身份与教学覆盖范围分开记录.

Mattuck 的 \textit{Notes and Exercises} 以官方阅读目录署名为依据, 不将所有文件的 PDF 元数据中统一出现的 Miller 名字误作逐件作者. 笔记16件涉及图解、复数、积分形式、算子、输入响应、Laplace、线性系统 LS.1--6、图示系统及极限环. 其中 S 原件的实际标题为 Stability, 对应第5、8章, 不依据网页的含糊描述臆测其主题. 第12章的幂级数、正则奇点与边值内容为编者补充; 原练习目录虽含 power series 主题, 不把它冒称为已完整翻译的课堂讲义. 14份原练习及解答在调查清单中标为范围外链接, 不归档、不宣称逐题翻译. Jeremy Orloff 的差分与 Z 变换补充为14页, 第13章只采用线性响应及稳定极点衔接, 不扩展成离散控制课程.
\originalsection{12.006J 的动力系统补充}
16份文件按原课程组合讲次登记. 含出版社保留图片或经许可使用的第三方图片者, 整份原件仅保留官方链接; 本册不提取、描摹或重新标注这些图片. 所用相图由明确方程独立构造.
\begin{longtable}{p{.16\textwidth}p{.48\textwidth}p{.24\textwidth}}
\toprule 原课讲次 & 原内容 & 本册采用边界 \\\midrule\endhead
1 & 非线性与混沌引论 & 第3、13章概念衔接 \\
2--3 & 一维流与分岔 & 第3章 \\
4--5 & 单摆、二维稳定性 & 第9、10章 \\
6--7 & 耗散与体积收缩 & 第10、13章 \\
8--9 & 受迫振子、极限环 & 第5、11章 \\
10--11 & 二维分岔 & 第11章正规模型; 一般定理未全文证明 \\
12--14 & 动力系统谱分析 & 第7章仅线性周期响应; 随机谱留延伸 \\
15--16 & Poincar\'e 截面 & 第11章 \\
17--18 & 流体、Rayleigh--B\'enard 对流 & 留延伸; 偏微分方程不纳入 \\
19 & 奇异吸引子 & 第13章机制提示; 不作全局分类 \\
20--21 & Lorenz 方程 & 第13章平衡、吸收界及线性失稳 \\
22 & H\'enon 吸引子 & 留延伸; 未把二维映射全部理论纳入 \\
23 & 分形维数 & 留延伸; 未建立分形测度理论 \\
24 & Lyapunov 指数 & 第13章定义与变分增长; 不证明遍历极限存在 \\
25--27 & 倍周期通往混沌 & 第13章 logistic 二周期; 不证明通用常数 \\
28 & 间歇性、准周期 & 留延伸 \\
\bottomrule
\end{longtable}
\originalsection{编者补充与引用工具}
本册补全局部积分迭代、最大延续、连续依赖、矩阵指数、线性化稳定性的 Lyapunov 证明、紧集 LaSalle 论证、平面周期乘子与 Dulac 判据. 幂级数的正系数支配证明、边值 Green 核与相容条件、Lorenz 吸收界和三次稳定判据也作为补充论证明确写出. 例题及练习为本册整理或另设, 不复用商业教材题面.

有限维 Jordan 标准形、一致多项式逼近、基本积分定理、隐函数定理以及平面 Jordan 曲线和 Green 公式作为已声明的先修工具使用. Hartman--Grobman 与 Poincar\'e--Bendixson 的一般定理准确列条件后引用, 不声称已给完整拓扑证明. Van der Pol 极限环的全局唯一性使用核实的 LC 第5页 Levinson--Smith 引用版本, 同样明确为引用工具. 一般 Hopf 定理、稳定流形定理和混沌的全局存在证明不作为本册已证明结果. 数值绘图或抽样轨道不承担这些证明.

\begin{thebibliography}{9}
\bibitem{primary} Haynes Miller, Arthur Mattuck. \textit{18.03 Differential Equations, Spring 2010}. MIT OpenCourseWare. \href{https://ocw.mit.edu/courses/18-03-differential-equations-spring-2010/pages/lecture-notes/}{课堂讲义}; \href{https://ocw.mit.edu/courses/18-03-differential-equations-spring-2010/pages/readings/}{阅读目录及补充讲义}.
\bibitem{dynamics} Daniel Rothman. \textit{12.006J Nonlinear Dynamics: Chaos, Fall 2022}. MIT OpenCourseWare. \href{https://ocw.mit.edu/courses/12-006j-nonlinear-dynamics-chaos-fall-2022/pages/lecture-notes/}{官方讲义目录}.
\bibitem{coursework} MIT OpenCourseWare, 18.03 Spring 2010. \href{https://ocw.mit.edu/courses/18-03-differential-equations-spring-2010/pages/assignments/}{原作业}, \href{https://ocw.mit.edu/courses/18-03-differential-equations-spring-2010/pages/recitations/}{习题课}, \href{https://ocw.mit.edu/courses/18-03-differential-equations-spring-2010/pages/exams/}{考试}. Mattuck 七个题库单元按原阅读目录署名. 全部题源版本、题面与官方答案匹配范围见 coursework-source-manifest.json 及逐题清单.
\bibitem{license} MIT OpenCourseWare. \href{https://ocw.mit.edu/pages/privacy-and-terms-of-use/}{使用条款}; \href{https://creativecommons.org/licenses/by-nc-sa/4.0/}{CC BY-NC-SA 4.0}. 逐文件例外见 source-manifest.json.
\end{thebibliography}
\end{document}
